\documentclass[openany]{book}

\usepackage{hyperref}
\usepackage{bookmark}
\usepackage[dvipsnames]{xcolor}
% \usepackage{graphicx}
\usepackage{geometry}
\geometry{
	paperwidth=4.25in, % Adjust the width of the page
	paperheight=7.5in, % Adjust the height of the page
	inner=0.1in, % Inner margin
	outer=0.1in, % Outer margin
	top=0.4in, % Top margin
	bottom=0.5in, % Bottom margin
	bindingoffset=0in % at least 0.75" to 0.875" inner margins.
% No text can exist within 0.25 inches of the physical edge of the book (The Safe Zone). If your margins are 0.1in, their automated uploader will reject
% 4.25" x 7.0" or 5" x 8". For an 800-page book, a 5" x 8" size is highly recommended
	% paperwidth=5in,       % Shift to a standard industry pocket trim
	% paperheight=8in,      % Standard pocket height
	% inner=0.4in,          % Gives you room to actually read without text spilling off the edge
	% outer=0.4in,          % Safe from the 0.25" printer rejection zone
	% top=0.4in, 
	% bottom=0.5in, 
	% bindingoffset=0in     % Keep at 0 for editing; adjust later for the final spine glue
% Or these:	
	% paperwidth=5in,       % Standard available pocket trim size
	% paperheight=8in,      
	% inner=0.375in,        % Gives the thick 800-page spine some breathing room
	% outer=0.3in,          % Looks incredibly tight and "Dessain-esque", but safely clear of the 0.25" failure limit
	% top=0.35in, 
	% bottom=0.4in, 
	% bindingoffset=0in     % Keep at 0 for editing; adjust later for the final spine glue
% Production margins
	% paperwidth=5in,
    % paperheight=8in,
    % twoside,                 % Crucial: alternates margins for left/right pages
    % inner=0.375in,           % Visual whitespace from text to the gutter curve
    % bindingoffset=0.5in,     % The "sacrifice zone" swallowed by the glue/curve
    % outer=0.3125in,          % Tight, old-school Dessain look (safely over the 0.25" POD limit)
    % top=0.375in,
    % bottom=0.4in
	}
% \usepackage{xparse}
\usepackage{titling}
\usepackage{lettrine}
\usepackage{multicol}
\usepackage{titlesec}
\usepackage{soul}
% \usepackage{pifont}
\usepackage{gregoriosyms}
% \usepackage{etoolbox}
\usepackage{xstring}% For string manipulation
\usepackage{fancyhdr}
\usepackage{setspace}
% \usepackage{lua-widow-control}
\usepackage{fontspec}
% \setmainfont{Georgia}
% \setmainfont{Times New Roman}
% \setmainfont{Libre Caslon Text}
% \setmainfont{Charis}
\setmainfont{Merriweather}
% \setmainfont{Libertinus Serif}
% \setmainfont{Gentium}% 93% of Charis, thinner
% \setmainfont{Source Serif 4 Regular}% too wide
% \setmainfont{EB Garamond}% Too thin
% \setmainfont{Junicode}
% [
	% Width=SemiCondensed, % or Condensed if desperate
	% Ligatures={Common, Rare, Historic}, % Turn on the medieval ligatures
	% Numbers=OldStyle
% ]
% \usepackage{lua-ul}
% \usepackage{luacolor}
\usepackage[latin]{babel}
\babelprovide[hyphenrules=latin]{latin}
% \usepackage{polyglossia}
% \setdefaultlanguage[variant=ecclesiastic,hyphenation=liturgical]{latin}
\usepackage[titles]{tocloft}
% \usepackage{imakeidx}

% Indices setup code
% \indexsetup{
	% level=\section*,
	% toclevel=section,
	% noclearpage,
	
	% othercode=\footnotesize\thispagestyle{empty}
    % othercode={%
		% This forces the index entries to start at the absolute edge of the column
		% \setlength{\leftskip}{0pt}%
		% Adjust this if you want a hanging indent for long entries
		% \everypar{\hangindent=1em \hangafter=1}% 
    % }
% }

% \makeindex[name=H, title=Index Hymnorum, columns=2, columnseprule, columnsep=0.1in, options=-s index.ist]
% \makeindex[name=R, title=Index Responsoriorum, columns=2, columnseprule, columnsep=0.1in]
% \makeindex[name=S, title=Index Responsoriorum Breve, columns=2, columnseprule, columnsep=0.1in]
% \makeindex[name=P, title=Index Psalmorum et Canticorum, columns=2, columnseprule, columnsep=0.1in, options=-s index.ist]
% \makeindex[name=I, title=Index Invitatorium, columns=2, columnseprule, columnsep=0.1in, options=-s index.ist]

% \newcommand{\cprintindex}[2]{%
    % #1 = index name (e.g., P, H, R)
    % #2 = display title (e.g., Index Psalmorum et Canticorum)
    % \phantomsection%
    % \addcontentsline{toc}{section}{#2}%
    % \hypertarget{#2}{}%
    % \renewcommand{\rightmark}{#2}%
    % \printindex[#1]%
% }

\usepackage[final]{microtype}

% \microtypecontext{spacing=nonfrench} % Slightly expands text spacing
\SetProtrusion{encoding=*,family=*}{. = {,500}}
\SetExpansion{encoding=*,family=*}{A = 500}
\SetExtraKerning[unit=space]{encoding={*},family={*}}{
  \textquoteleft = {500, }, % adjust quotes
  : = {100,100}
}

% https://tug.org/FontCatalogue/otherfonts.html#initials
% \usepackage{Acorn}% humanist
% \usepackage{AnnSton}% not appropriate
% \usepackage{ArtNouv}% more gothic
% \usepackage{ArtNouvc}% not appropriate
% \usepackage{Carrickc}% not medieval
% \usepackage{Eichenla}
% \usepackage{Eileen}% same as EileenBl but inverted, white lettering
% \usepackage{EileenBl}% identical to Romantik except for more spacing
% \usepackage{Elzevier}% humanist, very thin
% \usepackage{GotIn}% gothic blackletter
% \usepackage{GoudyIn}% humanist
% \usepackage{Kinigcap}% not appropriate
% \usepackage{Konanur}% not appropriate
% \usepackage{Kramer}% too modern block
% \usepackage{MorrisIn}% identical to Eichenla but slightly larger and rougher leaves
% \usepackage{Nouveaud}% too square
% \usepackage{Romantik}
% \usepackage{Rothdn}% looks like yfonts yinipar
% \usepackage{Royal}% gothic blackletter
% \usepackage{Sanremo}% not appropriate
% \usepackage{Starburst}% not appropriate
% \usepackage{Typocaps}% not appropriate
% \usepackage{Zallman}% humanist
% \input{Zallman.fd}
% \newcommand*\zallmancaps{\usefont{U}{Zallman}{xl}{n}}
% \newcommand*\zallmancaps{\usefont{TU}{Georgia}{b}{n}}
% third is ul/xl/l/m/sb/b/eb/bx fourth is n/it/sl/sc
\newcommand*\zallmancaps{\fontspec[Weight=Bold]{Libre Caslon Text}}
% \newcommand*\zallmancaps{\fontspec[Weight=700]{Georgia}}
% [Instance=Condensed Bold] % Name of size
% 300 = Light; 400 = Regular; 500 = Medium; 600 = SemiBold; 700 = Bold; 800 = ExtraBold; 900 = Black

% Automatically use zallmancaps in lettrine
% \renewcommand{\LettrineFontHook}{\zallmancaps}

% Adjust font spacing
\AtBeginEnvironment{multicols}{%
% \BeforeBeginEnvironment{multicols}{%
	% \hyphenpenalty=50
	% lower it
	\hyphenpenalty=20
	% \tolerance=200
	% increase it, but not noticeable here
	% \tolerance=500
	% shoudnt it be decreased?
	% \tolerance=50
	% \pretolerance=100
	% lower it to increase hyphens; increase to allow for wider gaps
	% \pretolerance=50
	% \emergencystretch=0pt
	% increase it
	% \emergencystretch=3em
	% Doesnt seem to make a difference
	% \sloppy
	\emergencystretch=1em  % Allow some stretch
	\hfuzz=0.5pt          % Don't complain about small overfulls
    \raggedcolumns  % Instead of balanced columns
    \clubpenalty=10000
    \widowpenalty=10000  
    \brokenpenalty=4991
    \predisplaypenalty=10000
    \postdisplaypenalty=1549
    \displaywidowpenalty=1602
}

% Adjust headheight and topmargin
% \setlength{\headheight}{13.6pt} % Set to at least 13.59999pt
\setlength{\headheight}{32pt} % Set to at least 13.59999pt
% \addtolength{\topmargin}{-1.6pt} % Adjust topmargin to compensate
% \addtolength{\topmargin}{-11.6pt} % Adjust topmargin to compensate
% \addtolength{\topmargin}{-14pt} % Adjust topmargin to compensate
\setlength{\headsep}{10pt}      % Space between header and text body
\setlength{\footskip}{12pt}     % Space between text body and footer

% Reduce padding and whitespace
\setlength{\parskip}{0.1em} % Reduce space between paragraphs
\setlength{\parindent}{1em} % Adjust paragraph indentation
% \setlength{\columnsep}{0.2in} % Adjust the space between columns
% \setlength{\columnsep}{0.1in} % Reduce space between columns
\setlength{\columnsep}{0.08in} % Reduce space between columns
\setlength{\columnseprule}{0.4pt} % Add a rule between columns and set its thickness

% Reduce line spacing
\setstretch{0.85} % Reduce line spacing slightly

\makeatletter
\def\leftmarginnote{%
	\lrmarginnote{\hskip -\marginparsep \hskip -12.25em}}

\def\rightmarginnote{%
	\lrmarginnote{\hskip\columnwidth \hskip -1em}}

\long\def\lrmarginnote#1#2{%
\vadjust{#1%
\smash{\vtop{{%
		\hsize\marginparwidth
        \@parboxrestore
        \@marginparreset
		\small
#2}}}}}
\makeatother

% \makeatletter
% \DeclareRobustCommand{\Vbar}{\vers@resp{-0.1em}{V}}
% \DeclareRobustCommand{\Rbar}{\vers@resp{0pt}{R}}
% \newcommand{\vers@resp}[2]{%
  % {\ooalign{\hidewidth\kern#1\raisebox{0.2ex}{\rotatebox[origin=c]{-20}{$\m@th\rceil$}}\hidewidth\cr#2\cr}}%
% }%
% \makeatother

% Defined by gregoriosyms
% \newcommand{\Vbar}{℣}
% \newcommand{\Rbar}{℟}

\frenchspacing

\def\V{\color{Red}\Vbar. \color{black}}
\def\R{\color{Red}\Rbar. \color{black}}

\newcommand{\doxology}{%
	{\bfseries \color{Red} G}lória Patri, et Fílio, * et Spirítui Sancto. \\
	{\bfseries \color{Blue} S}icut erat in princípio, et nunc, et semper, * et in sǽcula sæculórum. Amen. \\
}

\newcommand{\requiem}{%
	Réquiem ætérnam * dona eis, Dómine. \\
	Et lux perpétua * lúceat eis. \\
}

\input{psalms}

% Controls how to display Psalms.
\newcommand{\psalm}[1]{%
	% \noindent{\textit{\csname psalm#1Incipit\endcsname}}.%
	\noindent{\csname psalm#1Incipit\endcsname}.%
	% \vspace{-.75em}\noindent\begin{center}Ps. #1. \textit{\csname psalm#1Incipit\endcsname}\end{center}\vspace{-.75em}%
	% \csname psalm#1Text\endcsname%
}

% Define full Psalms with alternately colored, bold firstletters
% Track the verses

\newcommand{\fullpsalm}[1]{%
	% \index[P]{\csname psalm#1Incipit\endcsname}
	\begingroup % Start local group to limit scope
		\setlength{\parindent}{.2em} % Set paragraph indentation
		\let\psalmtext\empty % define command
		\expandafter\let\expandafter\psalmtext\csname psalm#1Text\endcsname % Display the psalm text
		\newif\ifblue
		\bluetrue
		\newcommand{\processline}[1]{%
			\StrLeft{##1}{1}[\firstletter]%
			\StrGobbleLeft{##1}{1}[\restofline]%
			\par\hspace*{.1em}%
			\ifblue
				{\color{Blue}\textbf{\firstletter}}%
				\bluefalse
			\else
				{\color{Red}\textbf{\firstletter}}%
				\bluetrue
			\fi
			\restofline%
		}
		\renewcommand{\\}{\processline} % Redefine \\ to use line processor
		\psalmtext
		\par % Add explicit paragraph break after psalm
	\endgroup
}

% Same but full Psalms
% \newcommand{\fullpsalm}[1]{%
	% % \noindent{\textit{\csname psalm#1Incipit\endcsname}}.%
	% % \vspace{-.75em}\noindent\begin{center}Ps. #1. \textit{\csname psalm#1Incipit\endcsname}\end{center}\vspace{-.75em}%
	% % \csname psalm#1Text\endcsname% % No indent
	% \begingroup % Start local group to limit scope
		% \setlength{\parindent}{.25em} % Set paragraph indentation
		% \renewcommand{\\}{\par\hspace*{.25em}} % Redefine \\ to end the line and add indentation
		% \csname psalm#1Text\endcsname % Display the psalm text
	% \endgroup % End local group
% }

% Same but no indent for credo litanie
\newcommand{\noindentpsalm}[1]{%
	\csname psalm#1Text\endcsname % No indent
}

% Same but Hymns
\newcommand{\hymn}[3]{%
	% \index[H]{#2}
	\hypertarget{#1}{\label{#1}}
		\setlength{\parindent}{.2em} % Set paragraph indentation
		% \setlength{\parindent}{1em} % Set paragraph indentation
	#3 %
}

% Define the responsory environment
\NewDocumentEnvironment{responsory}{O{} m m}{%
	\leavevmode% forces horizontal mode
	\unskip% Remove any previous spacing
	\hypertarget{#1}{\label{#1}}
	\def\firstline{#2}%
	\def\versicle{#3}%
	% \StrBefore{\firstline}{*}[\partone]%
	\StrBehind{\firstline}{*}[\rawparttwo]%
	\StrGobbleLeft{\rawparttwo}{1}[\parttwo] % Remove the leading space
	\StrBefore{\firstline}{ }[\firstword]%
	\StrBefore{\parttwo}{ }[\secondword]%
	% \unskip\ignorespaces% Prevent paragraph break
% The format of the responsory
	% \noindent \R \firstline \ \V \versicle \ \parttwo \ Gloria. \firstword. \ \secondword.
	% \noindent \R \firstline \ \V \versicle \ \secondword.
	% \newline \R \firstline \ \V \versicle \ \secondword.
	\oldnewline \R \firstline \ \V \versicle \ \secondword.%
	% \linebreak \R \firstline \ \V \versicle \ \secondword.
	% \index[R]{\firstline}
}{%
	% \unskip\ignorespaces% Prevent paragraph break
}

% Define the responsory-doxology environment
\NewDocumentEnvironment{responsory-doxology}{O{} m m}{%
	\leavevmode% forces horizontal mode
	\unskip% Remove any previous spacing
	\hypertarget{#1}{\label{#1}}
	\def\firstline{#2}%
	\def\versicle{#3}%
	% \StrBefore{\firstline}{*}[\partone]%
	\StrBehind{\firstline}{*}[\rawparttwo]%
	\StrGobbleLeft{\rawparttwo}{1}[\parttwo] % Remove the leading space
	\StrBefore{\firstline}{ }[\firstword]%
	\StrBefore{\parttwo}{ }[\secondword]%
	% \unskip\ignorespaces% Prevent paragraph break
% The format of the responsory
	% \noindent \R \firstline \ \V \versicle \ \parttwo \ Gloria. \firstword. \ \secondword.
	% \noindent \R \firstline \ \V \versicle \ \secondword. Gloria. \secondword.
	% \newline \R \firstline \ \V \versicle \ \secondword. Gloria. \secondword.
	\oldnewline \R \firstline \ \V \versicle \ \secondword. Gloria. \secondword.%
	% \index[R]{\firstline}
}{%
	% \unskip\ignorespaces% Prevent paragraph break
}

% Define the responsory-final environment allows us to change the word after doxology
\NewDocumentEnvironment{responsory-final}{O{} m m m}{%
	\leavevmode% forces horizontal mode
	\unskip% Remove any previous spacing
	\hypertarget{#1}{\label{#1}}
	\def\firstline{#2}%
	\def\versicle{#3}%
	% \StrBefore{\firstline}{*}[\partone]%
	\StrBehind{\firstline}{*}[\rawparttwo]%
	\StrGobbleLeft{\rawparttwo}{1}[\parttwo] % Remove the leading space
	\StrBefore{\firstline}{ }[\firstword]%
	\StrBefore{\parttwo}{ }[\secondword]%
	% \unskip\ignorespaces% Prevent paragraph break
% The format of the responsory
	% \noindent \R \firstline \ \V \versicle \ \parttwo \ Gloria. \firstword. \ \secondword.
	% \noindent \R \firstline \ \V \versicle \ \secondword. Gloria. #4.
	% \newline \R \firstline \ \V \versicle \ \secondword. Gloria. #4.
	\oldnewline \R \firstline \ \V \versicle \ \secondword. Gloria. #4.%
	% \index[R]{\firstline}
}{%
	% \unskip\ignorespaces% Prevent paragraph break
}

% Define the responsory-breve environment
\NewDocumentEnvironment{responsory-breve}{O{} m m}{%
	\leavevmode% forces horizontal mode
	\unskip% Remove any previous spacing
	\hypertarget{#1}{\label{#1}}
	\def\firstline{#2}%
	\def\versicle{#3}%
	% \StrBefore{\firstline}{*}[\partone]%
	\StrBehind{\firstline}{*}[\rawparttwo]%
	\StrGobbleLeft{\rawparttwo}{1}[\parttwo] % Remove the leading space
	\StrBefore{\firstline}{ }[\firstword]%
	\StrBefore{\parttwo}{ }[\secondword]%
	% \unskip\ignorespaces% Prevent paragraph break
% The format of the responsory
	% \noindent \R \firstline \ \V \versicle \ \parttwo \ Gloria. \firstword. \ \secondword.
	% \noindent \R \firstline \ \V \versicle \ \secondword. Gloria. \firstword.
	% \newline \R \firstline \ \V \versicle \ \secondword. Gloria. \firstword.
	\oldnewline \R \firstline \ \V \versicle \ \secondword. Gloria. \firstword.%
	% \index[S]{\firstline}
}{%
	% \unskip\ignorespaces% Prevent paragraph break
}

% Define the invitatory environment
\NewDocumentEnvironment{invitatory}{O{} m}{%
	\leavevmode% forces horizontal mode
	\unskip% Remove any previous spacing
	\hypertarget{#1}{\label{#1}}
	\def\firstline{#2}%
	\StrBehind{\firstline}{*}[\rawparttwo]%
	\StrGobbleLeft{\rawparttwo}{1}[\parttwo]% Remove the leading space
	\StrBefore{\firstline}{ }[\firstword]%
	\StrBefore{\parttwo}{ }[\secondword]%
	% \unskip\ignorespaces% Prevent paragraph break
	\space% Adds space
% The format of the invitatory: firstline or all
	\firstline
	% \index[I]{\firstline}
	% \par \noindent {\color{Red} Invit.} \firstline \\
	% {\color{Red} Invit.} \firstline
	% \lettrine[lines=2]{\zallmancaps \color{Red} V}{eníte}, exsultémus Dómino, iubilémus Deo, salutári nostro: præoccupémus fáciem eius in confessióne, et in psalmis iubilémus ei. \\
	% {\color{Red} Invit.} \firstline \\
	% {\bfseries \color{Blue} Q}uóniam Deus magnus Dóminus, et Rex magnus super omnes deos, quóniam non repéllet Dóminus plebem suam: quia in manu eius sunt omnes fines terræ, et altitúdines móntium ipse cónspicit. \\
	% {\color{Red} Invit.} \parttwo \\
	% {\bfseries \color{Red} Q}uóniam ipsíus est mare, et ipse fecit illud, et áridam fundavérunt manus eius veníte, adorémus, et procidámus ante Deum: plorémus coram Dómino, qui fecit nos, quia ipse est Dóminus, Deus noster; nos autem pópulus eius, et oves páscuæ eius. \\
	% {\color{Red} Invit.} \firstline \\
	% {\bfseries \color{Blue} H}ódie, si vocem eius audiéritis, nolíte obduráre corda vestra, sicut in exacerbatióne secúndum diem tentatiónis in desérto: ubi tentavérunt me patres vestri, probavérunt et vidérunt ópera mea. \\
	% {\color{Red} Invit.} \parttwo \\
	% {\bfseries \color{Red} Q}uadragínta annis próximus fui generatióni huic, et dixi; Semper hi errant corde, ipsi vero non cognovérunt vias meas: quibus iurávi in ira mea; Si introíbunt in réquiem meam. \\
	% {\color{Red} Invit.} \firstline \\
	% \doxology
	% {\color{Red} Invit.} \parttwo \\
	% {\color{Red} Invit.} \firstline \\
}{%
	\unskip\ignorespaces% Prevent paragraph break
}

% Define the invitatory-full environment
\NewDocumentEnvironment{invitatory-full}{O{} m}{%
	\leavevmode% forces horizontal mode
	\unskip% Remove any previous spacing
	\hypertarget{#1}{\label{#1}}
	\def\firstline{#2}%
	\StrBehind{\firstline}{*}[\rawparttwo]%
	\StrGobbleLeft{\rawparttwo}{1}[\parttwo]% Remove the leading space
	\StrBefore{\firstline}{ }[\firstword]%
	\StrBefore{\parttwo}{ }[\secondword]%
	% \unskip\ignorespaces% Prevent paragraph break
% The format of the invitatory: firstline or all
	% \firstline
	\par \noindent {\color{Red} Invitato.} \firstline
	% {\color{Red} Invit.} \firstline
	\par \noindent {\color{Red} Psalmus.}
	\lettrine[lines=2]{\zallmancaps \color{Blue} V}{eníte}, exsultémus Dómino, iubilémus Deo, salutári nostro: præoccupémus fáciem eius in confessióne, et in psalmis iubilémus ei. \\
	\firstword . \\
	{\bfseries \color{Red} Q}uóniam Deus magnus Dóminus, et Rex magnus super omnes deos, quóniam non repéllet Dóminus plebem suam: quia in manu eius sunt omnes fines terræ, et altitúdines móntium ipse cónspicit. \\
	\secondword . \\
	{\bfseries \color{Blue} Q}uóniam ipsíus est mare, et ipse fecit illud, et áridam fundavérunt manus eius veníte, adorémus, et procidámus ante Deum: plorémus coram Dómino, qui fecit nos, quia ipse est Dóminus, Deus noster; nos autem pópulus eius, et oves páscuæ eius. \\
	\firstword . \\
	{\bfseries \color{Red} H}ódie, si vocem eius audiéritis, nolíte obduráre corda vestra, sicut in exacerbatióne secúndum diem tentatiónis in desérto: ubi tentavérunt me patres vestri, probavérunt et vidérunt ópera mea. \\
	\secondword . \\
	{\bfseries \color{Blue} Q}uadragínta annis próximus fui generatióni huic, et dixi; Semper hi errant corde, ipsi vero non cognovérunt vias meas: quibus iurávi in ira mea; Si introíbunt in réquiem meam. \\
	\firstword . \\
	\doxology
	\secondword . \firstword .%
}{%
	\unskip\ignorespaces% Prevent paragraph break
}


% Remove space before chapter and section title: left, before, after
\titlespacing*{\chapter}{0pt}{-45pt}{0pt}
\titlespacing*{\section}{0pt}{-5pt}{0pt}
\titlespacing*{\subsection}{0pt}{-1pt}{0pt}

% Adjust section title spacing
% \titlespacing*{\section}{0pt}{0.25em}{0.25em} % Adjust the spacing before and after section titles

% Define chapter format
\titleformat{\chapter}[display]
{\normalfont\bfseries}{}{0pt}{}

% Define the bookmark text for unnumbered chapters
\newcommand{\ubchapter}[1]{%
	\chapter*{}% Empty chapter
	\addcontentsline{toc}{chapter}{#1}%
	\renewcommand{\leftmark}{#1}%
	\phantomsection%
	\hypertarget{#1}{}%
	\vspace{+\baselineskip}% Adjust the value as needed
}

% Define section format
\titleformat{\section}
% {\normalfont\Large\bfseries}{}{0pt}{}
{\normalfont\Large}{}{0pt}{}

% Define the bookmark text for unnumbered sections
\newcommand{\ubsection}[2]{%
	% \section*{}%
	\renewcommand{\rightmark}{#1}%
	\phantomsection%
	% \hypertarget{#1}{}%
	\addcontentsline{toc}{section}{#1}%
	\hypertarget{#2}{}%
	% \vspace{-\baselineskip}% Adjust the value as needed
}

% Define subsection format
\titleformat{\subsection}
% {\normalfont\large\bfseries}{}{0pt}{}
{\normalfont\centering}{}{0pt}{}

% Define the bookmark text for unnumbered subsections
\newcommand{\ubsubsection}[2]{%
	\subsection*{#1}%
	\phantomsection%
	% \hypertarget{#1}{}%
	\addcontentsline{toc}{subsection}{#1}%
	\hypertarget{#2}{}%
}

% Define the command for centered subsections, in particular the Kalendar
\newcommand{\csection}[2]{%
	\subsection*{\centering\makebox[\linewidth]{\normalfont\Large\bfseries #1}}%
	{\centering\large #2\par}%
	\renewcommand{\rightmark}{#1}%
	\phantomsection%
	% \hypertarget{#1}{}%
%	\refstepcounter{sectionctr}% Increment the section counter
	\addcontentsline{toc}{subsection}{#1}%
%	\hypertarget{\thesectionctr}{}%
	\hypertarget{#1}{}%
}

% Clear default fancyhdr settings
\fancyhead{}

\fancyhead[C]{\leftmark} % Display chapter in the center
\fancyhead[RO,LE]{\rightmark}
\fancyhead[RE,LO]{S. O. P.}

\pagestyle{fancy} % Activate the custom headers and footers

\renewcommand{\contentsname}{Index}

\pretitle{\begin{center}\bfseries\Large Excerpta De\\
\Huge Ecclesiasticum Officium
\end{center}}
\title{\begin{center}\bfseries\LARGE Secundum\\
\Huge Ordinem Fratrum Predicatorum
\end{center}}
\posttitle{\begin{center}\bfseries\Large In Hoc Volumine\\
\huge Per Quatuordecim Libros\\
\large Distinctum Hoc Ordine Continetur
\end{center}}


\author{\ \\ \bfseries\LARGE Humberti de Romanis}
\date{\ \\ \bfseries\Huge 1259}
%1256

% For soul package
\setulcolor{Red}

% Redefine the underline depth to adjust its position
\makeatletter
\def\SOUL@uldepth{.8pt}	% Adjust this value as needed
\makeatother

% For lua-ul underlining
% \newcommand{\ul}[1]{%
    % \underLine[color=Red]{#1}%
% }

% Change behvior of newline
\let\oldnewline\newline
% \renewcommand{\newline}{\linebreak}
\renewcommand{\newline}{\par \noindent}

% Format: \cftsetindents{entry}{indent}{numwidth}
% \cftsetindents{section}{0.5in}{0.5in}
\cftsetindents{section}{0in}{0in}
% \cftsetindents{subsection}{1.0in}{0.5in}
\cftsetindents{subsection}{.2in}{0in}

% Page colors
% \pagecolor[rgb]{0.98,0.94,0.82} % Soft cream
% \pagecolor[rgb]{0.99,0.95,0.84} % Antique cream, whitest
% \pagecolor[rgb]{0.93,0.89,0.75} % Warm old paper
\pagecolor[rgb]{0.90,0.85,0.70} % Aged book, more yellow than 1906
% \pagecolor[rgb]{0.85,0.81,0.69} % 1906

% Too bright
% \pagecolor[rgb]{0.95,0.91,0.80} % 1942
% \pagecolor[rgb]{1.0,0.98,0.90} % Light parchment
% \pagecolor[rgb]{0.96,0.96,0.86} % Soft beige

\begin{document}

% Baseline skip is the distance from one line's baseline to the next. It should typically be 1.2 to 1.3 times the font size
% \fontsize{14pt}{18pt}\selectfont
% \fontsize{12pt}{16pt}\selectfont % Comparable to Psalterium Tridentinum with Charis
% \fontsize{11pt}{14pt}\selectfont
\fontsize{10pt}{12pt}\selectfont %for Merriweather

\frontmatter

\maketitle

\phantomsection
\addcontentsline{toc}{section}{Index}
\tableofcontents
% \thispagestyle{empty}
\let\cleardoublepage\clearpage % Suppresses blank page

% \ubchapter{INCIPIT}

% \thispagestyle{fancy}

% \ubsection{Divine Office Responsories}{divine-office-responsories}

% \subsection*{Divine Office Responsories}

% Here is a sample responsory.

% \begin{responsory}[incipit1]
% {This is the main responsory verse * with a second part.}{This is the versicle.}
% \end{responsory}

% We can refer to the first responsory using its label. For example, see the responsory \hyperlink{incipit1}{\pageref{incipit1}}.


\ubchapter{Foreword}

\thispagestyle{fancy}

\ubsection{Introduction}{intro}



\mainmatter

\input{rubricae_humbert}

%pp013
\ubchapter{Martyrologium}

\thispagestyle{fancy}

\vspace{+.5em}

\ubsection{Anniversaria et Obitus}{anniversaria-et-obitus}

{\color{Red} \ubsubsection{Kalendarium in Quo Solum Anniversaria et Obitus Recitandi Per Annum Continentur}{kalendarium-in-quo-solum-anniversaria-et-obitus}}
% Obitus until XXX Magister inclusive

\begin{center}
\begin{tabular}{l | l | l | r | l l}
& & & & \color{Red} Januarius. \\
 % \color{Red} III & \color{Red} A & \textbf{\color{Red} K\color{Blue}L} & 1 & \color{Red} \\
 % & B & \color{Red} IIII & 2 & Oct. Sancti Stephani. & \color{Red} \\
% \color{Red} XI & C & \color{Red} III & 3 & Oct. Sancti Johannis. & \color{Red} Incipit Quartus Embolismus.\\
% & D & \color{Red} II & 4 & \color{Red} Oct. Sanctorum Innocentium. & \color{Red} \\
% \color{Red} XIX & E & \color{Red} Non. & 5 & & \color{Red} \\
\color{Red} VIII & F & \color{Red} VIII & 6 & Obiit Frater Raymundus Ordinis \\
& & & & \qquad Nostri Magister Tercius. \\
 % & G & \color{Red} VII & 7 & & \color{Red} Claves. LXX. \\
% \color{Red} XVI & \color{Red} A & \color{Red} VI & 8 & & \\
% \color{Red} V & B & \color{Red} V & 9 & & \\
 & C & \color{Red} IIII & 10 & Ob. F. Barnabas Mag. Ord. XV. \\
% \color{Red} XIII & D & \color{Red} III & 11 & & \color{Red} \\
% \color{Red} II & E & \color{Red} II. & 12 & & \color{Red} \\
 % & F & \color{Red} Idus & 13 & \multicolumn{2}{l}{Oct. Epiphanie. \color{Red} Simplex. \color{black} Hylarii et Remigii. \color{Red} Mem.} \\
% \color{Red} X & G & \color{Red} XIX & 14 & \multicolumn{2}{l}{Felicis Presbiteri et Confessoris. III. Lect.} \\
 % & \color{Red} A & \color{Red} XVIII & 15 & \color{Red} Mauri Abbatis. III. Lect. \\
% \color{Red} XVIII & B & \color{Red} XVII & 16 & Marcelli Pape et Martyris. III. Lect. \\
% \color{Red} VII & C & \color{Red} XVI & 17 & \color{Red} Anthonii Abbatis. III. Lect. \\
 % & D & \color{Red} XV & 18 & Prisce Virginis et Martyris. III. Lect. & \color{Red} Sol in Aquario. Initium LXX. \\
% \color{Red} XV & E & \color{Red} XIIII & 19 & \color{Red} \\
% \color{Red} IIII & F & \color{Red} XIII & 20 & \multicolumn{2}{l}{Fabiani et Sebastiani Martyrum. Simplex.} \\
 % & G & \color{Red} XII & 21 & Agnetis Virginis et Martyris. Simplex. \\
% \color{Red} XII & \color{Red} A & \color{Red} XI & 22 & Vincentii Martyris. Semid. \\
% \color{Red} I & B & \color{Red} X & 23 & \color{Red} Emerentiane Virginis et Martyris. Mem. \\
 % & C & \color{Red} IX & 24 & & \color{Red} \\
% \color{Red} IX & D & \color{Red} VIII & 25 & Conversio Sancti Pauli. Semid. \\
% & E & \color{Red} VII & 26 & & \color{Red} \\
% \color{Red} XVII & F & \color{Red} VI & 27 & \color{Red} Juliani Episcopi et Confessoris. Mem. \\
% \color{Red} VI & G & \color{Red} V & 28 & \color{Red} Agnetis Secundo. III. Lect. & Claves XL. \\
 % & \color{Red} A & \color{Red} IIII & 29 & & \color{Red} \\
% \color{Red} XIIII & B & \color{Red} III & 30 & & \color{Red} \\
% \color{Red} III & C & \color{Red} II & 31 & & \color{Red} \\


& & & & \color{Red} Februarius. \\
 % & D & \textbf{\color{Red} K\color{Blue}L} & 1 & \color{Red} Ignatii Episcopi et Martyris. Mem. & Finis. IV. embolismi. \\
% \color{Red} XI & E & \color{Red} IIII & 2 & \multicolumn{2}{l}{Purificatio Sancte Marie Virginis. Totum Dup.} \\
% \color{Red} XIX & F & \color{Red} III & 3 & \color{Red} Blasii Episcopi et Martyris. III. Lect. \\
\color{Red} VIII & G & \color{Red} II & 4 & Anniversarium Patrum et Matrum. \\
 % & \color{Red} A & \color{Red} Non. & 5 & \color{Red} Agathe Virginis et Martyris. Simplex. \\
% \color{Red} XVI & B & \color{Red} VIII & 6 &  Vedasti et Amandi Episcoporum. Memor. \\
% \color{Red} V & C & \color{Red} VII & 7 & \\
% \color{Red}  & D & \color{Red} VI & 8 & & \color{Red} Initium. XL. \\
% \color{Red} XIII & E & \color{Red} V & 9 & & \color{Red} \\
% \color{Red} II & F & \color{Red} IIII & 10 & Scolastice Virginis. Mem. & \color{Red} \\
% \color{Red}  & G & \color{Red} III & 11 & & \color{Red} \\
\color{Red} X & \color{Red} A & \color{Red} II & 12 & Ob. F. Iordanis, Ord. Nos. Mag. II. \\
% \color{Red}  & B & \color{Red} Idus & 13 & & \color{Red} \\
% \color{Red} XVIII & C & \color{Red} XVI & 14 & \color{Red} Valentini Martyris. III. Lect. \\
% \color{Red} VII & D & \color{Red} XV & 15 & & \color{Red} Sol in Piscibus. \\
% \color{Red}  & E & \color{Red} XIIII & 16 & & \color{Red} \\
% \color{Red} XV & F & \color{Red} XIII & 17 & & \color{Red} \\
% \color{Red} IIII & G & \color{Red} XII & 18 & & \color{Red} \\
% \color{Red}  & \color{Red} A & \color{Red} XI & 19 & & \color{Red} \\
% \color{Red} XII & B & \color{Red} X & 20 & & \color{Red} \\
% \color{Red} I & C & \color{Red} IX & 21 & & \color{Red} \\
% \color{Red}  & D & \color{Red} VIII & 22 & \color{Red} Cathedra Sancti Petri. Simplex. \\
% \color{Red} IX & E & \color{Red} VII & 23 & Ob. F. Iohannes Mag. Ord. XX. & \color{Red} \\
% \color{Red}  & F & \color{Red} VI & 24 & Mathie Apostoli. Semid. & \color{Red} Locus Byssexti. \\
% \color{Red} XVII & G & \color{Red} V & 25 & & \color{Red} \\
% \color{Red} VI & \color{Red} A & \color{Red} IIII & 26 & & \color{Red} \\
% \color{Red}  & B & \color{Red} III & 27 & & \color{Red} \\
% \color{Red} XVI & C & \color{Red} II & 28 & & \color{Red} \\


& & & & \color{Red} Martius. \\
\color{Red} III & D & \textbf{\color{Red} K\color{Blue}L} & 1 & Ob. F. Petrus Mag. Ord. XVIII. \\
% \color{Red}  & E & \color{Red} VI & 2 & & \color{Red} \\
% \color{Red} XI & F & \color{Red} V & 3 & \\
% \color{Red}  & G & \color{Red} IIII & 4 & & \color{Red} \\
% \color{Red} XIX & \color{Red} A & \color{Red} III & 5 & & \color{Red} Incipit. VII. Embolismus. \\
% \color{Red} VIII & B & \color{Red} II & 6 & & \color{Red} Incipit. III. Embolismus. \\
% \color{Red}  & C & \color{Red} Non. & 7 & & \color{Red} \\
% \color{Red} XVI & D & \color{Red} VIII & 8 & & \color{Red} \\
% \color{Red} V & E & \color{Red} VII & 9 & & \color{Red} \\
% \color{Red}  & F & \color{Red} VI & 10 & & \color{Red} \\
% \color{Red} XIII & G & \color{Red} V & 11 & & \color{Red} Claves Pasche. \\
% \color{Red} II & \color{Red} A & \color{Red} IV. & 12 & \color{Red} Gregorii Pape et Conf. Simplex. \\
% \color{Red}  & B & \color{Red} III & 13 & & \color{Red} \\
% \color{Red} X & C & \color{Red} II & 14 & & \color{Red} Ultima XXL. \\
% \color{Red}  & D & \color{Red} Idus & 15 & & \color{Red} Equinoctium Vernale. \\
% \color{Red} XVIII & E & \color{Red} XVII & 16 & Ob. F. Leonardus Mag. Ord. XXV. \\
% \color{Red} VII & F & \color{Red} XVI & 17 & & \color{Red} \\
% \color{Red}  & G & \color{Red} XV & 18 & & \color{Red} Sol in Ariete. \\
% \color{Red} XV & \color{Red} A & \color{Red} XIIII & 19 & & \color{Red} \\
% \color{Red} IIII & B & \color{Red} XIII & 20 & & \color{Red} \\
% \color{Red}  & C & \color{Red} XII & 21 & \color{Red} Benedicti Abbatis. Simplex. \\
% \color{Red} XII & D & \color{Red} XI & 22 & & \color{Red} Primum Pascha. \\
% \color{Red} I & E & \color{Red} X & 23 & & \color{Red} \\
% \color{Red}  & F & \color{Red} IX & 24 & & \color{Red} \\
% \color{Red} IX & G & \color{Red} VIII & 25 & Annunciatio Dominica. Totum Dupl.& \color{Red} \\
\color{Red}  & \color{Red} A & \color{Red} VII & 26 & Ob. F. Munio Mag. Ord. Septimus. & \color{Red} \\
% \color{Red} XVII & B & \color{Red} VI & 27 & & \color{Red} \\
% \color{Red} VI & C & \color{Red} V & 28 & & \color{Red} \\
% \color{Red}  & D & \color{Red} IIII & 29 & \color{Red} \\
% \color{Red} XIIII & E & \color{Red} III & 30 & & \color{Red} \\
% \color{Red} III & F & \color{Red} II & 31 & & \color{Red} \\


& & & & \color{Red} Aprilis. \\
% \color{Red}  & G & \textbf{\color{Red} K\color{Blue}L} & 1 & & \color{Red} \\
% \color{Red} XI & \color{Red} A & \color{Red} IIII & 2 & & \color{Red} \\
% \color{Red}  & B & \color{Red} III & 3 & & Finis. VII. Embolismus. \color{Red} \\
% \color{Red} XIX & C & \color{Red} II & 4 & Ambrosii Episcopi et Conf. Simplex. & \color{Red} Finis. VI. Embolismus. \\
% \color{Red} VIII & D & \color{Red} Non. & 5 & & \color{Red} \\
% \color{Red} XVI & E & \color{Red} VIII & 6 & & \color{Red} \\
% \color{Red} V & F & \color{Red} VII & 7 & & \color{Red} \\
% \color{Red}  & G & \color{Red} VI & 8 & Ob. F. Guarinus Mag. Ord. XIX. & \color{Red} \\
% \color{Red} XIII & \color{Red} A & \color{Red} V & 9 & & \color{Red} \\
% \color{Red} II & B & \color{Red} IIII & 10 & & \color{Red} \\
% \color{Red}  & C & \color{Red} III & 11 & & \color{Red} \\
% \color{Red} X & D & \color{Red} II & 12 & & \color{Red} \\
% \color{Red}  & E & \color{Red} Idus & 13 & & \color{Red} \\
% \color{Red} XVIII & F & \color{Red} XVIII & 14 & \multicolumn{2}{l}{Tyburtii Valeriani, et Maximi \color{Red} Martyrum. \color{black} III. Lect.} \color{Red} \\
% \color{Red} VII & G & \color{Red} XVII & 15 & & \color{Red} Claves Rogationum. \\
% \color{Red}  & \color{Red} A & \color{Red} XVI & 16 & \\
% \color{Red} XV & B & \color{Red} XV & 17 & & Sol in Tauro. \color{Red} \\
% \color{Red} IIII & C & \color{Red} XIIII & 18 & & \color{Red} \\
% \color{Red}  & D & \color{Red} XIII & 19 & & \color{Red} \\
% \color{Red} XII & E & \color{Red} XII & 20 & & \color{Red} \\
% \color{Red} I & F & \color{Red} XI & 21 & & \color{Red} \\
% \color{Red}  & G & \color{Red} X & 22 & & \color{Red} \\
% \color{Red} IX & \color{Red} A & \color{Red} IX & 23 & Georgii Martyris. Simplex. & \color{Red} \\
% \color{Red}  & B & \color{Red} VIII & 24 & & \color{Red} \\
% \color{Red} XVII & C & \color{Red} VII & 25 & \color{Red} Marci Evangeliste. Semid. & Ultimum Pascha. \\
% \color{Red} VI & D & \color{Red} VI & 26 & & \color{Red} \\
% \color{Red}  & E & \color{Red} V & 27 & & \color{Red} \\
% \color{Red} XIIII & F & \color{Red} IIII & 28 & Vitalis Martyris. III. Lect. & \color{Red} \\
% \color{Red} III & G & \color{Red} III & 29 & \multicolumn{2}{l}{Beati Petri Martyris de Ordines Predic. Totum Dup. \color{Red} Claves Pente.} \\
% \color{Red}  & \color{Red} A & \color{Red} II & 30 & & \color{Red} \\


& & & & \color{Red} Maius. \\
% \color{Red} XI & B & \textbf{\color{Red} K\color{Blue}L} & 1 & Philippi et Iacobi Apostolorum. Semid. \\
% \color{Red}  & C & \color{Red} VI & 2 & \\
% \color{Red} XIX & D & \color{Red} V & 3 & Inventio Sancte Crucis. Semid. \\
 % &  &  &  & \multicolumn{2}{l}{\qquad Alexandri Evencii, et Theodoli. \color{Red} Martyrum. \color{black} Mem.} \\
% \color{Red} VIII & E & \color{Red} IIII & 4 & Festum Corone Domini. Simplex. & \color{Red} \\
% \color{Red}  & F & \color{Red} III & 5 & & \color{Red} \\
% \color{Red} XVI & G & \color{Red} II & 6 & Ob. F. Petrus Mag. Ord. XXVII. \\
% \color{Red} V & \color{Red} A & \color{Red} Non. & 7 & & \color{Red} \\
% \color{Red}  & B & \color{Red} VIII & 8 & & \color{Red} \\
% \color{Red} XIII & C & \color{Red} VII & 9 & & \color{Red} \\
% \color{Red} II & D & \color{Red} VI & 10 & \multicolumn{2}{l}{Gordiani et Epimachi Martyrum. III. Lect.} \color{Red} \\
% \color{Red}  & E & \color{Red} V & 11 & & \color{Red} \\
% \color{Red} X & F & \color{Red} IIII & 12 & \multicolumn{2}{l}{Neri, Achillei, et Pancratii Martyrum. III. Lect.} \color{Red} \\
% \color{Red}  & G & \color{Red} III & 13 & & \color{Red} \\
% \color{Red} XVIII & \color{Red} A & \color{Red} II & 14 & & \color{Red} \\
% \color{Red} VII & B & \color{Red} Idus & 15 & & \color{Red} \\
% \color{Red}  & C & \color{Red} XVII & 16 &  & \color{Red} \\
% \color{Red} XV & D & \color{Red} XVI & 17 & & \color{Red} \\
% \color{Red} IIII & E & \color{Red} XV & 18 & & Sol in Geminis. \color{Red} \\
% \color{Red}  & F & \color{Red} XIIII & 19 & Potentiane Virginis. Mem. & \color{Red} \\
% \color{Red} XII & G & \color{Red} XIII & 20 & & \color{Red} \\
% \color{Red} I & \color{Red} A & \color{Red} XII & 21 & & \color{Red} \\
% \color{Red}  & B & \color{Red} XI & 22 & & \color{Red} \\
% \color{Red} IX & C & \color{Red} X & 23 & Translatio Beati Dominici. Totum Dup. & \color{Red} \\
% \color{Red}  & D & \color{Red} IX & 24 & & \color{Red} \\
% \color{Red} XVII & E & \color{Red} VIII & 25 & Urbani Pape et Martyris. III. Lect. & \color{Red} \\
% \color{Red} VI & F & \color{Red} VII & 26 & & \color{Red} \\
\color{Red}  & G & \color{Red} VI & 27 & Ob. F. Bernardus Mag. Ord. XI. & \color{Red} \\
% \color{Red} XIIII & \color{Red} A & \color{Red} V & 28 & & \color{Red} \\
% \color{Red} III & B & \color{Red} IIII & 29 & & \color{Red} \\
% \color{Red}  & C & \color{Red} III & 30 & & \color{Red} \\
% \color{Red} XI & D & \color{Red} II & 31 & Petronille Virginis. Mem. & \color{Red} \\


& & & & \color{Red} Iunius. \\
% \color{Red}  & E & \textbf{\color{Red} K\color{Blue}L} & 1 & & \color{Red} \\
% \color{Red} XIX & F & \color{Red} IIII & 2 & Marcellini et Petri Martyrum. III. Lect. & \color{Red} \\
% \color{Red} VIII & G & \color{Red} III & 3 & & \color{Red} \\
% \color{Red} XVI & \color{Red} A & \color{Red} II & 4 & & \color{Red} \\
% \color{Red} V & B & \color{Red} Non. & 5 & & \color{Red} \\
% \color{Red}  & C & \color{Red} VIII & 6 & & \color{Red} \\
% \color{Red} XIII & D & \color{Red} VII & 7 & Ob. F. Simonis Mag. Ord. XXI. & \color{Red} \\
% \color{Red} II & E & \color{Red} VI & 8 & Medardi Episcopi et Conf. Mem. & \color{Red} \\
% \color{Red}  & F & \color{Red} V & 9 & Primi et Feliciani Martyrum. III. Lect. & \color{Red} \\
% \color{Red} X & G & \color{Red} IIII & 10 & & \color{Red} \\
% \color{Red}  & \color{Red} A & \color{Red} III & 11 & \color{Red} Barnabe Apostoli. Semid. \\
% \color{Red} XVIII & B & \color{Red} II & 12 & \multicolumn{2}{l}{\color{Red} Basilidis Cyrini, Naboris, et Nazarii, Martyrum. III. Lect.} \\
% \color{Red} VII & C & \color{Red} Idus & 13 & & \color{Red} \\
% \color{Red}  & D & \color{Red} XVIII & 14 & & Solsticium Estivale. \color{Red} \\
% \color{Red} XV & E & \color{Red} XVII & 15 & \color{Red} Viti et Modest Martyrum. Mem. \\
% \color{Red} IIII & F & \color{Red} XVI & 16 & \color{Red} Cyrici et Iulice Martyrum. Mem. \\
% \color{Red}  & G & \color{Red} XV & 17 & & Sol in Cancio. \color{Red} \\
% \color{Red} XII & \color{Red} A & \color{Red} XIIII & 18 & \color{Red} Marci et Marcelliani Martyrum. III. Lect. \\
% \color{Red} I & B & \color{Red} XIII & 19 & \color{Red} Gervasii et Prothasii in Martyrum. Simplex. \\
% \color{Red}  & C & \color{Red} XII & 20 & Ob. F. Helias Mag. Ord. XXII. & \color{Red} \\% Or Dec 31?
% \color{Red} IX & D & \color{Red} XI & 21 & & \color{Red} \\
% \color{Red}  & E & \color{Red} X & 22 & & \color{Red} \\
% \color{Red} XVII & F & \color{Red} IX & 23 & & Vigilia. \color{Red} \\
% \color{Red} VI & G & \color{Red} VIII & 24 & Nativitas Sancti Iohannis Baptiste. Dupl. \\
% \color{Red}  & \color{Red} A & \color{Red} VII & 25 & & \color{Red} \\
% \color{Red} XIIII & B & \color{Red} VI & 26 & Iohannis et Pauli Martyrum. Simplex. & \color{Red} \\
% \color{Red} III & C & \color{Red} V & 27 & & \color{Red} \\
% \color{Red}  & D & \color{Red} IIII & 28 & Leonis Pape et Conf. Memoria. & Vigilia. \color{Red} \\
% \color{Red} XI & E & \color{Red} III & 29 & Apostolorum Petri et Pauli. Duplex. & \color{Red} \\
% \color{Red}  & F & \color{Red} II & 30 & Commemoratio Sancti Pauli. \color{Red} Semid. \\


%pp014
& & & & \color{Red} Iulius. \\
% \color{Red} XIX & G & \textbf{\color{Red} K\color{Blue}L} & 1 & Octava Sancti Iohannis Baptiste. Simplex. & \color{Red} \\
% \color{Red} VIII & \color{Red} A & \color{Red} VI & 2 & Processi et Martiniani Martyrum. Mem. & \color{Red} \\
% \color{Red}  & B & \color{Red} V & 3 & & \color{Red} \\
% \color{Red} XVI & C & \color{Red} IIII & 4 & & \color{Red} \\
\color{Red} V & D & \color{Red} III & 5 & Ob. Benedictus Papa \\
& & & & \qquad Prius F. Nycholaus Mag. Ord. IX. & \color{Red} \\
% \color{Red}  & E & \color{Red} II & 6 & \color{Red} Octava Apostolorum Petri et Pauli. Simplex. \\
\color{Red} XIII & F & \color{Red} Non. & 7 & Anniversarium Universorum \\
& & & & \qquad Sepultorum in Cimiteriis Nos. & \color{Red} \\
% \color{Red} II & G & \color{Red} VIII & 8 & & \color{Red} \\
% \color{Red}  & \color{Red} A & \color{Red} VII & 9 & & \color{Red} \\
% \color{Red} X & B & \color{Red} VI & 10 & Septem Fratrum. III. Lect. & \color{Red} \\
% \color{Red}  & C & \color{Red} V & 11 & & \color{Red} \\
% \color{Red} XVIII & D & \color{Red} IIII & 12 &  & \color{Red} \\
% \color{Red} VII & E & \color{Red} III & 13 & & \color{Red} \\
\color{Red}  & F & \color{Red} II & 14 & Ob. F. Humbertus Ord. Nos. Mag. V. & \color{Red} \\
% \color{Red} XV & G & \color{Red} Idus & 15 & & \color{Red} \\
% \color{Red} IIII & \color{Red} A & \color{Red} XVII & 16 & \\
% \color{Red}  & B & \color{Red} XVI & 17 & & \color{Red} \\
% \color{Red} XII & C & \color{Red} XV & 18 & & \color{Red} Sol in Leone. \\
% \color{Red} I & D & \color{Red} XIIII & 19 & & \color{Red} \\
% \color{Red}  & E & \color{Red} XIII & 20 & Margarete Virginis et Martyris. Simplex. & \color{Red} \\
% \color{Red} IX & F & \color{Red} XII & 21 & Praxedis Virginis. III. Lect. & \color{Red} \\
% \color{Red}  & G & \color{Red} XI & 22 & \color{Red} Marie Magdalene. Semid. \\
% \color{Red} XVII & \color{Red} A & \color{Red} X & 23 & \color{Red} Apollinaris Episcopi et Martyris. III. Lect. \\
% \color{Red} VI & B & \color{Red} IX & 24 & Christine Virginis et Martyris. Mem. \\
% \color{Red}  & C & \color{Red} VIII & 25 & \multicolumn{2}{l}{Iacobi Apostoli. Semi. Christofori et Cucufatis Martyrum. Mem.} \\
% \color{Red} XIIII & D & \color{Red} VII & 26 & & \color{Red} \\
% \color{Red} III & E & \color{Red} VI & 27 & & \color{Red} \\
% \color{Red}  & F & \color{Red} V & 28 & \multicolumn{2}{l}{Nazarii, Celsi, et Pantaleonis. Martyrum. III. Lect.} \color{Red} \\
% \color{Red} XI & G & \color{Red} IIII & 29 & \multicolumn{2}{l}{\color{Red} Felicis, Simplicii, Faustini, et Beatricis Martyrum. III. Lect.} \\
% \color{Red} XIX & \color{Red} A & \color{Red} III & 30 & Abdon et Sennen Martyrum. III. Lect. & \color{Red} \\
% \color{Red}  & B & \color{Red} II & 31 & Germani Episcopi et Conf. III. Lect. & \color{Red} \\


& & & & \color{Red} Augustus. \\
% \color{Red} VIII & C & \textbf{\color{Red} K\color{Blue}L} & 1 & Ad Vincula Sancti Petri. Simplex. & \color{Red} \\
 % &  & &  & \multicolumn{2}{l}{\qquad Sanctorum Machabeorum Martyrum. Mem. \color{Red}} \\
% \color{Red} XVI & D & \color{Red} IIII & 2 & \color{Red} Stephani Pape et Martyris. III. Lect. \\
% \color{Red} V & E & \color{Red} III & 3 & \color{Red} Inventio Sancti Stephani. Simplex. \\
% \color{Red}  & F & \color{Red} II & 4 & & \color{Red} \\
\color{Red} XIII & G & \color{Red} Non. & 5 & Transiit B. Dominicus Mag. Ord. I. & \color{Red} \\
\color{Red} II & \color{Red} A & \color{Red} VIII & 6 & Ob. F. Hugo Doctor in Theologia \\
& & & & \qquad Magister Ordinis. XVI. & \color{Red} \\
% \color{Red}  & B & \color{Red} VII & 7 & Ob. F. Thomas Mag. Ord. XXIV. & \color{Red} \\
\color{Red} X & C & \color{Red} VI & 8 & Ob. F. Herveus Doctor in Theologia \\
& & & & \qquad Magister Ordinis. XIIII. & \color{Red} \\
% \color{Red}  & D & \color{Red} V & 9 & & \color{Red} Vigilia. \\
% \color{Red} XVIII & E & \color{Red} IIII & 10 & \color{Red} Laurentii Martyris. Semid. & \color{Red} \\
% \color{Red} VII & F & \color{Red} III & 11 & Tiburtii Martyris. Mem. & \color{Red} \\
\color{Red}  & G & \color{Red} II & 12 & Ob. F. Hemericus Mag. Ord. XII. & \color{Red} \\
% \color{Red} XV & \color{Red} A & \color{Red} Idus & 13 & Ypoliti Sociorumque eius . Simplex. & \color{Red} \\
% \color{Red} IIII & B & \color{Red} XIX & 14 & Eusebii Presbiteri et Conf. Mem. & \color{Red} Vigilia. \\
% \color{Red}  & C & \color{Red} XVIII & 15 & \multicolumn{2}{l}{\color{Red} Assumptio Beate Marie Virginis. Totum Dupl.} \color{Red} \\
% \color{Red} XII & D & \color{Red} XVII & 16 & & \color{Red} \\
% \color{Red} I & E & \color{Red} XVI & 17 & Octava Sancti Laurentii. Simplex. & \color{Red} \\
% \color{Red}  & F & \color{Red} XV & 18 & Agapiti Martyris. Mem. & \color{Red} Sol in Virgine. \\
% \color{Red} IX & G & \color{Red} XIIII & 19 & & \color{Red} \\
% \color{Red}  & \color{Red} A & \color{Red} XIII & 20 & Bernardi Abbatis. Simplex. & \color{Red} \\
% \color{Red} XVII & B & \color{Red} XII & 21 & & \color{Red} \\
% \color{Red} VI & C & \color{Red} XI & 22 & \color{Red} Octava Sancte Marie. Simplex. \\
% & & & & \multicolumn{2}{l}{\qquad \color{Red} Tymothei et Symphoriani Martyrum. Mem.} \\
% \color{Red}  & D & \color{Red} X & 23 & & \color{Red} \\
% \color{Red} XIIII & E & \color{Red} IX & 24 & Bartholomei Apostoli. \color{Red} Semid. & \color{Red} \\
% \color{Red} III & F & \color{Red} VIII & 25 &  & \color{Red} \\
% \color{Red}  & G & \color{Red} VII & 26 &  & \color{Red} \\
\color{Red} XI & \color{Red} A & \color{Red} VI & 27 & Ob. F. Albertus Mag. Ord. X. \\
% \color{Red} XIX & B & \color{Red} V & 28 & \color{Red} Augustini Episcopi et Conf. Totum Dup. & \color{Red} \\
% \color{Red}  & C & \color{Red} IIII & 29 & \multicolumn{2}{l}{Decollatio Sancti Iohannis Baptiste. Simplex. Sabine Martyris. Mem.} \\
% \color{Red} VIII & D & \color{Red} III & 30 & \color{Red} Felicis, et Adaucti Martyrum. Mem. \\
% \color{Red}  & E & \color{Red} II & 31 & & \color{Red} Finis Quinti Embolismus. \\


& & & & \color{Red} Septembris. \\
% \color{Red} XVI & F & \textbf{\color{Red} K\color{Blue}L} & 1 & Egidii Abbatis. Mem. & \color{Red} \\
% \color{Red} V & G & \color{Red} IIII & 2 & & \color{Red} Incipit. II. Embolismus. \\
% \color{Red}  & \color{Red} A & \color{Red} III & 3 & & \color{Red} \\
% \color{Red} XIII & B & \color{Red} II & 4 & Octava Sancti Augustini. Simplex. & \color{Red} \\
% & & & & \qquad Marcelli Martyris. Mem. \\
\color{Red} II & C & \color{Red} Non. & 5 & Anni. Familiarium et Benefactorum. \\
% \color{Red}  & D & \color{Red} VIII & 6 & & \color{Red} \\
% \color{Red} X & E & \color{Red} VII & 7 & & \color{Red} \\
% \color{Red}  & F & \color{Red} VI & 8 & \multicolumn{2}{l}{\color{Red} Nativitas Sancte Marie Virginis. Totum Dup.} \\
% \color{Red} XVIII & G & \color{Red} V & 9 & Gorgonii Martyris. \color{Red} Mem. \\
% \color{Red} VII & \color{Red} A & \color{Red} IIII & 10 & & \color{Red} \\
% \color{Red}  & B & \color{Red} III & 11 & \color{Red} Prothi et Iacinti Martyrum. Mem. \\
% \color{Red} XV & C & \color{Red} II & 12 & & \color{Red} \\
% \color{Red} IIII & D & \color{Red} Idus & 13 & & \color{Red} \\
% \color{Red}  & E & \color{Red} XVIII & 14 & Exaltatio Sancte Crucis. Semid. & Equinoctium Autumpnale. \color{Red} \\
% & & & & \multicolumn{2}{l}{\qquad Cornelii, et Cipriani Martyrum. Mem.} \\
% \color{Red} XII & F & \color{Red} XVII & 15 & Octava Sancte Marie. Simplex. & \color{Red} \\
% & & & & \qquad Nichomedis Martyris. Mem. \\
% \color{Red} I & G & \color{Red} XVI & 16 & Eufemie Virginis et Martyris. III. Lect. & \color{Red} \\
% \color{Red}  & \color{Red} A & \color{Red} XV & 17 & Lamberti Episcopi et Martyris. Mem. & \color{Red} Sol in Libra. \\
\color{Red} IX & B & \color{Red} XIIII & 18 & Ob. Dom. Berangarius Mag. Ord. XIII. & \color{Red} \\
% There seems to be disagreement on the date of his death. There are three dates, Oct 20, Sep 18 and Sep 20. Cormier 1909 lists Sep 1 as compromise. Deutsch Wikipedia even notes Sep 30, but gives a footnote:
% Mit dem 18. bzw. 20. September 1330 geben mirabileweb bzw. Jan Prelog im Lexikon des Mittelalters abweichende Todesdaten an.
% There is good reason to believe Sep 18 is the correct date of death. If we take XIIII = XIIII and misread it as XX Kal Oct, then we easily understand the descrepancy, and 30 Sep is clearly a misreading of 20.
% \color{Red}  & C & \color{Red} XIII & 19 & & \color{Red} \\
% \color{Red} XVII & D & \color{Red} XII & 20 & \color{Red} & \color{Red} Vigilia. \\
% \color{Red} VI & E & \color{Red} XI & 21 & Ob. F. Martialus Mag. Ord. XXIX. \color{Red} \\
% Or Sep 11?
% \color{Red}  & F & \color{Red} X & 22 & \multicolumn{2}{l}{Mauritii Sociorumque eius Martyrum. Simplex.} \color{Red} \\
% \color{Red} XIIII & G & \color{Red} IX & 23 &  & \color{Red} \\
% \color{Red} III & \color{Red} A & \color{Red} VIII & 24 & & \color{Red} \\
% \color{Red}  & B & \color{Red} VII & 25 & & \color{Red} \\
% \color{Red} XI & C & \color{Red} VI & 26 & & \color{Red} \\
% \color{Red} XIX & D & \color{Red} V & 27 & \color{Red} Cosme et Damiani Martyrum. Simplex. \\
\color{Red}  & E & \color{Red} IIII & 28 & Ob. Dom. Geraldus Mag. Ord. XVII. & \color{Red} \\
% \color{Red} VIII & F & \color{Red} III & 29 & \color{Red} Michaelis Archangeli. Dup. & \color{Red} \\
% \color{Red}  & G & \color{Red} II & 30 & Iheronimi Presbiteri et Conf. Simplex. & \color{Red} Finis Secundi Embolismus. \\


& & & & \color{Red} Octobris. \\
% \color{Red} XVI & \color{Red} A & \textbf{\color{Red} K\color{Blue}L} & 1 & \color{Red} Remigii Episcopi et Conf. III. Lect. \\
% \color{Red} V & B & \color{Red} VI & 2 & \color{Red} Leodegarii Episcopi et Martyris. Mem. \\
% \color{Red} XIII & C & \color{Red} V & 3 & & \color{Red} \\
% \color{Red} II & D & \color{Red} IIII & 4 & Francisci Confessoris. Simplex. & \color{Red} \\
% \color{Red}  & E & \color{Red} III & 5 & Ob. F. Raymundus Mag. Ord. XXIII. & \color{Red} \\
% \color{Red} X & F & \color{Red} II & 6 & & \color{Red} \\
% \color{Red}  & G & \color{Red} Non. & 7 & Marci Pape et Conf. III. Lect. \\
 % &  & & & \multicolumn{2}{l}{\qquad Sergii, et Bachi, Marcelli, et Apulei. Martyrum. Mem.} \\
% \color{Red} XVIII & \color{Red} A & \color{Red} VIII & 8 & & \color{Red} \\
% \color{Red} VII & B & \color{Red} VII & 9 & \multicolumn{2}{l}{\color{Red} Dionisii Sociorumque eius Martyrum. Simplex.} \\
\color{Red}  & C & \color{Red} VI & 10 & Anniversarium Omnium Fratrum. \\
% \color{Red} XV & D & \color{Red} V & 11 & & \color{Red} \\
% \color{Red} IIII & E & \color{Red} IIII & 12 & & \color{Red} \\
% \color{Red}  & F & \color{Red} III & 13 &  & \color{Red} \\
% \color{Red} XII & G & \color{Red} II & 14 & Calyxti Pape et Martyris. Mem. & \color{Red} \\
% \color{Red} I & \color{Red} A & \color{Red} Idus & 15 & Ob. F. Bartholomeus Mag. Ord. XXVI. & \color{Red} \\
% \color{Red}  & B & \color{Red} XVII & 16 &  & \color{Red} \\
% \color{Red} IX & C & \color{Red} XVI & 17 &  & \color{Red} \\
% \color{Red}  & D & \color{Red} XV & 18 & Ob. F. Guidonus Mag. Ord. XXVIII. & \color{Red} Sol in Scorpione. \\
% \color{Red} XVII & E & \color{Red} XIIII & 19 & & \color{Red} \\
% \color{Red} VI & F & \color{Red} XIII & 20 & & \color{Red} \\
% \color{Red}  & G & \color{Red} XII & 21 & \multicolumn{2}{l}{\color{Red} Undecim Milium Virginum et Martyrum. Mem.} \\
% \color{Red} XIIII & \color{Red} A & \color{Red} XI & 22 & & \color{Red} \\
% \color{Red} III & B & \color{Red} X & 23 &  & \color{Red} \\
% \color{Red}  & C & \color{Red} IX & 24 & & \color{Red} \\
% \color{Red} XI & D & \color{Red} VIII & 25 & Crispini et Crispiani Martyrum. Mem. & \color{Red} \\
% \color{Red} XIX & E & \color{Red} VII & 26 & & \color{Red} \\
% \color{Red}  & F & \color{Red} VI & 27 &  & \color{Red} Vigilia. \\
% \color{Red} VIII & G & \color{Red} V & 28 & \color{Red} Symonis et Iude Apostolorum. Semid. & \color{Red} \\
% \color{Red}  & \color{Red} A & \color{Red} IIII & 29 & & \color{Red} \\
% \color{Red} XII & B & \color{Red} III & 30 &  & \color{Red} \\
% \color{Red} V & C & \color{Red} II & 31 & Quintini Martyris. Mem. & Vigilia. \color{Red} \\


& & & & \color{Red} Novembris. \\
% \color{Red}  & D & \textbf{\color{Red} K\color{Blue}L} & 1 & Festivitas Omnium Sanctorum. Totum Dup. & \color{Red} \\
% \color{Red} XIII & E & \color{Red} IIII & 2 & Commemoratio Omnium Fidelium Defunctorum. & \color{Red} Incipit. V. Embolismus. \\
% \color{Red} II & F & \color{Red} III & 3 & & \color{Red} \\
\color{Red}  & G & \color{Red} II & 4 & Ob. F. Iohannes Ord. Nos. Mag. IV. & \color{Red} \\
% \color{Red} X & \color{Red} A & \color{Red} Non. & 5 & & \color{Red} \\
% \color{Red}  & B & \color{Red} VIII & 6 & & \color{Red} \\
% \color{Red} XVIII & C & \color{Red} VII & 7 & & \color{Red} \\
% \color{Red} VII & D & \color{Red} VI & 8 & Quatuor Coronatorum Martyrum. III. Lect. & \color{Red} \\
% \color{Red}  & E & \color{Red} V & 9 & Theodori Martyris. III. Lect. & \color{Red} \\
% \color{Red} XV & F & \color{Red} IIII & 10 & & \color{Red} \\
% \color{Red} IIII & G & \color{Red} III & 11 & \multicolumn{2}{l}{\color{Red} Martini Episcopi et Conf. Semid. Menne Martyris. Mem.} \color{Red} \\
% \color{Red}  & \color{Red} A & \color{Red} II & 12 &  & \color{Red} \\
% \color{Red} XII & B & \color{Red} Idus & 13 & Brictii Episcopi et Conf. Mem. & \color{Red} \\
% \color{Red} I & C & \color{Red} XVIII & 14 & & \color{Red} \\
% \color{Red}  & D & \color{Red} XVII & 15 & & \color{Red} \\
% \color{Red} IX & E & \color{Red} XVI & 16 & & \color{Red} \\
% \color{Red}  & F & \color{Red} XV & 17 &  & \color{Red} Sol in Sagittario. \\
% \color{Red} XVII & G & \color{Red} XIIII & 18 & Octava Sancti Martini. \color{Red} Simplex. \\
% \color{Red} VI & \color{Red} A & \color{Red} XIII & 19 & Sancte Elysabeth. Mem. & \color{Red} \\
% \color{Red}  & B & \color{Red} XII & 20 & & \color{Red} \\
% \color{Red} XIIII & C & \color{Red} XI & 21 & & \color{Red} \\
\color{Red} III & D & \color{Red} X & 22 & Ob. F. Stephanus Mag. Ord. Octavus. & \color{Red} \\
% \color{Red}  & E & \color{Red} IX & 23 & Clementis Pape et Martyris. Simplex. & \color{Red} \\
% \color{Red} XI & F & \color{Red} VIII & 24 & Crisogoni Martyris. Mem. & \color{Red} \\
% \color{Red} XIX & G & \color{Red} VII & 25 & \color{Red} Katherine Virginis et Martyris. Semid. \\
% \color{Red}  & \color{Red} A & \color{Red} VI & 26 & & \color{Red} \\
% \color{Red} VIII & B & \color{Red} V & 27 & Vitalis et Agricole Martyrum. Mem. & \color{Red} \\
% \color{Red}  & C & \color{Red} IIII & 28 & & \color{Red} \\
\color{Red} XVI & D & \color{Red} III & 29 & Ob. F. Iohannes Mag. Ord. Sextus. & \\
% \color{Red} V & E & \color{Red} II & 30 & Andree Apostoli. Semid. & \color{Red} \\


& & & & \color{Red} Decembris. \\
% \color{Red} XIII & F & \textbf{\color{Red} K\color{Blue}L} & 1 & & \color{Red} Finis. V. Embolismus. \\
% \color{Red} II & G & \color{Red} IIII & 2 & & \color{Red} Incipit Primus Embolismus. \\
% \color{Red}  & \color{Red} A & \color{Red} III & 3 & & \color{Red} \\
% \color{Red} X & B & \color{Red} II & 4 & & \color{Red} \\
% \color{Red}  & C & \color{Red} Non. & 5 & & \color{Red} \\
% \color{Red} XVIII & D & \color{Red} VIII & 6 & \color{Red} Nicolai Episcopi. Semid. & \color{Red} \\
% \color{Red} VII & E & \color{Red} VII & 7 & \color{Red} Octava Sancti Andree. Mem. \\
% \color{Red}  & F & \color{Red} VI & 8 & & \color{Red} \\
% \color{Red} XV & G & \color{Red} V & 9 &  & \color{Red} \\
% \color{Red} IIII & \color{Red} A & \color{Red} IIII & 10 &  & \color{Red} \\
% \color{Red}  & B & \color{Red} III & 11 & Damasi Pape et Confess. Mem. & \color{Red} \\
% \color{Red} XII & C & \color{Red} II & 12 &  & \color{Red} \\
% \color{Red} I & D & \color{Red} Idus & 13 & Luce Virginis et Martyris. Simplex. & \color{Red} \\
% \color{Red}  & E & \color{Red} XIX & 14 &  & \color{Red} \\
% \color{Red} IX & F & \color{Red} XVIII & 15 & Ob. F. Conradus Mag. Ord. XXX. & \\
% \color{Red}  & G & \color{Red} XVII & 16 & & \color{Red} \\
% \color{Red} XVII & \color{Red} A & \color{Red} XVI & 17 & & O Sapiencia. \color{Red} \\
% \color{Red} VI & B & \color{Red} XV & 18 & & \color{Red} Sol in Capricorno. \\
% \color{Red}  & C & \color{Red} XIIII & 19 & & \color{Red} \\
% \color{Red} XIIII & D & \color{Red} XIII & 20 & & \color{Red} \\
% \color{Red} III & E & \color{Red} XII & 21 & Thome Apostoli. Semid. & \color{Red} \\
% \color{Red}  & F & \color{Red} XI & 22 & & \color{Red} \\
% \color{Red} XI & G & \color{Red} X & 23 & & \color{Red} \\
% \color{Red} XIX & \color{Red} A & \color{Red} IX & 24 & & \color{Red} Vigilia. \\
% \color{Red}  & B & \color{Red} VIII & 25 & \multicolumn{2}{l}{\color{Red} Nativitas Domini Nos. Ihesu Christi. Totum Dupl.} \color{Red} \\
% \color{Red} VIII & C & \color{Red} VII & 26 & Stephani Prothomartyris. Totum Dup. & \color{Red} \\
% \color{Red}  & D & \color{Red} VI & 27 & \multicolumn{2}{l}{\color{Red} Iohannis Apostoli et Evangeliste. Totum Dupl.} \color{Red} \\
% \color{Red} XVI & E & \color{Red} V & 28 & Sanctorum Innocentium. Simplex. & \color{Red} \\
% \color{Red} V & F & \color{Red} IIII & 29 & \color{Red} Thome Episcopi et Martyris. Simplex. & \color{Red} \\
% \color{Red}  & G & \color{Red} III & 30 &  & \color{Red} \\
% \color{Red} XIII & \color{Red} A & \color{Red} II & 31 & Silvestri Pape et Confessoris. Simplex. & \color{Red} Finis. I. Embolismus. \\
\end{tabular}
\end{center}

\newpage

\begin{multicols*}{2}

\ubsection{Rubrice}{martyrologium-rubrice}

{\color{Red} \ubsubsection{De Officio Legende De Kalendario et Luna et Martyrologyo et Aliis His Annexis}{de-officio-legende-de-kalendario-et-luna-et-martyrologyo-et-aliis-his-annexis}}
\lettrine[lines=2]{\zallmancaps \color{Blue} D}{}\ul{\textsc{um} dicitur psalmus} \psalm{116} \ul{in Matutinis: accedit qui lecturus est Kalendas ad prelatum, dicens ei summisse: Capitulum? Qui si dixerit Non: legende erunt Kalende in choro, finitis Matutinis. Si autem dixerit Post Primam: tunc reservande sunt usque ad Primam. Si autem dixerit Sic: legende erunt in capitulo post Matutinas.}
% \newline \ul{Cum ergo legende erunt in capitulo post matutinas: debet cum incipitur Benedictus pulsare capitulum, tam prolixe quod fratres possint commode advertere quod ad capitulum pulsetur. Qui modus pulsandi, semper est servandus in convocatione capituli. Debet uatem lumen in capitulo tam diligenter parare: quod durante capitulo non oporteat circa ideam iterum aliquem occupari. Debet etiam providere: ne iurge pro discplinis desint. Prato autem libro, et lumine pro se ad legendum, et tabula pro tempore, et aliis pertinentibus ad officium suum: stans ante pulpitum in quo debet legere, dum intrat conventus in capitulum, divertat aliquantulum transeunte prelato ob reverentiam eius. Deinde rediens ad locum illum, cum conventus intraverit: incipiat legere in voce mediocri que legenda fuerint, secundum  modum infrascriptum. Perlectis autem omnibus, et clauso libro, et aliis faciendis breviter expeditis: vadat ad locum suum, finito capitulo reportaturus que fuerint alibi reportanda.}
% \newline \ul{Cum autem reservantur usque post Primam, dum dicitur ultimus psalmis in Prima: accedat qui legere debet ad prelatum, interrogans eum ut supra. Etsi dixerit sic: tunc cito pulsetur ad capitulum ante terminationem hore, et omnia faciat ut supra, pretquam de luminibus, si sit tanta dies quod non sit lumen necessarium. Si autem dixerit Non: legende sunt in choro. Cum autem in choro leguntur sive post Matutinas sive post Primam: non pulsatur capitulum, sed alia fiunt modo supradicto, prout competunt.}
% \newline \ul{De supradictis autem: legendum est sic.}
% \newline Decimonono Kalendas \ul{vel} Sextodecimo Kalendas \ul{vel} Pridie Kalendas \ul{vel} Kalendis Ianuarii: \ul{vel} Quarto Nonas \ul{vel} Pridie Nonas \ul{vel}  \ul{vel} Nonis Ianuarii: \ul{vel} Octavo Idus \ul{vel} Pridie Idus \ul{vel} Idibus Ianuarii: Luna prima, \ul{vel} luna sextadecima; \ul{vel} tricesima. Rome: via apia: natalis sancti talis et cetera. Et alilorum plurimorum sanctorum: martyum, confessorum, atque virginum. \ul{Hec autem clausula semper addatur in fine lectionis de martyologio. Deinde dicta oratione} Dirigere et sanctificare etc: \ul{sequitur} Iube domne. \ul{Data autem benedictione: sequitur lectio de constitutionibus, vel de evangelio: sic} Secundum Matheum; \ul{vel} Secundum Marcum; \ul{vel} Secundum Lucam; \ul{vel} Secundum Iohannem; \ul{Et legatur de evangelio totum illud quod infrascriptum est in hoc libro. In fine vero lectionis tam de evangelio quam de constitutionibus dicatur semper} Tu autem.
\newline \ul{De evangelio autem legendum est omnibus diebus Dominicis, et festis simplicibus et supra. Item in vigilia Nativitatis Domini evangelium de vigilia: sive fuerit in Dominica: sive non. Et in omni vigilia que proprium habet officium in Missa, si in Dominica evenerit: evangelium de ipsa vigilia legendum est, preterquam in vigilia Beati Andree cum in Prima Dominica Adventus evenerit: qua tunc evangelium de Dominica est legendum. Item in crastino sancti Thome martyris. Item feria quarta in capite ieiunii, et in Cena Domini. Item in vigilia Pasche, et cotidie per Octavas Pasche et Pentecostes. Item in commemoratione omnium fidelium defunctorum. Ceteris vero diebus legendum est de constitutionibus. In varietatibus autem que occurrunt in huiusmodi lectionibus: servandus est modus qui in lectionario est notatus.}
\newline \ul{Dicto igitur} Tu autem \ul{et facta responsione sequitur} Commemoratio fratrum, familiarium, benefactorum, defunctorum: ordinis nostri.
\newline \ul{Si vero recitandus fuerit obitus alicuius, vel anniversarium aliquod fuerit in crastino: tunc antequam dicatur} Commemoratio \ul{dicendum est} Obiit talis, \ul{vel} Anniversarium Patrum et Matrum.
% \newline \ul{Solet autem tabula septimana semper legi cum Kalendis. Tabula vero pro diebus aliis in mensa, vel in choro post gratiarum actiones dicto} Fidelium anime \ul{et} Pater noster, \ul{et fratribus sedentibus in locis suis.} 

\end{multicols*}

%pp041
\ubchapter{Collectarium}
% Liber iste qui collectarius dicitur eo quo propter collectas sive orationes habendas in promptu principaliter scriptus est, vel quia in unum sunt in eo collecta ea que dicenda sunt a sacerdote extra officium Misse: contient ista per ordinem.

\ubsection{Kalendarium Sanctorum}{kalendarium-sanctorum}

\thispagestyle{fancy}

\vspace{+.5em}
\color{Red}\csection{Januarius}{\color{Red} habet dies. XXXI. luna. XXX.}\color{black}

\begin{center}
\begin{tabular}{l | l | l | r | l l}
 \color{Red} III & \color{Red} A & \textbf{\color{Red} K\color{Blue}L} & 1 & \color{Red} Circumcisio domini. \color{Red} Duplex. \\
 & B & \color{Red} IIII & 2 & Oct. Sancti Stephani. & \color{Red} \\
\color{Red} XI & C & \color{Red} III & 3 & Oct. Sancti Johannis. & \\
 &  & & & \multicolumn{2}{l}{\qquad \color{Red} Incipit Quartus Embolismus.} \\
& D & \color{Red} II & 4 & \color{Red} Oct. Sanctorum Innocentium. & \color{Red} \\
\color{Red} XIX & E & \color{Red} Non. & 5 & & \color{Red} \\
\color{Red} VIII & F & \color{Red} VIII & 6 & Epiphania Domini. Totum Dup. \\
 & G & \color{Red} VII & 7 & \qquad \color{Red} Claves. LXX. \\
\color{Red} XVI & \color{Red} A & \color{Red} VI & 8 & & \\
\color{Red} V & B & \color{Red} V & 9 & & \\
 & C & \color{Red} IIII & 10 & Pauli Heremite. \color{Red} Mem. \\
\color{Red} XIII & D & \color{Red} III & 11 & & \color{Red} \\
\color{Red} II & E & \color{Red} II. & 12 & & \color{Red} \\
 & F & \color{Red} Idus & 13 & \multicolumn{2}{l}{Oct. Epiphanie. \color{Red} Simplex.} \\
 &  & & & \multicolumn{2}{l}{\quad \color{black} Hylarii et Remigii. \color{Red} Mem.} \\
\color{Red} X & G & \color{Red} XIX & 14 & \multicolumn{2}{l}{Felicis Presbiteri et Conf. III. Lect.} \\
 & \color{Red} A & \color{Red} XVIII & 15 & \color{Red} Mauri Abbatis. III. Lect. \\
\color{Red} XVIII & B & \color{Red} XVII & 16 & Marcelli Pape et Martyris. III. Lect. \\
\color{Red} VII & C & \color{Red} XVI & 17 & \color{Red} Anthonii Abbatis. III. Lect. \\
 & D & \color{Red} XV & 18 & Prisce Virginis et Martyris. III. Lect. \\
 &  & & & \qquad \color{Red} Sol in Aquario. Initium LXX. \\
\color{Red} XV & E & \color{Red} XIIII & 19 & \color{Red} \\
\color{Red} IIII & F & \color{Red} XIII & 20 & \multicolumn{2}{l}{Fabiani et Sebastiani Mart. Simplex.} \\
 & G & \color{Red} XII & 21 & Agnetis Virginis et Mart. Simplex. \\
\color{Red} XII & \color{Red} A & \color{Red} XI & 22 & Vincentii Martyris. Semid. \\
\color{Red} I & B & \color{Red} X & 23 & \color{Red} Emerentiane Virg. et Mart. Mem. \\
 & C & \color{Red} IX & 24 & & \color{Red} \\
\color{Red} IX & D & \color{Red} VIII & 25 & Conversio Sancti Pauli. Semid. \\
& E & \color{Red} VII & 26 & & \color{Red} \\
\color{Red} XVII & F & \color{Red} VI & 27 & \color{Red} Juliani Episcopi et Conf. Mem. \\
\color{Red} VI & G & \color{Red} V & 28 & \color{Red} Agnetis Secundo. III. Lect. \\
 &  & & & \qquad Claves XL. \\
 & \color{Red} A & \color{Red} IIII & 29 & & \color{Red} \\
\color{Red} XIIII & B & \color{Red} III & 30 & & \color{Red} \\
\color{Red} III & C & \color{Red} II & 31 & & \color{Red} \\
\end{tabular}
\color{Red} Hore noctis. XVI. Diei. VIII.
\end{center}



\newpage
\color{Red}\csection{Februarius}{\color{Red} habet dies. XXVIII. luna. XXIX.}\color{black}

\begin{center}
\begin{tabular}{l | l | l | r | l l}
 & D & \textbf{\color{Red} K\color{Blue}L} & 1 & \color{Red} Ignatii Episcopi et Martyris. Mem. \\
 &  & & & \qquad Finis. IV. embolismi. \\
\color{Red} XI & E & \color{Red} IIII & 2 & \multicolumn{2}{l}{Purificatio S. Marie Virg. Tot. Dup.} \\
\color{Red} XIX & F & \color{Red} III & 3 & \color{Red} Blasii Episcopi et Martyris. III. Lect. \\
\color{Red} VIII & G & \color{Red} II & 4 & Anniversarium Patrum et Matrum. \\
 & \color{Red} A & \color{Red} Non. & 5 & \color{Red} Agathe Virginis et Mart. Simplex. \\
\color{Red} XVI & B & \color{Red} VIII & 6 &  Vedasti et Amandi Epis. Memor. \\
\color{Red} V & C & \color{Red} VII & 7 & \\
\color{Red}  & D & \color{Red} VI & 8 & \qquad \color{Red} Initium. XL. \\
\color{Red} XIII & E & \color{Red} V & 9 & & \color{Red} \\
\color{Red} II & F & \color{Red} IIII & 10 & Scolastice Virginis. Mem. & \color{Red} \\
\color{Red}  & G & \color{Red} III & 11 & & \color{Red} \\
\color{Red} X & \color{Red} A & \color{Red} II & 12 & \\
\color{Red}  & B & \color{Red} Idus & 13 & & \color{Red} \\
\color{Red} XVIII & C & \color{Red} XVI & 14 & \color{Red} Valentini Martyris. III. Lect. \\
\color{Red} VII & D & \color{Red} XV & 15 & \qquad \color{Red} Sol in Piscibus. \\
\color{Red}  & E & \color{Red} XIIII & 16 & & \color{Red} \\
\color{Red} XV & F & \color{Red} XIII & 17 & & \color{Red} \\
\color{Red} IIII & G & \color{Red} XII & 18 & & \color{Red} \\
\color{Red}  & \color{Red} A & \color{Red} XI & 19 & & \color{Red} \\
\color{Red} XII & B & \color{Red} X & 20 & & \color{Red} \\
\color{Red} I & C & \color{Red} IX & 21 & & \color{Red} \\
\color{Red}  & D & \color{Red} VIII & 22 & \color{Red} Cathedra Sancti Petri. Simplex. \\
\color{Red} IX & E & \color{Red} VII & 23 & & \color{Red} \\
\color{Red}  & F & \color{Red} VI & 24 & Mathie Apostoli. Semid. \\
 &  & & & \qquad \color{Red} Locus Byssexti. \\
\color{Red} XVII & G & \color{Red} V & 25 & & \color{Red} \\
\color{Red} VI & \color{Red} A & \color{Red} IIII & 26 & & \color{Red} \\
\color{Red}  & B & \color{Red} III & 27 & & \color{Red} \\
\color{Red} XVI & C & \color{Red} II & 28 & & \color{Red} \\
% \hline
\end{tabular}
\newline \color{Red} Hore noctis. XIIII. Diei. X.
\end{center}



\newpage
\color{Red}\csection{Martius}{\color{Red} habet dies. XXXI. luna. XXX.}\color{black}

\begin{center}
\begin{tabular}{l | l | l | r | l l}
\color{Red} III & D & \textbf{\color{Red} K\color{Blue}L} & 1 & \color{Red} Albini Episcopi et Conf. Mem. \\
\color{Red}  & E & \color{Red} VI & 2 & & \color{Red} \\
\color{Red} XI & F & \color{Red} V & 3 & \\
\color{Red}  & G & \color{Red} IIII & 4 & & \color{Red} \\
\color{Red} XIX & \color{Red} A & \color{Red} III & 5 & \qquad \color{Red} Incipit. VII. Embolismus. \\
\color{Red} VIII & B & \color{Red} II & 6 & \qquad \color{Red} Incipit. III. Embolismus. \\
\color{Red}  & C & \color{Red} Non. & 7 & & \color{Red} \\
\color{Red} XVI & D & \color{Red} VIII & 8 & & \color{Red} \\
\color{Red} V & E & \color{Red} VII & 9 & & \color{Red} \\
\color{Red}  & F & \color{Red} VI & 10 & & \color{Red} \\
\color{Red} XIII & G & \color{Red} V & 11 & \qquad \color{Red} Claves Pasche. \\
\color{Red} II & \color{Red} A & \color{Red} IV. & 12 & \color{Red} Gregorii Pape et Conf. Simplex. \\
\color{Red}  & B & \color{Red} III & 13 & & \color{Red} \\
\color{Red} X & C & \color{Red} II & 14 & \qquad \color{Red} Ultima XXL. \\
\color{Red}  & D & \color{Red} Idus & 15 & \qquad \color{Red} Equinoctium Vernale. \\
\color{Red} XVIII & E & \color{Red} XVII & 16 & \\
\color{Red} VII & F & \color{Red} XVI & 17 & & \color{Red} \\
\color{Red}  & G & \color{Red} XV & 18 & \qquad \color{Red} Sol in Ariete. \\
\color{Red} XV & \color{Red} A & \color{Red} XIIII & 19 & & \color{Red} \\
\color{Red} IIII & B & \color{Red} XIII & 20 & & \color{Red} \\
\color{Red}  & C & \color{Red} XII & 21 & \color{Red} Benedicti Abbatis. Simplex. \\
\color{Red} XII & D & \color{Red} XI & 22 & \qquad \color{Red} Primum Pascha. \\
\color{Red} I & E & \color{Red} X & 23 & & \color{Red} \\
\color{Red}  & F & \color{Red} IX & 24 & & \color{Red} \\
\color{Red} IX & G & \color{Red} VIII & 25 & Annunciatio Dominica. Totum Dupl.& \color{Red} \\
\color{Red}  & \color{Red} A & \color{Red} VII & 26 & & \color{Red} \\
\color{Red} XVII & B & \color{Red} VI & 27 & & \color{Red} \\
\color{Red} VI & C & \color{Red} V & 28 & & \color{Red} \\
\color{Red}  & D & \color{Red} IIII & 29 & \color{Red} \\
\color{Red} XIIII & E & \color{Red} III & 30 & & \color{Red} \\
\color{Red} III & F & \color{Red} II & 31 & & \color{Red} \\
\end{tabular}
\newline \color{Red} Hore noctis. XII. Diei. XII.
\end{center}


%pp042
\newpage
\color{Red}\csection{Aprilis}{\color{Red} habet dies. XXX. luna. XXIX.}\color{black}

\begin{center}
\begin{tabular}{l | l | l | r | l l}
\color{Red}  & G & \textbf{\color{Red} K\color{Blue}L} & 1 & & \color{Red} \\
\color{Red} XI & \color{Red} A & \color{Red} IIII & 2 & & \color{Red} \\
\color{Red}  & B & \color{Red} III & 3 & \qquad Finis. VII. Embolismus. \color{Red} \\
\color{Red} XIX & C & \color{Red} II & 4 & Ambrosii Episcopi et Conf. Simplex. \\
 & & & & \qquad \color{Red} Finis. VI. Embolismus. \\
\color{Red} VIII & D & \color{Red} Non. & 5 & & \color{Red} \\
\color{Red} XVI & E & \color{Red} VIII & 6 & & \color{Red} \\
\color{Red} V & F & \color{Red} VII & 7 & & \color{Red} \\
\color{Red}  & G & \color{Red} VI & 8 & & \color{Red} \\
\color{Red} XIII & \color{Red} A & \color{Red} V & 9 & & \color{Red} \\
\color{Red} II & B & \color{Red} IIII & 10 & & \color{Red} \\
\color{Red}  & C & \color{Red} III & 11 & & \color{Red} \\
\color{Red} X & D & \color{Red} II & 12 & & \color{Red} \\
\color{Red}  & E & \color{Red} Idus & 13 & & \color{Red} \\
\color{Red} XVIII & F & \color{Red} XVIII & 14 & \multicolumn{2}{l}{Tyburtii Valeriani, et Maximi} \\
 & & & & \quad \color{Red} Martyrum. \color{black} III. Lect. \color{Red} \\
\color{Red} VII & G & \color{Red} XVII & 15 & \qquad \color{Red} Claves Rogationum. \\
\color{Red}  & \color{Red} A & \color{Red} XVI & 16 & \\
\color{Red} XV & B & \color{Red} XV & 17 & \qquad Sol in Tauro. \color{Red} \\
\color{Red} IIII & C & \color{Red} XIIII & 18 & & \color{Red} \\
\color{Red}  & D & \color{Red} XIII & 19 & & \color{Red} \\
\color{Red} XII & E & \color{Red} XII & 20 & & \color{Red} \\
\color{Red} I & F & \color{Red} XI & 21 & & \color{Red} \\
\color{Red}  & G & \color{Red} X & 22 & & \color{Red} \\
\color{Red} IX & \color{Red} A & \color{Red} IX & 23 & Georgii Martyris. Simplex. & \color{Red} \\
\color{Red}  & B & \color{Red} VIII & 24 & & \color{Red} \\
\color{Red} XVII & C & \color{Red} VII & 25 & \color{Red} Marci Evangeliste. Semid. \\
 & & & & \qquad Ultimum Pascha. \\
\color{Red} VI & D & \color{Red} VI & 26 & & \color{Red} \\
\color{Red}  & E & \color{Red} V & 27 & & \color{Red} \\
\color{Red} XIIII & F & \color{Red} IIII & 28 & Vitalis Martyris. III. Lect. & \color{Red} \\
\color{Red} III & G & \color{Red} III & 29 & \multicolumn{2}{l}{Beati Petri Martyris de Ordinis } \\
 & & & & \quad Predicatorum. Totum Duplex. \\
 & & & & \qquad \color{Red} Claves Pente. \\
\color{Red}  & \color{Red} A & \color{Red} II & 30 & & \color{Red} \\
\end{tabular}
\newline \color{Red} Hore noctis. X. Diei. XIIII.
\end{center}



\newpage
\color{Red}\csection{Maius}{\color{Red} habet dies. XXXI. luna. XXIX.}\color{black}

\begin{center}
\begin{tabular}{l | l | l | r | l l}
\color{Red} XI & B & \textbf{\color{Red} K\color{Blue}L} & 1 & Philippi et Iacobi Apos. Semid. \\
\color{Red}  & C & \color{Red} VI & 2 & \\
\color{Red} XIX & D & \color{Red} V & 3 & Inventio Sancte Crucis. Semid. \\
 &  &  &  & \multicolumn{2}{l}{\quad Alexandri Evencii, et Theodoli.} \\
 &  &  &  & \quad \color{Red} Martyrum. \color{black} Mem. \\
\color{Red} VIII & E & \color{Red} IIII & 4 & Festum Corone Domini. Simplex. & \color{Red} \\
\color{Red}  & F & \color{Red} III & 5 & & \color{Red} \\
\color{Red} XVI & G & \color{Red} II & 6 & \color{Red} Iohannis Ante Portam Lati. Semid. \\
\color{Red} V & \color{Red} A & \color{Red} Non. & 7 & & \color{Red} \\
\color{Red}  & B & \color{Red} VIII & 8 & & \color{Red} \\
\color{Red} XIII & C & \color{Red} VII & 9 & & \color{Red} \\
\color{Red} II & D & \color{Red} VI & 10 & \multicolumn{2}{l}{Gordiani et Epimachi Mart. III. Lect.} \color{Red} \\
\color{Red}  & E & \color{Red} V & 11 & & \color{Red} \\
\color{Red} X & F & \color{Red} IIII & 12 & \multicolumn{2}{l}{Neri, Achillei, et Pancratii Mart.} \\
 &  &  &  & \quad III. Lect. \color{Red} \\
\color{Red}  & G & \color{Red} III & 13 & & \color{Red} \\
\color{Red} XVIII & \color{Red} A & \color{Red} II & 14 & & \color{Red} \\
\color{Red} VII & B & \color{Red} Idus & 15 & & \color{Red} \\
\color{Red}  & C & \color{Red} XVII & 16 &  & \color{Red} \\
\color{Red} XV & D & \color{Red} XVI & 17 & & \color{Red} \\
\color{Red} IIII & E & \color{Red} XV & 18 & \qquad Sol in Geminis. \color{Red} \\
\color{Red}  & F & \color{Red} XIIII & 19 & Potentiane Virginis. Mem. & \color{Red} \\
\color{Red} XII & G & \color{Red} XIII & 20 & & \color{Red} \\
\color{Red} I & \color{Red} A & \color{Red} XII & 21 & & \color{Red} \\
\color{Red}  & B & \color{Red} XI & 22 & & \color{Red} \\
\color{Red} IX & C & \color{Red} X & 23 & Translatio Beati Dominici. Tot. Dup. & \color{Red} \\
\color{Red}  & D & \color{Red} IX & 24 & & \color{Red} \\
\color{Red} XVII & E & \color{Red} VIII & 25 & Urbani Pape et Martyris. III. Lect. & \color{Red} \\
\color{Red} VI & F & \color{Red} VII & 26 & & \color{Red} \\
\color{Red}  & G & \color{Red} VI & 27 & & \color{Red} \\
\color{Red} XIIII & \color{Red} A & \color{Red} V & 28 & & \color{Red} \\
\color{Red} III & B & \color{Red} IIII & 29 & & \color{Red} \\
\color{Red}  & C & \color{Red} III & 30 & & \color{Red} \\
\color{Red} XI & D & \color{Red} II & 31 & Petronille Virginis. Mem. & \color{Red} \\
\end{tabular}
\newline \color{Red} Hore noctis. VIII. Diei. XVI.
\end{center}



\newpage
\color{Red}\csection{Iunius}{\color{Red} habet dies. XXX. luna. XXIX.}\color{black}

\begin{center}
\begin{tabular}{l | l | l | r | l l}
\color{Red}  & E & \textbf{\color{Red} K\color{Blue}L} & 1 & & \color{Red} \\
\color{Red} XIX & F & \color{Red} IIII & 2 & Marcellini et Petri Mart. III. Lect. & \color{Red} \\
\color{Red} VIII & G & \color{Red} III & 3 & & \color{Red} \\
\color{Red} XVI & \color{Red} A & \color{Red} II & 4 & & \color{Red} \\
\color{Red} V & B & \color{Red} Non. & 5 & & \color{Red} \\
\color{Red}  & C & \color{Red} VIII & 6 & & \color{Red} \\
\color{Red} XIII & D & \color{Red} VII & 7 & & \color{Red} \\
\color{Red} II & E & \color{Red} VI & 8 & Medardi Episcopi et Conf. Mem. & \color{Red} \\
\color{Red}  & F & \color{Red} V & 9 & Primi et Feliciani Mart. III. Lect. & \color{Red} \\
\color{Red} X & G & \color{Red} IIII & 10 & & \color{Red} \\
\color{Red}  & \color{Red} A & \color{Red} III & 11 & \color{Red} Barnabe Apostoli. Semid. \\
\color{Red} XVIII & B & \color{Red} II & 12 & \multicolumn{2}{l}{\color{Red} Basilidis Cyrini, Naboris, et } \\
 &  &  &  & \quad \color{Red} Nazarii, Martyrum. III. Lect. \\
\color{Red} VII & C & \color{Red} Idus & 13 & & \color{Red} \\
\color{Red}  & D & \color{Red} XVIII & 14 & \qquad Solsticium Estivale. \color{Red} \\
\color{Red} XV & E & \color{Red} XVII & 15 & \color{Red} Viti et Modest Martyrum. Mem. \\
\color{Red} IIII & F & \color{Red} XVI & 16 & \color{Red} Cyrici et Iulice Martyrum. Mem. \\
\color{Red}  & G & \color{Red} XV & 17 & \qquad Sol in Cancio. \color{Red} \\
\color{Red} XII & \color{Red} A & \color{Red} XIIII & 18 & \color{Red} Marci et Marcelliani Mart. III. Lect. \\
\color{Red} I & B & \color{Red} XIII & 19 & \color{Red} Gervasii et Prothasii Mart. Simplex. \\
\color{Red}  & C & \color{Red} XII & 20 & & \color{Red} \\
\color{Red} IX & D & \color{Red} XI & 21 & & \color{Red} \\
\color{Red}  & E & \color{Red} X & 22 & & \color{Red} \\
\color{Red} XVII & F & \color{Red} IX & 23 & \qquad Vigilia. \color{Red} \\
\color{Red} VI & G & \color{Red} VIII & 24 & Nativitas S. Iohannis Baptiste. Dupl. \\
\color{Red}  & \color{Red} A & \color{Red} VII & 25 & & \color{Red} \\
\color{Red} XIIII & B & \color{Red} VI & 26 & Iohannis et Pauli Mart. Simplex. & \color{Red} \\
\color{Red} III & C & \color{Red} V & 27 & & \color{Red} \\
\color{Red}  & D & \color{Red} IIII & 28 & Leonis Pape et Conf. Memoria. \\
 &  &  &  & \qquad Vigilia. \color{Red} \\
\color{Red} XI & E & \color{Red} III & 29 & Apostolorum Petri et Pauli. Duplex. & \color{Red} \\
\color{Red}  & F & \color{Red} II & 30 & Commemoratio Sancti Pauli. \color{Red} Semid. \\
\end{tabular}
\newline \color{Red} Hore noctis. VI. Diei. XVIII.
\end{center}



\newpage
\color{Red}\csection{Iulius}{\color{Red} habet dies. XXXI. luna. XXX.}\color{black}

\begin{center}
\begin{tabular}{l | l | l | r | l l}
\color{Red} XIX & G & \textbf{\color{Red} K\color{Blue}L} & 1 & Octava S. Iohannis Baptiste. Simplex. & \color{Red} \\
\color{Red} VIII & \color{Red} A & \color{Red} VI & 2 & Processi et Martiniani Mart. Mem. & \color{Red} \\
\color{Red}  & B & \color{Red} V & 3 & & \color{Red} \\
\color{Red} XVI & C & \color{Red} IIII & 4 & & \color{Red} \\
\color{Red} V & D & \color{Red} III & 5 & & \color{Red} \\
\color{Red}  & E & \color{Red} II & 6 & \color{Red} Oct. Apos. Petri et Pauli. Simplex. \\
\color{Red} XIII & F & \color{Red} Non. & 7 & & \color{Red} \\
\color{Red} II & G & \color{Red} VIII & 8 & & \color{Red} \\
\color{Red}  & \color{Red} A & \color{Red} VII & 9 & & \color{Red} \\
\color{Red} X & B & \color{Red} VI & 10 & Septem Fratrum. III. Lect. & \color{Red} \\
\color{Red}  & C & \color{Red} V & 11 & & \color{Red} \\
\color{Red} XVIII & D & \color{Red} IIII & 12 &  & \color{Red} \\
\color{Red} VII & E & \color{Red} III & 13 & & \color{Red} \\
\color{Red}  & F & \color{Red} II & 14 & \qquad \color{Red} Dies Caniculares. \\
\color{Red} XV & G & \color{Red} Idus & 15 & & \color{Red} \\
\color{Red} IIII & \color{Red} A & \color{Red} XVII & 16 & \\
\color{Red}  & B & \color{Red} XVI & 17 & & \color{Red} \\
\color{Red} XII & C & \color{Red} XV & 18 & \qquad \color{Red} Sol in Leone. \\
\color{Red} I & D & \color{Red} XIIII & 19 & & \color{Red} \\
\color{Red}  & E & \color{Red} XIII & 20 & Margarete Virg. et Mart. Simplex. & \color{Red} \\
\color{Red} IX & F & \color{Red} XII & 21 & Praxedis Virginis. III. Lect. & \color{Red} \\
\color{Red}  & G & \color{Red} XI & 22 & \color{Red} Marie Magdalene. Semid. \\
\color{Red} XVII & \color{Red} A & \color{Red} X & 23 & \color{Red} Apollinaris Epis. et Mart. III. Lect. \\
\color{Red} VI & B & \color{Red} IX & 24 & Christine Virginis et Mart. Mem. \\
\color{Red}  & C & \color{Red} VIII & 25 & \multicolumn{2}{l}{Iacobi Apostoli. Semi. \quad Christofori } \\
 &  &  &  & \quad et Cucufatis Martyrum. Mem. \\
\color{Red} XIIII & D & \color{Red} VII & 26 & & \color{Red} \\
\color{Red} III & E & \color{Red} VI & 27 & & \color{Red} \\
\color{Red}  & F & \color{Red} V & 28 & \multicolumn{2}{l}{Nazarii, Celsi, et Pantaleonis.} \\
 &  &  &  & \quad Martyrum. III. Lect. \color{Red} \\
\color{Red} XI & G & \color{Red} IIII & 29 & \multicolumn{2}{l}{\color{Red} Felicis, Simplicii, Faustini, et} \\
 &  &  &  & \quad \color{Red} Beatricis Martyrum. III. Lect. \\
\color{Red} XIX & \color{Red} A & \color{Red} III & 30 & Abdon et Sennen Mart. III. Lect. & \color{Red} \\
\color{Red}  & B & \color{Red} II & 31 & Germani Episcopi et Conf. III. Lect. & \color{Red} \\
\end{tabular}
\newline \color{Red} Hore noctis. VIII. Diei. XVI.
\end{center}



\newpage
\color{Red}\csection{Augustus}{\color{Red} habet dies. XXXI. luna. XXX.}\color{black}

\begin{center}
\begin{tabular}{l | l | l | r | l l}
\color{Red} VIII & C & \textbf{\color{Red} K\color{Blue}L} & 1 & Ad Vincula Sancti Petri. Simplex. & \color{Red} \\
 &  & &  & \multicolumn{2}{l}{\qquad S. Machabeorum Mart. Mem. \color{Red}} \\
\color{Red} XVI & D & \color{Red} IIII & 2 & \color{Red} Stephani Pape et Martyris. III. Lect. \\
\color{Red} V & E & \color{Red} III & 3 & \color{Red} Inventio Sancti Stephani. Simplex. \\
\color{Red}  & F & \color{Red} II & 4 & & \color{Red} \\
\color{Red} XIII & G & \color{Red} Non. & 5 & Beati Dominici Conf. Totum Dup. & \color{Red} \\
\color{Red} II & \color{Red} A & \color{Red} VIII & 6 & \multicolumn{2}{l}{Sixti Pape, Felicissimi et Agapiti}\\
 &  & &  & \quad Martyrum. Mem. \color{Red} \\
\color{Red}  & B & \color{Red} VII & 7 & Donati Episcopi et Martyris. Mem. & \color{Red} \\
\color{Red} X & C & \color{Red} VI & 8 & Cyriaci Socio. eius Mart. Mem. & \color{Red} \\
\color{Red}  & D & \color{Red} V & 9 & \qquad \color{Red} Vigilia. \\
\color{Red} XVIII & E & \color{Red} IIII & 10 & \color{Red} Laurentii Martyris. Semid. & \color{Red} \\
\color{Red} VII & F & \color{Red} III & 11 & Tiburtii Martyris. Mem. & \color{Red} \\
\color{Red}  & G & \color{Red} II & 12 & Octava Sancti Dominici. Simplex. & \color{Red} \\
\color{Red} XV & \color{Red} A & \color{Red} Idus & 13 & Ypoliti Sociorumque eius . Simplex. & \color{Red} \\
\color{Red} IIII & B & \color{Red} XIX & 14 & Eusebii Presbiteri et Conf. Mem. & \\
 &  & &  & \qquad \color{Red} Vigilia. \\
\color{Red}  & C & \color{Red} XVIII & 15 & \multicolumn{2}{l}{\color{Red} Assumptio B. Marie Virg. Tot. Dupl.} \color{Red} \\
\color{Red} XII & D & \color{Red} XVII & 16 & & \color{Red} \\
\color{Red} I & E & \color{Red} XVI & 17 & Octava Sancti Laurentii. Simplex. & \color{Red} \\
\color{Red}  & F & \color{Red} XV & 18 & Agapiti Martyris. Mem. & \\
 &  & &  & \qquad \color{Red} Sol in Virgine. \\
\color{Red} IX & G & \color{Red} XIIII & 19 & & \color{Red} \\
\color{Red}  & \color{Red} A & \color{Red} XIII & 20 & Bernardi Abbatis. Simplex. & \color{Red} \\
\color{Red} XVII & B & \color{Red} XII & 21 & & \color{Red} \\
\color{Red} VI & C & \color{Red} XI & 22 & \color{Red} Octava Sancte Marie. Simplex. \\
& & & & \multicolumn{2}{l}{\quad \color{Red} Tymothei et Symphoriani} \\
& & & & \multicolumn{2}{l}{\quad \color{Red} Martyrum. Mem.} \\
\color{Red}  & D & \color{Red} X & 23 & & \color{Red} \\
\color{Red} XIIII & E & \color{Red} IX & 24 & Bartholomei Apostoli. \color{Red} Semid. & \color{Red} \\
\color{Red} III & F & \color{Red} VIII & 25 &  & \color{Red} \\
\color{Red}  & G & \color{Red} VII & 26 &  & \color{Red} \\
\color{Red} XI & \color{Red} A & \color{Red} VI & 27 & \color{Red} Rufi Martyris. Mem. \\
\color{Red} XIX & B & \color{Red} V & 28 & \color{Red} Augustini Epis. et Conf. Totum Dup. & \color{Red} \\
\color{Red}  & C & \color{Red} IIII & 29 & \multicolumn{2}{l}{Decollatio S. Iohannis Bapt. Simplex.} \\
 &  & &  & \quad Sabine Martyris. Mem. \\
\color{Red} VIII & D & \color{Red} III & 30 & \color{Red} Felicis, et Adaucti Martyrum. Mem. \\
\color{Red}  & E & \color{Red} II & 31 & \qquad \color{Red} Finis Quinti Embolismus. \\
\end{tabular}
\color{Red} Hore noctis. X. Diei. XIIII.
\end{center}



\newpage
\color{Red}\csection{Septembris}{\color{Red} habet dies. XXX. luna. XXX.}\color{black}

\begin{center}
\begin{tabular}{l | l | l | r | l l}
\color{Red} XVI & F & \textbf{\color{Red} K\color{Blue}L} & 1 & Egidii Abbatis. Mem. & \color{Red} \\
\color{Red} V & G & \color{Red} IIII & 2 & \qquad \color{Red} Incipit. II. Embolismus. \\
\color{Red}  & \color{Red} A & \color{Red} III & 3 & & \color{Red} \\
\color{Red} XIII & B & \color{Red} II & 4 & Octava Sancti Augustini. Simplex. & \color{Red} \\
& & & & \qquad Marcelli Martyris. Mem. \\
\color{Red} II & C & \color{Red} Non. & 5 & Anniversarium Familiarium et & \color{Red} \\
& & & & \qquad Benefactorum Ordinis Nostri. \\
\color{Red}  & D & \color{Red} VIII & 6 & & \color{Red} \\
\color{Red} X & E & \color{Red} VII & 7 & & \color{Red} \\
\color{Red}  & F & \color{Red} VI & 8 & \multicolumn{2}{l}{\color{Red} Nativitas S. Marie Virg. Tot. Dup.} \\
\color{Red} XVIII & G & \color{Red} V & 9 & Gorgonii Martyris. \color{Red} Mem. \\
\color{Red} VII & \color{Red} A & \color{Red} IIII & 10 & & \color{Red} \\
\color{Red}  & B & \color{Red} III & 11 & \color{Red} Prothi et Iacinti Martyrum. Mem. \\
\color{Red} XV & C & \color{Red} II & 12 & & \color{Red} \\
\color{Red} IIII & D & \color{Red} Idus & 13 & & \color{Red} \\
\color{Red}  & E & \color{Red} XVIII & 14 & Exaltatio Sancte Crucis. Semid. \\
& & & & \multicolumn{2}{l}{\quad Cornelii, et Cipriani Mart. Mem.} \\
& & & & \qquad Equinoctium Autumpnale. \color{Red} \\
\color{Red} XII & F & \color{Red} XVII & 15 & Octava Sancte Marie. Simplex. & \color{Red} \\
& & & & \quad Nichomedis Martyris. Mem. \\
\color{Red} I & G & \color{Red} XVI & 16 & Eufemie Virginis et Mart. III. Lect. & \color{Red} \\
\color{Red}  & \color{Red} A & \color{Red} XV & 17 & Lamberti Episcopi et Mart. Mem. \\
& & & & \qquad  \color{Red} Sol in Libra. \\
\color{Red} IX & B & \color{Red} XIIII & 18 & & \color{Red} \\
\color{Red}  & C & \color{Red} XIII & 19 & & \color{Red} \\
\color{Red} XVII & D & \color{Red} XII & 20 & \color{Red} \qquad \color{Red} Vigilia. \\
\color{Red} VI & E & \color{Red} XI & 21 & \multicolumn{2}{l}{Matthei Apos. et Evan. Semid.} \color{Red} \\
\color{Red}  & F & \color{Red} X & 22 & \multicolumn{2}{l}{Mauritii Socio. eius Mart. Simplex.} \color{Red} \\
\color{Red} XIIII & G & \color{Red} IX & 23 &  & \color{Red} \\
\color{Red} III & \color{Red} A & \color{Red} VIII & 24 & & \color{Red} \\
\color{Red}  & B & \color{Red} VII & 25 & & \color{Red} \\
\color{Red} XI & C & \color{Red} VI & 26 & & \color{Red} \\
\color{Red} XIX & D & \color{Red} V & 27 & \color{Red} Cosme et Damiani Mart. Simplex. \\
\color{Red}  & E & \color{Red} IIII & 28 & & \color{Red} \\
\color{Red} VIII & F & \color{Red} III & 29 & \color{Red} Michaelis Archangeli. Dup. & \color{Red} \\
\color{Red}  & G & \color{Red} II & 30 & Iheronimi Pres. et Conf. Simplex. \\
& & & & \qquad  \color{Red} Finis Secundi Embolismus. \\
\end{tabular}
\newline \color{Red} Hore noctis. XII. Diei. XII.
\end{center}



\newpage
\color{Red}\csection{Octobris}{\color{Red} habet dies. XXXI. luna. XXIX.}\color{black}

\begin{center}
\begin{tabular}{l | l | l | r | l l}
\color{Red} XVI & \color{Red} A & \textbf{\color{Red} K\color{Blue}L} & 1 & \color{Red} Remigii Episcopi et Conf. III. Lect. \\
\color{Red} V & B & \color{Red} VI & 2 & \color{Red} Leodegarii Episcopi et Mart. Mem. \\
\color{Red} XIII & C & \color{Red} V & 3 & & \color{Red} \\
\color{Red} II & D & \color{Red} IIII & 4 & Francisci Confessoris. Simplex. & \color{Red} \\
\color{Red}  & E & \color{Red} III & 5 & & \color{Red} \\
\color{Red} X & F & \color{Red} II & 6 & & \color{Red} \\
\color{Red}  & G & \color{Red} Non. & 7 & Marci Pape et Conf. III. Lect. \\
 &  & & & \multicolumn{2}{l}{\quad Sergii, et Bachi, Marcelli, et} \\
 &  & & & \multicolumn{2}{l}{\quad Apulei. Martyrum. Mem.} \\
\color{Red} XVIII & \color{Red} A & \color{Red} VIII & 8 & & \color{Red} \\
\color{Red} VII & B & \color{Red} VII & 9 & \multicolumn{2}{l}{\color{Red} Dionisii Socio. eius Mart. Simplex.} \\
\color{Red}  & C & \color{Red} VI & 10 & \multicolumn{2}{l}{\color{Red} Anni. Omnium Fratrum Ord. Nos.} \\
\color{Red} XV & D & \color{Red} V & 11 & & \color{Red} \\
\color{Red} IIII & E & \color{Red} IIII & 12 & & \color{Red} \\
\color{Red}  & F & \color{Red} III & 13 &  & \color{Red} \\
\color{Red} XII & G & \color{Red} II & 14 & Calyxti Pape et Martyris. Mem. & \color{Red} \\
\color{Red} I & \color{Red} A & \color{Red} Idus & 15 & & \color{Red} \\
\color{Red}  & B & \color{Red} XVII & 16 &  & \color{Red} \\
\color{Red} IX & C & \color{Red} XVI & 17 &  & \color{Red} \\
\color{Red}  & D & \color{Red} XV & 18 & Luce Evangeliste. Semid. & \\
& & & & \qquad \color{Red} Sol in Scorpione. \\
\color{Red} XVII & E & \color{Red} XIIII & 19 & & \color{Red} \\
\color{Red} VI & F & \color{Red} XIII & 20 & & \color{Red} \\
\color{Red}  & G & \color{Red} XII & 21 & \multicolumn{2}{l}{\color{Red} XI. Milium Virg. et Mart. Mem.} \\
\color{Red} XIIII & \color{Red} A & \color{Red} XI & 22 & & \color{Red} \\
\color{Red} III & B & \color{Red} X & 23 &  & \color{Red} \\
\color{Red}  & C & \color{Red} IX & 24 & & \color{Red} \\
\color{Red} XI & D & \color{Red} VIII & 25 & Crispini et Crispiani Mart. Mem. & \color{Red} \\
\color{Red} XIX & E & \color{Red} VII & 26 & & \color{Red} \\
\color{Red}  & F & \color{Red} VI & 27 &  \qquad \color{Red} Vigilia. \\
\color{Red} VIII & G & \color{Red} V & 28 & \color{Red} Symonis et Iude Apos. Semid. & \color{Red} \\
\color{Red}  & \color{Red} A & \color{Red} IIII & 29 & & \color{Red} \\
\color{Red} XII & B & \color{Red} III & 30 &  & \color{Red} \\
\color{Red} V & C & \color{Red} II & 31 & Quintini Martyris. Mem. & \\
& & & & \qquad Vigilia. \color{Red} \\
\end{tabular}
\newline \color{Red} Hore noctis. XIIII. Diei. X.
\end{center}



\newpage
\color{Red}\csection{Novembris}{\color{Red} habet dies. XXX. luna. XXIX.}\color{black}

\begin{center}
\begin{tabular}{l | l | l | r | l l}
\color{Red}  & D & \textbf{\color{Red} K\color{Blue}L} & 1 & Fest. Omnium Sanctorum. Tot. Dup. & \color{Red} \\
\color{Red} XIII & E & \color{Red} IIII & 2 & Comm. Omni. Fidelium Defunc. & \\
& & & & \qquad \color{Red} Incipit. V. Embolismus. \\
\color{Red} II & F & \color{Red} III & 3 & & \color{Red} \\
\color{Red}  & G & \color{Red} II & 4 & & \color{Red} \\
\color{Red} X & \color{Red} A & \color{Red} Non. & 5 & & \color{Red} \\
\color{Red}  & B & \color{Red} VIII & 6 & & \color{Red} \\
\color{Red} XVIII & C & \color{Red} VII & 7 & & \color{Red} \\
\color{Red} VII & D & \color{Red} VI & 8 & XL. Coronatorum Mart. III. Lect. & \color{Red} \\
\color{Red}  & E & \color{Red} V & 9 & Theodori Martyris. III. Lect. & \color{Red} \\
\color{Red} XV & F & \color{Red} IIII & 10 & & \color{Red} \\
\color{Red} IIII & G & \color{Red} III & 11 & \multicolumn{2}{l}{\color{Red} Martini Episcopi et Conf. Semid. } \\
& & & & \quad \color{Red} Menne Martyris. Mem. \color{Red} \\
\color{Red}  & \color{Red} A & \color{Red} II & 12 &  & \color{Red} \\
\color{Red} XII & B & \color{Red} Idus & 13 & Brictii Episcopi et Conf. Mem. & \color{Red} \\
\color{Red} I & C & \color{Red} XVIII & 14 & & \color{Red} \\
\color{Red}  & D & \color{Red} XVII & 15 & & \color{Red} \\
\color{Red} IX & E & \color{Red} XVI & 16 & & \color{Red} \\
\color{Red}  & F & \color{Red} XV & 17 & \qquad \color{Red} Sol in Sagittario. \\
\color{Red} XVII & G & \color{Red} XIIII & 18 & Octava Sancti Martini. \color{Red} Simplex. \\
\color{Red} VI & \color{Red} A & \color{Red} XIII & 19 & Sancte Elysabeth. Mem. & \color{Red} \\
\color{Red}  & B & \color{Red} XII & 20 & & \color{Red} \\
\color{Red} XIIII & C & \color{Red} XI & 21 & & \color{Red} \\
\color{Red} III & D & \color{Red} X & 22 & \color{Red} Cecilie Virginis et Mart. Simplex. & \color{Red} \\
\color{Red}  & E & \color{Red} IX & 23 & Clementis Pape et Mart. Simplex. & \color{Red} \\
\color{Red} XI & F & \color{Red} VIII & 24 & Crisogoni Mart. Mem. & \color{Red} \\
\color{Red} XIX & G & \color{Red} VII & 25 & \color{Red} Katherine Virginis et Mart. Semid. \\
\color{Red}  & \color{Red} A & \color{Red} VI & 26 & & \color{Red} \\
\color{Red} VIII & B & \color{Red} V & 27 & Vitalis et Agricole Mart. Mem. & \color{Red} \\
\color{Red}  & C & \color{Red} IIII & 28 & & \color{Red} \\
\color{Red} XVI & D & \color{Red} III & 29 & \color{Red} Saturnini Martyris. Mem. & \\
& & & & \qquad \color{Red} Vigilia. \\
\color{Red} V & E & \color{Red} II & 30 & Andree Apostoli. Semid. & \color{Red} \\
\end{tabular}
\newline \color{Red} Hore noctis. XVI. Diei. VIII.
\end{center}



%pp043
\newpage
\color{Red}\csection{Decembris}{\color{Red} habet dies. XXXI. luna. XXX.}\color{black}

\begin{center}
\begin{tabular}{l | l | l | r | l l}
\color{Red} XIII & F & \textbf{\color{Red} K\color{Blue}L} & 1 & \qquad \color{Red} Finis. V. Embolismus. \\
\color{Red} II & G & \color{Red} IIII & 2 & \qquad \color{Red} Incipit Primus Embolismus. \\
\color{Red}  & \color{Red} A & \color{Red} III & 3 & & \color{Red} \\
\color{Red} X & B & \color{Red} II & 4 & & \color{Red} \\
\color{Red}  & C & \color{Red} Non. & 5 & & \color{Red} \\
\color{Red} XVIII & D & \color{Red} VIII & 6 & \color{Red} Nicolai Episcopi. Semid. & \color{Red} \\
\color{Red} VII & E & \color{Red} VII & 7 & \color{Red} Octava Sancti Andree. Mem. \\
\color{Red}  & F & \color{Red} VI & 8 & & \color{Red} \\
\color{Red} XV & G & \color{Red} V & 9 &  & \color{Red} \\
\color{Red} IIII & \color{Red} A & \color{Red} IIII & 10 &  & \color{Red} \\
\color{Red}  & B & \color{Red} III & 11 & Damasi Pape et Confess. Mem. & \color{Red} \\
\color{Red} XII & C & \color{Red} II & 12 &  & \color{Red} \\
\color{Red} I & D & \color{Red} Idus & 13 & Luce Virginis et Martyris. Simplex. & \color{Red} \\
\color{Red}  & E & \color{Red} XIX & 14 &  & \color{Red} \\
\color{Red} IX & F & \color{Red} XVIII & 15 & \qquad \color{Red} Solsticium Hyemale. \\
\color{Red}  & G & \color{Red} XVII & 16 & & \color{Red} \\
\color{Red} XVII & \color{Red} A & \color{Red} XVI & 17 & \qquad O Sapiencia. \color{Red} \\
\color{Red} VI & B & \color{Red} XV & 18 & \qquad \color{Red} Sol in Capricorno. \\
\color{Red}  & C & \color{Red} XIIII & 19 & & \color{Red} \\
\color{Red} XIIII & D & \color{Red} XIII & 20 & & \color{Red} \\
\color{Red} III & E & \color{Red} XII & 21 & Thome Apostoli. Semid. & \color{Red} \\
\color{Red}  & F & \color{Red} XI & 22 & & \color{Red} \\
\color{Red} XI & G & \color{Red} X & 23 & & \color{Red} \\
\color{Red} XIX & \color{Red} A & \color{Red} IX & 24 & \qquad \color{Red} Vigilia. \\
\color{Red}  & B & \color{Red} VIII & 25 & \multicolumn{2}{l}{\color{Red} Nativitas D.N.I.C. Totum Dup.} \color{Red} \\
\color{Red} VIII & C & \color{Red} VII & 26 & Stephani Prothomart. Tot. Dup. & \color{Red} \\
\color{Red}  & D & \color{Red} VI & 27 & \multicolumn{2}{l}{\color{Red} Iohannis Apostoli et Evan. Tot. Dup.} \color{Red} \\
\color{Red} XVI & E & \color{Red} V & 28 & Sanctorum Innocentium. Simplex. & \color{Red} \\
\color{Red} V & F & \color{Red} IIII & 29 & \color{Red} Thome Episcopi et Mart. Simplex. & \color{Red} \\
\color{Red}  & G & \color{Red} III & 30 &  & \color{Red} \\
\color{Red} XIII & \color{Red} A & \color{Red} II & 31 & Silvestri Pape et Conf. Simplex. & \\
& & & & \qquad \color{Red} Finis. I. Embolismus. \\
\end{tabular}
\newline \color{Red} Hore noctis. XVIII. Diei. VI.
\end{center}

%pp043
\newpage

\begin{multicols*}{2}

\ubsection{Rubrice}{rubrice-collectarium}

{\color{Red} \ubsubsection{De Varietatibus Terminationum in Orationibus}{de-varietatibus-terminationum-in-orationibus}}
\lettrine[lines=2]{\zallmancaps \color{Red} N}{}\ul{\textsc{otandum} quod orationes que ad Patrem sine persone alterius expressione diriguntur, sic concluduntur.}
\newline Per dominum nostrum Ihesum Christum filium tuum, Qui tecum vivit et regnat in unitate spiritus sancti deus: per omnia secula seculorum. \ul{Ut est illa de Beata Virgine,} Concede nos, \ul{et consimiles. Similiter et illa in qua fit mentio de Trinitate,} \hyperlink{omnipotens-trinitas}{Omnipotens sempiterne deus qui dedisti.}
\newline \ul{Item orationes que ad patrem cum expressione filii diriguntur: quedam sic concluduntur.}
\newline Per eumdem dominum nostrum Ihesum Christum filium tuum, Qui tecum. \ul{et cetera. Ut illa in die Natalis,} \hyperlink{concede-nativitas}{Concede quesumus omnipotens deus ut nos unigeniti.} \ul{Quedam per.} 
\newline Qui tecum vivit, \ul{etc ut prius sicut est illa in die Circumcisionis.} \hyperlink{deus-nobis-circumcisionis}{Deus qui nobis nati,} \ul{et consimiles. Sunt et alie orationes que ad patrem diriguntur, in fine quarum fit mentio de filio, et hoc vel per rectum hoc modo.}
\newline Ihesus Christus dominus noster, Qui tecum vivit, \ul{etc ut supra sicut in post communione in vigilia Ascensionis,} Tribue. \ul{Vel per obliquum sic.}
\newline Dominum nostrum Ihesum Christum filium tuum, Qui tecum, \ul{etc, sicut inoratione de sancto Stephano,} Da nobis quesumus domine imitari. \ul{etc, et in illa,} Deus qui salutis eterne, \ul{et consimilibus. Orationes vero que ad filium diriguntur: sic concludantur.}
\newline Qui vivis et regnas cum deo patre, in unitate spiritus sancti deus per omnia secula seculorum. \ul{Ut est illa, in prima Dominica Adventus,} Excita quesumus, \ul{et consimiles. Orationes vero que ad patrem vel filium cum expressione spiritus sancti diriguntur: sic concluduntur in fine}
\newline In unitate eiusdem spiritus sancti deus, Per omnia secula seculorum. \ul{Orationes autem que ad ipsam trinitatem diriguntur sic concluduntur,}
\newline Qui vivis et regnas deus: Per omnia secula seculorum. \ul{Ut est illa,} Placeat tibi sancta trinitas.
\newline \ul{Quando vero post aliquam memoriam alia facienda est, prior oratio si ad patrem dirigitur, terminetur sic}
\newline Per Christum dominum nostrum. \ul{Vel,}
\newline Per eundem Christum dominum nostrum. \ul{Vel,}
\newline Christum dominum nostrum. \ul{Horum exempla patent supra. Si autem ad filium oratio dirigitur: concludatur sic,}
\newline Qui vivis et regnas per omnia secula seculorum.
{\color{Red} \ubsubsection{De Confiteor Quomodo Sit Dicenda}{de-confiteor-quomodo-sit-dicenda-collectarium}}
\newline \ul{Quandocumque dicitur,} Confiteor. \ul{in choro a prelato vel ab incipiente officium, vel extra ab illo qui incepit horas, vel a sacerdote Missam celebraturo, sive in conventu, sive in Missa privata, eciam uno solo astante dicatur sic.}
\lettrine[lines=2]{\zallmancaps \color{Blue} C}{onfiteor} deo et beate Marie et omnibus sanctis, et vobis fratres, quia peccavi nimis cogitatione, locutione, opere et omissione, mea culpa precor vos orate pro me.
\newline \ul{Ille vero sive illi qui respondent in choro vel extra, vel in Missa conventuali, vel privata semper dicere debent eciam prelato.}
\lettrine[lines=2]{\zallmancaps \color{Red} M}{isereatur} tui omnipotens deus, et dimittat tibi omnia peccata tua, liberet te ab omni malo salvet et confirmet in omni opere bono, et perducat ad vitam eternam. \R Amen.
\newline \ul{Dicto miseratur subiungatur ab eodem, vel ab eisdem.}
\lettrine[lines=2]{\zallmancaps \color{Blue} C}{onfiteor} deo et beate Marie et omnibus sanctis et tibi pater, quia peccavi nimis, cogitatione, locutione, opere, et omissione, mea culpa, precor te ora pro me.
\newline \ul{Et ille qui primo dixit} Confiteor. \ul{subiungat, eciam si unus solus sit cui dicitur.}
\lettrine[lines=2]{\zallmancaps \color{Red} M}{isereatur} vestri omnipotens deus, et dimittat vobis omnia peccata vestra, liberet vos ab omni malo salvet, \ul{etc.} \R Amen.
\newline \ul{Quando vero qui incipit horas non est sacerdos, respondeatur,} Confiteor. \ul{etc, et tibi frater, etc. Quando vero frater solus dicit Primam vel Completorium: nec habet cum quo dicat} Confiteor: \ul{dicat sic.}
\lettrine[lines=2]{\zallmancaps \color{Blue} C}{onfiteor} deo et beate Marie et omnibus sanctis: quia peccavi nimis cogitatione, locutione, opere et omissione, mea culpa precor beatam Mariam et omnes sanctos orare pro me. \ul{Postea subiungat.}
\lettrine[lines=2]{\zallmancaps \color{Red} M}{isereatur} michi omnipotens deus et dimittat michi omnia peccata mea, liberet me, \ul{etc.}
{\color{Red} \ubsubsection{De Modo Dicendi Preces}{de-modo-dicendi-preces-collectarium}}
\newline \ul{Finito versiculo, Vel antiphona: cantore incipiente chorus presequatur.}
\lettrine[lines=2]{\zallmancaps \color{Red} K}{yrie} eleyson. Christe eleyson. Kyrie eleyson. \psalm{pater} \ul{Secreto. Deinde dicatur.}
\lettrine[lines=2]{\zallmancaps \color{Blue} E}{t} ne nos inducas in temptationem. \R Set libera nos a malo. \V In pace in idipsum. \R Dormiam et requiescam. \V Vivet anima mea et laudabit te. \R Et iudicia tua adiuvabunt me. \V
\lettrine[lines=2]{\zallmancaps \color{Red} E}{rravi} sicut ovis que periit. \R Quere servum tuum domine quia mandata tua non sum oblitus. \psalm{credo} \ul{Secreto. Deinde.}
\lettrine[lines=2]{\zallmancaps \color{Blue} C}{arnis} resurrectionem. \R Vitam eternam amen. \V Dignare domine nocta ista. {\color{Red} Vel.} die isto. \R sine peccato nos custodire.
\newline \ul{Preces omnes que dicuntur cum nota dicende sunt in eadem voce cum oratione que postea est dicenda.}
{\color{Red} \ubsubsection{Quomodo Terminande Sunt Hore a Sacerdote}{quomodo-terminande-sunt-hore-a-sacerdote-collectarium}}
\lettrine[lines=2]{\zallmancaps \color{Red} B}{enedicamus} domino. \R Deo gratias. {\color{Blue} R}equiescant in pace. \R Amen.% horas vel defunctorum
\lettrine[lines=2]{\zallmancaps \color{Blue} F}{idelium} anime per misericordiam dei requiscant in pace. \R Amen.
\newline \ul{Sine nota sonore et humili voce dicatur, et eodem modo respondeatur. Amen.}
%pp044
{\color{Red} \ubsubsection{Benedictiones Lectionum in Matutinis}{benedictiones-lectionum-in-matutinis-collectarium}}
\newline {\color{Red} In Primo Nocturno.}
\newline {\color{Red} B}enedictione perpetua benedicat nos pater eternis. Amen.
\newline {\color{Blue} U}nigenitus dei filius nos benedicere et adiuvare dignetur.
\newline {\color{Red} S}piritus sancti gratia illuminare dignetur sensus et corda nostra.
\newline {\color{Red} In Secundo Nocturno.}
\newline {\color{Blue} D}eus pater omnipotens sit nobis propitius et clemens.
\newline {\color{Red} A}d gaudia paradisi perducat nos misericordia Christi.
\newline {\color{Blue} I}gnem sui amoris accendat deus in cordibus nostris.
\newline {\color{Red} In Tercio Nocturno.}
\newline {\color{Red} I}lle nos benedicat qui sine fine vivit et regnat.
\newline {\color{Red} Quando omelia legitur.}
\newline {\color{Blue} E}vangelica lectio sit nobis salus et protectio.
\newline {\color{Red} D}ivinum auxilium maneat semper nobiscum.
\newline {\color{Blue} A}d societatem civium supernorum perducat nos rex angelorum.
\newline {\color{Red} In Natali Domini Octava Benedictio.}
\newline {\color{Red} P}er evangelica dicta deleantur nostra delicta. {\color{Red} Nona.}
\newline {\color{Blue} I}ntellectum sancti evangelii aperiat nobis virtus spiritus sancti.
{\color{Red} \ubsubsection{Benedictiones in Festis Beate Virginis}{benedictiones-in-festis-beate-virginis-collectarium}}
\newline {\color{Red} In Primo Nocturno.}
\newline {\color{Red} A}lma virgo virginum intercedat pro nobis ad dominum.
\newline {\color{Red} S}ancta dei genitrix sit nobis auxiliatrix.
\newline {\color{Blue} N}os cum prole pia benedicat virgo Maria.
\newline {\color{Red} In Secundo Nocturno.}
\newline {\color{Red} A}b hoste maligno eripiat nos dei genitrix virgo.
\newline {\color{Blue} I}hesus Marie filius sit nobis clemens et propitius.
\newline {\color{Red} S}tella Maria maris succurre piissima nobis.
\newline {\color{Red} In Tercio Nocturno.}
\lettrine[lines=2]{\zallmancaps \color{Blue} O}{ret} voce pia pro nobis virgo Maria. {\color{Red} Vel.} Evangelica lectio{\color{Red} :} \ul{Si fuerit omelia.}
\newline {\color{Red} I}n tribulatione et angusita subveniat nobis pia virgo Maria.
\newline {\color{Blue} A}d societatem civium supernorum perducat nos regina angelorum.
\newline \ul{In Sabbatis dicantur tres ultime quando de ea fit officium in conventu.}

%skip

%pp050
{\color{Red} \ubsubsection{De Disciplinis Post Completorium}{de-disciplinis-post-completorium-collectarium}}
\newline \ul{Dicto} Confiteor, et Misereatur, \ul{et psalmo,} \psalm{50} et \psalm{pater} \ul{Subiungatur,} Et ne nos. \V Salvos fac servos tuos. Dominus vobiscum. Oremus. {\color{Red} Oratio.}
\lettrine[lines=2]{\zallmancaps \color{Red} D}{eus} cui proprium est misereri smper et parcere, suscipe deprecationem nostram, et quos delictorum catena constringit, miseratio tue pietatis absolvat. Per Christum.
% {\color{Red} \ubsubsection{De Officio in Receptione Novitiorum}{-collectarium}}
% {\color{Red} \ubsubsection{De Officio in Electionibus}{-collectarium}}
% {\color{Red} \ubsubsection{De Oratione Pro Capitulo Generali, et Pro Pergentibus Ad Illud}{-collectarium}}
% {\color{Red} \ubsubsection{De Benedictione Itinerantium}{benedictione-itinerantiu--collectarium}}
% \lettrine[lines=2]{\zallmancaps \color{Red} P}{}\ul{\textsc{ro} fratribus egredientibus dicatur Ps.} \psalm{120} \ul{cum} Gloria patri. \ul{Deinde} Kyrie eleyson. Christe eleyson. Kyrie eleyson. Pater noster. \ul{Postea,} Et ne nos. \V Salvos fac. Dominus vobiscum. Oremus. {\color{Red} Oratio.} {\color{Blue} A}desto domine. Ut supra.
% \newline \ul{Circa venientes de via idem observetur: excepto quo pro eis dicatur Ps.} \psalm{122} \ul{et Oratio.}

{\color{Red} \ubsubsection{De Preciosa}{de-preciosa-collectarium}}
\newline \ul{Finitis que de Kalendario, luna, et martyologio pronuncianda sunt: subiungatur.} {\color{Red} Versus.}
\lettrine[lines=2]{\zallmancaps \color{Red} P}{reciosa} est in conspectu domini. \R Mors sanctorum eius. {\color{Red} Oratio.}
\lettrine[lines=2]{\zallmancaps \color{Blue} S}{ancta} Maria et omnes sancti intercedant pro nobis ad dominum, ut mereamur adiuvari ab eo. Qui vivit et regnat per omnia secula seculorum. \R Amen.
\newline \ul{Deinde dicatur ter.} \V Deus in adiutorium meum intende. \R Domine ad adiuvandum me festina.
\newline \ul{Quo ter dicto subiungatur.} Gloria patri et filio et spiritui sancto. Sicut erat in principio et nunc et semper et in secula seculorum amen. Kyrie eleyson. Christe eleyson. Kyrie eleyson. \psalm{pater} Et ne nos inducas  in temptationem. Sed libera nos a malo. \V Respice domine in servos tuos et in opera tua, et dirige filios eorum. \R Et sit splendor domini dei nostri super nos, et opera manuum nostrarum dirige super nos, et opus manuum nostrarum dirige.
\newline \ul{Deinde sequitur oratio dicenda eo modo, quo dicuntur orationes ad horas.}
\lettrine[lines=2]{\zallmancaps \color{Red} D}{irigere} et sanctificare digneris domine sancte pater omnipotens, eterne deus, hodie corda et corpora nostra in lege tua et in operibus mandatorum tuorum ut hic et in eternum te auxiliante semper salvi esse mereamur: Per Christum dominum nostrum amen.
\newline \ul{Postea dicto,} Iube. \ul{Si de evangelio legendum fuerit, detur benedictio.}
\newline Divinum auxilium maneat semper nobiscum. Amen.
\newline \ul{Si vero de constitutionibus: detur benedictio.}
\newline Regularibus disciplinis instruat nos magister celestis. Amen.
\newline \ul{Data benedictione sequitur lectio de evangelio, vel de constitutionibus: Fratres tamen qui extra conventu sunt: si evangelium proprium in promptu non habuerint: dicant.} {\color{Red} Secundum Iohannem.}
\lettrine[lines=2]{\zallmancaps \color{Blue} I}{n} principio erat verbum. Et verbum erat apud deum, et deus erat verbum. Hoc erat in principio apud deum. Omnia per ipsum facta sunt, et sine ipso factum est nichil. Tu autem.
\newline \ul{De constitutionibus vero dicatur.} {\color{Red} Lectio.}
\lettrine[lines=2]{\zallmancaps \color{Red} Q}{uoniam} ex precepto regule iubemur habere cor unum et animam unam in domino, iustum est ut qui sub una regula et unius professionis voto vivimus, uniformes in observantiis canonice religionis inveniamur. Tu autem.
\newline \ul{Finita lectione sequitur,} Commemoratio fratrum, familiarum, benefactorum, defunctorum ordinis nostri. - \ul{Prelatus} - Requiescant in pace. Amen. \ul{Deinde sequitur psalmus.}
\lettrine[lines=2]{\zallmancaps \color{Blue} L}{audate} dominum omnes gentes laudate eum omnes populi. Quoniam, etc. Gloria patri. Sicut erat. \V Ostende nobis domine misericordiam tuam. \R Et salutare tuum da nobis. Dominus vobiscum. Oremus. {\color{Red} Oratio.}
\lettrine[lines=2]{\zallmancaps \color{Red} A}{ctiones} nostras quesumus domine aspirando preveni, et adiuvando prosequere: ut cuncta nostra operatio a te semper incipiat et per te cepta finiatur. Per Christum dominum nostrum.
%pp051
\newline \ul{Postea dicatur} Adiutorium nostrum in nomine domini. Qui fecit celum et terram.
\newline \ul{In capitulo dicto} Requiescant in pace.\ul{et responso} Amen. \ul{dicat qui tenet capitulum.} Benedicite. \ul{Et responso.} Dominus. \ul{et recitatis benedicus dicatur.}
\lettrine[lines=2]{\zallmancaps \color{Red} R}{etribue} dignare domine omnibus nobis bona facientibus propter nomen sanctum tuum vitam eternam. \V Amen.
\newline \ul{Postea sequitur psalmus hoc modo.} \psalm{122} \ul{Et terminetur,} per Gloria. \ul{Deinde sequitur Ps.} \psalm{129} \ul{Dicendus eodem modo: et terminandus,} per Requiem. \ul{Postea sequitur.} Kyrie eleyson. Christe eleyson. Kyrie eleyson. \ul{ut supra.} \psalm{pater} Et ne nos. \ul{Deinde.} \V Oremus pro domino papa. \R Dominus conservet eum et vivificet eum, et beatum faciat eum in terra, et non tradat eum in animam inimicorum eius. \V Salvos fac servos tuos et ancillas tuas. \R Deus meus. Requiescant in pace. \R Amen. Dominus vobiscum. Oremus. {\color{Red} Oratio.}
\lettrine[lines=2]{\zallmancaps \color{Blue} O}{mnipotens} sempiterne deus qui facis mirabilia solus, pretende super famulum tuum papam nostram, et super cunctas congregationes illi comissas spiritum gratie salutaris, et ut in veritate tibi complaceant, perpetuum eis rorem tue benedictionis infunde. {\color{Red} Oratio.}
\lettrine[lines=2]{\zallmancaps \color{Red} P}{retende} domine famulis et familiabus tuis dexteram celestis auxilii, ut te toto corde perquirant, et que digne postulant assequantur: {\color{Red} Oratio.}
\lettrine[lines=2]{\zallmancaps \color{Blue} F}{idelium} deus omnium conditor \ul{et cetera.}
\newline \ul{Predicte tres orationes dicantur eo modo quo dicuntur Orationes ad horas et terminentur sub uno.} Qui vivis et regnas per omnia secula seculorum. Amen.
% {\color{Red} \ubsubsection{De Benedictione Mense}{de-benedictione-mense-collectarium}}
{\color{Red} \ubsubsection{De Collatione}{de-collatione-collectarium}}
\newline \ul{Tempore ieiunii in collatione detur lectori benedictio.} Noctem quietam et finem perfectum, tribuat nobis omnipotens et misericors dominus. \ul{Infra lectionem dicto:} Benedicite. \ul{Detur superiori potum benedico.} Largitor omnium bonorum, {\color{Red} \grecross } benedicat potum servorum suorum. \ul{Dicto} Tu autem. \ul{In fine lectionis subiugatur.} \V Adiutorium nostrum. %superiori from sr
\newline \ul{Quando vero non fit collatio: detur in choro benedictio.} Noctem quietam etc. \ul{Deinde subiugatur lectio.} Fratres sobrii. \ul{Qua finita et responso.} Deo gratias. \ul{subiungat qui prior est, vel si ipse defuerit ebdomadarius.} Adiutorium nostrum. \ul{Deinde dicto.} Pater noster. \ul{dicatur,} Confiteor. \ul{Postea Completorium inchoetur. Post Completorium detur benedictio sonore et humili voce}
\lettrine[lines=2]{\zallmancaps \color{Red} B}{enedictio} dei omnipotenteis, patris, et filii et spritu sancti descendat super vos et maneat semper. Amen.

%skip

%pp070
{\color{Red} \ubsubsection{Ante Completorium}{ante-completorium-psalterium}}
\newline {\color{Red} Lectio.}
\lettrine[lines=2]{\zallmancaps \color{Blue} F}{rates}, sobrii estote et vigilate, quia adversarius vester diabolus tanquam leo rugiens circuit querens quem devoret: cui resistite fortes in die. Tu autem.

% \newline {\color{Red} .}
% \newline {\color{Red} .}
% \newline {\color{Red} .}
% \newline {\color{Red} .}

% \lettrine[lines=2]{\zallmancaps \color{Red} O}{}
% \lettrine[lines=2]{\zallmancaps \color{Blue} I}{n}
% \lettrine[lines=2]{\zallmancaps \color{Red} O}{}
% \lettrine[lines=2]{\zallmancaps \color{Blue} I}{n}
% \lettrine[lines=2]{\zallmancaps \color{Red} O}{}
% \lettrine[lines=2]{\zallmancaps \color{Blue} I}{n}
% \lettrine[lines=2]{\zallmancaps \color{Red} O}{}

% {\color{Red} \ubsubsection{Ad Vesperas}{-collectarium}}
% {\color{Red} \ubsubsection{Ad Vesperas}{-collectarium}}

\end{multicols*}

%pp071
\newpage

\ubchapter{Psalterium}

\thispagestyle{fancy}
% \fancyhead[RO,LE]{}

\begin{multicols*}{2}

\ubsection{Dominica ad Matutinas}{dominica-ad-matutinas}

\noindent \ul{Dominica prima in Adventu et in sequentibus. In primo nocturno.}
\newline {\color{Red} Ant.} Scientes.
\newline \ul{Dominica prima post Octavas Epiphanie et in sequentibus usque ad Dominicam in Passione.}
\newline {\color{Red} Ant.} Servite.
\newline \ul{Dominica in Passione et sequenti.}
\newline {\color{Red} Ant.} Quid molesti.
\newline \ul{Dominica prima post Trinitatem et in sequentibus usque ad Adventum.}
\newline {\color{Red} Ant.} Pro fidei.
\fullpsalm{1}
\fullpsalm{2}
\fullpsalm{3}
\fullpsalm{4}
\fullpsalm{5}
\fullpsalm{6}
\newline {\color{Red} Ant.} Servite domino in timore, et exultate ei cum tremore.
\newline {\color{Red} Ant.} Domine.
\newline {\color{Red} Ant.} Pro fidei meritis vocitatur iure beatus legem qui domini meditatur nocte dieque.
\newline {\color{Red} Ant.} Iuste deus.
\fullpsalm{7}
\fullpsalm{8}
\fullpsalm{9}
%pp072
\fullpsalm{10}
\newline {\color{Red} Ant.} Domine deus meus in te speravi.
\newline {\color{Red} Ant.} Tu Domine.
\newline {\color{Red} Ant.} Iuste deus iudex fortis patiensque benignus in te sperantes muni miserando fideles.
\newline {\color{Red} Ant.} Surge.
\fullpsalm{11}
\fullpsalm{12}
\fullpsalm{13}
\fullpsalm{14}
\newline {\color{Red} Ant.} Scientes quia hora est iam nos de somno surgere nunc enim propior est nostra salus quam cum credidimus nox precessit dies autem appropinquavit.
\newline {\color{Red} Ant.} Hora est.
\newline {\color{Red} Ant.} Quid molesti estis huic mulieri opus enim bonum operata est in me.
\newline {\color{Red} Ant.} Mittens.
\newline {\color{Red} Ant.} Tu Domine servabis nos et custodies nos.
\newline {\color{Red} Ant.} Bonorum.
\newline {\color{Red} Ant.} Surge et in eternum serva munimine sacro custodique tuos astripotens famulos.
\newline {\color{Red} Ant.} Nature.
\fullpsalm{15}
\newline {\color{Red} Ant.} Bonorum meorum non indiges in te speravi conserva me Domine.
\newline {\color{Red} Ant.} Inclina.
\newline {\color{Red} Ant.} Nature genitor conserva a morte redemptos facque tuo dignos servitio famulos.
\newline {\color{Red} Ant.} Pectora.
\fullpsalm{16}
\newline {\color{Red} Ant.} Inclina Domine aurem tuam michi et exaudi verba mea.
\newline {\color{Red} Ant.} Dominus.
\newline {\color{Red} Ant.} Pectora nostra tibi tu conditor orbis adure igne pio purgans atque cremando probans.
\newline {\color{Red} Ant.} Tu populum.
\fullpsalm{17}
\newline {\color{Red} Ant.} Hora est iam nos de somno surgere et aperti sunt oculi nostri surgere ad Christum quia lux vera est fulgens in mundo.
\newline {\color{Red} Ant.} Nox precessit.
\newline {\color{Red} Ant.} Mittens hec mulier in corpus meum hoc unguentum ad sepeliendum me fecit.
\newline {\color{Red} Ant.} Magister dicit.
\newline {\color{Red} Ant.} Dominus firmamentum meum et refugium meum.
\newline {\color{Red} Ant.} Non sunt loquele.
\newline {\color{Red} Ant.} Tu populum humilem salvasti a morte redemptor atque superba tuo colla premis iaculo.
\newline {\color{Red} Ant.} Sponsus.
\fullpsalm{18}
%pp073
\newline {\color{Red} Ant.} Non sunt loquele neque sermones quorum non audiantur voces eorum.
\newline {\color{Red} Ant.} Exaudiat.
\newline {\color{Red} Ant.} Sponsus ut e thalamo processit Christus in orbem descendens celo iure salutifero.
\newline {\color{Red} Ant.} Auxilium.
\fullpsalm{19}
\newline {\color{Red} Ant.} Exaudiat te Dominus in die tribulationis.
\newline {\color{Red} Ant.} Domine.
\newline {\color{Red} Ant.} Auxilium nobis salvator mitte salutis et tribuas vitam tempore perpetuo.
\newline {\color{Red} Ant.} Rex sine fine.
\fullpsalm{20}
\newline {\color{Red} Ant.} Nox precessit dies autem appropinquavit abiciamus ergo opera tenebrarum et induamur arma lucis sicut in die honeste ambulemus.
\newline {\color{Red} Ant.} Domine in virtute tua letabitur rex.
\newline {\color{Red} Ant.} Magister dicit tempus meum prope est apud te facio Pascha cum discipulis meis.
\newline {\color{Red} Ant.} Rex sine fine manens miseris tu parce ruinis premia concedens et tua cuncta regens.
\newline {\color{Red} In Laudibus. Ant.} Regnavit Dominus precinctus fortitudine cum decore virtutum cuius sedes parata est in eternum.
\psalm{92}
\fullpsalm{21}
\fullpsalm{22}
\fullpsalm{23}
\fullpsalm{24}
\fullpsalm{25}
\ubsection{Feria II ad Matutinas}{feria-ii-ad-matutinas}
\newline {\color{Red} Ant.} Dominus.
\fullpsalm{26}
\fullpsalm{27}
\newline {\color{Red} Ant.} Dominus defensor vite mee.
\newline {\color{Red} Ant.} Adorate.
\fullpsalm{28}
\fullpsalm{29}
\newline {\color{Red} Ant.} Adorate Dominum in aula sancta eius.
\newline {\color{Red} Ant.} In tua.
\fullpsalm{30}
%pp074
\fullpsalm{31}
\newline {\color{Red} Ant.} In tua iusticia libera me domine
\newline {\color{Red} Ant.} Rectos.
\fullpsalm{32}
\fullpsalm{33}
\newline {\color{Red} Ant.} Rectos decet collaudatio.
\newline {\color{Red} Ant.} Expugna.
\fullpsalm{34}
\fullpsalm{35}
\newline {\color{Red} Ant.} Expugna impugnantes me.
\newline {\color{Red} Ant.} Revela.
\fullpsalm{36}
\fullpsalm{37}
\newline {\color{Red} Ant.} Revela Domino viam tuam.
\newline {\color{Red} In Laudibus. Ant.} Miserere mei deus.
\psalm{50}
\newline {\color{Red} Ant.} Intellige clamorem meum domine.
\psalm{5}
\newline {\color{Red} Ant.} Deus deus meus ad te de luce vigilo.
\psalm{62}
\newline {\color{Red} Ant.} Conversus est furor tuus domine et consolatus es me.
\psalm{isaias}
\newline {\color{Red} Ant.} Laudate domine de celis.
\psalm{148}
\newline {\color{Red} Ant.} Benedictus deus Israhel.
\psalm{benedictus}
\ubsection{Feria III ad Matutinas}{feria-iii-ad-matutinas}
\newline {\color{Red} Ant.} Ut non delinquam.
\fullpsalm{38}
\fullpsalm{39}
%pp075
\newline {\color{Red} Ant.} Ut non delinquam in lingua mea.
\newline {\color{Red} Ant.} Sana domine.
\fullpsalm{40}
\fullpsalm{41}
\fullpsalm{42}
\newline {\color{Red} Ant.} Sana domine animam meam quia peccavi tibi.
\newline {\color{Red} Ant.} Eructavit.
\fullpsalm{43}
\fullpsalm{44}
\newline {\color{Red} Ant.} Eructavit cor meum verbum bonum.
\newline {\color{Red} Ant.} Adiutor.
\fullpsalm{45}
\fullpsalm{46}
\newline {\color{Red} Ant.} Adiutor in tribulationibus.
\newline {\color{Red} Ant.} Auribus.
\fullpsalm{47}
\fullpsalm{48}
\newline {\color{Red} Ant.} Auribus percipite qui habitatis orbem.
\newline {\color{Red} Ant.} Deus.
\fullpsalm{49}
\fullpsalm{50}
\fullpsalm{51}
\newline {\color{Red} Ant.} Deus deorum locutus est.
\newline {\color{Red} In Laudibus. Ant.} Secundum magnam misericordiam tuam miserere mei deus.
\psalm{50}
\newline {\color{Red} Ant.} Salutare vultus mei deus meus.
\psalm{41}
\newline {\color{Red} Ant.} Ad te de luce vigilo deus.
\psalm{62}
\newline {\color{Red} Ant.} Cunctis diebus vite nostre salvos nos fac domine.
%pp076
\psalm{ezechias}
\newline {\color{Red} Ant.} Omnes angeli eius laudate dominum de celis.
\psalm{148}
\newline {\color{Red} Ant.} Erexit Dominus nobis cornu salutis in domo David pueri sui.
\psalm{benedictus}
\ubsection{Feria IV ad Matutinas}{feria-iv-ad-matutinas}
\newline {\color{Red} Ant.} Avertet.
\fullpsalm{52}
\fullpsalm{53}
\fullpsalm{54}
\newline {\color{Red} Ant.} Avertet Dominus captivitatem plebis sue.
\newline {\color{Red} Ant.} Quoniam.
\fullpsalm{55}
\fullpsalm{56}
\newline {\color{Red} Ant.} Quoniam in te confidit anima mea.
\newline {\color{Red} Ant.} Iuste.
\fullpsalm{57}
\fullpsalm{58}
\newline {\color{Red} Ant.} Iuste iudicate filii hominum.
\newline {\color{Red} Ant.} Da nobis.
\fullpsalm{59}
\fullpsalm{60}
\newline {\color{Red} Ant.} Da nobis Domine auxilium de tribulatione.
\newline {\color{Red} Ant.} A timore.
\fullpsalm{61}
\fullpsalm{62}
\fullpsalm{63}
\fullpsalm{64}
\newline {\color{Red} Ant.} A timore inimici eripe Domine animam meam.
\newline {\color{Red} Ant.} In ecclesiis.
\fullpsalm{65}
\fullpsalm{66}
\fullpsalm{67}
%pp077
\newline {\color{Red} Ant.} In ecclesiis benedicite Domino.
\newline {\color{Red} In Laudibus. Ant.} Amplius lava me Domine ab iniusticia mea.
\psalm{50}
\newline {\color{Red} Ant.} Te decet hymnus deus in Sion.
\psalm{64}
\newline {\color{Red} Ant.} Labia mea laudabit te in vita mea deus meus.
\psalm{62}
\newline {\color{Red} Ant.} Dominus iudicabit fines terre.
\psalm{annas}
\newline {\color{Red} Ant.} Celi celorum laudate deum.
\psalm{148}
\newline {\color{Red} Ant.} Salutem ex inimicis nostris et de manu omnium qui nos oderunt.
\psalm{benedictus}
\ubsection{Feria V ad Matutinas}{feria-v-ad-matutinas}
\newline {\color{Red} Ant.} Domine deus.
\fullpsalm{68}
\fullpsalm{69}
\newline {\color{Red} Ant.} Domine deus in adiutorium meum intende.
\newline {\color{Red} Ant.} Esto michi.
\fullpsalm{70}
\fullpsalm{71}
\newline {\color{Red} Ant.} Esto michi Domine in deum protectorem.
\newline {\color{Red} Ant.} Liberasti.
\fullpsalm{72}
\fullpsalm{73}
\newline {\color{Red} Ant.} Liberasti virgam hereditatis tue.
\newline {\color{Red} Ant.} In Israhel.
\fullpsalm{74}
\fullpsalm{75}
\newline {\color{Red} Ant.} Tu es deus.
\fullpsalm{76}
%pp078
\fullpsalm{77}
\newline {\color{Red} Ant.} Tu es deus qui facis mirabilia.
\newline {\color{Red} Ant.} Propitius.
\fullpsalm{78}
\fullpsalm{79}
\newline {\color{Red} Ant.} Propitius esto peccatis nostris Domine.
\newline {\color{Red} In Laudibus. Ant.} Tibi soli peccavi Domine miserere mei.
\psalm{50}
\newline {\color{Red} Ant.} Domine refugium factus es nobis.
\psalm{89}
\newline {\color{Red} Ant.} In matutinis domine meditabor in te.
\psalm{62}
\newline {\color{Red} Ant.} In eternum dominus regnabit et ultra.
\psalm{moysis}
\newline {\color{Red} Ant.} In sanctis eius laudate deum.
\psalm{148}
\newline {\color{Red} Ant.} In sanctitate serviamus Domino et liberabit nos ab inimicis nostris.
\psalm{benedictus}
\ubsection{Feria VI ad Matutinas}{feria-vi-ad-matutinas}
\newline {\color{Red} Ant.} Exultate.
\fullpsalm{80}
\fullpsalm{81}
\newline {\color{Red} Ant.} Exultate deo adiutori nostro.
\newline {\color{Red} Ant.} Tu solus.
\fullpsalm{82}
\fullpsalm{83}
\newline {\color{Red} Ant.} Tu solus altissimus super omnem terram.
\newline {\color{Red} Ant.} Benedixisti.
\fullpsalm{84}
\fullpsalm{85}
\newline {\color{Red} Ant.} Benedixisti Domine terram tuam.
\newline {\color{Red} Ant.} Fundamenta.
\fullpsalm{86}
%pp079
\fullpsalm{87}
\newline {\color{Red} Ant.} Fundamenta eius in montibus sanctis.
\newline {\color{Red} Ant.} Benedictus.
\fullpsalm{88}
\fullpsalm{89}
\fullpsalm{90}
\fullpsalm{91}
\fullpsalm{92}
\fullpsalm{93}
\fullpsalm{94}
\newline {\color{Red} Ant.} Benedictus Dominus in eternum.
\newline {\color{Red} Ant.} Cantate.
\fullpsalm{95}
\fullpsalm{96}
\newline {\color{Red} Ant.} Cantate Domino et benedicite nomini eius.
\newline {\color{Red} In Laudibus. Ant.} Spiritu principali confirma cor meum deus.
\psalm{50}
\newline {\color{Red} Ant.} In veritate tua exaudi me Domine.
\psalm{142}
\newline {\color{Red} Ant.} Illumina Domine vultum tuum super nos.
\psalm{62}
\newline {\color{Red} Ant.} Domine audivi auditum tuum et timui.
\psalm{habacuc}
\newline {\color{Red} Ant.} In tympano et choro in cordis et organo laudate deum.
\psalm{148}
\newline {\color{Red} Ant.} Per viscera misericordie dei nostri in quibus visitavit nos oriens ex alto.
\psalm{benedictus}
\ubsection{Sabbato ad Matutinas}{sabbato-ad-matutinas}
\newline {\color{Red} Ant.} Quia mirabilia.
\fullpsalm{97}
\fullpsalm{98}
\newline {\color{Red} Ant.} Quia mirabilia fecit Dominus.
\newline {\color{Red} Ant.} Iubilate.
\fullpsalm{99}
%pp080
\fullpsalm{100}
\newline {\color{Red} Ant.} Iubilate deo omnis terra.
\newline {\color{Red} Ant.} Clamor meus.
\fullpsalm{101}
\fullpsalm{102}
\newline {\color{Red} Ant.} Clamor meus ad te veniat deus.
\newline {\color{Red} Ant.} Benedic.
\fullpsalm{103}
\fullpsalm{104}
\newline {\color{Red} Ant.} Benedic anima mea Domino.
\newline {\color{Red} Ant.} Visita nos.
\fullpsalm{105}
\fullpsalm{106}
\newline {\color{Red} Ant.} Visita nos Domine in salutari tuo.
\newline {\color{Red} Ant.} Confitebor.
\fullpsalm{107}
%pp081
\fullpsalm{108}
\newline {\color{Red} Ant.} Confitebor Domino nimis in ore meo.
\newline {\color{Red} In Laudibus. Ant.} Benigne fac in bona voluntate tua Domine.
\psalm{50}
\newline {\color{Red} Ant.} Bonum est confiteri Domino.
\psalm{91}
\newline {\color{Red} Ant.} Metuant Dominum omnes fines terre.
\psalm{62}
\newline {\color{Red} Ant.} Et in servis suis Dominus miserebitur.
\psalm{audite}
\newline {\color{Red} Ant.} In cymbalis benesonantibus laudate Dominum.
\psalm{148}
\newline {\color{Red} Ant.} In viam pacis dirige nos Domine.
\psalm{benedictus}
\ubsection{Dominica ad Vesperas}{dominica-ad-vesperas}
\newline {\color{Red} Ant.} Dixit Dominus.
\fullpsalm{109}
\newline {\color{Red} Ant.} Dixit Dominus Domino meo sede a dextris meis.
\newline {\color{Red} Ant.} Fidelia.
\fullpsalm{110}
\newline {\color{Red} Ant.} Fidelia omnia mandata eius confirmata in seculum seculi.
\newline {\color{Red} Ant.} In mandatis.
\fullpsalm{111}
\newline {\color{Red} Ant.} In mandatis eius volet nimis.
\newline {\color{Red} Ant.} Sit nomen.
\fullpsalm{112}
\newline {\color{Red} Ant.} Sit nomen Domini benedictum in secula.
\newline {\color{Red} Ant.} Nos qui vivimus.
\fullpsalm{113}
\newline {\color{Red} Ant.} Nos qui vivimus benedicimus Domino.
\ubsection{Feria II ad Vesperas}{feria-ii-ad-vesperas}
\newline {\color{Red} Ant.} Inclinavit.
\fullpsalm{114}
\newline {\color{Red} Ant.} Inclinavit Dominus aurem suam michi.
\newline {\color{Red} Ant.} Credidi.
\fullpsalm{115}
\newline {\color{Red} Ant.} Laudate.
\fullpsalm{116}
\newline {\color{Red} Ant.} Laudate Dominum omnes gentes.
\fullpsalm{117}
\ubsection{Ad Primam}{ad-primam}
\newline {\color{Red} Ant.} Veni.
\psalm{53}
\newline {\color{Red} Ant.} Gloria.
\psalm{53}
\newline {\color{Red} Ant.} Deus.
\psalm{53}
\newline {\color{Red} Ant.} Vivo ego.
\psalm{53}
\newline {\color{Red} Ant.} Anime.
\psalm{53}
\newline {\color{Red} Ant.} Cognoverunt.
\psalm{53}
\fullpsalm{118.a}
\fullpsalm{118.b}
\fullpsalm{118.c}
\fullpsalm{118.d}
\newline {\color{Red} Ant.} Veni et libera nos deus noster.
\newline {\color{Red} Ant.} Gloria tibi trinitas equalis una deitas et ante omnia secula et nunc et in perpetuum.
\newline {\color{Red} Ant.} Deus exaudi orationem meam auribus percipe verba oris mei.
\newline {\color{Red} Ant.} Vivo ego dicit Dominus nolo mortem peccatoris sed ut magis convertatur et vivat.
\newline {\color{Red} Ant.} Anime impiorum fremebant adversum me et gravatum est cor meum super eos.
\newline {\color{Red} Ant.} Cognoverunt Dominum alleluya in fractione panis alleluya.
\ubsection{Ad Tertiam}{ad-tertiam}
\newline {\color{Red} Ant.} Veni Domine.
\newline {\color{Red} Ant.} Laus et perennis.
\newline {\color{Red} Ant.} Deduc me.
\newline {\color{Red} Ant.} Advenerunt.
\newline {\color{Red} Ant.} Iudicasti.
\newline {\color{Red} Ant.} Pax vobis.
\fullpsalm{118.e}
\fullpsalm{118.f}
%pp082
\fullpsalm{118.g}
\fullpsalm{118.h}
\fullpsalm{118.i}
\fullpsalm{118.j}
\newline {\color{Red} Ant.} Veni Domine et noli tardare relaxa facinora plebis tue Israhel.
\newline {\color{Red} Ant.} Laus et perennis gloria deo Patri et Filio Sancto simul Paraclito in secula seculorum.
\newline {\color{Red} Ant.} Deduc me Domine in semita mandatorum tuorum.
\newline {\color{Red} Ant.} Advenerunt nobis dies penitentie ad redimenda peccata et salvandas animas.
\newline {\color{Red} Ant.} Iudicasti Domine causam anime mee defensor vite mee Domine deus meus.
\newline {\color{Red} Ant.} Pax vobis ego sum alleluya nolite timere alleluya.
\ubsection{Ad Sextam}{ad-sextam}
\newline {\color{Red} Ant.} Tuam Domine.
\newline {\color{Red} Ant.} Gloria.
\newline {\color{Red} Ant.} Adiuva me.
\newline {\color{Red} Ant.} Commendemus.
\newline {\color{Red} Ant.} Popule meus.
\newline {\color{Red} Ant.} Ecce ego.
\fullpsalm{118.k}
\fullpsalm{118.l}
\fullpsalm{118.m}
\fullpsalm{118.n}
\fullpsalm{118.o}
\fullpsalm{118.p}
\newline {\color{Red} Ant.} Tuam Domine excita potentiam et veni ad salvandum nos.
\newline {\color{Red} Ant.} Gloria laudis resonet in ore omnium patri geniteque proli spiritui sancto pariter resultet laude perenni.
\newline {\color{Red} Ant.} Adiuva me et salvus ero Domine.
\newline {\color{Red} Ant.} Commendemus nosmetipsos in multa patientia in ieiuniis multis per arma iusticie.
\newline {\color{Red} Ant.} Popule meus quid feci tibi aut quid molestus fui responde michi.
\newline {\color{Red} Ant.} Ecce ego vobiscum sum alleluya omnibus diebus alleluya.
\ubsection{Ad Nonam}{ad-nonam}
\newline {\color{Red} Ant.} In tuo.
\newline {\color{Red} Ant.} Ex quo.
\newline {\color{Red} Ant.} Aspice.
\newline {\color{Red} Ant.} Per arma.
\newline {\color{Red} Ant.} Numquid redditur.
\newline {\color{Red} Ant.} Noli flere.
\fullpsalm{118.q}
\fullpsalm{118.r}
\fullpsalm{118.s}
\fullpsalm{118.t}
\fullpsalm{118.u}
\fullpsalm{118.v}
\newline {\color{Red} Ant.} In tuo adventu erue nos Domine.
\newline {\color{Red} Ant.} Ex quo omnia per quem omnia in quo omnia ipsi gloria in secula.
\newline {\color{Red} Ant.} Aspice in me Domine et miserere mei.
\newline {\color{Red} Ant.} Per arma iusticie virtutis dei commendemus nosmetipsos in multa patientia.
\newline {\color{Red} Ant.} Numquid redditur pro bono malum quia foderunt foveam anime mee.
\newline {\color{Red} Ant.} Noli flere Maria alleluya resurrexit Dominus alleluya.
\newline {\color{Red} Ant.} Clamavi.
\fullpsalm{119}
\newline {\color{Red} Ant.} Clamavi et exaudivit me.
\newline {\color{Red} Ant.} Auxilium.
\fullpsalm{120}
\newline {\color{Red} Ant.} Auxilium meum a Domino.
\newline {\color{Red} Ad Mag. Ant.} Magnificat anima mea Dominum.
\psalm{magnificat}
\ubsection{Feria III ad Vesperas}{feria-iii-ad-vesperas}
\newline {\color{Red} Ant.} In domum.
\fullpsalm{121}
\newline {\color{Red} Ant.} In domum Domini letantes ibimus.
\newline {\color{Red} Ant.} Qui habitas.
\fullpsalm{122}
\newline {\color{Red} Ant.} Qui habitas in celis miserere nobis.
\newline {\color{Red} Ant.} Adiutorium.
\fullpsalm{123}
\newline {\color{Red} Ant.} Adiutorium nostrum in nomine Domini.
\newline {\color{Red} Ant.} Bene fac.
\fullpsalm{124}
\newline {\color{Red} Ant.} Bene fac Domine bonis et rectis corde.
\newline {\color{Red} Ant.} Facti sumus.
\fullpsalm{125}
%pp083
\newline {\color{Red} Ant.} Facti sumus sicut consolati.
\newline {\color{Red} Ad Mag. Ant.} Exultavit spiritus meus in deo salutari meo.
\psalm{magnificat}
\ubsection{Feria IV ad Vesperas}{feria-iv-ad-vesperas}
\newline {\color{Red} Ant.} Beatus vir.
\fullpsalm{126}
\newline {\color{Red} Ant.} Beatus vir qui implevit desiderium suum.
\newline {\color{Red} Ant.} Beati.
\fullpsalm{127}
\newline {\color{Red} Ant.} Beati omnes qui timent Dominum.
\newline {\color{Red} Ant.} Benediximus.
\fullpsalm{128}
\newline {\color{Red} Ant.} Benediximus vobis in nomine Domini.
\newline {\color{Red} Ant.} De profundis.
\fullpsalm{129}
\newline {\color{Red} Ant.} De profundis clamavi ad te Domine.
\newline {\color{Red} Ant.} Speret.
\fullpsalm{130}
\newline {\color{Red} Ant.} Speret Israhel in Domino.
\newline {\color{Red} Ad Mag. Ant.} Respexit humilitatem meam Domine deus meus.
\psalm{magnificat}
\ubsection{Feria V ad Vesperas}{feria-v-ad-vesperas}
\newline {\color{Red} Ant.} Et omnis.
\fullpsalm{131}
\newline {\color{Red} Ant.} Et omnis mansuetudinis eius.
\newline {\color{Red} Ant.} Ecce quam bonum.
\fullpsalm{132}
\fullpsalm{133}
\newline {\color{Red} Ant.} Ecce quam bonum et quam iocundum.
\newline {\color{Red} Ant.} Omnia.
\fullpsalm{134}
\newline {\color{Red} Ant.} Omnia quecunque voluit Dominus fecit.
\newline {\color{Red} Ant.} Quoniam.
\fullpsalm{135}
\newline {\color{Red} Ant.} Quoniam in eternum misericordia eius.
\newline {\color{Red} Ant.} Ymnum.
\fullpsalm{136}
\newline {\color{Red} Ant.} Ymnum cantate nobis de canticis Syon.
\newline {\color{Red} Ad Mag. Ant.} Deposuit potentes sanctos persequentes et exaltavit humiles Christum confitentes.
\psalm{magnificat}
\ubsection{Feria VI ad Vesperas}{feria-vi-ad-vesperas}
\newline {\color{Red} Ant.} In conspectu.
\fullpsalm{137}
\newline {\color{Red} Ant.} In conspectu angelorum psallam tibi deus meus.
\newline {\color{Red} Ant.} Domine.
\fullpsalm{138}
\newline {\color{Red} Ant.} Domine probasti me et cognovisti me.
\newline {\color{Red} Ant.} A viro.
\fullpsalm{139}
\newline {\color{Red} Ant.} A viro iniquo libera me Domine.
\newline {\color{Red} Ant.} Domine.
\fullpsalm{140}
\newline {\color{Red} Ant.} Domine clamavi ad te exaudi me.
\newline {\color{Red} Ant.} Portio.
\fullpsalm{141}
\newline {\color{Red} Ant.} Portio mea Domine sit in terra viventium.
\newline {\color{Red} Ad Mag. Ant.} Suscepit deus Israhel puerum suum sicut locutus est Abraham et semini eius et exaltavit humiles usque in seculum.
\psalm{magnificat}
\ubsection{Sabbato ad Vesperas}{sabbato-ad-vesperas}
\fullpsalm{142}
\newline {\color{Red} Ant.} Benedictus.
\fullpsalm{143}
%pp084
\newline {\color{Red} Ant.} Benedictus Dominus deus meus.
\newline {\color{Red} Ant.} In eternum.
\fullpsalm{144}
\newline {\color{Red} Ant.} In eternum et in seculum seculi.
\newline {\color{Red} Ant.} Laudabo.
\fullpsalm{145}
\newline {\color{Red} Ant.} Laudabo deum meum in vita mea.
\newline {\color{Red} Ant.} Deo nostro.
\fullpsalm{146}
\newline {\color{Red} Ant.} Deo nostro iocunda sit laudatio.
\newline {\color{Red} Ant.} Lauda.
\fullpsalm{147}
\newline {\color{Red} Ant.} Lauda Iherusalem Dominum.
\fullpsalm{148}
\fullpsalm{149}
\fullpsalm{150}
\ubsection{Canticis}{canticis}
\fullpsalm{benedicite}
\fullpsalm{isaias}
\fullpsalm{ezechias}
\fullpsalm{annas}
\fullpsalm{moysis}
\fullpsalm{habacuc}
\fullpsalm{audite}
%pp085
\fullpsalm{magnificat}
\fullpsalm{benedictus}
\fullpsalm{nunc}
\fullpsalm{quicumque}
% \noindentpsalm{credo}
\fullpsalm{credo}
% \noindentpsalm{litanie}
\fullpsalm{litanie}
\fullpsalm{tedeum}
%pp086
\ubsection{De Officio Beate Virginis}{de-officio-bmv}
\ubsubsection{De Officio Beate Virginis Cotidiano}{psalterium-de-officio-beate-virginis-cotidiano}
\color{Red} \lettrine[lines=2]{\zallmancaps \color{Blue} O}{fficium} cotidianum beate virginis totum dicitur extra chorum, preter Completorium. Consuerunt autem dici inter duo signa horarum canonicarum. Attendendum tamen quod quando Missa cantatur immediate post Primam et Tertia continuatur Misse, non est dicenda Tertia de Beata Virgine ante Primam diei, sed prius Tertiam similiter in Dominicis quando Nona dicitur immediate post gratiarum actiones, si Vespere vel vigilie defunctorum continuande fuerint None diei, Nona de Beata Virgine dicatur post Vesperas, vel vigilias defunctorum. Similiter in Quadragesima non sunt dicende Vespere de Beata Virgine ante Nonam diei, sed post Vesperas sive Nona et Vespere continuentur, seu Vespere post Missam dicantur. Ante omnes autem horas dicatur {\color{black} Ave Maria gratia plena Dominus tecum} et qui incipit dicat usque {\color{black} Dominus tecum} ita alte quod audiri possit, et alii sive unus, sive plures prosequantur {\color{black} Benedicta tu in mulieribus et benedictus fructus ventris tui.} \color{black}%typo: in mediate, twice
{\ubsubsection{Ad Vesperas}{psalterium-de-officio-beate-virginis-cotidiano-ad-vesperas}}
\lettrine[lines=2]{\zallmancaps \color{Red} D}{eus} in adiutorium {\color{Red} etc.} Gloria Patri. {\color{Red} etc.} Sicut erat. {\color{Red} et cetera.} Alleluya. {\color{Red} vel.} Laus tibi Domine rex eterne glorie. {\color{Red} Psalmi.}
\psalm{109}
\psalm{112}
\psalm{121}
\psalm{126}
\psalm{147}
{\color{Red} Ant.} Sancta dei genitrix virgo semper Maria intercede pro nobis ad Dominum deum nostrum. {\color{Red} Non dicatur in huiusmodi officio aliquid de antiphona in principio psalmorum. Capitulum.}
\lettrine[lines=2]{\zallmancaps \color{Blue} S}{icut} cynamomum et balsamum aromatizans odorem dedi, quasi mirra electa dedi suavitatem odoris. Deo gratias. {\color{Red} Ymnus.}
% \hymn{ave-maris-stella}{Ave maris stella}{
\lettrine[lines=2]{\zallmancaps \color{Red} A}{ve} maris stella,
dei mater alma,
atque semper virgo,
felix celi porta.
\par {\bfseries \color{Blue} S}umens illud ave
gabrielis ore,
funda nos in pace
mutans nomen Eve.
\par {\bfseries \color{Red} S}olve vincla reis,
profer lumen cecis,
mala nostra pelle,
bona cuncta posce.
\par {\bfseries \color{Blue} M}onstra te esse matrem,
sumat per te precem,
qui pro nobis natus,
tulit esse tuus.
\par {\bfseries \color{Red} V}irgo singularis,
inter omnes mittis,
nos culpis solutos
mites fac et castos.
\par {\bfseries \color{Blue} V}itam presta puram,
iter para tutum,
ut videntes Ihesum,
semper colletemur.
\par {\bfseries \color{Red} S}it laus deo patri
summo christo decus,
spiritui sancto
tribus honor unus. Amen.
% }
\newline \V Ora pro nobis sancta dei genitrix. \R Ut digni efficiamur promissionibus Christi.
\newline {\color{Red} Ad Magnificat Ant.} Sancta Maria succurre miseris, iuva pusillanimes refove flebiles, ora pro populo, interveni pro clero, intercede pro devoto femineo sexu. Dominus vobiscum. Et cum spiritu tuo. {\color{Red} vel.} Domine exaudi orationem meam. Et clamor meus ad te veniat. {\color{Red} Si non sit sacerdos qui dicit orationem. Et hoc idem servetur semper in similibus casibus. Oremus. Oratio.}
\lettrine[lines=2]{\zallmancaps \color{Blue} C}{oncede} nos famulos tuos quesumus domine deus, perpetua mentis et corporis salute gaudere, et gloriosa beate Marie semper virginis intercessione, a presenti liberari tristicia, et eterna perfrui leticia. Per Christum dominum nostrum. Amen.
\newline {\color{Red} Ad memoriam de omnibus sanctis. Ant.} Sancti dei omnes intercedere dignemini pro nostra omniumque salute. \V Orate pro nobis omnes sancti dei. \R Ut digni efficiamur promissionibus Christi. Oremus. {\color{Red} Oratio.}
\lettrine[lines=2]{\zallmancaps \color{Red} T}{ribue} quesumus domine omnes sanctos tuos iugiter orare pro nobis, et eos clementer exaudire digneris. Per Christum dominum nostrum. Amen.
\newline {\color{Red} Ad memoriam de pace. Ant.} Da pacem domine in diebus nostris quia non est alius qui pugnet pro nobis nisi tu deus noster. \V Fiat pax in virtute tua. \R Et abundantia in turribus tuis. Oremus. {\color{Red} Oratio.}
\lettrine[lines=2]{\zallmancaps \color{Blue} D}{eus} a quo sancta desideria, recta consilia, et iusta sunt opera, da servis tuis illam quam mundus dare non potest pacem, ut et corda nostra mandatis tuis dedita, et hostium sublata formidine tempora sint tua protectione tranquilla. Per dominum nostrum, etc. Amen. Dominus vobiscum. {\color{Red} vel.} Domine exaudi. etc. Benedicamus domino. Deo gratias.
{\ubsubsection{Ad Completorium}{psalterium-de-officio-beate-virginis-cotidiano-ad-completorium}}
\lettrine[lines=2]{\zallmancaps \color{Red} C}{onverte} nos deus salutaris noster. Et averte iram tuam a nobis. Deus in adiutorium. etc. {\color{Red} Psalmi.}
\psalm{131}
\psalm{132}
\psalm{133}
{\color{Red} Ant.} Virgo Maria non est tibi similis nata in mundo inter mulieres, florens ut rosa, fragrans sicut lilium, ora pro nobis sancta dei genitrix. {\color{Red} Capitulum.}
\lettrine[lines=2]{\zallmancaps \color{Blue} E}{go} mater pulchre dilectionis et timoris et agnitionis et sancte spei.
\begin{responsory-breve}[intercede-pro-nobis]
{Intercede pro nobis sancta virgo virginum * Mater Dei Maria.}{Ut digni efficiamur promissionibus Christi.}
\end{responsory-breve}%
{\color{Red} Ymnus.}
\hymn{virgo-singularis}{Virgo singularis}{%second half of ave-maris-stella
\lettrine[lines=2]{\zallmancaps \color{Red} V}{irgo} singularis,
Inter omnes mitis,
Nos culpis solútos
Mites fac et castos.
\par {\bfseries \color{Blue} V}itam presta puram,
Iter para tutum,
Ut videntes Ihesum,
Semper colletemur.
\par {\bfseries \color{Red} S}it laus Deo Patri
Summo Christo decus,
Spiritui Sancto
Tribus honor unus. Amen.
}
\newline \V Post partum virgo inviolata permansisti. \R Dei genitrix intercede pro nobis.
\newline {\color{Red} Ad Nunc Dimittis. Ant.} Sub tuum presidium confugimus sancta dei genitrix, nostras deprecationes ne despicias in necessitatibus sed a periculis cunctis libera nos semper virgo benedicta. Dominus vobiscum. etc. Oremus.
\lettrine[lines=2]{\zallmancaps \color{Blue} C}{oncede} misericors deus fragilitati nostre presidium, ut qui sancte dei genitricis memoriam agimus, intercessionis eius auxilio, a nostris iniquitatibus resurgamus. Per eumdem dominum nostrum. Dominus vobiscum. etc.
{\ubsubsection{Ad Matutinas}{psalterium-de-officio-beate-virginis-cotidiano-ad-matutinas}}
\lettrine[lines=2]{\zallmancaps \color{Red} D}{omine} labia mea aperies. Et os meum annunciabit laudem tuam. Deus in adiutorium. etc.
\begin{invitatory-full}[regem-virginis]
{Regem virginis filium * Venite adoremus.}
\end{invitatory-full}
{\color{Red} Ymnus.}
% \hymn{quem-terra}{Quem terra pontus}{
\lettrine[lines=2]{\zallmancaps \color{Red} Q}{uem} terra pontus ethera
colunt adorant predicant,
trinam regentem machinam
claustrum Marie baiulat.
\par {\bfseries \color{Red} C}ui luna sol et omnia,
deserviunt per tempora
perfusa celi gratia
gestant puelle viscera.
\par {\bfseries \color{Blue} B}eata mater munere
cuius supernus artifex
mundum pugillo continens
ventris sub archa clausus est.
\par {\bfseries \color{Red} B}eata celi nuncio
fecunda sancto spiritu
desideratus gentibus
cuius per alvum fusus est.
\par {\bfseries \color{Blue} G}loria tibi domine
qui natus es de virgine
cum patre et sancto spiritu,
in sempiterna secula. Amen.
% }
\newline {\color{Red} Psalmi.}
\psalm{8}
\psalm{18}
\psalm{23}
{\color{Red} Ant.} Benedicta tu in mulieribus et benedictus fructus ventris tui.
\V Diffusa est gratia in labiis tuis. \R Propterea benedixit te deus in eternum. Pater noster. Et ne nos. Sed libera. Iube domne benedicere.
\newline {\color{Red} Benedictio.} Alma virgo virginum, intercedat pro nobis ad dominum. Amen. {\color{Red} Lectio Prima.}
\lettrine[lines=2]{\zallmancaps \color{Blue} S}{ancta} Maria virgo virginum mater et filia regis regum omnium, tuum nobis impende solatium, ut celestis regni per te mereamur habere premium, et cum electis dei regnare in perpetuum. Tu autem domine miserere nostri. Deo gratias.
\begin{responsory}[parvum-sancta-et-immaculata]
{Sancta et immaculata virginitas quibus te laudibus efferam nescio * Quia quem celi capere non poterant tuo gremio contulisti.}{Benedicta tu in mulieribus et benedictus fructus ventris tui.}
\end{responsory}%
Iube etc. {\color{Red} Benedictio.} Sancta dei genitrix sit nobis auxiliatrix. Amen. {\color{Red} Lectio Secunda.}
\lettrine[lines=2]{\zallmancaps \color{Red} S}{ancta} Maria piarum piissima intercede pro nobis sanctarum sanctissima, per te virgo, sumat nostra precamina, qui pro nobis ex te natus regnat super ethera ut sua caritate nostra deleantur peccamina. Tu autem, etc.
\begin{responsory}[beata-es-maria]
{Beata es Maria que domini portasti creatorem mundi * Genuisti qui te fecit et in eternum permanes virgo.}{Ave Maria gratia plena dominus tecum.}
\end{responsory}%
Iube do. ben. {\color{Red} Benedictio.} Nos cum prole pia benedicat virgo Maria. Amen. {\color{Red} Lectio Tertia.}
\lettrine[lines=2]{\zallmancaps \color{Blue} S}{ancta} dei genitrix que digne meruisti concipere quem totus orbis nequivit comprehendere, tuo pro interventu culpas nostras ablue, ut perennis sedem glorie per te redempti valeamus scandere, ubi regnas cum eodem filio tuo sine tempore. Tu autem.
\begin{responsory-final}[parvum-felix-namque]
{Felix namque es sacra virgo Maria et omni laude dignissima * Quia ex te ortus est sol iusticie * Christus deus noster.}{Ora pro populo interveni pro clero intercede pro devoto femineo sexu, sentiant omnes tuum iuvamen quicumque celebrant tuam commemorationem.}{Christus}%
\end{responsory-final}%
\psalm{tedeum} {\color{Red} dicatur quando in choro dicitur.} \V Ora pro nobis sancta dei genitrix. \R Ut digni.
{\ubsubsection{In Laudibus}{psalterium-de-officio-beate-virginis-cotidiano-in-laudes}}
\par \noindent Deus in adiutorium. etc. {\color{Red} Psalmi.}
\psalm{92}
\psalm{99}
\psalm{62}
\psalm{benedicite}
\psalm{148}
{\color{Red} Ant.} Post partum virgo inviolata permansisti dei genitrix intercede pro nobis. {\color{Red} Capitulum.}
\lettrine[lines=2]{\zallmancaps \color{Red} E}{go} quasi vitis fructificavi suavitatem odoris: et flores mei fructus honoris et honestatis. Deo gratias. {\color{Red} Ymnus.}
% \hymn{o-gloriosa}{O gloriosa domina}{
\lettrine[lines=2]{\zallmancaps \color{Blue} O}{}\ gloriosa domina
excelsa supra sydera,
qui te creavit provide
lactasti sacro ubere.
\par {\bfseries \color{Red} Q}uod Eva tristis abstulit,
tu reddis almo germine,
intrent ut astra flebiles
celi fenestra facta es.
\par {\bfseries \color{Blue} T}u regis alti ianua
et porta lucis fulgida,
vitam datam per virginem,
gentes redempte plaudite.
\par {\bfseries \color{Red} G}loria tibi domine
qui natus es de virgine
cum patre et sancto spiritu,
in sempiterna secula. Amen.
% }
\newline \V Elegit eam deus et preelegit eam. \R Et habitare eam facit in tabernaculo suo.
\newline {\color{Red} Ad Benedictus Ant.} O gloriosa dei genitrix virgo semper Maria que dominum omnium meruisti portare et regem angelorum sola virgo semper lactare nostri quesumus pia memorare, et pro nobis Christum deprecare, ut tuis fulti patrociniis ad celestia regna mereamur pervenire. Dominus vobiscum. Oremus.
{\bfseries \color{Red} C}{oncede} nos famulos tuos. etc.
{\color{Red} Antiphone, Versiculi, Orationes, et alia pro memoriis omnia ut supra in vesperis.}
{\ubsubsection{Ad Primam}{psalterium-de-officio-beate-virginis-cotidiano-ad-primam}}
\par \noindent Deus in adiutorium. {\color{Red} Ymnus.}
\hymn{memento-salutis}{Memento salutis actor}{
\lettrine[lines=2]{\zallmancaps \color{Red} M}{emento} salutis actor
quod nostri quondam corporis,
ex illibata virgine,
nascendo formam sumpseris.
\par {\bfseries \color{Blue} M}aria mater gratie,
mater misericordie,
tu nos ab hoste protege
et hora mortis suscipe.
\par {\bfseries \color{Red} G}loria tibi domine
qui natus es de virgine
cum patre et sancto spiritu,
in sempiterna secula. Amen.
}
\newline {\color{Red} Psalmi.}
\psalm{119}
\psalm{120}
\psalm{121}
{\color{Red} Ant.} Dignare me laudare te virgo sacrata, da michi virtutem contra hostes tuos. {\color{Red} Capitulum.}
\lettrine[lines=2]{\zallmancaps \color{Blue} A}{b} initio et ante secula creata sum, et usque ad futurum seculum non desinam: et in habitatione sancta coram ipso ministravi. Deo gratias.
\begin{responsory-breve}[parvum-post-partum]
{Post partum virgo, * Inviolata permansisti.}{Dei genitrix intercede pro nobis.}
\end{responsory-breve}%
\V Benedicta tu in mulieribus. \R Et benedictus fructus ventris tui. Dominus vobiscum. etc. {\color{Red} Oratio.}
\lettrine[lines=2]{\zallmancaps \color{Red} G}{ratiam} tuam quesumus domine mentibus nostris infunde, ut qui angelo nunciante Christi filii tui incarnationem cognovimus per passionem eius et crucem, ad resurrectionis gloriam perducamur. Per eumdem do. etc.
{\ubsubsection{Ad Tertiam}{psalterium-de-officio-beate-virginis-cotidiano-ad-tertiam}}
\par \noindent Deus in adiutorium. etc. {\color{Red} Ymnus.} \hyperlink{memento-salutis}{Memento salutis actor}. {\color{Red} Psalmi.}
\psalm{122}
\psalm{123}
\psalm{124}
{\color{Red} Ant.} Gaude Maria virgo cunctas hereses sola interemisti in universo mundo. {\color{Red} Capitulum.}
\lettrine[lines=2]{\zallmancaps \color{Blue} E}{t} sic in Syon firmata sum, et in civitate sanctificata similiter requievi, et in Hierusalem potestas mea. Deo gratias.
\begin{responsory-breve}%[sancta-maria-mater]
{Sancta Maria mater Christi * Audi rogantes servulos.}{Et impetratam nobis celitus tu defer indulgentiam.}
\end{responsory-breve}%
\V Ora pro nobis sancta dei genitrix. \R Ut digni. etc. Dominus vobiscum. etc. {\color{Red} Oratio.}
{\bfseries \color{Red} C}{oncede} nos famulos tuos. etc. Per do. etc.
{\ubsubsection{Ad Sextam}{psalterium-de-officio-beate-virginis-cotidiano-ad-sextam}}
\par \noindent Deus in adiutorium meum. etc. {\color{Red} Ymnus.} \hyperlink{memento-salutis}{{\bfseries \color{Red} M}emento salutis actor}. {\color{Red} Psalmi.}
\psalm{125}
\psalm{126}
\psalm{127}
{\color{Red} Ant.} In prole mater, in partu virgo, gaude et letare virgo mater domini. {\color{Red} Capitulum.}
\lettrine[lines=2]{\zallmancaps \color{Red} E}{t} radicavi in populo honorificato: et in partes dei mei hereditas illius, et in plenitudine sanctorum detentio mea. Deo gratias.
\begin{responsory-breve}%[ora-pro-nobis]
{Ora pro nobis * Sancta dei genitrix.}{Ut digni efficiamur promissionibus Christi.}
\end{responsory-breve}%
\V Elegit eam deus et preelegit eam. \R Et habitare eam facit in tabernaculo suo. Dominus vobiscum. Oremus. {\color{Red} Oratio.}
\lettrine[lines=2]{\zallmancaps \color{Blue} P}{rotege} domine famulos tuos subsidiis pacis et beate Marie semper virginis patrociniis confidentes: a cunctis hostibus redde securos. Per do. etc.
{\bfseries \color{Red} C}{oncede} nos famulos tuos. etc. Per do. etc.
{\ubsubsection{Ad Nonam}{psalterium-de-officio-beate-virginis-cotidiano-ad-nonam}}
\par \noindent Deus in adiutorium. {\color{Red} Ymnus.} \hyperlink{memento-salutis}{{\bfseries \color{Red} M}emento salutis actor}. etc. {\color{Red} Psalmi.}
\psalm{128}
\psalm{129}
\psalm{130}
{\color{Red} Ant.} Beata mater, et intacta virgo, gloriosa regina mundi intercede pro nobis ad dominum. {\color{Red} Capitulum.}
\lettrine[lines=2]{\zallmancaps \color{Red} Q}{uasi} cedrus exaltata sum in Libano, et quasi Cypressus in monte Syon, quasi palma exaltata sum in Cades, et quasi plantatio rose in Ihericho. Deo gratias.
\begin{responsory-breve}%[elegit-eam]
{Elegit eam deus * Et preelegit eam.}{Et habitare eam facit in tabernaculo suo.}
\end{responsory-breve}%
\V Sancta dei genitrix virgo semper Maria. \R Intercede pro nobis ad dominum deum nostrum. Dominus vobiscum. {\color{Red} Oremus. Oratio.}
\lettrine[lines=2]{\zallmancaps \color{Blue} F}{amulorum} tuorum quesumus domine delictis ignosce, ut qui tibi placere de actibus nostris non valemus, genitricis filii tui domini nostri Ihesu Christi intercessione salvemur. Qui tecum vivit. etc.
\newline {\color{Red} Predicto modo fiat per totum annum cotidianum officium beate virginis quando faciendum est, excepto quod in Adventu dicitur ad Magnificat Ant.} O virgo virginum quomodo fiet istud, quia nec primam similem visa es nec habere sequentem, filie Iherusalem quid me admiramini, divinum est mysterium hoc quod cernitis. {\color{Red} Et ad Vesperas et ad Matutinas et ad Tertiam. Oratio.}
\lettrine[lines=2]{\zallmancaps \color{Red} D}{eus} qui de beate Marie virginis utero verbum tuum angelo nunciante carnem suscipere voluisti, presta supplicibus tuis, ut qui vere eam genitricem dei credimus, eius apud te intercessionibus adiuvemur.
{\color{Red} Terminanda in Vesperis et Matutinas sic.} Per eumdem Christum dominum nostrum. {\color{Red} In Tertia vero sic.} Per eumdem dominum nostrum. etc. Amen. {\color{Red} Item ab Octava Epyphanie usque ad Purificationem in eisdem Horis. Oratio.}
\lettrine[lines=2]{\zallmancaps \color{Blue} D}{eus} qui salutis eterne beate Marie virginitate fecunda, humano generi premia prestitisti, tribue quesumus ut ipsam pro nobis intercedere sentiamus, per quam meruimus actorem vite suscipere.
{\color{Red} Terminanda in Vesperis et Matutinas sic.} Christum Dominum nostrum. {\color{Red} In Tertia vero sic.} Dominum nostrum Ihesum Christum. etc. Amen. {\color{Red} Item in tempore Paschali usque ad Penthecosten dicitur ad Magnificat Ant.} Regina celi letare alleluya quia quem meruisti portare alleluya, resurrexit sicut dixit, alleluya, ora pro nobis deum alleluya.
{\color{Red} Sed in tempore Ascensionis dicitur.} Iam ascendit sicut dixit. {\color{Red} Ubi scriptum est.} Resurrexit.


\end{multicols*}


%pp087
\ubchapter{Breviarium}

\ubsection{Rubrice Temporum}{breviarium}

\thispagestyle{fancy}
\fancyhead[RO,LE]{\rightmark}

\begin{multicols}{2}

{\ubsubsection{De Quo Officium Sit Agendum}{breviarium-de-quo-officium-sit-agendum}}
\lettrine[lines=3]{\zallmancaps \color{Blue} N}{}\ul{\textsc{otandum} quod per totum annum Officium tam diei quam noctis ad Horas faciendum est de Tempore nisi in festis Sanctorum et per octavas et in octavis eorum et nisi in Sabbatis quando de Beata Virgine agitur in conventu.
\newline Quando igitur de Tempore fuerit faciendum: omnibus Dominicis extra Tempus Pascale faciende sunt novem lectiones in Matutinis. In profestis autem diebus per totum annum fiant tantum tres nisi in Triduo ante Pasca. Tempus autem Paschale vocatur missa de vigilia Pasche et deinceps usque ad primas Vesperas Trinitatis exclusive.}

{\ubsubsection{De Responsorio in Primis Vesperis Dominice}{breviarium-de-responsorio-in-primis-vesperis-dominice}}
\lettrine[lines=2]{\zallmancaps \color{Red} O}{}\ul{\textsc{mni} tempore quando hystoria dominicalis in sua prima Dominica inchoatur, responsorium dicendum est in primis Vesperis. Quando vero propter aliquod impedimentum Dominica non habet primas Vesperas, vel hystoria dominicalis non cantatur in Dominica in qua primo debuit inchoari: responsorium in Vesperis de cetero non dicatur.}

{\ubsubsection{De Oratione Dominicali}{breviarium-de-oratione-dominicali}}
\lettrine[lines=2]{\zallmancaps \color{Blue} N}{}\ul{\textsc{otandum} est quod oratio Dominicalis semper dicenda est in Sabbato precedenti ad Vesperas, et in ipsa Dominica ad omnes Horas tam diei quam noctis quando de Dominica agitur, preterquam in Prima et in Completorio. Et hoc sciendum quod ubicumque dicitur hec oratio dicatur ad Horas nunquam intelligitur de Prima neque de Completorio nisi specialiter exprimatur. Dicenda est eciam oratio Dominicalis per ferias sequentis ebdomade usque ad Vesperas Sabbati exclusive ad omnes Horas tam diei quam noctis quando de Tempore agitur: nisi in ieiuniis Quatuor Temporum Adventus et in quarta feria Cinerum et deinceps usque ad Dominicam in Albis exclusive, et nisi in tribus diebus Rogationum, et nisi in ebdomada Pentecostes, et nisi due aut plures orationes Dominicales infra ebdomadam sint dicende. Quando vero due vel plures orationes infra unam ebdomadam dicende fuerint, postquam una infra ebdomadam fuerit incepta, dicatur quando de Tempore agitur usque dum alia inchoetur: et sic deinceps si plures dicende fuerint.
Ille autem orationes que infra ebdomadam inchoantur, semper sunt in Laudibus inchoande. Hoc eciam sciendum quod oratio que dicitur ad Laudes semper dicenda est ad Tertiam, Sextam et Nonam nisi in ieiuniis Quatuor Temporum Adventus et in duobus diebus Rogationum. Dicitur eciam ad Vesperas quando agitur de eodem de quo in Laudibus, nisi per ferias in Quadragesima et in duobus diebus Rogationum.}

{\ubsubsection{De Omelia Dominicali Quando Legenda vel Transferenda vel Officio Dimittenda Sit}{breviarium-de-omelia-dominicali-quando-legenda-vel-transferenda-vel-officio-dimittenda-sit}}
\lettrine[lines=2]{\zallmancaps \color{Red} I}{}\ul{\textsc{n} omni Dominica, nisi festum vel octava dies festi habentis octavam impediat, tres ultime lectiones de expositione evangelii dominicalis legantur preterquam in vigilia Natalis Domini, quando in Dominica evenerit, tunc enim omelia dominicalis omittitur illo anno. Et nisi quando aliqua vigilia habens propriam omeliam in Dominica evenerit. Tunc enim tres ultime lectiones erunt de omelia vigilie et omelia dominicalis in aliquam feriam transferetur, preterquam in vigilia Sancti Andree, quando in Dominica evenerit. Tunc enim omelia dominicalis legenda est, et omelia de vigilia illo anno penitus dimittenda. Quando vero propter aliquod impedimentum dominicalis omelia non legitur in Dominica, semper infra ebdomadam suam in aliqua feria est legenda, nisi quando festum Sancti Silvestri in Dominica evenerit, et quando festum Beati Petri Martiris, vel Apostolorum Philippi et Iacobi vel Sancte Crucis vel Corone Domini, in Dominica ante Ascensionem evenerint. Tunc enim omelia dominicalis dimittitur illo anno. Dimittuntur eciam due ultime omelie dominicales que sunt post,} \hyperlink{domine-ne-in-ira}{Domine ne in ira}, \ul{quando inter octavas Epiphanie et Septuagesimam non sunt ferie vacantes in quibus possint legi. Similiter et omelia evangelii} Loquente Ihesu \ul{dimittitur quando tantum viginti due septimane fuerint inter} \hyperlink{deus-omnium}{Deus omnium,} \ul{et Adventum, si in vicesima secunda septimana non fuerint tot ferie vacantes, quod omnia Evangelia legi possint.}
\newline \psalm{tedeum} \ul{dicatur omnibus Dominicis, preterquam in Adventu et in Septuagesima et deinceps usque ad Pascha exclusive, nisi quando festum Purificationis in Septuagesima vel Sexagesima vel Quinquagesima evenerit.}

{\ubsubsection{De Divisione Hystorie Dominicalis per Ferias}{breviarium-de-divisione-hystorie-dominicalis-per-ferias}}
\lettrine[lines=2]{\zallmancaps \color{Blue} S}{}\ul{\textsc{ciendum} quod per totum annum responsoria de hystoria dominicali per ebdomadam, quando de Tempore agitur hoc ordine sunt dicenda. Videlicet in secunda et quinta feria, tria prima. In tertia et sexta feria tria media. In quarta feria et Sabbato, tria ultima, et hoc observandum est nisi aliter sint signata. Quando autem festum aliquod infra ebdomadam celebratur, nichilominus ordo predictus in reliquis feriis observetur.}

{\ubsubsection{Ubi Adventus Domini Celebrandus Sit}{breviarium-ubi-adventus-domini-celebrandus-sit}}
\lettrine[lines=2]{\zallmancaps \color{Red} A}{}\ul{\textsc{dventus} Domini semper celebrari debet in Dominica que vicinior est diei Sancti Andree, sive precedat sive sequatur, vel in ipso die Sancti Andree quando in Dominica evenerit.
\newline Omnia que ad Horas tam diei quam noctis per totum Adventum dicenda sunt, que in Breviario non sunt notata, dicenda sunt secundum quod in Dominica,} \hyperlink{domine-ne-in-ira}{Domine ne in ira} \ul{et in eius feriis sunt signata.}

\ubsection{Proprium de Tempore}{proprium-de-tempore}

{\ubsubsection{Dominica Prima in Adventu Domini Sabbato Precedenti}{breviarium-dominica-prima-in-adventu-domini-sabbato-precedenti}}
\par \noindent {\color{Red} Ad Vesperas. Ant.} Benedictus. {\color{Red} Psalmi.} Super psalmos. {\color{Red} Cetere ad ceteros. Capitulum.}
\lettrine[lines=2]{\zallmancaps \color{Blue} E}{cce} \hypertarget{ecce-dies-capitulum}{\label{ecce-dies-capitulum}} dies veniunt dicit dominus et suscitabo David germen iustum, et regnabit rex et sapiens erit, et faciet iudicium et iusticiam in terra.
\R \hyperlink{missus-est}{Missus est.} {\color{Red} III. Ymnus.}
\hymn{conditor}{Conditor alme syderum}{
\lettrine[lines=2]{\zallmancaps \color{Red} C}{onditor} alme syderum,
eterna lux credentium,
Christe redemptor omnium,
exaudi preces supplicum.
\par {\bfseries \color{Blue} Q}ui condolens interitu
mortis perire seculum,
salvasti mundum languidum,
donans reis remedium.
\par {\bfseries \color{Red} V}ergente mundi vespere
uti sponsus de thalamo,
egressus honestissima
virginis matris clausula.
\par {\bfseries \color{Blue} C}uius forti potentie
genu curvantur omnia,
celestia terrestria
fatentur nutu subdita.%sic
\par {\bfseries \color{Red} T}e deprecamur Agyie
venture iudex seculi
conserva nos in tempore
hostis a telo perfidi.
\par {\bfseries \color{Blue} L}aus honor virtus gloria
Deo Patri et Filio,
sancto simul Paraclito
in sempiterna secula. Amen.
}
\newline \V Rorate celi desuper et nubes pluant iustum. \R Aperiatur terra et germinet salvatorem.
\newline {\color{Red} Ad Magnificat. Ant.} Ecce nomen domini venite de longinquo, et claritas eius replet orbem terrarum. {\color{Red} Oratio.}
\lettrine[lines=2]{\zallmancaps \color{Blue} E}{xcita} quesumus domine potentiam tuam et veni, ut ab imminentibus peccatorum nostrorum periculis te mereamur protegente eripi, te liberante salvari. Qui vivis.
\newline {\color{Red} Ad Matutinas. Invitat.} 
\begin{invitatory}[ecce-venite]
{Ecce venit rex * Occurramus obviam salvatori nostro.}
\end{invitatory}
{\color{Red} Ymnus.}
\phantomsection
\hymn{verbum}{Verbum supernum prodiens}{
\lettrine[lines=2]{\zallmancaps \color{Red} V}{erbum} supernum prodiens
a patre olim exiens,
qui natus orbi subvenis
cursu declivi temporis.
\par {\bfseries \color{Red} I}llumina nunc pectora
tuoque amore concrema,
audito ut preconio
sint pulsa tandem lubrica.
\par {\bfseries \color{Blue} I}udexque cum post aderis
rimari facta pectoris
reddens vicem pro abditis
iustisque regnum pro bonis.
\par {\bfseries \color{Red} N}on demum arctemur malis
pro qualitate criminis
sed cum beatis compotes
simus perennes celibes.
\par {\bfseries \color{Blue} L}aus honor virtus gloria
Deo Patri et Filio,
sancto simul Paraclito
in sempiterna secula. Amen.
}
\newline {\color{Red} In Primo Nocturno. Ant.} Scientes quia hora est iam nos de somno surgere, nunc enim propior est nostra salus quam cum credidimus, nox precessit dies autem appropinquavit. {\color{Red} Psalmi.}
\psalm{1}
\psalm{2}
\psalm{3}
\psalm{6}
Gloria patri. 
\psalm{7}
\psalm{8}
\psalm{9}
\psalm{10}
Gloria patri. 
\psalm{11}
\psalm{12}
\psalm{13}
\psalm{14}
Gloria. \V Ex Syon species decoris eius. \R Deus noster manifeste veniet.
\newline {\color{Red} Lectio Prima.}
\lettrine[lines=2]{\zallmancaps \color{Blue} V}{isio} Ysaie filii Amos, quam vidit super Iudam et Hierusalem, in diebus Ozie, Ioathan, Achaz, Ezechie, regum Iuda. Audite celi et auribus percipe terra: quoniam dominus locutus est. Filios enutrivi et exaltavi: ipsi autem spreverunt me. Hec dicit dominus deus, convertimini ad me et salvi eritis.
% Hec dicit is for Isaias
\hypertarget{aspiciens}{\label{aspiciens}}
\newline \R Aspiciens a longe, ecce video dei potentiam venientem et nebulam totam terram tegentem, * Ite obviam ei et dicite, * Nuntia nobis si tu es ipse, * Qui regnaturus es, * In populo Israhel. \V Quique terrigene et filii hominum simul in unum dives et pauper: Ite. \ul{Repetatur usque in finem.} \V Qui regis Israhel intende qui deducis velut ovem Ioseph. Nuntia. \ul{Usque in finem}. \V Excita domine potentiam tuam et veni et salvos facias nos. Qui regnaturus. \ul{Usque in finem.} Gloria Patri. In populo Israhel. \R Aspiciens.
\newline {\color{Red} Lectio Secunda.}
\lettrine[lines=2]{\zallmancaps \color{Red} C}{ognovit} bos possessorem suum, et asinus presepe domini sui: Israhel autem me non cognovit et populus meus non intellexit. Ve genti peccatrici, populo gravi iniquitate, semini nequam, filiis sceleratis. Dereliquerunt dominum: blasphemaverunt sanctum Israhel, abalienati sunt retrorsum. Hec dicit.
\begin{responsory}[aspiciebam]
{Aspiciebam in visu noctis, et ecce in nubibus celi filius hominis venit et datum est ei regnum et honor, * Et omnis populus tribus et lingue servient ei.}{Potestas eius potestas eterna que non auferetur, et regnum eius quod non corrumpetur.}
\end{responsory}%
\newline {\color{Red} Lectio Tertia.}
\lettrine[lines=2]{\zallmancaps \color{Blue} S}{uper} quo percutiam vos ultra addentes prevaricationem? Omne caput languidum, et omne cor merens, a planta pedis usque ad verticem non est in eo sanitas. Vulnus et livor et plaga tumens, non est circumligata, neque curata medicamine, neque fota oleo. Hec dicit.
\begin{responsory-doxology}[missus-est]
{Missus est Gabriel angelus ad Mariam virginem desponsatam Ioseph nuncians ei verbum et expavescit virgo de lumine, ne timeas Maria invenisti gratiam apud dominum, ecce concipies et paries, * Et vocabitur altissimi filius.}{Dabit ei dominus deus sedem David patris eius, et regnabit in domo Iacob in eternum.}
\end{responsory-doxology}%
\newline {\color{Red} In Secundo Nocturno. Ant.} Hora est iam nos de somno surgere, et aperti sunt oculi nostri surgere ad Christum, quia lux vera est fulgens in mundo. {\color{Red} Psalmi.}
\psalm{15} Gloria.
\psalm{16} Gloria.
\psalm{17} Gloria. \V Egredietur virga de radice Iesse. \R Et flos de radice eius ascendet.
\newline {\color{Red} Lectio Quarta.}
\lettrine[lines=2]{\zallmancaps \color{Red} T}{erra} vestra deserta, civitates vestre succense igni. Regionem vestram coram vobis alieni
%pp088
devorant, et desolabitur sicut in vastitate hostili. Et derelinquetur filia Syon ut umbraculum in vinea, et sicut tugurium in cucumerario, sicut civitas que vastatur. Hec dicit.
\begin{responsory}[ave-maria]
{Ave Maria gratia plena dominus tecum, * Spiritus sanctus superveniet in te et virtus altissimi obumbrabit tibi, quod enim ex te nascetur sanctum vocabitur filius dei.}{Quomodo fiet istud quoniam virum non cognosco, et respondens angelus dixit ei.}
\end{responsory}%
\newline {\color{Red} Lectio Quinta.}
\lettrine[lines=2]{\zallmancaps \color{Blue} N}{isi} dominus exercituum reliquisset nobis semen, quasi Sodoma fuissemus, et quasi Gomorra similes essemus. Audite verbum domini principes Sodomorum, auribus percipite legem dei nostri populus Gomorre. Quo michi multitudinem victimarum vestrarum dicit dominus? Plenus sum. Hec dicit.
\begin{responsory}[salvatorem-expectamus]
{Salvatorem expectamus dominum Ihesum Christum, * Qui reformabit corpus humilitatis nostre configuratum corpori claritatis sue.}{Sobrie et iuste et pie vivamus in hoc seculo expectantes beatam spem et adventum glorie magni dei.}
\end{responsory}%
\newline {\color{Red} Lectio Sexta.}
\lettrine[lines=2]{\zallmancaps \color{Red} H}{olocausta} arietum et adipem pinguium, et sanguinem vitulorum et agnorum et hyrcorum, nolui. Cum veniretis ante conspectum meum quis quesivit hec de manibus vestris, ut ambularetis in atriis meis? Ne offeratis ultra sacrificium frustra. Incensum abhominatio est michi. Neomeniam, et sabbatum, et festivitates alias non feram. Hec dicit.
\begin{responsory-doxology}[audite-verbum]
{Audite verbum domini gentes et annunciate illud in finibus terre et in insulis que procul sunt dicite, * Salvator noster adveniet.}{Annunciate et auditum facite loquimini et clamate.}
\end{responsory-doxology}%
\newline {\color{Red} In Tertio Nocturno. Ant.} Nox precessit dies autem appropinquavit, abiciamus ergo opera tenebrarum et induamur arma lucis sicut in die honeste ambulemus. {\color{Red} Psalmi.}
\psalm{18} Gloria.
\psalm{19} Gloria.
\psalm{20} Gloria. \V Egredietur dominus de loco sancto suo. \R Veniet ut salvet populum suum.
\newline {\color{Red} Lectio VII. Lectio sancti Evangelii Secundum Matheum.}
% https://la.wikisource.org/wiki/Expositio_evangelii_secundum_Lucam/IX
% Same as Roman Palmis
\lettrine[lines=2]{\zallmancaps \color{Blue} I}{n} illo tempore: Cum appropinquasset Ihesus Hierosolimis, et venissent Bethphage ad montem Oliveti, tunc misit duos discipulos dicens eis. Ite in castellum quod contra vos est. Et reliqua.
\newline {\color{Red} Omelia Beati Ambrosii Episcopi.}
\lettrine[lines=2]{\zallmancaps \color{Red} P}{ulchre} relictis Iudeis habitaturus in mentibus gentium, templum dominus ascendit. Hoc enim templum est: quod fidei series non lapidum structura fundavit. Deseruntur ergo qui adorant, eliguntur qui amaturi erant.
\begin{responsory}[ecce-virgo]
{Ecce virgo concipiet et pariet filium dicit dominus, et vocabitur nomen eius, * Admirabilis deus fortis.}{Super solium David et super regnum eius sedebit in eternum.}
\end{responsory}%
\newline {\color{Red} Lectio VIII.}
\lettrine[lines=2]{\zallmancaps \color{Blue} E}{t} ideo ad montem venit Oliveti, ut novellas olivas in sublimi virtute plantaret, quarum mater est illa que sursum est Hierusalem. In hoc monte est ille celestis agricola, ut plantati omnes in domo domini possint veraciter dicere. Ego sicut oliva fructifera in domo domini.
\begin{responsory}[obsecro-domine]
{Obsecro domine mitte quem missurus es, vide afflictionem populi tui sicut locutus es, * Veni et libera nos.}{A solis ortu et occasu ab aquilone et mari.}
\end{responsory}%
\newline {\color{Red} Lectio IX.}
\lettrine[lines=2]{\zallmancaps \color{Red} E}{t} fortasse ipse Christus mons est. Quis enim alius tales fructus ferret olivarum, non curvescentium ubertate baccarum, sed spiritus plenitudine gentium fecundarum? Ipse est per quem ascendimus. Ipse est ianua, ipse est via, qui aperitur et aperit, qui pulsatur ab ingredientibus et ab egredientibus adoratur.
\begin{responsory-final}[ecce-dies-veniunt]
{Ecce dies veniunt dicit dominus, et suscitabo David germen iustum, et regnabit rex et sapiens erit, et faciet iudicium et iusticiam in terra, * Et hoc est nomen quod vocabunt eum, * Dominus iustus noster.}{In diebus illis salvabitur Iuda et Israhel habitabit confidenter.}{Dominus}
\end{responsory-final}%
\newline \ul{Non dicatur,} \psalm{tedeum} \ul{per totum Adventum nisi in festis sanctorum, sed nonum responsorium in omnibus dominicis repetatur.}
\newline \V Emitte agnum domine dominatorem terre. \R De petra deserti ad montem filie Syon.
\newline \ul{Predictus versiculus dicatur per totum Adventum quando de tempore agitur.}
\newline {\color{Red} In Laudibus. Ant.} In illa die stillabunt montes dulcedinem et colles fluent lac et mel, alleluya. {\color{Red} Ps.} \psalm{92} {\color{Red} Ant.} Iocundare filia Syon exulta satis filia Hierusalem, alleluya. {\color{Red} Ps.} \psalm{99} {\color{Red} Ant.} Ecce dominus veniet et omnes sancti eius cum eo et erit in die illa lux magna, alleluya. {\color{Red} Ps.} \psalm{62} {\color{Red} Ant.} Omnes sitientes venite ad aquas, querite dominum dum inveniri potest, alleluya. {\color{Red} Can.} \psalm{benedicite} {\color{Red} Ant.} Ecce veniet propheta magnus et ipse renovabit Hierusalem, alleluya. {\color{Red} Ps.} \psalm{148} {\color{Red} Cap.} \hyperlink{ecce-dies-capitulum}{Ecce dies.} {\color{Red} Ymnus.}
\hymn{vox-clara}{Vox clara ecce}{
\lettrine[lines=2]{\zallmancaps \color{Blue} V}{ox} clara ecce intonat
obscura queque increpat
pellantur eminus somnia
ab ethre Christus promicat.
\par {\bfseries \color{Red} M}ens iam resurgat torpida
que sorde extat saucia,
sydus refulget iam novum,
ut tollat omne noxium.
\par {\bfseries \color{Blue} E} sursum agnus mittitur
laxare gratis debitum,
omnes pro indulgentia
vocem demus cum lacrimis.
\par {\bfseries \color{Red} S}ecundo ut cum fulserit,
mundumque horror cinxerit,
non pro reatu puniat
sed pius nos tunc protegat.
\par {\bfseries \color{Blue} L}aus honor virtus gloria
Deo Patri et Filio,
sancto simul Paraclito
in sempiterna secula. Amen.
}
\newline \V Vox clamantis in deserto. \R Parate viam domini.
\newline {\color{Red} Ad Benedictus. Ant.} Spiritus sanctus in te descendet Maria ne timeas, habebis in utero filium dei alleluya.
\newline \ul{Ad horas diei antiphone Laudum quarta pretermissa et hoc per totum Adventum in Dominicis observetur.}
\newline {\color{Red} Ad Primam.} \R Ihesu Christe fili dei vivi miserere nobis. \V Qui venturus es in mundum.
\newline \ul{Hic versus dicatur ad primam per totum Adventum eciam in festis sanctorum.}
\newline {\color{Red} Ad Tertiam. Cap.} Ecce dies.
\begin{responsory-breve}[veni-ad-liberandum]
{Veni ad liberandum nos, * Domine deus virtutum.}{Ostende faciem tuam et salvi erimus.}
\end{responsory-breve}%
\V Timebunt gentes nomen tuum domine. \R Et omnes reges terre gloriam tuam.
\newline {\color{Red} Ad Sextam. Capitulum.}
\lettrine[lines=2]{\zallmancaps \color{Red} I}{n} diebus illis salvabitur Iuda et Israhel habitabit confidenter, et hoc est nomen quod vocabunt eum dominus iustus noster.
\begin{responsory-breve}[ostende-nobis]
{Ostende nobis domine * Misericordiam tuam.}{Et salutare tuum da nobis.}
\end{responsory-breve}%
\V Memento nostri domine in beneplacito populi tui. \R Visita nos in salutari tuo.
\newline {\color{Red} Ad Nonam. Capitulum.}
\lettrine[lines=2]{\zallmancaps \color{Blue} E}{rit} in novissimis diebus preparatus mons domus domini in vertice montium, et elevabitur super colles et fluent ad eum omnes gentes.
\begin{responsory-breve}[super-te-hierusalem]
{Super te Hierusalem, * Orietur dominus.}{Et gloria eius in te videbitur.}
\end{responsory-breve}%
\V Domine deus virtutum converte nos. \R Et ostende faciem tuam et salvi erimus.
\newline {\color{Red} Ad Vesperas.} \ul{Capitulum, Ymnus, \Vbar, ut supra in primis Vesperis.} {\color{Red} Ad Magnificat. Ant.} Ne timeas Maria invenisti gratiam apud dominum, ecce concipies et paries filium alleluya. 
\newline \ul{Huius diei capitula, Ymni, versiculi, et responsoria horarum cum suis versiculis dicantur suis locis quando agitur de Adventu usque ad vigiliam Natalis Domini exclusive.}
\newline {\color{Red} Feria Secunda. Ad Matutinas. Invitat.}
\begin{invitatory}[regem-venturum]
{Regem venturum dominum * Venite adoremus.}
\end{invitatory}
\newline \ul{Hoc invitatorium dicatur cotidie per Adventum in profestis diebus usque ad Vigiliam Natalis Domini, nisi in ieiuniis quatuor temporum.} \V Ex Syon.
\newline \ul{Lectiones feriales a secundo capitulo usque ad quartum decimum accipe ubi volueris. Et nota quod illi qui non habent Biblias possunt repetere per ferias lectiones dominicales prenotatas. Qui vero habent possunt repetere vel quod melius est legere in capitulis predictis. Et hoc similiter in sequentibus observetur.}
\R \hyperlink{aspiciens}{Aspiciens}. \V Excita domine, sine Gloria Patri. \R \hyperlink{aspiciebam}{Aspiciebam}. \R \hyperlink{missus-est}{Missus est}. {\color{Red} Ad Ben. Ant.} Angelus domini nunciavit Marie et concepit de Spiritu Sancto alleluya. {\color{Red} Ad Primam. Ant.} Veni et libera nos deus noster. {\color{Red} Ad Tertiam. Ant.} Veni domine et noli tardare relaxa facinora plebis tue Israhel. {\color{Red} Ad Sextam. Ant.} Tuam domine excita potentiam et veni ad salvandum nos. {\color{Red} Ad Nonam. Ant.} In tuo adventu erue nos domine.
\newline \ul{Precedentes antiphone dicantur in profestis diebus ad horas per Adventum.} {\color{Red} Ad Mag. Ant.} Hierusalem respice ad orientem et vide alleluya.
\newline {\color{Red} Feria Tertia.} \V Egredietur virga. {\color{Red} Ad Ben. Ant.} Leva Hierusalem oculos et vide potentiam regis, ecce salvator venit solvere te a vinculo. {\color{Red} Ad Mag. Ant.} Querite dominum dum inveniri potest, invocate eum dum prope est alleluya.
\newline {\color{Red} Feria Quarta.} \V Egredietur dominus. 
\newline \ul{Predicti tres versiculi dicantur per Adventum in profestis diebus secundum ordinem nocturnorum.} {\color{Red} Ad Ben. Ant.} De Syon exibit lex et verbum domini de Hierusalem. {\color{Red} Ad Mag. Ant.} Veniet fortior me post me cuius non sum dignus solvere corrigiam calciamentorum eius.
\newline {\color{Red} Feria Quinta.} {\color{Red} Ad Ben. Ant.} Benedicta tu in mulieribus et benedictus fructus ventris tui. {\color{Red} Ad Mag. Ant.} Expectabo dominum salvatorem meum, et prestolabor eum dum prope est alleluya.
\newline {\color{Red} Feria Sexta.} {\color{Red} Ad Ben. Ant.} Ecce veniet deus et homo de domo David sedere in throno alleluya. {\color{Red} Ad Mag. Ant.} Ex Egypto vocavi filium meum veniet ut salvet populum suum.
\newline {\color{Red} Sabbato.} {\color{Red} Ad Ben. Ant.} Syon noli timere ecce deus tuus veniet alleluya.
{\ubsubsection{Dominica Secunda Sabbato Precedenti}{breviarium-dominica-secunda-sabbato-precedenti}}
\par \noindent {\color{Red} Ad Vesperas} \R \hyperlink{ecce-dominus-veniet}{Ecce dominus veniet.} {\color{Red} II. Ad Mag. Ant.} Veni domine visitare nos in pace ut letemur coram te corde perfecto. {\color{Red} Oratio.}
\lettrine[lines=2]{\zallmancaps \color{Red} E}{xcita} domine corda nostra ad preparandas unigeniti tui vias, ut per eius adventum purificatis tibi mentibus servite mereamur. Qui tecum.
\newline {\color{Red} Ad Matutinas. Invitat.} 
\begin{invitatory}[rex-noster-invitatorium]
{Rex noster adveniet Christus * Quem Iohannes predicavit agnum esse venturum.}
\end{invitatory}
\newline \ul{Antiphone, psalmi, et versiculi nocturnorum sicut in precedenti Dominica.}
\newline {\color{Red} Lectio Prima.}
\lettrine[lines=2]{\zallmancaps \color{Blue} P}{rope} est ut veniat tempus eius, et dies eius non elongabuntur. Miserebitur enim dominus Iacob, et eliget adhuc de Israhel, et requiescere eos faciet super humum suam. Adiungetur advena ad eos, et adherebit domui Iacob. Et tenebunt eos populi: et adducent eos in locum suum et possidebit eos domus Israhel super terram domini in servos et ancillas. Et erunt capientes eos qui se ceperant: et subiicient exactores suos. Hec dicit.
\begin{responsory}[hierusalem-cito-veniet]
{Hierusalem cito veniet salus tua, quare merore consumeris, nunquid consiliarius non est tibi quia innovavit te dolor, * Salvabo te et liberabo te noli timere.}{Israhel si me audieris non erit in te deus recens neque adorabis deum alienum ego enim dominus.}
\end{responsory}%
\newline {\color{Red} Lectio Secunda.}
\lettrine[lines=2]{\zallmancaps \color{Blue} E}{t} erit, in die illa cum requiem dederit tibi Deus a labore tuo, et a concussione tua et a servitute dura qua ante servisti: sumes parabolam istam contra regem Babylonis et dices. Quomodo cessavit exactor: quievit tributum? Contrivit dominus baculum impiorum, virgam dominantium cedentem populos in indignatione plaga insanabili, subicientem in furore gentes: persequentem crudeliter. Hec dicit.
\begin{responsory}[ecce-dominus-veniet]
{Ecce dominus veniet et omnes sancti eius cum eo, et erit in die illa lux magna, et exibunt de Hierusalem sicut aqua munda, et regnabit dominus in eternum, * Super omnes gentes.}{Ecce dominus cum virtute veniet et regnum in manu eius et potestas et imperium.}
\end{responsory}%
\newline {\color{Red} Lectio Tertia.}
\lettrine[lines=2]{\zallmancaps \color{Red} C}{onquievit}, et siluit omnis terra: gavisa est et exultavit. Abietes quoque letate sunt super te, et cedri Libani. Ex quo dormisti, non ascendit qui succidat nos. Infernus subter conturbatus est in occursum adventus tui: suscitavit tibi gigantes. Hec dicit.
\begin{responsory-doxology}[civitas-hierusalem]
{Civitas Hierusalem noli flere, quoniam doluit dominus super te, * Et auferet a te omnem tribulacionem.}{Ecce in fortitudine veniet et brachium eius dominabitur.}
\end{responsory-doxology}%
\newline {\color{Red} Lectio Quarta.}
\lettrine[lines=2]{\zallmancaps \color{Blue} O}{mnes} principes terre surrexerunt de soliis suis, omnes principes nationum. Universi respondebunt, et dicent tibi. Et tu vulneratus es sicut et nos, nostri similis effectus es. Detracta est ad inferos superbia tua: concidit cadaver tuum. Subter te sternetur tinea, et operimentum tuum erunt vermes. Hec dicit.
\begin{responsory}[ecce-veniet-dominus]
{Ecce veniet dominus protector noster sanctus Israhel, * Corona regni habens in capite suo.}{Et dominabitur a mari usque ad mare, et a flumine usque ad terminos orbis terrarum.}
\end{responsory}%
\newline {\color{Red} Lectio Quinta.}
\lettrine[lines=2]{\zallmancaps \color{Red} Q}{uomodo} cecidisti de celo Lucifer qui mane oriebaris: corruisti in terram qui vulnerabas gentes? Qui dicebas in corde tuo. In celum conscendam, super astra dei exaltabo solium meum. Sedebo in monte testamenti: in lateribus aquilonis. Ascendam super altitudinem nubium, similis ero altissimo. Veruntamen ad infernum detraheris, in profundum laci. Hec dicit.
\begin{responsory}[sicut-mater]
{Sicut mater consolatur filios suos ita consolabor vos dicit dominus, et de Hierusalem civitate quam elegi veniet vobis auxilium, et videbitis, * Et gaudebit cor vestrum.}{Dabo in Syon salutem et in Hierusalem gloriam meam.}
\end{responsory}%
\newline {\color{Red} Lectio Sexta.}
\lettrine[lines=2]{\zallmancaps \color{Blue} Q}{ui} te viderint ad te inclinabuntur, teque prospicient. Nunquid iste est vir qui conturbavit terram, qui concussit regna qui posuit orbem desertum, et urbes eius destruxit, vinctis eius non aperuit carcerem? Omnes reges gentium universi dormierunt in gloria: vir in domo sua. Hec dicit.
\begin{responsory-doxology}[hierusalem-plantabis]
{Hierusalem plantabis vineam in montibus tuis et exultabis quia dies domini veniet, surge Syon convertere ad deum tuum, gaude et letare Iacob, * Quia de medio gentium salvator tuus veniet.}{Exulta satis filia Syon, iubila filia Hierusalem.}
\end{responsory-doxology}%
\newline {\color{Red} Lectio VII. Secundum Lucam.}
% Same as Roman Prima Adventus
\lettrine[lines=2]{\zallmancaps \color{Red} I}{n} illo tempore: Dixit Ihesus discipulis suis. Erunt signa in sole et luna et stellis, et in terris pressura gentium, pre confusione sonitus maris et fluctuum, arescentibus hominibus pre timore et expectatione que supervenient universo orbi. Et reliqua.
\newline {\color{Red} Omelia Beati Gregorii Pape.} \grecross
\lettrine[lines=2]{\zallmancaps \color{Blue} D}{ominus} ac redemptor noster paratos nos invenire desiderans, senescentem mundum que mala sequantur denunciat: ut nos ab eius amore compescat. Appropinquantem eius terminum quante persecutiones preveniant innotescit ut si Deum metuere in tranquillitate nolumus: vicinum eius iudicium vel percussionibus attriti timeamus.
\begin{responsory}[egredietur-dominus-de-samaria]
{Egredietur dominus de Samaria ad portam que respicit ad orientem et veniet in Bethleem ambulans super aquas redemptionis vide, * Tunc salvus erit omnis homo quia ecce veniet.}{Et preparabitur in misericordia solium eius et sedebit super illud iudicans in equitate.}
\end{responsory}%
\newline {\color{Red} Lectio VIII.}
\lettrine[lines=2]{\zallmancaps \color{Red} H}{inc} lectioni sancti evangelii paulo superius dominus premisit dicens. Exurget gens contra gentem, et regnum terremotus magni per loca, et pestilentie et fames. Et quibusdam interpositis, hoc quod modo audistis adiunxit. Et erunt signa in sole et luna et stellis.
\begin{responsory}[docebit-nos]
{Docebit nos dominus vias suas, et ambulabimus in semitis eius, * Quia de Syon exibit lex et verbum domini de Hierusalem.}{Venite ascendamus ad montem domini, et ad domum dei Iacob.}
\end{responsory}%
\newline {\color{Red} Lectio IX.}
\lettrine[lines=2]{\zallmancaps \color{Blue} E}{x} his profecto omnibus alia iam facta cernimus, alia in proximo ventura formidamus. Nam gentem supra gentem exurgere, earumque pressuram terris insistere, plus iam in nostris tribulationibus quam in codicibus legimus. Quod terremotus urbes innumeras subruat ex aliis mundi partibus scitis quam frequenter audivimus. Pestilentias sine cessatione patimur.
\begin{responsory-doxology}[rex-noster]
{Rex noster adveniet Christus * Quem Iohannes predicavit agnum esse venturum.}{Super ipsum continebunt reges os suum ipsum gentes deprecabuntur.}Rex.
\end{responsory-doxology}%
\newline {\color{Red} In Laudibus. Ant.} Ecce in nubibus celi dominus veniet cum potestate magna alleluya. {\color{Red} Ps.} \psalm{92} {\color{Red} Ant.} Urbs fortitudinis nostre Syon, salvator ponetur in ea murus et ante murale, aperite portas quia nobiscum deus alleluya. {\color{Red} Ps.} \psalm{99} {\color{Red} Ant.} Ecce apparebit dominus et non mentietur, si moram fecerit expecta eum quia veniet et non tardabit, alleluya. {\color{Red} Ps.} \psalm{62} {\color{Red} Ant.} Montes et colles cantabunt coram deo laudem, et omnia ligna silvarum plaudent manibus, quia veniet dominus dominator in regnum eternum, alleluya alleluya. {\color{Red} Can.} \psalm{benedicite} {\color{Red} Ant.} Ecce dominus noster cum virtute veniet, et illuminabit oculos servorum suorum alleluya. {\color{Red} Ps.} \psalm{148}
\newline {\color{Red} Ad Ben. Ant.} Super solium David et super regnum eius sedebit in eternum alleluya.
\newline {\color{Red} Ad Mag. Ant.} Beata es Maria que credidisti, perficientur in te que dicta sunt tibi a domino alleluya.
\newline \ul{Lectiones feriales a decimo sexto capitulo usque ad tricesimum ubi volueris.}
\newline {\color{Red} Feria Secunda. Ad Benedictus. Ant.} De celo veniet dominator dominus et in manu eius honor et imperium. {\color{Red} Ad Mag. Ant.} Ecce rex veniet dominus terre, et ipse auferet iugum captivitatis nostre.
\newline {\color{Red} Feria Tertia. Ad Ben. Ant.} Super te Hierusalem orietur dominus, et gloria eius in te videbitur. {\color{Red} Ad Mag. Ant.} Vox clamantis in deserto parate viam domini rectas facite semitas dei nostri.
\newline {\color{Red} Feria Quarta. Ad Ben. Ant.} Ecce mitto angelum meum qui preparabit viam tuam ante faciem tuam. {\color{Red} Ad Mag. Ant.} Syon renovaberis et videbis iustum tuum qui venturus est in te.
\newline {\color{Red} Feria Quinta. Ad Ben. Ant.} Tu es qui venturus es domine quem expectamus ut salvum facias populum tuum. {\color{Red} Ad Mag. Ant.} Qui post me venit ante me factus est, cuius non sum dignus calciamenta solvere.
\newline {\color{Red} Feria Sexta. Ad Ben. Ant.} Dicite pusillanimes confortamini, ecce dominus deus noster veniet. {\color{Red} Ad Mag. Ant.} Cantate domino canticum novum laus eius ab extremis terre.
\newline {\color{Red} Sabbato. Ad Ben. Ant.} Levabit dominus signum in nationibus et congregabit dispersos Israhel.
{\ubsubsection{Dominica Tercia Sabbato Precedenti}{breviarium-dominica-tercia-sabbato-precedenti}}
\par \noindent {\color{Red} Ad Vesperas} \R \hyperlink{qui-venturus}{Qui venturus.} {\color{Red} III. Ad Mag. Ant.} Ante me non est formatus deus, et post me non erit quia michi curvabitur omne genu et confitebitur omnis lingua. {\color{Red} Oratio.}
\lettrine[lines=2]{\zallmancaps \color{Red} A}{urem} tuam quesumus domine precibus nostris accommoda et mentis nostre tenebras gratia tue visitationis illustra. Qui vivis.
\newline {\color{Red} Ad Matutinas. Invitat.} 
\begin{invitatory}[surgite-vigilemus]
{Surgite vigilemus, venite adoremus, * Quia nescitis horam quando veniet dominus.}
\end{invitatory}
\newline \ul{Antiphone, psalmi, et versiculi nocturnorum sicut in Prima Dominica.}
\newline {\color{Red} Lectio Prima.}
\lettrine[lines=2]{\zallmancaps \color{Blue} H}{ec} dicit dominus deus sanctus Israhel: Si revertamini et quiescatis: salvi eritis. In silentio et in spe: erit fortitudo vestra. Et noluistis et dixistis. Nequaquam: sed ad equos fugiemus. Ideo fugietis. Et super veloces ascendemus. Ideo veloces erunt qui persequentur vos. Hec dicit.
\begin{responsory}[ecce-apparebit]
{Ecce apparebit dominus super nubem candidam et cum eo sanctorum milia, habens in vestimento et in femore suo scriptum, * Rex regum et dominus dominantium.}{Apparebit in finem et non mentietur si moram fecerit expecta eum quia veniens veniet.}
\end{responsory}%
\newline {\color{Red} Lectio Secunda.}
\lettrine[lines=2]{\zallmancaps \color{Red} M}{ille} homines a facie terroris unius, et a facie terroris quinque fugietis, donec relinquamini quasi malus navis in vertice montis, et quasi signum super collem. Propterea expectat dominus ut misereatur vestri, et ideo exaltabitur parcens vobis, quia deus iudicii dominus. Beati omnes qui expectant eum. Hec dicit.
\begin{responsory}[bethleem-civitas]
{Bethleem civitas dei summi ex te exiet dominator Israhel, et egressus eius sicut a principio a diebus eternitatis, et magnificabitur in medio universe terre et pax erit in terra nostra, * Dum venerit.}{Loquetur pacem gentibus et potestas eius a mari usque ad mare.}
\end{responsory}%
\newline {\color{Red} Lectio Tertia.}
\lettrine[lines=2]{\zallmancaps \color{Blue} P}{opulus} enim Syon habitabit in Hierusalem. Plorans nequaquam plorabis. Miserans miserebitur tui. Ad vocem clamoris tui, statim ut audierit respondebit tibi. Et dabit dominus vobis panem artum et aquam brevem, et non faciet
%pp089
avolare a te ultra doctorem tuum. Hec dicit.
\begin{responsory-doxology}[qui-venturus]
{Qui venturus est veniet et non tardabit iam non erit timor in finibus nostris, * Quoniam ipse est salvator noster.}{Deponet omnes iniquitates nostras, et proiciet omnia peccata nostra.}
\end{responsory-doxology}%
\newline {\color{Red} Lectio Quarta.}
\lettrine[lines=2]{\zallmancaps \color{Red} E}{t} erunt oculi tui videntes preceptorem tuum, et aures tue audient verbum post tergum monentis. Hec via ambulate in ea, neque ad dexteram neque ad sinistram. Et contaminabis laminas sculptilium argenti tui: et vestimentum conflatilis auri tui, et disperges ea sicut immundiciam menstruate. Egredere, dices ei. Hec dicit.
\begin{responsory}[suscipe-verbum]
{Suscipe verbum virgo Maria quod tibi a domino per angelum transmissum est, concipies et paries deum pariter et hominem, * Ut benedicta dicaris inter omnes mulieres.}{Paries quidem filium et virginitatis non patieris detrimentum efficieris gravida et eris mater semper intacta.}
\end{responsory}%
\newline {\color{Red} Lectio Quinta.}
\lettrine[lines=2]{\zallmancaps \color{Blue} E}{t} dabitur pluvia semini tuo, ubicunque seminaveris in terra, et panis frugum terre erit uberrimus et pinguis. Pascetur in possessione tua in die illa agnus spaciose. Et tauri tui et pulli asinorum qui operantur terram, commixtum migma comedent, sicut in area ventilatum est. Hec dicit.
\begin{responsory}[egypte-noli-flere]
{Egypte noli flere, quia dominator tuus veniet tibi ante cuius conspectum movebuntur abyssi, * Liberare populum suum de manu potentie.}{Ecce dominator dominus cum virtute veniet.}
\end{responsory}%
\newline {\color{Red} Lectio Sexta.}
\lettrine[lines=2]{\zallmancaps \color{Red} E}{t} erunt super omnem montem excelsum et super omnem collem elevatum, rivi currentium aquarum in die interfectionis multorum, cum ceciderint turres. Et erit lux lune sicut lux solis, et lux solis erit septempliciter sicut lux septem dierum, in die qua alligaverit dominus vulnus populi sui, et percussuram plage eius sanaverit. Hec dicit.
\begin{responsory-final}[prope-est]
{Prope est ut veniat tempus eius, et dies eius non elongabuntur, * Miserebitur dominus Iacob, * Et Israhel salvabitur.}{Qui venturus est veniet et non tardabit, iam non erit timor in finibus nostris.}{Et Israhel}
\end{responsory-final}%
\newline {\color{Red} Lectio VII. Secundum Matheum.}
% Same as Roman Secunda Adventus
\lettrine[lines=2]{\zallmancaps \color{Blue} I}{n} illo tempore: Cum audisset Iohannes in vinculis opera Christi, mittens duos de discipulis suis ait illi. Tu es qui venturus es an alium expectamus? Et reliqua.
\newline {\color{Red} Omelia Beati Gregorii Pape.} \grecross
\lettrine[lines=2]{\zallmancaps \color{Red} Q}{uerendum} nobis est fratres karissimi, Iohannes propheta et plus quam propheta, qui venientem ad Iordanis baptisma dominum ostendit dicens, ecce agnus dei, ecce qui tollit peccatum mundi, cur in carcere positus mittens discipulos requirit, tu es qui venturus es an alium expectamus: tanquam si ignoret quem ostenderat, et an ipse sit nesciat quem ipsum esse prophetando, baptizando, ostendendo clamaverat.
\begin{responsory}[descendet-dominus]
{Descendet dominus sicut pluvia in vellus, * Orietur in diebus eius iusticia et abundantia pacis.}{Et adorabunt eum omnes reges omnes gentes servient ei.}
\end{responsory}%
\newline {\color{Red} Lectio VIII.}
\lettrine[lines=2]{\zallmancaps \color{Blue} S}{et} hec citius questio solvitur, si geste rei tempus et ordo pensetur. Ad Iordanis enim fluenta positus, quia ipse redemptor mundi esset asseruit, missus vero in carcere an ipse veniat requirit, non quia ipsum esse mundi redemptorem dubitet: sed querit ut sciat si is qui per se in mundum venerat, per se eciam ad inferni claustra descendat. Quem enim precurrens mundo nunciaverat, hunc moriendo ad inferos precurrebat.
\begin{responsory}[veni-domine]
{Veni domine et noli tardare relaxa facinora plebis tue, * Et revoca dispersos in terram suam.}{Excita domine potentiam tuam et veni ut salvos facias nos.}
\end{responsory}%
\newline {\color{Red} Lectio IX.}
\lettrine[lines=2]{\zallmancaps \color{Red} A}{it} ergo. Tu es qui venturus es an alium expectamus? Ac si aperte dicat. Sicut pro hominibus nasci dignatus es, an eciam pro hominibus mori et ad inferos descendere digneris insinua: ut qui nativitatis tue precursor extiti, mortis eciam precursor fiam, et venturum inferno te nuntiem, quem iam venisse mundo nuntiavi.
\begin{responsory-doxology}[ecce-radix-iesse]
{Ecce radix Iesse ascendet in salutem populorum ipsum gentes deprecabuntur, * Et erit nomen eius gloriosum.}{Erit iusticia cingulum lumborum eius et fides cinctorum renum eius.}Ecce.
\end{responsory-doxology}%
\newline {\color{Red} In Laudibus. Ant.} Veniet dominus et non tardabit, et illuminabit abscondita tenebrarum, et manifestabit se ad omnes gentes, alleluya. {\color{Red} Ps.} \psalm{92} {\color{Red} Ant.} Hierusalem gaude gaudio magno quia veniet tibi salvator alleluya. {\color{Red} Ps.} \psalm{99} {\color{Red} Ant.} Dabo in Syon salutem et in Hierusalem gloriam meam alleluya. {\color{Red} Ps.} \psalm{62} {\color{Red} Ant.} Montes et omnes colles humiliabuntur et erunt prava in directa, et aspera in vias planas veni domine et noli tardare alleluya. {\color{Red} Can.} \psalm{benedicite} {\color{Red} Ant.} Iuste et pie vivamus expectantes beatam spem et adventum domini. {\color{Red} Ps.} \psalm{148}
\newline {\color{Red} Ad Ben. Ant.} Iohannes autem cum audisset in vinculis opera Christi mittens duos de discipulis suis ait illi tu es qui venturus es an alium expectamus.
\newline {\color{Red} Ad Mag. Ant.} Ite dicite Iohanni ceci vident, et surdi audiunt, claudi curantur, leprosi mundantur. {\color{Red} vel} \hyperlink{o-antiphons}{O}.
\newline \ul{Quarta feria post istam Dominicam semper ieiunium quatuor temporum celebratur. Sciendum est quod antiphone Laudum que infra notate sunt ante antiphonas que incipiunt per} \hyperlink{o-antiphons}{O}, \ul{debent dici per sex ferias proximas ante Vigiliam Nativitatis Domini suis propriis diebus. Ille scilicet que pertinent ad secundam feriam in secunda feria et sic de aliis. Laudes autem de feria in qua celebrabitur festum Sancti Thome intermittentur illo anno. Lectiones pro feriis a trigesimo tertio capitulo usque ad quadragesimum quartum ad libitum.}
\newline {\color{Red} Feria Secunda. Ad Benedictus. Ant.} Egredietur virga de radice Iesse et replebitur omnis terra gloria domini et videbit omnis caro salutare dei. {\color{Red} Ad Mag. Ant. } Elevare elevare consurge Hierusalem solve vincla colli tui captiva filia Syon. {\color{Red} vel} \hyperlink{o-antiphons}{O}.
\newline {\color{Red} Feria Tertia. Ad Ben. Ant.} Tu Bethleem terra Iuda, non eris minima ex te enim exiet dux qui regat populum meum Israhel. {\color{Red} Ad Mag. Ant.} Erumpant montes iocunditatem et colles iusticiam quia lux mundi dominus cum potentia venit. {\color{Red} vel} \hyperlink{o-antiphons}{O}.
\newline {\color{Red} Feria Quarta. Ad Matutinas. Invitat.}
\begin{invitatory}[prope-est-iam-invitatorium]
{Prope est iam dominus, * Venite adoremus.}
\end{invitatory}
\newline {\color{Red} Lectio Prima. Secundum Lucam.}
\lettrine[lines=2]{\zallmancaps \color{Blue} I}{n} illo tempore: Missus est angelus Gabriel a Deo, in civitatem Galilee cui nomen Nazareth, ad virginem desponsatam viro cui nomen erat Ioseph de domo David, et nomen virginis Maria. Et reliqua.
\newline {\color{Red} Omelia Beati Ambrosii Episcopi.}
\lettrine[lines=2]{\zallmancaps \color{Red} L}{atent} quidem divina misteria, nec facile iuxta propheticum dictum, quisquam hominum potest scire consilium dei. Sed tamen ex ceteris factis atque preceptis domini salvatoris possumus intelligere, et propensioris fuisse consilii, quod ea potissimum electa est ut dominum pareret, que erat desponsata viro.%typo iusta
\begin{responsory}[clama-in-fortitudine]
{Clama in fortitudine qui annuncias pacem in Hierusalem, * Dic civitatibus vide et habitatoribus Syon, ecce deus noster quem expectabamus adveniet.}{Super montem excelsum ascende tu qui evangelizas Syon exalta in fortitudine vocem tuam.}
\end{responsory}%
\newline {\color{Red} Lectio Secunda.}
\lettrine[lines=2]{\zallmancaps \color{Blue} C}{ur} autem non ante quam desponsaretur impleta est? Fortasse ne diceretur quod conceperat ex adulterio. Et bene utrunque posuit scriptura, ut et desponsata esset et virgo. Virgo: ut expers virilis consortii videretur. Desponsata: ne temerate virginitatis adureretur infamia, cui gravis alvus corruptele videretur insigne preferre.
\begin{responsory}[orietur-stella]
{Orietur stella ex Iacob et exurget homo de Israhel, et confringet omnes duces alienigenarum, * Et erit omnis terra possessio eius.}{Et adorabunt eum omnes reges terre omnes gentes servient ei.}
\end{responsory}%
\newline {\color{Red} Lectio Tertia.}
\lettrine[lines=2]{\zallmancaps \color{Red} M}{aluit} autem dominus aliquos de suo ortu, quam de matris pudore dubitare. Sciebat enim teneram esse virginis verecundiam, et lubricam famam pudoris. Nec putavit ortus sui fidem matris iniuriis astruendam. Servatur itaque Sancte Marie sicut pudore integra, ita inviolabilis opinione virginitas. Oportet enim sanctos et ab his testimonium habere qui foris sunt.
\begin{responsory-doxology}[egredietur-dominus-et-preliabitur]
{Egredietur dominus et preliabitur contra gentes, * Et stabunt pedes eius supra montes olivarum ad orientem.}{Et elevabitur super omnes colles et fluent ad eum omnes gentes.}
\end{responsory-doxology}%
\newline \ul{Predicta tria responsoria et alia que sequuntur feria sexta et Sabbato, dici tantum debent in diebus in quibus notata sunt, nisi festum Sancti Thome in sequenti sexta feria vel Sabbato evenerit, et tunc responsoria illius diei in quinta feria dicantur. Omelia vero diei in quo festum evenerit dimittatur.}
\newline {\color{Red} Ad Ben. Ant.} Missus est Gabriel angelus ad Mariam virginem desponsatam Ioseph. {\color{Red} Oratio.}
\lettrine[lines=2]{\zallmancaps \color{Blue} F}{estina} quesumus domine ne tardaveris, et auxilium nobis superne virtutis impende ut adventus tui consolationibus subleventur, qui in tua pietate confidunt. Qui vivis.
\newline \ul{Hec oracio dicatur ad Laudes tantum, ad alias autem horas oracio Dominicalis: idem fiat de oracionibus que sequuntur feria sexta et Sabbato.}
\newline {\color{Red} Ad Mag. Ant.} Quomodo fiet istud angele dei, qui virum in concipiendo non pertuli, audi Maria virgo Christi, spiritus sanctus superveniet in te et virtus altissimi obumbrabit tibi. {\color{Red} vel} \hyperlink{o-antiphons}{O}.
\newline {\color{Red} Feria Quinta. Invitat.} \hyperlink{regem-venturum}{Regem venturum}.
\newline \ul{Responsoria tercii nocturni de hystoria Dominicali nisi festum Sancti Thome in sexta feria vel in Sabbato evenerit.} {\color{Red} Ad Ben. Ant.} Vigilate animo, in proximo est dominus deus noster. {\color{Red} Ad Mag. Ant.} Letamini cum Hierusalem et exultate in ea omnes qui diligitis eam in eternum. {\color{Red} vel} \hyperlink{o-antiphons}{O}.
\newline {\color{Red} Feria Sexta. Invitat.} \hyperlink{prope-est-iam-invitatorium}{Prope}.
\newline {\color{Red} Lectio Prima. Secundum Lucam.}
\lettrine[lines=2]{\zallmancaps \color{Blue} I}{n} illo tempore: Exurgens Maria abiit in montana cum festinatione in civitatem Iuda, et intravit in domum Zacharie et salutavit Elysabeth. Et reliqua.
\newline {\color{Red} Omelia Venerabilis Bede Presbiteri.}
\lettrine[lines=2]{\zallmancaps \color{Red} L}{ectio} quam audivimus sancti evangelii, et redemptionis nostre nobis semper veneranda primordia predicat, et salutaria semper imitande humilitatis remedia commendat. Nam quia peste superbie attactum genus humanum perierat, decebat ut medicamentum humilitatis quo sanaretur prima mox incipientis salutis tempora pretenderent.
\begin{responsory}[precursor-pro-nobis]
{Precursor pro nobis ingreditur agnus sine macula, secundum ordinem Melchisedech pontifex factus est, * In eternum et in seculum seculi.}{Ipse est rex iusticie cuius generacio non habet finem.}
\end{responsory}%
\newline {\color{Red} Lectio Secunda.}
\lettrine[lines=2]{\zallmancaps \color{Blue} E}{t} quia per temeritatem seducte mulieris mors in mundum introierat, congruum fuit ut in indicium vite revertentis, mulieres devote se humilitatis invicem ac pietatis prevenirent obsequiis.
\begin{responsory}[modo-veniet]
{Modo veniet dominator dominus * Et nomen eius Emmanuel vocabitur.}{Orietur in diebus eius iusticia et abundantia pacis.}
\end{responsory}%
\newline {\color{Red} Lectio Tertia.}
\lettrine[lines=2]{\zallmancaps \color{Red} P}{rior} ergo nobis beata dei genitrix ad sublimitatem patrie celestis iter ostendit humilitatis, non minus religionis quam castitatis exemplo venerabilis. Siquidem gloria virginei et intemerati corporis, qualis sit vita superne civitatis ad quam suspiramus insinuat, ubi neque nubent neque nubentur, sed sunt sicut angeli dei in celis, ac virtutem mentis eximiam qua ad hanc pertingere debeamus indicat.
\begin{responsory-doxology}[videbunt-gentes]
{Videbunt gentes iustum tuum, et cuncti reges inclitum tuum, * Et vocabitur tibi nomen novum quod os domini nominavit.}{Et erit corona glorie in manu domini, et diadema regni in manu dei tui.}
\end{responsory-doxology}%
{\color{Red} Ad Ben. Ant.} Ex quo facta est vox salutationis tue in auribus meis exultavit in gaudio infans in utero meo alleluya. {\color{Red} Oratio.}
\lettrine[lines=2]{\zallmancaps \color{Blue} E}{xcita} quesumus domine potentiam tuam et veni, ut hi qui in tua pietate confidunt, ab omni citius adversitate liberentur. Qui vivis.
\newline {\color{Red} Ad Mag. Ant.} Beata es Maria que credidisti perficientur in te que dicta sunt tibi a domino alleluya. {\color{Red} vel} \hyperlink{o-antiphons}{O}.
\newline {\color{Red} Sabbato. Invitat.} \hyperlink{prope-est-iam-invitatorium}{Prope est}.
\newline {\color{Red} Lectio Prima. Secundum Lucam.}
\lettrine[lines=2]{\zallmancaps \color{Red} A}{nno} quintodecimo imperii Tyberii Cesaris, procurante Pontio Pilato Iudeam, tetrarcha autem Galilee Herode, Philippo autem fratre eius tetrarcha Yturee, et Traconitidis regionis, et Lysania Abiline tetrarcha, sub principibus sacerdotum Anna et Caypha: factum est verbum domini super Iohannem Zacharie filium in deserto. Et reliqua.
\newline {\color{Red} Omelia Beati Gregorii Pape.}
\lettrine[lines=2]{\zallmancaps \color{Blue} R}{edemptoris} precursor quo tempore verbum predicationis acceperit, memorato romane rei publice principe, et Iudee regibus designatur. Quia enim illum predicare veniebat, qui ex Iudea quosdam et multos ex gentibus redempturus erat: per regem gentium et principes Iudeorum predicationis eius tempora designantur.
\begin{responsory}[emitte-agnum]
{Emitte agnum domine dominatorem terre de petra deserta * Ad montem filie Syon.}{Ex Syon species decoris eius deus noster manifeste veniet.}
\end{responsory}%
\newline {\color{Red} Lectio Secunda.}
\lettrine[lines=2]{\zallmancaps \color{Red} Q}{uia} autem gentilitas colligenda erat, et Iudea pro culpa perfidie dispergenda: ipsa quoque descriptio terreni principatus ostendit, quoniam et in Romana re publica unus prefuisse describitur, et in Iudee regno per quartam partem plurimi principabantur. Voce etenim nostri redemptoris dicitur. Omne regnum in semet ipsum divisum desolabitur.
\begin{responsory}[germinaverunt-campi]
{Germinaverunt campi heremi germen odoris Israhel, quia ecce deus noster cum virtute veniet, * Et splendor eius cum eo.}{Ecce dominator dominus cum virtute veniet.}
\end{responsory}%
\newline {\color{Red} Lectio Tertia.}
\lettrine[lines=2]{\zallmancaps \color{Blue} L}{iquet} ergo quod ad finem regni Iudea pervenerat, que tot regibus divisa subiacebat. Apte quoque non solum quibus regibus sed eciam sub quibus sacerdotibus actum sit, demonstratur: ut quia illum Iohannes Baptista predicabat, qui simul rex et sacerdos existeret: Lucas evangelista predicationis eius tempora per regnum et sacerdotium designavit.
\begin{responsory-final}[radix-iesse]
{Radix Iesse qui exurget iudicare gentes, in eum gentes sperabunt: Et erit nomen eius benedictum * In secula.}{Super quem continebunt reges os suum, ipsum gentes deprecabuntur.}{In secula}
\end{responsory-final}%
{\color{Red} Ad Ben. Ant.} Omnis vallis implebitur et omnis mons et collis humiliabitur et videbit omnis caro salutare dei. {\color{Red} Oracio.}
\lettrine[lines=2]{\zallmancaps \color{Red} D}{eus} qui conspicis quia ex nostra pravitate affligimur, concede propitius ut ex tua visitatione consolemur. Qui vivis.
{\ubsubsection{Dominica Quarta Sabbato Precedenti}{breviarium-dominica-quarta-sabbato-precedenti}}
\par \noindent {\color{Red} Ad Vesperas} \R \hyperlink{non-auferetur}{Non auferetur.} {\color{Red} II. Ad Mag. Ant.} \hyperlink{o-antiphons}{O}. {\color{Red} Oracio.}
\lettrine[lines=2]{\zallmancaps \color{Blue} E}{xcita} domine potentiam tuam et veni, et magna nobis virtute succurre, ut per auxilium gratie tue quod nostra peccata prepediunt indulgentia tue propitiationis acceleret. Qui vivis.
\newline {\color{Red} Ad Matutinas. Invitat.} 
\begin{invitatory}[ecce-iam-venit-invitatorium]
{Ecce iam venit plenitudo temporis in quo misit deus filium suum natum de virgine factum sub lege, * Venite adoremus.}
\end{invitatory}
\newline \ul{Antiphone, psalmi, et versiculi nocturnorum sicut in Prima Dominica, excepto quod tercius versiculus non dicetur si Vigilia Natalis Domini hac die evenerit.}
\newline {\color{Red} Lectio Prima.}
\lettrine[lines=2]{\zallmancaps \color{Red} A}{udite} me qui sequimini quod iustum est et queritis dominum. Attendite ad petram unde excisi estis, et ad cavernam laci de qua precisi estis. Attendite ad Abraham patrem vestrum, et ad Saram que peperit vos: Quia unum vocavi eum, et benedixi ei, et multiplicavi eum. Hec dicit.
\begin{responsory}[canite-tuba]
{Canite tuba in Syon vocate gentes, annunciate populis et dicite, * Ecce deus salvator noster adveniet.}{Annunciate in finibus terre et in insulis que procul sunt dicite.}
\end{responsory}%
\newline {\color{Red} Lectio Secunda.}
\lettrine[lines=2]{\zallmancaps \color{Blue} C}{onsolabitur} ergo dominus Syon, et consolabitur omnes ruinas eius. Et ponet desertum eius quasi delicias, et solitudinem eius quasi ortum domini. Gaudium et leticia invenietur in ea, gratiarum actio et vox laudis. Attendite ad me popule meus, et tribus mea me audite, quia lex a me exiet, et iudicium meum in lucem populorum requiescet. Hec dicit.
\begin{responsory}[non-auferetur]
{Non auferetur sceptrum de Iuda et dux de femore eius, donec veniat qui mittendus est, * Et ipse erit expectatio gentium.}{Pulchriores sunt oculi eius vino et dentes eius lacte candidiores.}
\end{responsory}%
\newline {\color{Red} Lectio Tertia.}
\lettrine[lines=2]{\zallmancaps \color{Red} P}{rope} est iustus meus, egressus est salvator meus, et brachia mea populos iudicabunt. Me insule expectabunt, et brachium meum sustinebunt. Levate in celum oculos vestros, et videte sub terra deorsum, quia celi sicut fumus liquescent, et terra sicut vestimentum atteretur, et habitatores eius sicut hec interibunt. Salus autem mea in sempiternum erit, et iusticia mea non deficiet. Hec dicit.
\begin{responsory-doxology}[me-oportet]
{Me oportet minui, illum autem crescere, qui post me venit ante me factus est, * Cuius non sum dignus corrigiam calciamenti eius solvere.}{Ego baptizo vos aqua ille autem baptizabit vos spiritu sancto.}
\end{responsory-doxology}%
\newline {\color{Red} Lectio Quarta.}
\lettrine[lines=2]{\zallmancaps \color{Blue} A}{udite} me qui scitis iustum, populus, lex mea in corde eorum. Nolite timere obprobrium hominum, et blasphemias eorum ne metuatis. Sicut enim vestimentum sic comedet eos vermis, et sicut lanam sic devorabit eos tinea. Salus autem mea in sempiternum erit, et iusticia mea in generationes generationum. Hec dicit.
\begin{responsory}[ecce-iam-venit]
{Ecce iam venit plenitudo temporis, in quo misit deus filium suum in terris natum de virgine factum sub lege, * Ut eos qui sub lege erant redimeret.}{Propter nimiam caritatem suam qua dilexit nos deus, filium suum misit in similitudinem carnis peccati.}
\end{responsory}%
\newline {\color{Red} Lectio Quinta.}
\lettrine[lines=2]{\zallmancaps \color{Red} C}{onsurge} consurge, induere fortitudinem brachium domini. Consurge sicut in diebus antiquis, in generationibus seculorum. Nunquid non tu percussisti superbum, vulnerasti draconem? Nunquid non tu siccasti mare, aquam abyssi vehementis? Qui posuisti profundum maris viam, ut transirent liberati. Hec dicit.
\begin{responsory}[virgo-israel]
{Virgo Israhel revertere ad civitates tuas, usquequo dolens averteris, generabis dominum salvatorem oblationem novam in terra, ambulabunt homines, * In salvationem.}{In caritate perpetua dilexi te, ideo attraxi te miserans.}
\end{responsory}%
\newline {\color{Red} Lectio Sexta.}
\lettrine[lines=2]{\zallmancaps \color{Blue} E}{t} nunc qui redempti sunt a domino revertentur, et venient in Syon laudantes, et leticia
%pp090
sempiterna super capita eorum. Gaudium et leticiam tenebunt, fugiet dolor et gemitus. Ego ego ipse consolabor vos. Hec dicit.
\begin{responsory-doxology}[iuravi-dicit]
{Iuravi dicit dominus ut ubera iam non irascar super terram, montes enim et colles suscipient iusticiam meam, et testamentum pacis erit, * In Hierusalem.}{Iuxta est salus mea ut veniat, et iusticia mea ut reveletur, in Hierusalem.}
\end{responsory-doxology}%
\newline \ul{Quando Vigilia Natalis Domini in hac Dominica evenerit tres ultime lectiones fiant de omelia vigilie, haec autem omelia non legitur illo anno.}
\newline {\color{Red} Lectio VII. Secundum Iohannem.}
% Same as Roman Tertia Adventus
\lettrine[lines=2]{\zallmancaps \color{Red} I}{n} illo tempore: Miserunt Iudei ab Iherosolimis sacerdotes et Levitas ad Iohannem, ut interrogarent eum, tu quis es. Et reliqua.
\newline {\color{Red} Omelia Beati Gregorii Pape.}
\lettrine[lines=2]{\zallmancaps \color{Blue} E}{x} huius nobis lectionis verbis fratres karissimi Iohannis humilitas commendatur. Qui cum tante virtutis esset ut Christus credi potuisset, elegit solide subsistere in se, ne humana opinione raperetur inaniter super se.
\begin{responsory}[non-discedimus]
{Non discedimus a te, vivificabis nos domine, et nomen tuum invocabimus, ostende nobis faciem tuam, * Et salvi erimus.}{Memento nostri domine in beneplacito populi tui, visita nos in salutari tuo.}
\end{responsory}%
\newline {\color{Red} Lectio VIII.}
\lettrine[lines=2]{\zallmancaps \color{Red} N}{am} confessus est et non negavit. Confessus est, quia non sum ego Christus. Sed qui dixit, non sum, negavit plane quod non erat, sed non negavit quod erat, ut veritatem loquens eius membrum fieret: cuius sibi nomen fallaciter non usurparet.
\begin{responsory}[intuemini-quantus]
{Intuemini quantus sit iste qui ingreditur ad salvandas gentes, ipse est rex iusticie, * Cuius generatio non habet finem.}{Precursor pro nobis ingreditur secundum ordinem Melchisedech pontifex factus est in eternum.}
\end{responsory}%
\newline {\color{Red} Lectio IX.}
\lettrine[lines=2]{\zallmancaps \color{Blue} C}{um} ergo non vult appetere nomen Christi, factus est membrum Christi, quia dum infirmitatem suam studuit humiliter agnoscere, illius celsitudinem meruit veraciter obtinere.
\begin{responsory-doxology}[montes-israel]
{Montes Israhel ramos vestros expandite, et florete, et fructus facite, * Prope est ut veniat dies domini.}{Rorate celi de super et nubes pluant iustum, aperiatur terra et germinet salvatorem.}
\end{responsory-doxology}%
Montes.
\newline {\color{Red} In Laudibus. Ant.} Canite tuba in Syon quia prope est dies domini ecce veniet ad salvandum nos alleluya alleluya. {\color{Red} Ps.} \psalm{92} {\color{Red} Ant.} Ecce veniet desideratus cunctis gentibus et replebitur gloria domus domini alleluya. {\color{Red} Ps.} \psalm{99} {\color{Red} Ant.} Erunt prava in directa et aspera in vias planas, veni domine et noli tardare alleluya. {\color{Red} Ps.} \psalm{62} {\color{Red} Ant.} Dominus veniet, occurrite illi, dicentes magnum principium et regni eius non erit finis, deus fortis dominator, princeps pacis alleluya alleluya. {\color{Red} Can.} \psalm{benedicite} {\color{Red} Ant.} Omnipotens sermo tuus domine a regalibus sedibus venit alleluya. {\color{Red} Ps.} \psalm{148}
\newline {\color{Red} Ad Ben. Ant.} Ave Maria gratia plena dominus tecum, benedicta tu in mulieribus alleluya. {\color{Red} Ad Mag. Ant.} \hyperlink{o-antiphons}{O}.
\newline \ul{Lectiones feriales a quinquagesimo quarto capitulo usque ad finem ad libitum.}
\newline {\color{Red} Feria Secunda. Ad Benedictus Ant.} \hypertarget{ant-dicit-dominus}{Dicit} dominus, penitentiam agite, appropinquabit enim regnum celorum alleluya. {\color{Red} Ad Mag. Ant.} \hyperlink{o-antiphons}{O}.
\newline {\color{Red} Feria Tertia. Ad Ben. Ant.} Consurge, consurge, induere fortitudinem brachium domini. {\color{Red} Ad Mag. Ant.} \hyperlink{o-antiphons}{O}.
\newline {\color{Red} Feria Quarta. Ad Ben. Ant.} Ponent domino gloriam, et laudem eius in insulis nunciabunt quia ecce veniet et non tardabit. {\color{Red} Ad Mag. Ant.} \hyperlink{o-antiphons}{O}.
\newline {\color{Red} Feria Quinta. Ad Ben. Ant.} Consolamini consolamini popule meus dicit deus vester. {\color{Red} Ad Mag. Ant.} \hyperlink{o-antiphons}{O}.
\newline {\color{Red} Feria Sexta. Ad Ben. Ant.} Dies domini sicut fur ita in nocte veniet, et vos estote parati quia qua hora non putatis filius hominis veniet. {\color{Red} Ad Mag. Ant.} \hyperlink{o-antiphons}{O}.
{\ubsubsection{Antiphone Dicende In Laudibus Per Sex Ferias Ante Vigiliam Natalis Domini Secundum Ordinem Feriarum}{breviarium-antiphone-dicende-in-laudibus-per-sex-ferias}}
\par \noindent {\color{Red} Feria Secunda. In Laudibus. Ant.} Ecce veniet dominus princeps regum terre beati qui parati sunt occurrere illi. {\color{Red} Ps.} \psalm{50} {\color{Red} Ant.} Dum venerit filius hominis putas inveniet fidem super terram. {\color{Red} Ps.} \psalm{5} {\color{Red} Ant.} Ecce iam venit plenitudo temporis, in quo misit deus filium suum in terris. {\color{Red} Ps.} \psalm{62} {\color{Red} Ant.} Haurietis aquas in gaudio de fontibus salvatoris. {\color{Red} Can.} \psalm{isaias} {\color{Red} Ant.} Egredietur dominus de loco sancto suo, veniet ut salvet populum suum. {\color{Red} Ps.} \psalm{148}
\newline {\color{Red} Feria Tercia. In Laudibus. Ant.} Rorate celi desuper et nubes pluant iustum, aperiatur terra et germinet salvatorem. {\color{Red} Ps.} \psalm{50} {\color{Red} Ant.} Emitte agnum domine dominatorem terre de petra deserti ad montem filie Syon. {\color{Red} Ps.} \psalm{42} {\color{Red} Ant.} Ut cognoscamus domine in terra viam tuam, in omnibus gentibus salutare tuum. {\color{Red} Ps.} \psalm{62} {\color{Red} Ant.} Da mercedem domine sustinentibus te, ut prophete tui fideles inveniantur. {\color{Red} Can.} \psalm{ezechias} {\color{Red} Ant.} Lex per Moysen data est, gratia et veritas per Ihesum Christum facta est. {\color{Red} Ps.} \psalm{148}
\newline {\color{Red} Feria Quarta. In Laudibus. Ant.} Prophete predicaverunt nasci salvatorem de virgine Maria. {\color{Red} Ps.} \psalm{50} {\color{Red} Ant.} Spiritus domini super me evangelizare pauperibus misit me. {\color{Red} Ps.} \psalm{64} {\color{Red} Ant.} Annunciate populis et dicite, ecce deus salvator noster veniet. {\color{Red} Ps.} \psalm{62} {\color{Red} Ant.} Ecce veniet dominus ut sedeat cum principibus, et solium glorie teneat. {\color{Red} Can.} \psalm{annas} {\color{Red} Ant.} Propter Syon non tacebo donec egrediatur ut splendor iustus eius. {\color{Red} Ps.} \psalm{148}
\newline {\color{Red} Feria Quinta. In Laudibus. Ant.} De Syon veniet dominus omnipotens, ut salvum faciat populum suum. {\color{Red} Ps.} \psalm{50} {\color{Red} Ant.} Convertere domine aliquantulum, et ne tardes venire ad servos tuos. {\color{Red} Ps.} \psalm{89} {\color{Red} Ant.} De Syon veniet qui regnaturus est, dominus Emmanuel magnum nomen eius. {\color{Red} Ps.} \psalm{62} {\color{Red} Ant.} Ecce deus meus et honorabo eum, deus patris mei et exaltabo eum. {\color{Red} Can.} \psalm{moysis} {\color{Red} Ant.} Dominus legifer noster dominus rex noster, ipse veniet et salvabit nos. {\color{Red} Ps.} \psalm{148}
\newline {\color{Red} Feria Sexta. In Laudibus. Ant.} Constantes estote videbitis auxilium domini super vos. {\color{Red} Ps.} \psalm{50} {\color{Red} Ant.} Ad te domine levavi animam meam, veni et eripe me domine ad te confugi. {\color{Red} Ps.} \psalm{142} {\color{Red} Ant.} Veni domine et noli tardare, relaxa facinora plebis tue Israhel. {\color{Red} Ps.} \psalm{62} {\color{Red} Ant.} Deus a Libano veniet, et splendor eius sicut lumen erit. {\color{Red} Can.} \psalm{habacuc} {\color{Red} Ant.} Ego autem ad dominum aspiciam et expectabo deum salvatorem meum. {\color{Red} Ps.} \psalm{148}
\newline {\color{Red} Sabbato. In Laudibus. Ant.} Intuemini quantus sit gloriosus iste, qui ingreditur ad salvandos populos. {\color{Red} Ps.} \psalm{50} {\color{Red} Ant.} Consurge consurge induere fortitudinem brachium domini. {\color{Red} Ps.} \psalm{91} {\color{Red} Ant.} Elevare elevare consurge Hierusalem solve vincla colli tui captiva filia Syon. {\color{Red} Ps.} \psalm{62} {\color{Red} Ant.} Expectetur sicut pluvia eloquium domini, et descendet sicut ros super nos deus noster. {\color{Red} Can.} \psalm{audite} {\color{Red} Ant.} Paratus esto Israhel in occursum domini, quoniam venit. {\color{Red} Ps.} \psalm{148}
{\ubsubsection{O Antiphone}{o-antiphons}}
\par \noindent \ul{Sequentes antiphone sunt dicende septem proximas diebus ante Vigiliam Natalis Domini. Ad Mag. In diebus vero in quibus cantantur nec preces nec prostrationes sunt in Vesperis faciende, sed statim post antiphonam subiungatur Oratio.} {\color{Red} Ant.}
\lettrine[lines=2]{\zallmancaps \color{Red} O}{}\ Sapientia que ex ore altissimi prodisti, attingens a fine usque ad finem fortiter, suaviterque disponens omnia, veni ad docendum nos, viam prudentie. {\color{Red} Ant.}
\lettrine[lines=2]{\zallmancaps \color{Blue} O}{}\ Adonay et dux domus Israhel qui Moysy in igne flamme rubi apparuisti et ei in Sina legem dedisti, veni ad redimendum nos in brachio extento. {\color{Red} Ant.}
\lettrine[lines=2]{\zallmancaps \color{Red} O}{}\ Radix Iesse qui stans in signum populorum super quem continebunt reges os suum, quem gentes deprecabuntur, veni ad liberandum nos, iam noli tardare. {\color{Red} Ant.}
\lettrine[lines=2]{\zallmancaps \color{Blue} O}{}\ Clavis David et sceptrum domus Israhel, qui aperis et nemo claudit, claudis et nemo aperit, veni et educ vinctum de domo carceris sedentem in tenebris et umbra mortis. {\color{Red} Ant.}
\lettrine[lines=2]{\zallmancaps \color{Red} O}{}\ Oriens splendor lucis eterne et sol iusticie, veni et illumina sedentes in tenebris et umbra mortis. {\color{Red} Ant.}
\lettrine[lines=2]{\zallmancaps \color{Blue} O}{}\ Rex gentium et desideratus earum lapisque angularis, qui facis utraque unum, veni salva hominem quem de limo formasti. {\color{Red} Ant.}
\lettrine[lines=2]{\zallmancaps \color{Red} O}{}\ Emmanuel rex et legifer noster, expectatio gentium, et salvator earum, veni ad salvandum nos domine deus noster.
\newline {\color{Red} Memoria de Adventu in die Sancti Thome. In Laudibus. Ant.} Nolite timere, quinta die veniet ad vos dominus noster.
\newline \ul{Quando autem festum transfertur in secundam feriam predicta Ant., Oratio dimittatur, et Antiphona ferialis scilicet,} \hyperlink{ant-dicit-dominus}{Dicit dominus}. \ul{ad memoriam dicatur.}
{\ubsubsection{In Vigilia Natalis Domini}{breviarium-in-vigilia-natalis-domini}}
\par \noindent {\color{Red} Ad Matutinas. Prostrationes non fiant. Invitat.}
\begin{invitatory}[hodie-scietis]
{Hodie scietis quia veniet dominus, * Et mane videbitis gloriam eius.}
\end{invitatory}
{\color{Red} Ymnus.} \hyperlink{verbum}{Verbum supernum}. Antiphone et psalmi ferialis. \V Hodie scietis quia veniet dominus. \R Et mane videbitis gloriam eius.
\newline {\color{Red} Lectio Prima vel Septima si Dominica fuerit. Secundum Matheum.}
\lettrine[lines=2]{\zallmancaps \color{Blue} I}{n} illo tempore: Cum esset desponsata mater eius Maria Ioseph: ante quam convenirent inventa est in utero habens de Spiritu Sancto. Et reliqua.
\newline {\color{Red} Omelia Origenis.}
\lettrine[lines=2]{\zallmancaps \color{Red} Q}{ue} fuit necessitas ut desponsata esset Maria Ioseph, nisi propterea quatenus hoc sacramentum diabolo celaretur, et ille malignus fraudis commenta adversus desponsatam virginem nulla penitus invenisset?
\begin{responsory}[sanctificamini-hodie]
{Sanctificamini hodie et estote parati quia die crastina videbitis * Maiestatem dei in vobis.}{Hodie scietis quia veniet dominus et mane videbitis.}
\end{responsory}%
\newline {\color{Red} Lectio Secunda vel Octava.}
\lettrine[lines=2]{\zallmancaps \color{Blue} V}{el} ideo desponsata fuit Ioseph, ut nato infanti vel ipsi Marie curam videretur gerere Ioseph, sive in Egyptum iens, vel inde denuo veniens. Ideo desponsata fuit Ioseph, non tamen in concupiscentia iuncta. Mater inquit eius. Mater immaculata, mater incorrupta, mater intacta.
\begin{responsory}[constantes-estote]
{Constantes estote, videbitis auxilium domini super vos, Iudea et Hierusalem nolite timere, * Cras egrediemini et dominus erit vobiscum.}{Sanctificamini filii Israhel et estote parati.}
\end{responsory}%
\newline {\color{Red} Lectio Tertia vel Nona.}
\lettrine[lines=2]{\zallmancaps \color{Red} M}{ater} eius. Cuius eius? Mater dei unigeniti, domini et regis omnium, plasmatoris et creatoris cunctorum. Illius qui in excelsis sine matre, et in terra est sine patre. Illius qui in celis secundum deitatem in sinu est patris, et in terris secundum corporis susceptionem in sinu est matris.
\begin{responsory-doxology}[sanctificamini-filii]
{Sanctificamini filii Israhel dicit dominus, in die enim crastina descendet dominus, * Et auferet a vobis omnem langorem.}{Ex Syon species decoris eius, deus noster manifeste veniet.}
\end{responsory-doxology}%
\newline \ul{Hac die in Laudibus et deinceps usque ad versus exclusive fiat officium, sicut in festo duplici.}
\newline \V Crastina die delebitur iniquitas terre. \R Et regnabit super vos salvator mundi.
\newline {\color{Red} In Laudibus. Ant.} Iudea et Hierusalem nolite timere, cras egrediemini et dominus erit vobiscum. {\color{Red} Ps.} \psalm{92} {\color{Red} Ant.} Crastina die erit vobis salus, dicit dominus deus exercituum. {\color{Red} Ps.} \psalm{99} {\color{Red} Ant.} Hodie scietis quia veniet dominus, et mane videbitis gloriam eius. {\color{Red} Ps.} \psalm{62} {\color{Red} Ant.} Propter Syon non tacebo, donec egrediatur ut splendor iustus eius. {\color{Red} Can.} \psalm{benedicite} {\color{Red} Ant.} Crastina die delebitur iniquitas terre, et regnabit super vos salvator mundi. {\color{Red} Ps.} \psalm{148} {\color{Red} Capitulum.}
\lettrine[lines=2]{\zallmancaps \color{Blue} P}{ropter} Syon non tacebo et propter Hierusalem non quiescam, donec egrediatur ut splendor iustus eius, et salvator eius ut lampas accendatur. {\color{Red} Ymnus.} \hyperlink{vox-clara}{Vox clara}. \V Constantes estote. \R Videbitis auxilium domini super vos. {\color{Red} Ad Ben. Ant.} Cum esset desponsata mater Ihesu Maria Ioseph, antequam convenirent inventa est in utero habens, quod enim in ea natum est de Spiritu Sancto est, alleluya. {\color{Red} Oratio.}
\lettrine[lines=2]{\zallmancaps \color{Red} D}{eus} \hypertarget{oratio-nativitatis-domini}{qui} nos redemptionis nostre annua expectatione letificas, presta ut unigenitum tuum, quem redemptorem leti suscipimus, venientem quoque iudicem securi videamus. Dominum nostrum.
\newline \ul{Hec Oratio dicatur ad Terciam, Sextam, Nonam et Vesperas. Ad horas diei, antiphone Laudum.}
\newline {\color{Red} Ad Terciam Cap.} Propter Syon.
\begin{responsory-breve}[hodie-scietis-breve]
{Hodie scietis * Quia veniet dominus.}{Et mane videbitis gloriam eius.}
\end{responsory-breve}%
\V Crastina die delebitur.
\newline {\color{Red} Ad Sextam Capitulum.}
\lettrine[lines=2]{\zallmancaps \color{Blue} V}{idebunt} gentes iustum tuum, et cuncti reges inclitum tuum, et vocabitur tibi nomen novum, quod os domini nominavit.
\begin{responsory-breve}[crastina-die-delebitur]
{Crastina die, * Delebitur iniquitas terre.}{Et regnabit super vos salvator mundi.}
\end{responsory-breve}%
\V Crastina die erit.
\newline {\color{Red} Ad Nonam Capitulum.}
\lettrine[lines=2]{\zallmancaps \color{Red} N}{on} vocaberis ultra derelicta, et terra tua non vocabitur amplius desolata, sed vocaberis voluntas mea in ea, et terra tua inhabitabitur.
\begin{responsory-breve}[crastina-die-erit]
{Crastina die, * Erit vobis salus.}{Dicit dominus exercituum.}
\end{responsory-breve}%
\V Constantes estote.
\newline \ul{Si hec Vigilia in Dominica evenerit hoc modo fiat officium. Ad Matutinas. Invitat.} \hyperlink{hodie-scietis}{Hodie scietis} \ul{Ymnus ut supra. Super psalmos primi nocturni Ant. Scientes. Ps.} \psalm{1} \ul{et cet. Cetere ad ceteros ut supra. In primo nocturno. \Vbar .} Ex Syon. \ul{In secundo nocturno. \Vbar .} Egredietur virga. \ul{Sex lectiones et sex responsoria de Dominica. In tertio nocturno. \Vbar .} Hodie scietis. \ul{Tres ultime lectiones de expositione evangelii.} Cum esset desponsata, \ul{et tria ultima responsoria de Vigilia. Expositio autem Dominicalis evangelii cum tribus ultimis responsoriis de hystoria.} \hyperlink{canite-tuba}{Canite,} \ul{et cum Laudibus: illo anno omittantur. \Vbar .} Crastina. \ul{In Laudibus. Ant. Cap. \Vbar . et Ant., Ad Ben. et Oratio de Vigilia erunt. Post orationem de Vigilia memoria fiat de Dominica. Antiphone, Cap., responsoria, versiculi, et Oracio ad horas de Vigilia erunt.}
\newline {\color{Red} Ad Vesperas. Ant.} Rex pacificus magnificatus est, cuius vultum desiderat universa terra. {\color{Red} Ps.} \psalm{112} {\color{Red} Ant.} Magnificatus est rex pacificus super omnes reges universe terre. {\color{Red} Ps.} \psalm{116} {\color{Red} Ant.} Scitote quia prope est regnum dei, amen dico vobis quia non tardabit. {\color{Red} Ps.} \psalm{145} {\color{Red} Ant.} Levate capita vestra, ecce appropinquabit redemptio vestra. {\color{Red} Ps.} \psalm{146} {\color{Red} Ant.} Completi sunt dies Marie ut pareret filium suum primogenitum. {\color{Red} Ps.} \psalm{147} {\color{Red} Capitulum.}
\lettrine[lines=2]{\zallmancaps \color{Blue} P}{opulus} qui ambulabat in tenebris vidit lucem magnam, habitantibus in regione umbre mortis lux orta est eis.
\begin{responsory-final}[iudea-et-hierusalem]
{Iudea et Hierusalem nolite timere, * Cras egrediemini, * Et dominus erit vobiscum.}{Constantes estote, videbitis auxilium domini super vos.}{Et dominus}
\end{responsory-final}%
{\color{Red} Ymnus.}
\hymn{veni-redemptor}{Veni redemptor gentium}{
\lettrine[lines=2]{\zallmancaps \color{Red} V}{eni} redemptor gentium,
ostende partum virginis,
miretur omne seculum,
talis decet partus deum.
\par {\bfseries \color{Blue} N}on ex virili semine
sed mistico spiramine,
verbum dei factum caro,
fructusque ventris floruit.
\par {\bfseries \color{Red} A}lvus tumescit virginis,
claustra pudoris permanent,
vexilla virtutum micant,
versatur in templo deus.
\par {\bfseries \color{Blue} P}rocedens de thalamo suo
pudoris aula regia,
gemine gigas substantie,
alacris ut currat viam.
\par {\bfseries \color{Red} E}gressus eius a patre,
regressus eius ad patrem,
excursus usque ad inferos,
recursus ad sedem dei.
\par {\bfseries \color{Blue} E}qualis eterno patri
carnis tropheo accingere,
infirma nostri corporis
virtute firmans perpeti.
\par {\bfseries \color{Red} P}resepe iam fulget tuum
lumenque nox spirat novum,
quod nulla nox interpolet
fideque iugi luceat.
\par {\bfseries \color{Blue} G}loria tibi domine
qui natus es de virgine,
cum patre et sancto spiritu
in sempiterna secula. Amen.
}
\newline \ul{Iste versus dicatur ad omnes ymnos eiusdem in metri tam de tempore quam de sanctis usque ad Primas Vesperas Epiphanie exclusive.}
\newline \V Tanquam sponsus. \R Dominus procedens de thalamo suo.
\newline {\color{Red} Ad Mag. Ant.} Dum ortus fuerit sol de celo videbitis regem regum procedentem a patre tanquam sponsum de thalamo suo. {\color{Red} Oratio.} \hyperlink{oratio-nativitatis-domini}{Deus qui nos redemp.}
\newline {\color{Red} Ad Completorium. Super psalmos. Ant.} Completi sunt. {\color{Red} Ad Nunc Dimittis. Ant.} Ecce completa sunt omnia que dicta sunt per angelum de virgine Maria.
{\ubsubsection{In Natali Domini Ad Matutinas}{breviarium-in-natali-domini-ad-matutinas}}
\par \noindent {\color{Red} Invitat.}
\begin{invitatory}[christus-natus-invitatorium]
{Christus natus est nobis, * Venite adoremus.}
\end{invitatory}
{\color{Red} Ymnus.}
\hymn{christe-redemptor}{Christe redemptor omnium}{
\lettrine[lines=2]{\zallmancaps \color{Blue} C}{hriste} redemptor omnium,
ex patre patris unice,
solus ante principium
natus ineffabiliter.
\par {\bfseries \color{Red} T}u lumen, tu splendor patris,
tu spes perennis omnium,
intende quas fundunt preces
tui per orbem famuli.
\par {\bfseries \color{Blue} M}emento salutis actor,
quod nostri quondam corporis
ex illibata virgine
nascendo formam sumpseris.
\par {\bfseries \color{Red} S}ic presens testatur dies
currens per anni circulum,
quod solus a sede patris
mundi salus adveneris.
\par {\bfseries \color{Blue} H}unc celum, terra, hunc mare,
hunc omne quod in eis est,
actorem adventus tui,
laudat exultans cantico.
\par {\bfseries \color{Red} N}os quoque qui sancto tuo
redempti sanguine sumus,
ob diem natalis tui
ymnum novum concinimus.
\par {\bfseries \color{Blue} G}loria tibi domine
qui natus es de virgine
cum patre et sancto spiritu,
in sempiterna secula. Amen.
}
\newline {\color{Red} In Primo Nocturno. Ant.} Dominus dixit ad me, filius meus es tu, ego hodie genui te. {\color{Red} Ps.} \psalm{2} {\color{Red} Ant.} Tanquam sponsus dominus procedens de thalamo suo. {\color{Red} Ps.} \psalm{18} {\color{Red} Ant.} Diffusa est gratia in labiis tuis propterea benedixit te deus in eternum. {\color{Red} Ps.} \psalm{44} \V Verbum caro factum est. \R Et habitavit in nobis.
\newline {\color{Red} Lectio Prima.}
\lettrine[lines=2]{\zallmancaps \color{Red} P}{rimo} tempore alleviata est terra Zabulon et terra Nephtali, et novissimo aggravata est via maris trans Iordanem Galilee gentium. Populus qui ambulabat in tenebris: vidit lucem magnam. Habitantibus in regione umbre mortis: lux orta est eis. Multiplicasti gentem, et non magnificasti leticiam. Letabuntur coram te sicut qui letantur in messe: sicut exultant victores capta preda, quando dividunt spolia. Iugum enim oneris eius, et virgam humeri eius, et sceptrum exactoris eius superasti: sicut in die Madian. Quia omnis violentia predatio cum tumultu, et vestimentum mixtum sanguine erit in combustionem, et cibus ignis. Parvulus enim natus est nobis: et filius datus est nobis. Hec dicit.
\begin{responsory}[hodie-nobis-celorum]
{Hodie nobis celorum rex de virgine nasci dignatus est, ut hominem perditum ad celestia regna revocaret, gaudet exercitus angelorum, * Quia salus eterna humano generi apparuit.}{Gloria in excelsis deo, et in terra pax hominibus bone voluntatis.}
\end{responsory}%
\newline {\color{Red} Lectio Secunda.}
\lettrine[lines=2]{\zallmancaps \color{Blue} C}{onsolamini} consolamini, popule meus, dicit Deus vester. Loquimini ad cor Hierusalem et advocate eam, quoniam completa est malicia eius: dimissa est iniquitas illius. Suscepit de manu domini duplicia: pro omnibus peccatis suis. Vox clamantis in deserto. Parate viam domini, rectas facite in solitudine semitas Dei nostri. Omnis vallis exaltabitur, et omnis mons et collis humiliabitur: et erunt prava in directa, et aspera in vias planas. Et revelabitur gloria domini: et videbit omnis caro pariter quod os Domini locutum est. Hec dicit.
\begin{responsory}[hodie-nobis-de-celo]
{Hodie nobis de celo pax vera descendit, * Hodie per totum mundum melliflui facti sunt celi.}{Hodie illuxit dies redemptionis nove, reparationis antique felicitatis, eterne.}
\end{responsory}%
\newline {\color{Red} Lectio Tertia.}
\lettrine[lines=2]{\zallmancaps \color{Red} C}{onsurge} consurge, induere fortitudini tua Syon. induere vestimentis glorie tue Hierusalem civitas sancti, quia non adiciet ultra ut pertranseat per te incircuncisus et immundus. Excutere de pulvere, consurge, sede, Hierusalem. Solve vincla colli tui captiva filia Syon. Quia hec dicit dominus. Gratis venundati estis, et sine argento redimemini. Quia hec dicit dominus Deus. In Egyptum descendit populus meus in principio, ut colonus esset ibi: et Assur absque ulla causa calumniatus est
%pp091
eum. Et nunc quid michi est hic dicit dominus? Quoniam ablatus est populus meus gratis. Dominatores eius inique agunt dicit dominus: et iugiter tota die nomen meum blasphematur. Propter hoc sciet populus meus nomen meum in die illa: quia ego ipse qui loquebar ecce adsum. Hec dicit.
\begin{responsory-doxology}[descendit-de-celis]
{Descendit de celis, deus verus a patre genitus, introivit in uterum virginis, nobis ut appareret visibilis, indutus carne humana prothoparente edita, et exivit per clausam portam deus et homo, * Lux et vita conditor mundi.}{Tanquam sponsus dominus procedens de thalamo suo.}
\end{responsory-doxology}%
\newline {\color{Red} In Secundo Nocturno. Ant.} Suscepimus deus misericordiam tuam in medio templi tui. {\color{Red} Ps.} \psalm{47} {\color{Red} Ant.} Orietur diebus domini abundantia pacis, et dominabitur. {\color{Red} Ps.} \psalm{71} {\color{Red} Ant.} Veritas de terra orta est, et iusticia de celo prospexit. {\color{Red} Ps.} \psalm{84} \V Notum fecit dominus. \R Salutare suum.
\newline \ul{Ysidorus in Sermone de Natali Domini.} Natalis Domini dies. \grecross \
{\color{Red} Lectio Quarta.}
\lettrine[lines=2]{\zallmancaps \color{Blue} P}{ostquam} invidia diaboli parens primus cecidit, confestim ille exul et perditus in omni genere suo radicem malicie transduxit, crescebatque in malum vehementius omne genus mortalium, diffusis ubique sceleribus. Volens ergo Deus terminare peccatum, consuluit, verbo, lege, prophetis, signis, plagis, prodigiis. Sed cum nec sic quidem errores suos admonitus agnosceret mundus, misit Deus filium suum ut carne indueretur, et hominibus appareret et peccatores sanaret. Qui ideo in hominem venit: quia per seipsum ab hominibus cognosci non potuit. Assumpsit autem humanitatem, non amisit divinitatem, manens quod erat, suscipiens quod non erat, ut liberaret quod fecerat. Hec est diei huius gloriosa festivitas.
\begin{responsory}[quem-vidistis]
{Quem vidistis pastores dicite, annunciate nobis in terris quis apparuit, * Natum vidimus, et choros angelorum collaudantes dominum.}{Dicite quidnam vidistis et annunciate Christi nativitatem.}
\end{responsory}%
\newline \ul{Augustinus in sermone de Natali Domini,} legimus sanctum Moysen.
{\color{Red} Lectio Quinta.}
% https://www.jstor.org/stable/pdf/42998637.pdf
\lettrine[lines=2]{\zallmancaps \color{Red} Q}{uid} autem Sibilla vaticinando de Christo clamaverit, in medium proferamus: ut ex uno lapide utrorumque frontes percutiantur, Iudeorum scilicet atque paganorum. Audite quid dixerit.
Iudicii signum, tellus sudore madescet.
E celo rex adveniet: per secla futurus.
Scilicet in carnem presens ut iudicet orbem.
Unde Deum cernent incredulus atque fidelis, celsum cum sanctis.
Evi iam termino in ipso, sic anime cum carne aderunt, quas iudicat ipse: cum iacet incultus densis in vepribus orbis.
Reicient simulacra viri, cunctam quoque gazam.
Exuret terras ignis: pontumque polumque.
Inquirens tetri portas effringet Averni.
Sanctorum sed enim cuncte lux libera carni tradetur: sontes eterna flamma cremabit.
Occultos actus retegens tunc quisque loquetur.
Secreta atque Deus reserabit pectora luci.
\begin{responsory}[o-magnum-mysterium]
{O magnum mysterium et admirabile sacramentum ut animalia viderent dominum natum, * Iacentem in presepio, beata virgo cuius viscera meruerunt, portare dominum Christum.}{Domine audivi auditum tuum et timui, consideravi opera tua et expavi in medio duum animalium.}% duum: duorum
\end{responsory}%
\newline \ul{Post hec iterum subiungit dicens.}
{\color{Red} Lectio Sexta.}
\lettrine[lines=2]{\zallmancaps \color{Blue} T}{unc} erit et luctus, stridebunt dentibus omnes.
Eripitur solis iubar et chorus interit astris.
Solvetur celum: lunaris splendor obibit.
Deiciet colles, valles extollet ab imo.
Non erit in rebus hominum sublime vel altum.
Iam equantur campis montes et cerula ponti.
Omnia cessabunt tellus confracta peribit.
Sic pariter fontes torrentur, fluminaque igni.
Et tuba tum sonitum tristem demittet ab alto.
Orbe gemens facinus miserum, variosque labores Tartareumque chaos monstrabit terra dehiscens.
Et coram hic domino reges sistentur ad unum.
Decidet e celo, ignisque et sulphuris amnis. %typo annis
Hec de Christi nativitate, passione, et resurrectione, atque secundo eius adventu ita dicta sunt ut siquis in Greco capita horum versuum discernere voluerit, inveniet Ihesus Christus Yos Theu Sother, quod et in Latinum translatis eisdem versibus apparet, preter quod Grecarum litterarum proprietas non adeo potuit observari.
\begin{responsory-doxology}[sancta-et-immaculata-doxology]
{Sancta et immaculata virginitas quibus te laudibus efferam nescio, * Quia quem celi capere non poterant tuo gremio contulisti.}{Benedicta tu in mulieribus et benedictus fructus ventris tui.}
\end{responsory-doxology}%
\newline {\color{Red} In Tertio Nocturno. Ant.} Ipse invocabit me alleluya, pater meus es tu alleluya. {\color{Red} Ps.} \psalm{88} {\color{Red} Ant.} Letentur celi et exultet terra ante faciem domini quoniam venit. {\color{Red} Ps.} \psalm{95} {\color{Red} Ant.} Notum fecit dominus alleluya, salutare suum, alleluya. {\color{Red} Ps.} \psalm{97} \V Ipse invocabit me. \R Pater meus es tu.
\newline {\color{Red} Lectio VII. Secundum Lucam.}
\lettrine[lines=2]{\zallmancaps \color{Blue} I}{n} illo tempore: Exiit edictum a Cesare Augusto ut describeretur universus orbis. Et reliqua.
\newline {\color{Red} Omelia Beati Gregorii Pape.}
\lettrine[lines=2]{\zallmancaps \color{Red} Q}{uia} largiente domino missarum sollemnia ter hodie celebraturi sumus: loqui diu de Evangelica lectione non possumus. Sed nos aliquid vel breviter dicere redemptoris nostri nativitas ipsa compellit. Quid est quod nascituro domino mundus describitur: nisi noc quod aperte monstratur quia ille veniebat in carne, qui electos suos ascriberet in eternitate? Quo contra de reprobis per prophetam dicitur. Deleantur de libro viventium, et cum iustis non scribantur. Qui bene eciam in Bethleem nascitur. Bethleem quippe: domus panis interpretatur. Ipse namque est qui ait. Ego sum panis vivus qui de celo descendi.
\begin{responsory}[beata-dei-genitrix]
{Beata dei genitrix Maria cuius viscera intacta permanent, * Hodie genuit salvatorem seculi.}{Beata que credidit quoniam perfecta sunt omnia que dicta sunt ei a domino.}
\end{responsory}%
\newline {\color{Red} Lectio VIII. Secundum Lucam.}
\lettrine[lines=2]{\zallmancaps \color{Blue} I}{n} illo tempore: Pastores loquebantur ad invicem. Tranesamus usque Bethleem, et videamus hoc verbum quod factum est, quod fecit dominus, et ostendit nobis. Et reliqua.
\newline {\color{Red} Omelia Venerabilis Bede Presbiteri.}
\lettrine[lines=2]{\zallmancaps \color{Red} N}{ato} in Bethleem domino salvatore, sicut sacra evangelii testatur hystoria, pastoribus qui in regione eadem erant vigilantes et custodientes vigilias noctis super gregem suum, angelus domini magna cum luce apparuit, exortumque mundo solem iusticie, non solum celestis voce sermonis, verum eciam claritate divine lucis astruxit. Nusquam enim in tota veteris instrumenti serie reperimus angelos qui tam sedulo apparuere patribus, cum luce apparuisse. Sed hoc privilegium recte hodierno tempori servatum est, quando exortum est in tenebris lumen rectis corde, misericors et miserator dominus.
\begin{responsory}[beata-viscera]
{Beata viscera Marie virginis que portaverunt eterni patris filium, et beata ubera que lactaverunt Christum dominum, * Quia hodie pro salute mundi de virgine nasci dignatus est.}{Dies sanctificatus illuxit nobis venite gentes et adorate dominum.}
\end{responsory}%
\newline {\color{Red} Lectio IX. Secundum Iohannem.}
\lettrine[lines=2]{\zallmancaps \color{Blue} I}{n} principio erat verbum. Et verbum erat apud deum, et deus erat verbum. Et reliqua.
\newline {\color{Red} Omelia Venerabilis Bede Presbiteri.}
\lettrine[lines=2]{\zallmancaps \color{Red} Q}{uia} temporalem mediatoris Dei et hominum hominis Ihesu Christi nativitatem que hodierna die facta est sanctorum verbis evangelistarum, Mathei videlicet et Luce manifeste descriptam esse cognovimus: libet eciam de verbi id est divinitatis eius eternitate in qua patri manet semper equalis beati Iohannis evangeliste dicta scrutari, qui singularis privilegio meruit castitatis, ut ceteris altius divinitatis ipsius caperet, simul et patefaceret archanum. Alii igitur evangeliste Christum ex tempore natum describunt: Iohannes in principio eundem fuisse testatur, dicens. In principio erat verbum.
\begin{responsory-doxology}[verbum-caro-factum]
{Verbum caro factum est et habitavit in nobis cuius gloriam vidimus quasi unigeniti a patre, * Plenum gratie et veritatis.}{In principio erat verbum, et verbum erat apud deum, et deus erat verbum.}
\end{responsory-doxology}%
\newline \ul{Postea sequitur.} Dominus vobiscum. Et cum spiritu tuo.
\newline Initium sancti Evangelium Secundum Matheum. Gloria tibi domine.
\lettrine[lines=2]{\zallmancaps \color{Blue} L}{iber} generationis Ihesu Christi filii David, filii Abraham. Abraham genuit Ysaac. Ysaac autem genuit Iacob. Iacob autem genuit Iudam, et fratres eius. Iudas autem genuit Phares et Zaram de Thamar. Phares autem genuit Esron. Esron autem genuit Aram. Aram autem genuit Aminadab. Aminadab autem genuit Naason. Naason autem genuit Salmon. Salmon autem genuit Booz de Raab. Booz autem genuit Obeth ex Ruth. Obeth autem genuit Iesse. Iesse autem genuit David regem. David autem rex genuit Salomonem ex ea que fuit Urie. Salomon autem genuit Roboam. Roboam autem genuit Abiam. Abia autem genuit Asa. Asa autem genuit Iosophath. Iosophath autem genuit Ioram. Ioram autem genuit Oziam. Ozias autem genuit Ioathan. Ioathan autem genuit Achaz. Achaz autem genuit Ezechiam. Ezechias autem genuit Manassen. Manasses autem genuit Amon. Amon autem genuit Iosiam. Iosias autem genuit Iechoniam et fratres eius, in transmigratione Babilonis. Et post transmigrationem Babilonis, Iechonias genuit Salathiel. Salathiel autem genuit Zorobabel. Zorobabel autem genuit Abiud. Abiud autem genuit Elyachim. Elyachim autem genuit Azor. Azor autem genuit Sadoch. Sadoch autem genuit Achim. Achim autem genuit Eliud. Eliud autem genuit Eleazar. Eleazar autem genuit Mathan. Mathan autem genuit Iacob. Iacob autem genuit Ioseph virum Marie. De qua natus est Ihesus: qui vocatur Christus.
\psalm{tedeum} \ul{Finito} \psalm{tedeum}, \ul{Si missa immediate non sequatur dicatur,} \V Puer natus est nobis. \R Et filius datus est nobis. Deus in adiutorium, \ul{et dicantur Laudibus more solito.}
\newline {\color{Red} In Laudibus. Ant.} Quem vidistis pastores dicite, annunciate nobis in terris quis apparuit: natum vidimus, et choros angelorum collaudantes dominum alleluya alleluya. {\color{Red} Ps.} \psalm{92} {\color{Red} Ant.} Genuit puerpera regem cui nomen eternum et gaudium matris habes cum virginitatis honore nec primam similem visa est nec habere sequentem alleluya. {\color{Red} Ps.} \psalm{99} {\color{Red} Ant.} Angelus ad pastores ait: annuncio vobis gaudium magnum quia natus est vobis hodie salvator mundi alleluya. {\color{Red} Ps.} \psalm{62} {\color{Red} Ant.} Facta est cum angelo multitudo celestis exercitus, laudantium et dicentium gloria in excelsis deo, et in terra pax hominibus bone voluntatis alleluya. {\color{Red} Can.} \psalm{benedicite} {\color{Red} Ant.} Parvulus filius hodie natus est nobis, et vocabitur deus fortis alleluya alleluya. {\color{Red} Ps.} \psalm{148}
\newline \ul{Non dicatur, Capitulum neque Ymnus, neque \Vbar , sed immediate post quintam antiphonam dicatur,} {\color{Red} Ad Ben. Ant.} Gloria in excelsis deo et in terra pax hominibus bone voluntatis alleluya alleluya. {\color{Red} Oracio.}
\lettrine[lines=2]{\zallmancaps \color{Blue} C}{oncede} \hypertarget{concede-nativitas}{quesumus} omnipotens deus, ut nos unigeniti tui nova per carnem nativitas liberet, quos sub peccati iugo vetusta servitus tenet. Per eundem.
\newline \ul{Deinde} Dominus vobiscum \ul{vel} Domine exaudi \ul{et} Benedicamus domino. \ul{Predicto modo dicantur Laudibus quando Missa non celebratur inter Matutinas et Laudibus, quando vero Missa fuerit celebranda statim post.} \psalm{tedeum} \ul{inchoetur, et dicta communione absque versiculo et absque} Deus in adiutorium, \ul{immediate inchoetur Ant. in Laudibus.} Quem vidistis, \ul{finita vero post} \psalm{benedictus} \ul{Ant., dicat sacerdos post communione, et subiuncto} Dominus vobiscum, \ul{dicatur,} Ite missa est, \ul{et sic Missa et Laudibus simul terminentur. Ad horas diei antiphone Laudum.}
\newline {\color{Red} Ad Primam.} \R Ihesu Christe fili dei vivi miserere nobis, * Alleluya alleluya. \V Qui natus es de virgine Maria, alleluya alleluya. \ul{Versus iste dicatur cotidie usque ad Epyphaniam exclusive.}
\newline {\color{Red} Ad Terciam. Capitulum.}
\lettrine[lines=2]{\zallmancaps \color{Red} M}{ultifarie} multisque modis olim deus loquens patribus in prophetis, novissime diebus istis locutus est nobis in filio quem constituit heredem universorum, per quem fecit et secula.
\begin{responsory-breve}[verbum-caro-factum-breve]
{Verbum caro factum est, * Alleluya alleluya.}{Et habitavit in nobis.}
\end{responsory-breve}%
\V Notum fecit dominus. \R Salutare suum. {\color{Red} Oracio.} \hyperlink{concede-nativitas}{Concede.}
\newline {\color{Red} Ad Sextam. Capitulum.}
\lettrine[lines=2]{\zallmancaps \color{Blue} A}{pparuit} benignitas, et humanitas salvatoris nostri dei, non ex operibus iusticie que fecimus nos, sed secundum suam misericordiam salvos nos fecit.
\begin{responsory-breve}[notum-fecit]
{Notum fecit dominus * Alleluya alleluya.}{Salutare suum.}
\end{responsory-breve}%
\V Ipse invocavit me. \R Pater meus es tu.
\newline {\color{Red} Ad Nonam. Capitulum.}
\lettrine[lines=2]{\zallmancaps \color{Red} A}{pparuit} gratia dei salvatoris nostri omnibus hominibus erudiens nos, ut abnegantes impietatem et secularia desideria sobrie, et iuste et pie vivamus in hoc seculo.
\begin{responsory-breve}[ipse-invocavit]
{Ipse invocavit me, * Alleluya alleluya.}{Pater meus es tu.}
\end{responsory-breve}%
\V Benedictus qui venit in nomine domini. \R Deus dominus et illuxit nobis.
\newline \ul{Ab hac die usque in crastinum Octavarum Epiphanie exclusive cotidie dicantur responsoria horarum cum alleluya, excepto festo Innocentum quando extra Dominicam evenerit. In Octava tamen eorum dicantur cum alleluya.}
\newline {\color{Red} Ad Vesperas. Ant.} Tecum principium in die virtutis tue, in splendoribus sanctorum ex utero ante luciferum genui te. {\color{Red} Ps.} \psalm{109} {\color{Red} Ant.} Redemptionem misit dominus populo suo, mandavit in eternum testamentum suum. {\color{Red} Ps.} \psalm{110} {\color{Red} Ant.} Exortum est in tenebris lumen rectis corde, misericors et miserator et iustus dominus. {\color{Red} Ps.} \psalm{111} {\color{Red} Ant.} Apud dominum misericordia et copiosa apud eum redemptio. {\color{Red} Ps.} \psalm{129} {\color{Red} Ant.} De fructu ventris tui ponam super sedem tuam. {\color{Red} Ps.} \psalm{131}
\newline \ul{Predicte antiphone cum suis psalmis dicantur cotidie ad Vesperas usque ad Octavas Epyphanie, et in ipsis Octavis.} {\color{Red} Cap.} Multifarie. \ul{Ymnus, \Vbar , ut supra in Primis Vesperis.} {\color{Red} Ad Mag. Ant.} Hodie Christus natus est, hodie salvator apparuit, hodie in terra canunt angeli, letantur archangeli, hodie exultant iusti dicentes: gloria in excelsis deo alleluya. {\color{Red} Oracio.} \hyperlink{concede-nativitas}{Concede.} \ul{Postea fiat memoria de Sancto Stephano.}
\newline {\color{Red} Ad Completorium. Super psalmos. Ant.} Natus est nobis hodie salvator qui est Christus dominus in civitate David.
\newline {\color{Red} Ad Nunc Dimittis. Super psalmos. Ant.} Alleluya verbum caro factum est alleluya, et habitavit in nobis alleluya alleluya.
\newline \ul{Predicte antiphone dicantur in Completorio usque ad Vigiliam Epyphanie exclusive. Subscripto modo fiat memoria de Natali. In Laudibus, et in Vesperas cotidie usque ad Primas Vesperas Circumcisionis exclusive.}
\newline {\color{Red} In Laudibus. Ant.} Nesciens mater virgo virum peperit sine dolore salvatorem seculorum ipsum regem angelorum sola virgo lactabat ubere de celo pleno. \V Benedictus qui venit. {\color{Red} Oracio.} \hyperlink{concede-nativitas}{Concede.} {\color{Red} In Vesperis. Ant.} Gaudeamus omnes fideles, salvator noster natus est in mundum, hodie processit proles magnifici germinis, et perseverat pudor virginitatis. \V Tanquam sponsus. {\color{Red} Oracio.} \hyperlink{concede-nativitas}{Concede.}
\newline \ul{Si festum Natalis Domini, vel sancti Stephani vel sancti Iohannis, vel sanctorum Innocentium, vel sancti Thome in Dominica evenerit illa die de Dominica nichil fiat, sed in crastino sancti Thome post memoria de sanctis. In Laudibus fiat memoria de Dominica per Ant.} Dum medium, \ul{et per \Vbar .} Notum fecit. \ul{et Orationem.} \hyperlink{omnipotens-nativitas}{Omnipotens sempiterne.} \ul{In Vesperis vero per Ant.} Erant Ioseph \ul{et \Vbar .} Verbum caro. \ul{et Orationem predictam. Proxima die post festum sancti Thome, Ad Matutinas. Invitat.} \hyperlink{christus-natus-invitatorium}{Christus natus.} \ul{Ymnus} \hyperlink{christe-redemptor}{Christe Redemptor.} \ul{Super psalmos Ant.} Dominus dixit. \ul{sola dicatur psalmi.} \psalm{2} \ul{et ceteri de Nativitate Domini, Versiculis et tria Responsoria de hystoria Nativitatis secundum feriam. Tres lectiones de expositione Dominicalis evangelii, que omittenda est quando festum sancti Silvestri in Dominica evenerit, et tunc in hac die, scilicet, post festum sancti Thome, fiant tres lectiones de aliquo sermone Natalis supra notato.}
\newline {\color{Red} Lectio Prima.}
\lettrine[lines=2]{\zallmancaps \color{Blue} I}{n} illo tempore: Erant Ioseph et Maria mater Ihesu mirantes super iis que dicebantur de illo. Et reliqua.
\newline {\color{Red} Omelia Origenis.}
% https://la.wikisource.org/wiki/Translatio_XXXIX_Homiliarum
\lettrine[lines=2]{\zallmancaps \color{Red} C}{ongregemus} in unum ea que in ortu Domini Ihesu dicta scriptaque sunt de eo: Et tunc scire poterimus singula queque digna esse miraculo. Quam ob rem mirabatur et pater, sic enim appellatus est Ioseph, quia nutritius fuit, mirabatur et mater super omnibus que dicebantur de eo.
\newline {\color{Red} Lectio Secunda.}
\lettrine[lines=2]{\zallmancaps \color{Blue} Q}{ue} nam ergo sunt que de parvulo Ihesu fama disperserat? Pastores erant in regione illa vigilantes et observantes vigilias noctis super gregem suum. Venit angelus sub ipsa hora nativitatis Ihesu, et ait ad eos. Annuntio vobis gaudium magnum. Ite, et invenietis infantem involutum pannis: et positum in presepio.
\newline {\color{Red} Lectio Tertia.}
\lettrine[lines=2]{\zallmancaps \color{Red} N}{ec} dum angelus verba finierat, et ecce multitudo celestis exercitus laudare cepit et benedicere Deum. Cum hoc pastores trepidi perspexissent, et angelus recessisset ab eis: dixerunt ad se invicem. Eamus in Bethleem, et videamus hoc factum quod Dominus ostendit nobis. Venerunt. et invenerunt parvulum.
\newline \psalm{tedeum} \ul{dicantur. \Vbar .} Puer natus. \ul{In Laudibus una Ant.} Quem vidistis. \ul{Psalmi.} \psalm{92} \ul{Cap.} Multifarie.
{\color{Red} Ymnus.}
\hymn{a-solis-ortus}{A solis ortus cardine}{
\lettrine[lines=2]{\zallmancaps \color{Blue} A}{}\ solis ortus cardine
ad usque terre limitem,
christum canamus principem,
natum Maria virgine.
\par {\bfseries \color{Red} B}eatus actor seculi
servile corpus induit,
ut carne carnem liberans
non perderet quos condidit.
\par {\bfseries \color{Blue} C}aste parentis viscera
celestis intrat gratia,
venter puelle baiulat
secreta que non noverat.
\par {\bfseries \color{Red} D}omus pudici pectoris,
templum repente fit dei,
intacta nesciens virum
verbo concepit filium.
\par {\bfseries \color{Blue} E}nixa est puerpera
quem Gabriel predixerat,
quem matris alvo gestiens
clausus Iohannes senserat.
\par {\bfseries \color{Red} F}eno iacere pertulit,
presepe non abhorruit,
parvoque lacte pastus est,
per quem nec ales esurit.
\par {\bfseries \color{Blue} G}audet chorus celestium,
et angeli canunt deo,
palamque fit pastoribus
pastor creator omnium.
\par {\bfseries \color{Red} G}loria tibi domine
Qui natus es de virgine
Cum patre et sancto spiritu,
in sempiterna secula. Amen.
}
\newline \V Benedictus qui venit.
\newline \ul{Ad Ben. Ant.} Nesciens mater. \ul{Oracio.} \hyperlink{concede-nativitas}{Concede.} \ul{Memoria de sanctis, ultimo de Dominica, nisi hec dies in Sabbato evenerit. Tunc enim memoria Dominicalis in crastino reservabitur.}
\newline {\color{Red} Ad memoriam de Dominica. Ant.} Dum medium silentium tenerent omnia et nox in suo cursu iter perageret omnipotens sermo tuus domine a regalibus sedibus venit alleluya. \V Notum fecit.
%pp092
\lettrine[lines=2]{\zallmancaps \color{Red} O}{mnipotens} \hypertarget{omnipotens-nativitas}{sempiterne} deus, dirige actus nostros in beneplacito tuo, ut in nomine dilecti filii tui mereamur bonis operibus abundare. Per eundem.
\newline \ul{Ad Primam. Ant.} Quem vidistis. \ul{Cetere ad ceteras horas. Capitula, responsoria, versiculi, et oracio de die Natalis Domini.}
\newline {\color{Red} Ad memoriam de Dominica in Vesperis. Ant.} Erant Ioseph et Maria mater Ihesu mirantes super his que dicebantur de illo. \V Verbum caro. Oracio. \hyperlink{omnipotens-nativitas}{Omnipotens sempiterne.}
\newline \ul{Si vero in hac die Dominica evenerit ad Matutinas, Invitat, Ymnus, Antiphone, psalmi, versiculi, responsoria, de Nativitate Domini, sex lectiones fiant de tribus mediis lectiones natalis tres ultime de omelia.} Erant Ioseph. \ul{In Laudibus. Ant., psalmi, capitulum, ymnus, \Vbar , de Natali. Ad Ben. Ant.} Nesciens. \ul{Oracio.} \hyperlink{concede-nativitas}{Concede.} \ul{Memoria de sanctis. Ultimo de Dominica, hore diei de Natali. Ad Vesperas. Capitulum, et cetera que sequuntur de sanctis Silvestro erunt, et memoria fiet de natali, deinde de sanctis ultimo de Dominica.}
{\ubsubsection{In Circumcisione Domini}{breviarium-circumcisione-domini-vesperas}}
\par \noindent {\color{Red} Ad Vesperas. Capitulum.} Multifarie. \R \hyperlink{verbum-caro-factum-doxology}{Verbum caro.} {\color{Red} Ymnus.} \hyperlink{veni-redemptor}{Veni Redemptor.} \V Tanquam sponsus. {\color{Red} Ad Mag. Ant.} Qui de terra est de terra loquitur, qui de celo venit super omnes est, et quod vidit et audivit hoc testatur, et testimonium eius nemo accipit. Qui autem acceperit eius testimonium signavit quia deus verax est. {\color{Red} Oracio.}
\lettrine[lines=2]{\zallmancaps \color{Blue} D}{eus} \hypertarget{deus-nobis-circumcisionis}{qui} nobis nati salvatoris diem celebrare concedis octavum, fac nos quesumus eius perpetua divinitate muniri, cuius sumus carnali commercio reparati. Qui tecum. \ul{Memoria de sancto Silvestro, deinde de sanctis. Postea de Dominica: si hac die evenerit.}
\newline {\color{Red} Ad Matutinas Invit.} \hyperlink{christus-natus-invitatorium}{Christus natus est.} {\color{Red} Ymnus.} \hyperlink{christe-redemptor}{Christe Redemptor.} {\color{Red} In Primo Nocturno. Ant.} \hypertarget{ant-dominus-dixit}{Dominus} dixit ad me filius meus es tu, ego hodie genui te. {\color{Red} Ps.} \psalm{2} {\color{Red} Ant.} In sole posuit tabernaculum suum, et ipse tanquam sponsus procedens de thalamo suo. {\color{Red} Ps.} \psalm{18} {\color{Red} Ant.} Elevamini porte eternales et introibit rex glorie. {\color{Red} Ps.} \psalm{23} \V Verbum caro.
\newline{\color{Red} Lectio Prima.} \ul{Sermo Ambrosii.} {\color{Red} \grecross}
\lettrine[lines=2]{\zallmancaps \color{Red} C}{ircumciditur} hodie puer ille qui natus est nobis. Factus est enim sub lege, ut eos qui sub lege erant lucrifaceret.
Circumcisus autem Dominico dignus iudicabatur obtutu, %typo: obtuitu
quia oculi Domini super iustos. Circumcisio autem, purgationem significat delictorum. Sed quoniam humane carnis et mentis fragilitas, inextricabilibus vitiis implicatur, culpe totius futura purgatio resurrectionis prefigurabatur etate.
\begin{responsory}[in-principio-erat]
{In principio erat verbum et verbum erat apud deum et deus erat verbum, hoc erat in principio apud deum, * Omnia per ipsum facta sunt et sine ipso factum est nichil.}{Quod factum est in ipso vita erat et vita erat lux hominum.}
\end{responsory}%
\newline {\color{Red} Lectio Secunda.} \ul{Origenes in sermone. Quod mortuus.} {\color{Red} \grecross}
% https://la.wikisource.org/wiki/Translatio_XXXIX_Homiliarum
\lettrine[lines=2]{\zallmancaps \color{Blue} S}{icut} autem Christo commortui sumus illo moriente, sic et cum eo circumcisi sumus, et sollemni purgatione mundati, nequaquam iam indigemus circumcisione carnali. Audi Paulum. %typo: Unum Paulus.
In quo et circumcisi estis circumcisione sine manibus in expoliatione corporis carnis. Et mors igitur et circumcisio eius pro nobis facte sunt.
\begin{responsory}[benedictus-qui-venit]
{Benedictus qui venit in nomine domini, deus dominus, * Et illuxit nobis.}{Lapidem quem reprobaverunt edificantes hic factus est in caput anguli.}
\end{responsory}%
\newline {\color{Red} Lectio Tertia.}
\lettrine[lines=2]{\zallmancaps \color{Red} E}{t} cum impleti essent dies circumcidendi puerum, vocatum est nomen eius Ihesus. Vocabulum Ihesu gloriosum, omni adoratu cultuque dignissimum, nomen quod est super omne nomen non debuit primum ab hominibus appellari, neque ab eis efferri in mundum, sed ab excellentiori quadam maiorique natura. Unde signanter evangelista addidit dicens. Quod vocatum fuerat ab angelo ante quam conciperetur in utero.
\begin{responsory-final}[ecce-agnus-dei]
{Ecce agnus dei qui tollit peccata mundi, ecce de quo dicebam vobis, * Qui post me venit ante me factus est, * Cuius non sum dignus corrigiam calciamenti solvere.}{Hoc est testimonium quod perhibuit Iohannes dicens.}{Cuius}
\end{responsory-final}%
\newline {\color{Red} In Secundo Nocturno. Ant.} Speciosus forma pre filiis hominum, diffusa est gratia in labiis tuis. {\color{Red} Ps.} \psalm{44} {\color{Red} Ant.} Suscepimus deus misericordiam tuam in medio templi tui. {\color{Red} Ps.} \psalm{47} {\color{Red} Ant.} Homo natus est in ea et ipse fundavit eam altissimus. {\color{Red} Ps.} \psalm{86} \V Notum fecit dominus.
\newline {\color{Red} Lectio Quarta.} \ul{Sermo Augustini in Kalendis Ianuarii.} {\color{Red} \grecross}
\lettrine[lines=2]{\zallmancaps \color{Blue} A}{dmonemus} caritatem vestram fratres, ut memineritis quomodo cantastis. Cantastis enim. Salva nos Domine Deus noster: et congrega nos de Gentibus. Nolite misceri Gentibus, similitudine morum atque factorum. Dant illi strenas: date vos elemosinas. Advocantur illi cantionibus luxuriarum, advocemini vos sermonibus Scripturarum. Currunt illi ad theatrum, vos ad ecclesiam. Inebriantur illi, vos ieiunate. Si hodie non potestis ieiunare, saltem cum sobrietate prandete.
\begin{responsory}[nesciens-mater]
{Nesciens mater virgo virum peperit sine dolore salvatorem seculorum, * Ipsum regem angelorum sola virgo lactabat ubere de celo pleno.}{Beata viscera Marie virginis que portaverunt eterni patris filium.}
\end{responsory}%
\newline {\color{Red} Lectio Quinta.}
\lettrine[lines=2]{\zallmancaps \color{Red} M}{ulti} luctabuntur hodie cum verbo quod audierunt. Diximus enim ut dicent. Nolite strenas dare pauperibus. Parum est ut tantum detis, eciam amplius date. Non vultis amplius, vel tantum date. Sed dicis michi. Quando strenas do, accipio et ego. Quid ergo? Quando das pauperi nichil accipies? Oblitus es illius promissionis ubi dicitur, qui dat pauperibus nunquam egebit? Oblitus es eciam illud quid dicturus est deus eis qui pauperibus dederunt, venite benedicti Patris mei, percipite regnum?
\begin{responsory}[sancta-et-immaculata]
{Sancta et immaculata virginitas quibus te laudibus efferam nescio, * Quia quem celi capere non poterant tuo gremio contulisti.}{Benedicta tu in mulieribus et benedictus fructus ventris tui.}
\end{responsory}%
\newline {\color{Red} Lectio Sexta.}
\lettrine[lines=2]{\zallmancaps \color{Blue} A}{udite} fratres apostolum qui dixit. Nolo vos socios fieri demoniorum. Demonia delectantur canticis vanitatis nugatorio spectaculo, turpitudinibus variis theatrorum, insania circi, crudelitate amphyteatri, certaminibus animosis eorum qui pro pestilentibus hominibus lites suscipiunt, pro mimo, pro ystrione pro pantomimo, pro auriga, pro venatore. Ista facientes quasi thura ponunt demoniis. Nolite inquit effici participes eorum.%typo: pantamimo
\begin{responsory-final}[confirmatum-est-cor]
{Confirmatum est cor virginis in quo divina mysteria angelo narrante suscepit, que speciosum forma pre filiis hominum castis concepit visceribus, * Et benedicta in eternum, * Deum nobis protulit et hominem.}{Domus pudici pectoris templum repente fit dei intacta nesciens virum, verbo concepit filium.}{Deum}
\end{responsory-final}%
\newline {\color{Red} In Tertio Nocturno. Ant.} Exultabunt omnia ligna silvarum ante faciem domini quoniam venit. {\color{Red} Ps.} \psalm{95} {\color{Red} Ant.} In principio et ante secula deus erat verbum, ipse natus est nobis salvator mundi. {\color{Red} Ps.} \psalm{96} {\color{Red} Ant.} Ante luciferum genitus et ante secula dominus, salvator noster hodie nasci dignatus est. {\color{Red} Ps.} \psalm{97} \V Ipse invocabit.
\newline {\color{Red} Lectio VII. Secundum Lucam.}
\lettrine[lines=2]{\zallmancaps \color{Red} I}{n} illo tempore: Postquam consummati sunt dies octo ut circumcideretur puer, vocatum est nomen eius Ihesus. Et reliqua.
\newline {\color{Red} Omelia Venerabilis Bede Presbiteri.} \grecross
\lettrine[lines=2]{\zallmancaps \color{Blue} S}{anctam} venerandamque presentis festi memoriam: paucis quidem verbis evangelista comprehendit, sed non pauca celestis mysterii virtute gravidam reliquit. Ut autem nobis necessariam obediendi virtutem precipuo commendaret exemplo, factum sub lege filium suum misit deus in mundum: non quia ipse legi quicquam deberet, qui unus magister noster, unus est legislator et iudex, sed ut eos qui sub lege positi legis onera portare nequiverant, sua compassione iuvaret: ac de servili conditione que sub lege erat ereptos, in adoptionem filiorum que per gratiam est sua largitate reduceret.
\begin{responsory}[o-regem-celi]
{O regem celi cui talia famulantur obsequia stabulo ponitur qui continet mundum, * Iacet in presepio et in nubibus tonat.}{Qui celum terramque regit in secula solus.}
\end{responsory}%
\newline {\color{Red} Lectio VIII.}
\lettrine[lines=2]{\zallmancaps \color{Red} S}{uscepit} igitur circumcisionem lege decretam in carne: qui absque omni prorsus labe pollutionis apparuit in carne. Et qui in similitudine carnis peccati, non autem in carne peccati advenit, remedium quo caro peccati consueverat mundari non respuit: sicut eciam undam baptismatis qua nove gratie populos a peccatorum sorde lavari voluit, ipse non necessitatis sed exempli causa subiit.
\begin{responsory}[congratulamini-michi-circumcisione]
{Congratulamini michi omnes qui diligitis dominum quia cum essem parvula placui altissimo, * Et de meis visceribus genui deum et hominem.}{Beatam me dicent omnes generationes qui ancillam humilem respexit deus.}
\end{responsory}%
\newline {\color{Red} Lectio IX.}
\lettrine[lines=2]{\zallmancaps \color{Blue} S}{cire} etenim debet vestra fraternitas quia idem salutifere curationis auxilium circumcisio in lege contra originalis peccati vulnus agebat, quod nunc baptismus agere revelate gratie tempore consuevit, excepto quod regni celestis ianuam necdum intrare poterant. Qui enim nunc per evangelium suum terribiliter ac salubriter clamat nisi qui renatus fuerit ex aqua et spiritu sancto non potest introire in regnum Dei, ipse dudum per legem suam clamabat: Masculus cuius preputii caro circumcisa non fuerit, peribit anima illa de populo suo, quia pactum meum irritum fecit.
\begin{responsory-doxology}[verbum-caro-factum-doxology]
{Verbum caro factum est et habitavit in nobis cuius gloriam vidimus quasi unigeniti a patre, * Plenum gratie et veritatis.}{In principio erat verbum et verbum erat apud deum et deus erat verbum.}
\end{responsory-doxology}%
\newline {\color{Red} In Laudibus. Ant.} O admirabile commercium creator generis humani animatum corpus sumens de virgine nasci dignatus est, et procedens homo sine semine, largitus est nobis suam deitatem. {\color{Red} Ps.} \psalm{92} {\color{Red} Ant.} Quando natus es ineffabiliter ex virgine, tunc implete sunt scripture, pluvia in vellus descendisti, ut salvum faceres genus humanum, te laudamus deus noster. {\color{Red} Ps.} \psalm{99} {\color{Red} Ant.} Rubum quem viderat Moyses incombustum, conservatam agnovimus tuam laudabilem virginitatem, dei genitrix intercede pro nobis. {\color{Red} Ps.} \psalm{62} {\color{Red} Ant.} Germinavit radix Iesse, orta est stella ex Iacob, virgo peperit salvatorem, te laudamus deus noster. {\color{Red} Can.} \psalm{benedicite} {\color{Red} Ant.} Ecce Maria genuit nobis salvatorem, quem Iohannes videns exclamavit dicens, ecce agnus dei ecce qui tollit peccata mundi alleluya. {\color{Red} Ps.} \psalm{148} {\color{Red} Cap.} Multifarie. {\color{Red} Ymnus.} \hyperlink{a-solis-ortus}{A Solis Ortus.} \V Benedictus qui venit.
\newline {\color{Red} Ad Ben. Ant.} Mirabile mysterium declaratur, hodie innovantur nature, deus homo factus est id quod fuit permansit et quod non erat assumpsit, non commixtionem passus neque divisionem. {\color{Red} Oracio.} \hyperlink{deus-nobis-circumcisionis}{Deus qui nobis.} \ul{Memoria de sanctis, Capitula, Responsoria, \Vbar , ad horas sicut in die Natalis cum Oratione.} \hyperlink{deus-nobis-circumcisionis}{Deus qui nobis.}
\newline {\color{Red} Ad Vesperas.} \ul{Capitulum, Ymnus, \Vbar , ut in primis vesperis.} {\color{Red} Ad Mag. Ant.} Magnum hereditatis mysterium templum dei factus est uterus nesciens virum, non est pollutus ex ea carnem assumens, omnes gentes venient dicentes gloria tibi domine.
\newline \ul{Memoria de sancto Stephano. Ant.} Tu principatum. \ul{deinde de aliis sanctis. In vigilia Epiphanie si in dominica contigerit fiant, novem lectiones. Sex lectiones de tribus lectionibus mediis in Natali et tres de omelia sequenti. Invitat, Ymnus, Ant., Psalmi, \Vbar , Responsoria, Antiphone Laudum, Capitulum, Ymnus, \Vbar , de Circumcisione.}
\newline {\color{Red} Ad Ben. Ant.} Revertere in terram Iuda mortui sunt enim qui querebant animam pueri. {\color{Red} Oracio.}
\lettrine[lines=2]{\zallmancaps \color{Blue} C}{orda} \hypertarget{corda-nostra-vigilia-epiphanie}{nostra} quesumus domine venture festivitatis splendor illustret, quo mundi huius tenebris carere valeamus, et perveniamus ad patriam claritatis eterne. Per.
\newline \ul{Per horas diei, Ant., Cap., \Rbar , \Vbar , de Circumcisione, cum Oracione.} \hyperlink{corda-nostra-vigilia-epiphanie}{Corda.} \ul{Si vero dominica non fuerit, ita fiat officium. Invitat.} \hyperlink{christus-natus-invitatorium}{Christus natus.} \ul{Ymnus.} \hyperlink{christe-redemptor}{Christe Redemptor.} \ul{Super psalmos. Ant.} \hyperlink{ant-dominus-dixit}{Dominus dixit.} \ul{sola dicatur, psalmi, ut in Circumcisione. \Vbar .} Notum fecit.
\newline {\color{Red} Lectio Prima. Secundum Matheum.}
\lettrine[lines=2]{\zallmancaps \color{Red} I}{n} illo tempore: Defuncto Herode, ecce angelus domini apparuit in somnis Ioseph in Egypto, dicens. Surge et accipe puerum et matrem eius: et vade in terram Israhel. Et reliqua.
\newline {\color{Red} Omelia Venerabilis Bede Presbiteri.}
\lettrine[lines=2]{\zallmancaps \color{Blue} Q}{uod} occisis pro domino pueris Herodes non longe post obiit, et Ioseph monente angelo dominum cum matre ad terram Israhel reduxit, significat omnes persecutiones que contra ecclesiam erant movende, ipsorum persecutorum morte vindicandas, eisdemque mulctatis persecutoribus pacem ecclesie denuo reddendam, et sanctos qui latuerant ad sua loca fuisse reversuros. \R \hyperlink{confirmatum-est-cor}{Confirmatum.}
\newline {\color{Red} Lectio Secunda.}
\lettrine[lines=2]{\zallmancaps \color{Red} P}{ossumus} quoque odium Herodis quo perdere Ihesum voluit, super persecutionibus que apostolorum sunt temporibus facte in Iudea specialiter accipere, quando invalescente invidia, predicatores verbi sunt pene omnes expulsi de provincia et in gentibus predicaturi, sunt longe lateque dispersi. Sic et factum est, ut gentilitas que per Egyptum figuratur, peccatis ante tenebrosa lumen verbi perciperet. \R \hyperlink{o-regem-celi}{O regem celi.}
\newline {\color{Red} Lectio Tertia.}
\lettrine[lines=2]{\zallmancaps \color{Blue} H}{oc} est autem puerum Ihesum et matrem eius per Ioseph in Egyptum transferri, dominice scilicet fidem incarnationis et ecclesie societatem per doctores sanctos gentibus committi. Quod erant in Egypto usque ad obitum Herodis: indicat figurate fidem Christi in gentibus mansuram donec plenitudo earum introeat, et sic omnis Israhel salvus fiat. \R \hyperlink{congratulamini-michi-circumcisione}{Congratulamini.} \psalm{tedeum} \V Puer natus.
\newline \ul{In Laudibus sola Ant.} O admirabile. \ul{Ps.} \psalm{92} \ul{Capitulum.} Multifarie. \ul{Ymnus.} \hyperlink{a-solis-ortus}{A Solis.} \ul{\Vbar .} Benedictus qui venit. \ul{Ad Ben. Ant.} Revertere. \ul{Oracio} \hyperlink{corda-nostra-vigilia-epiphanie}{Corda Nostra.} \ul{Ad Primam. Ant.} O admirabile. \ul{Cap.} Regi. \ul{Oracio.} In hac hora. \ul{Ad Terciam, Sextam, et Nonam. Ant., Cap., \Rbar , \Vbar , de Circumcisione cum oratione.} \hyperlink{corda-nostra-vigilia-epiphanie}{Corda.}
\newline {\color{Red} Ad Vesperas. Capitulum.}
\lettrine[lines=2]{\zallmancaps \color{Red} S}{urge} \hypertarget{surge-capitulum}{illuminare} Hierusalem, quia venit lumen tuum, et gloria domini super te orta est. \R \hyperlink{in-columbe}{In columbe.} {\color{Red} II. Ymnus.}
\hymn{hostis-herodes}{Hostis Herodes impie}{
\lettrine[lines=2]{\zallmancaps \color{Blue} H}{ostis} Herodes impie
christum venire quid times,
non eripit mortalia,
qui regna dat celestia.
\par {\bfseries \color{Red} I}bant magi quam viderant
stellam sequentes previam,
lumen requirunt lumine,
deum fatentur munere.
\par {\bfseries \color{Blue} L}avacra puri gurgitis
celestis agnus attigit,
peccata que non detulit
nos abluendo sustulit.
\par {\bfseries \color{Red} N}ovum genus potentie
aque rubescunt ydrie,
vinumque iussa fundere,
mutavit unda originem.
\par {\bfseries \color{Blue} G}loria tibi domine
qui apparuisti hodie
cum patre et sancto spiritu,
in sempiterna secula. Amen.
}
\newline \ul{Iste versus dicatur ad omnes ymnos per octavam Epiphanie.}
\newline \V Reges Tharsis et insule munera offerent. \R Reges Arabum et Saba dona adducent.
\newline {\color{Red} Ad Mag. Ant.} Magi videntes stellam, dixerunt ad invicem, hoc signum magni regis est, eamus et inquiramus eum et offeramus ei munera aurum thus et mirram. {\color{Red} Oracio.} \hyperlink{corda-nostra-vigilia-epiphanie}{Corda nostra.}
\newline {\color{Red} Ad Completorium. Super psalmos. Ant.} Lux de luce apparuisti Christe cui magi munera offerunt alleluya alleluya alleluya. {\color{Red} Ad Nunc Dimittis. Ant.} Alleluya, omnes de Saba venient alleluya aurum et thus deferentes alleluya, alleluya.
\newline \ul{Predicte antiphone dicantur ad Completorium usque ad Octavam inclusive.}
{\ubsubsection{In Epiphania Ad Matutinas}{breviarium-in-epiphania-ad-matutinas}}
\par \noindent \ul{Non dicatur invitatorium neque ymnus, sed statim post,} Domine labia \ul{et} Deus in adiutorium, \ul{incipiatur,} {\color{Red} Ant.} Afferte domino filii dei adorate dominum in aula sancta eius. {\color{Red} Ps.} \psalm{28} {\color{Red} Ant.} Fluminis impetus letificat alleluya, civitatem dei alleluya. {\color{Red} Ps.} \psalm{45} {\color{Red} Ant.} Psallite deo nostro psallite, psallite regi nostro psallite sapienter. {\color{Red} Ps.} \psalm{46} \V Reges Tharsis.
\newline {\color{Red} Lectio Prima.}
\lettrine[lines=2]{\zallmancaps \color{Red} O}{mnes} sitientes venite ad aquas: et qui non habetis argentum, properate, emite, et comedite. Venite emite absque argento, et absque ulla commutatione: vinum et lac. Quare appenditis argentum non in panibus, et laborem vestrum non in saturitate? Audite audientes me et comedite bonum: et delectabitur in crassitudine anima vestra. Hec dicit.
\begin{responsory}[hodie-in-iordane]
{Hodie in Iordane baptizato domino aperti sunt celi, et sicut columba super eum spiritus mansit, et vox patris intonuit * Hic est filius meus dilectus in quo michi complacui.}{Descendit Spiritus Sanctus corporali specie sicut columba in ipsum et vox de celo facta est.}
\end{responsory}%
\newline {\color{Red} Lectio Secunda.}
\lettrine[lines=2]{\zallmancaps \color{Blue} I}{nclinate} aurem vestram et venite ad me: audite, et vivet anima vestra. Et feriam vobiscum pactum sempiternum: misericordias David fideles. Ecce testem populis dedi eum, ducem ac preceptorem gentibus. Ecce gentem quam nesciebas vocabis, et gentes que te non cognoverunt ad te current: propter Dominum Deum tuum, et Sanctum Israhel, quia glorificavit te. Hec dicit.
\begin{responsory}[in-columbe]
{In columbe specie Spiritus Sanctus visus est, paterna vox audita est, * Hic est filius meus dilectus in quo michi bene complacui.}{Celi aperti sunt super eum, et vox patris intonuit.}
\end{responsory}%
\newline {\color{Red} Lectio Tertia.}
\lettrine[lines=2]{\zallmancaps \color{Red} Q}{uerite} Dominum dum inveniri potest, invocate eum dum prope est. Derelinquat impius viam suam, et vir iniquus cogitationes suas: et revertatur ad Dominum, et miserebitur eius, et ad Deum nostrum, quoniam multus est ad ignoscendum. Hec dicit.
\begin{responsory-doxology}[reges-tharsis]
{Reges Tharsis et insule munera offerent, * Reges Arabum et Saba dona domino deo adducent.}{Omnes de Saba venient aurum et thus deferentes.}
\end{responsory-doxology}%
\newline {\color{Red} In Secundo Nocturno. Ant.} Omnis terra adoret te et psallat tibi psalmum dicat nomini tuo domine. {\color{Red} Ps.} \psalm{65} {\color{Red} Ant.} Reges Tharsis et insule munera offerent regi domino. {\color{Red} Ps.} \psalm{71} {\color{Red} Ant.} Omnes gentes quascumque fecisti venient et adorabunt coram te domine. {\color{Red} Ps.} \psalm{85} \V Omnes de Saba venient. \R Aurum et thus deferentes.
\newline \ul{Maximus in sermone.} Cum plura. {\color{Red} Lectio Quarta.} \grecross
% https://la.wikisource.org/wiki/Homiliae_(Maximus)/1
\lettrine[lines=2]{\zallmancaps \color{Blue} S}{icut} posteritati sue fidelis mandavit antiquitas, hodie salvator celestibus ostensus indiciis, a Chaldeis est adoratus. Hodie Christus beati Iohannis ministerio: fluenta Iordanis benedictione proprii baptismatis consecravit. Hodie invitatus ad nuptias divinitatis sue potentiam manifestans: aquam mutavit in vinum. Aquas inquam quas sanctificavit in baptismo, nobilitavit in nuptiis. Tria quidem hec mysteria unius diei, per unum Dominum dignatio Trinitatis implevit.
\begin{responsory}[illuminare-illuminare]
{Illuminare illuminare Hierusalem, venit lux tua, * Et gloria domini super te orta est.}{Et ambulabunt gentes in lumine tuo et reges in splendore ortus tui.}
\end{responsory}%
\newline {\color{Red} Lectio Quinta.}
\lettrine[lines=2]{\zallmancaps \color{Red} H}{odie} magus per stellam reperit, quod Iudeus noluit credere per prophetas. Hodie gentilitas edocta de celo,
%pp093
redemptorem mundi inquirit ut regem, muneratur ut fortem, adorat ut Deum. Recte autem nove stella resplenduit cuius gens perfida nec radium occultare possit nec abscondere veritatem: ubi celum ipsum universorum oculis sydereo lumine fulgebat. Quam mirabile autem est hoc quod exiguus stelle radius alienigenarum corda permovit, cum illum Iudaicum populum cui mare divisum, cui prebitum manna de nubibus, nec ignea potuerit de celo micans columna convertere.
\begin{responsory}[omnes-de-saba]
{Omnes de Saba venient aurum et thus deferentes et * Laudem domino annunciantes.}{Reges Tharsis et insule munera offerent, reges Arabum et Saba dona adducent.}
\end{responsory}%
\newline {\color{Red} Lectio Sexta.}
\lettrine[lines=2]{\zallmancaps \color{Blue} Q}{uid} hac sollemnitate sacratius, in qua partum virginis qui latebat in cunis, stella fulgentior ceteris coruscavit, in qua unigenitum Dei in sanctificationem sui Iordanis unda suscepit, homo Spiritum sanctum vidit, Patrem mundus audivit. In qua eciam sobrio illo vini mirabilis poculo peccatis ebrium mortale genus visitasse se Christus ostendit? Et ideo karissimi in spiritualibus donis, lubrice carnis voluptate deposita, spiritaliter gaudeamus, ut tenebras cordis nostri, lumen stelle celestis irradiet, benedictio paterne vocis illustret, et salutare nos vinum, Christo propinante letificet. Amen.
\begin{responsory-doxology}[magi-veniunt]
{Magi veniunt ab oriente Hierosolimam querentes et dicentes, ubi est qui natus est cuius stellam vidimus, * Et venimus adorare eum.}{Vidimus stellam eius in oriente.}
\end{responsory-doxology}%
\newline {\color{Red} In Tertio Nocturno. Ant.} Venite adoremus deum quia ipse est dominus deus noster. {\color{Red} Ps.} \psalm{94} \ul{predicta antiphona cum psalmo non cantetur infra octavas, sed loco ipsius dicatur.} {\color{Red} Ant.} Homo natus est in ea et ipse fundavit eam altissimus. {\color{Red} Ps.} \psalm{86} {\color{Red} Ant.} Adorate dominum alleluya, in aula sancta eius alleluya. {\color{Red} Ps.} \psalm{95} {\color{Red} Ant.} Adorate deum alleluya omnes angeli eius alleluya. {\color{Red} Ps.} \psalm{96} \V Omnes terra adoret te deus et psallat tibi. \R Psalmum dicat nomini tuo.
\newline {\color{Red} Lectio VII. Secundum Matheum.}
\lettrine[lines=2]{\zallmancaps \color{Blue} C}{um} natus esset Ihesus in Bethleem vide in diebus Herodis regis: Ecce magi ab oriente venerunt Hierosolimam dicentes. Ubi est qui natus est rex Iudeorum? Et reliqua.
\newline {\color{Red} Omelia Beati Gregorii Pape.}
\lettrine[lines=2]{\zallmancaps \color{Red} S}{icut} ex lectione evangelica fratres audistis, celi rege nato rex terre turbatus est, quia nimirum terrena altitudo confunditur cum celsitudo celestis aperitur. Sed querendum nobis est, quidnam sit quod redemptore nato, pastoribus in Iudea angelus apparuit, atque ad adorandum hunc ab oriente magos non angelus sed stella perduxit. Quia videlicet Iudeis tanquam ratione utentibus rationale animal id est angelus predicare debuit, gentiles vero quia uti ratione nesciebant ad cognoscendum deum non per vocem sed per signa perducuntur: quia et illis prophetie tanquam fidelibus non infidelibus, et istis signa tanquam infidelibus non fidelibus data sunt.
\begin{responsory}[stella-quam]
{Stella quam viderant magi in oriente antecedebat eos, donec venirent ad locum ubi puer erat, videntes autem eam, * Gavisi sunt gaudio magno.}{Et intrantes domum invenerunt puerum cum Maria matre eius, et procidentes adoraverunt eum.}
\end{responsory}%
\newline {\color{Red} Lectio VIII.}
\lettrine[lines=2]{\zallmancaps \color{Blue} E}{t} notandum quod redemptorem nostrum, cum iam perfecte esset etatis eisdem gentilibus apostoli predicant, eumque parvulum et necdum per humani corporis officium loquentem stella gentibus denuntiat: quia nimirum rationis ordo poscebat, ut et loquentem iam dominum loquentes nobis predicatores innotescerent, et necdum loquentem elementa muta predicarent.
\begin{responsory}[videntes-stellam]
{Videntes stellam magi gavisi sunt gaudio magno, et intrantes domum invenerunt puerum cum Maria matre eius, et procidentes adoraverunt eum, * Et apertis thesauris suis obtulerunt ei munera, aurum thus et mirram.}{Stella quam viderant magi in oriente antecedebat eos, usque dum veniens staret supra ubi erat puer.}
\end{responsory}%
\newline {\color{Red} Lectio IX.}
\lettrine[lines=2]{\zallmancaps \color{Red} S}{ed} in omnibus signis que vel nascente domino vel moriente monstrata sunt, considerandum nobis est quanta fuerit in quorumdam Iudeorum corde duritia, que hunc nec per prophetie donum, nec per miracula agnovit. Omnia quippe elementa deum venisse testata sunt.
\begin{responsory-final}[tria-sunt]
{Tria sunt munera preciosa, que obtulerunt magi domino in die ista et habent in se divina mysteria, * In auro ostenditur regis potentia in thure sacerdotem magnum considera, * Et in mirra dominicam sepulturam.}{Salutis nostre actorem magi venerati sunt in cunabilis et de thesauris suis mysticas ei munerum species obtulerunt.}{Et in mirra}
\end{responsory-final}%
\newline \ul{Postea sequitur.} Dominus vobiscum. Et cum spiritu tuo.
\newline {\color{Red} Sequencia sancti Evangelium Secundum Lucam.}
\lettrine[lines=2]{\zallmancaps \color{Red} F}{actum} est autem cum baptizaretur omnis populus, et Ihesu baptizato et orante apertum est celum, et descendit Spiritus Sanctus corporali specie sicut columba in ipsum. Et vox de celo facta est. Tu es filius meus dilectus in te complacui michi. Et ipse Ihesus erat incipiens quasi annorum triginta, ut putabatur filius Ioseph. Qui fuit Hely. Qui fuit Mathan. Qui fuit Levi. Qui fuit Melchi. Qui fuit Iamne. Qui fuit Ioseph. Qui fuit Mathathie. Qui fuit Amos. Qui fuit Naum. Qui fuit Hesly. Qui fuit Nagge. Qui fuit Maath. Qui fuit Mathathie. Qui fuit Semei. Qui fuit Ioseph. qui fuit Iuda. Qui fuit Iohanna. Qui fuit Resa. Qui fuit Zorobabel. Qui fuit Salathiel. Qui fuit Neri. Qui fuit Melchi. Qui fuit Addi. Qui fuit Chosan. Qui fuit Elmadan. Qui fuit Her. Qui fuit Ihesu. Qui fuit Heliezer. Qui fuit Ioram. Qui fuit Mathath. Qui fuit Levi. Qui fuit Symeon. Qui fuit Iuda. Qui fuit Ioseph. Qui fuit Iona. Qui fuit Heliachim. Qui fuit Melcha. Qui fuit Menna. Qui fuit Mathathia. Qui fuit Nathan. Qui fuit David. Qui fuit Iesse. Qui fuit Obeth. Qui fuit Booz. Qui fuit Salmon. Qui fuit Naason. Qui fuit Aminadab. Qui fuit Aram. Qui fuit Esron. Qui fuit Phares. Qui fuit Iude. Qui fuit Iacob. Qui fuit Ysaac. Qui fuit Abrahe. Qui fuit Thare. Qui fuit Nachor. Qui fuit Saruch. Qui fuit Ragau. Qui fuit Phalech. Qui fuit Heber. Qui fuit Sale. Qui fuit Caynan. Qui fuit Arphaxad. Qui fuit Sem. Qui fuit Noe. Qui fuit Lamech. Qui fuit Matusale. Qui fuit Enoch. Qui fuit Iared. Qui fuit Malaleel. Qui fuit Caynan. Qui fuit Enos. Qui fuit Seth. Qui fuit Adam. Qui fuit Dei. Ihesus autem plenus Spiritu Sancto regressus est a Iordane.
\psalm{tedeum} \V Vidimus stellam eius in oriente. \R Et venimus adorare eum.
\newline {\color{Red} In Laudibus. Ant.} Ante luciferum genitus et ante secula dominus, salvator noster hodie mundo apparuit. {\color{Red} Ps.} \psalm{92} {\color{Red} Ant.} Venit lumen tuum Hierusalem, et gloria domini super te orta est, et ambulabunt gentes in lumine tuo alleluya. {\color{Red} Ps.} \psalm{99} {\color{Red} Ant.} Apertis thesauris suis obtulerunt magi domino aurum thus et mirram, alleluya. {\color{Red} Ps.} \psalm{62} {\color{Red} Ant.} Maria et flumina benedicite domino, ymnum dicite fontes domino alleluya. {\color{Red} Can.} \psalm{benedicite} {\color{Red} Ant.} Tria sunt munera que obtulerunt magi domino aurum thus et mirram, filio dei regi magno alleluya. {\color{Red} Ps.} \psalm{148} {\color{Red} Cap.} \hyperlink{surge-capitulum}{Surge illuminare.}
\hymn{a-patre-unigenitus}{A patre unigenitus}{
\lettrine[lines=2]{\zallmancaps \color{Blue} A}{}\ patre unigenitus
ad nos venit per virginem,
baptisma cruce consecrans,
cunctos fideles generans.
\par {\bfseries \color{Red} D}e celo celsus prodiit
excepit formam hominis,
facturam morte redimens
gaudia vite largiens.
\par {\bfseries \color{Blue} H}oc te redemptor quesumus.
illabere propitius
clarumque nostris sensibus
lumen prebe fidelibus.
\par {\bfseries \color{Red} M}ane vobiscum domine,
noctem obscuram remove,
omne delictum ablue,
piam medelam tribue.
\par {\bfseries \color{Blue} Q}uem iam venisse novimus
redire item credimus,
sub sceptro tuo inclito,
tuum defende populum.
\par {\bfseries \color{Red} G}loria tibi domine
qui apparuisti hodie
cum patre et sancto spiritu,
in sempiterna secula. Amen.
}
\newline \V Adorate dominum. \R In aula sancta eius.
\newline {\color{Red} Ad Ben. Ant.} Hodie celesti sponso iuncta est ecclesia, quoniam in Iordane lavit Christus eius crimina, currunt cum muneribus magi ad regales nuptias et ex aqua facto vino letantur convive, alleluya. {\color{Red} Oracio.}
\lettrine[lines=2]{\zallmancaps \color{Blue} D}{eus} \hypertarget{deus-qui-hodierna-epiphanie}{qui} hodierna die unigenitum tuum gentibus stella duce revelasti, concede propitius, ut qui iam te ex fide cognovimus, usque ad contemplandum speciem tue celsitudinis perducamur. Per eundem. \ul{Ad horas, antiphone Laudum.}
\newline {\color{Red} Ad Primam.} \R Ihesu Christe. \V Qui apparuisti hodie. \ul{Iste \Vbar . dicatur usque ad octavas inclusive.}
\newline {\color{Red} Ad Terciam. Cap.} \hyperlink{surge-capitulum}{Surge.}
\begin{responsory-breve}[reges-tharsis-breve]
{Reges Tharsis et insule munera offerent, * Alleluya alleluya.}{Reges Arabum et Saba dona adducent.}
\end{responsory-breve}%
\V Omnes de Saba venient.
\newline {\color{Red} Ad Sextam. Capitulum.}
\lettrine[lines=2]{\zallmancaps \color{Red} S}{uper} te Hierusalem orietur dominus, et gloria eius in te videbitur, et ambulabunt gentes in lumine tuo, et reges in splendore ortus tui.
\begin{responsory-breve}[omnes-de-saba-breve]
{Omnes de Saba venient, * Alleluya alleluya.}{Aurum et thus deferentes.}
\end{responsory-breve}%
\V Omnis terra adoret te deus et psallat tibi.
\newline {\color{Red} Ad Nonam. Capitulum.}
\lettrine[lines=2]{\zallmancaps \color{Blue} O}{mnes} de Saba venient aurum et thus deferentes et laudem domino annunciantes.
\begin{responsory-breve}[omnis-terra-adoret-breve]
{Omnis terra adoret te deus et psallat tibi, * Alleluya alleluya.}{Psalmum dicat nomini tuo.}
\end{responsory-breve}%
\V Adorate dominum. \R In aula sancta eius.
\newline {\color{Red} Ad Vesperas.} \ul{Capitulum, Ymnus, \Vbar , ut in Primis Vesperis} {\color{Red} Ad Mag. Ant.} Ab oriente venerunt magi in Bethleem adorare dominum et apertis thesauris suis preciosa munera obtulerunt aurum sicut regi magno thus sicut deo vero mirram sepulture eius, alleluya.
{\ubsubsection{Per Octavas Epiphanie}{breviarium-per-octavas-epiphanie}}
\par \noindent {\color{Red} Invitat.}
\begin{invitatory}[christus-apparuit-invitatorium]
{Christus apparuit nobis, * Venite adoremus.}
\end{invitatory}
{\color{Red} Ymnus.} \hyperlink{hostis-herodes}{Hostis Herodes.} \ul{cotidie dicatur. In ferialibus diebus: una antiphona de novem super psalmos. Prima die post festum dicatur prima, secunda die, secunda, et sic deinceps quot fuerint necessarie. Psalmi de festo, excepto psalmo,} \psalm{94} \ul{loco cuius dicatur psalmus,} \psalm{86} \ul{versiculi, antiphone, lectiones, et tria responsoria de hystoria secundum ordinem nocturnorum excepto quod quando dici contigerit Ant.} Omnis terra: \ul{dicendus est. \Vbar .} Reges Tharsis, \ul{et quando dicetur Ant.} Reges Tharsis, \ul{dicendus est \Vbar .} Omnes de Saba. \psalm{tedeum} \ul{cotidie dicatur, \Vbar .} Vidimus.
\newline \ul{In Laudibus prima antiphona laudibus super psalmos, Cap., Ymnus, \Vbar , de die.} {\color{Red} Ad Ben. Ant.} Descendit Spiritus Sanctus corporali specie sicut columba in ipsum, et vox de celo facta est hic est filius meus dilectus alleluya. \ul{Oratio sicut in die. Hore diei sicut in festo.}
\newline {\color{Red} Ad Vesperas.} \ul{Cap., Ymnus, \Vbar , de die.} {\color{Red} Ad Mag. Ant.} Videntes stellam magi gavisi sunt gaudio magno et intrantes domum obtulerunt domino aurum thus et mirram. \ul{Oracio ut supra. Predicte due antiphone dicantur cotidie infra octavas, eciam in dominica.} Descendit \ul{ad Ben. et} Videntes \ul{Ad Mag.}
\newline \ul{Si dies Epiphanie in Dominica evenerit in sequenti feria quarta legatur ad Matutinas Omelia, evangelii,} Cum factus. \ul{In laudibus vero post orationem de Epyphania dicatur, ad memoriam antiphone} Fili \ul{cum \Vbar ,} Omnis terra \ul{et oratione} \hyperlink{vota-quesumus-epiphanie}{Vota.} \ul{In Vesperis vero similiter post orationem de Epyphania dicatur ad memoriam Ant.} Puer Ihesus \ul{cum \Vbar .} Omnes de Saba \ul{et oratione} \hyperlink{vota-quesumus-epiphanie}{Vota.} \ul{Sabbato autem sequenti post orationem,} \hyperlink{deus-cuius-unigenitus-epiphanie}{Deus cuius unigenitus} \ul{fiat memoria de Dominica sequenti, per Ant.} Peccata \ul{et \Vbar .} Vespertina oratio. \ul{et per orationem} \hyperlink{omnipotens-sempiterne-dom-i-post}{Omnipotens sempiterne deus qui celestia.} \ul{Ad Laudes vero diei sequentis in qua totum officium est de octava Epiphanie fiat memoria de dominica, per Ant.} Nuptie \ul{et \Vbar .} Dominus regnavit. \ul{et orationem,} \hyperlink{omnipotens-sempiterne-dom-i-post}{Omnipotens sempiterne.} \ul{In Vesperis autem per antiphonam} Deficiente. \ul{et \Vbar .} Dirigatur \ul{et orationem predictam. In quarta vero feria sequenti ad Matutinas legatur omelia de nuptiis, et totum officium fiat de feria, et de sancto Marcello fiat tantum memoria. Ad Laudes antiphone et ad horas diei dicetur Oracio,} \hyperlink{omnipotens-sempiterne-dom-i-post}{Omnipotens sempiterne.}
\newline \ul{Quando vero infra octavas Epyphanie Dominica evenerit sic fiat officium. Ad Matutinas. Invitat.} \hyperlink{christus-apparuit-invitatorium}{Christus apparuit.} \ul{Ymnus,} \hyperlink{hostis-herodes}{Hostis.} \ul{Ant, Psalmi, \Vbar , Responsoria, de die Epiphanie excepta Ant.} Venite \ul{cum suo psalmo loco cuius dicatur Ant.} Homo natus \ul{cum psalmo} \psalm{86} \ul{In laudibus antiphone, Cap., Ymnus, \Vbar , de die. Ad Ben. Ant.} Descendit. \ul{Oracio,} \hyperlink{deus-qui-hodierna-epiphanie}{Deus qui hodierna.} \ul{Deinde fiat memoria de Dominica per Ant.} Fili \ul{cum \Vbar .} Omnis terra, \ul{et oratione,} \hyperlink{vota-quesumus-epiphanie}{Vota.} \ul{Hore sicut in die Epiphanie. Ad Vesperas. Cap., Ymnus, \Vbar , de Epiphania. Ad Mag. Ant.} Videntes. \ul{Oracio,} \hyperlink{deus-qui-hodierna-epiphanie}{Deus qui hodierna.} \ul{Deinde fiat memoria de Dominica, per Ant.} Puer \ul{cum \Vbar .} Omnes de Saba \ul{et oratione,} \hyperlink{vota-quesumus-epiphanie}{Vota.}
\newline \ul{Infra Octavas per ferias, et in Dominica et in Octava die repetantur tres lectiones de sermonibus supra notate in festo, et tres sequentes.}
\newline \ul{Maximus in sermone,} Salutare. \grecross \ {\color{Red} Lectio.}
% https://la.wikisource.org/wiki/Homiliae_(Maximus)/1
\lettrine[lines=2]{\zallmancaps \color{Red} S}{icut} nascendo gloriosus est filius dei: ita apparendo mirabilis. Mirum quod processit ex virgine: magnificum quod ostensus est e celo. Apud Iudeam in presepi vagiebat, in Chaldea inter sydera coruscabat. Apud Iudeos sordebat in pannis: apud gentiles fulgebat in gloria, ut celorum dominum testimonium celeste precederet: et actorem lucis signum luminis revelaret.
{\color{Red} Lectio.}
\lettrine[lines=2]{\zallmancaps \color{Blue} N}{ullus} ambigat quin Christus in terra degens fulgeret e celo: quia sic venit ad terras, ut non relinqueret celum. Inde est quod virginalem hunc puerum laudant angeli, mirantur pastores, celum attestatur, venerantur et magi, quia mystica eius nativitas sempiterno gaudio et terram letificavit et celum. Hunc ergo regem celestium tremit Herodes: sed obcecata tenebris Iudea, sola videre non potest, quod omnibus splendet e celo.
\newline \ul{Idem in sermone,} Ait prophetarum. \grecross \ {\color{Red} Lectio.}
\lettrine[lines=2]{\zallmancaps \color{Red} V}{idimus} iniquiunt magi stellam eius, bene eius: que specialiter eius nunciabat adventum. Et cetere quidem stelle facte sunt, ut mundi huius tempora cursusque distinguerent: hec vero prodire iussa est, ut ipsum mundi dominum, et regni celestis adesse tempus ostenderet. Stetit ait evangelista supra ubi erat puer. Nunquid non dicere videbatur, hic est puer quem natum testabar e celo? Hic est rex ille magnus, qui venit celesti imperio regnum sociare terrenum?
\newline {\color{Red} Dominica infra Octavas. Lectio VII. Secundum Lucam.}
\lettrine[lines=2]{\zallmancaps \color{Blue} I}{n} illo tempore: Cum factus esset Ihesus annorum duodecim, ascendentibus parentibus eius Hierosolimam secundum consuetudinem diei festi, consummatisque diebus cum redirent, remansit puer Ihesus in Hierusalem et non cognoverunt parentes eius. Et reliqua.
\newline {\color{Red} Omelia Venerabilis Bede Presbiteri.} \grecross
\lettrine[lines=2]{\zallmancaps \color{Red} Q}{uod} dominus per omnes annos cum parentibus in Pascha Hierosolimam venit, humane nimirum humilitatis est indicium. Hominis namque est ad offerenda Deo sacrificiorum spiritualium vota concurrere, et actorem suum orationibus lacrimisque sibi conciliare profusis.
\newline {\color{Red} Lectio VIII.}
\lettrine[lines=2]{\zallmancaps \color{Blue} F}{ecit} ergo dominus inter homines homo natus, quod faciendum hominibus per angelos imperaverat Deus. Servavit ipse legem quam dedit, ut nobis qui puri homines sumus, servandum per omnia quicquid Deus iubet ostenderet. Sequamur iter humane conversationis eius, si deitatis gloriam delectamur intueri: si optamus habitare in domo eius eterna in celis omnibus diebus vite nostre.
\newline {\color{Red} Lectio IX.}
\lettrine[lines=2]{\zallmancaps \color{Red} Q}{uod} autem duodenis in templo sedit in medio doctorum audiens et interrogans illos, humane est humilitatis indicium, immo eciam eximium discende humilitatis exemplum. Dei quippe virtus et Dei sapientia eternaque divinitas que loquitur, Ego sapientia habito in consiliis, ipsa hominem induta ad audiendos homines venire dignata est: ut nimirum hominibus quamvis summo ingenio preditis necessariam discendi verbi formam prerogaret, ne si qui discipuli veritatis fieri refugerent: magistri efficerentur erroris.
\newline {\color{Red} Ad memoriam de Dominica. Ant.} Fili quid fecisti nobis sic, ego et pater tuus dolentes querebamus te, quid est quod me querebatis nesciebatis quia in his que patris mei sunt oportet me esse alleluya. \V Omnis terra.
\lettrine[lines=2]{\zallmancaps \color{Blue} V}{ota} \hypertarget{vota-quesumus-epiphanie}{quesumus} domine supplicantes populi celesti pietate prosequere, ut et que agenda sunt videant et ad implenda que viderint convalescant. Per.
\newline {\color{Red} In Vesperis. Ad memoriam. Ant.} Puer Ihesus proficiebat etate et sapientia coram deo et hominibus. \V Omnes de Saba. {\color{Red} Oracio.} \hyperlink{vota-quesumus-epiphanie}{Vota.}
{\ubsubsection{In Octava Epyphanie}{breviarium-in-octava-epyphanie}}
\par \noindent {\color{Red} In Primis Vesperis. Cap.} \hyperlink{surge-capitulum}{Surge.} {\color{Red} Ymnus.} \hyperlink{hostis-herodes}{Hostis.} \V Reges. {\color{Red} Ad Mag. Ant.} Baptizat miles regem, servus dominum suum, Iohannes salvatorem, aqua Iordanis stupuit, columba protestata est, paterna vox audita est, hic est filius meus. {\color{Red} Oracio.}
\lettrine[lines=2]{\zallmancaps \color{Red} D}{eus} \hypertarget{deus-cuius-unigenitus-epiphanie}{cuius} unigenitus in substantia nostre carnis apparuit, presta quesumus ut per eum quem similem nobis foris agnovimus, intus reformari mereamur. Qui tecum.
\newline{\color{Red} Ad Matutinas. Invitat.} \hyperlink{christus-apparuit-invitatorium}{Christus apparuit.} \ul{Ymnus, Ant., psalmi, \Vbar , responsoria, sicut supra in die est notatum.}
\newline {\color{Red} Lectio VII. Secundum Matheum.}
\lettrine[lines=2]{\zallmancaps \color{Blue} I}{n} illo tempore: Venit Ihesus a Galilea in Iordanem ad Iohannem ut baptizaretur ab eo. Et reliqua.
\newline {\color{Red} Omelia Venerabilis Bede Presbiteri.} \grecross
\lettrine[lines=2]{\zallmancaps \color{Red} L}{ectio} sancti evangelii quam modo fratres audivimus, magnum nobis et in domino et in servo dat perfecte humilitatis exemplum. In domino quidem: quia cum sit dominus deus, non solum ab homine servo baptizari, verum eciam ipse ad hunc baptizandus venire dignatus est. In servo autem, quia cum sciret precursorem se ac baptistam sui salvatoris esse destinatum, memor tamen proprie fragilitatis: iniunctum sibi officium humiliter recusavit dicens. Ego a te debeo baptizari, et tu venis ad me?
\newline {\color{Red} Lectio VIII.}
\lettrine[lines=2]{\zallmancaps \color{Blue} V}{enit} filius Dei baptizari ab homine, non anxia necessitate abluendi alicuius sui peccati, qui peccatum non fecit ullum: sed pia dispensatione abluendi omnis nostri contagionem peccati, qui in multis offendimus omnes, et si dixerimus quia peccatum non habemus nos ipsos seducimus.
\newline {\color{Red} Lectio IX.}
\lettrine[lines=2]{\zallmancaps \color{Red} V}{enit} baptizari in aquis, ipsarum conditor aquarum, ut nobis qui in iniquitatibus concepti, et in delictis sumus generati, secunde nativitatis, que per aquam et spiritum celebratur, appetendum insinuaret mysterium. Dignatus est lavari
%pp094
aquis Iordanicis qui mundus erat a sordibus cunctis, ut ad diluendas nostrorum sordes scelerum: omnium fluenta sanctificaret aquarum.
\V Vidimus.
\newline {\color{Red} In Laudibus. Ant.} Veterem hominem renovans salvator venit ad baptismum, ut naturam que corrupta est per aquam recuperaret incorruptibili veste circumamictans nos. {\color{Red} Ps.} \psalm{92} {\color{Red} Ant.} Te qui in spiritu et igne purificas humana contagia, deum ac redemptorem omnes glorificamus. {\color{Red} Ps.} \psalm{99} {\color{Red} Ant.} Baptista contremuit, et non audet tangere sanctum dei verticem sed clamat cum tremore sanctifica me salvator. {\color{Red} Ps.} \psalm{62} {\color{Red} Ant.} Caput draconis salvator contrivit in Iordanis flumine, ab eius potestate omnes eripuit. {\color{Red} Can.} \psalm{benedicite} {\color{Red} Ant.} Magnum mysterium declaratur hodie quia creator omnium in Iordane expurgat nostra facinora. {\color{Red} Ps.} \psalm{148} \ul{Cap., Ymnus, \Vbar , de die.}
\newline {\color{Red} Ad Ben. Ant.} Precursor Iohannes exultat cum Iordane baptizato domino facta est orbis terrarum exultatio, facta est peccatorum nostrorum remissio, sanctificans aquas, ipsi omnes clamemus miserere nobis. {\color{Red} Oracio.} \hyperlink{deus-cuius-unigenitus-epiphanie}{Deus cuius unigenitus.}%typo: incoruptibili
\newline \ul{Ad horas diei Ant. Laudum, Capitula, \Rbar , \Vbar , de die Epiphanie, Oracio.} \hyperlink{deus-cuius-unigenitus-epiphanie}{Deus cuius.}
\newline \ul{Ad Vesperas. Cap., Ymnus, \Vbar , ut in primis vesperis.} {\color{Red} Ad Mag. Ant.} Fontes aquarum sanctificati sunt Christo apparente in gloria, orbis terrarum haurite aquas de fonte salvatoris, sanctificavit enim nunc omnem creaturam, Christus deus noster. {\color{Red} Oratio.} \hyperlink{deus-cuius-unigenitus-epiphanie}{Deus cuius.} \ul{Si hec octava in Sabbato evenerit, in vesperis tantum fiat memoria de Dominica sequenti.}
% ORDINARIUM
\lettrine[lines=2]{\zallmancaps \color{Red} N}{}\ul{\textsc{otandum} quod hystoria,} \hyperlink{domine-ne-in-ira}{Domine ne in ira} \ul{debet cantari in diebus dominicis usque ad Septuagesimam exclusive responsoria vero ferialia, nisi festivitates impediant omnibus diebus infra ebdomadam sunt cantanda, usque ad Septuagesimam, preterquam in Sabbatis in quibus de Beata Virgine est agendum.
Providendum tamen quod quando primum Sabbatum quod est post Octavam Epiphanie fuerit omnino vacans, ita quod non sit in eo festum trium lectionum vel maius, tunc in ipso Sabbato fiat Officium feriale et cantentur responsoria que pro Sabbato sunt signata.
Hoc autem contingit quandocumque Octava Epiphanie est in Dominica, nisi Septuagesima fuerit in proxima Dominica immediate, videlicet in festo Fabiani et Sebastiani, quia tunc festum predictorum martirum est in Sabbato celebrandum. Alias vero nunquam contingit.}
\newline \ul{Considerandum est autem diligenter utrum inter Octavas Epyphanie et Septuagesimam tot sint dies Dominici quod sunt Dominicalia Officia et orationes, ut in ipsis dicantur. Si vero non sint tot dies Dominici dicantur per ferias tali quidem die si fieri potest in quo nulla sit festivitas. Quod si necessitas coegerit ut in illa die dicatur Officium Dominicale in qua festum trium lectionum occurrerit, memoria tantum fiat de festo.
Ad Matutinas autem omelia, et antiphona ad Ben., et ad Mag., et Missa in die, de Dominica erunt.
Quando vero inter Octavam Epiphanie et Septuagesimam nulla est Dominica vacans in qua possit cantari, hystoria,} \hyperlink{domine-ne-in-ira}{Domine ne in ira} \ul{tunc responsoria de Psalmis que per ferias notata sunt vel eciam Officium de festis trium lectionum, si alie ferie non vacaverint sunt intermittenda, et hystoria,} \hyperlink{domine-ne-in-ira}{Domine ne in ira} \ul{per easdem ferias cantetur.
Propter orationes autem et Evangelia que post duo officia posita sunt, non debet intermitti festum trium lectionum, nec in Sabbato Officium de Beata Virgine.}
\newline \ul{Sciendum igitur quod si Septuagesima evenerit decimo quinto kalendas Februarii, scilicet in festo Prisce, tunc in quarta feria precedenti, que est immediate post Octavas Epyphanie, hoc est in die Felicis, legendum est omelia de nuptiis et cantanda sunt tria ultima responsoria de hystoria Dominicali, et sequenti die scilicet in festo Mauri quod erit quinta feria legende sunt tres prime lectiones de Paulo ad Romanos, et cantanda tria prima responsoria hystorie Dominicalis. In sexta vero feria, scilicet in die Marcelli legenda est omelia,} Cum descendisset \ul{et cantanda sunt tria media responsoria hystorie predicte. In Sabbato autem hoc est in die Antonii fiet de Beata Virgine.}
\newline \ul{Quando vero Septuagesima evenerit decimo quarto kalendas Februarii, tunc in tertia feria precedenti que erit immediate post Octavam Epiphanie, hoc est in die Felicis, legantur tres prime lectiones de Paulo, et cantentur tria media responsoria hystorie Dominicalis. In quarta autem feria sequenti hoc est in die Mauri, legatur omelia de nuptiis, et cantentur tria ultima responsoria hystorie Dominicalis. In quinta vero feria hoc est in die Marcelli, legatur omelia,} Cum descendisset, \ul{et cantentur tria prima responsoria hystorie predicte. In sexta feria fiat festum Antonii. In Sabbato vero fiat de Beata Virgine.
Quando autem Septuagesima evenerit decimo tertio kalendas Februarii, tunc in secunda feria precedenti que erit immediate post Octavam Epyphanie, legantur prime tres lectiones de epistola Pauli ad Romanos, et cantentur tria prima responsoria hystorie Dominicalis. In tertia feria fiat festum Mauri. In quarta feria hoc est in die Marcelli legatur omelia de nuptiis, et cantentur tria ultima responsoria hystorie Dominicalis. In quinta feria fiat festum Antonii. In sexta feria hoc est in die Prisce legatur omelia,} Cum descendisset \ul{et cantentur tria media responsoria hystorie predicte. In Sabbato fiet festum Fabiani et Sebastiani.}
{\ubsubsection{Dominica Prima Post Octavas Epyphanie, Sabbato Precedenti}{breviarium-dominica-prima-post-octavas-epyphanie-sabbato-precedenti}}
\lettrine[lines=2]{\zallmancaps \color{Blue} S}{\color{Red} ic} {\color{Red} inchoetur Vespere.} Deus in adiutorium meum intende. {\color{Red} Responsio.} Domine ad adiuvandum me festina. Gloria Patri. Sicut erat. Alleluya. \ul{Hoc modo inchoentur omnes hore et laudes, nisi quod ad Completorium premittendus est versiculus,} Converte nos deus salutaris noster. Et averte iram tuam a nobis, \ul{et ad Matutinas, \Vbar .} Domine labia mea aperies. Et os meum annunciabit laudem tuam, \ul{et ante Laudes dicitur versiculus. Excipitur totum triduum ante Pascha, et dies animarum et Vespere in die Pasche, et deinceps omnes Vespere usque ad Vesperas Sabbati in Albis exclusive, et quod Alleluya post Sicut Erat non dicitur in Completorio Sabbati ante Dominicam Septuagesime et deinceps usque ad Completorio Sabbati Pasche exclusive. Deinde inchoetur.} {\color{Red} Ant.} Benedictus dominus deus meus. {\color{Red} Ps.} \psalm{143} {\color{Red} Ant.} In eternum et in seculum seculi. {\color{Red} Ps.} \psalm{144} {\color{Red} Ant.} Laudabo deum meum in vita mea. \psalm{145} {\color{Red} Ant.} Deo nostro iocunda sit laudatio. {\color{Red} Ps.} \psalm{146} {\color{Red} Ant.} Lauda Hierusalem dominum. {\color{Red} Ps.} \psalm{147}%typo ira
\newline \ul{Predicte antiphone cum suis psalmis dicantur in omnibus Sabbatis per annum quando de dominica agitur nisi in tempore Paschali.} {\color{Red} Capitulum.}
\lettrine[lines=2]{\zallmancaps \color{Red} B}{enedictus} deus et pater domini nostri Ihesu Christi, pater misericordiarum et deus totius consolationis, qui consolatur nos in omni tribulatione nostra. {\color{Red} Responsio.} Deo gratias. \R \hyperlink{deus-qui-sedes}{Deus qui sedes.} {\color{Red} II. Ymnus.}
\hymn{o-lux-beata}{O lux beata trinitas}{
\lettrine[lines=2]{\zallmancaps \color{Blue} O}{}\ lux beata trinitas
et principalis unitas,
iam sol recedit igneus
infunde lumen cordibus.
\par {\bfseries \color{Red} T}e mane laudum carmine,
te deprecemur vesperi,
te nostra supplex gloria
per cuncta laudet secula.
\par {\bfseries \color{Blue} D}eo patri sit gloria
eiusque soli filio
cum spiritu paraclito
et nunc et in perpetuum. Amen.
}
\newline \V Vespertina oratio ascendat ad te domine. \R Et descendat super nos misericordia tua.
\newline \ul{Capitulum precedens, Ymnus, Versiculus, dicantur omnibus Sabbatis ad Vesperas usque ad Quadragesimam exclusive et in Sabbato post festum Trinitatis et deinceps in omnibus Sabbatis usque ad Adventum exclusive quando de Dominica agitur.}
\newline {\color{Red} Ad Mag. Ant.} Peccata mea domine sicut sagitte infixa sunt in me sed antequam vulnera generent in me, sana me domine medicamento penitentie deus.
\newline \ul{Predicta antiphone dicatur in Sabbatis usque ad Septuagesimam inclusive ad Magnificat vel ad memoriam de Dominica.}
\newline \ul{Finita antiphona post Mag. sequitur,} Dominus vobiscum. Oremus. {\color{Red} Oratio.}
\lettrine[lines=2]{\zallmancaps \color{Red} O}{mnipotens} \hypertarget{omnipotens-sempiterne-dom-i-post}{sempiterne} deus qui celestia simul et terrena moderaris supplicationes populi tui clementer exaudi, et pacem tuam nostris concede temporibus. Per dominum nostrum etc.
\newline \ul{Finita oratione et responso,} Amen, \ul{sequitur.} Dominus vobiscum, \ul{nisi memoria fierit facienda tunc enim responso,} Amen \ul{non dicatur,} Dominus vobiscum \ul{nec} Benedicamus \ul{sed immediate inchoetur antiphona ad memoriam, qua finita dicatur \Vbar , deinde absque,} Dominus vobiscum \ul{premisso,} Oremus, \ul{dicatur oracio, et si alia memoria fuerit faciendi terminetur oracio,} Per Christum, \ul{vel} Per eundem, \ul{vel} Qui vivis et regnas per omnia, etc. \ul{Secundum varietatem terminationum, si vero non fuerit alia memoria facienda ultima oracio terminetur sicut prima et ea finita subiungatur,} Dominus vobiscum, \ul{et,} Benedicamus. \ul{Hoc idem in Laudibus observetur, responso,} Deo gratias. \ul{Post} Benedicamus \ul{dicatur} Fidelium anime per misericordiam dei requiescant in pace. \ul{Quod semper dicendum est post quamlibet horam nisi in die Cene et deinceps usque ad Completorium Sabbati exclusive, et nisi post officium defunctorum. Quando vero plures, hore pariter sunt dicende, tunc solum in fine ultime hore dicatur} Fidelium \ul{nisi ut dictum est post officium defunctorum. Dicto} Fidelium \ul{vel} Requiescant in pace, \ul{dicatur,} Pater noster.
{\ubsubsection{Ad Completorium}{breviarium-ad-completorium}}
\par \noindent {\color{Red} Super psalmos Ant.} Miserere michi domine et exaudi orationem meam. {\color{Red} Ps.} \psalm{4} {\color{Red} Ps.} \psalm{30} \ul{usque ad versum,} Odisti observantes \ul{exclusive.} {\color{Red} Ps.} \psalm{90} {\color{Red} Ps.} \psalm{133}
\lettrine[lines=2]{\zallmancaps \color{Blue} T}{u} in nobis es domine et nomen sanctum tuum invocatum est super nos, ne dereliquas nos domine deus noster.
\begin{responsory-breve}[in-manus-tuas-breve]
{In manus tuas domine * Commendo spiritum meum.}{Redemisti me domine deus veritatis.}
\end{responsory-breve}%
{\color{Red} Ymnus.}
\hymn{te-lucis}{Te lucis ante terminum}{
\lettrine[lines=2]{\zallmancaps \color{Red} T}{e} lucis ante terminum
rerum creator poscimus
ut solita clementia
sis presul ac custodia.
\par {\bfseries \color{Blue} P}rocul recedant somnia
et noctium fantasmata,
hostemque nostrum comprime
ne polluantur corpora.
\par {\bfseries \color{Red} P}resta pater omnipotens
per Ihesum Christum dominum
qui tecum in perpetuum
regnat cum sancto spiritu. Amen.
}
\newline \V Custodi nos domine ut pupillam oculi. \R Sub umbra alarum tuarum protege nos.
\newline {\color{Red} Ad Nunc Dimittis. Ant.} Salva nos domine vigilantes, custodi nos dormientes, ut vigilemus cum Christo et requiescamus in pace.
\newline Kyrie eleyson. Christe eleyson. Kyrie eleyson. Pater noster. \ul{secreto. Quo dicto subiungatur alte.} Et ne nos inducas in temptationem. {\color{Red} Responsio.} Sed libera nos a malo. \V In pace in idipsum. {\color{Red} Responsio.} Dormiam et requiescam. \ul{Deinde dicatur.} \psalm{credo} \ul{secreto. Quo dicto subiungatur alte.} Carnis resurrectionem. {\color{Red} Responsio.} Vitam eternam. Amen. \ul{Deinde,} Dignare domine nocte ista. {\color{Red} Responsio.} Sine peccato nos custodire. Dominus vobiscum. {\color{Red} Oratio.} Oremus.
\lettrine[lines=2]{\zallmancaps \color{Blue} V}{isita} \hypertarget{visita-quesumus}{quesumus} domine habitationem istam, et omnes insidias inimici ab ea longe repelle, et angeli tui sancti habitantes in ea nos in pace custodiant, et benedictio tua sit super nos semper. Per. Dominus vobiscum. Benedicamus domino. \ul{Deinde detur benedictio.}
\newline \ul{Predicta Ant.} Miserere \ul{dicatur ad Completorium super psalmos per totum annum nisi in die Cene et deinceps usque ad Completorium Sabbati ante Trinitatem exclusive, et nisi in vigilia Natalis Domini et deinceps usque in crastinum octavarum Epyphanie exclusive, et nisi quando de Beata Virgine agitur in conventu.}
\newline \ul{Psalmi vero dicantur per totum annum, nisi quod in vigilia Pasche deinceps usque ad Completorium Sabbati in Albis exclusive, et in vigilia Pentecostes et deinceps usque ad Completorium Sabbati ante festum Trinitatis exclusive intermittitur psalmus,} \psalm{90}
\newline \ul{Capitulum,} Tu in nobis \ul{semper dicatur nisi in Cena Domini, et deinceps usque ad Completorium Sabbati in Albis exclusive.}
\newline \ul{Responsorium autem,} \hyperlink{in-manus-tuas-breve}{In manus,} \ul{per totum annum dicatur, nisi in Dominicis et in festis sanctorum ad utrunque Completorium in Quadragesima usque ad Completorium Sabbati ante Dominicam in Passione exclusive, et nisi in Cena Domini et deinceps usque ad Completorium Sabbati in Albis exclusive.}
\newline \ul{Ymnus,} \hyperlink{te-lucis}{Te lucis,} \ul{semper dicatur nisi in Sabbato ante prima Dominicam Quadragesime, et deinceps usque ad Sabbatum ante festum Trinitatis exclusive.}
\newline \ul{Versiculus,} Custodi nos, \ul{dicatur semper nisi in Cena Domini et deinceps usque ad Completorium Sabbati in Albis exclusive.}
\newline \ul{Canticum,} \psalm{nunc} \ul{semper dicatur.}
\newline \ul{Antiphona,} Salva nos, \ul{per totum annum dicatur, nisi in Sabbato ante primam Dominicam Quadragesime, et deinceps usque ad Sabbatum ante Trinitatem exclusive, et nisi in vigilia Natalis Domini, et deinceps usque in crastinum octavarum Epyphanie exclusive, et nisi quando de Beata Virgine agitur in conventu.}
\newline \ul{Preces vero cotidie dicantur nisi in Cena Domini et deinceps usque ad secundam feriam post Dominicam in Albis exclusive, et nisi in ebdomada Pentecostes, et nisi infra Octavam Nativitatis Domini, et nisi duplicibus et in totis duplicibus.}
\newline \ul{Versiculus,} Dignare domine \ul{ante} Dominus vobiscum, \ul{non dicatur quando preces non dicuntur, similiter nec in Prima.}
\newline \ul{Oracio,} \hyperlink{visita-quesumus}{Visita,} \ul{semper dicatur, nisi in Cena Domini et deinceps usque ad terciam feriam post Pascha inclusive.}
\newline \ul{Benedictio semper detur nisi in die Cene et Parascheves.}
\newline {\color{Red} Ad Matutinas. Invit.}
\begin{invitatory}[venite-exultemus-domino-invitatorium]
{Venite exultemus domino, * Iubilemus deo salutari nostro.}
\end{invitatory}
\ul{Quo dicto sequitur versus,} Preoccupemus. \ul{Predictum, Invit., dicatur Dominicis diebus usque ad Septuagesimam exclusive quando de Dominica agitur. Quandocumque autem Invit., dicitur, dicatur eciam psalmus,} \psalm{94} {\color{Red} Ymnus.}
\hymn{nocte-surgentes}{Nocte surgentes}{
\lettrine[lines=2]{\zallmancaps \color{Red} N}{octe} surgentes vigilemus omnes
semper in psalmis meditemur, atque
viribus totis Domino canamus
dulciter ymnos.
\par {\bfseries \color{Blue} U}t pio regi pariter canentes
cum suis sanctis mereamur aulam
ingredi celi, simul et beatam
ducere vitam.
\par {\bfseries \color{Red} P}restet hoc nobis deitas beata
patris ac nati, pariterque sancti
spiritus, cuius reboat in omni
gloria mundo. Amen.
}
\newline {\color{Red} In Primo Nocturno. Ant.} Servite domino in timore et exultate ei cum tremore. {\color{Red} Ps.} \psalm{1} \psalm{2} \psalm{3} \psalm{6} Gloria. {\color{Red} Ant.} Domine deus meus in te speravi. {\color{Red} Ps.} \psalm{7} \psalm{8} \psalm{9} \psalm{10} Gloria. {\color{Red} Ant.} Tu domine servabis nos et custodies nos. {\color{Red} Ps.} \psalm{11} \psalm{12} \psalm{13} \psalm{14} Gloria. \V Memor fui nocte nominis tui domine. \R Et custodi legem tuam. \ul{Deinde,} \psalm{pater} \ul{Quo dicto subiungatur alte,} Et ne nos. \ul{Et hoc generaliter observetur ante primam lectionem et quartam et septimam nisi in triduo ante Pascha.}
\newline {\color{Red} Lectio Prima.}
\lettrine[lines=2]{\zallmancaps \color{Blue} P}{aulus} servus Ihesu Christi, vocatus apostolus, segregatus in evangelium Dei, quod ante promiserat per prophetas suos in scripturis sanctis de filio suo, qui factus est ei ex semine David secundum carnem, qui predestinatus est filius Dei in virtute secundum spiritum sanctificationis ex resurrectione mortuorum Ihesu Christi Domini nostri, per quem accepimus gratiam et apostolatum ad obediendum fidei in omnibus gentibus pro nomine eius in quibus estis et vos vocati Ihesu Christi: omnibus qui sunt Rome, dilectis Dei, vocatis sanctis, gratia vobis et pax a Deo patre nostro, et domino Ihesu Christo. Tu autem domine miserere nostri. {\color{Red} Responsio.} Deo gratias. \ul{Per hanc clausulam.} Tu autem \ul{terminentur semper lectiones nisi in triduo ante Pascha et nisi ubi aliud est signatum.}
\lettrine[lines=2]{\zallmancaps \color{Red} D}{omine} \hypertarget{domine-ne-in-ira}{\label{domine-ne-in-ira}} ne in ira tua arguas me neque in furore tuo corripias me, * Miserere michi domine quoniam infirmus sum. \V Timor et tremor venerunt super me et contexerunt me tenebre, et dixi. Miserere.
\newline {\color{Red} Lectio Secunda.}
\lettrine[lines=2]{\zallmancaps \color{Blue} P}{rimum} quidem gratias ago Deo meo per Ihesum Christum pro omnibus vobis, quia fides vestra annunciatur in universo mundo. Testis enim michi est Deus cui servio in spiritu meo in evangelio filii eius, quod sine intermissione memoriam vestri facio semper in orationibus meis, obsecrans si quomodo tandem aliquando prosperum iter habeam in voluntate Dei veniendi ad vos.
\begin{responsory}[deus-qui-sedes]
{Deus qui sedes super thronum et iudicas equitatem, esto refugium pauperum in tribulatione, * Quia tu solus laborem et dolorem consideras.}{Tibi enim derelictus est pauper pupillo tu eris adiutor.}
\end{responsory}%
\newline {\color{Red} Lectio Tertia.}
\lettrine[lines=2]{\zallmancaps \color{Red} D}{esidero} enim videre vos, ut aliquid impertiar vobis gratie spiritualis ad confirmandos vos, id est simul consolari in vobis, per eam que invicem est fidem vestram atque meam. Nolo enim vos ignorare fratres quia sepe proposui venire ad vos, et prohibitus sum usque adhuc, ut aliquem fructum habeam et in vobis sicut et in ceteris gentibus.
\begin{responsory-final}[a-dextris-est]
{A dextris est michi dominus ne commovear, * Propter hoc dilatatum est cor meum, Et exultavit lingua mea.}{Conserva me domine quoniam in te speravi, dixi domino deus meus es tu.}{Et exultavit}
\end{responsory-final}%
\newline {\color{Red} In Secundo Nocturno. Ant.} Bonorum meorum non indiges in te speravi, conserva me domine. {\color{Red} Ps.} \psalm{15} {\color{Red} Ant.} Inclina domine aurem tuam michi et exaudi verba mea. {\color{Red} Ps.} \psalm{16} {\color{Red} Ant.} Dominus firmamentum meum et refugium meum. {\color{Red} Ps.} \psalm{17} \V Media nocte surgebam ad confitendum tibi. \R Super iudicia iustificationis tue.
\newline {\color{Red} Lectio Quarta.}
\lettrine[lines=2]{\zallmancaps \color{Blue} G}{recis} ac barbaris, sapientibus et insipientibus debitor sum, ita quod in me promptum est et vobis qui Rome estis evangelizare. Non enim erubesco evangelium. Virtus enim Dei est in salutem omni credenti, Iudeo primum et Greco. Iusticia enim Dei in eo revelatur ex fide in fidem, sicut scriptum est. Iustus autem ex fide vivit.
\begin{responsory}[notas-michi-fecisti]
{Notas michi fecisti domine vias vite, adimplebis me leticia cum vultu tuo, * Delectationes in dextera tua usque in finem.}{Tu es domine qui restitues hereditatem meam michi.}
\end{responsory}%
\newline {\color{Red} Lectio Quinta.}
\lettrine[lines=2]{\zallmancaps \color{Red} R}{evelatur} enim ira Dei de celo super omnem impietatem, et iniustitiam hominum eorum qui veritatem Dei in iniustitia detinent, quia quod notum est Dei manifestum est in illis. Deus enim illis revelavit.
\begin{responsory}[diligam-te-domine]
{Diligam te domine virtus mea, * Dominus firmamentum meum et refugium meum.}{Laudans invocabo dominum, et ab inimicis meis salvus ero.}
\end{responsory}%
\newline {\color{Red} Lectio Sexta.}
\lettrine[lines=2]{\zallmancaps \color{Blue} I}{nvisibilia} enim ipsius a creatura mundi per ea que facta sunt intellecta conspiciuntur, sempiterna quoque eius virtus et divinitas, ita ut sint inexcusabiles. Quia cum cognovissent Deum non sicut Deum glorificaverunt aut gratias egerunt, sed evanuerunt in cogitationibus suis, et obscuratum est insipiens cor eorum. Dicentes enim se esse sapientes, stulti facti sunt.
\begin{responsory-final}[domini-est-terra]
{Domini est terra et plenitudo eius, * Orbis terrarum et universi * Qui habitant in eo.}{In manu eius sunt omnes fines terre et altitudines montium ipsius sunt.}{Qui habitant}
\end{responsory-final}%
\newline {\color{Red} In Tertio Nocturno. Ant.} Non sunt loquele neque sermones quorum non audiantur voces eorum. {\color{Red} Ps.} \psalm{18} {\color{Red} Ant.} Exaudiet te dominus in die tribulationis. {\color{Red} Ps.} \psalm{19} {\color{Red} Ant.} Domine in virtute tua letabitur rex. {\color{Red} Ps.} \psalm{20} \V Exaltare domine in virtute tua. \R Cantabimus et psallemus virtutes tuas.
\newline {\color{Red} Lectio VII. Secundum Iohannem.}
\lettrine[lines=2]{\zallmancaps \color{Red} I}{n} illo tempore: Nuptie facte sunt in Chana Galilee, et erat mater Ihesu ibi. Et reliqua.
\newline {\color{Red} Omelia Venerabilis Bede Presbiteri.} \grecross
\lettrine[lines=2]{\zallmancaps \color{Blue} Q}{uod} Dominus noster atque salvator ad nuptias vocatus, non solum venire sed et miraculum ibidem quo convivas
%pp095
letificaret facere dignatus est, exceptis sacramentorum celestium figuris, eciam iuxta litteram fidem recte credentium confirmat.
\begin{responsory}[ad-domine-levavi]
{Ad domine levavi animam meam * Deus meus in te confido non erubescam.}{Neque irrideant me inimici mei etenim universum qui sustinent te non confundentur.}
\end{responsory}%
\newline {\color{Red} Lectio VIII.}
\lettrine[lines=2]{\zallmancaps \color{Red} P}{orro} Taciani et Marciani ceterorumque qui nuptiis detrahunt perfidia quam sit damnabilis insinuat. Si enim thoro immaculato et nuptiis debita castitate celebratis culpa inesset, nequaquam Dominus ad has venire, nequaquam eas signorum suorum initiis consecrare voluisset.
\begin{responsory}[audiam-domine]
{Audiam domine vocem laudis tue, * Ut enarrem universa mirabilia tua.}{Domine dilexi decorem domus tue et locum habitationis glorie tue.}
\end{responsory}%
\newline {\color{Red} Lectio IX.}
\lettrine[lines=2]{\zallmancaps \color{Blue} N}{unc} autem quia bona est castitas coniugalis, melior continencia vidualis, optima perfectio virginalis: ad comprobandam omnium electionem graduum, discernendum tamen meritum singulorum: ex intemerato Marie virginis utero nasci dignatus est, a prophetico Anne vidue ore mox natus benedicitur, a nuptiarum celebratoribus iam iuvenis invitatur: et hos sue presentia virtutis honorat.
\begin{responsory-final}[abscondi-tanquam]
{Abscondi tanquam aurum peccata mea et celavi in sinu meo iniquitatem meam, * Miserere mei deus, * Secundum magnam misericordiam tuam.}{Quoniam iniquitatem meam ego agnosco et delictum meum coram me est semper, tibi soli peccavi.}{Secundum}
\end{responsory-final}%
\newline \V Excelsus super omnes gentes dominus. \R Et super celos gloria eius.
\newline {\color{Red} In Laudibus. Super psalmos Ant.} Regnavit dominus precinctus fortitudine cum decore virtutum, cuius sedes parata est in eternum. {\color{Red} Ps.} \psalm{92} \ul{Cum ceteris.} {\color{Red} Capitulum.}
\lettrine[lines=2]{\zallmancaps \color{Red} B}{enedictio} et claritas et sapientia et gratiarum actio, honor virtus et fortitudo deo nostro, in secula seculorum. Amen. {\color{Red} Ymnus.}
\hymn{ecce-iam-noctis}{Ecce iam noctis}{
\lettrine[lines=2]{\zallmancaps \color{Blue} E}{cce} iam noctis tenuatur umbra,
lucis aurora rutilans coruscat,
nisibus totis rogitemus omnes
cunctipotentem.
\par {\bfseries \color{Red} U}t noster miseratus omnem
pellat languorem, tribuat salutem,
donec et nobis pietate patris
regna polorum.
\par {\bfseries \color{Blue} P}restet hoc nobis deitas beata
patris ac nati, pariterque sancti
spiritus, cuius reboat in omni
gloria mundo. Amen.
}
\newline \V Dominus regnavit. \R Decorem induit.
\newline {\color{Red} Ad Ben. Ant.} Nuptie facte sunt in Chana Galilee, et erat ibi Ihesus cum Maria matre eius.
\newline \ul{Predicti ymni,} \hyperlink{nocte-surgentes}{Nocte surgentes} \ul{et} \hyperlink{ecce-iam-noctis}{Ecce iam noctis,} \ul{dicantur omnibus dominicis et feriis per totum annum quando de tempore agitur ad Matutinas et ad Laudes nisi in Quadragesima, et in Paschali tempore et in Adventu Domini.}
\newline \ul{Antiphone nocturnorum cum psalmis, dicantur dominicis diebus usque ad Dominicam in Passione exclusive, quando de dominica agitur.}
\newline \ul{Versiculi nocturnorum, et versiculus ante Laudes dicantur in dominicis usque ad Quadragesimam exclusive.}
\newline \ul{Antiphona super psalmos, in Laudibus, et \Vbar , dicantur in Dominicis usque ad Septuagesimam exclusive quando de Dominica agitur.}
\newline \ul{Capitulum vero,} Benedictio, \ul{dicatur omnibus dominicis quando de Dominica agitur usque ad Quadragesimam exclusive.}
\newline \ul{Postquam predicta hystoria,} \hyperlink{domine-ne-in-ira}{Domine ne in ira} \ul{in Dominica fuerit cantata, in reliquis Dominicis si alique fuerint usque ad Septuagesimam exclusive, loco noni responsorii dicatur,} \Rbar . \hyperlink{honor-virtus}{Honor} \ul{de Trinitate.}
\newline {\color{Red} Ad Primam. Ymnus.}
\hymn{iam-lucis}{Iam lucis orto sydere}{
\lettrine[lines=2]{\zallmancaps \color{Red} I}{am} lucis orto sydere
deum precemur supplices,
ut in diurnis actibus
nos servet a nocentibus.
\par {\bfseries \color{Red} L}inguam refrenans temperet
ne litis horror insonet,
visum fovendo contegat,
ne vanitates hauriat.
\par {\bfseries \color{Blue} S}int pura cordis intima,
absistat et vecordia,
carnis terat superbiam
potus cibique parcitas.
\par {\bfseries \color{Red} U}t cum dies abscesserit,
noctemque sors reduxerit,
mundi per abstinentiam
ipsi canamus gloriam.
\par {\bfseries \color{Blue} D}eo patri sit gloria
eiusque soli filio
cum spiritu paraclito
et nunc et in perpetuum. Amen.
}
\newline {\color{Red} Ant.} Gloria tibi trinitas. {\color{Red} Ps.} \psalm{53} \psalm{118.1} \psalm{118.2} \psalm{quicumque} {\color{Red} Capitulum.}
\lettrine[lines=2]{\zallmancaps \color{Blue} R}{egi} seculorum immortali invisibili soli deo honor et gloria: in secula seculorum. Amen.
\begin{responsory-breve}[ihesu-christe-fili]
{Ihesu Christe fili dei vivi * Miserere nobis.}{Qui sedes ad dexteram patris.}
\end{responsory-breve}%
\V Exurge domine adiuva nos. \R Et libera nos propter nomen tuum.
\newline Kyrie eleyson. Christe eleyson. Kyrie eleyson. Pater noster. Et ne nos. Sed libera. \V Vivet anima mea et laudabit te. \R Et iudicia tua adiuvabunt me. \V Erravi sicut ovis que periit. \R Quere servum tuum domine quia mandata tua non sum oblitus. Credo in deum. Carnis resurrectionem. Vitam eternam. \ul{Deinde} Confiteor. Misereatur. \V Dignare domine die isto. \R Sine peccato. Dominus vobiscum. Oremus. {\color{Red} Oracio.}
\lettrine[lines=2]{\zallmancaps \color{Red} I}{n} \hypertarget{in-hac-hora}{\label{in-hac-hora}} hac hora huius diei tua nos domine reple misericordia ut per totum diem exultantes in tuis laudibus iugiter delectemur. Per. Dominus vobiscum. Benedicamus domino. \ul{Et hoc modo terminentur Tercia, Sexta et Nona.}
\newline \ul{Predictus ymnus,} \hyperlink{iam-lucis}{Iam lucis} \ul{dicatur per totum annum ad Primam, nisi in Cena domini et deinceps usque ad Dominicam in Albis exclusive et nisi die animarum.}
\newline \ul{Psalmi vero semper dicantur nisi quod in Septuagesima et deinceps in Dominicis usque ad Pascha exclusive adduntur alii secundum quod in Dominica Septuagesime sunt notati.}
\newline \psalm{quicumque} \ul{dicatur semper in Dominicis, nisi festum totum duplex evenerit.}
\newline \ul{Capitulum} Regi \ul{dicatur ad Primam omnibus Dominicis et festis eciam trium lectionum, et per omnes octavas, nisi in Cena Domini et in ebdomada Pasche quando nullum capitulum dicitur, et nisi in die animarum. Dicitur eciam in Sabbatis quando de Beata Virgine agitur in conventu.}
\newline \ul{Responsorium,} \hyperlink{ihesu-christe-fili}{Ihesu Christe} \ul{semper dicatur, nisi in Dominica in Passione et deinceps usque ad Dominican in Albis exclusive: quando de tempore agitur.}
\newline \ul{\Vbar .} Qui sedes, \ul{semper dicatur nisi in Dominica in Passione et deinceps usque ad diem Pentecostes exclusive, et nisi in Adventu Domini et deinceps usque in crastinum oct. Epiphanie exclusive, et nisi quando de Beata Virgine agitur in conventu et non in die animarum.}
\newline \ul{Versiculus,} Exurge \ul{semper dicatur nisi in die Cene et deinceps usque ad Dominicam in Albis exclusive, et nisi in die animarum.}
\newline \ul{Preces eciam cotidie dicende sunt nisi in Cena Domini et deinceps usque ad secundam feriam post Dominicam in Albis exclusive, et nisi in ebdomada Pentecostes et nisi infra Octavam Natalis Domini, et in duplicibus et totis duplicibus, et nisi in die animarum.}
\newline Confiteor \ul{in Prima semper dicatur, nisi in triduo ante Pascha.}
\newline \ul{Oratio} In hac hora \ul{semper dicatur quando dicitur Capitulum.} Regi seculorum. \ul{Dicitur eciam in quatuor diebus ultimis in ebdomada Pasche.}
\newline {\color{Red} Ad Terciam. Ymnus.}
\hymn{nunc-sancte}{Nunc sancte nobis spiritus}{
\lettrine[lines=2]{\zallmancaps \color{Blue} N}{unc} sancte nobis spiritus
unum patri cum filio,
dignare promptus ingeri,
nostro refusus pectori.
\par {\bfseries \color{Red} O}s, lingua, mens, sensus, vigor,
confessionem personent,
flammascat igne caritas
accendat ardor proximos.
\par {\bfseries \color{Blue} P}resta pater piissime,
patrique compar unice,
cum spiritu paraclito
regnans per omne seculum. Amen.
}
\newline {\color{Red} Ant.} Laus et perennis. {\color{Red} Ps.} \psalm{118.3} \psalm{118.4} \psalm{118.5} {\color{Red} Cap.}
\lettrine[lines=2]{\zallmancaps \color{Red} D}{eus} caritas est, et qui manet in caritate in deo manet et deus in eo.
\begin{responsory-breve}[inclina-cor-meum]
{Inclina cor meum deus * In testimonia tua.}{Averte oculos meos ne videant vanitatem, in via tua vivifica me.}
\end{responsory-breve}%
\V Adiutor meus esto domine ne derelinquas me. \R Neque despicias me deus salutaris meus.
\newline {\color{Red} Ad Sextam. Ymnus.}
\hymn{rector-potens}{Rector potens verax deus}{
\lettrine[lines=2]{\zallmancaps \color{Blue} R}{ector} potens verax deus
qui temperas rerum vices,
splendore mane instruis,
et ignibus meridiem.
\par {\bfseries \color{Red} E}xtingue flammas litium,
aufer calorem noxium,
confer salutem corporum,
veramque pacem cordium.
\par {\bfseries \color{Blue} P}resta pater piissime,
patrique compar unice,
cum spiritu paraclito
regnans per omne seculum. Amen.
}
\newline {\color{Red} Ant.} Gloria laudis. {\color{Red} Ps.} \psalm{118.6} \psalm{118.7} \psalm{118.8} {\color{Red} Capitulum.}
\lettrine[lines=2]{\zallmancaps \color{Red} G}{ratia} domini nostri Ihesu Christi et caritas dei et communicatio sancti spiritus, sit semper cum omnibus nobis. Amen.
\begin{responsory-breve}[in-eternum-domine]
{In eternum domine, * Permanet verbum tuum.}{In celo et in seculum seculi veritas tua.}
\end{responsory-breve}%
\V Dominus regit me et nichil michi deerit. \R In loco paschue ibi me collocavit.
\newline {\color{Red} Ad Nonam. Ymnus.}
\hymn{rerum-deus-tenax}{Rerum deus tenax vigor}{
\lettrine[lines=2]{\zallmancaps \color{Blue} R}{erum} deus tenax vigor,
immotus in te permanens,
lucis diurne tempora
successibus determinans.
\par {\bfseries \color{Red} L}argire clarum vespere,
quo vita nusquam decidat,
sed premium mortis sacre
perennis instet gloria.
\par {\bfseries \color{Blue} P}resta pater piissime,
patrique compar unice,
cum spiritu paraclito
regnans per omne seculum. Amen.
}
\newline {\color{Red} Ant.} Ex quo omnia. {\color{Red} Ps.} \psalm{118.9} \psalm{118.10} \psalm{118.11} {\color{Red} Capitulum.}
\lettrine[lines=2]{\zallmancaps \color{Red} T}{res} sunt qui testimonium dant in celo, Pater, Verbum et Spiritus Sanctus, et hi tres unum sunt.
\begin{responsory-breve}[clamavi-in-toto]
{Clamavi in toto corde, * Exaudi me domine.}{Iustificationes tuas requiram.}
\end{responsory-breve}%
\V Ab occultis meis munda me domine. \R Et ab alienis parce servo tuo.
\newline \ul{Predicti tres ymni ad Terciam, Sextam et Nonam semper dicantur, nisi in Cena Domini et deinceps usque ad Dominicam in Albis exclusive, et nisi quod in ebdomada Pentecostes ad Terciam loco de,} \hyperlink{nunc-sancte}{Nunc sancte} \ul{dicitur} \hyperlink{veni-creator-spiritus}{Veni creator.} \ul{Et preterquam in die animarum.}
\newline \ul{Psalmi vero semper dicantur.}
\newline \ul{Antiphona,} Gloria tibi trinitas \ul{ad Primam, et tres supra posite ad Terciam, Sextam et Nonam, dicantur omnibus Dominicis usque ad Septuagesimam exclusive, quando de Dominica agitur.}
\newline \ul{Tria vero Capitula precedentia ad Terciam, Sextam et Nonam dicantur in Dominicis usque ad Primam Dominicam Quadragesime exclusive, quando de Dominica agitur.}
\newline \ul{Responsoria autem horarum dicantur in Dominicis usque ad Septuagesimam exclusive, quando de Dominica agitur.}
\newline \ul{Versiculi autem dicantur in Dominicis et feriis quando de tempore agitur, usque ad Primam Dominicam in Quadragesima exclusive.}
\newline {\color{Red} Ad Vesperas. Ant.} Dixit dominus domino meo sede a dextris meis. {\color{Red} Ps.} \psalm{109} {\color{Red} Ant.} Fidelia omnia mandata eius confirmata in seculum seculi. {\color{Red} Ps.} \psalm{110} {\color{Red} Ant.} In mandatis eius volet nimis. {\color{Red} Ps.} \psalm{111} {\color{Red} Ant.} Sit nomen domini benedictum in secula. {\color{Red} Ps.} \psalm{112} {\color{Red} Ant.} Nos qui vivimus benedicimus domino. {\color{Red} Ps.} \psalm{113} {\color{Red} Capitulum.}
\lettrine[lines=2]{\zallmancaps \color{Red} D}{ominus} dirigat corda et corpora nostra in caritate dei et paciencia Christi. {\color{Red} Ymnus.}
\hymn{lucis-creator}{Lucis creator optime}{
\lettrine[lines=2]{\zallmancaps \color{Blue} L}{ucis} creator optime
lucem dierum proferens,
primordiis lucis nove
mundi parans originem.
\par {\bfseries \color{Red} Q}ui mane iunctum vesperi
diem vocari precipis,
tetrum chaos illabitur
audi preces cum fletibus.
\par {\bfseries \color{Blue} N}e mens gravata crimine,
vite sit exul munere,
dum nil perenne cogitat
seseque culpis illigat.
\par {\bfseries \color{Red} C}elorum pulset intimum,
vitale tollat premium,
vitemus omne noxium
purgemus omne pessimum.
\par {\bfseries \color{Blue} P}resta pater piissime,
patrique compar unice,
cum spiritu paraclito
regnans per omne seculum. Amen.
}
\newline \V Dirigatur domine oratio mea. \R Sicut incensum in conspectu tuo.
\newline {\color{Red} Ad Mag. Ant.} Deficiente vino iussit Ihesus impleri ydrias aqua, que in vinum conversa est, alleluya.
\newline \ul{Predicti psalmi dicantur omnibus Dominicis, in vesperis quando de simplici Dominica agitur.}
\newline \ul{Antiphone vero super psalmos dicantur in Dominicis usque ad Dominicam in Ramis inclusive quando de Dominica agitur.}
\newline \ul{Capitulum} Dominus dirigat, \ul{dicatur omnibus Dominicis in Secundis Vesperis et per ferias quando de tempore agitur usque ad diem Cinerum exclusive.}
\newline \ul{Ymnus} \hyperlink{lucis-creator}{Lucis creator} \ul{et Versiculus} Dirigatur, \ul{dicantur omnibus Dominicis in Secundis Vesperis, et per ferias quando de tempore agitur, usque ad Vesperas Sabbati ante Primam Dominicam Quadragesime exclusive.}
{\ubsubsection{Feria Secunda}{breviarium-feria-secunda}}
\par \noindent {\color{Red} Ad Matutinas. Invit.}
\begin{invitatory}[venite-exultemus-invitatorium]
{Venite, * Exultemus domino.}
\end{invitatory}
\newline \ul{Quo dicto sequitur,} Iubilemus etc.
\newline {\color{Red} In Nocturno. Ant.} Dominus defensor vite mee. {\color{Red} Ps.} \psalm{26} {\color{Red} Ps.} \psalm{27} Gloria. {\color{Red} Ant.} Adorate dominum in aula sancta eius. {\color{Red} Ps.} \psalm{28} {\color{Red} Ps.} \psalm{29} Gloria. {\color{Red} Ant.} In tua iusticia libera me domine. {\color{Red} Ps.} \psalm{30} {\color{Red} Ps.} \psalm{31} Gloria. {\color{Red} Ant.} Rectos decet collaudatio. {\color{Red} Ps.} \psalm{32} {\color{Red} Ps.} \psalm{33} Gloria. {\color{Red} Ant.} Expugna impugnantes me. {\color{Red} Ps.} \psalm{34} {\color{Red} Ps.} \psalm{35} Gloria. {\color{Red} Ant.} Revela domino viam tuam. {\color{Red} Ps.} \psalm{36} {\color{Red} Ps.} \psalm{37} Gloria. \V Domine in celo misericordia tua. \R Et veritas tua usque ad nubes.
\newline \ul{Lectiones feriales a secundo capitulo usque ad Epistolam ad Corinthios Primam, ad libitum.} \R
\lettrine[lines=2]{\zallmancaps \color{Red} Q}{uam} \hypertarget{quam-magna}{\label{quam-magna}} magna multitudo dulcedinis tue domine, * Quam abscondisti timentibus te. \V Perfecisti eis qui sperant in te in conspectu filiorum hominum. Quam ab.
\begin{responsory}[benedicam-dominum]
{Benedicam dominum in omni tempore, * Semper laus eius in ore meo.}{In domino laudabitur anima mea, audiant mansueti et letentur.}
\end{responsory}%
\begin{responsory-doxology}[delectare-in-domino]
{Delectare in domino, * Et dabit tibi petitiones cordis tui.}{Spera in domino et fac bonitatem, et inhabita terram et pasceris in divitiis eius.}
\end{responsory-doxology}%
\newline \V Fiat misericordia tua domine super nos. \R Quemadmodum speravimus in te.
\newline {\color{Red} In Laudibus. Ant.} Miserere mei deus. {\color{Red} Ps.} \psalm{50} {\color{Red} Ant.} Intellige clamorem meum domine. {\color{Red} Ps.} \psalm{5} {\color{Red} Ant.} Deus deus meus ad te de luce vigilo. {\color{Red} Ps.} \psalm{62} {\color{Red} Ant.} Conversus est furor tuus domine, et consolatus es me. {\color{Red} Can.} \psalm{isaias} {\color{Red} Ant.} Laudate dominum de celis. {\color{Red} Ps.} \psalm{148} {\color{Red} Capit.}
\lettrine[lines=2]{\zallmancaps \color{Blue} N}{ox} precessit dies autem appropinquavit, abiciamus ergo opera tenebrarum et induamur arma lucis, sicut in die honeste ambulemus.
\newline \V In matutinis domine meditabor in te. \R Quia fuisti adiutor meus.
\newline {\color{Red} Ad Ben. Ant.} Benedictus deus Israhel. \ul{Dicta antiphone subiungatur.} Kyrie eleyson. Christe eleyson. Kyrie eleyson. Pater noster. Et ne nos. Dominus vobiscum. \ul{Deinde Oracio.}
\newline \ul{Hoc modo dicantur preces ad Terciam, Sextam, Nonam, et Vesperas quando fiunt prostrationes, nisi in Parascheve et Sabbato Pasche.}
\newline {\color{Red} Ad Primam. Ant.} Deus exaudi orationem meam auribus percipe verba oris mei. {\color{Red} Capitulum.}
\lettrine[lines=2]{\zallmancaps \color{Red} P}{acem} et veritatem diligite, ait dominus omnipotens. {\color{Red} Oratio.}
\lettrine[lines=2]{\zallmancaps \color{Blue} D}{omine} deus omnipotens qui nos ad principium huius diei pervenire fecisti, tua nos salva virtute, ut in hac die ad nullum declinemus peccatum, sed semper ad tuam iusticiam faciendam nostra procedant eloquia, dirigantur cogitationes et opera. Per. \ul{Cetera ut supra in Prima de Dominica est notatum. Capitulum} Pacem \ul{et Oracio,} Domine deus, \ul{dicantur omnibus diebus profestis ad Primam, nisi in Parascheve et in Sabbato Sancto.}
\newline {\color{Red} Ad Terciam. Ant.} Deduc me domine in semita mandatorum tuorum. {\color{Red} Capitulum.}
\lettrine[lines=2]{\zallmancaps \color{Red} S}{ana} me domine et sanabor, salvum me fac et salvus ero, quoniam laus mea tu es.
\begin{responsory-breve}[sana-animam-meam-breve]
{Sana animam meam, * Quia peccavi tibi.}{Ego dixi domine miserere mei.}
\end{responsory-breve}%
\V Adiutor.
\newline {\color{Red} Ad Sextam. Ant.} Adiuva me et salvus ero domine. {\color{Red} Capitulum.}
\lettrine[lines=2]{\zallmancaps \color{Blue} A}{lter} alterius onera portate, et sic adimplebitis legem Christi.
\begin{responsory-breve}[benedicam-dominum-breve]
{Benedicam dominum, * In omni tempore.}{Semper laus eius in ore meo.}
\end{responsory-breve}%
\V Dominus regit.
\newline {\color{Red} Ad Nonam. Ant.} Aspice in me domine et miserere mei. {\color{Red} Capitulum.}
\lettrine[lines=2]{\zallmancaps \color{Red} E}{mpti} estis precio magno, glorificate et portate deum in corpore vestro.
\begin{responsory-breve}[redime-me-domine]
{Redime me domine, * Et miserere mei.}{Pes enim meus stetit in via recta in ecclesiis benedicam te domine.}
\end{responsory-breve}%
\V Ab occultis meis.
\newline {\color{Red} Ad Vesperas. Ant.} Inclinavit dominus aurem suam michi. {\color{Red} Ps.} \psalm{114} {\color{Red} Ant.} Credidi propter quod locutus sum. {\color{Red} Ps.} \psalm{115} {\color{Red} Ant.} Laudate dominum omnes gentes. {\color{Red} Ps.} \psalm{116} {\color{Red} Ant.} Clamavi et exaudivit me. {\color{Red} Ps.} \psalm{119} {\color{Red} Ant.} Auxilium meum a domino. {\color{Red} Ps.} \psalm{120}
\newline {\color{Red} Ad Mag. Ant.} Magnificat anima mea dominum.
{\ubsubsection{Feria Tercia}{breviarium-feria-tercia}}
\par \noindent {\color{Red} Ad Matutinas. Invit.}
\begin{invitatory}[iubilemus-deo-invitatorium]
{Iubilemus deo * Salutari nostro.}
\end{invitatory}
\newline {\color{Red} In Nocturno. Ant.} Ut non delinquam in lingua mea. {\color{Red} Ps.} \psalm{38} {\color{Red} Ps.} \psalm{39} Gloria. {\color{Red} Ant.} Sana domine animam meam quia peccavi tibi. {\color{Red} Ps.} \psalm{40} {\color{Red} Ps.} \psalm{41} Gloria. {\color{Red} Ant.} Eructavit cor meum verbum bonum. {\color{Red} Ps.} \psalm{43} {\color{Red} Ps.} \psalm{44} Gloria. {\color{Red} Ant.} Adiutor in tribulationibus. {\color{Red} Ps.} \psalm{45} {\color{Red} Ps.} \psalm{46} Gloria. {\color{Red} Ant.} Auribus percipite qui habitatis orbem. {\color{Red} Ps.} \psalm{47} {\color{Red} Ps.} \psalm{48} Gloria. {\color{Red} Ant.} Deus deorum dominus locutus est. {\color{Red} Ps.} \psalm{49} {\color{Red} Ps.} \psalm{51} Gloria. \V Immola deo sacrificium laudis. \R Et redde altissimo vota tua. {\color{Red} Responsorium.}
\lettrine[lines=2]{\zallmancaps \color{Red} A}{uribus} \hypertarget{auribus-percipe}{\label{auribus-percipe}} percipe domine lacrimas meas ne sileas a me, remitte michi, * Quoniam incola ego sum apud te et peregrinus. \V Dixi custodiam vias meas ut non delinquam in lingua mea. Quoniam.
\begin{responsory}[statuit-dominus]
{Statuit dominus supra petram pedes meos et direxit gressus meos deus meus, * Et immisit in os meum canticum novum.}{Expectans expectavi dominum et respexit me.}
\end{responsory}%
\begin{responsory-doxology}[ego-dixi-domine]
{Ego dixi domine miserere mei, * Sana animam meam quia peccavi tibi.}{Domine ne in ira tua arguas me, neque in furore tuo corripias me.}
\end{responsory-doxology}%
\newline {\color{Red} In Laudibus. Ant.} Secundum magnam misericordiam tuam. {\color{Red} Ps.} \psalm{50} {\color{Red} Ant.} Salutare vultus mei deus meus. {\color{Red} Ps.} \psalm{42} {\color{Red} Ant.} Ad te de luce vigilo deus. {\color{Red} Ps.} \psalm{62} {\color{Red} Ant.} Cunctis diebus vite nostre salvos nos fac domine. {\color{Red} Can.} \psalm{ezechias} {\color{Red} Ant.} Omnes angeli eius laudate dominum de celis. {\color{Red} Ps.} \psalm{148}
\newline {\color{Red} Ad Ben. Ant.} Erexit dominus nobis cornu salutis in domo David pueri sui.
\newline {\color{Red} Ad Vesperas. Ant.} In domum domini letantes ibimus. {\color{Red} Ps.} \psalm{121} {\color{Red} Ant.} Qui habitas in celis miserere nobis. {\color{Red} Ps.} \psalm{122} {\color{Red} Ant.} Adiutorium nostrum in nomine domini. {\color{Red} Ps.} \psalm{123} {\color{Red} Ant.} Bene fac domine bonis et rectis corde. {\color{Red} Ps.} \psalm{124} {\color{Red} Ant.} Facti sumus sicut consolati. {\color{Red} Ps.} \psalm{125}
\newline {\color{Red} Ad Mag. Ant.} Exultavit spiritus meus in deo salutari meo.
{\ubsubsection{Feria Quarta}{breviarium-feria-quarta}}
\par \noindent {\color{Red} Ad Matutinas. Invit.}
\begin{invitatory}[in-manu-tua-invitatorium]
{In manu tua domine, * Omnes fines terre.}
\end{invitatory}
\newline {\color{Red} In Nocturno. Ant.} Avertet dominus captivitatem plebis sue. {\color{Red} Ps.} \psalm{52} {\color{Red} Ps.} \psalm{54} Gloria. {\color{Red} Ant.} Quoniam in te confidit anima mea. {\color{Red} Ps.} \psalm{55} {\color{Red} Ps.} \psalm{56} Gloria. {\color{Red} Ant.} Iuste iudicate filii hominum. {\color{Red} Ps.} \psalm{57} {\color{Red} Ps.} \psalm{58} Gloria. {\color{Red} Ant.} Da nobis domine auxilium de tribulatione. {\color{Red} Ps.} \psalm{59} {\color{Red} Ps.} \psalm{60} Gloria. {\color{Red} Ant.} A timore inimici eripe domine animam meam. {\color{Red} Ps.} \psalm{61} {\color{Red} Ps.} \psalm{63} Gloria. {\color{Red} Ant.} In ecclesiis benedicite domino. {\color{Red} Ps.} \psalm{65} {\color{Red} Ps.} \psalm{67} Gloria. \V Deus vitam meam annunciavi tibi. \R Posuisti lacrimas meas in conspectu meo. \R
\lettrine[lines=2]{\zallmancaps \color{Red} N}{e} \hypertarget{ne-perdideris}{\label{ne-perdideris}} perdideris me domine cum iniquitatibus meis, * Neque in finem iratus reserves mala mea. \V Miserere mei deus miserere mei, quoniam in te confidit anima mea. Neque.
\begin{responsory}[paratum-cor-meum]
{Paratum cor meum deus paratum cor meum, * Cantabo et psalmum dicam domino.}{Exurge gloria mea exurge psalterium et cythara exurgam diluculo.}
\end{responsory}%
\begin{responsory-doxology}[adiutor-meus-tibi]
{Adiutor meus tibi psallam quia deus susceptor meus es, * Deus meus misericordia mea.}{Eripe me de inimicis meus deus meus et ab insurgentibus in me libera me.}
\end{responsory-doxology}%
\newline {\color{Red} In Laudibus. Ant.} Amplius lava me domine ab iniusticia mea. {\color{Red} Ps.} \psalm{50} {\color{Red} Ant.} Te decet ymnus deus in Syon. {\color{Red} Ps.} \psalm{64} {\color{Red} Ant.} Labia mea laudabunt te in vita mea deus meus. {\color{Red} Ps.} \psalm{62} {\color{Red} Ant.} Dominus iudicabit fines terre. {\color{Red} Can.} \psalm{annas} {\color{Red} Ant.} Celi celorum laudate deum. {\color{Red} Ps.} \psalm{148}
\newline {\color{Red} Ad Ben. Ant.} Salutem ex inimicis nostris et de manu omnium qui nos oderunt libera nos domine.
\newline {\color{Red} Ad Vesperas. Ant.} Beatus vir qui implevit desiderium suum. {\color{Red} Ps.} \psalm{126} {\color{Red} Ant.} Beati omnes qui timent dominum. {\color{Red} Ps.} \psalm{127} {\color{Red} Ant.} Benediximus vobis in nomine domini. {\color{Red} Ps.} \psalm{128} {\color{Red} Ant.} De profundis clamavi ad te domine. {\color{Red} Ps.} \psalm{129} {\color{Red} Ant.} Speret Israhel in domino. {\color{Red} Ps.} \psalm{130}
\newline {\color{Red} Ad Mag. Ant.} Respexisti humilitatem meam domine deus meus.
{\ubsubsection{Feria Quinta}{breviarium-feria-quinta}}
\par \noindent {\color{Red} Ad Matutinas. Invit.}
\begin{invitatory}[adoremus-dominum-invitatorium]
{Adoremus dominum, * Quoniam ipse fecit nos.}
\end{invitatory}
\newline {\color{Red} In Nocturno. Ant.} Domine deus in adiutorium meum intende. {\color{Red} Ps.}
%pp096
\psalm{68} {\color{Red} Ps.} \psalm{69} Gloria. {\color{Red} Ant.} Esto michi domine in deum protectorem. {\color{Red} Ps.} \psalm{70} {\color{Red} Ps.} \psalm{71} Gloria. {\color{Red} Ant.} Liberasti virgam hereditatis tue. {\color{Red} Ps.} \psalm{72} {\color{Red} Ps.} \psalm{73} Gloria. {\color{Red} Ant.} In Israhel magnum nomen eius. {\color{Red} Ps.} \psalm{74} {\color{Red} Ps.} \psalm{75} Gloria. {\color{Red} Ant.} Tu es deus qui facis mirabilia. {\color{Red} Ps.} \psalm{76} {\color{Red} Ps.} \psalm{77} Gloria. {\color{Red} Ant.} Propitius esto peccatis nostris domine. {\color{Red} Ps.} \psalm{78} {\color{Red} Ps.} \psalm{79} Gloria. \V Gaudebunt labia mea, cum cantavero tibi. \R Et anima mea quam redemisti. {\color{Red} Responsorium.}
\lettrine[lines=2]{\zallmancaps \color{Red} D}{eus} \hypertarget{deus-in-te}{\label{deus-in-te}} in te speravi domine non confundar in eternum in tua iusticia, * Libera me et eripe me. \V Esto michi in deum protectorem et in locum munitum ut salvum me facias. Libera.
\begin{responsory}[repleatur-os-meum]
{Repleatur os meum laude, ut ymnum dicam glorie tue, tota die magnificentie tue, noli me proicere in tempore senectutis, cum defecerit virtus mea, * Deus ne derelinquas me.}{Deus ne elongeris a me, deus meus in auxilium meum respice.}
\end{responsory}%
\begin{responsory-doxology}[gaudebunt-labia-mea]
{Gaudebunt labia mea dum cantavero tibi, * Et anima mea quam redemisti domine.}{Sed et lingua mea meditabitur iusticiam tuam, tota die laudem tuam.}
\end{responsory-doxology}%
\newline {\color{Red} In Laudibus. Ant.} Tibi soli peccavi domine miserere mei. {\color{Red} Ps.} \psalm{50} {\color{Red} Ant.} Domine refugium factus es nobis. {\color{Red} Ps.} \psalm{89} {\color{Red} Ant.} In matutinis domine meditabor in te. {\color{Red} Ps.} \psalm{62} {\color{Red} Ant.} In eternum dominus regnabit et ultra. {\color{Red} Can.} \psalm{moysis} {\color{Red} Ant.} In sanctis eius laudate deum. {\color{Red} Ps.} \psalm{148}
\newline {\color{Red} Ad Ben. Ant.} In sanctitate serviamus domino et liberabit nos ab inimicis nostris.
\newline {\color{Red} Ad Vesperas. Ant.} Et omnis mansuetudinis eius. {\color{Red} Ps.} \psalm{131} {\color{Red} Ant.} Ecce quam bonum et quam iocundum. {\color{Red} Ps.} \psalm{132} {\color{Red} Ant.} Omnia quecumque voluit dominus fecit. {\color{Red} Ps.} \psalm{134} {\color{Red} Ant.} Quoniam in eternum misericordia eius. {\color{Red} Ps.} \psalm{135} {\color{Red} Ant.} Ymnum cantate nobis de canticis Syon. {\color{Red} Ps.} \psalm{136}
\newline {\color{Red} Ad Mag. Ant.} Deposuit potentes sanctos persequentes et exaltavit humiles Christum confitentes.
{\ubsubsection{Feria Sexta}{breviarium-feria-sexta}}
\par \noindent {\color{Red} Ad Matutinas. Invit.}
\begin{invitatory}[dominum-qui-fecit-invitatorium]
{Dominum qui fecit nos, * Venite adoremus.}
\end{invitatory}
\newline {\color{Red} In Nocturno. Ant.} Exultate deo adiutori nostro. {\color{Red} Ps.} \psalm{80} {\color{Red} Ps.} \psalm{81} Gloria. {\color{Red} Ant.} Tu solus altissimus super omnem terram. {\color{Red} Ps.} \psalm{82} {\color{Red} Ps.} \psalm{83} Gloria. {\color{Red} Ant.} Benedixisti domine terram tuam. {\color{Red} Ps.} \psalm{84} {\color{Red} Ps.} \psalm{85} Gloria. {\color{Red} Ant.} Fundamenta eius in montibus sanctis. {\color{Red} Ps.} \psalm{86} {\color{Red} Ps.} \psalm{87} Gloria. {\color{Red} Ant.} Benedictus dominus in eternum. {\color{Red} Ps.} \psalm{88} {\color{Red} Ps.} \psalm{93} Gloria. {\color{Red} Ant.} Cantate domino et benedicite nomini eius. {\color{Red} Ps.} \psalm{95} {\color{Red} Ps.} \psalm{96} Gloria. \V Beatus homo quem tu erudieris domine. \R Et de lege tua docueris eum. {\color{Red} Responsorium.}
\lettrine[lines=2]{\zallmancaps \color{Red} C}{onfitebor} \hypertarget{confitebor-tibi}{\label{confitebor-tibi}} tibi domine deus in toto corde meo et honorificabo nomen tuum in eternum, * Quia misericordia tua domine magna est super me. \V Et eripuisti animam meam ex inferno inferiori. Quia.
\begin{responsory}[misericordia-tua-domine]
{Misericordia tua domine magna est super me, * Et liberasti animam meam ex inferno inferiori.}{Deus iniqui insurrexerunt in me, et fortes quesierunt animam meam.}
\end{responsory}%
\begin{responsory-doxology}[factus-est-michi]
{Factus est michi dominus in refugium et deus meus, * In auxilium spei mee.}{Deus ultionum dominus, deus ultionum libere egit.}
\end{responsory-doxology}%
\newline {\color{Red} In Laudibus. Ant.} Spiritu principali confirma cor meum deus. {\color{Red} Ps.} \psalm{50} {\color{Red} Ant.} In veritate tua exaudi me domine. {\color{Red} Ps.} \psalm{142} {\color{Red} Ant.} Illumina domine vultum tuum super nos. {\color{Red} Ps.} \psalm{62} {\color{Red} Ant.} Domine audivi auditum tuum et timui. {\color{Red} Can.} \psalm{habacuc} {\color{Red} Ant.} In tympano et choro, in cordis et organo laudate deum. {\color{Red} Ps.} \psalm{148}
\newline {\color{Red} Ad Ben. Ant.} Per viscera misericordie dei nostri in quibus visitavit nos oriens ex alto.
\newline {\color{Red} Ad Vesperas. Ant.} In conspectu angelorum psallam tibi deus meus. {\color{Red} Ps.} \psalm{137} {\color{Red} Ant.} Domine probasti me et cognovisti me. {\color{Red} Ps.} \psalm{138} {\color{Red} Ant.} A viro iniquo libera me domine. {\color{Red} Ps.} \psalm{139} {\color{Red} Ant.} Domine clamavi ad te exaudi me. {\color{Red} Ps.} \psalm{140} {\color{Red} Ant.} Porcio mea domine sit in terra vivencium. {\color{Red} Ps.} \psalm{141}
\newline {\color{Red} Ad Mag. Ant.} Suscepit deus Israhel puerum suum sicut locutus est Abraham et semini eius et exaltavit humiles usque in seculum.
{\ubsubsection{Sabbato}{breviarium-sabbato}}
\par \noindent {\color{Red} Ad Matutinas. Invit.}
\begin{invitatory}[dominum-deum-nostrorum-invitatorium]
{Dominum deum nostrorum, * Venite adoremus.}
\end{invitatory}
\newline {\color{Red} In Nocturno. Ant.} Quia mirabilia fecit dominus. {\color{Red} Ps.} \psalm{97} {\color{Red} Ps.} \psalm{98} Gloria. {\color{Red} Ant.} Iubilate deo omnis terra. {\color{Red} Ps.} \psalm{99} {\color{Red} Ps.} \psalm{100} Gloria. {\color{Red} Ant.} Clamor meus ad te veniat deus. {\color{Red} Ps.} \psalm{101} {\color{Red} Ps.} \psalm{102} Gloria. {\color{Red} Ant.} Benedic anima mea domino. {\color{Red} Ps.} \psalm{103} {\color{Red} Ps.} \psalm{104} Gloria. {\color{Red} Ant.} Visita nos domine in salutari tuo. {\color{Red} Ps.} \psalm{105} {\color{Red} Ps.} \psalm{106} Gloria. {\color{Red} Ant.} Confitebor domino nimis in ore meo. {\color{Red} Ps.} \psalm{107} {\color{Red} Ps.} \psalm{108} Gloria. \V Domine exaudi orationem meam. \R Et clamor meus ad te veniat. \R
\lettrine[lines=2]{\zallmancaps \color{Red} M}{isericordiam} \hypertarget{misericordiam-et-iudicium}{\label{misericordiam-et-iudicium}} et iudicium cantabo tibi domine, * Psallam et intelligam in via immaculata quando venies ad me. \V Perambulabam in innocencia cordis mei, in medio domus mee.
\begin{responsory}[domine-exaudi-orationem]
{Domine exaudi orationem meam et clamor meus ad te perveniat, * Quia non spernis deus preces pauperum.}{De profundis clamavi ad te domine domine exaudi vocem meam.}
\end{responsory}%
\begin{responsory-final}[velociter-exaudi]
{Velociter exaudi me domine, * Quia defecerunt sicut fumus * Dies mei.}{Dies mei sicut umbra declinaverunt et ego sicut fenum arui.}{Dies}
\end{responsory-final}%
\newline {\color{Red} In Laudibus. Ant.} Benigne fac in bona voluntate tua domine. {\color{Red} Ps.} \psalm{50} {\color{Red} Ant.} Bonum est confiteri domino. {\color{Red} Ps.} \psalm{91} {\color{Red} Ant.} Metuant dominum omnes fines terre. {\color{Red} Ps.} \psalm{62} {\color{Red} Ant.} Et in servis suis dominus miserebitur. {\color{Red} Can.} \psalm{audite} {\color{Red} Ant.} In cimbalis benesonantibus laudate dominum. {\color{Red} Ps.} \psalm{148}
\newline {\color{Red} Ad Ben. Ant.} In viam pacis dirige nos domine.
\lettrine[lines=2]{\zallmancaps \color{Blue} S}{}\ul{\textsc{ciendum} quod ordo precedentis ebdomade de Invitat., antiphonis, psalmis, et canticis. Capitulum, Versiculum, et horarum Responsoriis ita servandus est, quando de feria agitur. Videlicet ut Invitatoria dicantur ab Octava Epyphanie usque ad Dominicam in Passione. Capitula usque ad feriam quartam in capite ieiunii exclusive. Versiculi et Responsoria ad horas usque ad Dominicam Primam Quadragesime. Singula predictorum quatuor. A} \hyperlink{deus-omnium}{Deus omnium.} \ul{usque ad Adventum. Item antiphone et psalmi de nocturnis usque ad Cenam Domini exclusive, et a} \hyperlink{deus-omnium}{Deus omnium.} \ul{usque ad Nativitatem Domini. Item antiphone Laudum usque ad Ramos Palmarum et a} \hyperlink{deus-omnium}{Deus omnium.} \ul{usque ad feriam in qua incipiuntur antiphone Laudum notate ante Vigiliam Natalis Domini exclusive. Item psalmi et cantica Laudum usque ad Sabbatum Pasche exclusive, et a} \hyperlink{deus-omnium}{Deus omnium.} \ul{usque ad Vigiliam Natalis exclusive. Item antiphone et psalmi ad Vesperis usque ad Cenam Domini exclusive, et a} \hyperlink{deus-omnium}{Deus omnium.} \ul{usque ad Vigiliam Natalis domini. Item antiphone Ad} Ben. et ad Mag. \ul{dicuntur usque ad Septuagesimam et a} \hyperlink{deus-omnium}{Deus omnium.} \ul{usque ad Adventum.}
{\ubsubsection{Dominica Secunda Post Octavas Epyphanie}{breviarium-dominica-secunda-post-octavas-epyphanie}}
\par \noindent {\color{Red} Oracio.}
\lettrine[lines=2]{\zallmancaps \color{Red} O}{mnipotens} sempiterne deus infirmitatem nostram propicius respice, atque ad protegendum nos dexteram tue maiestatis extende. Per.
\newline {\color{Red} Ad Matutinas. Lectio Prima.}
\lettrine[lines=2]{\zallmancaps \color{Blue} P}{aulus} vocatus apostolus Ihesu Christi per voluntatem Dei et Sostenes frater, ecclesie Dei que est Corinthi, sanctificatis in Christo Ihesu, vocatis sanctis, cum omnibus qui invocant nomen domini nostri Ihesu Christi in omni loco ipsorum et nostro, gracia vobis et pax a Deo patre nostro et domino Ihesu Christo.
\newline {\color{Red} Lectio Secunda.}
\lettrine[lines=2]{\zallmancaps \color{Red} G}{racias} ago Deo meo semper pro vobis in gratia Dei que data est vobis in Christo Ihesu, quia in omnibus divites facti estis in illo, in omni verbo et in omni scientia: sicut testimonium Christi confirmatum est in vobis, ita ut nichil vobis desit in ulla gratia expectantibus revelationem domini nostri Ihesu Christi, qui et confirmabit vos usque in finem sine crimine in die adventus Domini nostri Ihesu Christi.
\newline {\color{Red} Lectio Tertia.}
\lettrine[lines=2]{\zallmancaps \color{Blue} F}{idelis} Deus per quem vocati estis in societatem filii eius Ihesu Christi Domini nostri. Obsecro autem vos fratres per nomen Domini nostri Ihesu Christi, ut idipsum dicatis omnes, et non sint in vobis scismata. Sitis autem perfecti in eodem sensu et in eadem sciencia.
\newline {\color{Red} Lectio Quarta.}
\lettrine[lines=2]{\zallmancaps \color{Red} S}{ignificatum} est enim michi fratres mei de vobis ab his qui sunt Cloes, quia contentiones sunt inter vos. Hoc autem dico: quod unusquisque vestrum dicit. Ego quidem sum Pauli, ego autem Appollo, ego vero Cephe, ego autem Christi. Divisus est Christus. Nunquid Paulus crucifixus est pro nobis? aut in nomine Pauli baptizati estis?
\newline {\color{Red} Lectio Quinta.}
\lettrine[lines=2]{\zallmancaps \color{Blue} G}{ratias} ago Deo quod neminem vestrum baptizavi nisi Crispum et Gaium, ne quis dicat quod in nomine meo baptizati estis. Baptizavi autem et Stephane domum. Ceterum nescio si quem alium baptizaverim. Non enim misit me Christus baptizare, sed evangelizare. Non in sapiencia verbi: ut non evacuetur crux Christi.
\newline {\color{Red} Lectio Sexta.}
\lettrine[lines=2]{\zallmancaps \color{Red} V}{erbum} enim crucis pereuntibus quidem stulticia est, his autem qui salvi fiunt id est nobis virtus Dei est. Scriptum est enim. Perdam sapientiam sapientium: et prudentiam prudentium reprobabo. Ubi sapiens? Ubi scriba? Ubi conquisitor huius seculi? Nonne stultam fecit Deus sapientiam huius mundi?
\newline {\color{Red} Lectio VII. Secundum Matheum.}
\lettrine[lines=2]{\zallmancaps \color{Blue} I}{n} illo tempore: Cum descendisset Ihesus de monte: secute sunt eum turbe multe. Et ecce leprosus veniens adorabat eum dicens. Domine si vis potes me mundare. Et reliqua.
\newline {\color{Red} Omelia Origenis.}
% https://la.wikisource.org/wiki/Homiliae_de_tempore
\lettrine[lines=2]{\zallmancaps \color{Red} D}{ocente} in monte domino, discipuli venerunt ad eum, sicut alacres, sicut domestici, sicut proximi, sicut amici vel fratres. Ideo et Dominus ait ad illos. Vos estis sal terre, vos estis lux mundi. Nunc vero descendente eo de monte turbe secute sunt eum, que in montem ascendere non potuerunt, ut pigri populi, ut negligentes, ut imperfecti.
\newline {\color{Red} Lectio VIII.}
\lettrine[lines=2]{\zallmancaps \color{Red} I}{ta} et filii Israhel primitus in montem ascendere non valuerunt, ad obviandum Deo pergere non potuerunt propter suam irreligiositatem et impietatem, sed solus Moyses ascendit, et pauci cum eo seniorum Israhel. Ita et cum Domino discipuli soli in montem ascenderunt, et tardiores deorsum steterunt.
\newline {\color{Red} Lectio IX.}
\lettrine[lines=2]{\zallmancaps \color{Blue} D}{escendente} nunc Domino hoc est inclinante se ad infirmitatem et impotentiam ceterorum, misertoque imperfectioni eorum vel infirmitati, secute sunt eum turbe multe: aliquanti propter caritatem, aliquanti propter doctrinam, aliquanti propter admirationem et curationem. Et ecce homo leprosus unus ex illis qui curam querebant qui levamen desiderabant: Ecce leprosus veniens adorabat eum dicens. Domine, si vis, potes me mundare.
\newline {\color{Red} Ad Ben. Ant.} Cum autem descendisset Ihesus de monte, ecce leprosus veniens adorabat eum dicens, domine si vis potes me mundare, et extendens manum tetigit eum dicens, volo mundare.
\newline {\color{Red} Ad Mag. Ant.} Domine si vis potes me mundare, et ait Ihesus, volo mundare.
\newline \ul{Lectiones feriales a secundo Capitulo usque ad Epistolam ad Galathas ad libitum.}
{\ubsubsection{Dominica Tertia Post Octavas Epyphanie}{breviarium-dominica-tertia-post-octavas-epyphanie}}
\par \noindent {\color{Red} Oracio.}
\lettrine[lines=2]{\zallmancaps \color{Red} D}{eus} qui nos in tantis periculis constitutos pro humana scis fragilitate non posse subsistere, da nobis salutem mentis et corporis, ut ea que pro peccatis nostris patimur, te adiuvante vincamus. Per.
\newline {\color{Red} Ad Matutinas. Lectio Prima.}
\lettrine[lines=2]{\zallmancaps \color{Blue} P}{aulus} apostolus non ab hominibus neque per hominem, sed per Ihesum Christum et Deum Patrem qui suscitavit eum a mortuis, et qui mecum sunt omnes fratres: ecclesiis Galathie. Gratia vobis et pax a Deo Patre nostro et domino nostro Ihesu Christo, qui dedit semetipsum pro peccatis nostris, ut eriperet nos de presenti seculo nequam secundum voluntatem Dei et Patris nostri, cui est gloria in secula seculorum. Amen.
\newline {\color{Red} Lectio Secunda.}
\lettrine[lines=2]{\zallmancaps \color{Red} M}{iror} quod sic tam cito transferimini ab eo qui vos vocavit in gratiam Christi in aliud Evangelium, quod non est aliud: nisi sint aliqui qui vos conturbant, et volunt convertere Evangelium Christi.
\newline {\color{Red} Lectio Tertia.}
\lettrine[lines=2]{\zallmancaps \color{Blue} S}{ed} licet nos aut angelus de celo evangelizet vobis preterquam quod evangelizavimus vobis, anathema sit. Sicut predixi, et nunc iterum dico. Siquis vobis evangelizaverit preter id quod accepistis, anathema sit. Modo enim hominibus suadeo, an Deo? An quero hominibus placere? Si adhuc hominibus placerem, Christi servus non essem.
\newline {\color{Red} Lectio Quarta.}
\lettrine[lines=2]{\zallmancaps \color{Red} N}{otum} enim vobis facio fratres evangelium quod evangelizatum est a me, quia non est secundum hominem. Neque enim ego ab homine accepi illud, neque didici: sed per revelationem Ihesu Christi. Audistis enim conversationem meam aliquando in Iudaismo quoniam supra modum persequebar ecclesiam dei, et expugnabam illam, et proficiebam in Iudaismo supra multos coetaneos meos in genere meo, abundantius emulator existens paternarum mearum traditionum.
\newline {\color{Red} Lectio Quinta.}
\lettrine[lines=2]{\zallmancaps \color{Blue} C}{um} autem placuit ei qui me segregavit ex utero matris mee, et vocavit per gratiam suam, ut revelaret filium suum in me, ut evangelizarem illum in gentibus, continuo non adquievi carni et sanguini, neque veni Hierosolimam ad antecessores meos apostolos, sed abii in Arabiam, et iterum reversus sum Damascum.
\newline {\color{Red} Lectio Sexta.}
\lettrine[lines=2]{\zallmancaps \color{Red} D}{einde} post annos tres veni Hierosolimam videre Petrum, et mansi apud eum diebus quindecim. Alium autem apostolorum vidi neminem: nisi Iacobum fratrem Domini. Que autem scribo vobis, ecce coram Deo, quia non mencior. Deinde veni in partes Syrie et Cilicie. Eram autem ignotus facie ecclesiis Iudee que erant in Christo. Tantum autem auditum habebant, quoniam qui persequebatur nos aliquando, nunc evangelizat fidem quam aliquando expugnabat, et in me clarificabant Deum.
\newline {\color{Red} Lectio VII. Secundum Matheum.}
\lettrine[lines=2]{\zallmancaps \color{Blue} I}{n} illo tempore: Ascendente Ihesu in naviculam, secuti sunt eum discipuli eius. Et ecce motus magnus factus est in mari, ita ut navicula operiretur fluctibus. Et reliqua.
\newline {\color{Red} Omelia Origenis.}
% https://la.wikisource.org/wiki/Homiliae_de_tempore
\lettrine[lines=2]{\zallmancaps \color{Red} I}{ngrediente} domino in naviculam, secuti sunt eum discipuli eius, non imbecilles, sed firmi et stabiles in fide, mansueti et pii, spernentes mundum, non duplici corde sed simplici. Hi ergo secuti sunt eum: non tantum gressus eius sequentes, sed magis sanctitatem comitantes, et iusticiam eius consectantes.
\newline {\color{Red} Lectio VIII.}
\lettrine[lines=2]{\zallmancaps \color{Blue} E}{t} ecce tempestas magna facta est in mari, ita ut navicula operiretur fluctibus. Cum enim multa magna et miranda ostendisset in terra transiit ad mare, ut et ibidem adhuc excellentiora opera demonstraret, quatinus terre marisque Dominum se esse cunctis ostenderet.
\newline {\color{Red} Lectio IX.}
\lettrine[lines=2]{\zallmancaps \color{Red} I}{ngressus} ergo naviculam fecit turbari mare, commovit ventos, concitavit fluctus. Cur hoc? Ideo ut discipulos mitteret in timorem, et suum auxilium postularent, suamque potenciam rogantibus manifestaret. Illa tempestas non ex sese oborta est, sed potestati paruit imperantis, eius qui educit ventos de thesauris suis, qui terminum mari harenam constituit.
\newline {\color{Red} Ad Ben. Ant.} Ascendente Ihesu in naviculam ecce motus magnus factus est in mari, et suscitaverunt eum discipuli eius dicentes domine salva nos perimus.
\newline {\color{Red} Ad Mag. Ant.} Domine salva nos perimus impera et fac deus tranquillitatem.
\newline \ul{Lectiones feriales a tercio capitulo usque ad Epistolam ad Thessalonicenses ad libitum.}
{\ubsubsection{Dominica Quarta Post Octavas Epyphanie}{breviarium-dominica-quarta-post-octavas-epyphanie}}
\par \noindent {\color{Red} Oracio.}
\lettrine[lines=2]{\zallmancaps \color{Blue} F}{amiliam} \hypertarget{familiam-tuam-dominica-quarta-post-octavas-epyphanie}{tuam} quesumus domine continua pietate custodi, ut que in sola spe gracie celestis innititur, tua semper protectione muniatur. Per.
\newline {\color{Red} Ad Matutinas. Lectio Prima.}
\lettrine[lines=2]{\zallmancaps \color{Red} P}{aulus} et Silvanus, et Timotheus: ecclesie Thessalonicensium in Deo Patre et Domino Ihesu Christo: gratia vobis et pax. Gracias agimus Deo semper pro omnibus vobis, memoriam vestri facientes in orationibus nostris sine intermissione memores operis fidei vestre, et laboris, et caritatis, et sustinencie spei Domini nostri Ihesu Christi ante Deum et Patrem nostrum. Scientes fratres dilecti a Deo electionem vestram, quia evangelium nostrum non fuit apud vos in sermone tantum, sed in virtute et in Spiritu Sancto et in plenitudine multa, sicut scitis quales fuerimus in vobis propter vos.
\newline {\color{Red} Lectio Secunda.}
\lettrine[lines=2]{\zallmancaps \color{Blue} E}{t} vos imitatores nostri facti estis et Domini, excipientes verbum in tribulatione multa cum gaudio Spiritus Sancti, ita ut facti sitis forma omnibus credentibus in Macedonia et in Achaia. A vobis enim diffamatus est sermo Domini non solum in Macedonia et in Achaia, sed in omni loco fides vestra que est ad Deum profecta est, ita ut non sit vobis necesse quicquam loqui.
\newline {\color{Red} Lectio Tertia.}
\lettrine[lines=2]{\zallmancaps \color{Red} I}{psi} enim de vobis annunciant qualem introitum habuerimus ad vos, et quomodo conversi estis ad Deum a simulacris servire Deo vivo et vero, et expectare filium eius de celis quem suscitavit a mortuis Ihesum, qui eripuit nos ab ira ventura.
\newline {\color{Red} Lectio Quarta.}
\lettrine[lines=2]{\zallmancaps \color{Blue} N}{am} ipsi scitis fratres introitum nostrum ad vos, quia non inanis fuit, sed ante passi et contumeliis affecti sicut scitis in Philippis: fiduciam habuimus in Deo nostro loqui ad vos evangelium Dei in multa sollicitudine.
\newline {\color{Red} Lectio Quinta.}
\lettrine[lines=2]{\zallmancaps \color{Red} E}{xortatio} enim nostra non de errore, neque de immundicia, neque in dolo, sed sicut probati sumus a Deo ut crederetur nobis Evangelium: ita loquimur, non quasi hominibus placentes, sed Deo qui probat corda nostra.
\newline {\color{Red} Lectio Sexta.}
\lettrine[lines=2]{\zallmancaps \color{Blue} N}{eque} enim aliquando fuimus in sermone adulationis sicut scitis, neque in occasione avaritie, Deus testis est, nec querentes ab hominibus gloriam, neque a vobis neque ab aliis, cum possemus oneri vobis esse ut Christi apostoli, sed facti sumus parvuli in medio vestrum tanquam si nutrix foveat filios suos. Ita desiderantes vos cupide, volebamus tradere vobis non solum Evangelium Dei, sed eciam animas nostras, quoniam karissimi nobis facti estis.
\newline {\color{Red} Lectio VII. Secundum Matheum.}
\lettrine[lines=2]{\zallmancaps \color{Red} I}{n} illo tempore: Dixit Ihesus discipulis suis, parabolam hanc. Simile factum est regnum celorum, homini qui seminavit bonum semen in agro suo. Et reliqua.
\newline {\color{Red} Omelia Beati Iheronimi Presbiteri.}
\lettrine[lines=2]{\zallmancaps \color{Blue} D}{imittit} turbas Ihesus, et domum revertitur, ut accedant ad eum discipuli, et interrogent que nec populus merebatur
%pp097
audire, nec poterat. Nam dimissis turbis venit in domum, et accesserunt ad eum discipuli eius dicentes. Edissere nobis parabolam zizaniorum agri. Qui respondens ait. Qui seminat bonum semen: est filius hominis.
\newline {\color{Red} Lectio VIII.}
\lettrine[lines=2]{\zallmancaps \color{Red} P}{erspicue} vero exposuit quod ager mundus sit, sator filius hominis, bonum semen filii regni, zizania filii pessimi, zizaniorum sator diabolus, messis consummacio mundi, messores angeli. Omnia ergo scandala referuntur ad zizania: iusti reputantur in filios regni.
\newline {\color{Red} Lectio IX.}
\lettrine[lines=2]{\zallmancaps \color{Blue} E}{rgo} que exposita sunt a Domino, his debemus accommodare fidem. Que autem tacita et nostre intelligentie derelicta, breviter sunt intelligenda. Nam homines dormientes significant ecclesiarum magistros: et servi patrisfamilias angelos, qui cotidie vident faciem Patris.
\newline {\color{Red} Ad Ben. Ant.} Domine nonne bonum semen seminasti in agro tuo, unde habet zizania, et ait, hoc fecit inimicus homo.
\newline {\color{Red} Ad Mag. Ant.} Colligite primum zizania, et alligate ea fasciculos ad comburendum, triticum autem congregate in orreum meum.
\newline \ul{Lectiones feriales a tercio capitulo usque ad epistolam ad Thimotheum. II.}
{\ubsubsection{Dominica Quinta Post Octavas Epyphanie}{breviarium-dominica-quinta-post-octavas-epyphanie}}
\par \noindent {\color{Red} Oracio.} \hyperlink{familiam-tuam-dominica-quarta-post-octavas-epyphanie}{Familiam}.
\newline {\color{Red} Ad Matutinas. Lectio Prima.}
\lettrine[lines=2]{\zallmancaps \color{Red} P}{aulus} Apostolus Ihesu Christi per voluntatem Dei, secundum promissionem vite que est in Christo Ihesu Thymotheo karissimo filio, gratia, misericordia, pax a Deo Patre, et Christo Ihesu Domino nostro. Gratias ago Deo meo cui servio a progenitoribus meis in conscientia pura, quod sine intermissione habeam tui memoriam in orationibus meis nocte ac die desiderans te videre, memor lacrimarum tuarum, ut gaudio implear, recordationem accipiens eius fidei que est in te non ficta, que et habitavit primum in avia tua Loyde, et matre tua Euniche, certus sum autem quod et in te.
\newline {\color{Red} Lectio Secunda.}
\lettrine[lines=2]{\zallmancaps \color{Blue} P}{ropter} quam causam admoneo te, ut resuscites gratiam Dei que est in te per impositionem manuum mearum. Non enim dedit nobis Deus spiritum timoris, sed virtutis et dilectionis et sobrietatis.
\newline {\color{Red} Lectio Tertia.}
\lettrine[lines=2]{\zallmancaps \color{Red} N}{oli} itaque erubescere testimonium Domini nostri, neque me vinctum eius, sed collabora Evangelio secundum virtutem Dei qui nos liberavit, et vocavit vocatione sua sancta, non secundum opera nostra, sed secundum propositum suum et graciam que data est nobis in Christo Ihesu ante tempora secularia. Manifestata est autem nunc per illuminationem Salvatoris nostri Ihesu Christi, qui destruxit quidem mortem, illuminavit autem vitam et incorruptionem, per evangelium in quo positus sum ego predicator et Apostolus et magister gentium.
\newline {\color{Red} Lectio Quarta.}
\lettrine[lines=2]{\zallmancaps \color{Blue} O}{b} quam causam eciam hec patior: sed non confundor. Scio enim cui credidi, et certus sum quia potens est depositum meum servare in illum diem. Formam habens sanorum verborum que a me audisti in fide et in dilectione in Christo Ihesu, bonum depositum custodi per Spiritum Sanctum qui habitat in nobis.
\newline {\color{Red} Lectio Quinta.}
\lettrine[lines=2]{\zallmancaps \color{Red} S}{cis} hoc quod aversi sunt a me omnes qui in Asia sunt, ex quibus est Phygelus et Hermogenes. Det misericordiam Dominus Onesiphori domui, quia sepe me refrigeravit et catenam meam non erubuit, sed cum Romam venisset sollicite me quesivit, et invenit. Det illi Dominus invenire misericordiam a Domino in illa die. Et quanta Ephesi ministravit michi: tu melius nosti.
\newline {\color{Red} Lectio Sexta.}
\lettrine[lines=2]{\zallmancaps \color{Blue} T}{u} ergo fili mi confortare in gratia que est in Christo Ihesu, et que audisti a me per multos testes hec commenda fidelibus hominibus, qui ydonei erant et alios docere. Labora sicut bonus miles Christi Ihesu. Nemo militans Deo implicat se negociis secularibus, ut ei placeat cui se probavit. Nam et qui certat in agone, non coronatur nisi legitime certaverit.
\newline \ul{Tres ultime lectiones de omelia precedentis Dominice. Ad Ben. et ad Mag. antiphone sicut in Dominica precedenti. Lectiones feriales a tercio capitulo usque ad finem epistolarum ad libitum.}
{\ubsubsection{Dominica in Septuagesima Sabbato Precedenti}{breviarium-dominica-in-septuagesima-sabbato-precedenti}}
\par \noindent {\color{Red} Ad Vesperas.} \R \hyperlink{igitur-perfecti}{Igitur}. {\color{Red} III. Ad Mag. Ant.} Peccata. {\color{Red} Oracio.} 
\lettrine[lines=2]{\zallmancaps \color{Red} P}{reces} populi tui quesumus domine clementer exaudi, ut qui iuste pro peccatis nostris affligimur, pro tui nominis gloria misericorditer liberemur. Per. Benedicamus domino alleluya alleluya.
\newline \ul{Ab hinc usque ad Missam Vigilie Pasche non dicatur alleluya, sed in principio horarum post} sicut erat \ul{dicatur} Laus tibi domine rex eterne glorie.
\newline {\color{Red} Ad Matutinas. Invit.}
\begin{invitatory}[preoccupemus-faciem-invitatorium]
{Preoccupemus faciem domini * Et in psalmis iubilemus ei.}
\end{invitatory}
\newline \ul{Predictum invitat. non resumatur post} salutari nostro \ul{sed post} iubilemus ei.
\newline {\color{Red} Lectio Prima.}
\lettrine[lines=2]{\zallmancaps \color{Blue} I}{n} principio creavit Deus celum et terram. Terra autem erat inanis et vacua, et tenebre erant super faciem abyssi, et spiritus Dei ferebatur super aquas. Dixitque Deus. Fiat lux. Et facta est lux. Et vidit Deus lucem quod esset bona: et divisit lucem a tenebris. Appellavitque lucem diem, et tenebras noctem. Factumque est vespere et mane dies unus. {\color{Red} Responsorium.}
\lettrine[lines=2]{\zallmancaps \color{Red} I}{n} \hypertarget{in-principio-fecit}{\label{in-principio-fecit}} principio fecit deus celum et terram, et creavit in ea hominem * Ad ymaginem et similitudinem suam. \V Formavit igitur dominus hominem de limo terre, et inspiravit in faciem eius spiraculum vite. Ad ymaginem.
\newline {\color{Red} Lectio Secunda.}
\lettrine[lines=2]{\zallmancaps \color{Blue} D}{ixit} quoque Deus. Fiat firmamentum in medio aquarum, et dividat aquas ab aquis. Et fecit Deus firmamentum, divisitque aquas que erant sub firmamento ab his que erant super firmamentum. Et factum est ita. Vocavitque Deus firmamentum celum. Et factum est vespere et mane dies secundus.
\begin{responsory}[in-principio-deus]
{In principio deus creavit celum et terram et spiritus domini ferebatur super aquas, et vidit deus cuncta que fecerat, * Et erant valde bona.}{Igitur perfecti sunt celi et terra et omnis ornatus eorum.}
\end{responsory}%
\newline {\color{Red} Lectio Tertia.}
\lettrine[lines=2]{\zallmancaps \color{Red} D}{ixit} vero Deus. Congregentur aque que sub celo sunt in locum unum, et appareat arida. Factumque est ita. Et vocavit Deus aridam Terram, congregationesque aquarum appellavit Maria. Et vidit Deus quod esset bonum, et ait. Germinet terra herbam virentem et facientem semen: et lignum pomiferum faciens fructum iuxta genus suum, cuius semen in semetipso sit super terram. Et factum est ita.
\begin{responsory-final}[igitur-perfecti]
{Igitur perfecti sunt celi et terra, et omnis ornatus eorum complevitque deus die septimo opus suum quod fecerat, * Et requievit, * Ab omni opere quod patrarat.}{Viditque deus cuncta que fecerat et erant valde bona.}{Ab omni}
\end{responsory-final}%
\newline {\color{Red} Lectio Quarta.}
\lettrine[lines=2]{\zallmancaps \color{Blue} E}{t} protulit terra herbam virentem, et facientem semen iuxta genus suum, lignumque faciens fructum et habens unumquodque sementem secundum speciem suam. Et vidit Deus quod esset bonum, factumque est vespere et mane dies tercius.
\begin{responsory}[formavit-igitur]
{Formavit igitur dominus hominem de limo terre et inspiravit in faciem eius spiraculum vite, et factus est homo, * In animam viventem.}{In principio fecit deus celum et terram et creavit in ea hominem.}
\end{responsory}%
\newline {\color{Red} Lectio Quinta.}
\lettrine[lines=2]{\zallmancaps \color{Red} D}{ixit} autem Deus. Fiant luminaria in firmamento celi, et dividant diem ac noctem, et sint in signa et tempora et dies et annos: ut luceant in firmamento celi, et illuminent terram. Et factum est ita.
\begin{responsory}[tulit-ergo]
{Tulit ergo dominus hominem et posuit eum in paradiso voluptatis, * Ut operaretur et custodiret illum.}{Plantaverat autem dominus deus paradisum voluptatis a principio in quo posuit hominem quem formaverat.}
\end{responsory}%
\newline {\color{Red} Lectio Sexta.}
\lettrine[lines=2]{\zallmancaps \color{Blue} F}{ecitque} Deus duo magna luminaria, luminare maius ut preesset diei, et luminare minus ut preesset nocti. Et stellas. Et posuit eas in firmamento celi, ut lucerent super terram, et preessent diei ac nocti, et dividerent lucem ac tenebras. Et vidit Deus quod esset bonum: et factum est vespere et mane dies quartus.
\begin{responsory-doxology}[dixit-dominus]
{Dixit dominus deus non est bonum esse hominem solum, * Faciamus ei adiutorium simile sibi.}{Ade vero non inveniebatur adiutor similis sibi dixit quoque deus.}
\end{responsory-doxology}%
\newline {\color{Red} Lectio VII. Secundum Matheum.}
\lettrine[lines=2]{\zallmancaps \color{Red} I}{n} illo tempore: Dixit Ihesus discipulis suis, parabolam hanc. Simile est regnum celorum homini patrifamilias: qui exiit primo mane conducere operarios in vineam suam. Et reliqua.
\newline {\color{Red} Omelia Beati Gregorii Pape.}
\lettrine[lines=2]{\zallmancaps \color{Blue} I}{n} explanatione sua multa ad loquendum sancti evangelii lectio postulat, quam volo si possum sub brevitate perstringere, ne vos et extensa processio et prolixa expositio videatur onerare. Regnum celorum patrifamilias simile dicitur, qui ad excolendam vineam suam operarios conducit.
\begin{responsory}[immisit-dominus]
{Immisit dominus soporem in Adam, et tulit unam de costis eius, et edificavit costam quam tulerat de Adam in mulierem, et adduxit eam ad Adam ut videret quid vocaret eam, * Et vocavit nomen eius virago quia de viro suo sumpta est.}{Cumque obdormisset tulit unam de costis eius et replevit carnem pro ea.}
\end{responsory}%
\newline {\color{Red} Lectio VIII.}
\lettrine[lines=2]{\zallmancaps \color{Red} Q}{uis} vero patrisfamilias similitudinem rectius tenet quam conditor noster, qui regit quos condidit, et electos suos sic in hoc mundo possidet, quasi subiectos dominus in domo? Qui habet vineam, universalem scilicet ecclesiam que ab Abel iusto usque ad ultimum electum qui in fine mundi nasciturus est, quot sanctos protulit: quasi tot palmites misit.
\begin{responsory}[ecce-adam]
{Ecce Adam quasi unus ex nobis factus est sciens bonum et malum, * Videte ne forte sumat de ligno vite et vivat in eternum.}{Cherubin collocavit dominus et flammeum gladium atque versatilem ad custodiendam viam ligni vite et ait.}
\end{responsory}%
\newline {\color{Red} Lectio IX.}
\lettrine[lines=2]{\zallmancaps \color{Red} H}{ic} itaque paterfamilias ad excolendam vineam suam mane, hora tercia, sexta, nona, et undecima operarios conducit, quia a mundi huius inicio usque in finem ad erudiendam plebem fidelium predicatores congregare non destitit. Mane etenim mundi: fuit ab Adam usque ad Noe.
\begin{responsory-doxology}[ubi-est-abel]
{Ubi est Abel frater tuus dixit dominus ad Chaim, nescio domine nunquid custos fratris mei sum ego et dixit ad eum, quid fecisti, * Ecce vox sanguinis fratris tui Abel clamat ad me de terra.}{Maledictus eris super terram que aperuit os suum, et suscepit sanguinem fratris tui de manu tua.}
\end{responsory-doxology}%
\newline \ul{Non dicatur} \psalm{tedeum} \ul{nec deinceps usque ad Pascha nisi in totis duplicibus, sed nonum responsorium repetatur.}%repetitur Ubi est Abel, etc
\newline {\color{Red} In Laudibus. Ant.} Miserere mei deus, et a delicto meo munda me, quia tibi soli peccavi. {\color{Red} Ps.} \psalm{50} {\color{Red} Ant.} Confitebor tibi domine quoniam exaudisti me. {\color{Red} Ps.} \psalm{117} {\color{Red} Ant.} Deus deus meus ad te de luce vigilo, quia factus es adiutor meus. {\color{Red} Ps.} \psalm{62} {\color{Red} Ant.} Benedictus es in firmamento celi et laudabilis in secula deus noster. {\color{Red} Can.} \psalm{benedicite} {\color{Red} Ant.} Laudate dominum de celis. {\color{Red} Ps.} \psalm{148} \V Domine refugium factus es nobis. \R A generatione et progenie. {\color{Red} Ad Ben. Ant.} Simile est regnum celorum homini patrifamilias qui exiit primo mane conducere operarios in vineam suam, dicit dominus.
\newline {\color{Red} Ad Primam. Ant.} Conventione autem facta cum operariis ex denario diurno misit eos in vineam suam. {\color{Red} Ps.} \psalm{21} \psalm{22} {\color{Red} Gloria.} \psalm{23} \psalm{24} {\color{Red} Gloria.} \psalm{25} \psalm{53} {\color{Red} Gloria.} \psalm{92} \psalm{118.1} {\color{Red} Gloria.} \psalm{118.2} {\color{Red} Gloria.} \psalm{quicumque} {\color{Red} Gloria.}
\ul{Predicti psalmi dicantur predicto modo Dominicis diebus ad Primam usque ad Ramos inclusive quando de Dominica agitur.}
\newline {\color{Red} Ad Terciam. Ant.} Ite et vos in vineam meam et quod iustum fuerit dabo vobis.
\begin{responsory-breve}[spes-mea]
{Spes mea domine * A iuventute mea.}{In te confirmatus sum ex utero de ventre matris mee tu es meus protector.}
\end{responsory-breve}%
\V Adiutor meus.
\newline {\color{Red} Ad Sextam. Ant.} Quid hic statis tota die ociosi responderunt et dixerunt nemo nos conduxit.
\begin{responsory-breve}[esto-nobis]
{Esto nobis domine, * Turris fortitudinis.}{A facie inimici.}
\end{responsory-breve}%
\V Dominus regit me.
\newline {\color{Red} Ad Nonam. Ant.} Voca operarios et redde illis mercedem suam dicit dominus.
\begin{responsory-breve}[educ-de-carcere]
{Educ de carcere animam meam, * Ut confiteatur nomini tuo domine.}{Periit a fuga a me et non est qui requirat animam meam.}
\end{responsory-breve}%
\V Ab occultis.
\newline {\color{Red} Ad Mag. Ant.} Dixit paterfamilias operariis suis, quid hic statis tota die ociosi at illi respondentes dixerunt, qui nemo nos conduxit, ite et vos in vineam meam et quod iustum fuerit dabo vobis.
\newline \ul{Per ebdomadam quando de tempore agitur.}
\newline {\color{Red} Ad Ben. Ant.} Hi novissimi una hora fecerunt et pares illos nobis fecisti, qui portavimus pondus diei et estus.
\newline {\color{Red} Ad Mag. Ant.} Erunt primi novissimi et novissimi primi, multi enim sunt vocati, pauci vero electi, dicit dominus.
\newline \ul{Lectiones feriales a secundo capitulo usque ad sextum ad libitum.}
{\ubsubsection{Dominica in Sexagesima Sabbato Precedenti}{breviarium-dominica-in-sexagesima-sabbato-precedenti}}
\par \noindent {\color{Red} Ad Vesperas.} \R \hyperlink{benedicens-ergo}{Benedicens}. {\color{Red} IX. Ad Mag. Ant.} Erunt primi. {\color{Red} Oracio.} 
\lettrine[lines=2]{\zallmancaps \color{Blue} D}{eus} qui conspicis quia ex nulla nostra actione confidimus, concede propicius ut contra adversa omnia doctoris gentium protectione muniamur. Per.
\newline {\color{Red} Ad Matutinas. Invit.} \hyperlink{preoccupemus-faciem-invitatorium}{Preoccupemus.}
\newline {\color{Red} Lectio Prima.}
\lettrine[lines=2]{\zallmancaps \color{Red} C}{um} cepissent homines multiplicari super terram, et filias procreassent, videntes filii Dei filias hominum quod essent pulchre, acceperunt uxores sibi ex omnibus quas elegerant. Dixitque Deus. Non permanebit spiritus meus in homine in eternum: quia caro est. Eruntque dies illius centum viginti annorum. \R
\lettrine[lines=2]{\zallmancaps \color{Blue} N}{oe} \hypertarget{noe-vir-iustus}{\label{noe-vir-iustus}} vir iustus atque perfectus cum deo ambulavit, * Et fecit omnia quecumque precepit ei deus. \V Noe autem invenit graciam ante dominum deum. Et fecit.
\newline {\color{Red} Lectio Secunda.}
\lettrine[lines=2]{\zallmancaps \color{Red} G}{igantes} autem erant super terram in diebus illis. Postquam enim ingressi sunt filii Dei ad filias hominum, illeque genuerunt: isti sunt potentes a seculo, viri famosi.
\begin{responsory}[quadraginta-dies]
{Quadraginta dies et noctes aperti sunt celi, et ex omni carne habente spiritum vite, * Ingressi sunt in archam et clausit a foris hostium dominus.}{Noe vero et filii eius uxor eius et uxores filiorum eius.}
\end{responsory}%
\newline {\color{Red} Lectio Tertia.}
\lettrine[lines=2]{\zallmancaps \color{Blue} V}{idens} autem Deus quod multa malicia hominum esset in terra, et cuncta cogitatio cordis intenta esset ad malum omni tempore, penituit eum quod hominum fecisset in terra, et precavens in futurum, et tactus dolore cordis intrinsecus, delebo inquit hominem quem creavi a facie terre ab homine usque ad animancia, a reptili usque ad volucres celi. Penitet enim me fecisse eos.
\begin{responsory-doxology}[requievit-archa]
{Requievit archa mense septimo super montes Armenie, * Et aque ibant et decrescebant usque ad decimum mensem.}{Decimo enim mense primo die mensis cacumina montium apparuerunt.}
\end{responsory-doxology}%
\newline {\color{Red} Lectio Quarta.}
\lettrine[lines=2]{\zallmancaps \color{Red} N}{oe} vero invenit gratiam coram Domino. He sunt generationes Noe. Noe vir iustus atque perfectus fuit in generationibus suis. Cum Deo ambulavit, et genuit tres filios, Sem, Cham, et Iaphet. Corrupta est autem terra coram Deo, et repleta est iniquitate.
\begin{responsory}[volens-noe]
{Volens Noe scire si iam cessassent aque emisit columbam que ramum virentis olive in ore suo deferens, * Ad archam reversa est.}{Deferens autem signum clementie dei columba in ore suo.}
\end{responsory}%
\newline {\color{Red} Lectio Quinta.}
\lettrine[lines=2]{\zallmancaps \color{Blue} C}{umque} vidisset Deus terram esse corruptam, omnis quippe caro corruperat viam suam super terram: dixit ad Noe. Finis universe carnis venit coram me. Repleta est terra iniquitate a facie eorum, et ego disperdam eos cum terra. Fac tibi archam de lignis levigatis.
\begin{responsory}[edificavit-noe]
{Edificavit Noe altare domino offerens super illud holocaustum, odoratus est dominus odorem suavitatis et benedixit eis, * Crescite et multiplicamini et replete terram.}{Ecce ego statuam pactum meum vobiscum et cum semine vestro post vos.}
\end{responsory}%
\newline {\color{Red} Lectio Sexta.}
\lettrine[lines=2]{\zallmancaps \color{Red} M}{ansiunculas} in archa facies, et bitumine linies intrinsecus et extrinsecus, et sic facies eam. Trecentorum cubitorum erit longitudo arche, quinquaginta cubitorum latitudo, et triginta cubitorum altitudo illius. Fenestram in archa facies, et in cubito consummabis summitatem eius. Ostium autem arche pones in latere deorsum. Cenacula et tristega facies in ea.
\begin{responsory-doxology}[ponam-arcum]
{Ponam arcum meum in nubibus celi dixit dominus ad Noe, * Et recordabor federis mei quod pepigi tecum.}{Cumque obduxero nubibus celum apparebit arcus meus in nubibus.}
\end{responsory-doxology}%
\newline {\color{Red} Lectio VII. Secundum Lucam.}
\lettrine[lines=2]{\zallmancaps \color{Blue} I}{n} illo tempore: Cum turba plurima conveniret, et de civitatibus properarent ad Ihesum: dixit per similitudinem. Exiit qui seminat seminare semen suum. Et reliqua.
\newline {\color{Red} Omelia Beati Gregorii Pape.}
\lettrine[lines=2]{\zallmancaps \color{Red} L}{ectio} sancti evangelii quam modo fratres karissimi audistis, expositione non indiget sed admonitione. Quam enim per semetipsam veritas exposuit, hanc discutere humana fragilitas non presumit.
\begin{responsory}[per-memetipsum]
{Per memetipsum iuravi dicit dominus, non adiciam ultra aquas diluvii super terram, pacti mei recordabor, * Ut non perdam aquis diluvii omnem carnem.}{Arcum meum ponam in nubibus, et erit signum federis inter me et inter terram.}
\end{responsory}%
\newline {\color{Red} Lectio VIII.}
\lettrine[lines=2]{\zallmancaps \color{Blue} S}{ed} est quod sollicite in hac ipsa expositione dominica pensare debeatis. Quia si nos vobis, semen verbum, agrum mundum, volucres demonia, spinas divitias significare diceremus, ad credendum nobis mens forsitan vestra dubitaret. Unde et idem Dominus per semetipsum dignatus est exponere quod dicebat, ut sciatis rerum significationes querere, in his eciam que per semetipsum noluit explanare.
\begin{responsory}[ecce-ego-statuam]
{Ecce ego statuam pactum meum vobiscum et cum semine vestro post vos, * Neque erit deinceps diluvium dissipans terram.}{Erit arcus meus in nubibus, et videbo illum, et recordabor federis mei sempiterni.}
\end{responsory}%
\newline {\color{Red} Lectio IX.}
\lettrine[lines=2]{\zallmancaps \color{Red} E}{xponendo} ergo quod dixit, figurate se loqui innotuit: quatenus certos vos redderet cum vobis nostra fragilitas verborum illius figuras aperiret. Quis enim michi unquam crederet si spinas divicias interpretari voluissem, maxime cum ille pungant, iste delectent? Et tamen spine sunt, quia cogitationum suarum punctionibus mentem lacerant, et cum usque ad peccatum pertrahunt, quasi inflicto vulnere cruentant.
\begin{responsory-doxology}[benedicens-ergo]
{Benedicens ergo deus Noe ait, nequaquam ultra maledicam terre propter hominem, * Ad ymaginem quippe dei factus est homo.}{Hoc erit signum federis inter me et terram, arcum meum ponam in nubibus celi.}
\end{responsory-doxology}%
Benedicens.
\newline {\color{Red} In Laudibus. Ant.} Secundum multitudinem miserationum tuarum domine dele iniquitatem meam. {\color{Red} Ps.} \psalm{50} {\color{Red} Ant.} Deus meus es tu et confitebor tibi, deus meus es tu et exaltabo te. {\color{Red} Ps.} \psalm{117} {\color{Red} Ant.} Ad te de luce vigilo deus ut videam virtutem tuam. {\color{Red} Ps.} \psalm{62} {\color{Red} Ant.} Ymnum dicite et superexaltate eum in secula. {\color{Red} Can.} \psalm{benedicite} {\color{Red} Ant.} Omnes angeli eius laudate dominum de celis. {\color{Red} Ps.} \psalm{148} \V Domine refugium. {\color{Red} Ad Ben. Ant.} Cum turba plurima conveniret ad Ihesum et de civitatibus properarent ad eum, dixit per similitudinem exiit qui seminat seminare semen suum.
\newline {\color{Red} Ad Primam. Ant.} Semen cecidit in terram bonam, et attulit fructum, aliud centesimum et aliud sexagesimum. {\color{Red} Ps.}
%pp098
\psalm{21}
\newline {\color{Red} Ad Terciam. Ant.} Semen cecidit in terram bonam et attulit fructum in patientia. \R \hyperlink{spes-mea}{Spes mea domine.}
\newline {\color{Red} Ad Sextam. Ant.} Ihesus hec dicens clamabat qui habet aures audiendi audiat. \R \hyperlink{esto-nobis}{Esto nobis.}
\newline {\color{Red} Ad Nonam. Ant.} Qui verbum dei retinent corde perfecto et optimo fructum afferunt in patiencia. \R \hyperlink{educ-de-carcere}{Educ de carcere.}
\newline {\color{Red} Ad Mag. Ant.} Vobis datum est nosse mysterium regni dei, ceteris autem in parabolis dixit Ihesus discipulis suis.
\newline \ul{Per ebdomadam quando de tempore agitur.}
\newline {\color{Red} Ad Ben. Ant.} Semen est verbum dei sator autem Christus omnis qui audit eum manebit in eternum.
\newline {\color{Red} Ad Mag. Ant.} Quod autem cecidit in terram bonam hi sunt qui in corde bono et optimo verbum retinent et fructum afferunt in patientia.
\newline \ul{Lectiones feriales ab octavo Capitulo usque ad duodecimum ad libitum.}
{\ubsubsection{Dominica in Quinquagesima Sabbato Precedenti}{breviarium-dominica-in-quinquagesima-sabbato-precedenti}}
\par \noindent {\color{Red} Ad Vesperas.} \R \hyperlink{dum-staret}{Dum staret}. {\color{Red} IIII. Ad Mag. Ant.} Quod autem. {\color{Red} Oratio.} 
\lettrine[lines=2]{\zallmancaps \color{Red} P}{reces} nostras quesumus domine clementer exaudi, atque a peccatorum vinculis absolutos, ab omni nos adversitate custodi. Per.
\newline {\color{Red} Ad Matutinas. Invit.} \hyperlink{preoccupemus-faciem-invitatorium}{Preoccupemus.}
\newline {\color{Red} Lectio Prima.}
\lettrine[lines=2]{\zallmancaps \color{Blue} D}{ixit} Dominus ad Abram. Egredere de terra tua, et de cognatione tua, et de domo patris tui, et veni in terram quam monstravero tibi. Faciamque te in gentem magnam, et benedicam tibi, et magnificabo nomen tuum: erisque benedictus. Benedicam benedicentibus tibi, et maledicam maledicentibus tibi: atque in te benedicentur universe cognationes terre. \R
\lettrine[lines=2]{\zallmancaps \color{Red} L}{ocutus} \hypertarget{locutus-est-quinquagesima}{\label{locutus-est-quinquagesima}} est dominus ad Abraham dicens egredere de terra et de cognatione tua, et veni in terram quam monstravero tibi, et faciam te crescere, * In gentem magnam. \V Benedicens benedicam tibi et multiplicabo te. In gentem.
\newline {\color{Red} Lectio Secunda.}
\lettrine[lines=2]{\zallmancaps \color{Blue} E}{gressus} est itaque Abram sicut preceperat ei Dominus, et ivit cum eo Loth. Septuaginta quinque annorum erat Abram, cum egrederetur de Haran. Tulitque Sarai uxorem suam, et Loth filium fratris sui, universamque substantiam quam possederant, et animas quas fecerant in Haran, et egressi sunt ut irent in terram Chanaan.
\begin{responsory}[revertenti-abraham]
{Revertenti Abraham a cede quatuor regum occurrit rex Salem Melchisedech, offerens panem et vinum erat enim dei sacerdos, * Et benedixit illi.}{Benedictus Abraham deo altissimo qui creavit celum et terram.}
\end{responsory}%
\newline {\color{Red} Lectio Tertia.}
\lettrine[lines=2]{\zallmancaps \color{Red} C}{umque} venissent in eam, pertransivit Abram terram usque ad locum Sichem, et usque ad convallem illustrem. Chananeus autem tunc erat in terra. Apparuit autem Dominus Abram, et dixit ei. Semini tuo dabo terram hanc. Qui edificavit ibi altare Domino qui apparuerat, ei.
\begin{responsory-doxology}[dixit-autem]
{Dixit autem dominus ad Abraham ego sum et pactum meum tecum, sanctificetur in vobis omne masculinum, * Et erit signum inter me et semen tuum.}{Ex te namque egredietur in quo omnes gentes benedicentur.}
\end{responsory-doxology}%
\newline {\color{Red} Lectio Quarta.}
\lettrine[lines=2]{\zallmancaps \color{Blue} E}{t} inde transgrediens ad montem qui erat contra orientem Bethel, tetendit ibi tabernaculum suum, ab occidente habens Bethel, et ab oriente Hay. Edificavit quoque ibi altare Domino, et invocavit nomen eius. Perrexitque Abram vadens et ultra progrediens ad meridiem.
\begin{responsory}[dum-staret]
{Dum staret Abraham ad ilicem Mambre, vidit tres viros descendentes per viam, tres vidit, * Et unum adoravit.}{Cumque elevasset oculos apparuerunt ei tres viri stantes iuxta eum.}
\end{responsory}%
\newline {\color{Red} Lectio Quinta.}
\lettrine[lines=2]{\zallmancaps \color{Red} F}{acta} est autem fames in terra. Descenditque Abram in Egyptum, ut peregrinaretur ibi. Prevaluerat enim fames in terra. Cumque prope esset ut ingrederetur Egyptum, dixit Sarai uxori sue. Novi quod pulchra sis mulier, et quod cum viderint te Egyptii, dicturi sunt, uxor ipsius est, et interficient me, et te reservabunt. Dic ergo obsecro te quod soror mea sis, ut bene sit michi propter te, et vivat anima mea ob gratiam tui.
\begin{responsory}[clamor-inquit]
{Clamor inquit dominus Sodomorum et Gomore venit ad me, * Descendam et videbo utrum clamorem opere compleverint.}{Abraham stabat coram deo et ait, absit a te domine ut perdas iustum cum impiis.}
\end{responsory}%
\newline {\color{Red} Lectio Sexta.}
\lettrine[lines=2]{\zallmancaps \color{Blue} C}{um} itaque ingressus esset Abram Egyptum, viderunt Egyptii mulierem quod esset pulchra nimis. Et nunciaverunt principes Pharaoni, et laudaverunt eam apud illum. Et sublata est mulier in domum Pharaonis. Abram vero bene usi sunt propter illam. Fueruntque ei oves et boves et asini, et servi, et famule, et asine, et cameli. Flagellavit autem Dominus Pharaonem plagis maximis, et domum eius, propter Sarai uxorem Abram.
\begin{responsory-doxology}[ascendens-ergo]
{Ascendens ergo deus ab Abraham pluit ignem et sulphur super Sodomam, * Abraham mane consurgens stet, et eversas urbes a longe conspexit.}{Recordatus est deus Abraham et liberavit Loth de subversione Sodome.}
\end{responsory-doxology}%
\newline {\color{Red} Lectio VII. Secundum Lucam.}
\lettrine[lines=2]{\zallmancaps \color{Red} I}{n} illo tempore: Assumpsit Ihesus duodecim discipulos suos: et ait illis. Ecce ascendimus Iherosolimam, et consummabuntur omnia que scripta sunt per prophetas de filio hominis. Et reliqua.
\newline {\color{Red} Omelia Beati Gregorii Pape.}
\lettrine[lines=2]{\zallmancaps \color{Blue} R}{edemptor} noster previdens ex passione sua discipulorum animos perturbandos, eis longe ante, et eiusdem passionis penam et resurrectionis sue gloriam predicit, ut cum eum morientem sicut predictum est cernerent, eciam resurrecturum non dubitarent.
\begin{responsory}[temptavit-deus]
{Temptavit deus Abraham et dixit ad eum, * Tolle filium tuum quem diligis Ysaac, et offeres illum ibi in holocaustum super unum montium quem dixero tibi.}{Vocatus quoque a domino respondit assum, et ait ei dominus.}
\end{responsory}%
\newline {\color{Red} Lectio VIII.}
\lettrine[lines=2]{\zallmancaps \color{Red} S}{ed} quia carnales adhuc discipuli nullo modo valebant capere verba mysterii, venitur ad miraculum. Ante eorum oculos cecus lumen recipit, ut qui celestis mysterii verba non caperent eos ad fidem celestia facta solidarent. Sed miracula Domini et salvatoris nostri sic accipienda sunt fratres mei, ut et in veritate credantur facta, et tamen per significationem nobis aliquid innuant. Opera quippe eius et per potentiam aliud ostendunt, et per mysterium aliud loquuntur.
\begin{responsory}[angelus-domini-vocavit]
{Angelus domini vocavit Abraham dicens, * Ne extendas manum tuam super puerum eo quod timeas dominum.}{Cumque extendisset manum ut immolaret filium, ecce angelus domini de celo clamavit dicens.}
\end{responsory}%
\newline {\color{Red} Lectio IX.}
\lettrine[lines=2]{\zallmancaps \color{Blue} E}{cce} quis iuxta hystoriam cecus iste fuerit ignoramus: sed tamen quem per mysterium significet, novimus. Cecus quippe est genus humanum quod in parente primo a paradisi gaudiis expulsum, claritatem superne lucis ignorans, damnationis sue tenebras patitur, sed tamen per redemptoris sui presentiam illuminatur, ut interne lucis gaudia iam per desiderium videat, atque in via vite boni operis gressus ponat.
\begin{responsory-final}[deus-domini-mei]
{Deus domini mei Abraham dirige viam meam, * Ut cum salute revertar, * In domum domini mei.}{Obsecro domine fac misericordiam cum servo tuo.}{In domum}
\end{responsory-final}%
\newline {\color{Red} In Laudibus. Ant.} Averte domine faciem tuam a peccatis meis, et omnes iniquitates meas dele. {\color{Red} Ps.} \psalm{50} {\color{Red} Ant.} Fortitudo mea et laus mea dominus, factus est michi in salutem. {\color{Red} Ps.} \psalm{117} {\color{Red} Ant.} In matutinis domine meditabor in te quia factus es adiutor meus. {\color{Red} Ps.} \psalm{62} {\color{Red} Ant.} Benedicamus patrem et filium in secula cum sancto spiritu. {\color{Red} Can.} \psalm{benedicite} {\color{Red} Ant.} Iuvenes et virgines senes cum iunioribus laudent nomen domini. {\color{Red} Ps.} \psalm{148} \V Domine refugium. {\color{Red} Ad Ben. Ant.} Ecce ascendimus Hierosolimam, et consummabuntur omnia que scripta sunt de filio hominis tradetur enim gentibus et illudetur et flagellabitur et conspuetur, et postquam flagellaverint occident eum et die tercia resurget.
\newline {\color{Red} Ad Primam. Ant.} Iter faciente Ihesu dum appropinquaret Hiericho cecus clamabat ad eum ut lumen recipere mereretur. {\color{Red} Ps.} \psalm{21}
\newline {\color{Red} Ad Terciam. Ant.} Transeunte domino clamabat cecus ad eum, miserere mei fili David. \R \hyperlink{spes-mea}{Spes mea.}
\newline {\color{Red} Ad Sextam. Ant.} Cecus sedebat secus viam et clamabat miserere mei fili David. \R \hyperlink{esto-nobis}{Esto nobis.}
\newline {\color{Red} Ad Nonam. Ant.} Cecus magis ac magis clamabat ut eum dominus illuminaret. \R \hyperlink{educ-de-carcere}{Educ de carcere.}
\newline {\color{Red} Ad Mag. Ant.} Stans autem Ihesus iussit cecum adduci ad se, et ait illi quid vis ut faciam tibi, domine ut videam, et Ihesus ait illi, respice fides tua te salvum fecit et confestim vidit et sequebatur illum magnificans deum.
\newline \ul{Lectiones feriales a terciodecimo capitulo ad sextumdecimum ad libitum.}
\newline \ul{Feria secunda et tercia quando de tempore agitur.}
\newline {\color{Red} Ad Ben. Ant.} Tradetur enim gentibus ad illudendum et flagellandum et crucifigendum.
\newline {\color{Red} Ad Mag. Ant.} Omnis plebs ut vidit dedit laudem deo.
{\ubsubsection{Feria Quarta In Capite Ieiunii}{breviarium-feria-quarta-in-capite-ieiunii}}
\par \noindent {\color{Red} Ad Matutinas. Lectio I. Secundum Matheum.}
\lettrine[lines=2]{\zallmancaps \color{Red} I}{n} illo tempore: Dixit Ihesus discipulis suis. Cum ieiunatis nolite fieri sicut ypocrite tristes. Et reliqua.
\newline {\color{Red} Omelia Beati Maximi Episcopi.}
% https://la.wikisource.org/wiki/Homiliae_(Maximus)/1
\lettrine[lines=2]{\zallmancaps \color{Blue} Q}{uia} nonnullorum est consuetudo advenientes quadragesime dies devotionis ieiunio prevenire, necessarie presens evangelii decursa est lectio, in qua Dominus noster virtutum spiritualium retributor sacratam nobis perfectamque dedit regulam ieiunandi dicens. Cum ieiunatis, nolite fieri sicut ypocrite tristes. Exterminant enim facies suas, ut appareant hominibus ieiunantes.
\newline {\color{Red} Lectio Secunda.}
\lettrine[lines=2]{\zallmancaps \color{Red} N}{on} enim Deo ieiunat sed hominibus, quicunque ostentatione sui ieiunii gloriam requirit humanam. Exterminant autem facies suas, qui merore simulato religiosum vultum populorum oculis mentiuntur. Non vult enim discipline celestis magister irreligiose agi, quod fieri pro religione precipitur, nec patitur invocantium se laborem quibus eternam parat mercedem: infructuose iactantie vitio deperire.
\newline {\color{Red} Lectio Tertia.}
\lettrine[lines=2]{\zallmancaps \color{Blue} N}{on} eritis inquit sicut ypocrite tristes. Si enim vere tristis et mestus es quia ieiunas, nullam tibi ut ingratus: apud Deum gratiam reservabis, quia quamvis opus facias bonum: pravitatem tamen degeneris animi invitus operaris. Si autem pro quadam sanctitatis ymagine tristiciam simulas: omnem fructum promissionis divine obsequii tui ostentator amittis.
\newline {\color{Red} In Laudibus. Cap.}
\lettrine[lines=2]{\zallmancaps \color{Red} C}{onvertimini} ad me in toto corde vestro, in ieiunio et fletu et planctu, et scindite corda vestra et non vestimenta vestra, ait dominus omnipotens. {\color{Red} Ad Ben. Ant.} Cum ieiunatis nolite fieri sicut ypocrite tristes. {\color{Red} Oracio.}
\lettrine[lines=2]{\zallmancaps \color{Blue} P}{resta} domine fidelibus tuis, ut ieiuniorum veneranda sollemnia, et congrua pietate suscipiant, et secura devotione percurrant. Per.
\newline \ul{Hac die et deinceps usque ad Cenam Domini exclusive orationes que in singulis feriis sunt notate ad Laudes, dicantur ad Terciam Sextam et Nonam quando de tempore agitur, ad Vesperas vero dicantur proprie que notate sunt per singulas ferias nisi in Sabbatis, in quibus semper dicitur Oratio sequentis Dominice.}
\newline {\color{Red} Ad Terciam. Cap.} Convertimini.
\newline {\color{Red} Ad Sextam. Cap.}
\lettrine[lines=2]{\zallmancaps \color{Red} D}{erelinquat} impius viam suam, et vir iniquus cogitationes suas, et revertatur ad dominum et miserebitur eius, et ad deum nostrum quoniam multus est ad ignoscendum.
\newline {\color{Red} Ad Nonam. Capitulum.}
\lettrine[lines=2]{\zallmancaps \color{Blue} F}{range} esurienti panem tuum, et egenos vagosque induc in domum tuam, cum videris nudum operi eum, et carnem tuam ne despexeris, ait dominus omnipotens.
\newline {\color{Red} Ad Vesperas. Capitulum.} Convertimini. {\color{Red} Ad Mag. Ant.} Thesaurizate vobis thesauros in celo ubi nec erugo nec tinea demolitur. {\color{Red} Oratio.}
\lettrine[lines=2]{\zallmancaps \color{Red} I}{nclinantes} se domine maiestati tue propitiatus intende ut qui divino munere sunt refecti, celestibus semper nutriantur auxiliis. Per.
\newline \ul{Predicta tria capitula dicantur cotidie ad horas in quibus notata sunt, diebus profestis usque ad Dominicam in Passione. Hac die et deinceps cotidie per totam Quadragesimam dicantur Vespere ante prandium nisi in Dominicis diebus.}
\newline {\color{Red} Feria Quinta. Ad Ben. Ant.} Domine puer meus iacet paraliticus in domo, et male torquetur, amen dico tibi ego veniam et curabo eum. {\color{Red} Oratio.}
\lettrine[lines=2]{\zallmancaps \color{Blue} D}{eus} qui culpa offenderis penitentia placaris, preces populi tui supplicantis propicius respice, et flagella tue iracundie que pro peccatis nostris meremur averte. Per. {\color{Red} Ad Mag. Ant.} Domine non sum dignus ut intres sub tectum meum, sed tantum dic verbo et sanabitur puer meus. {\color{Red} Oratio.}
\lettrine[lines=2]{\zallmancaps \color{Red} P}{arce} domine parce populo tuo, ut dignis flagellationibus castigatus in tua miseratione respiret. Per.
\newline {\color{Red} Feria Sexta. Ad Ben. Ant.} Cum facis elemosinam nesciat sinistra tua quid faciat dextera tua. {\color{Red} Oratio.}
\lettrine[lines=2]{\zallmancaps \color{Blue} I}{nchoata} ieiunia quesumus domine benigno favore prosequere, ut observantiam quam corporaliter exhibemus, mentibus eciam sinceris exercere valeamus. Per. {\color{Red} Ad Mag. Ant.} Tu autem cum oraveris intra in cubiculum et clauso ostio ora patrem tuum. {\color{Red} Oratio.}
\lettrine[lines=2]{\zallmancaps \color{Blue} T}{uere} domine populum tuum, et ab omnibus peccatis clementer emunda, quia nulla ei nocebit adversitas si nulla dominetur iniquitas. Per.
\newline {\color{Red} Sabbato. Ad Ben. Ant.} Quare ieiunavimus et non aspexisti humiliavimus animas nostras et nescisti. {\color{Red} Oratio.}
\lettrine[lines=2]{\zallmancaps \color{Blue} A}{desto} domine supplicationibus nostris, et concede ut hoc sollemne ieiunium quod animis corporibusque curandis salubriter institutum est, devoto servitio celebremus. Per.
{\ubsubsection{Dominica Prima in Quadragesima Sabbato Precedenti}{breviarium-dominica-prima-in-quadragesima-sabbato-precedenti}}
\par \noindent {\color{Red} Ad Vesperas. Capitulum.}
\lettrine[lines=2]{\zallmancaps \color{Red} H}{ortamur} \hypertarget{hortamur-vos}{\label{hortamur-vos}} vos ne in vacuum gratiam dei recipiatis, ait enim tempore accepto exaudivi te, et in die salutis adiuvi te. \R \hyperlink{emendemus-in-melius}{Emendemus.} {\color{Red} III. Ymnus.}
\hymn{audi-benigne-conditor}{Audi benigne conditor}{
\lettrine[lines=2]{\zallmancaps \color{Blue} A}{udi} benigne conditor
nostras preces cum fletibus
in hoc sacro ieiunio
fusas quadragenario.
\par {\bfseries \color{Red} S}crutator alme cordium
infirma tu scis virium,
ad te reversis exhibe
remissionis gratiam.
\par {\bfseries \color{Blue} M}ultum quidem peccavimus,
sed parce confitentibus
ad laudem tui nominis
confer medelam languidis.
\par {\bfseries \color{Red} S}ic corpus extra conteri
dona per abstinentiam,
ieiunet ut mens sobria
a labe prorsus criminum.
\par {\bfseries \color{Blue} P}resta beata trinitas
concede simplex unitas,
ut fructuosa sint tuis
ieiuniorum munera. Amen.
}
\newline \V Angelis suis deus mandavit de te. \R Ut custodiant te in omnibus viis tuis.
\newline {\color{Red} Ad Mag. Ant.} Ecce nunc tempus acceptabile ecce nunc dies salutis in his ergo diebus exhibeamus nos sicut dei ministros in multa patientia, in vigiliis, in ieiuniis, et in caritate non ficta. {\color{Red} Oratio.}
\lettrine[lines=2]{\zallmancaps \color{Red} D}{eus} qui ecclesiam tuam annua quadragesimali observatione purificas, presta familie tue, ut quod a te obtinere abstinendo nititur, hoc bonis operibus exequatur. Per.
\newline {\color{Red} Ad Completorium.}
\begin{responsory-breve}[in-pace]
{In pace in idipsum * Dormiam et requiescam.}{Si dedero somnum oculis meis et palpebris meis dormitationem.}
\end{responsory-breve}%
\newline \ul{Hoc responsorium dicatur tantum in Sabbatis et Dominicis diebus et festivis in utroque Completorio, usque ad Sabbatum ante Dominicam in Passione exclusive. Ceteris autem diebus dicatur responsorium.} \hyperlink{in-manus-tuas-breve}{In manus.} {\color{Red} Ymnus.}
\hymn{christe-qui-lux}{Christe qui lux es}{
\lettrine[lines=2]{\zallmancaps \color{Blue} C}{hriste} qui lux es et dies,
noctis tenebras detegis,
lucisque lumen crederis,
lumen beatum predicans.
\par {\bfseries \color{Red} P}recamur sancte domine,
defende nos in hac nocte,
sit nobis in te requies,
quietam noctem tribue.
\par {\bfseries \color{Blue} N}e gravis somnus irruat,
nec hostis nos subripiat,
nec caro illi consentiens
nos tibi reos statuat.
\par {\bfseries \color{Red} O}culi somnum capiant
cor ad te semper vigilet,
dextera tua protegat
famulos qui te diligunt.
\par {\bfseries \color{Blue} D}efensor noster aspice
insidiantes reprime,
guberna tuos famulos
quos sanguine mercatus es.
\par {\bfseries \color{Red} M}emento nostri domine
in gravi isto corpore,
qui es defensor anime
ad esto nobis domine.
\par {\bfseries \color{Blue} P}resta pater omnipotens
per Ihesum Christum dominum
qui tecum in perpetuum
regnat cum sancto spiritu. Amen.
}
\newline \ul{Predictus ymnus dicatur cotidie usque ad diem Cene exclusive.}
\newline {\color{Red} Ad Nunc Dimittis. Ant.} Evigila super nos eterne salvator ne nos apprehendat callidus temptator quia tu factus es nobis sempiternus adiutor.
\newline \ul{Hec Ant. dicatur super} \psalm{nunc} \ul{per quindecim dies, nisi in Annunciatione quando infra hoc tempus evenerit.}
\newline {\color{Red} Ad Matutinas. Invitat.}
\begin{invitatory}[non-sit-invitatorium]
{Non sit vobis vanum mane surgere ante lucem, * Quia promisit dominus coronam vigilantibus.}
\end{invitatory}
{\color{Red} Ymnus.}
\hymn{summi-largitor}{Summi largitor premii}{
\lettrine[lines=2]{\zallmancaps \color{Red} S}{ummi} largitor premii
spes qui es unica mundi,
preces intende servorum
ad te devote clamantum.
\par {\bfseries \color{Red} N}oster te conscientia
grave offendisse monstrat,
quam emundes supplicamus
ab omnibus piaculis.
\par {\bfseries \color{Blue} S}i renuis quis tribuet
indulge quia potens es,
te corde rogare mundo
fac nos precamur domine.
\par {\bfseries \color{Red} E}rgo nunc accepta nostrum
qui sacrasti ieiunium,
quo mistice paschalia
capiamus sacramenta.
\par {\bfseries \color{Blue} S}umma nobis hoc conserat
in deitate trinitas,
in qua gloriatur unus,
per cuncta secula deus. Amen.
}
\newline {\color{Red} In Primo Nocturno.} \V Dicet domino susceptor meus es tu. \R Et refugium meum deus meus.
\newline {\color{Red} Gregorii in Omelia super Evangelium.} Ductus est Ihesus in desertum.
\newline {\color{Red} Lectio Prima.}
\lettrine[lines=2]{\zallmancaps \color{Blue} Q}{uia} quadragesime tempus inchoamus, discutiendum est cur hec abstinentia custoditur. Moyses enim ut legem acciperet secundo: Helyas eciam in deserto, quadraginta diebus abstinuit. Ipse quoque actor hominum, quadraginta diebus nullum cibum sumpsit. Nos quoque annuo quadragesime tempore, carnem nostram per abstinentiam affligere conamur. Cur ergo in abstinentia hic numerus custoditur, nisi quia virtus Decalogi per libros quatuor sancti evangelii impletur? Denarius enim quater ductus in quadragenarium surgit. Quia tunc mandata Decalogi perficimus, cum quatuor libros sancti evangelii custodimus.
{\color{Red} Responsorium.}
\lettrine[lines=2]{\zallmancaps \color{Red} E}{cce} \hypertarget{ecce-nunc}{\label{ecce-nunc}} nunc tempus acceptabile ecce nunc dies salutis, commendemus nosmetipsos in multa patientia in ieiuniis multis, * Per arma iusticie virtutis dei. \V In omnibus exhibeamus nosmetipsos sicut dei ministris in multa patientia. Per arma.
\newline {\color{Red} Lectio Secunda.}
\lettrine[lines=2]{\zallmancaps \color{Blue} I}{n} hoc etenim mortali corpore ex quatuor elementis subsistimus, et per voluptates eiusdem corporis preceptis dominicis contraimus, que per Decalogum sunt accepta. Qui ergo per carnis desideria Decalogi mandata contempsimus: dignum est ut eandem carnem quaterdecies affligamus. A presenti eciam die usque ad Paschalis sollemnitatis gaudia sex ebdomade veniunt quarum dies quadraginta duo fiunt, ex quibus dum sex dies Dominici ab abstinentie subtrahuntur, non plus
%pp099
in abstinentia quam triginta et sex dies remanent. Dum vero per trecentos et sexaginta sex dies annus ducitur, nos autem per triginta et sex dies affligimur quasi anni nostri decimas Deo damus, ut qui nobis per annum viximus, actori nostro in eius decimis mortificemus.
\begin{responsory}[in-omnibus]
{In omnibus exhibeamus nos sicut dei ministros in multa patientia, * Ut non vituperetur ministerium nostrum.}{Ecce nunc tempus acceptabile ecce nunc dies salutis nemini dantes ullam offensionem.}
\end{responsory}%
\newline {\color{Red} Lectio Tertia.}
\lettrine[lines=2]{\zallmancaps \color{Red} S}{icut} igitur in lege fratres offere iubemini decimas rerum, ita offerre contendite et decimas dierum. Unusquisque carnem maceret, eiusque desideria affligat, ut hostia viva fiat. Hostia quippe et immolatur et viva est, quando et ab hac vita homo non deficit, et tamen se carnalibus desideriis occidit. Caro nos leta traxit ad culpam afflicta reducat ad veniam. A paradisi gaudiis per cibum cecidimus, ad hec per abstinentiam resurgamus.
\begin{responsory-doxology}[emendemus-in-melius]
{Emendemus in melius que ignorantur peccavimus ne subito preoccupati die mortis queramus spatium penitentie et invenire non possimus, * Attende domine et miserere quia peccavimus tibi.}{Peccavimus cum patribus nostris iniuste egimus iniquitatem fecimus.}
\end{responsory-doxology}%
\newline {\color{Red} In Secundo Nocturno.} \V Ipse liberavit me de laqueo venantium. \R Et a verbo aspero.
\newline {\color{Red} Lectio Quarta.}
\lettrine[lines=2]{\zallmancaps \color{Blue} N}{emo} sibi abstinentiam solam credat posse sufficere, cum dominus per prophetam dicat. Nonne hoc est ieiunium quod elegi? Frange esurienti panem tuum. Illud ergo ieiunium Deus approbat, quod ad eius oculos manus elemosinarum levat, quod cum proximi dilectione agitur, quod ex pietate conditur.
\begin{responsory}[paradisi-portas]
{Paradisi portas aperiat nobis ieiunii tempus suscipiamus illud orantes et deprecantes, * Ut in die resurrectionis cum domino gloriemur.}{Ecce nunc tempus acceptabile ecce nunc dies salutis, commendemus nosmetipsos in multa patientia.}
\end{responsory}%
\newline {\color{Red} Lectio Quinta.}
\lettrine[lines=2]{\zallmancaps \color{Red} H}{oc} ergo quod tibi subtrahis, alteri largire, et unde caro tua affligitur, inde egentis proximi caro recreetur. Hinc est quod per prophetam Dominus dicit. Cum ieiunaretis nunquid ieiunium ieiunastis michi? Sibi quippe quisque ieiunat, si ea que sibi ad tempus subtrahit, non pauperibus tribuit, sed ventri postmodum offerenda custodit. Hinc per Ioel dicitur. Sanctificate ieiunium. Ieiunium quippe sanctificare est, adiunctis aliis bonis dignam Deo abstinentiam carnis ostendere.
\begin{responsory}[scindite-corda]
{Scindite corda vestra et non vestimenta vestra, * Et convertimini ad dominum deum quia benignus et misericors est.}{Revertimini unusquisque a via sua mala et a pessimis cogitationibus vestris.}
\end{responsory}%
\newline {\color{Red} Lectio Sexta.}
\lettrine[lines=2]{\zallmancaps \color{Blue} C}{esset} ira, sopiantur iurgia. In cassum enim caro atteritur si a suis voluptatibus animus non refrenatur, cum per prophetam Dominus dicat. Ecce in die ieiunii vestri invenitur voluntas vestra, et omnes debitores vestros repetitis, neque enim qui debitorem se in penitentia macerat, eciam hoc quod sibi iuste competit interdicat. Sic sic nobis afflictis a Deo dimittitur quod iniuste egimus, si pro amore illius hoc quod iuste competit relaxemus.
\begin{responsory-doxology}[abscondite-elemosinam]
{Abscondite elemosinam in sinu pauperis, et ipsa orat pro vobis ad dominum, * Quia sicut aqua extinguit ignem, ita elemosina extinguit peccatum.}{Date elemosinam dicit dominus et ecce omnia munda sunt vobis.}
\end{responsory-doxology}%
\newline {\color{Red} In Tertio Nocturno.} \V Scapulis suis obumbrabit tibi. \R Et sub pennis eius sperabis.
\newline {\color{Red} Lectio VII. Secundum Matheum.}
\lettrine[lines=2]{\zallmancaps \color{Red} I}{n} illo tempore: Ductus est Ihesus in desertum a spiritu ut temptaretur a diabolo. Et reliqua.
\newline {\color{Red} Omelia Beati Gregorii Pape.}
\lettrine[lines=2]{\zallmancaps \color{Blue} D}{ubitari} a quibusdam solet, a quo spiritu ductus sit Ihesus in desertum, propter hoc quod subditur, assumpsit eum diabolus in sanctam civitatem, et rursus, assumpsit eum in montem excelsum. Sed vere et absque ulla questione convenienter accipitur, ut a spiritu sancto in desertum ductus credatur, ut illuc eum suus spiritus duceret, ubi nunc ad temptandum malignus spiritus inveniret.
\begin{responsory}[in-ieiunio]
{In ieiunio et fletu orabant sacerdotes dicentes, * Parce domine parce populo tuo et ne des hereditatem tuam in perditionem.}{Inter vestibulum et altare plorabant sacerdotes dicentes,}
\end{responsory}%
\newline {\color{Red} Lectio VIII.}
\lettrine[lines=2]{\zallmancaps \color{Red} S}{ed} ecce cum dicitur Deus homo vel in excelsum montem vel in sanctam civitatem a diabolo assumptus, mens refugit, humane hoc audire aures expavescunt. Qui tamen non esse incredibilia ista cognoscimus, si in illo et alia facta pensamus.
\begin{responsory}[tribularer-si]
{Tribularer si nescirem misericordias tuas domine tu dixisti, nolo mortem peccatoris, sed convertatur et vivat, * Qui Chananeam et publicanum vocasti ad penitentiam.}{Et Petrum lacrimantem suscepisti misericors domine.}
\end{responsory}%
\newline {\color{Red} Lectio IX.}
\lettrine[lines=2]{\zallmancaps \color{Blue} C}{erte} iniquorum omnium diabolus caput est, et huius capitis membra sunt omnes iniqui. An non diaboli membrum fuit Pilatus? An non diaboli membra Iudei persequentes et milites crucifigentes fuerunt? Quid ergo mirum, si se ab illo permisit in montem duci, qui se perculit eciam a membris illius crucifigi? Non est indignum redemptori nostro quod temptari voluit, qui venerat occidi.
\begin{responsory-final}[ductus-est]
{Ductus est Ihesus in desertum a spiritu ut temptaretur a diabolo, * Et accedens temptator dixit ei, * Si filius dei es, dic ut lapides isti panes fiant.}{Et cum ieiunasset quadraginta diebus et quadraginta noctibus postea esuriit.}{Si filius}
\end{responsory-final}%
{Ductus est.} \V In manibus portabunt te. \R Ne forte offendas ad lapidem pedem tuum.
\newline {\color{Red} In Laudibus. Ant.} Cor mundum crea in me deus et spiritum rectum innova in visceribus meis. {\color{Red} Ps.} \psalm{50} {\color{Red} Ant.} O domine salvum me fac, o domine bene prosperare. {\color{Red} Ps.} \psalm{117} {\color{Red} Ant.} Sic benedicam te in vita mea domine, et in nomine tuo levabo manus meas. {\color{Red} Ps.} \psalm{62} {\color{Red} Ant.} In spiritu humilitatis et in animo contrito suscipiamur domine a te, et sict fiat sacrificium nostrum ut a te suscipiatur hodie, et placeat tibi domine deus. {\color{Red} Can.} \psalm{benedicite} {\color{Red} Ant.} Laudate deum celi celorum et aque omnes. {\color{Red} Ps.} \psalm{148} {\color{Red} Cap.} \hyperlink{hortamur-vos}{Hortamur.} {\color{Red} Ymnus.}
\hymn{iam-christe}{Iam Christe sol iusticie}{
\lettrine[lines=2]{\zallmancaps \color{Red} I}{am} Christe sol iusticie
mentis dehiscant tenebre,
virtutum ut lux redeat
terris diem cum reparas.
\par {\bfseries \color{Blue} D}a tempus acceptabile
et penitens cor tribue,
convertat ut benignitas
quos longa suffert pietas.
\par {\bfseries \color{Red} Q}uiddamque penitentie
da ferre quamvis grandium
maiori tuo munere
quo demptio sit criminum.
\par {\bfseries \color{Blue} D}ies venit dies tua,
in qua relorent omnia,
letemur in hac ut tue
per hanc reducti gratie.
\par {\bfseries \color{Red} T}e rerum universitas
clemens adoret trinitas,
et nos novi per veniam
novum canamus canticum. Amen.
}
\newline \V Scuto circumdabit te veritas eius. \R Non timebis a timore nocturno.
\newline {\color{Red} Ad Ben. Ant.} Ductus est Ihesus in desertum a spiritu ut temptaretur a diabolo, et cum ieiunasset quadraginta diebus et quadraginta noctibus postea esuriit.
\newline {\color{Red} Ad Primam. Ant.} Ihesus autem cum ieiunasset quadraginta diebus et quadraginta noctibus postea esuriit.
\newline {\color{Red} Ad Terciam. Ant.} Non in solo pane vivit homo, sed in omni verbo quod procedit de ore dei. {\color{Red} Cap.} \hyperlink{hortamur-vos}{Hortamur.}
\begin{responsory-breve}[participem-me]
{Participem me fac deus omnium timentium te, * Et custodientium mandata tua.}{Aspice in me et miserere mei secundum iudicium diligentium nomen tuum.}
\end{responsory-breve}%
\V Dicet domino.
\newline {\color{Red} Ad Sextam. Ant.} Tunc assumpsit eum diabolus in sanctam civitatem, et statuit eum supra pinnaculum templi, et dixit ei, si filius dei es mitte te deorsum. {\color{Red} Capitulum.}
\lettrine[lines=2]{\zallmancaps \color{Blue} E}{cce} nunc tempus acceptabile, ecce nunc dies salutis, nemini dantes ullam offensionem ut non vituperetur ministerium nostrum.
\begin{responsory-breve}[ab-omni-via]
{Ab omni via mala prohibui pedes meos, * Ut custodiam mandata tua domine.}{A iudiciis tuis non declinavi quia tu legem posuisti michi.}
\end{responsory-breve}%
\V Ipse liberavit me.
\newline {\color{Red} Ad Nonam. Ant.} Vade Sathana non temptabis dominum deum tuum. {\color{Red} Capitulum.}
\lettrine[lines=2]{\zallmancaps \color{Red} I}{n} omnibus exhibeamus nosmetipsos sicut dei ministros in multa patientia in ieiuniis multis per arma iusticie.
\begin{responsory-breve}[septies-in-die]
{Septies in die laudem dixi tibi domine deus meus, * Ne perdas me.}{Erravi sicut ovis que perierat, require servum tuum domine quia mandata tua non sum oblitus.}
\end{responsory-breve}%
\V Scapulis suis.
\newline \ul{Predicta responsoria dicantur ad horas per quindecim dies, quando de tempore agitur.}
\newline {\color{Red} Ad Vesperas. Capitulum.} \hyperlink{hortamur-vos}{Hortamur.}
\newline \ul{Predicta tria Capitula dicantur ad horas in quibus notata sunt in Dominicis diebus usque ad Vesperas Sabbati ante Dominicam in Passione exclusive.}
\newline {\color{Red} Ymnus.} \hyperlink{audi-benigne-conditor}{Audi benigne.} \V Angelis suis.
\newline \ul{Predicta Ymni et Versiculi ad Matutinas et Vesperas, et ad alias horas in quibus notati sunt dicantur usque ad Vesperas Sabbati ante Dominicam in Passione exclusive, quando de tempore agitur.}
\newline {\color{Red} Ad Mag. Ant.} Reliquit eum temptator et accesserunt angeli et ministrabant ei.
\newline \ul{Lectiones supra posite de sermonibus in Dominica repetantur per ferias.}
\newline {\color{Red} Feria Secunda. Ad Benedictus. Antiphona.} Venite benedicti patris mei percipite regnum quod vobis paratum est ab origine mundi. {\color{Red} Oratio.}
\lettrine[lines=2]{\zallmancaps \color{Blue} C}{onverte} nos deus salutaris noster, et ut ieiunium nobis quadragesimale proficiat, mentes nostras celestibus instrue disciplinis. Per.
\newline {\color{Red} Ad Primam. Ant.} Vivo ego dicit dominus, nolo mortem peccatoris, sed ut magis convertatur et vivat.
\newline {\color{Red} Ad Terciam. Ant.} Advenerunt nobis dies penitentie ad redimenda peccata et salvandas animas.
\newline {\color{Red} Ad Sextam. Ant.} Commendemus nosmetipsos in multa patientia in ieiuniis multis per arma iusticie.
\newline {\color{Red} Ad Nonam. Ant.} Per arma iusticie virtutis dei commendemus nosmetipsos in multa patiencia.
\newline \ul{Antiphone predicte ad horas dicantur cotidie in profestis diebus usque ad Passionem Domini.}
\newline {\color{Red} Ad Mag. Ant.} Quod uni ex minimis meis fecistis michi fecistis dicit dominus. {\color{Red} Oratio.}
\lettrine[lines=2]{\zallmancaps \color{Red} A}{bsolve} quesumus domine nostrorum vincula peccatorum, et quicquid pro eis meremur propitiatus averte. Per.
\newline {\color{Red} Feria Tercia. Ad Ben. Ant.} Intravit Ihesus in templum dei, et eiciebat omnes vendentes et ementes, et mensas nummulariorum, cathedras vendentium columbas evertit. {\color{Red} Oratio.}
\lettrine[lines=2]{\zallmancaps \color{Blue} R}{espice} domine familiam tuam, et presta ut aput te mens nostra tuo desiderio fulgeat, que se carnis maceratione castigat. Per.
\newline {\color{Red} Ad Mag. Ant.} Abiit Ihesus foras extra civitatem in Bethaniam ibique docebat eos de regno dei. {\color{Red} Oratio.}
\lettrine[lines=2]{\zallmancaps \color{Red} A}{scendant} ad te domine preces nostre, et ab ecclesia tua cunctam repelle nequitiam. Per.
\newline {\color{Red} Feria Quarta. Ad Ben. Ant.} Generatio hec prava et perversa signum querit et signum non dabitur ei nisi signum Ione prophete. {\color{Red} Oratio.}
\lettrine[lines=2]{\zallmancaps \color{Blue} P}{reces} nostras quesumus domine clementer exaudi, et contra cuncta nobis adversantia, dexteram tue maiestatis extende. Per.
\newline {\color{Red} Ad Mag. Ant.} Sicut fuit Ionas in ventre ceti tribus diebus et tribus noctibus ita erit filius hominis in corde terre. {\color{Red} Oratio.}
\lettrine[lines=2]{\zallmancaps \color{Red} M}{entes} nostras quesumus domine lumine tue claritatis illustra, ut videre possimus que agenda sunt, et que recta sunt agere valeamus. Per.
\newline {\color{Red} Feria Quinta. Ad Ben. Ant.} Si vos manseritis in sermone meo vere discipuli mei eritis, et cognoscetis veritatem, et veritas liberabit vos. {\color{Red} Oratio.}
\lettrine[lines=2]{\zallmancaps \color{Blue} D}{evotionem} populi tui quesumus domine benignus intende, ut qui per abstinentiam macerantur in corpore, per fructum boni operis reficiantur in mente. Per.
\newline {\color{Red} Ad Mag. Ant.} Ego enim ex deo processi et veni, neque enim a me ipso veni, sed pater meus ille me misit. {\color{Red} Oratio.}
\lettrine[lines=2]{\zallmancaps \color{Red} D}{a} quesumus domine populis Christianis et que profitentur agnoscere, et celeste munus diligere quod frequentant. Per.
\newline {\color{Red} Feria Sexta. Ad Ben. Ant.} Angelus domini descendebat de celo, movebatur aqua et sanabatur unus. {\color{Red} Oratio.}
\lettrine[lines=2]{\zallmancaps \color{Blue} E}{sto} domine propitius plebi tue, et quam tibi facis esse devotam benigno refove miseratus auxilio. Per.
\newline {\color{Red} Ad Mag. Ant.} Qui me sanum fecit ille michi precepit tolle grabatum tuum et ambula in pace. {\color{Red} Oratio.}
\lettrine[lines=2]{\zallmancaps \color{Red} E}{xaudi} nos misericors deus et mentibus nostris gratie tue lumen ostende. Per.
\newline {\color{Red} Sabbato. Ad Ben. Ant.} Assumpsit Ihesus discipulos suos et ascendit in montem et transfiguratus est ante eos. {\color{Red} Oratio.}
\lettrine[lines=2]{\zallmancaps \color{Blue} P}{opulum} tuum quesumus domine propitius respice, atque ab eo flagella tue iracundie clementer averte. Per.
{\ubsubsection{Dominica Secunda in Quadragesima Sabbato Precedenti}{breviarium-dominica-secunda-in-quadragesima-sabbato-precedenti}}
\par \noindent {\color{Red} Ad Vesperas.} \R \hyperlink{det-tibi-deus}{Det tibi deus.} {\color{Red} III. Ad Mag. Ant.} Domine bonum est nos hic esse, si vis faciamus hic tria tabernacula, tibi unum, Moysi unum, et Helye unum. {\color{Red} Oratio.}
\lettrine[lines=2]{\zallmancaps \color{Red} D}{eus} qui conspicis omni nos virtute destitui, interius exteriusque custodi, ut et ab omnibus adversitatibus muniamur in corpore, et a pravis cogitationibus mundemur in mente. Per.
\newline {\color{Red} Ad Matutinas. Invitat.}
\begin{invitatory}[quoniam-deus-invitatorium]
{Quoniam deus magnus dominus, * Et rex magnus super omnes deos.}
\end{invitatory}
\newline \ul{Hic versus in Psalmo intermittatur postquam secunda vice dictum fuerit Invitatorium.}
\newline {\color{Red} Lectio Prima.}
\lettrine[lines=2]{\zallmancaps \color{Blue} S}{enuit} Ysaac et caligaverunt oculi eius, et videre non poterat. Vocavitque Esau filium suum maiorem, et dixit ad eum. Fili mi. Qui respondit. Assum. Cui pater. Vides inquit quod senuerim, et ignorem diem mortis mee? Sume arma tua, pharetram, et arcum, et egredere foras. Cumque venatu aliquid apprehenderis, fac michi inde pulmentum sicut velle me nosti, et affer ut comedam, et benedicat tibi anima mea antequam moriar. \R
\lettrine[lines=2]{\zallmancaps \color{Red} T}{olle} \hypertarget{tolle-arma}{\label{tolle-arma}} arma tua pharetram et arcum et affer de venatione tua ut comedam, * Et benedicat tibi anima mea. \V Cumque venatu aliquid attuleris fac michi inde pulmentum ut comedam. Et benedicat.
\newline {\color{Red} Lectio Secunda.}
\lettrine[lines=2]{\zallmancaps \color{Blue} Q}{uod} cum audisset Rebeca, et ille abisset in agrum ut iussionem patris expleret, dixit filio suo Iacob. Audivi patrem tuum loquentem cum Esau fratre tuo, et dicentem ei. Affer michi de venatione tua, et fac cibos ut comedam, et benedicam tibi coram Domino antequam moriar.
\begin{responsory}[ecce-odor-filii]
{Ecce odor filii mei sicut odor agri pleni quem benedixit dominus, crescere te faciat deus meus sicut arenam maris, * Et donet tibi de rore celi benedictionem.}{Qui maledixerit tibi sit ille maledictus, et qui benedixerit tibi benedictionibus repleatur.}
\end{responsory}%
\newline {\color{Red} Lectio Tertia.}
\lettrine[lines=2]{\zallmancaps \color{Red} N}{unc} ergo fili mi adquiesce consiliis meis, et pergens ad gregem affer michi duos edos optimos, ut faciam ex eis escas patri tuo quibus libenter vescitur. Quas cum intuleris et comederit, benedicat tibi priusquam moriatur.
\begin{responsory-doxology}[det-tibi-deus]
{Det tibi deus de rore celi et de pinguedine terre abundantiam, serviant tibi tribus populi, * Esto dominus fratrum tuorum.}{Et incurventur ante te filii matris tue.}
\end{responsory-doxology}%
\newline {\color{Red} Lectio Quarta.}
\lettrine[lines=2]{\zallmancaps \color{Blue} C}{ui} ille respondit. Nosti quod Esau frater meus homo pilosus sit, et ego lenis? Si attrectaverit me pater meus et senserit, timeo ne putet sibi me voluisse illudere, et inducat super me maledictionem pro benedictione.
\begin{responsory}[quis-igitur]%typo: ventionem
{Quis igitur ille est qui dudum captam venationem attulit michi, et comedi ex omnibus priusquam tu venires, * Benedixique ei et erit benedictus.}{Dominum tuum illum constitui, et omnes fratres eius servituti illius subiugavi.}
\end{responsory}%
\newline {\color{Red} Lectio Quinta.}
\lettrine[lines=2]{\zallmancaps \color{Red} A}{d} quem mater. In me sit ait ista maledictio fili mi. Tantum audi vocem meam, et pergens affer que dixi. Abiit et attulit, deditque matri. Paravit illa cibos sicut velle noverat patrem illius: et vestibus Esau valde bonis quas apud se habebat domi induit eum. Pelliculasque hedorum circumdedit manibus, et colli nuda protexit.
\begin{responsory}[dum-exiret]
{Dum exiret Iacob de terra sua, vidit gloriam dei et ait quam terribilis est locus iste, * Non est hic aliud nisi domus dei et porta celi.}{Vere dominus est in loco isto et ego nesciebam.}
\end{responsory}%
\newline {\color{Red} Lectio Sexta.}
\lettrine[lines=2]{\zallmancaps \color{Blue} D}{editque} pulmentum, et panes quos coxerat tradidit. Quibus illatis, dixit. Pater mi. Et ille respondit. Audio. Quis es tu fili mi? Dixitque Iacob. Ego sum Esau primogenitus tuus, feci sicut precepisti michi. Surge, sede, et comede de venatione mea, ut benedicat michi anima tua.
\begin{responsory-final}[si-dominus-deus]
{Si dominus deus meus fuerit mecum in via ista per quam ego ambulo et custodierit me, et dederit michi panem ad edendum et vestimentum quo operiar, et revocaverit me cum salute, * Erit michi dominus in refugium, * Et lapis iste in signum.}{Surgens mane Iacob tulit lapidem et erexit in titulum fundensque oleum desuper ait.}{Et lapis}
\end{responsory-final}%
\newline {\color{Red} Lectio VII. Secundum Matheum.}
\lettrine[lines=2]{\zallmancaps \color{Red} I}{n} illo tempore: Egressus Ihesus secessit in partes Tyri et Sydonis. Et ecce mulier Chananea a finibus illis egressa clamavit dicens ei. Miserere mei domine fili David: filia mea male a demonio vexatur. Et reliqua.
\newline {\color{Red} Omelia lectiones eiusdem.} \grecross
\lettrine[lines=2]{\zallmancaps \color{Blue} E}{t} ecce mulier Chananea. Mira res. Mulier caput peccati, arma diaboli, expulsio paradisi, mater dilecti, corruptio legis antique veniebat ad dominum Ihesum.
\begin{responsory}[erit-michi]
{Erit michi dominus in deum et lapis iste quem erexi in titulum vocabitur domus dei, et de universis que dederis michi, * Decimas et hostias pacificas offeram tibi.}{Si dominus deus meus fuerit mecum in via ista per quam ego ambulo et custodierit me.}
\end{responsory}%
\newline {\color{Red} Lectio VIII.}
\lettrine[lines=2]{\zallmancaps \color{Red} M}{irum} negocium. Iudei fugiunt, Chananea sequitur. Domestici derelinquunt, alienigena heret. Et ecce mulier. Cherubin aspectum eius tremunt in celo, angeli metuunt in excelsis, et mulier non veretur in terris, sed loquitur cum eo. O miranda res. Sursum tremor, deorsum fiducia.
\begin{responsory}[dixit-angelus]
{Dixit angelus ad Iacob, dimitte me aurora est, respondit ei, non dimittam te nisi benedixeris michi, * Et benedixit eum in eodem loco.}{Nequaquam vocaberis Iacob sed Israhel erit nomen tuum.}
\end{responsory}%
\newline {\color{Red} Lectio IX.}
\lettrine[lines=2]{\zallmancaps \color{Blue} P}{ropterea} inquit descendit, propterea carnem assumpsit, et homo factus est: ut et ego audeam ei loqui. Non perrexit ad divinos, non vocavit magos, non alligaturas quesivit, non impostatrices mulieres adduxit, non ad falsa confugit presidia, sed reliquit omnes diaboli falsitates, et venit ad dominum Ihesum omnium salvatorem.
\begin{responsory-doxology}[oravit-iacob]
{Oravit Iacob et dixit domine qui dixisti michi, revertere in terram nativitatis tue, * Erue me de manu fratris mei, quia valde eum timeo.}{Deus in cuius conspectu ambulaverunt patres nostri, domine qui pascis me a iuventute mea.}
\end{responsory-doxology}%
{\color{Red} Oravit.}
\newline {\color{Red} In Laudibus. Ant.} Domine labia mea aperies, et os meum annuntiabit laudem tuam. {\color{Red} Ps.} \psalm{50} {\color{Red} Ant.} Dextera domini fecit virtutem, dextera domini exaltavit me. {\color{Red} Ps.} \psalm{117} {\color{Red} Ant.} Factus est adiutor meus deus meus. {\color{Red} Ps.} \psalm{62} {\color{Red} Ant.} Trium puerorum cantemus ymnum quem cantabant in camino ignis benedicentes
%pp100
dominum. {\color{Red} Can.} \psalm{benedicite} {\color{Red} Ant.} Statuit ea in eternum et in seculum seculi preceptum posuit et non preteribit. {\color{Red} Ps.} \psalm{148}
\newline {\color{Red} Ad Ben. Ant.} Egressus Ihesus secessit in partes Tyri et Sidonis, et ecce mulier Chananea a finibus illis egressa clamabat dicens, miserere mei domine fili David.
\newline {\color{Red} Ad Primam. Ant.} Accedentes discipuli Ihesu ad eum, rogabant eum dicentes, dimitte eam quia clamat post nos.
\newline {\color{Red} Ad Terciam. Ant.} Non sum missus nisi ad oves que perierant domus Israhel dicit dominus.
\newline {\color{Red} Ad Sextam. Ant.} O mulier magna est fides tua fiat tibi sicut petisti.
\newline {\color{Red} Ad Nonam. Ant.} Vade mulier semel tibi dixi, si credideris videbis mirabilia.
\newline {\color{Red} Ad Mag. Ant.} Dixit dominus mulieri Chananee, non est bonum sumere panem filiorum et mittere canibus ad manducandum, utique domine, nam et catelli edunt de micis que cadunt de mensa dominorum suorum, ait illi Ihesus, mulier magna est fides tua fiat tibi sicut petisti.%typo: magna est est
\newline \ul{Lectiones feriales a duodetricesimo capitulo usque ad trigesimo secundo ad libitum.}
\newline {\color{Red} Feria Secunda. Ad Ben. Ant.} Ego principium qui et loquor vobis. {\color{Red} Oratio.}
\lettrine[lines=2]{\zallmancaps \color{Red} P}{resta} quesumus omnipotens deus, ut familia tua que se affligendo carne ab alimentis abstinet, sectando iusticiam a culpa ieiunet. Per.
\newline {\color{Red} Ad Mag. Ant.} Qui me misit mecum est, et non reliquit me solum, quia que placita sunt ei facio semper. {\color{Red} Oratio.}
\lettrine[lines=2]{\zallmancaps \color{Blue} A}{desto} supplicationibus nostris omnipotens deus et quibus fiduciam sperande pietatis indulges, consuete misericordie tribue benignus effectum. Per.
\newline {\color{Red} Feria Tercia. Ad Ben. Ant.} Unus est enim magister vester qui in celis est, dicit dominus. {\color{Red} Oratio.}
\lettrine[lines=2]{\zallmancaps \color{Red} P}{erfice} quesumus domine benignus in nobis observantie sancte subsidium, ut que te actore facienda cognovimus te operante impleamus. Per.
\newline {\color{Red} Ad Mag. Ant.} Qui maior est vestrum erit minister vester, quia omnis qui se exaltat humiliabitur, dicit dominus. {\color{Red} Oratio.}
\lettrine[lines=2]{\zallmancaps \color{Blue} P}{ropiciare} domine supplicationibus nostris, et animarum nostrarum medere languoribus ut remissione percepta, in tua semper benedictione letemur. Per.
\newline {\color{Red} Feria Quarta. Ad Ben. Ant.} Ecce ascendimus Iherosolimam et filius hominis tradetur ad crucifigendum. {\color{Red} Oratio.}
\lettrine[lines=2]{\zallmancaps \color{Red} P}{opuli} tuum domine propicius respice, et quos ab escis carnalibus precipis abstinere, a noxiis quoque vitiis cessare concede. Per.
\newline {\color{Red} Ad Mag. Ant.} Sedere autem mecum non est meum dare vobis sed quibus paratum est a patre meo. {\color{Red} Oratio.}
\lettrine[lines=2]{\zallmancaps \color{Blue} D}{eus} innocentie restitutor et amator, dirige ad te tuorum corda servorum, ut spiritus tui fervore concepto et in fide inveniantur stabiles et in opere efficaces. Per eiusdem.
\newline {\color{Red} Feria Quinta. Ad Ben. Ant.} Pater Abraham miserere mei et mitte Lazarum ut intingat extremum digiti sui in aquam ut refrigeret linguam meam. {\color{Red} Oratio.}
\lettrine[lines=2]{\zallmancaps \color{Red} P}{resta} nobis domine quesumus auxilium gratie tue, ut ieiuniis et orationibus convenienter intenti, liberemur ab hostibus mentis et corporis. Per.
\newline {\color{Red} Ad Mag. Ant.} Fili recordare quia recepisti bona in vita tua, et Lazarus similiter mala. {\color{Red} Oratio.}
\lettrine[lines=2]{\zallmancaps \color{Blue} A}{desto} domine famulis tuis, et perpetuam benignitatem largire poscentibus, ut hiis qui te actore et gubernatore gloriantur, et congregata restaures, et restaurata conserves. Per.
\newline {\color{Red} Feria Sexta. Ad Ben. Ant.} Malos male perdet, et vineam suam locabit aliis agricolis, qui reddant ei fructum temporibus suis. {\color{Red} Oratio.}
\lettrine[lines=2]{\zallmancaps \color{Red} D}{a} quesumus omnipotens deus, ut sacro nos purificante ieiunio sinceris mentibus ad sancta ventura facias pervenire. Per.
\newline {\color{Red} Ad Mag. Ant.} Querentes eum tenere timuerunt turbas, quia sicut prophetam eum habebant. {\color{Red} Oratio.}
\lettrine[lines=2]{\zallmancaps \color{Blue} D}{a} quesumus domine populo tuo salutem mentis et corporis, ut bonis operibus inherendo tue semper virtutis mereatur protectione defendi. Per.
\newline {\color{Red} Sabbato. Ad Ben. Ant.} Vadam ad patrem meum et dicam ei, pater, fac me sicut unum ex mercenariis tuis. {\color{Red} Oratio.}
\lettrine[lines=2]{\zallmancaps \color{Red} D}{a} quesumus domine nostris effectum ieiuniis salutarem, ut castigatio carnis assumpta ad nostrarum vegetationem transeat animarum. Per.
{\ubsubsection{Dominica Tertia in Quadragesima Sabbato Precedenti}{breviarium-dominica-tertia-in-quadragesima-sabbato-precedenti}}
\par \noindent {\color{Red} Ad Vesperas.} \R \hyperlink{loquens-ioseph}{Loquens Ioseph.} {\color{Red} IX. Ad Mag. Ant.} Dixit autem pater ad servos suos, cito proferte stolam primam et induite illum, et date anulum in manu eius et calciamenta in pedibus eius. {\color{Red} Oratio.}
\lettrine[lines=2]{\zallmancaps \color{Blue} Q}{uesumus} omnipotens deus, vota humilium respice, atque ad defensionem nostram dexteram tue maiestatis extende. Per.
\newline {\color{Red} In Completorio. Ad Nunc Dimittis. Ant.} Media vita in morte sumus, quem querimus adiutorem nisi te domine qui pro peccatis nostris iuste irasceris, * Sancte deus, sancte fortis, sancte et misericors salvator amare morti ne tradas nos. \V Ne proicias nos in tempore senectutis, cum defecerit virtus nostra ne derelinquas nos domine. Sancte deus.
\newline \ul{Precedens antiphona cum suo versu dicatur per quindecim dies super} \psalm{nunc} \ul{nisi in Annunciatione si infra hoc tempus evenerit.}
\newline {\color{Red} Ad Matutinas. Invitat.}
\begin{invitatory}[ploremus-coram-invitatorium]
{Ploremus coram domino deo nostro, * Qui fecit nos.}
\end{invitatory}
\newline {\color{Red} Lectio Prima.}
\lettrine[lines=2]{\zallmancaps \color{Red} I}{oseph} cum sexdecim esset annorum, pascebat gregem cum fratribus suis adhuc puer. Et erat cum filiis Bale et Zelphe uxorum patris sui. Accusavitque fratres suos apud patrem crimine pessimo. Israhel autem diligebat Ioseph super omnes filios suos, eo quod in senectute genuisset eum. Fecitque ei tunicam polimitam. \R
\lettrine[lines=2]{\zallmancaps \color{Blue} V}{identes} \hypertarget{videntes-ioseph}{\label{videntes-ioseph}} Ioseph a longe loquebantur mutuo fratres dicentes, ecce somniator venit, * Venite occidamus eum, et videamus si prosint illi somnia sua. \V Cumque vidissent Ioseph fratres sui quod a patre cunctis fratribus plus amaretur oderant eum nec poterant ei quicquam pacifice loqui unde et dicebant. Venite.
\newline {\color{Red} Lectio Secunda.}
\lettrine[lines=2]{\zallmancaps \color{Red} V}{identes} autem fratres eius quod a patre plus cunctis filiis amaretur, oderant eum, nec poterant ei quicquam pacifice loqui. Accidit quoque ut visum somnium referret fratribus suis. Que causa maioris odii seminarium fuit.
\begin{responsory}[dixit-iudas-fratribus]
{Dixit Iudas fratribus suis, ecce Ismaelite transeunt, venite venundetur et manus nostre non polluantur, * Caro enim et frater noster est.}{Quid enim prodest si occiderimus fratrem nostrum, et celaverimus sanguinem ipsius, melius est ut venundetur.}
\end{responsory}%
\newline {\color{Red} Lectio Tertia.}
\lettrine[lines=2]{\zallmancaps \color{Blue} D}{ixitque} ad eos. Audite somnium meum quod vidi. Putabam nos ligare manipulos in agro, et quasi consurgere manipulum meum et stare, vestrosque manipulos circunstantes adorare manipulum meum. Responderunt fratres eius. Nunquid rex noster eris, aut subiciemur ditioni tue? Hec ergo causa somniorum atque sermonum invidie et odii fomitem ministravit.
\begin{responsory-final}[videns-iacob]
{Videns Iacob vestimenta Ioseph, scidit vestimenta sua cum fletu et dixit, * Fera pessima devoravit, * Filium meum Ioseph.}{Congregatis autem cunctis liberis eius ut delinirent dolorem patris noluit eos audire sed dixit.}{Filium}
\end{responsory-final}%
\newline {\color{Red} Lectio Quarta.}
\lettrine[lines=2]{\zallmancaps \color{Red} A}{liud} quoque vidit somnium, quod narrans fratribus ait. Vidi per somnium quasi solem et lunam et stellas undecim adorare me. Quod cum patri suo et fratribus retulisset, increpavit eum pater et ait. Quid sibi vult hoc somnium quod vidisti? Nunquid ego et mater tua et fratres tui adorabimus te super terram? Invidebant ei igitur fratres sui, pater vero rem tacitus considerabat.
\begin{responsory}[ioseph-dum-intraret]
{Ioseph dum intraret in terram Egypti linguam quam non novit audivit, * Manus eius in laboribus servierunt, et lingua eius inter principes loquebatur sapientiam.}{Divertit ab oneribus dorsum eius.}
\end{responsory}%
\newline {\color{Red} Lectio Quinta.}
\lettrine[lines=2]{\zallmancaps \color{Blue} C}{umque} fratres illius in pascendis gregibus patris sui morarentur in Sichem, dixit ad eum Israhel. Fratres tui pascunt oves in Sichimis, veni mittam te ad eos. Quo respondente presto sum, ait. Vade et vide si cuncta prospera sint erga fratres tuos et pecora, et renuncia michi quid agatur.
\begin{responsory}[memento-mei-dum]
{Memento mei dum bene tibi fuerit, * Ut suggeras Pharaoni, ut educat me de isto carcere, quia furtim sublatus sum, et hic innocens in lacum missus sum.}{Tres enim adhuc dies sunt post quos recordabitur Pharao ministerii tui et restituet te in gradum pristinum, tunc memento mei.}
\end{responsory}%
\newline {\color{Red} Lectio Sexta.}
\lettrine[lines=2]{\zallmancaps \color{Red} M}{issus} de valle Ebron, venit in Sichem. Invenitque eum vir errantem in agro, et interrogavit quid quereret. At ille respondit. Fratres meos quero. Indica michi ubi pascant greges. Dixitque ei vir. Recesserunt de loco isto. Audivi autem eos dicentes, eamus in Dothaim. Perrexit ergo Ioseph post fratres suos, et invenit eos in Dothaim.%typo: Perexit
\begin{responsory-doxology}[dixit-ruben]
{Dixit Ruben fratribus suis, nunquid non dixi vobis, nolite peccare in puerum et non audistis me, * En sanguis eius exquiritur.}{Merito hec patimur quia peccavimus in fratrem nostrum videntes angustiam anime eius dum deprecaretur nos et non audivimus.}
\end{responsory-doxology}%
\newline {\color{Red} Lectio VII. Secundum Lucam.}
\lettrine[lines=2]{\zallmancaps \color{Blue} I}{n} illo tempore: Erat Ihesus eiciens demonium, et illud erat mutum. Et reliqua.
\newline {\color{Red} Omelia Venerabilis Bede Presbiteri.}
\lettrine[lines=2]{\zallmancaps \color{Red} D}{emoniacus} iste apud Matheum non solum mutus sed et cecus fuisse narratur, curatusque dicitur a domino, ita ut loqueretur et videret. Tria generalia signa simul in uno homine perpetrata sunt. Cecus videt? mutus loquitur, possessus a demone liberatur.
\begin{responsory}[merito-hec]
{Merito hec patimur quia peccavimus in fratrem nostrum, videntes angustiam anime eius dum deprecaretur nos et non audivimus, * Idcirco venit super nos tribulatio.}{Dixit Ruben fratribus suis nunquid non dixi vobis nolite peccare in puerum et non audistis me.}
\end{responsory}%
\newline {\color{Red} Lectio VIII.}
\lettrine[lines=2]{\zallmancaps \color{Blue} H}{oc} tunc quidem carnaliter factum est: sed et cotidie impletur in conversione credentium, ut expulso primum demone, fidei lumen aspiciant, deinde ad laudem dei tacentia prius ora laxentur. Quidam autem ex eis dixerunt. In Beelzebub principe demoniorum eicit demonia. Non hoc aliqui de turba, sed Pharisei calumniantes dicebant et scribe, sicut alii evangeliste testantur.%typo: callumniantes
\begin{responsory}[tollite-hinc]
{Tollite hinc nobiscum munera et ite ad dominum terre, et dum inveneritis adorate eum super terram, * Deus autem meus faciat vobis eum placabilem et remittat et hunc fratrem vestrum vobiscum, et eum quem tenet in vinculis.}{Sumite de optimis terre frugibus in vasis vestris et deferte viro munera.}
\end{responsory}%
\newline {\color{Red} Lectio IX.}
\lettrine[lines=2]{\zallmancaps \color{Red} T}{urbis} quippe que minus erudite videbantur domini semper facta mirantibus, illi contra vel negare hec, vel que negare nequiverant sinistra interpretatione pervertere laborabant, quasi hec non divinitatis sed immundi spiritus opera fuissent, idest Beelzebub, qui deus erat Acharon.
\begin{responsory-final}[loquens-ioseph]
{Loquens Ioseph fratribus suis dixit, * Pax vobis nolite timere, * Per salute ei vestra misit me dominus ante vos.}{Elevavitque vocem cum fletu, quam audierunt Egiptii omnisque domus Pharaonis, et dixit fratribus suis.}{Per salute}Loquens.
\end{responsory-final}%
\newline {\color{Red} In Laudibus. Ant.} Fac benigne in bona voluntate tua ut edificentur domine muri Hierusalem. {\color{Red} Ps.} \psalm{50} {\color{Red} Ant.} Dominus in adiutor est, non timebo quid faciat michi homo. {\color{Red} Ps.} \psalm{117} {\color{Red} Ant.} Deus miseratur nostri et benedicat nobis. {\color{Red} Ps.} \psalm{62} {\color{Red} Ant.} Vim virtutis sue oblitus est ignis ut pueri tui liberarentur illesi. {\color{Red} Can.} \psalm{benedicite} {\color{Red} Ant.} Sol et luna laudate deum quia exaltatum est nomen eius solius. {\color{Red} Ps.} \psalm{148}
\newline {\color{Red} Ad Ben. Ant.} Et cum eiecisset demonium locutus est mutus, et admirate sunt turbe.
\newline {\color{Red} Ad Primam. Ant.} Si in digito dei eicio demonia profecto venit in vos regnum dei.
\newline {\color{Red} Ad Terciam. Ant.} Dum fortis armatus custodit atrium suum in pace sunt omnia que possidet.
\begin{responsory-breve}[bonum-michi]
{Bonum michi domine quod humiliasti me, * Bonum michi lex oris tui super milia auri et argenti.}{Manus tue fecerunt me et plasmaverunt me da michi intellectum ut discam mandata tua.}
\end{responsory-breve}%
\V Dicet domino.
\newline {\color{Red} Ad Sextam. Ant.} Qui non colligit mecum dispergit et qui non est mecum contra me est.
\begin{responsory-breve}[servus-tuus]
{Servus tuus ego sum, * Da michi intellectum, domine.}{Ut discam mandata tua.}
\end{responsory-breve}%
\V Ipse libera.
\newline {\color{Red} Ad Nonam. Ant.} Cum immundus spiritus exierit ab homine ambulat per loca inaquosa querens requiem et non invenit.
\begin{responsory-breve}[declara-super]
{Declara super nos deus, * Tuam misericordiam.}{Declaratio sermonum tuorum dat intellectum domine.}
\end{responsory-breve}%
\V Scapulis suis.
\newline \ul{Responsoria prenotata dicantur per quindecim dies ad horas quando de tempore agitur.}
\newline {\color{Red} Ad Mag. Ant.} Extollens quedam mulier vocem de turba dixit beatus venter qui de portavit, et ubera que suxisti, at Ihesus dixit ei, quinimo beati qui audiunt verbum dei et custodiunt illud.
\newline \ul{Lectiones feriales a trigesimonono capitulo usque ad finem ad libitum.}
\newline {\color{Red} Feria Secunda. Ad Ben. Ant.} Amen dico vobis quia nemo propheta acceptus est in patria sua. {\color{Red} Oratio.}
\lettrine[lines=2]{\zallmancaps \color{Blue} C}{ordibus} nostris quesumus domine gratiam tuam benignus infunde: ut sicut ab escis corporalibus abstinemus, ita sensus quoque nostros a noxiis retrahamus excessibus. Per.%typo: benigus
\newline {\color{Red} Ad Mag. Ant.} Ihesus autem transiens per medium illorum ibat. {\color{Red} Oratio.}
\lettrine[lines=2]{\zallmancaps \color{Red} S}{ubveniat} nobis domine misericordia tua, ut ab imminentibus peccatorum nostrorum periculis te mereamur protegente salvari. Per.%sic, missing eripi te liberante
\newline {\color{Red} Feria Tercia. Ad Ben. Ant.} Si duo ex vobis consenserint super terram de omni re quamcumque pecierint fiet illis a patre meo qui in celis est, dicit dominus. {\color{Red} Oratio.}
\lettrine[lines=2]{\zallmancaps \color{Blue} E}{xaudi} nos omnipotens et misericors deus, et continentie salutaris propitius nobis dona concede. Per.
\newline {\color{Red} Ad Mag. Ant.} Ubi duo vel tres congregati fuerint in nomine meo, in medio eorum sum dicit dominus. {\color{Red} Oratio.}
\lettrine[lines=2]{\zallmancaps \color{Red} T}{ua} nos domine protectione defende, et ab omni semper iniquitate te custodi. Per.
\newline {\color{Red} Feria Quarta. Ad Ben. Ant.} Audite et intelligite traditiones quas dominus dedit vobis. {\color{Red} Oratio.}
\lettrine[lines=2]{\zallmancaps \color{Blue} P}{resta} nobis quesumus domine, ut salutaribus ieiuniis eruditi, a noxiis quoque vitiis abstinentes, propitiationem tuam facilius impetremus. Per.
\newline {\color{Red} Ad Mag. Ant.} Non lotis manibus manducare non coinquinat hominem. {\color{Red} Oratio.}%typo: conquinat
\lettrine[lines=2]{\zallmancaps \color{Red} C}{oncede} quesumus omnipotens deus, ut qui protectionis tue gratiam querimus, liberati a malis omnibus secura tibi mente serviamus. Per.
\newline {\color{Red} Feria Quinta. Ad Ben. Ant.} Exibant autem demonia a multis clamantia et dicentia quia tu es filius dei, et increpans non sinebat ea loqui quia sciebant ipsum esse Christum. {\color{Red} Oratio.}
\lettrine[lines=2]{\zallmancaps \color{Blue} C}{oncede} quesumus omnipotens deus, ut ieiuniorum nobis sancta devotio, et purificationem tribuat, et maiestati tue nos reddat acceptos. Per.
\newline {\color{Red} Ad Mag. Ant.} Omnes qui habebant infirmos ducebant illos ad Ihesum et curabantur. {\color{Red} Oratio.}
\lettrine[lines=2]{\zallmancaps \color{Red} S}{ubiectum} tibi populum quesumus domine propitatio celestis amplificet et tuis semper faciat servire mandatis. Per.
\newline {\color{Red} Feria Sexta. Ad Ben. Ant.} Domine ut video propheta es tu, patres nostri in monte hoc adoraverunt. {\color{Red} Oratio.}%notandum: abstinem at tu es
\lettrine[lines=2]{\zallmancaps \color{Red} I}{eiunia} nostra quesumus domine benigno favore prosequere, ut sicut ab alimentis in corpore, ita a vitiis ieiunemus in mente. Per.
\newline {\color{Red} Ad Mag. Ant.} Veri adoratores adorabunt patrem in spiritu et veritate. {\color{Red} Oratio.}
\lettrine[lines=2]{\zallmancaps \color{Blue} P}{resta} quesumus omnipotens deus, ut qui in tua protectione confidimus, cuncta nobis adversantia te adiuvante vincamus. Per.
\newline {\color{Red} Sabbato. Ad Ben. Ant.} Inclinavit se Ihesus et scribebat in terra siquis sine peccato est mittat in eam lapidem. {\color{Red} Oratio.}
\lettrine[lines=2]{\zallmancaps \color{Red} P}{resta} quesumus omnipotens deus, ut qui se affligendo carne ab alimentis abstinent, sectando iusticiam a culpa ieiunent. Per.
{\ubsubsection{Dominica Quarta in Quadragesima Sabbato Precedenti}{breviarium-dominica-quarta-in-quadragesima-sabbato-precedenti}}
\par \noindent {\color{Red} Ad Vesperas.} \R \hyperlink{audi-israel}{Audi Israhel.} {\color{Red} IX. Ad Mag. Ant.} Nemo te condemnavit mulier nemo domine, nec ego te condemnabo, iam amplius noli peccare. {\color{Red} Oratio.}
\lettrine[lines=2]{\zallmancaps \color{Blue} C}{oncede} quesumus omnipotens deus, ut qui ex merito nostre actionis affligimur, tue gratie consolatione respiremus. Per.
\newline {\color{Red} Ad Matutinas. Invitat.}
\begin{invitatory}[populus-domini-invitatorium]
{Populus domini et oves paschue eius, * Venite adoremus dominum.}
\end{invitatory}
%Ex 1:1
\newline {\color{Red} Lectio Prima.}
\lettrine[lines=2]{\zallmancaps \color{Red} H}{ec} sunt nomina filiorum Israhel, qui ingressi sunt in Egiptum cum Iacob. Singuli cum domibus suis introierunt. Ruben, Simeon, Levi, Iudas, Ysachar, Zabulon, et Beniamin, Dan et Neptalim, Gad et Aser. Erant igitur omnes anime eorum qui egressi sunt de femore Iacob septuaginta. Ioseph autem in Egipto erat. Quo mortuo et universis fratribus eius, omnique cognatione sua, filii Israhel creverunt, et quasi germinantes multiplicati sunt, ac roborati nimis impleverunt terram. \R
\lettrine[lines=2]{\zallmancaps \color{Blue} L}{ocutus} \hypertarget{locutus-est-quarta-quadragesima}{\label{locutus-est-quarta-quadragesima}} est dominus ad Moysen dicens descende in Egiptum, dic Pharaoni ut dimittat populum meum, induratum est cor Pharaonis non vult dimittere populum meum nisi, * In manu forti. \V Videns vidi afflictionem populi mei qui est in Egipto, et gemitum eius audivi, et descendi liberare eum. In manu.
\newline {\color{Red} Lectio Secunda.}
\lettrine[lines=2]{\zallmancaps \color{Red} S}{urrexit} interea rex novus super Egiptum, qui ignorabat Ioseph. Et ait ad populum suum. Ecce populus filiorum Israhel multus et fortior nobis est. Venite, sapienter opprimamus eum, ne forte multiplicetur, et si ingruerit contra nos bellum addatur inimicis nostris, expugnatisque nobis egrediatur e terra.
\begin{responsory}[stetit-moyses]
{Stetit Moyses coram Pharaone et dixit, hec dicit dominus, * Dimitte populum meum ut sacrificet michi in deserto.}{Dominus deus Hebreorum misit me ad te dicens.}
\end{responsory}%
\newline {\color{Red} Lectio Tertia.}
\lettrine[lines=2]{\zallmancaps \color{Blue} P}{reposuit} itaque eis magistros operum, ut affligerent eos oneribus. Edificaveruntque urbes tabernaculorum Pharaoni, Phiton et Ramesses. Quantoque opprimebant eos, tanto magis multiplicabantur et crescebant. Oderantque filios Israhel Egiptii, et affligebant illudentes eis, atque ad amaritudinem perducebant vitam eorum operibus duris luti et lateris, omnique famulatu quo in terre operibus premebantur.
\begin{responsory-doxology}[cantemus-domine]
{Cantemus domine gloriose enim honorificatus est, equum et ascensorem proiecit in mare, * Adiutor et protector factus est michi dominus in salutem.}{Dominus quasi vir pugnator omnipotens nomen eius.}
\end{responsory-doxology}%
\newline {\color{Red} Lectio Quarta.}
\lettrine[lines=2]{\zallmancaps \color{Red} D}{ixit} autem rex Egipti obstetricibus Hebreorum, quarum una vocabatur Sephora, altera Phua, precipiens eis. Quando obstetricabitis Hebreas et tempus partus advenerit, si masculus fuerit interficite illum, si femina reservate.
\begin{responsory}[in-mari-via]
{In mari via tua et semite tue in aquis multis, deduxisti sicut oves populum tuum, * In manu Moysi et Aaron.}{Transtulisti illos per mare rubrum et transvexisti eos per aquam nimiam.}
\end{responsory}%
\newline {\color{Red} Lectio Quinta.}
\lettrine[lines=2]{\zallmancaps \color{Blue} T}{imuerunt} autem obstetrices Deum, et non fecerunt iuxta preceptum regis Egipti, sed conservabant mares. Quibus ad se accersitis, rex ait. Quidnam est hoc quod facere voluistis ut pueros servaretis? Que responderunt. Non sunt Hebree sicut Egiptie mulieres. Ipse enim obstetricandi habent scientiam, et priusquam veniamus ad eas pariunt.
\begin{responsory}[qui-persequebantur]
{Qui persequebantur populum tuum domine demersisti eos in profundum, * Et in columna nubis ductor eorum fuisti domine.}{Deduxisti sicut oves populum tuum in manu Moysi et Aaron.}
\end{responsory}%
\newline {\color{Red} Lectio Sexta.}
\lettrine[lines=2]{\zallmancaps \color{Red} B}{ene} ergo fecit Deus obstetricibus, et crevit populus, confortatusque est nimis. Et quia timuerunt obstetrices
%pp101
Deum, edificavit illis domos. Precepit ergo Pharao omni populo suo dicens. Quicquid masculini sexus natum fuerit, in flumen proicite, quicquid feminei reservate.
\begin{responsory-doxology}[moyses-famulus]
{Moyses famulus domini ieiunavit quadraginta diebus et quadraginta noctibus, * Ut legem domini mereretur accipere.}{Ascendit Moyses in montem Sinai.}
\end{responsory-doxology}%
\newline {\color{Red} Lectio VII. Secundum Iohannem.}
\lettrine[lines=2]{\zallmancaps \color{Blue} I}{n} illo tempore: Abiit Ihesus transmare Galilee quod est Tyberiadis, et sequebatur eum multitudo magna quia videbant signa que faciebat super his qui infirmabantur. Et reliqua.
\newline {\color{Red} Omelia Beati Augustini Episcopi.}
\lettrine[lines=2]{\zallmancaps \color{Red} M}{iracula}, que fecit Dominus noster Ihesus Christus sunt quidem divina opera, et ad intelligendum Deum de visibilibus admonent humanam mentem. Quia enim ille non est talis substantia que videri oculis possit, et miracula eius quibus totum mundum regit, universamque creaturam administrat, assiduitate viluerunt, ita ut pene nemo dignetur attendere opera Dei mira et stupenda, in quolibet seminis grano secundum suam misericordiam servavit sibi quedam que faceret oportuno tempore preter usitatum cursum ordinemque nature, ut non maiora sed insolita videndo stuperent quibus cotidiana viluerunt.
\begin{responsory}[splendida-facta]
{Splendida facta est facies Moysi dum aspiceret in eum dominus, * Videntes seniores claritatem vultus eius admirantes timuerunt valde.}{Descendit Moyses de monte portans duas tabulas lapideas in manibus suis scriptas utrasque digito dei.}
\end{responsory}%
\newline {\color{Red} Lectio VIII.}
\lettrine[lines=2]{\zallmancaps \color{Blue} M}{aius} quippe miraculum est gubernatio totius mundi, quam saturatio quinque milium hominum de quinque panibus. Et tamen hoc nemo miratur? Illud mirantur homines: non quia maius est, sed quia rarum est. Quis enim et nunc pascit universum mundum, nisi ille qui de paucis granis segetes creat?
\begin{responsory}[ecce-mitto]
{Ecce mitto angelum meum qui precedat te et custodiat semper, observa et audi vocem meam: et inimicus ero inimicis tuis et affligentes te affligam, * Et precedet te angelus meus.}{Israhel si me audieris non erit in te deus recens neque adorabis deum alienum, ego enim dominus.}
\end{responsory}%
\newline {\color{Red} Lectio IX.}
\lettrine[lines=2]{\zallmancaps \color{Red} F}{ecit} ergo quomodo Deus. Unde enim multiplicat de paucis granis segetes, inde in manibus suis multiplicavit quinque panes. Potestas enim erat in manibus suis. Panes autem illi quinque, quasi semina erant, non quidem terre mandata, sed ab eo qui terram fecit multiplicata.
\begin{responsory-doxology}[audi-israel]
{Audi Israhel precepta domini et ea in corde tuo quasi in libro scribe, * Et dabo tibi terram fluentem lac et mel.}{Observa igitur et audi vocem meam et inimicus ero inimicis tuis.}Audi.
\end{responsory-doxology}%
\newline {\color{Red} In Laudibus. Ant.} Tunc acceptabis sacrificium iusticie si averteris faciem tuam a peccatis meis. {\color{Red} Ps.} \psalm{50} {\color{Red} Ant.} Bonum est sperare in domino quam sperare in principibus. {\color{Red} Ps.} \psalm{117} {\color{Red} Ant.} Benedicat nos deus deus noster, benedicat nos deus. {\color{Red} Ps.} \psalm{62} {\color{Red} Ant.} Potens es domine eripere nos de manu mortis, liberare nos deus noster. {\color{Red} Can.} \psalm{benedicite} {\color{Red} Ant.} Reges terre et omnes populi laudate deum. {\color{Red} Ps.} \psalm{148}
\newline {\color{Red} Ad Ben. Ant.} Abiit Ihesus transmare Galilee quod est Tyberiadis, et sequebatur eum multitudo magna, quia videbant signa super his qui infirmabantur, subiit ergo in montem Ihesus, et ibi sedebat cum discipulis suis, erat autem proximum pascha dies festus Iudeorum.
\newline {\color{Red} Ad Primam. Ant.} Accepit ergo Ihesus panes, et cum gratias egisset distribuit discumbentibus.
\newline {\color{Red} Ad Terciam. Ant.} De quinque panibus et duobus piscibus satiavit dominus quinque milia hominum.
\newline {\color{Red} Ad Sextam. Ant.} Satiavit dominus quinque milia hominum de quinque panibus et duobus piscibus.
\newline {\color{Red} Ad Nonam. Ant.} Illi homines cum signum vidissent quod factum fuerat glorificabant deum et dicebant, quia hic est salvator mundi.
\newline {\color{Red} Ad Mag. Ant.} Cum sublevasset oculos Ihesus et vidisset maximam multitudinem venientem ad se, dixit ad Philippum, unde ememus panes ut manducent hi? hoc autem dicebat temptans eum, ipse enim sciebat quid esset facturus.
\newline \ul{Lectiones feriales a secundo capitulo usque ad finem ad libitum.}
\newline {\color{Red} Feria Secunda. Ad Ben. Ant.} Auferte ista hinc dicit dominus, et nolite facerer domum patris mei domum negotiationis. {\color{Red} Oratio.}
\lettrine[lines=2]{\zallmancaps \color{Blue} P}{resta} quesumus omnipotens deus ut observationes sacras annua devotione recolentes, et corpore tibi placeamus et mente. Per.
\newline {\color{Red} Ad Mag. Ant.} Solvite templum hoc dicit dominus, et post triduum reedificabo illud, hoc autem dicebat de templo corporis sui. {\color{Red} Oratio.}
\lettrine[lines=2]{\zallmancaps \color{Red} D}{eprecationem} nostram quesumus domine benignus exaudi, et quibus supplicandi prestas affectum, tribue defensionis auxilium. Per.
\newline {\color{Red} Feria Tercia. Ad Ben. Ant.} Quid me queritis interficere hominem qui vera locutus sum vobis. {\color{Red} Oratio.}
\lettrine[lines=2]{\zallmancaps \color{Blue} S}{acre} nobis domine quesumus observationis ieiunia, et pie conversationis augmentum, et tue propitiationis continuum prestent auxilium. Per.
\newline {\color{Red} Ad Mag. Ant.} Unum opus feci et admiramini, quia totum hominem sanum feci in Sabbato. {\color{Red} Oratio.}
\lettrine[lines=2]{\zallmancaps \color{Red} M}{iserere} domine populo tuo, et continuis tribulationibus laborantem propitius respirare concede. Per.
\newline {\color{Red} Feria Quarta. Ad Ben. Ant.} Rabi quis peccavit, hic aut parentes eius quod cecus natus est, respondit Ihesus et dixit, neque hic peccavit, neque parentes eius, sed ut manifestentur opera dei in illo. {\color{Red} Oratio.}
\lettrine[lines=2]{\zallmancaps \color{Blue} D}{eus} qui et iustis premia meritorum, et peccatoribus per ieiunium veniam prebes, miserere supplicibus tuis, ut reatus nostri confessio, indulgentiam valeat percipere delictorum. Per.
\newline {\color{Red} Ad Mag. Ant.} Ille homo qui dicitur Ihesus lutum fecit ex sputo et linivit oculos meos et modo video. {\color{Red} Oratio.}
\lettrine[lines=2]{\zallmancaps \color{Red} P}{ateant} aures misericordie tue domine precibus supplicantium, et ut petentibus desiderata concedas, fac eos que tibi sunt placita postulare. Per.
\newline {\color{Red} Feria Quinta. Ad Ben. Ant.} Accessit Ihesus et tetigit loculum, ii autem qui portabant steterunt, et ait, adolescens tibi dico surge, et resedit qui erat mortuus et cepit loqui. {\color{Red} Oratio.}
\lettrine[lines=2]{\zallmancaps \color{Blue} P}{resta} quesumus omnipotens deus, ut quos ieiunia votiva castigant, ipsa quoque devotio sancta letificet, ut terrenis affectibus mitigatis, facilius celestia capiamus. Per.
\newline {\color{Red} Ad Mag. Ant.} Propheta magnus surrexit in nobis, et quia deus visitavit plebem suam. {\color{Red} Oracio.}
\lettrine[lines=2]{\zallmancaps \color{Red} P}{opuli} tui deus institutor et rector peccata quibus impugnatur expelle, ut semper tibi placitus, et tuo munimine sit securus. Per.
\newline {\color{Red} Feria Sexta. Ad Ben. Ant.} Lazarus amicus noster dormit, eamus et a somno excitemus eum. {\color{Red} Oratio.}
\lettrine[lines=2]{\zallmancaps \color{Blue} D}{eus} qui ineffabilibus mundum renovas sacramentis, presta quesumus ut ecclesia tua eternis proficiat institutis, et temporalibus non destituatur auxiliis. Per.
\newline {\color{Red} Ad Mag. Ant.} Domine si hic fuisses Lazarus non esset mortuus, ecce iam fetet quatriduanus in monumento. {\color{Red} Oratio.}
\lettrine[lines=2]{\zallmancaps \color{Red} D}{a} quesumus omnipotens deus, ut qui infirmitatis nostre conscii de tua virtute confidimus, sub tua semper pietate gaudeamus. Per.
\newline {\color{Red} Sabbato. Ad Ben. Ant.} Ego sum lux mundi qui sequitur me non ambulat in tenebris sed habebit lumen vite dicit dominus. {\color{Red} Oratio.}
\lettrine[lines=2]{\zallmancaps \color{Blue} F}{iat} quesumus domine per gratiam tuam fructuosus nostre devotionis affectus, quia tunc nobis proderunt suscepta ieiunia, si tue sint placita pietati. Per.
{\ubsubsection{Dominica in Passione Sabbato Precedenti}{breviarium-dominica-in-passione-sabbato-precedenti}}
\par \noindent {\color{Red} Ad Vesperas. Capitulum.}
\lettrine[lines=2]{\zallmancaps \color{Red} T}{u} \hypertarget{tu-autem-capitulum}{\label{tu-autem-capitulum}} autem domine Sabaoth, qui iudicas iuste, et probas renes et corda, videam ultionem tuam ex eis, tibi enim revelavi causam meam domine deus meus. \R \hyperlink{multiplicati-sunt}{Multiplicati.} {\color{Red} II.}
\newline \ul{Non dicatur,} Gloria \ul{sed repetatur responsorium. Et notandum quod ab hac die usque ad Pascha exclusive post versum responsorii} Gloria \ul{non dicitur nec ad Matutinas, nec ad alias horas canonicas, sed loco,} Gloria \ul{repetatur responsorium, nisi festum evenerit.} {\color{Red} Ymnus.}
\hymn{vexilla-regis}{Vexilla regis prodeunt}{
\lettrine[lines=2]{\zallmancaps \color{Blue} V}{exilla} regis prodeunt,
fulget crucis mysterium,
quo carne carnis conditor
suspensus est patibulo.
\par {\bfseries \color{Red} Q}uo vulneratus insuper
mucrone diro lancee,
ut nos lavaret crimine
manavit unda sanguine.
\par {\bfseries \color{Blue} I}mpleta sunt que concinit
David fideli carmine,
dicendo nationibus
regnavit a ligno deus.
\par {\bfseries \color{Red} A}rbor decora et fulgida
ornata regis purpura,
electa digno stipite
tam sancta membra tangere.
\par {\bfseries \color{Blue} B}eata cuius brachiis
secli pependit precium,
statera facta corporis,
predamque tulit Tartari.
\par {\bfseries \color{Red} O} crux ave spes unica,
hoc passionis tempore,
auge piis iusticiam,
reisque dona veniam.
\par {\bfseries \color{Blue} T}e summa deus trinitas
collaudet omnis spiritus,
quos per crucis mysterium
salvas rege per secula. Amen.
}
\newline \V Eripe me domine ab homine malo. \R A viro iniquo eripe me.
\newline {\color{Red} Ad Mag. Ant.} Ego sum qui testimonium perhibeo de me ipso et testimonium perhibet de me qui misit me pater. {\color{Red} Oratio.}
\lettrine[lines=2]{\zallmancaps \color{Red} Q}{uesumus} omnipotens deus familiam tuam propitius respice, ut te largiente regatur in corpore et te servante custodiatur in mente. Per.
\newline {\color{Red} Ad Completorium.} \R \hyperlink{in-manus-tuas-breve}{In manus.} {\color{Red} Ad Nunc Dimittis. Ant.} O rex gloriose inter sanctos tuos qui semper es laudabilis et tamen ineffabilis, tu in nobis es domine et nomen sanctum tuum invocatum est super nos, ne derelinquas nos deus noster, ut in die iudicii nos collocare digneris inter sanctos et electos tuos rex benedicte.
\newline \ul{Hec antiphona dicatur cotidie ad} \psalm{nunc}, \ul{usque ad Cenam Domini exclusive nisi in Annunciatione si infra hoc tempus evenerit.}
\newline {\color{Red} Ad Matutinas. Invitat.}
\begin{invitatory}[hodie-si-vocem-invitatorium]
{Hodie si vocem domini audieritis, * Nolite obdurare corda vestra.}
\end{invitatory}
\newline \ul{Hic versus intermittatur in psalmo, et eo dicto sequitur,} Sicut in exacerbatione, \ul{in finem,} Venite, Gloria \ul{non dicatur, sed post ultimum versum, invitatorium bis ex integro repetatur.} {\color{Red} Ymnus.}
\hymn{pange-lingua-fortunatus}{Pange lingua gloriosi}{
\lettrine[lines=2]{\zallmancaps \color{Blue} P}{ange} lingua gloriosi
prelium certaminis,
et super crucis tropheum
dic triumphum nobilem
qualiter redemptor orbis
immolatus vicerit.
\par {\bfseries \color{Red} D}e parentis prothoplasti
fraude facta condolens,
quando pomi noxialis
morte morsu corruit
ipse lignum tunc notavit
damna ligni ut solveret.
\par {\bfseries \color{Blue} H}oc opus nostre salutis
ordo depoposcerat,
multiformis proditoris
ars ut artem falleret,
et medelam ferret inde
hostis unde leserat.
\par {\bfseries \color{Red} Q}uando venit ergo sacri
plenitudo temporis,
missus est ab arce patris
natus orbis conditor,
ac de ventre virginali
caro factus prodiit.
\par {\bfseries \color{Blue} V}agit infans inter arta
conditus presepia:
membra pannis involuta
virgo mater alligat,
manus pedes atque crura
stricta cingit fascia.
\par {\bfseries \color{Red} G}loria et honor deo,
usquequo altissimo,
una patri filioque,
inclito paraclito,
cui laus est et potestas
per eterna secula. Amen.
}
\newline {\color{Red} In Primo Nocturno. Ant.} Quid molesti estis huic mulieri, opus enim bonum operata est in me.%typo: hinc
\newline \ul{Hec antiphona dicatur super psalmos primi nocturni cum triplici,} Gloria.
\newline \V Erue a framea deus animam meam. \R Et de manu canis unicam meam.
\newline {\color{Red} Lectio Prima.}
\lettrine[lines=2]{\zallmancaps \color{Red} V}{erba} Ieremie filii Helchie de sacerdotibus qui fuerunt in Anathot, in terra Beniamin. Quod factum est verbum Domini ad eum in diebus Iosie filii Amon regis Iuda, in tertiodecimo anno regni eius. Et factum est in diebus Ioachim filii Iosie regis Iuda, usque ad consummationem undecimi anni Sedechie filii Iosie regis Iuda, usque ad transmigrationem Hierusalem in mense quinto. Hec dicit. \R
\lettrine[lines=2]{\zallmancaps \color{Blue} I}{sti} \hypertarget{isti-sunt-dies}{\label{isti-sunt-dies}} sunt dies quos observare debetis temporibus suis, * Quarta decima die ad vesperum Pascha domini est et in quinta decima sollemnitatem celebrabitis altissimo domino. \V Locutus est dominus ad Moysen dicens, loquere filiis Israhel et dices ad eos. Quarta.
\newline {\color{Red} Lectio Secunda.}
\lettrine[lines=2]{\zallmancaps \color{Red} E}{t} factum est verbum Domini ad me dicens. Priusquam te formarem in utero novi te, et antequam exires de vulva sanctificavi te: et prophetam in gentibus dedi te. Et dixi. A a, a, Domine Deus, ecce nescio loqui, quia puer ego sum. Et dixit Dominus ad me. Noli dicere puer sum, quoniam ad omnia que mittam te ibis, et universa quecumque mandavero tibi loqueris. Ne timeas a facie eorum, quia ego tecum sum, ut eruam te dicit Dominus. Hec dicit.
\begin{responsory}[multiplicati-sunt]
{Multiplicati sunt qui tribulant me et dicunt, non est salus illi in deo eius, * Exurge domine salvum me fac deus meus.}{Nequando dicat inimicus meus prevalui adversus eum.}
\end{responsory}%
\newline {\color{Red} Lectio Tertia.}
\lettrine[lines=2]{\zallmancaps \color{Blue} E}{t} misit Dominus manum suam, et tetigit os meum. Et dixit Dominus ad me. Ecce dedi verba mea in ore tuo. Ecce constitui te hodie super gentes et super regna, ut evellas et destruas et disperdas, et dissipes, et edifices, et plantes. Hec dicit.
\begin{responsory}[qui-custodiebant]
{Qui custodiebant animam meam consilium fecerunt in unum dicentes, deus dereliquit eum, persequimini et comprehendite eum, quia non est qui liberet eum, deus meus ne elonges a me, * Deus meus in adiutorium meum intende.}{Omnes inimici mei adversum me cogitabant mala michi, verbum iniquum mandaverunt adversum me.}Qui custodiebant.
\end{responsory}%
\newline {\color{Red} Super psalmos Secundi Nocturni cum triplici,} Gloria. {\color{Red} Ant.} Mittens hec mulier in corpus meum hoc unguentum ad sepeliendum me fecit. \V De ore leonis libera me domine. \R Et a cornibus unicornium humilitatem meam.%typo: humilitem
\newline {\color{Red} Lectio Quarta.}
\lettrine[lines=2]{\zallmancaps \color{Red} E}{t} factum est verbum Domini ad me dicens. Quid tu vides Ieremia? Et dixi. Virgam vigilantem ego video. Et dixit Dominus ad me. Bene vidisti, quia vigilabo ego super verbo meo ut faciam illud. Et factum est verbum Domini secundo ad me dicens. Quid tu vides? Et dixi. Ollam succensam ego video, et faciem eius a facie aquilonis. Hec dicit.
\begin{responsory}[deus-meus-es]
{Deus meus es tu ne discedas a me, * Quoniam tribulatio proxima est et non est qui adiuvet.}{Deus deus meus respice in me.}
\end{responsory}%
\newline {\color{Red} Lectio Quinta.}
\lettrine[lines=2]{\zallmancaps \color{Blue} E}{t} dixit Dominus ad me. Ab aquilone pandetur malum super omnes habitatores terre. Quia ecce ego convocabo omnes cognationes regnorum aquilonis, ait Dominus, et venient, et ponent unusquisque solium suum in introitu portarum Hierusalem, et super omnes muros eius in circuitu, et super universas urbes Iuda, et loquar iudicia mea cum eis super omnem malicia eorum qui dereliquerunt me, et libaverunt diis alienis, et adoraverunt opus manuum suarum. Hec dicit.
\begin{responsory}[tota-die]
{Tota die contristatus ingrediebar domine, quoniam anima mea impleta est illusionibus, * Et vim faciebant qui querebant animam meam.}{Et qui inquirebant mala michi locuti sunt vanitates et dolos tota die meditabantur}
\end{responsory}%
\newline {\color{Red} Lectio Sexta.}
\lettrine[lines=2]{\zallmancaps \color{Red} T}{u} ergo accinge lumbos tuos, et surge, et loquere ad eos omnia que ego precipio tibi. Ne formides a facie eorum. Nec enim timere te faciam vultum eorum. Ego quippe dedi te hodie in civitatem munitam, et in columnam ferream, et in murum eneum super omnem terram, regibus Iuda, principibus eius, sacerdotibus et populo terre. Et bellabunt adversum te, et non prevalebunt, quia ego tecum sum dicit Dominus: ut liberem te. Hec dicit.
\begin{responsory}[deus-meus-eripe]
{Deus meus eripe me de manu peccatoris, et de manu contra legem agentis et iniqui, * Quoniam tu es patientia mea.}{Eripe me de inimicis meis deus meus et ab insurgentibus in me libera me.}Deus meus.
\end{responsory}%
\newline {\color{Red} Super psalmos Tertii Nocturni cum triplici,} Gloria. {\color{Red} Ant.} Magister dicit, tempus meum prope est apud te facio Pascha cum discipulis meis. \V Ne perdas cum impiis deus animam meam. \R Et cum viris sanguinum vitam meam.
\newline {\color{Red} Lectio VII. Secundum Iohannem.}
\lettrine[lines=2]{\zallmancaps \color{Blue} I}{n} illo tempore: Dicebat Ihesus turbis Iudeorum, et principibus sacerdotum. Quis ex vobis arguet me de peccato? Et reliqua.
\newline {\color{Red} Omelia Beati Gregorii Pape.}
\lettrine[lines=2]{\zallmancaps \color{Red} P}{ensate} fratres karissimi mansuetudinem Dei. Relaxare peccata venerat, et dicebat. Quis ex vobis arguet me de peccato? Non dedignatur ex ratione ostendere se peccatorem non esse, qui ex virtute divinitatis poterat peccatores iustificare.
\begin{responsory}[adiutor-et-susceptor]
{Adiutor et susceptor meus tu es domine et in verbum tuum speravi, * Declinate a me maligni, et scrutabor mandata dei mei.}{Iniquos odio habui et legem tuam dilexi.}
\end{responsory}%
\newline {\color{Red} Lectio VIII.}
\lettrine[lines=2]{\zallmancaps \color{Blue} S}{ed} terribile est valde quod subditur. Qui est ex Deo, verba Dei audit. Propterea vos non auditis, quia ex Deo non estis. Si enim ipse verba Dei audit qui ex Deo est, et audire verba dei non potest quisquis ex illo non est: interroget se unusquisque si verba Dei in aure cordis percipit, et intelliget unde sit.
\begin{responsory}[in-te-iactatus]
{In te iactatus sum ex utero de ventre matris mee deus meus es tu, ne discedas a me, * Quoniam tribulatio proxima est, et non est qui adiuvet.}{Salva me ex ore leonis et a cornibus unicornium humilitatem meam.}
\end{responsory}%
\newline {\color{Red} Lectio IX.}
\lettrine[lines=2]{\zallmancaps \color{Red} C}{elestem} patriam desiderare veritas iubet, carnis desideria conteri, mundi gloriam declinare, aliena non appeti, propria largiri. Penset ergo apud se unusquisque vestrum si hec vox Dei in cordis eius aure convaluit, et quia iam ex Deo sit agnoscit. Nam sunt nonnulli qui precepta Dei nec aure corporis percipere dignantur.
\begin{responsory}[in-proximo]
{In proximo est tribulatio mea domine et non est qui adiuvet ut fodiant manus meas et pedes meos libera me de ore leonis, * Ut enarrem nomen tuum fratribus meis.}{Erue a framea deus animam meam et de manu canis unicam meam.}In proximo.
\end{responsory}%
\newline \V Intende anime mee et libera eam. \R Propter inimicos meos eripe me.
\newline {\color{Red} In Laudibus. Ant.} Vide domine afflictionem meam quoniam erectus est inimicus meus. {\color{Red} Ps.} \psalm{50} {\color{Red} Ant.} In tribulatione invocavi dominum et exaudivit me in latitudine. {\color{Red} Ps.} \psalm{117} {\color{Red} Ant.} Iudicasti domine causam anime mee, defensor vite mee domine deus meus. {\color{Red} Ps.} \psalm{62} {\color{Red} Ant.} Popule meus quid feci tibi aut quid molestus fui responde michi. {\color{Red} Can.} \psalm{benedicite} {\color{Red} Ant.} Nunquid redditur pro bono malum, quia foderunt foveam anime mee. {\color{Red} Ps.} \psalm{148} {\color{Red} Cap.} \hyperlink{tu-autem-capitulum}{Tu autem domine Sabaoth.} {\color{Red} Ymnus.}
\hymn{lustra-sex}{Lustra sex qui iam}{
\lettrine[lines=2]{\zallmancaps \color{Blue} L}{ustra} sex qui iam peracta
tempus implens corporis,
se volente natus ad hoc
passioni deditus,
agnus in cruce levatur
immolandus stipite.
\par {\bfseries \color{Red} H}ic acetum, fel, arundo,
sputa, clavi, lancea
mite corpus perforatur:
sanguis, unda profluit,
terra, pontus, astra, mundus,
quo lavantur flumine.
\par {\bfseries \color{Blue} C}rux fidelis inter omnes
arbor una nobilis,
nulla silva talem profert
fronde, flore, germine,
dulce lignum dulces clavos
dulce pondus sustinens.
\par {\bfseries \color{Red} F}lecte ramos arbor alta
tensa laxa viscera,
et rigor lentescat ille
quem dedit nativitas,
ut superni membra regis
miti tendas stipite.
\par {\bfseries \color{Blue} S}ola digna tu fuisti
ferre secli precium
atque portum preparare
nauta mundo naufrago,
quem sacer cruor perunxit
fusus agni
%pp102
corpore.
\par {\bfseries \color{Red} G}loria et honor deo,
usquequo altissimo,
una patri filioque,
inclito paraclito,
cui laus est et potestas
per eterna secula. Amen.
}
\newline \V Eripe me de inimicis meis deus meus. \R Et ab insurgentibus in me libera me.
\newline {\color{Red} Ad Ben. Ant.} Quis ex vobis arguet me de peccato si veritatem dico quare vos non creditis michi, qui est ex deo verba dei audit, propterea vos non auditis quia ex deo non estis.
\newline {\color{Red} Ad Primam. Ant.} Ego demonium non habeo sed honorifico patrem meum dicit dominus.
\newline \ul{In hac die et deinceps usque ad Dominican in Albis exclusive non dicatur ad Primam \Rbar ,} Ihesu Christe, \ul{nisi in festis sanctorum si occurrerint. Sed in hac die et deinceps usque ad Cenam Domini exclusive quando de tempore agitur: post Cap. immediate subiungatur Versiculus:} Exurge domine adiuva nos.
\newline {\color{Red} Ad Terciam. Ant.} Ego gloriam meam non quero, est qui querat et iudicet. {\color{Red} Cap.} \hyperlink{tu-autem-capitulum}{Tu autem domine.}
\begin{responsory}[erue-a-framea]
{Erue a framea deus animam meam, * Et de manu canis unicam meam.}{Salva me ex ore leonis, et a cornibus unicornium humilitatem meam.}Erue.
\end{responsory}%
\V De ore leonis libera me domine.
\newline {\color{Red} Ad Sextam. Ant.} Amen amen dico vobis, siquis sermonem meum servaverit mortem non gustabit in eternum. {\color{Red} Cap.}
\lettrine[lines=2]{\zallmancaps \color{Red} F}{aciem} meam non averti ab increpantibus et conspuentibus in me, dominus deus auxiliator meus, et ideo non sum confusus.
\begin{responsory}[de-ore-leonis]
{De ore leonis libera me domine, * Et a cornibus unicornium humilitatem meam.}{Erue a framea deus animam meam et de manu canis unicam meam.}De ore.
\end{responsory}%
\V Ne perdas cum impiis.
\newline {\color{Red} Ad Nonam. Ant.} Abraham pater vester exultavit ut videret diem meum, vidit et gavisus est. {\color{Red} Capitulum.}
\lettrine[lines=2]{\zallmancaps \color{Blue} C}{onfundantur} qui me persequuntur et non confundar ego, paveant illi et non paveam ego, induc super eos diem afflictionis, et duplici contritione contere eos domine deus noster.
\begin{responsory}[principes-persecuti]
{Principes persecuti sunt me gratis et a verbis tuis formidavit cor meum, * Letabor ego super eloquia tua.}{Quasi qui invenit spolia multa.}Principes.
\end{responsory}%
\V Eripe me de inimicis meis deus meus.
\newline {\color{Red} Ad Vesperas.} \ul{Capitulum, Ymnus, \Vbar , ut in Primis Vesperis.} {\color{Red} Ad Mag. Ant.} Tulerunt lapides Iudei ut iacerent in Ihesum, Ihesus autem abscondit se et exivit de templo.
\newline \ul{Ordo huius diei Capitulorum, Ymnorum, Versiculorum, Responsorium ad horas servetur cotidie usque ad Cenam Domini exclusive: quando de tempore agitur.}
\newline \ul{Lectiones feriales a secundo capitulo usque ad duodecimum ad libitum.}
\newline {\color{Red} Feria Secunda. Ad Matutinas. Invitat.}
\begin{invitatory}[adoremus-dominum-qui-nos-invitatorium]
{Adoremus dominum, * Qui nos redemit per crucem.}
\end{invitatory}
\newline \ul{Hoc invitatorium dicatur cotidie usque ad Cenam Domini in ferialibus diebus quando de tempore agitur, et in fine Psalmi,} Venite, \ul{bis ex integro repetatur.}
\newline \R \hyperlink{isti-sunt-dies}{Isti sunt dies} \ul{cum duobus sequentibus, et dicantur similiter in quinta feria.}
\newline {\color{Red} Ad Ben. Ant.} In die magno festivitatis stabat Ihesus et clamabat dicens, siquis sitit veniat ad me et bibat. {\color{Red} Oracio.}
\lettrine[lines=2]{\zallmancaps \color{Red} S}{anctifica} quesumus domine nostra ieiunia, et cunctarum nobis propitius indulgentiam largire culparum. Per.
\newline {\color{Red} Ad Primam. Ant.} Anime impiorum fremebant adversum me, et gravatum est cor meum super eos.
\newline {\color{Red} Ad Terciam. Ant.} Iudicasti. {\color{Red} Ad Sextam. Ant.} Popule meus. {\color{Red} Ad Nonam. Ant.} Nunquid redditur.
\newline \ul{Predicte tres antiphone querantur in Laudibus precedentes dominice et dicantur ad horas in profestis diebus usque ad Cenam Domini.}
\newline {\color{Red} Ad Mag. Ant.} Siquis sitit veniat et bibat, et de ventre eius fluent aque vive. {\color{Red} Oratio.}
\lettrine[lines=2]{\zallmancaps \color{Blue} D}{a} quesumus domine populo tuo salutem mentis et corporis, ut bonis operibus inherendo, tua semper mereatur protectione defendi. Per.
\newline {\color{Red} Feria Tercia. Ad Ben. Ant.} Tempus meum nondum advenit, tempus autem vestrum semper est paratum. {\color{Red} Oratio.}
\lettrine[lines=2]{\zallmancaps \color{Red} N}{ostra} tibi quesumus domine sint accepta ieiunia, que nos et expiando gratia tua dignos efficiant, et ad remedia perducant eterna. Per.
\newline {\color{Red} Ad Mag. Ant.} Vos ascendite ad diem festum hunc, ego non ascendam, quia tempus meum nondum advenit. {\color{Red} Oratio.}
\lettrine[lines=2]{\zallmancaps \color{Blue} D}{a} nobis quesumus domine perseverantem in tua voluntate famulatum, ut in diebus nostris et merito et numero populus tibi serviens augeatur. Per.
\newline {\color{Red} Feria Quarta. Ad Ben. Ant.} Oves mee vocem meam audiunt, et ego dominus agnosco eas. {\color{Red} Oratio.}
\lettrine[lines=2]{\zallmancaps \color{Red} S}{anctificato} hoc ieiunio deus tuorum corda fidelium miserator illustra, et quibus devotionis prestas affectum, prebe supplicantibus pium benignus auditum. Per.
\newline {\color{Red} Ad Mag. Ant.} Multa bona opera operatus sum vobis propter quod opus vultis me occidere. {\color{Red} Oratio.}
\lettrine[lines=2]{\zallmancaps \color{Blue} A}{desto} supplicationibus nostris omnipotens deus, et quibus fiduciam sperande pietatis indulges, consuete misericordie tribuet benignus effectum. Per.
\newline {\color{Red} Feria Quinta. Ad Ben. Ant.} Quid molesti estis huic mulieri, opus enim bonum operata est in me. {\color{Red} Oratio.}
\lettrine[lines=2]{\zallmancaps \color{Red} P}{resta} quesumus omnipotens deus ut dignitas conditionis humane per immoderantiam sauciata medicinalis parsimonie studio reformetur. Per.
\newline {\color{Red} Ad Mag. Ant.} Mittens hec mulier in corpus meum hoc unguentum ad sepeliendum me fecit. {\color{Red} Oratio.}
\lettrine[lines=2]{\zallmancaps \color{Blue} E}{sto} quesumus domine propitius plebi tue, ut que tibi non placent respuens, tuorum potius repleatur delectationibus mandatorum. Per.
\newline {\color{Red} Feria Sexta. Ad Ben. Ant.} Appropinquabat autem dies festus, et querebant principes sacerdotum quomodo Ihesum interficerent, sed timebant plebem. {\color{Red} Oratio.}
\lettrine[lines=2]{\zallmancaps \color{Red} C}{ordibus} nostris quesumus domine auxilium gratie tue benignus infunde, ut peccata nostra castigatione voluntaria cohibentes, temporaliter potius maceremur, quam suppliciis deputemur eternis. Per.
\newline {\color{Red} Ad Mag. Ant.} Principes sacerdotum consilium fecerunt ut Ihesum occiderent, dicebant autem non in die festo, ne forte tumultus fieret in populo. {\color{Red} Oratio.}
\lettrine[lines=2]{\zallmancaps \color{Blue} C}{oncede} quesumus omnipotens deus, ut qui protectionis tue gratiam querimus, liberati a malis omnibus secura tibi mente serviamus. Per.
\newline {\color{Red} Sabbato. Ad Ben. Ant.} Desiderio desideravi hoc Pascha manducare vobiscum antequam patiar. {\color{Red} Oratio.}
\lettrine[lines=2]{\zallmancaps \color{Red} P}{roficiat} quesumus domine plebs tibi dicata pie devotionis affectu: ut sacris actionibus erudita, quanto maiestati tue fit gratior, tanto donis potioribus augeatur. Per.
{\ubsubsection{Dominica in Ramis Palmarum Sabbato Precedenti}{breviarium-dominica-in-ramis-palmarum-sabbato-precedenti}}
\par \noindent {\color{Red} Ad Vesperas.} \R \hyperlink{circumdederunt-me-viri}{Circumdederunt.} {\color{Red} IX. Ad Mag. Ant.} Clarifica me pater apud temet ipsum claritate quam habui priusquam mundus fieret. {\color{Red} Oracio.}
\lettrine[lines=2]{\zallmancaps \color{Blue} O}{mnipotens} sempiterne deus qui humano generi ad imitandum humilitatis exemplum, salvatorem nostrum carnem sumere et crucem subire fecisti, concede propitius ut et patientie ipsius habere documenta, et resurrectionis consortia mereamur. Per eundem.
\newline {\color{Red} Ad Matutinas. Invitat.}
\begin{invitatory}[ipsi-vero-invitatorium]
{Ipsi vero non cognoverunt vias meas, * Quibus iuravi in ira mea si introibunt in requiem meam.}
\end{invitatory}
\newline \ul{Dicto ultimo versu psalmi,} Venite, \ul{usque} hi errant corde. \ul{Invitatorium bis ex integro repetatur.}
\newline \ul{Antiphone in nocturnis ut in precedenti Dominica.}
\newline {\color{Red} Lectio Prima.}
\lettrine[lines=2]{\zallmancaps \color{Red} I}{ustus} quidem tu es Domine, si disputem tecum. Veruntamen iusta loquar ad te. Quare via impiorum prosperatur, bene est omnibus qui prevaricantur et inique agunt? Plantasti eos et radicem miserunt: proficiunt et faciunt fructum. Prope es tu ori eorum, et longe a renibus eorum. Hec dicit. \R
\lettrine[lines=2]{\zallmancaps \color{Blue} I}{n} \hypertarget{in-die-qua}{\label{in-die-qua}} die qua invocavi te domine, dixisti noli timere, iudicasti causam meam et liberasti me, * Deus meus. \V In die tribulationis mee clamavi ad te quia exaudisti me. Deus meus.
\newline {\color{Red} Lectio Secunda.}
\lettrine[lines=2]{\zallmancaps \color{Red} E}{t} tu Domine nosti me vidisti me, et probasti cor meum tecum. Congrega eos quasi gregem ad victimam, et sanctifica eos in die occisionis. Usquequo lugebit terra, et herba omnis regionis siccabitur? Propter malitiam habitantium in ea, consumptum est animal et volucre, quoniam dixerunt, non videbit novissima nostra. Hec dicit.
\begin{responsory}[fratres-mea]
{Fratres mei elongaverunt a me et noti mei, * Quasi alieni recesserunt a me.}{Dereliquerunt me proximi mei, et qui me noverunt.}%typo:elongaverut
\end{responsory}%
\newline {\color{Red} Lectio Tertia.}
\lettrine[lines=2]{\zallmancaps \color{Blue} S}{i} cum peditibus currens laborasti: quomodo poteris contendere cum equis? Cum autem in terra pacis securus fueris, quid facies in superbia Iordanis? Nam et fratres tui et domus patris tui, eciam ipsi pugnaverunt adversus te, et clamaverunt post te plena voce. Ne credas eis, cum locuti fuerint tibi bona. Hec dicit.
\begin{responsory}[attende-domine]
{Attende domine ad me et audi voces adversariorum meorum, * Nunquid redditur pro bono malum quia foderunt foveam anime mee.}{Recordare quod steterim in conspectu tuo ut loquerer pro eis bonum, et averterem indignationem tuam ab eis.}Attende.
\end{responsory}%
\newline {\color{Red} Lectio Quarta.}
\lettrine[lines=2]{\zallmancaps \color{Red} R}{eliqui} domum meam, dimisi hereditatem meam, dedi dilectam animam meam in manu inimicorum eius. Facta est michi hereditas mea quasi leo in silva, dedit contra me vocem, ideo odivi eam. Nunquid avis discolor hereditas mea michi? Nunquid avis tincta per totum? Hec dicit.
\begin{responsory}[conclusit-vias]
{Conclusit vias meas inimicus, insidiator factus est michi sicut leo in abscondito, replevit et inebriavit me amaritudine, deduxerunt in lacum mortis vitam meam, et posuerunt lapidem contra me, * Vide domine iniquitates illorum et iudica causam anime mee defensor vite mee.}{Factus sum in derisum omni populo meo canticum eorum tota die.}
\end{responsory}%
\newline {\color{Red} Lectio Quinta.}
\lettrine[lines=2]{\zallmancaps \color{Blue} V}{enite} congregamini omnes bestie terre, properate ad devorandum. Pastores multi demoliti sunt vineam meam, conculcaverunt partem meam, dederunt portionem meam desiderabilem in desertum solitudinis, posuerunt eam in dissipationem: luxitque super me. Hec dicit.
\begin{responsory}[salvum-me-fac]
{Salvum me fac deus quoniam intraverunt aque usque ad animam meam, ne avertas faciem tuam a me, quoniam tribulor exaudi me, * Deus meus.}{Intende anime mee et libera eam.}
\end{responsory}%
\newline {\color{Red} Lectio Sexta.}
% Jer. 15:10
\lettrine[lines=2]{\zallmancaps \color{Red} V}{e} michi mater mea. Quare genuisti me virum rixe, virum discordie in universa terra? Non feneravi, nec feneravit michi quisquam. Omnes maledicunt michi. Dicit Dominus. Si non reliquie tue in bonum, si non occurri tibi in tempore afflictionis, et in tempore tribulationis et angustie adversus inimicum. Hec dicit.
\begin{responsory}[noli-esse]
{Noli esse michi domine alienus, parce michi in die mala, confundantur omnes qui me persequuntur, * Et non confundar ego.}{Confundantur omnes inimici mei qui querunt animam meam ut auferant eam.}Noli esse.
\end{responsory}%
\newline {\color{Red} Lectio VII. Secundum Matheum.}
\lettrine[lines=2]{\zallmancaps \color{Blue} I}{n} illo tempore: Cum appropinquasset Ihesus Hierosolimis et venisset Bethphage ad montem Oliveti, tunc misit duos discipulos dicens eis. Ite in castellum quod contra vos est. Et reliqua.
\newline {\color{Red} Omelia Venerabilis Bede Presbiteri.}
\lettrine[lines=2]{\zallmancaps \color{Red} M}{ediator} Dei et hominum homo Christus Ihesus, qui pro humani generis salute passurus de celo descenderat ad terras, appropinquante hora passionis, appropinquare voluit loco passionis, ut eciam per hoc claresceret quia non invitus sed sponte pateretur.
\begin{responsory}[dominus-mecum]
{Dominus mecum est tanquam bellator fortis, propterea persecuti sunt me: et intelligere non potuerunt, domine probans renes et corda, * Tibi revelavi causam meam.}{Vidisti domine iniquitates eorum adversum me iudica iudicium meum.}
\end{responsory}%
\newline {\color{Red} Lectio VIII.}
\lettrine[lines=2]{\zallmancaps \color{Blue} I}{n} asino venire, et a turbis rex appellari ac laudari voluit, ut eciam per hoc eruditus quisque agnosceret ipsum esse Christum, quem sic illo venturum prophecia olim premissa signaverat. Ante quinque dies pasche venire voluit, sicut ex Iohannis evangelio didicimus, ut eciam per hoc ostenderet se esse agnum immaculatum qui peccata tolleret mundi.
\begin{responsory}[dixerunt-impii]
{Dixerunt impii apud se non recte cogitantes, circumveniamus iustum quoniam contrarius est operibus nostris, promittit se scientiam dei habere filium dei se nominat, et gloriatur patrem se habere deum, videamus si sermones illius veri sint, et si est verus filius dei liberet illum de manibus nostris, * Morte turpissima condemnemus eum.}{Tanquam nugaces estimati sumus ab illo, et abstinet se a viis nostris tanquam ab immundiciis, et prefert novissima iustorum.}
\end{responsory}%
\newline {\color{Red} Lectio IX.}
\lettrine[lines=2]{\zallmancaps \color{Red} A}{gnus} quippe paschalis cuius immolatione populus Israhel est ab Egiptia servitute liberatus, ante quinque dies pasche, idest quarta decima luna domum inferri, et ipsa die pasche, idest quartadecima luna ad vesperam iussus est immolari, significans eum qui nos suo sanguine redempturus, ante quinque dies pasche, idest hodierno die magno precedencium sequentiumque populorum gaudio ac laudatione deductus venit in templum Dei, et erat cotidie docens in eo.
\begin{responsory}[circumdederunt-me-viri]
{Circumdederunt me viri mendaces sine causa, flagellis ceciderunt me, * Sed tu domine defensor vindica me.}{Quoniam tribulatio proxima est et non est qui adiuvet.}Circumdederunt.
\end{responsory}%
\V Intende anime mee.
\newline {\color{Red} In Laudibus. Ant.} Dominus deus auxiliator meus et ideo non sum confusus. {\color{Red} Ps.} \psalm{50} {\color{Red} Ant.} Circumdantes circumdederunt me, et in nomine domini vindicabor in eis. {\color{Red} Ps.} \psalm{117} {\color{Red} Ant.} Iudica causam meam, defende quia potens es domine. {\color{Red} Ps.} \psalm{62} {\color{Red} Ant.} Cum angelis et pueris fideles inveniantur, triumphatori mortis clamantes osanna in excelsis. {\color{Red} Can.} \psalm{benedicite} {\color{Red} Ant.} Confundantur qui me persequuntur et non confundar ego domine deus meus. {\color{Red} Ps.} \psalm{148}
\newline {\color{Red} Ad Ben. Ant.} Turba multa que convenerat ad diem festum, clamabat domino, benedictus qui venit in nomine domini.
\newline {\color{Red} Ad Primam. Ant.} Osanna filio David, benedictus qui venit in nomine domini.
\newline {\color{Red} Ad Terciam. Ant.} Pueri Hebreorum vestimenta prosternebant in via, et clamabant dicentes osanna filio David benedictus qui venit in nomine domini. \R \hyperlink{erue-a-framea}{Erue.}
\newline {\color{Red} Ad Sextam. Ant.} Pueri Hebreorum tollentes ramos olivarum obviaverunt domino clamantes et dicentes osanna in excelsis. \R \hyperlink{de-ore-leonis}{De ore.}
\newline {\color{Red} Ad Nonam. Ant.} Occurrunt turbe cum floribus et palmis redemptori obviam, et victori triumphanti digna dant obsequia, filium dei ore gentes predicant, et in laude Christi voces tonant per nubila osanna. \R \hyperlink{principes-persecuti}{Principes.}
\newline {\color{Red} Ad Mag. Ant.} Ceperunt omnes turbe descendentium gaudentes laudare deum voce magna super omnibus quas viderant virtutibus, dicentes benedictus qui venit rex in nomine domini, pax in celo et gloria in excelsis.
\newline \ul{Lectiones per secunda et tercia et quarta feriales a octodecimo capitulo usque ad Trenos ad libitum.}
\newline {\color{Red} Feria Secunda. In Laudibus. Ant.} Faciem meam non averti ab increpantibus et conspuentibus in me. {\color{Red} Ps.} \psalm{50} {\color{Red} Ant.} Framea suscitare adversus eos qui dispergunt gregem meum. {\color{Red} Ps.} \psalm{5} {\color{Red} Ant.} Appenderunt mercedem meam triginta argenteis, quos appreciatus sum ab eis. {\color{Red} Ps.} \psalm{62} {\color{Red} Ant.} Inundaverunt aque super caput meum, dixi perii invocabo nomen tuum domine deus. {\color{Red} Can.} \psalm{isaias} {\color{Red} Ant.} Labia insurgentium et cogitationes eorum vide domine. {\color{Red} Ps.} \psalm{148}
\newline {\color{Red} Ad Ben. Ant.} Non haberes in me potestatem nisi desuper datum tibi fuisset. {\color{Red} Oracio.}
\lettrine[lines=2]{\zallmancaps \color{Blue} D}{a} quesumus omnipotens deus, ut qui in tot adversis ex nostra infirmitate deficimus, intercedente unigeniti filii tui passione respiremus. Per eundem.
\newline {\color{Red} Ad Mag. Ant.} Potestatem habeo ponendi animam meam et iterum sumendi eam. {\color{Red} Oratio.}
\lettrine[lines=2]{\zallmancaps \color{Red} A}{diuva} nos deus salutaris noster, et ad beneficia recolenda quibus nos instaurare dignatus es: tribue venire gaudentes. Per.
\newline {\color{Red} Feria Tercia. In Laudibus. Ant.} Vide domine et considera quoniam tribulor velociter exaudi me. {\color{Red} Ps.} \psalm{50} {\color{Red} Ant.} Discerne causam meam domine ab homine iniquo et doloso eripe me. {\color{Red} Ps.} \psalm{42} {\color{Red} Ant.} Dum tribularer clamavi ad dominum de ventre inferi et exaudivit me. {\color{Red} Ps.} \psalm{62} {\color{Red} Ant.} Domine vim pacior responde pro me, quia nescio quid dicam inimicis meis. {\color{Red} Can.} \psalm{ezechias} {\color{Red} Ant.} Dixerunt impii opprimamus virum iustum quoniam contrarius est operibus nostris. {\color{Red} Ps.} \psalm{148}
\newline {\color{Red} Ad Ben. Ant.} Nemo tollet a me animam meam sed ego pono eam et iterum sumo eam. {\color{Red} Oratio.}
\lettrine[lines=2]{\zallmancaps \color{Blue} O}{mnipotens} sempiterne deus, da nobis ita dominice passionis sacramenta peragere, ut indulgentiam percipere mereamur. Per eundem.
\newline {\color{Red} Ad Mag. Ant.} Cotidie apud vos eram in templo docens et non me tenuistis, et ecce flagellatum ducitis ad crucifigendum. {\color{Red} Oratio.}
\lettrine[lines=2]{\zallmancaps \color{Red} T}{ua} nos misericordia deus et ab omni subreptione vetustatis expurget, et capaces sancte novitatis efficiat. Per.
\newline {\color{Red} Feria Quarta. In Laudibus. Ant.} Libera me de sanguinibus deus deus meus et exultabit lingua mea iusticiam tuam. {\color{Red} Ps.} \psalm{50} {\color{Red} Ant.} Contumelias et terrores passus sum ab eis, et dominus mecum est tanquam bellator fortis. {\color{Red} Ps.} \psalm{64} {\color{Red} Ant.} Ipsi vero in vanum quesierunt animam meam, introibunt in inferiora terre. {\color{Red} Ps.} \psalm{62} {\color{Red} Ant.} Omnes inimici mei audierunt malum meum domine letati sunt quoniam tu fecisti. {\color{Red} Can.} \psalm{annas} {\color{Red} Ant.} Alliga domine in vinculis nationes gencium et reges earum in compedibus. {\color{Red} Ps.} \psalm{148}
\newline {\color{Red} Ad Ben. Ant.} Simon dormis non potuisti una hora vigilare mecum. {\color{Red} Oratio.}
\lettrine[lines=2]{\zallmancaps \color{Blue} P}{resta} quesumus omnipotens deus, ut qui nostris excessibus incessanter affligimur, per unigeniti tui passionem liberemur. Qui tecum.
\newline \ul{Ad Vesperas non dicantur preces, nec fiat prostratio nec aliqua memoria.}
\newline {\color{Red} Ad Mag. Ant.} Tanto tempore vobiscum eram docens vos in templo et non me tenuistis, modo flagellatum ducitis ad crucifigendum. {\color{Red} Oracio.}
\lettrine[lines=2]{\zallmancaps \color{Red} R}{espice} \hypertarget{respice-oratio}{\label{respice-oratio}} quesumus domine super hanc familiam tuam, pro qua dominus noster Ihesus Christus non dubitavit manibus tradi nocencium et crucis subire tormentum. Qui tecum.
\newline \ul{Conplectorium solite more absque prostratione dicatur, et benedictio post Conplectorium detur. In duobus vero diebus sequentibus non detur benedictio.}
\newline \ul{Antiphona de Beata Virgine hac die et duobus sequentibus non dicatur.}
{\ubsubsection{Feria Quinta in Cena Domini}{feria-quinta-in-cena-domini-breviarium}}
\lettrine[lines=2]{\zallmancaps \color{Blue} F}{}\ul{\textsc{eria} Quinta in Cena Domini, et duabus sequentibus ad Matutinas non dicatur} Domine labia, \ul{nec} Deus in adiutorium, \ul{neque invitatorium neque hympnus. Sed dicto} Pater noster, \ul{et} Credo, \ul{postquam signaverint se fratres, quod eciam fiat ad alias horas, statim incipiatur.}
\newline {\color{Red} In Primo Nocturno. Ant.} Zelus domus tue comedit me et obprobria exprobrantium tibi ceciderunt super me. {\color{Red} Ps.} \psalm{68} \ul{Non dicatur,} Gloria patri, \ul{in fine psalmorum, nec canticorum.} {\color{Red} Ant.} Avertantur retrorsum et erubescant qui cogitant michi mala. {\color{Red} Ps.} \psalm{69} {\color{Red} Ant.} Deus meus eripe me de manu peccatoris. {\color{Red} Ps.} \psalm{70} \V Exurge domine. \R Et iudica causam meam. \ul{Non pronuncietur.} Et ne nos, \ul{sed dicto,} Pater noster, \ul{lector incipiat legere, et hoc eciam fiat ante quartam et septimam lectionem.} Iube domne, \ul{et} Tu autem \ul{non dicantur, sed finita lectione sive Lamentatione immediate inchoetur responsorium. Iste modus ad Matutinas in duobus diebus sequentibus observetur.}
\newline {\color{Red} Lectio Prima. Aleph.}
\lettrine[lines=2]{\zallmancaps \color{Red} Q}{uomodo} sedet sola civitas plena populo, facta est quasi vidua domina gentium, princeps provintiarum facta est sub tributo. {\color{Red} Beth.} Plorans ploravit in nocte, et lacrime eius in maxillis eius. Non est qui consoletur eam, ex omnibus caris eius. Omnes amici eius spreverunt eam, et facti sunt ei inimici. Hierusalem Hierusalem convertere ad dominum deum tuum. {\color{Red} Responsorium.}
%pp103
\lettrine[lines=2]{\zallmancaps \color{Blue} I}{n} \hypertarget{in-monte-oliveti}{\label{in-monte-oliveti}} monte Oliveti oravi ad patrem pater si fieri potest transeat a me calix iste, spiritus quidem promptus est caro autem infirma, * Fiat voluntas tua. \V Veruntamen non sicut ego volo sed sicut tu vis. Fiat.
\newline {\color{Red} Lectio Secunda. Gimel.}
\lettrine[lines=2]{\zallmancaps \color{Red} M}{igravit} Iuda propter afflictionem, et multitudinem servitutis, habitavit inter gentes, nec invenit requiem. Omnes persecutores eius apprehenderunt eam inter angustias. {\color{Red} Deleth.} Vie Syon lugent, eo quod non sint qui veniant ad sollemnitatem. Omnes porte eius destructe, sacerdotes eius gementes: virgines eius squalide, et ipsa oppressa amaritudine. Hierusalem Hierusalem, etc.
\begin{responsory}[tristis-est]
{Tristis est anima mea usque ad mortem, sustinete hic et vigilate mecum nunc videbitis turbam que circumdabit me, * Vos fugam capietis et ego vadam immolari pro vobis.}{Ecce appropinquabit hora et filius hominis tradetur in manus peccatorum.}
\end{responsory}%
\newline {\color{Red} Lectio Tertia. He.}
\lettrine[lines=2]{\zallmancaps \color{Blue} F}{acti} sunt hostes eius in capite, inimici eius locupletati sunt, quia Dominus locutus est super eam propter multitudinem iniquitatum eius. Parvuli eius ducti sunt in captivitatem, ante faciem tribulantis. {\color{Red} Vau.} Et egressus est a filia Sion, omnis decor eius. Facti sunt principes eius velut arietes non invenientes pascua, et abierunt absque fortitudine ante faciem subsequentis. Hierusalem Hierusalem.
\begin{responsory}[ecce-vidimus]
{Ecce vidimus eum non habentem speciem neque decorem aspectus eius in eo non est, hic peccata nostra portavit et pro nobis dolens, ipse autem vulneratus est propter iniquitates nostras, * Cuius livore sanati sumus.}{Vere langores nostros ipse tulit et dolores nostros ipse portavit.}Ecce.
\end{responsory}%
\newline {\color{Red} In Secundo Nocturno. Ant.} Liberavit dominus pauperem a potente, et inopem cui non erat adiutor. {\color{Red} Ps.} \psalm{71} {\color{Red} Ant.} Cogitaverunt impii et locuti sunt nequiciam, iniquitatem in excelso locuti sunt. {\color{Red} Ps.} \psalm{72} {\color{Red} Ant.} Exurge domine et iudica causam meam. {\color{Red} Ps.} \psalm{73} \V Deus meus eripe me de manu peccatoris. \R Et de manu contra legem agentis et iniqui.
\newline \ul{Beda in Omelia super Evangelia ante diem festum Pasche.} \grecross
\newline {\color{Red} Lectio Quarta.}
\lettrine[lines=2]{\zallmancaps \color{Red} S}{cripturus} evangelista Iohannes, memorabile illud Domini ministerium, quo discipulis in Pascha priusquam ad passionem iret, pedes lavare dignatus est: primo ipsum nomen Pasche quid mystice signaret, aperire curavit, ita incipiens. Ante diem autem festum Pasche sciens Ihesus quia venit hora eius ut transeat ex hoc mundo ad Patrem. Pascha quippe transitus interpretatur, nomen ex eo vetus habens quod transierit in eo Dominus per Egiptum, percutiens primogenita, et filios Israhel liberans, et quod ipsi filii Israhel transierint illa nocte de Egiptia servitute, ut venirent ad terram promisse olim hereditatis et pacis.
\begin{responsory}[amicus-meus]
{Amicus meus osculi me tradidit signo, quem osculatus fuero ipse est tenere eum hoc malum fecit signum qui per osculum adimplevit homicidium, * Infelix pretermisit precium sanguinis et in fine laqueo se suspendit.}{Filius quidem hominis vadit sicut scriptum est de illo ve autem homini illi per quem tradetur.}
\end{responsory}%
\newline {\color{Red} Lectio Quinta.}
\lettrine[lines=2]{\zallmancaps \color{Blue} M}{ystice} autem significans quod in eo Dominus transiturus esset ex hoc mundo ad Patrem, et quod eius exemplo, fideles abiectis temporalibus desideriis, abiecta viciorum servitute, continuis virtutum studiis transire debent ad promissionem patrie celestis. Quomodo autem transierit Ihesus ex hoc mundo ad Patrem, pulchro sermone designat evangelista cum dicit. Cum dilexisset suos qui erant in mundo, in finem dilexit eos, id est in tantum dilexit, ut ipsa dilectione vitam ad tempus finiret corporalem, mox de morte ad vitam, de hoc mundo transiturus ad Patrem. Unde recte uterque transitus legalis videlicet et evangelicus sanguine consecratus est, iste sanguine fuso in cruce, ille in crucis modum asperso in medio limine.
\begin{responsory}[unus-ex-vobis]
{Unus ex vobis tradet me hodie ve ille per quem tradar ego, * Melius illi erat si natus non fuisset.}{Qui intingit mecum manum in parapside hic me traditurus est in manus peccatorum.}
\end{responsory}%
\newline {\color{Red} Lectio Sexta.}
\lettrine[lines=2]{\zallmancaps \color{Red} E}{t} cena inquit facta, cum diabolus misisset in cor ut traderet eum Iudas Simonis Scariothis, sciens quia omnia dedit ei Pater in manus, et quia a Deo exivit, et ad Deum vadit, surgit a cena, et ponit vestimenta sua. Locuturus de maxima humilitate assumpte humanitatis, prius commemorat eternitatem divine potestatis, ut ipsum verum Deum hominemque demonstraret, et ut nos illius precepti admoneat, ut quanto magni sumus humiliemur in omnibus. Homo erat verus qui hominum pedes lavare, ac per hominem tradi poterat? Deus erat verus, cui omnia dedit Pater in manus, qui a Deo exivit, et ad Deum rediit.
\begin{responsory}[eram-quasi]
{Eram quasi agnus innocens, ductus sum ad immolandum et nesciebam, consilium fecerunt inimici mei adversum me dicentes, * Venite mittamus lignum in panem eius, et eradamus eum de terra viventium.}{Omnes inimici mei adversum me cogitabant mala michi, verbum iniquum mandaverunt adversum me dicentes.}Eram.
\end{responsory}%
\newline {\color{Red} In Tertio Nocturno. Ant.} Dixi iniquis nolite loqui adversus deum iniquitatem. {\color{Red} Ps.} \psalm{74} {\color{Red} Ant.} Terra tremuit et quievit dum resurgeret in iudicium deus. {\color{Red} Ps.} \psalm{75} {\color{Red} Ant.} In die tribulationis mee deum exquisivi manibus meis. {\color{Red} Ps.} \psalm{76} \V Homo pacis mee in quo speravi. \R Ampliavit adversum me supplantationem.
\newline {\color{Red} Lectio VII.}
\lettrine[lines=2]{\zallmancaps \color{Blue} S}{ciebat} Dominus quia diabolus iam miserat in cor Iude ut traderet eum, sciebat quia omnia dedit ei Pater in manus, et inter omnia ipsum traditorem, et eos quibus tradendus erat, et mortem quam erat passurus, ut de his omnibus que vellet. Sciebat quia per humilitatem incarnationis a Deo exivit, et per victoriam resurrectionis ad Deum erat rediturus. Sciebat hec omnia, et tamen in magne exemplum humilitatis surgit a cena, et non Dei Domini sed hominis servi implet officium, eius pedes abluens, cuius manus in sua traditione noverat esse polluendas.
\begin{responsory}[una-hora-non]
{Una hora non potuistis vigilare mecum qui exortabamini mori per me, * Vel Iudam non videtis quomodo non dormit sed festinat tradere me Iudeis.}{Dormite iam et requiescite ecce appropinquabit qui me tradet.}
\end{responsory}%
\newline {\color{Red} Lectio VIII.}
\lettrine[lines=2]{\zallmancaps \color{Red} I}{n} sentencia vero qua Petro dictum est, si non lavero te non habebis partem mecum, aperte monstratum est, quod hec pedum lotio spiritualem purificationem sine qua ad consortium Christi nemo pervenire potest insinuat, non illam que semel in baptismate datur, sed illam qua cotidiani reatus sine quibus in hac vita non vivitur, cotidiana Dei gratia mundantur. Pedes namque quibus terram tangimus, ideoque eos a contagione pulveris custodire nequimus, terrene inhabitationis necessitatem designant, qua in tantum eciam magni viri a celestium contemplatione revocantur, ut si dixerimus quia peccatum non habemus, nosmetipsos seducamus.
\begin{responsory}[seniores-populi]
{Seniores populi consilium fecerunt, * Ut Ihesum dolo tenerent et occiderent, cum gladiis et fustibus exierunt tanquam ad latronem.}{Collegerunt ergo Pontifices et Pharisei concilium.}
\end{responsory}%
\newline {\color{Red} Lectio IX.}
\lettrine[lines=2]{\zallmancaps \color{Blue} O}{b} horum ergo sordidationem simul et emundationem cotidie dicimus, et dimitte nobis debita nostra. Verum hec de apostolis eorumque similibus dicta sunt. At nos qui divini amoris obliti, iter levum incedimus: non illa levi oratione possumus liberari. Sed maior necesse est inquinatio maiori orationum, vigiliarum, ieiuniorum, lacrimarum, atque elemosinarum exercitio purgetur, ipso nimirum interius vestigia cordis nostri mundante, qui sedens in dextra Dei per assumpte habitum humanitatis cotidie interpellat pro nobis.
\begin{responsory}[revelabunt-celi]
{Revelabunt celi iniquitatem Iude et terra adversus eum consurget et manifestum erit peccatum illius in die furoris domini, * Cum iis qui dixerunt domino deo recede a nobis, scienciam viarum tuarum nolumus.}{In diem perdicionis servabitur et ad diem ultionis ducetur.}Revelabunt.
\end{responsory}%
\newline \ul{Finito nono responsorio, non dicatur, \Vbar . ante Laudes, sed statim inchoetur.}
\newline {\color{Red} In Laudibus. Ant.} Iustificeris domine in sermonibus tuis et vincas cum iudicaris. {\color{Red} Ps.} \psalm{50} {\color{Red} Ant.} Dominus tanquam ovis ad victimam ductus est, et non aperuit os suum. {\color{Red} Ps.} \psalm{89} {\color{Red} Ant.} Contritum est cor meum in medio mei, contremuerunt omnia ossa mea. {\color{Red} Ps.} \psalm{62} {\color{Red} Ant.} Exhortatus es in virtute tua et in refectione sancta tua domine. {\color{Red} Can.} \psalm{moysis} {\color{Red} Ant.} Oblatus est quia ipse voluit et peccata nostra ipse portavit. {\color{Red} Ps.} \psalm{148} \ul{Capitulum non dicatur, neque Ymnus, neque Versiculus, sed finita quinta antiphona inchoetur.}
\newline {\color{Red} Ad Ben. Ant.} Traditor autem dedit eis signum dicens quem osculatus fuero ipse est tenete eum.
\newline \ul{Finita antiphona post} Benedictus, \ul{duo fratres stantes ante gradum altaris dicant,} Kyrie eleyson, \ul{et choro idem respondente iterum illi duo fratres dicant,} Kyrie eleyson. \ul{Deinde alii duo fratres stantes in medio chori dicant,} Domine miserere, \ul{et chorus respondeat,} Christus dominus factus est obediens usque ad mortem. \ul{Item fratres qui sunt ante gradus cantent versum,} Qui passurus advenisti propter nos. \ul{Et chorus dicat,} Christe eleyson. \ul{Item fratres ante gradus versum,} Qui expansis in cruce manibus traxisti omnia ad te secula. \ul{Item chorus,} Christe eleyson. \ul{Item fratres ante gradus, \Vbar ,} Qui prophetice prompsisti ero mors tua o mors. \ul{Item chorus,} Christe eleyson. \ul{Deinde fratres in medio chori stantes dicant,} Domine miserere, \ul{et chorus respondeat,} Christus dominus. \ul{Item fratres ante gradus dicant,} Kyrie eleyson, \ul{et chorus} Kyrie eleyson. \ul{Et fratres ante gradus} Kyrie eleyson. \ul{Deinde duo fratres in medio chori stantes dicant,} Domine miserere, \ul{et chorus respondeat,} Christus dominus. \ul{Item fratres ante gradus dicant alta voce, \Vbar ,} Mortem autem crucis. \ul{Finito versu Mortem autem, prosternant se et unusquisque dicat} Pater noster. \ul{Quo dicto duo et duo dicant privatim psalmum,} Miserere, \ul{cum oratione,} \hyperlink{respice-oratio}{Respice quesumus domine,} \ul{sine,} Dominus vobiscum, \ul{et sine} Oremus, \ul{et absque,} Qui tecum, \ul{et non dicatur,} Amen, \ul{nec} Benedicamus, \ul{neque} Fidelium. \ul{Sed dicto post orationem} Pater noster, \ul{surgant omnes.}
\newline \ul{Hac die non fiant prostrationes ad horas propter sollemnitatem Cene, nisi in hoc hora tantum.}
\newline \ul{Usque ad Completorium Sabbati exclusive non dicatur,} Deus in adiutorium, \ul{in principio horarum nec,} Converte nos, \ul{ad Completorium nec} Gloria \ul{in fine psalmorum nec canticorum, nec Capitulum, nec Ymnus, nec Versiculus, nec,} Dominus vobiscum, \ul{nec} Oremus, \ul{nec} Qui tecum, \ul{nec} Benedicamus, \ul{nec} Fidelium, \ul{post horas. More tamen solito dicatur} Pater noster.
\newline \ul{Ad Primam dicto} Pater \ul{et} Credo, \ul{inchoetur.}
\newline {\color{Red} Ant.} Christus factus est pro nobis obediens usque ad mortem, mortem autem crucis. {\color{Red} Ps.} \psalm{53} {\color{Red} Ps.} \psalm{118.1} {\color{Red} Ps.} \psalm{118.2} \ul{Dicta post psalmos antiphona absque,} Kyrie eleyson, \ul{et precibus, et sine psalmo,} \psalm{50} \ul{dicatur Oratio,} \hyperlink{respice-oratio}{Respice.} Confiteor \ul{non dicatur usque ad Completorium Sabbati exclusive. Preciosa solito more dicatur, sine} Gloria patri. \ul{Tercia, Sexta et Nona dicantur supradicto modo cum Ant.} Christus factus, \ul{et oratione,} \hyperlink{respice-oratio}{Respice.}
\newline \ul{Ad Vesperas si missa precesserit immediate finita communione incipiatur,} {\color{Red} Ant.} Calicem salutaris accipiam, et nomen domini invocabo. {\color{Red} Ps.} \psalm{115} {\color{Red} Ant.} Cum iis qui oderunt pacem eram pacificus, cum loquebar illis impugnabant me gratis. {\color{Red} Ps.} \psalm{119} {\color{Red} Ant.} Ab hominibus iniquis libera me domine. {\color{Red} Ps.} \psalm{139} {\color{Red} Ant.} Custodi me a laqueo quem statuerunt michi et a scandalis operantium iniquitatem. {\color{Red} Ps.} \psalm{140} {\color{Red} Ant.} Considerabam ad dexteram et videbam et non erat qui cognosceret me. {\color{Red} Ps.} \psalm{141} \ul{Finita quinta antiphona inchoetur}
\newline {\color{Red} Ad Mag. Ant.} Cenantibus autem accepit Ihesus panem, benedixit ac fregit, dedit discipulis suis.
\newline \ul{Finita post canticum antiphona sequitur,} Dominus vobiscum \ul{et post communio,} Refecti \ul{si missa precesserit. Deinde,} Benedicamus domino. \ul{Qui vero misse non interfuerint, dicant oratione} \hyperlink{respice-oratio}{Respice} \ul{eo modo quo ad alias horas.}
\newline \ul{Ad Completorium non dicatur Lectio,} Fratres sobrii, \ul{sed dicto} Pater noster \ul{inchoetur Ant.} Christus factus, \ul{psalmi} \psalm{4} \psalm{30a} \psalm{90} \psalm{133} \ul{Canti.} \psalm{nunc} \ul{Postea dicatur Ant.} Christus factus est. \ul{Post Ant. Oracio} \hyperlink{respice-oratio}{Respice.}
{\ubsubsection{Feria Sexta in Parascheve}{feria-sexta-in-parascheve-breviarium}}
\lettrine[lines=2]{\zallmancaps \color{Blue} F}{}\ul{\textsc{eria} Sexta in Parascheve ante Matutinas prosternant se fratres et dicto} Pater noster, \ul{et} Credo \ul{surgant. Et postquam signaverint se incipiatur.}
\newline {\color{Red} In Primo Nocturno. Ant.} Astiterunt reges terre et principes convenerunt in unum adversus dominum et adversus Christum eius. {\color{Red} Ps.} \psalm{2} {\color{Red} Ant.} Diviserunt sibi vestimenta mea, et super vestem meam miserunt sortem. {\color{Red} Ps.} \psalm{21} {\color{Red} Ant.} Insurrexerunt in me testes iniqui et mentita est iniquitas sibi. {\color{Red} Ps.} \psalm{26} \V Diviserunt sibi vestimenta mea. \R Et super vestem meam miserunt sortem.
\newline {\color{Red} Lectio Prima. Aleph.}
\lettrine[lines=2]{\zallmancaps \color{Red} Q}{uomodo} obtexit caligine in furore suo Dominus filiam Syon: proiecit de celo in terram inclitam Israhel, et non est recordatus scabelli pedum suorum in die furoris sui? {\color{Red} Beth.} Precipitavit Dominus nec pepercit omnia speciosa Iacob, destruxit in furore suo munitiones virginis Iuda, deiecit in terram. Polluit regnum et principes eius. Hierusalem Hierusalem. \R
\lettrine[lines=2]{\zallmancaps \color{Blue} O}{mnes} \hypertarget{omnes-amici}{\label{omnes-amici}} amici mei dereliquerunt me, et prevaluerunt insidiantes michi, tradidit me quem diligebam, et terribilibus oculis plaga crudeli percuciens, * Aceto potabunt me. \V Et dederunt in escam meam fel, et in siti mea. Aceto.%typo 2nd Acceto
\newline {\color{Red} Lectio Secunda. Gimel.}
\lettrine[lines=2]{\zallmancaps \color{Red} C}{onfregit} in ira furoris sui omne cornu Israhel, avertit retrorsum dexteram suam a facie inimici, et succendit in Iacob quasi ignem flamme devorantis in giro. {\color{Red} Deleth.} Tetendit arcum suum quasi inimicus, firmavit dexteram suam quasi hostis, et occidit omne quod pulchrum erat visu in tabernaculis filie Syon: effudit quasi ignem indignationem suam. Hierusalem Hierusalem.
\begin{responsory}[velum-templi]
{Velum templi scissum est et omnis terra tremuit, latro de cruce clamabat dicens, * Memento mei domine dum veneris in regnum tuum.}{Amen dico tibi, hodie mecum eris in paradiso.}
\end{responsory}%
\newline {\color{Red} Lectio Tertia. He.}
\lettrine[lines=2]{\zallmancaps \color{Blue} F}{actus} est Dominus velut inimicus, precipitavit Israhel. Precipitavit omnia menia eius, dissipavit munitiones eius, et replevit in filia Iuda humiliatum et humiliatam. {\color{Red} Vau.} Et dissipavit quasi ortum tentorium suum, demolitus est tabernaculum suum. Oblivioni tradidit Dominus in Syon festivitatem et sabbatum, et in obprobrium et indignationem furoris sui regem et sacerdotem. Hierusalem Hierusalem.
\begin{responsory}[vinea-mea]
{Vinea mea electa ego te plantavi, * Quomodo conversa es in amaritudinem, ut me crucifigeres et Barrabam dimitteres.}{Ego quidem plantavi te vineam meam electam, omne semen verum.}Vinea.
\end{responsory}%
\newline {\color{Red} In Secundo Nocturno. Ant.} Vim faciebant qui querebant animam meam. {\color{Red} Ps.} \psalm{37} {\color{Red} Ant.} Confundantur et revereantur qui querunt animam meam ut auferant eam. {\color{Red} Ps.} \psalm{39} {\color{Red} Ant.} Alieni insurrexerunt in me et fortes quesierunt animam meam. {\color{Red} Ps.} \psalm{53} \V Ab insurgentibus in me. \R Libera me domine.
\newline \ul{Iohannes episcopus in sermone hodie incipiamus.} \grecross
% https://books.google.com/books?id=H_0tzqg_9Z4C&pg=PA480#v=onepage&q&f=false
\newline {\color{Red} Lectio Quarta.}
\lettrine[lines=2]{\zallmancaps \color{Red} H}{odie} karissimi crux fixa est, et seculum sanctificatum est. Hodie crux fixa est, et mors submersa est. Hodie crux fixa est, et demones dispersi sunt. Hodie crux vicit, et mors devicta est. Hodie diabolus victus est, et homo absolutus est, et Deus glorificatus est. Hodie apud Iudeos terra contremuit, et apud nos orbis terre firmatus est. Hodie apud Iudeos petre scisse sunt: hodie amabilis Deo pietas, universum consuit orbem. Decebat enim ea que pietatis erant coniungi, ea que infidelitatis disiungi.
\begin{responsory}[tanquam-ad-latronem]
{Tanquam ad latronem existis cum gladiis comprehendere me, cotidie apud vos eram in templo docens et non me tenuistis, et ecce flagellatum ducitis, * Ad crucifigendum.}{Filius quidem hominis vadit sicut scriptum est de illo, ve autem homini illi per quem tradetur.}
\end{responsory}%
\newline {\color{Red} Lectio Quinta.}
\lettrine[lines=2]{\zallmancaps \color{Blue} H}{odie} apud Iudeos sol obscuratus est. Non enim ferre poterat creatura iniurias creatoris. Retraxit sol radios suos ne videret impiorum facinora: apud nos autem nox in diem conversa est. Fidelibus enim nox in diem mutatur, infidelibus autem eciam lux ipsa tenebrescit. Hodie Adam eiectus est de paradiso: et hodie latro paradisum ingreditur. Exiit fur, et introivit fur. exivit contemptor verbi, et introivit confitens verbum. Exivit contempnens salutem, et introivit de cruce mercans salutem.
\begin{responsory}[tenebre-facte]
{Tenebre facte sunt dum crucifixissent Ihesum Iudei, et circa horam nonam exclamavit Ihesus voce magna, deus meus ut quid me dereliquisti, * Et inclinato capite emisit spiritum.}{Cum ergo accepisset acetum dixit consummatum est.}
\end{responsory}%
\newline {\color{Red} Lectio Sexta.}
\lettrine[lines=2]{\zallmancaps \color{Red} V}{ide} iustum Dei iudicium. Propter unum peccatum damnatus est Adam, et propter unam fidei vocem, latro salvatus est. Unum peccatum deiecit illum, et una iusticia introduxit istum. Adam factus est exul, et latro habitator factus est paradisi. O admiranda rerum materie. Nullum ante latronem invenies paradisi repromissionem meruisse: non Abraham, non Ysaac, non Iacob, non Moysen, nec prophetas, nec apostolos, sed ante omnes reperies latronem. Hodie inquit dominus, mecum eris in paradiso. Nusquam autem repromissionem paradisi ante latronem invenies.
\begin{responsory}[barrabas-latro]
{Barrabas latro dimittitur, et innocens Christus occiditur, nam et Iudas armis doctus sceleris, qui per pacem didicit facere bellum, * Osculando tradidit dominum Ihesum Christum.}{Ecce turba et qui vocabatur Iudas venit, et dum appropinquaret ad Ihesum.}Barrabas.
\end{responsory}%
\newline {\color{Red} In Tertio Nocturno. Ant.} Ab insurgentibus in me libera me domine quia occupaverunt animam meam. {\color{Red} Ps.} \psalm{58} {\color{Red} Ant.} Longe fecisti notos meos a me traditus sum et non egrediebar. {\color{Red} Ps.} \psalm{87} {\color{Red} Ant.} Captabunt in animam iusti et sanguinem innocentem condempnabunt. {\color{Red} Ps.} \psalm{93} \V Locuti sunt adversum me lingua dolosa. \R Et sermonibus odii circumdederunt me.
\newline {\color{Red} Lectio VII.}
\lettrine[lines=2]{\zallmancaps \color{Blue} N}{ecessarium} est querere, cur latroni ante alios et tales viros, dominus paradisum repromisit. Credidit quidem illi, sed modo diverso Dominum glorie sicut possibile fuit intuentes. Iste autem vidit salvatorem, non super thronum regaliem, non adorari in templo, non loquentem de celis, non per angelos disponentem, sed in pena sociatum cum latrone. Videt in tormento et tanquam in gloria adorat, vidit in cruce, et rogat quasi in celis sedentem. Vidit condemnatum, et regem invocat dicens. Domine memento mei cum veneris in regnum tuum.
\begin{responsory}[tradiderunt-me]
{Tradiderunt me in manus impiorum et inter iniquos proiecerunt me et non pepercerunt anime mee congregati sunt adversum me fortes, * Et sicut gigantes steterunt contra me.}{Astiterunt reges terre et principes convenerunt in unum.}
\end{responsory}%
\newline {\color{Red} Lectio VIII.}
%pp104
\lettrine[lines=2]{\zallmancaps \color{Red} O}{}\ admiranda latronis confessio. Crucifixum vides, et regem predicas? In ligno pendere cernis et celorum regna meditaris? Nunquid nam scripturas legisti ab iniquitate non cessans? Iudei crucifixerunt qui noverunt legem et prophetas, et tu horum ignarus Deum vides crucifixum et adoras condemnatum? Quis te erudivit talia dicere? Non me inquit lex docuit, sed sol occultans lumen suum. Vidi quidem crucifixum, sed terre motum sensi, et propter patricidas elementa indignantia intellexi.
\begin{responsory}[ihesum-tradidit]
{Ihesum tradidit impius summis principibus sacerdotum et senioribus populi, Petrus autem sequebatur a longe, * Ut videret finem.}{Et ingressus est Petrus in atrium principis sacerdotum.}
\end{responsory}%
\newline {\color{Red} Lectio IX.}
\lettrine[lines=2]{\zallmancaps \color{Blue} E}{t} Petre quidem scindebantur, sed corda hominum non metuebant. Et surrexerunt mortui, et intraverunt in sanctam civitatem, ut inexcusabilem redderent causam eorum. Dicebant autem Iudei. Non habemus regem nisi Cesarem. Abnuerunt Dei regnum, et traditi sunt regno Romanorum quod pecierunt, et sanguis Christi factus est illis ad condemnationem: nobis autem ad iustificationem. Et decuit splendore Crucis obcecare infideles, et illuminari credentes. Sola enim Christi crux vere solvit tenebras, et regnum demonum dissipavit, et omnes terrorem malignantium abstulit. Crux sanctitatem prodit crux sol iustitie facta est, ut illuminati misericordia illius qui pependit in ea, perveniamus ad eum qui est lux eterna et vera. Amen.
\begin{responsory}[sicut-ovis]
{Sicut ovis ad occisionem ductus est et dum male tractaretur non aperuit os suum, traditus est ad mortem, * Ut vivificaret populum suum.}{Ipse vulneratus est propter iniquitates nostras, attritus est propter scelera nostra.}Sicut ovis.%typo apperuit
\end{responsory}%
\newline {\color{Red} In Laudibus. Ant.} Proprio filio suo non pepercit deus, sed pro nobis omnibus tradidit illum. {\color{Red} Ps.} \psalm{50} {\color{Red} Ant.} Anxiatus est in me spiritus meus, in me turbatum est cor meum. {\color{Red} Ps.} \psalm{142} {\color{Red} Ant.} Ait latro ad latronem, nos quidem digna factis recipimus, hic autem quid fecit, memento mei domine dum veneris in regnum tuum. {\color{Red} Ps.} \psalm{62} {\color{Red} Ant.} Dum conturbata fuerit anima mea domine misericordie memor eris. {\color{Red} Can.} \psalm{habacuc} {\color{Red} Ant.} Memento mei domine deus dum veneris in regnum tuum. {\color{Red} Ps.} \psalm{148}
\newline {\color{Red} Ad Ben. Ant.} Posuerunt super capud eius causam ipsius scriptam Ihesus Nazarenus rex Iudeorum.
\newline \ul{Finita antiphona sequitur,} Kyrie eleyson, \ul{et cetera sicut in precedenti nocte cum his versibus,} \V Agno miti basia cui lupus dedit venenosa. \V Vita in ligno moritur, infernus et mors lugens spoliatur. \V Te qui vinciri voluisti, nosque a mortis vinculis eripuisti.
\newline {\color{Red} Ad Primam. Ant.} Christus factus. {\color{Red} Ps.} \psalm{53} {\color{Red} Ps.} \psalm{118.1} {\color{Red} Ps.} \psalm{118.2} \ul{finita post psalmos antiphona dicatur,} Pater noster, \ul{et psalmus,} \psalm{50} \ul{Quo dicto statim subsequatur oratio} \hyperlink{respice-oratio}{Respice} \ul{supradicto modo. Hic modus servetur ad Terciam, Sextam, et Nonam, hac die et in Sabbato.}
\newline {\color{Red} Ad Vesperas. Ant.} Calicem. {\color{Red} Ps.} \psalm{115} \ul{Cetere antiphone cum suis psalmis sicut in die Cene.}
\newline {\color{Red} Ad Mag. Ant.} Cum accepisset acetum dixit, consummatum est, et inclinato capite tradidit spiritum.
\newline \ul{Finita post canticum antiphona dicatur} Pater noster, \ul{et psalmus,} \psalm{50} \ul{post psalmum dicatur oratio,} \hyperlink{respice-oratio}{Respice,} \ul{et sic Vespere finiantur.}
\newline \ul{Ad Completorium Ant., psalmi, et canticum sicut in die Cene. Post antiphonam dicatur} Pater noster, \ul{et psalmus,} \psalm{50} \ul{et oracio} \hyperlink{respice-oratio}{Respice.}
{\ubsubsection{In Sabbato Sancto Pasche}{in-sabbato-sancto-pasche-breviarium}}
\lettrine[lines=2]{\zallmancaps \color{Red} I}{}\ul{\textsc{n} Sabbato Sancto Pasche ante Matutinas fiat sicut in sexta feria dictum est, et postmodum incipiatur,}
\newline {\color{Red} In Primo Nocturno. Ant.} In pace in idipsum dormiam et requiescam. {\color{Red} Ps.} \psalm{4} {\color{Red} Ant.} Habitabit in tabernaculo tuo requiescet in monte sancto tuo. {\color{Red} Ps.} \psalm{14} {\color{Red} Ant.} Caro mea requiescet in spe. {\color{Red} Ps.} \psalm{15} \V In pace in idipsum. \R Dormiam et requiescam.
\newline {\color{Red} Lectio Prima. Aleph.}
\lettrine[lines=2]{\zallmancaps \color{Blue} Q}{uomodo} obscuratum est aurum, mutatus est color optimus, dispersi sunt lapides sanctuarii in capite omnium platearum? {\color{Red} Beth.} Filii Syon incliti et amicti auro primo, quomodo reputati sunt in vasa testea, opus manuum figuli? {\color{Red} Gimel.} Sed et lamie nudaverunt mammam, lactaverunt catulos suos. Filia populi mei crudelis, quasi strutio in deserto. Hierusalem Hierusalem. \R
\lettrine[lines=2]{\zallmancaps \color{Red} S}{epulto} \hypertarget{sepulto-domino}{\label{sepulto-domino}} domino signatum est monumentum, volventes lapidem ad ostium monumenti, * Ponentes milites qui custodirent illud. \V Ne forte veniant discipuli eius et furentur eum et dicant plebi surrexit a mortuis. Ponentes.
\newline {\color{Red} Lectio Secunda. Deleth.}
\lettrine[lines=2]{\zallmancaps \color{Blue} A}{dhesit} lingua lactentis ad palatum eius in siti. Parvuli pecierunt panem, et non erat qui frangeret eis. {\color{Red} He.} Qui vescebantur voluptuose, interierunt in viis: qui nutriebantur in croceis, amplexati sunt stercora. {\color{Red} Vau.} Et maior effecta est iniquitas filie populi mei peccato Sodomorum, que subversa est in momento, et non ceperunt in ea manus. Hierusalem Hierusalem.
\begin{responsory}[hierusalem-luge]
{Hierusalem luge et exue te vestibus iocunditatis, induere cinere et cilicio, * Quia in te est occisus salvator Israhel.}{Deduc quasi torrentem lacrimas per diem et noctem et non taceat pupilla oculi tui.}%typo: occuli
\end{responsory}%
\newline {\color{Red} Lectio Tertia. Zai.}
\lettrine[lines=2]{\zallmancaps \color{Red} C}{andidiores} Nazarei eius nive, nitidiores lacte, rubicundiores ebore antiquo, saphyro pulchriores. {\color{Red} Heth.} Denigrata est super carbones facies eorum, et non sunt cogniti in plateis. Adhesit cutis eorum ossibus, aruit, et facta est quasi lignum. {\color{Red} Theth.} Melius fuit occisis gladio quam interfectis fame, quoniam isti extabuerunt consumpti ab sterilitate terre. Hierusalem Hierusalem.
\begin{responsory}[plange-quasi]
{Plange quasi virgo plebs mea, ululate pastores in cinere et cilicio, * Quia veniet dies domini magna et amara valde.}{Ululate pastores et clamate aspergite vos cinere.}Plange.
\end{responsory}%
\newline {\color{Red} In Secundo Nocturno. Ant.} Elevamini porte eternales et introibit rex glorie. {\color{Red} Ps.} \psalm{23} {\color{Red} Ant.} Credo videre bona domini in terra vivencium. {\color{Red} Ps.} \psalm{26} {\color{Red} Ant.} Domine abstraxisti, ab inferis animam meam. {\color{Red} Ps.} \psalm{29} \V Tu autem domine miserere mei. \R Et resuscita me, et retribuam eis.
\newline \ul{Sermo Augustini in Sabbato Sancto.} \grecross
% https://la.wikisource.org/wiki/Sermones_(Augustinus_Incertus)
\newline {\color{Red} Lectio Quarta.}
\lettrine[lines=2]{\zallmancaps \color{Blue} E}{xulta} celum, et in leticia esto terra, quoniam iste dies nobis amplius ex sepulchro radiavit quam de sole refulsit. Dolet infernus quia resolutus est, gaudet quia visitatus est. Resultat quia ignotam lucem post secula longa vidit, et post profunde noctis caliginem respiravit. O pulchra lux que de candido celi fastigio promicasti, et post noctem perpetuam, sedentes in tenebris et umbra mortis subita claritate vestisti. Confestim igitur Christo descendente siluit stridor ille lugentium, dirupta cecidere vincula damnatorum? Attonite mentes obstupescunt tortorum, omnisque simul impia officina contremuit.
\begin{responsory}[recessit-pastor]
{Recessit pastor noster fons aque vive, ad cuius transitum sol obscuratus est, nam et ille captus est qui captivum tenebat primum hominem, * Hodie portas mortis et seras pariter salvator noster dirupit.}{Destruxit quidem claustra inferni, et subvertit potentias diaboli.}
\end{responsory}%
\newline {\color{Red} Lectio Quinta.}
\lettrine[lines=2]{\zallmancaps \color{Red} Q}{uisnam} est inquiunt iste terribilis et niveo splendore coruscus? Nunquam talem noster excepit Tartarus, nunquam in nostram cavernam talem evomit mundus. Invasor est iste non debitor. Effractor est non peccator. Iudicem videmus non supplicem. Pugnare venit non succumbere, eripere non manere. Periit potestas illa cunctis semper populis formidata. Nullus sub cede nostra captivus palpitat, sed quod est gemendum insultat. Nulli iam resonant eiulatus, omnis gemencium mugitus obmutuit. Mox igitur Christus in ipsos crudeles penarum ministros aciem dirigit, atque implacabiles turmas framea divina concidit.
\begin{responsory}[o-vos-omnes]
{O vos omnes qui transitis per viam attendite et videte, * Si est dolor similis sicut dolor meus.}{Attendite universi populi et videte dolorem meum.}
\end{responsory}%
\newline {\color{Red} Lectio Sexta.}
\lettrine[lines=2]{\zallmancaps \color{Blue} T}{ristes} igitur mox lugentesque diuturno squalore turbe, populique vexati concurrunt, et ad salvatoris sui vestigia volutantur, Succurre inquiunt fessis, miserandos resolve gemitus. Redemisti vivos cruce tua, eripe mortuos morte tua. Pari nobiscum sorte mundus perierat, et ad adventum tuum omnis creatura pendebat. Tibi nostra tormenta suspirabant, te semper infernus iste pavebat. Dum hic es absolve reos, dum ascenderis defende tuos. Dum hic es universa claustra solvantur, dum ascenderis non claudantur. Et si redieris ad corpus maiestate tua, non privetur infernus potentia tua. Post auditas preces, post compositas leges, post demersos alta voragine fossa tortores: iudex noster ab inferis laureatus exivit. Sed et leca cum principe suo omnis beatorum toga processit. Repedat igitur ad stadium suum triumphator iterum vivus, ut noverit omnis mundus quod redivit ab inferis dominus Christus.
\begin{responsory}[caligaverunt-oculi]
{Caligaverunt oculi mei a fletu meo quia elongatus est a me qui consolabatur me, videte omnes populi, * Si est dolor similis sicut dolor meus.}{Quos omnes qui transitis per viam attendite et videte.}Caligaverunt.
\end{responsory}%
\newline {\color{Red} In Tertio Nocturno. Ant.} Deus adiuvat me et dominus susceptor est anime mee. {\color{Red} Ps.} \psalm{53} {\color{Red} Ant.} In pace factus est locus eius et in Syon habitatio eius. {\color{Red} Ps.} \psalm{75} {\color{Red} Ant.} Factus sum sicut homo sine adiutorio inter mortuos liber. {\color{Red} Ps.} \psalm{87} \V In pace factus est locus eius. \R Et habitatio eius in Syon.
\newline \ul{Sequens omelia legatur sine titulo scilicet ut nec Secundum Matheum, nec Omelia Beati Hieronimi, dicatur,}
\newline {\color{Red} Lectio VII. Secundum Matheum.}
\lettrine[lines=2]{\zallmancaps \color{Blue} I}{n} illo tempore: Vespere sabbati que lucescit in prima sabbati, venit Maria Magdalene et altera Maria videre sepulchrum. Et reliqua.
\newline {\color{Red} Omelia Beati Hieronimi. Presbiteri.}
\lettrine[lines=2]{\zallmancaps \color{Red} Q}{uod} diversa tempora istarum mulierum in evangeliis describuntur, non mendacii signum est ut impii obiciunt, sed sedule visitationis officia, dum crebro abeunt ac recurrunt, et non paciuntur sepulchro Domini diu abesse vel longius.
\begin{responsory}[ecce-quomodo]
{Ecce quomodo moritur iustus et nemo percipit corde, et viri iusti tolluntur et nemo considerat, a facie iniquitatis sublatus est iustus, * Et erit in pace memoria eius.}{In pace factus est locus eius, et in Syon habitatio eius.}
\end{responsory}%
\newline {\color{Red} Lectio VIII.}
\lettrine[lines=2]{\zallmancaps \color{Blue} E}{t} ecce terremotus factus est magnus, angelus enim Domini descendit de celo. Dominus noster atque idem filius Dei et filius hominis iuxta utramque naturam divinitatis et carnis, nunc magnitudinis sue, nunc humilitatis signa demonstrat.
\begin{responsory}[estimatus-sum]
{Estimatus sum cum descendentibus in lacum, * Factus sum sicut homo sine adiutorio inter mortuos liber.}{Posuerunt me in lacu inferiori, in tenebrosis et in umbra mortis.}
\end{responsory}%
\newline {\color{Red} Lectio IX.}
\lettrine[lines=2]{\zallmancaps \color{Red} U}{nde} in presenti loco quamquam homo sit qui crucifixus est, qui sepultus, qui clausus tumulo, quem lapis oppositus cohibet, tamen que foris aguntur ostendunt filium Dei, sol fugiens, tenebre ingruentes, terra mota, velum scissum, saxa diruta, mortui suscitati, angelorum ministeria que et ab inicio nativitatis eius Deum probabant.
\begin{responsory}[agnus-dei-christus]
{Agnus dei Christus immolatus est pro salute mundi, nam de parentis prothoplasti fraude facta condolens quando pomi noxialis morte morsu corruit ipse lignum tunc notavit, * Damna ligni ut solveret.}{Lustra sex qui iam peracta tempus implens corporis se volente natus ad hoc passioni deditus, agnus in cruce lavatur immolandus stipite.}Agnus.
\end{responsory}%
\newline {\color{Red} In Laudibus. Ant.} O mors ero mors tua, morsus tuus ero inferne. {\color{Red} Ps.} \psalm{50} {\color{Red} Ant.} Plangent eum quasi unigenitum, quia innocens dominus occisus est. {\color{Red} Ps.} \psalm{42} {\color{Red} Ant.} Attendite universi populi et videte dolorem meum. {\color{Red} Ps.} \psalm{62} {\color{Red} Ant.} A porta inferi erue domine animam meam. {\color{Red} Can.} \psalm{ezechias} {\color{Red} Ant.} O vos omnes qui transitis per viam attendite et videte si est dolor sicut dolor meus. {\color{Red} Ps.} \psalm{148} {\color{Red} Ad Ben. Ant.} Mulieres sedentes ad monumentum lamentabantur flentes dominum.
\newline Kyrie eleyson, \ul{et cetera omnia sicut in nocte Cene. Hac die Prima et Tercia et Sexta et Nona, dicantur sicut in die hesterna.}
\newline {\color{Red} Ad Vesperas. Ant.} Alleluya alleluya alleluya, \ul{tota dicatur ante psalmum,} {\color{Red} Ps.} \psalm{116} Gloria patri. \ul{Repetita antiphona post psalmum, sequitur,}
\newline {\color{Red} Ad Mag. Ant.} Vespere autem sabbati que lucescit in prima sabbati venit Maria Magdalene et altera Maria videre sepulchrum, alleluya. \ul{Dicto Mag. et antiphona repetita: sequitur,} Dominus vobiscum. Oremus. {\color{Red} Oratio.}
\lettrine[lines=2]{\zallmancaps \color{Blue} S}{piritum} \hypertarget{spiritum-nobis}{\label{spiritum-nobis}} nobis domine tue caritatis infunde, ut quos sacramentis Paschalibus satiasti, tua facias pietate concordes. Per eiusdem. \ul{Deinde,} Dominus vobiscum, \ul{et} Benedicamus domino alleluya alleluya: \ul{nisi quando Vespere et Missa simul terminantur.}
\newline \ul{Completorium hac die dicatur hoc modo. Dicto,} Confiteor \ul{dicatur} Converte nos. Deus in adiutorium, \ul{postea incipiatur antiphona,} Alleluya. {\color{Red} Ps.} \psalm{4} {\color{Red} Ps.} \psalm{30} {\color{Red} Ps.} \psalm{133} {\color{Red} Can.} \psalm{nunc} \ul{Cum} Gloria patri, \ul{dicantur psalmi et canticum postea sequitur,} {\color{Red} Ant.} Alleluya alleluya alleluya alleluya. \ul{Deinde,} Dominus vobiscum. Oremus. \ul{Oratio} Spiritum nobis \ul{ut supra in Vesperis. Deinde,} Dominus vobiscum, \ul{et} Benedicamus. \ul{Postea sequitur antiphona de Beata Virgine.}
{\ubsubsection{In Die Pasche}{in-die-pasche-breviarium}}
\lettrine[lines=2]{\zallmancaps \color{Red} A}{}\ul{\textsc{d} Matutinas.} Domine labia. Deus in adiutorium. {\color{Red} Invitat.}
\begin{invitatory}[alleluya-invitatorium]
{Alleluya alleluya, * Alleluya.}
\end{invitatory}
\newline \ul{Hac die et per totam ebdomadam usque ad Vesperas sequentis sabbati exclusive ad horas tam diei quam noctis nullus dicatur Ymnus.}
\newline {\color{Red} In Nocturno. Ant.} Ego sum qui sum et consilium meum non est cum impiis sed in lege domini voluntas mea est alleluya. {\color{Red} Ps.} \psalm{1} {\color{Red} Ant.} Postulavi patrem meum alleluya, dedit michi gentes alleluya, in hereditatem alleluya. {\color{Red} Ps.} \psalm{2} {\color{Red} Ant.} Ego dormivi et somnum cepi, et resurrexi quia dominus suscepit me, alleluya, alleluya. {\color{Red} Ps.} \psalm{3} \V Resurrexit dominus alleluya. \R Sicut dixit vobis alleluya.
\newline \ul{Hic versiculus dicatur in nocturno, in Dominica et feria secunda et quinta usque ad Ascensionem quando de tempore agitur.}
\newline {\color{Red} Lectio Prima. Secundum Marcum.}
\lettrine[lines=2]{\zallmancaps \color{Blue} I}{n} illo tempore: Maria Magdalene et Maria Iacobi et Salome emerunt aromata: ut venientes ungerent Ihesum. Et valde mane una sabbatorum veniunt ad monumentum: orto iam sole. Et reliqua.
\newline {\color{Red} Omelia Beati Gregorii Pape.} \grecross
\lettrine[lines=2]{\zallmancaps \color{Red} A}{udistis} fratres karissimi quod sancte mulieres que Dominum fuerant secute cum aromatibus ad monumentum venerunt, et ei quem viventem dilexerant, eciam mortuo studio humanitatis obsequuntur. Sed res gesta aliquid in sancta ecclesia signat gerendum. Sic quippe necesse est, ut audiamus que facta sunt: quatenus cogitemus eciam que nobis sunt ex eorum imitatione facienda. {\color{Red} Responsorium.}
\lettrine[lines=2]{\zallmancaps \color{Blue} A}{ngelus} \hypertarget{angelus-domini-descendit}{\label{angelus-domini-descendit}} domini descendit de celo et accedens revolvit lapidem et super eum sedit, et dixit mulieribus, nolite timere, scio enim quia crucifixum queritis, * Iam surrexit, venite et videte locum ubi positus erat dominus, alleluya. \V Angelus domini locutus est mulieribus dicens quem queritis an Ihesum queritis. Iam.
\newline {\color{Red} Lectio Secunda.}
\lettrine[lines=2]{\zallmancaps \color{Red} E}{t} nos ergo in eum qui est mortuus credentes, si odore virtutum referti, cum opinione bonorum operum Dominum querimus: ad monumentum profecto illius cum aromatibus venimus. Ille autem mulieres angelos vident que cum aromatibus venerunt, quia videlicet ille mentes supernos cives aspiciunt, que cum virtutum odoribus ad Deum per sancta desideria proficiscuntur.
\begin{responsory}[angelus-domini-locutus]
{Angelus domini locutus est mulieribus dicens quem queritis, an Ihesum queritis iam surrexit, * Venite et videte alleluya alleluya.}{Ihesum queritis Nazarenum crucifixum surrexit non est hic.}
\end{responsory}%
\newline {\color{Red} Lectio Tertia.}
\lettrine[lines=2]{\zallmancaps \color{Blue} N}{otandum} vero nobis est, quidnam sit quod in dextris sedere angelus cernitur. Quid namque per sinistram nisi vita presens? quid vero per dexteram nisi perpetua vita signatur? Unde scriptum est. Leva eius sub capite meo, et dextera illius amplexabitur me. Quia igitur redemptor noster iam presentis vite corruptionem transierat, recte angelus qui nunciare perennem eius vitam venerat: in dextera sedebat.
\begin{responsory-final}[dum-transisset]
{Dum transisset Sabbatum Maria Magdalene et Maria Iacobi et Salome emerunt aromata, * Ut venientes ungerent Ihesum, * Alleluya alleluya.}{Et valde mane una sabbatorum veniunt ad monumentum orto iam sole.}{Alleluya}
\end{responsory-final}%
\newline \psalm{tedeum} \ul{et dicatur per totam ebdomadam.}
\newline \V In resurrectione tua Christe, alleluya. \R Celi et terra letentur alleluya.
\newline \ul{Hic versiculus dicatur cotidie ante Laudes usque ad Ascensionem exclusive nisi in festis sanctorum.}
\newline {\color{Red} In Laudibus. Ant.} Angelus autem domini descendit de celo et accedens revolvit lapidem, et sedebat super eum alleluya alleluya. {\color{Red} Ps.} \psalm{92} {\color{Red} Ant.} Et ecce terremotus factus est magnus angelus autem domini descendit de celo alleluya. {\color{Red} Ps.} \psalm{99} {\color{Red} Ant.} Erat autem aspectus eius sicut fulgur, vestimenta eius sicut nix, alleluya alleluya. {\color{Red} Ps.} \psalm{62} {\color{Red} Ant.} Pre timore autem eius exterriti sunt custodes et facti sunt velut mortui alleluya. {\color{Red} Can.} \psalm{benedicite} {\color{Red} Ant.} Respondens autem angelus dixit mulieribus nolite timere scio enim quod Ihesum queritis alleluya. {\color{Red} Ps.} \psalm{148}
\newline \ul{Post predictam antiphonam statim incipiatur,}
\newline {\color{Red} Ad Ben. Ant.} Et valde mane una sabbatorum veniunt ad monumentum orto iam sole, alleluya. {\color{Red} Oratio.}
\lettrine[lines=2]{\zallmancaps \color{Red} D}{eus} \hypertarget{deus-qui-hodierna-pasche}{\label{deus-qui-hodierna-pasche}} qui hodierna die per unigenitum tuum eternitatis nobis aditum devicta morte reserasti, vota nostra que preveniendo aspiras, eciam adiuvando prosequere. Per eundem. Benedicamus domino alleluya alleluya. Deo gratias alleluya alleluya.
\newline Benedicamus domino \ul{cum duplici} alleluya \ul{dicantur in Laudibus et in Vesperis per totam ebdomadam, et per ebdomadam Pentecostes, et in omnibus Dominicis et festis simplicibus et supra usque ad Vesperas sabbati ante Trinitatem exclusive. In profestis autem diebus et festis trium lectionum et infra octavas quando de octavis agitur preterquam in dominica, dicatur tantum cum uno} alleluya. \ul{Ad ceteras autem horas nunquam dicatur, cum} alleluya.
\newline \ul{Ad Primam dicto,} Deus in adiutorium, \ul{et} Gloria patri, \ul{statim inchoetur antiphona similiter, et ad alias horas. Et sciendum quod per totam ebdomadam, usque ad Vesperas sequentis sabbati exclusive, nec Capitula nec Responsoria ad horas nec Versiculi dicuntur nisi in nocturnis et ante Laudes, et in memoriis.} {\color{Red} Ant.} Angelus. \ul{Cetere ad ceteras horas.} {\color{Red} Ps.} \psalm{53} {\color{Red} Ps.} \psalm{118.1} {\color{Red} Ps.} \psalm{118.2} \ul{Finita antiphona post psalmos statim inchoetur} \R Hec dies quam fecit dominus exultemus et letemur in ea, \ul{sine versu. Et sic dicatur cotidie ad omnes horas per ebdomadam usque ad Sabbatum exclusive preterquam ad Matutinas. Ad Vesperas autem dicatur cum versu illius diei.}
\newline \ul{Finito Responsorio,} Hec dies \ul{Ad Primam, dicatur,} Confiteor. \ul{Deinde,} Dominus vobiscum, \ul{et Oratio,} \hyperlink{deus-qui-hodierna-pasche}{Deus qui hodierna,} Dominus vobiscum, Benedicamus domino. \ul{Predicta oratio dicatur ad Primam hac die et duobus sequentibus. Et hoc modo dicatur Prima per totam septimanam nisi quod in quarta feria et deinceps dicatur Oratio,} \hyperlink{in-hac-hora}{In hac hora,} \ul{et in Sabbato loco, Responsorii,} Hec dies, \ul{dicatur} Alleluya, hec dies.
\newline \ul{Ad Vesperas non dicatur,} Deus in adiutorium, \ul{sed dicto} Pater noster, \ul{sic incipiantur,} Kyrie eleyson, Kyrie eleyson, Kyrie eleyson, Christe eleyson, Christe eleyson,
%pp105
Christe eleyson, Kyrie eleyson, Kyrie eleyson, Kyrie eleyson. \ul{Et hoc modo incipiantur usque ad Sabbatum exclusive. Post ultimum,} Kyrie, \ul{statim inchoetur Ant.} Angelus autem. {\color{Red} Ps.} \psalm{109} {\color{Red} Ps.} \psalm{110} {\color{Red} Ps.} \psalm{111} \ul{Predicta antiphona et predicti tres psalmi dicantur ad Vesperas cotidie usque ad Vesperas sequentis sabbati exclusive. Finita antiphona post psalmos sequitur,} \R Hec dies. \V Confitemini domino quoniam bonus, quoniam in seculum misericordia eius. \ul{Responsorium non repetatur. Deinde dicatur,} Alleluya, \V Pascha nostrum immolatus est Christus.
\newline Alleluya \ul{post versum repetatur. Postea sequitur.} {\color{Red} Prosa.}
\lettrine[lines=2]{\zallmancaps \color{Red} V}{ictime} \hypertarget{victime-paschali-laudes}{\label{victime-paschali-laudes}} Paschali laudes immolent Christiani,
Agnus redemit oves, Christus innocens patri reconciliavit peccatores.
Mors et vita duello conflixere mirando, dux vite mortuus regnat vivus.
Dic nobis Maria quid vidisti in via.
Sepulchrum Christi viventis et gloriam vidi resurgentis.
Angelicos testes sudarium et vestes.
Surrexit Christus spes mea, precedet suos in Galilea.
Credendum est magis soli Marie veraci, quam Iudeorum tube fallaci.
Scimus Christum surrexisse ex mortuis vere, tu nobis victor rex miserere, alleluya.
\newline \ul{Precedens Prosa dicatur tantum hac die et duabus sequentibus. Postea sequitur,}
\newline {\color{Red} Ad Mag. Ant.} Et respicientes viderunt revolutum lapidem erat quippe magnus valde, alleluya. \ul{Oratio,} \hyperlink{deus-qui-hodierna-pasche}{Deus qui hodierna.}
\newline \ul{Dicta Oratione, sequitur.} {\color{Red} Ant.} Christus resurgens ex mortuis iam non moritur mors illi ultra non dominabitur, * Quod enim vivit vivit deo, alleluya alleluya. \V Dicant nunc Iudei quomodo milites custodientes sepulchrum perdiderunt regem ad lapidis positionem, quare non servabant Petram iusticie, aut sepultum reddant aut resurgentem adorent nobiscum dicentes. Quod enim. \V Dicite in nationibus alleluya. \R Quia dominus regnavit a ligno alleluya. {\color{Red} Oratio.}
\lettrine[lines=2]{\zallmancaps \color{Blue} D}{eus} \hypertarget{deus-qui-pro-cruce}{\label{deus-qui-pro-cruce}} qui pro nobis filium tuum crucis patibulum subire voluisti, ut inimici a nobis expelleres potestatem, concede nobis famulis tuis, ut resurrectionis gratiam consequamur. Per eundem Christum. \ul{Deinde dicatur,}
\newline {\color{Red} Ant.} Regina celi letare alleluya, quia quem meruisti portare alleluya resurrexit sicut dixit alleluya, ora pro nobis deum alleluya. \V Ora pro nobis sancta dei genitrix alleluya. \R Ut digni. {\color{Red} Oratio.}
\lettrine[lines=2]{\zallmancaps \color{Red} G}{ratiam} tuam quesumus domine mentibus nostris infunde, ut qui angelo nunciante Christi filii tui incarnationem cognovimus, per passionem eius et crucem ad resurrectionis gloriam perducamur. Per eundem dominum nostrum, etc. Dominus vobiscum. Benedicamus domino alleluya alleluya. Deo gratias alleluya alleluya.
\newline \ul{Et hoc modo his tribus diebus Vespere terminentur. Et sciendum, quod Invitatoria et antiphone et responsoria horarum cum suis versiculis, et omnes alii versiculi usque ad festum Sancte Trinitatis exclusive debent terminari cum alleluya.}
\newline {\color{Red} Ad Completorium. Ant.} Alleluya alleluya alleluya alleluya. {\color{Red} Ps.} \psalm{4} {\color{Red} Ps.} \psalm{30a} {\color{Red} Ps.} \psalm{133} \ul{Finita post psalmos antiphona dicatur \Rbar ,} Hec dies. {\color{Red} Ad Nunc Dimittis. Ant.} Alleluya resurrexit dominus alleluya, sicut dixit vobis alleluya alleluya. \ul{Post antiphona Oratio,} \hyperlink{spiritum-nobis}{Spiritum nobis,} \ul{cum,} Benedicamus domino, \ul{hoc modo dicatur Completorium per totam ebdomadam, usque ad Sabbatum exclusive excepto quod predicta oratio non dicitur nisi in duobus diebus sequentibus.}
\newline {\color{Red} Feria Secunda. Ad Matutinas. Invitat.}
\begin{invitatory}[surrexit-dominus-invitatorium]
{Surrexit dominus vere, * Alleluya.}
\end{invitatory}
\newline {\color{Red} Super Nocturnum. Ant.} Surrexit Christus et illuxit populo suo quem redemit sanguine suo alleluya. \ul{Hec antiphona dicatur per ebdomadam super nocturnum.} {\color{Red} Ps.} \psalm{4} {\color{Red} Ps.} \psalm{5} {\color{Red} Ps.} \psalm{6} \V Resurrexit dominus alleluya.
\newline {\color{Red} Lectio Prima. Secundum Lucam.}
\lettrine[lines=2]{\zallmancaps \color{Blue} I}{n} illo tempore: Duo ex discipulis Ihesu ibant ipsa die in castellum quod erat in spatio stadiorum sexaginta ab Hierusalem: nomine Emaus. Et reliqua.
\newline {\color{Red} Omelia Beati Gregorii Pape.}
\lettrine[lines=2]{\zallmancaps \color{Red} I}{n} cotidiana sollemnitate nobis laborantibus pauca loquenda sunt. Et fortasse hec utilius proderunt, quia sepe et alimenta que minus sufficiunt, avidius sumuntur. Ecce audistis fratres karissimi quia duobus discipulis ambulantibus in via, non quidem credentibus sed tamen de se loquentibus Dominus apparuit, sed eis speciem quam recognoscerent non ostendit. {\color{Red} Responsorium.}
\lettrine[lines=2]{\zallmancaps \color{Blue} M}{aria} \hypertarget{maria-magdalene}{\label{maria-magdalene}} Magdalene et altera Maria ibant diluculo ad monumentum, Ihesum quem queritis non est hic, surrexit sicut locutus est, * Precedet vos in Galileam ibi eum videbitis alleluya alleluya. \V Cito euntes dicite discipulis eius et Petro quia surrexit dominus. Precedet.
\newline {\color{Red} Lectio Secunda.}
\lettrine[lines=2]{\zallmancaps \color{Red} H}{oc} ergo egit foris Dominus in oculis corporis, quod apud ipsos agebatur intus in oculis cordis. Ipsi namque apud semetipsos intus et amabant et dubitabant. Eis autem Dominus foris et presens aderat, et quis esset non ostendebat. De se ergo loquentibus presentiam exhibuit, sed de se dubitantibus cognitionis sue speciem abscondit.
\begin{responsory}[surgens-ihesus]
{Surgens Ihesus dominus noster et stans in medio discipulorum suorum dixit, * Pax vobis alleluya gavisi sunt discipuli viso domino alleluya.}{Una ergo sabbatorum cum fores essent clause ubi erant discipuli congregati venit Ihesus et stetit in medio et dixit eis.}
\end{responsory}%
\newline {\color{Red} Lectio Tertia.}
\lettrine[lines=2]{\zallmancaps \color{Blue} V}{erba} quidem contulit, duritiam intellectus increpavit, scripture sacre misteria que de ipso erant aperuit, et tamen quia adhuc in eorum cordibus peregrinus erat a fide, ire se longius finxit. Fingere namque componere dicimus. Unde et compositores luti figulos vocamus. Nichil ergo simplex Veritas per duplicitatem fecit, sed talem se eis exhibuit in corpore: qualis apud illos erat in mente.
\begin{responsory-final}[congratulamini-michi-pasche]
{Congratulamini michi omnes qui diligitis dominum, quia quem querebam apparuit michi, * Et dum flerem ad monumentum vidi dominum meum, * Alleluya.}{Recedentibus discipulis non recedebam et amoris eius igne succensa ardebam desiderio.}{Alleluya}
\end{responsory-final}%
\newline {\color{Red} In Laudibus. Ant.} Angelus autem. \ul{Hec sola antiphona dicatur cotidie per octavas in Laudibus super psalmos. Dicta post psalmos antiphona, statim incipiatur antiphona ad Ben., et hoc per totam ebdomadam observetur.}
\newline {\color{Red} Ad Ben. Ant.} Qui sunt hi sermones quos confertis ad invicem ambulantes et estis tristes, alleluya alleluya. {\color{Red} Oratio.}
\lettrine[lines=2]{\zallmancaps \color{Red} D}{eus} qui sollemnitate Paschali mundo remedia contulisti, populum tuum quesumus celesti dono prosequere, ut et perfectam libertatem consequi mereatur, et ad vitam proficiat sempiternam. Per.
\newline \ul{Hac die et sequentibus diebus per ebdomadam nisi in Sabbato: Tercia et Sexta et Nona dicantur sicut in die Pasche cum propriis orationibus.}
\newline {\color{Red} Ad Vesperas.} \R Hec dies. \V Dicat nunc Israhel quoniam bonus, quoniam in seculum misericordia eius, Alleluya. \V Nonne cor nostrum ardens erat in nobis de Ihesu dum nobis loqueretur in via. {\color{Red} Sequentia.} \hyperlink{victime-paschali-laudes}{Victime.}
\newline {\color{Red} Ad Mag. Ant.} Nonne sic oportuit pati Christum et intrare in gloriam suam alleluya. \ul{Postea sequitur antiphona,} Christus \ul{et cetera sicut in precedenti die.}
\newline {\color{Red} Feria Tercia. In Nocturno. Ps.} \psalm{7} {\color{Red} Ps.} \psalm{8} {\color{Red} Ps.} \psalm{10} \V Surrexit dominus vere alleluya. \R Et apparuit Simoni alleluya.
\newline \ul{Hic versiculus dicatur in nocturno feria tercia et sexta usque ad Ascensionem quando de tempore agitur.}
\newline {\color{Red} Lectio Prima. Secundum Lucam.}
\lettrine[lines=2]{\zallmancaps \color{Blue} I}{n} illo tempore: Stetit Ihesus in medio discipulorum suorum, et dicit eis. Pax vobis. Et reliqua.
\newline {\color{Red} Omelia Venerabilis Bede Presbiteri.}
\lettrine[lines=2]{\zallmancaps \color{Red} G}{loriam} resurrectionis sue paulatim discipulis et per incrementa temporis dominus et redemptor noster ostendit, quia nimirum tanta erat virtus miraculi, ut hanc repente totam capere fragilia mortalium pectora non possent.
\begin{responsory}[surrexit-dominus-de]
{Surrexit dominus de sepulchro alleluya, * Qui pro nobis pependit in ligno alleluya.}{Letentur celi et exultet terra ante faciem domini.}
\end{responsory}%
\newline {\color{Red} Lectio Secunda.}
\lettrine[lines=2]{\zallmancaps \color{Blue} C}{onsulens} igitur ipse imbecillitati querentium, primo venientibus ad monumentum et feminis suo amore ferventibus et viris, revolutum lapidem, et ablato suo corpore linteamina quibus involutum erat, sola posita monstravit. Deinde curiosius inquirentibus feminis, et de eo quod invenerant mente consternatis, angelorum visionem ostendit, qui illum resurrexisse certa manifestatione patefacerent, ac sic precurrente fama patrate resurrectionis, tandem ipse Dominus virtutum et rex glorie apparens, quanta potentia mortem quam ad horam gustaverat, devicisset, aperuit.
\begin{responsory}[tulerunt-dominum]
{Tulerunt dominum meum et nescio ubi posuerunt eum, dicit ei angelus noli flere Maria, * Surrexit sicut dixit, precedet vos in Galileam ibi eum videbitis alleluya alleluya.}{Dum ergo fleret inclinavit se et prospexit in monumentum et vidit duos angelos in albis sedentes qui dicunt ei.}
\end{responsory}%
\newline {\color{Red} Lectio Tertia.}
\lettrine[lines=2]{\zallmancaps \color{Red} Q}{uantum} autem ex evangelice lectionis serie invenimus, quinquies ipsa qua resurrexit die visus est hominibus. Primo Marie Magdalene ad monumentum, quando ei desideranti pedes eius amplecti, dictum est, noli me tangere, nondum enim ascendi ad Patrem meum. Deinde apparuit duabus a monumento currentibus nunciare discipulis que ab angelis de perfecta eius resurrectione didicerant, de quibus scriptum est, quia accesserunt et tenuerunt pedes eius, et adoraverunt eum.
\begin{responsory-final}[expurgate-vetus]
{Expurgate vetus fermentum ut sitis nova conspersio, etenim Pascha nostrum immolatus est Christus, * Itaque epulemur in domino alleluya.}{Non in fermento malicie et nequitie sed in azimis sinceritatis et veritatis.}{Alleluya}
\end{responsory-final}%
\newline {\color{Red} Ad Ben. Ant.} Stetit Ihesus in medio discipulorum suorum et dixit eis, pax vobis alleluya alleluya. {\color{Red} Oratio.}
\lettrine[lines=2]{\zallmancaps \color{Blue} D}{eus} qui ecclesiam tuam novo semper fetu multiplicas, concede famulis tuis, ut sacramentum vivendo teneant quod fide perceperunt. Per.
\newline {\color{Red} Ad Vesperas. Ant.} \R Hec dies. \V Dicant nunc qui redempti sunt a domino quos redemit de manu inimici, de regionibus congregavit eos. Alleluya. \V Surrexit dominus et occurrens mulieribus ait, avete, tunc accesserunt et tenuerunt pedes eius. {\color{Red} Sequentia.} \hyperlink{victime-paschali-laudes}{Victime.}
\newline {\color{Red} Ad Mag. Ant.} Videte manus meas et pedes meos quia ego ipse sum alleluya alleluya.
\newline \ul{Post orationem sequitur Ant.} Christus \ul{et cetera sicut in die Pasche.}
\newline {\color{Red} Feria Quarta. In Nocturno. Ps.} \psalm{11} {\color{Red} Ps.} \psalm{12} {\color{Red} Ps.} \psalm{13} \V Surrexit dominus de sepulchro alleluya. \R Qui pro nobis pependit in ligno alleluya.
\newline \ul{Hic versiculus dicatur in nocturno feria quarta et Sabbato usque ad Ascensionem quando de tempore agitur.}
\newline {\color{Red} Lectio Prima. Secundum Iohannem.}
\lettrine[lines=2]{\zallmancaps \color{Blue} I}{n} illo tempore: Manifestavit se iterum Ihesus discipulis suis ad mare Tyberiadis: manifestavit autem sic. Erant simul Simon Petrus et Thomas qui dicitur Didimus et Nathanel qui erat Chana Galilee, et filii Zebedei, et alii ex discipulis eius duo. Et reliqua.
\newline {\color{Red} Omelia Beati Gregorii Pape.}
\lettrine[lines=2]{\zallmancaps \color{Red} L}{ectio} sancti evangelii que modo in auribus vestris lecta est fratres karissimi questione animum pulsat, sed pulsatione sua vim discretionis indicat. Queri etenim potest cur Petrus qui piscator ante conversionem fuit, post conversionem ad piscationem rediit. Et cum Veritas dicat nemo mittens manum suam in aratrum et aspiciens retro aptus est regno Dei, cur repetiit quod dereliquit? {\color{Red} Responsorium.}
\lettrine[lines=2]{\zallmancaps \color{Blue} E}{cce} \hypertarget{ecce-vicit-leo}{\label{ecce-vicit-leo}} vicit leo de tribu Iuda radix David aperire librum et solvere septem signacula eius, * Alleluya alleluya alleluya. \V Et unus de senioribus dixit michi ne fleveris, dignus est agnus qui occisus est accipere virtutem et fortitudinem. Alleluya.
\newline {\color{Red} Lectio Secunda.}
\lettrine[lines=2]{\zallmancaps \color{Red} S}{ed} si virtus discretionis inspicitur, citius videtur quia nimirum negotium quod ante conversionem sine peccato extitit, hoc eciam post conversionem repetere culpa non fuit. Nam piscatorem Petrum, Matheum, vero thelonearium scimus. Post conversionem suam Petrus ad piscationem rediit, Matheus vero ad thelonei negotium non resedit, quia aliud est victum per piscationem querere, aliud autem thelonei lucris pecunias augere.
\begin{responsory}[si-consurrexistis]
{Si consurrexistis cum Christo que sursum sunt querite ubi Christus est in dextera dei sedens, * Que sursum sunt sapite non que super terram, alleluya.}{Mortui enim estis et vita vestra abscondita est cum Christo.}
\end{responsory}%
\newline {\color{Red} Lectio Tertia.}
\lettrine[lines=2]{\zallmancaps \color{Blue} S}{unt} enim pleraque negotia que sine peccato exerceri aut vix aut nullatenus possunt. Que ergo ad peccatum implicant, ad hec necesse est, ut post conversionem animus non recurrat. Queri eciam potest cur discipulis in mari laborantibus post resurrectionem suam Dominus in littore stetit, qui ante resurrectionem suam coram discipulis suis in maris fluctibus ambulavit.
\begin{responsory-final}[surrexit-pastor]
{Surrexit pastor bonus qui posuit animam suam pro ovibus suis, * Et pro suo grege mori dignatus est, * Alleluya alleluya.}{Etenim Pascha nostrum immolatus est Christus.}{Alleluya}
\end{responsory-final}%
\newline {\color{Red} Ad Ben. Ant.} Mittite in dexteram navigii rete et invenietis, alleluya. {\color{Red} Oracio.}
\lettrine[lines=2]{\zallmancaps \color{Red} D}{eus} qui nos resurrectionis dominice annua sollemnitate letificas concede propicius, ut per temporalia festa que agimus, pervenire ad gaudia eterna mereamur. Per eundem.
\newline {\color{Red} Ad Vesperas.} \R Hec dies. \V Dextera domini fecit virtutem, dextera domini exaltavit me. {\color{Red} \bfseries A}lleluya. \V Christus resurgens ex mortuis iam non moritur mors illi ultra non dominabitur.
\newline {\color{Red} Ad Mag. Ant.} Hoc iam tercio manifestavit se Ihesus postquam resurrexit a mortuis alleluya.
\newline \ul{Post orationem fiat memoria de cruce sicut in die Pasche excepto quod versus antiphone scilicet} Dicant nunc Iudei, \ul{non dicatur, et hoc cotidie fiat in Vesperis usque ad Vesperas sabbati exclusive, memoria vero de Beata Virgine non fiat.}
\newline {\color{Red} Feria Quinta. In Nocturno. Ps.} \psalm{14} {\color{Red} Ps.} \psalm{15} {\color{Red} Ps.} \psalm{16} \V Resurrexit dominus.
\newline {\color{Red} Lectio Prima. Secundum Matheum.}
\lettrine[lines=2]{\zallmancaps \color{Blue} I}{n} illo tempore: Maria stabat ad monumentum foris plorans, et reliqua.
\newline {\color{Red} Omelia Beati Gregorii Pape.}
\lettrine[lines=2]{\zallmancaps \color{Red} M}{aria} Magdalene que fuerat in civitate peccatrix, amando veritatem lavit lacrimis maculas criminis, et vox Veritatis impletur qua dicitur. Dimissa sunt ei peccata multa, quoniam dilexit multum. Que enim prius frigida peccando remanserat, postmodum amando fortiter ardebat. \ul{Responsorium sicut in secunda feria.}
\newline {\color{Red} Lectio Secunda.}
\lettrine[lines=2]{\zallmancaps \color{Blue} N}{am} postquam venit ad monumentum, ibique corpus dominicum non invenit, sublatum credidit, atque discipulis nunciavit. Qui venientes viderunt atque ita esse ut mulier dixerat crediderunt. Et de eis protinus scriptum est. Abierunt ergo discipuli ad semetipsos. Deinde subiungitur. Maria stabat ad monumentum foris plorans. Qua in re pensandum est huius mulieris mentem quanta vis amoris accenderat, que a monumento Domini eciam discipulis recedentibus non recedebat.
\newline {\color{Red} Lectio Tertia.}
\lettrine[lines=2]{\zallmancaps \color{Red} E}{xquirebat} Maria quem non invenerat. Flebat inquirendo, et amoris sui igne succensa, eius quem ablatum credidit ardebat desiderio. Unde contigit, ut eum sola tunc videret que remansit ut quereret, quia nimirum virtus boni operis perseverantia est.
\newline {\color{Red} Ad Ben. Ant.} Maria stabat ad monumentum plorans, vidit duos angelos in albis sedentes, et sudarium quod fuerat super capud Ihesu alleluya.
\lettrine[lines=2]{\zallmancaps \color{Blue} D}{eus} qui diversitatem gentium in confessione tui nominis adunasti, da ut renatis fonte baptismatis, una sit fedes mentium et pietas actionum. Per.
\newline {\color{Red} Ad Vesperas.} \R Hec dies. \V Lapidem quem reprobaverunt edificantes hic factus est in caput anguli a domino factum est et est mirabile in oculis nostris. {\color{Red} \bfseries A}lleluya. \V In die resurrectionis mee dicit dominus precedam vos in Galileam.
\newline {\color{Red} Ad Mag. Ant.} Tulerunt dominum meum et nescio ubi posuerunt eum, si tu sustulisti eum dicito michi alleluya, et ego eum tollam alleluya.
\newline {\color{Red} Feria Sexta. In Nocturno. Ps.} \psalm{18} {\color{Red} Ps.} \psalm{19} {\color{Red} Ps.} \psalm{20} \V Surrexit dominus vere.
\newline {\color{Red} Lectio Prima. Secundum Matheum.}
\lettrine[lines=2]{\zallmancaps \color{Red} I}{n} illo tempore: Undecim discipuli abierunt in Galileam, in montem ubi constituerat illis Ihesus. Et reliqua.
\newline {\color{Red} Omelia Beati Hieronimi Presbiteri.}
\lettrine[lines=2]{\zallmancaps \color{Red} P}{ost} resurrectionem Ihesus in Galilea in monte conspicitur, ibique adoratur, licet quidam dubitent, et dubitatio eorum nostram augeat fidem. Tunc manifestius ostendit Thome et latus lancea vulneratum, et manus fixas demonstrat clavis. \ul{Responsoria ut in tercia feria.}
\newline {\color{Red} Lectio Secunda.}
\lettrine[lines=2]{\zallmancaps \color{Blue} E}{t} accedens Ihesus: locutus est eis dicens. Data est michi omnis potestas, in celo et in terra. Illi potestas data est, qui paulo ante crucifixus, qui sepultus in tumulo, qui mortuus iacuerat, qui postea resurrexit. In celo autem et in terra ei potestas data est, ut qui ante regnabat in celo: per fidem credentium regnet in terris.
\newline {\color{Red} Lectio Tertia.}
\lettrine[lines=2]{\zallmancaps \color{Red} E}{untes} ergo docete omnes gentes, baptizantes eos in nomine Patris et Filii et Spiritus sancti. Primum docent omnes gentes, deinde doctas intingunt aqua. Non enim potest fieri ut corpus baptismi recipiat sacramentum, nisi ante anima fidei susceperit veritatem.
\newline {\color{Red} Ad Ben. Ant.} Undecim discipuli in Galilea videntes dominum adoraverunt alleluya. {\color{Red} Oratio.}
\lettrine[lines=2]{\zallmancaps \color{Blue} O}{mnipotens} sempiterne deus qui Paschale sacramentum in reconciliationis humane federe contulisti, da mentibus nostris ut quod professione celebramus imitemur effectu. Per.
\newline {\color{Red} Ad Vesperas.} \R Hec dies. \V Benedictus qui venit in nomine domini, deus dominus et illuxit nobis. {\color{Red} \bfseries A}lleluya. \V Angelus domini descendit de celo et accedens revolvit lapidem et sedebat super eum.
\newline {\color{Red} Ad Mag. Ant.} Data est michi omnis potestas in celo et in terra alleluya alleluya.
\newline {\color{Red} Sabbato. In Nocturno. Ps.} \psalm{22} {\color{Red} Ps.} \psalm{23} {\color{Red} Ps.} \psalm{25} \V Surrexit dominus de sepulchro.
\newline {\color{Red} Lectio Prima. Secundum Iohannem.}
\lettrine[lines=2]{\zallmancaps \color{Red} I}{n} illo tempore: Una sabbati Maria Magdalene venit mane cum adhuc tenebre essent ad monumentum, et vidit lapidem sublatum a monumento. Et reliqua.
\newline {\color{Red} Omelia Beati Gregorii Pape.} \grecross
\lettrine[lines=2]{\zallmancaps \color{Blue} L}{ectio} sancti evangelii quam modo fratres audistis valde in superficie hystorica aperta est, sed eius nobis sunt mysteria sub brevitate requirenda. Maria Magdalene cum adhuc tenebre essent venit ad monumentum. \ul{Responsoria sicut in feria quarta.}
\newline {\color{Red} Lectio Secunda.}
\lettrine[lines=2]{\zallmancaps \color{Red} I}{uxta} hystoriam notatur hora: iuxta intellectum vero mysticum requirentis signatur intelligentia. Maria enim actorem omnium quem carne viderat mortuum, querebat in monumento. Et quia hunc minime invenit: furatum credidit.
\newline {\color{Red} Lectio Tertia.}
\lettrine[lines=2]{\zallmancaps \color{Blue} A}{dhuc} ergo erant tenebre, cum venit ad monumentum. Cucurrit citius, discipulis nuntiavit. Sed illi pre ceteris cucurrerunt, qui pre ceteris amaverunt: Petrus videlicet et Iohannes. Currebant autem duo simul: sed Iohannes precucurrit citius Petro. Venit prior ad monumentum, sed ingredi non presumpsit. Venit vero posterior Petrus, et intravit.
\newline {\color{Red} Ad Ben. Ant.} Currebant duo simul et ille alius discipulus precucurrit citius Petro et venit prior ad monumentum alleluya.
%pp106
\lettrine[lines=2]{\zallmancaps \color{Red} C}{oncede} quesumus omnipotens deus, ut qui festa Paschalia venerando egimus, per hec contingere ad gaudia eterna mereamur. Per.
\newline \ul{Hac die non dicatur Responsorium,} Hec dies \ul{ad horas sed ad Primam et alias horas dicatur post antiphonam loco Responsorii,} {\bfseries \color{Blue} A}lleluya. \V Hec dies quam fecit dominus exultemus et letemur in ea. \ul{Post versum,} Alleluya, \ul{non repetatur.}
{\ubsubsection{Dominica in Octavis Pasche Sabbato Precedenti}{dominica-in-octavis-pasche-sabbato-precedenti-breviarium}}
\newline {\color{Red} Ad Vesperas.} Deus in adiutorium. {\color{Red} Super psalmos Ant.} Alleluya alleluya alleluya. {\color{Red} Ps.} \psalm{143} {\color{Red} Ps.} \psalm{144} {\color{Red} Ps.} \psalm{145} {\color{Red} Ps.} \psalm{146} {\color{Red} Ps.} \psalm{147} {\color{Red} Capitulum.}
\lettrine[lines=2]{\zallmancaps \color{Blue} C}{hristus} \hypertarget{christus-resurgens-capitulum}{\label{christus-resurgens-capitulum}} resurgens ex mortuis iam non moritur, mors illi ultra non dominabitur, quod enim mortuus est peccato mortuus est semel quod autem vivit, vivit deo. {\color{Red} Ymnus.}
\hymn{ad-cenam-agni}{Ad cenam agni providi}{
\lettrine[lines=2]{\zallmancaps \color{Red} A}{d} cenam agni providi
et stolis albis candidi,
post transitum maris rubri
Christo canamus principi.
\par {\bfseries \color{Blue} C}uius corpus sanctissimum
in ara crucis torridum,
cruore eius roseo
gustando vivimus deo.
\par {\bfseries \color{Red} P}rotecti Pasche vespere
a devastante angelo,
erepti de durissimo
Pharaonis imperio.
\par {\bfseries \color{Blue} I}am Pascha nostrum Christus est,
qui immolatus agnus est
sinceritatis azima
caro eius oblata est.
\par {\bfseries \color{Red} O} vere digna hostia
per quam fracta sunt Tartara, %fca same as facta
redempta plebs captivata
reddita vite premia.
\par {\bfseries \color{Blue} C}onsurgit Christus tumulo,
victor redit de baratro,
tyrannum trudens vinculo %typo: tyramnum
et reserans paradisum.
\par {\bfseries \color{Red} Q}uesumus actor omnium
in hoc Paschali gaudio
ab omni mortis impetu
tuum defende populum.
\par {\bfseries \color{Blue} G}loria tibi domine
qui surrexisti a mortuis
cum patre et sancto spiritu,
in sempiterna secula. Amen.
}
\newline \ul{Isti duo versus scilicet,} Quesumus actor, \ul{et} Gloria tibi domine qui surrexisti, \ul{dicantur in fine omnium Ymnorum eiusdem metri usque ad Primas Vesperas Ascensionis, nisi in Ymnis,} \hyperlink{quem-terra}{Quem terra,} \ul{et} \hyperlink{o-gloriosa}{O Gloriosa,} \ul{quando Annunciatio post Pascha celebratur, et nisi in Ymnis} Beati Petri Martiris, \ul{et} Corone Domini, \ul{qui predicti omnes dicantur secundum quod in suis locis signatum est.}
\newline \V Mane nobiscum domine alleluya. \R Quoniam advesperascit alleluya.
\newline \ul{Hic ordo servetur omni die sabbati ad Vesperas usque ad Ascensionem scilicet de antiphona super psalmos, et de psalmis et de Capitulo, et de Ymno, et de Versiculo, quando de Dominica agitur.}
\newline {\color{Red} Ad Mag. Ant.} Cum esset sero die illa una sabbatorum et fores essent clause ubi erant discipuli congregati, stetit Ihesus in medio et dixit eis, pax vobis alleluya. {\color{Red} Oratio.}
\lettrine[lines=2]{\zallmancaps \color{Blue} P}{resta} quesumus omnipotens deus, ut qui Paschalia festa peregimus, hec te largiente moribus et vita teneamus. Per.
\newline{Postea fiat memoria de cruce.} {\color{Red} Ant.} Crucem sanctam subiit qui infernum confregit accinctus est potentia surrexit die tercia alleluya. \V Dicite in nationibus alleluya. \ul{Oratio,} \hyperlink{deus-qui-pro-cruce}{Deus qui pro nobis.}
\newline {\color{Red} Ad Completorium. Ant.} Alleluya alleluya alleluya alleluya. {\color{Red} Ps.} \psalm{4} \psalm{30a} \psalm{90} \psalm{133} \ul{Predicta antiphona dicatur ad Completorium super psalmos per totum tempus Paschale nisi in Annunciatione quando post Pascha celebrabitur.} {\color{Red} Cap.} Tu in nobis.
\begin{responsory-final}[in-manus-tuas-breve-pasche]
{In manus tuas domine commendo spiritum meum, * Alleluya alleluya.}{Redemisti me domine deus veritatis.}{In manus}
\end{responsory-final}%
{\color{Red} Ymnus.}
\hymn{ihesu-nostra-redemptio}{Ihesu nostra redemptio}{
\lettrine[lines=2]{\zallmancaps \color{Red} I}{hesu} nostra redemptio,
amor et desiderium,
deus creator omnium,
homo in fine temporum.
\par {\bfseries \color{Blue} Q}ue te vicit clementia
ut ferres nostra crimina,
crudelem mortem patiens,
ut nos a morte tolleres.
\par {\bfseries \color{Red} I}nferni claustra penetrans,
tuos captivos redimens,
victor triumpho nobili
ad dexteram patris residens.
\par {\bfseries \color{Blue} I}psa te cogat pietas,
ut mala nostra superes
parcendo, et voti compotes
nos tuo vultu saties.
}
Quesumus actor. Gloria tibi.
\newline \ul{Iste Ymnus dicatur usque ad festum Sancte Trinitatis exclusive.}
\newline \V Custodi. {\color{Red} Ad Nunc. Ant.} Alleluya resurrexit. \ul{Predicta antiphona dicatur ad} \psalm{nunc} \ul{usque ad Ascensionem exclusive nisi in Annuntiatione quando post Pascha celebrabitur. Oratio} \hyperlink{visita-quesumus}{Visita.}
\newline {\color{Red} Ad Matutinas. Invitat.}
\begin{invitatory}[alleluya-surrexit-invitatorium]
{Alleluya, surrexit domine vere, * Venite adoremus alleluya.}
\end{invitatory}
\newline \ul{Hoc invitatorium dicatur singulis Dominicis diebus usque ad Ascensionem quando de tempore agitur.} {\color{Red} Ymnus.}
\hymn{aurora-lucis}{Aurora lucis rutilat}{
\lettrine[lines=2]{\zallmancaps \color{Blue} A}{urora} lucis rutilat
celum laudibus intonat,
mundus exultans iubilat,
gemens infernus ululat.
\par {\bfseries \color{Red} C}um rex ille fortissimus
mortis confractis viribus
pede conculcans Tartara
soluit a pena miseros.
\par {\bfseries \color{Blue} I}lle qui clausus lapide
custoditur sub milite,
triumphans pompa nobili
victor surgit de funere.
\par {\bfseries \color{Red} S}olutis iam gemitibus
et inferni doloribus,
quia surrexit dominus
resplendens clamat angelus.
\par {\bfseries \color{Blue} T}ristes erant apostoli
de nece sui domini,
quem pena mortis crudeli
servi damnarant impii.
\par {\bfseries \color{Red} Q}uesumus actor omnium
in hoc Paschali gaudio
ab omni mortis impetu
tuum defende populum.
\par {\bfseries \color{Blue} G}loria tibi domine
qui surrexisti a mortuis
cum patre et sancto spiritu,
in sempiterna secula. Amen.
}
\newline \ul{Antiphone psalmi, \Vbar , Responsoria, sicut in die Pasche.}
\newline {\color{Red} Lectio Prima. Secundum Iohannem.}
\lettrine[lines=2]{\zallmancaps \color{Red} I}{n} illo tempore: Cum esset sero die illo una sabbatorum et fores essent clause, ubi erant discipuli congregati propter metum Iudeorum, venit Ihesus et stetit in medio et dicit eis. Pax vobis, et reliqua.
\newline {\color{Red} Omelia Beati Gregorii Pape.}
\lettrine[lines=2]{\zallmancaps \color{Blue} P}{rima} lectionis huius evangelice questio animum pulsat, quomodo post resurrectionem corpus dominicum verum fuit quod clausis ianuis ad discipulos ingredi potuit. Sed sciendum nobis est quod divina operatio si ratione comprehenditur non est admirabilis, nec fides habet meritum cui humana ratio prebet experimentum.
\newline {\color{Red} Lectio Secunda.}
\lettrine[lines=2]{\zallmancaps \color{Red} S}{ed} hec ipsa redemptoris nostri opera que ex semetipsis comprehendi nequaquam possunt, ex alia eius operatione pensanda sunt, ut rebus mirabilibus fidem prebeant facta mirabiliora. Illud enim corpus Domini intravit ad discipulos ianuis clausis, quod videlicet ad humanos oculos per nativitatem suam clauso exiit utero Virginis.
\newline {\color{Red} Lectio Tertia.}
\lettrine[lines=2]{\zallmancaps \color{Blue} Q}{uid} ergo mirum si ianuis clausis post resurrectionem suam in eternum iam victurus intravit, qui moriturus veniens non aperto utero Virginis exivit? Sed quia ad illud corpus quod videri poterat fides intuentium dubitabat, ostendit eis protinus manus et latus. Palpandam carnem prebuit, quam clausis ianuis introduxit.
\newline \V In resurrectione tua. \ul{In Laudibus antiphone sicut in die Pasche que etenim dicantur ad horas.} {\color{Red} Cap.} \hyperlink{christus-resurgens-capitulum}{Christus resurgens.} {\color{Red} Ymnus.}
\hymn{sermone-blando}{Sermone blando angelus}{
\lettrine[lines=2]{\zallmancaps \color{Red} S}{ermone} blando angelus
predixit mulieribus
in Galilea dominus
videndus est quam totius.
\par {\bfseries \color{Blue} I}lle dum pergunt concite
apostolis hoc dicere,
videntes eum vivere
osculantur pedes domini.
\par {\bfseries \color{Red} Q}uo agnito discipuli
in Galileam propere
pergunt videre faciem
desideratam domini.
\par {\bfseries \color{Blue} C}laro Paschali gaudio
sol mundo nitet radio,
cum Christum iam apostoli
visu cernunt corporeo.
\par {\bfseries \color{Red} O}stensa sibi vulnera
in Christi carne fulgida
resurrexisse dominum
voce fatentur publica.
\par {\bfseries \color{Blue} R}ex Christe clementissime
tu corda nostra posside,
ut tibi laudes debitas
reddamus omni tempore.
\par {\bfseries \color{Red} Q}uesumus actor omnium
in hoc Paschali gaudio
ab omni mortis impetu
tuum defende populum.
\par {\bfseries \color{Blue} G}loria tibi domine
qui surrexisti a mortuis
cum patre et sancto spiritu,
in sempiterna secula. Amen.
}
\newline \V Gavisi sunt discipuli alleluya. \R Viso domino alleluya.
\newline {\color{Red} Ad Ben. Ant.} Post dies octo ianuis clausis ingressus dominus dixit eis pax vobis alleluya alleluya. \ul{Oratio,} Presta quesumus, \ul{Memoria de cruce.} {\color{Red} Ant.} Crucifixus surrexit a mortuis redemit nos alleluya alleluya. \V Dicite in nationibus. \ul{Oratio,} \hyperlink{deus-qui-pro-cruce}{Deus qui pro nobis.}
\newline \ul{Predicto modo in Vesperis et in Laudibus, post alias memorias que occurrerint preterquam de Beato Dominico: fiat memoria de cruce cotidie usque ad Vigiliam Ascensionis in Vesperis exclusive, nisi in Inventione Sancte Crucis et in festis totis duplicibus.}
\newline {\color{Red} Ad Primam.} \R Ihesu Christe fili dei vivi miserere nobis, * Alleluya alleluya. \V Qui surrexisti a mortuis alleluya.
\newline \ul{Iste Versus dicatur cotidie usque ad Ascensionem exclusive nisi in festo Annuntiationis quando post Pascha celebrabitur.}
\newline \V Exurge domine.
\newline {\color{Red} Ad Terciam. Cap.} \hyperlink{christus-resurgens-capitulum}{Christus resurgens.}
\begin{responsory-breve}[resurrexit-dominus-breve]
{Resurrexit dominus, * Alleluya alleluya.}{Sicut dixit vobis.}
\end{responsory-breve}%
\V Surrexit dominus vere.
\newline {\color{Red} Ad Sextam. Capitulum.}
\lettrine[lines=2]{\zallmancaps \color{Blue} S}{i} consurrexistis cum Christo que sursum sunt querite ubi Christus est in dextera dei sedens, que sursum sunt sapite non que super terram.
\begin{responsory-breve}[surrexit-dominus-vere-breve]
{Surrexit dominus vere, * Alleluya alleluya.}{Et apparuit Simoni.}
\end{responsory-breve}%
\V Surrexit dominus de sepulchro.
\newline {\color{Red} Ad Nonam. Capitulum.}
\lettrine[lines=2]{\zallmancaps \color{Red} C}{um} autem Christus apparuerit vita vestra, tunc et vos apparebitis cum ipso in gloria.
\begin{responsory-breve}[surrexit-dominus-de-breve]
{Surrexit dominus de sepulchro, * Alleluya alleluya.}{Qui pro nobis pependit in ligno.}
\end{responsory-breve}%
\V Gavisi sunt.
\newline {\color{Red} Ad Vesperas. Super psalmos. Ant.} Alleluya alleluya alleluya alleluya. {\color{Red} Ps.} \psalm{109} \psalm{110} \psalm{111} \psalm{112} \psalm{113}
\newline \ul{Predicta antiphona cum psalmis dicatur omnibus Dominicis diebus usque ad Ascensionem quando de tempore agitur. Cap., Ymnus, \Vbar , ut in Primis Vesperis.}
\newline \ul{Ordo huius diei scilicet Ymnorum, Versiculorum, Capitulorum, Responsorium ad horas servetur usque ad Ascensionem exclusive tam in Dominicis quam in ferialibus diebus quando de tempore agitur.}
\newline {\color{Red} Ad Mag. Ant.} Quia vidisti me Thoma credidisti beati qui non viderunt et crediderunt alleluya.
\newline \ul{Per ebdomadam quando de tempore agitur.}
\newline {\color{Red} Ad Ben. Ant.} Multa quidem et alia signa fecit Ihesus in conspectu discipulorum suorum alleluya, que non sunt scripta in libro hoc alleluya.
\newline {\color{Red} Ad Mag. Ant.} Hec autem scripta sunt ut credatis quia Ihesus est Christus filius dei, et ut credentes vitam habeatis in nomine ipsius alleluya.
\lettrine[lines=2]{\zallmancaps \color{Red} F}{}\ul{\textsc{eria} Secunda et deinceps usque ad secundam feriam post Dominicam,} \hyperlink{deus-omnium}{Deus omnium:} \ul{non fiant prostrationes ad Matutinas nec ad alias horas.}
\newline {\color{Red} Ad Matutinas. Invitat.} \hyperlink{alleluya-invitatorium}{Alleluya alleluya, * Alleluya.} {\color{Red} Super psalmos Ant.} Alleluya.
\newline \ul{Invitatorium et antiphona predicta dicantur omnibus ferialibus diebus usque ad Ascensionem quando de feria agitur.}
\newline {\color{Red} Ebdomada Prima post Octavas Pasche Psalmi ad Matutinas sunt hi.}
\newline {\color{Red} Feria Secunda. Ps.} \psalm{26} {\color{Red} Ps.} \psalm{27} {\color{Red} Ps.} \psalm{28}
\newline {\color{Red} Feria Tercia. Ps.} \psalm{38} {\color{Red} Ps.} \psalm{39} {\color{Red} Ps.} \psalm{40}
\newline {\color{Red} Feria Quarta. Ps.} \psalm{52} {\color{Red} Ps.} \psalm{54} {\color{Red} Ps.} \psalm{55}
\newline {\color{Red} Feria Quinta. Ps.} \psalm{68} {\color{Red} Ps.} \psalm{69} {\color{Red} Ps.} \psalm{70}
\newline {\color{Red} Feria Sexta. Ps.} \psalm{80} {\color{Red} Ps.} \psalm{81} {\color{Red} Ps.} \psalm{82}
\newline {\color{Red} Sabbato. Ps.} \psalm{97} {\color{Red} Ps.} \psalm{98} {\color{Red} Ps.} \psalm{99}
\newline {\color{Red} Ebdomada Secunda. Feria Secunda. Ps.} \psalm{29} {\color{Red} Ps.} \psalm{30} {\color{Red} Ps.} \psalm{31}
\newline {\color{Red} Feria Tercia. Ps.} \psalm{41} {\color{Red} Ps.} \psalm{43} {\color{Red} Ps.} \psalm{44}
\newline {\color{Red} Feria Quarta. Ps.} \psalm{56} {\color{Red} Ps.} \psalm{57} {\color{Red} Ps.} \psalm{58}
\newline {\color{Red} Feria Quinta. Ps.} \psalm{71} {\color{Red} Ps.} \psalm{72} {\color{Red} Ps.} \psalm{73}
\newline {\color{Red} Feria Sexta. Ps.} \psalm{83} {\color{Red} Ps.} \psalm{84} {\color{Red} Ps.} \psalm{85}
\newline {\color{Red} Sabbato. Ps.} \psalm{100} {\color{Red} Ps.} \psalm{101} {\color{Red} Ps.} \psalm{102}
\newline {\color{Red} Ebdomada Tercia. Feria Secunda. Ps.} \psalm{32} {\color{Red} Ps.} \psalm{33} {\color{Red} Ps.} \psalm{34}
\newline {\color{Red} Feria Tercia. Ps.} \psalm{45} {\color{Red} Ps.} \psalm{46} {\color{Red} Ps.} \psalm{47}
\newline {\color{Red} Feria Quarta. Ps.} \psalm{59} {\color{Red} Ps.} \psalm{60} {\color{Red} Ps.} \psalm{61}
\newline {\color{Red} Feria Quinta. Ps.} \psalm{74} {\color{Red} Ps.} \psalm{75} {\color{Red} Ps.} \psalm{76}
\newline {\color{Red} Feria Sexta. Ps.} \psalm{86} {\color{Red} Ps.} \psalm{87} {\color{Red} Ps.} \psalm{88}
\newline {\color{Red} Sabbato. Ps.} \psalm{103} {\color{Red} Ps.} \psalm{104} {\color{Red} Ps.} \psalm{105}
\newline {\color{Red} Ebdomada Quarta. Feria Secunda. Ps.} \psalm{35} {\color{Red} Ps.} \psalm{36} {\color{Red} Ps.} \psalm{37}
\newline {\color{Red} Feria Tercia. Ps.} \psalm{48} {\color{Red} Ps.} \psalm{49} {\color{Red} Ps.} \psalm{51}
\newline {\color{Red} Feria Quarta. Ps.} \psalm{63} {\color{Red} Ps.} \psalm{65} {\color{Red} Ps.} \psalm{67}
\newline {\color{Red} Feria Quinta. Ps.} \psalm{77} {\color{Red} Ps.} \psalm{78} {\color{Red} Ps.} \psalm{79}
\newline {\color{Red} Feria Sexta. Ps.} \psalm{93} {\color{Red} Ps.} \psalm{95} {\color{Red} Ps.} \psalm{96}
\newline {\color{Red} Sabbato. Ps.} \psalm{106} {\color{Red} Ps.} \psalm{107} {\color{Red} Ps.} \psalm{108}
\newline \ul{Feria secunda et tercia et quarta ebdomade quinte dicantur psalmi sicut in primis tribus feriis prime ebdomade post Octavas Pasche, quando de tempore agitur.}
\newline \ul{Feria secunda post Octavas Pasche incipitur hystoria,} \hyperlink{dignus-es}{Dignus es,} \ul{nisi festum impediat, et cantatur in hac ebdomada et duabus sequentibus quando de tempore agitur, et legitur de Apocalipsi.}
\newline {\color{Red} Lectio Prima.}
\lettrine[lines=2]{\zallmancaps \color{Blue} A}{pocalipsis} Ihesu Christi quam dedit illi Deus palam facere servis suis que oportet fieri cito, et significavit mittens per angelum suum servo suo Iohanni, qui testimonium perhibuit verbo Dei, et testimonium Ihesu Christi quecumque vidit. Beatus qui legit et qui audiunt verba prophecie huius, et servant ea que in ea scripta sunt. Tempus enim prope est. {\color{Red} Responsorium.}
\lettrine[lines=2]{\zallmancaps \color{Red} D}{ignus} \hypertarget{dignus-es}{\label{dignus-es}} es domine accipere librum et aperire signacula eius, quoniam occisus es et redemisti nos deo, * In sanguine tuo alleluya. \V Fecisti enim nos deo nostro regnum et sacerdotes. In sanguine.
\newline {\color{Red} Lectio Secunda.}
\lettrine[lines=2]{\zallmancaps \color{Blue} I}{ohannes} septem ecclesiis que sunt in Asia, gratia vobis et pax ab eo qui est, et qui erat, et qui venturus est, et a septem spiritibus qui in conspectu throni eius sunt, et a Ihesu Christo qui est testis fidelis, primogenitus mortuorum, et princeps regum terre. Qui dilexit nos, et lavit nos a peccatis nostris in sanguine suo, et fecit nos regnum et sacerdotes Deo et Patri suo. Ipsi gloria et imperium in secula seculorum. Amen.
\begin{responsory}[ego-sicut]
{Ego sicut vitis fructificavi suavitatem odoris alleluya, * Transite ad me omnes qui concupiscitis me, et a generationibus meis adimplemini alleluya alleluya.}{Ego diligentes me diligo, et qui mane vigilaverint ad me invenient me.}
\end{responsory}%
\newline {\color{Red} Lectio Tertia.}
\lettrine[lines=2]{\zallmancaps \color{Red} E}{cce} venit cum nubibus, et videbit eum omnis oculus, et qui eum pupugerunt. Et plangent se super eum omnes tribus terre. eciam amen. Ego sum alpha et \omega , principium et finis dicit Dominus Deus: qui est, et qui erat, et qui venturus est omnipotens.
\begin{responsory-final}[audivi-vocem-tanquam]
{Audivi vocem in celo tanquam vocem tonitrui magni alleluya, regnabit deus noster in eternum alleluya, * Quia facta est salus et virtus et potetas Christi eius, * Alleluya alleluya.}{Et vox de throno exivit dicens, laudem dicite deo nostro omnes sancti eius, et qui timetis deum pusilli et magni.}{Alleluya}
\end{responsory-final}%
\newline {\color{Red} In Laudibus. Ant.} Alleluya alleluya. \ul{Hec antiphona dicatur in Laudibus usque ad Ascensionem tam in Dominicis diebus quam in ferialibus quando de tempore agitur.} {\color{Red} Ps.} \psalm{92} \ul{et ceteri. Et dicantur per totum tempus Paschale.}
\newline {\color{Red} Ad Primam. Ant.} Cognoverunt dominum alleluya, in fractione panis alleluya.
\newline {\color{Red} Ad Terciam. Ant.} Pax vobis ego sum alleluya nolite timere alleluya.
\newline {\color{Red} Ad Sextam. Ant.} Ecce ego vobiscum sum alleluya omnibus diebus alleluya.
\newline {\color{Red} Ad Nonam. Ant.} Noli flere Maria alleluya, resurrexit dominus alleluya.
\newline \ul{Predicte antiphone dicantur tam in Dominicis quam in ferialibus diebus ad horas usque ad Ascensionem quando de tempore agitur.}
\newline {\color{Red} Ad Vesperas. Super psalmos. Ant.} Alleluya.
\newline \ul{Predicta antiphona dicatur super psalmos ad Vesperas, in ferialibus diebus usque ad Ascensionem quando de tempore agitur.}
\newline {\color{Red} Feria Tercia. Lectio Prima.}
\lettrine[lines=2]{\zallmancaps \color{Red} E}{go} Iohannes frater vester et particeps in tribulatione et regno et patientia in Ihesu, fui in insula que appellatur Pathmos propter verbum Dei et testimonium Ihesu, fui in spiritu in dominica die. Et audivi post me vocem magnam tanquam tube dicentis. Quod vides scribe in libro, et mitte septem ecclesiis, Epheso et Smirne, et Pergamo et Thyatyre et Sardis et Phyladelphie et Laodicie. \R
\lettrine[lines=2]{\zallmancaps \color{Blue} L}{ocutus} \hypertarget{locutus-est-pasche}{\label{locutus-est-pasche}} est ad me unus de septem angelis dicens, veni ostendam tibi novam nuptam sponsam agni, * Et vidi Hierusalem descendentem de celo ornatam monilibus suis alleluya alleluya alleluya. \V Et sustulit me in spiritu, in montem magnum et altum. Et vidi.
\newline {\color{Red} Lectio Secunda.}
\lettrine[lines=2]{\zallmancaps \color{Red} E}{t} conversus sum ut viderem vocem que loquebatur mecum. Et conversus vidi septem candelabra aurea, et in medio septem candelabrorum aureorum similem filio hominis vestitum podere, et precinctum ad mammillas zona aurea.
\begin{responsory}[audivi-vocem-angelorum]
{Audivi vocem in celo angelorum multorum dicentium, * Timete dominum et date claritatem illi, et adorate eum qui fecit celum et terram, mare et fontes aquarum alleluya alleluya.}{Vidi angelum dei volantem per medium celi, voce magna clamantem et dicentem.}
\end{responsory}%
\newline {\color{Red} Lectio Tertia.}
\lettrine[lines=2]{\zallmancaps \color{Blue} C}{aput} autem eius et capilli erant candidi tanquam lana alba, et tanquam nix, et oculi eius tanquam flamma ignis: et pedes eius similes auricalco sicut in camino ardenti. Et vox illius tanquam vox aquarum multarum. Et habebat in dextera sua stellas septem, et de ore eius gladius utraque parte acutus exibat, et facies eius sicut sol lucet in virtute sua.
\begin{responsory-final}[ostendit-michi]
{Ostendit michi angelus fontem aque vive, et dixit ad me alleluya, * Hic deum adora, * Alleluya alleluya alleluya.}{Postquam audissem et vidissem cecidi ut adorarem ante pedes angeli qui michi hec ostendebat et dixit michi.}{Alleluya}
\end{responsory-final}%
\newline \ul{Sex premisse lectiones possunt resumi in aliis feriis. Vel potest legi a quarto capitulo usque ad sextum ad libitum.}
\newline {\color{Red} Feria Quarta. Responsorium.}
\lettrine[lines=2]{\zallmancaps \color{Red} V}{idi} \hypertarget{vidi-portam}{\label{vidi-portam}} portam civitatis ad orientem positam, et apostolorum nomina et agni super eam scripta, * Et super muros eius angelorum custodiam, alleluya. \V Vidi civitatem sanctam Hierusalem descendentem de celo, ornatam tanquam sponsam viro suo. Et super.
\begin{responsory}[vidi-hierusalem]
{Vidi Hierusalem descendentem de celo ornatam auro mundo, * Et lapidibus preciosis intextam, alleluya alleluya.}{Ab intus in fimbriis aureis circumamicta varietate.}%typo: circuamicta
\end{responsory}%
\begin{responsory-final}[in-diademate]
{In diademate capitis Aaron lapides preciosi fulgebant, * Dum perficeretur opus dei, * Alleluya alleluya alleluya}{Corona aurea super caput eius expressa signo sanctitatis.}{Alleluya}
\end{responsory-final}%
{\ubsubsection{Dominica Prima post Octavas Pasche Sabbato Precedenti}{dominica-prima-post-octavas-pasche-sabbato-precedenti-breviarium}}
\newline {\color{Red} Ad Magnificat. Antiphona.} Hec autem. {\color{Red} Oratio.}
\lettrine[lines=2]{\zallmancaps \color{Blue} D}{eus} qui in filii tui humilitate iacentem mundum erexisti, fidelibus tuis perpetuam concede leticiam, ut quos perpetue mortis eripuisti casibus, gaudiis facias sempiternis perfrui. Per eundem.
\newline {\color{Red} Ad Matutinas. Super psalmos. Ant.} Surrexit Christus et illuxit populo quo quem redemit sanguine suo alleluya. {\color{Red} Ps.} \psalm{1} {\color{Red} Ps.} \psalm{2} {\color{Red} Ps.} \psalm{3}
\newline \ul{Predicta antiphona cum suis psalmis dicatur in Dominicis diebus usque ad Ascensionem quando de tempore agitur.}
\newline {\color{Red} Lectio Prima. Secundum Iohannem.}
\lettrine[lines=2]{\zallmancaps \color{Red} I}{n} illo tempore: Dixit Ihesus discipulis suis. Ego sum pastor bonus. Bonus pastor animam suam dat pro ovibus suis. Et reliqua.
\newline {\color{Red} Omelia Beati Gregorii Pape.}
\lettrine[lines=2]{\zallmancaps \color{Blue} A}{udistis} fratres karissimi ex lectione evangelica eruditionem vestram, audistis et periculum nostrum. Ecce enim is qui non ex accidenti dono sed essentialiter bonus est, dicit. Ego sum pastor bonus. Atque eiusdem bonitatis formam quam nos imitemur, adiungit, dicens. Bonus pastor animam suam ponit pro ovibus suis. \R \hyperlink{dignus-es}{Dignus.}
\newline {\color{Red} Lectio Secunda.}
\lettrine[lines=2]{\zallmancaps \color{Red} F}{ecit} quod monuit: ostendit quod iussit. Bonus pastor animam suam pro ovibus suis posuit, ut in sacramento nostro corpus suum et sanguinem verteret, et oves quas redemerat carnis sue alimento satiaret. Ostensa nobis est de contemptu mortis, via quam sequamur, apposita forma cui imprimamur. \R \hyperlink{ego-sicut}{Ego sicut.}
\newline {\color{Red} Lectio Tertia.}
\lettrine[lines=2]{\zallmancaps \color{Blue} P}{rimum} nobis est exteriora nostra misericorditer ovibus impendere, postremum vero si necesse sit, eciam mortem
%pp107
nostram pro eisdem ovibus ministrare. A primo autem hoc minimo pervenitur ad postremum maius. Sed cum incomparabiliter longe sit melior anima qua vivimus, terrena substantia quam exterius possidemus: qui non dat pro ovibus substantiam suam, quando pro his daturus est animam suam? \R \hyperlink{audivi-vocem-tanquam}{Audivi.}
\newline {\color{Red} Ad Ben. Ant.} Ego sum pastor bonus qui pasco oves meas, et pro ovibus meis pono animam meam alleluya.
\newline {\color{Red} Ad Mag. Ant.} Ego sum pastor ovium, ego sum via et veritas, ego sum pastor bonus, et cognosco meas et cognoscunt me mee alleluya alleluya.
\newline \ul{Per ebdomadam quando de tempore agitur.}
\newline {\color{Red} Ad Ben. Ant.} Pastor bonus animam suam ponit pro ovibus suis alleluya.
\newline {\color{Red} Ad Mag. Ant.} Alias oves habeo que non sunt ex hoc ovili et illas oportet me adducere, et vocem meam audient, et fiet unum ovile et unus pastor alleluya.
\newline {\color{Red} Feria Secunda. Lectio Prima.}
\lettrine[lines=2]{\zallmancaps \color{Red} E}{t} vidi quod aperuisset agnus unum de septem signaculis, et audivi unum de quatuor animalibus dicens tanquam vocem tonitrui. Veni, et vide. Et vidi, et ecce equus albus, et qui sedebat super illum habebat arcum. Et data est ei corona: et exivit vincens ut vinceret.
\newline {\color{Red} Lectio Secunda.}
\lettrine[lines=2]{\zallmancaps \color{Blue} E}{t} cum aperuisset sigillum secundum, audivi secundum animal dicens. Veni, et vide. Et exivit alius equus rufus, et qui sedebat super eum datum est illi ut sumeret pacem de terra, et ut invicem se interficiant. Et datus est illi gladius magnus.
\newline {\color{Red} Lectio Tertia.}
\lettrine[lines=2]{\zallmancaps \color{Red} E}{t} cum aperuisset sigillum tercium, audivi tercium animal dicens, Veni, et vide. Et ecce equus niger, et qui sedebat super eum habebat stateram in manu sua. Et audivi tanquam vocem in medio quatuor animalium dicentium. Bilibris tritici denario, et tres bilibres ordei denario: et vinum et oleum. Ne leseris.
\newline {\color{Red} Feria Tercia. Lectio Prima.}
\lettrine[lines=2]{\zallmancaps \color{Blue} E}{t} cum aperuisset sigillum quartum, audivi vocem quarti animalis dicentis. Veni, et vide. Et ecce equus pallidus, et qui sedebat super eum nomen illi mors, et infernus sequebatur eum. Et data est illi potestas super quatuor partes terre, interficere gladio, fame, et morte, et bestiis terre.
\newline {\color{Red} Lectio Secunda.}
\lettrine[lines=2]{\zallmancaps \color{Red} E}{t} cum aperuisset sigillum quintum, vidi subtus altare animas interfectorum propter verbum Dei, et propter testimonium quod habebant. Et clamabant voce magna dicentes. Usquequo Domine sanctus et verus non iudicas et vindicas sanguinem nostrum de his qui habitant in terra?
\newline {\color{Red} Lectio Tertia.}
\lettrine[lines=2]{\zallmancaps \color{Blue} E}{t} date sunt illis singule stole albe, et dictum est illis ut requiescerent tempus adhuc modicum, donec compleantur conservi eorum, et fratres eorum, qui interficiendi sunt sicut et illi.
\newline \ul{Lectiones alie feriales a sexto capitulo usque ad decimum octavum ad libitum.}
{\ubsubsection{Dominica Secunda post Octavas Pasche Sabbato Precedenti}{dominica-secunda-post-octavas-pasche-sabbato-precedenti-breviarium}}
\newline {\color{Red} Ad Mag. Ant.} Alias. {\color{Red} Oratio.}
\lettrine[lines=2]{\zallmancaps \color{Red} D}{eus} qui errantibus ut in viam possint redire iustitie veritatis tue lumen ostendis, da cunctis qui Christiana professione censentur, et illa respuere que huic inimica sunt nomini, et ea que sunt apta sectari. Per.
\newline {\color{Red} Lectio Prima. Secundum Iohannem.}
\lettrine[lines=2]{\zallmancaps \color{Blue} I}{n} illo tempore: Dixit Ihesus discipulis suis. Modicum et iam non videbitis me, et iterum modicum et videbitis me, quia vado ad patrem. Et reliqua.
\newline {\color{Red} Omelia Beati Augustini Episcopi.}
\lettrine[lines=2]{\zallmancaps \color{Red} H}{ec} Domini verba ubi ait, modicum et iam non videbitis me, et iterum modicum et videbitis me, quia vado ad patrem, ita erant obscura discipulis antequam id quod dicit esset impletum, ut querentes inter se quid esset quod diceret omnino se nescire faterentur. \ul{Responsoria que in precedenti Dominica.}
\newline {\color{Red} Lectio Secunda.}
\lettrine[lines=2]{\zallmancaps \color{Blue} H}{oc} est enim quod eos movebat, quia dixit modicum et non videbitis me, et iterum modicum et videbitis me. Nam in precedentibus, quia non dixerat modicum, sed dixerat ad patrem vado, et iam non videbitis me, tanquam aperte illis visus est locutus, nec inter se de hoc aliquid quesierunt.
\newline {\color{Red} Lectio Tertia.}
\lettrine[lines=2]{\zallmancaps \color{Red} N}{unc} ergo quod illis tunc obscurum fuit, et mox manifestatum est, iam nobis utique manifestum est. Post paululum enim passus est, et non viderunt eum: rursus post paululum resurrexit, et viderunt eum.
\newline {\color{Red} Ad Ben. Ant.} Modicum et non videbitis me dicit dominus, iterum modicum et videbitis me quia vado ad patrem alleluya alleluya.
\newline {\color{Red} Ad Mag. Ant.} Iterum autem videbo vos et gaudebit cor vestrum alleluya, et gaudium vestrum nemo tollet a vobis alleluya.
\newline \ul{Per ebdomadam quando de tempore agitur.}
\newline {\color{Red} Ad Ben. Ant.} Amen amen dico vobis quia plorabitis et flebitis vos mundus autem gaudebit, vos autem contristabimini, sed tristicia vestra vertetur in gaudium alleluya.
\newline {\color{Red} Ad Mag. Ant.} Tristicia implevit cor vestrum et gaudium vestrum nemo tollet a vobis alleluya alleluya.
\newline {\color{Red} Feria Secunda. Lectio Prima.}
\lettrine[lines=2]{\zallmancaps \color{Blue} E}{t} vidi angelum descendentem de celo, habentem potestatem magnam, et terra illuminata est a gloria eius, et exclamavit in forti voce dicens. Cecidit, cecidit Babilon magna, et facta est habitatio demoniorum et custodia omnis spiritus immundi, et custodia omnis volucris immunde et odibilis, quia de ira fornicationis eius biberunt omnes gentes, et reges terre cum illa fornicati sunt, et mercatores terre de virtute deliciarum eius divites facti sunt.
\newline {\color{Red} Lectio Secunda.}
\lettrine[lines=2]{\zallmancaps \color{Red} E}{t} audivi aliam vocem de celo dicentem. Exite de illa populus meus, et ne participes sitis delictorum eius, et de plagis eius non accipiatis, quoniam pervenerunt peccata eius usque ad celum, et recordatus est Dominus iniquitatum eius.
\newline {\color{Red} Lectio Tertia.}
\lettrine[lines=2]{\zallmancaps \color{Blue} R}{eddite} illi sicut ipsa reddidit vobis, et duplicate duplicia secundum opera eius. In poculo quo miscuit vobis, miscete illi duplum. Quantum glorificavit se, et in deliciis fuit, tantum date illi tormentum et luctum, quia in corde suo dicit. Sedeo regina, et vidua non sum, et luctum non videbo. Ideo in una die venient plage eius, mors et luctus et fames, et igni comburetur, quia fortis est Deus qui iudicabit illam.
\newline {\color{Red} Feria Tercia. Lectio Prima.}
\lettrine[lines=2]{\zallmancaps \color{Red} E}{t} flebunt et plangent se super illam reges terre qui cum illa fornicati sunt et in deliciis vixerunt, cum viderint fumum incendii eius, longe stantes propter timorem tormentorum eius dicentes. Ve ve civitas illa magna Babilon, civitas illa fortis, quoniam una hora venit iudicium tuum.
\newline {\color{Red} Lectio Secunda.}
\lettrine[lines=2]{\zallmancaps \color{Blue} E}{t} negotiatores terre flebunt, et lugebunt super illam, quoniam merces eorum nemo emet amplius, merces auri, et argenti, lapidis preciosi et margariti, et bissi, et purpure, et serici, et cocci, et omne lignum thyinum, et omnia vasa eboris, et omnia vasa de lapide precioso, et eramento et ferro et marmore, et cynamomum, et amomum, odoramentorum et unguenti, et thuris et vini, et olei, et simile, et tritici, et iumentorum et ovium et equorum, et redarum, et mancipiorum et animarum hominum.
\newline {\color{Red} Lectio Tertia.}
\lettrine[lines=2]{\zallmancaps \color{Red} E}{t} poma desiderii anime tue discesserunt a te, et omnia pinguia et preclara perierunt a te. Et amplius illa non invenient mercatores horum. Qui divites facti sunt ab ea longe stabunt propter timorem tormentorum eius, flentes ac lugentes ac dicentes. Ve ve civitas illa magna que amicta erat bysso et purpura et cocco, et deaurata est auro et lapide precioso et margaritis, quoniam una hora destitute sunt tante divitie.
\newline \ul{Lectiones alie feriales a decimonono capitulo usque ad finem ad libitum.}
{\ubsubsection{Dominica Tercia post Octavas Pasche Sabbato Precedenti}{dominica-tercia-post-octavas-pasche-sabbato-precedenti-breviarium}}
\newline {\color{Red} Ad Vesperas.} \R \hyperlink{viderunt-te}{Viderunt. II.} {\color{Red} Ad Magnificat. Antiphona.} Tristicia. {\color{Red} Oratio.}
\lettrine[lines=2]{\zallmancaps \color{Blue} D}{eus} qui fidelium mentes unius efficis voluntatis, da populis tuis id amare quod precipis, id desiderare quod promittis ut inter mundanas varietates ibi nostra fixa sint corda, ubi vera sunt gaudia. Per.
\newline {\color{Red} Ad Matutinas. Lectio Prima. Secundum Iohannem.}
\lettrine[lines=2]{\zallmancaps \color{Red} I}{n} illo tempore: Dixit Ihesus discipulis suis. Vado ad eum qui me misit, et nemo ex vobis interrogat me quo vadis. Et reliqua.
\newline {\color{Red} Omelia Venerabilis Bede Presbiteri.}
\lettrine[lines=2]{\zallmancaps \color{Blue} V}{ado} inquit dominus ad eum qui me misit: et nemo ex vobis interrogat me quo vadis. Ac si aperte dicat. Revertar ascendendo ad eum qui me incarnari constituit, et tanta tamque manifesta erit eiusdem ascensionis claritas, ut nemini vestrum opus sit interrogare quo vadam, videntibus cunctis quia ad celos pergam. {\color{Red} Responsorium.}
\lettrine[lines=2]{\zallmancaps \color{Red} S}{i} \hypertarget{si-oblitus}{\label{si-oblitus}} oblitus fuero tui alleluya, obliviscatur me dextera mea, * Adhereat lingua mea faucibus meis si non meminero tui alleluya alleluya. \V Super flumina Babilonis illic sedimus et flevimus dum recordaremur tui Sion. Adhereat.
\newline {\color{Red} Lectio Secunda.}
\lettrine[lines=2]{\zallmancaps \color{Blue} B}{ene} autem cum de ascensione sua dixisset. vado ad eum qui me misit: addidit. Et nemo ex vobis interrogat me. quo vadis. Superius namque cum de sua passione protestaretur dicens, quo ego vado vos non potestis venire, interrogavit eum Petrus et ait. Domine quo vadis? Responsumque est ei. Quo ego vado non potes me modo sequi, sequeris autem postea. Quia nimirum passionis mortisque eius mysterium necdum intelligere: necdum poterant imitari. Maiestatem vero ascensionis statim ut viderunt cognoverunt, totisque animi votis, ut sequi mererentur optabant.
\begin{responsory}[viderunt-te]
{Viderunt te aque deus, viderunt te aque et timuerunt, * Multitudo sonitus aquarum vocem dederunt nubes alleluya alleluya alleluya.}{Illuxerunt coruscationes tue orbi terre, commota est et contremuit terra.}
\end{responsory}%
\newline {\color{Red} Lectio Tertia.}
\lettrine[lines=2]{\zallmancaps \color{Red} S}{ed} quia hec locutus sum vobis, tristicia implevit cor vestrum. Sciebat ipse Dominus quid hec sua verba in discipulorum cordibus agerent, quia videlicet tristiciam magis de abscessu quo eos desereret, quam de ascensu quo Patrem peteret, leticiam generarent. Unde benigne consolando subiunxit. Sed ego veritatem dico vobis. Expedit vobis ut ego vadam. Expedit ut forma servi vestris subtrahatur aspectibus, quatenus amor divinitatis artius vestris infigatur mentibus.
\begin{responsory-final}[narrabo-nomen]
{Narrabo nomen tuum fratribus meis, alleluya in medio ecclesie, * Laudabo te, * Alleluya alleluya.}{Confitebor tibi in ecclesia magna, in populo gravi.}{Alleluya}
\end{responsory-final}%
\newline {\color{Red} Ad Ben. Ant.} Vado ad eum qui misit me, sed quia hec locutus sum vobis tristicia implevit cor vestrum alleluya.
\newline {\color{Red} Ad Mag. Ant.} Ego veritatem dico vobis, expedit vobis ut ego vadam, si enim non abiero Paraclitus non veniet alleluya.
\newline \ul{Per ebdomadam quando de tempore agitur.}
\newline {\color{Red} Ad Ben. Ant.} Adhuc multa habeo vobis dicere sed non potestis portare modo, cum autem venerit ille spiritus veritatis docebit vos omnem veritatem, alleluya.
\newline {\color{Red} Ad Mag. Ant.} Ille me clarificabit quia de meo accipiet et annuntiabit vobis alleluya.
\newline {\color{Red} Feria Secunda. Lectio Prima.}
\lettrine[lines=2]{\zallmancaps \color{Blue} I}{acobus} Dei et Domini nostri Ihesu Christi servus, duodecim tribubus que sunt in dispersione, salutem. Omne gaudium existimate fratres mei, cum in temptationes varias incideritis, scientes quod probatio fidei vestre patientiam operatur. Patientia autem opus perfectum habeat, ut sitis perfecti et integri in nullo deficientes. \R \hyperlink{si-oblitus}{Si oblitus.}
\newline {\color{Red} Lectio Secunda.}
\lettrine[lines=2]{\zallmancaps \color{Red} S}{iquis} autem vestrum indiget sapientia postulet a Deo, qui dat omnibus affluenter et non improperat: et dabitur ei. Postulet autem in fide nichil hesitans. Qui enim hesitat, similis est fluctui maris qui a vento movetur et circumfertur. Non ergo estimet homo ille, quod accipiat aliquid a Domino. Vir duplex animo: inconstans est in omnibus viis suis. \R \hyperlink{viderunt-te}{Viderunt.}
\newline {\color{Red} Lectio Tertia.}
\lettrine[lines=2]{\zallmancaps \color{Blue} G}{lorietur} autem frater humilis in exaltatione sua, dives autem in humilitate sua, quoniam sicut flos feni transibit. Exortus est enim sol cum ardore, et arefecit fenum, et flos eius decidit, et decor vultus eius deperiit. Ita et dives in itineribus suis marcescet. \R \hyperlink{narrabo-nomen}{Narrabo.}
\newline {\color{Red} Feria Tercia. Lectio Prima.}
\lettrine[lines=2]{\zallmancaps \color{Red} B}{eatus} vir qui suffert temptationem, quoniam cum probatus fuerit, accipiet coronam vite, quam repromisit Deus diligentibus se. Nemo cum temptatur, dicat quoniam a Deo temptor. Deus enim intemptator malorum est. Ipse autem neminem temptat. Unusquisque vero temptatur a concupiscentia sua, abstractus et illectus. Deinde concupiscentia cum conceperit parit peccatum. Peccatum vero cum consummatum fuerit, generat mortem. \R
\lettrine[lines=2]{\zallmancaps \color{Blue} I}{n} \hypertarget{in-toto-corde}{\label{in-toto-corde}} toto corde meo alleluya, exquisivi te alleluya, * Ne repellas me a mandatis tuis alleluya alleluya. \V Benedictus es tu domine doce me iustificationes tuas. Ne repellas.
\newline {\color{Red} Lectio Secunda.}
\lettrine[lines=2]{\zallmancaps \color{Red} N}{olite} itaque errare fratres mei dilectissimi. Omne datum optimum et omne donum perfectum desursum est, descendens a Patre luminum, apud quem non est transmutatio, nec vicissitudinis obumbratio. Voluntarie enim genuit nos verbo veritatis, ut simus initium aliquod creature eius.
\begin{responsory}[in-ecclesiis]
{In ecclesiis benedicite deo domino de fontibus Israhel, * Alleluya alleluya.}{Cantate domino canticum novum cantate domino omnis terra.}
\end{responsory}%
\newline {\color{Red} Lectio Tertia.}
\lettrine[lines=2]{\zallmancaps \color{Blue} S}{citis} fratres mei dilectissimi. Sit autem omnis homo velox ad audiendum, tardus ad loquendum: et tardus ad iram. Ira enim viri iusticiam Dei non operatur. Propter quod abicientes omnem immundiciam et abundantiam malitie, in mansuetudine suscipite insitum verbum, quod potest salvare animas vestras.
\begin{responsory-final}[ymnum-cantate]
{Ymnum cantate nobis alleluya, * Quomodo cantabimus canticum domini in terra aliena, * Alleluya alleluya.}{Illic interrogaverunt nos qui captivos duxerunt nos verba cantionum.}{Alleluya}
\end{responsory-final}%
\newline \ul{Lectiones alie feriales usque in finem canonicarum ad libitum.}
\newline {\color{Red} Feria Quarta. Responsorium.}
\lettrine[lines=2]{\zallmancaps \color{Red} D}{eus} \hypertarget{deus-canticum}{\label{deus-canticum}} canticum novum cantabo tibi alleluya, * In psalterio decem cordarum psallam tibi alleluya alleluya. \V Deus meus es tu et confitebor tibi, deus meus es tu et exaltabo te. In psalterio.
\begin{responsory}[alleluya-audivimus]
{Alleluya audivimus eam in Effrata invenimus eam in campis silve, * Introibimus in tabernaculum eius, adorabimus in loco ubi steterunt pedes eius alleluya alleluya.}{Surge domine in requiem tuam, tu et archa sanctificationis tue.}
\end{responsory}%
\begin{responsory-final}[deduc-me]
{Deduc me in semita mandatorum tuorum alleluya, quoniam ipsam volui alleluya, * Inclina cor meum in testimonia tua, * Alleluya alleluya alleluya.}{Averte oculos meos ne videant vanitatem in via tua vivifica me.}{Alleluya}
\end{responsory-final}%
{\ubsubsection{Dominica Quarta post Octavas Pasche Sabbato Precedenti}{dominica-quarta-post-octavas-pasche-sabbato-precedenti-breviarium}}
\newline {\color{Red} Ad Mag. Ant.} Ille me clarificabit. {\color{Red} Oratio.}
\lettrine[lines=2]{\zallmancaps \color{Blue} D}{eus} a quo bona cuncta procedunt largire supplicibus, ut cogitemus te inspirante que recta sunt et te gubernante eadem faciamus. Per.
\newline {\color{Red} Ad Matutinas. Lectio Prima. Secundum Iohannem.}
\lettrine[lines=2]{\zallmancaps \color{Red} I}{n} illo tempore: Dixit Ihesus discipulis suis. Amen amen dico vobis, si quid pecieritis patrem in nomine meo, dabit vobis. Et reliqua.
\newline {\color{Red} Omelia Venerabilis Bede Presbiteri.} \grecross
\lettrine[lines=2]{\zallmancaps \color{Blue} P}{otest} movere infirmos auditores, quomodo in capite lectionis huius evangelice discipulis salvator promittat, si quid inquiens pecieritis Patrem in nomine meo dabit vobis: cum non solum nostri similes multa que Patrem in Christi nomine videntur petere non accipiant, verum eciam ipse apostolus Paulus tercio Dominum rogaverit ut a se angelus Sathane a quo tribulabatur abscederet, nec impetrare potuerit. \ul{Responsoria que in precedenti Dominica.}
\newline {\color{Red} Lectio Secunda.}
\lettrine[lines=2]{\zallmancaps \color{Red} S}{ed} huius motio questionis antiqua iam Patrum explanatione reserata est, qui veraciter intellexerunt illos solum in nomine Salvatoris petere, qui ea que ad perpetuam salutem pertinent, petunt, ideoque apostolum non in nomine salvatoris petisse ut temptatione careret quam ob custodiam humilitatis acceperat, quia si hac caruisset salvus esse non posset.
\newline {\color{Red} Lectio Tertia.}
\lettrine[lines=2]{\zallmancaps \color{Blue} Q}{uotiescunque} ergo petentes non exaudimur, ideo fit: quia vel contra auxilium nostre salutis petimus, ac propterea a misericorde Patre beneficii gratia nobis quod inepte petimus negatur, vel utilia quidem et que ad veram salutem respiciant petimus, sed ipsi male vivendo auditum a nobis iusti Iudicis avertimus, vel dum pro peccantibus quibusquam ut resipiscant oramus, et si salubriter petimus, atque ex nostro merito digni sumus auditu, ipsorum tamen perversitas ne impetremus: obsistit.
\newline {\color{Red} Ad Ben. Ant.} Usque modo non petistis quicquam in nomine meo, petite et accipietis alleluya.
\newline {\color{Red} Ad Mag. Ant.} Petite et accipietis ut gaudium vestrum plenum sit, ipse enim pater amat vos quia vos me amastis et credidistis alleluya.
\newline {\color{Red} Feria Secunda in Rogationibus. Ad Mat. Lectio Prima. Secundum Lucam.}
\lettrine[lines=2]{\zallmancaps \color{Red} I}{n} illo tempore: Dixit Ihesus discipulis suis. Quis vestrum habebit amicum, et ibit ad illum media nocte, et dicet illi. Amice commoda michi tres panes quoniam amicus meus veniet de via ad me, et non habeo quod ponam ante illum. Et reliqua.
\newline {\color{Red} Omelia Venerabilis Bede Presbiteri.}
\lettrine[lines=2]{\zallmancaps \color{Blue} R}{ogatus} a discipulis salvator non modo formam orationis, sed et instantiam frequentiamque tradit orandi. Amicus ergo ad quem media nocte venitur, ipse Deus intelligitur, cui in media tribulatione supplicare, et tres panes, id est intelligentiam Trinitatis qua presentis vite consolentur labores: efflagitare debemus. \R \hyperlink{si-oblitus}{Si oblitus.}
\newline {\color{Red} Lectio Secunda.}
\lettrine[lines=2]{\zallmancaps \color{Red} A}{micus} qui venit de via, ipse noster est animus, qui toties a nobis recedit, quoties ad appetenda terrena et temporalia foris vagatur. Redit ergo, celestique alimonia refici desiderat, cum in se reversus superna ceperit ac spiritalia meditari. De quo pulchre qui petierat adiungit, non se habere quod ponat ante illum, quoniam anime post seculi tenebras Deum suspiranti: nil preter eum cogitare, nil loqui, nil libet intueri. Solum quod recognovit summe Trinitatis gaudium contemplari, atque ad hoc plenius intuendum pervenire satagit. \R \hyperlink{viderunt-te}{Viderunt.}
\newline {\color{Red} Lectio Tertia.}
\lettrine[lines=2]{\zallmancaps \color{Blue} S}{equitur.} Et ille deintus respondens dicat. Noli
%pp108
michi molestus esse, iam ostium clausum est, et pueri mei mecum sunt in cubili, non possum surgere, et dare tibi. Ostium amici, divini est intelligentia sermonis, quod sibi Apostolus orat aperiri ad loquendum mysterium Christi. Quod clausum est tempore famis verbi, cum intelligentia non datur. \R \hyperlink{narrabo-nomen}{Narrabo.}
\newline {\color{Red} Ad Ben. Ant.} Petite et dabitur vobis, querite et invenietis pulsate et aperietur vobis alleluya. {\color{Red} Oratio.}
\lettrine[lines=2]{\zallmancaps \color{Blue} P}{resta} quesumus omnipotens deus: ut qui in afflictione nostra de tua pietate confidimus, contra adversa omnia tua semper protectione muniamur. Per.
\newline \ul{Hec oratio in hac die et sequenti ad Laudes dicantur, ad alias horas et ad Vesperas, oratio Dominicalis. Si festum simplex vel maius in hac feria secunda evenerit, predicta omelia legatur in tercia feria sequenti nisi aliud festum immediate occurrerit tunc enim omittitur.}
\newline {\color{Red} Ad Mag. Ant.} Omnis qui petit accipit, et qui querit invenit et pulsanti aperietur alleluya.
\newline {\color{Red} Feria Tercia. Lectio Prima.}
\newline \ul{Avitus archiepiscopus Viennensis quem Gingnadius nominat inter viros illustres, et fuit successor beati Mammerti qui tempore Clodovei regis invenit observantiam rogationum.} \grecross %Gennadius
% https://la.wikisource.org/wiki/Homilia_de_rogationibus_(Avitus_Viennensis)
\lettrine[lines=2]{\zallmancaps \color{Red} T}{emporibus} beati Mammerti archiepiscopi Viennensis apud Viennam contigerunt incendia crebra, terremotus assidui, nocturni sonitus, populosis hominum conventibus fere silvestres se ingerebant, in tantum ut pavidi cervi per angusta portarum usque ad fori loca penetrarent. Tracta sunt hec usque ad noctem vigiliarum dominice Resurrectionis. In qua cum sperarent homines hec finem habere, edes publica quam precelsa in civitatis vertice sublimitas in immensum pretulerat: flammis terribilibus conflagrare crepusculo cepit. \R \hyperlink{in-toto-corde}{In toto corde.}
\newline {\color{Red} Lectio Secunda.}
\lettrine[lines=2]{\zallmancaps \color{Blue} I}{nterpolatur} sollemnitas, pleno timoribus populo ecclesia vacuatur. Perstitit autem invictus antistes festivis altaribus, et calore fidei succensus: flumine lacrimarum permissam ignibus potestatem compescuit, et tunc concepit animo: quicquid hodie psalmis ac precibus mundus inclamat. \R \hyperlink{in-ecclesiis}{In ecclesiis.}
\newline {\color{Red} Lectio Tertia.}
\lettrine[lines=2]{\zallmancaps \color{Red} E}{voluta} autem sollemnitate, secreta collatione tractavit quomodo aut quando fieri deberet quod animo conceperat. Eligitur tempus inter Ascensionem Domini et Dominicam precedentem, et cum maxima compunctione et devotione et multitudine populi huiusmodi observantia inchoatur. Secute sunt tandem quedam ecclesie Galliarum rem tam commendabilis exempli, licet non eisdem temporibus. Tandem crescente sacerdotum concordia, ad idem tempus observantia est redacta. \R \hyperlink{ymnum-cantate}{Ymnum cantate.}
\newline \ul{Antiphona ad Ben. et oratio sicut in precedenti feria secunda.}
\newline {\color{Red} Ad Mag. Ant.} Omnis qui petit. \ul{Ut supra in secunda feria.}
{\ubsubsection{In Vigilia Ascensionis}{in-vigilia-ascensionis-breviarium}}
\newline {\color{Red} Ad Mat.} \ul{Invitatorium, Ymnus, \Vbar , Ut supra.}
\newline {\color{Red} Lectio Prima. Secundum Iohannem.}
\lettrine[lines=2]{\zallmancaps \color{Blue} I}{n} illo tempore: Sublevatus Ihesus oculis in celum: dixit. Pater venit hora: clarifica filium tuum ut filius tuus clarificet te. Et reliqua.
\newline {\color{Red} Omelia Beati Augustini Episcopi.}
\lettrine[lines=2]{\zallmancaps \color{Red} P}{oterat} Dominus noster unigenitus et coeternus Patri in forma servi et ex forma servi si hoc opus esset orare silentio, sed ita se Patri exhibere voluit precatorem, ut meminisset se esse nostrum doctorem. Proinde quam fecit orationem pro nobis, notam fecit et nobis: quoniam tanti magistri nostri non solum apud ipsos sermocinatio, sed eciam pro ipsis ad Patrem oratio discipulorum est edificatio. \R \hyperlink{deus-canticum}{Deus canticum.}
\newline {\color{Red} Lectio Secunda.}
\lettrine[lines=2]{\zallmancaps \color{Blue} I}{n} hoc autem quod dicit venit hora, ostendit omne tempus, et quid quando faceret vel fieri sineret, ab illo esse dispositum qui tempori subditus non est, quoniam que futura fuerant per singula tempora, in Dei sapientia causas efficientes habent in qua nulla sunt tempora. \R \hyperlink{alleluya-audivimus}{Alleluya audivimus.}
\newline {\color{Red} Lectio Tertia.}
\lettrine[lines=2]{\zallmancaps \color{Red} N}{on} ergo credatur hec hora fato urgente venisse, sed Deo potius ordinante, nec siderea necessitas Christi connexuit passionem. Absit enim ut sidera mori cogerent siderum conditorem. Non itaque Christum tempus ut moreretur impegit, sed tempus Christus quo moreretur elegit: qui eciam tempus quo de Virgine natus est cum Patre constituit de quo sine tempore natus est. \R \hyperlink{deduc-me}{Deduc me.}
\newline {\color{Red} Ad Ben. Ant.} Clarifica me pater apud temet ipsum claritate quam habui priusquam mundus fieret alleluya. {\color{Red} Oratio.}
\lettrine[lines=2]{\zallmancaps \color{Blue} P}{resta} \hypertarget{presta-quesumus-ascensionis}{\label{presta-quesumus-ascensionis}} quesumus omnipotens deus: ut nostre mentis intentio quo sollemnitatis venture gloriosus actor ingressus est semper intendat, et quo fide pergit: conversatione perveniat. Per eundem. \ul{Hore diei ut supra cum oratione de vigilia.}
\newline {\color{Red} Ad Vesperas. Super psalmos. Ant.} Ascendens Christus in altum captivam duxit captivitatem dedit dona hominibus alleluya. {\color{Red} Ps.} \psalm{112} \ul{cum ceteris.} {\color{Red} Capitulum.}
\lettrine[lines=2]{\zallmancaps \color{Red} P}{rimum} \hypertarget{primum-quidem-capitulum}{\label{primum-quidem-capitulum}} quidem sermonem feci de omnibus o Theophile que cepit Ihesus facere et docere, usque in diem qua precipiens apostolis per spiritum sanctum quos elegit, assumptus est. \R \hyperlink{omnis-pulchritudo}{Omnis pulchritudo.} {\color{Red} II. Ymnus.}
\hymn{eterne-rex-altissime}{Eterne rex altissime}{
\lettrine[lines=2]{\zallmancaps \color{Blue} E}{terne} rex altissime
redemptor et fidelium,
quo mors soluta deperit,
datur triumphus gratie.
\par {\bfseries \color{Red} S}candens tribunal dextere
patris, potestas omnium
collata est Ihesu celitus
que non erat humanitus.
\par {\bfseries \color{Blue} U}t trina rerum machina
celestium terrestrium,
et infernorum condita
flectat genu iam subdita.
\par {\bfseries \color{Red} T}remunt videntes angeli
versa vice mortalium,
culpat caro purgat caro,
regnat deus dei caro.
\par {\bfseries \color{Blue} T}u esto nostrum gaudium
qui es futurus premium,
sit nostra in te gloria
per cuncta semper secula.
\par {\bfseries \color{Red} G}loria tibi domine
qui scandis supra sidera
cum patre et sancto spiritu,
in sempiterna secula. Amen.
}
\newline \ul{Isti duo versus, scilicet,} Tu esto, \ul{et} Gloria tibi domine qui scandis, \ul{dicantur in fine omnium ymnorum eiusdem metri usque ad Vigiliam Pentecosten in Vesperis exclusive, nisi in festo Corone.}
\newline \V Elevata est magnificentia tua alleluya. \V Super celos deus alleluya.
\newline {\color{Red} Ad Mag. Ant.} Pater manifestavi nomen tuum hominibus quos dedisti michi, nunc autem pro eis rogo non pro mundo, quia ad te vado alleluya. \ul{Oratio.} \hyperlink{presta-quesumus-ascensionis}{Presta quesumus omnipotens deus ut nostre.}
\newline {\color{Red} Ad Nunc Dimittis. Ant.} Alleluya, ascendens Christus in altum alleluya, captivam duxit captivitatem alleluya alleluya.
\newline \ul{Hec antiphona dicatur super Nunc Dimittis usque ad Completorium Vigilie Pentecostes exclusive.}
{\ubsubsection{In Die Ascensionis Domini}{in-die-ascensionis-domini-breviarium}}
\newline {\color{Red} In die ad Matutinas. Invitat.}
\begin{invitatory}[alleluya-christum-invitatorium]
{Alleluya Christum dominum ascendentem in celum, * Venite adoremus alleluya.}
\end{invitatory}
\newline {\color{Red} Ymnus.} \hyperlink{eterne-rex-altissime}{Eterne.}
\newline {\color{Red} In Nocturno. Ant.} Elevata est magnificentia tua super celos deus alleluya. {\color{Red} Ps.} \psalm{8} {\color{Red} Ant.} Dominus in templo sancto suo, dominus in celo alleluya. {\color{Red} Ps.} \psalm{10} {\color{Red} Ant.} A summo celo egressio eius et occursus eius alleluya. %sic
{\color{Red} Ps.} \psalm{18} \V Elevata est.
\newline {\color{Red} Lectio Prima. Secundum Marcum.}
\lettrine[lines=2]{\zallmancaps \color{Red} I}{n} illo tempore: Recumbentibus undecim discipulis apparuit illis Ihesus, et exprobravit incredulitatem illorum et duritiam cordis, quia iis qui viderant eum resurrexisse non crediderant. Et reliqua.
\newline {\color{Red} Omelia Beati Gregorii Pape.}
\lettrine[lines=2]{\zallmancaps \color{Blue} Q}{uod} resurrectionem dominicam discipuli tarde crediderunt, non tam illorum infirmitas quam nostra ut ita dicam futura firmitas fuit. Ipsa namque resurrectio illis dubitantibus per multa argumenta monstrata est. Que dum nos legentes agnoscimus? quid aliud quam de eorum dubitatione solidamur?
\lettrine[lines=2]{\zallmancaps \color{Red} P}{ost} \hypertarget{post-passionem}{\label{post-passionem}} suam per dies quadraginta apparens eis et loquens de regno dei alleluya, et videntibus illis elevatus est alleluya, * Et nubes suscepit eum ab oculis eorum alleluya. \V Et convescens precepit eis ab Hierosolimis ne discederent sed expectarent promissionem patris. Et nubes.
\newline {\color{Red} Lectio Secunda.}
\lettrine[lines=2]{\zallmancaps \color{Blue} M}{inus} enim michi Maria Magdalene prestitit que citius credidit, quam Thomas qui diu dubitavit. Ille etenim dubitando vulnerum cicatrices tetigit, et de nostro pectore dubietatis vulnus amputavit. Ad insinuandam quoque veritatem dominice resurrectionis notandum nobis est, quid Lucas referat dicens. Convescens precepit eis a Hierosolimis ne discederent. Et post pauca. Videntibus illis elevatus est, et nubes suscepit eum ab oculis eorum.
\begin{responsory}[omnis-pulchritudo]
{Omnis pulchritudo domini exaltata est super sidera, species eius in nubibus celi, * Et nomen eius in eternum permanet alleluya.}{A summo celo egressio eius et occursus eius usque ad summum eius.}
\end{responsory}%
\newline {\color{Red} Lectio Tertia.}
\lettrine[lines=2]{\zallmancaps \color{Red} N}{otate} verba, signate mysteria. Convescens elevatus est. Comedit et ascendit, ut videlicet per effectum comestionis: veritas patesceret carnis. Marcus vero priusquam celum Dominus ascendat, eum de cordis atque infidelitatis duritia increpasse discipulos commemorat.
\begin{responsory-final}[viri-galilei]
{Viri Galilei quid admiramini aspicientes in celum, * Quemadmodum vidistis eum ascendentem in celum ita veniet, * Alleluya alleluya alleluya.}{Cumque intuerentur in celum euntem illum ecce duo viri astiterunt iuxta illos in vestibus albis qui et dixerunt.}{Alleluya}
\end{responsory-final}%
\newline \V Ascendo ad patrem meum et patrem vestrum alleluya. \R Deum meum et deum vestrum alleluya.
\newline \ul{Hic versiculus dicatur cotidie per octavas quando de ipsis agitur, et in Octava die.}
\newline {\color{Red} In Laudibus. Ant.} Viri Galilei quid aspicitis in celum, hic Ihesus qui assumptus est a vobis in celum sic veniet alleluya. {\color{Red} Ps.} \psalm{92} {\color{Red} Ant.} Cumque intuerentur in celum euntem illum dixerunt alleluya. {\color{Red} Ant.} Elevatus manibus ferebatur in celum et benedixit eis alleluya. {\color{Red} Ant.} Exaltate regem regum et ymnum dicite deo alleluya. {\color{Red} Ant.} Videntibus illis elevatus est et nubes suscepit eum in celum alleluya. {\color{Red} Cap.} \hyperlink{primum-quidem-capitulum}{Primum quidem.} {\color{Red} Ymnus.}
\hymn{tu-christe-nostrum}{Tu Christe nostrum gaudium}{
\lettrine[lines=2]{\zallmancaps \color{Red} T}{u} Christe nostrum gaudium
manens olimpo preditum,
mundi regis, qui fabricam
mundana vincens gaudia.
\par {\bfseries \color{Blue} H}inc te precantes quesumus
ignosce culpis omnibus,
et corda sursum subleva
ad te superna gratia.
\par {\bfseries \color{Red} U}t cum rubente ceperis
clarere nube Iudicis,
penas repellas debitas
reddas coronas perditas.
\par {\bfseries \color{Blue} T}u esto nostrum gaudium
qui es futurus premium,
sit nostra in te gloria
per cuncta semper secula.
\par {\bfseries \color{Red} G}loria tibi domine
qui scandis supra sidera
cum patre et sancto spiritu,
in sempiterna secula. Amen.
}
\newline \V A summo celo egressio eius alleluya. \R Et occursus eius usque ad summum eius, alleluya.
\newline {\color{Red} Ad Ben. Ant.} Ascendo ad patrem meum et patrem vestrum deum meum et deum vestrum, alleluya. {\color{Red} Oratio.}
\lettrine[lines=2]{\zallmancaps \color{Blue} C}{oncede} \hypertarget{concede-ascensionis}{\label{concede-ascensionis}} quesumus omnipotens deus, ut qui hodierna die unigenitum tuum redemptorem nostrum ad celos ascendisse credimus, ipsi quoque mente in celestibus habitemus. Per eundem. \ul{Ad horas antiphone Laudis.}
\newline {\color{Red} Ad Primam.} \R Ihesu Christe. \V Qui scandis supra sidera. \ul{Iste versus dicatur cotidie usque ad Pentecosten exclusive.}
\newline {\color{Red} Ad Terciam. Cap.} \hyperlink{primum-quidem-capitulum}{Primum quidem.}
\begin{responsory-breve}[elevata-est]
{Elevata est magnificentia tua, * Alleluya alleluya.}{Super celos deus.}
\end{responsory-breve}%
\V Ascendit deus in iubilatione alleluya. \R Et dominus in voce tube alleluya.
\newline {\color{Red} Ad Sextam. Capitulum.}
\lettrine[lines=2]{\zallmancaps \color{Red} E}{t} convescens precepit eis ab Hierosolimis ne discederent, sed expectarent promissionem patris quam audistis inquit per os meum.
\begin{responsory-breve}[ascendit-deus]
{Ascendit deus in iubilatione, * Alleluya alleluya.}{Et dominus in voce tube.}
\end{responsory-breve}%
\V Ascendens Christus in altum alleluya. \R Captivam duxit captivitatem alleluya.
\newline {\color{Red} Ad Nonam. Capitulum.}
\lettrine[lines=2]{\zallmancaps \color{Blue} V}{iri} Galilei quid statis aspicientes in celum, hic Ihesus qui assumptus est a vobis in celum, sic veniet quemadmodum vidistis eum euntem in celum.
\begin{responsory-breve}[ascendens-christus]
{Ascendens Christus in altum, * Alleluya alleluya.}{Captivam duxit captivitatem.}
\end{responsory-breve}%
\V A summo celo.
\newline \ul{Predicto modo dicantur hore per octavas quando de Ascensione agitur.}
\newline {\color{Red} Ad Vesperas. Super psalmos. Ant.} Viri Galilei. {\color{Red} Ps.} \psalm{109} {\color{Red} Ps.} \psalm{110} {\color{Red} Ps.} \psalm{111} {\color{Red} Ps.} \psalm{112} {\color{Red} Ps.} \psalm{116} \ul{Hi psalmi dicantur per octavas quando de octavis agitur cum predicta antiphona et capitulo,} \hyperlink{primum-quidem-capitulum}{Primum quidem,} \ul{et ymno,} \hyperlink{eterne-rex-altissime}{Eterne rex,} \ul{et \Vbar .} Elevata.
\newline {\color{Red} Ad Mag. Ant.} O rex glorie domine virtutem qui triumphator hodie super omnes celos ascendisti, ne derelinquas nos orphanos sed mitte promissum patris in nos, spiritum veritatis alleluya.
\newline {\color{Red} Feria Sexta. Ad Matutinas. Invitat.}
\begin{invitatory}[ascendens-christus-invitatorium]
{Ascendens Christus in altum, * Alleluya.}
\end{invitatory}
\newline \ul{Hoc invitatorium dicatur per octavas quando de octavis agitur, nisi in Dominica, Ymnus sicut in die Ascensionis.}
\newline {\color{Red} Super psalmos. Ant.} Exaltare domine in virtute tua cantabimus et psallemus, alleluya. {\color{Red} Ps.} \psalm{20} {\color{Red} Ps.} \psalm{29} {\color{Red} Ps.} \psalm{46} \V Ascendit deus in iubilatione alleluya.
% https://mlat.uzh.ch/browser/cps_2.BerCla.InAsDo#:2
\newline \ul{Sermo Beati Bernardi.} {\color{Red} Lectio Prima. \grecross}
\lettrine[lines=2]{\zallmancaps \color{Red} S}{ollempnitas} ista fratres karissimi gloriosa est et gaudiosa, in qua et Christo gloria et nobis leticia exhibetur. Consummatio enim et adimpletio est reliquarum sollemnitatum, et felix clausura totius itinerarii filii Dei. Qui enim descendit ipse est et qui ascendit hodierna die super omnes celos, ut adimpleret omnia.
\lettrine[lines=2]{\zallmancaps \color{Blue} S}{i} \hypertarget{si-enim-non}{\label{si-enim-non}} enim non abiero Paraclitus non veniet ad vos, si autem abiero mittam eum ad vos, cum autem venerit ille docebit vos omnem veritatem, * Alleluya. \V Non enim loquetur a semetipso, sed quecumque audiet loquetur, et que ventura sunt annuntabit vobis. Alleluya.
\newline {\color{Red} Lectio Secunda.}
\lettrine[lines=2]{\zallmancaps \color{Red} I}{am} cum se dominum universorum que sunt in terra et in mari et in inferno probasset filius Dei, non restabat nisi ut aeris et celorum se esse dominum comprobaret. Terra enim cognovit Dominum, quia ad vocem eius Lazarum mortuum reddidit. Cognovit mare, quia solidum se prebuit sub pedibus eius. Cognovit infernus, cuius ipse portas ereas confregit, ubi et ligavit insatiabilem homicidam.
\begin{responsory}[ascendens-in-altum]
{Ascendens in altum alleluya, captivam duxit captivitatem, * Dedit dona hominibus alleluya alleluya.}{Ascendit deus in iubilatione et dominus in voce tube.}
\end{responsory}%
\newline {\color{Red} Lectio Tertia.}
\lettrine[lines=2]{\zallmancaps \color{Blue} A}{d} claudendam igitur tunicam tuam inconsutilem Domine Ihesu, ad perficiendam fidei nostre integritatem: restat ut videntibus discipulis per medium aeris sicut aeris dominus ascendas super omnes celos, et tunc probabitur quia Dominus universorum tu es, et iam tibi profecto debebitur ut in nomine tuo omne genu flectatur, celestium, terrestrium, et infernorum, et omnis lingua confiteatur, quia tu es in gloria et in dextera Patris.
\begin{responsory-final}[non-relinquam]
{Non relinquam vos orphanos alleluya vado et venio ad vos, * Et gaudebit cor vestrum, * Alleluya alleluya.}{Ego rogabo patrem et alium Paraclytum dabit vobis.}{Alleluya}
\end{responsory-final}%
\newline \ul{Per octavas quando de octavis agitur, Ant.} Viri Galilei, \ul{sola dicatur in Laudibus, nisi in Dominica et in octava die. Et Cap., Ymnus, \Vbar , et Oratio de die Ascensionis. Sequentes vero antiphone dicantur per octavas eciam in Dominica.}
\newline {\color{Red} Ad Ben. vel ad memoriam in Laudibus. Ant.} Rogabo patrem meum et alium Paraclytum dabit vobis alleluya.
\newline {\color{Red} Ad Mag. vel ad memoriam in Vesperis. Ant.} Euntes in mundum universum predicate evangelium omni creature alleluya, qui crediderit et baptizatus fuerit salvus erit alleluya, qui vero non crediderit condemnabitur, alleluya.
\newline {\color{Red} Sabbato. Ad Matutinas. Super psalmos. Ant.} Nimis exaltatus es alleluya super celos deus, alleluya. {\color{Red} Ps.} \psalm{96} {\color{Red} Ps.} \psalm{98} {\color{Red} Ps.} \psalm{102} \V Ascendens Christus in altum alleluya. \ul{Idem in eodem.} \grecross
\newline {\color{Red} Lectio Prima.}
\lettrine[lines=2]{\zallmancaps \color{Red} Q}{uantus} putatis fratres dolor et timor irruperit apostolica pectora, cum viderunt dominum a se tolli, et attolli in aera, et si angelico comitatum obsequio non tamen fultum auxilio, sed gradientem in multitudine fortitudinis sue? Dolor nimius erat quia videbant illum propter quem omnia reliquerant a suis aspectibus tolli, ut non possent ablato sponso sponsi filii non lugere. Timor: quia orphani relinquebantur in medio Iudeorum necdum confirmati virtute ex alto. {\color{Red} Responsorium.}
\lettrine[lines=2]{\zallmancaps \color{Blue} E}{xaltare} \hypertarget{exaltare-domine}{\label{exaltare-domine}} domine alleluya, * In virtute tua alleluya. \V Cantabimus et psallemus. In virtute.
\newline {\color{Red} Lectio Secunda.}
\lettrine[lines=2]{\zallmancaps \color{Red} B}{enedicens} ergo eis ferebatur in celum, forte concussis illius singularis misericordie visceribus, cum miseros suos et pauperem suam scolam relinqueret. Quam digna ista processio ad quam nec ipsi apostoli digni fuerunt admitti, cum et animarum sanctarum et celestium virtutum triumphali pompa deductus est ad Patrem. Vere adimplevit omnia, quia natus inter homines et cum hominibus conversatus, ab hominibus, et pro hominibus passus et mortuus est, et resurrexit: ascendit et sedet ad dexteram Dei.
\begin{responsory}[ponis-nubem]
{Ponis nubem ascensum tuum domine, * Qui ambulas super pennas ventorum alleluya.}{Confessionem et decorem induisti amictus lumine sicut vestimento.}
\end{responsory}%
\newline {\color{Red} Lectio Tertia.}
\lettrine[lines=2]{\zallmancaps \color{Blue} V}{eruntamen} quid michi et sollemnitatibus istis? Quis me consolabitur Domine Ihesu, quia te non vidi cruce suspensum, plagis lividum, pallidum morte, quia non sum crucifixo compassus, obsecutus mortuo, ut saltem lacrimis meis loca illa vulnerum delinirem? Quomodo me dereliquisti insalutatum, cum te in alta celorum recepisti? Prorsus renuisset consolari anima mea: nisi me angeli in voce exultationis prevenissent, qui dixerunt. Hic Ihesus qui assumptus est a vobis in celum, sic veniet quemadmodum vidistis eum euntem in celum.
\begin{responsory-final}[non-conturbetur]
{Non conturbetur cor vestrum, ego vado ad patrem, et dum assumptus fuero a vobis: mittam vobis alleluya spiritum veritatis, * Et gaudebit cor vestrum, * Alleluya.}{Ego rogabo patrem et alium Paraclitum dabit vobis.}{Alleluya}
\end{responsory-final}%
{\ubsubsection{Dominica Infra Octavas}{dominica-infra-octavas-breviarium}}
\newline \ul{Ad Mat. Invitat, Ymnus, \Vbar , ac psalmi, Responsoria, sicut in die Ascensionis.}
\newline {\color{Red} Lectio Prima. Secundum Iohannem.}
\lettrine[lines=2]{\zallmancaps \color{Red} I}{n} illo tempore: Dixit Ihesus discipulis suis. Cum venerit Paraclytus quem ego mittam vobis a patre, spiritum veritatis qui a patre procedit, ille testimonium perhibebit de me. Et reliqua.
\newline {\color{Red} Omelia Venerabilis Bede Presbiteri.} \grecross
\lettrine[lines=2]{\zallmancaps \color{Blue} N}{otandum} in primis quod Dominus Spiritum veritatis et a se mittendum esse testatur, et eundem mox a Patre procedere subiungit, non quia idem Spiritus aliter a Patre procedit quam cum a Filio mittitur, sed ideo eum Filius a se mitti et a Patre dicit procedere, ut aliam Patris, aliam esse suam personam designet, atque ut in eadem distinctione personarum unam esse operationem ac voluntatem suam cum voluntate Patris et operatione denunciet.
\newline {\color{Red} Lectio Secunda.}
\lettrine[lines=2]{\zallmancaps \color{Red} C}{um} enim eiusdem Spiritus gratia datur hominibus, mittitur profecto Spiritus a Patre, mittitur et a Filio, procedit a Patre, procedit et a Filio, quia eius missio ipsa processio est qua ex Patre procedit et Filio. Venit et sua sponte, quia sicut coequalis est Patri et Filio, ita eandem habet voluntatem cum Patre et Filio communem. Spiritus enim ubi vult spirat, et sicut Apostolus enumeratis donis celestibus ait, hec omnia operatur unus atque idem Spiritus dividens singulis prout vult.
\newline {\color{Red} Lectio Tertia.}
\lettrine[lines=2]{\zallmancaps \color{Blue} A}{dveniens} autem Spiritus testimonium de Domino perhibuit, quia discipulorum cordibus aspirans, omnia que de illo erant sentienda mortalibus, clara illis luce revelavit, quod videlicet equalis et consubstantialis Patri erat ante secula, quod consubstantialis factus est nobis in fine seculorum: quod de
%pp109
Virgine natus sine peccato vixit in mundo, quod quando voluit, et qua morte voluit transivit de mundo, quod veraciter mortem resurgendo destruxit, veram in qua passus et resurrexit carnem ad celos ascendendo levavit, atque in paterne dextere gloria constituit: quod omnia prophetarum scripta illi testimonium perhibent, quod confessio nominis eius usque ad fines erat propaganda terrarum, ceteraque fidei illius mysteria Spiritus sancti testimonio sunt reserata discipulis.
\newline \ul{In Laudibus Ant.} Viri Galilei, \ul{cum ceteris.}
\newline {\color{Red} Ad Ben. Ant.} Rogabo. {\color{Red} Oratio.} \hyperlink{concede-ascensionis}{Concede.}
\newline {\color{Red} Ad Memoriam de Dominica. Ant.} Cum venerit Paraclitus quem ego mittam vobis spiritum veritatis qui a patre procedit ille testimonium perhibebit de me alleluya. \V Ascendit deus in iubilatione alleluya. {\color{Red} Oratio.}
\lettrine[lines=2]{\zallmancaps \color{Red} O}{mnipotens} sempiterne deus fac nos tibi semper et devotam gerere voluntatem, et maiestati tue sincero corde servire. Per.
\newline {\color{Red} In Vesperis. Ad Mag. Ant.} Euntes.
\newline {\color{Red} Ad Memoriam de Dominica. Ant.} Hec locutus sum vobis, ut cum venerit hora eorum reminiscamini quia ego dixi vobis alleluya. \V Ascendens Christus in altum alleluya. \ul{Oratio ut supra.}
\newline \ul{Feria secunda, tercia et quarta et in octava repetantur lectiones de sexta feria et de sabbato supra notate.}
\newline {\color{Red} Feria Secunda. At Mat. Super psalmos. Ant.} Dominus in templo sancto suo dominus in celo alleluya. \ul{Psalmi, \Vbar , Responsoria sicut in die Ascensionis.}
\newline {\color{Red} Feria Tercia. At Mat. Super psalmos. Ant.} Exaltabo te domine quoniam suscepisti me alleluya. \ul{Psalmi, \Vbar , Responsoria sicut in precedenti feria sexta.}
\newline {\color{Red} Feria Quarta. At Mat. Super psalmos. Ant.} Dominus in Syon alleluya magnus et excelsus alleluya. \ul{Psalmi, \Vbar , Responsoria sicut in precedenti sabbato.}%typo exelsus
{\ubsubsection{In Octava Ascensionis}{in-octava-ascensionis-breviarium}}
\newline {\color{Red} Fiat sicut in Festo Simplici.} \ul{In Vesperis precedentibus super psalmos Ant.} Viri, \ul{psalmi,} \psalm{109} \ul{et ceteri. Cetera omnia ad Vesperas et ad Matutinas et ad horas sicut in die Ascensionis excepto quod Responsorium in Primis Vesperis non dicitur, et preces dicuntur in Completorio et in Prima.}
\newline {\color{Red} Feria Sexta post Octavam Ascensionis. Ad Mat. Invitati} \hyperlink{alleluya-invitatorium}{Alleluya alleluya, * Alleluya.} {\color{Red} Ymnus.} \hyperlink{eterne-rex-altissime}{Eterne rex.} {\color{Red} Super psalmos. Ant.} Regnabit deus super gentes deus sedet super sedem sanctam suam alleluya. \ul{Psalmi, \Vbar , et Responsoria sicut in precedenti feria sexta.}
\newline {\color{Red} Lectio Prima.}
% Acts 1:12
\lettrine[lines=2]{\zallmancaps \color{Blue} I}{n} diebus illis reversi sunt apostoli Hierosolimam a monte qui vocatur Oliveti qui est iuxta Hierusalem, sabbati habens iter. Et cum introissent in cenaculum, ascenderunt ubi manebant Petrus et Iohannes: Iacobus et Andreas, Philippus et Thomas, Bartholomeus et Matheus, Iacobus Alphei et Simon Zelotes et Iudas Iacobi. Hi omnes erant perseverantes unanimiter in oratione cum mulieribus et Maria matre Ihesu et fratribus eius.
\newline {\color{Red} Lectio Secunda.}
\lettrine[lines=2]{\zallmancaps \color{Red} I}{n} diebus illis exurgens Petrus in medio fratrum, dixit. Erat autem turba hominum fere centum viginti. Viri fratres oportet impleri scripturam quam predixit Spiritus sanctus per os David de Iuda qui fuit dux eorum qui comprehenderunt Ihesum, qui connumeratus erat in nobis, et sortitus est sortem ministerii huius.
\newline {\color{Red} Lectio Tertia.}
\lettrine[lines=2]{\zallmancaps \color{Blue} E}{t} hic quidem possedit agrum de mercede iniquitatis, et suspensus crepuit medius, et diffusa sunt omnia viscera eius. Et notum factum est omnibus habitantibus Hierusalem, ita ut appellaretur ager ille lingua eorum Acheldemach, hoc est ager sanguinis. Scriptum est enim in libro psalmorum. Fiat commoratio eorum deserta, et non sit qui habitet in ea, et episcopatum eius accipiat alter.
\newline \ul{Non dicatur hac die nec sequenti,} \psalm{tedeum} \ul{nisi festum fuerit quia quasi de feria agitur.}
\newline \V Ascendo ad patrem. \ul{In Laudibus. Ant.} Viri Galilei, \ul{sola dicatur. Cap., Ymnus, \Vbar , de Ascensione.}
\newline {\color{Red} Ad Ben. Ant.} Cum venerit Paraclytus quem ego mittam vobis spiritum veritatis qui a patre procedit: ille testimonium perhibebit de me alleluya. \ul{Oratio de Dominica.}
\newline \ul{Ad Primam Ant.} Viri, \ul{Cap,} Pacem et veritatem. \ul{Ad ceteras horas diei antiphone, Capitula, Responsoria, \Vbar , de Ascensione oratio de Dominica.}
\newline \ul{Ad Vesperas. Antiphona, psalmi, Cap., Ymnus, \Vbar , de Ascensione.}
\newline {\color{Red} Ad Mag. Ant.} Hec locutus sum vobis, ut cum venerit hora eorum reminiscamini quia ego dixi vobis alleluya. \ul{Oratio de Dominica. Si festum in hac sexta feria evenerit nulla memoria fiat de Ascensione.}
{\ubsubsection{In Vigilia Pentecostes}{in-vigilia-pentecostes-breviarium}}
\newline {\color{Red} Ad Matutinas. Invitat.} \hyperlink{alleluya-invitatorium}{Alleluya alleluya, * Alleluya.} {\color{Red} Ymnus.} \hyperlink{eterne-rex-altissime}{Eterne rex.} {\color{Red} Super psalmos. Ant.} Dominus in celo alleluya paravit sedem suam alleluya. \ul{Psalmi, \Vbar , et Responsoria sicut in sabbato precedenti.}
\newline {\color{Red} Lectio Prima. Secundum Iohannem.}
\lettrine[lines=2]{\zallmancaps \color{Red} I}{n} illo tempore: Dixit Ihesus discipulis suis. Si diligitis me: mandata mea servate. Et ego rogabo patrem et alium Paraclytum dabit vobis. Et reliqua.
\newline {\color{Red} Omelia Venerabilis Bede Presbiteri.}
\lettrine[lines=2]{\zallmancaps \color{Blue} Q}{uia} sancti Spiritus hodie fratres karissimi celebramus adventum, debemus ipsi congruere sollemnitati quam colimus. Hoc etenim ordine tanta festivitatis huius digne gaudia celebramus, si nos quoque Domino opitulante aptos reddiderimus ad quos Spiritus sanctus advenire, et in quibus habitare dignetur.
\newline {\color{Red} Lectio Secunda.}
\lettrine[lines=2]{\zallmancaps \color{Red} E}{a} autem solummodo ratione adventu et illustratione sancti Spiritus apti existimus, si et corda nostra divino amore repleta, et corpora sint Dominicis mancipata preceptis. Unde in exordio lectionis huius evangelice discipulis veritas ait. Si diligitis me mandata mea servate, et ego rogabo Patrem et alium Paraclytum dabit vobis.
\newline {\color{Red} Lectio Tertia.}
\lettrine[lines=2]{\zallmancaps \color{Blue} P}{araclitus} quippe consolator interpretatur. Et Spiritus sanctus recte Paraclytus vocatur, quia corda fidelium ne inter huius seculi adversa deficiant: celestis vite desideriis sublevat ac reficit. Unde in actibus apostolorum crescente sancta ecclesia dicitur. Et edificabatur ambulans in timore Domini, et consolatione Spiritus sancti replebatur.
\newline {\color{Red} In Laudibus. Super psalmos. Ant.} Viri Galilei. {\color{Red} Ps.} \psalm{92} \ul{cum ceteris. Cap., Ymnus, \Vbar , ut in Ascensione.}
\newline {\color{Red} Ad Ben. Ant.} Si diligitis me mandata mea servate alleluya alleluya alleluya. {\color{Red} Oratio.}
\lettrine[lines=2]{\zallmancaps \color{Red} P}{resta} \hypertarget{presta-quesumus-pentecostes}{\label{presta-quesumus-pentecostes}} quesumus omnipotens deus ut claritatis tue super nos splendor effulgeat et lux tue lucis corda eorum qui per gratiam tuam renati sunt sancti spiritus illustratione confirmet. Per, eiusdem.
\newline \ul{Hore diei ut in precedenti sexta feria cum oratione de vigilia.}
\newline {\color{Red} Ad Vesperas. Super psalmos. Ant.} Veni sancte spiritus reple tuorum corda fidelium, et tui amoris in eis ignem accende, qui per diversitatem linguarum multarum gentes in unitatem fidei congregasti alleluya alleluya. {\color{Red} Ps.} \psalm{145} {\color{Red} Ps.} \psalm{146} {\color{Red} Ps.} \psalm{147} {\color{Red} Capitulum.}
\lettrine[lines=2]{\zallmancaps \color{Blue} F}{actus} est repente de celo sonus tanquam advenientis spiritus vehementis, et replevit totam domum ubi erant apostoli sedentes. \R \hyperlink{repleti-sunt-omnes}{Repleti sunt.} {\color{Red} II. Ymnus.}
\hymn{beata-nobis-gaudia}{Beata nobis gaudia}{
\lettrine[lines=2]{\zallmancaps \color{Red} B}{eata} nobis gaudia,
anni reduxit orbita,
cum spiritus Paraclytus
effulsit in discipulos.
\par {\bfseries \color{Blue} I}gnis vibrante lumine,
lingue figuram detulit,
verbis ut essent proflui
et caritate fervidi.
\par {\bfseries \color{Red} L}inguis loquuntur omnium,
turbe pavent gentilium,
musto madere deputant
quos spiritus repleverat.
\par {\bfseries \color{Blue} P}atrata sunt hec mystice
Pasche peracto tempore,
sacro dierum numero
quo lege fit remissio.
\par {\bfseries \color{Red} T}e nunc deus piissime
vultu precamur cernuo,
illapsa nobis celitus
largire dona spiritus.
\par {\bfseries \color{Blue} D}udum sacrata pectora
tua replesti gratia,
dimitte nunc peccamina
et da quieta tempora.
\par {\bfseries \color{Red} S}it laus patri cum filio
sancto simul Paraclyto
nobisque mittat filius
karisma sancti spiritus. Amen.
}
\newline \ul{Isti duo versus scilicet,} Dudum sacrata, \ul{et} Sit laus, \ul{dicantur in fine Ymnorum per totam ebdomadam nisi in Ymno,} \hyperlink{veni-creator-spiritus}{Veni creator.} \ul{In quo tantum dicitur,} Sit laus.
\newline \V Repleti sunt omnes spiritu sancto alleluya. \R Et ceperunt loqui alleluya.
\newline {\color{Red} Ad Mag. Ant.} Non vos relinquam orphanos alleluya, vado et venio ad vos alleluya, et gaudebit cor vestrum alleluya. {\color{Red} Oratio.} \hyperlink{presta-quesumus-pentecostes}{Presta quesumus.}
\newline {\color{Red} Ad Completorium. Ant.} Alleluya alleluya alleluya alleluya. {\color{Red} Ps.} \psalm{4} \psalm{30a} \psalm{133} \ul{et dicantur usque ad sabbatum exclusive.}
\newline {\color{Red} Ad Nunc Dimittis. Ant.} Alleluya, spiritus Paraclytus alleluya, docebit vos omnia alleluya alleluya. \ul{Hec antiphona dicatur usque ad Completorium sequentis sabbati exclusive.}
{\ubsubsection{In Die Pentecostes}{in-die-pentecostes-breviarium}}
\newline {\color{Red} In Die ad Matutinas. Invitatorium.}
\begin{invitatory}[alleluya-spiritus-invitatorium]
{Alleluya, Spiritus domini replevit orbem terrarum, * Venite adoremus alleluya.}
\end{invitatory}
\hymn{iam-christus-astra}{Iam Christus astra}{
\lettrine[lines=2]{\zallmancaps \color{Blue} I}{am} Christus astra ascenderat
regressus unde venerat,
promissum patris munere
sanctum daturus spiritum.
\par {\bfseries \color{Red} S}ollemnis surgebat dies
quo mystico septemplici
orbis volutus septies
signat beata tempora.
\par {\bfseries \color{Blue} D}um hora cunctis tercia
repente mundus intonat,
orantibus apostolis
deum venisse nunciat.
\par {\bfseries \color{Red} D}e patris ergo lumine
decorus ignis missus est,
qui fida Christi pectora
calore verbi complevit.
\par {\bfseries \color{Blue} D}udum sacrata pectora
tua replesti gratia,
dimitte nunc peccamina
et da quieta tempora.
\par {\bfseries \color{Red} S}it laus patri cum filio
sancto simul Paraclyto
nobisque mittat filius
karisma sancti spiritus. Amen.
}
\newline {\color{Red} Ant.} Factus est repente de celo sonus advenientis spiritus vehementis alleluya alleluya. {\color{Red} Ps.} \psalm{47} {\color{Red} Ant.} Confirma hoc deus quod operatus es in nobis a templo sancto tuo quod est in Hierusalem alleluya alleluya. {\color{Red} Ps.} \psalm{67} {\color{Red} Ant.} Emitte spiritum tuum et creabuntur et renovabis faciem terre alleluya alleluya. {\color{Red} Ps.} \psalm{103} \V Repleti sunt.
\newline {\color{Red} Lectio Prima. Secundum Iohannem.}
\lettrine[lines=2]{\zallmancaps \color{Red} I}{n} illo tempore: Dixit Ihesus discipulis suis. Si quis diligit me sermonem meum servabit, et pater meus diliget eum, et ad eum veniemus, et mansionem apud eum faciemus. Et reliqua.
\newline {\color{Red} Omelia Beati Gregorii Pape.}
\lettrine[lines=2]{\zallmancaps \color{Blue} L}{ibet} fratres karissimi evangelice verba lectionis sub brevitate transcurrere, ut post diutius liceat in contemplatione tante sollemnitatis immorari. Hodie namque Spiritus sanctus repentino sonitu super discipulos venit, mentesque carnalium in sui amorem permutavit. \R
\lettrine[lines=2]{\zallmancaps \color{Red} D}{um} \hypertarget{dum-complerentur-dies}{\label{dum-complerentur-dies}} complerentur dies Pentecostes erant omnes pariter dicentes alleluya, et subito factus est sonus de celo alleluya, * Tanquam spiritus vehemens replevit totam domum alleluya alleluya. \V Dum ergo essent discipuli in unum congregati propter metum Iudeorum sonus repente de celo venit super eos. Tanquam. 
\newline {\color{Red} Lectio Secunda.}
\lettrine[lines=2]{\zallmancaps \color{Blue} E}{t} foris apparentibus linguis igneis, intus facta sunt corda flammantia, quia dum Deum in ignis visione suscipiunt, per amorem suaviter arserunt. Ipse namque Spiritus sanctus amor est. Unde et Iohannes dicit. Deus caritas est. Qui ergo mente integra Deum desiderat, profecto iam habet quem amat. Neque enim posset Deum quisquam diligere: si eum quem diligit non haberet.
\begin{responsory}[repleti-sunt-omnes]
{Repleti sunt omnes spiritu sancto et ceperunt loqui prout spiritus sanctus dabat eloqui illis, * Et convenit multitudo dicentium alleluya.}{Loquebantur variis linguis apostoli magnalia dei.}
\end{responsory}%
\newline {\color{Red} Lectio Tertia.}
\lettrine[lines=2]{\zallmancaps \color{Red} S}{ed} ecce si unusquisque vestrum requiratur an diligat Deum, tota fiducia et secura mente respondet. Diligo. In ipso autem lectionis exordio audistis quid Veritas dicat. Siquis diligit me sermonem meum servabit. Probatio ergo dilectionis, exhibitio est operis. Hinc in epistola sua idem Iohannes dicit. Qui dicit quia diligo Deum, et mandata eius non custodit: mendax est. Vere enim Deum diligimus, si mandata eius servamus.
\begin{responsory-final}[spiritus-sanctus-procedens]
{Spiritus sanctus procedens a throno apostolorum pectora invisibiliter penetravit, novum sanctificationis signum, * Ut in ore eorum omnium genera nascerentur linguarum, * Alleluya.}{Advenit ignis divinus non comburens sed illuminans et tribuit eis karismatum dona.}{Alleluya}
\end{responsory-final}%
\psalm{tedeum} \V Emitte spiritum tuum et creabuntur alleluya. \R Et renovabis faciem terre, alleluya.
\newline \psalm{tedeum} \ul{et \Vbar .} Emitte. \ul{dicantur per totam ebdomadam.}
\newline {\color{Red} In Laudibus. Ant.} Dum complerentur dies Pentecostes erant omnes pariter dicentes alleluya. {\color{Red} Ps.} \psalm{92} {\color{Red} Ant.} Spiritus domini replevit orbem terrarum alleluya. {\color{Red} Ant.} Repleti sunt omnes spiritu sancto et ceperunt loqui alleluya. {\color{Red} Ant.} Fontes et omnia que moventur in aquis ymnum dicite deo alleluya. {\color{Red} Ant.} Loquebantur variis linguis apostoli magnalia dei alleluya alleluya. {\color{Red} Capitulum.} Factus est repente. {\color{Red} Ymnus.}
\hymn{impleta-gaudent-viscera}{Impleta gaudent viscera}{
\lettrine[lines=2]{\zallmancaps \color{Blue} I}{mpleta} gaudent viscera
afflata sancto spiritu,
voces diversas intonant,
fantur dei magnalia.
\par {\bfseries \color{Red} E}x omni gente cogniti,
Grecis, Latinis barbaris,
cunctisque admirantibus
linguis loquuntur omnibus.
\par {\bfseries \color{Blue} I}udea tunc incredulia,
vesana torvo spiritu,
ructare musti crapulam
alumnos Christi concrepat.
\par {\bfseries \color{Red} S}ed signis et virtutibus
occurrit et docet Petrus,
falsa profari perfidos
Iohelis testimonio.
\par {\bfseries \color{Blue} D}udum sacrata pectora
tua replesti gratia,
dimitte nunc peccamina
et da quieta tempora.
\par {\bfseries \color{Red} S}it laus patri cum filio
sancto simul Paraclyto
nobisque mittat filius
karisma sancti spiritus. Amen.
}
\newline \V Spiritus Paraclytus alleluya. \R Docebit vos omnia alleluya.
\newline {\color{Red} Ad Ben. Ant.} Accipite spiritum sanctum, quorum remiseritis peccata remittuntur eis alleluya. {\color{Red} Oratio.}
\lettrine[lines=2]{\zallmancaps \color{Red} D}{eus} qui hodierna die corda fidelium sancti spiritus illustratione docuisti, da nobis in eodem spiritu recta sapere, et de eius semper consolatione gaudere. Per, eiusdem.
\newline {\color{Red} Ad Primam.} \R \hyperlink{ihesu-christe-fili}{Ihesu Christe.} \V Qui sedes.
\newline {\color{Red} Ad Terciam. Ymnus.} 
\hymn{veni-creator-spiritus}{Veni creator spiritus}{
\lettrine[lines=2]{\zallmancaps \color{Blue} V}{eni} creator spiritus,
mentes tuorum visita,
imple superna gratia
que tu creasti pectora.
\par {\bfseries \color{Blue} Q}ui Paraclitus diceris,
donum dei altissimi,
fons vivus, ignis, caritas
et spiritalis unctio.
\par {\bfseries \color{Red} T}u septiformis munere,
dextre dei tu digitus,
tu rite promisso patris
sermone ditans guttura.
\par {\bfseries \color{Blue} A}ccende lumen sensibus,
infunde amorem cordibus,
infirma nostri corporis
virtute firmans perpeti.
\par {\bfseries \color{Red} H}ostem repellas longius,
pacemque dones protinus,
ductore sic te previo
vitemus omne noxium.
\par {\bfseries \color{Blue} P}er te sciamus da patrem,
noscamus atque filium,
te utriusque spiritum
credamus omni tempore.
\par {\bfseries \color{Red} S}it laus patri cum filio
sancto simul Paraclyto
nobisque mittat filius
karisma sancti spiritus. Amen.
}
\newline \ul{Iste Ymnus dicatur per totam ebdomadam ad Terciam.} {\color{Red} Capitulum.} Factus est.
\begin{responsory-breve}[repleti-sunt-breve]
{Repleti sunt omnes spiritu sancto, * Alleluya alleluya.}{Et ceperunt loqui.}
\end{responsory-breve}%
\V Loquebantur variis linguis apostoli alleluya. \R Magnalia dei alleluya.
\newline {\color{Red} Ad Sextam. Capitulum.} 
\lettrine[lines=2]{\zallmancaps \color{Red} A}{pparuerunt} apostolis dispertite lingue tanquam ignis, seditque supra singulos eorum spiritus sanctus.
\begin{responsory-breve}[loquebantur-variis-breve]
{Loquebantur variis linguis apostoli, * Alleluya alleluya.}{Magnalia dei.}
\end{responsory-breve}%
\V Spiritus domini replevit orbem terrarum alleluya. \R Et hoc quod continet omnia scientiam habet vocis alleluya.
\newline {\color{Red} Ad Nonam. Capitulum.} 
\lettrine[lines=2]{\zallmancaps \color{Blue} F}{acta} autem hac voce convenit multitudo et mente confusa est, quoniam audiebat unusquisque lingua sua illos loquentes.
\begin{responsory-breve}[spiritus-domini-breve]
{Spiritus domini replevit orbem terrarum, * Alleluya alleluya.}{Et hoc quod continet omnia scientiam habet vocis.}
\end{responsory-breve}%
\V Spiritus Paraclytus alleluya. 
\newline {\color{Red} Ad Vesperas. Ant.} Dum complerentur. {\color{Red} Ps.} \psalm{109} {\color{Red} Ps.} \psalm{110} {\color{Red} Ps.} \psalm{111}
\newline \ul{Predicta antiphona et hi tres psalmi dicantur cotidie ad Vesperas usque ad Vesperas sequentes Sabbati exclusive. Cap., Ymnus, \Vbar , ut in Primis Vesperis.}%typo: inprimis
\newline {\color{Red} Ad Mag. Ant.} Hodie completi sunt dies Pentecostes alleluya, hodie spiritus sanctus in igne discipulis apparuit, et tribuens eis karismatum dona misit eos in universum mundum predicare et testificari, qui crediderit et baptizatus fuerit, salvus erit, alleluya.
\newline {\color{Red} Feria Secunda et per ebdomadam. At Mat. Invit.}
\begin{invitatory}[repleti-sunt-invitatorium]
{Repleti sunt omnes spiritu sancto, * Alleluya.}
\end{invitatory}
\newline {\color{Red} Ymnus.} \hyperlink{iam-christus-astra}{Iam Christus} {\color{Red} In Nocturno. Ant.} Factus. 
\newline \ul{Ceteris diebus dicantur cetere sicut sunt in ordine et cum dicte fuerint reincipiantur. Psalmi sicut in precedenti die et dicantur per totam ebdomadam.} \V Repleti sunt.
\newline {\color{Red} Lectio Prima. Secundum Iohannem.}
\lettrine[lines=2]{\zallmancaps \color{Red} I}{n} illo tempore: Dixit Ihesus discipulis suis. Sic deus dilexit mundum, ut filium suum unigenitum daret, ut omnis qui credit in eum non pereat, sed habeat vitam eternam. Et reliqua.
% https://la.wikisource.org/wiki/Homiliae_(Haymo_Halberstatensis)/2
\newline {\color{Red} Omelia Beati Augustini Episcopi.} \grecross
\lettrine[lines=2]{\zallmancaps \color{Blue} N}{otandum} quod eadem de filio Dei unigenito replicat, que de filio hominis in cruce exaltato promiserat dicens: ut omnis qui credit in eum non pereat: sed habeat vitam eternam. Profecto idem conditor et redemptor noster, filius Dei ante secula existens, filius hominis factus est in fine seculorum, ut qui per divinitatis sue potentiam creaverat ad perfruendam vite perennis beatitudinem, ipse per fragilitatem humanitatis nostre nos restauraret ad recipiendam quam perdidimus vitam. \R
\lettrine[lines=2]{\zallmancaps \color{Red} I}{am} \hypertarget{iam-non-dicam}{\label{iam-non-dicam}} non dicam vos servos sed amicos meos, quia omnia cognovistis que operatus sum in medio vestri alleluya, * Accipite spiritum sanctum in vobis Paraclytum, ipse est quem pater mittet vobis alleluya. \V Quorum remiseritis peccata remittuntur eis. Accipite.
\newline {\color{Red} Lectio Secunda.}
\lettrine[lines=2]{\zallmancaps \color{Blue} S}{equitur} autem. Non enim misit Deus filium suum in mundum ut iudicet mundum: sed ut salvetur mundus per ipsum. Ergo quantum in medico est, sanare venit egrotum. Ipse autem se interimit, qui precepta medici servare non vult.
\begin{responsory}[disciplinam-et-sapientiam]
{Disciplinam et sapientiam docuit eos dominus alleluya, firmavit in illis gratiam spiritus sui, * Et intellectu adimplevit corda eorum alleluya alleluya.}{Repentino namque sonitu spiritus sanctus super eos venit.}
\end{responsory}%
\newline {\color{Red} Lectio Tertia.}
\lettrine[lines=2]{\zallmancaps \color{Red} V}{enit} salvator in mundum. Quare salvator dictus est mundi, nisi ut salvet mundum non ut iudicet mundum? Salvari non vis ab ipso: ex te iudicaberis. Vide quid ait. Qui credit in eum non iudicatur. Qui autem non credit, quid dicturum putas nisi iudicetur? Quid addidit? Iam inquit, iudicatus est. Nondum apparuit iudicium, sed iam factum est iudicium. Novit enim Dominus qui sunt eius.
\begin{responsory-doxology}[ultimo-festivitatis]
{Ultimo festivitatis die dicebat Ihesus qui in me credit flumina de ventre eius fluent aque vive, * Alleluya.}{Hoc autem dicebat de spiritu quem accepturi erant credentes in eum.}
\end{responsory-doxology}%
\newline \ul{In Laudibus per ebdomadam hec sola antiphona.} Dum complerentur. \ul{Cap.} Factus est. \ul{Ymnus.} \hyperlink{impleta-gaudent-viscera}{Impleta.} \ul{\Vbar .} Spiritus Paraclitus.
\newline {\color{Red} Ad Ben. Ant.} Sic deus dilexit mundum ut filium suum unigenitum daret, ut omnis qui credit in ipsum non pereat sed habeat vitam eternam alleluya. {\color{Red} Oratio.}
\lettrine[lines=2]{\zallmancaps \color{Blue} D}{eus} qui apostolis tuis sanctum dedisti spiritum: concede plebi tue pie petitionis effectum, ut quibus dedisti fidem, largiaris et pacem. Per, eiusdem.
\newline \ul{Ad horas antiphona, Cap., Responsoria, \Vbar . de die precedenti et dicantur per totam ebdomadam. Ad Vesperas antiphona, psalmi, Cap., Ymnus, \Vbar . sicut in die precedenti et dicantur per totam ebdomadam.}
\newline {\color{Red} Ad Mag. Ant.} Non enim misit deus filium suum in mundum ut iudicet mundum sed ut salvetur mundus per ipsum alleluya.
\newline {\color{Red} Feria Tercia.} \V Loquebantur.
\newline {\color{Red} Lectio Prima. Secundum Iohannem.}
\lettrine[lines=2]{\zallmancaps \color{Red} I}{n} illo tempore: Dixit Ihesus discipulis suis. Amen amen dico vobis, qui non intrat per ostium in ovile ovium sed ascendit aliunde, ille fur est et latro. Et reliqua.
\newline {\color{Red} Omelia Beati Augustini Episcopi.} \grecross
\lettrine[lines=2]{\zallmancaps \color{Blue} D}{ominus} de grege suo et de ostio quo intratur ad ovile similitudinem posuit in hodierna lectione. Dicant ergo Pagani. Bene vivimus. Si per ostium non intrant, quid eis prodest unde gloriantur? Ad hoc enim debet unicuique prodesse bene vivere, ut detur illi semper vivere. \R
\lettrine[lines=2]{\zallmancaps \color{Red} L}{oquebantur} \hypertarget{loquebantur-variis}{\label{loquebantur-variis}} variis linguis apostoli alleluya magnalia dei, * Prout spiritus sanctus dabat eloqui illis alleluya. \V Repleti sunt omnes spiritu sancto et ceperunt loqui. Prout.
\newline {\color{Red} Lectio Secunda.}
\lettrine[lines=2]{\zallmancaps \color{Blue} Q}{uid} enim prodest bene vivere, cui non datur semper vivere? Quia nec bene vivere dicendi sunt, qui finem bene vivendi, vel cecitate nesciunt, vel inflatione contempnunt. Non
%pp110
est autem cuiquam spes vera et certa semper vivendi nisi agnoscat vitam que est Christus, et per ianuam intret in ovile. Querunt ergo plerumque tales homines persuadere hominibus ut bene vivant, et christiani non sunt. Per aliam partem volunt ascendere, rapere et occidere, non ut pastor conservare atque salvare.
\begin{responsory}[apparuerunt-apostolis]
{Apparuerunt apostolis dispertite lingue tanquam ignis alleluya, * Seditque supra singulos eorum spiritus sanctus alleluya alleluya.}{Spiritus domini replevit totam domum ubi erant sedentes apostoli.}
\end{responsory}%
\newline {\color{Red} Lectio Tertia.}
\lettrine[lines=2]{\zallmancaps \color{Red} F}{uerunt} quidam philosophi de virtutibus atque vitiis subtilia multa tractantes, dividentes, diffinientes, ratiocinationes, acutissimas concludentes, libros implentes, sapientiam suam buccis crepantibus ventilantes: qui eciam dicere auderent hominibus. Nos sequimini, sectam nostram tenete, si vultis beate vivere. Sed quia non intrabant per ostium, perdere volebant, mactare, et occidere.
\begin{responsory-final}[spiritus-domini]
{Spiritus domini replevit orbem terrarum, * Et hoc quod continet omnia scientiam habet vocis * Alleluya alleluya.}{Omnium est enim artiphex omnem habens virtutem omnia prospiciens.}{Alleluya}
\end{responsory-final}%
\newline {\color{Red} Ad Ben. Ant.} Amen amen dico vobis qui non intrat per ostium in ovile ovium sed ascendit aliunde ille fur est et latro, qui autem intrat per ostium pastor est ovium, alleluya. {\color{Red} Oratio.}
\lettrine[lines=2]{\zallmancaps \color{Blue} A}{ssit} nobis quesumus domine virtus spiritus sancti, que et corda nostra clementer expurget, et ab omnibus tueatur adversis. Per, eiusdem.
\newline {\color{Red} Ad Mag. Ant.} Ego sum ostium dicit dominus, per me siquis introierit salvabitur et pascua inveniet, alleluya.
\newline {\color{Red} Feria Quarta.} \V Spiritus domini.
\newline {\color{Red} Lectio Prima. Secundum Iohannem.}
\lettrine[lines=2]{\zallmancaps \color{Red} I}{n} illo tempore: Dixit Ihesus turbis Iudeorum. Nemo potest venire ad me nisi pater qui misit me traxerit eum. Et reliqua.
\newline {\color{Red} Omelia Beati.}
\lettrine[lines=2]{\zallmancaps \color{Blue} M}{agna} gratie commendatio. Nemo venit nisi tractus. Quem trahat et quem non trahat, quare illum trahat et illum non trahat, noli velle iudicare, si non vis errare. Semel accipe et intellige. Nondum traheris: ora ut traharis. \R
\lettrine[lines=2]{\zallmancaps \color{Red} A}{dvenit} \hypertarget{advenit-ignis-divinus}{\label{advenit-ignis-divinus}} ignis divinus non comburens sed illuminans non consumens sed lucens et invenit corda discipulorum receptacula munda, * Et tribuit eis karismatum dona alleluya alleluya. \V Invenit eos concordes caritate et illustravit eos inundans divinitas deitatis. Et tribuit.
\newline {\color{Red} Lectio Secunda.}
\lettrine[lines=2]{\zallmancaps \color{Blue} Q}{uid} hic dicemus fratres? Si trahimur ad Christum: ergo inviti credimus. Ergo violentia adhibetur, non voluntas excitatur. Noli te cogitare invitum trahi. Invitus non trahitur animus sed amore. Est quedam voluptas cordis, cui dulcis est ille panis celestis.
\begin{responsory}[facta-autem-hac]
{Facta autem hac voce convenit multitudo et mente confusa est, quia audiebat unusquisque lingua sua, * Illos loquentes magnalia dei alleluya alleluya alleluya.}{Nonne ecce omnes isti qui loquuntur Galilei sunt, et quomodo nos audivimus unusquisque lingua sua.}
\end{responsory}%
\newline {\color{Red} Lectio Tertia.}
\lettrine[lines=2]{\zallmancaps \color{Red} P}{orro} si poete dicere licuit, trahit sua quemque voluptas, non obligatio sed delectatio, quanto fortius nos dicere debemus trahi hominem ad Christum qui delectatur veritate, delectatur beatitudine, delectatur iusticia, delectatur eterna vita, quod totum Christus est? Da amantem et sentit quod dico. Da desiderantem, da ferventem, da in ista solitudine peregrinantem atque sicientem, et ad fontem eterne patrie suspirantem, da talem, et scit quid dicam. Ramum viridem ostendis ovi, et trahis eam.
\begin{responsory-final}[spiritus-sanctus-replevit]
{Spiritus sanctus replevit totam domum ubi erant apostoli et apparuerunt illis dispertite lingue tanquam ignis seditque supra singulos eorum, * Et repleti sunt omnes spiritu sancto, * Alleluya alleluya alleluya.}{Dum essent discipuli in unum congregati propter mecum Iudeorum sonus repente de celo venit super eos.}{Alleluya}
\end{responsory-final}%
\newline {\color{Red} Ad Ben. Ant.} Amen amen dico vobis qui credit in me habet vitam eternam, alleluya, alleluya. {\color{Red} Oratio.}
\lettrine[lines=2]{\zallmancaps \color{Blue} M}{entes} nostras quesumus domine Paraclytus qui a te procedit illuminet, et inducat in omnem sicut tuus promisit filius veritatem. Qui tecum, eiusdem.
\newline {\color{Red} Ad Mag. Ant.} Ego sum panis vivus qui de celo descendi, siquis manducaverit ex hoc pane vivet in eternum, et panis quem ego dabo caro mea est pro mundi vita alleluya alleluya.
\newline {\color{Red} Feria Quinta.} \ul{Versiculus et Responsoria de Secunda Feria.}
\newline {\color{Red} Lectio Prima. Secundum Lucam.}
\lettrine[lines=2]{\zallmancaps \color{Red} I}{n} illo tempore: Convocatis Ihesus duodecim apostolis dedit illis virtutem et potestatem super omnia demonia, et ut languores curarent. Et reliqua.
\newline {\color{Red} Omelia Beati Ambrosii Episcopi.}
\lettrine[lines=2]{\zallmancaps \color{Blue} Q}{ualis} debeat esse qui evangelizat regnum Dei precipitur, preceptis evangelicis designatur, ut sine virga, sine pera, sine calciamento, sine pane, sine pecunia, hoc est subsidii secularis adminicula non requirens, fideique tutus putet sibi quo minus ea requirat magis posse suppetere.
\newline {\color{Red} Lectio Secunda.}
\lettrine[lines=2]{\zallmancaps \color{Red} Q}{ue} possunt qui volunt ad eum derivare tractatum: ut spiritalem tantummodo locus iste formare videatur affectum: qui velut indumentum quoddam videatur corporis exuisse non solum potestate reiecta contemptisque divitiis, sed eciam carnis ipsius illecebris abdicatis.%typo abddicatis
\newline {\color{Red} Lectio Tertia.}
\lettrine[lines=2]{\zallmancaps \color{Blue} H}{is} primo omnium datur pacis atque constantie generale mandatum, ut pacem ferant, constantiam servent, hospitalis necessitudinis iura custodiant, alienum a predicatore regni celestis astruens cursitare per domos, et inviolabilis hospitii iura mutare.
\newline {\color{Red} Ad Ben. Ant.} Convocatis Ihesus duodecim apostolis dedit illis virtutem, et potestatem super omnia demonia, et ut languores curarent et misit illos predicare regnum dei et sanare infirmos alleluya alleluya. {\color{Red} Oratio.}
\lettrine[lines=2]{\zallmancaps \color{Red} P}{resta} quesumus omnipotens et misericors deus, ut spiritus sanctus adveniens templum nos glorie sue dignanter inhabitando perficiat. Per, eiusdem.
\newline {\color{Red} Ad Mag. Ant.} Egressi duodecim apostoli circuibant per castella evangelizantes et curantes ubique, alleluya alleluya.
\newline {\color{Red} Feria Sexta.} \ul{Versiculus et Responsoria de Feria Tercia.}
\newline {\color{Red} Lectio Prima. Secundum Lucam.}
\lettrine[lines=2]{\zallmancaps \color{Blue} I}{n} illo tempore: Factum est in una dierum, et Ihesus sedebat docens. Et erant Pharisei sedentes et legis doctores qui venerant ex omni castello Galilee et Iudee et Hierusalem. Et reliqua.
\newline {\color{Red} Omelia Beati Ambrosii Episcopi.}
\lettrine[lines=2]{\zallmancaps \color{Red} N}{on} ociosa huius Paracliti nec angusta medicina est, quando dominus et orasse premittitur, non utique propter suffragium sed propter exemplum. Imitandi enim specimen dedit, non impetrandi ambitum requisivit.
\newline {\color{Red} Lectio Secunda.}
\lettrine[lines=2]{\zallmancaps \color{Blue} E}{t} convenientibus ex omni Galilea et Iudea et Hierusalem legis doctoribus, inter ceterorum remedia debilium, paralitici istius medicina describitur. Primum omnium quod ante diximus unusquisque eger petende precatores salutis debet adhibere, per quos vite nostre compago resoluta, actuumque nostrorum clauda vestigia verbi celestis remedio reformentur.
\newline {\color{Red} Lectio Tertia.}
\lettrine[lines=2]{\zallmancaps \color{Red} S}{int} igitur aliqui monitores mentis qui animum hominis quamvis exterioris corporis debilitate torpentem ad superiora erigant. quorum rursus adminiculis et attollere et humiliare se facilis, ante Ihesum locetur, dominico videri dignus aspectu. Humilitatem enim respicit dominus, quia respexit humilitatem ancille sue.
\newline {\color{Red} Ad Ben. Ant.} Factum est in una dierum, et Ihesus sedebat docens, et erant Pharisei sedentes et legis doctores qui venerant ex omni castello Galilee et Iudee et Hierusalem, et virtus erat domini ad sanandum eos, alleluya alleluya. {\color{Red} Oratio.}
\lettrine[lines=2]{\zallmancaps \color{Blue} D}{a} quesumus ecclesie tue misericors deus, ut spiritu sancto congregata hostili nullatenus incursione turbetur. Per, eiusdem.
\newline {\color{Red} Ad Mag. Ant.} Tulit ergo paraliticus lectum suum in quo iacebat, magnificans deum, et omnis plebs ut vidit dedit laudem deo alleluya.
\newline {\color{Red} Sabbato.} \ul{Versiculus et Responsoria de Quarta Feria.}
\newline {\color{Red} Lectio Prima. Secundum Lucam.}
\lettrine[lines=2]{\zallmancaps \color{Red} I}{n} illo tempore: Surgens Ihesus de sinagoga: introivit in domum Simonis. Socrus autem Simonis tenebatur magnis febribus. Et reliqua.
\newline {\color{Red} Omelia Venerabilis Bede Presbiteri.}
\lettrine[lines=2]{\zallmancaps \color{Blue} D}{omus} Petri quod dicitur, vetus lex intelligitur, Socrus autem Petri sinagoga, quia stimulis invidie agitabatur adversus ecclesiam gentium, ne locum eius preciperet credendo in Christum, et voluntati eius famulando.
\newline {\color{Red} Lectio Secunda.}
\lettrine[lines=2]{\zallmancaps \color{Red} S}{ive} aliter potest intelligi. Si virum a demonio liberatum moraliter animum ab immunda cogitatione purgatum significare dixerimus, consequenter femina febribus tenta sed ad imperium domini curata, carnem ostendit a concupiscentie sue fevore per continentie precepta frenatam.
\newline {\color{Red} Lectio Tertia.}
\lettrine[lines=2]{\zallmancaps \color{Blue} O}{mnis} enim amaritudo et ira et indignatio et clamor et blasfemia: spiritus immundi furor est. Fornicationem vero immunditiam, libidinem, concupiscentiam malam, et avaritiam que est simulacrorum servitus, hec omnia et his similia, febrem humane nature possumus accipere. A quibus quamdiu anime fidelium tenentur, egritudine tanta oppresse sunt, ut nequaquam digne deo quid boni operis valeant exhibere.
\newline {\color{Red} Ad Ben. Ant.} Vespere autem facto cum sol occidisset veniebant ad Ihesum omnes male habentes, at ille singulis manus imponens curabat eos alleluya. {\color{Red} Oratio.}
\lettrine[lines=2]{\zallmancaps \color{Red} M}{entibus} nostris quesumus domine spiritum sanctum benignus infunde, cuius et sapientia conditi sumus et providentia gubernamus. Per, eiusdem.
{\ubsubsection{In Festo Sancte Trinitatis}{in-festo-sancte-trinitatis-breviarium}}
\lettrine[lines=2]{\zallmancaps \color{Blue} A}{\color{Red} d} {\color{Red} Vesperas super Psalmos. Ant.} O beata et benedicta et gloriosa trinitas pater, et filius et spiritus sanctus. {\color{Red} Ps.} \psalm{112} \ul{Cum ceteris.} {\color{Red} Capitulum.}
\lettrine[lines=2]{\zallmancaps \color{Red} O}{}\ altitudo divitiarum sapientie et scientie dei quam incomprehensibilia sunt iudicia eius et investigabiles vie eius. \R \hyperlink{honor-virtus}{Honor.} {\color{Red} VI. Ymnus.}
\hymn{adesto-sancta-trinitas}{Adesto sancta trinitas}{
\lettrine[lines=2]{\zallmancaps \color{Blue} A}{desto} sancta trinitas,
par splendor una deitas,
que extas rerum omnium
sine fine principium.
\par {\bfseries \color{Red} T}e celorum militia
laudat, adorat, predicat,
triplexque mundi machina
benedicit per secula.
\par {\bfseries \color{Blue} A}ssumens et nos cernui
te adorantes famuli
vota precesque supplicum
ymnus iunge celestium.
\par {\bfseries \color{Red} U}num te lumen credimus,
quod et ter idem colimus,
alpha et \omega , quem dicimus,
te laudat omnis spiritus.
\par {\bfseries \color{Blue} L}aus patri sit ingenito,
laus eius unigenito,
laus sit sancto spiritui,
trino deo et simplici. Amen.
}
\newline \V Benedicamus patrem et filium cum sancto spiritu. \R Laudemus et superexaltemus eum in secula.
\newline {\color{Red} Ad Mag. Ant.} Gratias tibi deus, gratias tibi vera una trinitas, una et summa deitas, sancta et una unitas. {\color{Red} Oratio.}
\lettrine[lines=2]{\zallmancaps \color{Red} O}{mnipotens} \hypertarget{omnipotens-trinitas}{\label{omnipotens-trinitas}} sempiterne deus qui dedisti famulis tuis in confessione vere fidei eterne trinitatis gloriam agnoscere, et in potentia maiestatis adorare unitatem, quesumus ut eiusdem fidei firmitate ab omnibus semper muniamur adversis. Per.
\newline {\color{Red} Ad Completorium. Ant.} Miserere. {\color{Red} Ps.} \psalm{4} \ul{cum ceteris.} {\color{Red} Ad Nunc Dimittis. Ant.} Salva nos.
\newline {\color{Red} Ad Mat. Invit.}
\begin{invitatory}[deum-verum-invitatorium]
{Deum verum unum in trinitate et trinitatem in unitate, * Venite adoremus.}
\end{invitatory}
\newline {\color{Red} Ymnus.} \hyperlink{adesto-sancta-trinitas}{Adesto.}
\newline {\color{Red} In Primo Nocturno. Ant.} Adesto deus unus omnipotens pater et filius et spiritus sanctus. {\color{Red} Ps.} \psalm{8} {\color{Red} Ant.} Te unum in substantia, trinitatem in personis, confitemur. {\color{Red} Ps.} \psalm{18} {\color{Red} Ant.} Te semper idem esse vivere et intelligere profitemur. {\color{Red} Ps.} \psalm{23} \V Benedicamus patrem.
\newline \ul{Augustinus in sermone de sancta trinitate.} De unitate patris. \ul{et cetera.} \grecross
\newline {\color{Red} Lectio Prima.}
\lettrine[lines=2]{\zallmancaps \color{Blue} D}{iximus} alibi ea dici proprie in illa trinitate distincte ad singulas personas pertinentia, que relative dicuntur ad invicem, sicut pater et filius, et utriusque donum spiritus sanctus. Non enim pater trinitas, aut filius trinitas, aut donum trinitas. Quod vero ad se dicuntur singuli, non dici pluraliter tres, sed unum ipsam trinitatem, sicut deus pater, deus filius, deus spiritus sanctus, et bonus pater, bonus filius, bonus spiritus sanctus, et omnipotens pater, omnipotens filius, omnipotens spiritus sanctus. \R
\lettrine[lines=2]{\zallmancaps \color{Red} B}{enedicat} \hypertarget{benedicat-nos-deus}{\label{benedicat-nos-deus}} nos deus deus noster, benedicat nos deus, * Et metuant eum omnes fines terre. \V Deus misereatur nostri et benedicat nos deus. Et.
\newline {\color{Red} Lectio Secunda.}
\lettrine[lines=2]{\zallmancaps \color{Blue} N}{on} autem dicimus tres dii, aut tres boni, aut tres omnipotentes, sed unus deus bonus, omnipotens ipsa trinitas, et quicquid aliud non ad invicem relative, sed ad se singuli dicuntur. Hoc enim secundum essentiam dicuntur, quia hoc est ibi esse quod magnum esse, quod bonum, quod sapientem esse, et quicquid aliud ad se ibi unaqueque persona vel ipsa trinitas dicitur.
\begin{responsory}[benedictus-dominus-deus]
{Benedictus dominus deus Israhel qui facit mirabilia solus, * Et benedictum nomen maiestatis eius in eternum.}{Replebitur maiestate eius omnis terra fiat fiat.}
\end{responsory}%
\newline {\color{Red} Lectio Tertia.}
\lettrine[lines=2]{\zallmancaps \color{Red} I}{deo} autem dicimus tres personas vel tres substantias, non ut aliqua intelligatur diversitas essentie, sed ut vel uno aliquo vocabulo responderi possit cum dicitur quid tres vel quid tria. Tantamque esse equalitatem in ea trinitate ut non solum pater non sit maior, quam filius, quod attinet ad divinitatem, sed nec pater et filius simul maius aliquid sint, quam spiritus sanctus, aut singula queque persona quelibet trium minus aliquid quam ipsa trinitas.
\begin{responsory-final}[quis-deus-magnus]
{Quis deus magnus sicut deus noster, * Tu es deus, * Qui facis mirabilia.}{Notam fecisti in populis virtutem tuam redemisti in brachio tuo populum tuum.}{Qui}
\end{responsory-final}%
\newline {\color{Red} In Secundo Nocturno. Ant.} Te invocamus, te adoramus, te laudamus o beata trinitas. {\color{Red} Ps.} \psalm{46} {\color{Red} Ant.} Spes nostra, salus nostra, honor noster, o beata trinitas. {\color{Red} Ps.} \psalm{47} {\color{Red} Ant.} Libera nos, salva nos, iustifica nos, o beata trinitas. {\color{Red} Ps.} \psalm{71} \V Benedictus es domine in firmamento celi. \R Et laudabilis et gloriosus in secula.
\newline {\color{Red} Lectio Quarta.}
\lettrine[lines=2]{\zallmancaps \color{Blue} V}{erbum} ergo dei patris, unigenitus filius per omnia patri similis et equalis, deus de deo, lumen de lumine, sapientia de sapientia, essentia de essentia, hoc omnino est quod pater, non tamen pater, quia iste filius ille pater, ac per hoc novit omnia que novit pater, sed ei nosse de patre est sicut esse. Nosse enim et esse ibi unum est.
\begin{responsory}[magnus-dominus-et]
{Magnus dominus et magna virtus eius, * Et sapientie eius non est numerus.}{Magnus dominus et laudabilis nimis et magnitudinis eius non est finis.}
\end{responsory}%
\newline {\color{Red} Lectio Quinta.}
\lettrine[lines=2]{\zallmancaps \color{Red} S}{piritus} autem sanctus secundum scripturas sanctas, nec patris est solius nec filii solius, sed amborum: et ideo communem qua invicem se diligunt pater et filius, nobis insinuat caritatem. Verum cum nos exercet sermo divinus non dixit scriptura spiritus sanctus caritas est, sed dixit Deus caritas est, ut incertum sit, et ideo requirendum utrum deus pater sit caritas, an deus filius, an deus spiritus sanctus, an deus ipsa trinitas.
\begin{responsory}[gloria-patri-geniteque]
{Gloria patri geniteque proli et tibi compar utriusque semper spiritus alme deus unus, * Omni tempore secli.}{Da gaudiorum premia, da gratiarum munera, dissolve litis vincula, astringe pacis federa.}
\end{responsory}%
\newline {\color{Red} Lectio Sexta.}
\lettrine[lines=2]{\zallmancaps \color{Blue} N}{escio} autem cur non sicut sapientia et pater dicitur et filius et spiritus sanctus, et simul omnes non tres sed una sapientia: ita et pater caritas et filius caritas et spiritus sanctus caritas, et simul omnes una caritas dicantur. Non frustra autem in hac trinitate non dicitur verbum dei nisi filius, nec donum dei nisi spiritus sanctus, nec de quo genitum est verbum: nisi deus pater.
\begin{responsory-doxology}[honor-virtus]
{Honor virtus et potestas et imperium sit trinitati in unitate, unitati in trinitate, * In perenni seculorum tempore.}{Trinitati lux perennis unitati sit decus perpetim.}
\end{responsory-doxology}%
\newline {\color{Red} In Tertio Nocturno. Ant.} Caritas pater est, gratia Christus, communicacio spiritus, o beata trinitas. {\color{Red} Ps.} \psalm{95} {\color{Red} Ant.} Verax est pater, veritas filius, veritatis spiritus, o beata trinitas. {\color{Red} Ps.} \psalm{96} {\color{Red} Ant.} Una igitur pater logos Paraclytusque substantia est, o beata trinitas. {\color{Red} Ps.} \psalm{97} \V Verbo domini celi firmati sunt. \R Et spiritu oris eius omnis virtus eorum.
\newline {\color{Red} Lectio VII.}
\lettrine[lines=2]{\zallmancaps \color{Red} H}{inc} ergo factum est ut proprie dei verbum, eciam dei sapientia diceretur, cum sit sapientia et pater et spiritus sanctus. Si ergo proprie aliquid horum trium caritas nuncupanda est, quid aptius quam ut hoc sit spiritus sanctus, ut scilicet in illa simplici summaque natura, non sit aliud substantia et aliud caritas, sed substantia ipsa sit caritas, et caritas ipsa substantia, sive in patre, sive in filio, sive in spiritu sancto, et tamen proprie spiritus sanctus caritas nuncupetur?
\begin{responsory}[tibi-laus-tibi]
{Tibi laus tibi gloria tibi gratiarum actio, * In secula sempiterna o beata trinitas.}{Et benedictum nomen glorie tue sanctum et laudabile et superexaltatum.}
\end{responsory}%
\newline {\color{Red} Lectio VIII.}
\lettrine[lines=2]{\zallmancaps \color{Blue} C}{redamus} igitur patrem et filium et spiritum sanctum esse unum deum, universe creature conditorem atque rectorem, nec patrem esse filium, nec filium patrem esse vel spiritum sanctum, nec spiritum sanctum, vel patrem esse vel filium, sed trinitatem relatarum ad invicem personarum, et unitatem equalis essentie. Queramus autem hoc intelligere, ab eo ipso quem intelligere volumus auxilium precantes.
\begin{responsory}[benedicamus-patrem]
{Benedicamus patrem et filium cum sancto spiritu, laudemus et superexaltemus eum, * In secula.}{Benedictus es domine in firmamento celi et laudabilis et gloriosus.}
\end{responsory}%
\newline {\color{Red} Lectio IX.}
\lettrine[lines=2]{\zallmancaps \color{Red} Q}{uantum} autem nobis tribuitur, quod intelligimus explicare, tanta cura et sollicitudine pietatis cupiamus ut eciam si aliquid aliud pro alio dicimus, nichil tamen dicamus indignum. Ut siquid verbi gratia de patre dicimus quod patri proprie non conveniat, aut filio conveniat aut spiritui sancto, aut ipsi trinitati. Et siquid de filio quod filio proprie non conveniat, patri saltem congruat, aut spiritui sancto, aut trinitati. Item siquid de spiritu sancto quod proprietatem spiritus sancti non doceat, non tamen alienum sit a patre et a filio aut ab uno deo ipsa trinitate.
\begin{responsory-doxology}[summe-trinitati]
{Summe trinitati simplici deo, una divinitas, equalis gloria, coeterna maiestas, patri prolique sanctoque flamini, * Qui totum subdit suis orbem legibus.}{Prestet nobis gratiam deitas beata patris ac nati pariterque spiritus almi.}
\end{responsory-doxology}%
\newline \V Magnus dominus et laudabilis nimis. \R Et magnitudinis eius non est finis.
\newline {\color{Red} In Laudibus. Ant.} Gloria tibi trinitas equalis una deitas, et ante omnia secula, et nunc et in perpetuum. {\color{Red} Ps.} \psalm{92} {\color{Red} Ant.} Laus et perennis gloria deo patri et filio, sancto simul Paraclyto in secula seculorum. {\color{Red} Ant.} Gloria laudis resonet in ore omnium patri, geniteque proli, spiritui sancto pariter resultet laude perenni. {\color{Red} Ant.} Laus deo patri, parilique proli, et tibi semper studio perenni spiritus, nostro resonet ab ore omne per evum. {\color{Red} Ant.} Ex quo omnia, per quem omnia, in quo omnia, ipsi gloria in secula. {\color{Red} Cap.} O altitudo.
\hymn{o-trinitas-laudabilis}{O trinitas laudabilis}{
\lettrine[lines=2]{\zallmancaps \color{Blue} O}{}\ trinitas laudabilis,
et unitas mirabilis,
in simplici substantia
virtus manens intermina.
\par {\bfseries \color{Blue} T}u caritas, tu puritas,
tu pax et immortalitas,
patris, nati, Paraclyti,
decore pollens perpeti.
\par {\bfseries \color{Red} F}ides, corona supplicum
in te pie fidentium,
exterge sordes mentium,
sorti misertus pauperum.
\par {\bfseries \color{Blue} L}aus patri sit ingenito,
laus eius unigenito,
laus sit sancto spiritui,
trino deo et simplici. Amen.
}
%pp111
\newline \V Sit nomen domini benedictum. \R Ex hoc nunc et usque in seculum.
\newline {\color{Red} Ad Ben. Ant.} Benedicta sit creatrix et gubernatrix omnium sancta et individua trinitas, et nunc et semper et per infinita seculorum secula.
\newline \ul{Oratio.} \hyperlink{omnipotens-trinitas}{Omnipotens sempiterne etc.} \ul{Ad horas antiphone Laudum.}
\newline {\color{Red} Ad Terciam. Cap.} O altitudo.
\begin{responsory-breve}[benedicamus-patrem-breve]
{Benedicamus patrem et filium cum sancto spiritu, * Alleluya alleluya.}{Laudemus et superexaltemus eum in secula.}
\end{responsory-breve}%
\V Benedictus es domine.
\newline {\color{Red} Ad Sextam. Capitulum.}
\lettrine[lines=2]{\zallmancaps \color{Red} G}{ratia} domini nostri Ihesu Christi, et caritas dei, et communicatio sancti spiritus, sit semper cum omnibus nobis. Amen.
\begin{responsory-breve}[benedictus-es-domine]
{Benedictus es domine in firmamento celi, * Alleluya alleluya.}{Et laudabilis et gloriosus in secula.}
\end{responsory-breve}%
\V Verbo domini.
\newline {\color{Red} Ad Nonam. Capit.}
\lettrine[lines=2]{\zallmancaps \color{Blue} T}{res} sunt qui testimonium dant in celo, pater verbum et spiritus sanctus, et hi tres unum sunt.
\begin{responsory-breve}[verbo-domini-celi]
{Verbo domini celi firmati sunt, * Alleluya alleluya.}{Et spiritu oris eius omnis virtus eorum.}
\end{responsory-breve}%
\V Sit nomen domini.
\newline {\color{Red} Ad Vesperas. Super psalmos. Ant.} Gloria tibi trinitas. {\color{Red} Ps.} \psalm{109} {\color{Red} Ps.} \psalm{110} {\color{Red} Ps.} \psalm{111} {\color{Red} Ps.} \psalm{112} {\color{Red} Ps.} \psalm{116}
\newline \ul{Cap., Ymnus, \Vbar , ut in Primis Vesperis.}
\newline {\color{Red} Ad Mag. Ant.} Te deum patrem ingenitum, te filium unigenitum, te spiritum sanctum Paraclitum, sanctam et individuam trinitatem, toto corde et ore confitemur, laudamus atque benedicimus tibi gloria in secula.
\newline \ul{Per octavas quando de octava agitur.} {\color{Red} Invitat.}
\begin{invitatory}[deum-verum-octavam-invitatorium]
{Deum verum unum in trinitate, * Venite adoremus.}
\end{invitatory}
\newline \ul{Ad Ben. vel ad Mem. in Laudibus Ant.} Benedicta sit creatrix.
\newline \ul{Ad Mag. vel ad Mem. in Vesperis Ant.} Te deum patrem.
\newline \ul{Cetera per octavam fiant secundum ordinem qui infra in rubrica de octavis sanctorum notatus est, usque ad Vesperas sequentis Sabbati exclusive, et repetantur lectiones supraposite.}
{\ubsubsection{Dominica Prima post Festum Trinitatis}{dominica-prima-post-festum-trinitatis-breviarium}}
\newline \ul{Et deinceps usque ad Primam Dominicam Augusti exclusive quando de tempore agitur, cantetur hystoria,} \hyperlink{deus-omnium}{Deus omnium,} \ul{et legantur lectiones de Libris Regum. Sabbato vero precedenti ad Vesperas antiphone, Psalmi, Cap., Ymnus, \Vbar , sicut in Sabbato post Octavas Epyphanie.}
\newline \R \hyperlink{dominus-qui-eripuit}{Dominus qui eripuit.} {\color{Red} II. Ad Mag. Ant.} Loquere domine quia audit servus tuus.
\newline \ul{Orationes Dominicales usque ad Adventum exclusive sicut infra in ordine omeliarum notate sunt.}
\newline {\color{Red} Ad Matutinas. Invitat.} \hyperlink{venite-exultemus-domino-invitatorium}{Venite exultemus domino, * Iubilemus deo salutari nostro.}
\newline \ul{Hoc Invitatorium dicatur usque ad Adventum Domini exclusive in Dominicis quando de Dominica agitur.}
\newline {\color{Red} In Primo Nocturno. Ant.} Pro fidei meritis vocitatur iure beatus, legem qui domini meditatur nocte dieque. {\color{Red} Ps.} \psalm{1} \psalm{2} \psalm{3} \psalm{6} Gloria. {\color{Red} Ant.} Iuste deus iudex fortis patiensque benignus. In te sperantes muni miserando fideles. {\color{Red} Ps.} \psalm{7} \psalm{8} \psalm{9} \psalm{10} Gloria. {\color{Red} Ant.} Surge et in eternum serva munimine sacro, custodique tuos astripotens famulos. {\color{Red} Ps.} \psalm{11} \psalm{12} \psalm{13} \psalm{14} Gloria. \V Memor fui.
\newline {\color{Red} Lectio Prima.}
\lettrine[lines=2]{\zallmancaps \color{Red} F}{uit} vir unus de Ramathaimsophim de monte Ephraym, et nomen eius Helchana filius Ieroam, filii Heliu, filii Thou, filii Suph, Effrateus, et habuit duas uxores. Nomen uni Anna, et nomen secunde Phenenna. Fueruntque Phenenne filii: Anne autem non erant liberi. \R
\lettrine[lines=2]{\zallmancaps \color{Blue} D}{eus} \hypertarget{deus-omnium}{\label{deus-omnium}} omnium exauditor est, ipse misit angelum suum, et tulit me de ovibus patris mei, * Et unxit me unctione misericordie sue. \V Dominus qui eripuit me de ore leonis et de manu bestie liberavit me. Et unxit.
\newline {\color{Red} Lectio Secunda.}
\lettrine[lines=2]{\zallmancaps \color{Red} E}{t} ascendebat vir ille de civitate sua statutis diebus: ut adoraret et sacrificaret domino exercituum in Silo. Erant autem ibi duo filii Hely, Ophni et Phinees sacerdotes domini. Venit ergo dies et immolavit Helchana, deditque Phenenne uxori sue, et cunctis filiis eius et filiabus partes.
\begin{responsory}[dominus-qui-eripuit]
{Dominus qui eripuit me de ore leonis et de manu bestie liberavit me, * Ipse me eripiet de manibus inimicorum meorum.}{Misit deus misericordiam suam et veritatem suam animam meam eripuit de medio catulorum leonum.}
\end{responsory}%
\newline {\color{Red} Lectio Tertia.}
\lettrine[lines=2]{\zallmancaps \color{Blue} A}{nne} autem dedit partem unam tristis, quia Annam diligebat. Dominus autem concluserat vulvam eius. Affligebat quoque eam emula eius, et vehementer angebat, in tantum ut exprobraret quod conclusisset dominus vulvam eius. Sicque faciebat per singulos annos cum redeunte tempore ascenderent in domum domini, et sic provocabat eam.
\begin{responsory-final}[ego-te-tuli]
{Ego te tuli de domo patris tui dicit dominus, et posui te pascere gregem populi mei, et fui tecum in omnibus ubicumque ambulasti, * Firmans regnum tuum, * In eternum.}{Fecique tibi nomen grande iuxta nomen magnorum qui sunt in terra, et requiem dedi tibi ab omnibus inimicis tuis.}{In}
\end{responsory-final}%
\newline {\color{Red} In Secundo Nocturno. Ant.} Nature genitor conserva a morte redemptos. Facque tuo dignos servitio famulos. {\color{Red} Ps.} \psalm{15} {\color{Red} Ant.} Pectora nostra tibi tu conditor orbis adure, igne pio purgans atque cremando probans. {\color{Red} Ps.} \psalm{16} {\color{Red} Ant.} Tu populum humilem salvasti a morte redemptor atque superba tuo colla premis iaculo. {\color{Red} Ps.} \psalm{17} \V Media nocte.
\newline {\color{Red} Lectio Quarta.}
\lettrine[lines=2]{\zallmancaps \color{Red} P}{orro} Anna flebat, et non capiebat cibum. Dixit ergo ei Helchana vir suus. Anna cur fles, et quare non comedis, et quam ob rem affligitur cor tuum? Nunquid non ego melior sum tibi quam decem filii?
\begin{responsory}[preparate-corda-vestra]
{Preparate corda vestra domino et servite illi soli, * Et liberabit vos de manibus inimicorum vestrorum.}{Auferte deos alienos de medio vestri.}
\end{responsory}%
\newline {\color{Red} Lectio Quinta.}
\lettrine[lines=2]{\zallmancaps \color{Blue} S}{urrexit} autem Anna postquam comederat in Silo et biberat. Et Hely sacerdote sedente super sellam ante postes domus domini, cum esset Anna amaro animo, oravit dominum flens largiter, et votum vovit dicens. Domine exercituum, si respiciens videris afflictionem famule tue, et recordatus mei fueris, nec oblitus ancille tue, dederisque serve tue sexum virilem, dabo eum domino omnibus diebus vite eius, et novacula non ascendet super caput eius.
\begin{responsory}[percussit-saul]
{Percussit Saul mille et David decem milia, quia manus domini erat cum illo, percussit Philisteum, * Et abstulit obprobrium ex Israhel.}{Nonne iste est David de quo canebant in choro dicentes Saul percussit mille et David decem milia in milibus suis.}
\end{responsory}%
\newline {\color{Red} Lectio Sexta.}
\lettrine[lines=2]{\zallmancaps \color{Red} F}{actum} est autem cum illa multiplicaret preces coram domino, ut Hely observaret os eius. Porro Anna loquebatur in corde suo, tantumque labia illius movebantur, et vox penitus non audiebatur. Estimavit ergo eam Hely temulentam. Dixitque ei. Usquequo ebria eris? Digere paulisper vinum quo mades.
\begin{responsory-final}[montes-gelboe]
{Montes Gelboe nec ros nec pluvia veniant super vos, * Ubi ceciderunt * Fortes Israhel?}{Omnes montes in circuitu eius visitet dominus a Gelboe autem transeat.}{Fortes}
\end{responsory-final}%
\newline {\color{Red} In Tertio Nocturno. Ant.} Sponsus ut e thalamo processit Christus in orbem descendens celo iure salutifero. {\color{Red} Ps.} \psalm{18} {\color{Red} Ant.} Auxilium nobis salvator mitte salutis, et tribuas vitam tempore perpetuo. {\color{Red} Ps.} \psalm{19} {\color{Red} Ant.} Rex sine fine manens miseris tu parce ruinis. Premia concedens et tua cuncta regens. {\color{Red} Ps.} \psalm{20} \V Exaltare domine.
\newline \ul{Predicte antiphone, psalmi, et versiculi nocturnorum dicantur omnibus Dominicis quando de Dominica agitur usque ad Adventum Domini exclusive.}
\begin{responsory}[audi-domine]
{Audi domine ymnum et orationem quam servus tuus orat coram te hodie, ut sint oculi tui aperti, et aures tue intente, * Super domum istam die ac nocte.}{Respice domine de sanctuario tuo et de excelso celorum habitaculo.}
\end{responsory}%
\begin{responsory}[peccavi-super]
{Peccavi super numerum harene maris et multiplicata sunt peccata mea, et non sum dignus videre altitudinem celi pre multitudine iniquitatis mee, quoniam irritavi iram tuam, * Et malum coram te feci.}{Quoniam iniquitatam meam ego agnosco et delictum meum coram me est semper, tibi soli peccavi.}
\end{responsory}%
\begin{responsory-doxology}[factum-est-dum]
{Factum est dum tolleret dominus Helyam per turbinem in celum, Helyseus clamabat dicens, * Pater mi, pater mi, currus Israhel et auriga eius.}{Oro domine ut fiat spiritus tuus duplex in me, at ille inquit, si videris quando tollar a te erit quod petisti.}
\end{responsory-doxology}%
\newline \ul{Inqualibet Dominica a} \hyperlink{deus-omnium}{Deus omnium.} \ul{usque ad Adventum exclusive quando de Dominica agitur: Responsorium,} \hyperlink{honor-virtus}{Honor virtus.} \ul{loco noni responsorii dicatur: postquam Dominicales hystoria sive in feriis sive in Dominica fuerit percantanta.}
\newline \ul{In Laudibus antiphona, Psalmi, Cap., Ymnus, \Vbar , sicut in Dominica} \hyperlink{domine-ne-in-ira}{Domine ne in ira.} \ul{Antiphone vero ad Ben., et ad Mag., secundum quod infra in serie evangeliorum cum orationibus sunt notate dicantur. Ad Primam et ad alias horas antiphone, Cap., Responsoria, \Vbar , dicantur sicut sunt in hystoria} \hyperlink{domine-ne-in-ira}{Domine ne in ira.} \ul{Et hic ordo In Laudibus et in Horis servetur omnibus Dominicis quando de Dominica agitur usque ad Adventum Domini exclusive.}%typo noni for novi
\newline \ul{Sequentes antiphone sunt dicende diebus Sabbatorum suo ordine, ad Mag., vel ad memoriam usque ad Sabbatum ante primam Dominicam Augusti exclusive.}
\newline {\color{Red} Ant.} Cognoverunt omnes a Dan usque Bersabee, quod fidelis Samuel propheta esset domini. {\color{Red} Ant.} Prevaluit David in Philisteum in funda et lapide in nomine domini. {\color{Red} Ant.} Nonne iste est David de quo canebant in choro dicentes Saul percussit mille, et David decem milia in milibus suis. {\color{Red} Ant.} Iratus rex Saul dixit, michi mille dederunt et filio Isai dederunt decem milia. {\color{Red} Ant.} Quis enim in omnibus sicut David fidelis inventus est in regno tuo, egrediens et regrediens et pergens ad imperium regis. {\color{Red} Ant.} Doleo super te frater mi Ionatha, amabilis valde super amorem mulierum, sicut mater unicum amat filium ita te diligebam, sagitta Ionathe nunquam abiit retrorsum, nec declinavit clipeus eius in bello, et hasta eius non est aversa. {\color{Red} Ant.} Dixitque David ad dominum cum vidisset angelum cedentem populum, ego sum qui peccavi, ego inique egi, isti qui oves sunt quid fecerunt. {\color{Red} Ant.} Montes Gelboe nec ros nec pluvia veniant super vos, quia ibi abiectus est clipeus fortium, clipeus Saul quasi non esset unctus oleo, quomodo ceciderunt fortes in prelio: Ionathas in excelsis tuis interfectus est, Saul et Ionathas amabiles et decori valde in vita sua, in morte quoque non sunt separati. {\color{Red} Ant.} Rex autem David cooperto capite incedens lugebat filium dicens, Absalon fili me, fili me Absalon, quis michi det ut ego moriar pro te fili mi Absalon?
\newline \ul{Lectiones feriales a secundo capitulo usque ad septimum ad libitum.}
{\ubsubsection{Dominica Secunda}{dominica-secunda-post-festum-trinitatis-breviarium}}
\newline {\color{Red} Lectio Prima.}
\lettrine[lines=2]{\zallmancaps \color{Blue} V}{enerunt} viri Cariathiarim et reduxerunt archam domini, et intulerunt eam in domum Aminadab in Gabaa. Eleazarum autem filium eius sanctificaverunt, ut custodiret archam domini. Et factum est ex qua die mansit archa domini in Cariathiarim, multiplicati sunt dies. Erat quippe iam annus vigesimus. Et requievit omnis domus Israhel post dominum.%deleted ergo
\newline {\color{Red} Lectio Secunda.}
\lettrine[lines=2]{\zallmancaps \color{Red} A}{it} autem Samuel ad universam domum Israhel dicens. Si in toto corde vestro revertimini ad dominum: auferte deos alienos de medio vestrum et Astaroth, et preparate corda vestra domino, et servite illi soli, et eruet vos de manu Phylistiim. Abstulerunt ergo filii Israhel Baalim et Astaroth: et servierunt domino soli.
\newline {\color{Red} Lectio Tertia.}
\lettrine[lines=2]{\zallmancaps \color{Blue} D}{ixit} autem Samuel. Congregate universum Israhel in Masphat: ut orem pro vobis dominum. Et convenerunt in Masphat. Hauseruntque aquam et effuderunt in conspectu domini: et ieiunaverunt in die illa et dixerunt. Tibi peccavimus domino. Iudicavitque Samuel filios Israhel in Masphat.
\newline {\color{Red} Lectio Quarta.}
\lettrine[lines=2]{\zallmancaps \color{Red} E}{t} audierunt Philistiim quod congregati essent filii Israhel in Masphat, et ascenderunt satrape Philistinorum ad Israhel. Quod cum audissent filii Israhel: timuerunt a facie Philistinorum. Dixeruntque ad Samuelem. Ne cesses pro nobis clamare ad dominum Deum nostrum, ut salvet nos de manu Philistinorum.
\newline {\color{Red} Lectio Quinta.}
\lettrine[lines=2]{\zallmancaps \color{Blue} T}{ulit} autem Samuel agnum lactentem unum, et obtulit illum holocaustum integrum domino. Et clamavit Samuel ad dominum pro Israhel et exaudivit eum dominus. Factum est autem cum Samuel offerret holocaustum, Philisteos inire prelium contra Israhel. Intonuit autem dominus fragore magno in die illa super Philistiim, et exterruit eos, et cesi sunt a facie Israhel.
\newline {\color{Red} Lectio Sexta.}
\lettrine[lines=2]{\zallmancaps \color{Red} E}{gressique} viri Israhel de Masphat: persecuti sunt Philisteos, et percusserunt eos usque ad locum qui erat subter Bethchar. Tulit autem Samuel lapidem unum, et posuit eum inter Masphat et inter Sen, et vocavit nomen loci illius Lapis adiutorii. Dixitque. Huc usque auxiliatus est nobis Dominus. Et humiliati sunt Philistiim, nec apposuerunt ultra ut venirent in terminos Israhel.
\newline \ul{Lectiones feriales ab octavo capitulo usque ad decimum ad libitum.}
{\ubsubsection{Dominica Tercia}{dominica-tercia-post-festum-trinitatis-breviarium}}
\newline {\color{Red} Lectio Prima.}
\lettrine[lines=2]{\zallmancaps \color{Blue} L}{ocutus} est Samuel ad populum legem regni, et scripsit in libro, et reposuit coram domino. Et dimisit Samuel omnem populum, singulos in domum suam. Sed et Saul abiit in domum suam in Gabaath: et abiit cum eo pars exercitus quorum tetigerat Deus corda. Filii vero Belial dixerunt. Num salvare nos poterit iste? Et despexerunt eum et non attulerunt ei munera. Ille vero dissimulabat se audire.%deleted autem
\newline {\color{Red} Lectio Secunda.}
\lettrine[lines=2]{\zallmancaps \color{Red} A}{scendit} autem Naas Ammonites, et pugnare cepit adversum Iabes Galaad. Dixeruntque omnes viri Iabes ad Naas. Habeto nos federatos, et serviemus tibi. Et respondit ad eos Naas Ammonites. In hoc feriam vobiscum fedus ut eruam omnium vestrum oculos dextros, ponamque vos obprobrium in universo Israhel.
\newline {\color{Red} Lectio Tertia.}
\lettrine[lines=2]{\zallmancaps \color{Blue} E}{t} dixerunt ad eum seniores Iabes. Concede nobis septem dies, ut mittamus nuncios in universos terminos Israhel. Et si non fuerit qui defendat nos: egrediemur ad te. Venerunt ergo nuncii in Gabaath Saulis, et locuti sunt verba hec audiente populo. Et levavit omnis populus vocem suam: et flevit.
\newline {\color{Red} Lectio Quarta.}
\lettrine[lines=2]{\zallmancaps \color{Red} E}{t} ecce Saul veniebat sequens boves de agro, et ait. Quid habet populus quod plorat? Et narraverunt ei verba virorum Iabes. Et insilivit spiritus domini in Saul cum audisset verba hec: et iratus est furor eius nimis. Et assumens utrumque bovem et concidit in frusta, misitque in omnes terminos Israhel per manum nunciorum, dicens. Quicumque non exierit, secutusque fuerit Saul et Samuel: sic fiet bobus eius.
\newline {\color{Red} Lectio Quinta.}
\lettrine[lines=2]{\zallmancaps \color{Blue} I}{nvasit} ergo timor domini populum, et egressi sunt quasi vir unus: et recensuit eos in Bezech. Fueruntque filiorum Israhel trecenta milia: virorum autem Iuda triginta milia. Et dixerunt nunciis qui venerant. Sic dicetis viris qui sunt in Iabes Galaad. Cras erit vobis salus cum incaluerit sol.
\newline {\color{Red} Lectio Sexta.}
\lettrine[lines=2]{\zallmancaps \color{Red} V}{enerunt} ergo nuncii, et nunciaverunt viris Iabes. Qui letati sunt, et dixerunt. Mane exibimus ad vos, et facietis nobis omne quod placuerit vobis. Et factum est cum venisset dies crastinus: constituit Saul populum in tres partes. Et ingressus est media castra in vigilia matutina, et percussit Ammon usque dum incalesceret dies. Reliqui autem dispersi sunt, ita ut non relinquerentur in eis duo pariter.
\newline \ul{Lectiones feriales a sexto decimo capitulo usque ad vicesimum octavum ad libitum.}
{\ubsubsection{Dominica Quarta}{dominica-quarta-post-festum-trinitatis-breviarium}}
\newline {\color{Red} Lectio Prima.}
\lettrine[lines=2]{\zallmancaps \color{Blue} I}{n} diebus illis congregaverunt Philistiim agmina sua, ut prepararentur contra Israhel ad bellum. Dixitque Achis ad David. Sciens nunc scito, quoniam mecum egredieris in castris tu, et viri tui. Dixitque David ad Achis. Nunc scies que facturus est servus tuus. Et ait Achis ad David. Et ego custodem capitis mei ponam te cunctis diebus.
\newline {\color{Red} Lectio Secunda.}
\lettrine[lines=2]{\zallmancaps \color{Red} S}{amuel} autem mortuus est, planxitque eum omnis Israhel, et sepelierunt eum in Ramatha urbe sua. Et Saul abstulit magos et ariolos de terra. Congregatique sunt Philistiim: et venerunt et castrametati sunt in Sunam.
\newline {\color{Red} Lectio Tertia.}
\lettrine[lines=2]{\zallmancaps \color{Blue} C}{ongregavit} autem et Saul universum Israhel, et venit in Gelboe. Et vidit Saul castra Philistiim et timuit et expavit cor eius. Consuluitque dominum, et non respondit ei neque per somnia neque per sacerdotes, neque per prophetas. Dixitque Saul servis suis. Querite michi mulierem habentem phithonem, et vadam ad eam et sciscitabor per illam.
\newline {\color{Red} Lectio Quarta.}
\lettrine[lines=2]{\zallmancaps \color{Red} E}{t} dixerunt servi eius ad eum. Est mulier habens phithonem in Endor. Mutavit ergo habitum suum, vestitusque est aliis vestimentis. Et abiit ipse, et duo viri cum eo: veneruntque ad mulierem nocte. Et ait. Divina michi in phithone: et suscita michi quem dixero tibi.
\newline {\color{Red} Lectio Quinta.}
\lettrine[lines=2]{\zallmancaps \color{Blue} E}{t} ait mulier ad eum. Ecce tu nosti quanta fecerit Saul, et quomodo eraserit magos et ariolos de terra. Quare ergo insidiaris anime mee ut occidar? Et iuravit ei Saul in domino, dicens. Vivit dominus, quia non eveniet tibi quicquam mali propter hanc rem. Dixitque ei mulier. Quem suscitabo tibi? Qui ait. Samuelem suscita michi.
\newline {\color{Red} Lectio Sexta.}
\lettrine[lines=2]{\zallmancaps \color{Red} C}{um} autem vidisset mulier Samuelem, exclamavit voce magna et dixit ad Saul. Quare imposuisti michi? Tu es enim Saul. Dixitque ei rex. Noli timere. Quid vidisti? Et ait mulier ad Saul. Deos vidi ascendentes de terra. Dixitque ei. Qualis est forma eius? Que ait. Vir senex ascendit: et ipse amictus est pallio.
\newline \ul{Lectiones feriales a vicesimo nono capitulo usque ad secundum regum ad libitum.}
{\ubsubsection{Dominica Quinta}{dominica-quinta-post-festum-trinitatis-breviarium}}
\newline {\color{Red} Lectio Prima.}
\lettrine[lines=2]{\zallmancaps \color{Blue} F}{actum} est postquam mortuus est Saul, ut David reverteretur a cede Amalech: et maneret in Siceleg duos dies. In die autem tercia apparuit homo veniens de castris Saul veste conscissa: et pulvere aspersus caput. Et ut venit ad David: cecidit super faciem suam, et adoravit.%deleted autem
\newline {\color{Red} Lectio Secunda.}
\lettrine[lines=2]{\zallmancaps \color{Red} D}{ixitque} ad eum David. Unde venis? Qui ait ad eum. De castris Israhel fugi. Et dixit ad eum David. Quod est verbum quod factum est? Indica michi. Qui ait. Fugit populus e prelio, et multi corruentes ex populo mortui sunt. Sed et Saul et Ionathas filius eius interierunt. Dixitque David ad adolescentem qui nunciabat ei. Unde scis quia mortuus est Saul et Ionathas filius eius?
\newline {\color{Red} Lectio Tertia.}
\lettrine[lines=2]{\zallmancaps \color{Blue} E}{t} ait adolescens qui nunciabat ei. Casu veni in montem Gelboe: et Saul incumbebat super hastam suam. Porro currus et equites appropinquabant ei. Et conversus post tergum suum, vidensque me: vocavit. Cui cum respondissem, assum: dixit michi. Quisnam es tu? Et aio ad eum. Amalchites sum.
\newline {\color{Red} Lectio Quarta.}
\lettrine[lines=2]{\zallmancaps \color{Red} E}{t} locutus est michi. Sta super me, et interfice me: quoniam tenent me angustie, et adhuc tota anima mea in me est. Stansque super eum, occidi illum. Sciebam enim quod vivere non poterat post ruinam. Et tuli diadema quod erat in capite eius et armillam de brachio illius, et attuli ad te dominum meum huc.
\newline {\color{Red} Lectio Quinta.}
\lettrine[lines=2]{\zallmancaps \color{Blue} A}{pprehendens} autem David vestimenta sua scidit, omnesque viri qui erant cum eo: et planxerunt et fleverunt et ieiunaverunt usque ad vesperam super Saul et super Ionatham filium eius, et super populum domini, et super domum Israhel, quod corruissent
%pp112
gladio. Dixitque David ad iuvenem qui nunciaverat ei. Unde es tu? Qui respondit. Filius hominis advene Amalechite ego sum.
\newline {\color{Red} Lectio Sexta.}
\lettrine[lines=2]{\zallmancaps \color{Red} E}{t} ait ad eum David. Quare non timuisti mittere manum tuam, ut occideres christum domini? Vocansque David unum de pueris suis: ait. Accedens: irrue in eum. Qui percussit illum: et mortuus est. Et ait ad eum David. Sanguis tuus super caput tuum. Os enim tuum locutum est adversum te, dicens: ego interfeci christum domini.
\newline \ul{Lectiones feriales a secundo capitulo secundi regum usque ad undecimum ad libitum.}
{\ubsubsection{Dominica Sexta}{dominica-sexta-post-festum-trinitatis-breviarium}}
\newline {\color{Red} Lectio Prima.}
\lettrine[lines=2]{\zallmancaps \color{Blue} F}{actum} est vertente anno, eo tempore quo solent reges ad bella procedere, misit David Ioab et servos suos cum eo, et universum Israhel, et vastaverunt filios Amon, et obsederunt Rabath: David autem remansit in Hierusalem.%deleted autem, Kabath?
\newline {\color{Red} Lectio Secunda.}
\lettrine[lines=2]{\zallmancaps \color{Red} D}{um} hec agerentur, accidit ut surgeret David de strato suo post meridiem, et deambularet in solario domus regie. Viditque mulierem se lavantem ex adverso super solarium suum. Erat autem mulier pulchra valde. Misit ergo rex, et requisivit que esset mulier. Nunciatumque est ei, quod ipsa esset Bethsabee filia Heliam: uxor Urie Ethei.
\newline {\color{Red} Lectio Tertia.}
\lettrine[lines=2]{\zallmancaps \color{Blue} M}{issis} itaque David nunciis, tulit eam. Que cum ingressa esset ad illum: dormivit cum ea. Statimque sanctificata est ab immunditia sua: et reversa est domum suam concepto fetu. Mittensque nuntiavit David: et ait. Concepi. Misit autem David ad Ioab, dicens. Mitte ad me Uriam Etheum. Misitque Ioab Uriam ad David.
\newline {\color{Red} Lectio Quarta.}
\lettrine[lines=2]{\zallmancaps \color{Red} E}{t} venit Urias ad David. Quesivitque David, quam recte ageret Ioab et populus: et quomodo administraretur bellum. Et dixit David ad Uriam. Vade in domum tuam et lava pedes tuos. Egressusque est Urias de domo regis: secutusque est eum cibus regius.
\newline {\color{Red} Lectio Quinta.}
\lettrine[lines=2]{\zallmancaps \color{Blue} D}{ormivit} autem Urias ante portam domus regie cum aliis servis suis domini sui, et non descendit ad domum suam. Nuntiatumque est David a dicentibus. Non ivit Urias in domum suam. Et ait David ad Uriam. Nunquid non de via venisti? Quare non descendisti in domum tuam?
\newline {\color{Red} Lectio Sexta.}
\lettrine[lines=2]{\zallmancaps \color{Red} E}{t} ait Urias ad David. Archa Dei et Israhel et Iuda habitant in papilionibus, et dominus meus Ioab et domini servi mei super faciem terre manent, et ego ingrediar domum meam ut comedam et bibam et dormiam cum uxore mea? Per salutem tuam et per salutem anime tue, quod non faciam rem hanc.
\newline \ul{Lectiones feriales a duodecimo capitulo usque ad tercium regum ad libitum.}
{\ubsubsection{Dominica Septima}{dominica-septima-post-festum-trinitatis-breviarium}}
\newline {\color{Red} Lectio Prima.}
\lettrine[lines=2]{\zallmancaps \color{Blue} E}{t} rex David senuerat, habebatque etatis plurimos dies. Cumque operiretur vestibus: non calefiebat. Dixerunt ergo ei servi sui. Queramus domino nostro regi adolescentulam virginem, et stet coram rege, et foveat eum, dormiatque in sinu suo, et calefaciat dominum nostrum regem.
\newline {\color{Red} Lectio Secunda.}
\lettrine[lines=2]{\zallmancaps \color{Red} Q}{uesierunt} igitur adolescentulam speciosam in omnibus finibus Israhel, et invenerunt Abysag Sunanitem: et adduxerunt eam ad regem. Erat autem puella pulchra nimis. Dormiebatque cum rege et ministrabat ei. Rex vero non cognovit eam.
\newline {\color{Red} Lectio Tertia.}
\lettrine[lines=2]{\zallmancaps \color{Blue} A}{donias} autem filius Agith, elevabatur dicens. Ego regnabo. Fecitque sibi currum et equites, et quinquaginta viros qui ante eum currerent. Nec corripuit eum pater suus aliquando dicens: quare hoc fecisti. Erat autem et ipse pulcher valde, secundus natus post Absalon. Et sermo eius fuit cum Ioab filio Sarvie, et cum Abiathar sacerdote, qui adiuvabant partes Adonie.
\newline {\color{Red} Lectio Quarta.}
\lettrine[lines=2]{\zallmancaps \color{Red} S}{adoch} vero sacerdos, et Banaias filius Ioiade, et Nathan propheta, et Semei, et Cerethi et Phelethi, et robur exercitus David non erat cum Adonia. Immolatis ergo Adonias arietibus et vitulis et universis pinguibus iuxta lapidem Zoelech qui erat vicinus fonti Rogel, vocavit universos fratres suos filios regis, et omnes viros Iuda servos regis.
\newline {\color{Red} Lectio Quinta.}
\lettrine[lines=2]{\zallmancaps \color{Blue} N}{athani} autem prophetam, et Banaiam, et robustos quosque, et Salomonem fratrem suum: non vocavit. Dixit itaque Nathan ad Bethsabee matrem Salomonis. Num audisti quod regnaverit Adonias filius Agith et dominus noster David hoc ignorat? Nunc ergo veni et accipe consilium a me, et salva animam tuam filiique tui Salomonis.
\newline {\color{Red} Lectio Sexta.}
\lettrine[lines=2]{\zallmancaps \color{Red} V}{ade} et ingredere ad regem David: et dic ei. Nonne tu domine mi rex iurasti michi ancille tue dicens, quod Salomon filius tuus regnabit post me, et ipse sedebit in solio meo? Quare ergo regnat Adonias? Et adhuc ibi te loquente cum rege: ego veniam post te, et complebo sermones tuos.
\newline \ul{Lectiones feriales a secundo capitulo tertii regum usque ad decimum ad libitum.}
{\ubsubsection{Dominica Octava}{dominica-octava-post-festum-trinitatis-breviarium}}
\newline {\color{Red} Lectio Prima.}
\lettrine[lines=2]{\zallmancaps \color{Blue} R}{egina} Saba audita fama Salomonis, in nomine domini venit temptare eum in enigmatibus. Et ingressa Hierusalem cum multo comitatu et divitiis, camelis portantibus aromata, et aurum infinitum nimis et gemmas preciosas, venit ad regem Salomonem, et locuta est ei universa que habebat in corde suo. Et docuit eam Salomon omnia verba que proposuit. Non fuit sermo qui regem posset latere, et non responderet ei.
\newline {\color{Red} Lectio Secunda.}
\lettrine[lines=2]{\zallmancaps \color{Red} V}{idens} autem regina Saba omnem sapientiam Salomonis, et domum quam edificaverat, et cibos mense eius, et habitacula servorum, et ordines ministrantium, vestesque eorum et pincernas, et holocausta que offerebat in domo domini, non habebat ultra spiritum.
\newline {\color{Red} Lectio Tertia.}
\lettrine[lines=2]{\zallmancaps \color{Blue} D}{ixitque} ad regem. Verus est sermo quem audivi in terra mea, super sermonibus tuis, et super sapientia tua, et non credebam narrantibus michi donec ipsa veni, et vidi oculis meis et probavi quod media pars michi nunciata non fuerat. Maior est sapientia tua et opera tua, quam rumor quem audivi.
\newline {\color{Red} Lectio Quarta.}
\lettrine[lines=2]{\zallmancaps \color{Red} B}{eati} viri tui, et beati servi tui hi qui stant coram te semper et audiunt sapientiam tuam. Sit dominus Deus tuus benedictus cui complacuisti, et posuit te super thronum Israhel, eo quod dilexerit dominus Israhel in sempiternum, et constituit te regem, ut faceres iudicium et iusticiam.
\newline {\color{Red} Lectio Quinta.}
\lettrine[lines=2]{\zallmancaps \color{Blue} D}{edit} ergo regi centum viginti talenta auri, et aromata multa nimis, et gemmas preciosas. Non sunt allata ultra aromata tam multa quam ea que dedit regina Saba regi Salomoni. Sed et classis Hiram que portabat aurum de Ophyr, attulit de Ophyr ligna thiina multa nimis, et gemmas preciosas.
\newline {\color{Red} Lectio Sexta.}
\lettrine[lines=2]{\zallmancaps \color{Red} F}{ecitque} rex de lignis thiinis fulcra domus domini et domus regie: et cytharas lyrasque cantoribus. Non sunt allata huiusmodi ligna thiina neque visa usque in presentem diem. Rex autem Salomon dedit regine Saba omnia que voluit et petivit ab eo, exceptis his que ultro obtulerat ei munere regio. Que reversa est, et abiit in terram suam cum servis suis.
\newline \ul{Lectiones feriales ab undecimo capitulo usque ad quartum regum ad libitum.}
{\ubsubsection{Dominica Nona}{dominica-nona-post-festum-trinitatis-breviarium}}
\newline {\color{Red} Lectio Prima.}
\lettrine[lines=2]{\zallmancaps \color{Blue} P}{revaricatus} est Moab in Israhel: postquam mortuus est Achab, ceciditque Ochozias per cancellos cenaculi sui quod habebat in Samaria, et egrotavit. Misitque nuncios, dicens ad eos. Ite, consulite Beelzebub deum Acharon utrum vivere queam de infirmitate mea hac.%deleted autem
\newline {\color{Red} Lectio Secunda.}
\lettrine[lines=2]{\zallmancaps \color{Red} A}{ngelus} autem domini locutus est ad Heliam Thesbiten dicens. Surge ascende in occursum nunciorum regis Samarie: et dices ad eos. Nunquid non est Deus in Israhel, ut eatis ad consulendum Beelzebub deum Acharon? Quam ob rem hec dicit dominus. De lectulo super quem ascendisti non descendes: sed morte morieris. Et abiit Helyas.
\newline {\color{Red} Lectio Tertia.}
\lettrine[lines=2]{\zallmancaps \color{Blue} R}{eversique} sunt nuncii ad Ochoziam. Qui dixit eis. Quare reversi estis? At illi responderunt ei. Vir occurrit nobis, et dixit ad nos. Ite et revertimini ad regem qui misit vos: et dicetis ei. Hec dicit dominus. Nunquid quia non erat Deus in Israhel, mittis ut consulatur Beelzebub deus Acharon? Idcirco de lectulo super quem ascendisti non descendes: sed morte morieris.
\newline {\color{Red} Lectio Quarta.}
\lettrine[lines=2]{\zallmancaps \color{Red} Q}{ui} dixit eis. Cuius figure et habitus est vir qui occurrit vobis: et locutus est verba hec? At illi dixerunt. Vir pilosus, et zona pellicea accinctus renibus. Qui ait. Helyas Thesbites est. Misitque ad eum quinquagenarium principem, et quinquaginta qui erant sub eo. Qui ascendit ad eum.
\newline {\color{Red} Lectio Quinta.}
\lettrine[lines=2]{\zallmancaps \color{Blue} S}{edentique} in vertice montis: ait. Homo Dei, rex precepit ut descendas. Respondensque Helyas, dixit quinquagenario. Si homo Dei sum: descendat ignis de celo et devoret te et quinquaginta tuos. Descendit itaque ignis de celo, et devoravit eum, et quinquaginta qui erant cum eo.
\newline {\color{Red} Lectio Sexta.}
\lettrine[lines=2]{\zallmancaps \color{Red} R}{ursum} misit ad eum principem quinquagenarium alterum, et quinquaginta cum eo. Qui locutus est illi. Homo Dei: hec dicit rex. Festina, descende. Respondens Helyas: ait. Si homo Dei sum: descendat ignis de celo, et devoret te et quinquaginta tuos. Descendit ergo ignis de celo, et devoravit illum, et quinquaginta eius.
\newline \ul{Lectiones feriales a secundo capitulo quarti regum usque ad octavum ad libitum.}
{\ubsubsection{Dominica Decima}{dominica-decima-post-festum-trinitatis-breviarium}}
\newline {\color{Red} Lectio Prima.}
\lettrine[lines=2]{\zallmancaps \color{Blue} H}{elyseus} locutus est ad mulierem cuius vivere fecerat filium, dicens. Surge, vade tu et domus tua, et peregrinare ubicumque repereris, vocavit enim dominus famem, et veniet super terram septem annis. Que surrexit, et fecit iuxta verbum hominis Dei: et vadens cum domo sua peregrinata est in terra Philistiim diebus multis.
\newline {\color{Red} Lectio Secunda.}
\lettrine[lines=2]{\zallmancaps \color{Red} C}{umque} finiti essent anni septem reversa est mulier de terra Philistiim. Et egressa est: ut interpellaret regem pro domo sua et pro agris suis. Rex autem loquebatur cum Gyezi puero viri Dei: dicens. Narra michi omnia magnalia que fecit Helyseus. Cumque ille narraret regi quomodo mortuum suscitasset: apparuit mulier cuius vivificaverat filium, clamans ad regem pro domo sua et pro agris suis.
\newline {\color{Red} Lectio Tertia.}
\lettrine[lines=2]{\zallmancaps \color{Blue} D}{ixitque} Gyezi. Domine mi rex. Hec est mulier: et hic est filius eius quem suscitavit Helyseus. Et interrogavit rex mulierem. Que narravit ei.
\newline {\color{Red} Lectio Quarta.}
\lettrine[lines=2]{\zallmancaps \color{Red} D}{editque} ei rex eunuchum unum dicens. Restitue ei omnia que sua sunt, et universos reditus agrorum a die qua reliquit terram, usque ad presens.
\newline {\color{Red} Lectio Quinta.}
\lettrine[lines=2]{\zallmancaps \color{Blue} V}{enit} quoque Helyseus Damascum, et Benadad rex Sirie egrotabat. Nunciaveruntque ei dicentes. Venit vir Dei huc. Et ait rex ad Azael. Tolle tecum munera, et vade in occursum viri Dei: et consule dominum per eum dicens. Si evadere potero de infirmitate mea hac?
\newline {\color{Red} Lectio Sexta.}
\lettrine[lines=2]{\zallmancaps \color{Red} I}{vit} igitur Azael in occursum eius habens secum munera, et omnia bona Damasci, onera quadraginta camelorum. Cumque stetisset coram eo ait. Filius tuus Benadad rex Sirie, misit me ad te dicens. Si sanari potero de infirmitate mea hac? Dixitque ei Helyseus. Vade, dic ei. Sanaberis. Porro ostendit michi dominus: quia morte morietur.
\newline \ul{Lectiones feriales a quarto decimo capitulo quarti regum usque ad finem libri ad libitum.}
\lettrine[lines=2]{\zallmancaps \color{Blue} D}{}\ul{\textsc{ominica} prima Augusti et deinceps usque ad primam Dominicam Septembris exclusive cantetur hystoria,} \hyperlink{in-principio}{In principio.} \ul{quando de tempore agitur, et legantur lectiones de Libris Salomonis et de Libris Sapientie Et notandum quod Dominica prima mensis appellatur illa que propinquior est Kalendis, sive precedat Kalendas, sive sequatur, sive sit in ipsis Kalendis. Si igitur Kalende fuerunt in Dominica ibi est hystoria que illi Dominice assignatur nisi festum impediat inchoanda. Si vero Kalende fuerunt secunda, tercia, vel quarta feria, tunc in precedenti Dominica hystoria inchoetur. Si autem quinta vel sexta feria vel Sabbato Kalende fuerint, inchoari debet in sequenti Dominica.}
{\ubsubsection{Sabbato ante Primam Dominicam Augusti}{sabbato-ante-primam-dominicam-augusti-breviarium}}
\newline {\color{Red} Ad Vesperas.} \R \hyperlink{girum-celi}{Girum.} {\color{Red} II. Ad Mag. Ant.} Omnis sapientia a domino deo est, et cum illi fuit semper et est ante evum.
\newline {\color{Red} Ad Matutinas. Lectio Prima.}
\lettrine[lines=2]{\zallmancaps \color{Red} P}{arabole} Salomonis filii David regis Israhel ad sciendam sapientiam et disciplinam, ad intelligenda verba prudentie, et suscipiendam eruditionem doctrine, iusticiam et iudicium et equitatem: ut detur parvulis astutia et adolescenti scientia et intellectus. \R
\lettrine[lines=2]{\zallmancaps \color{Blue} I}{n} \hypertarget{in-principio}{\label{in-principio}} principio deus antequam terram faceret, priusquam abissos constitueret, priusquam produceret fontes aquarum, antequam montes collocarentur, * Ante omnes colles generavit me dominus. \V Quando preparabat celos aderam, cum eo cuncta componens. Ante.
\newline {\color{Red} Lectio Secunda.}
\lettrine[lines=2]{\zallmancaps \color{Red} A}{udiens} sapiens sapientior erit, et intelligens gubernacula possidebit. Animadvertet parabolam et interpretationem: verba sapientum et enigmata eorum. Timor domini principium sapientie. Sapientiam atque doctrinam stulti despiciunt.
\begin{responsory}[girum-celi]
{Girum celi circuivi sola et in fluctibus maris ambulavi, in omni gente et in omni populo primatum tenui, * Superborum et sublimium colla propria virtute calcavi.}{Ego in altissimis habito et thronus meus in colunna nubis.}%typo tenuii
\end{responsory}%
\newline {\color{Red} Lectio Tertia.}
\lettrine[lines=2]{\zallmancaps \color{Blue} A}{udi} fili mi disciplinam patris tui, et ne dimittas legem matris tue: ut addatur gratia capiti tuo, et torques collo tuo. Fili mi si te lactaverint peccatores, ne acquiescas eis. Si dixerint, veni nobiscum, insidiemur sanguini abscondamus tendiculas contra insontem frustra, deglutiamus eum sicut infernus viventem et integrum, quasi descendentem in lacum, omnem preciosam substantiam reperiemus, implebimus domos nostras spoliis, sortem mitte nobiscum, marsupium unum sit omnium nostrum: fili mi ne ambules cum eis.
\begin{responsory-final}[emitte-domine]
{Emitte domine sapientiam de sede magnitudinis tue, ut mecum sit et mecum laboret, * Ut sciam quid acceptum sit coram te * Omni tempore.}{Da michi domine sedium tuarum assistricem sapientiam.}{Omni}
\end{responsory-final}%
\newline {\color{Red} Lectio Quarta.}
\lettrine[lines=2]{\zallmancaps \color{Red} P}{rohibe} pedem tuum a semitis eorum. Pedes enim illorum ad malum currunt, et festinant ut effundant sanguinem. Frustra autem iacitur rete ante oculos pennatorum. Ipsi quoque contra sanguinem suum insidiantur, et moliuntur fraudes contra animas suas. Sic semite omnis avari: animas possidentium rapiunt.
\begin{responsory}[da-michi-domine]
{Da michi domine sedium tuarum assistricem sapientiam, et noli me reprobare a pueris tuis, * Quoniam servus tuus sum ego et filius ancille tue.}{Mitte illam a sede magnitudinis tue ut mecum sit et mecum laboret.}
\end{responsory}%
\newline {\color{Red} Lectio Quinta.}
\lettrine[lines=2]{\zallmancaps \color{Blue} S}{apientia} foris predicat: in plateis dat vocem suam. In capite turbarum clamitat, in foribus portarum urbis profert verba sua dicens. Usquequo parvuli diligitis infantiam, et stulti ea que sibi sunt noxia cupient, et imprudentes odibunt scientiam?
\begin{responsory}[domine-pater]
{Domine pater et deus vite mee, ne derelinquas me in cogitatu maligno, extollentiam oculorum meorum ne dederis michi, et desiderium malignum averte a me domine, aufer a me concupiscentiam, * Et animo irreverenti et infrunito ne tradas me domine.}{Ne dereliquas me domine ne accrescant ignorantie mee et multiplicentur delicta mea.}
\end{responsory}%
\newline {\color{Red} Lectio Sexta.}
\lettrine[lines=2]{\zallmancaps \color{Red} C}{onvertimini} ad correptionem meam. En proferam vobis spiritum meum et ostendam verba mea. Quia vocavi et renuistis, extendi manum meam, et non fuit qui aspiceret. Despexistis omne consilium meum, et increpationes meas neglexistis. Ego quoque in interitu vestro ridebo, et subsannabo cum vobis quod timebatis advenerit.
\begin{responsory-doxology}[super-salutem-augusti]
{Super salutem et omnem pulchritudinem dilexi sapientiam et proposui pro luce habere illam, * Venerunt michi omnia bona pariter cum illa.}{Dixi sapientie soror mea es, et prudentiam vocavi amicam meam.}
\end{responsory-doxology}%
\begin{responsory}[initium-sapientie]
{Initium sapientie timor domini, intellectus bonus omnibus facientibus eum, * Laudatio eius manet in seculum seculi.}{Dilectio illius custodia legum est, quia omnis sapientia timor domini.}
\end{responsory}%
\begin{responsory}[magna-enim-sunt]
{Magna enim sunt iudicia tua domine et inenarrabilia verba tua, * Magnificasti populum tuum et honorasti.}{Deduxisti sicut oves populum tuum in manu Moysi et Aaron.}
\end{responsory}%
\begin{responsory-final}[verbum-iniquum]
{Verbum iniquum et dolosum longe fac a me domine, * Divitias et paupertates ne dederis michi, * Sed tantum victui meo tribue necessaria.}{Duo rogavi te ne deneges michi antequam moriar.}{Sed}
\end{responsory-final}%
\newline \ul{Sequentes antiphone sunt dicende in Sabbatis ad Mag. vel ad memoriam quamdiu cantabitur hystoria,} \hyperlink{in-principio}{In principio.}
\newline {\color{Red} Ant.} Sapientia edificavit sibi domum, scidit columnas septem, subdidit sibi gentes, superborumque et sublimium colla, propria virtute calcavit. {\color{Red} Ant.} Sapientia clamitat in plateis, siquis diligit sapientiam ad me declinet, et eam inveniet, et eam cum invenerit beatus erit si tenuerit eam. {\color{Red} Ant.} Dominus possedit me in initio viarum suarum antequam quicquam faceret a principio, necdum erant abyssi et ego parturiebar, quando preparabat celos aderam cum eo componens omnia. {\color{Red} Ant.} Ego in altissimis habitavi et thronus meus in columna nubis.
\newline \ul{Lectiones alie per totum mensem usque ad Septembrem per quinque Libros Salomonis ad libitum.}
\lettrine[lines=2]{\zallmancaps \color{Blue} D}{}\ul{\textsc{ominica} prima Septembris et deinceps usque ad terciam Dominicam eiusdem mensis exclusive idest per quindecim dies tantum quando de tempore agitur cantetur hystoria,} \hyperlink{si-bona}{Si bona.} \ul{et legatur de Libro Iob.}
{\ubsubsection{Sabbato ante Primam Dominicam Septembris}{sabbato-ante-primam-dominicam-septembris-breviarium}}
\newline \ul{Sabbato vero precedenti Ad Vesperas.} \R \hyperlink{antequam-comedam}{Antequam.} {\color{Red} II. Ad Mag. Ant.} Cum audisset Iob nuntiorum verba sustinuit patienter et ait, si bona suscepimus de manu domini, mala autem quare non sustineamus, in omnibus his non peccavit Iob labiis suis, neque stultum aliquid contra deum locutus est.
\newline \ul{Si hec Dominica in Kalendis occurrerit fiat sicut infra octavam Beati Augustini in festo sancti Egidii est notatum.}
\newline Sequenti Sabbato. {\color{Red} Ad Mag. Ant.} In omnibus his non peccavit Iob labiis suis neque stultum quid contra deum locutus est.
\newline {\color{Red} Ad Matutinas. Lectio Prima.}
%pp113
\lettrine[lines=2]{\zallmancaps \color{Red} V}{ir} erat in terra Hus nomine Iob. Et erat vir ille simplex et rectus ac timens Deum: et recedens a malo. Natique sunt ei septem filii et tres filie. Et fuit possessio eius, septem milia ovium, et tria milia camelorum, quingenta quoque iuga boum, et quingente asine: ac familia multa nimis. \R
\lettrine[lines=2]{\zallmancaps \color{Blue} S}{i} \hypertarget{si-bona}{\label{si-bona}} bona suscepimus de manu domini mala autem quare non sustineamus, dominus dedit dominus abstulit, * Sicut domino placuit ita factum est, sit nomen domini benedictum. \V Nudus egressus sum de utero matris mee nudus revertar illuc. Sicut.
\newline {\color{Red} Lectio Secunda.}
\lettrine[lines=2]{\zallmancaps \color{Red} E}{ratque} vir ille magnus inter omnes orientales. Et ibant filii eius et faciebant convivium per domos unusquisque in die suo. Et mittentes vocabant tres sorores suas, ut comederent et biberent cum eis. Cumque in orbem transissent dies convivii: mittebat ad eos Iob: et sanctificabat illos.
\begin{responsory}[antequam-comedam]
{Antequam comedam suspiro et tanquam inundantes aque sic rugitus meus, quia timor quem timebam evenit michi, et quod verebar accidit, non ne dissimulavi, nonne silui nonne quievi, * Et venit super me indignatio tua domine.}{Ecce non est auxilium michi in me: necessarii quoque mei recesserunt a me.}
\end{responsory}%
\newline {\color{Red} Lectio Tertia.}
\lettrine[lines=2]{\zallmancaps \color{Blue} C}{onsurgensque} diluculo offerebat holocausta per singulos. Dicebat enim. Ne forte peccaverint filii mei: et benedixerint Deo in cordibus suis. Sic faciebat Iob cunctis diebus. Quadam autem die cum venissent filii Dei ut assisterent coram domino: affuit inter eos eciam Sathan. Cui dixit dominus. Unde venis? Qui respondens ait. Circuivi terram: et perambulavi eam.
\begin{responsory-doxology}[utinam-appenderentur]
{Utinam appenderentur peccata mea quibus iram merui, et calamitas quam patior, * In statera.}{Quasi arena maris hec gravior appareret.}
\end{responsory-doxology}%
\newline {\color{Red} Lectio Quarta.}
\lettrine[lines=2]{\zallmancaps \color{Red} D}{ixitque} dominus ad eum. Nunquid considerasti servum meum Iob, quod non sit ei similis in terra homo simplex et rectus ac timens Deum et recedens a malo?
\begin{responsory}[quare-detraxistis]
{Quare detraxistis sermonibus veritatis, ad increpandum verba componitis, et subvertere nitimini amicum vestrum, * Veruntamen que cogitastis explete.}{Quod iustum est iudicate et non invenietis in lingua mea iniquitatem.}
\end{responsory}%
\newline {\color{Red} Lectio Quinta.}
\lettrine[lines=2]{\zallmancaps \color{Blue} C}{ui} respondens Sathan ait. Nunquid frustra Iob timet Deum? Nonne tu vallasti eum, ac domum eius, universamque substantiam per circuitum? Operibus manuum eius benedixisti: et possessio eius crevit in terra. Sed extende paululum manum tuam et tange cuncta que possidet: nisi in facie benedixerit tibi.
\begin{responsory}[induta-est-caro]
{Induta est caro mea putredine et sordibus pulveris, cutis mea aruit et contracta est, * Memento mei domine quia ventus est vita mea.}{Dies mei sicut umbra declinaverunt et ego sicut fenum arui.}
\end{responsory}%
\newline {\color{Red} Lectio Sexta.}
\lettrine[lines=2]{\zallmancaps \color{Red} D}{ixit} ergo dominus ad Sathan. Ecce universa que habet in manu tua sunt, tantum in eum ne extendas manum tuam. Egressusque est Sathan a facie domini. Cum autem quadam die filii et filie eius comederent et biberent vinum in domo fratris sui primogeniti, nuntius venit ad Iob qui diceret. Boves arabant, et asine pascebantur iuxta eos, et irruerunt Sabei tuleruntque omnia, et pueros percusserunt gladio, et evasi ego solus ut nuntiarem tibi.
\begin{responsory-doxology}[paucitas-dierum]
{Paucitas dierum meorum finietur brevi, dimitte me domine ut plangam paululum dolorem meum antequam vadam, * Ad terram tenebrosam et opertam mortis caligine.}{Remitte michi ut refrigerer priusquam abeam.}
\end{responsory-doxology}%
\begin{responsory}[memento-mei-deus-breve]
{Memento mei deus quia ventus est vita mea, * Nec aspiciet me visus hominis.}{Et non revertetur oculus meus ut videat bona.}
\end{responsory}%
\begin{responsory}[ne-abscondas-me]
{Ne abscondas me domine a facie tua, manum tuam longe fac a me, * Et formido tua non me terreat.}{Corripe me domine in misericordia et non in furore tuo, ne forte ad nichilum redigas me.}
\end{responsory}%
\begin{responsory-final}[nocte-os-meum]
{Nocte os meum perforatur doloribus, et qui me comedunt non dormiunt, a multitudine eorum consumitur vestimentum meum, * Comparatus sum luto, * Et assimilatus sum faville et cineri.}{O custos hominum quare me posuisti contrarium tibi, et factus sum michimetipsi gravis, parce michi domine nichil enim sunt dies mei.}{Et}
\end{responsory-final}%
\newline \ul{Lectiones feriales usque ad Dominicam terciam Septembris per residuum Libri.}
\lettrine[lines=2]{\zallmancaps \color{Blue} D}{}\ul{\textsc{ominica} tercia Septembris et deinceps usque ad primam Dominicam Octobris exclusive cantetur hystoria de Thobia, Iudith et Hester quando de tempore agitur et legantur lectiones de eisdem Libris.}
{\ubsubsection{Sabbato ante Terciam Dominicam Septembris}{sabbato-ante-terciam-dominicam-septembris-breviarium}}
\newline \ul{Sabbato vero precedenti Ad Vesperas} \R \hyperlink{omni-tempore}{Omni tempore.} {\color{Red} II. Ad Mag. Ant.} Ingressus Raphael angelus ad Thobiam salutavit eum dicens, gaudium tibi semper sit, cui Thobias ait, quale gaudium michi erit qui in tenebris sedeo et lumen celi non video, cui angelus inquit, forti animo esto, in proximo est ut a deo cureris.
\newline {\color{Red} Ad Matutinas. Lectio Prima.}
\lettrine[lines=2]{\zallmancaps \color{Red} T}{hobias} ex tribu et civitate Neptalim que est in superioribus Galilee supra Naason post viam que ducit ad occidentem in sinistro habens civitatem Sephet cum captus esset in diebus Salmanasar regis Assiriorum: in captivitate tamen positus viam veritatis non deseruit, ita ut omnia que habere poterat cotidie concaptivis fratribus qui erant ex eius genere impertiret. {\color{Red} Responsorium.}
\lettrine[lines=2]{\zallmancaps \color{Blue} P}{eto} \hypertarget{peto-domine}{\label{peto-domine}} domine ut de vinculo improperii huius absolvas me, aut certe desuper terram eripias me, ne reminiscaris delicta mea vel parentum meorum, neque vindictam sumas de peccatis meis, * Quia eruis sustinentes te domine. \V Omnia iudicia tua iusta sunt et omnes vie tue misericordia et veritas et nunc domine memento mei. Quia.
\newline {\color{Red} Lectio Secunda.}
\lettrine[lines=2]{\zallmancaps \color{Red} C}{umque} esset iunior omnibus in tribu Neptalim: nichil tamen puerile gessit in opere. Denique cum irent omnes ad vitulos aureos quos Ieroboam fecerat rex Israhel: hic solus fugiebat consortia omnium, et pergebat in Hierusalem ad templum domini, et ibi adorabat dominum Deum Israhel, omnia primitiva sua et decimas suas fideliter offerens, ita ut in tercio anno proselitis et advenis ministraret omnem decimationem.
\begin{responsory}[omni-tempore]
{Omni tempore benedic deum et pete ab eo ut vias tuas dirigat, * Et omni tempore consilia tua in ipso permaneant.}{Inquire ut facias que placita sunt ei in veritate et in tota virtute tua.}
\end{responsory}%
\newline {\color{Red} Lectio Tertia.}
\lettrine[lines=2]{\zallmancaps \color{Blue} H}{ec} et his similia secundum legem Dei celi puerulus observabat. Cum vero factus fuisset vir accepit uxorem Annam de tribu sua: genuitque ex ea filium, nomen suum imponens ei. Quem ab infantia timere Deum docuit, et abstinere ab omni peccato.
\begin{responsory-doxology}[tempus-est-ut]
{Tempus est ut revertar ad eum qui me misit, vos autem benedicite deum, * Et enarrate omnia mirabilia eius.}{Confitemini ei coram omnibus viventibus quia fecit vobiscum misericordiam suam.}
\end{responsory-doxology}%
\newline {\color{Red} Lectio Quarta.}
\lettrine[lines=2]{\zallmancaps \color{Red} I}{gitur} cum per captivitatem devenisset cum uxore sua et filio in civitatem Niniven, cum omni tribu sua, et omnes ederent ex cibis gentilium, iste custodivit animam suam, et nunquam contaminatus est in escis eorum. Et quoniam memor fuit domini in toto corde suo, dedit illi Deus gratiam in conspectu Salmanasar regis, et dedit illi potestatem quocumque vellet ire, habens libertatem quecumque facere voluisset. Pergebat ergo per omnes qui erant in captivitate, et monita salutis dabat eis.
\begin{responsory}[adonay-domine]
{Adonay domine deus magne et mirabilis, qui dedisti salutem in manu femine, * Exaudi preces servorum tuorum.}{Benedictus es domine qui non derelinquas presumentes de te, et de sua virtute gloriantes humilias.}
\end{responsory}%
\newline {\color{Red} Lectio Quinta.}
\lettrine[lines=2]{\zallmancaps \color{Blue} C}{um} autem venisset in Rages civitatem Medorum, et ex his quibus honoratus fuerat a rege habuisset decem talenta argenti, et cum in multa turba generis sui Gabelum egentem videret, qui erat ex tribu eius, sub cyrographo dedit illi memoratum pondus argenti.
\begin{responsory}[dominator-domine]
{Dominator domine celorum et terre, creator aquarum, rex universe creature tue, * Exaudi orationes servorum tuorum.}{Tu domine cui humilium semper et mansuetorum placuit deprecatio.}
\end{responsory}%
\newline {\color{Red} Lectio Sexta.}
\lettrine[lines=2]{\zallmancaps \color{Red} P}{ost} multum vero temporis: mortuo Salmanasar rege cum regnaret Sennacherib filius eius pro eo, et filios Israhel exosos haberet in conspectu suo, Thobias pergebat per omnem cognationem suam, et consolabatur eos.
\begin{responsory-final}[spem-in-alium]
{Spem in alium nunquam habui preter in te deus Israhel, * Qui irasceris et propitius eris, * Et omnia peccata hominum in tribulatione dimittis.}{Domine deus celi et terre respice ad humilitatem nostram.}{Et omnia}
\end{responsory-final}%
\begin{responsory}[conforta-me-rex]
{Conforta me rex sanctorum principatum tenens, * Et da sermonem rectum et bene sonantem in os meum.}{Da nobis domine locum penitentie, et ne claudas ora canentium te domine.}
\end{responsory}%
\begin{responsory}[laudate-dominum]
{Laudate dominum deum nostrum, * Qui non deseruit sperantes in se, et in me adimplevit misericordiam suam, quam promisit domui Israhel.}{Laudate dominum omnes gentes et collaudate eum omnes populi.}
\end{responsory}%
\begin{responsory-doxology}[domine-rex-omnipotens]
{Domine rex omnipotens in ditione tua cuncta sunt posita, et non est qui possit resistere voluntati tue, libera nos, * Propter nomen tuum.}{Exaudi domine orationem nostram, et converte luctum nostrum in gaudium.}
\end{responsory-doxology}%
\newline \ul{Due sequentes antiphone dicantur in Sabbatis Ad Mag. vel ad memoriam.}
\newline {\color{Red} Ant.} Ne reminiscaris domine delicta mea vel parentum meorum, neque vindictam suam de peccatis meis. {\color{Red} Ant.} Adonay domine deus magne et mirabilis qui dedisti salutem in manu femine exaudi preces servorum tuorum.
\newline \ul{Si unum tantum Sabbatum fuerit dimittatur antiphona.} Adonay.
\newline {\color{Red} Lectiones de Iudith. Lectio.}
\lettrine[lines=2]{\zallmancaps \color{Blue} A}{rphaxat} itaque rex Medorum subiugaverat multas gentes imperio suo, et ipse edificavit civitatem potentissimam quam appellavit Egbathanis. Ex lapidibus quadratis et sectis fecit muros eius in altitudinem cubitorum septuaginta, et in latitudinem cubitorum triginta.
\newline {\color{Red} Lectio.}
\lettrine[lines=2]{\zallmancaps \color{Red} T}{urres} vero eius posuit in altitudinem cubitorum centum. Per quadrum vero earum latus utrumque vicenorum pedum spatio tendebatur. Posuitque portas eius in altitudinem turrium. Et gloriabatur quasi potens in potentia exercitus sui, et in gloria quadrigarum suarum.
\newline {\color{Red} Lectio.}
\lettrine[lines=2]{\zallmancaps \color{Blue} A}{nno} igitur duodecimo regni sui Nabuchodonosor rex Assiriorum, qui regnabat in Ninive civitate magna: pugnavit contra Arphaxat, et obtinuit eum in campo magno qui appellatur Ragau circa Eufraten et Tigrim et Iadasan in campo Erioch regis Elichorum.
\newline {\color{Red} Lectio.}
\lettrine[lines=2]{\zallmancaps \color{Red} T}{unc} exaltatum est regnum Nabuchodonosor, et cor eius elevatum est. Et misit ad omnes qui habitabant in Cilicia, et Damasco, et Libano, et ad gentes que sunt in Carmelo et Cedar, et inhabitantes Galileam, in campo magno Esdrelon, et ad omnes qui erant in Samaria et trans flumen Iordanem usque ad Hierusalem, et omnem terram Iesse, quousque perveniatur ad montes Ethiopie.
\newline {\color{Red} Lectio.}
\lettrine[lines=2]{\zallmancaps \color{Blue} A}{d} hos omnes misit nuntios Nabuchodonosor rex Assiriorum. Qui omnes uno animo contradixerunt: et remiserunt eos vacuos, ac sine honore abiecerunt. Tunc indignatus Nabuchodonosor rex ad omnem terram illam: iuravit per thronum suum et regnum quod defenderet se de omnibus his.
\newline {\color{Red} Lectio.}
\lettrine[lines=2]{\zallmancaps \color{Red} A}{nno} tertiodecimo Nabuchodonosor regis, vicesima et secunda die mensis primi, factum est verbum in domo Nabuchodonosor regis Assiriorum ut defenderet se. Vocavitque omnes maiores omnesque duces bellatores suos: et habuit cum eis misterium consilii sui. Dixitque cogitationem suam in eo esse: ut omnem terram suo subiugaret imperio.
\newline {\color{Red} Lectiones de Hester. Lectio.}
\lettrine[lines=2]{\zallmancaps \color{Blue} I}{n} diebus Assueri qui regnavit ab India usque Ethiopiam super centum viginti septem provincias, quando sedit in solio regni sui, Susa civitas regni eius exordium fuit. Tertio igitur anno imperii sui fecit grande convivium cunctis principibus et pueris suis fortissimis Persarum et Medorum inclitis, et prefectis provinciarum coram se, ut ostenderet divitias glorie regni sui, ac magnitudinem atque iactantiam potentie sue multo tempore, centum videlicet et octoginta diebus.
\newline {\color{Red} Lectio.}
\lettrine[lines=2]{\zallmancaps \color{Red} C}{umque} implerentur dies convivii invitavit omnem populum qui inventus est Susis a maximo usque ad minimum, et iussit septem diebus convivium preparari in vestibulo orti et nemoris quod regio cultu et manu consitum erat. Et pendebant ex omni parte tentoria aerei coloris et carbasini ac iacintini, sustentata funibus bissinis ac purpureis, qui eburneis circulis inserti erant, et columnis marmoreis fulciebantur.
\newline {\color{Red} Lectio.}
\lettrine[lines=2]{\zallmancaps \color{Blue} L}{ectuli} quoque aurei et argentei super pavimentum smaragdo et pario, stratum lapide dispositi erant: quod mira varietate pictura decorabat. Bibebant autem qui invitati erant aureis poculis, et aliis atque aliis vasis cibi inferebantur.
\newline {\color{Red} Lectio.}
\lettrine[lines=2]{\zallmancaps \color{Red} V}{inum} quoque ut magnificentia regia dignum erat, abundans et precipuum ponebatur. Nec erat qui nolentes cogeret ad bibendum, sed sic rex statuerat preponens mensis singulos de principibus suis, ut sumeret unusquisque quod vellet.
\newline {\color{Red} Lectio.}
\lettrine[lines=2]{\zallmancaps \color{Blue} V}{asthi} quoque regina fecit convivium feminarum in palatio ubi rex Assuerus manere consueverat. Itaque die septimo cum rex esset hilarior, et post nimiam potationem incaluisset mero, precepit Mauman, et Bazatha, et Arbona, et Bagatha, et Abgatha, et Zarath, et Carthas septem eunuchis qui in conspectu eius ministrabant, ut introducerent reginam Vasthi coram rege, posito super caput eius diademate, ut ostenderet cunctis populis et principibus illius pulchritudinem. Erat enim pulchra valde. Que renuit, et ad regis imperium quod per eunuchos mandaverat venire contempsit.
\newline {\color{Red} Lectio.}
\lettrine[lines=2]{\zallmancaps \color{Red} U}{nde} iratus rex, et nimio furore succensus, interrogavit sapientes qui ex more regio semper ei aderant, et illorum faciebat cuncta consilio, scientium leges ac iura maiorum, erant autem primi et proximi, Carsena, et Sechar, et Admatha, et Tharsis, et Mares, et Marsana et Mamucha, septem duces Persarum atque Medorum, qui videbant faciem regis, et primi post eum residere soliti erant, cui sententie Vasthi regina subiaceret que Assueri regis imperium quod per eunuchos mandaverat facere noluisset.
\lettrine[lines=2]{\zallmancaps \color{Blue} D}{}\ul{\textsc{ominica} prima Octobris et deinceps usque ad primam Dominicam Novembris exclusive cantetur hystoria,} \hyperlink{adaperiat-dominus}{Adaperiat.} \ul{et legatur de Libris Machabeorum quando de tempore agitur.}
{\ubsubsection{Sabbato ante Primam Dominicam Octobris}{sabbato-ante-primam-dominicam-octobris-breviarium}}
\newline \ul{Sabbato vero precedenti Ad Vesperas.} \R \hyperlink{tua-est-potentia}{Tua est.} {\color{Red} III. Ad Mag. Ant.} Adaperiat dominus cor vestrum in lege sua et in preceptis suis et faciat pacem.
\newline {\color{Red} Ad Matutinas. Lectio Prima.}
\lettrine[lines=2]{\zallmancaps \color{Red} E}{t} factum est, postquam percussit Alexander Philippi Macedo qui primus regnavit in Grecia egressus de terra Cethin Darium regem Persarum et Medorum, constituit prelia multa, et omnium obtinuit munitiones, et interfecit reges terre, et pertransiit usque ad fines terre. \R
\lettrine[lines=2]{\zallmancaps \color{Blue} A}{daperiat} \hypertarget{adaperiat-dominus}{\label{adaperiat-dominus}} dominus cor vestrum in lege sua et in preceptis suis et faciat pacem in diebus vestris, * Concedat vobis salutem et redimat vos a malis. \V Exaudiat dominus orationes vestras, et reconcilietur vobis, nec vos deserat in tempore malo. Concedat.
\newline {\color{Red} Lectio Secunda.}
\lettrine[lines=2]{\zallmancaps \color{Red} E}{t} accepit spolia multitudinis gentium: et siluit terra in conspectu eius. Et congregavit virtutem et exercitum fortem nimis, et exaltatum est et elevatum cor eius, et obtinuit regiones gentium, et tyrannos: et facti sunt illi in tributum.
\begin{responsory}[exaudiat-dominus]
{Exaudiat dominus orationes vestras et reconcilietur vobis, * Nec vos deserat in tempore malo.}{Det vobis cor omnibus ut colatis eum et faciatis eius voluntatem.}
\end{responsory}%
\newline {\color{Red} Lectio Tertia.}
\lettrine[lines=2]{\zallmancaps \color{Blue} E}{t} post hec decidit in lectum, et cognovit quia moreretur. Et vocavit pueros suos nobiles qui secum erant nutriti a iuventute sua, et divisit eis regnum suum cum adhuc viveret. Et regnavit Alexander annis duodecim: et mortuus est.
\begin{responsory-final}[tua-est-potentia]
{Tua est potentia, tuum regnum domine, tu es super omnes gentes, * Da pacem domine * In diebus nostris.}{Creator omnium deus, terribilis et fortis iustus et misericors.}{In diebus}
\end{responsory-final}%
\newline {\color{Red} Lectio Quarta.}
\lettrine[lines=2]{\zallmancaps \color{Red} E}{t} obtinuerunt pueri eius regnum unusquisque in loco suo. Et imposuerunt sibi omnes diademata post mortem eius, et filii eorum post eos annis multis, et multiplicata sunt mala in terra. Et exiit ex eis radix peccatrix Anthiocus illustris, filius Anthiochi regis, qui fuerat Rome obses: et regnavit in anno centesimo tricesimo et septimo regni Grecorum.
\begin{responsory}[congregata-sunt]
{Congregata sunt inimici nostri et gloriantur in virtute sua, contere fortitudinem illorum domine, et disperge illos, * Ut cognoscant quia non est alius qui pugnet pro nobis nisi tu deus noster.}{Disperge illos in virtute tua et depone eos protector noster domine.}
\end{responsory}%
\newline {\color{Red} Lectio Quinta.}
\lettrine[lines=2]{\zallmancaps \color{Blue} I}{n} diebus illis exierunt ex Israhel filii iniqui, et suaserunt multis dicentes. Eamus et disponamus testamentum cum gentibus que circa nos sunt, quia ex quo recessimus ab eis, invenerunt nos mala multa. Et visus est bonus sermo in oculis eorum.
\begin{responsory}[impetum-inimicorum]
{Impetum inimicorum ne timueritis, memores estote quomodo salvi facti sunt patres nostri, * Et nunc clamemus in celum et miserebitur nostri deus noster.}{Mementote mirabilium eius que fecit Pharaoni et exercitui eius in mari rubro.}
\end{responsory}%
\newline {\color{Red} Lectio Sexta.}
\lettrine[lines=2]{\zallmancaps \color{Red} E}{t} destinaverunt aliqui de populo, et abierunt ad regem. Et dedit illis potestatem: ut facerent iusticias gentium. Et edificaverunt ginnasium in Hierosolimis secundum leges nationum, et fecerunt sibi preputia, et recesserunt a testamento sancto, et iuncti sunt nationibus: et venundati sunt ut facerent malum.
\begin{responsory-doxology}[refulsit-sol]
{Refulsit sol in clipeos aureos et resplenduerunt montes ab eis, * Et fortitudo gentium dissipata est.}{Erat enim exercitus magnus valde et fortis, et appropiavit Iudas et exercitus eius in prelium.}
\end{responsory-doxology}%
\begin{responsory}[dixit-iudas-simoni]
{Dixit Iudas Simoni fratri suo, elige tibi viros et vade libera fratres tuos in Galileam, ego autem et Ionathas frater meus ibimus in Galadithim, * Sicut fuerit voluntas in celi sic fiat.}{Accingimini et estote filii potentes, quoniam melius est nobis mori in bello, quam videre mala gentis nostre et sanctorum.}
\end{responsory}%
\begin{responsory}[in-ymnis-et]
{In ymnis et confessionibus benedicebant dominum, * Qui magna fecit in Israhel, et victoriam dedit illis dominus omnipotens.}{Ornaverunt faciem templi coronis aureis, et dedicaverunt altare domino.}
\end{responsory}%
\begin{responsory-final}[ornaverunt-faciem]
{Ornaverunt faciem templi coronis aureis et dedicaverunt altare domino, * Et facta est letitia magna * In populo.}{In ymnis et confessionibus benedicebant dominum.}{In populo}
\end{responsory-final}%
\newline \ul{Sequentes antiphone dicende sunt in Sabbatis ad Mag. vel ad memoriam quamdiu cantabitur hystoria,} \hyperlink{adaperiat-dominus}{Adaperiat.}
\newline {\color{Red} Ant.} Da pacem domine in diebus nostris, quia non est alius qui pugnet pro nobis nisi tu deus noster. {\color{Red} Ant.} Tua est potentia tuum regnum domine, tu es super omnes gentes, da pacem domine in diebus nostris. {\color{Red} Ant.} Accingimini filii potentes, et estote parati pugnare in prelio, quoniam melius est nobis mori in bello, quam videre mala gentis nostre et sanctorum, sicut fuerit voluntas in celo sic fiat. {\color{Red} Ant.} Exaudiat dominus orationes vestras et reconcilietur vobis, nec vos deserat in tempore malo.
\newline \ul{Lectiones feriales ab illo loco,} Et veritus est. \ul{in primo capitulo primi Machabeorum usque ad illum locum,} Tunc congregata est synagoga. \ul{in secundo capitulo.}
{\ubsubsection{Dominica Secunda}{sabbato-ante-secundam-dominicam-octobris-breviarium}}
\newline {\color{Red} Lectio Prima.}
\lettrine[lines=2]{\zallmancaps \color{Blue} T}{unc} congregata est synagoga fortis viribus ex Israhel omnis voluntarius in lege. Et omnes qui fugiebant a malis additi sunt ad eos, et facti sunt illis ad firmamentum. Et collegerunt exercitum, et percusserunt peccatores in ira sua, et viros iniquos in indignatione sua, et ceteri fugerunt ad nationes ut evaderent.
\newline {\color{Red} Lectio Secunda.}
\lettrine[lines=2]{\zallmancaps \color{Red} E}{t} circuivit Mathathias et amici eius et destruxerunt aras, et circumciderunt pueros incircumcisos quotquot invenerunt
%pp114
in finibus Israhel in fortitudine, et persecuti sunt filios superbie, et prosperatum est opus in manu eorum. Et obtinuerunt legem de manibus gentium et de manibus regum, et non dederunt cornu peccatori.
\newline {\color{Red} Lectio Tertia.}
\lettrine[lines=2]{\zallmancaps \color{Blue} E}{t} appropinquaverunt dies Mathathie moriendi, et dixit filiis suis. Nunc confortata est superbia et castigatio, et tempus eversionis, et ira indignationis. Nunc ergo o filii emulatores estote legis, et date animas vestras pro testamento patrum. Et mementote operum patrum que fecerunt in generationibus suis, et accipietis gloriam magnam et nomen eternum.
\newline {\color{Red} Lectio Quarta.}
\lettrine[lines=2]{\zallmancaps \color{Red} A}{braham} nonne in temptatione inventus est fidelis, et reputatum est ei ad iusticiam? Ioseph in tempore angustie sue custodivit mandatum, et factus est dominus Egipti. Phinees pater noster zelando zelum Dei, accepit testamentum sacerdotii eterni.
\newline {\color{Red} Lectio Quinta.}
\lettrine[lines=2]{\zallmancaps \color{Blue} I}{hesus} dum implet verbum: factus est dux in Israhel. Caleph dum testificatur in ecclesia, accepit hereditatem. David in sua misericordia: consecutus est sedem regni in secula. Helias dum zelat zelum legis, receptus est in celum.
\newline {\color{Red} Lectio Sexta.}
\lettrine[lines=2]{\zallmancaps \color{Red} A}{nanias}, Azarias, Misael credentes: liberati sunt de flamma. Daniel in sua simplicitate: liberatus est de ore leonum. Et ita cogitate per generationem et generationem, quia omnes qui sperant in deo, non infirmantur. Et a verbis viri peccatoris ne timueritis: quia gloria eius stercus et vermis est. Hodie extollitur: et cras non invenietur.
\newline \ul{Lectiones feriales a tercio capitulo usque ad secundum Liber Machabeorum ad libitum.}
{\ubsubsection{Dominica Tercia}{sabbato-ante-terciam-dominicam-octobris-breviarium}}
\newline {\color{Red} Lectio Prima.}
\lettrine[lines=2]{\zallmancaps \color{Blue} F}{ratribus} qui sunt per Egiptum Iudeis, salutem dicunt fratres qui sunt in Hierosolimis Iudei, et qui in regione Iudea: et pacem bonam. Benefaciat vobis Deus, et meminerit testamenti sui quod ad Abraham et Ysaac et Iacob locutus est servorum suorum fidelium: et det vobis cor omnibus ut colatis eum et faciatis eius voluntatem corde magno et animo volenti.
\newline {\color{Red} Lectio Secunda.}
\lettrine[lines=2]{\zallmancaps \color{Red} A}{daperiat} cor vestrum in lege sua et in preceptis suis: et faciat pacem. Exaudiat orationes vestras et reconcilietur vobis, nec vos deserat in tempore malo.
\newline {\color{Red} Lectio Tertia.}
\lettrine[lines=2]{\zallmancaps \color{Blue} E}{t} nunc hic sumus, orantes pro vobis regnante Demetrio anno centesimo sexagesimo nono. Nos Iudei scripsimus vobis in tribulatione et impetu qui supervenit nobis in istis annis. Ex quo recessit Iason a sancta terra et a regno. Portam succenderunt, et effuderunt sanguinem innocentem.
\newline {\color{Red} Lectio Quarta.}
\lettrine[lines=2]{\zallmancaps \color{Red} E}{t} oravimus ad dominum et exauditi sumus, et obtulimus sacrificium, et similaginem, et accendimus lucernas, et proposuimus panes. Et nunc frequentate dies scenophegie mensis Casleu.
\newline {\color{Red} Lectio Quinta.}
\lettrine[lines=2]{\zallmancaps \color{Blue} A}{nno} centesimo octogesimo octavo, populus qui est Hierosolimis et in Iudea, senatusque et Iudas, Aristobolo magistro Ptolemei regis, qui est de genere christorum sacerdotum, et his qui in Egipto sunt Iudeis: salutem et sanitatem. De magnis periculis a Deo liberati, magnifice gratias agimus ipsi: utpote qui adversus talem regem dimicavimus.
\newline {\color{Red} Lectio Sexta.}
\lettrine[lines=2]{\zallmancaps \color{Red} I}{pse} enim ebullire fecit de Perside, eos qui pugnaverunt contra nos et sanctam civitatem. Nam cum in Perside esset dux ipse, et cum ipso immensus exercitus, cecidit in templo Nanee, consilio deceptus sacerdotis Nanee. Ut enim cum eo habitaturus venit ad locum Anthiochus, et amici eius: et ut acciperet pecunias multas dotis nomine.
\newline \ul{Lectiones feriales a secundo capitulo secundi Machabeorum usque ad sextum capitulum ad libitum.}
{\ubsubsection{Dominica Quarta}{sabbato-ante-quartam-dominicam-octobris-breviarium}}
\newline {\color{Red} Lectio Prima.}
\lettrine[lines=2]{\zallmancaps \color{Blue} C}{um} venisset Apollonius Hierosolimam, pacem simulans, quievit usque ad diem sanctum sabbati. Et tunc feriatis Iudeis, arma capere suis precepit, omnesque qui ad spectaculum processerant, trucidavit, et civitatem cum armatis discurrens: ingentem multitudinem peremit.
\newline {\color{Red} Lectio Secunda.}
\lettrine[lines=2]{\zallmancaps \color{Red} I}{udas} autem Machabeus qui decimus fuerat, secesserat in desertum locum, ibique inter feras vitam in montibus cum suis agebat. Et feni cibo vescentes demorabantur, ne participes essent coinquinationis.
\newline {\color{Red} Lectio Tertia.}
\lettrine[lines=2]{\zallmancaps \color{Blue} N}{on} post multum temporis misit rex senem quendam Anthiochenum, qui compelleret Iudeos ut se transferrent a patriis et Dei legibus, contaminare eciam quod in Hierosolimis erat templum, et cognominare Iovis Olimpii et in Garizim, prout erant hi qui locum inhabitabant Iovis hospitalis.
\newline {\color{Red} Lectio Quarta.}
\lettrine[lines=2]{\zallmancaps \color{Red} P}{essima} autem et universis gravis, malorum erat incursio. Nam templum luxuria et comessationibus erat plenum, et scortantium cum meretricibus: sacratisque edibus mulieres se ultro ingerebant, intro ferentes ea que non licebat.
\newline {\color{Red} Lectio Quinta.}
\lettrine[lines=2]{\zallmancaps \color{Blue} A}{ltare} eciam plenum erat illicitis que legibus prohibebantur. Neque autem sabbata custodiebantur, neque dies sollemnes patrii servabantur: nec simpliciter Iudeum quisquam se confitebatur. Ducebantur autem cum amara necessitate in die natalis regis ad sacrificia.
\newline {\color{Red} Lectio Sexta.}
\lettrine[lines=2]{\zallmancaps \color{Red} E}{t} cum Liberi sacra celebrarentur, cogebantur hedera coronati Libero circuire. Decretum autem exiit in proximas gentilium civitates suggerentibus Ptolemeis, ut pari modo et ipsi adversus Iudeos agerent ut sacrificarent: eosque qui nollent transire ad instituta gentilium: interficerent.
\newline \ul{Lectiones feriales a septimo capitulo usque ad nonum ad libitum.}
{\ubsubsection{Dominica Quinta}{sabbato-ante-quintam-dominicam-octobris-breviarium}}
\newline {\color{Red} Lectio Prima.}
\lettrine[lines=2]{\zallmancaps \color{Blue} A}{nthiochus} inhoneste revertebatur de Perside. Intraverat enim in eam que dicitur Persepolis. Et temptavit expoliare templa, et civitatem opprimere. Sed multitudine ad arma concurrente: in fugam versi sunt. Et ita contigit, ut Anthiochus post fugam turpiter rediret.
\newline {\color{Red} Lectio Secunda.}
\lettrine[lines=2]{\zallmancaps \color{Red} E}{t} cum venisset circa Egbathanan: recognovit que erga Nichanorem et Thimotheum gesta sunt. Elatus autem in ira, arbitrabatur se iniuriam illorum qui se fugaverant posse in Iudeos retorquere. Ideoque iussit agitari currum suum, sine intermissione agens iter, celesti eum iudicio perurgente: eo quod ita superbe locutus est venturum se Hierosolimam, et congeriem sepulchri Iudeorum eam facturum.
\newline {\color{Red} Lectio Tertia.}
\lettrine[lines=2]{\zallmancaps \color{Blue} S}{ed} qui universa conspicit dominus Deus Israhel: percussit eum insanabili et invisibili plaga. Ut enim finivit hunc ipsum sermonem, apprehendit eum dolor dirus viscerum, et amara internorum tormenta. Et quidem satis iuste: quippe qui multis et novis cruciatibus aliorum torserat viscera: licet ille nullo modo a sua malitia cessaret.
\newline {\color{Red} Lectio Quarta.}
\lettrine[lines=2]{\zallmancaps \color{Red} S}{uper} hec autem superbia repletus, ignem spirans animo in Iudeos, et precipiens accelerari negotium, contigit illum impetu euntem de curru cadere, et gravi corporis collisione membra vexari. Isque qui sibi videbatur eciam fluctibus maris imperare, supra humanum modum superbia repletus, et montium altitudines in statera appendere, nunc humiliatus ad terram in gestatorio portabatur, manifestam Dei virtutem in semetipso contestans, ita ut de corpore impii vermes scaturirent, ac viventis in doloribus carnes eius effluerent, odore eciam illius et fetore exercitus gravaretur.
\newline {\color{Red} Lectio Quinta.}
\lettrine[lines=2]{\zallmancaps \color{Blue} E}{t} qui paulo ante sidera celi contingere se arbitrabatur: eum nemo poterat propter intolerantiam fetoris portare. Hinc igitur cepit ex gravi superbia deductus ad agnitionem sui venire, divina admonitus plaga, per momenta singula doloribus suis augmenta capientibus.
\newline {\color{Red} Lectio Sexta.}
\lettrine[lines=2]{\zallmancaps \color{Red} E}{t} cum nec ipse iam fetorem suum ferre posset, sic ait. Iustum est subditum esse Deo: et mortalem non paria Deo sentire. Orabat autem hic scelestus dominum: a quo non esset misericordiam consecuturus. Et civitatem ad quam festinans veniebat, ut eam ad solum deduceret, ac sepulchrum congestorum faceret: nunc optat liberam reddere.
\newline \ul{Lectiones feriales a nono capitulo usque ad finem ad libitum.}
\lettrine[lines=2]{\zallmancaps \color{Blue} D}{}\ul{\textsc{ominica} prima Novembris et deinceps usque ad primam Dominicam Adventus exclusive quando de tempore agitur cantetur hystoria,} \hyperlink{vidi-dominum}{Vidi dominum.} \ul{et legantur Ezechiel et Daniel et duodecim prophete. Et notandum quod cum legitur Liber Ezechielis lectiones de primo capitulo per.} Tu autem. \ul{a secundo et deinceps per.} Hec dicit dominus. \ul{finiuntur.}
{\ubsubsection{Sabbato ante Primam Dominicam Novembris}{sabbato-ante-primam-dominicam-novembris-breviarium}}
\newline \ul{Sabbato vero precedenti ad Vesperas.} \R \hyperlink{aspice-domine-de}{Aspice domine de sede.} {\color{Red} II. Ad Mag. Ant.} Vidi dominum sedentem super solium excelsum, et plena erat omnis terra maiestate eius, et ea que sub ipso erant replebant templum.
\newline {\color{Red} Ad Matutinas. Lectio Prima.}
\lettrine[lines=2]{\zallmancaps \color{Red} E}{t} factum est in tricesimo anno, in quarto in quinta mensis: cum essem in medio captivorum iuxta fluvium Chobar: aperti sunt celi, et vidi visiones Dei. In quinta mensis, ipse est annus quintus transmigrationis regis Ioachim, factum est verbum domini ad Ezechielem filium Buzi sacerdotem in terra Chaldeorum, secus flumen Chobar. Tu autem. {\color{Red} Responsorium.}
\lettrine[lines=2]{\zallmancaps \color{Blue} V}{idi} \hypertarget{vidi-dominum}{\label{vidi-dominum}} dominum sedentem super solium excelsum et elevatum, et plena erat omnis terra maiestate eius, * Et ea que sub ipso erant replebant templum. \V Seraphim stabant super illud, sex ale uni et sex ale alteri. Et ea.
\newline {\color{Red} Lectio Secunda.}
\lettrine[lines=2]{\zallmancaps \color{Blue} E}{t} facta est super eum ibi manus domini. Et vidi, et ecce ventus turbinis veniebat ab aquilone, et nubes magna et ignis involvens: et splendor in circuitu eius. Et de medio eius quasi species electri: id est de medio ignis.
\begin{responsory}[aspice-domine-de]
{Aspice domine de sede sancta tua et cogita de nobis, inclina deus meus aurem tuam et audi, * Aperi oculos tuos et vide tribulationem nostram.}{Qui regis Israhel intende, qui deducis velut ovem Ioseph.}
\end{responsory}%
\newline {\color{Red} Lectio Tertia.}
\lettrine[lines=2]{\zallmancaps \color{Red} E}{t} in medio eius similitudo quatuor animalium, et hic aspectus eorum. Similitudo hominis in eis. Et quatuor facies uni, et quatuor penne uni. Et pedes eorum pedes recti: et planta, pedis eorum quasi planta pedis vituli, et scintille quasi aspectus eris candentis. Et manus hominis sub pennis eorum in quatuor partibus.
\begin{responsory-doxology}[aspice-domine-quia]
{Aspice domine quia facta est desolata civitas plena divitiis, sedet in tristicia domina gentium, * Non est qui consoletur eam nisi tu deus.}{Plorans ploravit in nocte et lacrime eius in maxillis eius.}
\end{responsory-doxology}%
\newline {\color{Red} Lectio Quarta.}
\lettrine[lines=2]{\zallmancaps \color{Blue} E}{t} facies et pennas per quatuor partes habebant, iuncteque erant penne eorum alterius ad alterum. Non revertebantur cum incederent, sed unumquodque ante faciem suam gradiebatur. Similitudo autem vultus eorum facies hominis, et facies leonis a dextris ipsorum quatuor, facies autem bovis a sinistris ipsorum quatuor: et facies aquile desuper ipsorum quatuor.
\begin{responsory}[super-muros]
{Super muros tuos Hierusalem constitui custodes tota die et nocte non tacebunt, * Laudantes nomen domini.}{Predicabunt populis fortitudinem meam et annuntiabunt in gentibus gloriam meam.}
\end{responsory}%
\newline {\color{Red} Lectio Quinta.}
\lettrine[lines=2]{\zallmancaps \color{Red} E}{t} facies eorum et penne eorum extente desuper. Due penne singulorum iungebantur: et due tegebant corpora eorum. Et unumquodque eorum coram facie sua ambulabat. Ubi erat impetus spiritus, illuc gradiebantur: nec revertebantur cum ambularent.
\begin{responsory}[muro-tuo]
{Muro tuo inexpugnabili circumcinge nos domine, et armis tue potentie protege nos semper, * Libera domine deus Israhel clamantes ad te.}{Erue nos in mirabilibus tuis et da gloriam nomini tuo.}
\end{responsory}%
\newline {\color{Red} Lectio Sexta.}
\lettrine[lines=2]{\zallmancaps \color{Blue} S}{imilitudo} animalium et aspectus eorum quasi carbonum ignis ardentium et quasi aspectus lampadarum. Hec erat visio discurrens in medio animalium, splendor ignis, et de igne fulgur egrediens. Et animalia ibant et revertebantur, in similitudinem fulguris choruscantis.
\begin{responsory-doxology}[sustinuimus-pacem]
{Sustinuimus pacem et non venit, quesivimus bona et ecce turbatio, cognovimus domine peccata nostra, * Non in perpetuum obliviscaris nos.}{Peccavimus cum patribus nostris, iniuste egimus, iniquitatem fecimus.}
\end{responsory-doxology}%
\begin{responsory}[angustie-michi]
{Angustie michi sunt undique, et quid eligam ignoro, * Melius est michi incidere in manus hominum quam derelinquere legem dei mei.}{Si enim hoc egero mors michi est, si autem non egero non effugiam manus vestras.}
\end{responsory}%
\begin{responsory}[misit-dominus]
{Misit dominus angelum suum et conclusit ora leonum et non me contaminaverunt, * Quia coram eo iusticia inventa est in me.}{Misit deus misericordiam suam et veritatem suam, animam meam eripuit de medio catulorum leonum.}
\end{responsory}%
\begin{responsory-doxology}[laudabilis-populis]
{Laudabilis populis quem dominus exercituum benedixit dicens, opus manuum mearum tu es, * Hereditas mea Israhel.}{Beata gens cuius est dominus deus eius populus electus in hereditatem.}
\end{responsory-doxology}%
\newline \ul{Sequentes antiphone dicende sunt in Sabbatis, Ad Mag. vel ad memoriam quamdiu cantabitur hystoria,} \hyperlink{vidi-dominum}{Vidi dominum.} {\color{Red} Ant.} Muro tuo inexpugnabili circumcinge nos domine et armis tue potentie protege nos semper deus noster. {\color{Red} Ant.} Qui celorum contines thronos et abyssos intueris domine rex regum, montes ponderas, terram palmo concludis, exaudi nos domine in gemitibus nostris. {\color{Red} Ant.} Sustinuimus pacem et non venit domine, quesivimus bona et ecce turbatio, cognovimus domine peccata nostra ne in eternum obliviscaris nos deus Israhel. {\color{Red} Ant.} Aspice domine quia facta est desolata civitas plena divitiis, sedet in tristicia domina gentium, non est qui consoletur eam nisi tu deus noster.
\newline \ul{Lectiones feriales a secundo capitulo Ezechielis usque ad septimum ad libitum.}
{\ubsubsection{Dominica Secunda}{sabbato-ante-secundam-dominicam-novembris-breviarium}}
\newline {\color{Red} Lectio Prima.}
\lettrine[lines=2]{\zallmancaps \color{Red} E}{t} factus est sermo domini ad me dicens. Et tu fili hominis, hec dicit dominus Deus terre Israhel. Finis venit, venit finis super quatuor plagas terre. Nunc finis super te, immittam furorem meum in te. Et iudicabo te iuxta vias tuas, et ponam contra te omnes abominationes tuas. Hec dicit.
\newline {\color{Red} Lectio Secunda.}
\lettrine[lines=2]{\zallmancaps \color{Blue} E}{t} non parcet oculus meus super te, et non miserebor, sed vias tuas ponam super te, et abominationes tue in medio tui erunt. Et scietis: quia ego dominus. Hec dicit dominus Deus. Afflictio una afflictio: ecce venit. Finis venit, venit finis. Evigilavit adversum te. Ecce venit. Venit contractio super te, qui habitas in terra. Hec dicit.
\newline {\color{Red} Lectio Tertia.}
\lettrine[lines=2]{\zallmancaps \color{Red} V}{enit} tempus, prope est dies occisionis et non glorie montium. Nunc de propinquo effundam iram meam super te, et complebo furorem meum in te. Et iudicabo te iuxta vias tuas, et imponam tibi omnia scelera tua: et non parcet oculus meus nec miserebor, sed vias tuas imponam tibi, et abhominationes tue in medio tui erunt: et scietis quia ego sum dominus percutiens. Hec dicit.
\newline {\color{Red} Lectio Quarta.}
\lettrine[lines=2]{\zallmancaps \color{Blue} E}{cce} dies: ecce venit. Egressa est contractio. Floruit virga, germinavit superbia: iniquitas surrexit in virga impietatis. Non ex eis, et non ex populo, neque ex sonitu eorum: et non erit requies in eis. Hec dicit.
\newline {\color{Red} Lectio Quinta.}
\lettrine[lines=2]{\zallmancaps \color{Red} V}{enit} tempus: appropinquavit dies. Qui emit non letetur, et qui vendit non lugeat, quia ira super omnem populum eius. Quia qui vendit ad id quod vendidit non revertetur: et adhuc in viventibus vita eorum. Visio enim ad omnem multitudinem eius non regredietur, et vir in iniquitate vite sue non confortabitur. Hec dicit.
\newline {\color{Red} Lectio Sexta.}
\lettrine[lines=2]{\zallmancaps \color{Blue} C}{anite} tuba, preparentur omnes, et non est qui vadat ad prelium. Ira enim mea super universum populum eius. Gladium foris, pestis et fames intrinsecus. Qui in agro est gladio morietur, et qui in civitate, pestilentia et fame devorabuntur. Et salvabuntur qui fugerint ex eis: et erunt in montibus quasi columbe convallium, omnes trepidi, unusquisque in iniquitate sua. Hec dicit.
\newline \ul{Lectiones feriales ab octavo capitulo usque ad finem ad libitum.}
{\ubsubsection{Dominica Tercia}{sabbato-ante-terciam-dominicam-novembris-breviarium}}
\newline {\color{Red} Lectio Prima.}
\lettrine[lines=2]{\zallmancaps \color{Red} A}{nno} tertio regni Ioachim regis Iuda venit Nabucchodonosor rex Babilonis Hierusalem, et obsedit eam, tradidit dominus in manu eius Ioachim regem Iuda, et partem vasorum domus Dei. Et asportavit ea in terram Sennaar in domum dei sui, et vasa intulit in domum thesauri dei sui. Tu autem.
\newline {\color{Red} Lectio Secunda.}
\lettrine[lines=2]{\zallmancaps \color{Blue} E}{t} ait rex Asfanaz preposito eunuchorum suorum, ut introduceret de filiis Israhel, et de semine regio et tyrannorum pueros, in quibus nulla esset macula, decoros forma: et eruditos omni sapientia, cautos scientia et doctos disciplina, et qui possent stare in palatio regis, ut doceret eos litteras et linguam Chaldeorum. Tu.
\newline {\color{Red} Lectio Tertia.}
\lettrine[lines=2]{\zallmancaps \color{Red} E}{t} constituit eis rex annonam per singulos dies de cibis suis, et de vino unde bibebat ipse: ut enutriti tribus annis, postea starent in conspectu regis. Fuerunt ergo inter eos de filiis Iuda, Daniel, Ananias, Misael, et Azarias. Et imposuit eis prepositus eunuchorum nomina, Danieli Balthasar, Ananie Sidrac, Misaeli Misach, et Azarie Abdenago. Tu autem.
\newline {\color{Red} Lectio Quarta.}
\lettrine[lines=2]{\zallmancaps \color{Blue} P}{roposuit} autem Daniel in corde suo ne pollueretur de mensa regis, neque de vino potus eius. Et rogavit eunuchorum prepositum ne contaminaretur. Dedit autem Deus Danieli gratiam et misericordiam in conspectu principis eunuchorum. Tu autem.
\newline {\color{Red} Lectio Quinta.}
\lettrine[lines=2]{\zallmancaps \color{Red} E}{t} ait princeps eunuchorum ad Danielem. Timeo ego dominum meum regem: qui constituit vobis cibum et potum. Qui si viderit vultus vestros macilentiores pre ceteris adolescentibus coevis vestris condemnabitis caput meum regi. Tu autem.
\newline {\color{Red} Lectio Sexta.}
\lettrine[lines=2]{\zallmancaps \color{Blue} E}{t} dixit Daniel ad Malasar, quem constituerat princeps eunuchorum super Danielem Ananiam, Misaelem, et Azariam. Tempta nos obsecro servos tuos diebus decem, et dentur nobis legumina ad vescendum et aqua ad bibendum, et contemplare vultus nostros et vultus puerorum qui vescuntur cibo regio: et sicut videris sic facias cum servis tuis. Tu autem.
\newline \ul{Lectiones feriales a secundo capitulo Danielis usque ad finem ad libitum.}
{\ubsubsection{Dominica Quarta}{sabbato-ante-quartam-dominicam-novembris-breviarium}}
\newline {\color{Red} Lectio Prima.}
\lettrine[lines=2]{\zallmancaps \color{Red} V}{erbum} domini quod factum est ad Osee filium Beeri in diebus Ozie, Ioathan, Achaz, Ezechie regum Iuda: et in diebus Ieroboam filii Ioas regis Israhel. Principium loquendi domino in Osee. Et dixit dominus ad Osee. Vade sume tibi uxorem fornicationum, et filios fornicationum, quia fornicans fornicabitur terra a domino. Et abiit, et accepit Gomer filiam Debelaim. Et concepit: et peperit ei filium. Hec dicit.
\newline {\color{Red} Lectio Secunda.}
\lettrine[lines=2]{\zallmancaps \color{Blue} E}{t} dixit dominus ad eum. Voca nomen eius Iezrael: quoniam adhuc modicum et visitabo sanguinem Iezrael super domum Hieu: et quiescere faciam regnum domus Israhel. Et in illa die conteram arcum Israhel in valle Iezrael. Et concepit adhuc, et peperit filiam. Et dixit ei. Voca nomen eius absque misericordia, quia non addam ultra misereri domui Israhel: sed oblivione obliviscar eorum. Et domui Iuda miserebor, et salvabo eos in domino Deo suo, et non salvabo eos in arcu et gladio et in bello, et in equis et in equitibus. Hec dicit.
\newline {\color{Red} Lectio Tertia.}
\lettrine[lines=2]{\zallmancaps \color{Red} E}{t} ablactavit eam que erat absque misericordia. Et concepit, et peperit filium. Et dixit. Voca nomen eius non populus meus: quia vos non populus meus, et ego non ero vester. Et erit numerus filiorum Israhel quasi harena maris que sine mensura est: et non numerabitur. Et erit, in loco ubi dicetur eis, non populus
%pp115
meus vos: dicetur eis, filii Dei viventis. Et congregabuntur filii Iuda et filii Israhel pariter, et ponent sibimet caput unum, et ascendent de terra: quia magnus dies Iezrael. Hec dicit.
\newline {\color{Red} Lectio Quarta.}
\lettrine[lines=2]{\zallmancaps \color{Blue} D}{icite} fratribus vestris populus meus, et sorori vestre misericordiam consecuta. Iudicate matrem vestram iudicate: quoniam ipsa non uxor mea, et ego non vir eius. Auferat fornicationes suas a facie sua, et adulteria sua de medio uberum suorum, ne forte expoliem eam nudam, et statuam eam secundum diem nativitatis sue. Hec dicit.
\newline {\color{Red} Lectio Quinta.}
\lettrine[lines=2]{\zallmancaps \color{Red} E}{t} ponam eam quasi solitudinem, et statuam eam velut terram inviam, et interficiam eam siti. Et filiorum illius non miserebor, quoniam filii fornicationum sunt: quia fornicata est mater eorum. Confusa est que concepit eos: quia dixit. Vadam post amatores meos, qui dant panes michi, et aquas meas, lanam meam, et linum meum, oleum meum et potum meum. Hec dicit.
\newline {\color{Red} Lectio Sexta.}
\lettrine[lines=2]{\zallmancaps \color{Blue} P}{ropter} hoc, ecce ego sepiam viam tuam spinis, et sepiam eam maceria, et semitas suas non inveniet. Et sequetur amatores suos et non apprehendet eos: et queret eos et non inveniet. Et dicet. Vadam et revertar ad virum meum priorem, quia bene michi erat tunc magis quam nunc. Hec dicit.
\newline \ul{Lectiones feriales a tercio capitulo Osee usque ad Micheam ad libitum.}
{\ubsubsection{Dominica Quinta}{sabbato-ante-quintam-dominicam-novembris-breviarium}}
\newline {\color{Red} Lectio Prima.}
\lettrine[lines=2]{\zallmancaps \color{Red} V}{erbum} domini quod factum est ad Micheam Morastiten in diebus Ioathan, Achaz, Ezechie regum Iuda, quod vidit super Samariam et Hierusalem. Audite populi omnes, et attendat terra et plenitudo eius, et sit dominus Deus vobis in testem, dominus de templo sancto suo. Quia ecce dominus egredietur de loco suo: et descendet et calcabit super excelsa terre. Hec dicit.
\newline {\color{Red} Lectio Secunda.}
\lettrine[lines=2]{\zallmancaps \color{Blue} E}{t} consumentur montes subtus eum, et valles scindentur sicut cera a facie ignis: et sicut aque que decurrunt in preceps. In scelere Iacob omne istud, et in peccatis domus Israhel. Quod scelus Iacob? Nonne Samaria? Et que excelsa Iude? Nonne Hierusalem? Hec dicit.
\newline {\color{Red} Lectio Tertia.}
\lettrine[lines=2]{\zallmancaps \color{Red} E}{t} ponam Samariam quasi acervum lapidum in agro cum plantatur vinea. Et detraham lapides eius in vallem: et fundamenta eius revelabo. Et omnia sculptilia eius concidentur: et omnes mercedes eius comburentur igni. Et omnia ydola eius ponam in perditionem, quia de mercedibus meretricis congregata sunt: et usque ad mercedem meretricis revertentur. Hec dicit.
\newline {\color{Red} Lectio Quarta.}
\lettrine[lines=2]{\zallmancaps \color{Blue} S}{uper} hoc plangam et ululabo. Vadam spoliatus et nudus. Faciam planctum velut drachonum, et luctum quasi strutionum: quia desperata est plaga eius, quia venit usque ad Iudam, tetigit portam populi mei usque ad Hierusalem. Hec dicit.
\newline {\color{Red} Lectio Quinta.}
\lettrine[lines=2]{\zallmancaps \color{Red} I}{n} Geth nolite annuntiare: lacrimis ne ploretis. In domo pulveris, pulvere vos conspergite. Et transite vobis habitatio pulchra, confusa ignominia. Non est egressa que habitat in exitu. Planctum domus vicina accipiet ex vobis que stetit sibimet, quia infirmata est in bonum que habitat in amaritudinibus, quia descendit malum a domino in portam Hierusalem, tumultus quadrige stuporis habitanti Lachis. Hec dicit.
\newline {\color{Red} Lectio Sexta.}
\lettrine[lines=2]{\zallmancaps \color{Blue} P}{rincipium} peccati est filie Syon, quia in te inventa sunt scelera Israhel. Propterea dabit emissarios super hereditatem Geth, domus mendacii in deceptionem regibus Israhel. Adhuc heredem adducam tibi que habitas in Maresa: usque ad Odollam veniet gloria Israhel. Hec dicit.
\newline \ul{Lectiones feriales a primo capitulo Zacharie usque ad finem Malachie ad libitum.}
{\ubsubsection{Serie Evangeliorum Cum Orationibus}{serie-evangeliorum-cum-orationibus-breviarium}}
\newline {\color{Red} Dominica Prima post festum Trinitatis. Oratio.}
\lettrine[lines=2]{\zallmancaps \color{Red} D}{eus} in te sperantium fortitudo, adesto propitius invocationibus nostris, et quia sine te nichil potest mortalis infirmitas, presta auxilium gratie tue, ut in exequendis mandatis tuis et voluntate tibi et actione placeamus. Per.
\newline {\color{Red} Lectio VII. Secundum Lucam.}
\lettrine[lines=2]{\zallmancaps \color{Blue} I}{n} illo tempore: Dicebat Ihesus turbis Iudeorum et Phariseorum. Homo quidam erat dives, et induebatur purpura et bisso, et epulabatur cotidie splendide. Et reliqua.
\newline {\color{Red} Omelia Beati Gregorii Pape.} \grecross
\lettrine[lines=2]{\zallmancaps \color{Red} Q}{uem} fratres karissimi, quem dives iste qui induebatur purpura et bisso et epulabatur cotidie splendide, nisi Iudaicum populum significat, qui cultum vite exterius habuit, qui accepte legis deliciis ad nitorem usus est non ad utilitatem? Quem vero Lazarus ulceribus plenus, nisi gentilem populum figuraliter exprimit, qui dum conversus ad Deum confiteri peccata sua non erubuit, huic vulnus in cute fuit.
\newline {\color{Red} Lectio VIII.}
\lettrine[lines=2]{\zallmancaps \color{Blue} I}{n} cutis quippe vulnere virus a visceribus trahitur, et foras erumpit. Quid est ergo peccatorum confessio, nisi quedam vulnerum ruptio? Quia peccati virus salubriter aperitur in confessione, quod pestifere latebat in mente. Vulnera etenim cutis in superficie trahunt humorem putredinis. Et confitendo peccata quid aliud agimus, nisi malum quod in nobis latebat aperimus?
\newline {\color{Red} Lectio IX.}
\lettrine[lines=2]{\zallmancaps \color{Red} S}{ed} Lazarus vulneratus cupiebat saturari de micis que cadebant de mensa divitis, et nemo illi dabat: quia gentilem quemque ad cognitionem legis admittere, superbus ille populus despiciebat. Qui dum doctrinam legis non ad caritatem habuit, sed ad elationem, quasi de acceptis opibus tumuit. Et quia ei verba defluebant de scientia: quasi mice cadebant de mensa.
\newline {\color{Red} Ad Ben. Ant.} Homo quidam erat dives et induebatur purpura et bisso et epulabatur cotidie splendide, et erat quidam mendicus nomine Lazarus, qui iacebat ad ianuam eius ulceribus plenus cupiens saturari de micis que cadebant de mensa divitis, et nemo illi dabat, sed et canes veniebant et lingebant ulcera eius alleluya.
\newline {\color{Red} Ad Mag. Ant.} Pater Abraham miserere mei et mitte Lazarum ut intingat extremum digiti sui in aquam, ut refrigeret linguam meam.
\newline {\color{Red} Dominica Secunda. Oratio.}
\lettrine[lines=2]{\zallmancaps \color{Blue} S}{ancti} nominis tui domine timorem pariter et amorem fac nos habere perpetuum, quia nunquam tua gubernatione destituis, quos in solidate tue dilectionis instituis. Per.
\newline {\color{Red} Lectio VII. Secundum Lucam.}
\lettrine[lines=2]{\zallmancaps \color{Red} I}{n} illo tempore: Dixit Ihesus discipulis suis similitudinem hanc. Homo quidam fecit cenam magnam: et vocavit multos. Et reliqua.
\newline {\color{Red} Omelia Beati Gregorii Pape.} \grecross
\lettrine[lines=2]{\zallmancaps \color{Blue} Q}{uis} est iste homo fratres karissimi, nisi ille de quo per prophetam dicitur: homo est, et quis cognoscet eum? Qui fecit cenam magnam: quia satietatem nobis interne dulcedinis preparavit. Qui vocavit multos, sed pauci veniunt, quia nonnunquam ipsi qui ei per fidem subiecti sunt: eterno eius convivio male vivendo contradicunt.
\newline {\color{Red} Lectio VIII.}
\lettrine[lines=2]{\zallmancaps \color{Red} M}{isit} autem servum suum hora cene, dicere invitatis ut venirent. Quid hora cene, nisi finis est mundi? In quo nimirum nos sumus, sicut iam dudum Paulus testatur dicens. Nos sumus in quos fines seculorum devenerunt. Si ergo hora cene iam est cum vocamur, tanto minus nos excusare debemus a convivio Dei, quanto propinquasse iam cernimus finem mundi. Quo enim pensamus quia nichil est quod restat, eo debemus pertimescere, ne tempus gratie quod presto est pereat.
\newline {\color{Red} Lectio IX.}
\lettrine[lines=2]{\zallmancaps \color{Blue} I}{dcirco} autem hoc convivium Dei non prandium sed cena nominatur, quia post prandium cena, post cenam vero convivium nullum restat. Et quia eternum Dei convivium nobis in extremo preparabitur, rectum fuit ut hoc non prandium sed cena vocaretur. Sed quis per hunc servum qui a patrefamilias ad invitandum mittitur, nisi predicatorum ordo designatur?
\newline {\color{Red} Ad Ben. Ant.} Homo quidam fecit cenam magnam et vocavit multos, et misit servum suum hora cene dicere invitatis ut venirent, quia omnia parata sunt alleluya.
\newline {\color{Red} Ad Mag. Ant.} Exi cito in plateas et vicos civitatis et pauperes ac debiles, et cecos et claudos compelle intrare ut impleatur domus mea alleluya.
\newline {\color{Red} Dominica Tercia. Oratio.}
\lettrine[lines=2]{\zallmancaps \color{Red} D}{eprecationem} nostram quesumus domine benignus exaudi et quibus supplicandi prestas affectum, tribue defensionis auxilium. Per.
\newline {\color{Red} Lectio VII. Secundum Lucam.}
\lettrine[lines=2]{\zallmancaps \color{Blue} I}{n} illo tempore: Erant appropinquantes ad Ihesum publicani et peccatores ut audirent illum. Et murmurabant Pharisei et scribe dicentes, quia hic peccatores recipit et manducat cum illis. Et reliqua.
\newline {\color{Red} Omelia Beati Gregorii Pape.} \grecross
\lettrine[lines=2]{\zallmancaps \color{Red} A}{udistis} ex lectione evangelica fratres mei, quia peccatores et publicani accesserunt ad redemptorem nostrum, et non solum ad colloquendum sed eciam ad convescendum recepti sunt. Quod videntes Pharisei dedignati sunt. Ex qua re colligite, quia vera iusticia compassionem habet, falsa iusticia dedignationem, quamvis et iusti recte soleant peccatoribus indignari. Sed aliud est quod agitur typo superbie: aliud quod zelo discipline.
\newline {\color{Red} Lectio VIII.}
\lettrine[lines=2]{\zallmancaps \color{Blue} D}{edignantur} iusti, sed non dedignantes. Desperant: sed non desperantes. Persecutionem commovent: sed amantes. Quia et si foris increpationes per disciplinam exaggerant: intus tamen dulcedinem per caritatem servant. Preponunt sibi in animo iusti, ipsos plerumque quos corrigunt, meliores existimant eos quoque quos iudicant. Quod videlicet agentes: et per disciplinam subditos, et per humilitatem custodiunt semetipsos.
\newline {\color{Red} Lectio IX.}
\lettrine[lines=2]{\zallmancaps \color{Red} A}{t} contra hi qui de falsa iusticia superbire solent, ceteros quosque despiciunt, nulla infirmantibus misericordia condescendunt. Et quo se peccatores esse non credunt, eo deterius peccatores fiunt. De quorum profecto numero Pharisei extiterant: qui diiudicantes dominum quod peccatores susciperet: arenti corde ipsum fontem misericordie reprehendebant.
\newline {\color{Red} Ad Ben. Ant.} Quis ex vobis homo qui habet centum oves, et si perdiderit unam ex illis nonne dimittit nonaginta novem in deserto, et vadit ad illam que perierat donec inveniat illam alleluya.
\newline {\color{Red} Ad Mag. Ant.} Que mulier habens dragmas decem, et si perdiderit dragmam unam, nonne accendit lucernam, et everrit domum, et querit diligenter donec inveniat?
\newline {\color{Red} Dominica Quarta. Oratio.}
\lettrine[lines=2]{\zallmancaps \color{Blue} P}{rotector} in te sperantium deus, sine quo nichil est validum nichil sanctum, multiplica super nos misericordiam tuam, ut te rectore te duce sic transeamus per bona temporalia, ut non amittamus eterna. Per.
\newline {\color{Red} Lectio VII. Secundum Lucam.}
\lettrine[lines=2]{\zallmancaps \color{Red} I}{n} illo tempore: Dixit Ihesus discipulis suis. Estote misericordes: sicut et pater vester misericors est. Nolite iudicare, et non iudicabimini. Et reliqua.
\newline {\color{Red} Omelia Venerabilis Bede Presbiteri.}
\lettrine[lines=2]{\zallmancaps \color{Blue} B}{enignus} est Deus super ingratos et malos, vel multiplici scilicet sua misericordia, qua eciam iumenta salvat temporalia bona largiendo, vel celestia dona singulari gratia qua electos suos solum glorificat inspirando. Sed sive hoc, sive illud, sive utrumque intelligas, magna Dei bonitate fit, que nobis imitanda precipitur si filii Dei esse volumus.
\newline {\color{Red} Lectio VIII.}
\lettrine[lines=2]{\zallmancaps \color{Red} N}{olite} inquit iudicare, et non iudicabimini. Nolite condemnare et non condemnabimini. Hoc loco nichil aliud precipi existimo, nisi ut ea facta que dubium est quo animo fiant: in meliorem partem interpretemur. Quod autem scriptum est a fructibus eorum cognoscetis eos: de manifestis dictum est, que non possunt bono animo fieri, sicuti sunt stupra, vel blasphemie, vel furta, vel ebrietates, et si qua sunt talia de quibus nobis iudicare permittitur.
\newline {\color{Red} Lectio IX.}
\lettrine[lines=2]{\zallmancaps \color{Blue} D}{e} genere autem ciborum, quia possunt bono animo et simplici corde sine vitio concupiscientie quicunque humani cibi indifferenter sumi, prohibet apostolus iudicari eos qui carnibus vescebantur, et bibebant vinum: ab eis qui se ab huiusmodi alimentis temperabant. Qui enim manducat inquit, non manducantem non spernat, et qui non manducat, manducantem non iudicet.
\newline {\color{Red} Ad Ben. Ant.} Estote ergo misericordes sicut et pater vester misericors est dicit dominus.
\newline {\color{Red} Ad Mag. Ant.} Nolite iudicare et non iudicabimini, in quo enim iudicio iudicaveritis iudicabimini dicit dominus.
\newline {\color{Red} Dominica Quinta. Oratio.}
\lettrine[lines=2]{\zallmancaps \color{Red} D}{a} nobis quesumus domine, ut et mundi cursus pacifice nobis tuo ordine dirigatur, et ecclesia tua tranquilla devotione letetur. Per.
\newline {\color{Red} Lectio VII. Secundum Lucam.}
\lettrine[lines=2]{\zallmancaps \color{Blue} I}{n} illo tempore: Cum turbe irruerunt ad Ihesum, ut audirent verbum dei: et ipse stabat secus stagnum Genesareth. Et reliqua.
\newline {\color{Red} Omelia Venerabilis Bede Presbiteri.}
\lettrine[lines=2]{\zallmancaps \color{Red} S}{tagnum} Genesareth idem dicunt esse quod mare Galilee vel mare Tiberiadis. Sed mare Galilee ab adiacente provincia dictum est, mare Tiberiadis a proxima civitate, que olim Cenereth vocata, sed ab Herode tetrarcha instaurata, in honorem Tiberii Cesaris: Tiberias est appellata.
\newline {\color{Red} Lectio VIII.}
\lettrine[lines=2]{\zallmancaps \color{Blue} P}{orro} Genesar a laci ipsius natura, qua crispantibus aquis de seipso excitare auram perhibetur, Greco vocabulo quasi generans sibi auram dicitur. Neque enim in stagni morem sternitur aqua, sed frequentibus auris spirantibus agitatur, haustu dulcis, et ad potandum habilis. Sed Hebraice lingue consuetudine, omnis aquarum congregatio, sive dulcis sive salsa, mare nuncupatur.
\newline {\color{Red} Lectio IX.}
\lettrine[lines=2]{\zallmancaps \color{Red} Q}{uia} ergo stagnum sive mare presens seculum designat, dominus secus mare stat, postquam vite labentis mortalitatem devincens in ea qua passus est carne stabilitatem perpetue quietis adiit. Turbarum conventus ad eum: gentium in fide concurrentium typus est, de quibus Ysaias, et fluent inquit ad eum omnes gentes, et ibunt populi multi.
\newline {\color{Red} Ad Ben. Ant.} Ascendit Ihesus in navim et sedens docebat turbas alleluya.
\newline {\color{Red} Ad Mag. Ant.} Preceptor per totam noctem laborantes nichil cepimus, in verbo autem tuo laxabo rete.
\newline {\color{Red} Dominica Sexta. Oratio.}
\lettrine[lines=2]{\zallmancaps \color{Blue} D}{eus} qui diligentibus te bona invisibilia preparasti, infunde cordibus nostris tui amoris affectum: ut te in omnibus et super omnia diligentes, promissiones tuas que omne desiderium superant consequamur. Per.
\newline {\color{Red} Lectio VII. Secundum Matheum.}
\lettrine[lines=2]{\zallmancaps \color{Red} I}{n} illo tempore: Dixit Ihesus discipulis suis. Amen dico vobis: nisi abundaverit iusticia vestra plus quam scribarum et Phariseorum: non intrabitis in regnum celorum. Et reliqua.
\newline {\color{Red} Omelia Venerabilis Bede Presbiteri.} \grecross
% la.wikisource.org/wiki/De_Sermone_Domini_in_monte
\lettrine[lines=2]{\zallmancaps \color{Blue} A}{udistis} fratres karissimi dominum in evangelio discipulos admonentem, et in illis omnes nos, nisi inquientem abundaverit iustitia vestra plus quam scribarum et phariseorum: non intrabitis in regnum celorum. Iustitia namque scribarum et phariseorum est, ut non occidant: iustitia vero illorum qui intraturi sunt in regnum celorum, ut non irascantur sine causa.
\newline {\color{Red} Lectio VIII.}
\lettrine[lines=2]{\zallmancaps \color{Red} Q}{uapropter} qui docet ut non irascamur non solvit legem ne occidamus, sed implet potius: ut et foris dum non occidimus, et in corde dum non irascimur, innocentiam custodiamus. Audistis inquit, quia dictum est antiquis. Non occides. Qui autem occiderit: reus erit iudicio. Ego autem dico vobis: quia omnis qui irascitur fratri suo, reus erit iudicio. Sed qui interest inter reum iudicio, et reum concilio, et reum gehenne ignis? Non enim reatus haberet gradus, nisi eciam gradatim peccata commemorarentur.
\newline {\color{Red} Lectio IX.}
\lettrine[lines=2]{\zallmancaps \color{Blue} G}{radus} igitur sunt in istis peccatis. In primo unum est, idest ira sola. In secundo duo: ira scilicet et vox que iram significat. In tertio vero tria, ira scilicet, et vox que iram significat: et in voce ipsa vituperationis expressio. Unde eciam trium reatuum, iudicii, concilii, gehenne ignis, diversitatis aliquid interesse, fateri cogit ipsa distinctio.
\newline {\color{Red} Ad Ben. Ant.} Audistis quia dictum est antiquis, non occides, qui autem occiderit reus erit iudicio.
\newline {\color{Red} Ad Mag. Ant.} Si offers munus tuum ad altare et recordatus fueris quia frater tuus habet aliquid adversum te, relinque ibi munus tuum ante altare et vade prius reconciliari fratri tuo, et tunc veniens offeres munus tuum alleluya.
\newline {\color{Red} Dominica Septima. Oratio.}
\lettrine[lines=2]{\zallmancaps \color{Red} D}{eus} virtutem cuius est totum quod est optimum, insere pectoribus nostris amorem tui nominis: et presta in nobis religionis augmentum, ut que sunt bona nutrias, ac pietatis studio que sunt nutrita custodias. Per.
\newline {\color{Red} Lectio VII. Secundum Marcum.}
\lettrine[lines=2]{\zallmancaps \color{Blue} I}{n} illo tempore: Cum turba multa esset cum Ihesu, nec haberent quod manducarent: convocatis discipulis, ait illis. Misereor super turbam quia ecce iam triduo sustinent me, nec habent quod manducent. Et reliqua.
\newline {\color{Red} Omelia Venerabilis Bede Presbiteri.} \grecross
\lettrine[lines=2]{\zallmancaps \color{Red} I}{n} hac lectione consideranda est in uno eodemque redemptore nostro distincta operatio divinitatis et humanitatis, atque Euthycetis error, qui unam tantum in Christo operationem dogmatizare presumpsit, procul a Christianis finibus expellendus. Quis enim non videat hoc quod super turbam miseretur dominus, ne vel inedia vel vie longioris labore deficiat: affectum esse et compassionem humane fragilitatis, quod vero de septem panibus et pisciculis paucis quatuor milia hominum saturavit, divine opus esse virtutis?
\newline {\color{Red} Lectio VIII.}
\lettrine[lines=2]{\zallmancaps \color{Blue} M}{ystice} autem hoc miraculo designatur, quod viam presentis seculi aliter incolumes transire nequimus: nisi nos gratia redemptoris nostri alimento verbi sui reficiat. Hoc vero typice inter hanc refectionem et illam quinque panum ac duorum piscium distat: quod ibi littera veteris instrumenti spiritali gratia plena esse signata est, hic autem novi veritas ac gratia testamenti fidelibus ministranda monstrata est.
\newline {\color{Red} Lectio IX.}
\lettrine[lines=2]{\zallmancaps \color{Red} S}{ane} utraque refectio in monte celebrata est, ut aliorum evangelistarum narratio declarat, quia utriusque scriptura testamenti recte intellecta: altitudinem nobis celestium preceptorum mandat et premiorum: utraque altitudinem Christi qui est mons domus domini in vertice montium, consona voce predicat.
\newline {\color{Red} Ad Ben. Ant.} Misereor super turbam, quia ecce iam triduo sustinent me, nec habent quod manducent, et si dimisero eos ieiunos deficient in via, alleluya.
\newline {\color{Red} Ad Mag. Ant.} Et accipiens Ihesus septem panes gratias agens benedixit, et fregit et dedit discipulis suis ut apponerent, et apposuerunt turbe alleluya.
\newline {\color{Red} Dominica Octava. Oratio.}
\lettrine[lines=2]{\zallmancaps \color{Blue} D}{eus} cuius providentia in sui dispostione non fallitur, te suplices exoramus, ut noxia cuncta summoveas, et omnia nobis profutura concedas. Per.
\newline {\color{Red} Lectio VII. Secundum Matheum.}
\lettrine[lines=2]{\zallmancaps \color{Red} I}{n} illo tempore: Dixit Ihesus discipulis suis. Attendite a falsis prophetis qui veniunt ad vos in vestimentis ovium, intrinsecus autem sunt lupi rapaces. Et reliqua.
\newline {\color{Red} Omelia Origenis.}
% https://la.wikisource.org/wiki/Commentarium_in_Matthaeum_(Rabanus_Maurus)
\lettrine[lines=2]{\zallmancaps \color{Blue} Q}{uod} paulo superius spatiosam et latam viam dominus nominavit, hoc nunc apertius falsos prophetas ostendit, per quos multi in perditionem abhominabilem abierunt: quia primo in Iudea multi apparuerunt, et modo perfidia sua totum repleverunt mundum.
\newline {\color{Red} Lectio VIII.}
\lettrine[lines=2]{\zallmancaps \color{Red} S}{ed} illi qui prius falsi prophete fuerunt, verissimos domini prophetas usque ad mortem persecuti sunt, sicut Hieremiam, et Micheam aliosque multos. Isti autem nunc falsi prophete et falsi Christiani qui sunt, veraces Christianos sine misericordia persequuntur, et
%pp116
opprimunt aliquando si detur copia eciam gladiis, sine intermissione autem suis pravis moribus et exemplis corrumpunt. Idcirco omnes preveniens dominus adhortatus est dicens. Attendite a falsis prophetis.
\newline {\color{Red} Lectio IX.}
\lettrine[lines=2]{\zallmancaps \color{Blue} A}{ttendite} diligentius, observate cautius, ut non seducamini, ut non circumveniamini, ut non fallamini. Attendite ergo, hoc est considerate, quia non sunt oves sed lupi in vestimento ovium, quia non sunt religiosi sed irreligiosi in figura religiositatis, quia nomine solum sunt Christiani, sed veritate vacui Christianorum persecutores.
\newline {\color{Red} Ad Ben. Ant.} Attendite a falsis prophetis qui veniunt ad vos in vestimentis ovium intrinsecus autem sunt lupi rapaces, a fructibus eorum cognoscetis eos.
\newline {\color{Red} Ad Mag. Ant.} Non omnis qui dicit michi domine domine intrabit in regnum celorum, sed qui facit voluntatem patris mei qui in celis est ipse intrabit in regnum celorum, alleluya.
\newline {\color{Red} Dominica Nona. Oratio.}
\lettrine[lines=2]{\zallmancaps \color{Red} L}{argire} nobis quesumus domine semper spiritum cogitandi que recta sunt propicius et agendi, ut qui sine te esse non possumus, secundum te vivere valeamus. Per.
\newline {\color{Red} Lectio VII. Secundum Lucam.}
\lettrine[lines=2]{\zallmancaps \color{Blue} I}{n} illo tempore: Dixit Ihesus discipulis suis parabolam hanc. Homo quidam erat dives qui habebat villicum. Et hic diffamatus est apud illum quasi dissipasset bona illius. Et reliqua.
\newline {\color{Red} Omelia Beati Hieronimi Presbiteri. \grecross}
\lettrine[lines=2]{\zallmancaps \color{Red} V}{illicus} proprie: ville gubernator est. Yconomus autem tam pecunie quam frugum et omnium que dominus possidet dispensator est. Iste igitur dispensator accusatus est ad dominum suum, quod dissiparet substantiam eius. Quo vocato: dixit. Quid hoc audio de te? Redde rationem dispensationis, neque enim in ea ultra dispensabis. Qui dixit in semetipso. Quid faciam quia dominus meus aufert a me dispensationem? Fodere non valeo, mendicare erubesco.
\newline {\color{Red} Lectio VIII.}
\lettrine[lines=2]{\zallmancaps \color{Blue} S}{i} dispensator iniqui mamone, domini voce laudatur, quod de re iniqua sibi iusticiam prepararit, et passus dispendium dominus laudat dispensatoris prudentiam, quod adversus dominum quidem fraudulenter sed pro se prudenter egerit: quanto magis Christus qui nullum potest sustinere damnum, et pronus est ad clementiam laudabit discipulos suos si in eos qui sibi credituri sunt misericordes fuerint?
\newline {\color{Red} Lectio IX.}
\lettrine[lines=2]{\zallmancaps \color{Red} D}{enique} post parabolam intulit. Et ego dico vobis. Facite vobis amicos de iniquo mamona. Mamona autem non Hebreorum sed Sirorum lingua divitie nuncupantur, quod de iniquitate collecte sint. Si ergo iniquitas bene dispensata vertitur in iusticiam: quanto magis sermo divinus in quo nulla est iniquitas, qui apostolis creditus est, si bene dispensatus fuerit dispensatores suos levabit in celum?
\newline {\color{Red} Ad Ben. Ant.} Dixit dominus villico, quid hoc audio de te, redde rationem villicationis tue alleluya.
\newline {\color{Red} Ad Mag. Ant.} Quid faciam quia dominus meus aufert a me villicationem, fodere non valeo, mendicare erubesco, scio quid faciam ut cum amotus fuero a villicatione recipiant me in domos suas.
\newline {\color{Red} Dominica Decima. Oratio.}
\lettrine[lines=2]{\zallmancaps \color{Blue} P}{ateant} aures misericordie tue domine precibus supplicantium, et ut petentibus desiderata concedas, fac eos que tibi sunt placita postulare. Per.
\newline {\color{Red} Lectio VII. Secundum Lucam.}
\lettrine[lines=2]{\zallmancaps \color{Red} I}{n} illo tempore: Cum appropinquaret Ihesus Hierusalem, videns civitatem flevit super illam dicens, quia si cognovisses et tu. Et reliqua.
\newline {\color{Red} Omelia Beati Gregorii Pape.} \grecross
\lettrine[lines=2]{\zallmancaps \color{Blue} Q}{uod} flente domino illa Hierosolimorum subversio describatur, que a Vaspasiano et Tito Romanis principibus facta est: nullus qui hystoriam eversionis eiusdem legit ignorat. Romani enim principes denunciantur cum dicitur. Quia venient dies in te et circundabunt te inimici tui vallo, et circundabunt te, et coangustabunt te undique et ad terram prosternent te, et filios tuos qui in te sunt.
\newline {\color{Red} Lectio VIII.}
\lettrine[lines=2]{\zallmancaps \color{Red} H}{oc} quoque quod additur: et non relinquent in te lapidem super lapidem, eciam ipsa iam eiusdem civitatis transmigratio testatur: quia dum nunc in eo loco constructa est ubi extra portam dominus fuerat crucifixus: prior illa Hierusalem ut dicitur, funditus est eversa. Cui ex qua culpa eversionis sue pena fuerit illata, subiungitur. Eo quod non cognoveris tempus visitationis tue.
\newline {\color{Red} Lectio IX.}
\lettrine[lines=2]{\zallmancaps \color{Blue} C}{reator} quippe omnium hanc per incarnationis sue mysterium visitare dignatus est, sed ipsa timoris et amoris illius recordata non est. Unde eciam per prophetam in increpationem cordis humani aves celi ad testimonium deducuntur, cum dicitur. Milvus in celo cognovit tempus suum, turtur et hyrundo et ciconia custodierunt tempus adventus sui, populus autem meus non cognovit iudicium domini.
\newline {\color{Red} Ad Ben. Ant.} Cum appropinquavit dominus Hierusalem videns civitatem flevit super illam et dixit, quia si cognovisses et tu, quia venient dies in te, et circundabunt te, et coangustabunt te undique, et ad terram prosternent te, eo quod non cognovisti tempus visitationis tue alleluya.
\newline {\color{Red} Ad Mag. Ant.} Scriptum est enim, quia domus mea domus orationis est cunctis gentibus, vos autem fecistis illam speluncam latronum, et erat cotidie docens in templo.
\newline {\color{Red} Dominica Undecima. Oratio.}
\lettrine[lines=2]{\zallmancaps \color{Red} D}{eus} qui omnipotentiam tuam parcendo maxime et miserando manifestas, multiplica super nos gratiam tuam, ut ad tua promissa currentes, celestium bonorum facias esse consortes. Per.
\newline {\color{Red} Lectio VII. Secundum Lucam.}
\lettrine[lines=2]{\zallmancaps \color{Blue} I}{n} illo tempore: Dixit Ihesus ad quosdam qui in se confidebant tanquam iusti, et aspernabantur ceteros parabolam istam. Duo homines ascenderunt in templum ut orarent, unus phariseus et alter publicanus, et reliqua.
\newline {\color{Red} Omelia Venerabilis Bede Presbiteri.} \grecross
\lettrine[lines=2]{\zallmancaps \color{Red} S}{uperbi} cum nequaquam omnia, sed modicum quid eorum que precepta sunt faciunt, non solum mox de sua iusticia presumunt, sed et infirmos quosque despiciunt: atque ideo quasi fide vacui cum oraverint non exaudiuntur.
\newline {\color{Red} Lectio VIII.}
\lettrine[lines=2]{\zallmancaps \color{Blue} P}{ublicanus} autem humiliter orans, ad illa prefate vidue hoc est ecclesie membra pertinet, de quibus supra dicitur. Deus autem non faciet vindictam electorum suorum clamantium ad se. Phariseus autem merita iactans: ad ea de quibus terribilis in conclusione sententia subditur. Veruntamen filius hominis veniens, putas inveniet fidem in terra?
\newline {\color{Red} Lectio IX.}
\lettrine[lines=2]{\zallmancaps \color{Red} P}{hariseus} autem stans: hec apud se orabat. Deus gratias ago tibi: quia non sum sicut ceteri hominum, raptores, iniusti, adulteri, velut eciam hic publicanus. Quatuor sunt species quibus omnis tumor arrogantium demonstratur: cum bonum aut a semetipsis habere se estimant, aut sibi datum desuper credunt, sed pro suis se hoc accepisse meritis putant: aut certe, cum iactant se habere quod non habent, aut despectis ceteris singulariter videri appetunt habere quod habent.
\newline {\color{Red} Ad Ben. Ant.} Duo homines ascenderunt in templum ut orarent, unus phariseus et alter publicanus, descendit hic iustificatus in domum suam ab illo alleluya.
\newline {\color{Red} Ad Mag. Ant.} Stans a longe publicanus nolebat oculos ad celum levare, sed percutiebat pectus suum dicens, Deus propitius esto michi peccatori.
\newline {\color{Red} Dominica Duodecima. Oratio.}
\lettrine[lines=2]{\zallmancaps \color{Blue} O}{mnipotens} sempiterne deus, qui abundantia pietatis tue et merita supplicum excedis et vota, effunde super nos misericordiam tuam, ut dimittas que conscientia metuit, et adicias quod oratio non presumit. Per.
\newline {\color{Red} Lectio VII. Secundum Marcum.}
\lettrine[lines=2]{\zallmancaps \color{Red} I}{n} illo tempore: Exiens Ihesus de finibus Tyri venit per Sidonem ad mare Galilee inter medios fines Decapoleos. Et reliqua.
\newline {\color{Red} Omelia Venerabilis Bede Presbiteri.}
\lettrine[lines=2]{\zallmancaps \color{Blue} S}{urdus} ille et mutus quem mirabiliter curatum a domino, modo cum evangelium legeretur audivimus: genus designat humanum in his qui ab errore diabolice deceptionis divina merentur gratia liberari. Obsurduit namque homo ab audiendo vite verbo: postquam mortifera serpentis verba contra Deum tumidus audivit.
\newline {\color{Red} Lectio VIII.}
\lettrine[lines=2]{\zallmancaps \color{Red} M}{utus} a laude conditoris effectus est: ex quo cum seductore colloquium habere presumpsit. Et merito clausit aures ab audienda inter angelos laude creatoris, quas ad audiendam eiusdem creatoris vituperationem, sermonibus hostis incautus aperuit. Merito clausit os a predicanda cum angelis laude creatoris, quod velut ad meliorandum eiusdem opus creatoris, cibi vetiti prevaricatione superbus implevit.
\newline {\color{Red} Lectio IX.}
\lettrine[lines=2]{\zallmancaps \color{Blue} H}{eu} miser generis humani defectus, quod in radice vitiosum germinavit: vitiosius multo dilatari cepit in propagine ramorum: ita ut veniente in carne domino exceptis paucis de Iudea fidelibus, totus pene mundus ab agnitione et confessione veritatis surdus erraret et mutus. Sed ubi abundavit delictum: superabundavit et gratia.
\newline {\color{Red} Ad Ben. Ant.} Exiens Ihesus de finibus Tyri venit per Sidonem ad mare Galilee inter medios fines Decapoleos alleluya.
\newline {\color{Red} Ad Mag. Ant.} Bene omnia fecit, surdos fecit audire et mutos loqui.
\newline {\color{Red} Dominica Terciadecima. Oratio.}
\lettrine[lines=2]{\zallmancaps \color{Red} O}{mnipotens} et misericors deus de cuius munere venit ut tibi a fidelibus tuis digne et laudabiliter serviatur, tribue quesumus nobis: ut ad promissiones tuas sine offensione curramus. Per.
\newline {\color{Red} Lectio VII. Secundum Lucam.}
\lettrine[lines=2]{\zallmancaps \color{Blue} I}{n} illo tempore: Dixit Ihesus discipulis suis. Beati oculi qui vident que vos videtis. Et reliqua.
\newline {\color{Red} Omelia Venerabilis Bede Presbiteri.} \grecross
\lettrine[lines=2]{\zallmancaps \color{Red} N}{on} oculi Scribarum et Phariseorum qui corpus tantum domini videre, sed illi beati oculi qui eius possunt cognoscere sacramenta: de quibus dicitur. Et revelasti ea parvulis. Beati oculi parvulorum quibus et se et patrem filius revelare dignatur. Dico enim vobis quod multi reges et prophete voluerunt videre que vos videtis, et non viderunt, et audire que auditis, et non audierunt.
\newline {\color{Red} Lectio VIII.}
\lettrine[lines=2]{\zallmancaps \color{Blue} A}{braham} exultavit ut videret diem Christi: et vidit et gavisus est. Ysaias quoque et Micheas et multi alii prophete viderunt gloriam domini, qui propterea videntes sunt appellati: sed hi omnes a longe aspicientes et salutantes, per speculum et in enigmate viderunt. Apostoli autem impresentiarum habentes dominum, convescentesque ei, et que voluissent interrogando discentes: nequaquam per angelos aut varias visionum species opus habebant doceri.
\newline {\color{Red} Lectio IX.}
\lettrine[lines=2]{\zallmancaps \color{Red} E}{t} ecce quidam legis peritus surrexit, temptans eum et dicens. Magister, quid faciendo vitam eternam possidebo? Legis peritus qui de vita eterna dominum temptans interrogat: occasionem ut reor temptandi de ipsis domini sermonibus scripsit, ubi ait. Gaudete autem quod nomina vestra scripta sunt in celis. Sed ipsa sua temptatione declarat, quam vera sit illa domini confessio qua patri loquitur, quod abscondisti hec a sapientibus et prudentibus: et revelasti ea parvulis.
\newline {\color{Red} Ad Ben. Ant.} Homo quidam descendebat ab Hierusalem in Iericho et incidit in latrones, qui eciam despoliaverunt eum, et plagis impositis abierunt semivivo relicto.
\newline {\color{Red} Ad Mag. Ant.} Quis tibi videtur proximus fuisse illi qui incidit in latrones, et ait, ille qui fecit misericordiam in illum, vade et tu fac similiter alleluya.
\newline {\color{Red} Dominica Quartadecima. Oratio.}
\lettrine[lines=2]{\zallmancaps \color{Blue} O}{mnipotens} sempiterne deus, da nobis fidei, spei et caritatis augmentum, et ut mereamur assequi quod promittis, fac nos amare quod precipis. Per.
\newline {\color{Red} Lectio VII. Secundum Lucam.}
\lettrine[lines=2]{\zallmancaps \color{Red} I}{n} illo tempore: Dum iret Ihesus in Hierusalem, transibat per mediam Samariam et Galileam. Et cum ingrederetur quoddam castellum: occurrerunt ei decem viri leprosi. Et reliqua.
\newline {\color{Red} Omelia Venerabilis Bede Presbiteri.}
\lettrine[lines=2]{\zallmancaps \color{Blue} L}{eprosi} non absurde intelligi possunt, qui scientiam vere fidei non habentes: varias doctrinas profitentur erroris. Non enim vel abscondunt imperitiam suam, sed pro summa peritia proferunt in lucem: et iactantiam sermonis ostendunt. Nulla porro falsa doctrina est: que non aliqua vera intermisceat.
\newline {\color{Red} Lectio VIII.}
\lettrine[lines=2]{\zallmancaps \color{Red} V}{era} ergo falsis inordinate permista in una disputatione vel narratione hominis, tanquam in unius corporis colore apparentia, significant lepram: tanquam veris falsisque colorum locis, humana corpora variantem atque maculantem. Hi autem tam vitandi sunt ecclesie: ut si fieri potest longius remoti magno clamore Christum interpellent. Unde et apte subiungitur. Qui steterunt a longe: et levaverunt vocem dicentes. Ihesu preceptor miserere nostri.
\newline {\color{Red} Lectio IX.}
\lettrine[lines=2]{\zallmancaps \color{Blue} E}{t} bene ut salventur Ihesum preceptorem nominant. Quia enim in eius verbis se errasse significant: hunc salvandi humiliter preceptorem vocant. Cumque ad cognitionem preceptoris redeunt, mox ad formam salutis recurrunt. Nam sequitur. Quos ut vidit dixit. Ite, ostendite vos sacerdotibus. Et factum est dum irent: mundati sunt.
\newline {\color{Red} Ad Ben. Ant.} Cum ingrederetur Ihesus quoddam castellum occurrerunt ei decem viri leprosi, qui steterunt a longe et levaverunt vocem dicentes, Ihesu preceptor miserere nostri.
\newline {\color{Red} Ad Mag. Ant.} Nonne decem mundati sunt, et novem ubi sunt, non est inventus qui rediret et daret gloriam Deo nisi hic alienigena, vade quia fides tua te salvum fecit, alleluya.
\newline {\color{Red} Dominica Quintadecima. Oratio.}
\lettrine[lines=2]{\zallmancaps \color{Red} C}{ustodi} domine quesumus ecclesiam tuam propitiatione perpetua et quia sine te labitur humana mortalitas, tuis semper auxiliis et abstrahatur a noxiis, et ad salutaria dirigatur. Per.
\newline {\color{Red} Lectio VII. Secundum Matheum.}
\lettrine[lines=2]{\zallmancaps \color{Blue} I}{n} illo tempore: Dixit Ihesus discipulis suis. Nemo potest duobus dominis servire. Et reliqua.
\newline {\color{Red} Omelia Beati Hieronimi Presbiteri.}
\lettrine[lines=2]{\zallmancaps \color{Red} M}{ammona} sermone Siriaco divitie nuncupantur. Non potestis Deo servire et mammone. Audiat hoc avarus, audiat qui censetur vocabulo Christiano, non posse simul divitiis Christoque serviri. Et tamen non dicit, qui habet divitias, sed qui servit divitiis. Qui enim divitiarum servus est, divitias custodit ut servus. Qui autem servitutis excussit iugum, distribuit eas ut dominus.
\newline {\color{Red} Lectio VIII.}
\lettrine[lines=2]{\zallmancaps \color{Blue} N}{e} solliciti sitis anime vestre quid manducetis, neque corpori vestro quid induamini. In nonnullis codicibus additum est. Neque quid bibatis. Ergo quod natura omnibus tribuit et iumentis ac bestiis hominibus quoque commune est: huius cura penitus liberamur. Et precipitur nobis ne solliciti simus quid comedamus. Et quia in sudore vultus preparamus nobis panem: labor exercendus est, sollicitudo tollenda.
\newline {\color{Red} Lectio IX.}
\lettrine[lines=2]{\zallmancaps \color{Red} H}{oc} quod dicitur ne solliciti sitis anime vestre quid comedatis, neque corpori vestro quid induamini: de carnali cibo ac vestimento accipiamus. Ceterum de spiritualibus cibis et vestimentis: semper debemus esse solliciti. Nonne anima plus est quam esca et corpus plus quam vestimentum? Quod dicit: istius modi est. Qui maiora prestitit: utique et minora prestabit.
\newline {\color{Red} Ad Ben. Ant.} Nolite solliciti esse dicentes, quid manducabimus aut quid bibemus, scit enim pater vester celestis quid vobis necesse sit alleluya.
\newline {\color{Red} Ad Mag. Ant.} Querite primum regnum Dei et iusticiam eius et hec omnia adicientur vobis.
\newline {\color{Red} Dominica Sextadecima. Oratio.}
\lettrine[lines=2]{\zallmancaps \color{Blue} E}{cclesiam} tuam domine miseratio continuata mundet et muniat, et quia sine te non potest salva consistere, tuo semper munere gubernetur. Per.
\newline {\color{Red} Lectio VII. Secundum Lucam.}
\lettrine[lines=2]{\zallmancaps \color{Red} I}{n} illo tempore: Ibat Ihesus in civitatem que vocatur Naym, et ibant cum illo discipuli eius, et turba copiosa. Cum autem appropinquaret porte civitatis: ecce defunctus efferebatur filius unicus matris sue. Et reliqua.
\newline {\color{Red} Omelia Venerabilis Bede Presbiteri.} \grecross
\lettrine[lines=2]{\zallmancaps \color{Blue} D}{efunctus} hic qui extra portam civitatis multis est intuentibus elatus, significat hominem letali criminum funere soporatum, eandemque insuper anime mortem, non cordis adhuc cubili tegentem, sed ad multorum notitiam per locutionis operisve indicium quasi per sue civitatis ostia propalantem.
\newline {\color{Red} Lectio VIII.}
\lettrine[lines=2]{\zallmancaps \color{Red} E}{t} bene filius unicus matri sue fuisse perhibetur, quia licet e multis collecta personis una sit perfecta et immaculata virgo mater ecclesia, singuli quique tamen fidelium universalis se ecclesie filios rectissime fatentur. Electus quilibet quando ad fidem imbuitur, filius est, quando alios imbuit: mater.
\newline {\color{Red} Lectio IX.}
\lettrine[lines=2]{\zallmancaps \color{Blue} A}{n} non materno erga parvulos angebatur affectu, qui ait. Filioli mei quos iterum parturio donec formetur Christus in vobis? Portam civitatis qua defunctus efferebatur: puto aliquem de sensibus esse corporeis. Qui enim seminat inter fratres discordias, qui iniquitatem in excelso loquitur: per oris portam extrahitur mortuus.
\newline {\color{Red} Ad Ben. Ant.} Ibat Ihesus in civitatem que vocatur Naym, et ecce defunctus efferebatur filius unicus matris sue.
\newline {\color{Red} Ad Mag. Ant.} Accepit autem omnes timor, et magnificabant deum, dicentes quia propheta magnus surrexit in nobis, et quia deus visitavit plebem suam.%typo Acepit
\newline {\color{Red} Dominica Decimaseptima. Oratio.}
\lettrine[lines=2]{\zallmancaps \color{Red} T}{ua} nos domine quesumus gratia semper et preveniat et sequatur, ac bonis operibus iugiter prestet esse intentos. Per.
\newline {\color{Red} Lectio VII. Secundum Lucam.}
\lettrine[lines=2]{\zallmancaps \color{Blue} I}{n} illo tempore: Cum intraret Ihesus in domum cuiusdam principis Phariseorum sabbato manducare panem: et ipsi observabant eum. Et ecce homo quidam ydropicus erat ante illum. Et reliqua.
\newline {\color{Red} Omelia Venerabilis Bede Presbiteri.}
\lettrine[lines=2]{\zallmancaps \color{Red} Y}{dropis} morbus: ab aquoso humore vocabulum trahit. Grece enim ydor: aqua vocatur. Est autem humor subcutaneus de vitio vesice natus, cum inflatione turgente, et hanelitu fetido. Propriumque est ydropici, quanto magis abundat humore inordinato: tanto amplius sitire. Et ideo recte comparatur ei quem fluxus carnalium voluptatum exuberans, aggravat.
\newline {\color{Red} Lectio VIII.}
\lettrine[lines=2]{\zallmancaps \color{Blue} C}{omparatur} diviti avaro, qui quanto est copiosior divitiis quibus non bene utitur: tanto ardentius talia concupiscit. Et respondens Ihesus, dixit ad legis peritos et Phariseos. Si licet Sabbato curare? At illi tacuerunt. Quod dicitur respondisse Ihesus, ad hoc respicit quod premissum est: et ipsi observabant eum. Dominus enim novit cogitationes hominum.
\newline {\color{Red} Lectio IX.}
\lettrine[lines=2]{\zallmancaps \color{Red} S}{ed} merito interrogati tacent, qui contra se dictum quicquid dixerint vident. Nam si licet Sabbato curare, quare salvatorem an curet observant? Si non licet, quare ipsi Sabbato pecora curant? Ipse vero apprehensum sanavit eum ac dimisit.
\newline {\color{Red} Ad Ben. Ant.} Dixit Ihesus ad legis peritos et phariseos, si licet Sabbato curare, at illi tacuerunt ipse vero apprehendens ydropicum sanavit eum ac dimisit.
\newline {\color{Red} Ad Mag. Ant.} Cum vocatus fueris ad nuptias recumbe inovissimo loco, ut dicat tibi amice ascende superius, tunc erit tibi gloria coram simul discumbentibus alleluya.%sic
\newline {\color{Red} Dominica Decimaoctava. Oratio.}
\lettrine[lines=2]{\zallmancaps \color{Blue} D}{a} quesumus domine populo tuo diabolica vitare contagia, et te solum dominum pura mente
%pp117
sectari. Per.
\newline {\color{Red} Lectio VII. Secundum Matheum.}
\lettrine[lines=2]{\zallmancaps \color{Red} I}{n} illo tempore: Pharisei audientes quod Ihesus silentium imposuisset Saduceis: convenerunt in unum. Et interrogavit eum unus ex eis, legis doctor: temptans eum. Magister, quod est mandatum magnum in lege? Et reliqua.
\newline {\color{Red} Omelia Beati Iohannis Episcopi.}
% https://la.wikisource.org/wiki/Commentarium_in_Matthaeum_(Rabanus_Maurus)
\lettrine[lines=2]{\zallmancaps \color{Blue} C}{onvenerunt} ut multitudine vincerent, quem ratione superare non poterant. A veritate nudos se esse professi sunt, qui multitudine se armaverant. Dicebant enim apud se. Unus loquatur pro omnibus, et omnes loquamur per unum, ut si quidem vicerit, omnes videamur victores: si autem victus fuerit, vel solus videatur confusus.
\newline {\color{Red} Lectio VIII.}
\lettrine[lines=2]{\zallmancaps \color{Red} O}{}\ Pharisei qui omnia propter homines cogitatis et facitis, primum quidem venientes cum uno, vincendi estis per unum. Tamen pone, quia uno victo homines non intelligunt omnes vos esse victos, nunquid conscientie vestre non sentiunt se esse confusas? Levis enim est ei consolatio qui in seipso confusus est: quod ab aliis ignoratur.
\newline {\color{Red} Lectio IX.}
\lettrine[lines=2]{\zallmancaps \color{Blue} M}{agister} quod est mandatum magnum in lege? Magistrum vocat cuius non vult esse discipulus. Simplicissimus interrogator, et malignissimus insidiator, de magno mandato interrogat: qui nec minimum observat. Ille enim debet interrogare de maiore iusticia, qui iam minorem complevit. Dominus autem sic ei respondit: ut interrogationis eius fictam conscientiam statim primo responso percuteret.
\newline {\color{Red} Ad Ben. Ant.} Magister, quod est mandatum magnum in lege, ait illi Ihesus, diliges dominum Deum tuum ex toto corde tuo, alleluya.
\newline {\color{Red} Ad Mag. Ant.} Quid nobis videtur de Christo cuius filius est, dicunt ei omnes, David, dixit eis Ihesus quomodo David in spiritu vocat eum dominum dicens, dixit dominus domino meo sede a dextris meis?
\newline {\color{Red} Dominica Decimanona. Oratio.}
\lettrine[lines=2]{\zallmancaps \color{Red} D}{irigat} corda nostra quesumus domine tue miserationes operatio, quia tibi sine te placere non possumus. Per.
\newline {\color{Red} Lectio VII. Secundum Matheum.}
\lettrine[lines=2]{\zallmancaps \color{Blue} I}{n} illo tempore: Ascendens Ihesus in naviculam, transfretavit et venit in civitatem suam. Et reliqua.
\newline {\color{Red} Omelia Beati Iohannis Episcopi.} \grecross
% https://archive.org/stream/patrologiaecursu186mign/patrologiaecursu186mign_djvu.txt
\lettrine[lines=2]{\zallmancaps \color{Red} C}{hristus} ecclesie navem seculi fluctum semper mitigaturus ascendit, ut credentes in se, ad celestem patriam tranquilla navigatione perducat, et municipes civitatis sue faciat, quos humanitatis sue fecit esse consortes. Non ergo Christus indiget navi, sed navis indiget Christo, quia sine celesti gubernatore navis ecclesie per mundanum pelagus, per talia et tanta discrimina ad celestem portum non valet pervenire.
\newline {\color{Red} Lectio VIII.}
\lettrine[lines=2]{\zallmancaps \color{Blue} A}{scendens} inquit in naviculam: transfretavit et venit in civitatem suam. Creator rerum, orbis terre dominus: postea quam se propter nos nostra angustavit in carne: cepit habere humanam patriam, cepit civitatis Iudaice esse civis, parentes habere cepit parentum omnium ipse parens, ut invitaret amor, attraheret caritas, vinciret affectio, suaderet humanitas, quos fugarat dominatio, metus disperserat, et fecerat vis potestatis extorres.%typo: inquid
\newline {\color{Red} Lectio IX.}
\lettrine[lines=2]{\zallmancaps \color{Red} E}{t} videns Ihesus fidem illorum: dixit paralitico. Confide fili, remittuntur tibi peccata tua. O mira humilitas. Despectum et debilem totisque membrorum compagibus dissolutum filium vocat: quem sacerdotes non dignabantur contingere. Aut certe ideo filium: quia dimittuntur ei peccata sua. Audit veniam, et tacet paraliticus, nec ullam respondet gratiam: quia plus corporis quam anime tendebat ad curam.
\newline {\color{Red} Ad Ben. Ant.} Dixit dominus paralitico, confide fili remittuntur tibi peccata tua.
\newline {\color{Red} Ad Mag. Ant.} Videntes autem turbe timuerunt, et glorificaverunt Deum qui dedit potestatem talem hominibus.
\newline {\color{Red} Dominica Vicesima. Oratio.}
\lettrine[lines=2]{\zallmancaps \color{Blue} O}{mnipotens} et misericors deus universa nobis adversantia propitiatus exclude, ut mente et corpore pariter expediti, que tua sunt liberis mentibus exequamur. Per.
\newline {\color{Red} Lectio VII. Secundum Matheum.}
\lettrine[lines=2]{\zallmancaps \color{Red} I}{n} illo tempore: Dicebat Ihesus turbis parabolam hanc. Simile factum est regnum celorum homini regi, qui fecit nuptias filio suo. Et misit servos suos vocare invitatos ad nuptias, et nolebant venire. Et reliqua.
\newline {\color{Red} Omelia Beati Gregorii Pape.} \grecross
\lettrine[lines=2]{\zallmancaps \color{Blue} Q}{uerendum} prius est: an hec apud Matheum ipsa sit lectio, que apud Lucam sub appellatione cene describitur. Et quidem sunt nonnulla que sibi dissona esse videntur: quia hic prandium illic cena memoratur. Hic qui ad nuptias non dignis vestibus intravit repulsus est, illic nullus qui intrasse dicitur, repulsus esse perhibetur.
\newline {\color{Red} Lectio VIII.}
\lettrine[lines=2]{\zallmancaps \color{Red} Q}{ua} ex re recte colligitur, quod et hic per nuptias presens ecclesia, et illic per cenam eternum et ultimum convivium designatur, quia et hanc nonnulli exituri intrant, et ad illud quisquis semel intraverit, ulterius non exibit. At siquis forte contendat, hanc eandem esse lectionem, ego melius puto salva fide alieno intellectui cedere: quam contentionibus deservire: quoniam et intelligi congrue forsitan potest, quia de proiecto eo qui cum nuptiali veste non venerat: quod Lucas tacuit, Matheus dixit.
\newline {\color{Red} Lectio IX.}
\lettrine[lines=2]{\zallmancaps \color{Blue} Q}{uod} vero per illum cena, per hunc autem prandium dicitur, nequaquam vel hoc nostre intelligentie obsistit, quia cum ad horam nonam apud antiquos cotidie prandium fieret: ipsum quoque prandium cena vocabatur. Sepe autem iam me dixisse memini: quia plerumque in sancto evangelio, regnum celorum presens ecclesia nominatur. Congregatio quippe iustorum regnum celorum dicitur.
\newline {\color{Red} Ad Ben. Ant.} Dicite invitatis, ecce prandium meum paravi, venite ad nuptias dicit dominus.
\newline {\color{Red} Ad Mag. Ant.} Nuptie quidem parate sunt sed qui invitati erant non fuerunt digni, ite ergo ad exitus viarum et quoscumque inveneritis vocate ad nuptias dicit dominus.
\newline {\color{Red} Dominica Vicesima Prima. Oratio.}
\lettrine[lines=2]{\zallmancaps \color{Red} L}{argire} quesumus domine fidelibus tuis indulgentiam placatus et pacem, ut pariter ab omnibus mundentur offensis, et secura tibi mente deserviant. Per.
\newline {\color{Red} Lectio VII. Secundum Iohannem.}
\lettrine[lines=2]{\zallmancaps \color{Blue} I}{n} illo tempore: Erat quidam regulus cuius filius infirmabatur Capharnaum. Et reliqua.
\newline {\color{Red} Omelia Beati Gregorii Pape.} \grecross
\lettrine[lines=2]{\zallmancaps \color{Red} H}{oc} nobis solummodo fratres de sancti evangelii expositione video esse requirendum, cur is qui ad salutem petendam venerat, audivit, nisi signa et prodigia videritis, non creditis. Qui enim querebat salutem filio: procul dubio credebat. Neque enim ab eo quereret salutem, quem non crederet salvatorem. Quare ergo dicitur, nisi signa et prodigia videritis non creditis cum ante crederet quam signum videret?
\newline {\color{Red} Lectio VIII.}
\lettrine[lines=2]{\zallmancaps \color{Blue} S}{ed} mementote quid petiit, et aperte cognoscetis, quia in fide dubitavit. Poposcit namque ut descenderet et sanaret filium eius. Corporalem ergo presentiam domini regulus querebat, qui per spiritum nusquam deerat. Minus itaque in illum credidit, quem non putavit posse salutem dare nisi presens esset et corpore.
\newline {\color{Red} Lectio IX.}
\lettrine[lines=2]{\zallmancaps \color{Red} S}{i} perfecte credidisset, procul dubio sciret quia non esset locus ubi non esset Deus. Ex magna ergo parte diffisus est, qui honorem non dedit maiestati, sed presentie corporali. Salutem itaque filio petiit, et tamen in fide dubitavit, quia eum ad quem venerat et potentem ad curandum credidit, et tamen morienti filio esse absentem Deum putavit.
\newline {\color{Red} Ad Ben. Ant.} Erat quidam regulus cuius filius infirmabatur Capharnaum, hic cum audisset quod Ihesus adveniret in Galileam rogabat eum ut descenderet, et sanaret filium eius.
\newline {\color{Red} Ad Mag. Ant.} Cognovit ergo pater quia illa hora erat, in qua dixit Ihesus filius tuus vivit, et credidit ipse et domus eius tota.
\newline {\color{Red} Dominica Vicesima Secunda. Oratio.}
\lettrine[lines=2]{\zallmancaps \color{Blue} F}{amiliam} tuam quesumus domine continua pietate custodi: ut a cunctis adversitatibus te protegente sit libera, et in bonis actibus tuo nomini sit devota. Per.
\newline {\color{Red} Lectio VII. Secundum Matheum.}
\lettrine[lines=2]{\zallmancaps \color{Red} I}{n} illo tempore: Dixit Ihesus discipulis suis. Simile est regnum celorum homini regi, qui volunt rationem ponere cum servis suis. Et reliqua.
\newline {\color{Red} Omelia Beati Hieronimi Presbiteri.}
\lettrine[lines=2]{\zallmancaps \color{Blue} F}{amiliare} est Siris et maxime Palestinis, ad omnem sermonem suum parabolas iungere, ut quod per simplex preceptum teneri ab auditoribus non potest: per similitudinem exempli teneatur. Precepit itaque dominus Petro sub comparatione regis et domini, et servi qui debitor decem milium talentorum a domino rogans: veniam impetraverat, ut ipse quoque dimittat conservis suis minora peccantibus.
\newline {\color{Red} Lectio VIII.}
\lettrine[lines=2]{\zallmancaps \color{Red} S}{i} enim ille rex et dominus servo suo debitori, decem milia talentorum tam facile dimisit, quanto magis servi debent conservis suis minora dimittere? Quod ut manifestius fiat: dicamus sub exemplo. Siquis nostrum commiserit adulterium, homicidium, sacrilegium, maiora crimina decem milia talentorum rogantibus dimittuntur, si et ipsi dimittamus minora peccantibus.
\newline {\color{Red} Lectio IX.}
\lettrine[lines=2]{\zallmancaps \color{Blue} S}{i} autem ob factam contumeliam simul implacabiles, et propter amarius verbum perpetes habeamus discordias, nonne nobis videmur recte redigendi in carcerem, et sub exemplo operis nostri hoc agere, ut maiorum nobis delictorum venia non relaxetur? Oblatus est ei unus, qui debebat decem milia talenta. Scio quosdam istum qui debebat decem milia talentorum, diabolum interpretari.
\newline {\color{Red} Ad Ben. Ant.} Dixit autem dominus servo, redde quod debes, procidens autem servus ille rogabat eum dicens, patientiam habe in me et omnia reddam tibi.
\newline {\color{Red} Ad Mag. Ant.} Serve nequam omne debitum dimisi tibi quoniam rogasti me, nonne ergo oportuit et te misereri conservi tui, sicut et ego tui misertus sum dicit dominus?
\newline {\color{Red} Dominica Vicesima Tercia. Oratio.}
\lettrine[lines=2]{\zallmancaps \color{Red} D}{eus} refugium nostrum et virtus, adesto piis ecclesie tue precibus actor ipse pietatis, et presta ut quod fideliter petimus efficaciter consequamur. Per.
\newline {\color{Red} Lectio VII. Secundum Matheum.}
\lettrine[lines=2]{\zallmancaps \color{Blue} I}{n} illo tempore: Abeuntes Pharisei consilium inierunt, ut caperent Ihesum in sermone. Et reliqua.
\newline {\color{Red} Omelia Beati Iohannis Crisostomi.} \grecross
\lettrine[lines=2]{\zallmancaps \color{Red} A}{bierunt} et consilium acceperunt. Quo abierunt? Ad Herodianos. Nam ex eo quod non dicit consiliati sunt, sed consilium acceperunt, et ex eo quod pariter cum Herodianis venerunt: apparet quod cum illis, eiusmodi circumventionis consilium tractaverunt. Homines coloni cuius terram possident: eius auxilio opus habent. Qui iusticiam tenet: nullius patrocinio indiget, nisi Dei.
\newline {\color{Red} Lectio VIII.}
\lettrine[lines=2]{\zallmancaps \color{Blue} Q}{ui} in diaboli iniquitatibus ambulat: diaboli adiutorium necessarium habet. Coloni enim dei, diaboli auxilium non requirit. Colonus autem diaboli auxilium dei, et si querit non inveniet. Vidisti aliquando euntem ad furtum Deum orare, ut bene prosperetur in furto? Aut qui vadit ad fornicationem, nunquam signum crucis ponit sibi in fronte, ut non comprehendatur in crimine? Quod si fecerit non iuvat, quia nescit iusticia Dei, patrocinium dare criminibus.
\newline {\color{Red} Lectio IX.}
\lettrine[lines=2]{\zallmancaps \color{Red} S}{ic} et isti Christum expugnare cupientes, convenienter non ad servos Dei viros religiosos, sed ad gentiles, idest Herodianos, se contulerunt. Quale consilium: tales et consiliatores. Quis autem poterat dare consilium contra Christum, nisi diabolus qui erat adversarius Christi? Cogitabant enim apud se sacerdotes. Si nos soli euntes interrogaverimus Christum, quamvis dixerit quia non licet Cesari dari tributum, tamen nemo nobis credet dicentibus contra eum. Iam enim omnes sciunt quia inimici eius sumus.
\newline {\color{Red} Ad Ben. Ant.} Magister scimus quia verax es et viam Dei in veritate doces alleluya.
\newline {\color{Red} Ad Mag. Ant.} Reddite ergo que sunt Cesaris Cesari, et que sunt Dei Deo, alleluya.
\newline {\color{Red} Dominica Vicesima Quarta. Oratio.}
\lettrine[lines=2]{\zallmancaps \color{Blue} A}{bsolve} quesumus domine tuorum delicta populorum, ut a peccatorum nostrorum nexibus, que pro nostra fragilitate contraximus, tua benignitate liberemur. Per.
\newline {\color{Red} Lectio VII. Secundum Matheum.}
\lettrine[lines=2]{\zallmancaps \color{Red} I}{n} illo tempore: Loquente Ihesu ad turbas, ecce princeps unus accessit et adorabat eum dicens. Domine filia mea modo defuncta est. Et reliqua.
\newline {\color{Red} Omelia Beati Hieronimi Presbiteri.}
\lettrine[lines=2]{\zallmancaps \color{Blue} O}{ctavum} signum est, in quo princeps suscitari postulat filiam suam, nolens de mysterio circumcisionis excludi. Sed subintrat mulier sanguine fluens, et octavo sanatur loco, ut principis filia de hoc exclusa numero veniat ad nonum, iuxta illud quod in psalmis dicitur. Ethiopia preveniat manus eius Deo. Et cum intraverit plenitudo gentium: tunc omnis Israhel salvus erit.
\newline {\color{Red} Lectio VIII.}
\lettrine[lines=2]{\zallmancaps \color{Red} E}{t} ecce mulier que sanguinis fluxum patiebatur duodecim annis, accessit retro, et tetigit fimbriam vestimenti eius. In evangelio secundum Lucam scribitur, quod filia principis duodecim annos haberet etatis. Nota ergo quod eo tempore: hec mulier, idest gentium populus ceperit egrotare: quo gens crediderit Iudeorum. Nisi enim ex comparatione virtutum: vitium non ostenditur.
\newline {\color{Red} Lectio IX.}
\lettrine[lines=2]{\zallmancaps \color{Blue} H}{ec} mulier sanguine fluens, non in domo, non in urbe accedit ad dominum, quia iuxta legem urbibus excludebatur: sed in itinere ambulante domino, ut dum pergit ad aliam: alia curaretur. Unde dicunt apostoli. Vobis quidem prius oportebat predicari verbum Dei, sed quoniam vos indignos iudicastis salute: transgredimur ad gentes.
\newline {\color{Red} Ad Ben. Ant.} Loquente Ihesu ad turbas ecce princeps unus accessit et adoravit eum dicens, domine filia mea modo defuncta est, sed veni impone manum tuam super eam et vivet alleluya.
\newline {\color{Red} Ad Mag. Ant.} Confide filia, fides tua te salvam fecit, et sanata est mulier ex illa hora, alleluya.
\newline {\color{Red} Dominica Vicesima Quinta. Oratio.}
\lettrine[lines=2]{\zallmancaps \color{Red} E}{xcita} quesumus domine tuorum fidelium voluntates, ut divini operis fructum propensius exequentes, pietatis tue remedia maiora percipiant. Per.
\newline {\color{Red} Lectio VII. Secundum Iohannem.}
\lettrine[lines=2]{\zallmancaps \color{Blue} I}{n} illo tempore: Cum sublevasset oculos Ihesus, et vidisset quia multitudo maxima venit ad eum: dixit ad Philippum. Unde ememus panes, ut manducent hi? Et reliqua.
\newline {\color{Red} Omelia Venerabilis Bede Presbiteri.} \grecross
\lettrine[lines=2]{\zallmancaps \color{Red} Q}{uod} sublevasse oculos Ihesus et venientem ad se multitudinem vidisse perhibetur, divine pietatis indicium est, quia videlicet cunctis ad se venire querentibus, gratia misericordie celestis occurrere consuevit. Et ne querendo errare possint: lucem sui spiritus aperire currentibus solitus est. Nam quod oculi Ihesu, dona spiritus eius mystice designent, testatur in Apocalypsi Iohannes, qui figurate loquens de illo ait. Vidi agnum stantem tanquam occisum habentem cornua septem et oculos septem, qui sunt septem spiritus Dei, missi in omnem terram.
\newline {\color{Red} Lectio VIII.}
\lettrine[lines=2]{\zallmancaps \color{Blue} Q}{uod} autem Philippum temptat dominus: provida utique dispensatione facit, non ut ipse que non noverat discat, sed ut Philippus tarditatem fidei sue, quam in magistro sciente ipse nesciebat, temptatus agnoscat, et miraculo facto castiget. Neque enim dubitare debuerat, presente rerum creatore qui educit panem de terra, et vino letificat cor hominis, paucorum denariorum panes sufficere, turbarum milibus non paucis, ut unusquisque sufficienter acciperet, et iam saturatus abiret.
\newline {\color{Red} Lectio IX.}
\lettrine[lines=2]{\zallmancaps \color{Red} Q}{uinque} autem panes quibus multitudinem populi saturavit quinque sunt libri Moysi: quibus spirituali intellectu patefactis, et abundantiore iam sensu multiplicatis, auditorum fidelium cotidie corda reficit. Qui bene ordeacei fuisse referuntur propter austeriora legis edicta, et tegumenta littere grossiora, que interiora spiritalis sensus: quasi medullam celabant.
\newline {\color{Red} Ad Ben. Ant.} Cum sublevasset oculos Ihesus et vidisset maximam multitudinem venientem ad se, dixit ad Philippum, unde ememus panes ut manducent hi, hoc autem dicebat temptans eum, ipse enim sciebat quid esset facturus.
\newline {\color{Red} Ad Mag. Ant.} Illi homines cum signum vidissent quod factum fuerat, glorificabant deum et dicebant, quia hic est salvator mundi.
\lettrine[lines=2]{\zallmancaps \color{Blue} S}{}\ul{\textsc{ecundum} quod a Dominica,} \hyperlink{deus-omnium}{Deus omnium.} \ul{usque ad Adventum exclusive: sunt viginti septem ebdomade quando prolixius est temporis spatium, quando vero brevius: viginti due. Evangelia vero Dominicalia cum orationibus: viginti quinque, et totidem antiphone ad Ben. et totidem ad Mag. Quando igitur tantum viginti due fuerint ebdomade, ultime omelie trium evangeliorum legantur per ferias vicesime secunde ebdomade si vacaverint, et dicantur orationes cum antiphonis ad Ben. et ad Mag. Quod si tot ferie non vacaverint, quot sunt omelie, omelia evangelii,} Loquente Ihesu. \ul{cum oratione sua et antiphonis dimittatur. Si autem viginti tres fuerint ebdomade, due residue omelie legantur infra ultimam ebdomadam predicto modo. Si vero viginti quatuor: omelia evangelii,} Cum sublevasset. \ul{legatur in aliqua feria vacante: eodem modo. Si autem viginti quinque: quelibet omelia legatur in sua Dominica. Si vero viginti sex: omelia evangelii,} Loquente Ihesu. \ul{repetatur in vicesima quinta Dominica, quando de dominica agetur, et omelia evangelii,} Cum sublevasset. \ul{Legatur in vicesima sexta. Si autem viginti septem: omelia evangelii,} Loquente Ihesu. \ul{legatur in vicesima quarta et vicesima quinta Dominica: et omelia evangelii,} Cum sublevasset. \ul{in vicesima sexta et vicesima septima. Hoc tamen sciendum quod quando sunt viginti sex vel viginti septem ebdomade, postquam omelia lecta fuerit in sua prima ebdomada, sive in Dominica sive in aliqua feria: si in sequenti Dominica in qua repeti debuerat festum occurrerit, vel octava: non est ulterius repetenda.}

\ubsection{Dedicatione Ecclesie}{dedicatione-ecclesie}

\lettrine[lines=2]{\zallmancaps \color{Red} N}{}\ul{\textsc{otandum} quod in primo anno dedicationis ecclesie, officium inchoatur in Missa. In hora vero que sequitur Missam et deinceps in omnibus Horis usque ad Completorium sequentis diei inclusive fiat officium totum duplex, et dicatur oratio,} Deus qui invisibiliter. \ul{ad omnes horas preterquam ad Completorium et Primam. In Vesperis diei quo ecclesia dedicata est, cantetur, responsorium, et die sequenti in Matutinis fiat novem lectiones, vel tres tempore Paschali. Octava vero si celebranda fuerit celebrabitur revoluta ebdomada eo die quo dedicatio facta fuerit, et dicetur oratio,} Deus qui invisibiliter. \ul{cotidie per octavas, quando de ipsis agendum fuerit. Secundo autem anno et deinceps singulis annis in anniversario fiat officium sicut in toto duplici. Et octave celebrentur, preterquam in Adventu et deinceps usque ad Octavas Epyphanie, et in Septuagesima et deinceps usque ad Octavas Pasche, et in diebus Rogationum et deinceps usque ad Dominica,} \hyperlink{deus-omnium}{Deus omnium.} \ul{Quando autem octave non fiunt: infra octavas memoria non fiat de dedicatione. Quod si hoc festum infra octavas alicuius sancti evenerit durantibus octavis de sancto, fiat memoria de octavis dedicationis. Fratres de conventu in quo festum huiusmodi celebratur tam intus quam extra faciant officium de dedicatione. Fratres eciam hospites dum sunt presentes in conventu aliquo: faciant officium sicut fit in conventu, tam de festo dedicationis quam de aliis, quamvis per ordinem generaliter non fiant. Postquam vero recesserint de conventu: redeant ad officium consuetum.}
{\ubsubsection{In Dedicatione Ecclesie et in Anniversario Eiusdem}{in-dedicatione-ecclesie-breviarium}}
\lettrine[lines=2]{\zallmancaps \color{Blue} A}{\color{Red} d} {\color{Red} Vesperas. Super psalmos. Ant.} Sanctificavit dominus tabernaculum suum, hec est domus domini in qua invocetur nomen eius, de qua scriptum est, erit nomen meum ibi, dicit dominus. {\color{Red} Ps.} \psalm{112} \ul{et ceteri.} {\color{Red} Capitulum.}
\lettrine[lines=2]{\zallmancaps \color{Red} V}{idi} civitatem sanctam Hierusalem, novam descendentem de celo a deo paratam, sicut sponsam ornatam viro suo. \R \hyperlink{terribilis-est-locus}{Terribilis.} {\color{Red} IX. Ymnus.}
\hymn{urbs-beata-hierusalem}{Urbs beata Hierusalem}{
\lettrine[lines=2]{\zallmancaps \color{Blue} U}{rbs} beata Hierusalem
dicta pacis visio,
que construitur in celis
vivis ex lapidibus,
% et angelis coronata %typo: coornata
et angelis coornata
ut sponsata comite.
\par {\bfseries \color{Red} N}ova veniens e celo
nuptiali thalamo,
%pp118
preparata ut sponsata,
copulata domino,
platee et muri eius
ex auro purissimo.
\par {\bfseries \color{Blue} P}orte nitent margaritis
aditis patentibus,
et virtute meritorum
illuc introducitur,
omnes qui pro Christi nomine
hic in mundo premitur.
\par {\bfseries \color{Red} T}unsionibus pressuris
expoliti lapides,
suis coaptantur locis
per manus artificis,
disponuntur permansuri
sacris edificiis.
\par {\bfseries \color{Blue} G}loria et honor deo
usquequo altissimo, %typo: altisimo
una patri filioque,
inclito Paraclyto,
cui laus est et potestas
per eterna secula. Amen.
}
\newline \V Domum tuam domine decet sanctitudo. \R In longitudinem dierum.
\newline {\color{Red} Ad Mag. Ant.} O quam metuendus est locus iste, vere non est hic aliud nisi domus dei et porta celi. {\color{Red} In primo Anno. Oratio.}
\lettrine[lines=2]{\zallmancaps \color{Red} D}{eus} qui invisibiliter omnia contines, et tamen pro salute generis humani signa tue potentie visibiliter ostendis, templum hoc potentia tue inhabitationis illustra, ut omnes qui huc deprecaturi conveniunt, ex quacumque tribulatione ad te clamaverint, consolationis tue beneficia consequantur. Per. %typo benebicia
{\color{Red} In secundo anno et deinceps. Oratio.}
\lettrine[lines=2]{\zallmancaps \color{Blue} D}{eus} qui nobis per singulos annos huius sancti templi tui consecrationis reparas diem, et sacris semper mysteriis representas incolumes, exaudi preces populi tui, et presta ut quisquis hoc templum beneficia petiturus ingreditur, cuncta se impretrasse letetur. Per.
\newline {\color{Red} Ad Matutinas. Invitat.}
\begin{invitatory}[exultemus-domino-regi-invitatorium]
{Exultemus domino regi summo, * Qui suum sanctificavit tabernaculum.}
\end{invitatory}
\newline {\color{Red} Ymnus.} \hyperlink{urbs-beata-hierusalem}{Urbs beata.}
\newline {\color{Red} In Primo Nocturno. Ant.} Tollite portas principes vestras et elevamini porte eternales. {\color{Red} Ps.} \psalm{23} {\color{Red} Ant.} Vidit Iacob scalem, summitas eius celos tangebat, et descendentes angelos, et dixit, vere locus iste sanctus est. {\color{Red} Ps.} \psalm{45} {\color{Red} Ant.} Cum evigilasset Iacob de summo ait, vere locus iste sanctus est. {\color{Red} Ps.} \psalm{47} \V Domum tuam domine.
% https://artflsrv04.uchicago.edu/philologic4.7/PLD/navigate/1924/2
\newline \ul{Ricardus in exceptionibus, parte secunda, libro decimo.} {\color{Red} Lectio Prima.}
\lettrine[lines=2]{\zallmancaps \color{Red} S}{anctificavit} tabernaculum suum altissimus. Habet karissimi tabernaculum domini idest ecclesia sancta, lapides, cementum, et alia multa que sacramentis sunt plena, et spiritualia preferunt documenta. Singuli lapides sunt singuli fideles, quadri stabilitate fidei, et firmi virtute patiendi. Cementum est caritas, que singulos coaptat atque vivificat: et ne per aliquam discordiam invicem discrepent, uniformiter equat. \R
\lettrine[lines=2]{\zallmancaps \color{Blue} I}{n} \hypertarget{in-dedicatione-templi}{\label{in-dedicatione-templi}} dedicatione templi decantabat populus laudem, * Et in ore eorum dulcis resonabat sonus. \V Obtulerunt sacrificium super altare domini et ceciderunt in facias suas et adoraverunt eum. Et in.
\newline {\color{Red} Lectio Secunda.}
\lettrine[lines=2]{\zallmancaps \color{Red} F}{undamentum} apostoli sunt et prophete, sicut scriptum est. Superedificati, super fundamentum apostolorum et prophetarum. Parietes sunt contemplativi: fundamento vicini, terrena deserentes, celestibus adherentes. Tectum in edificio spirituali non ut immateriali sursum eminet sed deorsum pendet. Tectum ergo sunt activi terrenis actionibus proximi, propter inperfectionem suam minus celestibus intendentes, ac necessitati proximorum terrenas res administrantes.
\begin{responsory}[fundata-est-domus]
{Fundata est domus domini super verticem montium et exaltata est super omnes colles, et venient ad eam omnes gentes, * Et dicent gloria tibi domine.}{Venientes autem venient cum exultatione portantes manipulos suos.}
\end{responsory}%
\newline {\color{Red} Lectio Tertia.}
\lettrine[lines=2]{\zallmancaps \color{Blue} L}{ongitudo} ecclesie consideratur secundum diuturnitatem populorum, altitudo secundum differentiam meritorum. In longitudine namque porrigitur, dum secundum tria tempora, scilicet naturalis,et scripte legis, et gratie: a primo iusto usque ad ultimum extenditur. In latitudinem vero distenditur, dum multorum populorum numerositate dilatatur. In altitudinem erigitur, dum per minorum et maiorum differentiam elevatur: dum super laicos ordo sacerdotalis, et super sacerdotalem episcopalis, et super hunc archiepiscopalis: super omnes denique papa collocatur.
\begin{responsory-final}[mane-surgens-iacob]
{Mane surgens Iacob erigebat lapidem in titulum, fundens oleum desuper votum vovit domino, * Vere locus iste sanctus est, * Et ego nesciebam.}{Vidit Iacob scalem, summitas eius celos tangebat, et descendentes angelos et ait.}{Et ego}
\end{responsory-final}%
\newline {\color{Red} In Secundo Nocturno. Ant.} Non est hic aliud nisi domus dei et porta celi. {\color{Red} Ps.} \psalm{83} {\color{Red} Ant.} Erexit Iacob lapidem in titulum, fundens oleum desuper. {\color{Red} Ps.} \psalm{84} {\color{Red} Ant.} Erit michi dominus in deum, et lapis iste vocabitur domus dei. {\color{Red} Ps.} \psalm{86} \V Hec est domus domini firmiter edificata. \R Bene fundata est supra firmam petram.
\newline {\color{Red} Lectio Quarta.}
\lettrine[lines=2]{\zallmancaps \color{Red} S}{acrarium} ecclesie, significat ordinem virginum, chorus: ordinem continentium, navis: ordinem coniugatorum. Strictius enim est sanctuarium quam chorus, et chorus quam navis, sic pauciores sunt virgines quam continentes, et isti quam coniugati. Sacratior eciam est locus sanctuarium quam chorus, et chorus quam navis. Sic et dignior est chorus virginum quam ordo continentium, et ordo continentium dignior est quam ordo coniugatorum.
\begin{responsory}[benedic-domine-domum]
{Benedic domine domum istam quam edificavi nomini tuo, venientium in locum istum, * Exaudi preces in excelso solio glorie tue.}{Domine si conversus fuerit populus tuus et oraverit ad sanctuarium tuum.}
\end{responsory}%
\newline {\color{Red} Lectio Quinta.}
\lettrine[lines=2]{\zallmancaps \color{Blue} A}{trium} sunt falsi Christiani, qui sanctificati sunt et baptizati, sed pleni putredine cadaverum, idest corruptione vitiorum. De hoc atrio scriptum est. Atrium quod foris est: noli metiri, quia proiectum est, et datum gentibus. Falsi namque Christiani damnabuntur cum gentibus, et dabuntur ad puniendum demonibus. Altare Christus est, super quem non solum sacrificia bonorum operum, sed eciam orationum offerimus: dicentes. Per dominum nostrum Ihesum Christum filium tuum.
\begin{responsory}[o-quam-metuendus]
{O quam metuendus est locus iste, vere, * Non est hic aliud nisi domus dei et porta celi.}{Mane surgens Iacob votum vovit domino et dixit.}
\end{responsory}%
\newline {\color{Red} Lectio Sexta.}
\lettrine[lines=2]{\zallmancaps \color{Red} T}{urris} est nomen domini, sicut scriptum est. Turris fortissima nomen domini: ad eam fugiet iustus et salvabitur. Cimbala vero predicatores intelliguntur: qui verba Dei loquuntur. Fenestre vitree sunt viri spirituales, per quos divina cognitio nobis illucet. Interior dealbatura significat mundiciam cordium, exterior, vero mundiciam corporum. Duodecim cerei, duodecim sunt apostoli, qui vexillum crucis et fidem passionis Christi predicaverunt per quatuor partes mundi.
\begin{responsory-doxology}[orantibus-in-loco]
{Orantibus in loco isto dimitte domine peccata populi tui, et ostende eis viam bonam per quam ambulent, * Et da gloriam nomini tuo.}{Domine exaudi orationem meam et clamor meus ad te veniat.}
\end{responsory-doxology}%
\newline {\color{Red} In Tertio Nocturno. Ant.} Qui habitat in adiutorio altissimi, in protectione dei celi commorabitur. {\color{Red} Ps.} \psalm{90} {\color{Red} Ant.} Domum istam protege domine, et angeli tui custodiant muros eius. {\color{Red} Ps.} \psalm{95} {\color{Red} Ant.} Fundata est domus domini super verticem montium, et exaltata est super omnes colles. {\color{Red} Ps.} \psalm{96} \V Beati qui habitant in domo tua domine. \R In secula seculorum laudabunt te.
\newline {\color{Red} Lectio VII. Secundum Lucam.}
\lettrine[lines=2]{\zallmancaps \color{Blue} I}{n} illo tempore: Ingressus Ihesus perambulabat Hiericho. Et ecce vir nomine Zacheus, et hic erat princeps publicanorum, et ipse dives. Et reliqua.
\newline {\color{Red} Omelia Venerabilis Bede Presbiteri.}
\lettrine[lines=2]{\zallmancaps \color{Red} Q}{ue} impossibilia sunt apud homines: possibilia sunt apud Deum. Ecce enim camelus deposita gibbi sarcina, per foramen acus transit, hoc est dives et publicanus relicto onere divitiarum contemptoque censu fraudium: angustam portam, arctamque viam que ad vitam ducit, ascendit.
\begin{responsory}[lapides-preciosi]
{Lapides preciosi omnes muri tui, * Et turres Hierusalem gemmis edificabuntur.}{Hec est domus domini firmiter edificata bene fundata est supra firmam petram.}
\end{responsory}%
\newline {\color{Red} Lectio VIII.}
\lettrine[lines=2]{\zallmancaps \color{Blue} Q}{ui} mira devotione fidei ad videndum salvatorem, quod natura minus habuerat ascensu supplet arboris, atque ideo iuste, quamvis ipse rogare non audeat, benedictionem dominice susceptionis quam desiderabat accipit. Mistice autem Zacheus qui interpretatur iustificatus, credentem ex gentibus populum significat: qui quanto curis secularibus occupatior, tanto flagitiis deprimentibus erat factus humilior. Sed ablutus est, sed sanctificatus, sed iustificatus in nomine domini nostri Ihesu Christi, et in spiritu Dei nostri.
\begin{responsory}[domus-mea-domus]
{Domus mea domus orationis vocabitur dicit dominus, * In ea omnis qui petit accipit et qui querit invenit, et pulsanti aperietur.}{Domum tuam domine decet sanctitudo in longitudinem dierum.}
\end{responsory}%
\newline {\color{Red} Lectio IX.}
\lettrine[lines=2]{\zallmancaps \color{Red} Q}{ui} intrantem Hiericho salvatorem videre querebat, sed pre turba non poterat, quia gratie fidei quam mundo salvator attulit participare cupiebat, sed inolita vitiorum consuetudo ne ad votum perveniret obstiterat. Eadem namque turba noxie consuetudinis, que supra cecum clamantem ne lumen peteret increpabat: eciam suspicientem publicanum, ne Ihesum videat, tardat. Sed sicut cecus turbarum voces magis ac magis clamando devicit: ita pusillus necesse est turbe nocentis obstaculum altiora petendo transcendat, terrena relinquat, arborem crucis ascendat.
\begin{responsory-final}[terribilis-est-locus]
{Terribilis est locus iste, non est hic aliud nisi domus dei et porta celi, * Vere etenim dominus est in loco isto, * Et ego nesciebam.}{Cumque evigilasset Iacob quasi de gravi somno ait.}{Et ego}
\end{responsory-final}%
\newline \V Domine dilexi decorem domus tue. \R Et locum habitationis glorie tue.
\newline {\color{Red} In Laudibus. Ant.} Domum tuam domine decet sanctitudo in longitudinem dierum. {\color{Red} Ps.} \psalm{92} {\color{Red} Ant.} Hec est domus domini firmiter edificata, bene fundata est supra firmam petram. {\color{Red} Ant.} Domus mea domus orationis vocabitur. {\color{Red} Ant.} Bene fundata est domus domini supra firmam petram. {\color{Red} Ant.} Lapides preciosi omnes muri tui et turres Hierusalem gemmis edificabuntur. {\color{Red} Capitulum.} Vidi civitatem. {\color{Red} Ymnus.}
\hymn{angulare-fundamentum}{Angulare fundamentum}{
\lettrine[lines=2]{\zallmancaps \color{Blue} A}{ngulare} fundamentum
lapis Christus missus est,
qui compage parientum
in utroque nectitur,
quem Syon sancta suscepit,
in quo credens permanet.
\par {\bfseries \color{Red} O}mnes illa deo grata
et dilecta civitas:
plena modulis in laude
et canoro iubilo,
trinum deum unicumque
cum favore predicat.
\par {\bfseries \color{Blue} H}oc in templo summe deus
exoratus adveni,
et clementi bonitate
precum vota suscipe,
largam benedictionem
hic infunde iugiter.
\par {\bfseries \color{Red} H}ic promereantur omnes
petita accipere,
et adepta possidere
cum sanctis perenniter,
paradisum introire
translati in requiem.
\par {\bfseries \color{Blue} G}loria et honor deo
usquequo altissimo,
una patri filioque,
inclito Paraclyto,
cui laus est et potestas
per eterna secula. Amen.
}
\newline \V Domus mea. \R Domus orationis vocabitur.
\newline {\color{Red} Ad Ben. Ant.} Mane surgens Iacob erigebat lapidem in titulum fundens oleum desuper, votum vovit domino, vere locus iste sanctus est et ego nesciebam. \ul{Ad horas antiphone Laudum.}
\newline {\color{Red} Ad Terciam. Cap.} Vidi civitatem.
\begin{responsory-breve}[domum-tuam]
{Domum tuam domine decet sanctitudo, * Alleluya alleluya.}{In longitudinem dierum.}
\end{responsory-breve}%
\V Hec est domus.
\newline {\color{Red} Ad Sextam. Capitulum.}
\lettrine[lines=2]{\zallmancaps \color{Blue} E}{cce} tabernaculum dei cum hominibus et habitabit cum eis et ipsi populus eius erunt, et ipse deus cum eis erit eorum deus.
\begin{responsory-breve}[hec-est-domus]
{Hec est domus domini firmiter edificata, * Alleluya alleluya.}{Bene fundata supra firmam petram.}
\end{responsory-breve}%
\V Beati qui habitant.
\newline {\color{Red} Ad Nonam. Capitulum.}
\lettrine[lines=2]{\zallmancaps \color{Red} F}{undamentum} aliud nemo potest ponere, preter id quod positum est, quod est Christus Ihesus.
\begin{responsory-breve}[beati-qui-habitant]
{Beati qui habitant In domo tua domine, * Alleluya alleluya.}{In secula seculorum laudabunt te.}
\end{responsory-breve}%
\V Domus mea.
\newline {\color{Red} Ad Vesperas. Super psalmos. Ant.} Domum tuam. {\color{Red} Ps.} \psalm{115} {\color{Red} Ps.} \psalm{121} {\color{Red} Ps.} \psalm{126} {\color{Red} Ps.} \psalm{146} {\color{Red} Ps.} \psalm{147}
\newline \ul{Cap., Ymnus, \Vbar , ut in Primis Vesperis.}
\newline {\color{Red} Ad Mag. Ant.} Zachee festinans descende, quia hodie in domo tua oportet me manere, at ille festinans descendit, et suscepit illum gaudens in domum suam, hodie huic domui salus a deo facta est.
\newline \ul{Per octavas et in octava, Invitat. de die, Ad Ben. et ad Mag. antiphone de die. Cetera fiant secundum modum octavarum de sanctis. Et legantur tres sequentes lectiones cum sex supra notatis de sermone. Ex supradicto sermone, Ricardus in exceptionibus parte secunda, libro decimo.}
\newline {\color{Red} Lectio.} \grecross
\lettrine[lines=2]{\zallmancaps \color{Blue} H}{abet} igitur ecclesia singulos lapides per singulos fideles, cementum per caritatem, fundamentum per prophetas et apostolos, parietes per contemplativos, tectum per activos. Habet longitudinem, per diuturnitatem temporum: latitudinem per multitudinem populorum: altitudinem per differentiam meritorum. Habet sacrarium per ordinem virginum, chorum per ordinem continentium, navem per ordinem coniugatorum. Habet altare per Christum: turrim per nomen domini, Cymbala per predicatores: Vitreas fenestras per viros spirituales. Habet interiorem dealbaturam per mundiciam cordium, exteriorem per mundiciam corporum: duodecim cereos per duodecim apostolos.
\newline {\color{Red} Lectio.}
\lettrine[lines=2]{\zallmancaps \color{Red} S}{equitur} pontifex Christum significans, qui circuivit ecclesiam suam primum in tempore naturalis legis, eam admonendo per patriarchas: de hinc in tempore scripte legis, admonendo per prophetas. Tandem in tempore gratie per semetipsum eandem circuivit et introivit, foris scilicet erudiendo per doctrinam, et intus sanctificando per gratiam. In primo circuitu dixit carnem cum sanguine non comedetis, in secundo non occidetis: in tercio vero, qui irascitur fratri suo reus est iudicio.
\newline {\color{Red} Lectio.}
\lettrine[lines=2]{\zallmancaps \color{Blue} M}{ulti} autem estimantur in hoc edificio et spirituali per sanctitatem contineri, qui tamen ab ipso longe per pravitatem sunt remoti. Extra ipsum quippe sunt, immundi, fornicarii, ebriosi, feneratores, avari, rapaces, odiosi, mendaces, et qui fratri suo fatue dicunt, et qui mulierem ad concupiscendum aspiciunt. Non enim sunt illius edificii lapides: qui caritatis non sunt participes. Talem ergo vitam fratres agamus: ut Dei lapides esse possimus.
\newline \ul{Tempore Paschali sic fiat officium. Ad Vesperas ut supra. Ad Matutinas, Invitat., Ymnus, ut supra. Tres psalmi cum antiphonis et versiculo et tribus responsorii, secundum feriam et ordinem nocturnorum. In Laudibus et ad horas ut supra. Post Gloria vero responsorium repetatur, Alleluya, et antiphone omnis et versiculi et responsoria terminentur cum Alleluya.}

\ubsection{Rubrice Sanctorum}{rubrice-sanctorum}

{\ubsubsection{De Festis in Communi}{de-festis-in-communi-breviarium}}
\lettrine[lines=2]{\zallmancaps \color{Red} I}{}\ul{\textsc{n} quacumque die festum simplex vel maius occurrerit: in ea celebrabitur, preterquam in Dominicis Adventus, et preterquam in Septuagesima et Dominics sequentibus usque ad Ramos inclusive, et preterquam in die Cinerum, et preterquam in secunda feria post Ramos et deinceps usque ad Dominicam in Albis inclusive. Et preterquam in die Ascensionis et in Octava eius, et in vigilia Pentecostes et deinceps usque ad festum Trinitatis inclusive. Excipitur festum Purificationis quod nunquam transfertur, et festum Annunciationis, de quo fiat sicut infra in rubrica \textit{De Translatione Festivitatum} notatum est.}
\newline \ul{Festum eciam trium lectionum fiat in quacumque die occurrerit nisi in Dominica, vel nisi festum simplex vel maius in eo fuerit celebrandum. Et nisi quando hystoria} \hyperlink{domine-ne-in-ira}{Domine ne in ira.} \ul{propter brevitatem temporis inter Octavas Epiphanie et Septuagesima in festo trium lectionum fuerit cantanda. Et nisi in quarta feria Cinerum et deinceps usque ad Octavas Pasche inclusive. Et nisi in prima die Rogationum, et deinceps usque ad Octavas Ascensionis inclusive, et in vigilia Pentecostes et deinceps usque ad,} \hyperlink{deus-omnium}{Deus omnium.} \ul{Et nisi in Sabbatis quando de Beata Virgine agitur in conventu et in ieiuniis Quatuor Temporum Septembris.}

{\ubsubsection{De Concomitantia Festivitatum}{de-concomitantia-festivitatum-breviarium}}
\lettrine[lines=2]{\zallmancaps \color{Blue} S}{}\ul{\textsc{i} festum simplex vel maius post festum totum duplex immediate evenerit, primum festum scilicet totum duplex in suo die habebit totas Vesperas, et de festo sequenti memoria tantum fiet. Excipitur festum Beati Petri Martiris et festum Translationis Beati Dominici, quando aliquod eorum in vigilia Ascensionis occurrerit, quia tunc secundas Vesperas non habebit.}
\newline \ul{Si vero festum totum duplex in crastino festi simplicis vel semiduplicis aut duplicis occurrerit celebrandum: festum totum duplex in primis Vesperis a principio inchoabitur, et de festo illius diei memoria tantum fiet.}
\newline \ul{Si autem festum semiduplex occurrerit post festum semiduplex vel simplex, in die primi festi in Vesperis psalmi de ipso primo festo erunt, antiphona vero super psalmos et capitulum et que sequuntur erunt de festo sequenti et de festo illius diei fiet memoria.}
\newline \ul{Quando vero festum simplex post festum semiduplex evenerit: festum semiduplex habebit in die suo totas Vesperas, et de festo simplici sequenti memoria tantum fiet.}
\newline \ul{Quando autem unum festum simplex post aliud simplex evenerit: primum festum in die suo habebit in Vesperis antiphonam et psalmos. Capitulum vero et que sequuntur de sequenti festo erunt, et memoria fiet de primo festo.}
\newline \ul{Quod de festo simplici et semiduplici supradictum est: intelligendum est quando primum festum habuit primas Vesperas. Si enim propter aliquod impedimentum, festum semiduplex vel simplex, non habuisset primas Vesperas: habere debet secundas, nisi festum duplex vel totum duplex sequatur.}
\newline \ul{Si festum simplex vel maius in quocumque Sabbato fuerit celebrandum: habebit festum in ipso Sabbato totas Vesperas, eciam si Dominica secundas Vesperas non habeat, nisi in Sabbatis Adventus, et in Sabbato ante Septuagesimam et deinceps in Sabbatis usque ad Pascha. In quibus festum simplex vel semiduplex in Vesperis habebit psalmos et antiphonam, a capitulo vero inchoabitur de Dominica, et memoria postea de festo fiet. Festum vero duplex vel totum duplex omni tempore habebit totas Vesperas, et de Dominica sequenti memoria tantum fiet.}
\newline \ul{Quando autem festum simplex post Dominicam occurrerit: psalmi et antiphone de Dominica erunt. Capitulum et que sequuntur de festo. Si vero semiduplex vel duplex fuerit: psalmi de Dominica erunt, antiphona vero super psalmos et capitulum et que sequuntur de festo sequenti erunt. Festum vero totum duplex a principio inchoabitur, et habebit totas Vesperas, et de Dominica fiet memoria.}
\newline \ul{Sciendum autem quod Dominica infra octavas in omnibus censetur sicut festum simplex, nisi quantum ad primas vesperas quando festum aliquod simplex infra octavas in Sabbato evenerit.}

{\ubsubsection{De Translatione Festivitatum}{de-translatione-festivitatum-breviarium}}
\lettrine[lines=2]{\zallmancaps \color{Red} S}{}\ul{\textsc{i} festum simplex vel maius in aliqua Dominica Adventus evenerit: vel in Dominica Septuagesime et deinceps in aliqua Dominica usque ad Ramos exclusive in secundam feriam transferatur, excepto festo Purificationis quod non transfertur: et festo Sanctorum Fabiani et Sebastiani quod in precedenti Sabbato celebrandum est, et festo Sancte Agnetis quod in crastinum Sancti Vincentii transferatur, si aliquod ipsorum in Dominica Septuagesime evenerit. Quodcumque autem festum simplex vel maius in die Cinerum occurrerit: in precedenti tertia feria celebretur.}
\newline \ul{Quando vero festum simplex vel supra in Dominica de Ramis Palmarum vel deinceps usque ad octavas Pasche inclusive evenerit: si simplex fuerit: illo anno nichil de eo fiat nisi memoria in Dominica et in secunda et tercia feria in utrisque Vesperis et in Laudibus, et in quarta feria in Laudibus. Semiduplex autem et maius in ebdomadam post octavas Pasche transferatur, preter festum Annunciationis quod si certo die generaliter ab omnibus fuerit in Patria celebrandum, celebretur eciam a fratribus ipso die, alioquin sicut dictum est de aliis: transferatur.}
\newline \ul{Festivitates vero que in predictam ebdomadam transferentur, celebrentur per ferias vacantes ipsius septimane eo ordine quo fuerint intermisse, ita tamen quod si festum aliquod nisi trium fuerit lectionum, aliqua die predicte ebdomade occurrerit, in sua die celebretur, nec propter aliud festum extra suam diem transferatur. Quodlibet autem festum quod in secunda feria et deinceps per hanc ebdomadam fuerit celebrandum: habebit primas Vesperas et in secundis fiet de eo memoria, si aliud immediate fuerit celebrandum nisi sit totum duplex quod habebit eciam secundas Vesperas, in quibus fiet tantum memoria de festo celebrando post ipsum immediate. Festum vero sequens post totum duplex habebit secundas Vesperas et sic deinceps: si plura successive et immediate occurrerint celebranda.}
\newline \ul{Quod si festum simplex vel maius in die Ascensionis vel in Octava eius evenerit: in sequentem
%pp119
sextam feriam transferatur. Si vero in vigilia Pentecostes: in precedenti sexta feria celebretur. Si autem in die Pentecostes vel deinceps usque ad festum Trinitatis inclusive: festum semiduplex vel maius occurrerit: transferatur in sequentem feriam post festum Trinitatis. Quod si simplex fuerit illo anno nichil de eo fiat.}

{\ubsubsection{De Festo Trium Lectionum}{de-festo-trium-lectionum-breviarium}}
\lettrine[lines=2]{\zallmancaps \color{Red} S}{}\ul{\textsc{ciendum} quod de festo trium lectionum nichil fit in Cena Domini et deinceps usque ad Dominicam in Albis inclusive, et in die Pentecostes et deinceps usque ad diem Trinitatis inclusive.
Item quandocumque festum totum duplex in die festi trium lectionum fuerit celebrandum.
Omnibus autem diebus Dominicis, et quando tempus ita breve fuerit inter Octavas Epyphanie et Septuagesimam: quod in feriis intervenientibus oportet cantari hystoriam} \hyperlink{domine-ne-in-ira}{Domine ne in ira:} \ul{et in quarta feria Cinerum et deinceps usque ad quartam feriam post Ramos inclusive et in tribus diebus Rogationum, et infra octavas Ascensionis et in Octava, et in vigilia Pentecostes et infra octavas Trinitatis, et infra octavas dedicationis et in octava die, tantum fit memoria de festo trium lectionum in Vesperas et in Laudibus, nisi quando in predictis temporibus festum totum duplex in die festi trium lectionum fuerit celebrandum.
Fit eciam tantum memoria de festo trium lectionum in Sabbatis quando de Beata Virgine agitur in conventu, et in ieiuniis quatuor temporum Septembris, et quando festum semiduplex propter translationem in die ipsius fuerit celebrandum.
Quando vero festum trium lectionum in crastino Dominice vel festi simplicis aut semiduplicis debuerit celebrari, in primis Vesperis fiet tantum de ipso memoria, et sequenti die Officium de ipso erit in Matutinis, et Laudibus et Horis diei. Ad Nonam autem semper terminatur.
Si autem in crastino totius duplicis celebrari debuerit, in precedenti die nichil fiet de eo, sed sequenti die Officium de ipso agatur ut predictum est.
Quando vero in crastino ferie vacantis occurrerit: in Vesperis precedentibus ad capitulum inchoetur. Sic autem fiat. Ad Matutinas Invitatorium sicut est notatum ab uno dicatur. Antiphona una super novem psalmos de festo. Versiculus ante lectiones et tria responsoria dicantur secundum feriam de hystoria que cantaretur si festum haberet novem lectiones extra tempus Paschale, excepto festo Praxedis in quo semper dicantur tria ultima responsoria. Tempore vero Paschali dicantur tantum tres psalmi cum antiphona de festo Paschalis temporis: Versiculus vero secundum ordinem nocturni. Tria responsoria sicut in festo simplici Paschalis temporis, oratio sicut est notata. Cetera fiant sicut in festo simplici excepto quod non dicatur} Te Deum \ul{et prima antiphona Laudum sola dicitur in Laudibus super psalmos. Prostrationes autem non fiant quamdiu de eo agitur, et nunquam transferatur.}

{\ubsubsection{De Festo Simplici}{de-festo-simplici-breviarium}}
\lettrine[lines=2]{\zallmancaps \color{Blue} F}{}\ul{\textsc{estum} simplex subscripto modo celebretur, nisi maius festum precedens impediat. In primis Vesperis antiphone et psalmi secundam feriam qui dicerentur si de feria ageretur, vel psalmi de festo si precesserit, vel de Octavis quando infra Octavas evenerit, cum antiphona de festo precedenti vel de Octavis. Capitulum vero et que sequuntur de festo erunt. Responsorium in Vesperis non dicitur. Invitatorium sicut est notatum. Novem lectiones fiant extra tempus Paschale. In tempore vero Paschali, tantum tres.} Te Deum laudamus, \ul{dicatur nisi a Septuagesima usque ad Pascha, et nisi in festo Innocentum quando extra Dominicam fuerit. In Laudibus dicantur hi psalmi,} \psalm{92} \psalm{99} \psalm{62} \psalm{66} \ul{Canticum} \psalm{benedicite} \ul{psalmis} \psalm{148} \ul{cum quinque antiphonis que eciam dicantur ad Horas diei suo ordine quarta pretermissa. Et hoc generaliter observetur, scilicet quod quandocumque antiphone Laudum ad Horas dicuntur: quarta intermittatur, sive de tempore sive de sanctis agatur.
Prima vero earum in festo simplici dicatur in secundis Vesperis super psalmos preterquam in festis que a Natali Domini usque ad Octavas Epyphanie occurrerunt.}

{\ubsubsection{De Festo Semiduplici}{de-festo-semiduplici-breviarium}}
\lettrine[lines=2]{\zallmancaps \color{Red} I}{}\ul{\textsc{n} festo semiduplici in primis Vesperis una tantum dicatur antiphona de ipso festo super psalmos, psalmi feriales vel de festo si precesserit vel de Octavis si infra Octavas evenerit. Responsorium in primis Vesperis sicut est signatum. Cetera fiant sicut in festo simplici super notatum est.}

{\ubsubsection{De Festo Duplici}{de-festo-duplici-breviarium}}
\lettrine[lines=2]{\zallmancaps \color{Blue} I}{}\ul{\textsc{n} omni festo duplici, in primis Vesperis una tantum antiphona de ipso festo super psalmos dicatur, et psalmi sicut in semiduplicibus, nisi in festo Circumcisionis. Responsorium in primis Vesperis sicut est signatum. In secundis vero Vesperis dicantur psalmi sub una antiphona: nisi in festo Circumcisionis.
Cetera fiant sicut in festo simplici est notatum, excepto quod in Completorio et in Prima preces non dicantur, sed in Completorio dicta antiphona post} Nunc dimittis \ul{et in Prima dicto,} Confiteor: \ul{statim subiungatur} Dominus vobiscum, \ul{et Oratio.}

{\ubsubsection{De Festo Toto Duplici}{de-festo-toto-duplici-breviarium}}
\lettrine[lines=2]{\zallmancaps \color{Red} F}{}\ul{\textsc{estum} totum duplex fiat in die Pasche et duobus sequentibus, In Ascensione Domini, In die Pentecostes et duobus sequentibus, In festo Trinitatis, in die Consecrationis Ecclesie et in Anniversario eiusdem, et in festo Patroni Ecclesie Nostre, et in ceteris que in Kalendario sunt notata.
In omnibus autem festivitatibus totis duplicibus, ad primas Vesperas dicantur hi psalmi scilicet,} \psalm{112} \psalm{116} \psalm{145} \psalm{146} \psalm{147} \ul{nisi in festo sancti Stephani et sancti Iohannis et Epyphanie et nisi in festo Pasche et duobus diebus sequentibus et nisi in festo Pentecostes et duobus diebus sequentibus, et nisi in festo Translationis Beati Dominici, quando post Ascensionem immediate celebrabitur. Responsorium sicut et quando est signatum. Antiphona ad Mag. et ad Ben. tota dicatur ante cantici inceptionem.} Te Deum, \ul{in Matutinis omni tempore dicatur. Responsoria horarum cum Alleluya dicantur, nisi in Completorio Sabbati ante Septuagesimam et deinceps usque ad Pascha. Versiculi vero nunquam dicantur cum Alleluya, nisi in tempore Paschali.
Cetera omnia fiant sicut supra in festo duplici est notatum, preterquam in die Natali et sancti Stephani et sancti Iohannis, et Epyphanie in quibus in secundis Vesperis dicuntur quinque antiphone super psalmos et quod secunda feria et tercia post Pascha et Pentecosten, una tantum antiphona dicitur in Laudibus super psalmos.}

{\ubsubsection{De Octavis Sanctorum}{de-octavis-sanctorum-breviarium}}
\lettrine[lines=2]{\zallmancaps \color{Blue} F}{}\ul{\textsc{estivitates} sanctorum habentium Octavas sunt he: Sancti Andree, Sancti Stephani, Sancti Iohannis Evangeliste, Sanctorum Innocentium, Sancti Iohannis Baptiste, Apostolorum Petri et Pauli, Beati Dominici, Sancti Laurentii, Assumptionis Beate Virginis, Sancti Augustini, Nativitatis Beate Marie, Sancti Martini.
Preter istas Octavas nulle alie Octave sanctorum fiant sive patronorum ecclesie, sive aliorum quorumcumque. Sic autem fiat Officium per Octavas sanctorum.
Per omnes ferias infra Octavas sanctorum quando de Octavis agitur. Ad matutinas Invitatorium sicut in suis locis notatum est, Ymnus sicut in die. Una tantum antiphona de festo super novem psalmos. Prima die post festum dicatur prima antiphona de nocturno, secunda die: secunda, et sic deinceps quot fuerint necessarie. Versiculus et tria responsoria de festo secundum feriam. Est autem dividenda hystoria per ferias eo modo quo supra de dominicali hystoria est notatum.} Te Deum. \ul{cotidie dicatur. Versiculus ante Laudes sicut in die.
In Laudibus super psalmos prima antiphona Laudum. Capitulum, Ymnus, Versus, de die. Ad Ben. antiphona, sicut est notata. Oratio sicut in die. Ad horas diei, antiphona, Capitula, responsoria, versiculus, Oratio, de festo, excepto quod si festum fuit totum duplex responsoria tamen non dicuntur cum Alleluya infra Octavas. Ad Vesperas antiphona, psalmi, Capitulum, hymnus, versiculus de festo. Ad Mag. antiphona, sicut est notata. Oratio de festo. Excipiuntur de predictis, Octave Apostolorum Petri et Pauli et sancti Andree.}
\newline \ul{In Octava autem die ad Vesperas precedentes psalmi de Octava, super psalmos prima antiphona Laudum de festo. Cetera fiant sicut in festo servato modo festi simplicis. Excipiuntur Octava sancti Laurentii propter, versus antiphonarum qui non dicuntur, et Octava Apostolorum Petri et Pauli, et Octave sancti Andree, sancti Stephani, sancti Iohannis, sanctorum Innocentum.
Octava eciam sancti Iohannis Baptiste in secundis Vesperis habet propriam antiphonam Ad Mag. Dominica vero infra Octavas fiat sicut de Octava die notatum est, nisi quod antiphone Ad Ben. et Ad Mag. de die festi non dicantur si alie per Octavas notate sint. Excipitur Dominica infra Octavas Beati Augustini, quando in Kalendis Septembris evenerit.
Quando autem omelia Dominicali legitur in aliqua feria infra Octavas, fiat totum aliud Officium de Octavis.}

{\ubsubsection{De Memoriis Faciendis}{de-memoriis-faciendis-breviarium}}
\lettrine[lines=2]{\zallmancaps \color{Red} P}{}\ul{\textsc{er} totum annum quando propter festum vel Octavas: de Dominica non fit plenum Officium: semper facienda est de ipsa memoria Sabbato precedenti in Vesperis, et in ipsa die in Laudibus et in Vesperis: preterquam in vigilia Natalis Domini in Vesperis, et in die Natalis, et Circumcisionis, et deinceps quandocumque Dominica fuerit usque in crastinum Epyphanie exclusive. Excipitur Dominica infra Octavas Natalis Domini de qua nichil faciendum est nisi in crastino sancti Thome, vel in die sancti Silvestri. Et Dominice infra Octavas Epyphanie et Ascensionis, de quibus in Sabbato precedenti nulla est memoria facienda.
Fiat eciam memoria de tempore per omnes ferias Adventus et Quadragesime, et tribus diebus Rogationum, quando festum simplex vel maius in predictis temporibus occurrerit celebrandum, aliis autem temporibus per totum annum de tempore per ferias nulla fiat memoria: quando festum infra ebdomadam contigerit celebrari.
Item de omnibus festis simplicibus et supra: si non habuerint integre primas Vesperas, vel secundas, memoria in utrisque Vesperis semper fiat. Excipitur festum commemorationis Beati Pauli: de quo in primis Vesperis nichil sit specialiter. Item de festo trium lectionum: quando non habet plenum Officium, facienda est memoria sicut supra in rubrica de ipso notatum est, facienda est eciam memoria de sanctis de quibus in Kalendario est signata: nisi quando de festo trium lectionum nichil fieret si eveniret. Quando autem infra quascumque Octavas: non fit Officium de Octavas: facienda est de ipsis memoria nisi in utrisque Vesperis et in Laudibus Assumptionis Beate Marie.
Omnibus vero Sabbatis per totum annum fiat memoria de Beata Virgine, preterquam in Sabbato infra Octavas Natalis Domini, et in Sabbato Sancto Pasche et sequenti, et Sabbato ante festum Trinitatis, et preterquam in festis duplicibus et totis duplicibus et in die animarum cum in Sabbato evenerint, et preterquam in Sabbatis in quibus de ipsa agitur in conventu, et in Sabbatis infra Octavas Assumptionis et Nativitatis eiusdem quando agitur de Octavis.
Fiat eciam de ipsa memoria in Octava sancti Stephani in Laudibus et deinceps tam in Vesperis quam in Laudibus, usque ad Vesperas Octavarum Epiphanie inclusive, nisi in vigilia et in festo Epyphanie.
De Cruce vero facienda est memoria cotidie per ebdomadam Pasche in Vesperis et Sabbato in Albis ad Vesperas et deinceps cotidie in Laudibus et in Vesperis, usque ad Vesperas vigilie Ascensionis exclusive, nisi in Inventione ipsius, et in festis totis duplicibus.
De Beato autem Dominico fiat cotidie memoria per totum annum nisi in vigilia Natalis Domini in Laudibus, et deinceps usque ad Octavam Epyphanie inclusive, et nisi in quarta feria post Ramos in Vesperis et deinceps usque ad Dominicam in Albis inclusive, et in vigilia Pentecostes in Vesperis, et deinceps usque ad festum Trinitatis, et nisi in festis duplicibus et totis duplicibus, et infra Octavas eiusdem quando de ipsis agitur et nisi in die animarum in Laudibus.}

{\ubsubsection{De Ordine Memoriarum}{de-ordine-memoriarum-breviarium}}
\lettrine[lines=2]{\zallmancaps \color{Blue} M}{}\ul{\textsc{emoria} de festo simplici et supra: precedit memoriam de Octavis et de Dominica, memoria vero de Octavis precedit memoriam de Dominica, precedit eciam memoriam de Beata Virgine in sexta feria et Sabbato.
Memoria vero de Dominica, et memoria de Beata Virgine in sexta feria et Sabbato: precedit memoriam trium lectionum, et festi habentis solam memoriam. Item memoria de Adventu et Quadragesima: precedit memoria de Beata Virgine in sexta feria et Sabbato, memoria autem de Beata Virgine in Adventu: precedit memoria de sancto Andrea in sexta feria et Sabbato et memoriam Crucis in Paschali tempore. Memoria vero de Cruce in Paschali tempore: fiat post omnes alias nisi post memoriam Beati Dominici, que omni tempore est post omnes alias memorias facienda nisi infra Octavam ipsius.
Quandocumque autem fuerit memoria facienda, in Vesperis et in Laudibus, per antiphonam et versiculum et Orationem fiat.}

{\ubsubsection{De Festivitatibus Extraordinariis}{de-festivitatibus-extraordinariis-breviarium}}
\lettrine[lines=2]{\zallmancaps \color{Red} F}{}\ul{\textsc{estivitates} sanctorum que sollemnes habentur in locis in quibus fratres commorantur possunt ab eis sollemniter celebrari et cantari hystorie et legende legi.}

{\ubsubsection{De Matutinis In Sero Post Completorium Dicendis}{de-matutinis-in-sero-post-completorium-dicendis-breviarium}}
\lettrine[lines=2]{\zallmancaps \color{Blue} I}{}\ul{\textsc{n} festo sancte Trinitatis et deinceps in omnibus festis duplicibus et totis duplicibus usque ad festum sancti Augustini inclusive, et in festo sancte Marie Magdalene cantentur in sero post Completorium Matutine.}

{\ubsubsection{De Orationibus et Vigiliis Sanctorum}{de-orationibus-et-vigiliis-sanctorum-breviarium}}
\lettrine[lines=2]{\zallmancaps \color{Blue} I}{}\ul{\textsc{n} omnibus festis sanctorum nisi vigiliam habuerint: Oratio que dicitur in Vesperis precedentibus diem festum: dicenda est in ipso die Ad Laudes, Terciam, Sextam et Nonam, et ad secundas Vesperas diei, nisi festum fuerit trium lectionum quod non habet secundas Vesperas.
In festis vero sanctorum habentium vigiliam, Oratio de vigilia si propria fuerit: dicenda est ad primas Vesperas, nisi in festo sancti Andree quando transfertur.
Oratio vero que dicitur in Laudibus diei: dicenda est ad Terciam, Sextam et Nonam et ad Vesperas diei. Sciendum autem quod nulla vigilia sancti alicuius habet aliquid in Matutinis nec in Horas ante Vesperas: nisi solam omeliam, propter quam tamen cum legenda fuerit: aliud Officium diei illius non mutatur nisi quandoque in benedictionibus lectionum.}

{\ubsubsection{Que Requirenda Sint in Communi Sanctorum}{que-requirenda-sint-in-communi-sanctorum-breviarium}}
\lettrine[lines=2]{\zallmancaps \color{Red} N}{}\ul{\textsc{otandum} quod per totum annum omnia que in festis sanctorum in suis locis non inveniuntur signata: de communi sanctorum Paschalis temporis sive alterius sunt assumenda, et dicenda suis locis secundum quod ibi notata sunt, et competunt festivitatibus de quibus agitur. Quandocumque autem non invenitur aliquid proprium de lectionibus in festo quocumque: recurrendum est ad commune sive propribus sive pro novem lectionibus sive pro sex quando est omelia propria. Cum autem fuerint Octave: repetende sunt lectiones predicte et in Octavis et infra quotiens fuerit necesse.}

\ubsection{Proprium Sanctorum}{proprium-sanctorum}

{\ubsubsection{In Vigilia Sancti Andree}{in-vigilia-sancti-andree-breviarium}}
\newline {\color{Red} Si Dominica fuerit legatur omelia de Dominica, alioquin de Vigilia. Lectio Prima. Secundum Iohannem.}
\lettrine[lines=2]{\zallmancaps \color{Blue} I}{n} illo tempore: Stabat Iohannes et ex discipulis eius duo. Et respiciens Ihesum ambulantem, dixit. Ecce agnus Dei. Et reliqua.
\newline {\color{Red} Omelia Venerabilis Bede Presbiteri.} \grecross
\lettrine[lines=2]{\zallmancaps \color{Red} S}{tabat} Iohannes, et respexit Ihesum: quia in magno iam culmine perfectionis animi gressum fixerat, unde iure celsitudinem divine maiestatis mereretur intueri. Stabat: quia illam virtutum arcem conscenderat, a qua nullis temptationum improbitatibus posset deici.
\newline {\color{Red} Lectio Secunda.}
\lettrine[lines=2]{\zallmancaps \color{Blue} S}{tabant} cum illo et discipuli: quia magisterium eius corde sequebantur immobili. Et respiciens inquit Ihesum ambulantem: dixit. Ecce agnus Dei. Ambulatio Ihesu dispensationem incarnationis qua ad nos venire, ac nobis exempla vivendi prebere dignatus est insinuat, sicut mansio eius eternitatem divinitatis in qua semper est patri equalis: aptissime designat.
\newline {\color{Red} Lectio Tertia.}
\lettrine[lines=2]{\zallmancaps \color{Red} E}{cce} agnus Dei. Cognovimus quare Ihesum precursor suus agnum Dei vocaverit quia videlicet eum pre ceteris mortalibus singulariter innocentem, hoc est ab omni peccato mundum novit. Et quia nos suo sanguine redempturum, atque ideo quia dona sui velleris sponte largiturum, ex quo vestem nobis nuptialem facere possimus, idest exempla vivendi nobis relicturum previdit, quibus in dilectione calefieri deberemus.
\newline {\color{Red} Ad Vesperas. Super psalmos. Ant.} Unus ex duobus qui secuti sunt dominum erat Andreas frater Symonis Petri. Alleluya. \R \hyperlink{homo-dei-ducebatur}{Homo dei.} {\color{Red} III. Ad Mag. Ant.} Ambulans Ihesus iuxta mare Galliee vidit Petrum et Andream fratrem eius et ait illis venite post me faciam vos fieri piscatores hominum, at illi relictis retibus et navi secuti sunt eum.
\lettrine[lines=2]{\zallmancaps \color{Blue} Q}{uesumus} omnipotens deus ut beatus Andreas apostolus, tuum pro nobis imploret auxilium ut a nostris reatibus absoluti a cunctis eciam periculis exuamur. Per.
\newline \ul{Si festum apostoli propter Dominicam transferatur: ad Primas Vesperas dicatur, Oratio,} \hyperlink{maiestatem-andree}{Maiestatem.}
{\ubsubsection{In Die Sancti Andree}{in-die-sancti-andree-breviarium}}
\newline {\color{Red} Ad Matutinas. Invitat.}
\begin{invitatory}[andree-invitatorium]
{Adoremus victoriosissimum regem Christum, * Qui victorem per crucis tropheum coronavit beatum Andream apostolum.}
\newline {\color{Red} In Primo Nocturno. Ant.} Vidit dominus Petrum et Andream et vocavit eos. {\color{Red} Ps.} \psalm{18} {\color{Red} Ant.} Venite post me dicit dominus faciam vos fieri piscatores hominum. {\color{Red} Ps.} \psalm{33} {\color{Red} Ant.} Relictis retibus suis, secuti sunt dominum redemptorem. {\color{Red} Ps.} \psalm{44} \V In omnem terram. \R
\end{invitatory}
\lettrine[lines=2]{\zallmancaps \color{Red} D}{um} \hypertarget{dum-perambularet}{\label{dum-perambularet}} perambularet dominus iuxta mare secus litus Galilee vidit Petrum et Andream retia mittentes in mare, vocavit eos dicens, * Venite post me faciam vos piscatores hominum. \V Erant enim piscatores et ait illis. Venite.
\begin{responsory}[mox-ut-vocem]
{Mox ut vocem domini predicantis audivit beatus Andreas: relictis retibus, quorum usu actuque vivebat eterne vite secutus est, * Premia largientem.}{Ad unius iussionis vocem Petrus et Andreas relictis retibus secuti sunt redemptorem.}
\end{responsory}%
\begin{responsory-doxology}[homo-dei-ducebatur]
{Homo dei ducebatur ut crucifigerent eum, populus autem clamabat voce magna dicens, * Innocens eius sanguis sine causa damnatur.}{Cumque carnifices ducerent eum ut crucifigeretur, factus est concursus populorum clamantium et dicentium.}
\end{responsory-doxology}%
\newline {\color{Red} In Secundo Nocturno. Ant.} Dilexit Andream dominus in odorem suavitatis alleluya. {\color{Red} Ps.} \psalm{46} {\color{Red} Ant.} Dignum sibi dominus computavit martirem, quem vocavit apostolum dum esset in mari alleluya. {\color{Red} Ps.} \psalm{60} {\color{Red} Ant.} Christus me misit ad istam provinciam, ubi non parvum populum adquisivi. {\color{Red} Ps.} \psalm{63} \V Constitues.
\begin{responsory}[doctor-bonus-et]
{Doctor bonus et amicus dei Andreas ducitur ad crucem, aspiciens a longe vidit crucem, salve crux, * Suscipe discipulum eius qui pependit in te, magister meus Christus.}{Salve crux que in corpore Christi dedicata es, et ex membris eius tanquam margaritis ornata.}
\end{responsory}%
\begin{responsory}[o-bona-crux]
{O bona crux que decorem et pulchritudinem de membris domini suscepisti, accipe me ab hominibus, et redde me magistro meo, * Ut per te me recipiat: qui per te me redemit.}{Securus et gaudens venio ad te, ita ut et tu exultans suscipias me.}
\end{responsory}%
\begin{responsory-doxology}[salve-crux-que]
{Salve crux que in corpore Christi dedicata es, et ex membris eius tanquam margaritis ornata, * Suscipe discipulum eius qui pependit in te.}{O bona crux diu desiderata et iam concupiscenti animo preparata.}
\end{responsory-doxology}%
\newline {\color{Red} In Tertio Nocturno. Ant.} Videns Andreas crucem: cum gaudio dicebat, quia amator tuus semper fui, et desideravi amplecti te, o bona crux. {\color{Red} Ps.} \psalm{74} {\color{Red} Ant.} Recipe me ab hominibus et redde me magistro meo, ut per te me recipiat qui per te me redemit. {\color{Red} Ps.} \psalm{96} {\color{Red} Ant.} Biduo vivens pendebat in cruce pro Christi nomine beatus Andreas et docebat populum. {\color{Red} Ps.} \psalm{98} \V Nimis honorati.
\newline {\color{Red} Lectio VII. Secundum Matheum.}
\lettrine[lines=2]{\zallmancaps \color{Blue} I}{n} illo tempore: Ambulans Ihesus iuxta mare Galilee vidit duos fratres, Symonem qui vocatur Petrus, et Andream fratrem eius: mittentes rete in mare. Erant enim piscatores. Et reliqua.
\newline {\color{Red} Omelia Beati Gregorii Pape.}
\lettrine[lines=2]{\zallmancaps \color{Red} A}{udistis} fratres karissimi, quia ad unius iussionis vocem, Petrus et Andreas relictis retibus secuti sunt redemptorem. Nulla vero hunc facere adhuc miracula viderant, nichil ab eo de premio eterne retributionis audierant: et tamen ad unum domini preceptum hoc quod possidere videbantur obliti sunt.
\begin{responsory}[expandi-manus-meas]
{Expandi manus meas tota die in cruce ad populum non credentem sed contradicentem michi, * Qui ambulant viam non bonam sed post peccata sua.}{Exaltare qui iudicas terram redde retributionem superbis.}
\end{responsory}%
\newline {\color{Red} Lectio VIII.}
\lettrine[lines=2]{\zallmancaps \color{Blue} Q}{uanta} nos eius miracula videmus, quot flagellis atterimur, quantis minarum asperitatibus deterremur: et tamen vocantem sequi contemnimus. In celo iam sedet, qui de conversione nos admonet, iam iugo fidei colla gentium subdidit: iam mundi gloriam stravit, iam ruinis eius crebrescentibus, districti sui iudicii diem propinquantem denuntiat, et tamen superba mens nostra non vult hoc sponte deserere quod cotidie perdit invita.
\begin{responsory}[dilexit-andream]
{Dilexit Andream dominus in odorem suavitatis, dum penderet in cruce, dignum sibi computavit martyrem quem vocavit apostolum dum esset in mari, * Et ideo amicus dei appellatus est.}{Vir iste in populo suo mittissimus apparuit, sanctitate autem et gratia plenus.}
\end{responsory}%
\newline {\color{Red} Lectio IX.}
\lettrine[lines=2]{\zallmancaps \color{Red} Q}{uid} ergo karissimi: quid in eius iudicio dicturi sumus, qui ab amore presentis seculi nec preceptis flectimur, nec verberibus emendamur? Sed fortasse aliquis tacitis sibi cogitationibus dicat. Ad vocem dominicam uterque iste piscator quid aut quantum dimisit, qui pene nichil habuit? Sed hac in re fratres karissimi affectum debemus potius pensare quam censum. Multum reliquit, qui sibi nichil retinuit. Multum reliquit, qui quamlibet parum: totum deseruit.
\begin{responsory-final}[videns-crucem-andreas]
{Videns crucem Andreas exclamavit dicens, * O crux inenarrabilis, o crux inestimabilis, o crux que per totum mundum fulges, * Non me dimittas errantem sicut ovem non habentem pastorem.}{Biduo vivens pendebat in cruce pro Christi nomine beatus Andreas et docebat populum et dicebat.}{Non me}
\end{responsory-final}%
\newline \V Dilexit Andream dominus. \R In odorem suavitatis.
\newline {\color{Red} In Laudibus. Ant.} Salve crux preciosa, suscipe discipulum eius qui pependit in te magister meus Christus. {\color{Red} Ant.} Beatus Andreas orabat dicens, domine rex eterne glorie suscipe me pendentem in patibulo. {\color{Red} Ant.} Andreas Christi famulus, dignus deo apostolus, germanus Petri et in passione socius. {\color{Red} Ant.} Maximilla Christo amabilis tulit corpus apostoli optimo loco cum aromatibus sepelivit. {\color{Red} Ant.} Qui persequebantur iustum demersisti eos domine in profundum, et in ligno crucis dux iusti fuisti.
\newline {\color{Red} Ad Ben. Ant.} Concede nobis hominem iustum, redde nobis hominem sanctum, ne interficias hominem deo carum, iustum, mansuetum et pium. {\color{Red} Oratio.}
\lettrine[lines=2]{\zallmancaps \color{Blue} M}{aiestatem} \hypertarget{maiestatem-andree}{\label{maiestatem-andree}} tuam domine suppliciter exoramus: ut sicut ecclesie tue beatus Andreas apostolis extitit predicator et rector: ita apud te sit pro nobis perpetuus intercessor. Per.
\newline {\color{Red} Ad Vesperas. Super psalmos. Ant.} Salve. {\color{Red} Ad Mag. Ant.} Cum pervenisset beatus Andreas ad locum ubi crux parata erat, exclamavit et dixit, o bona crux diu desiderata et iam concupiscenti animo preparata, securus et gaudens venio ad te, ita ut et tu exultans suscipias me discipulum eius qui pependit in te.
\newline \ul{Cotidie usque ad Octavam diem fiat memoria de sancto Andrea. In Laudibus per antiphonam.} Biduo. \ul{et \Vbar .} Annuntiaverunt. \ul{In Vesperis per antiphonam.} Beatus Andreas. \ul{et \Vbar .} Dilexit Andream. \ul{Oratio.} \hyperlink{maiestatem-andree}{Maiestatem.}
{\ubsubsection{Sancti Nicholai Episcopi}{sancti-nicholai-episcopi-breviarium}}
\newline {\color{Red} Ad Vesperas. Super psalmos. Ant.} Amicus dei Nicholaus pontificali decoratus infula, omnibus se amabilem exhibuit. \R \hyperlink{ex-eius-tumba}{Ex eius tumba.} {\color{Red} IX. Ad Mag. Ant.} O pastor eterne o clemens et bone custos, qui dum devoti gregis preces attenderes, voce lapsa de celo presuli sanctissimo, dignum episcopatu Nicholaum ostendisti tuum famulum. {\color{Red} Oratio.}
\lettrine[lines=2]{\zallmancaps \color{Red} D}{eus} qui beatum Nicholaum pontificem tuum innumeris decorasti miraculis, tribue quesumus nobis, ut eius meritis et precibus a Gehenne incendiis liberemur. Per.
\newline {\color{Red} Ad Matutinas. Invitat.}
\begin{invitatory}[adoremus-regem-nicholai-invitatorium]
{Adoremus regem seculorum, * In quo vivit Nicholaus honor sacerdotum.}
\end{invitatory}
\newline {\color{Red} In Primo Nocturno. Ant.} Nobilissimis siquidem natalibus ortus, velut lucifer Nicholaus emicuit. {\color{Red} Ps.} \psalm{1} {\color{Red} Ant.} Postquam domi puerilem decurrit etatem, cunctis mundi huius spretis oblectationibus: Christi se iugo subiciens, documentis sanctis suum prebuit auditum. {\color{Red} Ps.} \psalm{2} {\color{Red} Ant.} Pudore bono repletus dei famulus, sumptibus datis stupri nefas prohibuit. {\color{Red} Ps.} \psalm{3} \V Amavit eum. \R
\lettrine[lines=2]{\zallmancaps \color{Blue} C}{onfessor} \hypertarget{confessor-dei-nicholaus}{\label{confessor-dei-nicholaus}} dei Nicholaus nobilis progenie sed nobilior moribus, ab ipso puerili evo secutus dominum, * Meruit divina revelatione ad summum provehi sacerdotium. \V Erat enim valde compatiens et super afflictos pia gestans viscera. Meruit.
\begin{responsory}[operibus-sanctis]
{Operibus sanctis Nicholaus humiliter insistens, revelatione divina provectus est, * Ad summum sacerdotii gradum.}{Voce quippe de celo lapsa cuidam insinuatur
%pp120
presuli: dignum episcopatu Nicholaum.}
\end{responsory}%
\begin{responsory-doxology}[quadam-die]
{Quadam die tempestate sevissima quassati naute ceperunt sanctum vocare Nicholaum, * Et statim cessavit tempestas.}{Mox illis clamantibus apparuit quidam dicens eis, ecce assum, quid vocastis me.}
\end{responsory-doxology}%
\newline {\color{Red} In Secundo Nocturno. Ant.} Auro virginum incestus, auro patris earum inopiam, auro prorsus utroque detestabilem infamiam dei servus ademit Nicholaus. {\color{Red} Ps.} \psalm{4} {\color{Red} Ant.} Innocenter puerilia iura transcendens, evangelice institutionis discipulus effectus est. {\color{Red} Ps.} \psalm{5} {\color{Red} Ant.} Gloriam mundi sprevit cum suis oblectationibus, et ideo meruit provehi ad summum sacerdotii gradum. {\color{Red} Ps.} \psalm{8}
\begin{responsory}[audiens-christi]
{Audiens Christi confessor trium iuvenum innocentum necem, precucurrit quam totius ad locum quo fuerant plectendi, * Et liberavit eos.}{Statimque solutos a vinculis usque ad pretorium consulis secum adduxit.}
\end{responsory}%
\begin{responsory}[qui-cum-audisset]
{Qui cum audisset sancti Nicholai nomen statim expandunt manus utrasque ad celum, * Salvatoris laudantes clementiam.}{Clara quippe voce coram hominibus dignum referebant illum dei famulum.}
\end{responsory}%
\begin{responsory-doxology}[beatus-nicholaus]
{Beatus Nicholaus iam triumpho potitus, novit suis famulis prebere celestia comoda, qui toto corde poscunt eius largitiones, * Illi nimirum tota nos devotione oportet committere.}{Ut apud Christum eius patrociniis adiuvemur semper.}
\end{responsory-doxology}%
\newline {\color{Red} In Tertio Nocturno. Ant.} Pontifices almi divina revelatione letificati Nicholaum presulem devotissime consecraverunt. {\color{Red} Ps.} \psalm{14} {\color{Red} Ant.} Sanctus quidem triticum quod a nautis postulaverat acceptum et sagacitate distribuere, et augeri precibus impetravit. {\color{Red} Ps.} \psalm{20} {\color{Red} Ant.} Muneribus datis neci sunt iuvenes innocentes addicti, quibus domini servus fuit vite presidium festinanter. {\color{Red} Ps.} \psalm{23}
\begin{responsory}[summe-dei-confessor]
{Summe dei confessor Nicholae te venerantes protege, * Namque credimus tuis precibus nos posse salvari.}{Qui tres pueros morti addictos illesos abire fecisti, tuis laudibus instantem conserva plebem.}
\end{responsory}%
\begin{responsory}[servus-dei-nicholaus]
{Servus dei Nicholaus auri pondo trium virginum redemit pudorem, * Earumque patris impudicam remenso auro fugavit inopiam.}{Affluens itaque, misericordie visceribus metallo duplicato propulsavit earum infamiam.}
\end{responsory}%
\begin{responsory-final}[ex-eius-tumba]
{Ex eius tumba marmorea sacrum resudat oleum, quo liniti sanantur ceci, * Surdis auditus redditur, * Et debilis quisque sospes regreditur.}{Catervatim ruunt populi cernere cupientes que per eum fiunt mirabilia.}{Et debilis}%typo: Et eius
\end{responsory-final}%
\newline {\color{Red} In Laudibus. Ant.} Beatus Nicholaus adhuc puerulus multo ieiunio macerabat corpus. {\color{Red} Ant.} Ecclesie sancte frequentans limina, sacra pectori condebat mandata. {\color{Red} Ant.} Iuste et sancte vivendo ad honorem sacerdotii meruit promoveri divinitus. {\color{Red} Ant.} Amicus dei Nicholaus pontificali decoratus infula omnibus se amabilem exhibuit. {\color{Red} Ant.} O per omnia laudabilem virum cuius meritis ab omni clade liberantur qui ex toto corde querunt illum.
\newline {\color{Red} Ad Ben. Ant.} Copiose caritatis Nicholae pontifex qui cum deo gloriaris in celi palatio, condescende supplicamus ad te suspirantibus, ut exutos gravi carne pertrahas ad superos.
\newline {\color{Red} Ad Mag. Ant.} O Christi pietas omni prosequenda laude qui sui famuli Nicholai merita longe lateque declarat, nam ex tumba eius oleum manat, cunctosque languidos sanat.%typo: pro sequentia
\newline \ul{Memoria de sancto Andrea.} {\color{Red} Ant.} Ambulans. \V Dilexit Andream. {\color{Red} Oratio.}
\lettrine[lines=2]{\zallmancaps \color{Red} P}{rotegat} nos domine sepius beati Andree apostoli tui repetita sollemnitas, ut cuius patrocinia sine intermissione recolimus, perpetuam defensionem sentiamus. Per.
\newline \ul{In crastino. In Laudibus memoria de sancto Andrea.} {\color{Red} Ant.} Concede nobis. \V Annunciaverunt. {\color{Red} Oratio.} Protegat. \ul{In Vesperis nulla fiat de eo memoria.}
{\ubsubsection{Sancti Damasi Pape}{sancti-damasi-pape-breviarium}}
\newline {\color{Red} Oratio.}
\lettrine[lines=2]{\zallmancaps \color{Blue} D}{a} quesumus omnipotens deus, ut beati Damasi confessoris tui atque pontificis veneranda sollemnitas, et devotionem nobis augeat et salutem. Per.
{\ubsubsection{Sancte Lucie Virginis et Martiris}{sancte-lucie-virginis-et-martiris-breviarium}}
\newline {\color{Red} Ad Mag. Ant.} In tua patientia possedisti animam tuam Lucia sponsa Christi odisti que in mundo sunt, et choruscas cum angelis, sanguine proprio inimicum vicisti. {\color{Red} Oratio.}
\lettrine[lines=2]{\zallmancaps \color{Red} E}{xaudi} nos deus salutaris noster, ut sicut de beate Lucie virginis et martiris tue festivitate gaudemus, ita pie devotionis erudiamur effectu. Per.
\newline \ul{Ad Matutinas antiphone et responsoria, omnia preter nonum de communi.} {\color{Red} Nonum.}%typo: novum
\begin{responsory-doxology}[rogavi-dominum-meum]
{Rogavi dominum meum Ihesum Christum ut ignis iste non dominetur michi, et * Impetravi a domino inducias martyrii mei.}{Ut credentibus timorem auferam passionis, et non credentibus vocem insultationis.}
\end{responsory-doxology}%
\newline {\color{Red} In Laudibus. Ant.} Orante sancta Lucia apparuit ei beata Agatha, et consolabatur ancillam Christi. {\color{Red} Ant.} Lucia virgo quid a me petis quod ipsa poteris prestare continuo matri tue. {\color{Red} Ant.} Per te Lucia virgo civitas Siracusana decorabitur a domino Ihesu Christo. {\color{Red} Ant.} Benedico te pater domini mei Ihesu Christi, quia per filium tuum, ignis extinctus est a latere meo. {\color{Red} Ant.} Soror mea Lucia virgo deo devota, quid a me petis quod ipsa poteris prestare continuo matri tue.%typo: aparuit
\newline {\color{Red} Ad Ben. Ant.} Colunna es immobilis Lucia sponsa Christi quia omnis plebs te expectat, ut accipias coronam regni alleluya.
\newline {\color{Red} Ad Mag. Ant.} Tanto pondere eam fixit spiritus sanctus, ut virgo domini immobilis permaneret.
{\ubsubsection{Sancti Thome Apostoli}{sancti-thome-apostoli-breviarium}}
\newline {\color{Red} Ad Mag. Ant.} O Thoma Didime qui Christum meruisti tangere, te precibus rogamus altisonis, succurre nobis miseris ne damnemur cum impiis in adventu iudicis. {\color{Red} Oratio.}
\lettrine[lines=2]{\zallmancaps \color{Blue} D}{a} nobis quesumus domine beati apostoli tui Thome sollemnitatibus gloriari, ut eius semper et patrociniis sublevemur, et fidem congrua devotione sectemur. Per.
\newline {\color{Red} Lectio VII. Secundum Iohannem.}
\lettrine[lines=2]{\zallmancaps \color{Red} I}{n} illo tempore: Thomas unus ex duodecim qui dicitur Didimus, non erat cum eis quando venit Ihesus. Et reliqua.
\newline {\color{Red} Omelia Beati Gregorii Pape.} \grecross
\lettrine[lines=2]{\zallmancaps \color{Blue} I}{ste} unus discipulus defuit, reversus quod gestum est audivit: audita credere noluit. Venit iterum dominus, et non credenti discipulo latus palpandum prebuit, manus ostendit, et ostensa suorum cicatrice vulnerum, infidelitatis illius vulnus sanavit.
\newline {\color{Red} Lectio VIII.}
\lettrine[lines=2]{\zallmancaps \color{Red} Q}{uid} fratres karissimi quid inter hec animadvertitis? Non hoc casu sed divina dispensatione gestum est. Egit namque miro modo divina clementia, ut discipulus dubitans dum in magistro suo vulnera palparet carnis, in nobis vulnera sanaret infidelitatis.
\newline {\color{Red} Lectio IX.}
\lettrine[lines=2]{\zallmancaps \color{Blue} P}{lus} enim nobis Thome infidelitas ad fidem, quam fides credentium discipulorum profuit, quia dum ille ad fidem palpando reducitur: nostra mens omni dubitatione postposita in fide solidatur. Sic quippe discipulum post resurrectionem suam dubitare permisit, nec tamen in dubitatione deseruit: sicut ante nativitatem suam habere Mariam sponsum voluit: qui tamen ad eius nuptias non pervenit.
\newline {\color{Red} Ad Ben. Ant.} Quia vidisti me Thoma credidisti, beati qui non viderunt et crediderunt, alleluya.
\newline \ul{Memoria de Adventu per Ant.} Nolite timere. \ul{nisi festum transferatur.}
\newline {\color{Red} Ad Mag. Ant.} In regeneratione.
{\ubsubsection{Sancti Stephani Prothomartiris}{sancti-stephani-prothomartiris-breviarium}}%typo stephhani
\newline {\color{Red} Ad Memoriam in Vesperis. Ant.} Tu principatum tenes in choro martyrum similis angelo, qui pro te lapidantibus deum deprecatus es, beate Stephane intercede pro nobis ad dominum. {\color{Red} Oratio.}
\lettrine[lines=2]{\zallmancaps \color{Red} D}{a} nobis quesumus domine imitari quod colimus, ut discamus et inimicos diligere: quia eius natalicia celebramus, qui novit eciam pro persecutoribus exorare. Dominum nostrum.
\newline \ul{Quando predicta oratio dicitur ad memoriam et sequitur alia memoria immediate: finiatur sic.} Exorare Christum dominum nostrum.
\newline {\color{Red} Ad Matutinas. Invitat.}
\begin{invitatory}[christum-natum-invitatorium]
{Christum natum, qui beatum hodie coronavit Stephanum, * Venite adoremus.}
\end{invitatory}
\newline {\color{Red} In Primo Nocturno. Ant.} Beatus Stephanus iugi legis dei meditatione roboratus: tanquam lignum fructiferum secus salutarium aquarum plantatum decursus, fructum martyrii in tempore suo dedit primus. {\color{Red} Ps.} \psalm{1} {\color{Red} Ant.} Constitutus a deo predicator preceptorum eius, in timore sancto illi servire studiut, officioque fideliter peracto, in montem sanctum eius ascendere dignus fuit. {\color{Red} Ps.} \psalm{2} {\color{Red} Ant.} In tribulatione lapidum se prementium positus milia populi se circumdantis non timuit, quia susceptorem suum Ihesum ut eum salvum faceret, exurgere in celo vidit. {\color{Red} Ps.} \psalm{3} \R
\lettrine[lines=2]{\zallmancaps \color{Blue} S}{tephanus} \hypertarget{stephanus-autem-plenus}{\label{stephanus-autem-plenus}} autem plenus gratia et fortitudine, * Faciebat prodigia et signa magna in populo. \V Cum autem esset plenus spiritu sancto beatus Stephanus. Faciebat.
\begin{responsory}[videbant-omnes-stephanum]
{Videbant omnes Stephanum qui erant in concilio, * Et intuebantur vultum eius tanquam vultum angeli stantis inter illos.}{Stephanus autem plenus gratia et fortitudine faciebat prodigia et signa magna in populo.}
\end{responsory}%
\begin{responsory-final}[intuens-in-celum]
{Intuens in celum beatus Stephanus vidit gloriam dei et ait, * Ecce video celos apertos, et filium hominis stantem * A dextris virtutis dei.}{Cum autem esset Stephanus plenus spiritu sancto, intendens in celum ait.}{A dextris}
\end{responsory-final}%
\newline {\color{Red} In Secundo Nocturno. Ant.} Lumine vultus tui domine insignitus prothomartyr Stephanus sacrificium iusticie seipsum tibi sacrificavit, ideoque in leticia cordis, in pace obdormiens requiescit. {\color{Red} Ps.} \psalm{4} {\color{Red} Ant.} Benedictionis tue domine munere iustificatus, et scuto tue protectionis in passione munitus, nominis sui coronam Stephanus a te percipere meruit. {\color{Red} Ps.} \psalm{5} {\color{Red} Ant.} O quam admirabile est nomen tuum domine deus noster, pro quo beatus Stephanus passus, gloria et honore a te est coronatus, et super celos dono tue magnificentie exaltatus. {\color{Red} Ps.} \psalm{8} \V 
\begin{responsory}[lapidabant-stephanum]
{Lapidabant Stephanum invocantem et dicentem, domine Ihesu Christe accipe spiritum meum et * Ne statuas illis hoc peccatum.}{Positis autem genibus beatus Stephanus orabat dicens.}
\end{responsory}%
\begin{responsory}[impetum-fecerunt]
{Impetum fecerunt unanimiter in eum, et eiecerunt eum extra civitatem, * Invocantem et dicentem domine accipe spiritum meum.}{Et testes deposuerunt vestimenta sua secus pedes adolescentis, qui vocabatur Saulus, et lapidabant Stephanum.}
\end{responsory}%
\begin{responsory-doxology}[stephanus-servus-dei]
{Stephanus servus dei quem lapidabant Iudei, vidit celos apertos, vidit et introivit, * Beatus homo cui celi patebant.}{Cum igitur saxorum crepitantium turbine quateretur, inter ethereos aule celestis sinus divina ei claritas fulsit.}
\end{responsory-doxology}%
\newline {\color{Red} In Tertio Nocturno. Ant.} In domino deo suo confisus fortis athleta Stephanus lapidum fortiter sustinuit ictus, et idcirco ad montem virtutum transmigravit victoriosus. {\color{Red} Ps.} \psalm{10} {\color{Red} Ant.} Sine macula beatus Stephanus ingressus est domine in tabernaculum tuum, et quia operatus est iusticiam, requiescit in monte sancto tuo. {\color{Red} Ps.} \psalm{14} {\color{Red} Ant.} Domine virtus et leticia rectorum, quoniam tu prevenisti dilectum tibi Stephanum, dono gratuite benedictionis, te primum secutus est morte gloriose passionis, unde cum corona martyrii dedisti ei vitam in seculum seculi. {\color{Red} Ps.} \psalm{20}%typo: athelta
\newline {\color{Red} Lectio VII. Secundum Matheum.}
\lettrine[lines=2]{\zallmancaps \color{Red} I}{n} illo tempore: Dicebat Ihesus turbis Iudeorum et principibus sacerdotum. Ecce ego mitto ad vos prophetas et sapientes et scribas: et ex illis occidetis, et crucifigetis, et ex eis flagellabitis in synagogis vestris. Et reliqua.
\newline {\color{Red} Omelia Beati Hieronimi Presbiteri.} \grecross
\lettrine[lines=2]{\zallmancaps \color{Blue} E}{cce} ego mitto ad vos prophetas, et sapientes et scribas. Observa iuxta apostolum scribentem ad Corinthios, varia dona esse discipulorum Christi, alios prophetas qui ventura predicant, alios sapientes, qui noverint quando debeant proferre sermonem, alios scribas in lege doctissimos. Ex quibus lapidatus est Stephanus, Paulus occisus, crucifixus Petrus, flagellati in Actibus apostolorum discipuli.
\begin{responsory}[impii-super-iustum]
{Impii super iustum iacturam fecerunt ut eum morti traderent, * At ille gaudens suscepit lapides ut mereretur accipere coronam glorie.}{Continuerunt aures suas, et impetum fecerunt unanimes in eum.}
\end{responsory}%
\newline {\color{Red} Lectio VIII.}
\lettrine[lines=2]{\zallmancaps \color{Red} U}{t} veniat super vos omnis sanguis iustus, qui effusus est super terram a sanguine Abel iusti, usque ad sanguinem Zacharie filii Barachie, quem occidistis inter templum et altare. De Abel nulla est ambiguitas quin hic sit quem Cayn frater suus occiderit. Iustus autem non solum ex domini nunc sententia, sed eciam ex Genesis testimonio comprobatur, ubi accepta eius a Deo narrantur munera.
\begin{responsory}[patefacte-sunt-ianue]
{Patefacte sunt ianue celi Christi martyri beato Stephano, qui in numero martyrum inventus est primus, * Et ideo triumphat in celis coronatus.}{Mortem enim quam salvator dignatus est pro omnibus pati, hanc ille primus reddidit salvatori.}
\end{responsory}%
\newline {\color{Red} Lectio IX.}
\lettrine[lines=2]{\zallmancaps \color{Blue} Q}{uerimus} igitur, quis sit iste Zacharias filius Barachie, quia multos legimus Zacharias. In diversis diversa legi, et debeo singulorum opiniones ponere. Alii Zachariam filium Barachie dicunt, qui in duodecim prophetis undecimus est: patrisque in eo nomen consentit. Sed ubi occisus est inter templum et altare Scriptura non loquitur: maxime cum temporibus eius vix ruine templi fuerint.
\begin{responsory-final}[sancte-dei-preciose]
{Sancte dei preciose prothomartyr Stephane, qui virtute caritatis circumfultus undique dominum pro inimico exorasti populo, * Funde preces * Pro devoto tibi nunc collegio.}{Ut tuo propiciatus interventu dominus, nos purgatos a peccatis iungat celi civibus.}{Pro devoto}
\end{responsory-final}%
\newline {\color{Red} In Laudibus. Ant.} Lapidaverunt Stephanum, et ipse invocabat dominum dicens, ne statuas illis hoc peccatum. {\color{Red} Ant.} Lapides torrentes illi dulces fuerunt, ipsum sequuntur omnes anime iuste. {\color{Red} Ant.} Adhesit anima mea post te, quia caro mea lapidata est pro te deus meus. {\color{Red} Ant.} Stephanus vidit celos apertos, vidit et introivit, beatus homo cui celi patebant. {\color{Red} Ant.} Ecce video celos apertos et Ihesum stantem a dextris dei. {\color{Red} Capitulum.}
\lettrine[lines=2]{\zallmancaps \color{Red} S}{tephanus} plenus gratia et fortitudine, faciebat prodigia et signa magna in populo.
\newline {\color{Red} Ad Ben. Ant.} Intuens in celum beatus Stephanus vidit gloriam dei et ait, ecce video celos apertos et filium hominis stantem a dextris dei. {\color{Red} Oratio.} Da nobis. \ul{Memoria de Nativitate.}
\newline {\color{Red} Ad Terciam. Capitulum.} Stephanus plenus.
\newline {\color{Red} Ad Sextam. Cap.}
\lettrine[lines=2]{\zallmancaps \color{Blue} C}{um} esset Stephanus plenus spiritu sancto intendens in celum vidit gloriam dei, et Ihesum stantem a dextris dei.
\newline {\color{Red} Ad Nonam. Cap.}
\lettrine[lines=2]{\zallmancaps \color{Red} P}{ositis} genibus beatus Stephanus clamavit voce magna dicens, domine ne statuas illis hoc peccatum.
\newline {\color{Red} Ad Vesperas. Cap.} Stephanus. {\color{Red} Ad Mag. Ant.} O quam gloriosus est beatus Stephanus martyr et Levita, qui ante apostolos regna celestia possidere meruit, et ad patris dexteram filium videre.
\newline \ul{Memoria de sancto Iohanne.} {\color{Red} Ant.} Valde honorandus est beatus Iohannes qui supra pectus domini in cena recubuit. {\color{Red} Oratio.}
\lettrine[lines=2]{\zallmancaps \color{Blue} E}{cclesiam} tuam domine benignus illustra, ut beati Iohannis apostoli tui et evangeliste illuminata doctrinis ad dona perveniat sempiterna. Per. \ul{Memoria de Nativitate.}
{\ubsubsection{Sancti Iohannis Evangeliste}{sancti-iohannis-evangeliste-breviarium}}
\newline {\color{Red} Ad Matutinas. Invit.}
\begin{invitatory}[adoremus-regem-apostolorum-iohannem-invitatorium]
{Adoremus regem apostolorum, * Qui privilegio amoris Iohannem dilexit apostolum.}
\end{invitatory}
\newline {\color{Red} In Primo Nocturno. Ant.} Iohannes apostolus et evangelista virgo est electus a domino: atque inter ceteros magis dilectus. {\color{Red} Ps.} \psalm{18} {\color{Red} Ant.} Supra pectus domini Ihesu recumbens evangelii fluenta, de ipso sacro dominici pectoris fonte potavit. {\color{Red} Ps.} \psalm{33} {\color{Red} Ant.} Quasi unus de paradisi fluminibus evangelista Iohannes verbi dei gratiam in toto terrarum orbe diffudit. {\color{Red} Ps.} \psalm{44} \R
\lettrine[lines=2]{\zallmancaps \color{Red} V}{alde} \hypertarget{valde-honorandus-est}{\label{valde-honorandus-est}} honorandus est beatus Iohannes qui supra pectus domini in cena recubuit, * Cui Christus in cruce matrem virginem virgini commendavit. \V Virgo est electus a domino atque inter ceteros magis dilectus. Cui.
\begin{responsory}[hic-est-discipulus]
{Hic est discipulus qui testimonium perhibet de his et scripsit hec, * Et scimus quia verum est testimonium eius.}{Fluenta evangelii de ipso sacro dominici pectoris fonte potavit.}
\end{responsory}%
\begin{responsory-final}[hic-est-beatissimus]
{Hic est beatissimus evangelista et apostolus Iohannes, * Qui privilegio amoris precipui ceteris altius * A domino meruit honorari.}{Hic est discipulus ille quem diligebat Ihesus, qui supra pectus domini in cena recubuit.}{A domino}
\end{responsory-final}%
\newline {\color{Red} In Secundo Nocturno. Ant.} In ferventis olei dolium missus Iohannes apostolus divina se protegente gratia illesus exivit. {\color{Red} Ps.} \psalm{46} {\color{Red} Ant.} Propter insuperabilem evangelizandi constantiam exilio relegatus, divine visionis et allocutionis meruit crebra consolatione relevari. {\color{Red} Ps.} \psalm{60} {\color{Red} Ant.} Occurrit beato Iohanni ab exilio revertenti omnis populus virorum ac mulierum clamantium et dicentium benedictus qui venit in nomine domini. {\color{Red} Ps.} \psalm{63}
\begin{responsory}[qui-vicerit-faciam]
{Qui vicerit faciam illum columnam in templo meo dicit dominus, * Et scribam super eum nomen meum, et nomen civitatis nove Hierusalem.}{Vincenti dabo edere de ligno vite quod est in paradiso dei mei.}
\end{responsory}%
\begin{responsory}[diligebat-autem-eum]
{Diligebat autem eum Ihesus, quem specialis prerogativa castitatis ampliori dilectione fecerat dignum, * Quia virgo electus ab ipso virgo in evum permansit.}{In cruce denique moriturus, huic matrem virginem virgini commendavit.}
\end{responsory}%
\begin{responsory-doxology}[iste-est-iohannes]
{Iste est Iohannes qui supra pectus domini in cena recubuit, * Beatus apostolus cui revelata sunt secreta celestia.}{Iste est Iohannes cui Christus in cruce matrem virginem virgini commendavit.}
\end{responsory-doxology}%
\newline {\color{Red} In Tertio Nocturno. Ant.} Apparuit caro suo Iohanni dominus Ihesus Christus cum discipulis suis, et ait illi, veni dilecte meus ad me, quia tempus est ut epuleris in convivio meo cum fratribus tuis. {\color{Red} Ps.} \psalm{74} {\color{Red} Ant.} Expandens manus suas ad deum dixit, invitatus ad convivium tuum venio gratias agens, quia me dignatus es domine Ihesu Christe ad tuas epulas invitare, sciens quod ex toto corde meo desiderabam te. {\color{Red} Ps.} \psalm{96} {\color{Red} Ant.} Domine suscipe me ut cum fratribus meis sim, cum quibus veniens invitasti me, aperi michi ianuam vite et perduc me ad convivium epularum tuarum, tu es enim Christus filius dei vivi: qui precepto patris mundum salvasti, tibi gratias referimus per infinita seculorum secula. {\color{Red} Ps.} \psalm{98}
\newline {\color{Red} Lectio VII. Secundum Iohannem.}
\lettrine[lines=2]{\zallmancaps \color{Blue} I}{n} illo tempore: Dixit Ihesus Petro. Sequere me. Conversus Petrus vidit illum discipulum quem diligebat Ihesus sequentem, qui et recubuit in cena super pectus eius, et dixit, domine quis est qui tradet te. Et reliqua.
\newline {\color{Red} Omelia Venerabilis Bede Presbiteri.} \grecross
\lettrine[lines=2]{\zallmancaps \color{Red} C}{ommendat} nobis beatissimus evangelista et apostolus Iohannes privilegium amoris precipui, quo ceteris amplius meruit honorari a domino, commendat testimonium evangelice descriptionis, quod veritate divina subnixum, nullus fidelium dubitare permittitur.
\begin{responsory}[in-illa-die]
{In illa die suscipiam te servum meum et ponam te sicut signaculum in conspectu meo, * Quoniam ego elegi te dicit dominus.}{Esto fidelis usque ad mortem et dabo tibi coronam vite.}
\end{responsory}%
\newline {\color{Red} Lectio VIII.}
\lettrine[lines=2]{\zallmancaps \color{Blue} C}{ommendat} placidam sue carnis absolutionem, quam domino specialiter se visitante percepit. Notum autem novi vestre fraternitati, quis sit ille discipulus quem diligebat Ihesus, Iohannes videlicet, ipse cuius hodie natalicia festa celebramus.
\begin{responsory}[cibavit-illum-dominus]
{Cibavit illum dominus pane vite et intellectus et aqua sapientie salutaris potavit eum, * Et exaltavit eum apud proximos suos.}{In medio ecclesie aperuit os eius, et implevit eum spiritu sapientie et intellectus.}
\end{responsory}%
\newline {\color{Red} Lectio IX.}
\lettrine[lines=2]{\zallmancaps \color{Red} D}{iligebat} autem eum Ihesus, non exceptis ceteris singulariter solum, sed pre ceteris quos diligebat familiarius unum, quem specialis prerogativa castitatis ampliori dilectione fecerat dignum. Omnes quippe se diligere probat, quibus ante passionem loquitur, sicut dilexit me pater: et ego dilexi vos.
\begin{responsory-doxology}[in-medio-ecclesie]
{In medio ecclesie aperuit os eius, * Et implevit eum dominus spiritu sapientie et intellectus.}{Iocunditatem et exultationem thesaurizavit super eum.}
\end{responsory-doxology}%
\newline \V Valde honorandus est beatus Iohannes. \R Qui supra pectus domini in cena recubuit.
\newline {\color{Red} In Laudibus. Ant.} Hic est discipulus ille qui testimonium perhibet de his, et scimus quia verum est testimonium eius. {\color{Red} Ant.} Hic est discipulus meus sic eum volo manere donec veniam. {\color{Red} Ant.} Ecce puer meus electus quem elegi, posui super eum spiritum meum. {\color{Red} Ant.} Sunt de hic stantibus qui non gustabunt mortem donec videant filium hominis in regno suo. {\color{Red} Ant.} Si eum volo manere donec veniam, tu me sequere. {\color{Red} Cap.} Qui timet.
\newline {\color{Red} Ad Ben. Ant.} Iste est Iohannes qui supra pectus domini in cena recubuit, beatus apostolus cui revelata sunt secreta celestia. {\color{Red} Oratio.} Ecclesiam. \ul{Memoria de Nativitate. Deinde de sancto Stephano.} {\color{Red} Ant.} Ecce video. \V Corona aurea. {\color{Red} Oratio.} Da nobis.
\newline \ul{Ad Terciam et ad alias horas, Capitula de communi evangelistarum. Ad Vesperas. Capitulum.} Qui timet.
\newline {\color{Red} Ad Mag. Ant.} In medio ecclesie aperuit os eius, et implevit eum dominus spiritu sapientie et intellectus, stolaque glorie induit eum, alleluya alleluya alleluya.
\newline \ul{Memoria de Innocentibus.} {\color{Red} Ant.} Ambulabunt mecum in albis quoniam digni sunt, et non delebo nomina eorum de libro vite. \V Letamini in domino. {\color{Red} Oratio.}
\lettrine[lines=2]{\zallmancaps \color{Blue} D}{eus} cuius hodierna die preconium Innocentes martyres non loquendo sed moriendo confessi sunt, omnia in nobis vitiorum mala mortifica, ut fidem tuam quam lingua nostra loquitur eciam moribus vita fatetur. Per.
\newline \ul{Memoria de Nativitate, deinde de sancto Stephano.} {\color{Red} Ant.} Stephanus vidit.
{\ubsubsection{Sanctorum Innocentum}{sanctorum-innocentum-breviarium}}
\newline {\color{Red} Ad Matutinas. Invit.}
\begin{invitatory}[venite-adoremus-innocentum-invitatorium]
{Venite adoremus regem regum, * Qui beatis martiribus coronam dedit glorie.}
\end{invitatory}
{\color{Red} Ymnus.} \hyperlink{que-vox-que-poterit}{Que vox que poterit. Te summa.}
\newline {\color{Red} In Primo Nocturno. Ant.} Novit dominus vias Innocentium qui non steterunt in viis peccatorum. {\color{Red} Ps.} \psalm{1} {\color{Red} Ant.} Rex terre infremuit adversus Christum, quia rex in Syon constitutus super Innocentum milia regnat. {\color{Red} Ps.} \psalm{2} {\color{Red} Ant.} Deus iudex iustus iudica nos secundum innocentiam nostram. {\color{Red} Ps.} \psalm{7} \V Letamini.
\newline \ul{Sermo Beati Bernardi.} {\color{Red} Lectio Prima.}
\lettrine[lines=2]{\zallmancaps \color{Red} B}{enedictus} qui venit in nomine domini, Deus dominus et illuxit nobis. Benedictum nomen glorie eius quod est sanctum. Neque enim ociose venit quod ex Maria natum est Sanctum: sed copiose diffundit et nomen et gratiam sanctitatis. Nimirum inde Stephanus sanctus: inde Iohannes sanctus: inde sancti eciam Innocentes. \R
\lettrine[lines=2]{\zallmancaps \color{Blue} S}{ub} \hypertarget{sub-altare-dei}{\label{sub-altare-dei}} altare dei audivi voces occisorum dicentium, quare non defendis sanguinem nostrum, et acceperunt divinum responsum, * Adhuc sustinete modicum tempus donec impleatur numerus fratrum vestrorum. \V Et date sunt illis singule stole albe, et dictum est illis. Adhuc.
\newline {\color{Red} Lectio Secunda.}
\lettrine[lines=2]{\zallmancaps \color{Red} U}{tili} proinde dispositione triplex ista sollemnitas, Natale domimi comitatur, ut non modo inter contiguas sollemnitates, devotio continua perseveret: sed et fructus dominice Nativitatis, exinde nobis evidentius innotescat. Siquidem advertere est in his tribus sollemnitatibus, triplicem quandam speciem sanctitatis: nec facile preter hec tria sanctorum genera, quartum aliud in hominibus posse arbitror reperiri.
\begin{responsory}[effuderunt-sanguinem]
{Effuderunt sanguinem sanctorum velut aquam in circuitu Hierusalem, * Et non erat qui sepeliret.}{Vindica domine sanguinem sanctorum tuorum qui effusus est.}
\end{responsory}%
%pp121
\newline {\color{Red} Lectio Tertia.}
\lettrine[lines=2]{\zallmancaps \color{Red} H}{abemus} in beato Stephano martyrii simul et operis voluntatem, habemus solam voluntatem in beato Iohanne: solum in beatis Innocentibus opus. Biberunt omnes hi calicem salutaris, aut corpore simul et spiritu, aut solo spiritu, aut corpore solo. Calicem quidem meum bibetis: ait dominus Iacobo et Iohanni. Bibit ergo Iohannes calicem salutaris. De Innocentium vero coronis quis dubitat?
\begin{responsory-doxology}[adoraverunt-viventem]
{Adoraverunt viventem in secula seculorum, * Mittentes coronas suas ante thronum domini dei sui.}{Et ceciderunt in conspectu thorni in facies suas et benedixerunt viventem in secula.}
\end{responsory-doxology}%
\newline {\color{Red} In Secundo Nocturno. Ant.} Ex ore infantium deus et lactentium laude perfecta destruis inimicum. {\color{Red} Ps.} \psalm{8} {\color{Red} Ant.} Iudicabit dominus pupillum et viduam, et pluet super peccatores laqueos ignis. {\color{Red} Ps.} \psalm{10} {\color{Red} Ant.} Mortis usuras rex impius exegit super Innocentes, sed illi requiescunt in monte sancto dei, malignus ad nichilum est deductus. {\color{Red} Ps.} \psalm{14} \V Exultent iusti.
\newline {\color{Red} Lectio Quarta.}
\lettrine[lines=2]{\zallmancaps \color{Blue} I}{lle} pro Christo trucidatos infantes dubitet inter martyres coronari: qui regeneratos in Christo non credit inter adoptionis filios numerari. Alioquin quando coevos sibi pueros puer ille qui nobis natus est propter se pateretur occidi, qui utique solo nutu poterat cohibere nisi melius eis aliquid provideret?
\begin{responsory}[isti-sunt-sancti-qui-passi]
{Isti sunt sancti qui passi sunt propter te domine, vindica eos, * Quia clamant ad te cotidie.}{Vindica domine sanguinem sanctorum tuorum qui effusus est.}
\end{responsory}%
\newline {\color{Red} Lectio Quinta.}
\lettrine[lines=2]{\zallmancaps \color{Red} Q}{uemadmodum} igitur ceteris infantibus, tunc quidem circumcisio, nunc vero baptismus sine ullo proprie voluntatis usu sufficit ad salutem, sic nichilominus susceptum pro eo martyrium, illis sufficit ad sanctitatem. Si queris eorum apud Deum meritum, ut coronarentur: quere et apud Herodem crimina, ut trucidarentur. An forte minor Christi pietas quam Herodis impietas, ut ille quidem potuerit innoxios neci dare: Christus non potuerit propter se occisos coronare?
\begin{responsory}[cantabant-sancti-canticum]
{Cantabant sancti canticum novum ante sedem dei et agni, * Et resonabat terra in voces eorum.}{Sub throno dei sancti clamant vindica sanguinem nostrum deus noster.}
\end{responsory}%
\newline {\color{Red} Lectio Sexta.}
\lettrine[lines=2]{\zallmancaps \color{Blue} S}{it} ergo Stephanus martyr apud homines, cuius voluntaria passio evidenter apparuit in eo: vel maxime quod in ipso mortis articulo pro persequentibus quam pro se ipso sollicitudinem gereret ampliorem, ut illorum scelera magis quam sua vulnera plangeret. Sit Iohannes apud angelos martyr, quibus tanquam spiritualibus creaturis spiritualia devotionis eius signa certius innotuerunt. Ceterum hi sunt plane martyres tui Deus: ut in quibus nec homo nec angelus meritum invenit: singularis tue prerogativa gratie evidentius commendetur.
\begin{responsory-final}[isti-sunt-sancti-qui-non]
{Isti sunt sancti qui non inquinaverunt vestimenta sua, * Ambulabunt mecum in albis, * Quia digni sunt.}{Hi sunt qui cum mulieribus non sunt coinquinati virgines enim sunt.}{Quia}
\end{responsory-final}%
\newline {\color{Red} In Tertio Nocturno. Ant.} Quis ascendet aut quis stabit in loco sancto dei innocentes manibus et mundo corde. {\color{Red} Ps.} \psalm{23} {\color{Red} Ant.} Innocentes adheserunt michi, quorum pedes avulsi sunt a laqueo mortis. {\color{Red} Ps.} \psalm{24} {\color{Red} Ant.} Filio regis datum est iudicium ut salvet innocentes filios pauperum, et humiliabit Herodem calumniatorem. {\color{Red} Ps.} \psalm{71} \V Iusti autem.%typo: avilsi
\newline {\color{Red} Lectio VII. Secundum Matheum.}
\lettrine[lines=2]{\zallmancaps \color{Red} I}{n} illo tempore: Angelus domini apparuit in somnis Ioseph dicens. Surge et accipe puerum et matrem eius: et fuge in Egyptum: et esto ibi usque dum dicam tibi. Et reliqua.
\newline {\color{Red} Omelia Venerabilis Bede Presbiteri.}
\lettrine[lines=2]{\zallmancaps \color{Blue} D}{e} morte preciosa martyrum Christi innocentum, sacra nobis est fratres karissimi evangelii lectio recitata, in qua tamen omnium Christi martyrum: preciosa est mors designata. Quod enim parvuli occisi sunt, significat per humilitatis meritum ad martyrii perveniendum gloriam, et quia nisi conversus fuerit quis, et effectus ut parvulus: pro Christo animam non possit dare.
\begin{responsory}[vidi-sub-altare]
{Vidi sub altare dei animas interfectorum propter verbum dei quod habebant et clara voce dicebant, vindica domine sanguinem sanctorum tuorum, * Qui effusus est.}{Sub throno dei sancti clamant, vindica domine sanguinem nostrum.}
\end{responsory}%
\newline {\color{Red} Lectio VIII.}
\lettrine[lines=2]{\zallmancaps \color{Red} Q}{uod} in Bethleem et in omnibus finibus eius occisi sunt: ostendit non solum in Iudea unde ecclesie cepit origo, sed et in omnibus eiusdem ecclesie finibus, quacunque per orbem diffusa est, persecutionem sevituram perfidorum, et piorum patientiam esse coronandam.
\begin{responsory}[sub-throno-dei]
{Sub throno dei sancti clamant, * Vindica sanguinem nostrum deus noster.}{Sub altare dei audivi voces occisorum clamantium et dicentium.}
\end{responsory}%
\newline {\color{Red} Lectio IX.}
\lettrine[lines=2]{\zallmancaps \color{Blue} Q}{ui} bimi occisi sunt doctrina et operatione perfectos indicant. Qui vero infra: simplices vel ydiotas fidei tamen non ficte constantia fortes: eque denuntiant. Quod illi quidem occisi sunt, sed Christus qui querebatur illesus evasit, insinuat corpora quidem ab impiis posse perimi, sed Christum pro quo persecutio tota sevit nullatenus eis vel viventibus vel occisis posse auferri.
\begin{responsory-final}[centum-quadraginta]
{Centum quadraginta quatuor milia qui empti sunt de terra, hi sunt qui cum mulieribus non sunt coinquinati virgines enim permanserunt, * Ideo regnant cum deo, et agnus dei, * In illis.}{Hi empti sunt ex omnibus primitie deo et agno et in ore ipsorum non est inventum mendacium.}{In illis}
\end{responsory-final}%
\newline \ul{Non dicatur,} Te deum, \ul{nisi Dominica fuerit sed reincipiatur nonum, \Rbar .}%typo: novum
\newline \V Exultabunt sancti in gloria.
\newline {\color{Red} In Laudibus. Ant.} Herodes iratus occidit multos pueros in Bethleem Iude civitate David. {\color{Red} Ant.} A bimatu et infra, occidit multos pueros Herodes propter dominum. {\color{Red} Ant.} Vox in Rama audita est, ploratus et ululatus, Rachel plorans filios suos. {\color{Red} Ant.} Sub throno dei sancti clamant vindica sanguinem nostrum deus noster. {\color{Red} Ant.} Cantabant sancti canticum novum ante sedem dei et agni et resonabat terra in voces eorum. {\color{Red} Capitulum.}
\lettrine[lines=2]{\zallmancaps \color{Red} V}{idi} supra montem Syon agnum stantem, et cum eo centum quadraginta quatuor milia habentes nomen eius et nomen patris eius scriptum in frontibus suis. {\color{Red} Ymnus.}
\hymn{caterva-matrum-personat}{Caterva matrum personat}{
\lettrine[lines=2]{\zallmancaps \color{Blue} C}{aterva} matrum personat,
collisa deflens pignora,
quorum tyrannus milia
Christo sacravit victima.
\par {\bfseries \color{Red} G}loria tibi domine
qui natus es de virgine,
cum patre et sancto spiritu
in sempiterna secula. Amen.
}
\newline \V Mirabilis deus.
\newline {\color{Red} Ad Ben. Ant.} Hi sunt qui cum mulieribus non sunt coinquinati, virgines enim sunt et sequuntur agnum quocumque ierit. {\color{Red} Oratio.} Deus cuius.
\newline \ul{Memoria de Nativitate et de sancto Stephano, ut supra. Memoria de sancto Iohanne.} {\color{Red} Ant.} Hic est discipulus ille. \V Annunciaverunt. {\color{Red} Oratio.} Ecclesiam.
\newline {\color{Red} Ad Terciam. Cap.} Vidi supra. \ul{Responsoria et Versiculi plurimorum martirum ad omnes horas.}
\newline {\color{Red} Ad Sextam. Cap.}
\lettrine[lines=2]{\zallmancaps \color{Red} H}{i} sunt qui cum mulieribus non sunt coinquinati, virgines enim sunt, et sequuntur agnum quocumque ierit.
\newline {\color{Red} Ad Nonam. Capit.}
\lettrine[lines=2]{\zallmancaps \color{Blue} H}{i} empti sunt ex omnibus primitie deo et agno et in ore eorum non est inventum mendacium.
\newline {\color{Red} Ad Vesperas. Cap.} Vidi supra. {\color{Red} Ymnus.} \hyperlink{que-vox-que-poterit}{Que vox. Te summa.} \V Letamini.
\newline {\color{Red} Ad Mag. Ant.} Innocentes pro Christo infantes occisi sunt, ab iniquo rege lactentes interfecti sunt, ipsum sequuntur agnum sine macula et dicunt semper gloria tibi domine.
\newline \ul{Memoria de sancto Thoma.} {\color{Red} Ant.} Hic est vere martyr. {\color{Red} Oratio.}
\lettrine[lines=2]{\zallmancaps \color{Red} D}{eus} pro cuius ecclesia gloriosus pontifex Thomas gladiis impiorum occubuit, presta quesumus, ut omnes qui eius implorant auxilium, petitionis sue salutarem consequantur effectum. Per.
\newline \ul{Memoria de Natali et de sancto Stephano ut supra. De sancto Iohanne.} {\color{Red} Ant.} Valde honorandus. \V In omnem. \ul{Oratio ut supra.}
{\ubsubsection{Sancti Thome Episcopi et Martiris}{sancti-thome-episcopi-et-martiris-breviarium}}
\newline {\color{Red} Ad Mat.} \ul{Sicut unius martiris, post orationem de festo fiat memoria de Natali et de sancto Stephano et de sancto Iohanne ut supra. De Innocentibus.} {\color{Red} Ant.} Herodes iratus. \V Mirabilis. {\color{Red} Oratio.} Deus cuius hodierna.
\newline \ul{Ad Vesperas. Memoria de Nativitate. De sancto Stephano et de sancto Iohanne ut supra. De Innocentibus.} {\color{Red} Ant.} A bimatu. \V Letamini. {\color{Red} Oratio.} Deus cuius.
\newline \ul{Per Octavas Nativitas et sanctorum fiat de ipsis memoria eo modo et ordine qui in tribus diebus precedentibus est notatus.}
{\ubsubsection{Sancti Silvestri Pape}{sancti-silvestri-pape-breviarium}}
\newline {\color{Red} Oratio.}
\lettrine[lines=2]{\zallmancaps \color{Blue} D}{a} quesumus omnipotens deus, ut beati Silvestri confessoris tui atque pontificis, veneranda sollemnitas, et devotionem nobis augeat et salutem. Per.
\newline \ul{Memoria de Nativitate postea de sanctis, ultimo de Dominica si hac die evenerit. Ad Matutinas sicut de uno confessore. Post orationem de sancto Silvestro memoria de Natali postea de sanctis ultimo de Dominica si hac die evenerit.}
\newline \ul{In Octava sancti Stephani ad Matutinas, Invit.} \hyperlink{christum-natum-invitatorium}{Christum natum.} \ul{Ymnus de uno martire, prima antiphona, primi nocturni sola dicatur super psalmos de festo, versiculus de tercio nocturno unius martiris. Tria ultima responsoria de festo.} Te deum. \ul{In Laudibus una sola antiphona, scilicet,} Lapidaverunt. \ul{Cap., Ymnus, \Vbar , antiphona ad Ben., et Oratio sicut in festo. In hac die et deinceps usque ad Octavas Epyphanie inclusive cotidie nisi in vigilia et in die Epyphanie post alias memorias fiat memoria de beata Virgine in Laudibus et in Vesperis. In Laudibus per antiphonam.} Rubum. \ul{\Vbar .} Benedictus. \ul{Oratio.} Deus qui salutis. \ul{In Vesperis per antiphonam.} Germinavit. \ul{\Vbar .} Tanquam sponsus. \ul{Oratio ut supra. Ad horas diei, et ad Vesperas sicut in festo, eciam si in crastino fuerit Dominica. Quod si Octava sancti Stephani in Dominica evenerit omnia fiant sicut in festo eiusdem servato modo festi simplicis.}
\newline \ul{In Octava sancti Iohannis et Innocentum omnia fiant de festis, servato ordine qui in Octava sancti Stephani est notatus.}
\newline \ul{In Octava sanctorum Innocentum si Dominica fuerit repetantur sex lectiones de festo supra posite, et tres sequentes sint ultime. Si vero non fuerit Dominica: legantur tres sequentes.}
\newline \ul{Idem in eodem.} {\color{Red} Lectio.}
\lettrine[lines=2]{\zallmancaps \color{Red} E}{x} ore infantium et lactentium perfecisti laudem. Gloria in excelsis Deo angeli dicunt, et in terra pax hominibus bone voluntatis. Magna quidem sed audeo dicere necdum perfecta laus donec veniat qui dicat. Sinite parvulos venire ad me: quia talium est regnum celorum, et pax hominibus eciam sine voluntatis usu, in sacramento pietatis.
\newline {\color{Red} Lectio.}
\lettrine[lines=2]{\zallmancaps \color{Blue} C}{onsiderent} hec qui de voluntate et opere contentiosis solent disputationibus corrixari, et advertant neutrum negligi oportere. Ubi tamen facultas deest, non modo salutem conferre possit sed eciam sanctitatem. Sed et hoc teneant, prodesse quidem opus sine voluntate, non autem contra voluntatem, ut unde salvantur infantes: inde magis damnentur ficte accedentes. Nichilominus sane in quibusdam voluntas sine opere sufficiens est: non autem contra opus.
\newline {\color{Red} Lectio.}
\lettrine[lines=2]{\zallmancaps \color{Red} E}{mulemur} proinde caritatem et sectemur bona opera fratres mei: nec infirmitatis nec ignorantie peccata ullo modo parvipendamus. Magis autem solliciti et timorati agamus gratias benignissimo et largissimo salvatori, qui humane salutis occasiones tam copiosa caritate perquirit, ut in his voluntatem et opus, in his sine opere voluntatem, in his eciam sine voluntate opus invenire letetur, qui vult omnes homines salvos fieri.
{\ubsubsection{Sancti Pauli Heremite}{sancti-pauli-heremite-breviarium}}
\newline {\color{Red} Oratio.}
\lettrine[lines=2]{\zallmancaps \color{Red} D}{eus} qui nos beati Pauli confessoris tui annua sollemnitate letificas, concede propicius ut cuius natalicia colimus per eius ad te exempla gradiamur. Per.
{\ubsubsection{Sanctorum Hylarii et Remigii Episcoporum}{sanctorum-hylarii-et-remigii-episcoporum-breviarium}}
\newline {\color{Red} Oratio.}
\lettrine[lines=2]{\zallmancaps \color{Blue} D}{eus} qui nos sanctorum tuorum Hylarii et Remigii confessionibus gloriosis circundas et protegis, da nobis et eorum imitatione proficere et intercessione gaudere. Per.
{\ubsubsection{Sancti Felicis Presbiteri et Confessoris}{sancti-felicis-presbiteri-et-confessoris-breviarium}}
\newline {\color{Red} Oratio.}
\lettrine[lines=2]{\zallmancaps \color{Red} C}{oncede} quesumus omnipotens deus ut ad meliorem vitam sanctorum tuorum exempla nos provocent, quatinus quorum sollemnia agimus eciam actus imitemur. Per.
{\ubsubsection{Sancti Mauri Abbatis}{sancti-mauri-abbatis-breviarium}}
\newline {\color{Red} Oratio.}
\lettrine[lines=2]{\zallmancaps \color{Blue} I}{ntercessio} nos quesumus domine beati Mauri abbatis commendet, ut quod nostris meritis non valemus eius patrocinio assequamur. Per.
{\ubsubsection{Sancti Marcelli Pape et Martiris}{sancti-marcelli-pape-et-martiris-breviarium}}
\newline {\color{Red} Oratio.}
\lettrine[lines=2]{\zallmancaps \color{Red} P}{reces} populi tui quesumus domine clementer exaudi, ut beati Marcelli martyris tui atque pontificis meritis adiuvemur, cuius passione letamur. Per.
{\ubsubsection{Sancti Antoni Abbatis}{sancti-antoni-abbatis-breviarium}}
\newline {\color{Red} Oratio.}
{\zallmancaps \bfseries \color{Blue} I}{ntercessio}. \ul{Ut supra.}
{\ubsubsection{Sancte Prisce Virginis et Martyris}{sancte-prisce-virginis-et-martyris-breviarium}}
\newline {\color{Red} Oratio.}
\lettrine[lines=2]{\zallmancaps \color{Red} D}{a} quesumus omnipotens deus, ut qui beate Prisce martyris tue natalicia colimus, et annua sollemnitate letemur, et tante fidei proficiamus exemplo. Per.
{\ubsubsection{Sanctorum Fabiani et Sebastiani Martirum}{sanctorum-fabiani-et-sebastiani-martirum-breviarium}}
\newline {\color{Red} Oratio.}
\lettrine[lines=2]{\zallmancaps \color{Blue} D}{eus} qui beatos martyres tuos Fabianum atque Sebastianum virtute constantie in passione roborasti, ex eorum nobis imitatione tribue pro amore suo prospera mundi despicere, et nulla eius adversa formidare. Per.
{\ubsubsection{Sancte Agnetis Virginis et Martyris}{sancte-agnetis-virginis-et-martyris-breviarium}}
\newline {\color{Red} Ad Mag. Ant.} Beata Agnes in medio flammarum expansis manibus orabat, te deprecor venerande, colende, pater metuende, quia per sanctum filium tuum minas evasi sacrilegi tyranni, et carnis spurcicias immaculato calle transivi et ecce venio ad te, quem amavi, quem quesivi quem semper optavi. {\color{Red} Oratio.}
\lettrine[lines=2]{\zallmancaps \color{Red} O}{mnipotens} sempiterne deus qui infirma mundi eligis ut fortia queque confundas, concede propitius, ut qui beate Agnetis virginis et martyris tue sollemnia colimus, eius apud te patrocinia sentiamus. Per. \ul{Memoria de martyribus nisi festa transferantur.}
\newline {\color{Red} Ad Matutinas. In Primo Nocturno. Ant.} Discede a me pabulum mortis quia iam ab alio amatore preventa sum. {\color{Red} Ps.} \psalm{8} {\color{Red} Ant.} Anulo suo subarravit me dominus meus Ihesus Christus et tanquam sponsam decoravit me corona. {\color{Red} Ps.} \psalm{18} {\color{Red} Ant.} Dextram meam et collum meum cinxit lapidibus preciosis, tradidit auribus meis inestimabiles margaritas. {\color{Red} Ps.} \psalm{23} \R
\lettrine[lines=2]{\zallmancaps \color{Blue} D}{iem} \hypertarget{diem-festum-sacratissime} festum sacratissime virginis celebremus, qualiter passa sit beata Agnes ad memoriam revocemus, terciodecimo etatis sue anno mortem perdidit et vitam invenit, * Quia solum vite dilexit actorem. \V Infantia quidem computabatur in annis sed erat senectus mentis immensa. Quia.
\begin{responsory}[dextram-meam]
{Dextram meam et collum meum cinxit lapidibus preciosis, tradidit auribus meis inestimabiles margaritas, * Et circumdedit me vernantibus atque coruscantibus gemmis.}{Induit me dominus ciclade auro texta, et immensis monilibus ornavit me.}
\end{responsory}%
\begin{responsory-doxology}[induit-me-dominus]
{Induit me dominus vestimento salutis et indumento leticie circumdedit me, * Et tanquam sponsam decoravit me corona.}{Tradidit auribus meis inestimabiles margaritas, et circumdedit me vernantibus atque coruscantibus gemmis.}
\end{responsory-doxology}%
\newline {\color{Red} In Secundo Nocturno. Ant.} Christus circumdedit me vernantibus atque coruscantibus gemmis preciosis. {\color{Red} Ps.} \psalm{44} {\color{Red} Ant.} Induit me dominus ciclade auro texta, et immensis monilibus ornavit me. {\color{Red} Ps.} \psalm{45} {\color{Red} Ant.} Posuit signum in faciem meam ut nullum preter eum amatorem admittam. {\color{Red} Ps.} \psalm{86}
\begin{responsory}[amo-christum-in]
{Amo Christum in cuius thalamum introivi, cuius mater virgo est, cuius pater feminam nescit, cuius michi organa modulatis vocibus cantant, * Quem cum amavero casta sum, cum tetigero munda sum, cum accepero virgo sum.}{Anulo suo subarravit me, et immensis monilibus ornavit me.}
\end{responsory}%
\begin{responsory}[mel-et-lac]
{Mel et lac ex eius ore suscepi, * Et sanguis eius ornavit genas meas.}{Cuius pulchritudinem sol et luna mirantur: ipsi soli servo fidem.}
\end{responsory}%
\begin{responsory-doxology}[ipsi-sum-desponsata]
{Ipsi sum desponsata cui angeli serviunt, cuius pulchritudinem sol et luna mirantur, Ipsi soli servo fidem, * Ipsi me tota devotione comitto.}{Dextram meam et collum meum cinxit lapidibus preciosis, tradidit auribus meis inestimabiles margaritas.}
\end{responsory-doxology}%
\newline {\color{Red} In Tertio Nocturno. Ant.} Mel et lac ex eius ore suscepi, et sanguis eius ornavit genas meas. {\color{Red} Ps.} \psalm{95} {\color{Red} Ant.} Cuius pulchritudinem sol et luna mirantur, ipsi soli servo fidem. {\color{Red} Ps.} \psalm{96} {\color{Red} Ant.} Ipsi soli servo fidem, ipsi me tota devotione comitto. {\color{Red} Ps.} \psalm{97}
\begin{responsory}[iam-corpus-eius]
{Iam corpus eius corpori meo sociatum est, et sanguis eius ornavit genas meas, * Cuius mater virgo est, cuius pater feminam nescit.}{Ipsi sum desponsata cui angeli serviunt, cuius pulchritudinem sol, et luna mirantur.}
\end{responsory}%
\begin{responsory}[pulchra-facie]
{Pulchra facie, sed pulchrior fide beata es Agnes, respuens mundum letaberis cum angelis, * Intercede pro omnibus nobis.}{Specie tua et pulchritudine tua intende prospere et regna.}%typo: pulchior
\end{responsory}%
\begin{responsory-doxology}[omnipotens-adorande]
{Omnipotens adorande, colende, tremende, benedico te, * Quia per filium tuum unigenitum: evasi minas hominum impiorum, et spurcitias diaboli impolluto calle transivi.}{Te confiteor labiis te corde, te totis visceribus concupisco.}
\end{responsory-doxology}%
\newline {\color{Red} In Laudibus. Ant.} Ingressa Agnes turpitudinis locum, angelum domini preparatum invenit. {\color{Red} Ant.} Mecum enim habeo custodem corporis mei angelum domini. {\color{Red} Ant.} Ipsi sum desponsata cui angeli serviunt, cuius pulchritudinem sol et luna mirantur. {\color{Red} Ant.} Benedico te pater domini mei Ihesu Christi, quia per filium tuum ignis extinctus est a latere meo. {\color{Red} Ant.} Congaudete mecum et congratulamini, quia cum his omnibus lucidas sedes accepi.
\newline {\color{Red} Ad Ben. Ant.} Stans beata Agnes in medio flammarum expansis manibus orabat ad dominum, omnipotens adorande, colende, tremende, benedico te, et glorifico nomen tuum in eternum.
{\ubsubsection{Sancti Vincentii Martyris}{sancti-vincentii-martyris-breviarium}}
\newline {\color{Red} Ad Vesperas. Super psalmos. Ant.} Iste est qui pro lege dei sui morti se tradidit non dubitavit mori, ab iniquis interfectus est et in eternum vivit cum Christo, agnum secutus est, et accepit palmam.
\newline \ul{Psalmi unius virginis, Vel de Dominica, si in secunda feria propter Septuagesima festum evenerit vel transferatur.}
\newline \R \hyperlink{christi-miles-preciosus}{Christi miles.} {\color{Red} VI. Ad Mag. Ant.} Sacram presentis diei sollemnitatem humili celebremus devotione, qua invictus Christi martyr Vincentius tyranno devicto: insignem victorie palmam celo gaudens intulit. {\color{Red} Oratio.}
\lettrine[lines=2]{\zallmancaps \color{Red} A}{desto} domine supplicationibus nostris, ut qui ex iniquitate nostra reos nos esse cognoscimus, beati Vincentii martyris tui intercessione liberemur. Per.
\newline \ul{Memoria de sancta Agnete, nisi festa transferantur.} {\color{Red} Ant.} Ecce quod cupivi iam video, quod speravi iam teneo, illi sum iuncta in celis quem in terris posita tota devotione dilexi.
\newline {\color{Red} Ad Matutinas. Invit.}
\begin{invitatory}[vincentem-mundum-invitatorium]
{Vincentem mundum adoremus dominum, * In quo pressuram mundi martyr vicit Vincentius.}
\end{invitatory}
\newline {\color{Red} In Primo Nocturno. Ant.} Sanctus Vincentius a pueritia studiis litterarum traditus, superna sibi providente clementia: gemina scientia efficacissime claruit. {\color{Red} Ps.} \psalm{1} {\color{Red} Ant.} Sanctitate quoque insignis diaconii arce suscepta: vices pontificis diligenter exequebatur. {\color{Red} Ps.} \psalm{2} {\color{Red} Ant.} Valerius igitur episcopus, et levita Vincentius, spe fruendi victoria divinitus subnixi, in confessione deitatis alacriter cucurrerunt. {\color{Red} Ps.} \psalm{3} \R
\lettrine[lines=2]{\zallmancaps \color{Blue} S}{acram} \hypertarget{sacram-presentis-diei}{\label{sacram-presentis-diei}} presentis diei sollemnitatem humili celebremus devotione, qua invictus Christi martyr Vincentius tyranno devicto, * Insignem victorie palmam, celo gaudens intulit. \V Peracto passionis sue venerando triumpho, angelorum cetibus comitatur. Insignem.
\begin{responsory}[sanctus-vincentius]
{Sanctus Vincentius Christi martyr providente superna clementia, que sibi eum previdebat vas electionis futurum, * Gemina scientia efficacissime claruit.}{Sanctitate quoque insignis diaconii arce suscepta: vices pontificis diligenter exequebatur.}
\end{responsory}%
\begin{responsory-final}[levita-vincentius]
{Levita Vincentius dixit beato Valerio, si iubes pater sancte responsis iudicem aggrediar, iam tibi fili karissime divini verbi curam commiseram, * Nunc quoque, * Pro fide qua astamus responsa committo.}{Tibi enim gemina scientia pollenti, ac superni amoris igne ferventi: celestis olim doctrine ministerium delegavi.}{Pro fide}
\end{responsory-final}%
\newline {\color{Red} In Secundo Nocturno. Ant.} Tanto namque feliciores se esse credebant: quanto acriora tyranni supplicia pia longanimitate certassent evincere. {\color{Red} Ps.} \psalm{4} {\color{Red} Ant.} Levita Vincentius dixit beato Valerio, si iubes pater sancte responsis iudicem aggrediar. {\color{Red} Ps.} \psalm{5} {\color{Red} Ant.} Iam tibi fili karissime divini verbi curam commiseram, nunc quoque pro fide qua astamus responsa committo. {\color{Red} Ps.} \psalm{8}
\begin{responsory}[ecce-iam-in]
{Ecce iam in sublime agor, et omnes principes tuos seculo altior tyranne, despicio, nolo gloriam meam minuas, nec damna inferas laudi, * Paratus sum enim ad omnia tormenta pro salvatoris nomine sustinenda.}{Insurge ergo, et toto malignitatis spiritu debachare, videbis me dei virtute plus posse dum torqueor: quam possis ipse qui torques.}
\end{responsory}%
\begin{responsory}[assumptus-ex-eculeo]
{Assumptus ex eculeo levita Vincentius, atque ad patibulum raptus, tortores suos pene preveniens, * Moras carnificum arguendo ad penam alacriter properabat.}{Intrepidus itaque dei athleta candentis ferri machinam ultro conscendit.}
\end{responsory}%
\begin{responsory-final}[christi-miles-preciosus]
{Christi miles preciosus levita Vincentius, ut tribunal sponte rogum conscendit intrepidus, * Cuius salis crepitantis per corpus minutie sparsim ibant, * Atque prune vernabantur sanguine.}{Inter hec manet immotus ille dei servulus, orans Christum in sublime erectis luminibus.}{Atque}
\end{responsory-final}%
\newline {\color{Red} In Tertio Nocturno. Ant.} Beatus Vincentius cuius iam mens conscia corone dixit Datiano, hactenus a te habitus sermo de neganda fide peroravit. {\color{Red} Ps.} \psalm{10} {\color{Red} Ant.} Nefarium tamen apud Christianorum prudentiam esse cognosce: deitatis cultum abnegando blasphemare. {\color{Red} Ps.} \psalm{14} {\color{Red} Ant.} Profitemur enim Christiane religionis nos esse cultores, ac unius veri dei permanentis in secula famulos et testes. {\color{Red} Ps.} \psalm{20}
\begin{responsory}[beatus-dei-athleta]
{Beatus dei athleta Vincentius, horrendo clausus ergastulo carceris, psalmum deo et ymnum, angelorum frequentia relevatus, letus dicebat, * Quorum vallatus caterva venerando fovebatur obsequio et mulcebatur alloquio.}{Dantur ergo laudes deo altissimo, et resonante organo vocis angelice modulata suavitas procul diffunditur.}
\end{responsory}%
\begin{responsory}[agnosce-o-vincenti]
{Agnosce o Vincenti invictissime, pro cuius nomine fideliter decertasti, * Ipse tibi coronam preparatam servat in celis qui te fecit victorem in penis.}{Esto igitur iam securus de premio, quia mox deposito onere carnis nostro eris addendus collegio.}
\end{responsory}%
\begin{responsory-final}[vir-inclitus]
{Vir inclitus Vincentius martyr domini preciosus, succensus igne divini amoris: constanter sustinuit supplicia passionis, * Et per inmanitatem tormentorum: * Pervenit ad societatem angelorum.}{Cuius intercessio nobis obtineat veniam, qui per tormenta passionis eternam meruit palmam et sempiternam coronam.}{Pervenit}
\end{responsory-final}%
\newline {\color{Red} In Laudibus. Ant.} Assumptus ex eculeo levita Vincentius, atque ad patibulum raptus, moras carnificum arguendo ad penam alacriter properabat. {\color{Red} Ant.} Intrepidus itaque candentis ferri machinam ultro conscendit ac manens immotus, erectis in celum luminibus dominum precabatur. {\color{Red} Ant.} Hinc horrendo carceris clausus ergastulo dei athleta, angelorum venerando fovebatur obsequio: et mulcebatur alloquio. {\color{Red} Ant.} Agnosce o Vincenti invictissime, quia pro cuius nomine fideliter decertasti: ipse tibi coronam preparatam servat in celestibus. %typo: athelta
%pp122
{\color{Red} Ant.} Dantur ergo laudes deo altissimo, et resonante organo vocis angelice modulata suavitas procul diffunditur.
\newline {\color{Red} Ad Ben. Ant.} Egregius dei martyr Vincentius diris tormentorum suppliciis pro Christo alacriter superatis, ac felicis pugne agone constanter expleto, tandem preciosam resolutus in mortem celo triumphatis spiritum reddidit.
\newline {\color{Red} Ad Mag. Ant.} Beatus Vincentius applicatus tormentis alacri vultu dei presentia roboratus, dixit, hoc est quod semper optavi et votis omnibus exquisivi.
{\ubsubsection{Sancte Emerentiane Virginis et Martyris}{sancte-emerentiane-virginis-et-martyris-breviarium}}
\newline {\color{Red} Oratio.}
\lettrine[lines=2]{\zallmancaps \color{Red} I}{ndulgentiam} nobis domine beata Emerentiana virgo et martyr imploret, que tibi grata semper extitit et merito castitatis, et tue professione virtutis. Per.
{\ubsubsection{In Conversione Sancti Pauli}{in-conversione-sancti-pauli-breviarium}}
\newline {\color{Red} Ad Vesperas. Super psalmos. Ant.} Sancte Paule apostole predicator veritatis et doctor gentium intercede pro nobis ad deum qui te elegit. {\color{Red} Capitulum.}
\lettrine[lines=2]{\zallmancaps \color{Blue} S}{aulus} adhuc spirans minarum et cedis in discipulos domini, accessit ad principem sacerdotum, et petiit ab eo epistolas in Damascum ad synagogas, ut siquos inveniret huius vie viros ac mulieres, vinctos perduceret Hierusalem.
\newline \R \hyperlink{magnus-sanctus-paulus}{Magnus sanctus Paulus.} {\color{Red} VI. Ymnus.}
\hymn{doctor-egregie-paule}{Doctor egregie Paule}{
\lettrine[lines=2]{\zallmancaps \color{Red} D}{octor} egregie Paule, mores instrue,
et mente polum nos transferre satage,
donec perfectum largiatur plenius %typo perfetum
evacuato quod ex parte gerimus.
\par {\bfseries \color{Blue} S}it trinitati sempiterna gloria,
honor potestas atque iubilatio,
in unitate cui manet imperium
ex tunc et modo per eterna secula. Amen.
}
\newline {\color{Red} Ad Mag. Ant.} Elegit dominus virum de plebe, et claritatem visionis eterne dedit illi, celebremus conversionem sancti Pauli apostoli, gaudium sit in celo et in terra pax hominibus bone voluntatis. {\color{Red} Oratio.}
\lettrine[lines=2]{\zallmancaps \color{Blue} D}{eus} qui universum mundum beati Pauli apostoli predicatione docuisti, da nobis quesumus, ut qui eius hodie conversionem colimus, per eius ad te exempla gradiamur. Per.
\newline {\color{Red} Ad Matutinas. Invitat.}
\begin{invitatory}[laudemus-deum-conversione-invitatorium]
{Laudemus deum nostrum, * In conversione apostoli Pauli.}
\end{invitatory}
\newline {\color{Red} Ymnus.} \hyperlink{doctor-egregie-paule}{Doctor egregie.}
\newline {\color{Red} In Primo Nocturno. Ant.} Saulus adhuc spirans minarum et cedis in discipulos domini: abiit ad principem sacerdotum, et peciit ab eo ut ubicumque inveniret huius vie viros: vinctos perduceret Hierusalem. {\color{Red} Ps.} \psalm{18} \V Paternarum traditionum amplius emulator existens.
\newline \ul{Post psalmum \Vbar . ante antiphona dicatur et hoc in consimilibus observetur.}
{\color{Red} Ant.} Ibat igitur Saulus furiis invectus, dirum toto pectore virus efflabat, et sanctorum sanguinem sine intermissione sitiebat. {\color{Red} Ps.} \psalm{33} \V Per totam Iudeam insania ferebatur, ut Christi membra laniaret in terris. {\color{Red} Ant.} Et subito circumfulsit eum lux de celo, et cecidit in terram nichilque videbat. {\color{Red} Ps.} \psalm{44} \V Audivit autem vocem dicentem sibi, Saule Saule quid me persequeris. \R
\lettrine[lines=2]{\zallmancaps \color{Red} S}{aulus} \hypertarget{saulus-adhuc}{\label{saulus-adhuc}} adhuc spirans minarum et cedis in discipulos domini, abiit ad principem sacerdotum et petiit illum, * Ut ubicumque inveniret huius vie viros: vinctos perduceret Hierusalem. \V Ibat igitur Saulus furiis invectus, dirumque toto pectore virus efflabat. Ut ubicumque.
\begin{responsory}[ibat-igitur-saulus]
{Ibat igitur Saulus furiis invectus, dirumque toto pectore virus efflabat, * Et sanctorum sanguinem sine intermissione sitiebat.}{Per totam Iudeam insania ferebatur ut Christi membra laniaret in terris.}%typo: virtus
\end{responsory}%
\begin{responsory-doxology}[a-christo-de]
{A Christo de celo vocatus et in terram prostratus: ex persecutore effectus est vas electionis, et plus omnibus laborans multo latius inter omnes verbi gratiam seminavit, * Atque doctrinam evangelicam sua predicatione complevit.}{Inter apostolos vocatione novissimus predicatione primus, nomen Christi multarum manifestavit gentium populis.}
\end{responsory-doxology}%
\newline {\color{Red} In Secundo Nocturno. Ant.} Saule Saule quid me persequeris, quis es domine, ego sum Ihesus quem tu persequeris, durum est tibi contra stimulum calcitrare. {\color{Red} Ps.} \psalm{46} \V Circumfulsit eum lux de celo et audivit vocem dicentem sibi. {\color{Red} Ant.} Saulus autem tremens ac stupens dixit ad Ihesum. Domine quid me vis facere, ait autem dominus ad eum, surge et ingredere civitatem et dicetur tibi quid te oporteat facere. {\color{Red} Ps.} \psalm{60} \V Viri autem qui comitabantur cum eo stabant stupefacti. {\color{Red} Ant.} A Christo de celo vocatus et in terram prostratus ex persecutore factus est vas electionis. {\color{Red} Ps.} \psalm{63} \V Prostratus est sevissimus persecutor sed erectus est fidelissimus predicator.
\begin{responsory}[vade-anania]
{Vade Anania et quere Saulum, quia vas electionis est michi, ut portet nomen meum, * Coram gentibus et regibus et filiis Israhel.}{Ego enim ostendam illi quanta oporteat eum, pro nomine meo pati.}
\end{responsory}%
\begin{responsory}[ingressus-paulus]
{Ingressus Paulus in synagogas predicabat Ihesum Iudeis, affirmans * Quia hic est Christus.}{Stupebant autem omnes qui eum audiebant dicentem.}
\end{responsory}%
\begin{responsory-doxology}[magnus-sanctus-paulus]
{Magnus sanctus Paulus vas electionis vere digne est glorificandus, * Qui et meruit thronum duodecimum possidere.}{A Christo de celo vocatus et in terram prostratus: ex persecutore effectus est vas electionis.}
\end{responsory-doxology}%
\newline {\color{Red} In Tertio Nocturno. Ant.} Viri autem qui comitabantur cum eo stabant stupefacti, audientes quidem vocem, sed neminem videntes. {\color{Red} Ps.} \psalm{74} \V Saulus autem cadens in terram: apertis oculis nichil videbat. {\color{Red} Ant.} Ad manus autem illum trahentes introduxerunt Damascum, et erat tribus diebus non videns, et non manducavit neque bibit. {\color{Red} Ps.} \psalm{96} \V Surrexit autem Saulus de terra, apertisque oculis nichil videbat. {\color{Red} Ant.} Vade Anania et quere Saulum ecce enim orat, quia vas electionis est michi, ut portet nomen meum coram gentibus et regibus et filiis Israhel. {\color{Red} Ps.} \psalm{98} \V Dixit autem dominus in visu Ananie.
\begin{responsory}[tu-es-vas]
{Tu es vas electionis sancte Paule apostole predicator veritatis in universo mundo, * Per quem omnes gentes cognoverunt gratiam dei.}{Intercede pro nobis ad deum qui te elegit.}
\end{responsory}%
\begin{responsory}[sancte-paule-apostole]
{Sancte Paule apostole predicator veritatis et doctor gentium intercede pro nobis ad deum, * Qui te elegit.}{Ut digni efficiamur gratia dei.}
\end{responsory}%
\begin{responsory-doxology}[celebremus-conversionem]
{Celebremus conversionem sancti Pauli apostoli, * Quia hodie ex persecutore effectus est vas electionis.}{Gaudent angeli et exultant archangeli, et collaudant in celis filium dei.}
\end{responsory-doxology}%
\V Dedisti. 
\newline {\color{Red} In Laudibus. Ant.} Saulus qui et Paulus magnus predicator, a deo confortatus, convalescebat et confundebat Iudeos. \V Ostendens quia hic est Christus filius dei vivi. {\color{Red} Ant.} Saule frater dominus misit me Ihesus qui apparuit tibi in via qua veniebas, ut videas et implearis spiritu sancto. \V Et abiit Ananias et intravit in domum et imponens ei manum dixit. {\color{Red} Ant.} Sub manu continuo Ananie ceciderunt squame ab oculis eius, et surgens baptizatus est, et accipiens cibum confortatus est. \V Fuit autem cum discipulis qui erant Damasci per dies aliquot. {\color{Red} Ant.} Prostratus est sevissimus persecutor, sed erectus est fidelissimus predicator, quos domine verbis edocet de te meritis ducat ad te. \V Inter apostolos vocatione novissimus predicatione primus. {\color{Red} Ant.} Ingressus Paulus in synagogas predicabat Iudeis Ihesum, affirmans quia hic est Christus. \V Stupebant autem omnes qui eum audiebant. {\color{Red} Cap.} Saulus adhuc. {\color{Red} Ymnus.} \hyperlink{eterna-christi-munera-apostolorum}{Eterna Christi.}%typo Stupebat
\newline {\color{Red} Ad Ben. Ant.} Celebremus conversionem sancti Pauli apostoli quia hodie, ex persecutore effectus est vas electionis, ideo gaudent angeli, letantur archangeli, et collaudant in celis filium dei.
\newline \ul{Ad horas diei antiphone Laudum sine versibus.}
\newline {\color{Red} Ad Terciam. Cap.} Saulus.
\newline {\color{Red} Ad Sextam. Capitulum.}
\lettrine[lines=2]{\zallmancaps \color{Blue} V}{ade} Anania quoniam vas electionis michi est iste, ut portet nomen meum coram gentibus et regibus et filiis Israhel.
\newline {\color{Red} Ad Nonam. Cap.}
\lettrine[lines=2]{\zallmancaps \color{Red} E}{t} abiit Ananias et introivit in domum et imponens ei manum dixit, Saule frater, dominus misit me Ihesus qui apparuit tibi in via qua veniebas, ut videas et implearis spiritu sancto.
\newline {\color{Red} Ad Vesperas. Super psalmos. Ant.} Saulus qui et Paulus. \ul{sine versu. Cap., Ymnus, \Vbar , ut in Primis Vesperis.}
\newline {\color{Red} Ad Mag. Ant.} Cum autem complacuit ei qui me segregavit ex utero matris mee, et vocavit per gratiam suam, ut revelaret in me filium suum in gentibus, continuo non adquievi carni et sanguini, sed abii ad antecessores meos apostolos.
{\ubsubsection{Sancti Iuliani Episcopi et Confessoris}{sancti-iuliani-episcopi-et-confessoris-breviarium}}
\newline {\color{Red} Oratio.}
\lettrine[lines=2]{\zallmancaps \color{Blue} E}{xaudi} domine preces nostras quas in sancti Iuliani confessoris tui atque pontificis sollemnitate deferimus, ut qui tibi digne meruit famulari, eius intercedentibus meritis ab omnibus nos absolvas peccatis. Per.
{\ubsubsection{Sancte Agnetis Secundo}{sancte-agnetis-secundo-breviarium}}
\newline {\color{Red} Oratio.}
\lettrine[lines=2]{\zallmancaps \color{Red} D}{eus} qui nos annua beate Agnetis virginis et martyris tue sollemnitate letificas, da ut quam veneramur officio eciam pie conversationis sequamur exemplo. Per.
{\ubsubsection{Sancti Ignatii Episcopi et Martyris}{sancti-ignatii-episcopi-et-martyris-breviarium}}
\newline {\color{Red} Oratio.}
\lettrine[lines=2]{\zallmancaps \color{Blue} D}{a} quesumus omnipotens deus, ut qui beati Ignatii martyris tui atque pontificis sollemnia colimus, eius apud te intercessionibus adiuvemur. Per.
{\ubsubsection{In Festo Purificationis Beate Marie}{in-festo-purificationis-beate-marie-breviarium}}
\newline {\color{Red} Ad Vesperas. Super psalmos. Ant.} O admirabile commercium creator generis humani, animatum corpus sumens de virgine nasci dignatus est, et procedens homo sine semine largitus est nobis suam deitatem. {\color{Red} Ps.} \psalm{112} \ul{cum ceteris.} {\color{Red} Cap.}
\lettrine[lines=2]{\zallmancaps \color{Red} E}{cce} ego mittam angelum meum, qui preparabit viam ante faciem meam, et statim veniet ad templum suum dominator quem vos queritis, et angelus testamenti quem vos vultis. \R \hyperlink{gaude-maria-virgo}{Gaude Maria.} {\color{Red} IX.} \ul{cum prosa.} \hyperlink{inviolata-intacta-prosa}{Inviolata.} {\color{Red} Ymnus.}
\hymn{ave-maris-stella}{Ave maris stella}{
\lettrine[lines=2]{\zallmancaps \color{Blue} A}{ve} maris stella
dei Mater alma,
atque semper virgo
felix celi porta.
\par {\bfseries \color{Red} S}umens illud ave
gabrielis ore,
funda nos in pace
mutans nomen Eve.
\par {\bfseries \color{Blue} S}olve vincla reis,
profer lumen cecis,
mala nostra pelle,
bona cuncta posce.
\par {\bfseries \color{Red} M}onstra te esse matrem,
sumat per te precem
qui pro nobis natus
tulit esse tuus.
\par {\bfseries \color{Blue} V}irgo singularis
inter omnes mitis,
nos culpis solutos
mites fac et castos.
\par {\bfseries \color{Red} V}itam presta puram,
iter para tutum,
ut videntes Ihesum
semper colletemur.
\par {\bfseries \color{Blue} S}it laus deo patri
summo Christo decus
spiritui sancto,
tribus honor unus. Amen.
}
\newline \V Responsum accepit Symeon a spiritu sancto. \R Non visurum se mortem nisi videret Christum domini.
\newline {\color{Red} Ad Mag. Ant.} Cum inducerent puerum Ihesum parentes eius accepit eum Symeon in ulnas suas et benedixit deum dicens, nunc dimittis domine servum tuum in pace. {\color{Red} Oratio.}
\lettrine[lines=2]{\zallmancaps \color{Red} O}{mnipotens} sempiterne deus maiestatem tua supplices exoramus, ut sicut unigenitus filius tuus hodierna die cum nostre carnis substantia in templo est presentatus, ita nos facias purificatis tibi mentibus presentari. Per eundem.
\newline {\color{Red} Ad Completorium. Super psalmos. Ant.} Sancta dei genitrix virgo semper Maria intercede pro nobis ad dominum deum nostrum.
\newline {\color{Red} Ad Nunc Dimittis. Ant.} Nunc dimittis domine servum tuum in pace, quia viderunt oculi mei salutare tuum.
\newline {\color{Red} Ad Matutinas. Invit.}
\begin{invitatory}[ecce-venit-ad-invitatorium]
{Ecce venit ad templum sanctum suum dominator dominus, * Gaude et letare Syon occurens deo tuo.}%typo occcurens
\end{invitatory}
\newline {\color{Red} Ymnus.}
\hymn{quem-terra}{Quem terra pontus}{
\lettrine[lines=2]{\zallmancaps \color{Blue} Q}{uem} terra, pontus, ethera,
colunt, adorant, predicant,
trinam regentem machinam,
claustrum Marie baiulat.
\par {\bfseries \color{Red} C}ui luna, sol et omnia
deserviunt per tempora,
perfusa celi gratia,
gestant puelle viscera.
\par {\bfseries \color{Blue} B}eata mater munere
cuius supernus artifex
mundum pugillo continens
ventris sub archa clausus est.
\par {\bfseries \color{Red} B}eata celi nuntio,
fecunda sancto spiritu,
desideratus gentibus
cuius per alvum fusus est. {\color{Red} Tempore Resurrectionis.} Quesumus actor.
\par {\bfseries \color{Blue} G}loria tibi domine
qui natus es de virgine
cum patre et sancto spiritu,
in sempiterna secula. Amen.
}
\newline \ul{Predictus versus,} Gloria tibi domine qui natus, \ul{etc. dicatur ad omnes ymni eiusdem metri, quandocumque agitur de Beata Virgine.}%typo ym
\newline {\color{Red} In Primo Nocturno. Ant.} Benedicta tu in mulieribus et benedictus fructus ventris tui. {\color{Red} Ps.} \psalm{8} {\color{Red} Ant.} Sicut mirra electa odorem dedisti suavitatis sancta dei genitrix. {\color{Red} Ps.} \psalm{18} {\color{Red} Ant.} Post partum virgo inviolata permansisti, dei genitrix intercede pro nobis. {\color{Red} Ps.} \psalm{23} \V Sancta dei genitrix virgo semper Maria. \R Intercede pro nobis ad dominum deum nostrum.
\newline \ul{Sermo Beati Augustini} {\color{Red} \grecross \ Lectio Prima.}
\lettrine[lines=2]{\zallmancaps \color{Red} S}{icut} olim predictum est, mater Syon dicet homo, et homo factus est in ea. O omnipotentia nascentis, o magnificentia de celo ad terram descendentis. Adhuc in utero portabatur, et a Iohanne Baptista salutabatur, in templo presentabatur, et a Symeone sene, famoso, annoso, probato, cognoscebatur. Tunc cognovit, tunc adoravit: tunc dixit. Et nunc domine dimitte servum tuum in pace: quia viderunt oculi mei salutare tuum. Differebatur exire de seculo, ut videret natum per quem conditum est seculum. \R
\lettrine[lines=2]{\zallmancaps \color{Blue} A}{dorna} \hypertarget{adorna-thalamum}{\label{adorna-thalamum}} thalamum tuum Syon et suscipe regem Christum, * Quem virgo concepit, virgo peperit virgo post partum quem genuit adoravit. \V Accipiens Symeon puerum in manibus gratias agens benedixit dominum. Quem.
\newline {\color{Red} Lectio Secunda.}
\lettrine[lines=2]{\zallmancaps \color{Red} C}{ognovit} infantem senex: factus est in puero puer. Innovatus est in etate, qui plenus erat pietate. Rogo vos fratres: iste senex quam magnitudinem sentiebat in parvo? Quod videbat: mater portabat. Quod intelligebat: mundum regebat. Ipse ferebat, et Christus regebat. Symeon senex ferebat Christum infantem: Christus regebat Symeonem in senectute. Quid ergo dixit Symeon? Audite. Nunc domine dimitte servum tuum in pace. In pace inquam, quam commendaturus es discipulis tuis.
\begin{responsory}[postquam-impleti]
{Postquam impleti sunt dies purgationis eius secundum legem Moysi, tulerunt illum in Hierusalem, ut sisterent eum domino, * Sicut scriptum est in lege domini quia omne masculinum adaperiens vulvam sanctum domino vocabitur.}{Obtulerunt pro eo domino par turturum aut duos pullos columbarum.}
\end{responsory}%
\newline {\color{Red} Lectio Tertia.}
\lettrine[lines=2]{\zallmancaps \color{Blue} D}{ictum} illi fuerat a domino: quod non gustaret mortem, nisi videret Christum domini natum. Natus est Christus. Ille veniebat, et iste ibat. Sed nisi veniret ille: iste ire nolebat. Iam senectus naturam excludebat: sed sincera pietas detinebat. At ubi natus est Christus, at ubi eum manibus matris portatum vidit et divinam infantiam pia senectus agnovit: accepit eum in manibus suis et dixit. Nunc dimittis servum tuum domine in pace.
\begin{responsory-final}[obtulerunt-pro-eo]
{Obtulerunt pro eo domino par turturum aut duos pullos columbarum, * Sicut scriptum est, * In lege domini.}{Postquam impleti sunt dies purgationis eius secundum legem Moysi tulerunt illum in Hierusalem.}{In lege}
\end{responsory-final}%
\newline {\color{Red} In Secundo Nocturno. Ant.} Emissiones tue paradisus malorum punicorum cum pomorum fructibus. {\color{Red} Ps.} \psalm{44} {\color{Red} Ant.} Fons ortorum, puteus aquarum viventium, que fluunt impetu de Libano. {\color{Red} Ps.} \psalm{45} {\color{Red} Ant.} Veniat dilectus meus in ortum suum, ut comedat fructum pomorum suorum. {\color{Red} Ps.} \psalm{86} \V Post partum virgo inviolata permansisti. \R Dei genitrix intercede pro nobis.
\newline {\color{Red} Lectio Quarta.}
\lettrine[lines=2]{\zallmancaps \color{Red} I}{mpletum} est desiderium senis, in mundi ipsius senectute. Ipse ad senem venit: qui mundum inveteratum invenit. In isto quidem seculo diu esse nolebat, et Christum in hoc seculo videre cupiebat: cantans cum propheta et dicens. Ostende nobis domine misericordiam tuam: et salutare tuum da nobis. Denique ut noveritis ita esse: istius leticia conclusit ita dicens. Nunc dimittis servum tuum in pace: Quia viderunt oculi mei salutare tuum.
\begin{responsory}[symeon-iustus]
{Symeon iustus et timoratus expectabat redemptionem Israhel, * Et spiritus sanctus erat in eo.}{Responsum accepit Symeon a spiritu sancto non visurum se mortem nisi videret Christum domini.}
\end{responsory}%
\newline {\color{Red} Lectio Quinta.}
\lettrine[lines=2]{\zallmancaps \color{Blue} H}{ec} sunt testimonia tua domine Ihesu, antequam tibi fluctus maris cederent imperanti, antequam ventus te iubente siluisset, mortuus te vocante resurrexisset, sol te moriente palluisset, terra te resurgente tremuisset, celum te ascendente patuisset. Adhuc enim in manibus matris portabaris: et iam dominus orbis agnoscebaris. Ubique Christus dominus pulcher occurrebat, pulcher in celis, pulcher in terris: pulcher in patre verbum, pulcher in matre caro et verbum.
\begin{responsory}[responsum-acceperat]
{Responsum acceperat Symeon a spiritu sancto non visurum se mortem nisi videret Christum domini, et benedixit deum et dixit, * Nunc dimittis servum tuum in pace, quia viderunt oculi mei salutare tuum domine.}{Cum inducerent puerum Ihesum parentes eius, accepit eum Symeon in ulnas suas et dixit.}
\end{responsory}%
\newline {\color{Red} Lectio Sexta.}
\lettrine[lines=2]{\zallmancaps \color{Red} C}{um} hic fuisset dominus, et manibus portaretur, locuti sunt celi, gavisi sunt angeli. Magos stella perduxit, adoratus est in presepio, adoratus est in templo. Puer Christus a sene Symeone gestatus est in ulnas: a quo ipse regebatur. Pulcher ergo Christus in miraculis, in flagellis, in sermonibus, in verberibus. Pulcher non curans mortem, et mortuos suscitans, pulcher in ligno, pulcher in celo, pulcher in sepulcro. Quis igitur fratres non adoret primum et novissimum, initium et finem. Alpha et \omega ?
\begin{responsory-doxology}[cum-inducerent]
{Cum inducerent puerum Ihesum parentes eius, ut facerent secundum consuetudinem legis pro eo, accepit eum Symeon in ulnas suas, et benedixit deum et dixit, * Nunc dimittis domine servum tuum in pace.}{Suscipiens Symeon puerum in manibus exclamavit dicens.}
\end{responsory-doxology}%
\newline {\color{Red} In Tertio Nocturno. Ant.} Veni in ortum meum soror mea sponsa, messui mirram meam cum aromatibus meis. {\color{Red} Ps.} \psalm{95} {\color{Red} Ant.} Comedi favum cum melle meo, bibi vinum meum cum lacte meo. {\color{Red} Ps.} \psalm{96} {\color{Red} Ant.} Talis est dilectus meus, et ipse est amicus meus filie Hierusalem. {\color{Red} Ps.} \psalm{97} \V Speciosa facta es et suavis. \R In delitiis tuis sancta dei genitrix.
\newline {\color{Red} Lectio VII. Secundum Lucam.}
\lettrine[lines=2]{\zallmancaps \color{Blue} I}{n} illo tempore: Postquam impleti sunt dies purgationis Marie secundum legem Moysi, tulerunt Ihesum in Hierusalem, ut sisterent eum domino sicut scriptum est in lege domini, quia omne masculinum adaperiens vulvam, sanctum domino vocabitur. Et reliqua.
\newline {\color{Red} Omelia Beati Ambrosii Episcopi.}
\lettrine[lines=2]{\zallmancaps \color{Red} N}{on} solum ab angelis et prophetis, a pastoribus et parentibus: sed eciam a senioribus et iustis, generatio domini accipit testimonium. Omnis etas et uterque sexus, eventuumque miracula, fidem astruunt. Virgo generat, sterilis parit, mutus loquitur, Elyzabeth prophetat: Magus adorat.
\begin{responsory}[suscipiens-ihesum]
{Suscipiens Ihesum in ulnas suas Symeon exclamavit et dixit, * Tu es vere lumen ad illuminationem gentium, et gloriam plebis tue Israhel.}{Symeon in manibus infantem accepit, sed maiestatem intus agnovit et dixit.}
\end{responsory}%
\newline {\color{Red} Lectio VIII.}
\lettrine[lines=2]{\zallmancaps \color{Blue} U}{tero} clausus Iohannes esultat vidua confitetur: iustus expectat. Et bene iustus: qui non suam sed populi gratiam requirebat, cupiens ipse corporee vinculis fragilitatis absolvi, sed expectans videre promissum. Sciebat enim: quia beati oculi qui eum viderent.
\begin{responsory}[senex-puerum]
{Senex puerum portabat, puer autem senem regebat, * Quem virgo concepit, virgo peperit virgo post partum quem genuit adoravit.}{Responsum accepit Symeon a spiritu sancto, non visurum se mortem nisi videret Christum domini.}
\end{responsory}%
\newline {\color{Red} Lectio IX.}
\lettrine[lines=2]{\zallmancaps \color{Red} N}{unc} inquit domine servum tuum. Videte iustum velut corporee carcere molis inclusum, velle dissolvi: ut incipiat esse cum Christo. Dissolvi enim et esse cum Christo: multo melius est. Sed qui vult dimitti veniat in templum, veniat in Hierusalem, expectet Christum domini. Accipiat in manibus verbum Dei, complectatur velut quibusdam fidei sue brachiis, tunc dimittetur ut non videat mortem: quia viderit vitam.
\begin{responsory-final}[gaude-maria-virgo]
{Gaude Maria virgo cunctas hereses sola interemisti, que Gabrielis archangeli dictis credidisti, * Dum virgo deum et hominem genuisti, * Et post partum virgo inviolata permansisti.}{Gabrielem archangelum scimus divinitus te esse affatum, uterum tuum de spiritu sancto credimus impregnatum, erubescat Iudeus infelix, qui dicit Christum ex Ioseph semine esse natum.}{Et post}
\end{responsory-final}%
\newline \ul{Post gloria predicti responsorii, in Vesperis, dicatur immediate Prosa. In Matutinis autem non dicatur.} {\color{Red} Prosa.}
\lettrine[lines=2]{\zallmancaps \color{Blue} I}{nviolata} \hypertarget{inviolata-intacta-prosa}{\label{inviolata-intacta-prosa}} intacta et casta es Maria.
Que es effecta fulgida celi porta.
O mater alma Christi karisma.
Suscipe pie laudum preconia.
Nostra ut pura pectora sint et corpora.
Te nunc flagitant devota corda et ora.
Tua per precata dulcissona.
Nobis concedas veniam per secula.
O benigna que sola inviolata permansisti.
\newline \V Ora pro nobis sancta dei genitrix. \R Ut digni efficiamur promissionibus Christi.
\newline {\color{Red} In Laudibus. Ant.} Symeon iustus et timoratus expectabat redemptionem Israhel, et spiritus sanctus erat in eo. {\color{Red} Ant.} Responsum accepit Symeon a spiritu sancto non visurum se mortem nisi videret dominum. {\color{Red} Ant.} Accipiens Symeon puerum in manibus gratias agens benedixit deum. {\color{Red} Ant.} Obtulerunt pro eo domino par turturum, aut duos pullos columbarum. {\color{Red} Ant.} Lumen ad revelationem gentium, et gloriam plebis tue Israhel. {\color{Red} Cap.} Ecce ego. {\color{Red} Ymnus.}
\hymn{o-gloriosa}{O gloriosa domina}{
\lettrine[lines=2]{\zallmancaps \color{Red} O}{}\ gloriosa domina
excelsa supra sydera,
qui te creavit provide
lactasti sacro ubere.
\par {\bfseries \color{Blue} Q}uod Eva tristis abstulit
tu reddis almo germine,
intrent ut astra flebiles
celi fenestra facta es.
\par {\bfseries \color{Red} T}u regis alti ianua,
et porta lucis fulgida,
vitam datam per virginem
gentes redempte plaudite. {\color{Red} Tempore Resurrectionis.} Quesumus actor.
\par {\bfseries \color{Blue} G}loria tibi domine
qui natus es de virgine
cum patre et sancto spiritu,
in sempiterna secula. Amen.
}
\newline \V Elegit eam deus et preelegit eam. \R Et habitare eam facit in tabernaculo suo.
\newline {\color{Red} Ad Ben. Ant.} Senex puerum portabat, puer autem senem regebat, quem virgo peperit, et post partum virgo permansit, ipsum quem genuit, adoravit.
\newline {\color{Red} Ad Primam.} \R \hyperlink{ihesu-christe-fili}{Ihesu Christe.} \V Qui natus es de virgine Maria.
\newline \ul{Iste versus dicatur in omnibus sollemnitatibus Beate Virginis et per octavas quando de ipsis agitur et in Sabbatis quando de ipsa agitur in conventu.}
\newline {\color{Red} Ad Terciam. Cap.} Ecce ego.
\begin{responsory-breve}[sancta-dei-genitrix]
{Sancta dei genitrix virgo semper Maria, * Alleluya alleluya.}{Intercede pro nobis ad dominum deum nostrum.}
\end{responsory-breve}%
\V Post partum.
\newline {\color{Red} Ad Sextam. Cap.}
\lettrine[lines=2]{\zallmancaps \color{Blue} I}{pse} enim quasi ignis conflans et quasi herba fullonum, et sedebit conflans et emundans argentum: et purgabit filios Levi.
\begin{responsory-breve}[post-partum]
{Post partum virgo inviolata permansisti, * Alleluya alleluya.}{Dei genitrix intercede pro nobis.}
\end{responsory-breve}%
\V Speciosa facta es.
\newline {\color{Red} Ad Nonam. Capitulum.}
\lettrine[lines=2]{\zallmancaps \color{Red} E}{t} colabit eos quasi aurum, et quasi argentum, et erunt
%pp123
domino offerentes sacrificia in iusticia.
\begin{responsory-breve}[speciosa-facta-es]
{Speciosa facta es et suavis, * Alleluya alleluya.}{In deliciis tuis sancta dei genitrix.}
\end{responsory-breve}%
\V Elegit eam.
\newline \ul{Infra Septuagesimam non dicantur responsoria cum alleluya.}
\newline {\color{Red} Ad Vesperas. Super psalmos. Ant.} Symeon. {\color{Red} Psalmi.} \psalm{109} \psalm{112} \psalm{121} \psalm{126} \psalm{147}
\newline \ul{Predicti Psalmi dicantur in Vesperis in omnibus festivitatibus Beate Marie. Cap., Ymnus, \Vbar , ut in Primis Vesperis.}
\newline {\color{Red} Ad Mag. Ant.} Hodie beata virgo Maria puerum Ihesum presentavit in templo, et Symeon repletus spiritu sancto accepit eum in ulnas suas et benedixit deum.
{\ubsubsection{Sancti Blasii Episcopi et Martiris}{sancti-blasii-episcopi-et-martiris-breviarium}}
\newline {\color{Red} Oratio.}
\lettrine[lines=2]{\zallmancaps \color{Red} D}{a} quesumus omnipotens deus ut qui beati Blasii martiris tui atque pontificis sollemnia colimus, eius apud te intercessionibus adiuvemur. Per.
{\ubsubsection{Sancte Agathe Virginis et Martiris}{sancte-agathe-virginis-et-martiris-breviarium}}
\newline {\color{Red} Ad Mag. Ant.} Agatha virgo sacra nobili orta genere gloriosam propter Christum pertulit passionem. {\color{Red} Oratio.}
\lettrine[lines=2]{\zallmancaps \color{Blue} D}{eus} qui inter cetera potentie tue miracula eciam in sexu fragili victoriam martyrii contulisti, concede propicius, ut cuius natalicia colimus, per eius ad te exempla gradiamur. Per.
\newline {\color{Red} Ad Matutinas. In Primo Nocturno. Ant.} Ingenua sum et expectabilis genere ut omnis parentela mea testatur. {\color{Red} Ps.} \psalm{8} {\color{Red} Ant.} Ancilla Christi sum, ideo me ostendo servilem personam. {\color{Red} Ps.} \psalm{18} {\color{Red} Ant.} Summa ingenuitas ista est in qua servitus Christi comprobatur. {\color{Red} Ps.} \psalm{23} \R
\lettrine[lines=2]{\zallmancaps \color{Red} D}{um} \hypertarget{dum-ingrederetur-beata}{\label{dum-ingrederetur-beata}} ingrederetur beata Agatha in carcerem dixit ad iudicem, * Impie crudelis et dire tyranne, non es confusus amputare in femina quod ipse in matre suxisti. \V Ego habeo mammillas integras intus in anima mea quas ab infantia domino consecravi. Impie.
\begin{responsory}[vidisti-domine]
{Vidisti domine et expectasti agonem meum, quomodo pugnavi in stadio, sed quia nolui obedire mandatis principum, * Iussa sum in mammilla torqueri.}{Propter veritatem et mansuetudinem et iusticiam.}
\end{responsory}%
\begin{responsory-doxology}[quis-es-tu]
{Quis es tu qui venisti ad me curare vulnera mea, ego sum apostolus Christi nichil in me dubites filia, ipse me misit ad te, * Quem dilexisti mente et puro corde.}{Nam et ego apostolus eius sum, et in nomine eius scias te esse salvandam.}
\end{responsory-doxology}%
\newline {\color{Red} In Secundo Nocturno. Ant.} Agatha letissime et glorianter ibat ad carcerem, et quasi ad epulas invitata, agonem suum domino precibus commendabat. {\color{Red} Ps.} \psalm{44} {\color{Red} Ant.} Si ignem adhibeas rorem michi salvificum de celo angeli ministrabunt. {\color{Red} Ps.} \psalm{45} {\color{Red} Ant.} Nisi diligenter perfeceris corpus meum a carnificibus attrectari, non potest anima mea in paradisum domini cum palma intrare martyrii. {\color{Red} Ps.} \psalm{86}
\begin{responsory}[ipse-me-coronavit]
{Ipse me coronavit qui per apostolum suum in custodia me confortavit, pro eo quod iussa sum suspendi in eculeo propter fidem castitatis, * Adiuva me domine deus meus in tortura mammillarum mearum.}{Ipse me dignatus est ab omni plaga curare, et mammillam meam meo pectori restituere.}
\end{responsory}%
\begin{responsory}[agatha-letissime]
{Agatha letissime et glorianter ibat ad carcerem, et quasi ad epulas invitata, * Agonem suum domino precibus commendabat.}{Nobilissimis orta natalibus ab ignobili gaudens trahebatur ad carcerem.}
\end{responsory}%
\begin{responsory-doxology}[ego-autem-adiuvata]
{Ego autem adiuvata a domino perseverabo in confessione eius, * Qui me salvam fecit et consolatus est me.}{Mens mea solidata est in Christo fundata.}
\end{responsory-doxology}%
\newline {\color{Red} In Tertio Nocturno. Ant.} Agatha sancta dixit, si feras michi promittis, audito Christi nomine mansuescunt. {\color{Red} Ps.} \psalm{95} {\color{Red} Ant.} Vidisti domine agonem meum quomodo pugnavi in stadio, sed quia nolui obedire mandatis principum iussa sum in mammilla torqueri. {\color{Red} Ps.} \psalm{96} {\color{Red} Ant.} Propter fidem castitatis iussa sum suspendi in eculeo adiuva me domine deus meus in tortura mammillarum mearum. {\color{Red} Ps.} \psalm{97}
\begin{responsory}[gaudeamus-omnes]
{Gaudeamus omnes in domino diem festum celebrantes sub honore Agathe virginis, * De cuius passione gaudent angeli et collaudant filium dei.}{Immaculatus dominus immaculatam sibi famulam in hoc fragilitatis corpore positam misericorditer consecravit.}
\end{responsory}%
\begin{responsory}[qui-me-dignatus]
{Qui me dignatus est ab omni plaga curare et mammillam meam meo pectori restituere, * Ipsum invoco deum vivum.}{Medicinam carnalem corpori meo nunquam exhibui, sed habeo dominum Ihesum Christum qui solo sermone restaurat universa.}
\end{responsory}%
\begin{responsory-doxology}[beata-agatha]
{Beata Agatha ingressa carcerem expandit manus suas ad deum et dixit, domine qui me fecisti vincere tormenta carnificum, * Iube me domine ad tuam misericordiam pervenire.}{Domine qui me creasti, et tulisti a me amorem seculi, qui corpus meum a pollutione separasti.}
\end{responsory-doxology}%
\newline {\color{Red} In Laudibus. Ant.} Quis es tu qui venisti ad me curare vulnera mea, ego sum apostolus Christi nichil in me dubites filia. {\color{Red} Ant.} Medicinam carnalem corpori meo nunquam exhibui, sed habeo dominum Ihesum Christum, qui solo sermone restaurat universa. {\color{Red} Ant.} Gratias tibi ago domine, quia memor es mei, et misisti ad me apostolum tuum curare vulnera mea. {\color{Red} Ant.} Benedico te pater domini mei Ihesu Christi, quia per apostolum tuum mammillam meam meo pectori restituisti. {\color{Red} Ant.} Qui me dignatus est ab omni plaga curare et mammillam meam meo pectori restituere, ipsum invoco deum vivum.
\newline {\color{Red} Ad Ben. Ant.} Paganorum multitudo fugiens ad sepulchrum virginis, tulerunt velum eius contra ignem, ut comprobaret dominus quod a periculis incendii meritis Agathe martyris sue eos liberaret.
\newline {\color{Red} Ad Mag. Ant.} Mentem sanctam, spontaneam, honorem deo, et patrie liberationem.
{\ubsubsection{Sanctorum Vedasti et Amandi Episcoporum}{sanctorum-vedast-et-amandi-episcoporum-breviarium}}
\newline {\color{Red} Oratio.}
\lettrine[lines=2]{\zallmancaps \color{Blue} D}{eus} qui nos sanctorum confessorum tuorum Vedasti et Amandi confessionibus gloriosis circundas et protegis, da nobis et eorum imitatione proficere et intercessione gaudere. Per.
{\ubsubsection{Sancte Scolastice Virginis}{sancte-scolastice-virginis-breviarium}}
\newline {\color{Red} Oratio.}
\lettrine[lines=2]{\zallmancaps \color{Red} E}{xaudi} nos deus salutaris noster: ut sicut de beate Scolastice virginis tue festivitate gaudemus, ita pie devotionis erudiamur affectu. Per.
{\ubsubsection{Sancti Valentini Martiris}{sancti-valentini-martiris-breviarium}}
\newline {\color{Red} Oratio.}
\lettrine[lines=2]{\zallmancaps \color{Blue} P}{resta} quesumus omnipotens deus, ut qui beati Valentini martyris tui natalicia colimus, a cunctis malis imminentibus eius intercessione liberemur. Per.
{\ubsubsection{In Cathedra Sancti Petri}{in-cathedra-sancti-petri-breviarium}}
\newline {\color{Red} Ad Vesperas. Capitulum.} Dilectus deo. {\color{Red} Ymnus.}
\hymn{iam-bone-pastor}{Iam bone pastor}{
\lettrine[lines=2]{\zallmancaps \color{Red} I}{am} bone pastor Petre clemens accipe,
vota precantum, et peccati vincula,
resolve tibi potestate tradita,
qua cunctis celum, verbo claudis, aperis.
\par {\bfseries \color{Blue} S}it trinitati sempiterna gloria,
honor potestas atque iubilatio,
in unitate cui manet imperium
ex tunc et modo per eterna secula. Amen.
}
\newline \V Tu es Petrus. \R Et super hanc Petram edificabo ecclesiam meam.
\newline {\color{Red} Ad Mag. Ant.} Tu es pastor ovium princeps apostolorum tibi tradite sunt claves regni celorum. {\color{Red} Oratio.}
\lettrine[lines=2]{\zallmancaps \color{Blue} D}{eus} qui beato Petro apostolo tuo collatis clavibus regni celestis ligandi atque solvendi pontificium tradidisti, concede ut intercessionis eius auxilio, a peccatorum nostrorum nexibus liberemur. Qui vivis.
\newline {\color{Red} Ad Matutinas. Invit.}
\begin{invitatory}[tu-es-pastor-invitatorium]
{Tu es pastor ovium princeps apostolorum, * Tibi tradidit deus claves regni celorum.}
\end{invitatory}
\newline {\color{Red} Ymnus.} \hyperlink{iam-bone-pastor}{Iam bone.} \ul{Lectiones unius confessoris pontificis.}
\newline {\color{Red} Nonum.} \hypertarget{quodcumque-ligaveris}{\label{quodcumque-ligaveris}} \R Quodcumque ligaveris super terram erit ligatum et in celis, * Et quodcumque solveris super terram erit solutum et * In celis. \V Et ego dico tibi quia tu es Petrus et super hanc petram edificabo ecclesiam meam. Gloria. In celis.
\newline {\color{Red} Ante Laudes.} \V Exaltent eum in ecclesia plebis. \R Et in cathedra seniorum laudent eum. {\color{Red} In Laudibus.} \V Iustus germinabit.
\newline {\color{Red} Ad Ben. Antiphona.} Quodcumque ligaveris super terram erit ligatum et in celis et quodcumque solveris super terram erit solutum et in celis dixit dominus Symoni Petro.
\newline \ul{Ad Vesperas. Cap., Ymnus, \Vbar , ut in Primis Vesperis.}
\newline {\color{Red} Ad Mag. Ant.} Symon Bariona, tu vocaberis Cephas quod interpretatur Petrus, ianitor celi pulsantibus aperi, supra modum peccavimus omnes dimitte septuagies septies. \ul{Cetera omnia fiant de communi unius confessoris.}
{\ubsubsection{Sancti Mathie Apostoli}{sancti-mathie-apostoli-breviarium}}
\newline {\color{Red} Oratio.}
\lettrine[lines=2]{\zallmancaps \color{Red} D}{eus} qui beatum Mathiam apostolorum tuorum collegio sociasti, tribue quesumus, ut eius interventione tue circa nos pietatis semper viscera sentiamus. Per.
\newline \ul{In anno bissextili celebretur predicta festivitas in secunda die bissexti, nisi quarta feria Cinerum ipsa die evenerit, tunc enim in precedenti feria tercia celebrabitur.}
{\ubsubsection{Sancti Albini Episcopi et Confessoris}{sancti-albini-episcopi-et-confessoris-breviarium}}
\newline {\color{Red} Oratio.}
\lettrine[lines=2]{\zallmancaps \color{Blue} D}{a} quesumus omnipotens deus, ut beati Albini confessoris tui atque pontificis veneranda sollemnitas et devotionem nobis augeat et salutem. Per.
{\ubsubsection{Sancti Gregorii Pape}{sancti-gregorii-pape-breviarium}}
\newline {\color{Red} Oratio.}
\lettrine[lines=2]{\zallmancaps \color{Red} D}{eus} qui nos beati Gregorii confessoris tui atque pontificis annua sollemnitate letificas, concede propicius, ut cuius opere atque doctrina scripture tue mysteria multa cognovimus eius apud te semper patrocinia sentiamus. Per.%typo: annu
{\ubsubsection{Sancti Benedicti Abbatis}{sancti-benedicti-abbatis-breviarium}}
\newline {\color{Red} Oratio.}
\lettrine[lines=2]{\zallmancaps \color{Blue} I}{ntercessio} nos quesumus domine beati Benedicti abbatis commendet, ut quod nostris meritis non valemus, eius patrocinio assequamur. Per.
{\ubsubsection{In Annunciatione Dominica}{in-annunciatione-dominica-breviarium}}
\newline {\color{Red} Ad Vesperas. Super psalmos. Ant.} Ave Maria gratia plena dominus tecum, benedicta tu in mulieribus. {\color{Red} Tempore Resurrectionis.} Alleluya. \ul{Idem fiat de aliis alleluya que infra ponuntur post antiphonas et responsoria, scilicet quod tantum dicantur tempore resurrectionis.} {\color{Red} Ps.} \psalm{112} \ul{cum ceteris.} {\color{Red} Cap.}
\lettrine[lines=2]{\zallmancaps \color{Red} E}{cce} virgo concipiet et pariet filium, et vocabitur nomen eius Emmanuel, butirum et mel comedet, ut sciat reprobare malum, et eligere bonum. \R \hyperlink{quomodo-fiet-istud}{Quomodo.} {\color{Red} IX. Ymnus.} \hyperlink{ave-maris-stella}{Ave maris stella.} \V Rorate celi desuper et nubes pluant iustum. \R Aperiatur terra et germinet salvatorem.
\newline {\color{Red} Ad Mag. Ant.} Orietur sicut sol salvator mundi et descendet in uterum virginis sicut ymber super gramen alleluya alleluya alleluya. {\color{Red} Oratio.}
\lettrine[lines=2]{\zallmancaps \color{Blue} D}{eus} qui de beate Marie virginis utero, verbum tuum angelo nuntiante carnem suscipere voluisti, presta supplicibus tuis, ut qui vere eam genitricem dei credimus, eius apud te intercessionibus adiuvemur. Per eundem.
\newline {\color{Red} Ad Completorium. Super psalmos. Ant.} Ecce virgo concipiet et pariet filium et vocabitur nomen eius Emmanuel, alleluya alleluya.
\newline {\color{Red} Ad Nunc Dimittis. Ant.} Ecce ancilla domini, fiat michi secundum verbum tuum, alleluya.
\newline {\color{Red} Ad Mat. Invitat.}
\begin{invitatory}[ave-maria-invitatorium]
{Ave Maria gratia plena, * Dominus tecum, * Alleluya.}
\end{invitatory}
\newline {\color{Red} Ymnus.} \hyperlink{quem-terra}{Quem terra.}
\newline {\color{Red} In Primo Nocturno. Ant.} Ecce virgo concipiet et pariet filium et vocabitur nomen eius Emmanuel, Alleluya alleluya. {\color{Red} Ps.} \psalm{8} {\color{Red} Ant.} Beata es Maria que credidisti, perficientur in te que dicta sunt tibi a domino Alleluya. {\color{Red} Ps.} \psalm{18} {\color{Red} Ant.} Angelus domini nunciavit Marie et concepit de spiritu sancto, Alleluya. {\color{Red} Ps.} \psalm{23} \V Ex Syon species decoris eius.
\newline \ul{Sermo qui legitur in multis ecclesiis et videtur esse Leonis pape.} \grecross \ {\color{Red} Lectio Prima.}
\lettrine[lines=2]{\zallmancaps \color{Red} G}{audeamus} in domino dilectissimi, et licet non quantum debemus, tamen quales possumus gratias nostro referamus actori, ne superabundanti eius in nobis gratie reperiamur ingrati. Novum enim virginis conceptum nobis festivitas hodierna commendat: in qua nostre reparationis celebratur exordium. \R
\lettrine[lines=2]{\zallmancaps \color{Blue} M}{issus} est Gabriel angelus ad Mariam virginem desponsatam Ioseph nuncians ei verbum et expavescit virgo de lumine, ne timeas Maria, invenisti gratiam apud dominum, ecce concipies et paries, * Et vocabitur altissimi filius, Alleluya. \V Dabit ei dominus deus sedem David patris eius et regnabit in domo Iacob in eternum. Et vocabitur. %already defined missus-est
\newline {\color{Red} Lectio Secunda.}
\lettrine[lines=2]{\zallmancaps \color{Blue} C}{ertum} proponitur nobis divine pietatis et potestatis indicium. Si enim rerum dominus fugitivos servos requirens: iudicium exercere et non pietatem exhibere veniret, nequaquam lutei vasis huius fragilitatem, in qua nobis compati et pro nobis pati posset: indueret.
\begin{responsory}%already defined [non-auferetur]
{Non auferetur sceptrum de Iuda, et dux de femore eius, donec veniat qui mittendus est, * Et ipse erit expectatio gentium.}{Pulchriores sunt oculi eius vino et dentes eius lacte candidiores.}
\end{responsory}%
\newline {\color{Red} Lectio Tertia.}
\lettrine[lines=2]{\zallmancaps \color{Red} E}{t} quamvis hoc stultum et infirmum videatur phylosophie nomen retinentibus: quid tamen potentius quam contra iura nature virgini conceptum dare, et per mortem carnis mortalem superbiam ad immortalitatis gloriam revocare? Hec stulticia facta est nobis a deo magna sapiencia, quoniam vite perdite reparatrix nobis extitit medicina.%typo: physosophie
\begin{responsory-final}%already defined [ave-maria]
{Ave Maria gratia plena dominus tecum, * Spiritus sanctus superveniet in te et virtus altissimi obumbrabit tibi, * Quod enim ex te nascetur sanctum, vocabitur filius dei, Alleluya.}{Quomodo fiet istud quoniam virum non cognosco et respondens angelus dixit ei.}{Quod enim}
\end{responsory-final}%
\newline {\color{Red} In Secundo Nocturno. Ant.} Emissiones tue paradisus malorum punicorum cum pomorum fructibus. {\color{Red} Ps.} \psalm{44} {\color{Red} Ant.} Fons ortorum puteus aquarum viventium, que fluunt impetu de Libano. {\color{Red} Ps.} \psalm{45} {\color{Red} Ant.} Veniat dilectus meus in ortum suum ut comedat fructum pomorum suorum. {\color{Red} Ps.} \psalm{86} \V Egredietur virga de radice Iesse.
\newline {\color{Red} Lectio Quarta.}
\lettrine[lines=2]{\zallmancaps \color{Blue} S}{i} nobis terrena hereditate privatis et in obscurissimo carcere positis, regalis aliqua persona regio decore deposito condescenderet, ut tanquam unus ex nobis nostre captivitatis erumpnas toleraret, quatinus oportunitate inventa, nos ab ipsa liberaret, qua ratione posset melius suam dilectionem intimare, et nostram erga se excitare?
\begin{responsory}%already defined[suscipe-verbum]
{Suscipe verbum virgo Maria quod tibi a domino per angelum transmissum est, concipies et paries deum pariter et hominem, * Ut benedicta dicaris inter omnes mulieres.}{Paries quidem filium et virginitatis non patieris detrimentum efficieris gravida et eris mater semper intacta.}
\end{responsory}%
\newline {\color{Red} Lectio Quinta.}
\lettrine[lines=2]{\zallmancaps \color{Red} H}{uius} ergo tam necessarie liberationis admiranda initia nobis ad memoriam revocata, toto mentis affectu diligamus, debitis obsequiis honoremus, quatinus de iocundis principiis iocundos exitus expectare mereamur. Fit hodie porta celi uterus virginis, per quam deus descendit ad homines, ut eis ascensum preberet ad celum.
\begin{responsory}%already defined [ecce-virgo]
{Ecce virgo concipiet et pariet filium dicit dominus, et vocabitur nomen eius, * Admirabilis deus fortis.}{Super solium David et super regnum eius sedebit in eternum.}
\end{responsory}%
\newline {\color{Red} Lectio Sexta.}
\lettrine[lines=2]{\zallmancaps \color{Blue} D}{iligamus} karissimi misericordiam per quam redempti sumus cum perditi essemus, conservemus castitatem, cuius amatorem se esse monstravit qui de caste mulieris visceribus, immaculatum sibi corpus aptavit. Conformemur ei, qui vitam suam in terris regulam nobis proposuit, qui primo adventu suo voluit nos intus ad ymaginem suam reformare, secundo reformabit corpus humilitatis nostre configuratum corpori claritatis sue.
\begin{responsory-final}[egredietur-virga]
{Egredietur virga de radice Iesse et flos de radice eius ascendet, * Et erit iusticia cingulum lumborum eius, * Et fides cinctorum renum eius.}{Et requiescet semper eum spiritus domini, spiritus sapientie et intellectus.}{Et fides}
\end{responsory-final}%
\newline {\color{Red} In Tertio Nocturno. Ant.} Veni in ortum meum soror mea sponsa, messui mirram meam cum aromatibus meis. {\color{Red} Ps.} \psalm{95} {\color{Red} Ant.} Comedi favum cum melle meo, bibi vinum meum cum lacte meo. {\color{Red} Ps.} \psalm{96} {\color{Red} Ant.} Talis est dilectus meus, et ipse est amicus meus filie Hierusalem. {\color{Red} Ps.} \psalm{97} \V Egredietur dominus de loco sancto suo.
\newline {\color{Red} Lectio VII. Secundum Lucam.}
\lettrine[lines=2]{\zallmancaps \color{Red} I}{n} illo tempore: Missus est angelus Gabriel a Deo in civitatem Galilee cui nomen Nazareth ad virginem desponsatam viro cui nomen erat Ioseph de domo David: et nomen virginis Maria. Et reliqua.
\newline {\color{Red} Omelia Venerabilis Bede Presbiteri.}
\lettrine[lines=2]{\zallmancaps \color{Blue} E}{xordium} redemptionis nostre fratres karissimi hodierna nobis sancti evangelii lectio commendat, que angelum de celis a Deo missum narrat ad virginem, ut novam in carne nativitatem filii Dei predicaret, per quam nos abiecta vetustate noxia renovari, atque inter filios Dei computari possimus.
\begin{responsory}[rorate-celi]
{Rorate celi desuper et nubes pluant iustum, * Aperiatur terra et germinet salvatorem.}{Emitte agnum domine dominatorem terre, de petra deserti ad montem filie Syon.}
\end{responsory}%
\newline {\color{Red} Lectio VIII.}
\lettrine[lines=2]{\zallmancaps \color{Red} U}{t} ergo ad promisse salutis mereamur dona pertingere, primordium eius intenta curemus aure percipere. Missus est inquit angelus Gabriel a Deo in civitatem Galilee cui nomen Nazareth, ad virginem desponsatam viro cui nomen erat Ioseph.
\begin{responsory}[nascetur-nobis]
{Nascetur nobis parvulus, et vocabitur deus fortis, ipse sedebit super thronum David patris sui et imperabit, * Cuius potestas super humerum eius.}{In ipso benedicentur omnes tribus terre.}
\end{responsory}%
\newline {\color{Red} Lectio IX.}
\lettrine[lines=2]{\zallmancaps \color{Blue} A}{ptum} profecto humane restaurationis principium, ut angelus a Deo mitteretur ad virginem partu consecrandam divino, quia prima perditionis humane fuit causa cum serpens a diabolo mittebatur ad mulierem spiritu superbie decipiendam. immo ipse in serpente diabolus veniebat, qui genus humanum deceptis parentibus primis immortalitatis gloria nudaret.
\begin{responsory-final}[quomodo-fiet-istud]
{Quomodo fiet istud respondens ait Maria quia virum non cognosco, angelus ad hec, * Spiritus sanctus superveniet in te, * Et virtus altissimi obumbrabit tibi, * Alleluya.}{Ideoque quod nascetur ex te sanctum, vocabitur filius dei.}{Et virtus}{\color{Red} Vel} Alleluya.
\end{responsory-final}%
\newline \V Emitte agnum domine dominatorem terre.
\newline {\color{Red} In Laudibus. Ant.} Missus est Gabriel angelus ad Mariam virginem desponsatam Ioseph, Alleluya alleluya. {\color{Red} Ant.} Ave Maria gratia plena dominus tecum, benedicta tu in mulieribus, alleluya. {\color{Red} Ant.} Ne timeas Maria invenisti gratiam apud dominum, ecce concipies et paries filium, Alleluya. {\color{Red} Ant.} Dabit ei dominus sedem David patris eius et regnavit in eternum, Alleluya. {\color{Red} Ant.} Ecce ancilla domini fiat michi secundum verbum tuum, Alleluya. {\color{Red} Cap.} Ecce virgo. {\color{Red} Ymnus.} \hyperlink{o-gloriosa}{O gloriosa.} \V Elegit eam.
\newline {\color{Red} Ad Ben. Ant.} Quomodo fiet istud angele dei quia virum in concipiendo non pertuli: audi Maria virgo Christi, spiritus sanctus superveniet in te et virtus altissimi obumbrabit tibi, Alleluya.
\newline {\color{Red} Ad Terciam. Cap.} Ecce virgo. \R \hyperlink{sancta-dei-genitrix}{Sancta dei genitrix.} \V Post partum.
\newline {\color{Red} Ad Sextam. Cap.}
\lettrine[lines=2]{\zallmancaps \color{Red} E}{gredietur} virga de radice Iesse, et flos de radice eius ascendet, et requiescet super eum spiritus domini. \R \hyperlink{post-partum}{Post partum.} \V Speciosa.
\newline {\color{Red} Ad Nonam. Capitulum.}
\lettrine[lines=2]{\zallmancaps \color{Blue} R}{orate} celi desuper et nubes pluant iustum, aperiatur terra et germinet salvatorem, et iusticia oriatur simul: ego dominus creavi eum. \R \hyperlink{speciosa-facta-es}{Speciosa.} \V Elegit eam.
\newline {\color{Red} Ad Vesperas. Super psalmos. Ant.} Missus est. \ul{Cap., Ymnus, \Vbar , ut in Primis Vesperis.}
\newline {\color{Red} Ad Mag. Ant.} Hec est dies quam fecit dominus, hodie dominus afflictionem populi sui respexit et redemptionem misit, hodie mortem quam femina intulit femina fugavit, hodie deus homo factus, id quod fuit permansit, et quod non erat assumpsit, ergo exordium nostre redemptionis devote recolamus et exultemus dicentes, gloria tibi domine, Alleluya alleluya. \ul{Completorium ut supra.}
\newline \ul{In Paschali vero tempore, hoc modo fiat officium. Ad omnes horas preterquam ad Matutinas, omnia fiant sicut prenotatum est, nisi quod in fine ymnorum dicetur,} Quesumus actor, \ul{et} Gloria tibi domine qui natus, \ul{in ymnis eiusdem metri et responsoria ad horas diei et \Vbar , dicentur cum alleluya. Ad Matutinas autem invitatorium, ymnus ut supra. Tres prime antiphone cum suis psalmis, \Vbar .} Ex Syon. \ul{Tres lectiones de omelia evangelii,} Missus est. \ul{Responsorium primum,} \hyperlink{missus-est}{Missus.} \ul{Responsorium secundum:} \hyperlink{ave-maria}{Ave Maria.} \ul{Responsorium tercium:} \hyperlink{quomodo-fiet-istud}{Quomodo.} \ul{\Vbar .} Emitte agnum. \ul{Cetera ut supra. Completorium vero Paschalis temporis non mutatur nisi in antiphonis.}
{\ubsubsection{Sancti Ambrosii Episcopi et Confessoris}{sancti-ambrosii-episcopi-et-confessoris-breviarium}}
\newline {\color{Red} Oratio.}
\lettrine[lines=2]{\zallmancaps \color{Red} D}{eus} qui populo tuo eterne salutis beatum Ambrosium ministrum tribuisti, presta quesumus, ut quem doctorem vite habuimus in terris, intercessorem habere mereamur in celis. Per.
{\ubsubsection{In Communi Unius Apostoli Seu Plurimorum, Sive Unius Evangeliste in Paschali Tempore}{in-communi-unius-apostoli-seu-plurimorum-sive-unius-evangeliste-in-paschali-tempore-breviarium}}
\newline {\color{Red} Ad Vesperas. Super psalmos. Ant.} Estote. \R \hyperlink{candidi-facti-sunt}{Candidi.} {\color{Red} Ad Mag. Ant.} Filie Hierusalem, venite et videte martyrem cum corona qua coronavit eum dominus in die sollemnitatis et leticie alleluya alleluya.
\newline {\color{Red} Ad Matutinas. Invitat.}
\begin{invitatory}[exultent-in-domino-invitatorium]
{Exultent in domino sancti, * Alleluya.}
\end{invitatory}
\newline {\color{Red} Super psalmos. Ant.} Tristicia nostra alleluya vertetur in gaudium alleluya. \ul{Tres psalmi et versiculus ante lectiones de apostolis, dicantur secundum ordinem nocturnorum.} \R
\lettrine[lines=2]{\zallmancaps \color{Blue} V}{irtute} \hypertarget{virtute-magna}{\label{virtute-magna}} magna reddebant apostoli * Testimonium resurrectionis Ihesu Christi domini nostri, alleluya alleluya. \V Repleti quidem spiritu sancto loquebantur cum fiducia verbum dei. Testimonium.
\begin{responsory}[isti-sunt-agni]
{Isti sunt agni novelli qui annunciaverunt alleluya, modo venerunt ad fontes, * Repleti sunt claritate alleluya alleluya.}{In conspectu agni amicti stolis albis et palme in manibus eorum.}
\end{responsory}%
\begin{responsory-final}[candidi-facti-sunt]
{Candidi facti sunt Nazarei eius alleluya, splendorem deo dederunt alleluya, * Et sicut lac coagulati sunt * Alleluya alleluya.}{In omnem terram exivit sonus eorum, et in fines orbis terre verba eorum.}{Alleluya}
\end{responsory-final}%
\newline {\color{Red} In Laudibus. Ant.} Sancti tui domine florebunt sicut lilium alleluya, et sicut odor balsami erunt ante te alleluya. {\color{Red} Ant.} Sancti et iusti in domino gaudete alleluya vos elegit deus in hereditatem sibi, alleluya. {\color{Red} Ant.} In velamento clamabunt sancti tui domine, alleluya alleluya alleluya. {\color{Red} Ant.} Spiritus et anime iustorum ymnum dicite deo nostro alleluya alleluya. {\color{Red} Ant.} In celestibus regnis sanctorum habitatio est alleluya et in eternum requies eorum alleluya.
\newline {\color{Red} Ad Ben. Ant.} Lux perpetua lucebit sanctis tuis domine alleluya, et eternitas temporum alleluya alleluya alleluya. \ul{Ad horas diei antiphone Laudum.}
\newline {\color{Red} Ad Vesperas. Super psalmos. Ant.} Sancti tui. {\color{Red} Ps.} \psalm{109} \ul{et ceteri de Apostolis.}
\newline {\color{Red} Ad Mag. Ant.} In civitate domini clare sonant iugiter organa sanctorum, ibi angeli et archangeli ymnum decantant ante thronum dei, alleluya.
\ul{Cetera omnia sumantur de communi apostolorum et evangelistarum alterius temporis.}
{\ubsubsection{In Communi Unius Martyris Sive Confessoris, Sive Trium Lectionum Fuerit Sive Supra, in Tempore Paschali}{in-communi-unius-martyris-sive-confessoris-sive-trium-lectionum-fuerit-sive-supra-in-tempore-paschali-breviarium}}
\newline {\color{Red} Ad Mag. Ant.} Filie Hierusalem. {\color{Red} Ad Matutinas. Invitat.} \hyperlink{exultent-in-domino-invitatorium}{Exultent.} {\color{Red} Super psalmos. Ant.} Tristicia. \ul{Tres psalmi et versiculus dicantur de martyre vel de confessore secundum ordinem nocturnorum.} \R
\lettrine[lines=2]{\zallmancaps \color{Red} B}{eatus} \hypertarget{beatus-vir-qui-metuit}{\label{beatus-vir-qui-metuit}} vir qui metuit dominum alleluya, * In mandatis eius cupit nimis, alleluya alleluya alleluya. \V Potens in terra erit semen eius, generatio rectorum benedicetur. In mandatis.
\begin{responsory}[de-ore-prudentis]
{De ore prudentis procedit mel alleluya, dulcedo mellis est lingua eius, Alleluya, * Favus distillans labia eius alleluya alleluya.}{Sapientia requiescit in corde eius, et prudentia in sermone oris illius.}
\end{responsory}%
\begin{responsory-final}[filie-hierusalem]
{Filie Hierusalem venite et videte martyrem cum corona qua coronavit eum dominus, * In die sollemnitatis et leticie, * Alleluya.}{Induit illum dominus stola glorie et coronavit eum.}{Alleluya}
\end{responsory-final}%
\newline {\color{Red} In Laudibus. Ant.} Sancti tui. \ul{et alie ut supra. %
%pp124
In festo vero trium lectionum una tantum dicatur.} {\color{Red} Ad Ben. Ant.} Lux perpetua. \ul{Ad horas antiphone Laudum.}
\newline {\color{Red} Ad Vesperas. Super psalmos. Ant.} Sancti tui. {\color{Red} Ad Mag. Ant.} In civitate. \ul{Cetera omnia sumantur de communi unius martiris sive confessoris alterius temporis.}
{\ubsubsection{In Communi Plurimorum Martyrum, in Tempore Paschali}{in-communi-plurimorum-martyrum-in-tempore-paschali-breviarium}}
\lettrine[lines=2]{\zallmancaps \color{Blue} A}{\color{Red} d} {\color{Red} Mag. Ant.} Filie Hierusalem. {\color{Red} Ad Matutinas. Invitat.} \hyperlink{exultent-in-domino-invitatorium}{Exultent.} {\color{Red} Super psalmos. Ant.} Tristicia. \ul{Tres psalmi et versiculus dicantur de martyribus secundum ordinem nocturnorum.} \R
\lettrine[lines=2]{\zallmancaps \color{Red} P}{reciosa} \hypertarget{preciosa-in-conspectu}{\label{preciosa-in-conspectu}} in conspectu domini alleluya, * Mors sanctorum eius alleluya. \V In conspectu omnis populi eius. Mors.
\begin{responsory}[tristicia-vestra]
{Tristicia vestra alleluya, * Convertetur in gaudium alleluya alleluya.}{Mundus autem gaudebit vos autem contristabimini, sed tristicia vestra.}
\end{responsory}%
\begin{responsory-doxology}[lux-perpetua-lucebit]
{Lux perpetua lucebit sanctis tuis domine et eternitas temporum, * Alleluya alleluya.}{Leticia sempiterna super capita eorum gaudium et leticiam obtinebunt.}
\end{responsory-doxology}%
\newline {\color{Red} In Laudibus. Ant.} Sancti tui. {\color{Red} Ad Ben. Ant.} Lux perpetua. \ul{Cetera omnia sumantur de communi plurimorum martirum alterius temporis.}
{\ubsubsection{Ad Memoriam Unius Vel Plurimorum Sanctorum in Tempore Paschali}{ad-memoriam-unius-vel-plurimorum-sanctorum-in-tempore-paschali-breviarium}}
\newline {\color{Red} In Vesperis. Ant.} Filie Hierusalem. {\color{Red} In Laudibus. Ant.} Lux perpetua. \ul{Versiculi sicut in alio tempore.}
\newline \ul{De virginibus vero fiat memoria: sicut in alio tempore.}
{\ubsubsection{Sanctorum Tyburcii et Valeriani et Maximi Martyrum}{sanctorum-tyburcii-et-valeriani-et-maximi-martyrum-breviarium}}
\newline {\color{Red} Oratio.}
\lettrine[lines=2]{\zallmancaps \color{Blue} P}{resta} quesumus omnipotens deus: ut qui sanctorum tuorum Tiburcii et Valeriani et Maximi sollemnia colimus, eorum eciam virtutes imitemur. Per.
{\ubsubsection{Sancti Georgii Martiris}{sancti-georgii-martiris-breviarium}}
\newline {\color{Red} Oratio.}
\lettrine[lines=2]{\zallmancaps \color{Red} D}{eus} qui nos beati Georgii martyris tui meritis et intercessione letificas, concede propicius, ut qui tua per eum beneficia poscimus dono tue gratie consequamur. Per.
{\ubsubsection{Sancti Marci Evangeliste}{sancti-marci-evangeliste-breviarium}}
\newline {\color{Red} Oratio.}
\lettrine[lines=2]{\zallmancaps \color{Blue} D}{eus} qui beatum Marcum evangelistam tuum, evangelice predicationis gratia sublimasti, tribue quesumus, eius nos semper, et erudicione proficere et oratione defendi. Per.
{\ubsubsection{Sancti Vitalis Martiris}{sancti-vitalis-martiris-breviarium}}
\newline {\color{Red} Oratio.}
\lettrine[lines=2]{\zallmancaps \color{Red} P}{resta} quesumus omnipotens deus, ut intercedente beato Vitale martire tuo et a cunctis adversitatibus liberemur in corpore, et a pravis cogitationibus mundemur in mente. Per.
{\ubsubsection{In Festo Beati Petri Martiris}{in-festo-beati-petri-martiris-breviarium}}
\lettrine[lines=2]{\zallmancaps \color{Blue} A}{\color{Red} d} {\color{Red} Vesperas. Super psalmos. Ant. C}olletetur turba fidelium, triumphantis athlete gaudio, qui conservans pudoris lilium et coruscans doctrine radio, dum pro fide subit martyrum, trino felix potitur bravio alleluya. {\color{Red} Ps.} \psalm{112} \R \hyperlink{dum-sansonis}{Dum Sansonis.} {\color{Red} III. Ymnus.}
\hymn{magne-dies-leticie}{Magne dies leticie}{
\lettrine[lines=2]{\zallmancaps \color{Red} M}{agne} dies leticie
nobis illuxit celitus,
Petrus ad thornuim glorie
martyr pervenit inclitus.
\par {\bfseries \color{Blue} P}uer in fide claruit
parentum carens nebula,
deo servire studuit
sub paupertatis regula.
\par {\bfseries \color{Red} C}arnem afflixit iugiter
in labore multiplici,
viam sequens humiliter
patris sui Dominici.
\par {\bfseries \color{Blue} V}ita, mors, signa varia,
celum frequenti lumine,
dant Petro testimonia
de sanctitatis culmine.
\par {\bfseries \color{Red} Q}uesumus actor omnium
in hoc Paschali gaudio,
per ipsius suffragium
crescat nostra devotio.
\par {\bfseries \color{Blue} G}loria tibi domine
qui surrexisti a mortuis,
et fortes in certamine
sertis ornas perpetuis. Amen.
}
\newline \V Ora pro nobis beate Petre alleluya.
\newline {\color{Red} Ad Mag. Ant. O} Petre martyr inclite predicatorum gloria, virginitate predite, verbo, signis et grata, concessa nobis solite pietatis clementia, transacto mundi tramite nos transfer ad celestia, alleluya. {\color{Red} Oratio.}
\lettrine[lines=2]{\zallmancaps \color{Blue} P}{resta} quesumus omnipotens deus ut beati Petri martyris tui fidem congrua devotione sectemur, qui pro eiusdem fidei dilatione martyrii plamam meruit obtinere. Per.
\newline {\color{Red} Ad Matutinas. Invit.}
\begin{invitatory}[celestis-exercitus-invitatorium]
{Celestis exercitus regem adoremus, * Et pro palma militis eius iubilemus, alleluya.}
\end{invitatory}
\newline {\color{Red} Ymnus.}
\hymn{adest-triumphus-nobilis}{Adest triumphus nobilis}{
\lettrine[lines=2]{\zallmancaps \color{Red} A}{dest} triumphus nobilis
festumque celi curie,
quo rosa delectabilis
offertur regi glorie.
\par {\bfseries \color{Blue} P}etrus flos pulchritudinis,
et virtutum sacrarium,
nullum mortalis criminis
sensit unquam contagium.
\par {\bfseries \color{Red} R}oborare cum nititur
fidem verbi preconio,
pro fide tandem ceditur
hereticorum gladio.
\par {\bfseries \color{Blue} D}um sic in petra fidei
Petri tenet vestigia,
ad Petram Christum provehi
meretur cum victoria.
\par {\bfseries \color{Red} Q}uesumus actor omnium
in hoc Paschali gaudio,
per ipsius suffragium
crescat nostra devotio.
\par {\bfseries \color{Blue} G}loria tibi domine
qui surrexisti a mortuis,
et fortes in certamine
sertis ornas perpetuis. Amen.
}
\newline {\color{Red} In Nocturno. Ant.} De fumo lumen oritur et rose flos de sentibus, doctor et martyr nascitur Petrus de infidelibus, alleluya. {\color{Red} Ps.} \psalm{1} {\color{Red} Ant.} Predicatorum ordinis militans in acie, nunc coniunctus est agminis celestis militie, alleluya. {\color{Red} Ps.} \psalm{2} {\color{Red} Ant.} Mens fuit angelica lingua fructuosa, vita apostolica, mors quam preciosa, alleluya. {\color{Red} Ps.} \psalm{3} \V Gloria et honore. \R
\lettrine[lines=2]{\zallmancaps \color{Blue} O}{}\ \hypertarget{o-petre-fidus}{\label{o-petre-fidus}} Petre sidus aureum summi vas honoris, candorem servans niveum dono creatoris, * Decus habes virgineum in supernis choris, alleluya. \V Tu speculum mundicie, tu custos innocentie, tu sancti flos pudoris. Decus.
\begin{responsory}[predicator-fervidus]
{Predicator fervidus fidei zelator, circuit intrepidus heresum damnator, vitiorum rigidus semper expugnator, * Veritatis lucidus doctor et amator, alleluya}{Verbo, vita, predicat, signis crebris emicat, legis emulator.}
\end{responsory}%
\begin{responsory-final}[dum-sansonis]
{Dum Sansonis vulpes querit ab iniquis emitur, caput sacrum lictor ferit, iusti sanguis funditur, * Sic triumphi palmam gerit, dum pro fide moritur, * Alleluya.}{Stat invictus pugil fortis, constans profert hora mortis fidem pro qua patitur.}{Alleluya}
\end{responsory-final}%
\newline {\color{Red} In Laudibus. Ant. P}etrus novus incola celos laureatus, ascendit aureola triplici dotatus, alleluya. {\color{Red} Ant.} Turbe currunt languentium, signa coruscant varia, et in Petri preconium crebra crescunt prodigia, alleluya. {\color{Red} Ant.} Bolus digne suffocat guttur detrahentis, sed mox ipsum revocat votum penitentis alleluya. {\color{Red} Ant.} Motu, sensu corporis iuvenis privatur, tactu sacri pulveris vite restauratur, alleluya. {\color{Red} Ant.} Ad sancti Petri tumulum frequens lux descendit, in cuius laudis titulum lampades accendit, alleluya. {\color{Red} Ymnus.}
\hymn{exultet-claro-sidere}{Exultet claro sidere}{
\lettrine[lines=2]{\zallmancaps \color{Blue} E}{xultet} claro sidere
fulgens mater ecclesia,
Petrus martyr in ethere
nova profudit gaudia.
\par {\bfseries \color{Blue} P}auper, pudicus, humilis,
Christo se totum dedicat,
in lege dei docilis,
verbis, exemplis predicat.
\par {\bfseries \color{Red} T}riumphat per martyrium
Christi fortis in acie,
conservans semper lilium
virginali mundicie.
\par {\bfseries \color{Blue} L}ux celi, vite meritum,
cum signorum frequentia,
Petri commendant exitum,
et predicant magnalia.
\par {\bfseries \color{Red} Q}uesumus actor omnium
in hoc Paschali gaudio,
per ipsius suffragium
crescat nostra devotio.
\par {\bfseries \color{Blue} G}loria tibi domine
qui surrexisti a mortuis,
et fortes in certamine
sertis ornas perpetuis. Amen.
}
\newline {\color{Red} Ad Ben. Ant.} {\color{Blue} S}umma pollens Petrus mundicia, et prefulgens doctrine gratia, martyrii clarus victoria, trine fulget corone gloria, alleluya.
\newline \ul{Ad Vesperas. Ymnus, \Vbar , ut in Primis Vesperis.}
\newline {\color{Red} Ad Mag. Ant.} {\color{Blue} O} martyr egregie, doctor veritatis, puritatis vasculum, norma sanctitatis, tua per suffragia veniam peccatis, et vitam in gloria presta cum beatis, alleluya.
\newline \ul{Si hoc festum in vigilia Ascensionis evenerit, in Secundis Vesperis fiat tantum de eo memoria.}
{\ubsubsection{Sanctorum Philippi et Iacobi Apostolorum}{sanctorum-philippi-et-iacobi-apostolorum-breviarium}}
\newline {\color{Red} Ad Vesperas. Super psalmos. Ant.} Estote. {\color{Red} Cap.} Iam non estis. \R \hyperlink{tanto-tempore}{Tanto tempore.} {\color{Red} Ad Mag. Ant.} Filie Hierusalem. {\color{Red} Oratio.}
\lettrine[lines=2]{\zallmancaps \color{Blue} D}{eus} qui nos annua apostolorum tuorum Philippi et Iacobi sollemnitate letificas, presta quesumus, ut quorum gaudemus meritis instruamur exemplis. Per.
\newline \ul{Ad Matutinas. Invitat., et antiphone de communi Paschalis temporis.} \R \hyperlink{virtute-magna}{Virtute.} \R \hyperlink{isti-sunt-agni}{Isti sunt.}
\begin{responsory-final}[tanto-tempore]
{Tanto tempore vobiscum sum, et non cognovistis me, * Philippe qui videt me videt et patrem alleluya, non credis quia ego in patre et pater in me est, * Alleluya alleluya.}{Domine ostende nobis patrem et sufficit nobis.}{Alleluya}
\end{responsory-final}%
\newline {\color{Red} In Laudibus. Ant.} Domine ostende nobis patrem et sufficit nobis alleluya. {\color{Red} Ant.} Phylippe qui videt me videt et patrem meum alleluya. {\color{Red} Ant.} Tanto tempore vobiscum sum et non cognovistis me: Phylippe qui videt me videt et patrem meum alleluya. {\color{Red} Ant.} Si cognovissetis me et patrem meum utique cognovissetis, et amodo cognoscetis eum et vidistis eum, alleluya alleluya alleluya. {\color{Red} Ant.} Si diligeretis me gauderetis utique quia vado ad patrem, quia pater maior me est alleluya.
\newline {\color{Red} Ad Ben. Ant.} Si manseritis in me et verba mea in vobis manserint, quodcumque pecieritis fiet vobis, alleluya alleluya alleluya.
\newline {\color{Red} Ad Vesperas. Super psalmos. Ant.} Domine ostende.
\newline {\color{Red} Ad Mag. Ant.} Non turbetur cor vestrum neque formidet creditis in deum et in me credite, in domo patris mei, mansiones multe sunt alleluya alleluya.
{\ubsubsection{In Inventione Sancte Crucis}{in-inventione-sancte-crucis-breviarium}}
\newline {\color{Red} Ad Vesperas. Super psalmos. Ant.} Crucem sanctam subiit qui infernum confregit, accinctus est potentia surrexit die tercia alleluya. {\color{Red} Capitulum.}
\lettrine[lines=2]{\zallmancaps \color{Red} M}{ichi} absit gloriari nisi in cruce domini nostri Ihesu Christi per quem michi mundus crucifixus est et ego mundo. \R \hyperlink{per-tuam-crucem}{Per tuam crucem.} {\color{Red} III. Ymnus.}
\hymn{salve-crux-sancta}{Salve crux sancta}{
\lettrine[lines=2]{\zallmancaps \color{Blue} S}{alve} crux sancta, salve mundi gloria,
vera spes nostra, vera ferens gaudia,
signum salutis salus in periculis,
vitale lignum, vitam ferens omnium.
\par {\bfseries \color{Red} T}e adorandam te crucem vivificam,
per te redempti, dulce decus seculi,
semper laudamus semper tibi canimus,
per lignum servi per te sumus liberi.
\par {\bfseries \color{Blue} S}it deo patri laus in cruce filii
sit coequali laus sancto spiritui
civibus summis, gaudium sit angelis,
honor sit mundo crucis hec inventio. Amen. {\color{Red} Vel} crucis exaltatio. Amen.
}
\newline \V Hoc signum crucis erit in celo alleluya. \R Cum dominus ad iudicandum venerit alleluya.
\newline {\color{Red} Ad Mag. Ant.} O crux benedicta que sola fuisti digna portare regem celorum et dominum alleluya. {\color{Red} Oratio.}
\lettrine[lines=2]{\zallmancaps \color{Red} D}{eus} qui in preclara salutifere crucis inventione passionis tue miracula suscitasti, concede ut vitalis ligni precio eterne vite suffragia consequamur. Qui vivis.
\newline {\color{Red} Deinde Sancti Alexandri sociorumque eius. Memoria. Oratio.}
\lettrine[lines=2]{\zallmancaps \color{Blue} P}{resta} quesumus omnipotens deus, ut qui sanctorum tuorum Alexandri, Eventii et Theodoli natalicia colimus a cunctis malis imminentibus eorum intercessionibus liberemur. Per.
\newline \ul{Non fiat alia memoria de cruce, similiter nec in Laudibus, nec in Secundis Vesperis.}
\newline {\color{Red} Ad Matutinas. Invitat.}
\begin{invitatory}[alleluya-salve-sancta-invitatorium]
{Alleluya, salve sancta crux salve, * Alleluya.}
\end{invitatory}
\newline {\color{Red} Ymnus.} \hyperlink{salve-crux-sancta}{Salve sancta crux.} %sic
{\color{Red} In Nocturno. Super psalmos. Ant.} Adoramus te Christe et benedicimus tibi, quia per crucem tuam redemisti mundum alleluya. {\color{Red} Ps.} \psalm{8} \psalm{20} \psalm{23} \V Hoc signum.
\newline \ul{Crisostonus in Omelia. XIII. de cruce. \grecross } {\color{Red} Lectio Prima.}
% https://mlat.uzh.ch/browser/cps_2.JonAur.DeCuIm#:4;2
\lettrine[lines=2]{\zallmancaps \color{Red} S}{i} nosse desideras karissime crucis virtutem: audi quanta dicere possum ad ipsius laudem . Crux est spes Christianorum, crux resurrectio mortuorum, dux cecorum, baculus claudorum. Crux est consolatio pauperum, refrenatio divitum, destructio superborum. Crux est male viventium pena, crux adversus demonas victoria. Crux est pedagogus adolescentium, gubernator navigantium, portus periclitantium. \R
\lettrine[lines=2]{\zallmancaps \color{Blue} D}{ulce} \hypertarget{dulce-lignum-dulces}{\label{dulce-lignum-dulces}} lignum dulces clavos dulce pondus sustinuit, * Que digna fuit portare precium huius secula alleluya. \V Hoc signum crucis erit in celo cum dominus ad iudicandum venerit. Que.%typo: evenerit
\newline {\color{Red} Lectio Secunda.}
\lettrine[lines=2]{\zallmancaps \color{Red} C}{rux} est murus obsessorum, spes desperatorum, requies tribulatorum. Crux est pater orphanorum, defensor viduarum, consiliarius iustorum. Crux est custos parvulorum, caput virorum, finis senum. Crux est lumen in tenebris sedentium, magnificentia regum, scutumque perpetuum. Crux est sapientia stultorum, libertas servorum, phylosophya imperatorum: lex impiorum. Crux est preconizatio prophetarum, annunciatio apostolorum: gloriatio martyrum, abstinentia monachorum: castitas virginum.
\begin{responsory}[ecce-crucem-domini]
{Ecce crucem domini fugite partes adverse vicit leo de tribu Iuda, * Radix David alleluya.}{Crux benedicta in qua triumphavit rex angelorum.}
\end{responsory}%
\newline {\color{Red} Lectio Tertia.}
\lettrine[lines=2]{\zallmancaps \color{Blue} C}{rux} est gaudium sacerdotum, ecclesie fundamentum, orbis terre tutela, ydolorum repulsio, templorum destructio. Crux est scandalum Iudeorum: perditio impiorum. Crux est infirmorum virtus, egrotantium medicus. Crux est panis esurientium, fons sitientium: protectio nudorum. Crux Christi: clavis est paradisi.
\begin{responsory-final}[per-tuam-crucem]
{Per tuam crucem salva nos Christe redemptor, qui mortem nostram moriendo destruxisti, * Et vitam resurgendo reparasti, * Alleluya.}{Miserere nostri Ihesu benigne, qui passus es clementer pro nobis.}{Alleluya}
\end{responsory-final}%
\newline \V Dicite in nationibus alleluya.
\newline {\color{Red} In Laudibus. Ant.} Helena Constantini mater Hierosolimam petiit alleluya. {\color{Red} Ant.} Tunc precepit eos omnes igne cremari, at illi timentes tradiderunt Iudam alleluya. {\color{Red} Ant.} Cumque ascendisset Iudas de lacu perrexit ad locum ubi iacebat sancta crux alleluya. {\color{Red} Ant.} Orabat Iudas deus deus meus ostende michi lignum sancte crucis alleluya. {\color{Red} Ant.} Cum orasset Iudas commotus est locus ille in quo sancta crux iacebat alleluya. {\color{Red} Cap.} Michi absit. {\color{Red} Ymnus.}
\hymn{originale-crimen}{Originale crimen necans}{
\lettrine[lines=2]{\zallmancaps \color{Red} O}{riginale} crimen necans in cruce,
nos a privatis Christe munda maculis
humanitatem miseratus fragilem,
per crucem sanctam lapsis dona veniam.
\par {\bfseries \color{Blue} P}rotege, salva, benedic, sanctifica,
populum cunctum crucis per signaculum,
morbos averte corporis et anime,
hoc contra signum nullum stet periculum.
\par {\bfseries \color{Red} S}it laus in cruce filii
sit coequali laus sancto spiritui
civibus summis, gaudium sit angelis,
honor sit mundo crucis hec inventio. Amen. {\color{Red} Vel} crucis exaltatio. Amen.
}
\newline \V Tuam crucem adoremus domine alleluya. \R Tuam gloriosam recolimus passionem alleluya.
\newline {\color{Red} Ad Ben. Ant.} Helena sancta dixit ad Iudam, comple desiderium meum et vive super terram, ut ostendas michi qui dicitur Calvarie locus, ubi absconditum est preciosum lignum dominicum alleluya.
\newline {\color{Red} Ad Terciam. Cap.} Michi absit.
\begin{responsory-breve}[hoc-signum-breve]
{Hoc signum crucis erit in celo, * Alleluya alleluya.}{Cum dominus ad iudicandum venerit.}
\end{responsory-breve}%
\V Per signum crucis.
\newline {\color{Red} Ad Sextam. Capitulum.}
\lettrine[lines=2]{\zallmancaps \color{Blue} V}{erbum} crucis pereuntibus stulticia est, iis autem qui salvi fiunt, idest nobis: virtus dei est.
\begin{responsory-breve}[per-signum-crucis]
{Per signum crucis de inimicis nostris, * Alleluya alleluya.}{Libera nos deus noster.}
\end{responsory-breve}%
\V Adoramus te Christe.
\newline {\color{Red} Ad Nonam. Cap.}
\lettrine[lines=2]{\zallmancaps \color{Red} N}{os} autem predicamus Christum crucifixum, Iudeis quidem scandalum gentibus autem stulticiam: ipsis autem vocatis Iudeis atque Grecis, Christum dei virtutem et dei sapientiam.
\begin{responsory-breve}[adoramus-te-christe-breve]
{Adoramus te Christe et benedicimus tibi, * Alleluya alleluya.}{Quia per crucem tuam redemisti mundum.}
\end{responsory-breve}%
\V Tuam crucem.
\newline {\color{Red} Ad Vesperas. Super psalmos. Ant.} Helena. {\color{Red} Ps.} \psalm{109} \psalm{110} \psalm{111} \psalm{112} \psalm{116} \ul{Cap., Ymnus, \Vbar , ut in Primis Vesperis.}
\newline {\color{Red} Ad Mag. Ant.} O crux splendidior cunctis astris, mundo celebris, hominibus multum amabilis, sanctior universis, que sola fuisti digna portare talentum mundi, dulce lignum, dulces clavos, dulcia ferens pondera, salva presentem catervam in tuis hodie laudibus congregatam, alleluya, alleluya, alleluya. \ul{Post orationem de cruce fiat Memoria de Corona.}
\newline {\color{Red} Ant.} Gaude felix mater ecclesia, assunt tibi nova sollemnia nam corona quondam in gloria, nunc per orbem refulget gloria, alleluya. \V Tuam coronam adoramus domine alleluya. \R Tuum gloriosum recolimus triumphum alleluya. {\color{Red} Oratio.}
\lettrine[lines=2]{\zallmancaps \color{Blue} P}{resta} quesumus omnipotens deus, ut qui in memoriam passionis domini nostri Ihesu Christi coronam eius spineam veneramur in terris, ab ipso gloria et honore coronari mereamur in celis. Qui tecum.
\newline \ul{Si festum Sancte Crucis in die Ascensionis evenerit in sexta feria celebretur, et festum Corone in Sabbatum transferatur.}
{\ubsubsection{In Festo Corone}{in-festo-corone-breviarium}}
\newline {\color{Red} Ad Mat. Invitat.}
\begin{invitatory}[assunt-dominici-serti-invitatorium]
{Assunt dominici serti sollemnia, * Laude multiplici plaudat ecclesia, alleluya.}
\end{invitatory}
\newline {\color{Red} Ymnus.}
\hymn{eterne-regi-glorie}{Eterne regi glorie}{
\lettrine[lines=2]{\zallmancaps \color{Red} E}{terne} regi glorie
devota laudum cantica,
fideles solvant hodie
pro corona dominica.
\par {\bfseries \color{Blue} C}oronat regem omnium
corona contumelie,
cuius nobis obprobrium
coronam confert glorie.
\par {\bfseries \color{Red} D}e spinarum aculeis
Christi corona plectitur,
qua ministris Tartareis,
mundi potestas tollitur.
\par {\bfseries \color{Blue} C}orona Christi capitis
sacro perfusa sanguine:
penis solutis debitis
reos purgat a crimine.
\par {\bfseries \color{Red} Q}uesumus actor omnium, {\color{Red} Vel} Tu esto nostrum gaudium.
\par {\bfseries \color{Blue} L}aus Christo regi glorie
pro corone virtutibus
qua nos reformans gratie
coronat in celestibus. Amen.
}
\newline {\color{Red} Super psalmos. Ant.} Christum sub serto spineo deridet plebs perfidie, cuius cruore roseo sertum confertur glorie alleluya. {\color{Red} Psalmi.} \psalm{1} \psalm{2} \psalm{3} \V Tuam coronam adoramus domine.
\newline \ul{Ex quidam sermone de Corona Domini.} \grecross \ {\color{Red} Lectio Prima.}
\lettrine[lines=2]{\zallmancaps \color{Red} N}{on} miretur ortodoxorum quispiam, si iocunda corone dominice sollemnitas diem hunc expendat in laudibus redemptoris. Quia et si hec corona capiti salvatoris ad penam et ludibrium in die parascheves fuerit applicata, tamen quia dies illa non est gaudii sed meroris, quando membra compatiuntur capiti: differtur interim hec gratulabunda festivitas, in qua recolligimus salutis nostre messem de spinarum semine propagatam. \R
\lettrine[lines=2]{\zallmancaps \color{Blue} S}{pina} \hypertarget{spina-carens-flos}{\label{spina-carens-flos}} carens flos spina pungitur per quam culpe spina confringitur spina mortis, * Spinis retunditur dum vita moritur alleluya. \V Per hoc ludibrium hostis deluditur, mortis dominium per mortem tollitur. Spinis.
\newline {\color{Red} Lectio Secunda.}
\lettrine[lines=2]{\zallmancaps \color{Red} M}{iro} etenim modo deiectio capitis obtinuit veniam et gratiam corpori: et percusso vertice, solidata sunt membra. Sinagoga siquidem mater Christi secundum carnem, novercam se exhibens affectu crudelitatis et effectu: nostrum Salomonem corona spinea coronavit. Pudeat ergo sectari gloriam membra, quibus caput suum tam inglorium exhibetur peccatorum nostrorum spinis circundatum: pudeat sub spinato capite, membrum fieri delicatum.
\begin{responsory}[coronat-regem]
{Coronat regem omnium Iudea serto spineo, stat inter spinas lilium vernans cruore roseo, * Spinarum culpe nescium spine punctum aculeo alleluya.}{Sub decore fulget purpureo corpus nitens candore niveo.}
\end{responsory}%
\newline {\color{Red} Lectio Tertia.}
\lettrine[lines=2]{\zallmancaps \color{Blue} S}{cripsit} namque Iohannes evangelista. Exivit Ihesus portans spineam coronam et purpureum vestimentum: ut quasi pugil noster in armis et vestimentis rubeis appareret. Exeamus ergo ad eum egredientem extra castra improperium eius portantes, erumnam eius nostram reputantes: et dicamus cum psalmista. Confixus sum in erumna mea dum configitur spina.
\begin{responsory}[felix-spina-cuius]
{Felix spina cuius aculei guttis rubent roris sanguinei, * Vires frangunt regis Tartarei seras pandunt regni siderei * Alleluya.}{O spinarum immensa gloria que tot nobis prestant remedia.}{Alleluya}
\end{responsory}%
\newline \V Plectentes coronam de spinis, alleluya. \R Posuerunt super caput domini, alleluya.
\newline {\color{Red} In Laudibus. Ant.} Adest dies leticie quo diadema spinem, commendatur memorie Christi cruore roseum, alleluya. {\color{Red} Ant.} Summum regem glorie spinis coronatum, ridet plebs perfidie morti condemnatum, alleluya. {\color{Red} Ant.} O quam felix punctio quam beata spina, de qua fluit unctio mundi medicina alleluya. {\color{Red} Ant.} Pungens spina vulnerat Christum patientem, et a morte liberat populum credentem alleluya. {\color{Red} Ant.} Spine rubent sanguine Christum cruentates, mundum lavant crimine celum reserantes alleluya. {\color{Red} Capitulum.}
\lettrine[lines=2]{\zallmancaps \color{Red} E}{gredimini} et videte filie Syon regem Salomonem, in diademate quo coronavit eum mater sua. {\color{Red} Ymnus.}
\hymn{lauda-fidelis-contio}{Lauda fidelis contio}{
\lettrine[lines=2]{\zallmancaps \color{Blue} L}{auda} fidelis contio,
spine tropheum inclitum
per quam perit perditio,
viteque datur meritum.
\par {\bfseries \color{Red} N}os a puncturis liberat
eterni patris filius,
dum spinis pungi tolerat
spinarum culpe nescius.
\par {\bfseries \color{Blue} D}um spinarum aculeum
Christus pro nobis pertulit,
per diadema spineum
vite coronam contulit.
\par {\bfseries \color{Red} P}laudat turba fidelium
quod per spine ludibrium
purgat creator omnium
spineti nostri vitium.
\par Quesumus actor omnium, {\color{Red} Vel} Tu esto nostrum gaudium.
\par {\bfseries \color{Blue} L}aus Christo regi glorie
pro corone virtutibus
qua nos reformans gratie
coronat in celestibus. Amen.
}
\newline \V Eris corona glorie in manu domini alleluya. \R Et diadema regni in manu dei tui alleluya.
\newline {\color{Red} Ad Ben. Ant.} Ave spina pene remedium servi decus regis obprobrium tua plaga, dolor, ludibrium, vite nobis mercantur premium alleluya.
\newline {\color{Red} Ad Terciam. Cap.} Egredimini.
\begin{responsory-breve}[tuam-coronam-adoramus]
{Tuam coronam adoramus domine, * Alleluya alleluya.}{Tuum gloriosum recolimus triumphum.}
\end{responsory-breve}%
\V Gloria et honore coronasti eum domine alleluya.
\newline {\color{Red} Ad Sextam. Cap.}
\lettrine[lines=2]{\zallmancaps \color{Red} V}{idi} et ecce equus albus et qui sedebat super eum habebat arcum, et data est ei corona, et exivit vincens ut vinceret.
\begin{responsory-breve}%[gloria-et-honore-breve]
{Gloria et honore coronasti eum domine, * Alleluya alleluya.}{Et constituisti eum super opera manuum tuarum.}
\end{responsory-breve}%
\V Posuisti domine super caput eius alleluya.
\newline {\color{Red} Ad Nonam. Cap.}
\lettrine[lines=2]{\zallmancaps \color{Blue} I}{n} die illa erit dominus exercituum corona glorie et sertum exultationis residuo populi sui.
\begin{responsory-breve}%[posuisti-domine-breve]
{Posuisti domine super caput eius, * Alleluya alleluya.}{Coronam de lapide precioso.}
\end{responsory-breve}%
\V Eris corona glorie in manu domini alleluya.
\newline {\color{Red} Ad Vesperas. Super psalmos. Ant.} Adest. {\color{Red} Ps.} \psalm{109} \psalm{110} \psalm{111} \psalm{112} \psalm{116} {\color{Red} Cap.} Egredimini. {\color{Red} Ymnus.} \hyperlink{eterne-regi-glorie}{Eterne regi glorie.} \V Tuam coronam.
\newline {\color{Red} Ad Mag. Ant.} O decus ecclesie gloriosa spina sertum regis glorie, mundi medicina, presentis angustie dulcor et resina, te laudantes hodie serves a ruina alleluya.
{\ubsubsection{Sancti Iohannis Ante Portam Latinam}{sancti-iohannis-ante-portam-latinam-breviarium}}
\newline {\color{Red} Ad Vesperas.} \ul{Antiphona super psalmos, Cap., \Rbar , Ymnus, \Vbar , de communi evangelistarum.}
\newline {\color{Red} Ad Mag. Ant.} In ferventis olei dolium missus Iohannes apostolus divina se protegente gratia illesus exivit alleluya. {\color{Red} Oratio.}
\lettrine[lines=2]{\zallmancaps \color{Red} D}{eus} qui conspicis quia nos undique mala nostra perturbant, presta quesumus, ut beati Iohannis apostoli tui et evangeliste intercessio gloriosa nos protegat. Per.
\newline {\color{Red} Ad Matutinas.} \ul{Invitatorium, antiphona, psalmi, responsoria de communi evangelistarum Paschalis tempore.}
\newline {\color{Red} Lectio Prima. Secundum Iohannem.}
\lettrine[lines=2]{\zallmancaps \color{Blue} I}{n} illo tempore: Dixit Ihesus Petro. Sequere me. Et reliqua.
\newline {\color{Red} Omelia Venerabilis Bede Presbiteri.} \grecross
\lettrine[lines=2]{\zallmancaps \color{Red} N}{on} est putandum quia discipulus ille non sit mortuus in carne, quia nec dominus hoc de ipso futurum predixit. Sed ita potius intelligendum: quod ceteris discipulis Christi per passionem consummatis, ipse in pace ecclesie adventum vocationis expectaverit. Et hoc est quod ait dominus. Sic eum volo manere donec veniam. Non quia non et ipse multos antea labores pro domino pressurasque malorum toleraverit, sed quia ultimum in pace senium finierit, utpote ecclesiis Christi per Asiam quam regebat iam longe lateque fundatis.
\newline {\color{Red} Lectio Secunda.}
\lettrine[lines=2]{\zallmancaps \color{Blue} N}{am} et in actibus apostolorum cum ceteris apostolis flagellatus invenitur: quando ibant gaudentes a conspectu concilii, quoniam digni habiti sunt, pro nomine Ihesu contumeliam pati. Et a Domitiano Cesare in ferventis olei dolium missus in ecclesiastica narratur hystoria, ex quo tamen divina se protegente gratia tam intactus exivit, quam fuerat a corruptione concupiscentie carnalis extraneus.
\newline {\color{Red} Lectio Tertia.}
\lettrine[lines=2]{\zallmancaps \color{Red} N}{ec} multo post ab eodem principe propter insuperabilem evangelizandi constantiam in Pathmos insulam est exilio relegatus, ubi humano licet destitutus solatio, divine tamen visionis meruit et allocutionis crebra consolatione revelari.
\newline \ul{Ad Laudes et ad horas, Capitula, Ymni, \Vbar , Responsoria de communi evangelistarum.}
{\ubsubsection{Sanctorum Gordiani et Epimachi Martyrum}{sanctorum-gordiani-et-epimachi-martyrum-breviarium}}
\newline {\color{Red} Oratio.}
\lettrine[lines=2]{\zallmancaps \color{Blue} D}{a} quesumus omnipotens deus, ut qui beatorum martyrum tuorum Gordiani atque Epymachi sollemnia colimus, eorum apud te intercessionibus adiuvemur. Per.
{\ubsubsection{Sanctorum Nerei Achillei et Panchratii Martyrum}{sanctorum-nerei-achillei-et-panchratii-martyrum-breviarium}}
\newline {\color{Red} Oratio.}
\lettrine[lines=2]{\zallmancaps \color{Red} S}{emper} nos domine martyrum tuorum Nerei, Achillei atque Panchratii foveat quesumus beata sollemnitas, et tuo dignos reddat obsequio. Per.
{\ubsubsection{Sancte Potentiane Virginis}{sancte-potentiane-virginis-breviarium}}
\newline {\color{Red} Oratio.}
\lettrine[lines=2]{\zallmancaps \color{Blue} E}{xaudi} nos deus salutaris noster, ut sicut de beate Potentiane virginis tue festivitate gaudemus, ita pie devotionis erudiamur affectu. Per.
{\ubsubsection{In Translatione Beati Dominici}{in-translatione-beati-dominici-breviarium}}
\newline {\color{Red} Ad Vesperas. Super psalmos. Ant.} Adest dies leticie quo beatum Dominicum exaltavit rex glorie per odorem mirificum, alleluya. {\color{Red} Ps.} \psalm{112} \ul{cum ceteris, Cap., Ymnus, \Vbar , sicut infra in festo ipsius.} \R \hyperlink{in-odoris-mira}{In odoris.} {\color{Red} III. Ad Mag. Ant.} Magnificat Dominicum Christus novo miraculo, dum odorem mirificum cum eius profert ex tumulo alleluya. {\color{Red} Oratio.} \hyperlink{deus-qui-ecclesiam}{Deus qui ecclesiam.}
\newline {\color{Red} Ad Mat. Invit.}
\begin{invitatory}[assunt-dominici-leta-invitatorium]
{Assunt Dominici leta sollemnia, * Laude multiplici plaudat ecclesia alleluya.}
\end{invitatory}
\newline {\color{Red} Ymnus.} \hyperlink{novus-athleta-domini}{Novus athleta.}
\newline {\color{Red} In Nocturno. Ant.} Preco novus et celicus missus in fine seculi, pauper fulsit Dominicus forma previsus catuli alleluya. {\color{Red} Ps.} \psalm{1} {\color{Red} Ant.} Florem pudicicie servans illibatum, attigit eximie vite celibatum alleluya. {\color{Red} Ps.} \psalm{2} {\color{Red} Ant.} Documentis arcium eruditus satis transiit ad studium summe veritatis alleluya. {\color{Red} Ps.} \psalm{3} \V Amavit. \ul{Lectiones de communi unius confessoris doctoris trium lectionum. Et fiant de qualibet lectiones tres, si festum post Trinitatem evenerit.} \R
\lettrine[lines=2]{\zallmancaps \color{Red} F}{ulget} \hypertarget{fulget-decus-ecclesie}{\label{fulget-decus-ecclesie}} decus ecclesie novo clarum prodigio lucerna splendet hodie diu latens sub modio, * Quam odore rex glorie revelavit eximio, alleluya. \V O mira suavitas inaudite fragrantie per quam patris dignitas declaratur ecclesia. Quam.
\begin{responsory}[virgo-pugil]
{Virgo pugil Christi virtutum forma fuisti verbo fulsisti, dum transferri meruisti, * Fragrat odor dulces cantant celi agmina laudes alleluya.}{Regna tuis pande celi, doctor venerande.}
\end{responsory}%
\begin{responsory-final}[in-odoris-mira]
{In odoris mira fragrantia populorum currit frequentia, * Ubi passim Christi clementia * Sana reddit membra languentia * Alleluya.}{Sidus quondam oppressum nebula veri solis pandunt miracula.}{Alleluya}\ul{Vel} Sana. \ul{post Trinitatem}
\end{responsory-final}%
\newline \V Ora pro nobis.
\newline {\color{Red} In Laudibus. Ant.} Gemma sub terra latuit despecto iacens loculo, cuius virtus apparuit multiplici miraculo, alleluya. {\color{Red} Ant.} Corpus sacrum quod fuerat apotheca carismatum, universam exuperat fragrantiam aromatum alleluya. {\color{Red} Ant.} Glebam sacri corporis odor propalavit, quam neglecti temporis cursus occultavit alleluya. {\color{Red} Ant.} Tumba mira fragrantia profundens vim odoris corda reddit ardentia in laudem salvatoris alleluya. {\color{Red} Ant.} Superans fragrantiam odor pigmentorum cumulum et gloriam monstrat premiorum alleluya. \ul{Cap., Ymnus, \Vbar , sicut infra in festo ipsius.}
\newline {\color{Red} Ad Ben. Ant.} O quantus stupor populi quanta fratrum leticia dum loco sacri tumuli mira prodit fragrantia currunt senes et parvuli ad sancti beneficia alleluya.
\newline \ul{Ad horas capitula sicut in alio festo. Ad Vesperas. Cap., Ymnus, \Vbar , ut in Primis Vesperis.}
\newline {\color{Red} Ad Mag. Ant.} O speculum mundicie carnis carens spurcicia tue colentes hodie translationis gaudia transfer ad regnum glorie post huius vite stadia alleluya.%typo: transfert
\newline \ul{Si hoc festum in vigilia Ascensionis evenerit in Secundis Vesperis, antiphona, psalmi, et omnia de Ascensione erunt, et memoria fiet de beato Dominico. Si autem in crastino Ascensionis fuerit celebrandum: in die Ascensionis Vespere erunt de Ascensione, et memoria fiet de beato Dominico.}
\newline \ul{Quando autem hoc festum post Trinitatem occurrerit celebrandum: ad Vesperas fiat ut supra notatum est excepto quod alleluya post antiphonas et responsorium posita: non dicantur; Ad Matutinas, Invitat., Ymnus, antiphone super psalmos nocturnorum: sicut in alio festo. Responsorium primum.} \hyperlink{mundum-vocans}{Mundum.}
%pp125
\ul{Responsorium secundum.} \hyperlink{datum-mundo-pro}{Datum mundo.} \ul{Responsorium tercium.} \hyperlink{fulget-decus-ecclesie}{Fulget.} \ul{Responsorium quartum.} \hyperlink{paupertatis-ascendens}{Paupertatis.} \ul{Responsorium quintum.} \hyperlink{panis-oblatus}{Panis.} \ul{Responsorium sextum.} \hyperlink{virgo-pugil}{Virgo pugil.} \ul{Responsorium septimum.} \hyperlink{felix-vitis-de}{Felix.} \ul{Responsorium octavum.} \hyperlink{ascendenti-de-valle}{Ascendenti.} \ul{Responsorium nonum.} \hyperlink{in-odoris-mira}{In odoris.} \ul{Ad Laudes et alias horas, ut supra notatum est. Alleluya tamen post responsoria et antiphonas posita non dicantur.}
{\ubsubsection{Sancti Urbani Pape et Martiris}{sancti-urbani-pape-et-martiris-breviarium}}
\newline {\color{Red} Oratio.}
\lettrine[lines=2]{\zallmancaps \color{Blue} D}{a} quesumus omnipotens deus, ut beati Urbani martyris tui atque pontificis sollemnia colimus, eius apud te intercessionibus adiuvemur. Per.
{\ubsubsection{Sancte Petronille, Virginis}{sancte-petronille-virginis-breviarium}}
\newline {\color{Red} Oratio.}
\lettrine[lines=2]{\zallmancaps \color{Red} E}{xaudi} nos deus salutaris noster, ut sicut de beate Petronille virginis tue festivitate gaudemus, ita pie devotionis erudiamur affectu. Per.
{\ubsubsection{Sanctorum Marcellini et Petri Martirum}{sanctorum-marcellini-et-petri-martirum-breviarium}}
\newline {\color{Red} Oratio.}
\lettrine[lines=2]{\zallmancaps \color{Blue} D}{eus} qui nos annua beatorum martyrum tuorum Marcellini et Petri sollemnitate letificas, presta quesumus, ut quorum gaudemus meritis provocemur exemplis. Per.
{\ubsubsection{Sancti Medardi Episcopi et Confessoris}{sancti-medardi-episcopi-et-confessoris-breviarium}}
\newline {\color{Red} Oratio.}
\lettrine[lines=2]{\zallmancaps \color{Red} E}{xaudi} domine preces nostras, quas in sancti Medardi confessoris tui atque pontificis sollemnitate deferimus, ut qui tibi digne meruit famulari, eius intercedentibus meritis ab omnibus nos absolvas peccatis. Per.
{\ubsubsection{Sanctorum Primi et Feliciani Martirum}{sanctorum-primi-et-feliciani-martirum-breviarium}}
\newline {\color{Red} Oratio.}
\lettrine[lines=2]{\zallmancaps \color{Blue} F}{ac} nos quesumus domine sanctorum tuorum Primi et Feliciani semper festa sectari, et eorum suffragiis, protectionis tue dona sentire. Per.
{\ubsubsection{Sancti Barnabe Apostoli}{sancti-barnabe-apostoli-breviarium}}
\newline {\color{Red} Oratio.}
\lettrine[lines=2]{\zallmancaps \color{Red} Q}{uesumus} omnipotens deus ut beatus Barnabas apostolus tuum pro nobis imploret auxilium, ut a nostris reatibus absoluti a cunctis eciam periculis exuamur. Per.
{\ubsubsection{Sanctorum Basilidis, Cyrini Naboris et Nazarii, Martyrum}{sanctorum-basilidis-cyrini-naboris-et-nazarii-martyrum-breviarium}}
\newline {\color{Red} Oratio.}
\lettrine[lines=2]{\zallmancaps \color{Blue} S}{anctorum} tuorum Basilidis, Cyrini, Naboris et Nazarii quesumus domine natalicia nobis votiva resplendeant, et quod illis contulit excellentiam sempiternam, fructibus nostre devotionis accrescat. Per.
{\ubsubsection{Sanctorum Viti et Modesti Martirum}{sanctorum-viti-et-modesti-martirum-breviarium}}
\newline {\color{Red} Oratio.}
\lettrine[lines=2]{\zallmancaps \color{Red} D}{eus} qui nos annua sanctorum tuorum Viti et Modesti sollemnitate letificas, concede propicius, ut quorum gaudemus meritis accendamur exemplis. Per.
{\ubsubsection{Sanctorum Cyrici et Iulite Martirum}{sanctorum-cyrici-et-iulite-martirum-breviarium}}
\newline {\color{Red} Oratio.}
\lettrine[lines=2]{\zallmancaps \color{Blue} D}{eus} qui nos concedis sanctorum martyrum tuorum Cyrici et Iulite natalicia colere, da nobis in eterna beatitudine de eorum societate gaudere. Per.
{\ubsubsection{Sanctorum Marci et Marcelliani Martirum}{sanctorum-marci-et-marcelliani-martirum-breviarium}}
\newline {\color{Red} Oratio.}
\lettrine[lines=2]{\zallmancaps \color{Red} P}{resta} quesumus omnipotens deus, ut qui sanctorum Marci et Marcelliani natalicia colimus, a cunctis malis imminentibus eorum intercessionibus liberemur. Per.
{\ubsubsection{Sanctorum Gervasii et Prothasii Martirum}{sanctorum-gervasii-et-prothasii-martirum-breviarium}}
\newline {\color{Red} Oratio.}
\lettrine[lines=2]{\zallmancaps \color{Blue} D}{eus} qui nos annua sanctorum tuorum Gervasii et Prothasii sollemnitate letificas, concede propicius, ut quorum gaudemus meritis accendamur exemplis. Per.
{\ubsubsection{In Vigilia Sancti Iohannis Baptiste}{in-vigilia-sancti-iohannis-baptiste-breviarium}}
\newline {\color{Red} Lectio I. vel Septima si Dominica fuerit. Secundum Lucam.}
\lettrine[lines=2]{\zallmancaps \color{Red} F}{uit} in diebus Herodis regis Iudee sacerdos quidam nomine Zacharias de vice Abia, et uxor illi de filiabus Aaron, et nomen eius Elisabeth. Et reliqua.
\newline {\color{Red} Omelia Venerabilis Bede Presbiteri. \grecross }
\lettrine[lines=2]{\zallmancaps \color{Blue} E}{x} iustis parentibus Iohannes est genitus, ut eo confidentius iusticie precepta populis daret, quo hec ipse non quasi novitia didicisset, sed velut hereditario iure a progenitoribus accepta servaret. De sacerdotali prosapia ortus est, ut eo potentius mutationem sacerdotii preconizaret: quo ipsum ad sacerdotale genus pertinere claresceret.
\newline {\color{Red} Lectio II. vel VIII.}
\lettrine[lines=2]{\zallmancaps \color{Red} R}{edemptor} etenim noster in carne apparens sicut rex nobis fieri dignatus est regnum nobis celeste tribuendo: ita eciam pontifex factus est, semetipsum pro nobis offerendo hostiam deo in odorem suavitatis. Unde scriptum est. Tu es sacerdos in eternum secundum ordinem Melchisedech.
\newline {\color{Red} Lectio III. vel IX.}
\lettrine[lines=2]{\zallmancaps \color{Blue} M}{elchisedech} quippe ut legimus sacerdos Dei summi, longe legalis sacerdotii tempora precessit, offerens tunc panem et vinum. Et ideo redemptor sacerdos esse dicitur secundum ordinem Melchisedech, quia ablatis victimis legalibus: idem sacrificii genus in mysterium sui corporis et sanguinis in novo testamento offerendum instituit.
\newline {\color{Red} Ad Vesperas. Super psalmos. Ant.} Ingresso Zacharia templum domini, apparuit ei Gabriel angelus stans a dextris altaris incensi. {\color{Red} Capitulum.}
\lettrine[lines=2]{\zallmancaps \color{Red} P}{riusquam} te formarem in utero novi te, et antequam exires de vulva sanctificavi te, et prophetam in gentibus dedi te. \R \hyperlink{inter-natos-mulierum}{Inter natos.} {\color{Red} IX. Ymnus.}
\hymn{ut-queant-laxis}{Ut queant laxis}{
\lettrine[lines=2]{\zallmancaps \color{Blue} U}{t} queant laxis resonare fibris
mira gestorum famuli tuorum,
solve polluti labii reatum
sancte Iohannes.
\par {\bfseries \color{Red} N}untius celso veniens Olimpo,
te patri magnum fore nasciturum
nomen, et vite seriem gerende
ordine promit.
\par {\bfseries \color{Blue} I}lle promissi dubius superni,
perdidit prompte modulos loquele,
sed reformasti genitus perempte
organa vocis.
\par {\bfseries \color{Red} V}entris obstruso positus cubili:
senseras regem thalamo manentem,
hinc parens nati meritis uterque
abdita pandit.
\par {\bfseries \color{Blue} L}audibus cives celebrant superni,
te deus simplex pariterque trine,
supplices et nos veniam precamur,
parce redemptis. Amen.
}
\newline \V Fuit homo missus a deo. \R Cui nomen erat Iohannes.
\newline {\color{Red} Ad Mag. Ant.} Pro eo quod non credidisti verbis meis eris tacens et non poteris loqui usque in diem nativitatis eius. {\color{Red} Oratio.}
\lettrine[lines=2]{\zallmancaps \color{Red} P}{resta} quesumus omnipotens deus ut familia tua per viam salutis incedat, et beati Iohannis precursoris hortamenta sectando, ad eum quem predixit secura perveniat. Dominum nostrum Ihesum Christum filium tuum. Qui tecum.
{\ubsubsection{In Die Sancti Iohannis Baptiste}{in-die-sancti-iohannis-baptiste-breviarium}}
\newline {\color{Red} In die. Ad Mat. Invitat.}
\begin{invitatory}[regem-precursoris-dominum-invitatorium]
{Regem precursoris dominum, * Venite adoremus.}
\end{invitatory}
\newline {\color{Red} Ymnus.}
\hymn{antra-deserti-teneris}{Antra deserti teneris}{
\lettrine[lines=2]{\zallmancaps \color{Blue} A}{ntra} deserti teneris sub annis
civium turmas fugiens petisti,
ne levi saltem maculare vitam
famine posses.
\par {\bfseries \color{Red} P}rebuit hyrtum tegimen camelus
artubus sacris, strophium bidentes,
cui latex haustum, sociata pastum
mella locustis.
\par {\bfseries \color{Blue} C}eteri tantum cecinere vatum
corde presago iubar affuturum,
tu quidem mundi scelus auferentem
indice prodis.
\par {\bfseries \color{Red} N}on fuit vasti spatium per orbis,
sanctior quisquam genitus Iohanne,
qui nefas secli meruit lavantem
tingere limphis.
\par {\bfseries \color{Blue} L}audibus cives celebrant superni,
te deus simplex pariterque trine,
supplices et nos veniam precamur,
parce redemptis. Amen.
}
\newline {\color{Red} In Primo Nocturno. Ant.} Priusquam te formarem in utero novi te et antequam exires sanctificavi te. {\color{Red} Ps.} \psalm{1} {\color{Red} Ant.} Ad omnia que mittam te dicit dominus ibis, ne timeas et que mandavero tibi loqueris ad eos. {\color{Red} Ps.} \psalm{2} {\color{Red} Ant.} Ne timeas a facie eorum, quia ego tecum sum dicit dominus. {\color{Red} Ps.} \psalm{3} \V Gloria et honore.
\newline \ul{Oratio Anselmi ad beatum Iohannem.} \grecross \ {\color{Red} Lectio Prima.}
\lettrine[lines=2]{\zallmancaps \color{Red} S}{ancte} Iohannes, tu ille Iohannes qui deum baptizasti. Tu prius ab archangelo laudatus, quam genitus a patre, prius plenus deo quam natus ex matre: prius noscens deum quam notus in mundo. Tu ante monstrans matri gravidam matrem dei, quam gravida mater te diei. Tu de quo deus dixit. Inter natos mulierum, non surrexit maior. \R
\lettrine[lines=2]{\zallmancaps \color{Blue} F}{uit} \hypertarget{fuit-homo-missus}{\label{fuit-homo-missus}} homo missus a deo cui nomen Iohannes erat, * Hic venit ut testimonium perhiberet de lumine et pararet domino plebem perfectam. \V Erat Iohannes in deserto predicans baptismum penitentie. Hic venit.
\newline {\color{Red} Lectio Secunda.}
\lettrine[lines=2]{\zallmancaps \color{Red} A}{d} te domine tam magnum, tam sanctum, tam beatum, ad te venit scelerosus vermis, erumnosus homuntio, iam mortuo sensu vixisse se dolens, sed mortua anima sibi nimis dolendus peccator: ad te tam magne amice dei venit valde timens, dubius de salute sua, quia certus de magna culpa sua: sed sperans de maiori gratia tua.
\begin{responsory}[elisabeth-zacharie]
{Elisabeth Zacharie magnum virum genuit Iohannem baptistam precursorem domini, * Qui viam domino preparavit in heremo.}{Erat quidem infecunda corpore, sed servata mysterio, caruitque sterilitatis obprobrio gignendo filium de celo promissum.}
\end{responsory}%
\newline {\color{Red} Lectio Tertia.}
\lettrine[lines=2]{\zallmancaps \color{Blue} M}{aior} est domine gratia tua quam culpa mea, quia plus potes apud deum quam delere scelera mea. Ad te ergo quem gratia fecit tam amicum dei, ad te confugit anxius, quem iniquitas fecit tam reum dei: ad te quem tam beatum fecit gratia, ego quem tam miserum fecit nequitia. Vere domine fateor, iniquitas mea me fecit talem, sed te talem, non tu sed gratia dei tecum.
\begin{responsory-doxology}[priusquam-te-formarem]
{Priusquam te formarem in utero novi te, et antequam exires de ventre sanctificavi te, * Et prophetam in gentibus dedi te.}{Ecce dei verba mea in ore tuo, ecce constitui te super gentes et regna.}
\end{responsory-doxology}%
\newline {\color{Red} In Secundo Nocturno. Ant.} Misit dominus manum suam et tetigit os meum, et prophetam in gentibus dedit me dominus. {\color{Red} Ps.} \psalm{4} {\color{Red} Ant.} Ecce dedi verba mea in ore tuo, ecce constitui te super gentes et regna. {\color{Red} Ps.} \psalm{5} {\color{Red} Ant.} Dominus ab utero vocavit me de ventre matris mee recordatus est nominis mei. {\color{Red} Ps.} \psalm{8} \V Posuisti domine.
\newline {\color{Red} Lectio Quarta.}
\lettrine[lines=2]{\zallmancaps \color{Red} M}{emento} ergo domine, ut sicut gratia dei te sic sublimavit, sic misericordia tua erigat quem culpa sua sic humiliavit. Heu me qualem me feci? In peccatis eram conceptus, et natus, sed abluisti me domine: et ego peioribus sordidavi me. Reformasti in me imaginem amabilem tuam, et ego super induxi odibilem imaginem. Heu heu cuius? O cur non odi eius imitationem: cuius sic horreo nomen?
\begin{responsory}[descendit-angelus-domini]
{Descendit angelus domini ad Zachariam dicens, accipe puerum in senectute tua, * Et habebit nomen Iohannes baptista.}{Ne timeas Zacharia quoniam exaudita est oratio tua et Elysabeth uxor tua pariet tibi filium.}
\end{responsory}%
\newline {\color{Red} Lectio Quinta.}
\lettrine[lines=2]{\zallmancaps \color{Blue} I}{lle} sponte cecidit: ego volens sordui, Sed ille nulla precedente delicti vindicta superbiens peccavit: ego visa eius pena contempnens ad peccatum properavi. Ille semel in innocentia constitutus: ego restitutus. Ille contra eum qui fecit se: ego contra eum qui me fecit et refecit. Ille dereliquit deum permittentem: ego fugi deum prosequentem. Ille perstat in malicia deo reprobante, ego cucurri in illam deo revocante.
\begin{responsory}[hic-precursor-directus]
{Hic precursor directus et lucerna lucens ante dominum, ipse est enim Iohannes qui viam domino preparavit in heremo, sed et agnum dei demonstrabat, * Et illuminabat mentes hominum.}{Erat Iohannes in deserto baptizans et predicans baptismum penitentie.}%typo: precusor
\end{responsory}%
\newline {\color{Red} Lectio Sexta.}
\lettrine[lines=2]{\zallmancaps \color{Blue} I}{lle} obduratus ad punientem, ego obduratus ad blandientem. Et si ambo contra deum, ille contra nec requirentem se, ego contra morientem pro me. O infelix, ecce cuius ymaginis horrebam horrorem: in multis me aspicio horribiliorem. Reforma domine faciem quam fedavi, restaura innocentiam quam violavi. Qui abstulisti peccata que nascendo attuli, tolle peccata que vivendo contraxi. Tolle qui tollis peccata mundi, per merita illius qui hoc testimonio te ostendit mundo, tolle peccata que contraxi in mundo.
\begin{responsory-final}[innuebat-patri-eius]
{Innuebant patri eius quem vellet vocari eum, * Et postulans pugillarem scripsit dicens, * Iohannes est nomen eius.}{Vicini quoque et cognati eius venerunt circumcidere puerum, et sciscitati sunt eum quo nomine vocaretur.}{Iohannes}
\end{responsory-final}%
\newline {\color{Red} In Tertio Nocturno. Ant.} Posuit os meum dominus quasi gladium acutum, sub umbra manus sue protexit me. {\color{Red} Ps.} \psalm{10} {\color{Red} Ant.} Formans me ex utero servum sibi dominus in lucem gentium, ut sim salus eius ab extremis terre. {\color{Red} Ps.} \psalm{14} {\color{Red} Ant.} Reges videbunt et consurgent principes, et adorabunt dominum deum tuum qui elegit te. {\color{Red} Ps.} \psalm{20} \V Magna est gloria.
\newline {\color{Red} Lectio VII. Secundum Lucam.}
\lettrine[lines=2]{\zallmancaps \color{Red} I}{n} illo tempore: Elysabeth impletum est tempus pariendi, et peperit filium. Et reliqua.
\newline {\color{Red} Omelia Venerabilis Bede Presbiteri. \grecross .}
\lettrine[lines=2]{\zallmancaps \color{Blue} P}{recursoris} domini nativitas sicut sacratissima lectionis evangelice prodit hystoria, multa miraculorum sublimitate refulget. Quia nimirum decebat ut ille quo maior inter natos mulierum nemo surrexit: maiore pre ceteris sanctis in ipso mox ortu, virtutum iubare claresceret.
\begin{responsory}[precursor-domini-venit]
{Precursor domini venit, de quo ipse testatur, * Nullus maior inter natos mulierum Iohanne baptista.}{Hic est enim propheta et plusquam propheta de quo dominus ait.}%typo: nator
\end{responsory}%
\newline {\color{Red} Lectio VIII.}
\lettrine[lines=2]{\zallmancaps \color{Red} S}{enes} ac diu infecundi parentes dono nobilissime prolis exultant, ipsi patri quem incredulitas mutum reddiderat, ad salutandum nove preconem gratie, os et lingua reseratur. Nec solum facultas Deum benedicendi restituitur: sed de eo eciam prophetandi virtus augetur.
\begin{responsory}[gabriel-angelus-apparuit]
{Gabriel angelus apparuit Zacharie dicens nascetur tibi filius, nomen eius Iohannes vocabitur, * Et multi in nativitate eius gaudebunt.}{Erit enim magnus coram domino, vinum et siceram non bibet.}
\end{responsory}%
\newline {\color{Red} Lectio IX.}
\lettrine[lines=2]{\zallmancaps \color{Blue} E}{xcitati} fama facti omnes vicini, admiratione ac metu percelluntur, omnium qui audiere circumquaque ad adventum novi prophete corda preparantur. Unde merito sancta per orbem ecclesia, huius tantummodo post dominum eciam nativitatis diem celebrare consuevit.
\begin{responsory-doxology}[inter-natos-mulierum]
{Inter natos mulierum non surrexit maior Iohanne baptista, * Qui viam domino preparavit in heremo.}{Fuit homo missus a deo cui nomen Iohannes erat.}
\end{responsory-doxology}%
\V Fuit homo.
\newline {\color{Red} In Laudibus. Ant.} Elysabeth Zacharie magnum virum genuit Iohannem baptistam precursorem domini. {\color{Red} Ant.} Innuebant patri eius quem vellet vocari eum, et scripsit dicens Iohannes est nomen eius. {\color{Red} Ant.} Iohannes est nomen eius, vinum et siceram non bibet, et multi in nativitate eius gaudebunt. {\color{Red} Ant.} Iohannes vocabitur nomen eius et in nativitate eius multi gaudebunt. {\color{Red} Ant.} Iste puer magnus coram domino, nam et manus eius cum ipso est. {\color{Red} Cap.} Priusquam. {\color{Red} Ymnus.}%typo: vivum
\hymn{o-nimis-felix}{O nimis felix}{
\lettrine[lines=2]{\zallmancaps \color{Red} O}{}\ nimis felix meritique celsi,
nesciens labem nivei pudoris,
prepotens martyr, heremique cultor,
maxime vatum.
\par {\bfseries \color{Blue} S}erta ter denis alios coronant
aucta crementis, duplicata quosdam,
trina centeno cumulata fructu
te sacer ornant.
\par {\bfseries \color{Red} N}unc potens nostri meritis opimis,
pectoris duros lapides repelle,
asperum planans iter, et reflexos
dirige calles.
\par {\bfseries \color{Blue} U}t pius mundi sator et redemptor
mentibus pulsa livione puris,
rite dignetur veniens sacratos
ponere gressus.
\par {\bfseries \color{Red} L}audibus cives celebrant superni,
te deus simplex pariterque trine,
supplices et nos veniam precamur,
parce redemptis. Amen.
}
\newline \V Corona aurea.
\newline {\color{Red} Ad Ben. Ant.} Apertum est os Zacharie, et prophetavit dicens benedictus deus Israel. {\color{Red} Oratio.}
\lettrine[lines=2]{\zallmancaps \color{Blue} D}{eus} qui presentem diem honorabilem nobis in beati Iohannis nativitate fecisti, da populis tuis spiritualium gratiam gaudiorum, et omnium fidelium mentes dirige in viam salutis eterne. Per.
\newline {\color{Red} Ad Terciam. Cap.} Priusquam. \ul{\Rbar , \Vbar , unius martiris, similiter et ad alias horas.}
\newline {\color{Red} Ad Sextam. Cap.}
\lettrine[lines=2]{\zallmancaps \color{Red} A}{udite} insule, et attendite populi de longe, dominus ab utero vocavit me de ventre matris mee recordatus est nominis mei.
\newline {\color{Red} Ad Nonam. Cap.}
\lettrine[lines=2]{\zallmancaps \color{Blue} P}{osuit} os meum dominus quasi gladium acutum, in umbra manus sue protexit me, et posuit me sicut sagittam electam.
\newline {\color{Red} Ad Vesperas. Super psalmos. Ant.} Elysabeth. \ul{Psalmi unius martiris. Cap., Ymnus, \Vbar , ut in Primis Vesperis.}
\newline {\color{Red} Ad Mag. Ant.} Et posuerunt omnes qui audierant in corde suo dicentes, quis putas puer iste erit, etenim manus domini erat cum illo, et Zacharias pater eius impletus est spiritu sancto et prophetavit dicens, benedictus deus Israhel, quia visitavit et fecit redemptionem plebis sue.
\newline \ul{Per Octavam Invitat.} \hyperlink{regem-precursoris-dominum-invitatorium}{Regem precursoris.} \ul{Sequentes vero antiphone dicantur ad Ben. et ad Mag. cotidie vel ad Memoriam.}
\newline {\color{Red} Ad Ben. Ant.} Tu puer propheta altissimi vocaberis, preibis ante dominum parare vias eius.
\newline {\color{Red} Ad Mag. Ant.} Inter natos mulierum non surrexit maior Iohanne baptista.
\newline \ul{Si festum sancti Iohannis quinta vel sexta feria post Trinitatem evenerit, in crastino sancti Iohannis fiat officium de octavis ipsius, et memoria de Trinitate.}%later rubric for Corpus Christi
\newline \ul{Sequentes tres lectiones cum aliis sex de festo: legantur in octava, et infra octavas quotienscumque necesse fuerit.}
\newline \ul{Idem Anselmus in eadem oratione supradicta.} \grecross \ {\color{Red} Lectio.}
\lettrine[lines=2]{\zallmancaps \color{Red} T}{u} sancte Iohannes qui ostendisti mundo tollentem peccata sua, per gratiam tibi datam fac michi hanc misericordiam ut tollat peccata mea. Tu deus tollis peccata mundi, tu amice eius dicis, hic tollit peccata mundi. Ecce ante vos onustus peccatis mundi. Tu tollis, et tu dicis. Ecce ego cui tu tollas quod tu dicis. Ecce medicus, ecce testis eius, et ecce eger servus medici et opus eius, rogans medicum et testem eius.
\newline {\color{Red} Lectio.}
\lettrine[lines=2]{\zallmancaps \color{Blue} V}{ere} medice, oro te, sana me. Verax testis eius, oro te, ora pro me. Proba te in me michi, tu actum tuum, tu dictum tuum. Experiar quod audio, sentiam quod credo. Ihesu bone domine, si tu operaris quod ille testatur, fiat in me opus tuum. Iohannes monstrator dei, si tu testaris quod ille operatur, fiat in me verbum tuum.
\newline {\color{Red} Lectio.}
\lettrine[lines=2]{\zallmancaps \color{Red} S}{ana} me domine tu cuius est sanare, impetra hoc michi domine tu qui potes impetrare. Tu enim magnus dominus, et tu magnus coram domino. Tu summe potens per te ipsum: tu valde potens apud ipsum. Tu summe bonus deus, et ut valde amicus eius qui est in eternum misericors et benedictus deus.
{\ubsubsection{Sanctorum Iohannis et Pauli Martirum}{sanctorum-iohannis-et-pauli-martirum-breviarium}}
\newline {\color{Red} Oratio.}
\lettrine[lines=2]{\zallmancaps \color{Blue} Q}{uesumus} omnipotens deus, ut nos geminata leticia hodierne festivitatis excipiat que de beatorum Iohannis et Pauli glorificatione procedit, quos eadem fides et passio vere fecit esse germanos. Per.
\newline \ul{Ad Matutinas. \Rbar .} \hyperlink{absterget-deus-omnem}{Absterget.} \ul{\Rbar .} \hyperlink{viri-sancti-gloriosum}{Viri sancti.}
\newline \R {\color{Red} III.} \hypertarget{beati-martyres-christi}{\label{beati-martyres-christi}} Beati martyres Christi Iohannes et Paulus, dixerunt Iuliano, nobis alius non est, nisi unus deus, * Pater et filius et spiritus sanctus. \V Una fides, unum baptisma, unus spiritus erat in eis. Pater. Gloria. Pater.
\newline \ul{\Rbar .} \hyperlink{sancti-tui-domine}{Sancti tui.} \ul{\Rbar .} \hyperlink{verbera-carnificum}{Verbera.}
\newline \R {\color{Red} VI.} \hypertarget{hec-est-vera}{\label{hec-est-vera}} Hec est vera fraternitas que nunquam potuit violari certamine, qui effuso sanguine secuti sunt dominum, * Contempnentes aulam regiam, * Pervenerunt ad regna celestia. \V Ecce quam bonum et quam iocundum habitare fratres in unum. Contempnentes. Gloria. Pervenerunt.
\newline \ul{\Rbar .} \hyperlink{in-circuitu-tuo}{In circuitu.} \ul{\Rbar .} \hyperlink{sancti-mei-qui}{Sancti mei.}
\newline \R {\color{Red} IX.} \hypertarget{isti-sunt-due}{\label{isti-sunt-due}} Isti sunt due olive et duo candelabra lucentia ante dominum, * Habent potestatem claudere celum nubibus, et aperire portas eius, * Quia lingue eorum claves celi facte sunt. \V Isti sunt duo filii splendoris, qui assistunt dominatori universe terre. Habent. Gloria. Quia.
\newline {\color{Red} In Laudibus. Ant.} Paulus et Iohannes dixerunt Iuliano, nos unum deum colimus qui fecit celum et terram. {\color{Red} Ant.} Paulus et Iohannes dixerunt ad Terentianum si tuus dominus est Iulianus habeto pacem cum illo, nobis alius non est nisi dominus Ihesus Christus. {\color{Red} Ant.} Iohannes et Paulus cognoscentes tirannidem Iuliani, facultates suas pauperibus erogare ceperunt. {\color{Red} Ant.} Sancti spiritus et anime iustorum ymnum dicite deo in eternum. {\color{Red} Ant.} Iohannes et Paulus dixerunt ad Gallicanum, fac votum deo celi et eris victor melior quam fuisti.
\newline {\color{Red} Ad Ben. Ant.} Beati martyres Christi Iohannes et Paulus dixerunt Iuliano nobis alius non est nisi unus deus pater et filius et spiritus sanctus alleluya.
\newline {\color{Red} Ad Mag. Ant.} Isti sunt due olive et duo candelabra lucentia ante dominum, habent potestatem claudere celum nubibus, et aperire portas eius, quia lingue eorum claves celi facte sunt alleluya.
{\ubsubsection{Sancti Leonis Pape}{sancti-leonis-pape-breviarium}}
\newline {\color{Red} Oratio.}
\lettrine[lines=2]{\zallmancaps \color{Red} D}{eus} qui beatum Leonem pontificem, sanctorum tuorum meritis coequasti, concede propitius, ut qui commemorationis eius festa percolimus, vite quoque imitemur exempla. Per.
{\ubsubsection{In Vigilia Apostolorum Petri et Pauli}{in-vigilia-apostolorum-petri-et-pauli-breviarium}}
\newline {\color{Red} Lectio I. vel VII. si Dominica fuerit. Secundum Iohannem.}
\lettrine[lines=2]{\zallmancaps \color{Blue} I}{n} illo tempore: Dixit Simoni Petro Ihesus. Simon Iohannis diligis me plus his? Dicit ei. Eciam domine: tu scis quia amo te. Et reliqua.
\newline {\color{Red} Omelia Venerabilis Bede Presbiteri.}
\lettrine[lines=2]{\zallmancaps \color{Red} V}{irtutem} nobis perfecte dilectionis presens sancti evangelii lectio commendat. Perfecta namque dilectio est, qua Deum ex toto corde, ex tota anima, ex tota virtute: proximum autem tanquam nos ipsos, diligere iubemur.
\newline {\color{Red} Lectio II. vel VIII.}
\lettrine[lines=2]{\zallmancaps \color{Blue} E}{t} neutra harum dilectio sine altera valet esse perfecta, quia nec Deus vere sine proximo, nec sine Deo vere potest proximus amari. Unde dominus toties interrogato Petro an se diligeret et illo respondente quod eum ipso teste diligeret, adiungebat per singula ita concludens. Pasce oves meas, sive pasce agnos meos. Ac si aperte diceret. Hec sola et vera est probatio integri in dominum amoris, si erga fratres studueris curam solliciti exercere laboris.
\newline {\color{Red} Lectio III. vel IX.}
\lettrine[lines=2]{\zallmancaps \color{Red} Q}{uicunque} enim fratri opus pietatis quod valet impendere negligit, minus se conditorem diligere ostendit, cuius mandatum in sustentanda proximi necessitate contempnit. Hec profecto caritas quoniam sine divine gratie inspiratione minime possit haberi, tacite quodammodo dominus
%pp126
insinuat, quia Petrum de illa interrogans Symonem eum Iohannis quod nusquam alias cognominat.
\newline {\color{Red} Ad Vesperas. Super psalmos. Ant.} Quem dicunt homines esse filium hominis, dixit Ihesus discipulis suis, respondens Petrus dixit, tu es Christus filius dei vivi, et ego dico tibi, quia tu es Petrus et super hanc petram edificabo ecclesiam meam. \ul{Psalmi de Octava.} {\color{Red} Cap.}
\lettrine[lines=2]{\zallmancaps \color{Blue} P}{etrus} quidem servabatur in carcere, oratio autem fiebat sine intermissione ab ecclesia ad deum pro eo. \R \hyperlink{cornelius-centurio}{Cornelius.} {\color{Red} III. Ymnus.}
\hymn{aurea-luce-et}{Aurea luce et decore}{
\lettrine[lines=2]{\zallmancaps \color{Red} A}{urea} luce et decore roseo
lux lucis omne perfudisti seculum,
decorans celos inclito martyrio
hac sacra die, que dat reis veniam.
\par {\bfseries \color{Blue} I}anitor celi, doctor orbis pariter,
iudices secli, vera mundi lumina,
per crucem alter, alter ense triumphans,
vite senatum, laureati possident.
\par {\bfseries \color{Red} O}live bine pietatis unice,
fide devotos, spe robustos, maxime
fonte repletos caritatis gemine,
post mortem carnis impetrate vivere.
\par {\bfseries \color{Blue} S}it trinitati sempiterna gloria,
honor potestas atque iubilatio,
in unitate cui manet imperium,
ex tunc et modo per eterna secula. Amen.
}
\newline \V Tu es Petrus. \R Et super hanc petram edificabo ecclesiam meam.
\newline {\color{Red} Ad Mag. Ant.} Beatus Petrus apostolus vidit sibi Christum occurrere et adorans eum ait, domine quo vadis, venio Romam iterum crucifigi. {\color{Red} Oratio.}
\lettrine[lines=2]{\zallmancaps \color{Blue} D}{eus} qui nos beatorum apostolorum tuorum Petri et Pauli gloriosa natalicia prevenire concedis, tribue quesumus, eorum nos semper et beneficiis preveniri et orationibus adiuvari. Per. \ul{Memoria de sancto Iohanne.}
{\ubsubsection{In Die Apostolorum Petri et Pauli}{in-die--breviarium}}
\newline {\color{Red} In die. Ad Mat. Invitat.}
\begin{invitatory}[christum-regem-regum-invitatorium]
{Christum regem regum adoremus dominum, * Qui martyrio crucis beatum glorificavit Petrum apostolum.}
\end{invitatory}
\newline {\color{Red} Ymnus.} \hyperlink{aurea-luce-et}{Aurea luce. Ianitor. Sit trinitati.}
\newline {\color{Red} In Primo Nocturno. Ant.} In plateis ponebantur infirmi in lectulis, ut veniente Petro saltem umbra illius obumbraret quemquam illorum et liberarentur ab infirmitatibus suis. {\color{Red} Ps.} \psalm{18} {\color{Red} Ant.} Ait Petrus principibus sacerdotum, Ihesum quem vos interemistis suspendentes in ligno, hunc deus suscitavit et principem et salvatorem exaltavit ad dandam penitentiam in remissionem peccatorum. {\color{Red} Ps.} \psalm{33} {\color{Red} Ant.} Petrus apostolus dixit paralitico, Enea sanet te dominus Ihesus Christus, surge et sterne tibi, qui continuo surrexit et omnes qui viderunt conversi sunt ad dominum. {\color{Red} Ps.} \psalm{44} \V In omnem terram. \R
\lettrine[lines=2]{\zallmancaps \color{Red} S}{imon} \hypertarget{simon-petre-antequam}{\label{simon-petre-antequam}} Petre antequam de navi vocarem te novi te, et super plebem meam principem te constitui, * Et claves regni celorum tradidi tibi. \V Quodcumque ligaveris super terram erit ligatum et in celis, et quodcumque solveris super terram erit solutum in celis.
\begin{responsory}[si-diligis-me]
{Si diligis me Symon Petre pasce oves meas, domine tu nosti quia amo te, * Et animam meam pono pro te.}{Si oportuerit me mori tecum, non te negabo.}
\end{responsory}%
\begin{responsory-final}[cornelius-centurio]
{Cornelius centurio vir religiosus ac timens deum, vidit manifeste angelum dei dicentem sibi, * Corneli mitte et accersi Symonem qui cognominatur Petrus, * Hic dicet tibi quid te oporteat facere.}{Cum orasset Cornelius nondum in Christo renatus, apparuit ei angelus dicens.}{Hic dicet}
\end{responsory-final}%
\newline {\color{Red} In Secundo Nocturno. Ant.} Factum est autem ut quedam discipula nomine Tabita plena operibus bonis et elemosinis infirmata moreretur, miserunt autem discipuli ad Petrum rogantes, ne pigriteris venire usque ad nos. {\color{Red} Ps.} \psalm{46} {\color{Red} Ant.} Adveniente Petro circumsteterunt illum omnes vidue flentes et ostendentes tunicas et vestes quas faciebat illis Dorcas. {\color{Red} Ps.} \psalm{60} {\color{Red} Ant.} Ponens Petrus genua sua oravit, et conversus ad corpus dixit, Tabita surge, at illa aperuit oculos, et viso Petro resedit, dansque illi manum erexit eam et sanctis ac viduis assignavit vivam. {\color{Red} Ps.} \psalm{63} \V Constitues.
\begin{responsory}[domine-si-tu]
{Domine si tu es iube me venire ad te super aquas, et extendens manum apprehendit eum, * Et dixit Ihesus modice fidei quare dubitasti.}{Cumque vidisset ventum validum venientem timuit et cum cepisset mergi clamavit dicens, domine salvum me fac.}
\end{responsory}%
\begin{responsory}[tu-es-pastor]
{Tu es pastor ovium princeps apostolorum, tibi tradidit deus omnia regna mundi, * Et ideo tradite sunt tibi claves regni celorum.}{Tibi enim a domino collata est potestas ligandi atque solvendi.}
\end{responsory}%
\begin{responsory-doxology}[surge-petre]
{Surge Petre et indue vestimenta tua, accipe fortitudinem ad salvandas gentes, * Quia ceciderunt catene de manibus tuis.}{Angelus domini astitit et lumen refulsit in habitaculo carceris, percussoque latere Petri excitavit eum dicens, surge velociter.}
\end{responsory-doxology}%
\newline {\color{Red} In Tertio Nocturno. Ant.} Cornelius centurio vir religiosus ac timens deum, vidit manifeste angelum dei dicentem sibi, Corneli mitte et accersi Symonem qui cognominatur Petrus, hic dicet tibi, quid te oporteat facere. {\color{Red} Ps.} \psalm{74} {\color{Red} Ant.} Aperiens Petrus os suum dixit, in veritate comperi quoniam non est personarum acceptor deus, sed in omni gente qui timet deum et operatur iusticiam acceptus est illi. {\color{Red} Ps.} \psalm{96} {\color{Red} Ant.} Adhuc loquente Petro cecidit spiritus sanctus super omnes qui audiebant verbum et erant loquentes linguis et magnificantes deum, tunc Petrus iussit eos baptizari in nomine domini Ihesu Christi. {\color{Red} Ps.} \psalm{98} \V Nimis honorati.
\newline {\color{Red} Lectio VII. Secundum Matheum.}
\lettrine[lines=2]{\zallmancaps \color{Blue} I}{n} illo tempore: Venit Ihesus in partes Cesaree Philippi, et interrogabat discipulos suos dicens. Quem dicunt homines esse filius hominis? Et reliqua.
\newline {\color{Red} Omelia Beati Ambrosii Episcopi.}
\lettrine[lines=2]{\zallmancaps \color{Red} N}{ec} turbe quidem ociosa opinio est qua alii Helyam quem venturum putabant, alii Iohannem quem decollatum sciebant, aut unum de prophetis prioribus surrexisse credebant. Sed hoc supra nos est querere. Alterius scientie: alterius prudentie est.
\begin{responsory}[quem-dicunt-homines]
{Quem dicunt homines esse filium hominis dixit Ihesus discipulis suis, respondens Petrus dixit tu es Christus filius vivi, * Et ego dico tibi quia tu es Petrus et super hanc petram edificabo ecclesiam meam.}{Beatus es Symon Bariona quia caro et sanguis non revelavit tibi, sed pater meus qui est in celis.}
\end{responsory}%
\newline {\color{Red} Lectio VIII.}
\lettrine[lines=2]{\zallmancaps \color{Blue} N}{am} si apostolo Paulo satis est nichil scire nisi Ihesum Christum et hunc crucifixum quid amplius michi desiderandum est scire quam Christum? In uno enim hoc nomine et divinitatis et incarnationis expressio: et fides est passionis. Et ideo licet ceteri apostoli siluerint, Petrus tamen respondit pro ceteris. Tu es Christus filius Dei vivi.
\begin{responsory}[ego-pro-te]
{Ego pro te rogavi Petre, * Ut non deficiat fides tua, et tu aliquando conversus confirma fratres tuos.}{Symon, ecce Sathan expetivit vos ut cribraret sicut triticum, ego autem rogavi pro te.}%typo ro-rogavi
\end{responsory}%
\newline {\color{Red} Lectio IX.}
\lettrine[lines=2]{\zallmancaps \color{Red} C}{omplexus} est itaque omnia, qui et naturam et nomen expressit in quo summa virtutum est. Scivit enim Petrus quod in filio Dei omnia sunt. Omnia enim dedit pater filio. Si omnia dedit, eternitatem quam habet maiestatemque transfudit. Sed quo protelabor longius? Finis fidei mee Christus est: finis fidei mee filius Dei est.
\begin{responsory-final}
{Quodcumque ligaveris super terram erit ligatum et in celis, * Et quodcumque solveris super terram erit solutum et * In celis.}{Et ego dico tibi quia tu es Petrus et super hanc petram edificabo ecclesiam meam.}{In celis}%typo so-solutum
\end{responsory-final}%
\V Tu es Petrus.
\newline {\color{Red} In Laudibus. Ant.} Petrus et Iohannes ascendebant in templum ad horam orationis nonam. {\color{Red} Ant.} Argentum et aurum non est michi, quod autem habeo hoc tibi do. {\color{Red} Ant.} Dixit angelus ad Petrum circumda tibi vestimentum tuum et sequere me. {\color{Red} Ant.} Petre amas me pasce oves meas, tu scis domine quia amo te. {\color{Red} Ant.} Tu es Petrus et super hanc petram edificabo ecclesiam meam. {\color{Red} Cap.} Petrus quidem. {\color{Red} Ymnus.} \hyperlink{aurea-luce-et}{Olive bine. Sit trinitati.} \V Annunciaverunt.
\newline {\color{Red} Ad Ben. Ant.} Quodcumque ligaveris super terram erit ligatum et in celis et quodcumque solveris super terram erit solutum et in celis dixit dominus Symoni Petro. {\color{Red} Oratio.}
\lettrine[lines=2]{\zallmancaps \color{Blue} D}{eus} qui hodiernam diem apostolorum tuorum Petri et Pauli martyrio consecrasti, da ecclesie tue eorum in omnibus sequi preceptum per quos religionis sumpsit exordium. Per. \ul{Memoria de sancto Iohanne.}
\newline {\color{Red} Ad Terciam. Cap.} Petrus quidem. \ul{Responsoria et \Vbar . horarum de communi apostolorum.}
\newline {\color{Red} Ad Sextam. Cap.}
\lettrine[lines=2]{\zallmancaps \color{Red} A}{ngelus} domini astitit et lumen refulsit in habitaculo, percussoque latere Petri excitavit eum dicens, surge velociter, et ceciderunt catene de manibus eius.
\newline {\color{Red} Ad Nonam. Capitulum.}
\lettrine[lines=2]{\zallmancaps \color{Blue} N}{unc} scio vere quia misit dominus angelum suum, et eripuit me de manu Herodis, et de omni expectatione plebis Iudeorum.
\newline {\color{Red} Ad Vesperas. Super psalmos. Ant.} Petrus et Iohannes. {\color{Red} Ps.} \psalm{109} \ul{et ceteri de apostolis.} {\color{Red} Cap.} Iam non estis. \ul{Ymnus ut in Primis Vesperis.} \V In omnem.
\newline {\color{Red} Ad Mag. Ant.} Gloriosi principes terre quomodo in vita sua dilexerunt se, ita et in morte non sunt separati. {\color{Red} Oratio.} Deus qui hodiernam. \ul{Memoria de sancto Iohanne.}
{\ubsubsection{In Commemoratione Sancti Pauli}{in-commemoratione-sancti-pauli-breviarium}}
\newline {\color{Red} Ad Mat. Invit.}
\begin{invitatory}[laudemus-deum-commemoratione-invitatorium]
{Laudemus deum nostrum, * In sollemnitate apostoli Pauli.}
\end{invitatory}
\newline {\color{Red} Ymnus.}
\lettrine[lines=2]{\zallmancaps \color{Red} D}{octor} egregie Paule mores instrue
et mente polum nos transferre satage,
donec perfectum largiatur plenius,
evacuato quod ex parte gerimus.
\par {\bfseries \color{Blue} S}it trinitati sempiterna gloria,
honor potestas atque iubilatio,
in unitate cui manet imperium
ex tunc et modo per eterna secula. Amen.
\newline {\color{Red} In Primo Nocturno. Ant.} Qui operatus est Petro in apostolatum operatus est et michi inter gentes, et cognoverunt gratiam que data est michi a Christo domino. {\color{Red} Ps.} \psalm{18} \V Qui me segregavit ex utero matris mee et vocavit per gratiam suam. {\color{Red} Ant.} Scio cui credidi et certus sum quia potens est depositum meum servare in illum diem iustus iudex. {\color{Red} Ps.} \psalm{33} \V De reliquo reposita est michi corona iusticie. {\color{Red} Ant.} Michi vivere Christus est et mori lucrum, gloriari me oportet in cruce domini mei Ihesu Christi. {\color{Red} Ps.} \psalm{44} \V Per quem michi mundus crucifixus est et ego mundo. \R
\lettrine[lines=2]{\zallmancaps \color{Blue} Q}{ui} \hypertarget{qui-operatus-est}{\label{qui-operatus-est}} operatus est Petro in apostolatum, operatus est et michi inter gentes, et cognoverunt gratiam dei, * Que data est michi. \V Gratia dei in me vacua non fuit, et gratia eius semper in me manet. Que.
\begin{responsory}[scio-cui-credidi]
{Scio cui credidi, et certus sum quia potens est depositum meum servare, * In illum diem.}{Resposita est michi corona iusticie quam reddet michi dominus.}
\end{responsory}%
\begin{responsory-final}[bonum-certamen]
{Bonum certamen certavi, cursum consummavi, fidem servavi, * Ideoque reposita est michi, * Corona iusticie.}{Gratia dei sum id quod sum, et gratia eius in me vacua non fuit.}{Corona}
\end{responsory-final}%
\newline {\color{Red} In Secundo Nocturno. Ant.} Bonum certamen certavi, cursum consummavi, fidem servavi. {\color{Red} Ps.} \psalm{46} \V De reliquo reposita est michi corona iusticie. {\color{Red} Ant.} Tu es vas electionis sancte Paule apostole predicator veritatis in universo mundo. {\color{Red} Ps.} \psalm{60} \V Per quem omnes gentes cognoverunt gratiam dei. {\color{Red} Ant.} Magnus sanctus Paulus vas electionis vere digne est glorificandus, qui et meruit thronum duodecimum possidere. {\color{Red} Ps.} \psalm{63} \V In regeneratione cum sederit filius hominis in sede maiestatis sue.
\begin{responsory}[reposita-est-michi]
{Reposita est michi corona iusticie quam reddet michi dominus, * In illum diem iustus iudex.}{Scio cui credidi, et certus sum quia potens est depositum meum servare.}
\end{responsory}%
\begin{responsory}[gratia-dei-sum]
{Gratia dei sum id quod sum, * Et gratia eius in me vacua non fuit, sed semper in me manet.}{Qui operatus est Petro in apostolatum operatus est et michi inter gentes.}
\end{responsory}%
\begin{responsory-final}[michi-vivere-christus]
{Michi vivere Christus est et mori lucrum, * Gloriari me oportet, * In cruce domini mei Ihesu Christi.}{Per quem michi mundus crucifixus est et ego mundo.}{In cruce}
\end{responsory-final}%
\newline {\color{Red} In Tertio Nocturno. Ant.} Ter virgis cesus sum, semel lapidatus sum, ter naufragium pertuli pro Christo domino. {\color{Red} Ps.} \psalm{74} \V Nocte ac die in profundo maris fui. {\color{Red} Ant.} Reposita est michi corona iusticie quam reddet michi dominus in illum diem iustus iudex. {\color{Red} Ps.} \psalm{96} \V Cooperante gratia spiritus sancti. {\color{Red} Ant.} Ne magnitudo revelationum extollat me datus est michi stimulus angelus Sathane qui me colafizet, propter quod ter dominum rogavi ut auferretur a me, et dixit michi dominus, sufficit tibi Paule gratia mea. {\color{Red} Ps.} \psalm{98} \V Nam virtus in infirmitate perficitur.%typo regavi
\begin{responsory}[damasci-prepositus]
{Damasci prepositus gentis Arethe regis voluit me comprehendere, a fratribus per murum, * Summissus sum in sporta, et sic evasi manus eius in nomine domini.}{Deus et pater domini nostri Ihesu Christi scit quia non mentior.}
\end{responsory}%
\begin{responsory}
{Tu es vas electionis sancte Paule apostole predicator veritatis in universo mundo, * Per quem omnes gentes cognoverunt gratiam dei.}{Intercede pro nobis ad deum qui te elegit.}
\end{responsory}%
\begin{responsory-doxology}
{Sancte Paule apostole predicator veritatis et doctor gentium intercede pro nobis ad deum, * Qui te elegit.}{Ut digni efficiamur gratia dei.}
\end{responsory-doxology}%
\newline {\color{Red} In Laudibus. Ant.} Ego plantavi Apollo rigavit, deus autem incrementum dedit alleluya. \V Unusquisque propriam mercedem accipiet secundum suum laborem. {\color{Red} Ant.} Libenter gloriabor in infirmitatibus meis ut inhabitet in me virtus Christi. \V Quando enim infirmor, tunc fortior sum et potens. {\color{Red} Ant.} Gratia dei in me vacua non fuit, sed gratia eius semper in me manet. \V Gratia dei sum id quod sum. {\color{Red} Ant.} Damasci prepositus gentis Arethe regis voluit me comprehendere, a fratribus per murum summissus sum in sporta, et sic evasi manus eius in nomine domini. \V Deus et pater domini nostri Ihesu Christi scit quia non mentior. {\color{Red} Ant.} Sancte Paule apostole predicator veritatis et doctor gentium intercede pro nobis ad deum qui te elegit. \V Ut digni efficiamur gratia dei. {\color{Red} Capitulum.}
\lettrine[lines=2]{\zallmancaps \color{Blue} M}{agnificabitur} Christus in corpore meo, sive per vitam sive per mortem, michi enim vivere Christus est, et mori lucrum. {\color{Red} Ymnus.} \hyperlink{eterna-christi-munera-apostolorum}{Eterna Christi munera apostolorum.}
\newline {\color{Red} Ad Ben. Ant.} Ego enim iam delibor et tempus mee resolutionis instat, bonum certamen certavi, cursum consummavi, fidem servavi, super est michi corona iusticie quam reddet michi dominus in illa die iustus iudex. {\color{Red} Oratio.}
\lettrine[lines=2]{\zallmancaps \color{Red} D}{eus} qui universum mundum beati Pauli apostoli predicatione docuisti, da nobis quesumus, ut qui cuius natalicia colimus, eius apud te patrocinia sentiamus. Per.
\newline \ul{Memoria de sancto Petro.} {\color{Red} Ant.} Petre amas me pasce oves meas, tu scis domine quia amo te. \V Tu es Petrus. {\color{Red} Oratio.}
\lettrine[lines=2]{\zallmancaps \color{Blue} D}{eus} qui beato Petro apostolo tuo collatis clavibus regni celestis ligandi atque solvendi pontificium tradidisti, concede ut intercessionis eius auxilio a peccatorum nostrorum nexibus liberemur. Qui vivis.
\newline \ul{Ad horas, antiphone Laudum sine versibus.}
\newline {\color{Red} Ad Terciam. Cap.} Magnificabitur.
\newline {\color{Red} Ad Sextam. Cap.}
\lettrine[lines=2]{\zallmancaps \color{Red} S}{cio} cui credidi, et certus sum quia potens est depositum meum servare in illum diem.
\newline {\color{Red} Ad Nonam. Cap.}
\lettrine[lines=2]{\zallmancaps \color{Blue} B}{onum} certamen certavi, cursum consummavi, fidem servavi, in reliquo reposita est michi corona iusticie, quam reddet michi dominus in illa die iustus iudex.
\newline {\color{Red} Ad Vesperas. Super psalmos. Ant.} Ego plantavi. \ul{sine versu.} {\color{Red} Cap.} Iam non estis. {\color{Red} Ymnus.} \hyperlink{aurea-luce-et}{Aurea.} \V In omnem.
\newline {\color{Red} Ad Mag. Ant.} Isti sunt due olive et duo candelabra lucentia ante dominum, habent potestatem claudere celum nubibus et aperire portas eius quia lingue eorum claves celi facte sunt, alleluya. {\color{Red} Oratio.} Deus qui hodiernam. \ul{Memoria de sancto Iohanne.} {\color{Red} Ant.} Pro eo. \V Fuit homo. {\color{Red} Oratio.} Deus qui presentem.
{\ubsubsection{In Octava Sancti Iohannis}{in-octava-sancti-iohannis-breviarium}}
\newline \ul{In Laudibus post primam orationem fiat memoria de apostolis.} {\color{Red} Ant.} Petrus apostolus, et Paulus doctor gentium, ipsi nos docuerunt legem tuam domine. \V Annunciaverunt. {\color{Red} Oratio.} Deus qui hodiernam.
\newline {\color{Red} In Vesperis. Ad Mag. Ant.} Et factum est in die octavo venerunt pater et mater et cognati eius ad circumcisionem pueri, et vocabant eum nomine patris sui Zachariam, et dixit mater eius nequaquam, sed vocabitur Iohannes. \ul{Post orationem de sancto Iohanne fiat memoria de apostolis Ant.} Gloriosi. \ul{\Vbar .} In omnem. \ul{Oratio} Deus qui hodiernam. \ul{Postea, memoria de sanctis Processo et Martiniano.} {\color{Red} Oratio.}
\lettrine[lines=2]{\zallmancaps \color{Red} D}{eus} qui nos sanctorum tuorum Processi et Martiniani confessionibus gloriosis circundas et protegis, da nobis et eorum imitatione proficere et intercessione gaudere. Per.
\newline \ul{Per Octavam apostolorum nisi in Dominica, ad Matutinas, Invitat.}
\begin{invitatory}[regem-apostolorum-dominum-invitatorium]
{Regem apostolorum dominum, * Venite adoremus.}
\end{invitatory}
\ul{Ymnus sicut in die, super psalmos una tantum antiphona de communi apostolorum. Psalmi novem de apostolis, \Vbar , et Responsoria de communi apostolorum secundum feriam. Quod si versiculus sit idem cum antiphona: sequens \Vbar , dicatur.} Te deum. \ul{\Vbar .} Dedisti. \ul{In Laudibus} habet sola \ul{antiphonam.} Hoc est. \ul{Cap., \Vbar , de communi apostolorum. Ymnus.} \hyperlink{aurea-luce-et}{Olive.} \ul{Ad Ben. Ant.} Petrus apostolus. \ul{Oratio.} Deus qui hodiernam. \ul{Ad horas diei de communi apostolorum. Ad Vesperas super psalmos antiphona.} Hoc est preceptum. \ul{Psalmi de apostolis. Cap.} Iam non estis. \ul{Ymnus.} \hyperlink{aurea-luce-et}{Aurea.} \ul{sicut in die. \Vbar .} In omnem. \ul{Ad Mag. Ant.} Gloriosi. \ul{Oratio.} Deus qui hodiernam. \ul{Predicte antiphone scilicet.} Petrus apostolus. \ul{et} Gloriosi. \ul{dicantur cotidie per octavam.} Petrus. \ul{ad Ben. et} Gloriosi. \ul{Ad Mag.}
{\ubsubsection{Dominica infra Octavam Apostolorum}{dominica-infra-octavam-apostolorum-breviarium}}
\newline \ul{Ad Mat. Invit., antiphone, psalmi, versiculi, Responsoria de communi apostolorum. Ymnus.} \hyperlink{aurea-luce-et}{Aurea.} \ul{sicut in die. Tres ultime lectiones de omelia dominicali. In Laudibus Ant.} Hoc est preceptum. \ul{Cum ceteris. Ymnus.} \hyperlink{aurea-luce-et}{Olive.} \ul{Ad horas diei de communi apostolorum cum oratione.} Deus qui hodiernam. \ul{Ad Vesperas super psalmos. Ant.} Hoc est. \ul{Ymnus.} \hyperlink{aurea-luce-et}{Aurea.} \ul{sicut in die. \Vbar .} In omnem. \ul{Ad Mag. Ant.} Gloriosi.
{\ubsubsection{In Octava Apostolorum}{in-octava-apostolorum-breviarium}}
\newline \ul{Ad Vesperas. Super psalmos. Ant.} Hoc est. \ul{Psalmi de apostolis. Cap.} Iam non estis. \ul{Ymnus.} \hyperlink{aurea-luce-et}{Aurea.} \ul{\Vbar .} In omnem. \ul{Ad Mag. Ant.} Gloriosi. {\color{Red} Oratio.}
\lettrine[lines=2]{\zallmancaps \color{Blue} D}{eus} cuius dextera beatum Petrum ambulantem in fluctibus ne mergeretur erexit et coapostolum eius Paulum tercio naufragantem de profundo pelagi liberavit, exaudi nos propicius et concede, ut amborum meritis eternitatis gloriam consequamur. Qui vivis.
\newline \ul{Ad Matutinas. Invit.} \hyperlink{gaudete-et-exultate-invitatorium}{Gaudete.} \ul{Ymnus.} \hyperlink{aurea-luce-et}{Aurea.} \ul{antiphone, psalmi, \Vbar , Responsoria de communi apostolorum.}
\newline {\color{Red} Lectio VII. Secundum Matheum.}
\lettrine[lines=2]{\zallmancaps \color{Red} I}{n} illo tempore: Iussit Ihesus discipulos suos ascendere in naviculam, et precedere eum trans fretum, donec dimitteret turbas. Et reliqua.
\newline {\color{Red} Omelia Beati Augustini Episcopi. \grecross .}
\lettrine[lines=2]{\zallmancaps \color{Blue} C}{um} sanctum evangelium legeretur audivimus naviculam periclitantem, Christum periclitantibus subvenientem, Petrum venienti Christo occurrentem. In quibus omnibus miraculum spectavimus, mysterium requiramus.
\newline {\color{Red} Lectio VIII.}
\lettrine[lines=2]{\zallmancaps \color{Red} Q}{uando} enim possumus de divinis operibus que leguntur intellectum alicuius mystice significationis exculpere, quasi de abstrusis favorum cellis mella producimus: vel Christi discipulos imitantes spicas manibus confricamus, ut ad latentia grana perveniamus, et in eis vitam inveniamus.
\newline {\color{Red} Lectio IX.}
\lettrine[lines=2]{\zallmancaps \color{Blue} A}{scendit} ergo dominus noster Ihesus Christus, ut lectum est in montem solus orare. Mons altitudo est. Quid enim in hoc mundo altius celo? Quis vero in celum ascendit, novit optime fides vestra. Cur autem solus ascendit? Quia nemo ascendit in celum nisi qui descendit de celo: filius hominis qui est in celo.
\newline \ul{In Laudibus antiphone, Cap., de communi apostolorum. Ymnus.} \hyperlink{aurea-luce-et}{Olive.} \ul{\Vbar .} Annunciaverunt. \ul{Ad Ben. Ant.} Petrus apostolus. \ul{Oratio.} Deus cuius dextera. \ul{Ad horas diei de communi apostolorum. Ad Vesperas super psalmos Ant.} Hoc est. \ul{Psalmi de apostolis. Cap., Ymnus, \Vbar , ut in Primis Vesperis. Ad Mag. Ant.} Isti sunt due.
{\ubsubsection{Sanctorum Septem Fratrum Martirum}{sanctorum-septem-fratrum-martirum-breviarium}}
\newline {\color{Red} Oratio.}
\lettrine[lines=2]{\zallmancaps \color{Red} P}{resta} quesumus omnipotens deus ut qui gloriosos martyres fortes in sua confessione cognovimus, pios apud te in nostra intercessione sentiamus. Per.
{\ubsubsection{Sancte Margarete Virginis et Martiris}{sancte-margarete-virginis-et-martiris-breviarium}}
\newline {\color{Red} Oratio.}
\lettrine[lines=2]{\zallmancaps \color{Blue} I}{ndulgentiam} nobis domine beata Margareta virgo et martyr imploret, que tibi grata semper extitit et merito castitatis, et tue professione virtutis. Per.
{\ubsubsection{Sancte Praxedis Virginis}{sancte-praxedis-virginis-breviarium}}
\newline {\color{Red} Oratio.}
\lettrine[lines=2]{\zallmancaps \color{Red} E}{xaudi} nos deus salutaris noster, ut sicut de beate Praxedis virginis tue festivitate gaudemus, ita pie devotionis erudiamur affectu. Per.
\newline \ul{Ad Matutinas Ymnus.} \hyperlink{huius-obtentu}{Huius obtentu. Gloria patri geniteque.} \ul{Tria ultima Responsoria de communi unius virginis.}
{\ubsubsection{Sancte Marie Magdalene}{sancte-marie-magdalene-breviarium}}
\newline {\color{Red} Ad Vesperas. Super psalmos. Ant.} Recumbente Ihesu in domo pharisei Symonis, accessit ad eum Maria Magdalene afferens preciosi libram unguenti. {\color{Red} Capitulum.}
\lettrine[lines=2]{\zallmancaps \color{Blue} M}{ulierem} fortem quis inveniet, procul et de ultimis finibus precium eius, confidit in ea cor viri sui et spoliis non indigebit. \R \hyperlink{flavit-auster}{Flavit auster.} {\color{Red} VI. Ymnus.}
\hymn{lauda-mater-ecclesia}{Lauda mater ecclesia}{
\lettrine[lines=2]{\zallmancaps \color{Red} L}{auda} mater ecclesia,
lauda Christi clementiam,
qui septem purgat vitia
per septiformam gratiam.
\par {\bfseries \color{Blue} M}aria soror Lazari
que tot commisit crimina,
ab ipsa fauce Tartari
redit ad vite limina.
\par {\bfseries \color{Red} P}ost fluxe carnis scandala
fit ex lebete phyala
in vas translata glorie
de vase contumelie.
\par {\bfseries \color{Blue} E}gra currit ad medicum
vas ferens aromaticum,
et a morbo multiplici
verbo curatur medici.
\par {\bfseries \color{Red} S}urgentem cum victoria
Ihesum videt ab inferis,
prima meretur gaudia
que plus ardebat ceteris.
\par {\bfseries \color{Blue} U}ni deo sit gloria
pro multiformi gratia,
qui culpas et supplicia
remittit, et dat premia. Amen.
}
\newline \V Optimam partem elegit sibi Maria. \R Que non auferetur ab ea.
\newline {\color{Red} Ad Mag. Ant.} Celsi meriti Maria que solem verum resurgentem videre meruisti mortalium prima, obtine ut nos visu glorie sue tecum letificet in celis. {\color{Red} Oratio.}
\lettrine[lines=2]{\zallmancaps \color{Blue} L}{argire} nobis clementissime pater, quod sicut beata Maria Magdalena dominum nostrum Ihesum Christum super omnia diligendo suorum obtinuit veniam peccatorum, ita nobis apud misericordiam tuam sempiternam impetret beatitudinem. Per eundem.
\newline {\color{Red} Ad Matutinas. Invit.}
\begin{invitatory}[eternum-trinumque-invitatorium]
{Eternum trinumque deum laudemus et unum, * Qui sibi Mariam transvexit in ethera sanctam.}
\end{invitatory}
\newline {\color{Red} Ymnus.} \hyperlink{lauda-mater-ecclesia}{Lauda mater.}
\newline {\color{Red} In Primo Nocturno. Ant.} Cum discubuisset in domo Symonis dominus Ihesus mundi gloria, mulier secum que erat peccatrix nardi odorifera exhibuit unguenta. {\color{Red} Ps.} \psalm{8} {\color{Red} Ant.} Secus pedes domini astans incessanter crine soluto diva vestigia osculatur lacrimis plena. {\color{Red} Ps.} \psalm{18} {\color{Red} Ant.} Irrigabat igitur dominicos Ihesu pedes, fracto quoque alabastro precipuis odoribus fragrabat domus omnis, atque capillorum officio detergebat. {\color{Red} Ps.} \psalm{23} \V Diffusa est.
\newline \ul{Ex gestis eius. \textbf{T}.} {\color{Red} Lectio Prima.}
\lettrine[lines=2]{\zallmancaps \color{Red} B}{eata} Maria Magdalene et Martha et Lazarus ex nobili et regali progenie Syro patre et Eukaria matre nati, a beato Maximino discipulo domini baptizati fuerunt, et Magdalum castrum quod est secundo miliario a Genesareth, et Bethaniam que est iuxta Iherusalem, et magnam Hierosolimarum partem possidentes, ita sibi invicem diviserunt, quod Maria Magdalum a quo et Magdalena nuncupata est, et Lazarus partem urbis Hierusalem, et
%pp127
Martha Bethaniam propria possiderent. \R
\lettrine[lines=2]{\zallmancaps \color{Blue} L}{etetur} \hypertarget{letetur-omne-seculum}{\label{letetur-omne-seculum}} omne seculum in sollemnitate sancte Marie, * Quam Ihesus eternus amor dilexit plurimum. \V Hec Maria fuit illa domino gratissima que unguento precioso pedes unxit domini. Quam.
\newline {\color{Red} Lectio Secunda.}
\lettrine[lines=2]{\zallmancaps \color{Red} C}{um} autem Maria deliciis corporis se totam exponeret, et Lazarus militie plus vacaret, Martha prodens et eciam partes fratris et sororis sue suo sensu et probitate gubernans: militibus et familiis suis copiose necessaria ministrabat, et dominum cum discipulis et alios pauperes sustentavit, et post ascensionem domini residuum ante pedes apostolorum in commune divisit.
\begin{responsory}[hec-est-illa]
{Hec est illa Maria que unxit dominum unguento et extersit pedes eius capillis suis, * Cuius frater erat Lazarus quem a mortuis resuscitavit Ihesus.}{Hec illa est Maria cui se dominus resurgens a mortuis primum ostendit.}
\end{responsory}%
\newline {\color{Red} Lectio Tertia.}
\lettrine[lines=2]{\zallmancaps \color{Red} I}{llo} vero tempore quo dominus predicabat, Maria sic erat exposita corporee voluptati quod iam non solum nomen sed eciam cognomen suum perdiderat, et peccatrix consueverat appellari. Cum vero post missionem spiritus Paracliti crescente numero fidelium cresceret infidelis crudelitas Iudeorum, in tantum ut eos multis modis affligerent, et a terra per mare sine remis et velis et ceteris alimentis expellerent: beatus Maximinus cum Lazaro et sororibus eius et multis aliis transfretantes: regente domino Massiliam devenerunt.
\begin{responsory-final}[felix-maria-unxit]
{Felix Maria unxit pedes Ihesu et extersit capillis suis, * Et domus impleta est, * Ex odore unguenti.}{Mixtum rorem balsami fracto fudit alabastro, que unguento precioso pedes unxit domini.}{Ex odore}
\end{responsory-final}%
\newline {\color{Red} In Secundo Nocturno. Ant.} Symon autem intra se inquit, si esset propheta sciret profecto que et qualis esset que tangit eum quia peccatrix est. {\color{Red} Ps.} \psalm{44} {\color{Red} Ant.} Et conversus dominus dixit, postquam venit in domum tuam non cessat crine fuso hec meos osculari pedes. {\color{Red} Ps.} \psalm{45} {\color{Red} Ant.} Quoniam multum dilexeras dimissa scito tibi peccamina, et fide qua polles sacra in pace remitto. {\color{Red} Ps.} \psalm{86} \V Specie tua.
\newline {\color{Red} Lectio Quarta.}
\lettrine[lines=2]{\zallmancaps \color{Blue} P}{ostmodum} cum Aquensem civitatem cum adiacenti provincia verbis et miraculis convertissent: beata Magdalena soli deo vacare cupiens, in quadam rupe excelsa quatuordecim fere milibus a Massilia plusquam triginta duobus annis hominibus ignota permansit. Horis vero septem canonicis cotidie ibi manibus angelicis in ethera ferebatur: et sic post angelicas melodias dei laudibus plenissime satiata ad locum illum ab angelis reportabatur.
\begin{responsory}[videns-ihesus]
{Videns Ihesus flentem Mariam et sororem eius Martham lacrimatus est, et accedens ad monumentum, * Quatriduanum iam Lazarum resuscitavit.}{O felicis soror utraque meriti quarum lacrimis motus fons ipse pietatis.}
\end{responsory}%
\newline {\color{Red} Lectio Quinta.}
\lettrine[lines=2]{\zallmancaps \color{Red} T}{unc} temporis cum quidem heremita penitentie datus iuxta rupem predictam cellulam collocasset, vidit angelos ad rupem illam descendentes, et beatam Magdalenam in altum levatam post hore spacium referentes. Cum autem illuc admirans accessisset: invenit eam. Ipsa autem quenam esset eidem indicans: rogavit eum ut enarraret beato Maximino, et rogaret ut sequenti dominica hora matutinarum ecclesiam suam solus intraret, et ibi eam inveniret subvectam manibus angelorum.
\begin{responsory}[laudemus-opus-dei]
{Laudemus opus dei in Maria genitrice sed virgine, laudemus in Maria peccatrice sed penitente, Maria hec est data speculum innocentie, * Maria hec est exemplum penitentie.}{Ut confiteamur nomini sancto tuo deus et gloriemur in laude tua.}
\end{responsory}%
\newline {\color{Red} Lectio Sexta.}
\lettrine[lines=2]{\zallmancaps \color{Blue} V}{eniens} ille: narravit omnia Maximino. Ille mirans et gaudens: hora predicta ingreditur ecclesiam et beatam Magdalenam a terra duobus cubitis elevatam, inter angelos intuetur. Cum autem accedere trepidaret: iussu eius alios convocavit, et accipiens sacram communionem ab episcopo: ante altare orans cum lacrimis, cunctis videntibus expiravit.
\begin{responsory-final}[flavit-auster]
{Flavit auster et fugavit aquilonem quando lavit cor Marie penitentis ymber sancti spiritus, * Liquefecit et refecit resolutam in lamentis * Verbum missum celitus.}{Fluminis impetus civitatem dei letificavit.}{Verbum}
\end{responsory-final}%
\newline {\color{Red} In Tertio Nocturno. Ant.} Satagebat igitur Martha soror Marie Magdalene circa ministeria, que sedens secus pedes domini audiebat verba oris eius. {\color{Red} Ps.} \psalm{95} {\color{Red} Ant.} Non est Martha inquit tibi cure quod soror mea me reliquit solam ministrare, iube illi ut me adiuvet. {\color{Red} Ps.} \psalm{96} {\color{Red} Ant.} Et respondens dixit illi dominus unum est necesse, Maria optimam sibi partem elegit que non auferetur. {\color{Red} Ps.} \psalm{97} \V Adiuvabit eam.%typo: \V for Ant.
\newline {\color{Red} Lectio VII. Secundum Lucam.}
\lettrine[lines=2]{\zallmancaps \color{Red} I}{n} illo tempore: Rogabat Ihesum quidam Phariseus, ut manducaret cum illo. Et ingressus domum Pharisei: discubuit. Et reliqua.
\newline {\color{Red} Omelia Beati Gregorii pape. \grecross }
\lettrine[lines=2]{\zallmancaps \color{Blue} C}{ogitanti} michi de Marie penitentia: flere magis libet quam aliquid dicere. Cuius enim vel saxeum pectus ille huius peccatricis lacrime ad exemplum penitendi non emolliant? Consideravit namque quid fecit et noluit moderari quid faceret.
\begin{responsory}[o-certe-precipuus]
{O certe precipuus Marie Magdalene amor, * Que a monumento dominico discipulis recedentibus non recessit.}{Ardore caritatis succensa quem sepultum noverat sublatum credidit de sepulchro.}
\end{responsory}%
\newline {\color{Red} Lectio VIII.}
\lettrine[lines=2]{\zallmancaps \color{Red} S}{uper} convivantes ingressa est: non iussa venit, inter epulas lacrimas obtulit. Discite quo dolore ardet: que flere et inter epulas non erubescit. Sed ecce quia turpitudinis sue maculas aspexit, lavanda ad fontem misericordie cucurrit, convivantes non erubuit.
\begin{responsory}[o-felix-felicis]
{O felix felicis meriti Maria que resurgentem a mortuis dei filium videre meruisti prima, pro cuius amore: seculi contempsisti blandimenta, * Sedula nos apud ipsum quesumus prece commenda.}{Ut tecum mereamur o domina perfrui felicissima ipsius presentia.}
\end{responsory}%
\newline {\color{Red} Lectio IX.}
\lettrine[lines=2]{\zallmancaps \color{Blue} N}{am} quia semetipsam graviter erubescebat intus: nichil esse credidit quod verecundaretur foris. Quid igitur miramur fratres? Mariam venientem, an dominum suscipientem? Suscipientem dicam an trahentem? Dicam melius trahentem et suscipientem.
\begin{responsory-doxology}[gloriosa-diei]
{Gloriosa diei huius sollemnia devota celebremus omnes veneratione, qua beata Maria Magdalene, * Celos meruit ascendere.}{Consummato cursu presentis vite, celesti evocata pietate.}
\end{responsory-doxology}%
\V Optimam partem.
\newline {\color{Red} In Laudibus. Ant.} Laudibus excelsis omnis mundus exultet in sollemnitate sancte Marie Magdalene. {\color{Red} Ant.} Hec corde compuncta venit, ad pedes Ihesu accessit, illos lacrimis rigavit, capillis capitis tersit. {\color{Red} Ant.} Intercede supplicans assidue pro nobis Ihesu domino Maria Magdalene. {\color{Red} Ant.} O bone Ihesu laus tibi, remisisti peccatrici multa peccamina quia te dilexit multum. {\color{Red} Ant.} Maria ergo unxit pedes Ihesu et extersit capillis suis, et domus impleta est ex odore unguenti. {\color{Red} Capitulum.} Mulierem. {\color{Red} Ymnus.}
\hymn{eterni-patris-unice}{Eterni patris unice}{
\lettrine[lines=2]{\zallmancaps \color{Red} E}{terni} patris unice
nos pio vultu respice,
qui Magdalenam hodie
vocas ad thronum glorie.
\par {\bfseries \color{Blue} I}n thesauro reposita
regis est dragma perdita,
gemmaque lucet inclita
de luto luci reddita.
\par {\bfseries \color{Red} I}hesu dulce refugium,
spes una penitentium,
per peccatricis meritum
peccati solve debitum.
\par {\bfseries \color{Blue} P}ia mater et humilis,
nature memor fragilis,
in huius vite fluctibus
nos rege tuis precibus.
\par {\bfseries \color{Red} U}ni deo sit gloria
pro multiformi gratia,
qui culpas et supplicia
remittit, et dat premia. Amen.
}
\newline \V Dimissa sunt ei peccata multa. \R Quoniam dilexit multum.
\newline {\color{Red} Ad Ben. Ant.} O mundi lampas et margarita prefulgida que resurrectionem Christi nunciando apostolorum apostola fieri meruisti, Maria Magdalene semper pia exoratrix pro nobis assis ad deum, qui te elegit.
\newline {\color{Red} Ad Terciam. Cap.} Mulierem. \R \hyperlink{diffusa-est-gratia}{Diffusa est gratia.} \V Specie tua.
\newline {\color{Red} Ad Sextam. Cap.}
\lettrine[lines=2]{\zallmancaps \color{Blue} A}{ccinxit} fortitudine lumbos suos et roboravit brachium suum, gustavit et vidit quia bona est negociatio eius: non extinguetur in nocte lucerna eius.
\R \hyperlink{specie-tua-et}{Specie tua.} \V Adiuvabit eam.
\newline {\color{Red} Ad Nonam. Cap.}
\lettrine[lines=2]{\zallmancaps \color{Red} M}{ulte} filie congregaverunt divitias, tu supergressa es universas.
\R \hyperlink{adiuvabit-eam-deus}{Adiuvabit eam.} \V Dimissa sunt ei.
\newline {\color{Red} Ad Vesperas. Super psalmos. Ant.} Laudibus. \ul{Psalmi unius virginis. Cap., Ymnus, \Vbar , ut in Primis Vesperis.}
\newline {\color{Red} Ad Mag. Ant.} In diebus illis mulier que erat in civitate peccatrix ut cognovit quod Ihesus accubuisset in domo Symonis leprosi attulit alabastrum unguenti, et stans retro secus pedes domini Ihesu lacrimis cepit rigare pedes eius et capillis capitis sui tergebat et osculabatur pedes eius, et unguento ungebat.
{\ubsubsection{Sancti Apollinaris Episcopi et Martiris}{sancti-apollinaris-episcopi-et-martiris-breviarium}}
\newline {\color{Red} Oratio.}
\lettrine[lines=2]{\zallmancaps \color{Blue} D}{eus} qui nos beati Apollinaris martyris tui atque pontificis annua sollemnitate letificas, concede propicius ut cuius natalicia colimus, de eiusdem eciam protectione gaudeamus. Per.
{\ubsubsection{Sancte Christine Virginis et Martiris}{sancte-christine-virginis-et-martiris-breviarium}}
\newline {\color{Red} Oratio.}
\lettrine[lines=2]{\zallmancaps \color{Red} I}{ndulgentiam} nobis domine beata Christina virgo et martyr imploret, que tibi grata semper extitit, et merito castitatis, et tue professione virtutis. Per.
{\ubsubsection{Sancti Iacobi Apostoli}{sancti-iacobi-apostoli-breviarium}}
\newline {\color{Red} Ad Vesperas. Super psalmos. Ant.} O beate Iacobe omnium corde, ore, laudande, o patrone singularis, amabilis, intercede pro nobis ad dominum.
\newline {\color{Red} Ad Mag. Ant.} O lux et decus Hyspanie sanctissime Iacobe qui inter apostolos primatum tenes primus eorum martyrio laureatus, o singulare presidium qui meruisti videre redemptorem nostrum adhuc mortalem in deitate transformatum, exaudi preces servorum tuorum et intercede pro nostra salute omniumque populorum.
\newline {\color{Red} Oratio.}
\lettrine[lines=2]{\zallmancaps \color{Blue} E}{sto} domine plebi tue sanctificator et custos, ut apostoli tui Iacobi munita presidiis, et conversatione tibi placeat et secura deserviat. Per.
\newline \ul{Memoria de sanctis Christoforo et Cucufate.} {\color{Red} Oratio.}
\lettrine[lines=2]{\zallmancaps \color{Red} P}{resta} quesumus omnipotens deus ut qui gloriosos martyres Christoforum et Cucufatem fortes in sua confessione cognovimus, pios apud te in nostra intercessione sentiamus. Per.
\newline {\color{Red} Ad Matutinas. Lectio VII. Secundum Matheum.}
\lettrine[lines=2]{\zallmancaps \color{Blue} I}{n} illo tempore: Accessit ad Ihesum mater filiorum Zebedei cum filiis suis: adorans et petens aliquid ab eo. Et reliqua.
\newline {\color{Red} Omelia Beati Hieronimi Presbiteri.} \grecross
\lettrine[lines=2]{\zallmancaps \color{Red} U}{nde} opinionem regni habet mater filiorum Zebedei, ut cum dominus dixerit filius hominis tradetur principibus sacerdotum et scribis, et condemnabunt eum morte, et tradent gentibus ad illudendum, et flagellandum, et crucifigendum: et ignominiam passionis timentibus discipulis nuntiaverit, illa gloriam postulet triumphantis?
\newline {\color{Red} Lectio VIII.}
\lettrine[lines=2]{\zallmancaps \color{Blue} H}{ac} ut reor causa, quia post omnia dixerat dominus, et tercia die resurget, putavit eum mulier post resurrectionem ilico regnaturum, et quod in secundo adventu promittitur, primo esse complendum, et aviditate feminea presentia cupit, immemor futurorum.
\newline {\color{Red} Lectio IX.}
\lettrine[lines=2]{\zallmancaps \color{Red} P}{ostulat} autem mater filiorum Zebedei errore muliebri, et pietatis affectu, nesciens quid peteret. Nec mirum si ista arguatur imperitie, cum de Petro dicatur quando tria vult facere tabernacula, nesciens quid diceret. Respondens autem Ihesus: dixit eis. Nescitis quid petatis. Mater postulat, sed dominus discipulis loquitur, intelligens preces eius ex filiorum descendere voluntate.
{\ubsubsection{Sanctorum Nazarii, Celsi et Pantaleonis Martirum}{sanctorum-nazarii-celsi-et-pantaleonis-martirum-breviarium}}
\newline {\color{Red} Oratio.}
\lettrine[lines=2]{\zallmancaps \color{Blue} D}{eus} qui nos concedis sanctorum martyrum tuorum Nazarii Celsi, et Pantaleonis natalicia colere, da nobis in eterna beatitudine de eorum societate gaudere. Per.
{\ubsubsection{Sanctorum Felicis Simplicii Faustini et Beatricis Martirum}{sanctorum-felicis-simplicii-faustini-et-beatricis-martirum-breviarium}}
\newline {\color{Red} Oratio.}
\lettrine[lines=2]{\zallmancaps \color{Red} P}{resta} quesumus omnipotens deus, ut sicut populus Christianus martyrum tuorum Felicis, Simplicii, Faustini et Beatricis temporali sollemnitate congaudet, ita perfruatur eterna, et quod votis celebrat comprehendat effectu. Per.
{\ubsubsection{Sanctorum Abdon et Sennen Martirum}{sanctorum-abdon-et-sennen-martirum-breviarium}}
\newline {\color{Red} Oratio.}
\lettrine[lines=2]{\zallmancaps \color{Blue} D}{eus} qui sanctis tuis Abdon et Sennen ad hanc gloriam veniendi, copiosum munus gratie contulisti, da famulis tuis suorum veniam peccatorum, ut sanctorum tuorum intercedentibus meritis ab omnibus mereamur adversitatibus liberari. Per.
{\ubsubsection{Sancti Germani Episcopi et Confessoris}{sancti-germani-episcopi-et-confessoris-breviarium}}
\newline {\color{Red} Oratio.}
\lettrine[lines=2]{\zallmancaps \color{Red} E}{xaudi} domine preces nostras quas in sancti Germani confessoris tui atque pontificis sollemnitate deferimus, ut qui tibi digne meruit famulari eius intercedentibus meritis ab omnibus nos absolvas peccatis. Per.
{\ubsubsection{Ad Vincula Sancti Petri}{ad-vincula-sancti-petri-breviarium}}
\newline {\color{Red} Ad Vesperas. Cap.} Petrus quidem. {\color{Red} Ymnus.} \hyperlink{iam-bone-pastor}{Iam bone.} \ul{quere in Cathedra eiusdem.} \V Tu es Petrus.
\newline {\color{Red} Ad Mag. Ant.} Tu es pastor ovium princeps apostolorum, tibi tradite sunt claves regni celorum. {\color{Red} Oratio.}
\lettrine[lines=2]{\zallmancaps \color{Blue} D}{eus} qui beatum Petrum apostolum a vinculis absolutum illesum abire fecisti, nostrorum quesumus absolve vincula peccatorum, et omnia mala a nobis propiciatus exclude. Per.
{\color{Red} Eodem die Sanctorum Machabeorum Martirum. Oratio.}
\lettrine[lines=2]{\zallmancaps \color{Red} F}{raterna} nos domine martyrum tuorum corona letificet, que et fidei nostre prebeat incrementa virtutum, et multiplici nos suffragio consoletur. Per.
\newline {\color{Red} Ad Mat. Invit.}
\begin{invitatory}
{Tu es pastor ovium princeps apostolorum, * Tibi tradidit deus claves regni celorum.}
\end{invitatory}
\newline {\color{Red} Ymnus.} \hyperlink{iam-bone-pastor}{Iam bone.}
\newline \ul{Antiphone, Psalmi, Versiculi nocturnorum de communi apostolorum, Responsoria vero et versiculus ante Laudes, antiphone Laudibus, Cap., et antiphona Ad Ben. sicut in passione eiusdem. In Laudibus: Ymnus.} \hyperlink{eterna-christi-munera-apostolorum}{Eterna.} \ul{\Vbar .} Annunciaverunt. \ul{Antiphone, ad horas et capitula sicut in passione eiusdem. Ad Terciam Sextam et Nonam, Responsoria, et \Vbar , de communi apostolorum. Ad Vesperas super psalmos antiphona.} Petrus et Iohannes. \ul{Psalmi de apostolis. Cap., Ymnus, \Vbar , ut in Primis Vesperis.}
\newline {\color{Red} Ad Mag. Ant.} Angelus domini astitit, et lumen refulsit in habitaculo carceris percussoque latere Petri excitavit eum dicens, surge velociter, et continuo ceciderunt catene de manibus eius.
{\ubsubsection{Sancti Stephani Pape et Martiris}{sancti-stephani-pape-et-martiris-breviarium}}
\newline {\color{Red} Oratio.}
\lettrine[lines=2]{\zallmancaps \color{Blue} D}{eus} qui nos beati Stephani martyris tui atque pontificis annua sollemnitate letificas concede propicius, ut cuius natalicia colimus de eiusdem eciam protectione gaudeamus. Per.
{\ubsubsection{In Inventione Sancti Stephani Prothomartiris}{in-inventione-sancti-stephani-prothomartiris-breviarium}}
\newline \ul{Capitula sicut in alio festo ipsius.} {\color{Red} Ad Mag. Ant.} Ostendit sanctus Gamaliel per visum Luciano sacerdoti tres calathos aureos rosis refertos et quartum argenteum croco plenum et dixit, hi sunt nostri loculi et nostre reliquie, hic autem sanguineas habens rosas loculus est sancti Stephani qui solus ex nobis martyrio meruit coronari. {\color{Red} Oratio.}
\lettrine[lines=2]{\zallmancaps \color{Red} D}{a} nobis quesumus domine imitari quod colimus, ut discamus et inimicos diligere quia eius corporis inventionem celebramus, qui novit eciam pro persecutoribus exorare: Dominum nostrum Ihesum.
\newline {\color{Red} Ad Matutinas. Invit.}
\begin{invitatory}[adoremus-regem-magnum-invitatorium]
{Adoremus regem magnum dominum, * Qui in sanctis suis semper est mirabilis.}
\end{invitatory}
\newline \ul{Lectiones de communi unius martyris trium lectionum, et fiant de qua habet lectiones tres.}
\newline {\color{Red} Nonum.} \R Sancte dei preciose prothomartyr Stephane, qui virtute caritatis circumfultus undique dominum pro inimico exorasti populo, * Funde preces * Pro devoto tibi nunc collegio. \V Ut tuo propiciatus interventu dominus nos purgatos a peccatis iungat celi civibus. Funde. Gloria. Pro devoto.
\newline {\color{Red} Ad Ben. Ant.} Ex odoris mira fragantia sanitas egrotis emanavit maxima, septuaginta namque et tres tunc a variis langoribus curati sunt homines in inventione corporis Stephani dilectissimi martyris deum benedicentes.
\newline {\color{Red} In Secundis Vesperis. Ad Mag. Ant.} Hodie sanctus Iohannes pontifex cum maximo cleri plebisque tripudio preciosas prothomartyris Stephani reliquias, in sanctam transtulit ecclesiam Syon, ubi quondam archidiaconi functus est officio, cuius pia apud deum sit pro nobis quesumus intercessio. \ul{Cetera omnia de communi unius martyris.}
{\ubsubsection{In Festo Beati Dominici}{in-festo-beati-dominici-breviarium}}
\newline {\color{Red} Ad Vesperas. Super psalmos. Ant.} Gaude felix parens Hyspania nove prolis dans mundo gaudia, sed tu magis plaude Bononia tanti patris dotata gloria, lauda tota mater ecclesia nove laudis agens sollemnia. {\color{Red} Psalmi.} \psalm{112} \ul{cum ceteris.} {\color{Red} Cap.}
\lettrine[lines=2]{\zallmancaps \color{Blue} Q}{uasi} stella matutina in medio nebule, et quasi luna plena in diebus suis et quasi sol refulgens: sic iste effulsit in templo dei. \R \hyperlink{granum-excussum}{Granum.} {\color{Red} VI. Ymnus.}
\hymn{gaude-mater-ecclesia}{Gaude mater ecclesia}{
\lettrine[lines=2]{\zallmancaps \color{Red} G}{aude} mater ecclesia
letam agens memoriam,
que nove prolis gaudia
mittis ad celi curiam.
\par {\bfseries \color{Blue} P}redicatorum ordinis
dux et pater Dominicus
mundi iam fulget terminis
civis effectus celicus.
\par {\bfseries \color{Red} C}arnis liber ergastulo
celi potitur gloria,
pro paupertatis cingulo
stola dotatur regia.
\par {\bfseries \color{Blue} F}ragrans odor de tumulo
cum virtutum frequentia,
clamat pro Christi famulo
summi regis magnalia.
\par {\bfseries \color{Red} T}rino deo et simplici,
laus, honor, virtus, gloria:
qui nos prece Dominici
ducat ad celi gaudia. Amen.
}
\newline \V Ora pro nobis beate Dominice.
\newline {\color{Red} Ad Mag. Ant.} Transit pauper ad regni solium dux ad sceptrum victor ad premium, mors in vitam labor in ocium, presens cedit luctus in gaudium. {\color{Red} Oratio.}
\lettrine[lines=2]{\zallmancaps \color{Blue} D}{eus} \hypertarget{deus-qui-ecclesiam}{\label{deus-qui-ecclesiam}} qui ecclesiam tuam beati Dominici confessoris tui illuminare dignatus es meritis et doctrinis, concede ut eius intercessione temporalibus non destituatur auxiliis et spiritualibus semper proficiat incrementis. Per.
\newline {\color{Red} Ad Mat. Invit.}
\begin{invitatory}
{Assunt Dominici leta sollemnia, * Laude multiplici plaudat ecclesia.}
\end{invitatory}
\newline {\color{Red} Ymnus.}
\hymn{novus-athleta-domini}{Novus athleta domini}{
\lettrine[lines=2]{\zallmancaps \color{Red} N}{ovus} athleta domini
collaudetur Dominicus,
qui rem conformat nomini
vir factus evangelicus.
\par {\bfseries \color{Blue} C}onservans sine macula
virginitatis lilium,
ardebat quasi facula
pro zelo pereuntium.
\par {\bfseries \color{Red} M}undum calcans sub pedibus
manum misit ad fortia:
nudus occurrens hostibus,
Christi suffultus gratia.
\par {\bfseries \color{Blue} P}ugnat verbo, miraculis,
missis per orbem fratribus,
crebros adiungens sedulis
fletus orationibus.
\par {\bfseries \color{Red} T}rino deo et simplici,
laus, honor, virtus, gloria:
qui nos prece Dominici
ducat ad celi gaudia. Amen.
}
\newline {\color{Red} In Primo Nocturno. Ant.} Preco novus et celicus missus in fine seculi, pauper fulsit Dominicus forma previsus catuli. {\color{Red} Ps.} \psalm{1} {\color{Red} Ant.} Florem pudicicie servans illibatum attigit eximie vite celibatum. {\color{Red} Ps.} \psalm{2} {\color{Red} Ant.} Documentis artium eruditus satis transiit ad studium summe veritatis. {\color{Red} Ps.} \psalm{3} \ul{Lectiones de communi unius confessoris doctoris novem lectionum.} \R
\lettrine[lines=2]{\zallmancaps \color{Blue} M}{undum} \hypertarget{mundum-vocans}{\label{mundum-vocans}} vocans ad agni nuptias hora cene paterfamilias servum mittit, * Promittens varias vite delicias. \V Ad hoc convivium tam permagnificum elegit nuntium, sanctum Dominicum. Promittens.
\begin{responsory}[datum-mundo-pro]
{Datum mundo pro mundi gloria, mira Christi presignat gratia, * Cuius ortum precurrunt nuncia veri presagia.}{Stella micans in fronte parvuli novum iubar premonstrat seculi.}
\end{responsory}%
\begin{responsory-final}[verbum-vite]
{Verbum vite dum palam promitur surgunt hostes liber conscribitur favent omnes, * Sic error vincitur, * Fides extollitur.}{Ter in flammas libellus traditus ter exivit illesus penitus.}{Fides}
\end{responsory-final}%
\newline {\color{Red} In Secundo Nocturno. Ant.} Sub Augustini regula mente profecit sedula tandem virum canonicum auget in apostolicum. {\color{Red} Ps.} \psalm{4} {\color{Red} Ant.} Agonizans pro Christi nomine, mundum replet divino semine paupertatis degens sub tegmine. {\color{Red} Ps.} \psalm{5} {\color{Red} Ant.} Pernox cum Christo proprium non possidebat lectulum, post lacrimarum fluvium vix humi dans corpusculum. {\color{Red} Ps.} \psalm{8}
\begin{responsory}[paupertatis-ascendens]
{Paupertatis ascendens culmina clamat mundi detestans crimina, frangit hostes et fugat agmina, * Nulla sanctum frangunt discrimina.}{Nocte celi perlustrans limina, die terris dat verbi semina.}
\end{responsory}%
\begin{responsory}[panis-oblatus]
{Panis oblatus celitus fratrum supplet inopiam, * Viteque natus redditus matris pellit tristiciam.}{Signo crucis obedit pluvia lingua verba transformat varia.}%typo maria for varia
\end{responsory}%
\begin{responsory-final}[granum-excussum]
{Granum excussum palea nexu soluto luteo, de paupertatis area celi locatur horreo, * Cum mercede virginea, * Doctorum fulgens cuneo.}{Flos in florum virens areola, bina gaudet sortis aureola.}{Doctorum}
\end{responsory-final}%
\newline {\color{Red} In Tertio Nocturno. Ant.} Sitiebat servus Christi martyrium sicut sitit cervus ad aque fluvium. {\color{Red} Ps.} \psalm{14} {\color{Red} Ant.} Migrans pater filiis vite firmamentum paupertatis humilis condit testamentum. {\color{Red} Ps.} \psalm{20} {\color{Red} Ant.} Liber carnis vinculo celum introivit, ubi pleno poculo gustat quod sitivit. {\color{Red} Ps.} \psalm{23}
\newline {\color{Red} Lectio VII. Secundum Matheum.}
\lettrine[lines=2]{\zallmancaps \color{Red} I}{n} illo tempore: dixit Ihesus discipulis suis. Vos estis sal terre. Quod si sal evanuerit in quo salietur? Et reliqua.
\newline {\color{Red} Omelia Beati Augustini Episcopi.}
% https://la.wikisource.org/wiki/De_Sermone_Domini_in_monte
\lettrine[lines=2]{\zallmancaps \color{Blue} H}{ic} ostendit dominus fatuos esse iudicandos qui temporalium bonorum vel copiam sectantes, vel inopiam metuentes, amittunt eterna, que nec dari possunt ab hominibus nec auferri. Itaque si sal infatuatum fuerit in quo salietur? Idest si vos per quos condiendi sunt quodammodo populi, metu temporalium persecutionum amiseritis regna celestia: qui erunt homines per quos vobis error auferatur, cum vos elegerit Deus per quos errorem auferat ceterorum?
\begin{responsory}[felix-vitis-de]
{Felix vitis de cuius surculo tantum germen redundat seculo, * Celi vinum propinans populo vitali poculo.}{Ex ubertate palmitum mundi iam cinxit ambitum.}
\end{responsory}%
\newline {\color{Red} Lectio VIII.}
\lettrine[lines=2]{\zallmancaps \color{Red} E}{rgo} ad nichilum valet sal infatuatum nisi ut mittatur foras, et conculcetur ab hominibus. Non itaque calcatur ab hominibus qui patitur persecutionem, sed qui persecutionem timendo infatuatur. Calcari enim non potest nisi inferior. Sed
%pp128
inferior non est, qui quamvis corpore multa in terra sustineat, corde tamen in celo fixus est.
\begin{responsory}[ascendenti-de-valle]
{Ascendenti de valle lubrici mundi chori plaudunt angelici, * Ihesu bone prece Dominici, tibi presta nos gratos effici.}{Per quem multos a morte suscitas, penas nobis relaxa debitas.}
\end{responsory}%
\newline {\color{Red} Lectio IX.}
\lettrine[lines=2]{\zallmancaps \color{Blue} V}{os} estis lumen mundi. Quomodo dixit superius sal terre, sic nunc dicit lumen mundi. Nam neque superius ista terra accipienda est quam pedibus calcamus corporeis, sed homines qui in terra habitant vel eciam peccatores, quorum condiendis et extinguendis fetoribus apostolicum salem dominus misit.
\begin{responsory-doxology}[o-spem-miram]
{O spem miram quam dedisti mortis hora te flentibus dum post mortem promisiti te profuturum fratribus, * Imple pater quod dixisti nos tuis iuvans precibus.}{Qui tot signis claruisti in egrorum corporibus nobis opem ferens Christi egreis medere moribus.}
\end{responsory-doxology}%
\newline {\color{Red} In Laudibus. Ant.} Adest dies leticie quo beatus Dominicus aulam celestis curie civis intrat magnificus. {\color{Red} Ant.} Pauper in peculio dives vita pura paupertatis precio celi tenet iura. {\color{Red} Ant.} Scala celo prominens fratri revelatur per quam pater transiens sursum ferebatur. {\color{Red} Ant.} Lingue manus studio scalam hanc erexit, quam virgo cum filio mater sursum vexit. {\color{Red} Ant.} Fulget in choro virginum doctor veritatis, sertum honoris geminum gerens cum beatis. {\color{Red} Capitulum.} Quasi stella. {\color{Red} Ymnus.}
\hymn{ymnum-nove-leticie}{Ymnum nove leticie}{
\lettrine[lines=2]{\zallmancaps \color{Red} Y}{mnum} nove leticie
dulci productum cantico,
noster depromat hodie
chorus sancto Dominico.
\par {\bfseries \color{Blue} V}ergente mundi vespere
novum sidus exoritur,
et clausis culpe carcere
preco salutis mittitur.
\par {\bfseries \color{Red} D}octrinam evangelicam
spargens per orbis cardinem,
pestem fugat hereticam,
novum producens ordinem.
\par {\bfseries \color{Blue} H}ic est fons ille modicus
crescens in flumen maximum,
qui mundo iam mirificus
potum largitur optimum.
\par {\bfseries \color{Red} T}rino deo et simplici,
laus, honor, virtus, gloria:
qui nos prece Dominici
ducat ad celi gaudia. Amen.
}
\newline \V Iustus germinabit.
\newline {\color{Red} Ad Ben. Ant.} Carnis vigor mentis martyrium, mundum replens lingue preconium, pauper Christi post cursus stadium vite tibi mercantur premium.
\newline {\color{Red} Ad Terciam. Cap.} Quasi stella.
\newline {\color{Red} Ad Sextam. Cap.}
\lettrine[lines=2]{\zallmancaps \color{Blue} S}{piritus} meus qui est in te, et verba mea que posui in ore tuo non recedent de ore tuo et de ore seminis tui dicit dominus, amodo et usque in sempiternum.
\newline {\color{Red} Ad Nonam. Cap.}
\lettrine[lines=2]{\zallmancaps \color{Red} L}{ex} veritatis fuit in ore eius, et iniquitas non est inventa in labiis eius, in pace et equitate ambulavit mecum et multos avertit ab iniquitate.
\newline \ul{Ad Vesperas. Cap., Ymnus, \Vbar , ut in Primis Vesperis.} {\color{Red} Ad Mag. Ant.} O lumen ecclesie doctor veritatis, rosa patientie ebur castitatis aquam sapientie propinasti gratis, predicator gratie nos iunge beatis.
\newline \ul{Per octavas quando de ipsis agitur, et in octava die. Invitat. sicut in die. Sequentes vero antiphone dicantur cotidie infra octavas, Ad Ben. et ad Mag. vel ad memoriam.}
\newline {\color{Red} Ad Ben. Ant.} Benedictus redemptor omnium qui saluti providens hominum mundo dedit sanctum Dominicum.
\newline {\color{Red} Ad Mag. Ant.} Magne pater sancte Dominice mortis hora nos tecum suscipe et hic semper nos pie respice.
{\ubsubsection{Sanctorum Sixti Felicissimi Et Agapiti Martirum}{sanctorum-sixti-felicissimi-et-agapiti-martirum-breviarium}}
\newline {\color{Red} Oratio.}
\lettrine[lines=2]{\zallmancaps \color{Blue} D}{eus} qui nos concedis sanctorum martirum tuorum Sixti Felicissimi et Agapiti natalicia colere, da nobis in eterna beatitudine de eorum societate gaudere. Per.
{\ubsubsection{Sancti Donati Episcopi et Martiris}{sancti-donati-episcopi-et-martiris-breviarium}}
\newline {\color{Red} Oratio.}
\lettrine[lines=2]{\zallmancaps \color{Red} D}{eus} tuorum gloria sacerdotum presta quesumus, ut sancti martyris tui et episcopi Donati cuius festa gerimus sentiamus auxilium. Per.
{\ubsubsection{Sancti Cyriaci Martiris Sociorumque Eius}{sancti-cyriaci-martiris-sociorumque-eius-breviarium}}
\newline {\color{Red} Oratio.}
\lettrine[lines=2]{\zallmancaps \color{Blue} D}{eus} qui nos annua beatorum martyrum tuorum Cyriaci Largi et Smaragdi sollemnitate letificas, concede propicius, ut quorum natalicia colimus, virtutem quoque passionis imitemur. Per.
{\ubsubsection{Sancti Laurentii Martiris}{sancti-laurentii-martiris-breviarium}}%typo: Martirum
\newline {\color{Red} Ad Vesperas. Super psalmos. Ant.} Laurentius bonum opus operatus est qui per signum crucis cecos illuminavit. \ul{Psalmi de Octava.} \R \hyperlink{beatus-laurentius-clamavit}{Beatus Laurentius clamavit.} {\color{Red} IIII. Ad Mag. Ant.} Beatus Laurentius dum in craticula suprapositus ureretur, ad impiissimum tyrannum dixit, assatum est, iam versa et manduca, nam facultates ecclesie quas requiris in celestes thesauros manus pauperum deportaverunt. {\color{Red} Oratio.}
\lettrine[lines=2]{\zallmancaps \color{Red} A}{desto} domine supplicationibus nostris, et intercessione beati Laurentii martiris tui, perpetuam nobis misericordiam benignus impende. Per.%typo benigus
\newline {\color{Red} Ad Matutinas. Invit.}
\begin{invitatory}%[regem-sempiternum-invitatorium]
{Regem sempiternum venite adoremus, * Qui martyrem suum digne pro meritis coronavit in celis.}
\end{invitatory}
\newline {\color{Red} In Primo Nocturno. Ant.} Quo progrederis sine filio pater, quo sacerdos sancte sine ministro properas. {\color{Red} Ps.} \psalm{1} \V Beatus Laurentius dixit. {\color{Red} Ant.} Noli me derelinquere pater sancte, quia thesauros tuos iam expendi quos tradidisti michi. {\color{Red} Ps.} \psalm{2} \V Quid in me ergo displicuit paternitati tue. {\color{Red} Ant.} Non ego te desero fili, neque derelinquo sed maiora tibi debentur pro fide Christi certamina. {\color{Red} Ps.} \psalm{3} \V Beatus Sixtus dixit. \R
\lettrine[lines=2]{\zallmancaps \color{Blue} L}{evita} \hypertarget{levita-laurentius-bonum}{\label{levita-laurentius-bonum}} Laurentius bonum opus operatus est, qui per signum crucis cecos illuminavit, * Et thesauros ecclesie dedit pauperibus. \V Dispersit dedit pauperibus iusticia eius manet in seculum seculi. Et.
\begin{responsory}[quo-progrederis]
{Quo progrederis sine filio pater, quo sacerdos sancte sine diacono properas, * Tu nunquam sine ministro sacrificium offerre consueveras.}{Quid in me ergo displicuit paternitati tue, nunquid degenerem me probasti, experire certe utrum ydoneum ministrum elegeris cui commisisti dominici sanguinis dispensationem.}
\end{responsory}%
\begin{responsory-final}[noli-me-derelinquere]
{Noli me derelinquere pater sancte, quia thesauros tuos iam expendi, * Non ego te desero fili neque derelinquo * Sed maiora tibi debentur pro fide Christi certamina.}{Nos quasi senes levioris pugne cursum recipimus, te autem quasi iuvenem manet gloriosior de tiranno triumphus, post triduum me sequeris sacerdotem levita.}{Sed maiora}
\end{responsory-final}%
\newline {\color{Red} In Secundo Nocturno. Ant.} Beatus Laurentius orabat dicens, domine Ihesu Christe deus de deo miserere michi servo tuo. {\color{Red} Ps.} \psalm{4} \V Quia accusatus non negavi nomen sanctum tuum, interrogatus te Christum confessus sum. {\color{Red} Ant.} Beatus Laurentius dixit, mea nox obscurum non habet sed omnia in luce clarescunt. {\color{Red} Ps.} \psalm{5} \V Quia ipse dominus novit quod accusatus non negavi nomen sanctum eius. {\color{Red} Ant.} Dixit Romanus ad beatum Laurentium, video ante te iuvenem pulcherrimum, festina me baptizare. {\color{Red} Ps.} \psalm{8} \V Afferens autem urceum cum aqua misit se ad pedes eius.
\begin{responsory}[beatus-laurentius-clamavit]
{Beatus Laurentius clamavit et dixit deum meum colo et ipsi soli servio, * Et ideo non timeo tormenta tua.}{Mea nox obscurum non habet sed omnia in luce clarescunt.}
\end{responsory}%
\begin{responsory}[strinxerunt-corporis]
{Strinxerunt corporis membra posita in craticula subicientibus primas insultat levita Christi, * Beate Laurenti martyr Christi intercede pro nobis.}{Carnifices vero urgentes ministrabant carbones subter cratem ferream.}
\end{responsory}%
\begin{responsory-doxology}[beatus-laurentius-dixit]
{Beatus Laurentius dixit, ego me obtuli sacrificium deo, * In odorem suavitatis.}{Gaudeo plane quia hostia Christi effici merui.}
\end{responsory-doxology}%
\newline {\color{Red} In Tertio Nocturno. Ant.} Strinxerunt corporis membra posita in craticula, subicientibus prunas insultat levita Christi. {\color{Red} Ps.} \psalm{10} \V Carnifices vero urgentes ministrabant carbones subter cratem ferream. {\color{Red} Ant.} Igne me examinasti, et non est inventa in me iniquitas. {\color{Red} Ps.} \psalm{16} \V Probasti domine cor meum et visitasti nocte. {\color{Red} Ant.} Interrogatus te dominum confessus sum assatus gratias ago. {\color{Red} Ps.} \psalm{20} \V Gratias tibi ago domine Ihesu Christe quia ianuas tuas ingredi merui.
\begin{responsory}[in-craticula-te]
{In craticula te deum non negavi, et ad ignem applicatus, * Te dominum Ihesum Christum confessus sum.}{Accusatus non negavi nomen sanctum tuum et interrogatus.}
\end{responsory}%
\begin{responsory}[gaudeo-plane-quia]
{Gaudeo plane quia hostia Christi effici merui, accusatus non negavi nomen sanctum eius, interrogatus Christum confessus sum, * Assatus gratias ago.}{Ego me obtuli sacrificium deo in odorem suavitatis.}
\end{responsory}%
\begin{responsory-doxology}[meruit-esse-hostia]
{Meruit esse hostia Christi levita Laurentius, qui dum assaretur non negavit dominum, * Et ideo inventus est sacrificium laudis.}{In craticula positus deum non negavit, et ad ignem applicatus Christum confessus est.}
\end{responsory-doxology}%
\newline \V Ora pro nobis beate Laurenti.
\newline {\color{Red} In Laudibus. Ant.} Laurentius ingressus est martyr et confessus est nomen domini Ihesu Christi. {\color{Red} Ant.} Laurentius bonum opus operatus est, qui per signum crucis cecos illuminavit. {\color{Red} Ant.} Adhesit anima mea post te quia caro mea igne cremata est pro te deus meus. {\color{Red} Ant.} Misit dominus angelum suum et liberavit me de medio ignis, et non sum estuatus. {\color{Red} Ant.} Beatus Laurentius orabat dicens gratias tibi ago domine quia ianuas tuas ingredi merui.
\newline {\color{Red} Ad Ben. Ant.} In craticula te deum non negavi et ad ignem applicatus te Christum confessus sum, probasti cor meum et visitasti nocte igne me examinasti et non est inventa in me iniquitas. {\color{Red} Oratio.}
\lettrine[lines=2]{\zallmancaps \color{Red} D}{a} nobis quesumus omnipotens deus vitiorum nostrorum flammas extinguere, qui beato Laurentio tribuisti tormentorum suorum incendia superare. Per.
\newline {\color{Red} Ad Mag. Ant.} Levita Laurentius bonum opus operatus est qui per signum crucis cecos illuminavit et thesauros ecclesie dedit pauperibus.
\newline \ul{Memoria de sancto Tyburcio.} {\color{Red} Ant.} Inclitus martyr Tyburcius cum nudatis plantis super prunas ardentes incederet, Fabiano prefecto dixit, videtur michi quod super roseos flores incedam in nomine domini mei Ihesu Christi. \V Posuisti. {\color{Red} Oratio.}
\lettrine[lines=2]{\zallmancaps \color{Blue} B}{eati} Tyburcii nos domine foveant continuata presidia, quia non desinis propicius intueri quos talibus auxiliis concesseris adiuvari. Per.
\newline \ul{Per octavas sancti Laurentii preterquam in vigilia et in festo Assumptionis fiat de eo memoria. In Laudibus per antiphonam.} Laurentius ingressus. \ul{\Vbar .} Corona aurea. \ul{In Vesperis per antiphonam.} Laurentius bonum. \ul{\Vbar .} Gloria et honore. \ul{In crastino sancti Laurentii officium fiat de octavis beati Dominici, et in Laudibus post memoriam de sancto Laurentio fiat memoria de sancto Tyburcio antiphona.} Qui vult. \ul{\Vbar .} Magna est.
{\ubsubsection{Sancti Ypoliti Martiris Sociorumque Eius}{sancti-ypoliti-martiris-sociorumque-eius-breviarium}}
\newline {\color{Red} Oratio.}
\lettrine[lines=2]{\zallmancaps \color{Red} D}{a} quesumus omnipotens deus, ut beati Ypoliti martyris tui sociorumque eius veneranda sollemnitas, et devotionem nobis augeat et salutem. Per.
{\ubsubsection{Sancti Eusebii Presbiteri et Confessoris}{sancti-eusebii-presbiteri-et-confessoris-breviarium}}
\newline {\color{Red} Oratio.}
\lettrine[lines=2]{\zallmancaps \color{Blue} D}{eus} qui nos beati Eusebii confessoris tui annua sollemnitate letificas: concede propicius, ut cuius natalicia colimus, per eius ad te exempla gradiamur. Per.
{\ubsubsection{In Vigilia Assumptionis Beate Marie}{in-vigilia-assumptionis-beate-marie-breviarium}}
\newline {\color{Red} Ad Mat. Invit.} \hyperlink{regem-martyrum-invitatorium}{Regem martyrum.} \ul{Super psalmos antiphona.} Quo progrederis. \ul{sine versu, psalmi de die sancti Laurentii, \Vbar , et tria responsoria de hystoria beati Laurenti secundum feriam. Tres lectiones de omelia sequenti.} \psalm{tedeum} \ul{\Vbar .} Ora pro nobis. \ul{In Laudibus habet sola antiphona.} Laurentius ingressus. \ul{Psalmi.} \psalm{92} \ul{Cap., Ymnus, \Vbar , et antiphona ad Ben. et Oratio et hore diei sicut in festo sancti Laurenti. Si autem hac die Dominica contigerit totum officium de sancto Laurentio fiat sicut in festo eius exceptis versibus antiphonarum qui non dicantur, et tres ultime lectiones fiant de omelia sequenti.}
\newline {\color{Red} Lectio Prima. Secundum Lucam.}
\lettrine[lines=2]{\zallmancaps \color{Red} I}{n} illo tempore: Loquente Ihesu ad turbas extollens vocem quedam mulier de turba: dixit illi. Beatus venter qui te portavit, et ubera que suxisti, et reliqua.
\newline {\color{Red} Omelia Beati Bede Presbiteri. \grecross }
\lettrine[lines=2]{\zallmancaps \color{Blue} U}{nigenitus} dei karissimi ante assumptam de intemerata virgine nostre mortalitatis substantiam, locutus est ad turbas per prophetas suos, locutus est per semetipsum in adventu suo, et post gloriosissimam in celum ascensionem: locutus est per discipulos. Factum est igitur in adventu suo cum ad turbas visibilis loqueretur, quedam mulier ventrem esse beatum qui eum portaverat et ubera que suxerat predicavit.
\newline {\color{Red} Lectio Secunda.}
\lettrine[lines=2]{\zallmancaps \color{Blue} A}{t} ille dixit. Quin immo beati qui audiunt verbum dei et custodiunt illud. Hoc in loco beate dei genitricis excellentia commendatur, quia in tanta speculationis custodia domino famulabatur, quod ipsius humilitati et sanctitati nullus sanctorum comparatur. Beatam igitur esse matrem suam salvator noster predicavit, que verbum dei audivit, et supra omnes huius vite mortales custodivit.
\newline {\color{Red} Lectio Tercia.}
\lettrine[lines=2]{\zallmancaps \color{Red} O}{}\ veneranda domina electa et preelecta, que mater est, et virgo esse non desinit. O admiranda puella que salvo virginitatis pudore suum genuit genitorem. O felix ancilla que creatorem et gubernatorem genuit universorum, et virgo permansit in eternum.
\newline {\color{Red} Ad Vesperas. Super psalmos. Ant.} Tota pulchra es amica mea et macula non est in te, favus distillans labia tua, mel et lac sub lingua tua, odor unguentorum tuorum super omnia aromata, iam enim hyems transiit, ymber abiit et recessit, flores apparuerunt, vinee florentes odorem dederunt, et vox turturis audita est in terra nostra, surge propera amica mea, veni de Libano veni coronaberis. {\color{Red} Psalmi.} \psalm{112} \ul{cum ceteris.} {\color{Red} Cap.}
\lettrine[lines=2]{\zallmancaps \color{Blue} I}{n} omnibus requiem quesivi et in hereditate domini morabor, tunc precepit et dixit michi creator omnium, et qui creavit me requievit in tabernaculo meo. \R \hyperlink{beata-es-virgo}{Beata es.} {\color{Red} IV. Ymnus.} \hyperlink{ave-maris-stella}{Ave maris stella.} \V Exaltata es sancta dei genitrix. \R Super choros angelorum ad celestia regna.
\newline {\color{Red} Ad Mag. Ant.} Ascendit Christus super celos et preparavit sue castissime matri immortalitatis locum, et hec est illa preclara festivitas omnium sanctorum festivitatibus incomparabilis, in qua gloriosa et felix mirantibus celestis curie ordinibus ad ethereum pervenit thalamum, quo pia sui memorum immemor nequaquam existat. {\color{Red} Oratio.}
\lettrine[lines=2]{\zallmancaps \color{Red} D}{eus} qui virginalem aulam beate Marie in qua habitares eligere dignatus es, da quesumus ut sua nos defensione munitos, iocundos faciat sue interesse festivitati. Qui vivis.
\newline {\color{Red} Ad Completorium. Super psalmos. Ant.} Virgo Maria non est tibi similis nata in mundo inter mulieres, florens ut rosa, fragrans sicut lilium, ora pro nobis sancta dei genitrix.
\newline {\color{Red} Ad Nunc Dimittis. Ant.} Sub tuum presidium confugimus sancta dei genitrix, nostras deprecationes ne despicias in necessitatibus, sed a periculis cunctis libera nos semper virgo benedicta.
\newline \ul{Predicte antiphone dicantur in Completorio per totas octavas, nisi in octava sancti Laurentii et in festo sancti Bernardi.}
{\ubsubsection{In Die Assumptionis Beate Marie}{in-die-assumptionis-beate-marie-breviarium}}
\newline {\color{Red} Ad Mat. Invitat.}
\begin{invitatory}[venite-adoremus-regem-cuius-invitatorium]
{Venite adoremus regem regum, * Cuius hodie ad ethereum mater virgo assumpta est celum.}
\end{invitatory}
\newline {\color{Red} Ymnus.} \hyperlink{quem-terra}{Quem terra.}
\newline {\color{Red} In Primo Nocturno. Ant.} Exaltata es sancta dei genitrix super choros angelorum ad celestia regna. {\color{Red} Ps.} \psalm{8} {\color{Red} Ant.} Paradisi porte per te nobis aperte sunt, quas hodie gloriosa cum angelis triumphas. {\color{Red} Ps.} \psalm{18} {\color{Red} Ant.} Ortus conclusus es dei genitrix, ortus conclusus, fons signatus, surge propera amica mea et veni. {\color{Red} Ps.} \psalm{23} \V Sancta dei genitrix.
\newline \ul{Ex sermone,} Inter precipuas. \ul{qui legitur in aliquibus ecclesiis, et intitulatur a quibusdam beati Gregorii.} {\color{Red} Lectio Prima.} \grecross
\lettrine[lines=2]{\zallmancaps \color{Blue} H}{odierna} karissimi nobis refulsit insignis sollemnitas, in qua beata virgo ad ethereum thalamum assumpta est, in quo rex regum ineffabili sedet solio, ubi milia milium ministrant angelorum, et decies milies centena assistunt milia, quorum cum laudibus, et exequiis, suscepta est in celo. Hec quidem sollemnitas ante mundi constitutionem preordinata, sed hodierna die expleta, nobis est annua sed superius civibus perpetua est et continua. \R
\lettrine[lines=2]{\zallmancaps \color{Red} V}{idi} \hypertarget{vidi-speciosam-sicut}{\label{vidi-speciosam-sicut}} speciosam sicut columbam ascendentem desuper rivos aquarum, cuius inestimabilis odor erat nimis in vestimentis eius, * Et sicut dies verni circumdabant eam flores rosarum et lilia convallium. \V Que est ista que ascendit per desertum sicut virgula fumi ex aromatibus mirre et thuris. Et sicut.
\newline {\color{Red} Lectio Secunda.}
\lettrine[lines=2]{\zallmancaps \color{Blue} H}{ec} est dies in qua gloriosa dei genitrix inter filias Hierusalem speciosa resplenduit, in qua quasi aurora processit, et vadit post agnum gloriosior quocumque ille pergit, in qua clarior sole refulsit in throno claritatis. Quoniam et si duodecim leguntur throni apostolorum non desunt tamen alii quam plurimi throni, in quibus viginti quatuor residere dicuntur seniores, unde constat quod thronus beatissime virginis hodie celsior glorificatur, et veneratur eciam ab angelis.
\begin{responsory}[sicut-cedrus]
{Sicut cedrus exaltata sum in Libano et sicut cypressus in monte Syon quasi mirra electa, * Dedi suavitatem odoris.}{Sicut cinamomum et balsamum aromatizans.}
\end{responsory}%
\newline {\color{Red} Lectio Tertia.}
\lettrine[lines=2]{\zallmancaps \color{Red} G}{lorificemus} igitur virginem quam hodie paradisus excepit gaudens, quam angeli cum laudibus prosequuntur, quam apostolorum chorus veneratur, quam martyres candidati beatificant, quam sanctorum confessorum stolatus concelebrat numerus, cui hodie sanctarum virginum cum suis palmis victricibus exultans occurrit exercitus, quoniam hec est per quam omnis maledictio soluta est, et celestis benedictio in totum venit mundum. Hec est in cuius utero ecclesia dei coniuncta est federe sempiterno.
\begin{responsory-doxology}[que-est-ista]
{Que est ista que processit sicut sol et formosa tanquam Hierusalem, viderunt eam filie Syon, et beatam dixerunt, * Et regine laudaverunt eam.}{Et sicut dies verni circumdabant eam flores rosarum et lilia convallium.}
\end{responsory-doxology}%
\newline {\color{Red} In Secundo Nocturno. Ant.} Emissiones tue paradisus malorum punicorum cum pomorum fructibus. {\color{Red} Ps.} \psalm{44} {\color{Red} Ant.} Fons ortorum puteus aquarum viventium que fluunt impetu de Libano. {\color{Red} Ps.} \psalm{45} {\color{Red} Ant.} Veniat dilectus meus in ortum suum, ut comedat fructum pomorum suorum. {\color{Red} Ps.} \psalm{86} \V Post partum virgo.
\newline {\color{Red} Lectio Quarta.}
\lettrine[lines=2]{\zallmancaps \color{Blue} V}{irgo} hec refulget clarissima inter virgines, candidior inter martyres. Iure enim plusquam martyr est que nimio amore vulnerata testis extitit salvatoris, et pre merore in animo cruciatum sustinuit passionis. Quia igitur ineffabilibus est donis impresentiarum munerata, et privilegiis sublimata divinis, credere oportet quod hodie omnimoda celesti ita illustratur gloria, ut nullis digne possit in terris laudibus venerari.
\begin{responsory}[super-salutem-assumptionis]
{Super salutem et omnem pulchritudinem dilecta es a domino, et regina celorum vocari digna es, * Gaudent chori angelorum consortes et concives nostri.}{Exaltata es sancta dei genitrix super choros angelorum ad celestia regna.}
\end{responsory}%
\newline \ul{Ex sermone,} Adest nobis dilectissimi. \ul{qui in quibusdam libris attribuitur Augustino.} {\color{Red} Lectio Quinta.}
\lettrine[lines=2]{\zallmancaps \color{Red} E}{t} quidem sunt nonnulla sine autoris nomine de eius assumptione conscripta, que ita caventur, ut ad confirmandam rei veritatem legi minime permittantur. Hinc sane pulsantur non nulli, quia nec corpus eius in terra invenitur, nec assumptio eius cum carne ut in apocrifa legitur, in catholica hystoria reperitur.
\begin{responsory}[ista-est-speciosa]
{Ista est speciosa inter filias Hierusalem, sicut vidistis eam plenam caritate et dilectione, * In cubilibus et in ortis aromatum.}{Ista est que ascendit de deserto deliciis affluens.}
\end{responsory}%
\newline {\color{Red} Lectio Sexta.}
\lettrine[lines=2]{\zallmancaps \color{Blue} D}{icendum} est istis, quia si Moysi corpus ab homine non invenitur in terris cum quo locutus est deus facie ad faciem, illius querere dementie est: per quam idem maiestatis deus incarnatus effulsit in terris. Neque enim dignum est de corporis eius noticia sollicitum quempiam esse, quam non dubitat super angelos elevatam cum Christo regnare. Sufficere debet tantum noticie humane hanc vere fateri reginam celorum, pro eo quod regem peperit angelorum.
\begin{responsory-doxology}[beata-es-virgo]
{Beata es virgo Maria dei genitrix que credidisti domino perfecta sunt in te que dicta sunt tibi, ecce exaltata es super choros angelorum, * Intercede pro nobis ad dominum Ihesum Christum.}{Benedicta et venerabilis es virgo Maria, que sine tactu pudoris inventa es mater salvatoris.}%per facta
\end{responsory-doxology}%
\newline {\color{Red} In Tertio Nocturno. Ant.} Veni in ortum meum soror mea sponsa messui mirram meam cum aromatibus meis. {\color{Red} Ps.} \psalm{95} {\color{Red} Ant.} Comedi favum cum melle meo, bibi vinum meum cum lacte meo. {\color{Red} Ps.} \psalm{96} {\color{Red} Ant.} Talis est dilectus meus et ipse est amicus meus filie Hierusalem. {\color{Red} Ps.} \psalm{97} \V Speciosa facta es.
\newline {\color{Red} Lectio VII. Secundum Matheum.}
\lettrine[lines=2]{\zallmancaps \color{Red} I}{n} illo tempore: Intravit Ihesus in quoddam castellum, et mulier quedam Martha nomine excepit illum in domum suam. Et huic erat soror nomine Maria. Et reliqua.
\newline {\color{Red} Omelia Venerabilis Bede Presbiteri. \grecross .}
\lettrine[lines=2]{\zallmancaps \color{Blue} D}{ue} iste domino dilecte sorores duas vitas spiritales quibus in presenti sancta exercetur ecclesia designant, Martha quidem actualem qua proximo in caritate sociamur: Maria vero contemplativam qua in Dei amore suspiramus.
\begin{responsory}[hodie-maria-virgo]
{Hodie Maria virgo celos ascendit, * Gaudere quia cum Christo regnat in eternum.}{Regina mundi hodie de seculo nequam eripitur.}
\end{responsory}%
\newline {\color{Red} Lectio VIII.}
\lettrine[lines=2]{\zallmancaps \color{Red} A}{ctiva} enim vita est: panem esurienti tribuere, verbo sapientie nescientem docere, errantem corrigere, ad humilitatis viam superbientem proximum revocare: infirmantis curam gerere, que singulis quibusque expediant dispensare, et commissis nobis qualiter subsistere valeant providere.
\begin{responsory}[beatam-me-dicent]
{Beatam me dicent omnes generationes, quia fecit michi dominus magna, * Quia potens est et sanctum nomen eius.}{Magnificat anima mea dominum et exultavit spiritus meus in deo salutari meo.}
\end{responsory}%
\newline {\color{Red} Lectio IX.}
\lettrine[lines=2]{\zallmancaps \color{Blue} C}{ontemplativa} vero vita est, caritatem quidem Dei et proximi mente retinere, sed ab exteriori actione quiescere, soli desiderio conditoris inherere, ut nil iam agere libeat, sed calcatis curis omnibus ad videndam faciem sui creatoris animus inardescat, ita ut iam noverit carnis corruptibilis pondus cum merore portare, totisque desideriis appetere illis ymnidicis angelorum choris interesse, admisceri civibus celestibus, de eterna in conspectu Dei incorruptione gaudere.
\begin{responsory-final}[felix-namque]
{Felix namque es sacra virgo Maria et omni laude dignissima, * Quia ex te ortus est sol iusticie * Christus deus noster.}{Ora pro populo interveni pro clero, intercede pro devoto femineo sexu, sentiant omnes tuum iuvamen quicumque celebrant tuam assumptionem.}{Christus}
\end{responsory-final}%
\V Exaltata es sancta dei genitrix.
\newline {\color{Red} In Laudibus. Ant.} Assumpta est
%pp129
Maria in celum, gaudent angeli, laudantes benedicunt dominum. {\color{Red} Ant.} Maria virgo assumpta est ad ethereum thalamum, in quo rex regum stellato sedet solio. {\color{Red} Ant.} In odore unguentorum tuorum currimus, adolescentule dilexerunt te nimis. {\color{Red} Ant.} Benedicta filia tu a domino, quia per te fructum vite communicavimus. {\color{Red} Ant.} Pulchra es et decora filia Hierusalem, terribilis ut castrorum acies ordinata. {\color{Red} Cap.} In omnibus. {\color{Red} Ymnus.} \hyperlink{o-gloriosa}{O gloriosa.} \V Elegit eam deus.
\newline {\color{Red} Ad Ben. Ant.} Que est ista que ascendit sicut aurora consurgens, pulchra ut luna, electa ut sol, terribilis ut castrorum acies ordinata. {\color{Red} Oratio.}
\lettrine[lines=2]{\zallmancaps \color{Red} V}{eneranda} nobis domine huius diei festivitas opem conferat salutarem, in qua sancta dei genitrix mortem subiit temporalem, nec tamen mortis nexibus deprimi potuit, que filium tuum dominum nostrum de se genuit incarnatum. Qui tecum.%typo: die for diei
\newline \ul{Quando predicta oratio dicitur ad memoriam, et sequitur alia memoria immediate terminetur, sic.} Incarnatum. Qui tecum vivit et regnat, per omnia secula seculorum.
\newline {\color{Red} Ad Terciam. Cap.} In omnibus requiem. \R \hyperlink{sancta-dei-genitrix}{Sancta dei genitrix.} \V Post partum.
\newline {\color{Red} Ad Sextam. Cap.}
\lettrine[lines=2]{\zallmancaps \color{Blue} E}{t} sic in Syon firmata sum, et in civitate sanctificata similiter requievi, et in Hierusalem potestas mea.
\R \hyperlink{post-partum}{Post partum virgo.} \V Speciosa.
\newline {\color{Red} Ad Nonam. Capitulum.}
\lettrine[lines=2]{\zallmancaps \color{Red} E}{t} radicavi in populo honorificato, et in partes dei mei hereditas illius: et in plenitudine sanctorum detentio mea.
\R \hyperlink{speciosa-facta-es}{Speciosa facta es.} \V Elegit eam.
\newline \ul{Ad Vesperas. Cap., Ymnus, \Vbar , ut in Primis Vesperis.}
\newline {\color{Red} Ad Mag. Ant.} Hodie Maria virgo celos ascendit gaudete, quia cum Christo regnat in eternum.
\newline \ul{Per Octavas quando de ipsis agitur et in octava die, Invit. sicut in die. Sequentes vero antiphone dicantur cotidie per octavas ad Ben. et ad Mag. vel ad memoriam.}
\newline {\color{Red} Ad Ben. Ant.} Virgo prudentissima quo progrederis quasi aurora valde rutilans filia Syon, tota formosa et suavis es, pulchra ut luna, electa ut sol.
\newline {\color{Red} Ad Mag. Ant.} Speciosa facta est et suavis in deliciis virginitatis sancta dei genitrix quam videntes filie Syon vernantem in floribus rosarum et liliis convallium, beatissimam predicaverunt, et regine laudaverunt eam.
\newline \ul{Tres lectiones sequentes cum aliis sex de festo legantur in octava, et quotiens necesse fuerit infra octavas.}
\newline \ul{De sermone supradicto,} Adest. {\color{Red} Lectio.}
\lettrine[lines=2]{\zallmancaps \color{Blue} N}{on} invenitur apud Latinos aliquis tractatorum, de beate Marie morte, aliquid aperte dixisse. Nam cum illum versiculum evangelicum quo Symeon dicit ad domini Matrem, et tuam ipsius animam pertransibit gladius, beate recordationis tractaret Ambrosius: ait. Nec hystoria nec littere docent beatam Mariam gladio vitam finisse. Hinc et Ysidorus. Incertum est per hoc inquit dictum: utrum gladium spiritus an gladium dixerit persecutionis.
\newline {\color{Red} Lectio.}
\lettrine[lines=2]{\zallmancaps \color{Red} S}{ed} quid de his quibus loquor dicam, cum nec ipse qui hanc ante crucem domini accepit in sua, idest Iohannes evangelista, de hoc posteris aliquid retinendum scriptis mandavit? Nullus enim hoc fidelius enarrare potuit, si illud deus manifestari voluisset: quam ille utique qui hanc nutriendam suscepit, nec contra morem, filius matrem reliquit.
\newline {\color{Red} Lectio.}
\lettrine[lines=2]{\zallmancaps \color{Blue} R}{estat} ergo ut homo mendaciter non fingat apertum, quod deus voluit manere occultum. Vera autem de eius assumptione sententia hec esse probatur, ut secundum apostolum sive in corpore sive extra corpus ignorantes: assumptam super angelos credamus.
\newline \ul{In Octava sancti Laurentii. Ad Primas Vesperas, antiphona et psalmi de beata virgine. Cetera omnia sicut in festo, exceptis responsorio in Vesperis, et versibus antiphonarum.}
{\ubsubsection{Sancti Agapiti Martyris}{sancti-agapiti-martyris-breviarium}}
\newline {\color{Red} Oratio.}
\lettrine[lines=2]{\zallmancaps \color{Red} L}{etetur} ecclesia tua deus beati Agapiti martyris tui confisa suffragiis, atque eius precibus gloriosis et devota permaneat, et secura consistat. Per.
{\ubsubsection{Sancti Bernardi Abbatis}{sancti-bernardi-abbatis-breviarium}}
\newline {\color{Red} Oratio.}
\lettrine[lines=2]{\zallmancaps \color{Blue} I}{ntercessio} nos quesumus domine beati Bernardi abbatis commendet, ut quod nostris meritis non valemus, eius patrocinio assequamur. Per.
{\ubsubsection{Sanctorum Thymothei et Symphoriani Martyrum}{sanctorum-thymothei-et-symphoriani-martyrum-breviarium}}
\newline {\color{Red} Oratio.}
\lettrine[lines=2]{\zallmancaps \color{Red} A}{uxilium} tuum nobis domine quesumus placatus impende, et intercedentibus beatis martyribus tuis Thymotheo et Simphoriano dexteram super nos tue propiciationis extende. Per.
{\ubsubsection{Sancti Bartholomei Apostoli}{sancti-bartholomei-apostoli-breviarium}}
\newline {\color{Red} Oratio.}
\lettrine[lines=2]{\zallmancaps \color{Blue} O}{mnipotens} sempiterne deus qui huius diei venerandam sanctamque leticiam in beati apostoli tui Bartholomei festivitate tribuisti, da ecclesie tue quesumus, et amare quod credidit, et predicare quod docuit. Per.%typo tribusti
{\ubsubsection{Sancti Rufi Martyris}{sancti-rufi-martyris-breviarium}}
\newline {\color{Red} Oratio.}
\lettrine[lines=2]{\zallmancaps \color{Red} A}{desto} domine supplicationibus nostris, ut qui ex iniquitate nostra reos nos esse cognoscimus, beati Rufi martyris tui intercessione liberemur. Per.
{\ubsubsection{Sancti Augustini Episcopi et Confessoris}{sancti-augustini-episcopi-et-confessoris-breviarium}}
\newline {\color{Red} Ad Vesperas. Super psalmos. Ant.} Letare mater nostra Hierusalem quia rex tuus dispensatorem strenuum et civem fidelissimum, de servitute Babilonis tibi redemit Augustinum. {\color{Red} Ps.} \psalm{112} \ul{cum ceteris.} \R \hyperlink{volebat-enim}{Volebat.} {\color{Red} VI. Ymnus.}
\hymn{magne-pater-augustine}{Magne pater Augustine}{
\lettrine[lines=2]{\zallmancaps \color{Blue} M}{agne} pater Augustine
preces nostras suscipe,
et per eas conditori
nos placare satage,
atque rege gregem tuum,
summum decus presulum.
\par {\bfseries \color{Red} A}matorem paupertatis
te collaudant pauperes,
assertorem veritatis
amant veri iudices,
frangis nobis favos mellis
de scripturis disserens.
\par {\bfseries \color{Blue} Q}ue obscura prius erant
nobis plana faciens,
tu de verbis salvatoris
dulcem panem conficis,
et propinas potum vite
de psalmorum nectare.
\par {\bfseries \color{Red} T}u de vita clericorum
sanctam scribis regulam,
quam qui amant et sequuntur
viam tenent regiam,
atque tuo sancto ductu
redeunt ad patriam.
\par {\bfseries \color{Blue} R}egi regum salus vita,
decus et imperium,
trinitati laus et honor
sit per omne seculum,
qui concives nos ascribat
supernorum civium. Amen.
}
\newline {\color{Red} Ad Mag. Ant.} Adest dies celebris quo solutus nexu carnis sanctus presul Augustinus assumptus est cum angelis, ubi gaudet cum prophetis letatur cum apostolis, quorum plenus spiritu que predixerunt mystica fecit nobis pervia, post quos secunda dispensandi verbi dei primus refulsit gratia. {\color{Red} Oratio.}
\lettrine[lines=2]{\zallmancaps \color{Red} A}{desto} supplicationibus nostris omnipotens deus, et quibus fiduciam sperande pietatis indulges, intercedente beato Augustino confessore tuo atque pontifice consuete misericordie tribue benignus effectum. Per.
\newline {\color{Red} Ad Mat. Invit.}
\begin{invitatory}[magnus-dominus-invitatorium]
{Magnus dominus et laudabilis valde, * Qui de tenebris gentium lumen ecclesie sue vocavit Augustinum.}
\end{invitatory}
\newline {\color{Red} Ymnus.} \hyperlink{magne-pater-augustine}{Magne pater.}
\newline {\color{Red} In Primo Nocturno. Ant.} Aperuit Augustinus codicem apostolicum et coniectis oculis ad primum capitulum legit, induimini dominum Ihesum Christum, et statim quasi infusa luce securitatis ab eo omnes dubietatis tenebre diffugerunt. {\color{Red} Ps.} \psalm{1} {\color{Red} Ant.} Insinuavit ergo per litteras sancto viro Ambrosio presens votum suum, ut moneret, quid sibi de libris sanctis legendum esset quo percipiende Christiane gratie aptior fieret atque paratior. {\color{Red} Ps.} \psalm{2} {\color{Red} Ant.} At ille iussit Ysaiam prophetam, quod pre ceteris evangelii vocationisque gentium sit prenunciator apertior. {\color{Red} Ps.} \psalm{3} \R
\lettrine[lines=2]{\zallmancaps \color{Blue} I}{nvenit} \hypertarget{invenit-se-augustinus}{\label{invenit-se-augustinus}} se Augustinus longe esse a deo in regione dissimilitudinis tanquam audiret vocem dei de excelso, cibus sum grandium, * Cresce et manducabis me. \V Nec tu me mutabis in te sicut cibum carnis tue sed tu mutaberis in me. Cresce.
\begin{responsory}[sensit-igitur]
{Sensit igitur et expertus est non esse mirum quod palato non sano pena est panis qui sano est suavis, * Et oculis egris odiosa lux que puris est amabilis.}{Propter iniquitatem corripuisti hominem, et tabescere fecisti sicut araneam animam eius.}
\end{responsory}%
\begin{responsory-final}[tum-vero]
{Tum vero invisibilia dei per ea que facta sunt intellecta conspexit, sed aciem figere non valuit, * Quia non secum ferebat nisi amantem memoriam, * Et quasi olefacta desiderantem que comedere nondum posset.}{Quid autem sacramenti haberet verbum caro factum est nec suspicari quidem poterat.}{Et quasi}
\end{responsory-final}%
\newline {\color{Red} In Secundo Nocturno. Ant.} Veruntamen primam huius lectionem non intelligens totumque talem arbitrans, distulit repetendum exercitatior in dominico eloquio. {\color{Red} Ps.} \psalm{4} {\color{Red} Ant.} Inde ubi tempus advenit quo nomen eum dare oporteret relicto rure Mediolanum remeans baptizatus est a beato Ambrosio fugitque ab eo sollicitudo vite preterite. {\color{Red} Ps.} \psalm{5} {\color{Red} Ant.} Nec saciabatur illis diebus dulcedine mirabili considerare altitudinem consilii divini super salutem generis humani. {\color{Red} Ps.} \psalm{8}
\begin{responsory}[itaque-avidissime]
{Itaque avidissime legere cepit pre ceteris apostolum Paulum, * Et perierunt ille questiones in quibus aliquando visus est ei adversari testimoniis legis et prophetarum textus sermonis eius.}{Et apparuit ei una facies eloquiorum castorum et exultare cum tremore didicit.}
\end{responsory}%
\begin{responsory}[misit-ergo]
{Misit ergo dominus in mentem eius, visumque est bonum in conspectu eius pergere ad Simplicianum, qui ei bonus apparebat servus dei, * Et lucebat in eo gratia divina.}{Audierat enim quod a iuventute sua devotissime deo viveret, et vere sic erat.}%typo Simplianum
\end{responsory}%
\begin{responsory-doxology}[volebat-enim]
{Volebat enim conferenti estus suos ut proferret quis esset aptus modus vivendi sic affecto ut ipse erat ad ambulandum in via dei, * In qua alius sic, alius sic ibat.}{Displicebat ei quicquid agebat in seculo pre dulcedine dei et decore domus eius quam dilexit.}
\end{responsory-doxology}%
\newline {\color{Red} In Tertio Nocturno. Ant.} Flebat autem uberrime in ymnis et canticis suave sonantis ecclesie vocibus vehementer affectus. {\color{Red} Ps.} \psalm{14} {\color{Red} Ant.} Voces igitur ille influebant auribus eius et eliquabatur veritas in cor eius, et fluebant lacrime et bene illi erat cum eis. {\color{Red} Ps.} \psalm{20} {\color{Red} Ant.} Adiunctus inde Nebridio et Evodio remeabat ad Affricam, querens quis eos ad bene vivendum locus haberet utilius, sed cum essent apud Ostia Tiberina pia mater eius defuncta est. {\color{Red} Ps.} \psalm{23}
\begin{responsory}[vulneraverat-caritas]
{Vulneraverat caritas Christi cor eius et gestabat verba eius in visceribus quasi sagittas acutas, * Et exempla servorum dei quos de mortuis vivos fecerat tanquam carbones vastatores.}{Ascendenti a convalle plorationis et cantanti canticum graduum dederat sagittas accutas.}%sic accutas
\end{responsory}%
\begin{responsory}[accepta-baptismi]
{Accepta baptismi gratia positus apud Ostia Tiberina cum matre sua colloquebatur soli valde dulciter, et inhiabant ore cordis pariter in superna fluenta fontis vite, * Vilescebatque mundus iste inter verba cum delectationibus suis.}{Tunc ait illa filio nulla re iam delector in hac vita cum te contempta felicitate terrena videam servum Christi.}
\end{responsory}%
\begin{responsory-final}[verbum-dei-usque]
{Verbum dei usque ad ipsam suam egritudinem impretermisse alacriter et fortiter sana mente sanoque consilio in sancta ecclesia predicavit, membris omnibus sui corporis incolumis integro aspectu atque auditu, * Coram positis fratribus et orantibus * Dormivit cum patribus suis.}{Testamentum nullum fecit quia unde faceret pauper Christi non habuit.}{Dormivit}
\end{responsory-final}%
\newline {\color{Red} In Laudibus. Ant.} Post mortem matris reversus est Augustinus ad agros proprios, ubi cum amicis ieiuniis et orationibus vacans scribebat libros et docebat indoctos. {\color{Red} Ant.} Comperta autem eius fama beatus Valerius Yponensis episcopus eum ad se accersiri fecit, et licet invitum presbiterum ordinavit. {\color{Red} Ant.} Factus ergo presbiter monasterium clericorum mox instituit, et cepit vivere secundum regulam sub sanctis apostolis constitutam. {\color{Red} Ant.} Sanctus autem Valerius ordinator eius exultabat uberius hominem sibi talem datum divinitus, qui in doctrina sana edificare ecclesiam esset Ydoneus. {\color{Red} Ant.} Eodem tempore Fortunatus presbiter Manicheorum versutia plurimos seducebat, quem sanctus Augustinus in conventu omnium disputans publice superavit. {\color{Red} Ymnus.}
\hymn{celi-cives-applaudite}{Celi cives applaudite}{
\lettrine[lines=2]{\zallmancaps \color{Red} C}{eli} cives applaudite
et vos fratres concinite,
patris nostri sollemnia
solis reduxit orbita.
\par {\bfseries \color{Blue} Q}uod lingua foris personat
intus affectus sentiat,
nec imitari pigeat,
quod laudare mens approbat.
\par {\bfseries \color{Red} H}unc post mundi curricula
celi suscepit curia,
quem cum suis fidelibus
iam coronavit dominus.
\par {\bfseries \color{Blue} C}onemur totis viribus
iungamus preces precibus,
ut Augustini, meritis
celi fruamur gaudiis.
\par {\bfseries \color{Red} P}resta pater piissime,
patrique compar unice,
cum spiritu paraclito
regnans per omne seculum. Amen.
}
\newline {\color{Red} Ad Ben. Ant.} In diebus eius obsessa est civitas Yponensis ab exercitu barbarorum, inter que mala fuerunt Augustino lacrime sue panes die ac nocte, atque sub hoc eventu ad extremam horam veniens obdormivit in pace.
\newline {\color{Red} Ad Mag. Ant.} Hodie gloriosus pater Augustinus dissoluta huius habitationis domo, domum non manufactam accepit in celis, quam sibi cooperante dei gratia, manu lingua fabrefecit in terris, ubi iam quod sitivit internum gustat eternum, decoratus una stola securusque de reliqua.
\newline \ul{Post orationem de sancto Augustino fiat memoria de sancto Iohanne.} {\color{Red} Ant.} Arguebat Herodem Iohannes propter Herodiadem quam tulerat fratri suo Philippo uxorem. \V Gloria et honore. {\color{Red} Oratio.}
\lettrine[lines=2]{\zallmancaps \color{Blue} S}{ancti} Iohannis baptiste et martyris tui domine quesumus veneranda festivitas, salutaris auxilii nobis prestet effectum. Per.
\newline \ul{Per octavas sancti Augustini, quando de octava agitur nisi in Dominica, Invitat.} \hyperlink{regem-confessorum-invitatorium}{Regem confessorem.} \ul{Sequentes vero antiphone dicantur ad Ben. et ad Mag. vel ad memoriam.}
\newline {\color{Red} Ad Ben. Ant.} Huius mater devotissima quem carne prius pepererat mundo caritatis visceribus postmodum multo semine lacrimarum genuit Christo.
\newline {\color{Red} Ad Mag. Ant.} \hyperlink{sancti-augustini-episcopi-et-confessoris-breviarium}{Letare mater.} \ul{quere supra in Primis Vesperis super psalmos.}
{\ubsubsection{In Decollatione Sancti Iohannis Baptiste}{in-decollatione-sancti-iohannis-baptiste-breviarium}}
\newline {\color{Red} Ad Matutinas.} \ul{Invitat., Ymnus, antiphone, psalmi, versiculi nocturnorum de communi unius martyris.}
\newline \ul{Ex hystoria de capitis Iohannis inventione.} \ul{\grecross } {\color{Red} Lectio Prima.}
\lettrine[lines=2]{\zallmancaps \color{Red} D}{ecollato} Iohanne apud Macheronta, Herodias iubente Herode non est passa caput eius cum reliquo corpore sepeliri, metuens ne caput existens cum corpore facilius posset resurgere. Corpus siquidem eius apud Sebasten Palestine urbem humatum legitur, et a barbaris intactum usque ad apostatam Iulianum. Huius siquidem apostate tempore ad suggestionem eius pagani sepulchrum Iohannis invadunt, ossa dispergunt, et rursus collecta et cremata latius dispergunt. \R \hyperlink{iste-sanctus-pro}{Iste sanctus.}
\newline {\color{Red} Lectio Secunda.}
\lettrine[lines=2]{\zallmancaps \color{Blue} S}{ed} fuerunt quidam ex Hierosolimis monachi qui mixti colligentibus quecumque potuerunt colligere, ad Philippum Hierosolimorum episcopum pertulerunt. Ille vero tantum thesaurum se posse servare supra se estimans misit illud ad pontificem Alexandrie Athanasium. \R \hyperlink{posuisti-domine-super}{Posuisti domine.}
\newline {\color{Red} Lectio Tertia.}
\lettrine[lines=2]{\zallmancaps \color{Red} R}{egnante} vero Theodosio, qui quintus fuit post Iulianum, destructa sunt de eius mandato Alexandrie delubria ydolorum, et edificata ecclesia mire magnitudinis in honore Iohannis baptiste, et ibidem cum honore maximo reposite reliquie supradicte. Tunc temporis caput precursoris adhuc repositum erat in parte palatii: in qua Herodias illud infoderat.
\begin{responsory-final}[misit-herodes-rex]
{Misit Herodes rex manus ac tenuit Iohannem et vinxit eum in carcerem, quia metuebat eum propter Herodiadem. * Quam tulerat * Fratri suo Philippo uxorem.}{Arguebat Herodem Iohannes propter Herodiadem.}{Fratri}
\end{responsory-final}%
\newline {\color{Red} Lectio Quarta.}
\lettrine[lines=2]{\zallmancaps \color{Blue} T}{empore} vero Marciani et Valentiani qui quarto loco post Theodosium imperarunt, venerunt quidam monachi Hierosolimam, loca veneranda conspicere cupientes. Hi revelatione sancti Iohannis certa et multiplici sibi facta properantes ad predictum palatium, caput ipsius invenerunt sacris cilicinis involutum, vestibus ut estimo quibus in deserto fuerat obvolutus. \R \hyperlink{beatus-vir-qui-suffert}{Beatus vir qui suffert.}
\newline {\color{Red} Lectio Quinta.}
\lettrine[lines=2]{\zallmancaps \color{Red} C}{um} autem ipso accepto ad propria properarent, quidam Emissene civitatis figulus paupertatem fugiens eis sese exhibuit comitem. Hic dum peram sibi creditam cum sacro capite portaret, admonitus nocte a sancto Iohanne ipsos fugiens Emissenam urbem cum sancto capite ingressus est. \R \hyperlink{domine-prevenisti-eum}{Domine prevenisti.}
\newline {\color{Red} Lectio Sexta.}
\lettrine[lines=2]{\zallmancaps \color{Blue} I}{bique} quamdiu vixit in quodam specu sanctum caput venerans prosperitatem adeptus est, paupertate fugata. Moriens autem illud revelavit sub fide sorori sue. Et secundum eundem modum sibi invicem successores.
\begin{responsory-final}[in-medio-carceris]
{In medio carceris stabat beatus Iohannes, voce magna clamavit et dixit, * Domine deus meus * Tibi commendo spiritum meum.}{Misit rex et decollari iussit Iohannem in carcere orantem et dicentem.}{Tibi}
\end{responsory-final}%
\newline {\color{Red} Lectio VII.}
\lettrine[lines=2]{\zallmancaps \color{Red} C}{um} autem post multum tempus Marcellus presbiter in eodem specu sanctam vitam duceret, apparens ei beatus Iohannes et rem revelans precepit ut caput eius una cum episcopo Alexandrie deferret apud Alexandriam ubi erant alie reliquie corporis sui. Quo veniente episcopo ita factum est. \R \hyperlink{stola-iocunditatis}{Stola iocunditatis.}
\newline {\color{Red} Lectio VIII.}
\lettrine[lines=2]{\zallmancaps \color{Blue} C}{um} autem quidam presbiter qui cum episcopo venerat manum extenderet ad ydriam in qua caput erat repositum, dubitans an illud vere esset caput Iohannis, adhesit manus eius ad os ydrie, nec eam extrahere valuit donec orantibus aliis ipsam retraxit, sed invalidam. \R \hyperlink{corona-aurea-super}{Corona aurea.}
\newline {\color{Red} Lectio IX.}
\lettrine[lines=2]{\zallmancaps \color{Red} E}{piscopus} ergo ydriam cum sancto Thesauro collocavit in diaconio donec instaret sollemnitas decollationis, quam Theophilus quarto kalendas Septembris ordinaverat ut ipsa festivitas non decollatio sed collectio vel capitis eius inventio nominaretur. Translatum est autem caput sanctum, regnante Pipino rege Francorum ad Gallias in regionem Aquitanie.
\begin{responsory-final}[accedentes-discipuli]
{Accedentes discipuli sancti Iohannis baptiste, * Tulerunt corpus eius * Et sepelierunt illud.}{Iussu Herodis amputatum est caput Iohannis in carcere quo audito discipuli eius.}{Et sepelierunt}
\end{responsory-final}%
\V Ora pro nobis.
\newline {\color{Red} In Laudibus. Ant.} Herodes enim tenuit et ligavit Iohannem et posuit in carcerem propter Herodiadem. {\color{Red} Ant.} Puelle saltanti imperavit mater, nichil aliud petas nisi caput Iohannis. {\color{Red} Ant.} Domine mi rex da michi in disco caput Iohannis baptiste. {\color{Red} Ant.} Da michi in disco caput Iohannis baptiste, et contristatus est rex propter iusiurandum. {\color{Red} Ant.} Misit rex incredulus ministrum detestabilem et amputari iussit caput Iohannis baptiste. \ul{Cap., Ymnus, \Vbar , unius martiris.}
\newline {\color{Red} Ad Ben. Ant.} Misso Herodes spiculatore precepit amputari caput Iohannis in carcere, quo audito discipuli eius venerunt et tulerunt corpus eius, et posuerunt illud in monumento. \ul{Memoria de sancta Sabina martyre.} {\color{Red} Ant.} Veni electa. \V Specie tua. {\color{Red} Oratio.}
\lettrine[lines=2]{\zallmancaps \color{Blue} D}{eus} qui inter cetera potentie tue miracula eciam in sexu fragili victoriam martyrii contulisti, concede propicius, ut cuius natalicia colimus, per eius ad te exempla gradiamur. Per.
\newline \ul{Ad horas Cap., Responsoria, \Vbar , de communi unius martiris. Ad Vesperas super psalmos antiphona.} Herodes. \ul{Psalmi, Cap., Hymnus, \Vbar , unius martiris.}
\newline {\color{Red} Ad Mag. Ant.} Perpetuis nos domine sancti Iohannis baptiste tuere presidiis, et quanto fragiliores sumus tanto magis necessariis attolle suffragiis. \ul{Memoria de sanctis Felice et Adaucto martiribus.} {\color{Red} Oratio.}
\lettrine[lines=2]{\zallmancaps \color{Red} M}{aiestatem} tuam domine supplices deprecamur ut sicut nos iugiter sanctorum tuorum commemoratione letificas, ita semper supplicatione defendas. Per.
{\ubsubsection{Sancti Egidii Abbatis}{sancti-egidii-abbatis-breviarium}}
\newline {\color{Red} Oratio.}
\lettrine[lines=2]{\zallmancaps \color{Blue} I}{ntercessio} nos quesumus domine beati Egidii abbatis commendet, ut quod nostris meritis non valemus, eius patrocinio assequamur. Per.
\newline \ul{Si in hac die Dominica evenerit: tunc in Sabbato a capitulo Vesperarum fiat de Dominica et deinceps usque ad Matutinas secunde ferie exclusive et de octavis fiat tantum memoria.}
{\ubsubsection{Sancti Marcelli Martyris}{sancti-marcelli-martyris-breviarium}}
\newline {\color{Red} Oratio.}
\lettrine[lines=2]{\zallmancaps \color{Red} P}{resta} quesumus omnipotens deus ut qui beati Marcelli martyris tui natalicia colimus, intercessione eius in tui nominis amore roboremur. Per.
{\ubsubsection{In Nativitate Beate Marie Virginis}{in-nativitate-beate-marie-virginis-breviarium}}
\newline {\color{Red} Ad Vesperas. Super psalmos. Ant.} Hec est regina virginum que genuit regem velut rosa decora, virgo dei genitrix per quam reperimus deum et hominem, alma virgo intercedat pro nobis omnibus. {\color{Red} Ps.} \ul{Cum ceteris.} {\color{Red} Cap.}
\lettrine[lines=2]{\zallmancaps \color{Blue} E}{go} quasi vitis fructificavi suavitatem odoris, et flores mei fructus honoris et honestatis. \R \hyperlink{stirps-iesse}{Stirps Iesse.} {\color{Red} III. Ymnus.} \hyperlink{ave-maris-stella}{Ave maris stella.} \V Ora pro nobis sancta dei genitrix.
\newline {\color{Red} Ad Mag. Ant.} Beatissime virginis Marie nativitatem devotissime celebremus, ut ipsa pro nobis intercedat ad dominum Ihesum Christum. {\color{Red} Oratio.}
\lettrine[lines=2]{\zallmancaps \color{Red} S}{upplicationem} servorum tuorum deus miserator exaudi, ut qui in nativitate dei genitricis et virginis congregamur, eius intercessionibus a te de instantibus periculis eruamur. Per eundem.
\newline {\color{Red} Ad Completorium. Super psalmos. Ant.} Virgo Maria. {\color{Red} Ad Nunc Dimittis. Ant.} Corde et animo Christo canamus gloriam in hac sacra sollemnitate precelse genitricis dei Marie. \ul{Predicte antiphone dicantur ad Completorio per octavas nisi in festo sancte crucis.}
\newline {\color{Red} Ad Mat. Invitat.}
\begin{invitatory}[nativitatem-virginis-invitatorium]
{Nativitatem virginis Marie celebremus, * Christum eius filium adoremus dominum.}
\end{invitatory}
\newline {\color{Red} Ymnus.} \hyperlink{quem-terra}{Quem terra.}
\newline {\color{Red} In Primo Nocturno. Ant.} Ecce tu pulchra es amica mea ecce tu pulchra oculi tui columbarum. {\color{Red} Ps.} \psalm{8} {\color{Red} Ant.} Sicut lilium inter spinas, sic amica mea inter filias. {\color{Red} Ps.} \psalm{18} {\color{Red} Ant.} Favus distillans labia tua sponsa et odor vestimentorum tuorum sicut odor thuris. {\color{Red} Ps.} \psalm{23} \V Sancta dei genitrix.
\newline \ul{Bernardus in sermone de Assumptione.} {\color{Red} Lectio Prima.}
\lettrine[lines=2]{\zallmancaps \color{Blue} B}{eatam} virginem omnimodis constat ab originali contagio sola dei gratia priusquam nasceretur esse mundatam, quam nimirum pietas Christiana credere prohibet, vel minus a Hieremia in utero sanctificatam, vel non magis a Iohanne spiritu sancto repletam. Neque vero festis laudibus nascens honoraretur, si non sancta nasceretur. \R
\lettrine[lines=2]{\zallmancaps \color{Red} H}{odie} \hypertarget{hodie-nata-est}{\label{hodie-nata-est}} nata est beata virgo Maria ex progenie David per quam salus mundi credentibus apparuit, * Cuius vita gloriosa lucem dedit seculo. \V Beatissime virginis Marie nativitatem devotissime celebremus. Cuius.
\newline \ul{Idem in epistola ad canonicas Lugdunenses.} %typo Lugdimenses
%pp130
{\color{Red} Lectio Secunda.}
\lettrine[lines=2]{\zallmancaps \color{Red} E}{st} itaque virgo regia veris honorum titulis cumulata, que procul dubio sancta fuit antequam nata. Ego puto quod et copiosior sanctificationis benedictio quam in aliis sanctis ab utero sanctificatis in eam descenderit, que non solum ipsius ortum sanctificaret, sed et vitam eius ab omni peccato, deinceps immunem conservaret, quod alteri nemini in natis quidem mulierum creditur esse donatum. Decuit nimirum reginam virginum ut singulari previlegio vitam absque omni peccato duceret, que dum peccati mortisque peremptorem pareret: vite iusticieque munus omnibus obtineret.
\begin{responsory}[gloriose-virginis]
{Gloriose virginis Marie ortum dignissumum recolamus, * Cuius dominus humilitatem respiciens angelo nunciante concepit redemptorem mundi.}{Hodie nata est beata virgo Maria ex progenie David.}
\end{responsory}%
\newline \ul{Augustinus in libro de natura et gratia.} {\color{Red} Lectio Tertia.}
\lettrine[lines=2]{\zallmancaps \color{Blue} E}{xcepta} ergo sancta virgine Maria, si omnes sancti et sancte cum hic viverent interrogati fuissent utrum sine peccato essent, omnes una voce clamassent, si dixerimus quia peccatum non habemus ipsi nos seducimus, et veritas in nobis non est: excepta inquam hac sancta virgine, de qua propter honorem domini cum de peccatis agitur, nullam prorsus volo questionem habere. Scimus enim quod ei plus gratie collatum fuerit ad peccatum ex omni parte vincendum, que illum concipere et parere meruit quem constat nullum habuisse peccatum.
\begin{responsory-final}[stirps-iesse]
{Stirps Iesse virgam produxit virgaque florem, * Et super hunc florem * Requiescit spiritus almus.}{Virgo dei genitrix virga est, flos filius eius.}{Requiescit}
\end{responsory-final}%
\newline {\color{Red} In Secundo Nocturno. Ant.} Emissiones tue paradisus malorum punicorum cum pomorum fructibus. {\color{Red} Ps.} \psalm{44} {\color{Red} Ant.} Fons ortorum puteus aquarum viventium que fluunt impetu de Libano. {\color{Red} Ps.} \psalm{45} {\color{Red} Ant.} Veniat dilectus meus in ortum suum, ut comedat fructum pomorum suorum. {\color{Red} Ps.} \psalm{86} \V Post partum.
\newline \ul{Alcuinus in libro secundo de trinitate.} {\color{Red} Lectio Quarta.}%typo Alcumus
\lettrine[lines=2]{\zallmancaps \color{Red} B}{eata} nempe Maria lana mundissima fuit, et virginitate clarissima, que sola digna esset in se filii dei recipere dignitatem, sicut lana conchilii sanguinem, ut ex eadem fieret purpura imperiali maiestati tantummodo digna, quamque nullus esset dignus induere nisi augusta preditus dignitate.
\begin{responsory}[nativitas-gloriose]
{Nativitas gloriose virginis Marie ex semine Abrahe orta de tribu Iuda clara ex stirpe David, * Cuius vita inclita cunctas illustrat ecclesias.}{Regali ex progenie Maria exorta refulget.}
\end{responsory}%
\newline \ul{Odilio Cluniacensis in sermone quodam.} {\color{Red} Lectio Quinta.}
\lettrine[lines=2]{\zallmancaps \color{Blue} U}{num} quippe scimus pro certo quod omnis Marie vita et actio semper intenta fuit in domino. Ipsa siquidem optimam partem elegit, ideoque promeruit quod ex ea dei filius angelo nunciante corpus redemptionis nostre suscepit.
\begin{responsory}[nativitas-tua-dei]
{Nativitas tua dei genitrix virgo gaudium annunciavit universo mundo, * Ex te enim ortus est sol iusticie Christus deus noster, qui solvens maledictionem dedit benedictionem, et confundens mortem donavit nobis vitam sempiternam.}{Felix namque es que sola genitricis dignitatem obtinuisti, et virginalem pudiciciam non amisisti.}
\end{responsory}%
\newline \ul{Bernardus super missus est, omelia secunda.} {\color{Red} Lectio Sexta.}
\lettrine[lines=2]{\zallmancaps \color{Red} P}{roinde} factor hominum ut homo fieret nasciturus de homine, talem sibi debuit matrem eligere quin potius condere, qualem et decere sciebat et sibi placituram noverat. Voluit igitur esse virginem de qua immaculata immaculatus procederet omnium maculas purgaturus, voluit, et humilem de qua mitis et humilis corde prodiret, harum in se virtutum hominibus exemplum saluberrimum ostensurus.
\begin{responsory-doxology}[ad-nutum-domini]
{Ad nutum domini nostrum ditantis honorem, * Sicut spina rosam genuit Iudea Mariam.}{Ut vitium virtus operiret gratia culpam.}
\end{responsory-doxology}%
\newline {\color{Red} In Tertio Nocturno. Ant.} Veni in ortum meum soror mea sponsa, messui mirram meam cum aromatibus meis. {\color{Red} Ps.} \psalm{95} {\color{Red} Ant.} Comedi favum cum melle meo bibi vinum meum cum lacte meo. {\color{Red} Ps.} \psalm{96} {\color{Red} Ant.} Talis est dilectus meus et ipse est amicus meus filie Hierusalem. {\color{Red} Ps.} \psalm{97} \V Speciosa facta es.
\newline {\color{Red} Lectio VII. Secundum Matheum.}
\lettrine[lines=2]{\zallmancaps \color{Blue} L}{iber} generationis Ihesu Christi filii David: filii Abraham. Abraham genuit Ysaac. Ysaac autem genuit Iacob. Et reliqua.
\newline {\color{Red} Omelia Beati Hieronimi Presbiteri. \grecross .}
\lettrine[lines=2]{\zallmancaps \color{Red} I}{n} Ysaya legimus. Generationem eius quis enarrabit? Non putemus evangelistam prophete esse contrarium, ut quod ille impossibile dixit effatu, hic narrare incipiat, quia ibi de generatione divinitatis, hic vero de incarnatione dictum est. A carnalibus autem cepit, ut per hominem Deum dicere incipiamus.
\begin{responsory}[regali-ex-progenie]
{Regali ex progenie Maria exorta refulget, * Cuius precibus nos adiuvari mente et spiritu devotissime poscimus.}{Corde et animo Christo canamus gloriam in hac sacra sollemnitate precelse genitricis dei Marie.}
\end{responsory}%
\newline {\color{Red} Lectio VIII.}
\lettrine[lines=2]{\zallmancaps \color{Blue} F}{ilii} David, filii Abraham. Ordo preposterus sed necessario commutatus. Si enim primum posuisset Abraham, et postea David: rursus ei repetendus fuerat Abraham, ut generationis series texeretur. Ideo autem ceteris pretermissis horum filium nuncupavit, quia ad hos tantum facta est de Christo repromissio.
\begin{responsory}[corde-et-animo]
{Corde et animo Christo canamus gloriam, * In hac sacra sollemnitate precelse genitricis dei Marie.}{Cum iocunditate nativitatem sancte Marie celebremus.}
\end{responsory}%
\newline {\color{Red} Lectio IX.}
\lettrine[lines=2]{\zallmancaps \color{Red} I}{udas} autem genuit Phares et Zaram de Thamar. Notandum in genealogia salvatoris nullam sanctarum assumi mulierum, sed eas quas scriptura reprehendit ut qui propter peccatores venerat, de peccatricibus nascens, omnium peccata deleret. Unde et in consequentibus Ruth Moabitis ponitur, et Bethsabee uxor Urie.
\begin{responsory-final}[solem-iusticie]
{Solem iusticie regem paritura supremum, * Stella Maria maris * Hodie processit ad ortum.}{Cernere divinum lumen gaudete fideles.}{Hodie}
\end{responsory-final}%
\V Ora pro nobis.
\newline {\color{Red} In Laudibus. Ant.} Nativitas gloriose virginis Marie ex semine Abrahe orta de tribu Iuda clara ex stirpe David. {\color{Red} Ant.} Nativitas est hodie sancte Marie virginis, cuius vita inclita cunctas illustrat ecclesias. {\color{Red} Ant.} Regali ex progenie Maria exorta refulget, cuius precibus nos adiuvari mente et spiritu devotissime poscimus. {\color{Red} Ant.} Corde et animo Christo canamus gloriam in hac sacra sollemnitate precelse genitricis dei Marie. {\color{Red} Ant.} Cum iocunditate nativitatem sancte Marie celebremus ut ipsa pro nobis intercedat ad dominum Ihesum Christum. {\color{Red} Cap.} Ego quasi vitis. {\color{Red} Ymnus.} \hyperlink{o-gloriosa}{O gloriosa domina.} \V Elegit eam.
\newline {\color{Red} Ad Ben. Ant.} Nativitatem hodiernam perpetue virginis genitricis dei Marie sollemniter celebremus qua celsitudo throni processit alleluya.
\newline {\color{Red} Ad Terciam. Cap.} Ego quasi. \R \hyperlink{sancta-dei-genitrix}{Sancta dei genitrix.} \V Post partum.
\newline {\color{Red} Ad Sextam. Cap.}
\lettrine[lines=2]{\zallmancaps \color{Blue} T}{ransite} ad me omnes qui concupiscitis me, et a generationibus meis implemini. Spiritus enim meus super mel dulcis, et hereditas mea super mel et favum.
\R \hyperlink{post-partum}{Post partum.} \V Speciosa.
\newline {\color{Red} Ad Nonam. Cap.}
\lettrine[lines=2]{\zallmancaps \color{Red} Q}{ui} audit me non confundetur, et qui operantur in me non peccabunt, et qui elucidant me vitam eternam habebunt.
\R \hyperlink{speciosa-facta-es}{Speciosa.} \V Elegit eam.
\newline {\color{Red} Ad Vesperas. Ps.} \psalm{109} \ul{cum ceteris sicut in Purificatione. Cap., Ymnus, \Vbar , ut in Primis Vesperis.}
\newline {\color{Red} Ad Mag. Ant.} Nativitas tua dei genitrix virgo gaudium annunciavit universo mundo, ex te enim ortus est sol iusticie Christus deus noster, qui solvens maledictionem dedit benedictionem, et confundens mortem donavit nobis vitam sempiternam.
\newline \ul{Per octavas quando de ipsis agitur, et in octava die, Invitat. sicut in festo. Sequentes vero antiphone dicantur cotidie per octavas Ad Ben. et ad Mag. vel ad memoriam.}
\newline {\color{Red} Ad Ben. Ant.} Maria virgo semper letare que meruisti Christum portare, celi et terre conditorem, quia de tuo utero protulisti mundi salvatorem.
\newline {\color{Red} Ad Mag. Ant.} Sancta Maria virginum piissima suscipe vota servulorum assidua, lapsos erige, errantes corrige, trementes corrobora, pusillanimes conforta, ut tibi semper referamus laudes quam dei summi colimus genitricem.
\newline \ul{Tres lectiones sequentes cum aliis sex de festo legantur in octava, et quotiens necesse fuerit infra octavas.}
\newline \ul{Ysidorus in libro de ortu etc.} {\color{Red} Lectio.}
\lettrine[lines=2]{\zallmancaps \color{Blue} M}{aria} que interpretatur domina sive illuminatrix, clara stirps David, virga Iesse, ortus conclusus, fons signatus, mater domini, templum dei, sacrarium spiritus sancti, virgo sancta, virgo feta, virgo ante partum, virgo post partum, salutationem ab angelo accepit et mysterium conceptionis agnovit, partus qualitatem requirens, et contra legem nature, obsequii fidem non renuit.
\newline {\color{Red} Lectio.}
\lettrine[lines=2]{\zallmancaps \color{Red} H}{anc} dominus ipse in cruce positus per sanguinis testamentum virgini commendavit discipulo, ut ipsum mater haberet vite comitem quem filius noverat integritatis esse custodem. Hanc quidam corporalis necis passione asserunt ab hac vita migrasse, pro eo quod iustus Symeon complectens brachiis suis Christum, prophetaverit dicens. Et tuam ipsius animam pertransiet gladius.
\newline {\color{Red} Lectio.}
\lettrine[lines=2]{\zallmancaps \color{Blue} I}{ncertum} quidem est utrum pro materiali gladio dixerit an pro verbo dei valido et acutiore omni gladio ancipiti. Specialiter tamen nulla docet hystoria Mariam gladii animadversione peremptam, quod nec obitus eius usquam legitur quamvis reperiatur sepulchrum.
{\ubsubsection{Sancti Gorgonii Martyris}{sancti-gorgonii-martyris-breviarium}}
\newline {\color{Red} Oratio.}
\lettrine[lines=2]{\zallmancaps \color{Red} S}{anctus} martyr tuus Gorgonius sua nos domine intercessione letificet et pia faciat sollemnitate gaudere. Per.
{\ubsubsection{Sanctorum Prothy et Iacincti, Martyrum}{sanctorum-prothy-et-iacincti-martyrum-breviarium}}
\newline {\color{Red} Oratio.}
\lettrine[lines=2]{\zallmancaps \color{Blue} B}{eatorum} martyrum tuorum Prothy et Iacincti nos domine foveat preciosa confessio, et pia iugiter intercessio tueatur. Per.
{\ubsubsection{In Exaltatione Sancte Crucis}{in-exaltatione-sancte-crucis-breviarium}}
\newline {\color{Red} Ad Vesperas. Super psalmos. Ant.} O crux admirabilis evacuatio vulneris restitutio sanitatis. \ul{Psalmi de octava. Cap., Ymnus, sicut in Inventione.} \R \hyperlink{per-tuam-crucem}{Per tuam crucem.} {\color{Red} IX.} \V Hoc signum crucis erit in celo.
\newline {\color{Red} Ad Mag. Ant.} Crux fidelis inter omnes arbor una nobilis, nulla silva talem profert, fronde, flore, germine, dulce lignum dulces clavos dulce pondus sustinet. {\color{Red} Oratio.}
\lettrine[lines=2]{\zallmancaps \color{Red} D}{eus} qui nos hodierna die exaltationis sancte crucis annua sollemnitate letificas, presta, ut cuius mysterium in terra cognovimus, eius redemptionis in celo premia consequamur. Per.
\newline \ul{Memoria de Octava. Postea de martiribus.} {\color{Red} Oratio.}
\lettrine[lines=2]{\zallmancaps \color{Blue} I}{nfirmitatem} nostram quesumus domine propicius respice, et mala omnia que iuste meremur sanctorum tuorum Cornelii et Cypriani intercessione averte. Per.
\newline {\color{Red} Ad Matutinas. Invit.}
\begin{invitatory}[regem-dominum-invitatorium]
{Regem dominum qui per crucis mysterium redemit mundum, * Venite adoremus.}
\end{invitatory}
\newline {\color{Red} Ymnus.} \hyperlink{salve-crux-sancta}{Salve crux sancta.}
\newline {\color{Red} In Primo Nocturno. Ant.} Adoramus te Christe et benedicimus tibi, quia per crucem tuam redemisti mundum. {\color{Red} Ps.} \psalm{8} {\color{Red} Ant.} Hoc signum crucis erit in celo cum dominus ad iudicandum venerit. {\color{Red} Ps.} \psalm{20} {\color{Red} Ant.} Per signum crucis de inimicis nostris libera nos deus noster. {\color{Red} Ps.} \psalm{23} \V Hoc signum crucis erit in celo.
\newline \ul{Rabanus de laude crucis, libro primo.} {\color{Red} Lectio Prima.}
\lettrine[lines=2]{\zallmancaps \color{Red} O}{}\ vere bona et sancta crux Christi quis te rite totam enarrare potest aut condigne laudare, que celestium archanorum pia es revelatrix, mysteriorum dei sacra conservatrix, sacramentorum Christi ydonea dispensatrix? In te angeli gaudia sua accumulata conspiciunt, in te homines iura salutis sue cognoscunt, in te inferi iustam retributionem fraudis sue percipiunt. \R
\lettrine[lines=2]{\zallmancaps \color{Blue} D}{ulce} lignum dulces clavos, dulce pondus sustinuit, * Que digna fuit portare precium huius seculi. \V Hoc signum crucis erit in celo cum dominus ad iudicandum venerit. Que.
\newline {\color{Red} Lectio Secunda.}
\lettrine[lines=2]{\zallmancaps \color{Red} O}{}\ bona crux, preterita renovas, presentia illustras, futura premonstras, perdita requiris, inventa custodis, lapsa restituis, et in viam pacis dirigis. Tu eterni regis es victoria, celestis militie leticia, terrigenarum quoque potentia. O bona crux, tu peccatorum remissio, pietatis exhibitio, meritorum augmentatio. Tu infirmorum remedium, laborantium auxilium, lassorum refrigerium. Tu cura egrotos medicans, tu gaudium mestos consolans, tu sanitas dolentes letificans.
\begin{responsory}[hoc-signum-crucis]
{Hoc signum crucis erit in celo cum dominus ad iudicandum venerit, * Tunc manifesta erunt abscondita cordis nostri.}{Cum sederit filius hominis in sede maiestatis sue et ceperit iudicare seculum per ignem.}
\end{responsory}%
\newline {\color{Red} Lectio Tertia.}
\lettrine[lines=2]{\zallmancaps \color{Blue} O}{}\ bona crux, tu status recte credentium, tu firmitas bene operantium, tu beatitudo rite perseverantium. Et quicquid de redemptione mundi potest cogitari quicquid lingua rite loqui, ad laudem tuam decentissime potest adaptari: quia quicquid in te laudatur, crucifixo in te Christo regi deputatur.
\begin{responsory-final}[o-crux-gloriosa]
{O crux gloriosa, o crux adoranda, o lignum preciosum, et admirabile signum, * Per quod et diabolus est victus, * Et mundus Christi sanguine redemptus.}{Adoremus crucis signaculum per quod salutis sumpsimus sacramentum.}{Et mundus}%typo iuctus
\end{responsory-final}%
\newline {\color{Red} In Secundo Nocturno. Ant.} Salva nos Christe salvator per virtutem sancte crucis, qui salvasti Petrum in mari miserere nobis. {\color{Red} Ps.} \psalm{29} {\color{Red} Ant.} Salvator mundi salva nos qui per crucem et sanguinem redemisti nos auxiliare nobis te deprecamur deus noster. {\color{Red} Ps.} \psalm{46} {\color{Red} Ant.} Adoremus crucis signaculum per quod salutis sumpsimus sacramentum. {\color{Red} Ps.} \psalm{65} \V Per signum crucis de inimicis nostris. \R Libera nos deus noster.
\newline \ul{Idem in libro secundo.} \grecross . {\color{Red} Lectio Quarta.}
\lettrine[lines=2]{\zallmancaps \color{Red} S}{alve} crux dei veneranda, que sapientia, lumen et doctrix orbis terrarum es, que et vera laus, et amica virtus, et clara phylosophia apud celicolas terrigenasque indesinenter viges, quam magis decet imperialem vocari thronum, quam servile tormentum, quia imperator, et rex noster Christus regnum sibi in te, et potestatem in celo et in terra conquisivit, hostes superavit, et mundum deo reconciliavit.
\begin{responsory}
{Ecce crucem domini fugite partes adverse, vicit leo de tribu Iuda, * Radix David.}{Crux benedicta in qua triumphavit rex angelorum.}
\end{responsory}%
\newline {\color{Red} Lectio Quinta.}
\lettrine[lines=2]{\zallmancaps \color{Blue} O}{}\ Christiani populi vexillum, framea qua cum hoste sortimur bellum, te obsecro ut tua virtute pectus meum infirmum benedicas, ut dignis laudibus eterni regis in te triumphum valeam decantare, quomodo scilicet celestia simul et terrena uno federe coniungas, pactum confirmes, et mortis vincula dissolvas.
\begin{responsory}[crux-fidelis-inter]
{Crux fidelis inter omnes arbor una nobilis, nulla silva talem profert fronde, flore, germine, * Dulce lignum, dulces clavos, dulce pondus sustinuit.}{Super omnia ligna cedrorum tu sola excelsior.}%typo lingua for ligna
\end{responsory}%
\newline {\color{Red} Lectio Sexta.}
\lettrine[lines=2]{\zallmancaps \color{Red} V}{irtus} est animi habitus, nature decus, vite ratio, morum nobilitas, et vite moderatio, que cuncta o sancta crux religiosis hominibus tradis, et cuncta simul beneplacita deo, et salubria depromis. Arbor odoris suavissimi, et expansione pulchrarum frondium latissima, ortus deliciarum incomparabilis, affluens ubertate largissima, floribus virtutum et foliis verborum iocundissima.
\begin{responsory-doxology}[tuam-crucem-adoramus]
{Tuam crucem adoramus domine tuam gloriosam recolimus passionem, * Miserere nostri qui passus es pro nobis.}{Adoramus te Christe et benedicmus tibi, quia per crucem tuam redemisti mundum.}
\end{responsory-doxology}%
\newline {\color{Red} In Tertio Nocturno. Ant.} Tuam crucem adoramus domine tuam gloriosam recolimus passionem, miserere nostri qui passus es pro nobis. {\color{Red} Ps.} \psalm{75} {\color{Red} Ant.} Crucem tuam adoramus domine et sanctam resurrectionem tuam laudamus et glorificamus, ecce enim propter crucem venit gaudium in universo mundo. {\color{Red} Ps.} \psalm{95} {\color{Red} Ant.} Crux alma fulget per quam salus reddita est mundo, crux vincit, crux regnat, crux repellit omne crimen alleluya. {\color{Red} Ps.} \psalm{96} \V Adoramus te Christe et benedicimus tibi. \R Quia per crucem tuam redemisti mundum.
\newline \ul{Anselmus in oratione de cruce,} Sancta crux. \grecross \ {\color{Red} Lectio VII.}
\lettrine[lines=2]{\zallmancaps \color{Blue} O}{}\ crux amabilis, o lignum preciosum, o signum venerandum, o crux gloriosa, non es suscipienda secundum crudelium qui te mitissimo paraverunt insipientissimam impietatem, sed eius qui te sponte suscepit sapientissimam pietatem. Nam illi te elegerunt, ut per te scelus impietatis perpetrarent, ille te elegit, ut per te opus pietatis consummaret.
\begin{responsory}[o-crux-benedicta]
{O crux benedicta, * Que sola fuisti digna portare regem celorum et dominum.}{O crux admirabilis evacuatio vulneris restitutio sanitatis.}
\end{responsory}%
\newline {\color{Red} Lectio VIII.}
\lettrine[lines=2]{\zallmancaps \color{Red} I}{lli} te elegerunt ut iustum morti traderent, ille ut peccatores erueret a morte. Illi ut vitam interficerent, ille ut mortem perimeret. Illi ut salvatorem damnarent, ille ut damnatos salvaret. Illi ut vivificantem mortificarent, ille ut mortuos vivificaret. Illi insipienter et crudeliter, ille sapienter et misericorditer.
\begin{responsory}[adoramus-te-christe]
{Adoramus te Christe et benedicimus tibi, * Quia per crucem tuam redemisti mundum.}{Tuam crucem adoramus domine tuam gloriosam recolimus passionem.}
\end{responsory}%
\newline {\color{Red} Lectio IX.}
\lettrine[lines=2]{\zallmancaps \color{Blue} O}{}\ crux admirabilis, per te infernus spoliatur, et per te redemptis obturatur. Per te demones terrentur, comprimuntur, vincuntur, conculcantur. Per te mundus renovatur, atque veritate lucente et iusticia regnante decoratur. Per te humana natura peccatrix, est iustificata, damnata salvata, ancilla Tartari liberata, mortua, resuscitata. Per te civitas illa in celis restauratur. Sit itaque in te et per te gloria mea, sit in te et per te vera spes mea. Per te anima mea a vita veteri mortificetur, et in novam resuscitetur.
\begin{responsory-doxology}
{Per tuam crucem salva nos Christe redemptor qui mortem nostram moriendo destruxisti, * Et vitam resurgendo reparasti.}{Miserere nostri Ihesu benigne qui passus es clementer pro nobis.}
\end{responsory-doxology}%
\newline \V Dicite in nationibus. \R Quia dominus regnavit a ligno.
\newline {\color{Red} In Laudibus. Ant.} O magnum pietatis opus, mors mortua tunc est quando in ligno mortua vita fuit. {\color{Red} Ant.} Ecce crucem domini fugite partes adverse, vicit leo de tribu Iuda radix David alleluya. {\color{Red} Ant.} Nos autem gloriari oportet in cruce domini nostri Ihesu Christi. {\color{Red} Ant.} O crux admirabilis, evacuatio vulneris, restitutio sanitatis. {\color{Red} Ant.} Crux benedicta nitet, dominus quia carne pependit, atque cruore suo vulnera nostra lavit.%typo Crus, ant. 5
\newline \ul{Capitulum, Ymnus, \Vbar , sicut in Inventione eiusdem.}
\newline {\color{Red} Ad Ben. Ant.} Super omnia ligna cedrorum tu sola excelsior in qua vita mundi pendit, in qua Christus triumphavit, et mors mortem superavit alleluya.
\newline \ul{Memoria de octavis, postea de martyribus. Cap., Responsoria, Versiculi ad horas sicut in Inventione eiusdem excepto quod responsoria et versiculi dicantur sine alleluya.}
\newline {\color{Red} Ad Vesperas. Super psalmos. Ant.} O magnum. {\color{Red} Ps.} \psalm{109} \psalm{110} \psalm{111} \psalm{112} \psalm{116} \ul{Cap., Ymnus, \Vbar , ut in Primis Vesperis.}
\newline {\color{Red} Ad Mag. Ant.} O crux benedicta quia in te pependit salvator mundi et in te triumphavit rex angelorum alleluya.
\newline \ul{Memoria de octava. Ant.} Beatissime. \ul{Deinde de sancto Nichomede martyre.} {\color{Red} Oratio.}
\lettrine[lines=2]{\zallmancaps \color{Red} A}{desto} domine populo tuo, ut beati Nichomedis martyris tui merita preclara suscipiens, ad impetrandam misericordiam tuam semper eius patrociniis adiuvetur. Per.
{\ubsubsection{Sancte Eufemie Virginis et Martiris}{sancte-eufemie-virginis-et-martiris-breviarium}}
\newline {\color{Red} Oratio.}
\lettrine[lines=2]{\zallmancaps \color{Blue} O}{mnipotens} sempiterne deus qui infirma mundi eligis ut fortia queque confundas, concede propitius ut qui beate Eufemie virginis et martyris tue sollemnia colimus, eius apud te patrocinia sentiamus. Per.
{\ubsubsection{Sancti Lamberti Episcopi et Martiris}{sancti-lamberti-episcopi-et-martiris-breviarium}}
\newline {\color{Red} Oratio.}
\lettrine[lines=2]{\zallmancaps \color{Red} D}{a} quesumus omnipotens deus, ut qui beati Lamberti martyris tui atque pontificis sollemnia colimus, eius apud te intercessionibus adiuvemur. Per.
{\ubsubsection{In Vigilia Sancti Mathei}{in-vigilia-sancti-mathei-breviarium}}
\newline {\color{Red} Lectio I. vel VII. si Dominica fuerit. Secundum Lucam.}
\lettrine[lines=2]{\zallmancaps \color{Blue} I}{n} illo tempore: Vidit Ihesus publicanum nomine Levi sedentem ad theloneum, et ait illi. Sequere me. Et reliqua.
\newline {\color{Red} Omelia Venerabilis Bede Presbiteri.}
\lettrine[lines=2]{\zallmancaps \color{Red} A}{d} theloneum, ad curam dispensationemque vectigalium dicit. Thelos enim Grece, Latine vectigal nominatur. Idem autem Levi, qui et Matheus est, sed Lucas Marcusque propter verecundiam et honorem evangeliste, nomen ponere noluerunt vulgatum.
\newline {\color{Red} Lectio II. vel VIII.}
\lettrine[lines=2]{\zallmancaps \color{Blue} I}{pse} autem Matheus iuxta illud quod scriptum est, iustus accusator sui est in principio sermonis, Matheum se et publicanum nominat, ut ostendat legentibus nullum
%pp131
debere conversum de salute diffidere, cum ipse de publicano in apostolum, de theloneario in evangelistam sit repente mutatus.
\newline {\color{Red} Lectio III. vel IX.}
\lettrine[lines=2]{\zallmancaps \color{Red} E}{t} relictis omnibus surgens secutus est eum. Intelligens ergo Matheus quid sit veraciter dominum sequi, relictis omnibus, sequitur. Sequi enim imitari est. Ideoque ut pauperem Christum non tam gressu quam affectu consectari potuisset: reliquit propria qui rapere solebat aliena.
\newline {\color{Red} Ad Vesperas. Sicut unius evangeliste. Oratio.}
\lettrine[lines=2]{\zallmancaps \color{Blue} D}{a} nobis quesumus omnipotens deus, ut beati Mathei apostoli tui et evangeliste quam prevenimus veneranda sollemnitas, et devotionem nobis augeat et salutem. Per.
{\ubsubsection{In Die Sancti Mathei}{in-die-sancti-mathei-breviarium}}
\newline {\color{Red} Ad Mat. Lectio VII. Secundum Matheum.}
\lettrine[lines=2]{\zallmancaps \color{Red} I}{n} illo tempore: Cum transiret Ihesus, vidit hominem sedentem in theloneo Matheum nomine et ait illi. Sequere me. Et reliqua.
\newline {\color{Red} Omelia Venerabilis Bede Presbiteri. \grecross .}
\lettrine[lines=2]{\zallmancaps \color{Blue} E}{x} lectione evangelica audivimus, quia sedentem in theloneo Matheum, ac temporalibus curis intentum, repente dominus miseratus vocavit, et de publicano iustum, de theloneario discipulum fecit. Quem eciam proficientibus eiusdem gratie incrementis, de communi discipulorum numero ad apostolatus gradum promovit.
\newline {\color{Red} Lectio VIII.}
\lettrine[lines=2]{\zallmancaps \color{Red} N}{ec} solum predicandi, verum eciam scribendi evangelium illi ministerium commisit, ut qui terrestrium curam omiserat, celestium esse talentorum dispensator inciperet. Quod ob id utique providentia superna fieri disposuit, ut nullum enormitas peccatorum, nullum numerositas scelerum, a speranda venia revocaret, cum hunc tantis mundi nexibus absolutum, tam celestem factum conspiceret, ut communicato cum angelicis spiritibus vocabulo: evangelista nominaretur et esset.
\newline {\color{Red} Lectio IX.}
\lettrine[lines=2]{\zallmancaps \color{Blue} V}{idit} inquit Ihesus hominem sedentem in theloneo: Matheum nomine. Vidit autem non tam corporeo intuitu quam interne miserationis aspectibus, quibus et Petrum negantem ut reatum suum agnoscere ac deflere posset, aspicere dignatus est: quibus populum suum Egyptiaca servitute depressum ut eriperet, respexit, dicens ad Moysen. Videns vidi afflictionem populi mei qui est in Egipto, et gemitum eorum audivi, et descendi liberare eos.
\newline {\color{Red} Ad Laudes et alias horas. Oratio.}
\lettrine[lines=2]{\zallmancaps \color{Red} B}{eati} Mathei apostoli et evangeliste domine precibus adiuvemur, ut quod possibilitas nostra non obtinet, eius nobis intercessione donetur. Per. \ul{Cetera sicut in communi unius evangeliste.}
{\ubsubsection{Sanctorum Mauritii Sociorumque Eius Martirum}{sanctorum-mauritii-sociorumque-eius-martirum-breviarium}}
\newline {\color{Red} Oratio.}
\lettrine[lines=2]{\zallmancaps \color{Blue} A}{nnue} quesumus omnipotens deus, ut nos sanctorum martyrum tuorum, Mauritii Exuperii, Candidi, Vitalis, Innocentii ac sociorum eorumdem letificet festiva sollemnitas, ut quorum suffragiis nitimur, nataliciis gloriemur. Per.
{\ubsubsection{Sanctorum Cosme et Damiani Martirum}{sanctorum-cosme-et-damiani-martirum-breviarium}}
\newline {\color{Red} Oratio.}
\lettrine[lines=2]{\zallmancaps \color{Red} P}{resta} quesumus omnipotens deus, ut qui sanctorum martyrum tuorum Cosme et Damiani natalicia colimus, a cunctis malis imminentibus eorum intercessionibus liberemur. Per.
{\ubsubsection{Sancti Michaelis Archangeli}{sancti-michaelis-archangeli-breviarium}}
\newline {\color{Red} Ad Vesperas. Super psalmos. Ant.} Dum committeret bellum draco cum Michaele archangelo audita est vox milia milium dicentium salus deo nostro. {\color{Red} Capitulum.}%typo: Michale
\lettrine[lines=2]{\zallmancaps \color{Blue} S}{ignificavit} deus que oportet fieri cito, loquens per angelum suum servo suo Iohanni, qui testimonium perhibuit verbo dei, et testimonium Ihesu Christi quecumque vidit. \R \hyperlink{te-sanctum-dominum}{Te sanctum.} {\color{Red} IX. Ymnus.}
\hymn{tibi-christe-splendor}{Tibi Christe splendor}{
\lettrine[lines=2]{\zallmancaps \color{Red} T}{ibi} Christe splendor patris,
vita, virtus cordium,
in conspectu angelorum
votis, voce, psallimus,
alternantes concrepando
melos damus vocibus.
\par {\bfseries \color{Blue} C}ollaudamus venerantes
omnes celi milites,
sed precipue primatem
celestis exercitus
Michaelem in virtute
conterentem Zabulon.
\par {\bfseries \color{Red} Q}uo custode procul pelle
rex Christe piissime
omne nefas inimici
mundos corde et corpore,
paradiso redde tuo
nos sola clementia.
\par {\bfseries \color{Blue} G}loriam patri melodis
personemus vocibus,
gloriam Christo canamus,
gloriam paraclito,
qui trinus et unus deus
extat ante secula. Amen.
}
\newline \V Stetit angelus iuxta aram templi. \R Habens thuribulum aureum in manu sua.
\newline {\color{Red} Ad Mag. Ant.} Dum sacrum mysterium cerneret Iohannes, archangelus Michael tuba cecinit dignus es domine deus noster accipere librum et solvere signacula eius, alleluya. {\color{Red} Oratio.}
\lettrine[lines=2]{\zallmancaps \color{Blue} D}{eus} qui miro ordine angelorum ministeria hominumque dispensas, concede propitius, ut quibus tibi ministrantibus in celo semper assistitur ab his in terra vita nostra muniatur. Per.
\newline {\color{Red} Ad Mat. Invit.}
\begin{invitatory}[regem-angelorum-invitatorium]
{Regem angelorum dominum, * Venite adoremus.}
\end{invitatory}
\newline {\color{Red} Ymnus.} \hyperlink{tibi-christe-splendor}{Tibi Christe.}
\newline {\color{Red} In Primo Nocturno. Ant.} Stetit angelus iuxta aram templi habens thuribulum aureum in manu sua. {\color{Red} Ps.} \psalm{5} {\color{Red} Ant.} Ascendit fumus aromatum in conspectu domini de manu angeli. {\color{Red} Ps.} \psalm{8} {\color{Red} Ant.} Michael archangele veni in adiutorium populo dei. {\color{Red} Ps.} \psalm{10} \V Stetit angelus.
\newline \ul{Ex hystoria dedicationis ecclesie beati Michaelis.} \grecross \ {\color{Red} Lectio Prima.}
\lettrine[lines=2]{\zallmancaps \color{Red} E}{st} locus in Campanie finibus, ubi inter sinum Adriaticum et montem Garganum civitas Sepontus posita est, qui in supremo cacumine beati archangeli Michaelis gestat ecclesiam. Hanc mortalibus hoc modo cognitam, libellus in ecclesia eadem positus indicat. \R
\lettrine[lines=2]{\zallmancaps \color{Blue} F}{actum} \hypertarget{factum-est-silentium}{\label{factum-est-silentium}} est silentium in celo dum committeret bellum draco cum Michaele archangelo, audita est vox milia milium dicentium, salus, honor et virtus, * Omnipotenti deo. \V Milia milium ministrabant ei, et decies centena milia assistebant ei. Omnipotenti.
\newline {\color{Red} Lectio Secunda.}
\lettrine[lines=2]{\zallmancaps \color{Red} E}{rat} in illa civitate quidam predives nomine Garganus. Huius dum pecora per devia montis latera pascerentur, contigit taurum gregis consortia spernentem, singularem incedere et ad ultimum redeuntibus pecudibus domum non esse regressum.
\begin{responsory}[stetit-angelus-iuxta]
{Stetit angelus iuxta aram templi, habens thuribulum aureum in manu sua, et data sunt ei incensa multa, * Et ascendit fumus aromatum de manu angeli in conspectu domini.}{Factum est silentium in celo et accepit angelus thuribulum et implevit illud de igne altaris.}
\end{responsory}%
\newline {\color{Red} Lectio Tertia.}
\lettrine[lines=2]{\zallmancaps \color{Blue} C}{um} igitur dominus cum multitudine servorum hunc requireret, invenit tandem in vertice montis, in foribus cuiusdam spelunce iraque permotus arrepto arcu appetit illum sagita toxicata. Que velut venti flamine retorta, eum mox reversa percussit.
\begin{responsory-final}[hic-est-michael]
{Hic est Michael archangelus princeps militie angelorum, * Cuius honor prestat beneficia populorum, * Et oratio perducit ad regna celorum.}{Archangele Christi per gratiam quam meruisti te deprecamur ut eripias nos a laqueo mortis.}{Et oratio}
\end{responsory-final}%
\newline {\color{Red} In Secundo Nocturno. Ant.} Michael prepositus paradisi quem honorificant angelorum cives. {\color{Red} Ps.} \psalm{14} {\color{Red} Ant.} Gloriosus apparuisti in conspectu domini, propterea decorem induit te dominus. {\color{Red} Ps.} \psalm{29} {\color{Red} Ant.} Angelus, archangelus, Michael dei nuntius, de animabus iustis suscipiendis. {\color{Red} Ps.} \psalm{46} \V Ascendit fumus aromatum in conspectu domini. \R De manu angeli.
\newline {\color{Red} Lectio Quarta.}
\lettrine[lines=2]{\zallmancaps \color{Red} T}{urbati} super hoc cives consulunt episcopum. Qui indicto ieiunio triduano, quidnam esset novit a deo esse querendum. Quo peracto: archangelus episcopum per visionem alloquitur dicens. Scias hoc mea gestum voluntate. Ego enim sum Michael, qui locum hunc in terris incolere instituens hoc volui probari indicio, me loci esse custodem.
\begin{responsory}[venit-michael]
{Venit Michael archangelus cum multitudine angelorum cui tradidit deus animas sanctorum, * Ut perducat eas in paradisum exultationis.}{Data est potestas archangelo Michaeli super animas sanctorum.}
\end{responsory}%
\newline {\color{Red} Lectio Quinta.}
\lettrine[lines=2]{\zallmancaps \color{Blue} T}{unc} consuetudinem fecerunt cives, hic dominum precibus sanctumque deposcere Michaelem. Sed nec intra criptam intrare ausi, pro foribus cotidie orationi vacabant. Interea pagani Sepontinos bello lacessere temptant. Qui antistitis monitis triduo petunt inducias, ut triduano ieiunio liceret sancti Michaelis presidium implorare.
\begin{responsory}[archangeli-michaelis]
{Archangeli Michaelis interventione suffulti te domine deprecamur, * Ut quos honore prosequimur contingamus et mente.}{Perpetuum nobis domine tue miserationis presta subsidium quibus et angelica prestitisti suffragia non deesse.}
\end{responsory}%
\newline {\color{Red} Lectio Sexta.}
\lettrine[lines=2]{\zallmancaps \color{Red} E}{cce} autem ipsa nocte que belli precedebat diem, adest Michael, antistitis preces dicit exauditas. Mane ergo obviant Christiani paganis, atque in primo belli apparatu mons Garganus immenso tremore concutitur, fulgura crebra volant, et caligo tenebrosa totum montis cacumen obduxit.
\begin{responsory-final}[fidelis-sermo-et]
{Fidelis sermo et omni acceptione dignus, Michael archangelus pugnavit cum diabolo, * Gratia dei ille victor in celis extitit, * Et illi hosti antiquo patuit ruina magna.}{Gaudent angeli, letantur archangeli, et collaudant in celis filium dei.}{Et illi}
\end{responsory-final}%
\newline {\color{Red} In Tertio Nocturno. Ant.} Concussum est mare et contremuit terra, ubi archangelus Michael descendebat de celo. {\color{Red} Ps.} \psalm{96} {\color{Red} Ant.} Data sunt ei incensa multa ut adoleret ea ante altare aureum quod est ante oculos domini. {\color{Red} Ps.} \psalm{98} {\color{Red} Ant.} Laudemus dominum quem laudant angeli, cui cherubim et seraphim sanctus sanctus sanctus proclamant. {\color{Red} Ps.} \psalm{102} \V In conspectu angelorum psallam tibi deus meus. \R Adorabo ad templum sanctum tuum et confitebor nomini tuo.%typo confitobor
\newline {\color{Red} Lectio VII.}
\lettrine[lines=2]{\zallmancaps \color{Blue} F}{ugiunt} pagani et qui evasere videntes sexcentos ferme suorum fulmine interemptos, Christo continuo colla summittunt. Cumque victores reversi, vota domino ad templum referunt archangeli, vident iuxta ianuam septemtrionalem instar pusille portelle et quasi hominis vestigia marmori impressa. Agnoscuntque beatum Michaelem per hoc presentie sue signum voluisse monstrare.
\begin{responsory}[in-conspectu-gentium]
{In conspectu gentium nolite timere, vos enim in cordibus vestris adorate et timete dominum, * Angelus enim eius vobiscum est.}{Cantate ei canticum novum bene psallite ei in vociferatione.}
\end{responsory}%
\newline {\color{Red} Lectio VIII.}
\lettrine[lines=2]{\zallmancaps \color{Red} M}{ota} est dubitatio inter Sepontinos, quid de loco agant, an intrare debeant vel dedicare illic ecclesiam. Nocte ergo constituti ieiunii, ad hoc a deo revelandum, Michael Sepontino episcopo apparens, non est inquit opus hanc dedicare ecclesiam. Ipse enim qui condidi: dedicavi.
\begin{responsory}[in-conspectu-angelorum]
{In conspectu angelorum psallam tibi et adorabo ad templum sanctum tuum, * Et confitebor nomini tuo domine.}{Deus meus es tu et confitebor tibi deus meus es tu et exaltabo te.}
\end{responsory}%
\newline {\color{Red} Lectio IX.}
\lettrine[lines=2]{\zallmancaps \color{Blue} E}{t} addidit. Vos tantum intrate, et me astante iuxta morem patrono: precibus locum frequentate. Adveniunt illi mane et intrant, missarumque celebratione completa redierunt in sua. Hic est ergo locus in quo crebri sanantur egroti, et multa angelica potestate fieri miracula comprobantur.
\begin{responsory-final}[te-sanctum-dominum]
{Te sanctum dominum in excelsis laudant omnes angeli dicentes, * Te decet laus * Et honor domine.}{Cherubin quoque et seraphin sanctus proclamant, et omnis celicus ordo dicentes.}{Et honor}
\end{responsory-final}%
\newline \V Gloriosus apparuisti in conspectu domini. \R Propterea decorem induit te dominus.
\newline {\color{Red} In Laudibus. Ant.} Dum preliaretur Michael archangelus cum dracone audita est vox dicentium salus deo nostro alleluya. {\color{Red} Ant.} Dum committeret bellum draco cum Michaele archangelo audita est vox milia milium dicentium salus deo nostro. {\color{Red} Ant.} Archangele Michael constitui te principem super omnes animas suscipiendas. {\color{Red} Ant.} Angeli domini dominum benedicite in eternum. {\color{Red} Ant.} Angeli archangeli, throni et dominationes, principatus et potestates, virtutes celorum laudate dominum de celis alleluya. {\color{Red} Capitulum.} Significavit. {\color{Red} Ymnus.}
\hymn{christe-sanctorum-decus}{Christe sanctorum decus}{
\lettrine[lines=2]{\zallmancaps \color{Red} C}{hriste} sanctorum decus angelorum,
rector humani generis et rector,
nobis eternum tribue benignus
scandere regnum.
\par {\bfseries \color{Blue} A}ngelum pacis Michael ad istam,
celitus mitti rogitamus aulam,
nobis ut crebro veniente crescant
prospera cuncta.
\par {\bfseries \color{Red} A}ngelus fortis Gabriel ut hostem
pellat antiquum volitet ab alto
sepius templum veniens in istud
visere nostrum.
\par {\bfseries \color{Blue} A}ngelum nobis medicum salutis
mitte de celis Raphael ut omnes
sanet egrotos pariterque nostros
dirigat actus.
\par {\bfseries \color{Red} H}inc dei nostri genitrix Maria
totus et nobis chorus angelorum
semper assistat simul et beata
contio tota.
\par {\bfseries \color{Blue} P}restet hoc nobis deitas beata
patris ac nati, pariterque sancti
spiritus, cuius reboat in omni
gloria mundo. Amen.
}
\newline \V Adorate deum. \R Omnes angeli eius.
\newline {\color{Red} Ad Ben. Ant.} Factum est silentium in celo dum draco committeret bellum et Michael pugnavit cum eo et fecit victoriam alleluya.
\newline {\color{Red} Ad Terciam. Cap.} Significavit.
\begin{responsory-breve}[stetit-angelus-breve]
{Stetit angelus * Iuxta aram templi.}{Habens thuribulum aureum in manu sua.}
\end{responsory-breve}%
\V Ascendit fumus aromatum.
\newline {\color{Red} Ad Sextam. Cap.}
\lettrine[lines=2]{\zallmancaps \color{Blue} F}{actum} est prelium in celo, Michael et angeli eius preliabantur cum drachone, et draco pugnabat et angeli eius, et non valuerunt neque locus inventus est eorum amplius in celo.
\begin{responsory-breve}[ascendit-fumus-breve]
{Ascendit fumus aromatum * In conspectu domini.}{De manu angeli.}
\end{responsory-breve}%
\V In conspectu angelorum.
\newline {\color{Red} Ad Nonam. Capitulum.}
\lettrine[lines=2]{\zallmancaps \color{Red} A}{udivi} vocem angelorum multorum in circuitu throni et animalium, et seniorum, et erat numerus eorum milia milium voce magna dicentium salus deo nostro.
\begin{responsory-breve}[in-conspectu-angelorum-breve]
{In conspectu angelorum * Psallam tibi deus meus.}{Adorabo ad templum sanctum tuum et confitebor nomini tuo.}
\end{responsory-breve}%
\V Adorate deum.
\newline {\color{Red} Ad Vesperas. Super psalmos. Ant.} Dum preliaretur. {\color{Red} Ps.} \psalm{109} \psalm{110} \psalm{111} \psalm{112} \psalm{137} \ul{Cap., Ymnus, \Vbar , ut in Primis Vesperis.}
\newline {\color{Red} Ad Mag. Ant.} Archangeli Michaelis interventione suffulti te domine deprecamur, ut quos honore prosequimur contingamus et mente, alleluya.
{\ubsubsection{Sancti Hieronimi Presbiteri et Confessoris}{sancti-hieronimi-presbiteri-et-confessoris-breviarium}}
\newline {\color{Red} Oratio.}
\lettrine[lines=2]{\zallmancaps \color{Blue} D}{eus} qui nos beati Hieronimi confessoris tui annua sollemnitate letificas concede propicius, ut cuius natalicia colimus, per eius ad te exempla gradiamur. Per.
{\ubsubsection{Sancti Remigii Episcopi et Confessoris}{sancti-remigii-episcopi-et-confessoris-breviarium}}
\newline {\color{Red} Oratio.}
\lettrine[lines=2]{\zallmancaps \color{Red} D}{eus} qui populo tuo eterne salutis beatum Remigium ministrum tribuisti, presta quesumus, ut quem doctorem vite habuimus in terris, intercessorem habere mereamur in celis. Per.
\newline {\color{Red} Ad Matutinas. Ymnus dicatur sic.}
\lettrine[lines=2]{\zallmancaps \color{Blue} I}{ste} confessor domini sacratus,
sobrius castus fuit et quietus,
vita dum presens vegetavit eius
corporis artus.
\par {\bfseries \color{Blue} A}d sacrum, etc. \ul{Cetera de communi.}
{\ubsubsection{Sancti Leodegarii Episcopi et Martiris}{sancti-leodegarii-episcopi-et-martiris-breviarium}}
\newline {\color{Red} Oratio.}
\lettrine[lines=2]{\zallmancaps \color{Red} L}{etetur} ecclesia tua deus, beati Leodegarii martyris tui atque pontificis confisa suffragiis, atque eius precibus gloriosis et devota permaneat et secura consistat. Per.
{\ubsubsection{Sancti Francisci Confessoris}{sancti-francisci-confessoris-breviarium}}
\newline {\color{Red} Oratio.}
\lettrine[lines=2]{\zallmancaps \color{Blue} D}{eus} qui ecclesiam tuam beati Francisci meritis, fetu nove prolis amplificas, tribue nobis ex eius imitatione terrena despicere, et celestium donorum semper participatione gaudere. Per.
{\ubsubsection{Sancti Marci Pape}{sancti-marci-pape-breviarium}}
\newline {\color{Red} Oratio.}
\lettrine[lines=2]{\zallmancaps \color{Red} E}{xaudi} domine preces nostras, et interveniente beato Marco confessore tuo atque pontifice, supplicationes nostras placatus intende. Per.
{\ubsubsection{Sanctorum Sergii Bachi, Marcelli et Apulei Martirum}{sanctorum-sergii-bachi-marcelli-et-apulei martirum-breviarium}}
\newline {\color{Red} Oratio.}
\lettrine[lines=2]{\zallmancaps \color{Blue} S}{anctorum} tuorum nos domine Sergii et Bachi, Marcelli, et Apulei, beata merita prosequantur et tuo semper faciant amore ferventes. Per.
{\ubsubsection{Sanctorum Dionisii Sociorumque Eius}{sanctorum-dionisii-sociorumque-eius-breviarium}}
\newline {\color{Red} Oratio.}
\lettrine[lines=2]{\zallmancaps \color{Red} D}{eus} qui hodierna die beatum Dyonisium virtute constantie in passione roborasti, quique illi ad predicandum gentibus gloriam tuam Rusticum et Eleutherium sociare dignatus es, tribue nobis quesumus ex eorum imitatione pro amore tuo prospera mundi despicere, et nulla eius adversa formidare. Per.
{\ubsubsection{Sancti Calixti Pape et Martiris}{sancti-calixti-pape-et-martiris-breviarium}}%typo Martirum
\newline {\color{Red} Oratio.}
\lettrine[lines=2]{\zallmancaps \color{Blue} D}{eus} qui conspicis nos ex nostra infirmitate deficere ad amorem tuum nos misericorditer per sanctorum tuorum exempla restaura. Per.
{\ubsubsection{Sancti Luce Evangeliste}{sancti-luce-evangeliste-breviarium}}
\newline {\color{Red} Oratio.}
\lettrine[lines=2]{\zallmancaps \color{Red} I}{nterveniat} pro nobis domine quesumus sanctus tuus Lucas evangelista, qui crucis mortificationem iugiter in suo corpore pro tui nominis amore portavit. Per.
{\ubsubsection{Sanctarum Undecim Milium Virginum et Martirum}{sanctarum-undecim-milium-virginum-et-martirum-breviarium}}
\newline {\color{Red} Oratio.}
\lettrine[lines=2]{\zallmancaps \color{Blue} O}{mnipotens} sempiterne deus qui infirma mundi eligis ut fortia queque confundas, concede propitius, ut qui beatarum virginum et martyrum tuarum sollemnia colimus, earum apud te patrocinia sentiamus. Per.
{\ubsubsection{Sanctorum Crispini et Crispiniani}{sanctorum-crispini-et-crispiniani-breviarium}}
\newline {\color{Red} Oratio.}
\lettrine[lines=2]{\zallmancaps \color{Red} D}{eus} qui nos concedis sanctorum martyrum tuorum Crispini et Crispiniani natalicia colere, da nobis in eterna beatitudine de eorum societate gaudere. Per.
{\ubsubsection{In Vigilia Apostolorum Symonis et Iude}{in-vigilia-apostolorum-symonis-et-iude-breviarium}}
\newline {\color{Red} Ad Vesperas. Oratio.}
\lettrine[lines=2]{\zallmancaps \color{Blue} C}{oncede} quesumus omnipotens deus, ut sicut apostolorum tuorum Symonis et Iude gloriosa natalicia prevenimus, sic ad tua beneficia promerenda, maiestatem tuam pro nobis ipsi preveniant. Per.
{\ubsubsection{In Die Apostolorum Symonis et Iude}{in-die-apostolorum-symonis-et-iude-breviarium}}
\newline {\color{Red} In Die. Oratio.}
\lettrine[lines=2]{\zallmancaps \color{Red} D}{eus} qui nos per beatos apostolos tuos Symonem et Iudam, ad agnitionem tui nominis venire tribuisti, da nobis eorum gloriam sempiternam, et proficiendo celebrare et celebrando proficere. Per.
{\ubsubsection{Sancti Quintini Martyris}{sancti-quintini-martyris-breviarium}}
\newline {\color{Red} Oratio.}
\lettrine[lines=2]{\zallmancaps \color{Blue} A}{desto} domine supplicationibus nostris, ut qui ex iniquitate nostra reos nos esse cognoscimus, beati Quintini martyris tui intercessione liberemur. Per.
{\ubsubsection{In Festivitate Omnium Sanctorum}{in-festivitate-omnium-sanctorum-breviarium}}
\newline {\color{Red} Ad Vesperas. Super psalmos. Ant.} O quam gloriosum est regnum, in quo cum Christo gaudent omnes sancti, amicti stolis albis sequuntur agnum quocumque ierit. {\color{Red} Ps.} \psalm{112} \ul{cum ceteris.} {\color{Red} Cap.}
\lettrine[lines=2]{\zallmancaps \color{Red} E}{cce} ego Iohannes vidi alterum angelum ascendentem ab ortu solis, habentem signum dei vivi, et clamavit voce magna quatuor angelis quibus datum est nocere terre et mari dicens, nolite nocere terre et mari neque arboribus, quoadusque signemus servos dei nostri in frontibus eorum. \R \hyperlink{concede-nobis-domine}{Concede.} {\color{Red} IX. Ymnus.}
\hymn{ihesu-salvator-seculi}{Ihesu salvator seculi}{
\lettrine[lines=2]{\zallmancaps \color{Blue} I}{hesu} salvator seculi,
redemptis ope subveni,
et pia dei genetrix
salutem posce miseris.
\par {\bfseries \color{Red} C}etus omnes angelici,
patriarcharum cunei,
ac prophetarum merita:
nobis precentur veniam.
\par {\bfseries \color{Blue} B}aptista Christi previus,
et claviger ethereus,
cum ceteris apostolis
nos solvant nexu criminis.
\par {\bfseries \color{Red} C}horus sacratus martyrum,
confessio sacerdotum,
et virginalis castitas:
nos a peccatis abluant.
\par {\bfseries \color{Blue} E}lectorum suffragia,
omnesque cives celici,
annuant votis supplicum:
et vite poscant premium.
\par {\bfseries \color{Red} L}aus honor virtus gloria
deo patri et filio
sancto spiritu Paraclito
in sempiterna secula. Amen.
}
\newline \V Letamini in domino.
\newline {\color{Red} Ad Mag. Ant.} Beati estis sancti dei omnes qui meruistis consortes fieri celestium virtutum et perfrui claritatis gloria, ideoque precamur ut memores nostri intercedere dignemini pro nobis ad dominum deum nostrum. {\color{Red} Oratio.}
\lettrine[lines=2]{\zallmancaps \color{Red} D}{omine} deus noster multiplica super nos gratiam tuam, et quorum prevenimus gloriosa sollemnia, tribue subsequi in sancta professione leticiam. Per.
\newline \ul{Ad Completorium antiphone.} Miserere \ul{et} Salve nos.
\newline {\color{Red} Ad Mat. Invitat.}
\begin{invitatory}[regem-regum-dominum-invitatorium]
{Regem regum dominum venite adoremus, * Quia ipse est corona sanctorum omnium.}
\end{invitatory}
\newline {\color{Red} Ymnus.} \hyperlink{ihesu-salvator-seculi}{Ihesu salvator.}
\newline {\color{Red} In Primo Nocturno. Ant.} Adesto deus unus omnipotens pater et filius et spiritus sanctus. {\color{Red} Ps.} \psalm{8} {\color{Red} Ant.} Sicut lilium inter spinas, sic amica mea inter filias. {\color{Red} Ps.} \psalm{18} {\color{Red} Ant.} Laudemus dominum quem laudant angeli, cui cherubin et seraphin sanctus sanctus sanctus proclamant. {\color{Red} Ps.} \psalm{102} \V Letamini in domino.
\newline \ul{Sermo Alcuini vel magistri Albini secundum quosdam.} \grecross \ {\color{Red} Lectio Prima.}
% https://mlat.uzh.ch/browser/cps_2.AuInAuH.SeSuDeS2#:22.1;2
\lettrine[lines=2]{\zallmancaps \color{Blue} H}{odie} dilectissimi omnium sanctorum sub una sollemnitatis leticia celebramus festivitatem. O vere beata mater ecclesia, quam sic honor divine dignationis illuminat, quam vincentium martyrum gloriosus sanguis exornat, quam inviolate confessionis candida induit virginitas. Floribus eius nec rose nec lilia desunt. Certe nunc karissimi singuli ad utrosque honores amplissimas accipiunt dignitatum coronas, vel de
%pp132
virginitate candidas, vel de passione purpureas. \R
\lettrine[lines=2]{\zallmancaps \color{Red} S}{umme} trinitati, simplici deo, una divinitas, equalis gloria, coeterna maiestas, patri prolique sanctoque flamini, * Qui totum subdit suis orbem legibus. \V Prestet nobis gratiam deitas beata patris ac nati pariterque spiritus almi. Qui.
\newline {\color{Red} Lectio Secunda.}
\lettrine[lines=2]{\zallmancaps \color{Blue} I}{n} celestibus castris pax et acies habent flores suos, quibus milites Christi coronantur. Dei enim ineffabilis et immensa bonitas eciam hoc previdit ut laborum quidem tempus et agonis non extenderet, nec longum faceret aut eternum, sed breve et ut ita dicam momentaneum, ut in hac brevi et exigua vita: agones essent et labores, in illa vero celesti vita, eterna esset corona et premia meritorum.
\begin{responsory}
{Felix namque es sacra virgo Maria et omni laude dignissima, * Quia ex te ortus est sol iusticie Christus deus noster.}{Ora pro populo interveni pro clero intercede pro devoto femineo sexu, sentiant omnes tuum iuvamen quicumque celebrant tuam sollemnitatem.}
\end{responsory}%
\newline {\color{Red} Lectio Tertia.}
\lettrine[lines=2]{\zallmancaps \color{Red} P}{revidit} quidem, ut labores cito finirentur, meritorum vero premia sine fine durarent, ut post huius mundi tenebras visuri essent candidissimam lucem et accepturi maiorem cunctarum passionum acerbitatibus beatitudinem testante hoc apostolo ubi ait. Non sunt condigne passiones huius temporis ad superventuram gloriam que revelabitur in nobis.
\begin{responsory-final}
{Te sanctum dominum in excelsis laudant omnes angeli dicentes, * Te decet laus * Et honor domine.}{Cherubin quoque et seraphin, sanctus proclamant, et omnis celicus ordo dicentes.}{Et honor}
\end{responsory-final}%
\newline {\color{Red} In Secundo Nocturno. Ant.} Inter natos mulierum non surrexit maior Iohanne baptista. {\color{Red} Ps.} \psalm{91} {\color{Red} Ant.} Estote fortes in bello et pugnate cum antiquo serpente et accipietis regnum eternum alleluya. {\color{Red} Ps.} \psalm{33} {\color{Red} Ant.} Isti sunt sancti qui pro dei amore minas hominum contempserunt, sancti martyres in regno celorum exultant cum angelis, o quam preciosa est mors sanctorum qui assidue assistunt ante dominum et ab invicem non sunt separati. {\color{Red} Ps.} \psalm{45} \V Exultent iusti.
\newline {\color{Red} Lectio Quarta.}
\lettrine[lines=2]{\zallmancaps \color{Blue} Q}{uam} leto sinu de prelio revertentes, celestis excipit leticia, de hoste prostrato trophea ferentibus, occurrit. Cum triumphantibus viris et femine veniunt, que cum seculo sexum quoque vicerunt, et geminata gloria militie, virgines cum pueris teneros annos virtutibus transeuntes, sed et cetera fidelium turba aule perpetue regiam intravit, qui sinceritatem fidei inconcussis celestium preceptorum disciplinis: unica pace observaverunt.
\begin{responsory}
{Inter natos mulierum non surrexit maior Iohanne baptista, * Qui viam domino preparavit in heremo.}{Fuit homo missus a deo cui nomen Iohannes erat.}
\end{responsory}%
\newline {\color{Red} Lectio Quinta.}
\lettrine[lines=2]{\zallmancaps \color{Red} E}{rgo} agite fratres, aggrediamur iter vite, revertamur ad civitatem celestem in qua scripti sumus et cives decreti. Non sumus hospites sed cives sanctorum et domestici Dei, eciam illius heredes, coheredes autem Christi. Huius urbis nobis ianuas aperiet fortitudo, et fiducia latum prebebit ingressum.
\begin{responsory}%[qui-sunt-isti]
{Qui sunt isti qui ut nubes volant, * Et quasi columbe ad fenestras suas.}{Candidiores nive nitidiores lacte, rubicundiores ebore antiquo.}%typo atiquo, vive
\end{responsory}%
\newline {\color{Red} Lectio Sexta.}
\lettrine[lines=2]{\zallmancaps \color{Blue} Q}{uid} hac vita beatius, ubi non est paupertatis metus, non egritudinis imbecillitas? Nemo ibi leditur, nemo irascitur, nemo invidet, cupiditas nulla exardescit, nullum cibi desiderium, nulla honoris pulsat aut potestatis ambitio? nullus ibi diaboli metus, insidie demonum nulle, terror Gehenne procul. Mors neque corporis neque anime sed immortalitatis munere vita iocunda. Nulla erit tunc usquam discordia, sed cuncta consona. Tranquilla sunt omnia, iugis splendor.
\begin{responsory-doxology}%[laverunt-stolas]
{Laverunt stolas suas et candidas eas fecerunt * In sanguine agni.}{Isti sunt qui venerunt ex magna tribulatione et laverunt stolas suas.}
\end{responsory-doxology}%
\newline {\color{Red} In Tertio Nocturno. Ant.} Sint lumbi vestri precincti et lucerne ardentes in manibus vestris, et vos similes hominibus expectantibus dominum suum quando revertatur a nuptiis. {\color{Red} Ps.} \psalm{32} {\color{Red} Ant.} Simile est regnum celorum decem virginibus que accipientes lampades suas exierunt obviam sponso et sponse. {\color{Red} Ps.} \psalm{44} {\color{Red} Ant.} Laudem dicite deo nostro omnes sancti eius et qui timetis deum pusilli et magni quoniam regnavit dominus deus noster omnipotens gaudemus et exultemus et demus gloriam ei. {\color{Red} Ps.} \psalm{83} \V Iusti autem.
\newline {\color{Red} Lectio VII.}
\lettrine[lines=2]{\zallmancaps \color{Red} V}{erum} super hec omnia est consociari angelorum et archangelorum cetibus, omniumque celestium supernarum virtutum contuberniis perfrui, et intueri agmina sanctorum splendidius syderibus micantia, patriarcharum fide fulgentia, prophetarum spe letantia, apostolorum in duodecim tribus Israhel orbem iudicantia, martyrum purpureis coronis victorie lucentia, virginum choros candentia serta gestantes inspicere.
\begin{responsory}%[sint-lumbi-vestri]
{Sint lumbi vestri precincti et lucerne ardentes in manibus vestris, * Et vos similes hominibus expectantibus dominum suum quando revertatur a nuptiis.}{Vigilate ergo quia nescitis qua hora dominus vester venturus sit.}
\end{responsory}%
\newline {\color{Red} Lectio VIII.}
\lettrine[lines=2]{\zallmancaps \color{Blue} D}{e} rege autem qui horum medius residet dicere, vox nulla sufficit. Si enim cotidie nos oporteret tormenta perferre, si ipsam Gehennam parvo tempore tolerare, ut Christum videre possemus in gloria venientem, et sanctorum eius numero sociari, nonne dignum erat pati omne quod triste est, ut tanti boni tanteque glorie participes haberemur?
\begin{responsory}%[audivi-vocem-de]
{Audivi vocem de celo venientem venite omnes virgines sapientissime, * Oleum recondite in vasis vestris dum sponsus advenerit.}{Media nocte clamor factus est ecce sponsus venit.}%typo: recundite
\end{responsory}%
\newline {\color{Red} Lectio IX.}
\lettrine[lines=2]{\zallmancaps \color{Red} Q}{ue} erit fratres karissimi illa iustorum gloria, quam grandis leticia: cum dominus ceperit pro terrenis celestia, pro temporalibus sempiterna, pro modicis magna prestare, et adducere secum sanctos in visione paterne glorie, et facere secum in celestibus consedere, suisque immortalitatem largiri, ad quam eos sui sanguinis vivificatione reparavit?
\begin{responsory-doxology}[concede-nobis-domine]
{Concede nobis domine quesumus veniam delictorum, et intercedentibus sanctis quorum hodie sollemnia celebramus talem nobis tribue devotionem, * Ut ad eorum pervenire mereamur societatem.}{Adiuvent nos eorum merita quos propria impediunt scelera, excuset intercessio accusat, quos actio, et qui eis tribuisti celestis palmam triumphi nobis veniam non deneges peccati.}
\end{responsory-doxology}%
\newline \V Orate pro nobis omnes sancti dei. \R Ut.
\newline {\color{Red} In Laudibus. Ant.} Vidi turbam magnam quam dinumerare nemo poterat ex omnibus gentibus stantes ante thronum. {\color{Red} Ant.} Et omnes angeli stabant in circuitu throni, et ceciderunt in conspectu eius in facies suas, et adoraverunt eum. {\color{Red} Ant.} Redemisti nos domine deus in sanguine tuo ex omni tribu et lingua et populo et natione et fecisti nos deo nostro regnum. {\color{Red} Ant.} Benedicite dominum omnes electi eius agite diem leticie et confitemini ei. {\color{Red} Ant.} Ymnus omnibus sanctis eius filiis Israhel populo appropinquanti sibi, gloria hec est omnibus sanctis eius. {\color{Red} Capitulum.} Ecce ego Iohannes. {\color{Red} Ymnus.}
\hymn{christe-redemptor-omnium}{Christe redemptor omnium}{
\lettrine[lines=2]{\zallmancaps \color{Blue} C}{hriste} redemptor omnium
conserva tuos famulos,
beate semper virginis
placatus sanctis precibus.
\par {\bfseries \color{Red} B}eata quoque agmina
celestium spirituum,
preterita, presentia,
futura mala pellite.
\par {\bfseries \color{Blue} V}ates eterni iudicis,
apostolique domini,
suppliciter exposcimus
salvari vestris precibus.
\par {\bfseries \color{Red} M}artyres dei incliti,
confessoresque lucidi,
vestris orationibus
nos ferte in celestibus.
\par {\bfseries \color{Blue} C}horus sacrarum virginum,
electorumque omnium,
simul cum sanctis omnibus
consortes Christi facite.
\par {\bfseries \color{Red} G}entem auferte perfidam
credentium de finibus,
ut Christi laudes debitas
persolvamus alacriter.
\par {\bfseries \color{Blue} L}aus honor virtus gloria
deo patri et filio
sancto spiritu Paraclito
in sempiterna secula. Amen.
}
\newline \V Mirabilis deus.
\newline {\color{Red} Ad Ben. Ant.} Te gloriosus apostolorum chorus te prophetarum laudabilis numerus, te martyrum candidatus laudat exercitus, te omnes electi voce confitentur unanimes beata trinitas, unus deus. {\color{Red} Oratio.}
\lettrine[lines=2]{\zallmancaps \color{Red} O}{mnipotens} sempiterne deus qui nos omnium sanctorum merita, sub una tribuisti celebritate venerari, quesumus, ut desideratam nobis tue propitiationis abundantiam, multiplicatis intercessoribus largiaris. Per.
\newline {\color{Red} Ad Terciam. Cap.} Ecce ego Iohannes.
\newline \ul{Responsoria et versiculi plurimorum martyrum dicantur ad horas.}
\newline {\color{Red} Ad Sextam. Cap.}
\lettrine[lines=2]{\zallmancaps \color{Blue} A}{udivi} numerum signatorum centum quadraginta quatuor milia signati, ex omni tribu filiorum Israhel.
\newline {\color{Red} Ad Nonam. Cap.}
\lettrine[lines=2]{\zallmancaps \color{Red} V}{idi} turbam magnam quam dinumerare nemo poterat ex omnibus gentibus et tribubus et populis et linguis, stantes ante thronum.
\newline {\color{Red} Ad Vesperas. Super psalmos. Ant.} Vidi turbam. {\color{Red} Ps.} \psalm{109} \psalm{112} \psalm{115} \psalm{125} \psalm{139} \ul{Cap., Ymnus, \Vbar , ut in Primis Vesperis.}
\newline {\color{Red} Ad Mag. Ant.} Salvator mundi salva nos omnes, sancta dei genitrix virgo semper Maria ora pro nobis, precibus quoque sanctorum, apostolorum, martyrum et confessorum, atque sanctarum virginum, suppliciter petimus, ut a malis omnibus eruamur, bonisque omnibus nunc et semper perfrui mereamur.
\newline \ul{Finitis Vesperis, statim dicantur Vespere Defunctorum, nisi in crastino fuerit Dominica.}
{\ubsubsection{In Commemoratione Omnium Fidelium Defunctorum}{in-commemoratione-omnium-fidelium-defunctorum-breviarium}}
\newline {\color{Red} Ad Vesperas. Ant.} Placebo domino in regione vivorum. {\color{Red} Ps.} \psalm{114} {\color{Red} Ant.} Heu me quia incolatus meus prolongatus est. {\color{Red} Ps.} \psalm{119} {\color{Red} Ant.} Dominus custodit te ab omni malo custodiat animam tuam dominus. {\color{Red} Ps.} \psalm{120} {\color{Red} Ant.} Si iniquitates observaveris domine, domine quis sustinebit. {\color{Red} Ps.} \psalm{129} {\color{Red} Ant.} Opera manuum tuarum domine ne despicias. {\color{Red} Ps.} \psalm{137} {\color{Red} Ant.} Audivi vocem de celo dicentem beati mortui qui in domino moriuntur. {\color{Red} Canticum.} \psalm{magnificat} \ul{Finita antiphona post} Magnificat, \ul{dicatur,} Pater noster. \ul{Quo finito non pronuncietur,} Et ne nos, \ul{sed dicatur psalmus,} \psalm{145} \ul{et per,} Requiem, \ul{terminetur. Deinde sacerdos,} \V A porta inferi. \R Erue domine animas eorum. Dominus vobiscum. Oremus. {\color{Red} Oratio.}
\lettrine[lines=2]{\zallmancaps \color{Blue} F}{idelium} deus omnium conditor et redemptor, animabus famulorum famularumque tuarum remissionem cunctorum tribue peccatorum, ut indulgentiam quam semper optaverunt piis supplicationibus consequantur. Qui vivis et regnas cum deo.
\newline \ul{Non dicatur post orationem,} Dominus vobiscum. \ul{Nec,} Benedicamus, \ul{sed tantummodo,} Requiescant in pace. \R Amen.
\newline \ul{In crastino Ad Matutinas et ad ceteras horas defunctorum non fiant prostrationes.}
\newline \ul{Si festivitas Omnium Sanctorum Sabbato evenerit, in crastino totum officium tam diei quam noctis de Dominica erit. Et tunc in eadem die post Vesperas de Dominica: Vespere Defunctorum dicantur ut supra.}
\newline \ul{Ad Matutinas non dicatur,} Domine labia. \ul{Nec,} Deus in adiutorium. \ul{neque Invitat., neque,} Venite, \ul{sed dicto} Pater noster. \ul{Et} Credo, \ul{statim incipiatur.}
\newline {\color{Red} In Primo Nocturno. Ant.} Dirige domine deus meus in conspectu tuo viam meam. {\color{Red} Ps.} \psalm{5} {\color{Red} Ant.} Convertere domine et eripe animam meam quoniam non est in morte qui memor sit tui. {\color{Red} Ps.} \psalm{6} {\color{Red} Ant.} Nequando rapiat ut leo animam meam dum non est qui redimat neque qui salvum faciat. {\color{Red} Ps.} \psalm{7} \V A porta inferi. \ul{Lector non petat benedictionem, sed dicto,} Pater noster, \ul{et subiuncto,} Et ne nos, \ul{ac responso,} Sed libera, \ul{incipiat legere.}
\newline {\color{Red} Lectio Prima.}
\lettrine[lines=2]{\zallmancaps \color{Blue} P}{arce} michi domine, nichil enim sunt dies mei. Quid est homo quia magnificas eum, aut quid apponis erga eum cor tuum? Visitas eum diluculo, et subito probas illum. Usquequo non parcis michi, nec dimittis me ut glutiam salivam meam? Peccavi. Quid faciam tibi o custos hominum? Quare posuisti me contrarium tibi, et factus sum michimetipsi gravis? Cur non tollis peccatum meum, et quare non aufers iniquitatem meam? Ecce nunc in pulvere dormio et si mane me quesieris non subsistam. \ul{In fine lectionis, non dicatur,} Tu autem: \ul{sed finita lectione, statim incipiatur.} \R
\lettrine[lines=2]{\zallmancaps \color{Red} C}{redo} \hypertarget{credo-quod-redemptor}{\label{credo-quod-redemptor}} quod redemptor meus vivit et in novissimo die de terra surrecturus sum, et in carne mea videbo deum, * Salvatorem meum. \V Quem visurus sum ego ipse et non alius, et oculi mei conspecturi sunt. Salvatorem.
\newline {\color{Red} Lectio Secunda.}
\lettrine[lines=2]{\zallmancaps \color{Blue} T}{edet} animam meam vite mee, dimittam me adversum eloquium meum. Loquar in amaritudine anime mee, dicam Deo. Noli me condemnare. Indica michi: cur me ita iudices. Nunquid bonum tibi videtur si calumnieris, et opprimas me opus manuum tuarum, et consilium impiorum adiuves? Nunquid oculi carnei tibi sunt, aut sicut videt homo et tu videbis? Nunquid sicut dies hominis dies tui, et anni tui sicut humana sunt tempora, ut queras iniquitatem meam, et peccatum meum scruteris? Et scias quia nichil impium fecerim: cum sit nemo qui de manu tua possit eruere.
\begin{responsory}[qui-lazarum-resuscitasti]
{Qui Lazarum resuscitasti a monumento fetidum, * Tu eis domine dona requiem et locum indulgentie.}{Qui venturus es iudicare vivos et mortuos et seculum per ignem.}
\end{responsory}%
\newline {\color{Red} Lectio Tertia.}
\lettrine[lines=2]{\zallmancaps \color{Red} M}{anus} tue domine fecerunt me et plasmaverunt me totum in circuitu, et sic repente precipitas me? Memento queso quod sicut lutum feceris me, et in pulverem reduces me. Nonne sicut lac mulsisti me, et sicut caseum me coagulasti? Pelle et carnibus vestisti me: ossibus et nervis compegisti me. Vitam et misericordiam tribuisti michi, et visitatio tua custodivit spiritum meum.
\begin{responsory}[domine-quando]
{Domine quando veneris iudicare terram, ubi me abscondam a vultu ire tue, * Quia peccavi nimis in vita mea.}{Commissa mea pavesco et ante te erubesco, dum veneris iudicare noli me condemnare.}
\end{responsory}%
\newline \ul{Post resumptionem responsorii non dicatur,} Requiem, \ul{sed statim incipiatur.}
\newline {\color{Red} In Secundo Nocturno. Ant.} In loco paschue ibi me collocavit. {\color{Red} Ps.} \psalm{22} {\color{Red} Ant.} Delicta iuventutis mee et ignorantias meas ne memineris domine. {\color{Red} Ps.} \psalm{24} {\color{Red} Ant.} Credo videre bona domini in terra viventium. {\color{Red} Ps.} \psalm{26} \V In memoria eterna erunt iusti. \R Ab auditione mala non timebunt.
\newline {\color{Red} Lectio Quarta.}
\lettrine[lines=2]{\zallmancaps \color{Blue} R}{esponde} michi. Quantas habeo iniquitates et peccata: scelera mea et delicta ostende michi. Cur faciem tuam abscondis, et arbitraris me inimicum tuum? Contra folium quod vento rapitur ostendis potentiam tuam, et stipulam siccam persequeris. Scribis enim contra me amaritudines, et consumere me vis peccatis adolescentie mee. Posuisti in nervo pedem meum, et observasti omnes semitas meas: et vestigia pedum meorum considerasti. Qui quasi putredo consumendus sum: et quasi vestimentum quod comeditur a tinea.
\begin{responsory}[heu-michi-domine]
{Heu michi domine quia peccavi nimis in vita mea, quid faciam miser ubi fugiam nisi ad te deus meus, miserere mei dum veneris * In novissimo die.}{Anima mea turbata est valde sed tu domine succurre ei.}
\end{responsory}%
\newline {\color{Red} Lectio Quinta.}
\lettrine[lines=2]{\zallmancaps \color{Red} H}{omo} natus de muliere brevi vivens tempore, repletur multis miseriis. Qui quasi flos egreditur et conteritur, et fugit velut umbra, et nunquam in eodem statu permanet. Et dignum ducis super huiuscemodi aperire oculos tuos, et adducere tecum in iudicium? Quis potest facere mundum de immundo conceptum semine? Nonne tu qui solus es? Breves dies hominis sunt: numerus mensium eius apud te est. Constituisti terminos eius, qui preteriri non poterunt. Recede paululum ab eo ut quiescat, donec optata veniat sicut mercenarii dies eius.
\begin{responsory}[ne-recorderis]
{Ne recorderis peccata mea domine, * Dum veneris iudicare seculum per ignem.}{Dirige domine deus meus in conspectu tuo viam meam.}
\end{responsory}%
\newline {\color{Red} Lectio Sexta.}
\lettrine[lines=2]{\zallmancaps \color{Blue} Q}{uis} michi hoc tribuat, ut in inferno protegas me, et abscondas me donec pertranseat furor tuus, et constituas michi tempus in quo recorderis mei? Putasne mortuus homo rursum vivat? Cunctis diebus quibus nunc milito, expecto donec veniat immutatio mea. Vocabis me, et ego respondebo tibi. Operi manuum tuarum porriges dexteram. Tu quidem gressus meos dinumerasti, sed parce peccatis meis.
\begin{responsory}[domine-secundum-actum]
{Domine secundum actum meum noli me iudicare, nichil dignum in conspectu tuo egi, * Ideo deprecor maiestatem tuam, ut tu deus deleas iniquitatem meam.}{Amplius lava me domine ab iniusticia mea et a delicto meo munda me, quia tibi soli peccavi.}
\end{responsory}%
\newline {\color{Red} In Tertio Nocturno. Ant.} Complaceat tibi domine ut eripias me ad adiuvandum me respice. {\color{Red} Ps.} \psalm{39} {\color{Red} Ant.} Sana domine animam meam quia peccavi tibi. {\color{Red} Ps.} \psalm{40} {\color{Red} Ant.} Sitivit anima mea ad deum vivum quando veniam et apparebo ante faciem domini. {\color{Red} Ps.} \psalm{41} \V Ne tradas bestiis animas confitentes tibi. \R Et animas pauperum tuorum ne obliviscaris in finem.
\newline {\color{Red} Lectio VII.}
\lettrine[lines=2]{\zallmancaps \color{Red} S}{piritus} meus attenuabitur, dies mei breviabuntur, et solum michi superest sepulchrum. Non peccavi: et in amaritudinibus moratur oculus meus. Libera me domine, et pone me iuxta te: et cuiusvis manus pugnet contra me. Dies mei transierunt, cogitationes mee dissipate sunt: torquentes cor meum. Noctem verterunt in diem: et rursum post tenebras spero lucem. Si sustinuero, infernus domus mea est, et in tenebris stravi lectulum meum. Putredini dixi, pater meus es: mater mea et soror mea, vermibus. Ubi est ergo nunc prestolatio mea, et patientia mea? Tu es domine Deus meus.
\begin{responsory}[peccante-me]
{Peccante me cotidie et non me penitente timor mortis conturbat me, * Quia in inferno nulla est redemptio, miserere mei deus, et salva me.}{Deus in nomine tuo salvum me fac et in virtute tua libera me.}
\end{responsory}%
\newline {\color{Red} Lectio VIII.}
\lettrine[lines=2]{\zallmancaps \color{Blue} P}{elli} mee consumptis carnibus adhesit os meum et derelicta sunt tantummodo labia circa dentes meos. Miseremini mei miseremini mei, saltem vos amici mei: quia manus domini tetigit me. Quare persequimini me sicut Deus, et carnibus meis saturamini? Quis michi tribuat ut scribantur sermones mei? Quis michi det ut exarentur in libro, stilo ferreo et plumbi lamina, vel certe sculpantur in silice? Scio enim quod redemptor meus vivit, et in novissimo die de terra surrecturus sum, et rursum circumdabor pelle mea, et in carne mea videbo Deum. Quem visurus sum ego ipse: et oculi mei conspecturi sunt, et non alius. Reposita est hec spes mea in sinu meo.
\begin{responsory}[memento-mei-deus]
{Memento mei deus quia ventus est vita mea, * Nec aspiciet me visus hominis.}{Et non revertetur oculus meus ut videat bona.}
\end{responsory}%
\newline {\color{Red} Lectio IX.}
\lettrine[lines=2]{\zallmancaps \color{Red} Q}{uare} de vulva eduxisti me? Qui utinam consumptus essem ne oculus me videret. Fuissem quasi non essem: de utero translatus ad tumulum. Nunquid non paucitas dierum meorum finietur brevi? Dimitte ergo me, ut plangam paululum dolorem meum, antequam vadam et non revertar ad terram tenebrosam et opertam mortis caligine. Terram miserie et tenebrarum, ubi umbra mortis, et nullus ordo: sed sempiternus horor inhabitans.
\newline \hypertarget{libera-me-domine}{\label{libera-me-domine}} \R Libera me domine de morte eterna in die illa tremenda, * Quando celi movendi sunt et terra, * Dum veneris iudicare seculum per ignem. \V Dies illa dies ire calamitatis et miserie, dies magna et amara valde. Quando. \ul{usque} Dum. \V Tremens factus sum ego et timeo dum discussio venerit atque ventura ira. Dum. \V Quid ego miserrimus, quid dicam vel quid faciam dum nil boni perferam ante tantum iudicem. Quando. \V Nunc Christe te petimus miserere quesumus, qui venisti redimere perditos: noli damnare redemptos. Dum. \V Creator omnium rerum deus, qui me de limo terre formasti et mirabiliter proprio sanguine redemisti, corpusque meum licet modo putrescat de sepulchro facies in die iudicii resuscitari exaudi exaudi me, ut animam meam in sinu Abrahe patriarche tui iubeas collocari. Libera me. \ul{Repetito responsorio inchoetur.}
\newline {\color{Red} In Laudibus. Ant.} Exultabunt domino ossa humiliata. {\color{Red} Ps.} \psalm{50} {\color{Red} Ant.} Exaudi domine orationem meam ad te omnis caro veniet. {\color{Red} Ps.} \psalm{64} {\color{Red} Ant.} Me suscepit dextera tua domine. {\color{Red} Ps.} \psalm{62} {\color{Red} Ant.} A porta inferi erue domine animas eorum. {\color{Red} Can.} \psalm{ezechias} {\color{Red} Ant.} Omnis spiritus laudet dominum. {\color{Red} Ps.} \psalm{148} {\color{Red} Ant.} Ego sum resurrectio et vita, qui credit in me eciam si mortuus fuerit vivet, et omnis qui vivit et credit in me non morietur in eternum. {\color{Red} Can.} \psalm{benedictus}
\newline \ul{Finita post canticum antiphona dicatur,} Pater noster, \ul{et psalmus} \psalm{129} \ul{secundum quod supra scriptum est de psalmo,} \psalm{145} \ul{Deinde \Vbar , et oratio ut supra in Vesperis. Ad Primam dicto,} Pater noster, \ul{et} Credo, \ul{Incipiatur antiphona} Requiem, \ul{psalmus} \psalm{53} \ul{psalmus} \psalm{118.1} \ul{psalmus} \psalm{118.2} \ul{Post singulos psalmos non dicatur} Requiem: \ul{sed in fine omnium dicatur,} {\color{Red} Ant.} Requiem eternam dona eis domine et lux perpetua luceat eis. \ul{Deinde dicatur,} Confiteor. \ul{Quo %typo: Confitebor
%pp133
dicto sequatur, \Vbar .} A porta inferi. Dominus vobiscum. Oremus. Fidelium deus. Requiescant in pace. Amen. \ul{Hic ordo servetur ad alias horas diei nisi quod,} Confiteor \ul{non dicatur.} Preciosa \ul{more solito dicatur. Dicta Nona terminatur Officium Mortuorum.}
{\ubsubsection{Sanctorum Quatuor Coronatorum Martirum}{sanctorum-quatuor-coronatorum-martirum-breviarium}}
\newline {\color{Red} Oratio.}
\lettrine[lines=2]{\zallmancaps \color{Red} P}{resta} quesumus omnipotens deus, ut qui gloriosos martyres fortes in sua confessione cognovimus, pios apud te in nostra intercessione sentiamus. Per.
{\ubsubsection{Sancti Theodori Martiris}{sancti-theodori-martiris-breviarium}}
\newline {\color{Red} Oratio.}
\lettrine[lines=2]{\zallmancaps \color{Blue} D}{eus} qui nos beati Theodori martyris tui confessione gloriosa circundas et protegis, presta nobis eius imitatione proficere et oratione fulciri. Per.
{\ubsubsection{Sancti Martini Episcopi et Confessoris}{sancti-martini-episcopi-et-confessoris-breviarium}}
\newline {\color{Red} Ad Vesperas. Super psalmos. Ant.} O beatum pontificem qui totis visceribus diligebat Christum regem et non formidabat imperii principatum. O Martine dulcedo medicamentum et medice, o sanctissima anima quam si gladius persecutoris non abstulit tamen palmam martyrii non amisit. \R \hyperlink{beatus-martinus}{Beatus Martinus.} {\color{Red} III. Ad Mag. Ant.} O Martine o pie quod pium est gaudere de te, o Martine, prophetis compar, apostolis consertus, presulum gemma, fide et meritis egregie, pietate misericordia caritate ineffabilis, succurre nobis nunc et ante deum. {\color{Red} Oratio.}
\lettrine[lines=2]{\zallmancaps \color{Red} D}{eus} qui conspicis quia ex nulla nostra virtute subsistimus, concede propicius ut intercessione beati Martini confessoris tui atque pontificis, contra omnia adversa muniamur. Per.
\newline \ul{Eodem die sancti Menne martyris.} {\color{Red} Oratio.}
\lettrine[lines=2]{\zallmancaps \color{Blue} P}{resta} quesumus omnipotens deus, ut qui beati Menne martyris tui natalicia colimus, intercessione eius in tui nominis amore roboremur. Per.
\newline {\color{Red} Ad Mat. Invitat.}
\begin{invitatory}[adoremus-christum-invitatorium]
{Adoremus Christum regem, * Qui eterna sanctum Martynum coronavit gloria.}
\end{invitatory}
\newline {\color{Red} In Primo Nocturno. Ant.} Sanctus Martinus obitum suum longe ante prescivit, qui et discipulos in unum congregavit, atque sui dissolutionem imminere predixit. {\color{Red} Ps.} \psalm{1} {\color{Red} Ant.} Cum repente viribus corporis cepit destitui tunc meror et luctus omnium vox plangentium, unaque precantium, ne pastor oves deseras. {\color{Red} Ps.} \psalm{2} {\color{Red} Ant.} Scimus quidem te pater desiderare Christum, sed salva sunt tibi premia tua, nostri potius miserere quos deseris pater. {\color{Red} Ps.} \psalm{3}
\begin{responsory}[hic-est-martinus]
{Hic est Martinus electus dei pontifex, cui dominus post apostolos tantam gratiam conferre dignatus est, * Ut in virtute trinitatis deifice mereretur fieri trium mortuorum suscitator magnificus.}{Sancte trinitatis fidem Martinus confessus est.}
\end{responsory}%
\begin{responsory}[dum-sacramenta]
{Dum sacramenta offerret beatus Martinus, * Globus igneus apparuit super caput eius.}{Dum enim ad deum preces emitteret, et pro universo grege sibi commisso deprecaretur.}
\end{responsory}%
\begin{responsory-final}[beatus-martinus]
{Beatus Martinus obitum suum longe ante prescivit, dixitque fratribus, * Dissolutionem sui corporis imminere, * Quia indicavit se iam resolvi.}{Viribus corporis cepit repente destitui, convocatisque discipulis in unum dixit.}{Quia}
\end{responsory-final}%
\newline {\color{Red} In Secundo Nocturno. Ant.} Domine iam satis est quod huc usque certavi, sed quoad ipse iusseris militabo, at si parcis etati bonum est michi, hos quibus timeo ipse custodias. {\color{Red} Ps.} \psalm{4} {\color{Red} Ant.} Artus febre fatiscentes spiritui servire cogebat, stratuque suo illo nobili cilicio recubans, nec terram videre iam dignatus celo toto inhyabat. {\color{Red} Ps.} \psalm{5} {\color{Red} Ant.} Sinite me inquit celum videre, ut spiritus dirigatur ad dominum, nichil in me reperiet inimicus sed sinus Abrahe me suscipiet. {\color{Red} Ps.} \psalm{8}
\begin{responsory}[dixerunt-discipuli]
{Dixerunt discipuli ad beatum Martinum cur nos pater deseris aut cui nos desolatos relinquis, * Invadent enim gregem tuum lupi rapaces.}{Scimus quidem desiderare te Christum sed salva sunt tibi tua premia nostri potius miserere quos deseris.}
\end{responsory}%
\begin{responsory}[domine-si-adhuc]
{Domine si adhuc populo tuo sum necessarius non recuso subire propter eos laborem, * Fiat voluntas tua.}{Satis est domine quod hucusque certavi, nec deficientem causabor etatem, munia tua devotus implebo.}
\end{responsory}%
\begin{responsory-doxology}[oculis-ac-manibus]
{Oculis ac manibus in celum semper intentus, * Invictum ab oratione spiritum non relaxabat.}{Nobili illo stratu suo in cinere et cilicio recubans.}
\end{responsory-doxology}%
\newline {\color{Red} In Tertio Nocturno. Ant.} Media nocte dominica sanctus discessit, cui celi cives mox obviant, multi audierunt voces in sublime, qui et vitro purior, lacte candidior, carne quoque monstratus est gemma sacerdotum. {\color{Red} Ps.} \psalm{14} {\color{Red} Ant.} Glorificati hominis viderunt gloriam qui affuerunt, nam caro quam cinis semper obtexerat ita resplenduit, ut quoddam resurrectionis decus et extincta proferret. {\color{Red} Ps.} \psalm{20} {\color{Red} Ant.} Adest multitudo monachorum ac virginum, hi speciali gloria precipue flebant cum sentirent magis esse gaudendum si rationem vis doloris admitteret. {\color{Red} Ps.} \psalm{23}
\begin{responsory}[o-beatum-virum]
{O beatum virum in cuius transitu sanctorum canit numerus, angelorum exultant chori, * Omniumque celestium virtutum occurrit psallentium exercitus.}{Ecclesia illius virtute roboratur, sacerdotes dei revelatione glorificantur, quem Michael assumpsit cum angelis.}
\end{responsory}%
\begin{responsory}[o-quantus-erat]
{O quantus erat luctus omnium quanta precipue merentium lamenta monachorum, * Quia pium est gaudere Martino et pium est flere Martinum.}{Beati viri corpus, usque ad locum sepulchri ymnis canora celestibus turba prosequitur.}
\end{responsory}%
\begin{responsory-final}[martinus-abrahe]
{Martinus Abrahe sinu letus excipitur, Martinus hic pauper et modicus, * Celum dives ingreditur, * Ymnis celestibus honoratur.}{Martinus episcopus migravit a seculo vivit in Christo gemma sacerdotum.}{Ymnis}
\end{responsory-final}%
\newline {\color{Red} In Laudibus. Ant.} Dixerunt discipuli ad beatum Martinum cur nos pater deseris, aut cui nos desolatos relinquis invadent, enim gregem tuum lupi rapaces. {\color{Red} Ant.} Domine si adhuc populo tuo sum necessarius non recuso laborem, fiat voluntas tua. {\color{Red} Ant.} O virum ineffabilem nec labore victum nec morte vincendum, qui nec mori timuit nec vivere recusavit. {\color{Red} Ant.} Oculis ac manibus in celum semper intentus, invictum ab oratione spiritum non relaxabat alleluya alleluya. {\color{Red} Ant.} Martinus Abrahe sinu letus excipitur, Martinus hic pauper et modicus celum dives ingreditur ymnis celestibus honoratur.
\newline {\color{Red} Ad Ben. Ant.} O quantus luctus omnium, quanta precipue lamenta monachorum et virginum chori, quia pium est gaudere Martino et pium est flere Martinum.
\newline {\color{Red} Ad Mag. Ant.} O beatum virum cuius anima paradisum possidet, unde exultant angeli, letantur archangeli, chorus sanctorum proclamat turba virginum invitat, mane nobiscum in eternum.
\newline \ul{Per octavas beati Martini nisi in Dominica Invitat.} \hyperlink{regem-confessorum-invitatorium}{Regem confessorum.} \ul{Ad Ben. vero et ad Mag. dicantur antiphone cotidie sicut in festo.}
{\ubsubsection{Sancti Brictii Episcopi et Confessoris}{sancti-brictii-episcopi-et-confessoris-breviarium}}
\newline {\color{Red} Oratio.}
\lettrine[lines=2]{\zallmancaps \color{Red} D}{a} quesumus omnipotens deus ut beati Brictii confessoris tui atque pontificis veneranda sollemnitas et devotionem nobis augeat et salutem. Per.
{\ubsubsection{Sancte Elysabeth}{sancte-elysabeth-breviarium}}
\newline {\color{Red} Memoria in Vesperis. Ant.} Veni electa mea et ponam in te thronum meum quia concupivit rex speciem tuam. \V Diffusa. {\color{Red} Oratio.}
\lettrine[lines=2]{\zallmancaps \color{Blue} T}{uorum} corda fidelium deus miserator illustra, et beate Elysabeth precibus gloriosis, fac nos prospera mundi despicere, et celesti semper consolatione gaudere. Per.
\newline {\color{Red} In Laudibus. Ant.} Accinxit fortitudine lumbos suos, et roboravit brachium suum, ideoque lucerna eius non extinguetur in sempiternum. \V Specie tua. {\color{Red} Oratio.} Tuorum.
{\ubsubsection{Sancte Cecilie Virginis et Martiris}{sancte-cecilie-virginis-et-martiris-breviarium}}
\newline {\color{Red} Ad Mag. Ant.} Virgo gloriosa semper evangelium Christi gerebat in pectore suo, non diebus neque noctibus a colloquiis divinis et oratione cessabat. {\color{Red} Oratio.}
\lettrine[lines=2]{\zallmancaps \color{Red} D}{eus} qui nos annua beate Cecilie virginis et martyris tue sollemnitate letificas, da ut quam veneramur officio, eciam pie conversationis sequamur exemplo. Per.
\newline {\color{Red} Ad Mat. In Primo Nocturno. Ant.} Cecilia virgo Almachium exuperabat, Tyburtium et Valerianum ad coronas vocabat. {\color{Red} Ps.} \psalm{8} {\color{Red} Ant.} Expansis manibus orabat ad dominum, ut eam eriperet de inimicis. {\color{Red} Ps.} \psalm{18} {\color{Red} Ant.} Cilicio Cecilia membra domabat, deum gemitibus exorabat. {\color{Red} Ps.} \psalm{23}
\begin{responsory}[cantantibus-organis]
{Cantantibus organis Cecilia virgo in corde suo soli domino decantabat dicens, * Fiat domine cor meum et corpus meum immaculatum ut non confundar.}{Biduanis ac triduanis ieiuniis orans, suam domino pudicitiam commendabat dicens.}
\end{responsory}%
\begin{responsory}[o-beata-cecilia]
{O beata Cecilia que duos fratres convertisti, Almachium iudicem superasti, * Urbanum episcopum in vultu angelico demonstrasti.}{Beata es virgo et gloriosa, et benedictus sermo oris tui.}
\end{responsory}%
\begin{responsory-final}[virgo-gloriosa]
{Virgo gloriosa semper evangelium Christi gerebat in pectore, * Et non diebus neque noctibus vacabat, * A colloquiis divinis et oratione.}{Cilicio Cecilia membra domabat deum gemitibus exorabat.}{A colloquiis}
\end{responsory-final}%
\newline {\color{Red} In Secundo Nocturno. Ant.} Domine Ihesu Christe seminator casti consilii, suscipe seminum fructus quos in Cecilia seminasti. {\color{Red} Ps.} \psalm{44} {\color{Red} Ant.} Beata Cecilia dixit ad Tyburcium, hodie te fateor meum esse cognatum, quia amor dei te fecit esse contemptorem ydolorum. {\color{Red} Ps.} \psalm{45} {\color{Red} Ant.} Fiat domine cor meum et corpus meum immaculatum ut non confundar. {\color{Red} Ps.} \psalm{86}
\begin{responsory}[cilicio-cecilia]
{Cilicio Cecilia membra domabat, deum gemitibus exorabat, * Almachium exuperabat, Tyburcium et Valerianum ad coronas vocabat.}{Non diebus neque noctibus vacabat a colloquiis divinis et oratione.}
\end{responsory}%
\begin{responsory}[cecilia-me-misit]
{Cecilia me misit ad vos, ut ostendatis michi sanctum Urbanum, * Quia ad ipsum habeo secreta que perferam.}{Cumque pro signo acceperitis benedictionem eius, ducite me ad sanctum antistitem.}
\end{responsory}%
\begin{responsory-final}[ceciliam-intra]
{Ceciliam intra cubiculum orantem invenit, et iuxta eam stantem angelum domini, * Quem videns Valerianus * Nimio terrore correptus est.}{Angelus domini descendit de celo et lumen refulsit in habitaculo.}{Nimio}
\end{responsory-final}%
\newline {\color{Red} In Tertio Nocturno. Ant.} Credimus Christum filium dei verum deum esse, qui sibi talem elegit famulam. {\color{Red} Ps.} \psalm{95} {\color{Red} Ant.} Nos scientes sanctum nomen omnino negare non possumus. {\color{Red} Ps.} \psalm{96} {\color{Red} Ant.} Tunc Valerianus perrexit ad antistitem, et signo quod acceperat invenit sanctum Urbanum. {\color{Red} Ps.} \psalm{97}
\begin{responsory}[domine-ihesu]
{Domine Ihesu Christe pastor bone seminator casti consilii, suscipe seminum fructus quos in Cecilia seminasti, * Cecilia famula tua quasi apis tibi argumentosa deservit.}{Nam sponsum quem quasi leonem ferocem accepit, ad te quasi agnum mansuetissimum destinavit.}
\end{responsory}%
\begin{responsory}[beata-cecilia-dixit]
{Beata Cecilia dixit Tyburcio hodie te fateor meum esse cognatum, * Quia amor dei te fecit esse contemptorem ydolorum.}{Sicut enim amor dei michi fratrem tuum coniugem fecit: ita te michi cognatum faciet.}
\end{responsory}%
\begin{responsory-final}[dum-aurora-finem]
{Dum aurora finem daret, Cecilia dixit, * Eya milites Christi abicite opera tenebrarum, * Et induimini armis lucis.}{Cecilia vale dicens fratribus et exhortans eos ait.}{Et induimini}
\end{responsory-final}%
\newline {\color{Red} In Laudibus. Ant.} Cantantibus organis Cecilia domino decantabat dicens, fiat cor meum immaculatum ut non confundar. {\color{Red} Ant.} Est secretum Valeriane quod tibi volo dicere, angelum dei habeo amatorem, qui nimio zelo custodit corpus meum. {\color{Red} Ant.} Valerianus in cubiculo Ceciliam cum angelo orantem invenit. {\color{Red} Ant.} Benedico te pater domini mei Ihesu Christi, quia per filium tuum ignis extinctus est a latere meo. {\color{Red} Ant.} Cecilia famula tua domine quasi apis tibi argumentosa deservit.
\newline {\color{Red} Ad Ben. Ant.} Dum aurora finem daret Cecilia exclamavit dicens, eya milites Christi abicite opera tenebrarum et induimini armis lucis.
{\ubsubsection{Sancti Clementis Pape et Martiris}{sancti-clementis-pape-et-martiris-breviarium}}
\newline {\color{Red} Ad Mag. Ant.} Dedisti domine habitaculum martyri tuo Clementi in mari in modum templi marmorei, angelicis manibus preparatum iter prebens populo terre ut enarrent mirabilia tua. {\color{Red} Oratio.}
\lettrine[lines=2]{\zallmancaps \color{Blue} D}{eus} qui nos annua beati Clementis martyris tui atque pontificis sollemnitate letificas, concede propicius ut cuius natalicia colimus, virtutem quoque passionis imitemur. Per.
\newline \ul{Memoria de sancta Cecilia.} {\color{Red} Ant.} Triduanas a domino poposci indutias ut domum meam domino ecclesiam consecrarem.
\newline \ul{Ad Matutinas. \Rbar .} \hyperlink{iste-sanctus-pro}{Iste sanctus.} \ul{\Rbar .} \hyperlink{posuisti-domine-super}{Posuisti.}
\newline \R {\color{Red} III.} \hypertarget{oremus-omnes}{\label{oremus-omnes}} Oremus omnes ad dominum Ihesum Christum dixit, beatus Clemens, ut confessoribus suis fontis venas aperiat, * Ut de beneficiis eius gratulemur. \R Qui percussit in deserto Syna petram, et fluxerunt aque in abundantiam ipse nobis laticem effluentem imperciatur. Ut de. Gloria. Ut de.
\newline \ul{\Rbar .} \hyperlink{beatus-vir-qui-suffert}{Beatus vir.} \ul{\Rbar .} \hyperlink{domine-prevenisti-eum}{Domine prevenisti.}
\newline \R {\color{Red} VI.} \hypertarget{orante-sancto-clemente}{\label{orante-sancto-clemente}} Orante sancto Clemente apparuit ei agnus dei, * De sub cuius pede fons vivus emanat, * Fluminis impetus letificat civitatem dei. \V Cumque orationem complesset et hinc inde circumspiceret: vidit supra montem agnum stantem. De sub. Gloria. Fluminis.
\newline \ul{\Rbar .} \hyperlink{stola-iocunditatis}{Stola iocunditatis.} \ul{\Rbar .} \hyperlink{corona-aurea-super}{Corona aurea.}
\newline \R {\color{Red} IX.} \hypertarget{dedisti-domine}{\label{dedisti-domine}} Dedisti domine habitaculum martyri tuo Clementi in mari in modum templi marmorei angelicis manibus preparatum, iter prebens populo terre, * Ut enarrent mirabilia tua. \V Hoc domine ad laudem et gloriam nominis tui fieri voluisti. Ut enarrent. Gloria. Ut enarrent.
\newline {\color{Red} In Laudibus. Ant.} Orante sancto Clemente apparuit ei agnus dei. {\color{Red} Ant.} Non meis meritis ad vos me misit dominus: vestris coronis participem me fieri. {\color{Red} Ant.} Vidi supra montem agnum stantem, de sub cuius pede fons vivus emanat. {\color{Red} Ant.} De sub cuius pede fons vivus emanat, fluminis impetus letificat civitatem dei. {\color{Red} Ant.} Omnes gentes per girum crediderunt Christo domino.%typo apperuit, perticipem
\newline {\color{Red} Ad Ben. Ant.} Oremus omnes ad dominum Ihesum Christum ut confessoribus suis fontis venas aperiat.
\newline {\color{Red} Ad Mag. Ant.} Invenerunt in modum templi marmorei habitaculum a deo paratum.
{\ubsubsection{Sancti Grisogoni Martyris}{sancti-grisogoni-martyris-breviarium}}
\newline {\color{Red} Oratio.}
\lettrine[lines=2]{\zallmancaps \color{Red} A}{desto} domine supplicationibus nostris, ut qui ex iniquitate nostra reos nos esse cognoscimus, beati Grisogoni martyris tui intercessione liberemur. Per.
{\ubsubsection{Sancte Katherine Virginis et Martyris}{sancte-katherine-virginis-et-martyris-breviarium}}
\newline {\color{Red} Ad Vesperas. Super psalmos. Ant.} Virginis eximie Katherine martyris alme, festa celebrare da nobis rex pie Christe. \R \hyperlink{virgo-flagellatur}{Virgo flagellatur.} {\color{Red} VI. Ymnus.}
\hymn{katherine-collaudemus}{Katherine collaudemus}{
\lettrine[lines=2]{\zallmancaps \color{Blue} K}{atherine} collaudemus
virtutum insignia,
cordis ei presentemus
et oris obsequia,
ut ab ipsa reportemus
equa laudis premia.
\par {\bfseries \color{Red} F}ulta fide Katherina
iudicem Maxentium
non formidat, lex divina
sic firmat eloquium:
quod confutat ex doctrina
doctores gentilium.
\par {\bfseries \color{Blue} V}icti Christum confitentur
relictis erroribus,
iudex iubet ut crementur,
nec pilis aut vestibus
ignis nocet, sed torrentur
inustis corporibus.
\par {\bfseries \color{Red} G}loria et honor deo
usquequo altissimo
una patri filioque
inclito Paraclito,
cui laus est et potestas
per eterna secula. Amen.
}
\newline {\color{Red} Ad Mag. Ant.} Inclita sancte virginis Katherine sollemnia, suscipiat alacriter pia mater ecclesia, ave virgo deo digna, ave dulcis et benigna, obtine nobis gaudia que possides cum gloria. {\color{Red} Oratio.}
\lettrine[lines=2]{\zallmancaps \color{Red} D}{eus} qui dedisti legem Moysi in summitate montis Synay, et in eodem loco corpus beate Katherine virginis et martyris tue per sanctos angelos tuos mirabiliter collocasti: concede propitius, ut eius meritis et intercessione, ad montem qui Christus est valeamus pervenire. Per eundem.
\newline {\color{Red} Ad Mat. Invit.}
\begin{invitatory}[adoretur-virginum-invitatorium]
{Adoretur virginum rex in seculorum secula, * Virgini qui Katherine contulit celestia.}
\end{invitatory}
\newline {\color{Red} Ymnus.}
\hymn{pange-lingua-gloriose}{Pange lingua gloriose}{
\lettrine[lines=2]{\zallmancaps \color{Blue} P}{ange} lingua gloriose
virginis martyrium,
gemme iubar preciose
descendat in medium,
ut illustret tenebrose
mentis domicilium.
\par {\bfseries \color{Red} B}landimentis rex molitur
virginem seducere,
nec promissis emollitur,
nec terretur verbere,
compeditur, custoditur,
tetro clausa carcere.
\par {\bfseries \color{Blue} C}lause lumen ne claudatur
illucet Porphyrio,
qui regine federatur
fidei collegio,
quorum fidem imitatur
ducentena contio.
\par {\bfseries \color{Red} G}loria et honor deo
usquequo altissimo
una patri filioque
inclito Paraclito,
cui laus est et potestas
per eterna secula. Amen.
}
\newline \ul{\Rbar .} \hyperlink{veni-sponsa}{Veni sponsa.} \ul{\Rbar .} \hyperlink{propter-veritatem}{Propter veritatem.}
\newline \R {\color{Red} III.} \hypertarget{nobilis-et-pulchra}{\label{nobilis-et-pulchra}} Nobilis et pulchra prudens Katherina puella, * Flagrat amore dei, * Spernit vaga gaudia mundi. \V Cui rex carne pater fuerat, reginaque mater. Flagrat. Gloria. Spernit.
\newline \ul{\Rbar .} \hyperlink{grata-facta-est}{Grata.} \ul{\Rbar .} \hyperlink{hec-est-virgo}{Hec est virgo.}
\newline \R {\color{Red} VI.} \hypertarget{virgo-flagellatur}{\label{virgo-flagellatur}} Virgo flagellatur, crucianda fame religatur, carcere clausa manet, lux celica fusa refulget, * Fragrat odor, dulces * Cantant celi agmina laudes. \V Sponsus amat sponsam salvator visitat illam. Fragrat. Gloria. Cantant.
\newline \ul{\Rbar .} \hyperlink{audivi-vocem-de}{Audivi.} \ul{\Rbar .} \hyperlink{veni-electa-mea}{Veni electa.}
\newline \R {\color{Red} IX.} \hypertarget{percussa-gladio}{\label{percussa-gladio}} Percussa gladio dat lac pro sanguine collo, * Quam manus angelica * Sepelivit vertice Syna. \V Membris virgineis olei fluit unda salubris. Quam. Gloria. Sepelivit.%typo: virgines
\newline {\color{Red} In Laudibus. Ant.} Passionem gloriose virginis Katherine devote plebs celebret fidelis, que sui memores deo commendat meritis: et iuvat beneficiis. {\color{Red} Ant.} Post plurima supplicia martyr alma, ad decollandum est ducta, ad celum tendens oculos, collum summittit gladio, orans dat gloriam deo. {\color{Red} Ant.} Expecto pro te gladium Ihesu rex bone, tu meum colloca in paradiso spiritum, et fac misericordiam: meam agentibus memoriam. {\color{Red} Ant.} Vox de celis insonuit, veni electa mea, veni intra thalamum sponsi tui, que postulas impetrasti, qui te laudant salvi fient. {\color{Red} Ant.} Quia devotis laudibus tui memoriam virgo recolimus, o beata Katherina ora pro nobis quesumus. {\color{Red} Ymnus.}
\hymn{presens-dies-expendatur}{Presens dies expendatur}{
\lettrine[lines=2]{\zallmancaps \color{Red} P}{resens} dies expendatur
in eius preconium,
cuius virtus dilatatur
in ore laudantium,
si gestorum teneatur
finis et initium.
\par {\bfseries \color{Blue} I}mminente passione
virgo hec interserit
assequatur Ihesu bone
quod a te pecierit,
suo quisquis in agone
memor mei fuerit.
\par {\bfseries \color{Red} I}n hoc caput amputatur
fluit lac cum sanguine:
angelorum sublevatur
corpus multitudine,
et Synay collocatur
in supremo culmine.
\par {\bfseries \color{Blue} G}loria et honor deo
usquequo altissimo
una patri filioque
inclito Paraclito,
cui laus est et potestas
per eterna secula. Amen.
}
\newline {\color{Red} Ad Ben. Ant.} Prudens et vigilans virgo qualis es cum sponso illo qui te elegit de mundo, quam pulchra quam mirabilis, quanta luce spectabilis inter Syon iuvenculas et Iherusalem filias, thalamo gaudes regio coniuncta dei filio.
\newline {\color{Red} Ad Vesperas. Ymnus.} \hyperlink{katherine-collaudemus}{Katherine.} {\color{Red} Ad Mag. Ant.} Ave virginum gemma Katherina, ave sponsa regis regum gloriosa, ave viva Christi hostia, tua venerantibus patrocinia: impetrata non deneges suffragia.
{\ubsubsection{Sanctorum Vitalis et Agricole Martyrum}{sanctorum-vitalis-et-agricole-martyrum-breviarium}}
\newline {\color{Red} Oratio.}
\lettrine[lines=2]{\zallmancaps \color{Blue} D}{a} quesumus omnipotens deus, ut qui beatorum martyrum tuorum Vitalis et Agricole sollemnia colimus, eorum apud te intercessionibus adiuvemur. Per.
{\ubsubsection{Sancti Saturnini Martiris}{sancti-saturnini-martiris-breviarium}}
\newline {\color{Red} Oratio.}
\lettrine[lines=2]{\zallmancaps \color{Red} D}{eus} qui nos beati Saturnini martyris tui concedis natalitio perfrui, eius nos tribue meritis adiuvari. Per.

\ubsection{Commune Sanctorum}{commune-sanctorum}

{\ubsubsection{Incipit Officium, In Communi Sanctorum}{in-communi-sanctorum-breviarium}}
{\ubsubsection{In Communi Unius Vel Plurimorum Apostolorum Extra Tempus Paschale}{in-communi-unius-vel-plurimorum-apostolorum-extra-tempus-paschale-breviarium}}
\newline {\color{Red} Ad Vesperas. Super psalmos. Ant.} Estote fortes in bello et pugnate cum antiquo serpente, et accipietis regnum eternum alleluya. {\color{Red} Infra Septuagesimam.} Dicit dominus. {\color{Red} Capitulum.}
\lettrine[lines=2]{\zallmancaps \color{Blue} I}{am} non estis hospites et advene, sed estis cives sanctorum et domestici dei, superedificati super fundamentum apostolorum et prophetarum, ipso summo angulari lapide Christo Ihesu. \R \hyperlink{vos-estis-lux}{Vos estis.} {\color{Red} VI. Ymnus.}
\hymn{exultet-celum-laudibus}{Exultet celum laudibus}{
\lettrine[lines=2]{\zallmancaps \color{Red} E}{xultet} celum laudibus,
resultet terra gaudiis,
apostolorum gloriam
sacra canunt sollemnia.
\par {\bfseries \color{Red} V}os secli iusti iudices,
et vera mundi lumina,
votis precamur cordium
audite preces supplicum.
\par {\bfseries \color{Blue} Q}ui celum verbo clauditis,
serasque eius solvitis,
nos a peccatis omnibus,
solvite iussu quesumus.
\par {\bfseries \color{Red} Q}uorum precepto subditur
salus et languor omnium:
sanate egros moribus
nos reddentes virtutibus.
\par {\bfseries \color{Blue} U}t cum iudex advenerit
Christus in fine seculi,
nos sempiterni gaudii
faciat esse compotes.
\par {\bfseries \color{Red} D}eo patri sit gloria
eiusque soli filio
cum spiritu paraclito
et nunc et in perpetuum. Amen.
}
\newline \V In omnem terram exivit sonus eorum. \R Et in fines orbis terre verba eorum.
\newline {\color{Red} Ad Mag. Ant.} Beati eritis cum vos oderint homines et cum separaverint vos et exprobraverint, et eiecerint nomen vestrum tanquam malum propter filium hominis, gaudete et exultate, ecce enim merces vestra multa est in celo.
\newline {\color{Red} Ad Mat. Invit.}
\begin{invitatory}[gaudete-et-exultate-invitatorium]
{Gaudete et exultate, * Quia nomina vestra scripta sunt in celis.}
\end{invitatory}
\newline {\color{Red} Ymnus.} \hyperlink{exultet-celum-laudibus}{Exultet.}
\newline {\color{Red} In Primo Nocturno. Ant.} In omnem terram exivit sonus eorum et in fines orbis terre verba eorum. {\color{Red} Ps.} \psalm{8} {\color{Red} Ant.} Clamaverunt iusti et dominus exaudivit eos. {\color{Red} Ps.} \psalm{33} {\color{Red} Ant.} Constitues eos principes super omnem terram memores erunt nominis tui domine. {\color{Red} Ps.} \psalm{44} \V In omnem terram. \R
\lettrine[lines=2]{\zallmancaps \color{Blue} E}{cce} \hypertarget{ecce-ego-mitto}{\label{ecce-ego-mitto}} ego mitto vos sicut oves in medio luporum dicit dominus, * Estote ergo prudentes sicut serpentes, et simplices sicut columbe. \V Tradent enim vos in conciliis, et in synagogis suis flagellabunt vos. Estote.
\begin{responsory}[tollite-iugum]
{Tollite iugum meum super vos dicit dominus, et discite quia mitis sum et humilis corde, * Iugum enim meum suave est et onus meum leve.}{Et invenietis requiem animabus vestris.}%typo: humiles
\end{responsory}%
\begin{responsory-doxology}[dum-steteritis]
{Dum steteritis ante reges et presides nolite cogitare quomodo aut quid loquamini, * Dabitur enim vobis in illa hora quid loquamini.}{Non enim vos estis qui loquimini, sed spiritus patris vestri qui loquitur in vobis.}
\end{responsory-doxology}%
\newline {\color{Red} In Secundo Nocturno. Ant.} Principes populorum congregati sunt cum deo Abraham. {\color{Red} Ps.} \psalm{46} {\color{Red} Ant.} Dedisti hereditatem timentibus nomen tuum domine. {\color{Red} Ps.} \psalm{60} {\color{Red} Ant.} Annuntiaverunt opera dei et facta eius intellexerunt. {\color{Red} Ps.} \psalm{63} \V Constitues eos principes super omnem terram. \R Memores erunt nominis tui domine.
\begin{responsory}[vidi-coniunctos]
{Vidi coniunctos viros habentes splendidas vestes, et angelus domini locutus est ad me dicens, * Isti sunt viri sancti facti amici dei.}{Vidi angelum dei volantem per medium celi, voce magna clamantem et dicentem.}
\end{responsory}%
\begin{responsory}[isti-sunt-triumphatores]
{Isti sunt triumphatores et amici dei qui contempnentes iussa principum meruerunt premia eterna, * Modo coronantur et accipiunt palmam.}{Isti sunt qui venerunt ex magna tribulatione, et laverunt stolas suas in sanguine agni.}
\end{responsory}%
\begin{responsory-final}[vos-estis-lux]
{Vos estis lux mundi dicit dominus, * Qui in patientia vestra, * Possidebitis animas vestras.}{Et ego dispono vobis sicut disposuit michi pater meus regnum, ut edatis et bibatis super mensam meam in regno meo.}{Possidebitis}
\end{responsory-final}%
\newline {\color{Red} In Tertio Nocturno. Ant.} Peccatorum confringam cornua et exaltabuntur cornua iusti. {\color{Red} Ps.} \psalm{74} {\color{Red} Ant.} Lux orta est iusto et rectis corde leticia. {\color{Red} Ps.} \psalm{96} {\color{Red} Ant.} Custodiebant testimonia eius et precepta eius que dedit illis. {\color{Red} Ps.} \psalm{98} \V Nimis honorati sunt amici tui deus. \R Nimis confortatus est principatus eorum.
\begin{responsory}[isti-sunt-viri]
{Isti sunt viri sancti quos elegit dominus in caritate non ficta, et dedit illis gloriam sempiternam, * Quorum doctrina fulget ecclesia ut sole luna.}{In omnem terram exivit sonus eorum, et in fines orbis terre
%pp134
verba eorum.}
\end{responsory}% 
\begin{responsory}[qui-sunt-isti]
{Qui sunt isti qui ut nubes volant, * Et quasi columbe ad fenestras suas.}{Candidiores nive, nitidiores lacte, rubicundiores ebore antiquo.}
\end{responsory}%
\begin{responsory-doxology}[gregem-tuum-pastor]
{Gregem tuum pastor eterne non deseras, sed per beatos apostolos tuos * Perpetua defensione custodias.}{Protege domine populum tuum et apostolorum tuorum patrocinio confidentem.}
\end{responsory-doxology}%
\newline \V Dedisti hereditatem. \R Timentibus nomen tuum domine.
\newline {\color{Red} In Laudibus. Ant.} Hoc est preceptum meum, ut diligatis invicem sicut dilexi vos. {\color{Red} Ps.} \psalm{92} {\color{Red} Ant.} Maiorem caritatem nemo habet, ut animam suam ponat quis pro amicis suis. {\color{Red} Ant.} Vos amici mei estis si feceritis que precipio vobis, dicit dominus. {\color{Red} Ant.} Beati pacifici, beati mundo corde, quoniam ipsi deum videbunt. {\color{Red} Ant.} In patientia vestra possidebitis animas vestras. {\color{Red} Cap.} Iam non estis. {\color{Red} Ymnus.}
\hymn{eterna-christi-munera-apostolorum}{Eterna Christi munera}{
\lettrine[lines=2]{\zallmancaps \color{Red} E}{terna} Christi munera,
apostolorum gloriam,
laudes canentes debitas
letis canamus mentibus.
\par {\bfseries \color{Blue} E}cclesiarum principes,
belli triumphales duces,
celestis aule milites
et vera mundi lumina.
\par {\bfseries \color{Red} D}evota sanctorum fides,
invicta spes credentium,
perfecta Christi caritas,
mundi triumphat principem.
\par {\bfseries \color{Blue} I}n his paterna gloria,
in his voluntas spiritus,
exultat in his filius,
celum repletur gaudiis.
\par {\bfseries \color{Red} T}e nunc redemptor quesumus,
ut ipsorum consortio
iungas precantes servulos
in sempiterna secula. Amen.
}
\newline \ul{Iste versus ultimus non dicatur in festo sancti Iohannis evangeliste nec in tempore Paschali.}
\newline \V Annunciaverunt opera dei. \R Et facta eius intellexerunt.
\newline {\color{Red} Ad Ben. Ant.} Vos qui secuti estis me sedebitis super sedes iudicantes duodecim tribus Israhel, dicit dominus.
\newline {\color{Red} Ad Terciam. Cap.} Iam non estis.
\begin{responsory-breve}[in-omnem-terram]
{In omnem terram * Exivit sonus eorum.}{Et in fines orbis terre verba eorum.}
\end{responsory-breve}%
\V Constitues.
\newline {\color{Red} Ad Sextam. Capitulum.}
\lettrine[lines=2]{\zallmancaps \color{Blue} P}{er} manus autem apostolorum fiebant signa et prodigia multa in plebe.
\begin{responsory-breve}[constitues-eos-principes]
{Constitues eos principes * Super omnem terram.}{Memores erunt nominis tui domine.}
\end{responsory-breve}%
\V Nimis honorati.
\newline {\color{Red} Ad Nonam. Cap.}
\lettrine[lines=2]{\zallmancaps \color{Red} I}{bant} apostoli gaudentes a conspectu concilii, quoniam digni habiti sunt pro nomine Ihesu contumeliam pati.
\begin{responsory-breve}[nimis-honorati]
{Nimis honorati sunt * Amici tui deus.}{Nimis confortatus est principatus eorum.}
\end{responsory-breve}%
\V Annunciaverunt.
\newline {\color{Red} Ad Vesperas. Super psalmos. Ant.} Hoc est. {\color{Red} Ps.} \psalm{109} \psalm{112} \psalm{115} \psalm{125} \psalm{138} \ul{Cap., Ymnus, \Vbar , ut in Primis Vesperis.}
\newline {\color{Red} Ad Mag. Ant.} In regeneratione cum sederit filius hominis in sede maiestatis sue sedebitis et vos iudicantes duodecim tribus Israhel dicit dominus.
{\ubsubsection{In Communi Unius Evangeliste Extra Tempus Paschale}{in-communi-unius-evangeliste-extra-tempus-paschale-breviarium}}
\newline {\color{Red} Ad Vesperas. Capitulum.}
\lettrine[lines=2]{\zallmancaps \color{Blue} Q}{ui} timet deum faciet bona, et qui continens est iusticie apprehendet illam, et obviabit illi quasi mater honorificata.
\newline {\color{Red} Ad Mag. Ant.} Ecce ego Iohannes vidi ostium apertum in celo, et ecce sedes posita erat in eo, et in medio sedis et in circuitu eius quatuor animalia plena oculis ante et retro, et dabant gloriam et honorem et benedictionem sedenti super thronum, viventi in secula seculorum.%typo: inventi for viventi
\newline {\color{Red} Ad Mat. Invit.}
\begin{invitatory}[regem-evangelistarum-invitatorium]
{Regem evangelistarum dominum, * Venite adoremus.}
\end{invitatory}
\newline {\color{Red} In Laudibus. Cap.} Qui timet. {\color{Red} Ad Ben. Ant.} In medio et in circuitu sedis dei quatuor animalia senas alas habentia oculis undique plena, non cessant nocte ac die dicere, sanctus sanctus sanctus, dominus deus omnipotens, qui erat et qui est et qui venturus est.
\newline {\color{Red} Ad Terciam. Cap.} Qui timet.
\newline {\color{Red} Ad Sextam. Capitulum.}
\lettrine[lines=2]{\zallmancaps \color{Red} C}{ibavit} illum dominus pane vite et intellectus, et aqua sapientie salutaris potavit illum.
\newline {\color{Red} Ad Nonam. Capitulum.}
\lettrine[lines=2]{\zallmancaps \color{Blue} I}{n} medio ecclesie aperuit os eius, et implevit eum dominus spiritu sapientie et intellectus, et stola glorie induit eum.
\newline {\color{Red} Ad Vesperas. Cap.} Qui timet. {\color{Red} Ad Mag. Ant.} Tua sunt hec Christe opera qui sanctos tuos ita glorificas, ut eciam dignitatis gratia in eis future preire miracula facias, tu insignes evangelii predicatores animalium celestium admirabili figura presignasti, his namque celeste munus collatum gloriosis indiciis es dignatus ostendere, hinc laus, hinc gloria tibi resonet in secula.
{\ubsubsection{In Communi Unius Martiris Extra Tempus Paschale}{in-communi-unius-martiris-extra-tempus-paschale-breviarium}}
\newline {\color{Red} Ad Vesperas. Cap.}
\lettrine[lines=2]{\zallmancaps \color{Red} B}{eatus} vir qui suffert temptationem, quoniam cum probatus fuerit accipiet coronam vite, quam repromisit deus diligentibus se. {\color{Red} Ymnus.}
\hymn{deus-tuorum-militum}{Deus tuorum militum}{
\lettrine[lines=2]{\zallmancaps \color{Blue} D}{eus} tuorum militum
sors et corona premium,
laudes canentes martyris
absolve nexu criminis.
\par {\bfseries \color{Red} H}ic nempe mundi gaudia,
et blandimenta noxia
caduca rite deputans
pervenit ad celestia.
\par {\bfseries \color{Blue} P}enas cucurrit fortiter,
et sustulit viriliter,
pro te effundens sanguinem
eterna dona possidet.
\par {\bfseries \color{Red} O}b hoc precatu supplici
te poscimus piissime,
in hoc triumpho martyris
dimitte noxam servulis. \ul{In Inventione sancti Stephani dicatur sic.} In tanti festo martyris, \ul{Et cetera.}
\par {\bfseries \color{Blue} L}aus et perennis gloria,
deo patri et filio,
sancto simul paraclito,
in sempiterna secula. Amen.
}
\newline \V Gloria et honore coronasti eum domine. \R Et constituisti eum super opera manuum tuarum.%typo: supera, missing oepra
\newline {\color{Red} Ad Mag. Ant.} Hic est vere martyr qui pro Christi nomine sanguinem suum fudit, qui minas iudicum non timuit, nec terrene dignitatis gloriam quesivit, sed ad celestia regna feliciter pervenit.
\newline {\color{Red} Ad Mat. Invit.}
\begin{invitatory}[regem-sempiternum-invitatorium]
{Regem sempiternum venite adoremus, * Qui martyrem suum digne pro meritis coronavit in celis.}
\end{invitatory}
\newline {\color{Red} Ymnus.} \hyperlink{deus-tuorum-militum}{Deus tuorum.}
\newline {\color{Red} In Primo Nocturno. Ant.} In lege domini fuit voluntas eius die ac nocte. {\color{Red} Ps.} \psalm{1} {\color{Red} Ant.} Predicans preceptum domini constitutus est in monte sancto eius. {\color{Red} Ps.} \psalm{2} {\color{Red} Ant.} Voce mea ad dominum clamavi et exaudivit me de monte sancto suo. {\color{Red} Ps.} \psalm{3} \V Gloria et honore. \R
\lettrine[lines=2]{\zallmancaps \color{Red} I}{ste} \hypertarget{iste-sanctus-pro}{\label{iste-sanctus-pro}} sanctus pro lege dei sui certavit usque ad mortem et a verbis impiorum non timuit, * Fundatus enim erat supra firmam petram. \V Munimine regio septus nullatenus est ab adversariis superatus. Fundatus.
\begin{responsory}[posuisti-domine-super]
{Posuisti domine super caput eius coronam de lapide precioso, vitam peciit a te et tribuisti ei, * In seculum seculi.}{Magna est gloria eius in salutari tuo, gloriam et magnum decorem impones super eum.}
\end{responsory}%
\begin{responsory-doxology}[gloria-et-honore-coronasti]
{Gloria et honore coronasti eum domine, et constituisti eum super opera manuum tuarum, * Omnia subiecisti sub pedibus eius.}{Quoniam elevata est magnificentia tua super celos deus.}
\end{responsory-doxology}%
\newline {\color{Red} In Secundo Nocturno. Ant.} Filii hominum scitote quia dominus sanctum suum mirificavit. {\color{Red} Ps.} \psalm{4} {\color{Red} Ant.} Scuto bone voluntatis tue coronasti eum domine. {\color{Red} Ps.} \psalm{5} {\color{Red} Ant.} In universa terra gloria et honore coronasti eum. {\color{Red} Ps.} \psalm{8} \V Posuisti domine super caput eius. \R Coronam de lapide precioso.
\begin{responsory}[beatus-vir-qui-suffert]
{Beatus vir qui suffert temptationem quia cum probatus fuerit accipiet coronam vite, * Quam repromisit deus diligentibus se.}{Hic accipiet benedictionem a domino et misericordiam.}
\end{responsory}%
\begin{responsory}[domine-prevenisti-eum]
{Domine prevenisti eum in benedictionibus dulcedinis posuisti in capite eius * Coronam de lapide precioso.}{Vitam peciit a te et tribuisti ei domine.}
\end{responsory}%
\begin{responsory-final}[hic-est-vere]
{Hic est vere martyr, qui pro Christi nomine sanguinem suum fudit, * Qui minas iudicum non timuit, nec terrene dignitatis gloriam quesivit, * Sed ad celestia regna pervenit.}{Iustum deduxit dominus per vias rectas et ostendit illi regnum dei.}{Sed ad}
\end{responsory-final}%
\newline {\color{Red} In Tertio Nocturno. Ant.} Iustus dominus et iusticias dilexit, equitatem vidit vultus eius. {\color{Red} Ps.} \psalm{10} {\color{Red} Ant.} Habitabit in tabernaculo tuo requiescet in monte sancto tuo. {\color{Red} Ps.} \psalm{14} {\color{Red} Ant.} Posuisti domine super caput eius coronam de lapide precioso. {\color{Red} Ps.} \psalm{20} \V Magna est gloria eius in salutari tuo. \R Gloria et magnum decorem impones super eum.
\begin{responsory}[stola-iocunditatis]
{Stola iocunditatis induit eum dominus, * Et coronam pulchritudinis posuit super caput eius.}{Cibavit illum dominus pane vite et intellectus et aqua sapientie salutaris potavit eum.}
\end{responsory}%
\begin{responsory}[corona-aurea-super]
{Corona aurea super caput eius expressa signo sanctitatis gloria honoris, * Et opus fortitudinis.}{Quoniam prevenisti eum in benedictionibus dulcedinis posuisti in capite eius coronam de lapide precioso.}
\end{responsory}%
\begin{responsory-final}[gloriosus-dei-amicus]
{Gloriosus dei amicus, N. Ihesu domini confessione fundatus inter tormentorum supplicia stetit imperterritus, * Ac demum triumphali morte insignis * Celos petivit sanguine laureatus.}{Felici commercio pro terrenis celestia, pro perituris eterna commutans.}{Celos}
\end{responsory-final}%
\newline \V Ora pro nobis beate, N. \R Ut digni efficiamur promissionibus Christi.
\newline {\color{Red} In Laudibus. Ant.} Qui me confessus fuerit coram hominibus confitebor et ego eum coram patre meo. {\color{Red} Ps.} \psalm{92} {\color{Red} Ant.} Qui sequitur me non ambulat in tenebris sed habebit lumen vite dicit dominus. {\color{Red} Ant.} Siquis michi ministraverit honorificabit eum pater meus qui est in celis, dicit dominus. {\color{Red} Ant.} Qui michi ministrat me sequatur, et ubi sum ego illic erit et minister meus. {\color{Red} Ant.} Volo pater ut ubi ego sum illic sit et minister meus. {\color{Red} Cap.} Beatus vir qui suffert. {\color{Red} Ymnus.}
\hymn{martyr-dei-qui}{Martyr dei qui unicum}{
\lettrine[lines=2]{\zallmancaps \color{Blue} M}{artyr} dei qui unicum
patris sequendo filium,
victis triumphas hostibus
victor fruens celestibus.
\par {\bfseries \color{Red} T}ui precatus munere
nostrum reatum dilue,
arcens mali contagium
vite removens tedium.
\par {\bfseries \color{Blue} S}oluta sunt iam vincula
tui sacrati corporis,
nos solve vinclis seculi
amore filii dei.
\par {\bfseries \color{Red} D}eo patri sit gloria
eiusque soli filio
cum spiritu paraclito
et nunc et in perpetuum. Amen.
}
\newline \V Corona aurea super caput eius. \R Expressa signo sanctitatis.
\newline {\color{Red} Ad Ben. Ant.} Qui vult venire post me abneget semetipsum et tollat crucem suam et sequatur me.
\newline {\color{Red} Ad Terciam. Cap.} Beatus vir qui suffert.
\begin{responsory-breve}[gloria-et-honore-breve]
{Gloria et honore * Coronasti eum domine.}{Et constituisti eum super opera manuum tuarum.}
\end{responsory-breve}%
\V Posuisti domine super caput eius.
\newline {\color{Red} Ad Sextam. Cap.}
\lettrine[lines=2]{\zallmancaps \color{Red} C}{ustodivit} illum dominus ab inimicis, et a seductoribus tutavit eum, et certamen forte dedit illi ut vinceret.
\begin{responsory-breve}[posuisti-domine-breve]
{Posuisti domine, * Super caput eius.}{Coronam de lapide precioso.}
\end{responsory-breve}%
\V Magna est gloria.
\newline {\color{Red} Ad Nonam. Cap.}
\lettrine[lines=2]{\zallmancaps \color{Blue} P}{ro} iusticia agonizare pro anima tua, et usque ad mortem certa pro iusticia.
\begin{responsory-breve}[magna-est-gloria]
{Magna est gloria eius, * In salutari tuo.}{Gloriam et magnum decorem impones super eum.}
\end{responsory-breve}%
\V Corona aurea.
\newline {\color{Red} Ad Vesperas. Super psalmos. Ant.} Qui me confessus. {\color{Red} Ps.} \psalm{109} \psalm{111} \psalm{112} \psalm{115} \psalm{125} \ul{Cap., Ymnus, \Vbar , ut in Primis Vesperis.}
\newline {\color{Red} Ad Mag. Ant.} Iste est qui pro lege dei sui morti se tradidit, non dubitavit mori, ab iniquis interfectus est, et in eternum vivit cum Christo, agnum secutus est: et accepit palmam.
{\ubsubsection{In Communi Plurimorum Martyrum Extra Tempus Paschale}{in-communi-plurimorum-martyrum-extra-tempus-paschale-breviarium}}
\newline {\color{Red} Ad Vesperas. Capitulum.}
\lettrine[lines=2]{\zallmancaps \color{Red} S}{ancti} per fidem vicerunt regna, operati sunt iusticiam adepti sunt repromissiones, in Christo Ihesu domino nostro. {\color{Red} Ymnus.}
\hymn{sanctorum-meritis-inclita}{Sanctorum meritis inclita}{
\lettrine[lines=2]{\zallmancaps \color{Blue} S}{anctorum} meritis inclita gaudia,
pangamus socii gestaque fortia,
nam gliscit animus promere cantibus
victorum genus optimum.
\par {\bfseries \color{Blue} H}i sunt quos retinens mundus inhorruit,
ipsum nam sterili flore peraridum
sprevere penitus, teque secuti sunt
rex Christe bone celitus.
\par {\bfseries \color{Red} H}i pro te furias atque ferocias
calcarunt hominum sevaque verbera
cessit his lacerans fortiter ungula
nec carpsit penetralia.
\par {\bfseries \color{Red} C}eduntur gladiis more bidentium,
non murmur resonat non querimonia,
sed corde tacito mens bene conscia
conservat patientiam. \hypertarget{que-vox-que-poterit}{\label{que-vox-que-poterit}}
\par {\bfseries \color{Blue} Q}ue vox que poterit lingua retexere
que tu martyribus munera preparas,
rubri nam fluido sanguine laureis
ditantur bene fulgidis.
\par {\bfseries \color{Red} T}e summa deitas unaque poscimus
ut culpas abluas noxia subtrahas,
des pacem famulis, nos quoque gloriam,
per cuncta tibi secula. Amen.
}
\newline \V Letamini in domino et exultate iusti. \R Et gloriamini omnes recti corde.
\newline {\color{Red} Ad Mag. Ant.} Gaudent in celis anime sanctorum qui Christi vestigia sunt secuti, et quia pro eius amore sanguinem suum funderunt, ideo cum Christo regnabunt in eternum.
\newline {\color{Red} Ad Mat. Invit.}
\begin{invitatory}[venite-adoremus-regem-qui-invitatorium]
{Venite adoremus regem regum, * Qui beatis martyribus coronam dedit glorie.}
\end{invitatory}
\newline {\color{Red} Ymnus.} \hyperlink{sanctorum-meritis-inclita}{Sanctorum.}
\newline {\color{Red} In Primo Nocturno. Ant.} Secus decursus aquarum plantavit vineam iustorum, sed in lege domini fuit voluntas eorum. {\color{Red} Ps.} \psalm{1} {\color{Red} Ant.} Tanquam aurum in fornace probavit electos dominus, et quasi holocausta accepit eos in eternum. {\color{Red} Ps.} \psalm{2} {\color{Red} Ant.} Si coram hominibus tormenta passi sunt spes electorum est immortalis in eternum. {\color{Red} Ps.} \psalm{10} \V Letamini in domino. \R
\lettrine[lines=2]{\zallmancaps \color{Blue} A}{bsterget} \hypertarget{absterget-deus-omnem}{\label{absterget-deus-omnem}} deus omnem lacrimam ab oculis sanctorum, et iam non erit amplius neque luctus neque clamor, sed nec ullus dolor, * Quoniam priora transierunt. \V Non esurient neque sitient amplius, neque cadet super illos sol neque ullus estus. Quoniam.
\begin{responsory}[viri-sancti-gloriosum]
{Viri sancti gloriosum sanguinem fuderunt pro domino, amaverunt Christum in vita sua, imitati sunt eum in morte sua, * Et ideo coronas triumphales meruerunt.}{Unus spiritus et una fides erat in eis.}
\end{responsory}%
\begin{responsory-doxology}[tradiderunt-corpora]
{Tradiderunt corpora sua propter deum ad supplicia, * Et meruerunt habere coronas perpetuas.}{Isti sunt qui venerunt ex magna tribulatione, et laverunt stolas suas in sanguine agni.}
\end{responsory-doxology}%
\newline {\color{Red} In Secundo Nocturno. Ant.} Dabo sanctis meis locum nominatum in regno patris mei, dicit dominus. {\color{Red} Ps.} \psalm{14} {\color{Red} Ant.} Sanctis qui in terra sunt eius, mirificavit omnes voluntates meas inter illos. {\color{Red} Ps.} \psalm{15} {\color{Red} Ant.} Sancti qui sperant in domino habebunt fortitudinem, assument pennas ut aquile, volabunt et non deficient. {\color{Red} Ps.} \psalm{23} \V Exultent iusti in conspectu dei. \R Et delectentur in leticia.
\begin{responsory}[sancti-tui-domine]
{Sancti tui domine mirabile consecuti sunt iter, servientes preceptis tuis, * Ut invenirentur illesi in aquis validis, terra apparuit arida et in mari rubro via sine impedimento.}{Victricem manum tuam laudaverunt pariter, et decantaverunt domine nomen sanctum tuum.}
\end{responsory}%
\begin{responsory}[verbera-carnificum]
{Verbera carnificum non timuerunt sancti dei, morientes pro Christi nomine, * Ut heredes fierent in domo domini.}{Tradiderunt corpora sua propter deum ad supplicia.}
\end{responsory}%
\begin{responsory-doxology}[o-veneranda]
{O veneranda martyrum gloriosa certamina, qui in suis corporibus pro Christo immania pertulerunt tormenta, * Et ideo percipere meruerunt immarcessibilem eterne glorie coronam.}{Despecta namque presentis vite luce, contemptoque suorum corporum cruciatu: sevientem mundum dei pro amore vicerunt.}
\end{responsory-doxology}%
\newline {\color{Red} In Tertio Nocturno. Ant.} Iusti autem in perpetuum vivent, et apud dominum est merces eorum. {\color{Red} Ps.} \psalm{32} {\color{Red} Ant.} Tradiderunt corpora sua propter deum ad supplicia: ut heredes fierent in domo domini. {\color{Red} Ps.} \psalm{33} {\color{Red} Ant.} Ecce merces sanctorum copiosa est apud deum, ipsi vero mortui sunt pro Christo, et vivunt in eternum. {\color{Red} Ps.} \psalm{45} \V Iusti autem in perpetuum vivent. \R Et apud dominum est merces eorum.
\begin{responsory}[in-circuitu-tuo]
{In circuitu tuo domine lumen est quod nunquam deficiet, ubi constituisti lucidissimas mansiones, * Ibi requiescunt sanctorum anime.}{Magnus dominus et laudabilis nimis in civitate dei nostri in monte sancto eius.}
\end{responsory}%
\begin{responsory}[sancti-mei-qui]
{Sancti mei qui in isto seculo certamen habuistis, * Mercedem laborum ego reddam vobis.}{Venite benedicti patris mei percipite regnum.}
\end{responsory}%
\begin{responsory-doxology}[laverunt-stolas]
{Laverunt stolas suas et candidas eas fecerunt * In sanguine agni.}{Isti sunt qui venerunt ex magna tribulatione, et laverunt stolas suas.}
\end{responsory-doxology}%
\newline \V Exultabunt sancti in gloria. \R Letabuntur in cubilibus suis.
\newline {\color{Red} In Laudibus. Ant.} Iustorum autem anime in manu dei sunt et non tanget illos tormentum malicie. {\color{Red} Ant.} Cum palma ad regna pervenerunt sancti, coronas decoris meruerunt de manu domini. {\color{Red} Ant.} Corpora sanctorum in pace sepulta sunt, et vivent nomina eorum in eternum. {\color{Red} Ant.} Martyres domini, dominum benedicite in eternum. {\color{Red} Ant.} Martyrum chorus laudate dominum in excelsis. {\color{Red} Cap.} Sancti per fidem. {\color{Red} Ymnus.}
\hymn{eterna-christi-munera-et}{Eterna Christi munera}{
\lettrine[lines=2]{\zallmancaps \color{Red} E}{terna} Christi munera
et martyrum victorias,
laudes canentes debitas:
letis canamus mentibus.
\par {\bfseries \color{Blue} T}errore victo seculi,
penisque spretis corporis,
mortis sacre compendio:
vitam beatam possident.
\par {\bfseries \color{Red} T}raduntur igni martyres,
et bestiarum dentibus,
armata sevit ungulis
tortoris insani manus.
\par {\bfseries \color{Blue} N}udata pendent viscera,
sanguis sacratus funditur,
sed permanent immobiles:
vite perennis gratia.
\par {\bfseries \color{Red} T}e nunc redemptor quesumus
ut ipsorum consortio,
iungas precantes servulos
in sempiterna secula. Amen.
}
\newline \V Mirabilis deus. \R In sanctis suis.
\newline {\color{Red} Ad Ben. Ant.} Isti sunt sancti qui pro dei amore minas hominum contempserunt, sancti martyres in regno celorum exultant cum angelis, o quam preciosa est mors sanctorum qui assidue assistunt ante dominum, et ab invicem non sunt separati.
\newline {\color{Red} Ad Terciam. Cap.} Sancti per fidem.
\begin{responsory-breve}[letamini-in-domino]
{Letamini in domino * Et exultate iusti.}{Et gloriamini omnes recti corde.}
\end{responsory-breve}%
\V Exultent iusti in conspectu dei.
\newline {\color{Red} Ad Sextam. Cap.}
\lettrine[lines=2]{\zallmancaps \color{Blue} R}{eddidit} deus mercedem laborum sanctorum suorum, et deduxit illos in via mirabili, et fuit illis in velamento diei, et in luce stellarum nocte.
\begin{responsory-breve}[exultent-iusti]
{Exultent iusti, * In conspectu dei.}{Et delectentur in leticia.}
\end{responsory-breve}%
\V Iusti autem in perpetuum vivent.
\newline {\color{Red} Ad Nonam. Cap.}
\lettrine[lines=2]{\zallmancaps \color{Red} I}{ustorum} anime in manu dei sunt, et non tanget illos tormentum malicie, visi sunt oculis insipientium, mori, illi autem sunt in pace.
\begin{responsory-breve}[iusti-autem-in]
{Iusti autem, * In perpetuum vivent.}{Et apud dominum est merces eorum.}
\end{responsory-breve}%
\V Mirabilis deus.
\newline {\color{Red} Ad Vesperas. Super psalmos. Ant.} Iustorum. {\color{Red} Ps.} \psalm{109} \psalm{112} \psalm{115} \psalm{125} \psalm{139} \ul{Cap., Ymnus, \Vbar , ut in Primis Vesperis.}
\newline {\color{Red} Ad Mag. Ant.} Absterget deus omnem lacrimam ab oculis sanctorum, et iam non erit amplius neque luctus neque clamor, sed nec ullus dolor quoniam priora transierunt.
{\ubsubsection{In Communi Unius Confessoris Extra Tempus Paschale}{in-communi-unius-confessoris-extra-tempus-paschale-breviarium}}
\newline {\color{Red} Ad Vesperas. Capitulum.}
\lettrine[lines=2]{\zallmancaps \color{Blue} D}{ilectus} deo et hominibus, cuius memoria in benedictione est, similem illum fecit in gloria sanctorum. {\color{Red} Ymnus.}
\hymn{iste-confessor-domini}{Iste confessor domini}{
\lettrine[lines=2]{\zallmancaps \color{Red} I}{ste} confessor domini sacratus,
festa plebs cuius celebrat per orbem,
hodie letus meruit secreta
scandere celi.
\par {\bfseries \color{Blue} Q}ui pius, prudens, humilis, pudicus,
sobrius, castus fuit et quietus,
vita dum presens vegetavit eius,
corporis artus.
\par {\bfseries \color{Red} A}d sacrum cuius tumulum, frequenter
membra languentum modo sanitati
quolibet morbo fuerint gravati
restituuntur.
\par {\bfseries \color{Blue} U}nde nunc noster chorus in honore
ipsius, ymnum canit hunc libenter,
ut piis eius meritis iuvemur
omne per evum.
\par {\bfseries \color{Red} S}it salus illi decus atque virtus,
qui supra celi residens cacumen
totius mundi machinam gubernat
trinus et unus. Amen.
}
\newline \V Amavit eum dominus et ornavit eum. \R Stola glorie induit eum.
\newline {\color{Red} Ad Mag. Ant.} Confessor domini, N. astantem plebem corrobora sancta intercessione, ut qui vitiorum pondere premimur beatitudinis tue gloria sublevemur, et te duce eterna premia consequamur.
\newline {\color{Red} Ad Mat. Invit.}
\begin{invitatory}[confessorum-regem-invitatorium]
{Confessorum regem adoremus, * Qui celestis regni meritum et gloriam contulit sancto suo, N.}
\end{invitatory}
\newline {\color{Red} Ymnus.} \hyperlink{iste-confessor-domini}{Iste confessor.}
\newline {\color{Red} In Primo Nocturno. Ant.} Beatus vir qui in lege domini meditatur, voluntas eius permanet die ac nocte, et omnia quecumque faciet semper prosperabuntur. {\color{Red} Ps.} \psalm{1} {\color{Red} Ant.} Beatus iste sanctus qui confisus est in domino, predicavit preceptum domini, constitutus est in monte sancto eius. {\color{Red} Ps.} \psalm{2} {\color{Red} Ant.} Tu es gloria mea, tu es susceptor meus domine, tu exaltas caput meum, et exaudisti me de monte sancto tuo. {\color{Red} Ps.} \psalm{3} \V Amavit. \R
\lettrine[lines=2]{\zallmancaps \color{Blue} E}{uge} \hypertarget{euge-serve-bone}{\label{euge-serve-bone}} serve bone et fidelis quia in pauca fuisti fidelis supra multa te constituam, * Intra in gaudium domini tui. \V Domine quinque talenta tradidisti michi, ecce alia quinque superlucratus sum. Intra.
\begin{responsory}[iustus-germinabit]
{Iustus germinabit sicut lilium et * Florebit in eternum ante dominum.}{Plantatus in domo domini in atriis domus dei nostri.}
\end{responsory}%
\begin{responsory-final}[iste-sanctus-digne]
{Iste sanctus digne in memoriam vertitur hominum: qui ad gaudium transiit angelorum, * Quoniam in hac peregrinatione solo corpore constitutus: cogitatione et aviditate * In illa eterna paterna conversatus est.}{Vinculis carnis absolutus, talentum sibi creditum domino suo duplicatum reportavit.}{In illa}
\end{responsory-final}%
\newline {\color{Red} In Secundo Nocturno. Ant.} Invocantem exaudivit dominus sanctum suum, dominus exaudivit eum et constituit eum in pace. {\color{Red} Ps.} \psalm{4} {\color{Red} Ant.} Letentur omnes qui sperant in te domine quoniam tu benedixisti iusto, scuto bone voluntatis coronasti eum. {\color{Red} Ps.} \psalm{5} {\color{Red} Ant.} Domine dominus noster quam admirabile est nomen tuum in universa terra, quia gloria et honore coronasti sanctum tuum, et constituisti eum super opera manuum tuarum. {\color{Red} Ps.} \psalm{8} \V Iustum deduxit dominus per vias rectas. \R Et ostendit illi regnum dei.
\begin{responsory}[magnificavit-eum-in]
{Magnificavit eum in conspectu regum, * Et dedit illi coronam glorie.}{Statuit illi testamentum sempiternum, et beatificavit illum in gloria.}
\end{responsory}%
\begin{responsory}[amavit-eum-dominus]
{Amavit eum dominus et ornavit eum stola glorie induit eum, * Et ad portas paradisi coronavit eum.}{Induit eum dominus lorica fidei et ornavit eum.}
\end{responsory}%
\begin{responsory-final}[sancte-n-christi]
{Sancte, N. Christi confessor audi rogantes servulos * Et impetratam celitus * Tu defer indulgentiam.}{O sancte, N. sydus aureum domini gratia servorum gemitus solita suscipe clementia.}{Tu defer}
\end{responsory-final}%
\newline {\color{Red} In Tertio Nocturno. Ant.} Domine iste sanctus habitabit in tabernaculo tuo, operatus est iusticiam, requiescet in monte sancto tuo. {\color{Red} Ps.} \psalm{14} {\color{Red} Ant.} Vitam petiit a te tribuisti ei domine, gloriam et magnum decorem imposuisti super eum, posuisti in capite eius coronam de lapide precioso. {\color{Red} Ps.} \psalm{20} {\color{Red} Ant.} Hic accipiet benedictionem a domino et misericordiam a deo salutari suo, quia hec est generatio querentium dominum. {\color{Red} Ps.} \psalm{23} \V Iustus ut palma florebit in domo domini. \R Sicut cedrus Libani multiplicabitur.
\begin{responsory}[iustus-ut-palma]
{Iustus ut palma florebit sicut cedrus Libani multiplicabitur plantatus in domo domini, * In atriis domus dei nostri florebit.}{Gloria et divitie in domo eius et iusticia eius manet in seculum seculi.}
\end{responsory}%
\begin{responsory}[iustum-deduxit]
{Iustum deduxit dominus per vias rectas et ostendit illi regnum dei, et dedit illi scientiam sanctorum, * Honestavit illum in laboribus et complevit labores illius.}{Immortalis est enim memoria illius quia apud deum nota est et apud homines.}
\end{responsory}%
\begin{responsory-final}[sint-lumbi-vestri]
{Sint lumbi vestri precincti, et lucerne ardentes in manibus vestris, * Et vos similes hominibus expectantibus dominum suum, * Quando revertatur a nuptiis.}{Vigilate ergo quia nescitis qua hora dominus vester venturus sit.}{Quando}
\end{responsory-final}%
\newline \V Ora pro nobis beate, N. \R Ut digni efficiamur.
\newline {\color{Red} In Laudibus. Ant.} Domine quinque talenta tradidisti michi, ecce alia quinque superlucratus sum. {\color{Red} Ps.} \psalm{92} {\color{Red} Ant.} Beatus ille servus quem cum venerit dominus et pulsaverit ianuam invenerit vigilantem. {\color{Red} Ant.} Fidelis servus et
%pp135
prudens quem constituit dominus super familiam suam. {\color{Red} Ant.} Ecce vere Israhelita in quo dolus non est. {\color{Red} Ant.} Serve bone et fidelis intra in gaudium domini tui. {\color{Red} Cap.} Dilectus deo. {\color{Red} Ymnus.}
\hymn{ihesu-redemptor-omnium}{Ihesu redemptor omnium}{
\lettrine[lines=2]{\zallmancaps \color{Red} I}{hesu} redemptor omnium
corona confitentium,
in hac die clementius
nostris faveto vocibus.
\par {\bfseries \color{Blue} T}ui sacri qua nominis
confessor almus claruit,
cuius celebrat annua
devota plebs sollemnia.
\par {\bfseries \color{Red} Q}ui rite mundi gaudia
huius caduca respuens,
cum angelis celestibus
letus potitur premiis.
\par {\bfseries \color{Blue} H}uius benignus annue
nobis sequi vestigia,
huius precatu servulis
dimitte noxam criminis.
\par {\bfseries \color{Red} S}it Christe rex piissime
tibi patrique gloria,
cum spiritu paraclito
in sempiterna secula. Amen.
}
\newline \V Iustus germinabit sicut lilium. \R Et florebit in eternum ante dominum.
\newline {\color{Red} Ad Ben. Ant.} Euge serve bone et fidelis, quia in pauca fuisti fidelis supra multa te constituam dicit dominus.
\newline {\color{Red} Ad Terciam. Cap.} Dilectus.
\begin{responsory-breve}[amavit-eum-dominus-breve]
{Amavit eum dominus, * Et ornavit eum.}{Stola glorie induit eum.}
\end{responsory-breve}%
\V Iustum deduxit.
\newline {\color{Red} Ad Sextam. Capitulum.}
\lettrine[lines=2]{\zallmancaps \color{Blue} I}{ustus} cor suum tradet ad vigilandum diluculo ad deum qui fecit illum, et in conspectu altissimi deprecabitur.
\begin{responsory-breve}[iustum-deduxit-breve]
{Iustum deduxit dominus * Per vias rectas.}{Et ostendit illi regnum dei.}
\end{responsory-breve}%
\V Iustus ut palma.
\newline {\color{Red} Ad Nonam. Cap.}
\lettrine[lines=2]{\zallmancaps \color{Red} I}{ustum} deduxit dominus per vias rectas, et ostendit illi regnum dei, et dedit illi scientiam sanctorum, honestavit illum in laboribus, et complevit labores illius.
\begin{responsory-breve}[iustus-ut-palma-breve]
{Iustus ut palma florebit * In domo domini.}{Sicut cedrus Libani multiplicabitur.}
\end{responsory-breve}%
\V Iustus germinabit sicut lilium.
\newline {\color{Red} Ad Vesperas. Super psalmos. Ant.} Domine quinque. {\color{Red} Ps.} \psalm{109} \psalm{111} \psalm{112} \psalm{115} \psalm{125} \ul{Capitulum, Ymnus, \Vbar , ut in Primis Vesperis.}
\newline {\color{Red} Ad Mag. Ant.} Similabo eum viro sapienti qui edificavit domum suam supra petram.
{\ubsubsection{In Communi Unius Virginis Extra Tempus Paschale}{in-communi-unius-virginis-extra-tempus-paschale-breviarium}}
\newline {\color{Red} Ad Vesperas. Capitulum.}
\lettrine[lines=2]{\zallmancaps \color{Blue} C}{onfitebor} tibi domine rex, et collaudabo te deum salvatorem meum, confitebor nomini tuo quoniam adiutor et protector factus es michi, et liberasti corpus meum a perditione. {\color{Red} Ymnus.}
\hymn{virginis-proles-opifexque}{Virginis proles opifexque}{
\lettrine[lines=2]{\zallmancaps \color{Red} V}{irginis} proles, opifexque matris,
virgo quem gessit peperitque virgo,
virginis festum canimus tropheum
accipe votum.
\par {\bfseries \color{Blue} H}ec tua virgo duplici beata
sorte, dum gestit fragilem domare
corporis sexum domuit cruentum
corporis seclum.
\par {\bfseries \color{Red} U}nde nec mortem, nec amica mortis
seva penarum genera pavescens,
sanguine fuso meruit sacrata
scandere celum. \hypertarget{huius-obtentu}{\label{huius-obtentu}}
\par {\bfseries \color{Blue} H}uius obtentu deus alme nostris
parce iam culpis vitia remittens,
quo tibi puri resonemus almum
pectoris hymnum.
\par {\bfseries \color{Red} G}loria patri geniteque proli
et tibi compar utriusque semper
super alme, deus unus omni
tempore secli. Amen.
}
\newline \V Diffusa est gratia in labiis tuis. \R Propterea benedixit te deus in eternum.
\newline {\color{Red} Ad Mag. Ant.} Ista est virgo sapiens quam dominus vigilantem invenit, que accepta lampade sumpsit secum oleum, et veniente domino introivit cum eo ad nuptias.
\newline {\color{Red} Ad Mat. Invit.}
\begin{invitatory}[agnum-sponsum-invitatorium]
{Agnum sponsum virginum, * Venite adoremus dominum Ihesum Christum.}
\end{invitatory}
\newline {\color{Red} Ymnus.} \hyperlink{virginis-proles-opifexque}{Virginis proles.}
\newline {\color{Red} In Primo Nocturno. Ant.} Ante thorum huius virginis frequentate nobis dulcia cantica dramatis. {\color{Red} Ps.} \psalm{8} {\color{Red} Ant.} Unguentum effusum nomen tuum, ideo adolescentule dilexerunt te nimis. {\color{Red} Ps.} \psalm{18} {\color{Red} Ant.} O quam pulchra est casta generatio cum caritate. {\color{Red} Ps.} \psalm{23} \V Diffusa est gratia. \R
\lettrine[lines=2]{\zallmancaps \color{Blue} V}{eni} \hypertarget{veni-sponsa}{\label{veni-sponsa}} sponsa Christi accipe coronam quam tibi dominus preparavit, * Pro cuius amore sanguinem tuum fudisti, et cum angelis in paradisum introisti. \V Veni electa mea et ponam in te thronum meum, quia concupivit rex speciem tuam. Pro cuius.
\begin{responsory}[propter-veritatem]
{Propter veritatem et mansuetudinem et iusticiam, * Et deducet te mirabiliter dextera tua.}{Audi filia et vide et inclina aurem tuam.}
\end{responsory}%
\begin{responsory-final}[specie-tua-et]
{Specie tua et pulchritudine tua * Intende prospere, * Procede et regna.}{Diffusa est gratia in labiis tuis propterea benedixit te deus in eternum.}{Procede}
\end{responsory-final}%
\newline {\color{Red} In Secundo Nocturno. Ant.} Specie tua et pulchritudine tua intende prospere procede et regna. {\color{Red} Ps.} \psalm{44} {\color{Red} Ant.} Adiuvabit eam deus vultu suo, deus in medio eius non commovebitur. {\color{Red} Ps.} \psalm{45} {\color{Red} Ant.} Leva eius sub capite meo, et dextera illius amplexabitur me. {\color{Red} Ps.} \psalm{86} \V Specie tua et pulchritudine. \R Intende prospere procede et regna.
\begin{responsory}[grata-facta-est]
{Grata facta est a domino in certamine quia apud deum et apud homines glorificata est, in conspectu principum loquebatur sapientiam, * Et dominus omnium dilexit eam.}{Erecta namque in virtutis culmine tormenta derisit premia calcavit.}
\end{responsory}%
\begin{responsory}[hec-est-virgo]
{Hec est virgo sapiens quam dominus vigilantem invenit que acceptis lampadibus sumpsit secum oleum, * Et veniente domino introivit cum eo ad nuptias.}{Ista est una de numero prudentum que aptavit lampades suas.}
\end{responsory}%
\begin{responsory-doxology}[offerentur-regi]
{Offerentur regi virgines domino post eam proxime eius offerentur tibi * In leticia et exultatione.}{Prudentes virgines aptate lampades vestras ecce sponsus venit exite obviam ei.}
\end{responsory-doxology}%
\newline {\color{Red} In Tertio Nocturno. Ant.} Dum esset rex in accubitu suo nardus mea dedit odorem suavitatis. {\color{Red} Ps.} \psalm{95} {\color{Red} Ant.} Hec est que nescivit, thorum in delicto habebit fructum in respectione animarum sanctarum. {\color{Red} Ps.} \psalm{96} {\color{Red} Ant.} Hec est virgo sancta atque gloriosa, quia dominus omnium dilexit eam. {\color{Red} Ps.} \psalm{97} \V Adiuvabit eam deus vultu suo. \R Deus in medio eius non commovebitur.%missing eius in versiculus
\begin{responsory}[audivi-vocem-de]
{Audivi vocem de celo venientem, venite omnes virgines sapientissime, * Oleum recondite in vasis vestris dum sponsus advenerit.}{Media nocte clamor factus est ecce sponsus venit.}
\end{responsory}%
\begin{responsory}[veni-electa-mea]
{Veni electa mea et ponam in te thronum meum, * Quia concupivit rex speciem tuam.}{Specie tua et pulchritudine tua intende prospere procede et regna.}
\end{responsory}%
\begin{responsory-final}[regnum-mundi]
{Regnum mundi et omnem ornatum seculi contempsi propter amorem domini mei Ihesu Christi, * Quem vidi quem amavi, * In quem credidi quem dilexi.}{Eructavit cor meum verbum bonum dico ego opera mea regi.}{In quem}
\end{responsory-final}%
\newline \V Ora pro nobis beata, N. \R Ut digni.
\newline {\color{Red} In Laudibus. Ant.} Hec est virgo sapiens et una de numero prudentum. {\color{Red} Ps.} \psalm{92} {\color{Red} Ant.} Hec est virgo sapiens quam dominus vigilantem invenit. {\color{Red} Ant.} Veni sponsa Christi accipe coronam quam tibi dominus preparavit in eternum. {\color{Red} Ant.} Media nocte clamor factus est, ecce sponsus venit, exite obviam ei. {\color{Red} Ant.} Tunc surrexerunt omnes virgines ille et ornaverunt lampades suas. {\color{Red} Cap.} Confitebor. {\color{Red} Ymnus.}
\hymn{ihesu-corona-virginum}{Ihesu corona virginum}{
\lettrine[lines=2]{\zallmancaps \color{Red} I}{hesu} corona virginum
quem mater illa concipit,
que sola virgo peperit,
hec vota clemens accipe.
\par {\bfseries \color{Blue} Q}ui pascis inter lilia
septus choreis virginum,
sponsas decorans gloria,
sponsisque reddens premia.
\par {\bfseries \color{Red} Q}uocumque pergis virgines
sequuntur, atque laudibus
post te canentes cursitant,
ymnosque dulces personant.
\par {\bfseries \color{Blue} T}e deprecamur largius
nostris adauge sensibus
nescire prorsus omnia
corruptionis vulnera.
\par {\bfseries \color{Red} S}it Christe rex piissime,
tibi patrique gloria,
cum spiritu paraclyto
in sempiterna secula. Amen.
}
\newline \V Adducentur regi virgines post eam. \R Proxime eius afferentur tibi.%typo: Aducentur
\newline {\color{Red} Ad Ben. Ant.} Veni electa mea et ponam in te thronum meum, quia concupivit rex speciem tuam.
\newline {\color{Red} Ad Terciam. Cap.} Confitebor.
\begin{responsory-breve}[diffusa-est-gratia]
{Diffusa est gratia * In labiis tuis.}{Propterea benedixit te deus in eternum.}
\end{responsory-breve}%
\V Specie tua et pulchritudine.
\newline {\color{Red} Ad Sextam. Capitulum.}
\lettrine[lines=2]{\zallmancaps \color{Red} E}{mulor} vos dei emulatione, despondi enim vos uni viro virginem castam exhibere Christo.
\begin{responsory-breve}[specie-tua-et-breve]
{Specie tua * Et pulchritudine tua.}{Intende prospere procede et regna.}
\end{responsory-breve}%
\V Adiuvabit eam.
\newline {\color{Red} Ad Nonam. Cap.}
\lettrine[lines=2]{\zallmancaps \color{Blue} O}{}\ quam pulchra est casta generatio cum caritate, immortalis est enim memoria illius, quoniam et apud deum nota est et apud homines.
\begin{responsory-breve}[adiuvabit-eam-deus]
{Adiuvabit eam * Deus vultu suo.}{Deus in medio eius non commovebitur.}
\end{responsory-breve}%
\V Adducentur.
\newline {\color{Red} Ad Vesperas. Super psalmos. Ant.} Hec est virgo sapiens et una. {\color{Red} Ps.} \psalm{109} \psalm{112} \psalm{121} \psalm{126} \psalm{147} \ul{Capitulum, Ymnus, \Vbar , ut in Primis Vesperis.}
\newline {\color{Red} Ad Mag. Ant.} Accinxit fortitudine lumbos suos, et roboravit brachium suum. Ideoque lucerna eius non extinguetur in sempiternum.
{\ubsubsection{In Communi Sanctorum Trium Lectionum Extra Tempus Paschale}{in-communi-sanctorum-trium-lectionum-extra-tempus-paschale-breviarium}}
{\ubsubsection{Unius Martyris}{unius-martyris-breviarium}}
\newline {\color{Red} Ad Mat. Invit.}
\begin{invitatory}[regem-martyrum-invitatorium]
{Regem martyrum dominum, * Venite adoremus.}
\end{invitatory}
\newline {\color{Red} Super psalmos. Ant.} Letabitur iustus in domino et sperabit in eo et laudabuntur omnes recti corde.
{\ubsubsection{Plurimorum Martyrum}{plurimorum-martyrum-breviarium}}
\newline {\color{Red} Ad Mat. Invit.} \hyperlink{regem-martyrum-invitatorium}{Regem martyrum.} \ul{ut supra.} {\color{Red} Super psalmos. Ant.} Hi sunt qui venerunt ex magna tribulatione, et laverunt stolas suas in sanguine agni.%typo: venerut
{\ubsubsection{Unius Confessoris}{unius-confessoris-breviarium}}
\newline {\color{Red} Ad Mat. Invit.}
\begin{invitatory}[regem-confessorum-invitatorium]
{Regem confessorum dominum, * Venite adoremus.}
\end{invitatory}
\newline {\color{Red} Super psalmos. Ant.} Iste sanctus digne in memoriam vertitur hominum, qui ad gaudium transiit angelorum.
{\ubsubsection{Unius Virginis}{unius-virginis-breviarium}}
\newline {\color{Red} Ad Mat. Invit.}
\begin{invitatory}[regem-virginum-invitatorium]
{Regem virginum dominum, * Venite adoremus.}%typo: ad adoremus
\end{invitatory}
\newline {\color{Red} Super psalmos. Ant.} Veniente sponso prudens virgo preparata introivit cum eo ad nuptias.
\newline \ul{Cetera omnia fiant sicut supra in rubrica de festo trium lectionum: notatum est.}
{\ubsubsection{Quomodo Sint Memorie Faciende}{quomodo-sint-memorie-faciende-breviarium}}
\newline \ul{Ad faciendam memoriam Unius Martyris. In Vesperis. Ant.} Hic est vere. \ul{\Vbar .} Gloria et honore. \ul{In Laudibus. Ant.} Qui vult. \ul{\Vbar .} Corona.
\newline \ul{Ad faciendam memoriam Plurimorum Martyrum. In Vesperis. Ant.} Gaudent. \ul{\Vbar .} Letamini. \ul{In Laudibus. Ant.} Isti sunt sancti. \ul{\Vbar .} Mirabilis.
\newline \ul{Ad faciendam memoriam Unius Confessoris: In Vesperis. Ant.} Confessor. \ul{\Vbar .} Amavit. \ul{In Laudibus. Ant.} Euge. \ul{\Vbar .} Iustus germinabit.
\newline \ul{Ad faciendam memoriam Plurimorum Confessorum. In Vesperis. Ant.} Fulgebunt iusti et tanquam scintille in arundineto discurrent, iudicabunt nationes et regnabunt in eternum. \ul{\Vbar .} Letamini. \ul{In Laudibus. Ant.} Sint lumbi vestri precincti et lucerne ardentes in manibus vestris, et vos similes hominibus expectantibus dominum suum quando revertatur a nuptiis. \ul{\Vbar .} Mirabilis.
\newline \ul{Ad faciendam memoriam Unius Virginis. In Vesperis. Ant.} Ista est. \ul{\Vbar .} Diffusa. \ul{In Laudibus. Ant.} Veni electa. \ul{\Vbar .} Adducentur regi.
\newline \ul{Ad faciendam memoriam Plurimorum Virginum. In Vesperis. Ant.} Prudentes virgines aptate lampades vestras, ecce sponsus venit exite obviam ei. \ul{\Vbar .} Adducentur regi. \ul{In Laudibus. Ant.} Simile est regnum celorum decem virginibus que accipientes lampades suas exierunt obviam sponso et sponse. \ul{\Vbar .} Media nocte clamor factus est. \ul{\Rbar .} Ecce sponsus venit exite obviam ei.
\newline \ul{Orationes vero ad memorias dicantur sicut in suis locis notate sunt.}
\newline \ul{Ad faciendam beati Dominici per totum annum quando fuerit facienda, preterquam infra octavas ipsius. In Vesperis. Ant.} Magne pater. \ul{\Vbar .} Os iusti meditabitur sapientiam. \ul{\Rbar .} Et lingua eius loquetur iudicium. \ul{In Laudibus. Ant.} Benedictus redemptor. \ul{\Vbar .} Lex dei eius in corde ipsius. \ul{\Rbar .} Et non supplantabuntur gressus eius. \ul{Ad utramque memoriam, Oratio,} \hyperlink{deus-qui-ecclesiam}{Deus qui ecclesiam.}
\newline \ul{Si in festo alicuius martyris vel infra octavas fiat memoria de alio martyre, tunc in Vesperis dicatur ad memoriam de secundo sancto: versiculus secundi nocturni. In Laudibus, vero dicatur versiculus tercii nocturni. Eodem modo fiat de aliis sanctis.}
{\ubsubsection{Lectiones De Communi Sanctorum}{lectiones-de-communi-sanctorum-breviarium}}
\lettrine[lines=2]{\zallmancaps \color{Blue} N}{}\ul{\textsc{ota} quod lectiones infra scripte de communi sanctorum ad hoc sunt facte prolixe, ut qui volunt legant eas plene. Qui vero non delectantur in prolixis lectionibus faciant de una lectione duas vel tres, et sic legant modo istas modo illas ad fastidium tollendum, quod posset oriri ex repetitione eiusdem, cum non habeantur lectiones de sanctis in parvis breviariis nisi de communi sanctorum.}
{\ubsubsection{In Festo Unius Vel Plurimorum Apostolorum Novem Lectionum}{in-festo-unius-vel-plurimorum-apostolorum-novem-lectionum-breviarium}}
\newline \ul{Maximus in sermone de apostolis Petro et Paulo,} Gloriosissimos. \grecross \ {\color{Red} Lectio Prima.}
% https://mlat.uzh.ch/browser/cps_2.MaxTau.Homili3#:69.1;4
\lettrine[lines=2]{\zallmancaps \color{Red} A}{postoli} Latino sermone dicuntur missi. Qui ergo honorant missos: manifestum est eos honorare mittentem, quoniam dignitas que defertur ministris, illi cuius ministri sunt exhibetur, ut ait ipse salvator ad discipulos. Qui vos audit me audit, et qui vos recipit me recipit. Vere beata apostolorum merita, in quibus se Christus et recipi predicat et audiri. Beati sunt et illi: quorum devotio delata apostolis, recurrit in Christum. Tenentes itaque fratres huius promissionis fidem, de suppliciis patrum nostrorum pro Christi confessione susceptis, exultemus, quoniam qui de martyrum morte letatur: martyres non dubitat cum Christo regnare post mortem.
\newline \ul{Augustinus in libro qui dicitur catholice fidei sive manuale sive speculum.} \grecross \ {\color{Red} Lectio Secunda.}
% https://mlat.uzh.ch/browser/cps_2.Alcuin.ConFid5#:2.1;2
\lettrine[lines=2]{\zallmancaps \color{Blue} H}{i} sunt quibus mediator Dei et hominum manifestavit nomen tuum domine potissimum: et annunciavit eis omnia. Hi sunt, qui constituti sunt principes super omnem terram: quorum tuba predicationis ad totius aures mundi pervenit. Isti namque sicut ab initio viderunt et ministri fuerunt sermonis: tradiderunt nobis regni tui mysteria nunquam ante comperta, te nimirum verba illorum credibilia faciente, signis et portentis et virtutibus variis: mirisque sancti spiritus distributionibus. Hec sane causa est qua freti sine ulla ambiguitate suscipimus, quicquid mater ecclesia universalis et apostolica docet, usque adeo ut pro tantorum documentorum assertione per tuam gratiam omni genere tormentorum mori non timeamus: scientes pro certo vitam gloriosam restare nobis.
\newline \ul{Idem ad Volusianum.} \grecross \ {\color{Red} Lectio Tertia.}
\lettrine[lines=2]{\zallmancaps \color{Red} H}{i} sunt qui impleti spiritu sancto linguis omnium gentium loquuntur, errores fidenter arguunt, veritatem saluberrimam predicant: homines ad penitentiam hortantur, preterite vite culpabilis, indulgentiam de divina gratia pollicentur, et predicationem veritatis miracula consequuntur. Excitata est tandem adversus eos infidelitas seva, tolerant predicta, sperant promissa: docent precepta. Numero exigui per mundum discurrunt: populos mirabili facilitate convertunt. Inter inimicos augentur, persecutionibus crescunt: ac per afflictionum angustias usque in extrema terrarum dilatati, alternis adversitatibus et prosperitatibus rerum, patientiam et temperantiam vigilanter exercent.
\newline \ul{Crisostomus in omelia decima quinta de Pentecostes. \grecross} \ {\color{Red} Lectio Quarta.}
\lettrine[lines=2]{\zallmancaps \color{Blue} I}{sti} non solum stelle, sed supra stellas sunt: quia stelle in celo, apostoli super celos sunt. Stelle de igne sensibili: apostoli de igne intelligibili. Stelle in die obscurantur: apostoli die ac nocte suis radiis effulgent. Stelle orto sole tenebrescunt: apostoli sole iusticie resplendente lucescunt. Stelle in resurrectione cadent ut folia: apostoli vero rapientur in nubibus obviam Christo in aera. In istis sideribus aliud antifer: aliud apostolatur lucifer. In apostolis autem luciferi sunt omnes: et ideo stellis maiores. Quicumque eos vocaverit luminaria mundi non peccavit: non solum dum essent in corpore, sed eciam magis nunc. Gratie enim et merita sanctorum non minuuntur morte, non retardantur die, non dissolvuntur in terra: nec nocte obscurantur, nec tenebris obumbrantur.
\newline {\color{Red} Lectio Quinta.}
\lettrine[lines=2]{\zallmancaps \color{Red} P}{iscatores} quidem erant apostoli, et illis dormientibus retia nunc operantur. Et tunc quidem pisces capiebant ad mortem: nunc autem homines capiunt ad salutem. Vinitores erant: quoniam labruscas eradicaverunt: et pietatis semina plantaverunt. Et illis absentibus corpore: vinea nunc florescit et botros affert. Fuerunt et columne: quoniam virtute sua ecclesie robur firmatum est. Fuerunt et fundamentum: quoniam in confessione eorum fundata est ecclesia. Portus quoque: quoniam tempestates impias represserunt. Gubernatores eciam: quoniam orbem terrarum in viam rectam direxerunt. Pastores: quoniam lupos abiecerunt, et oves conservaverunt. Aratores: quoniam spinas eradicaverunt. Medici: quoniam vulnera nostra curaverunt.
\newline {\color{Red} Lectio Sexta.}
\lettrine[lines=2]{\zallmancaps \color{Blue} F}{uerunt} et plantatores, et fructores, et cursores, et athlete, et duces, et certatores: et coronas gestantes. Adducam paulum in medium, hec omnia facientem. Vis eum plantatorem videre? Audi. Ego plantavi: Apollo rigavit. Vis structorem videre? Sicut sapiens inquit architectus fundamentum posui. Vis dimicantem videre? Sic pugno ait: non ut aerem cedens. Vis videre cursorem? Ab Hierosolima et in circuitu usque ad Iliricum et inde ad Hispanias: repletum est evangelium. Vis athletam videre? Non est nobis colluctatio ait, adversus carnem et sanguinem. Vis videre ducem? Assumite ait arma dei. Vis videre certantem? Bonum inquit certamen certavi. Vis videre coronatum? Decetero reposita est michi ait corona iusticie.
\newline {\color{Red} Lectio VII.}
\lettrine[lines=2]{\zallmancaps \color{Red} P}{eragrarunt} piscatores orbem terrarum: et infirmum eum invenientes ad sanitatem reduxerunt, et in ruina positum ad stabilitatem revocaverunt: non scuta moventes, non arcus tendentes, non sagittas mittentes, non pecuniam largientes, non in eloquentia confidentes: pecunia indigentes, sed regnum celorum possidentes, non habentes humana solatia: habentes autem dominum secum. Ego enim inquit vobiscum sum: usque ad consummationem seculi. Peragrarunt universum orbem oves simul cum lupis, et vulnerate a bestiis non sunt oves: sed magis lupi ad ovium mansuetudinem sunt conversi.
\newline \ul{Gregorius in ultima parte Ezechielis, Omelia sexta.} \grecross \ {\color{Red} Lectio VIII.}
\lettrine[lines=2]{\zallmancaps \color{Blue} H}{orum} vitam intuens Ysaias: ait. Qui sunt isti qui ut nubes volant? Qui volare ut nubes dicuntur: quia in terra viventes extra terram fuerunt per omne quod egerunt. Unde et per quandam nubem dicitur. In carne ambulantes: non secundum carnem militamus. Priores etenim patres coniugiis utebantur, filios procreabant, substantias possidebant: curis rei familiaris intendebant. Istos autem prophecie spiritu videns substantias deserere, coniugia non appetere, filios non procreare, nichil in terra querere, nichil possidere: non eos per terram velut homines ambulantes, sed ut nubes volantes nominavit. Volant enim: quia mente ea que sunt celestia contemplantur. Qui terram quasi non tangunt: quia in ea nichil appetunt.
\newline \ul{Rabanus super Matheum.} \grecross \ {\color{Red} Lectio IX.}
% https://mlat.uzh.ch/browser/cps_2.RabMau.CoInMa5#:4.3;27
\lettrine[lines=2]{\zallmancaps \color{Red} I}{sti} sunt quibus dominus dicit. Ecce ego mitto vos. Hoc est, ego qui vos prescivi, qui vos de mundo elegi, qui ea que ventura sunt vobis omnia novi, qui eciam laborantibus vobis condignam mercedem restituam in futuro: ipse vos secundum nomen vestrum ex discipulis missos faciam.
% \newline \ul{Origenes super Numeros, Omelia XXVIII. \grecross }
% https://mlat.uzh.ch/browser/cps_2.RabMau.CoInEz2#:20.1;7
This is a note to who Rabanus was citing for the following line, Quis fratres, etc.
% \newline 
Quis fratres karissimi dignus habebitur tam beate sortis electione, non cui dicatur a domino intra in gaudium domini tui, sed cui dicat, sedete et vos mecum super duodecim thronos iudicantes duodecim tribus Israhel? Et de quibus dicat. Pater volo: ut ubi ego sum, et isti sint mecum. Volo eciam istos esse reges: ut ego sum rex regum. Volo et istos habere dominationem, ut ego sum dominus dominantium. Beati qui ad ista conscenderint fastigia meritorum.
{\ubsubsection{In Festo Unius Vel Plurimorum Apostolorum Trium Lectionum}{in-festo-unius-vel-plurimorum-apostolorum-trium-lectionum-breviarium}}
\newline \ul{Crisostomus super Matheum, Omelia ambulans Ihesus secus mare.} \grecross \ {\color{Red} Lectio Prima.}
% https://la.wikisource.org/wiki/Catena_aurea_in_quatuor_Evangelia/Catena_in_Matthaeum#Lectio_8_4
\lettrine[lines=2]{\zallmancaps \color{Blue} O}{mnis} rex pugnaturus contra adversarium regni, prius congregat exercitum: et sic vadit ad pugnam. Sic et dominus contra diabolum pugnaturus: prius congregavit apostolos verbi dei ministros. Expugnatio enim diaboli: est predicatio veritatis. Sagitta mortifera in visceribus eius est sermo iusticie: deiectio eius, est opus signorum. Spoliatio potestatis eius: est conversio fidelium. Et rex quidem terrenus congregat exercitum, ut illorum laboribus, ipse gloriam assequatur, iste autem congregat apostolos, non ut eorum laboribus ipse gloriam consequatur, sed ut suo labore illis victoriam prestet.
\newline \ul{Augustinus de civitate dei, libro vicesimo secundo.} \grecross \ {\color{Red} Lectio Secunda.}
\lettrine[lines=2]{\zallmancaps \color{Red} N}{on} autem peritos grammatica, non armatos dyalectica, non inflatos rethorica, sed disciplinis liberalibus ineruditos, et quantum ad phylosophorum doctrinam impolitos, piscatores cum fidei retibus ad mare huius seculi paucissimos misit: atque ita et ex omni genere tam multos pisces, et tanto mirabiliores quanto rariores, eciam ipsos philosophos cepit.
\newline \ul{Idem super Iohannem Omelia Septima.}
\newline Volens enim superborum frangere cervicem, non quesivit, per oratorem piscatorem: sed de piscatore lucratus est imperatorem.
\newline \ul{Idem super primam ad Corinthios. \grecross }
\newline Si enim sapientes et divites ac nobiles primitus eligerentur, merito talium rerum sibi eligi viderentur. Propter hoc ergo non elegit reges vel senatores, vel philosophos vel oratores: sed elegit plebeios pauperes,
%pp136
indoctos piscatores: Vnde et Nathanael doctus non est in apostolum electus.
\newline \ul{Gregorius in Omelia vicesimo sexto.} \grecross \ {\color{Red} Lectio Tertia.}
\lettrine[lines=2]{\zallmancaps \color{Blue} H}{i} sunt qui prius spiritum sanctum habuerunt ut innocenter viverent: et quibusdam in predicatione prodessent. Post resurrectionem vero patenter acceperunt: ut prodesse pluribus possent. Unde et in hac datione spiritus dicitur. Quorum remiseritis peccata remittuntur eis: et quorum retinueritis retenta sunt. Libet intueri, discipuli ad tanta humilitatis onera vocati, ad quantum glorie culmen perducti sunt. Ecce non solum de semetipsis securi fiunt: sed eciam obligationis aliene, potestatem relaxationis accipiunt. Principatum superni iudicii sortiuntur: ut vice Dei quibusdam peccata retineant, quibusdam relaxent. Ecce qui districtum Dei iudicium metuunt, animarum iudices fiunt, et alios damnant vel liberant: qui semetipsos damnari metuebant.
{\ubsubsection{In Festo Unius Evangeliste Novem Lectionum}{in-festo-unius-evangeliste-novem-lectionum-breviarium}}
\newline \ul{Gregorius super primam partem Ezechielis.} \grecross \ {\color{Red} Lectio Prima.}
\lettrine[lines=2]{\zallmancaps \color{Red} I}{n} medio eius similitudo quatuor animalium. Quod in medio eius dicitur sive electri scilicet Christi, sive ignis: nil obstat intelligi, quia quatuor hec animalia, sancti scilicet evangeliste et ex eiusdem domini incarnatione ad fidei virtutem solidati sunt: et in igne persecutionis multis tribulationibus afflicti. Et hic aspectus eorum: similitudo hominis in eis. Quis hic homo describitur nisi ille de quo scriptum est et habitu inventus est ut homo? Hec igitur animalia ad huius similitudinem non haberent. Quod egregius predicator ostendit dicens. Imitatores mei estote: sicut et ego Christi. Et quatuor facies uni: et quatuor penne uni. Quid per faciem nisi noticia, quid per pennam nisi volatus designatur? Per faciem quippe unusquisque cognoscitur: per pennam vero in altum, avium corpora sublevantur.
\newline {\color{Red} Lectio Secunda.}
\lettrine[lines=2]{\zallmancaps \color{Blue} F}{acies} itaque ad fidem pertinet, penna ad contemplationem. Per fidem namque a domino cognoscimur, sicut ipse dicit. Cognosco oves meas. Per contemplationem vero: super nosmetipsos tollimur. Quatuor ergo facies uni sunt: quia noticia fidei qua cognoscimur a domino, ipsa est in uno que est simul in quatuor. Et quatuor penne uni, quia filium dei simul omnes concorditer predicant: et ad divinitatem eius penna contemplationis volant. Evangelistarum vero facies ad humanitatem pertinent, penne ad divinitatem, quia in eo quem corporeum inspiciunt: quasi facies intendunt. Sed dum hunc esse incircumscriptum ex divinitate annunciant: per contemplationis pennam levantur. Et pedes eorum pedes recti: quia sanctorum evangelistarum atque omnium perfectorum opera ad sequendam iniquitatem non sunt retorta.
\newline {\color{Red} Lectio Tertia.}
\lettrine[lines=2]{\zallmancaps \color{Red} U}{t} vero in eisdem vite gravitas, fortitudo atque discretio monstraretur: recte subiungitur. Planta pedis eorum planta pedis vituli, scilicet mature incedens et fortis et divisa: quia unusquisque predicator et venerationem habet in maturitate, et fortitudinem in opere, et divisionem ungule in discretione. Et scintille quasi aspectus eris candentis. Quis es metallum valde sonorum est: et in omnem terram exivit sonus eorum. Bene autem es candens dicitur: quia vita predicantium sonat et ardet. Es ergo candens: est predicatio accensa. Sed de candente ere scintille prodeunt, quia de eorum exhortationibus verba flammantia procedunt: que eos quos in corde tetigerint incendunt. Sed scintille iste subtiles valde sunt: quia cum de celesti patria loquuntur, non tantum valent aperire verbo: quantum possunt ardere desiderio. Candens ergo es scintillas proicit, quando vix tenuiter predicator loqui sufficit hoc unde ipse fortiter ignescit.
\newline {\color{Red} Lectio Quarta.}
\lettrine[lines=2]{\zallmancaps \color{Blue} E}{t} manus hominis sub pennis eorum in quatuor partes. Possunt in hoc loco quatuor partes, quatuor mundi regiones accipi: quia sanctorum predicatio actore Deo in cunctis mundi partibus est egressa. Manus autem hominis sub pennis eorum est: quia nisi Deus homo fieret qui mentes predicantium ad celestia sublevasset, illa que apparent animalia: non volarent. Et facies et pennas per quatuor partes habebant: quia videlicet in cunctis mundi regionibus demonstrant, quicquid de humanitate, quicquid de divinitate nostri redemptoris sentiunt. Iuncteque erant penne eorum alterius ad alterum, quia omnis eorum virtus, omnis sapientia qua ceteros transcendunt: vicissim sibi in caritatis atque concordie pace sociatur.
\newline {\color{Red} Lectio Quinta.}
\lettrine[lines=2]{\zallmancaps \color{Red} N}{on} revertebantur cum incederent: sed unumquodque ante faciem suam gradiebatur. Quia scilicet ea que reliquerunt nullo iam appetitu respiciunt: et in eternis que appetunt sub contemplationis sue oculis boni operis pedem ponunt. Similitudo autem vultus eorum, facies hominis et facies leonis a dextris ipsorum quatuor. Facies autem bovis a sinistris ipsorum quatuor: et facies aquile desuper ipsorum quatuor. Nam quia ab humana generatione cepit, iure per hominem Mattheus: quia a clamore in deserto, recte per leonem Marcus, quia a sacrificio, recte per vitulum Lucas: quia a divinitate verbi cepit, digne per aquilam signatur Iohannes. Ipse eciam unigenitus Dei totum simul est, qui et nascendo homo, moriendo vitulus, resurgendo leo: et ascendendo ad celos aquila factus est.
\newline {\color{Red} Lectio Sexta.}
\lettrine[lines=2]{\zallmancaps \color{Blue} H}{omo} autem et leo a dextris: vitulus vero a sinistris esse perhibetur. A dextris enim leta: a sinistris vero tristia habemus. Nativitas autem eius et resurrectio, leticiam discipulis prebuit, quos passio contristavit. Iure autem locus aquile non iuxta, sed desuper esse describitur, quia per hoc quod eius ascensionem signat, seu quia verbum patris apud patrem esse denuntiat: super ceteros virtute contemplationis excrevit. Desuper autem ipsorum quatuor esse describitur, quia per hoc quod in principio verbum vidit: eciam super semetipsum transiit. Et facies eorum et penne extente desuper. Quia videlicet omnis intentio omnisque contemplatio sanctorum super se tendit, ut illud possit adipisci quod in celestibus appetit, sive bono operi, sive invigilet contemplationi.
\newline {\color{Red} Lectio VII.}
\lettrine[lines=2]{\zallmancaps \color{Red} D}{ue} penne singulorum iungebantur: et due tegebant corpora eorum. Dictum fuerat, et penne eorum extente desuper, atque mox subiunctum est: due penne singulorum iungebantur. Ubi aperte intelligitur: quia et extendebantur desuper et iungebantur. Due vero tegebant corpora eorum. Quatuor vero esse virtutes invenimus, que a terrenis actibus omne pennatum animal levant, in futuris scilicet amor et spes: de preteritis timor et penitentia. Penne ergo sibimet iuncte superius extenduntur: quia sanctorum mentem amor et spes ad superna levant. Apte quoque coniuncte nominantur, quia amant celestia que sperant, et sperant que amant.
\newline {\color{Red} Lectio VIII.}
\lettrine[lines=2]{\zallmancaps \color{Blue} D}{ue} vero corpora contegunt: quia timor et penitentia ab omnipotentis oculis eorum mala preterita abscondunt. Due vero que extente sursum sunt, dicuntur esse coniuncte, et due que corpus contegunt non dicuntur esse coniuncte, quia per amorem et spem unum est quod desiderant, sed per timorem et penitentiam diversum est quod deplorant. Et unumquodque eorum coram facie sua ambulabat. Dictum superius fuerat unumquodque ante faciem suam gradiebatur: nunc autem dicitur coram facie sua ambulabat. Coram autem quod in presenti est dicimus. Omnis autem iustus, vitam suam sollicitus aspicit, et diligenter considerat quantum cotidie in bonis crescat, aut a bonis decrescat. Quia ergo se ante se ponit: coram se ambulat.
\newline {\color{Red} Lectio IX.}
\lettrine[lines=2]{\zallmancaps \color{Red} U}{bi} erat impetus spiritus: illuc gradiebantur. In electis et reprobis diversi sunt impetus. In electis impetus spiritus: in reprobis impetus carnis. Impetus carnis ad voluptates et huiusmodi: impetus spiritus ad caritatem et virtutes huiusmodi. Quia ergo perfecti se semper virtutibus exercent: recte dicitur. Ubi erat impetus spiritus illuc gradiebantur. Hoc autem quod dictum superius fuerat, ut altius confirmetur, iterum replicatur, cum subiungitur. Non revertebantur cum ambularent. Sequitur. Et similitudo animalium et aspectus eorum quasi carbonum ignis ardentium, et quasi aspectus lampadarum. Aspectus animalium carbonibus ignis ardentibus comparatur. Quisquis enim carbonem ignis tangit: accenditur. Qui autem sancto viro adheret: ex eius assiduitate visionis, usu locutionis, exemplo operis, accipit ut accendatur in amorem veritatis. Lampades vero lucem suam longius spargunt. Quia ergo sancti viri quibusdam longe positis lucent: lampades fiunt.
{\ubsubsection{In Festo Unius Evangeliste Trium Lectionum}{in-festo-unius-evangeliste-trium-lectionum-breviarium}}
\newline {\color{Red} Lectio Prima.}
\lettrine[lines=2]{\zallmancaps \color{Blue} H}{ec} erat visio discurrens, in medio animalium: splendor ignis, et de igne fulgur egrediens. Splendor ignis, et de igne fulgur egrediens in medio animalium discurrens videtur, quia sanctus spiritus qui ignis dicitur universam replens ecclesiam, in electorum cordibus ex se ipso flammas amoris proicit, ut corusci more per terrorem feriat, et ad amorem suum corda torpentia accendat. Homo autem discurrens ubique obviam venit: et repente ubi non creditur: invenitur. Spiritus ergo discurrens dicitur: quia ubique eciam nescientibus occurrit. Et animalia ibant et revertebantur, in similitudinem fulguris coruscantis. Superius dictum est, ibant et non revertebantur, nunc dicitur ibant et revertebantur, quia sancti viri ab activa vita quam apprehendunt ad iniquitatem non corruunt: et a contemplativa vita quam tenere iugiter non possunt, in activam relabuntur.
\newline {\color{Red} Lectio Secunda.}
\lettrine[lines=2]{\zallmancaps \color{Blue} B}{ene} autem revertentia fulguri comparantur: quia cum ad superna contemplanda evolant, bona celestia que saltem per speculum contemplari potuerunt fratribus denuntiant: cordaque audientium feriunt et incendunt. Alio quoque modo vadunt et redeunt. Vadunt enim cum ad insinuandam celestis doni gratiam: in predicationem mittuntur, atque ut ad fidem trahant mira coram infidelibus faciunt, sed redeunt: quia hec omnipotentis virtuti tribuentes, sibimetipsis que fecerint, non ascribunt. Vadunt ergo et revertuntur in similitudinem fulguris coruscantis, quia postquam mira inter homines faciunt: ad dandam actori suo gloriam redeunt, ut illi laudem reddant.
\newline \ul{Crisostomus super Matheum, Omelia prima.} \grecross \ {\color{Red} Lectio Tertia.}
\lettrine[lines=2]{\zallmancaps \color{Red} P}{otest} autem queri, cur quatuor evangeliorum scriptores fuere: cum unus sufficeret cuncta que in eis sunt scripta memorare. Sufficeret utique, sed fit hoc sine dubio veritatis documentum pergrande, quod sunt quatuor qui hec ipsa nec eisdem temporibus, nec eisdem locis conscribant: et tamen omnia quasi ex uno ore pronunciant. Verum econtrario inquis accidit res: in multis enim inter se dissonant plures. At vero hoc ipsum: maximum est veritatis testimonium. Si enim ex toto usque ad singula verba concinerent: inimici veritatis eos ad decipiendum homines communi consilio congregatos scripsisse crederent.
{\ubsubsection{In Festo Unius Martyris, Novem Lectionum}{in-festo-unius-martyris-novem-lectionum-breviarium}}
\newline \ul{Augustinus in sermone de sancto Cipriano,} Gloriosissimi. \grecross \ {\color{Red} Lectio Prima.}
\lettrine[lines=2]{\zallmancaps \color{Blue} H}{odie} natalem martyris ecclesia celebrat. Quid est hoc fratres? Quando sanctus iste natus est ignoramus, et quia hodie passus est: natalem huius hodie celebramus. Sed quando natus sit, eciam si sciremus: nequaquam diem illum celebramus. Illo namque die peccatum originale traxit: isto die peccatum omne finivit. Illo die ex fastidioso matris utero, in hanc lucem processit, que oculos carnis illecebrat: isto autem die ex profundissimo carcere in illam lucem discessit que visum mentis illustrat. Ad preciosam mortem iuste vivendo venit: ad gloriosam vero vitam iniuste moriendo pervenit. Adeptus est triumphale martyris nomen: quia perduxit usque ad sanguinem pro veritate certamen. Cuius ergo in terra triumphum honoramus: in celo meritum cognoscamus.
\newline \ul{Idem in sermone de Genesio,} Cum frequenter. \grecross \ {\color{Red} Lectio Secunda.}
\lettrine[lines=2]{\zallmancaps \color{Red} M}{agna} ex sanctorum natalium nomine, demonstratur spes ac securitas ecclesie. Non enim talem debemus timere mortem: in qua renascimur ad salutem. Sanctus iste mortem corporis non metuit dum anime timuit: nec tyrannum timuit temporarium, timens iudicem sempiternum. Mira fides, mira devotio. Ante persecutorem assistebat, sed animo ante deum erat. Carnem exponebat supplicio: sed mentem habebat in celo. Negociator iste beatus dedit minus ut acciperet amplius, emit commercio sanguinis, regnum salutis: factaque est ei brevis passio perpetua beatitudo. Hinc est ergo quod illius natalem colimus, non quo est in terra natus, sed quo est in celo renatus, non quo pannis est vilibus involutus, sed quo divinis est floribus coronatus: non quo plorans in mundum venit, sed quo e mundo cum gloria triumphavit.
\newline \ul{Idem in sermone alio de sancto Cypriano. \grecross \ } {\color{Red} Lectio Tertia.}
\lettrine[lines=2]{\zallmancaps \color{Blue} D}{uplex} in sancto isto gemma prefulsit, dum vita clara et mors extitit preciosa. Qui antea alia offerebat, modo se ipsum offert. Assolens angelus venire de thuribulo precum eius, ante sedentem in throno offerebat incensum: modo autem levavit ipsum in aula celorum. O dedicata passio: et gloriosa confessio. Absolvitur per martyrum anima a vinculo carnis: et in celo cum triumpho levatur ab angelis. Interrogatus crudeliter respondit sub patientia fortiter. Terrena contempsit, ut inveniret celestia: temporalia reliquit ut possideret eterna. Fidei armis indutus Christi miles in prelio stetit, fortiter dimicavit, pugnavit et vicit, non feriendo sed patiendo: vitamque moriendo quam optabat, accepit a deo eterno.
\newline \ul{Cyprianus in sermone de Celerino. \grecross \ } {\color{Red} Lectio Quarta.}
% https://mlat.uzh.ch/browser/cps_2.CypCar.Episto5#:33
\lettrine[lines=2]{\zallmancaps \color{Red} B}{eatus} iste martyr, dum inexpugnabili firmitate sui certaminis adversarium vicit, vincendi ceteris viam fecit, non brevi vulnerum compendio victor: sed diu adherentibus penis longe colluctationis miraculo triumphator. In carceris custodia septus fuit: sed in vinculis corpore posito, solutus ac liber spiritus mansit. Caro famis ac sitis diuturnitate contabuit: sed animam fide et virtute viventem, nutrimentis spiritalibus Deus pavit. Iacuit inter penas penis suis fortior, inclusus includentibus maior: iacens stantibus celsior, vinctus vincientibus firmior, iudicantibus iudicandus sublimior. Lucent in corpore glorioso clara vulnerum signa: apparent in membris longa tabe consumptis, expressa Christi vestigia. Magna sunt et mira: que de virtutibus eius posset audire fraternitas vestra.
\newline \ul{Sermo Augustini de sancto Cypriano. \grecross \ } {\color{Red} Lectio Quinta.}
\lettrine[lines=2]{\zallmancaps \color{Blue} I}{n} domino laudabitur anima huius martyris: qui aliis fuit odor vite in vitam: aliis odor mortis in mortem. Alii siquidem illum imitando vixerunt: alii vero eidem invidendo perierunt. Ipse vero moriendo vixit: iudicatus iudicem superavit. Adversarium percussus vicit, et mortem occisus occidit. Illi ergo laus, illi gloria qui dignatus est hunc virum inter sanctos suos predestinare ante tempora, inter homines creare oportuno tempore: mundare sordentem, formare credentem, adiuvare pugnantem, coronare vincentem, in quo ostendit quantis malis esset opponenda quantisque bonis preponenda caritas, quam sic illi dilexit, ut et malis pro caritate non parceret, et malos pro pace toleraret: liber in dicendo: et pacificus in audiendo.
\newline \ul{Idem in alio sermone de eodem. \grecross \ } {\color{Red} Lectio Sexta.}
\lettrine[lines=2]{\zallmancaps \color{Red} F}{ramea} Dei est anima iusti in manu Dei: de qua in Psalmo dicitur. Erue animam meam ab impiis: frameam tuam ab inimicis. Quod enim dixit animam meam, hoc repetit frameam tuam. Hanc frameam effudit dominus: dum martyres suos usquequaque sparsit, et conclusit adversus eos qui ecclesiam persequebantur, ut qui predicantium vocibus non flectebantur: morientium virtutibus frangerentur. Fortia quippe arma sibi fabricat adversus inimicos: eos ipsos quos facit amicos. Magna itaque framea Dei est anima martyris, splendida caritate, acuta veritate, Dei pugnantis vibrata virtute: que bella fecit, quae contradicentium catervas redarguendo superavit, percussit infensos: prostravit aversos.
\newline {\color{Red} Lectio VII.}
\lettrine[lines=2]{\zallmancaps \color{Blue} H}{ec} framea in multorum cordibus inimicicias quibus oppugnabatur, occidit: eosque amicos quibus adversus alios Deus copiosius pugnaret, effecit. Ubi autem venit tempus ut tanquam prevalentibus hostibus prenderetur, tunc ne oppressus et victus ab impiis cederet: affuit ille per quem invictus prestaretur. Suscepit victoriam postquam nullum certamen ulterius remaneret: quam scilicet de hoc mundo et de huius mundi principe reportaret. Affuit omnino Christus fidelissimo suo testi usque ad mortem pro veritate certanti: fecit quod ille rogavit, dum animam eius ab impiis frameam suam ab inimicis eruit, cuius victricis anime caro sancta tanquam illius framee vagina ab hominibus honoratur in terra: eiusdem scilicet anime triumphali resurrectione reddenda, nullaque deinceps morte ponenda.
\newline \ul{Origenes super iudicum, Omelia septima.} \grecross \ {\color{Red} Lectio VIII.}
\lettrine[lines=2]{\zallmancaps \color{Red} Q}{uis} karissimi martyris huius animam sequi possit, que supergressa omnes aerias potestates, ad altare celeste tendit? Beata certe anima illa, que sibi occurrentes aerias demonum turmas, sanguinis in martyrio profusi cruore deturbat: que sic sequitur Christum: quomodo Christus eam precesserat.
\newline \ul{Sermo beati Augustini de sancto Cypriano. \grecross}
\newline Celebratio sollemnitatis martyrum: imitatio debet esse virtutum. Facile siquidem est martyrem celebrando venerari: magnum vero, fidem eius et patientiam imitari. Hoc ergo sic agamus: ut illud optemus. Hoc ita celebremus: ut illud potius diligamus, quod laudamus in martyris fide: quod usque ad mortem certavit pro veritate. Vicit autem: quia blandientem mundum contempsit,
%pp137
sevienti non cessit: et ob hoc ad Deum victor accessit.
\newline {\color{Red} Lectio IX.}
\lettrine[lines=2]{\zallmancaps \color{Blue} B}{eatus} hic martyr, vicit errores sapientia: terrores patientia. Agnum secutus: leonem vicit. Quando persecutor seviebat: leo fremebat. Sed deorsum leo calcabatur: quia sursum agnus attendebatur: qui in ligno pependit, sanguinem fudit: mundum redemit: mortemque sua morte destruxit. Contempnite ergo seculum Christiani: sicut contempsit martyr iste. Quid obsecro tantum amatis quod ille contempsit quem sic honoratis, et nisi contempsisset, non utique honoraretis? Tu ergo noli hoc amare, sed martyris exemplo contempne. Non enim intravit: et ostium contra te clausit. Patet ianua Christus, qua ingrediaris. Noli amare impedimentum: si non vis invenire tormentum. Quod amas in terra est impedimentum: viscus est pennarum spiritualium, quibus volatur ad Deum. Capi non vis, et viscum diligis? Nunquid non ideo caperis, quia dulciter caperis? Certe quanto magis delectat: tanto fortius strangulat.
{\ubsubsection{In Festo Unius Martiris, Trium Lectionum}{in-festo-unius-martiris-trium-lectionum-breviarium}}
\newline \ul{Maximus in sermone, Ad exhibendum. \grecross \ } {\color{Red} Lectio Prima.}
% https://mlat.uzh.ch/browser/7648:2.8
\lettrine[lines=2]{\zallmancaps \color{Red} B}{eatus} hic martyr, virtute caritatis indutus, bonum certamen certavit, cursum consummavit: fidem servavit, passus superavit: occisus triumphavit. Aliud est enim occidi pro mundo: aliud occidi pro Christo. Qui occiditur pro mundo, vincitur cum occiditur: quia nisi vinceretur non occideretur. Qui autem occiditur pro Christo victus non dicitur: quia cum occiditur victor efficitur. Corona quippe victoribus datur. Quomodo ergo occisus ille victor negatur, qui pro eo quod occiditur a iusto iudice coronatur? Ecce beatus hic martyr permansit victor occisus: quia pro invicto invenitur occisus, quia si occidi pro invicto metueret, ad victoriam non veniret: immo si non occideretur miserabiliter vinceretur. Vicit ergo pro fide moriendo: qui vinceretur sine fide vivendo.
\newline {\color{Red} Lectio Secunda.}
\lettrine[lines=2]{\zallmancaps \color{Blue} S}{i} enim amissa fide martyr hic vivere delectaretur: in anima moreretur. Corpus enim ex anima vivit, anima autem ex fide vivit: Christus autem in cordibus fidelium per fidem inhabitare dignatur. Discedente itaque fide ab homine, discedit Christus, discedente Christo discedit vita: quia Christus est vita nostra. Quid ergo proderit quia videtur in corpore vivus, qui fuerit in anima morte infidelitatis occisus? Iste beatus martyr plenus spiritu sancto cogitavit: et ut semper in anima viveret mortem corporis nullatenus formidavit. Noverat quidem anima non moriente: corpus cum ea in eterna felicitate victurum.
\newline {\color{Red} Lectio Tertia.}
\lettrine[lines=2]{\zallmancaps \color{Red} H}{inc} est quod ad subeundam corporis mortem indubitanter venit: quia in ipsa temporali morte eternam vitam anime simul et carnis invenit. Non est ergo victus occisus cui mors corporis in tantum facta est triumphalis: ut per illam fieret immortalis. Traditur itaque iudici: deputatur cruciatibus occidendus. Sed quia fervet ardore caritatis: exhibuit corpus suum hostiam viventem, sanctam, Deo placentem, et caritas holocaustum exhibuit, quem persecutoris crudelitas laceravit: et impietas iudicis interfecit, quem holocausti hostiam, passio gloriosa perfecit. Omnes igitur in hac festivitate letemur: et festivitatis gaudium cunctis proficiat ad virtutis exemplum.
{\ubsubsection{In Festo Plurimorum Martirum Novem Lectionum}{in-festo-plurimorum-martirum-novem-lectionum-breviarium}}
\newline \ul{Crisostomus in Omelia de Machabeis. \grecross \ } {\color{Red} Lectio Prima.}
% https://mlat.uzh.ch/browser/cps_2.JonAur.DeCuIm#:2;18
\lettrine[lines=2]{\zallmancaps \color{Blue} Q}{uam} speciosa festivitas hodie nobis enituit, et splendida dies hodierna prefulsit: non de usitatis solis radiis illuminata, sed de lumine martyrum illustrata. Hi namque radii Christo emittente iaciuntur in terra de superno circulo solis. Et sicut speciem ecclesie gaudentis illustrant: ita diaboli faciem invidentis obcecant. Nam ut audaces latronum duces, agnitis regis insignibus deterrentur ut fugiant, ita profecto demones ubi coronatorum martyrum corpora posita conspiciunt: ilico ab eorum longe conspectu pavidi fugiunt. Horum itaque corpora preciosa sunt: quia plagas pro domino susceperunt, et stigmata propter Christum impressa circumferunt. Et sicut ex corona regali undique fulgores emittuntur: ita martyrum corpora pro Christo vulneribus distincta, omni regum diademate delectabiliora redduntur.
\newline \ul{Cyprianus in epistola ad martyres et confessores. \grecross \ } {\color{Red} Lectio Secunda.}
\lettrine[lines=2]{\zallmancaps \color{Red} O}{}\ beati martyres quibus vos laudibus predicem, o milites fortissimi robur corporis vestri quo preconio vocis explicem? Usque ad glorie consummationem: tolerastis durissimam questionem. Nec cessistis suppliciis: sed potius supplicia cesserunt vobis. Finem doloribus quem non dabant tormenta dedit corona. Persecutio gravior ad hoc diu perseveravit, non ut stantem fidem deiceret: sed ut homines dei ad deum velocius mitteret. Vidit admirans presentium multitudo celeste certamen, et in prelio stetisse servos Christi, voce libera, mente incorrupta, virtute divina: telis quidem secularibus nudos, sed armis fidei armatos. Torti denique torquentibus fortiores steterunt: ac laniantes ungulas, laniata membra vicerunt. Fluebat sanguis qui persecutionis incendium extingueret: flammas Gehenne sopiret.
\newline {\color{Red} Lectio Tertia.}
\lettrine[lines=2]{\zallmancaps \color{Blue} O}{}\ quale deo fuit illud spectaculum, quam sublime quam magnum, quam preciosum et acceptum? Quam letus illic Christus affuit, quam libens in talibus pugnavit, et vicit preliatores, et assertores sui nominis erexit, corroboravit, animavit? Nam qui mortem semel vicit pro nobis: semper vincit in nobis. Hic est agon fidei nostre qua congredimur, qua vincimus, qua coronamur, quem per prophetas predictum, per dominum commissum, per apostolos gestum, sancti martyres exhibuerunt et palmam acceperunt: in fide stabiles, in dolore patientes, in questione victores. O martyres sancti: membra vestra felicia vinculis infamibus ligaverunt: sed ornamenta sunt ista non vincula. O pedes feliciter iuncti, qui non a fabro sed a domino resolvuntur, et itinere salutari ad paradisum diriguntur, ad presens in seculo ligati: et celeriter ad Christum itinere glorioso cursuri.
\newline \ul{Idem in epistola de lapsis.} \grecross \ {\color{Red} Lectio Quarta.}
\lettrine[lines=2]{\zallmancaps \color{Red} R}{epugnastis} o martyres sancti fortiter seculo, spectaculum gloriosum prebuistis deo: secuturis fratribus fuistis exemplo. Religiosa vox Christum locuta est: illustres manus sacrificiis sacrilegis restiterunt. Ora sanctificata celestibus cibis, post corpus domini ydolorum reliquias respuerunt, frons cum signo dei, diaboli coronam ferre non potuit: sed corone se domini reservavit. Quam leto sinu vos mater ecclesia de prelio revertentes excepit: quam beata quam gaudens portas suas aperuit, ut intraretis de hoste prostrato trophea referentes. Multitudo vestram gloriam sequitur: et vestra vestigia comitatur. Eadem in illis cordis sinceritas, eadem fidei integritas: et celestium preceptorum radicibus nisos, et evangelicis rationibus roboratos, non exilia prescripta, non destinata tormenta, non rei familiaris aut corporis terruerunt supplicia.
\newline \ul{Beda in Omelia de Innocentibus.} \grecross \ {\color{Red} Lectio Quinta.}
\lettrine[lines=2]{\zallmancaps \color{Blue} H}{i} sunt quos cum paterentur lacrimis prosecuta est ecclesia: ipsa que genuit, sed mox hinc expulsos superna Hierusalem obviis leticie manibus excepit: et in gaudiis domini sui perpetuo coronandos introduxit. Stant enim modo ante thronum Dei coronati: qui quondam ante thronum iudicum terrenorum iacebant penis attriti. Stant in conspectu agni nulla ratione a contemplanda eius gloria separandi: a cuius hic amore nec per supplicia potuerunt separari. Stolis albis fulgent et palmas in manibus habent, quia premia in operibus habent: dum corpora que pro domino dissipari passi sunt, per resurrectionem glorificata recipiunt. Magna autem voce salutem Deo canunt: quia cum magna gratiarum actione recolunt, non sua se virtute, sed ipso auxiliante superasse.
\newline \ul{Gregorius de exhortatione ad martyrium. \grecross \ } {\color{Red} Lectio Sexta.}
% https://mlat.uzh.ch/browser/cps_2.CypCar.EpDeExM#:3.13;2
\lettrine[lines=2]{\zallmancaps \color{Red} Q}{uis} non omnibus viribus elaboret ut post terrena tormenta ad premia divina perveniat? statimque cum Christo gaudeat? Si militibus gloriosum est ut redeant in patriam victores, quanto maior est gloria ad paradisum redire triumphantes? Gloriosum prorsus unde Adam eiectus est, illic hoste prostrato qui eum eiecerat, trophea victricia reportare, et acceptissimum munus Deo, fidem scilicet incorruptam offerre: cumque venerit vindictam recepturus eundem comitari, cum sederit iudicaturus, eius assistere lateri: Christi coheredem fieri, angelis coequari: cum patriarchis, et prophetis, et apostolis, regni celestis possessione letari. Huiuscemodi meditationibus fundata mens, durat in persecutione. Clauduntur oculi terre: sed patet celum. Minatur Antichristus: sed tuetur Christus. Mors infertur: sed immortalitas sequitur. Mundus eripitur: sed paradisus exhibetur. Vita temporalis extinguitur, sed eterna reparatur.
\newline \ul{Iohannes episcopus in sermone de martyribus,} Qui sanctorum. \grecross \ {\color{Red} Lectio VII.}
\lettrine[lines=2]{\zallmancaps \color{Blue} Q}{ui} sanctorum glorias, frequenti laude colloquitur: eorum mores sanctos imitetur. Aut enim imitari debet: aut laudare non debet. Propterea se ipsos nobis prebent exemplum: ut ex ipsis proficiamus. Delicatus es tu Christiane miles, si putas te posse sine certamine triumphare. Exere vires, dimica atrociter: in prelio concerta. Considera pactum, condictionem attende: militiam nosce, pactum quod spopondisti, condictionem qua accessisti: militiam cuius nomen dedisti. Denique considera te, presente Deo cum hoste pugnare. Favet ille ut vincas, favet ut obtineas: cum pugnas presens est. Quantum virium de eius presentia concipis: tantum inimico debilitatis accedit. Tu erigeris in virtutem: hostis incidit in debilitatem. Tibi arma divinitus ministrantur: illi malitie tela franguntur.
\newline {\color{Red} Lectio VIII.}
\lettrine[lines=2]{\zallmancaps \color{Red} T}{ibi} conspectus Dei virium adhibet incrementa: hosti adimit nocendi venena. Tibi angeli plaudunt: illi formidinem adhibent. Tibi fortitudo tribuitur: hostis malitia infirmatur. In tua pugna dominus congreditur: et victoria tibi ascribitur. Quid formidas, quasi tua virtute dimicans? Apprehende arma, procede in bellum: fortiter dimica. Non potest imperatorem promereri: qui hostem a se noluerit superari. Nec hostem poterit vincere: nisi qui atrociter voluerit dimicare. Victoria militis peremptio est hostis: et peremptio hostis gloria est imperatoris. Denique in certamine Christiano, aut superstes miles hoste prostrato, beatus de prelio redit: aut hostem victor moriendo devincit.
\newline {\color{Red} Lectio IX.}
\lettrine[lines=2]{\zallmancaps \color{Blue} F}{elicter} quidem vincit: qui post victoriam vinci non novit. Feliciter vincit: qui et diabolum et mundum post victoriam dereliquit. Feliciter vincit qui proficiscens e mundo, diabolum quem in presenti devicit, in futuro cum domino iudicabit. Alios post unum certamen ad alias palmas dominus servat: alios consummato martyrio iam coronat. Alios victores retinet in exemplum: alios iam perfectos transmittit ad celum. Alios frequenter desiderat videre certantes, alios passione iam meritos imponit regnis celestibus triumphantes: ut imperatoris Christi ab omnibus dignatio collaudetur. Sic dimicaverunt quos diligere Christiane consuevisti: sic preliati sunt quos semper amasti.
{\ubsubsection{In Festo Plurimorum Martirum, Trium Lectionum}{in-festo-plurimorum-martirum-trium-lectionum-breviarium}}
\newline \ul{Sermo beati Maximi. \grecross \ } {\color{Red} Lectio Prima.}
% https://mlat.uzh.ch/browser/cps_2.AuInMaT.Sermon11#:2.5;4
\lettrine[lines=2]{\zallmancaps \color{Red} R}{ecte} beatorum martyrum celebrantur insignia, qui illuminantes virtute fidei sue innumerabiles populos, soli quidem illatam sibi senserunt mortem: sed non soli sue mortis beneficia possederunt. Immensa enim Dei pietas multiplex ad bonitatem, et artifex ad salutem, illorum merita: nostra vult esse suffragia. Dumque nobis fidem veram duro martyrii agone commendat: afflictionem precedentium, instructionem efficit posterorum. Illos exanimat, ut nos erudiat, illos conterit, ut nos adquirat: eorumque cruciatus nostros contulit esse profectus. Illi periculis suis pugnant: et nostris utilitatibus militant. Magni itaque periculi res est si post prophetarum oracula, post apostolorum testimonia, post martyrum vulnera veterem fidem quasi novellam discutere presumas: et post morientium sudores, ociosa disputatione de veritate religionis contendas.
\newline {\color{Red} Lectio Secunda.}
\lettrine[lines=2]{\zallmancaps \color{Blue} V}{eneremur} ergo in sanctorum martyrum gloria fidem nostram, congratulemur dilectissimi magne fidei nostre: per quam dum exules proflua caritate suscipitis, ipsos intercessores habere mereamini, quorum et si seminetur in cineribus portio: tamen manet integra in virtutibus plenitudo. Mater martyrii fides catholica est: in qua isti illustres athlete, sanguine subscripserunt, quam morte sua iuverunt. Ab illis enim dum mortis tolerantia indubitanter recipitur, spes immortalitatis evidenter asseritur. In sanctis itaque preconibus veritatis, et passionis magnanimitas, et resurrectionis autoritas est. Puniti enim pereunt in melius reparandi: quibus per angustias tribulationum, aperitur exitus ad plenitudinem gaudiorum. Cruentas impio ex ore sententias respuunt manu, repudiant auditu: refugiunt officio. Itaque impetum passionis magnanimitate provocant: humilitate declinant.
\newline {\color{Red} Lectio Tertia.}
% https://mlat.uzh.ch/browser/cps_2.MaxTau.Sermon10#:88.1;4
\lettrine[lines=2]{\zallmancaps \color{Red} I}{ncumbente} persecutionis furore: adhuc cathecumeni in celo animum preparant, regenerantur morte, damnatione salvantur: gladio consecrantur. Nos post pie regenerationis candorem in peccatorum ceno revoluti, innovata fedamus: repetentes opera non iam undis purganda sed flammis. Felix baptismus sanguinis, quo omne delictum ita exiit: ut iam non possit redire. Honoremus karissimi beatos martyres, principes fidei, intercessores mundi: precones regni, coheredes Dei. Quod si dicas michi. Quid honoras in carne iam consumpta: honoro in carne martyris exceptas pro Christi nomine cicatrices, honoro pro confessione domini sacratos cineres, honoro in cineribus semina eternitatis. Honoro corpus quod michi Deum meum ostendit diligere: quod me propter dominum mortem docuit non timere, quod reverentur et demones: quod afflixerunt in supplicio: sed glorificant in sepulcro. Honoro itaque corpus quod Christum honoravit in gladio: quod et cum Christo regnabit in celo.
{\ubsubsection{In Festo Unius Confessoris Pontificis, Novem Lectionum}{in-festo-unius-confessoris-pontificis-novem-lectionum-breviarium}}
\newline {\color{Red} Lectio Prima.}
\lettrine[lines=2]{\zallmancaps \color{Blue} C}{elebrantes} karissimi sancti huius pontificis sollemnia, qui eo quod digne prefuit, honore duplici in celo scilicet et in terra dignus est habitus: de pastoralis cura regiminis quod traditum tenemus a sanctis pastoribus ad memoriam revocemus.
\newline \ul{Gregorius in pastorali. \grecross \ }
\newline Pensandum quippe valde est, ad culmen quisque regiminis qualiter veniat: atque ad hoc rite perveniens qualiter vivat, et bene vivens qualiter doceat: et recte docens infirmitatem suam cotidie quanta consideratione cognoscat. Nulla ars doceri presumitur, nisi intenta prius meditatione discatur. Ab imperitis ergo pastorale magisterium qua temeritate suscipitur, quando ars est arcium regimen animarum? Quis autem cogitationum vulnera occultiora esse nesciat vulneribus viscerum? Et tamen sepe qui nequaquam spiritalia precepta cognoverunt, cordis se medicos profiteri non metuunt: dum qui pigmentorum vim nesciunt, videri medici carnis erubescunt.
\newline {\color{Red} Lectio Secunda.}
\lettrine[lines=2]{\zallmancaps \color{Red} S}{unt} eciam nonnulli, qui per speciem regiminis gloriam affectant honoris, transcendere ceteros concupiscunt: atque ideo susceptum cure pastoralis officium ministrare digne tanto magis nequeunt, quanto ad humilitatis magisterium ex sola elatione pervenerunt. Sunt autem nonnulli, qui solerti cura spiritalia precepta perscrutantur: sed que intelligendo penetrant, vivendo conculcant. Repente docent que non opere sed meditatione didicerant: et quod verbis predicant, moribus impugnant. Unde fit, ut cum pastor per abrupta graditur: ad precipicium grex sequatur. Nemo quippe amplius in ecclesia nocet: quam qui perverse agens, nomen vel ordinem sanctitatis habet. Delinquentem namque eum redarguere nullus presumit, et in exemplum culpa vehementer extenditur: quando pro reverentia ordinis peccator honoratur.
\newline {\color{Red} Lectio Tertia.}
\lettrine[lines=2]{\zallmancaps \color{Blue} P}{lerumque} quoque si ad regiminis culmen quis eruperit, in elationem protinus permutatur. Sic Saul qui indignum se prius considerans fugerat: mox ut regni gubernacula suscepit intumuit. Sepe quoque suscepta cura regiminis, cor per diversa diverberat, et impar quisque invenitur ad singula: dum confusa mente dividitur ad multa. Cumque foris mens per insolentem curam trahitur: a timoris intimi soliditate vacuatur. Fit in exteriorum dispositione sollicita: et sui solummodo ignara, scit multa cogitare se nesciens. Nam cum plusquam necesse est se exterioribus implicat: quasi occupata in itinere obliviscitur quo tendebat, ita ut a studio sue inquisitionis aliena, nec ipsa que patitur damna consideret: et per quanta delinquat ignoret. Hec itaque proferentes, non potestatem reprehendimus, sed ab appetitu illius cordis infirmitatem munimus: ne imperfecti quique, qui in planis stantes titubant: in precipiti pedem ponant.
\newline {\color{Red} Lectio Quarta.}
\lettrine[lines=2]{\zallmancaps \color{Red} S}{unt} eciam nonnulli qui pro exercitatione ceterorum, magnis muneribus exaltantur, qui culmen regiminis, si vocati suscipere renuunt: ipsa sibi plerumque dona adimunt, que non pro se tantummodo, sed eciam pro aliis acceperunt. Dum enim solius contemplationis studiis inardescunt: parere utilitati proximorum in predicatione refugiunt. Secretum quietis diligunt: secessum speculationis petunt. De quo si districte iudicentur, ex tantis procul dubio rei sunt: quantis venientes ad publicum prodesse potuerunt. Sunt autem nonnulli, qui ex sola humilitate refugiunt: ne eis quibus impares se estimant
%pp138
preferantur. Quorum profecto humilitas vera est: cum ad refugiendum hoc quod utiliter subire precipitur, pertinax non est. Cum enim regiminis culmen imperatur, si iam donis preventus est quibus aliis prosit: debet quisque et ex corde fugere, et invitus obedire.
\newline {\color{Red} Lectio Quinta.}
\lettrine[lines=2]{\zallmancaps \color{Blue} N}{onnunquam} quoque predicationis officium, et nonnulli laudabiliter appetunt: et ad hoc nonnulli coacti laudabiliter pertrahuntur. Ysaias quippe domino querente quem mitteret: ultro se obtulit dicens. Ecce ego, mitte me. Hieremias autem mittitur: et tamen ne mitti debeat humiliter reluctatur. Sed in utrisque subtiliter est intuendum: quia et is qui recusavit plene non restitit, et is qui mitti voluit, ante se per altaris calculum purgatum vidit, ne aut non purgatus adire quisque sacra misteria audeat, aut quem superna gratia elegit: sub humilitatis specie superbe contradicat. Quod Moyses utrumque miro opere explevit: qui preesse tante multitudini et noluit et obedivit. Superbus enim fortasse esset: si ducatum plebis innumere sine trepidatione susciperet, et rursum superbus existeret: si actoris imperio obedire recusaret.
\newline {\color{Red} Lectio Sexta.}
\lettrine[lines=2]{\zallmancaps \color{Red} P}{lerumque} vero qui preesse concupiscunt: ad usum sue libidinis instrumentum apostolici sermonis arripiunt, quo ait. Siquis episcopatum desiderat: bonum opus desiderat. Qui tamen laudans desiderium in pavorem vertit protinus quod laudavit: cum repente subiungit. Oportet autem episcopum irreprehensibilem esse. Notandum quoque quod illo in tempore dicitur: quo quisquis plebibus preerat, prius ad martyrii tormenta ducebatur. Tunc ergo fuit laudabile episcopatum querere: quando per hunc dubium non erat ad supplicia graviora pervenire. Unde ipsum episcopatus officium, boni operis expressione diffinitur cum dicitur. Siquis episcopatum desiderat bonum opus desiderat. Ipse ergo sibi testis est, quia episcopatum non appetit: qui non per hunc boni operis ministerium, sed honoris gloriam querit, qui subiectione ceterorum pascitur, laude propria letatur, ad honorem cor elevat, rerum affluentium abundantia exultat.
\newline {\color{Red} Lectio VII.}
\lettrine[lines=2]{\zallmancaps \color{Blue} P}{lerumque} eciam hi qui subire pastorale magisterium cupiunt: nonnulla bona opera animo proponunt. Et quamvis hoc elationis intentione appetant: operaturos se tamen magna pertractant. Fitque ut aliud intentio supprimat: aliud superficies cogitationis ostendat. Nam sepe mens sibi mentitur, et fingit de bono opere amare quod non amat: de mundi autem gloria non amare quod amat. Cumque percepti principatus officio perfrui ceperit: obliviscitur quicquid religiose cogitavit. Virtutibus igitur pollens, coactus ad regimen veniat: ne acceptam pecuniam in sudario ligans de eius occultatione iudicetur. At contra, qui regimen appetit: attendat ne per exemplum pravi operis ad ingressum regni tendentibus obstaculum fiat. Considerandum quoque est: quia cum causam populi presul suscipit, quasi ad egrum medicus accedit. Qua ergo presumptione percussum mederi properat, qui in facie vulnus portat.%typo: ostaculum
\newline {\color{Red} Lectio VIII.}
\lettrine[lines=2]{\zallmancaps \color{Red} I}{lle} igitur ille modis omnibus debet ad exemplum vivendi pertrahi: qui cunctis passionibus carnis moriens iam spiritaliter vivit. Qui prospera mundi postposuit: qui nulla adversa pertimescit. Qui sola eterna desiderat, cuius intentioni nec omnino per imbecillitatem deficiens corpus: nec valde per contumeliam repugnat spiritus. Qui ad aliena cupienda non ducitur: sed propria largitur. Qui ad ignoscendum citius flectitur: sed nunquam plusquam deceat ignoscens, ab arce rectitudinis inclinatur. Qui nulla illicita perpetrat: sed perpetrata ab aliis deplorat. Qui ex affectu aliene infirmitati compatitur: sicque in bonis proximi sicut in suis letatur. Qui ita se imitabilem ceteris in cunctis que agit insinuat: ut inter eos non habeat quod saltem de transactis erubescat. Qui sic studet vivere: ut proximorum corda arentia doctrine valeat fluentis irrigare.
\newline {\color{Red} Lectio IX.}
\lettrine[lines=2]{\zallmancaps \color{Blue} I}{lle} debet ad statum huiusmodi pertrahi: qui orationis experimento iam didicit, quod obtinere a domino que poposcerit possit. Cui per effectus vocem specialiter dicitur. Adhuc loquente te: dicam. Ecce assum. Qua enim mente apud Deum intercessionis locum pro populo arripit, qui familiarem eius se esse nescit? Tantum quoque debet actionem populi actio transcendere presulis: quantum distare solet a grege vita pastoris. Sit ergo necesse est rector cogitatione mundus: actione precipuus. Discretus in silentio: utilis in verbo. Singulis compassione proximus: pre cunctis contemplatione suspensus. Bene agentibus per humilitatem socius: contra delinquentium vitia per zelum iusticie erectus. Internorum curam exteriorum occupatione non minuens: exteriorum providentiam internorum sollicitudine non relinquens.
{\ubsubsection{In Festo Unius Confessoris Pontificis Trium Lectionum}{in-festo-unius-confessoris-pontificis-trium-lectionum-breviarium}}
\newline \ul{Sermo beati Ambrosii secundum quosdam libros, vel Maximi secundum alios. \grecross \ } {\color{Red} Lectio Prima.}
% https://mlat.uzh.ch/browser/cps_2.AucVar.SeSAmHa#:56.3;2
\lettrine[lines=2]{\zallmancaps \color{Red} A}{d} sancti ac beatissimi patris, cuius hodie festa celebramus laudes addidisse: aliquid decerpsisse est. Quicquid enim in sancta plebe eius virtutis fuit et gratie: de hoc quasi de quodam fonte lucidissimo omnium rivulorum hec puritas emanavit. Etenim qui castitatis pollebat vigore: qui abstinentie gloriabatur angustiis, qui blandimentis erat preditus lenitatis: omnium civium in Deum provocavit affectum. Qui pontificis administratione fulgebat: plures ex discipulis reliquit sui sacerdotii successores. Habuimus in eo medicum qui e celo lapsus est: cum quo sanitas fines nostros intravit. %typo: Quiquid
% https://mlat.uzh.ch/browser/cps_2.AucInc.SeInNaS2#:1;5
Quantis hic cecis a via veritatis errantibus, et de summa iam in profundum rupe pendentibus, amissum reddidit visum, et illum quo Christus videretur reparavit intuitum?
\newline {\color{Red} Lectio Secunda.}
\lettrine[lines=2]{\zallmancaps \color{Blue} Q}{uantorum} hic auribus surdis infidelitatis obturatione ad percipiendam vocem celestium mandatorum preciosum infudit auditum, ut vocanti Deo per misericordiam, responderent per obedientiam? Quantos iste intrinsecus vulneratos angelici oris arte et orationum firmitate curavit? Quantos per longam incuriam peccati labe solutos, et quadam lepre contagione perfusos, castigationibus et exhortationibus expiando mundavit? Quantorum animas viventium in corpore iam defunctas, et delictorum mole sepultas, Deo resuscitavit? De quantorum cordibus fugavit luxuriam, deiecit iram, et extinxit invidiam, et velut barbaris hospitibus perturbatis, illic fidem, iusticiam, castitatem, misericordiam, habitatores pacificos intromisit?
\newline {\color{Red} Lectio Tertia.}
\lettrine[lines=2]{\zallmancaps \color{Red} S}{umme} alacritatis instantia fuit vir iste cultor et custos animarum, cupiditates resecans iras comprimens, malicias extirpans. Ideo autem Deus illum nobis providit: ut ad actuum illius informemur exemplum. Sicut enim paterfamilias splendidissime cupiens depingere domum suam, perquirit quasque ymagines quas ante studiosi pictoris proponat aspectum, ut ille docta manu transferat in parietem opus alienum: ita Deus noster ad exornandas animas nostras, videtur quedam vivendi exemplaria viri huius gesta scilicet et merita demonstrasse in cordibus nostris. Secundum hoc ergo: animarum nostrarum vultibus viri huius virtutum superducamus colores. Rapiat sibi de eo, alter in gratia speciem compunctionis, alter candorem nitide castitatis: ille pallorem abstinentie, iste ruborem verecundie, ille ardentem fidem velut auri splendore radiantem: atque hoc modo ad unam illius faciem omnem exornemus domum, ut videns nos de futura beatitudine cogitantes pater iste: nos in sinum suum velut peculiares accipiat in illo die, et oves suas letus pastor agnoscat.%typo: in gratia
{\ubsubsection{In Festo Unius Confessoris Non Pontificis Novem Lectionum}{in-festo-unius-confessoris-non-pontificis-novem-lectionum-breviarium}}
\newline \ul{Ex sermone predicto.} \grecross \ {\color{Red} Lectio Prima.}
% https://mlat.uzh.ch/browser/cps_2.AucInc.SeInNaS2#:1;2
\lettrine[lines=2]{\zallmancaps \color{Blue} B}{ene} et congrue in hac die quam nobis beati viri huius ad paradisum transitus exultabilem reddit: hunc versiculum decantamus. In memoria eterna erit iustus. Digne enim in memoriam vertitur hominum qui ad gaudium transiit angelorum, qui iam gratiam Christi clarificatus invenit, quia mundi gloriam non quesivit: cavens illud quod dicit sermo divinus. Ne laudaveris hominem in vita sua. Tanquam si diceret. Lauda post vitam: magnifica post consummationem. Duplici enim ex causa utilius est hominum magis memorie dare laudem quam vite, ut illo potissimum tempore merita sanctitatis extollas: quando nec laudantem adulatio nocet, nec laudatum temptet elatio. Lauda ergo post periculum, predica securum. Lauda navigantis felicitatem: sed cum venerit ad portum. Lauda ducis virtutem: sed cum perductus est ad triumphum.
\newline {\color{Red} Lectio Secunda.}
\lettrine[lines=2]{\zallmancaps \color{Red} Q}{uis} autem vivens tute possit laudari, qui et de preterito meminit se habere quod deleat, et de futuro videt sibi superesse quod timeat? Quis vero in hoc corpusculo positus debeat sibi quicquam vendicare de meritis, quorum iter vite occulta demonum infestatio, velut latronum obsidet multitudo? Navigantibus nobis per hoc mare magnum, valde pertimescendum est, ne navem nostram aut procella tempestatis arripiat, aut fluctus absorbeat: aut in eternam predam pyrata crudelis abducat. Adversus ergo hec multiformia mala: castigatis et castis actibus resistamus. Beati autem viri huius merita iam in tuto posita, securi magnificemus: qui gubernaculum fidei viriliter tenens, anchoram spei tranquilla iam in statione composuit, et celestibus divitiis onustam navem optato in litore collocavit: qui contra omnes adversarios scutum timoris Dei tandiu infatigabiliter tenuit, donec ad victoriam perveniret.
\newline \ul{Gregorius in moralibus libro decimo. \grecross \ } {\color{Red} Lectio Tertia.}
\lettrine[lines=2]{\zallmancaps \color{Blue} F}{uit} iste de numero iustorum, de quibus dicitur. Deridetur iusti simplicitas. animus quippe viri iusti nonnunquam operi recto constanter innititur: et tamen humanis irrisionibus urgetur. Miranda agit, et obprobria recipit: quia ab huius mundi sapientibus, puritatis virtus fatuitas creditur. Quid namque stultius mundo videtur, quam mentem verbis ostendere, nil callida machinatione simulare, nullas iniuriis contumelias reddere, pro maledicentibus orare, paupertatem querere, possessa relinquere, rapienti non resistere: percutienti unam maxillam, alteram prebere? Sed quod abhominantur Egyptii, hoc Israhelite Deo offerunt, quia simplicitatem ac mansuetudinem Deo immolant: quam abhominantes reprobi: fatuitatem putant.
\newline {\color{Red} Lectio Quarta.}
\lettrine[lines=2]{\zallmancaps \color{Red} H}{ec} iusti simplicitas breviter, et sufficienter exprimitur: cum protinus subinfertur. Lampas contempta apud cogitationes divitum. Quid hoc loco signatur nomine divitum, nisi elatio superborum? Qui cum vitam simplicium in hoc mundo humilem abiectamque despiciunt: elatis protinus despectibus illudunt. Sepe namque contingit, ut electus quisque, quia ad eternam felicitatem ducitur, continua hic adversitate deprimatur: non hunc rerum abundantia fulciat, non dignitatum honorabilem gloria reddat, a cunctis despicabilis cernitur: et huius mundi gratia indignus estimatur. Sed tamen ante occulti iudicis oculos virtutibus emicat, honorari metuit, despici non refugit, corpus continentia afficit: sola in animo dilectione pinguescit. Bene itaque, et lampas esse dicitur et contempta: lampas quia interius lucet, contempta: quia exterius non lucet.
\newline {\color{Red} Lectio Quinta.}
\lettrine[lines=2]{\zallmancaps \color{Blue} S}{ed} usquequo lampas ista contempnitur? Usquequo despicabilis habetur? Nunquid nam fulgorem suum nullatenus exerit, et quanta claritate candeat nunquam ostendit? Ostendit plane. Nam cum lampas contempta apud cogitationes divitum diceretur: protinus additur. Parata ad tempus statutum. Statutum quippe contempte lampadis tempus: est extremi iudicii predestinatus dies, quo iustus quisque non despicitur, sed quanta potestate fulgeat demonstratur. Tunc enim cum Deo iudices veniunt: qui nunc iniuste pro Deo iudicantur. Tunc eorum lux tanto latius emicat: quanto nunc eos manus persequentium durius angustat. Tunc reproborum oculis patescit, quod celesti potestate subnixi sunt: qui terrena omnia sponte reliquerunt. Quot enim nunc pro amore veritatis sese libenter humiliant: tot tunc in iudicio lampades coruscant.
\newline \ul{Idem in eodem libro vicesimo nono. \grecross .} {\color{Red} Lectio Sexta.}
\lettrine[lines=2]{\zallmancaps \color{Red} H}{orum} iustorum vita alibi palme comparatur, que videlicet inferius est aspera: superius vero et visu et fructibus pulchra, inferius angusta: superius amplitudine pulchre viriditatis expansa. Sic eciam electorum vita despecta inferius, pulchra superius, in imis multis tribulationibus angustata: in summis est amplitudine summe retributionis expansa. Habet eciam preter alias arbores hoc palma: quod tenuis ab imis incipit: sed vastior ad summa succrescit. Arbusta igitur alia, mundi huius dilectoribus sunt similia: qui in terrenis sunt fortes, in celestibus debiles. At contra, iusti in terrenis studiis sunt similes, qui cum ad conversionem veniunt: non tales perseverant, quales ceperunt. Palme vero iustorum conversio comparatur: que scilicet plus in finiendo peragit, quam inchoando proponit, et si tepidius prima inchoat: ferventius extrema consummat.
\newline \ul{Ricardus in exceptionibus secunda parte, libro decimo. \grecross \ } {\color{Red} Lectio VII.}
% https://mlat.uzh.ch/browser/cps_2.AucInc.Sermon37#:9;2
\lettrine[lines=2]{\zallmancaps \color{Blue} D}{e} iustis huiusmodi: dicitur alibi. Iustus germinabit sicut lilium. Multe quidem lilii proprietates iusto conveniunt. Viror designat fidem: eo quod prima sit in virtutibus, sicut lilium primum cernitur in germinantibus. Erectio significat spem: per quam erigimur. Medicinalis est iustus instar lilii per predicationem. Habet quoque ut lilium fibras aptissime intus in radice sibi coniunctas: dum in ipso concordat cogitatio cogitationi, sensus sensui, voluntas voluntati, affectio affectioni. Habet hoc lilium folia, bona scilicet opera, quibus foris circumdatur et ornatur. Et hec folia eriguntur: quia quicquid exercet, sursum levatur per intentionem. Per iusticiam rectus est: tribuit enim unicuique quod suum est. Per ramusculos sursum invenitur pluraliter discretum, dum diversis spiritus sancti donis multipliciter perfunditur. Florem habet candidum, per mundiciam: sexifidum per vite perfectionem: intus quibusdam granulis croceum que significant caritatem: suaviter odoriferum, quod exprimit famam.%typo: fibras for dossas
\newline \ul{Idem in eodem. \grecross \ } {\color{Red} Lectio VIII.}
% https://mlat.uzh.ch/browser/cps_2.AucInc.Sermon37#:7;2
\lettrine[lines=2]{\zallmancaps \color{Red} L}{audem} huiusmodi iusti decantans scriptura: pronunciat dicens. Erit tanquam lignum, quod plantatum est secus decursus aquarum. Aque iste dona sunt gratie, quibus Deus iustum irrigat: ut crescat et fructum faciat. Decursus aquarum: vices sunt aspirationum. Affluit namque spiritus ad nostram iustificationem, defluit ad nostram humiliationem. Accedit ut crescamus in virtute: recedit ne propter abundantem gratiam superbiamus de virtute. Fons ergo Christus: flumina dona quibus iustus vegetatur, augmentatur et consummatur. Vegetatur per veram fidem, augmentatur per bonam operationem: consummatur per vite perfectionem. Radix huius arboris est fides: que est origo ac fundamentum aliarum virtutum et radix bonorum operum. Stipes designat spem per quam erigimur.
\newline {\color{Red} Lectio IX.}
\lettrine[lines=2]{\zallmancaps \color{Blue} R}{ami} huius ligni caritatem insinuant: que per dilectionem Dei et proximi extenditur. Folia quibus arbor vestitur: bonam significant actionem, viror actionis perseverantiam: medulla cordis intentionem, cortex conversationis superficiem: flores propter redolentiam, bonam famam. Habet quoque lignum hoc fructum, et intra et extra et supra. Intra: in animo per bonam conscientiam, extra: coram proximo per doctrinam: supra: in Deo per gloriam. Talis fuit karissimi iustus iste cuius hodie festa celebramus, cuius virtutes recitamus, cuius sanctitatem commendamus: cuius suffragium imploramus, cuius illuminamur verbo, formamur exemplo: munimur patrocinio. Tales igitur et nos esse karissimi laboremus: et imitemur patrem filii, magistrum discipuli, satellites militem, oves pastorem.
{\ubsubsection{In Festo Unius Confessoris Non Pontificis Trium Lectionum}{in-festo-unius-confessoris-non-pontificis-trium-lectionum-breviarium}}
\newline \ul{Bernardus in sermone de obitu sancti Humberti. \grecross \ } {\color{Red} Lectio Prima.}
% https://mlat.uzh.ch/browser/cps_2.BerCla.InObDHu#:1.1;2
\lettrine[lines=2]{\zallmancaps \color{Red} V}{ir} iste sanctus karissimi facticium nobis sermonem in omni forma sanctitatis exhibuit: nec oportet me os amplius aperire, si bene sermonem eius retinuistis. Nam sicut advena et peregrinus viam ac vitam istam pertransiit, quantum minus potuit de mundi rebus accipiens: utpote se non esse de hoc mundo cognoscens. Non habebat hic manentem civitatem sicut nec patres sui, sed extentus in anteriora sequebatur ad palmam superne vocationis. Nichil habet mundus quod clamet in eo, vel de ipso: quia nec mundus ei placuit, nec ipse mundo. Victum et vestitum habens his contentus fuit, non ad superfluitatem sed ad necessitatem: nisi quod ipsam necessitatem sepius causabatur esse superfluitatem. Erat autem vere mitis et humilis corde: ideoque sese amabilem et affabilem omnibus exhibebat.
\newline {\color{Red} Lectio Secunda.}
\lettrine[lines=2]{\zallmancaps \color{Blue} Q}{uis} unquam ex eius ore sonum detractionis, verbum scurrilitatis, glorie sermonem, invidie vocem audivit? Quis eum vel alios iudicantem, vel iudicanti consentientem aliquando deprehendit? Quis eum inania loquentem audire potuit, immo quis non ab eo si forte loqueretur talia audiri timuit? Nimirum sollicite custodiebat vias suas: ut non delinqueret in lingua sua. Nunquid aliquis eum ridentem eciam inter multos ridentes invenit? Vultum quidem suum assidentium gratia ne fieret onerosus serenabat: sed risum integrum non admittebat. Porro quanti fervoris in opere Dei diebus ac noctibus usque
%pp139
ad mortem fuerit: admirati estis. Sic se viscera pietatis induerat, ut omnes excusaret, pro omnibus eciam nescientibus intercederet, non acceptor utique personarum sed necessitatum.
\newline {\color{Red} Lectio Tertia.}
\lettrine[lines=2]{\zallmancaps \color{Red} E}{rat} igitur humilis corde, dulcis sermone, strenuus opere, fervens caritate, in commisso fidelis, in consilio circumspectus, supra ceteros compositus: idem omni tempore permanens. Ecce iam coram te pater dulcissime fons ille puritatis est: quem tanto ardore sitiebas. Ecce immersus es in illam divine pietatis abyssum: cuius tam devote memoriam abundantie suavitatis eructare solebas. Cui enim aliquando quinque verba locutus es, in quibus non vera puritas, non sancti Dei pietas resonaret? De reliquo fratres, si eius vestigia sequeremini, non tam facile in vanis cogitationibus et ociosis sermonibus, in iocis et scurrilitatibus laberemini. In his quippe multum et de vita vestra, et de tempore perditis. Scio quidem quia durum est dissoluto disciplinam apprehendere: verboso silentium tenere, vagari solito stabilem permanere: sed multo durius erit futuras molestias tolerare. Ergo karissimi, in doctrina ista vos exercete, et formam attendite quam in sancto audistis: ut ad eum ad quem ipse pervenit, et vos perveniatis.
{\ubsubsection{In Festo Unius Confessoris Doctoris Novem Lectionum}{in-festo-unius-confessoris-doctoris-novem-lectionum-breviarium}}
\newline \ul{Gregorius, in moralibus libro vicesimo septimo. \grecross \ } {\color{Red} Lectio Prima.}
\lettrine[lines=2]{\zallmancaps \color{Blue} D}{uo} in hac vita sunt genera iustorum: unum bene viventium, sed nulla docentium, aliud recte viventium, et recta docentium: sicut et in celi facie alie stelle prodeunt quas nulle pluvie subsequuntur, alie prodeunt: que terram magnis ymbribus infundunt. Possunt per stellas pluvie, sancti apostoli designari. De quibus recte dicitur in Iob. Qui aufert stellas pluvie: et effundit ymbres ad instar gurgitum. Quia videlicet reductis ad superna apostolis, per expositorum sequentium linguas: fluenta divine scientie diu abscondita, largiore effusione dominus patefecit. Nam quod illi sub brevitate locuti sunt: hoc exponendo isti multipliciter auxerunt. Unde expositorum predicamenta gurgitibus comparantur, quia dum multorum precedentium dicta colligunt: quasi ex guttis gurgites faciunt.
\newline {\color{Red} Lectio Secunda.}
\lettrine[lines=2]{\zallmancaps \color{Red} N}{equaquam} autem apostolis expositores in scientia se preferant: cum exponendo latius loquuntur. Meminisse quippe incessanter debent: per quos eiusdem scientie inventiones acceperunt. Unde et apte subiungitur. Qui de nubibus fluunt. Ymbrium quippe gurgites de nubibus fluunt: quia profunde sequentium intelligentie, originem a sanctis apostolis acceperunt. De quibus adhuc nubibus, apte subiungitur. Que pretexunt cuncta desuper. Nubes cum desuper aerem pretexunt, si oculos attollimus: non celum sed ipsas intuemur. Et cum de celo sol rutilat, prius interfuso aere spargitur: ut post in celo solis radios contemplemur. Quia vero carnales sumus: intellectus noster ad divina transire non sinitur. Sed cum celum suspicit nubes videt, quia comprehendere que Dei sunt appetit: sed vix mirari ea que hominibus sunt collata potest. Unde et alias dicitur. Illuminans tu mirabiliter a montibus eternis. Qui enim orientem solem contemplari non valet, radiatos montes aspicit: et quia ortus sit sol deprehendit.
\newline \ul{Idem in eodem.} \grecross \ {\color{Red} Lectio Tertia.}
\lettrine[lines=2]{\zallmancaps \color{Blue} D}{e} his nubibus alibi scriptum est. Nunquid nosti semitas nubium eius magnas et perfectas scientias? Habent iste nubes subtilissimas sancte predicationis vias, angusta quippe est porta que ducit ad vitam: habent et perfectas scientias. Perfecta scientia est scire omnia: et tamen iuxta quendam modum, scientem se esse nescire. Quia etsi iam Dei precepta novimus, etsi iam que intellexisse credimus agimus: adhuc tamen acta eadem qua districtione examinis sint discutienda nescimus, necdum faciem Dei cernimus: necdum occulta Dei iudicia videmus. Quanta ergo est nostra scientia, que quousque mortalitatis pondere premimur, ipso sue incertitudinis nubilo cecatur? In hoc itaque mundo dum vivimus, tunc perfecte que scienda sunt scimus: cum proficientes per intelligentiam, nil nos perfecte scire cognoscimus. Dicatur ergo. Nunquid nosti semitas nubium eius magnas, et perfectas scientias? Ac si aperte dicat. Nunquid summa acta predicantium conspicis, qui postquam per scientiam ad alta se elevant, deorsum se humiliter per cognitionem sue ignorantie inclinant?
\newline \ul{Idem in eodem decimo octavo. \grecross \ } {\color{Red} Lectio Quarta.}
\lettrine[lines=2]{\zallmancaps \color{Red} Q}{uia} vero testamenti novi doctores ad hoc usque perducti sunt ut occultas allegoriarum caligines in veteri testamento rimentur: recte alibi dicitur. Profunda quoque fluviorum scrutatus est: et abscondita produxit in lucem. Quid namque hic aliud flumina nisi antiquorum patrum dicta nuncupantur? Quis enim pensare sufficiat, quam vehemens fluvius dum legem conderet, ex pectore Moysi erupit, quam vehemens fluvius ex David corde defluxit, quanta Salomonis ore, aut que omnium prophetarum fluminum fluenta manarunt? Sed horum fluminum Iudea superficiem tenuit: dum littere superficiem servans, eorum profunda nescivit. Nos vero qui veniente domino in eis interna spiritualia querimus: eorum profunda rimamur. Quod ipse facere dicitur: quia id agere ipso donante prevalemus.
\newline {\color{Red} Lectio Quinta.}
\lettrine[lines=2]{\zallmancaps \color{Blue} P}{rofunda} ergo fluviorum dominus scrutatur, et abscondita in lucem producit: quia dicta legis caliginosa, nunc expositio spiritualis illuminat. Aperta namque dicta exponentium: fecerunt nobis esse iam conspicuas sententias antiquorum. Unde et Ysaias plana sancte ecclesie verba expositionis conspiciens: exclamavit dicens. Locus fluviorum rivi latissimi et patentes. Testamenti enim veteris dicta quasi angusti et clausi rivi fuerant, qui immensas sue scientie sententias collectione obscurissima constringebant. At contra: doctrina sancte ecclesie, rivi et lati et patentes sunt: quia eius dicta et invenientibus multa sunt, et querentibus plana, quia dum suis expositoribus intelligentie spiritum dominus infudit: prophetantium obscuritates aperuit. Nullatenus autem sine eius sapientia, dicta Dei penetrantur. Nisi enim quis spiritum eius acceperit: eius nullo modo verba cognoscit.
\newline \ul{Idem in eodem.} \grecross \ {\color{Red} Lectio Sexta.}
\lettrine[lines=2]{\zallmancaps \color{Red} N}{otandum} autem: quod in argento eloquium designari solet. Heretici autem sic de eloquii sui nitore superbiunt, ut nulla sacrorum librorum autoritate solidentur: qui nobis ad loquendum quasi quedam vene argenti sunt. Et ideo dominus eos ad sacre autoritatis paginas revocat, ut si vere loqui desiderant, inde sumere debeant, quid loquantur: cum ait alibi. Habet argentum venarum suarum principia. Ac si aperte dicat. Qui ad vere predicationis verba se preparat: necesse est ut causarum origines a sacris paginis sumat, ut omne quod loquitur ad divine auctoritatis fundamentum revocet: atque in eo edificium sue locutionis firmet. Quia dum heretici quasi nova quedam proferunt que in antiquorum patrum libris veteribus non tenentur, fit ut dum videri sapientes desiderant: miseris suis auditoribus stulticie semina spargant.
\newline \ul{Idem in eodem libro vicesimo octavo. \grecross \ } {\color{Red} Lectio VII.}
\lettrine[lines=2]{\zallmancaps \color{Blue} N}{on} immerito quoque doctores sancti basium nomine designantur: cum alibi dicitur. Super quo bases illius solidate sunt? Quia dum recta predicant, et predicationi sue vivendo concordant: omne pondus ecclesie morum suorum fixa gravitate sustentant. Unde et bene cum in typo ecclesie tabernaculum figeretur: ad Moysen dicitur. Facies columnas quatuor et bases earum vestitas argento. Bases quippe argento vestite, quatuor columnas tabernaculi sustinent: quia predicatores ecclesie divino eloquio decorati, ut in cunctis, se exemplum prebeant: quatuor evangelistarum dicta et ore et operibus portant. Dicat ergo dominus. Super quo bases illius solidate sunt, subaudi nisi super me, qui cuncta mirabiliter teneo, et bonis exterioribus intus originem presto?
\newline \ul{Idem in eodem libro tricesimo. \grecross \ } {\color{Red} Lectio VIII.}
\lettrine[lines=2]{\zallmancaps \color{Red} D}{e} huiusmodi doctoribus alibi dicitur. Quis posuit in visceribus hominis sapientiam? vel quis dedit gallo intelligentiam? Ac si diceret dominus. In cor hominis humana sapientis superne sapientie gratiam quis infudit, vel ipsis predicatoribus, quis nisi ego, intelligentiam dedit? Sapientia vero divinitus inspirata in visceribus ponitur, quia non in solis vocibus sed eciam in sensibus datur, ut iuxta quod loquitur lingua: vivat et conscientia. Gallo quoque intelligentia desuper tribuitur, quia doctori veritatis discretionis virtus ut noverit quibus, quid, quanto, vel quomodo inferat divinitus ministratur. Non enim una eademque cunctis exhortatio convenit: quia nec cunctos par morum qualitas astringit. Sepe autem aliis officiunt: que aliis prosunt. Nam et plerumque herbe que hec animalia reficiunt: alia occidunt. Et levis sibilus equos mitigat: catulos instigat. Et medicamentum quod hunc morbum imminuit: alteri vires iungit. Et panis qui vitam fortium roborat: parvulorum necat.
\newline {\color{Red} Lectio IX.}
\lettrine[lines=2]{\zallmancaps \color{Blue} P}{ro} qualitate igitur audientium formari debet sermo doctorum, ut ad sua singulis congruat: et tamen a communis edificationis arte nunquam recedat. Quid enim sunt intente mentes auditorum, nisi quedam in cythara tensiones strate cordarum, quas tangendo artifex ut non sibimetipsi dissimile canticum faciant, dissimiliter pulsat? Et idcirco corde consonam modulationem reddunt, quia uno quidem plectro: sed non uno impulsu feriuntur. Unde et doctor quisque ut in una cunctos virtute caritatis edificet: ex una doctrina non una eademque exhortatione tangere corda audientium debet. Quia igitur laus tante intelligentie non predicatoris virtus est sed actoris: recte per eiusdem actorem dicitur. Vel quis dedit gallo intelligentiam? Ac si diceret. Nisi ego qui doctorum mentes quas mire ex nichilo condidi, ad intelligenda que occulta sunt, mirabilius instruxi?
{\ubsubsection{In Festo Unius Confessoris Doctoris Trium Lectionum}{in-festo-unius-confessoris-doctoris-trium-lectionum-breviarium}}
\newline \ul{Idem in eodem libro decimo. \grecross \ } {\color{Red} Lectio Prima.}
\lettrine[lines=2]{\zallmancaps \color{Red} Y}{adum} nomine que imminente verno, ad celi faciem prodeunt cum sol iam caloris sui vires exerit, atque ad extendenda diei spatia ferventior exurgit, doctores ecclesie sancte designantur, qui eo tempore ad mundi noticiam venerunt, quo fides clarius elucet, et altius sol veritatis calet, expletisque longis noctibus infidelitatis per credulitatis vernum lucidior annus aperitur. Bene autem Yadum nuncupatione signantur. Greco quippe eloquio ydos pluvia vocatur: et ipsi sancte predicationis ymbres fuderunt. Dum ergo Yades cum pluviis veniunt: ad celi spatia altiora sol ducitur, quia apparente doctorum scientia, fidei calor augetur, et perfusa terra ad fructum proficit: quia uberiorem frugem boni operis reddimus, dum per sacram eruditionis flammam in corde clarius ardemus. Dumque per eos diebus singulis scientia celestis ostenditur, quasi interni nobis luminis vernum tempus aperitur, ut novus sol nostris mentibus rutilet: et eorum verbis nobis cognitus, se ipso cotidie clarior micet. Vergente etenim mundi fine superna scientia proficit: et largius cum tempore excrescit.
\newline {\color{Red} Lectio Secunda.}
\lettrine[lines=2]{\zallmancaps \color{Blue} N}{onnulli} vero a Greca littera que, y, dicitur, Yadas nuncupatas arbitrantur. Quod si ita est doctores his stellis non inconvenienter expressi sunt: qui a litteris nomen trahunt. Vir igitur sanctus redemptionis nostre ordinem contemplatus admiretur: atque admirans exclamet dicens. Qui extendit celos solus et graditur super fluctus maris: qui facit Arcturum, et Orionas et Yadas. Extensis enim celis dominus formavit Arcturum: quia in honorem deductis apostolis in celesti conversatione fundavit ecclesiam. Formato quoque Arcturo fecit Orionas: quia roborata fide universalis ecclesie contra procellas mundi edidit martyres. Editis quoque Orionibus protulit Yadas: quia convalescentibus contra adversa martyribus ad infundendam ariditatem humanorum cordium, doctrinam contulit magistrorum.
\newline \ul{Gregorius in secunda parte super Ezechielem. Omelia sexta. \grecross \ } {\color{Red} Lectio Tertia.}
\lettrine[lines=2]{\zallmancaps \color{Red} I}{sti} sunt gazophilatia templi: quia gazophilatium appellari solet locus quo divitie servantur. Quid itaque per gazophilatium signatur, nisi corda doctorum sapientie atque scientie divitiis plena? His divitiis abundare discipulos magister gentium viderat: cum dicebat. Divites facti estis in illo: in omni verbo et in omni scientia. Has veras esse divitias veritas denunciat: cum de transitoriis dicit. Fallacia divitiarum suffocat verbum. Vere enim divitie sunt: in quarum comparatione que transire possunt false nominantur. De his divitiis per Salomonem dicitur. Thesaurus desiderabilis, requiescit in ore sapientis. Excepto autem eo quod ad eternam patriam divitie spirituales ducunt, est eis a terrenis divitiis magna distantia, quia spirituales erogate proficiunt, terrene autem aut erogantur et deficiunt: aut retinentur et utiles non sunt. Qui ergo has in se veras divitias continent: gazophilatia recte vocantur.
{\ubsubsection{In Festo Unius Virginis et Martyris Novem Lectionum}{in-festo-unius-virginis-et-martyris-novem-lectionum-breviarium}}
\newline \ul{Sermo beati Augustini.} \grecross \ {\color{Red} Lectio Prima.}
\lettrine[lines=2]{\zallmancaps \color{Blue} H}{odie} nobis fratres karissimi virginale decus illuxit: quod nostre memorie virtutis divine victoriam renovavit. Porro quanto fragilior sexus, quanto infirmius vasculum, quod reportat ex hoste triumphum: tanto maiori diabolus obprobrio confusionis induitur, tanto mirabilior deus in sanctis suis agnoscitur: tanto rex martyrum Christus pugnatricibus suis iocundius delectatur. Hinc est enim quod salvator noster in huius mundi campum potestates aerias debellaturus adveniens, non elegit sapientes secundum carnem, non nobiles, non potentes: sed que stulta sunt huius mundi elegit: ut confundat fortia, et ignobilia et contemptibilia elegit et ea que non sunt ut ea que sunt destrueret: ut non glorietur omnis caro coram illo.
\newline \ul{Item sermo eiusdem secundum quosdam libros, licet magis videantur esse verba Fulgentii vel Maximi.} \grecross \ {\color{Red} Lectio Secunda.}%typo: quostam
\lettrine[lines=2]{\zallmancaps \color{Red} R}{efulget} et preminet inter martyres meritum sancte huius domini famule: nam ibi est corona gloriosior ubi sexus infirmior, quia profecto virilis animus in femina maius aliquid fecit: quando sub tanto pondere fragilitas non defecit. Bene inheserat uni viro a quo virtutem traxerat qua resisteret diabolo: ut femina prosterneret inimicum, qui per feminam prostraverat virum. Ille in ea apparuit invictus: qui pro ea factus est infirmus. Ille fecit hanc feminam viriliter mori: qui pro ea dignatus est de femina misericorditer nasci. Merito autem ille vetus inimicus, ne ullas preteriret insidias, qui per feminam vicerat virum: per virum hanc feminam superare temptavit. Sed illa interioris hominis cautissimo firmissimoque robore, omnes eius insidias obtudit impetusque confregit.
\newline {\color{Red} Lectio Tertia.}
\lettrine[lines=2]{\zallmancaps \color{Blue} Q}{uid} gloriosius hac femina quam viri mirantur facilius quam imitantur? Sed hoc illius potissimum laudis est, in quem nec masculus nec femina invenitur: qui eciam in his que sunt femine corpore virtutem mentis in sexu carnis abscondit. Calcatus est draco pede casto: et caput serpentis antiqui quod fuit precipitium femine cadenti, gradus factum est ascendenti. Quid hoc spectaculo suavius, quid hoc certamine fortius, quid hac victoria gloriosius? Cum sanctum corpus diversis suppliciis affligeretur: toto amphiteatro fremebant gentes, et populi meditabantur inania. Sed qui habitat in celis irridebit eos: et dominus subsannabit eos. Nunc autem posteri illorum quorum voces in carne martyris impie seviebant: martyrem piis vocibus laudant. Neque tunc tanto concursu hominum ad eos occidendos cavea crudelitatis impleta est: quanto nunc ad eos honorandos ecclesia pietatis impletur. Omni anno nunc spectat cum religione caritas: quod uno die cum sacrilegio tunc commisit impietas.
\newline \ul{Sermo Ambrosii. \grecross \ } {\color{Red} Lectio Quarta.}
% https://mlat.uzh.ch/browser/cps_2.AmbMed.DeVir#:1.2;2
\lettrine[lines=2]{\zallmancaps \color{Red} N}{atalis} est virginis: integritatem sectamur. Natalis est martyris, hostias immolemus. Stupeant nupte: imitentur innupte. Sed quid de ista virgine dignum dicere possumus? Facescant ingenia: conticescat eloquentia. Vox preconum est una: quot homines tot precones, qui martyrem dum loquuntur predicant. Hanc senes, iuvenes, et pueri laudant. Nemo laudabilior est: quam qui ab omnibus laudari potest. Cetere quidem puelle torvos vultus parentum non possunt ferre: et acu puncta districta solent quasi vulnera flere. Hec autem virgo cruentas carnificum impavida, manus expectavit, hec stridentium gravibus immobilis cathenarum tractibus cruento mucroni militis totum obtulit corpus: mori parata si ad aras raperetur invita. Processit ad locum supplicii leta successu, festina gradu, non in torto crine caput compta sed Christo: non flosculis sed moribus redimita.
\newline {\color{Red} Lectio Quinta.}
\lettrine[lines=2]{\zallmancaps \color{Blue} Q}{uanto} terrore carnifex egit, ut illa timeretur, quantis blandiciis ut ei persuaderet? At illa sponsi sui cavens iniuria cui se optabat placituram: ille me accipiet inquit qui me sibi prior elegit. Pereat corpus, quod oculis amari potest carnalibus. Denique stetit et oravit: ac mortis supplicio cervicem inflexit.
\newline \ul{Idem in alio sermone.} \grecross \
\newline Habetis in una hostia duplex martyrium: pudoris et religionis. Beata quippe hec: et virgo permansit, et martyrium obtinuit.
% https://mlat.uzh.ch/browser/cps_2.AmbMed.DeVir#:1.3.1;2
Invitati ergo integritatis amore: aliquid de virginitate, ne velut in transitu perstricta videatur, que principalis virtus esse comprobatur. Non enim ideo laudabilis est virginitas quod et in
%pp140
martyribus reperta fit: sed quoniam ipsa martyres facit.
\newline {\color{Red} Lectio Sexta.}
\lettrine[lines=2]{\zallmancaps \color{Red} Q}{uid} autem mirum si virginitatem pretulerunt vite, virgines Christi, quando similia fecerunt virgines seculi?
\newline \ul{Hieronimus contra Iovinianum libro primo.} \grecross \
% https://mlat.uzh.ch/browser/cps_2.HieStr.AdvJov#:3.39;2
\newline Quis enim valeat silentio preterire septem millesias virgines que Gallorum impetu cuncta vastante, ne quid indecens ab hostibus sustinerent turpitudinis mortem sumpserunt, exemplum sui cunctis virginibus relinquentes magis pudiciciam cure esse quam vitam? Nichanor victis Thebis atque subversis, unius captive virginis amore superatus est. Cuius coniugium expetens et voluntarios amplexus, quod scilicet captiva optare debuerat: sensit pudicis mentibus plus virginitatem esse quam regnum, interfectamque propria manu flens et lugens amator tenuit. Narrant scriptores Grecie, et aliam Thebanam virginem quam hostis Macedo corruperat, dissimulasse paulisper dolorem, et violatorem virginitatis sue iugulasse postea dormientem, seque interfecisse cum gaudio, ut nec vivere voluerit post perditam castitatem: nec ante mori quam sui ultrix existeret.
\newline \ul{Gregorius in omelia loquente Ihesu. \grecross \ } {\color{Red} Lectio VII.}
\lettrine[lines=2]{\zallmancaps \color{Blue} C}{onsideremus} nos qui viri sumus, in huius femine cuius sollemnia agimus comparatione, quid extimabimur. Sepe agenda aliqua bona proponimus, sed si levissimus sermo ab ore irridentis irruperit, fracti protinus resilimus: hanc vero frangere nec tormenta potuerunt. Nos ad precepta dominica largiri nostra superflua nolumus: hec pro illo eciam propriam carnem dedit. Cum ergo ad illud terribile examen districtus iudex venerit, quid nos viri dicemus cum huius femine gloriam viderimus? Que tunc erit viris excusatio, quando hec ostendetur, que cum seculo et sexum vicit?
\newline \ul{Idem in omelia simile est regnum celorum thesaurum.} \grecross \ 
\newline Erectus namque in virtutis culmine eius animus tormenta derisit, premia calcavit: ante armatos duces et presides ducta stetit, feriente robustior, iudicante sublimior. Quid igitur nos barbati et debiles dicemus, qui ire ad regna celestia puellas per ferrum videmus, quos ira superat, superbia inflat, ambitio perturbat, luxuria inquinat?
\newline \ul{Maximus in sermone, O beatissima. \grecross \ } {\color{Red} Lectio VIII.}
% https://mlat.uzh.ch/browser/cps_2.MaxTau.Sermon10#
\lettrine[lines=2]{\zallmancaps \color{Red} V}{irgo} hec beatissima magnum de se excitavit exemplum pudicicie: cum virginitatem pretulit vite. Omnem quoque terrenum amorem a se per amorem celestem expellens, crucis Christi tropheum arripuit: et tam contra libidines quam eciam contra universa supplicia fide potius armata quam ferro dimicavit, ac vicit, contempsit blandientem, renuitque minantem: supplicium attendens risit, nesciens aliquid amare quod transit. Denique virginitatem nimis diligens, nec ludibria nec supplicia nec carnifices metuit: postremo nec gladium furentis percussoris expavit. Turbatur tyrannus a puella contemptus, sacratissima virginitas usque ad sanguinis effusionem in suo gradu perseverat, et inconcussa Christo semel oblata perdurat: quam ille nimirum suo sanguine subarraverat.%typo: tyrannius
\newline {\color{Red} Lectio IX.}
\lettrine[lines=2]{\zallmancaps \color{Blue} V}{eni} iam beatissima virgo ad preciosum thalamum quem quesisti, et gladium percussoris admitte pro celestis amore sponsi, ut te sponsam, immaculata Christi virginitas comprobet: et martyrem eius confessio probata consignet. Felix certe virginum caterva: quibus studium est confessionis tue sequi vestigia. Habebunt enim tecum coronam in celo: que tecum spreverint inimicum in mundo. Nimirum sicut eterne remunerationis apparuit participatio tibi prosequenti Marie vestigia: sic et ceteris virginibus te imitantibus, eterna credimus gaudia non neganda. Igitur Christo splendida, Dei filio pulchra, omnibus angelis et archangelis grata, precibus quibus possumus ut nostri meminisse digneris exoramus: quatinus ille nobis peccatorum tribuat indulgentiam, qui tibi tuorum omnium laborum tradidit palmam.
{\ubsubsection{In Festo Unius Virginis Martyris Trium Lectionum}{in-festo-unius-virginis-martyris-trium-lectionum-breviarium}}
\newline \ul{Sermo beati Maximi. \grecross \ } {\color{Red} Lectio Prima.}
% https://mlat.uzh.ch/browser/cps_2.AucVar.SeSAmHa#:48.1;2
\lettrine[lines=2]{\zallmancaps \color{Red} C}{um} in toto mundo virgineus Marie flos immarcessibiles innectat coronas et immaculato conservet affectu sceptrigeram pudoris aulam, eo usque integritas perseveravit ad palmam: ut in puellis tropheum sanctitatis arriperet, ac ad celestem thalamum perveniret.
\newline \ul{Crisostomus in omelia de Machabeis. \grecross \ }
\newline Non solum enim senes ac iuvenes, sed eciam mulieres expedite ad premia celestia cucurrerunt: et illustres coronas adepte sunt. In mundanis quippe agonibus etatis ac forme generisque dignitas requiritur: et ideo servis ac mulieribus et senibus ac pueris ad eos aditus denegatur. In celestibus autem omni persone et etati ac sui indiscreta facultate stadium pater, et agonis adeundi commune ius et una libertas est: ut propensa sine invidia omnibus bonitate, et ineffabilem salvatoris virtutem agnoscas, et apostolicum sermonem, quia virtus in infirmitate perficitur, opere confirmari perspicias.
\newline \ul{Crisostomus in omelia de militia spirituali.} \grecross \ {\color{Red} Lectio Secunda.}
\lettrine[lines=2]{\zallmancaps \color{Blue} A}{pud} deum ergo femineus eciam militat sexus, Multe namque femine animo virili spiritalem militiam gesserunt. Quedam enim interioris hominis virtute viros equaverunt in agonibus martyrum: quedam eciam fortiores viris extiterunt. Nam quibusdam virorum ex ignavia muliebri degeneratibus, ille sexum infirmum virili mente vicerunt: ac sic in testimonio virtutis victrices penarum, vires labentibus constituerunt. Ex ipsis ergo sunt: que angelicum sacre virginitatis replent chorum. Ex ipsis et ille sunt: que gloriosis preliis et confessorum et martyrum triumphis coruscant. Sequebantur enim et ipsum dominum in eius presentia sicut scriptura designat, non solum viri sed et mulieres et obsequebantur et ministrabant ei. His ergo virtutibus quondam certatim, et patres filios, et matres Christo filias offerebant: suamque progeniem ad militandum regi celesti probare festinabant. Talia vero pignora vos devoti parentes habebitis non solum in consortium glorie, sed eciam in ornamentum et patrocinium perenne.
\newline \ul{Idem in omelia ad Eustropium vicesimo nono. \grecross \ } {\color{Red} Lectio Tertia.}
\lettrine[lines=2]{\zallmancaps \color{Red} O}{}\ diabole: nescis que tibi fecerint martyres? Ingressa est puella etate tenera, annis immatura, et inventa est ferro fortior cum latera eius scinderes: fidem tamen eius movere non posses. Frequenter in tormentis caro defecit: et robur fidei non defecit. Consumptum est corpus: et mens non potuit inclinari. Interiit substantia: et perstitit patientia.
\newline \ul{Ciprianus in epistola de habitu virginum.} \grecross \ 
\newline Neminem certe Christianum et maxime virginem decet claritatem aliquam carnis et honorem reputare: sed solum bona in eternum mansura diligere, aut si gloriadnum sit in carne: tunc plane cum in nominis Christi confessione cruciatur, cum fortior femina viris torquentibus invenitur: cum ignes aut cruces aut ferrum aut bestias patitur, ut coronetur. Ista sunt preciosa carnis monilia: ista corporis ornamenta meliora.
\newline \ul{Idem in epistola ad confessores. \grecross \ }
\newline Beate femine que dominicam fidem tenentes ac sexu suo fortiores, non solum ad coronam proximaverunt: sed et ceteris exemplum de sui constantia prebuerunt.
{\ubsubsection{In Festo Unius Virginis Non Martiris Si Fiant Novem Lectionum Ratione Patrie Vel Patronatus, Vel Alia De Causa}{in-festo-unius-virginis-non-martiris-si-fiant-novem-lectionum-ratione-patrie-vel-patronatus-vel-alia-de-causa-breviarium}}
\newline \ul{Augustinus in libro de virginitate. \grecross \ } {\color{Red} Lectio Prima.}
\lettrine[lines=2]{\zallmancaps \color{Blue} I}{n} festo virginis predicanda est virginitas. Hoc in isto sermone suscepimus. Adiuvet ad hoc Christus virginis filius et virginum sponsus, virginali utero corporaliter natus, virginali connubio spiritualiter coniugatus. Cum igitur universa ecclesia virgo sit sicut dicit apostolus, quanto digna sunt honore membra que hoc custodiunt in carne quod ipsa tota custodit in fide? Nulla ergo fecunditas carnis sancte virginitati eciam carnis comparari potest. Illa autem virgo coniugate merito preponitur, que nec multitudini se amandam preponit: sed speciosum forma pre filiis hominum sic amavit, ut quia eum sicut Maria concipere carne non posset: ei corde concepto, carnem eciam integram custodiret. Virginalis autem integritas et ab omni concubitu immunitas portio est angelica et in carne corruptibili incorruptionis perpetue meditatio.
\newline {\color{Red} Lectio Secunda.}
\lettrine[lines=2]{\zallmancaps \color{Red} C}{edat} virginitati omnis fecunditas carnis: omnis pudicicia coniugalis. Illa non est in potestate: illa non est in eternitate. Fecunditatem carnalem non habet liberum arbitrium: pudiciciam coniugalem non habet celum. Profecto habet virginitas magnum aliquid preter ceteros in illa communi immortalitate: que habet aliquid iam non carnis in carne. Cum enim dixisset dominus de spadonibus, dabo eis in domo mea locum nominatum multo meliorem quam filiorum atque filiarum, continuo subiecit, nomen eternum dabo eis: quod utique gloriam quandam propriam excellentemque significat. Nam ideo fortassis et nomen dictum est: quod eos quibus datur distinguit a ceteris. Pergite ergo sancti pueri dei ac puelle, mares ac femine, celibes et innupte: pergite perseverantes in finem. Laudate dulcius, quem cogitatis uberius, sperate felicius cui servitis instantius: amate ardentius cui placetis attentius.
\newline {\color{Red} Lectio Tertia.}
\lettrine[lines=2]{\zallmancaps \color{Red} L}{umbis} accinctis et lucernis accensis expectate dominum virgines quando veniat, afferetis ad nuptias agni canticum novum, quale nemo poterit dicere: nisi duodena milia sanctorum cytharedorum illibate virginitatis. Et quia sequi agnum quocumque ierit scripsit de vobis Iohannes, quo ire putamus hunc agnum quo nemo eum sequi debeat vel audeat nisi vos, nisi de Christo, in Christum, post Christum, per Christum, propter Christum? Gaudia propria virginum Christi sunt: nam sunt aliis alia sed nullis talia. Ite, in hec sequimini agnum: quia et agni caro utique virgo est. Sequimini eum tenendo perseveranter quod vovistis, ardenter facite quod potestis ne virginitatis bonum pereat: cum nichil facere possitis ut redeat. Videbit vos cetera multitudo nec invidebit: et illud canticum novum vestrum dicere non poterit, audire tamen poterit et delectari, sed vos qui dicetis felicius exultabitis, iocundiusque regnabitis.
\newline {\color{Red} Lectio Quarta.}
\lettrine[lines=2]{\zallmancaps \color{Blue} S}{equimini} virgines agnum quocumque ierit: sed ad humilem humiliter venite ne cadatis. Qui enim timet cadere: rogat et dicit. Non veniat michi pes superbie. Pergite ergo viam sublimitatis, pede humilitatis. Perseverantes prebeant vobis exemplum: cadentes autem augeant timorem vestrum. Illud amate ut imitemini: hoc lugete ne inflemini. Non inveniantur in vobis improbus vultus, non vagi oculi, non infrenis lingua, non petulans risus, non scurrilis iocus, non indecens habitus: non tumidus aut fluxus incessus. Toto corde amate speciosum forma pre filiis hominum, Inspicite pulchritudinem amatoris vestri. Cogitate equalem patri, subditum matri, in celis dominantem, in terris servientem: creantem omnia, creatum inter omnia. Inspicite vulnera pendentis, cicatrices resurgentis, sanguinem morientis, precium credentis, commercium redimentis. Toto vobis figatur in corde: qui pro vobis fixus est in cruce. Totus maneat in animo vestro: quem noluistis occupari connubio.
\newline \ul{Hieronimus contra Iovinianum, libro primo. \grecross \ } {\color{Red} Lectio Quinta.}
\lettrine[lines=2]{\zallmancaps \color{Red} C}{astitas} semper operi nuptiarum fuit prelata. Nam et primi parentes ante offensam in paradiso virgines fuerunt: post peccatum autem extra paradisum mox nuptias celebraverunt. Itaque nuptie replent mundum: virginitas paradisum. Unde dictum est, crescite et multiplicamini et replete terram. Prius enim oportuit silvam plantari et crescere: ut esset quod deinceps posset incidi. Quantumque interest inter radicem et fructum arboris: tantum inter nuptias et culmen virginitatis. Preterea Moyses moriens a populo plangitur: Iosue vero qui nec uxorem nec filios habuisse legitur, in morte non lugetur. Nuptie quidem finiuntur in morte: virginitas vero incipit coronari post mortem.
\newline \ul{Idem ad Eustochium. \grecross }
\newline Laudo nuptias, laudo coniugium. Sed quia michi virgines generant: lego de spinis rosam, de terra aurum, de conca margaritam.
\newline \ul{Idem contra Iovinianum libro primo. \grecross \ } {\color{Red} Lectio Sexta.}
\lettrine[lines=2]{\zallmancaps \color{Blue} P}{ercurram} breviter Grecas et Latinas barbaras; hystorias: et docebo virginitatem semper tenuisse pudicicie principatum. Quid referam Sibillam Eritheam atque Cumanam, et octo reliquas? Nam varro decem fuisse autumat: quarum insigne virginitas est, et virginitatis premium divinatio. Certe Romanus populus, in quanto honore virgines semper habuerit, hinc apparet: quod consules et imperatores in curribus triumphantes, qui de superatis gentibus trophea referebant et omnis dignitatis gradus, eis de via cedere solitus sit. Nam quod eciam inter duodecim signa celi quibus mundum volui putant, virginem collocarunt: nulli dubium est. Unde magna iniuria nuptiarum: ut ne inter Scorpios quidem et Centauros, et Cancros et Pisces, uxorem maritumque contruserint.%typo: trumphantes
\newline \ul{Cyprianus in epistola de habitu et disciplina virginum. \grecross \ } {\color{Red} Lectio VII.}
% https://mlat.uzh.ch/browser/cps_19.CypCar.DeHaVirCSEL#:3;2
\lettrine[lines=2]{\zallmancaps \color{Red} V}{irginitas} est flos ecclesiastici germinis, decus atque ornamentum gratie spiritalis: leta indoles laudis et honoris, opus integrum et incorruptum: ymago Dei respondens ad sanctimoniam domini: illustrior portio gregis Christi. Gaudet per illas et largiter floret in illis gloriosa fecunditas ecclesie matris. Quantoque plus numero suo copiosa virginitas addit: tanto gaudium matris augescit. Denique quod futuri sumus, vos virgines iam esse cepistis: gloriam resurrectionis in hoc seculo iam tenetis, dum per seculum sine seculi contagione transitis. Virgines magna vos merces manet virtutis: grande premium castitatis. Vos a maledictionis sententia quondam in mulierem lata libere estis. Nullus vobis de partu metus: nec maritus est dominus. Dominus vester et caput: Christus est. Ad instar masculi: vobis communis est conditio. Portemus ait apostolus, ymaginem eius qui de celo est. Hanc ymaginem portat virginitas, integritas, sanctitas.
\newline \ul{Ambrosius in sermone de sancta Agnete. \grecross \ } {\color{Red} Lectio VIII.}
% https://mlat.uzh.ch/browser/cps_2.AmbMed.DeVir#:1.3.2;2
\lettrine[lines=2]{\zallmancaps \color{Blue} Q}{uis} autem humano ingenio virginitatem comprehendere possit, quam nec natura suis legibus inclusit? Aut quis naturali voce complecti noverit, quod supra usum nature sit? Accersivit e celis hec virgo quod imitaretur in terris. Nec immerito usum vivendi quesivit e celo: que sibi sponsum invenit in celo. Hec nubes aeris et angelos sideraque transgrediens, Dei verbum in ipso sinu patris invenit: ac toto pectore hausit. Verbum inquam de sinu patris: virginitas hausit. Quis enim tantum relinquat bonum cum invenerit? Unguentum inquit effusum nomen tuum: ideo te adolescentule dilexerunt et attraxerunt. Postremo non meum illud quod dicitur in evangelio: quia que non nubunt neque nubentur, erunt sicut angeli in celo. Nemoque miretur si angelis comparantur: que angelorum domino copulantur.
\newline \ul{Hieronimus in sermone de Assumptione Beate Virginis, Cogitis. \grecross \ } {\color{Red} Lectio IX.}
% https://mlat.uzh.ch/browser/cps_2.AucVar.HoDeSa#:46;2
\lettrine[lines=2]{\zallmancaps \color{Red} C}{ognata} virginitas est angelis. Profecto in carne preter carnem vivere: non terrena vita est sed celestis. Unde in carne angelicam gloriam adquirere: maioris est meriti quam habere. Esse enim angelum felicitatis est: esse vero virginem virtutis, dum hoc obtinere viribus nititur homo cum gratia: quod habet angelus ex natura.
\newline \ul{Origenis in libro de singularitate clericorum. \grecross \ }
\newline Castitas seu virginitas est munimen sanctimonie, expugnatio infamie, refrenatio lascivie: anime quoque victoria, corporis preda. Ipsa est libertas gloriarum, captivitas criminum, abolitio scandalorum: pax virtutum, et inquietantium debellatio bellorum. Ipsa culmen puritatis, carcer libidinis: portus honestatis. Ipsa denique vita spiritus, carnis interitus, status qualitatis angelice, terminus humane substantie.
{\ubsubsection{In Festo Unius Virginis Non Martyris Trium Lectionum}{in-festo-unius-virginis-non-martyris-trium-lectionum-breviarium}}
\newline \ul{Fulgentius ad probam de virginitate. \grecross \ } {\color{Red} Lectio Prima.}
% https://mlat.uzh.ch/browser/cps_2.FulRus.Episto74#:3.9;2
\lettrine[lines=2]{\zallmancaps \color{Blue} V}{irgo} quasi virago dicitur. Virago autem a viro: et vir a virtute derivatur. Claret igitur inter ceteras ecclesie donationes illud precipuum munus spiritualis esse karismatis: ut nimirum ipsa virtus virginitatis, meretur vocabulo censeri virtutis. Nec in eo quod virginalis integritatis agnoscimus culmen: pudicicie coniugalis asserimus crimen. Sed virginitatem sanctam dicimus merito potiore distare a coniugali vita: quantum distant a bonis meliora, a terrenis celestia, a carne spiritus, ab infirmitate virtus. Immo nec dubitamus tantum a virginitate carnis et spiritus distare concubitum: quantum similitudo pecorum ab imitatione discernitur angelorum. In uno siquidem ad terram spiritus terrena carnis voluptate deprimitur: in altero vero terrena caro celesti delectatione spiritus ad celestia sublevatur.
\newline \ul{Maximus in sermone, O beatissima. \grecross \ } {\color{Red} Lectio Secunda.}
% https://mlat.uzh.ch/browser/cps_2.AucVar.SeSAmHa#:48.1
\lettrine[lines=2]{\zallmancaps \color{Red} Q}{uam} imitabile virginibus virgo hec cuius festum agimus tradidit amoris exemplum, carnis concupiscentias execrando: solamque Christi pulchritudinem adamando? Discite virgines quales amoris Christi flammas accendit in suo pectore: que eum concupivit solum, qui pro amore hominum mortis non recusavit interitum. Hec dicentes nequaquam coniugum casta commercia vituperamus: sed virginum perseverantiam adornamus. Itaque sancta coniugia patriarcharum coniuges, videlicet Saram, Rebeccam, Rachelem, Annam, Susannam, sequantur. Virgines autem Christi: solam virginem matrem luminis, subsequantur. Huius virginis matris pedissequa fuit hec virgo: que magnum de se excitavit exemplum pudicitie.
\newline {\color{Red} Lectio Tertia.}
\lettrine[lines=2]{\zallmancaps \color{Blue} V}{irgo} hec, quia non de corporis humani elegantia placere studebat, dum de anime feditate domino displicere timebat: dabat intus per fidem sue menti candorem, et excolebat anime pulchritudinem. Tantoque studebat in carnis aspectu carnalibus displicere: quanto instabat Christo et eius angelis complacere.
%pp141
Itaque virgines Christi discite normam puellaris exempli, oblata sub specie pietatis virorum munuscula tota intentione respuite: et quasi rabidi canis morsus, omni cum gaudio recusate. Accipe virgo: quod amori proficiat eterno.
\newline \ul{Cyprianus de habitu virginum. \grecross \ }
% https://mlat.uzh.ch/browser/cps_2.CypCar.DeHaVi#:22;2
\newline Igitur o virgines, estote tales: quales vos artifex Deus fecit. Maneat in vobis facies incorrupta, cervix pura: forma sincera. Non inferantur auribus vulnera: nec includat brachia vel colla de armillis et monilibus preciosa catena. Sint a compedibus aureis pedes liberi, crines nullo colore fucati: oculi conspiciendo Deo digni.
{\ubsubsection{De Officio Beate Virginis in Sabbatis in Conventu}{de-officio-beate-virginis-in-sabbatis-in-conventu-breviarium}}
\lettrine[lines=2]{\zallmancaps \color{Red} A}{}\ul{\textsc{b} octavis Epiphanie usque ad Septuagesimam, in omnibus Sabbatis nisi festum simplex vel maius occurrerit: fiat in conventu officium de beata virgine, hoc excepto quod quando Sabbatum primum post octavas Epiphanie omnino vacaverit: faciendum est officium feriale secundum quod supra in rubrica ante Dominicam,} \hyperlink{domine-ne-in-ira}{Domine ne in ira.} \ul{notatum est. Sic autem fiat.}
\newline {\color{Red} Feria Sexta ad Vesperas. Antiphone et Psalmi Feriales. Capitulum.}
\lettrine[lines=2]{\zallmancaps \color{Blue} S}{icut} cynamomum et balsamum aromatizans odorem dedi, quasi mirra electa dedi suavitatem odoris. {\color{Red} Ymnus.} \hyperlink{ave-maris-stella}{Ave maris stella.} \V Ora pro nobis sancta dei genitrix.
\newline {\color{Red} Ad Mag. Ant.} O admirabile commercium creator generis humani animatum corpus sumens de virgine nasci dignatus est et procedens homo sine semine largitus est nobis suam deitatem. {\color{Red} Oratio.}
\lettrine[lines=2]{\zallmancaps \color{Red} D}{eus} qui salutis eterne beate Marie virginitate fecunda humano generi premia prestitisti, tribue quesumus, ut ipsam pro nobis intercedere sentiamus, per quam meruimus actorem vite suscipere. Dominum nostrum Ihesum.
\newline \ul{Quando predicam oratio dicitur ad memoriam et sequitur alia memoria immediate, terminetur sic.} Suscipere: Christum dominum nostrum.
\newline {\color{Red} Ad Completorium. Super psalmos. Ant.} Virgo Maria non est tibi similis nata in mundo inter mulieres, florens ut rosa fragrans sicut lilium, ora pro nobis sancta dei genitrix.
\newline {\color{Red} Ad Nunc Dimittis. Ant.} Sub tuum presidium confugimus sancta dei genitrix nostras deprecationes ne despicias in necessitatibus sed a periculis cunctis libera nos semper virgo benedicta. {\color{Red} Oratio.} \hyperlink{visita-quesumus}{Visita.}
\newline {\color{Red} Ad Mat. Invitat.}
\begin{invitatory}
{Ave Maria gratia plena, * Dominus tecum.}
\end{invitatory}
\newline {\color{Red} Ymnus.} \hyperlink{quem-terra}{Quem terra.}
\newline {\color{Red} Super psalmos. Ant.} Paradisi porta per Evam cunctis clausa est et per Mariam virginem iterum patefacta est. {\color{Red} Ps.} \psalm{8} \ul{et ceteri sicut in festis eiusdem.} \V Speciosa facta es et suavis. \R In delitiis tuis sancta dei genitrix.
\newline \ul{Ex sermone,} Adest nobis, \ul{qui attribuitur Augustino secundum quosdam libros vel Fulgentio secundum alios.} {\color{Red} Lectio Prima.}
% https://mlat.uzh.ch/browser/cps_2.AuInAuH.SeSuDeS2#:7.5;2
\lettrine[lines=2]{\zallmancaps \color{Blue} O}{}\ beata Maria quis tibi digne valeat iura gratiarum et laudum preconia impendere, que tuo singulari merito mundo succurristi perdito? Quas tibi laudes fragilitas humani generis persolvet, que solo tuo commercio recuperandi aditum invenit? Accipe itaque exiles et quascumque meritis tuis impares gratiarum actiones, et cum susceperis vota: culpas nostras orando excusa. \R
\lettrine[lines=2]{\zallmancaps \color{Red} C}{um} \hypertarget{cum-esset-rex}{\label{cum-esset-rex}} esset rex in accubitu suo * Nardus mea dedit odorem suum. \V Pulchra sum et decora filie Hierusalem. Nardus.
\newline {\color{Red} Lectio Secunda.}
\lettrine[lines=2]{\zallmancaps \color{Blue} A}{dmitte} piissima Dei genitrix Maria nostras preces intra sacrarium exauditionis: et reporta nobis antidotum reconciliationis. Sit per te excusabile quod per te ingerimus: sit impetrabile quod fideli mente poscimus. Accipe quod offerimus, redona quod rogamus: et excusa quod timemus.
\begin{responsory}[ave-maria-gratia]
{Ave Maria gratia plena, * Dominus tecum.}{Benedicta tu in mulieribus et benedictus fructus ventris tui.}
\end{responsory}%
\newline {\color{Red} Lectio Tertia.}
\lettrine[lines=2]{\zallmancaps \color{Red} S}{ancta} Maria succurre miseris, iuva pusillanimes: refove flebiles. Ora pro populo, interveni pro clero: intercede pro devoto femineo sexu. Sentiant omnes tuum iuvamen: quicumque celebrant tuam commemorationem. Assiste parata votis poscentium: et repende nobis omnibus optatum effectum. Sit tibi studium assidue exorare pro populo Dei: que meruisti benedicta proferre precium mundi.
\begin{responsory-final}[virgo-parens]
{Virgo parens Christi paritura deum genuisti, fulgida stella maris nos protege nos tuearis, * Dum paris et gaudes, * Cantant celi agmina laudes.}{Intercede pia pro nobis virgo Maria.}{Cantant}
\end{responsory-final}%
\psalm{tedeum} \V Ora pro nobis.
\newline {\color{Red} In Laudibus. Ant.} Post partum virgo inviolata permansisti dei genitrix intercede pro nobis. {\color{Red} Ps.} \psalm{92} \ul{cum ceteris.} {\color{Red} Capitulum.} Sicut cynamomum. {\color{Red} Ymnus.} \hyperlink{o-gloriosa}{O gloriosa domina.} \V Elegit eam deus.
\newline {\color{Red} Ad Ben. Ant.} Genuit puerpera regem cui nomen eternum et gaudium matris habes cum virginitatis honore nec primam similem visa est nec habere sequentem, alleluya. {\color{Red} Oratio.} Deus qui salutis. \ul{et dicatur ad Terciam, Sextam, et Nonam.}
\newline {\color{Red} Ad Primam. Ant.} Dignare me laudare te virgo sacrata da michi virtutem contra hostes tuos. \R \hyperlink{ihesu-christe-fili}{Ihesu Christe.} \V Qui natus es de virgine Maria.
\newline {\color{Red} Ad Terciam. Ant.} Gaude Maria virgo cunctas hereses sola interemisti in universo mundo. {\color{Red} Cap.} Sicut cynamomum.
\begin{responsory-breve}[sancta-maria-mater]
{Sancta Maria mater Christi, * Audi rogantes servulos.}{Et impetratam nobis celitus tu defer indulgentiam.}
\end{responsory-breve}%
\V Ora pro nobis.
\newline {\color{Red} Ad Sextam. Ant.} In prole mater in partu virgo gaude et letare virgo mater domini. {\color{Red} Capitulum.}
\lettrine[lines=2]{\zallmancaps \color{Blue} E}{t} radicavi in populo honorificato, et in partes dei mei hereditas illius, et in plenitudine sanctorum detentio mea.
\begin{responsory-breve}[ora-pro-nobis]
{Ora pro nobis, * Sancta dei genitrix.}{Ut digni efficiamur promissionibus Christi.}
\end{responsory-breve}%
\V Elegit eam.
\newline {\color{Red} Ad Nonam. Ant.} Beata mater et intacta virgo gloriosa regina mundi intercede pro nobis ad dominum. {\color{Red} Capitulum.}
\lettrine[lines=2]{\zallmancaps \color{Red} Q}{uasi} cedrus exaltata sum in Libano, et quasi Cypressus in monte Sion, quasi palma exaltata sum in Cades: et quasi plantatio rose in Hierico.
\begin{responsory-breve}[elegit-eam]
{Elegit eam deus * Et preelegit eam.}{Et habitare eam facit in tabernaculo suo.}
\end{responsory-breve}%
\V Sancta dei genitrix virgo semper Maria. \R Intercede pro nobis ad dominum deum nostrum.
\newline \ul{Dicta Nona terminatur officium.}
\newline \ul{Si inter Purificationem et Septuagesimam Sabbatum aliquod intervenerit: fiat officium predicto modo, excepto quod ad Magnificat dicatur,} {\color{Red} Ant.} Alma redemptoris mater que pervia celi porta manens et stella maris succurre cadenti surgere qui curat populo, tu que genuisti, natura mirante tuum sanctum genitorem, virgo prius ac posterius Gabrielis ab ore sumens illud ave peccatorum miserere.
\newline \ul{Et ad Benedictus,} {\color{Red} Ant.} Ave stella matutina peccatorum medicina, mundi princeps et regina, sola virgo digna dici contra tela inimici clipeum pone salutis tue titulum virtutis, o sponsa dei elata esto nobis via recta ad eterna gaudia.
\newline \ul{Et Oratio,} Concede nos \ul{dicatur loco illius,} Deus qui salutis. {\color{Red} Oratio.}
\lettrine[lines=2]{\zallmancaps \color{Blue} C}{oncede} nos famulos tuos quesumus domine deus perpetua mentis et corporis salute gaudere, et gloriosa beate Marie semper virginis intercessione, a presenti liberari tristicia et eterna perfrui leticia. Per.
\newline \ul{Eodem eciam modo fiat officium in omnibus Sabbatis a} \hyperlink{deus-omnium}{Deus omnium,} \ul{usque ad Adventum nisi festum simplex aut maius, vel octave vel ieiunia quatuor temporum aut vigilie sanctorum que propriam Missam habeant occurrerint. Ita tamen quod ad Mag. dicatur Ant.} Alma redemptoris: \ul{et ad Ben. Ant.} Ave stella \ul{et oratio,} Concede, \ul{ut predictum est. Quando vero propter aliquod predictorum intermittitur: tunc fiat de ea memoria in Vesperis et in Laudibus, preterquam in festis duplicibus et in totis duplicibus.}
{\ubsubsection{Canticum Graduum}{canticum-graduum-breviarium}}
\newline \ul{Quando autem predictum officium fit in choro in Sabbatis, surgendo dicatur canticum graduum hoc modo. Isti quique psalmi, scilicet,} \psalm{119} \psalm{120} \psalm{121} \psalm{122} \psalm{123} \ul{dicantur sub uno.} Requiem. \ul{Deinde} Pater noster. Et ne nos. A porta inferi. Dominus vobiscum. \ul{vel} Domine exaudi, \ul{si non sit sacerdos, qui dicit,} Oremus. {\color{Red} Oratio.}
\lettrine[lines=2]{\zallmancaps \color{Red} A}{bsolve} quesumus domine animas famulorum famularumque tuarum ab omni vinculo delictorum, ut in resurrectionis gloria inter sanctos tuos resuscitati respirent. Per Christum.
\newline \ul{Isti alii quinque psalmi scilicet,} \psalm{124} \psalm{125} \psalm{126} \psalm{127} \psalm{128} \ul{dicantur singuli cum Gloria patri. Deinde.} Kyrie eleyson. Christe eleyson. Kyrie eleyson. Pater noster. Et ne nos. \V Memor esto domine congregationis tue. \R Quam possedisti ab initio. Dominus vobiscum. Oremus. {\color{Red} Oratio.}
\lettrine[lines=2]{\zallmancaps \color{Blue} D}{eus} cui proprium est misereri semper et parcere suscipe deprecationem nostram, et quos delictorum catena constringit, miseracio tue pietatis absolvat. Per Christum.
\newline \ul{Isti alii quinque psalmi scilicet,} \psalm{129} \psalm{130} \psalm{131} \psalm{132} \psalm{133} \ul{dicantur singuli cum gloria patri. Deinde.} Kyrie eleyson. Christe eleyson. Kyrie eleyson. Pater noster. Et ne nos. \V Salvos fac servos tuos et ancillas tuas. \R Deus meus sperantes in te. Dominus vobiscum. Oremus. {\color{Red} Oratio.}
\lettrine[lines=2]{\zallmancaps \color{Red} P}{retende} domine famulis et famulabus tuis dexteram celestis auxilii: ut te toto corde perquirant et que digne postulant assequantur. Per Christum.
\newline \ul{Ad faciendam memoriam in Sabbatis de beata virgine.}
\newline {\color{Red} Per Adventum. In Vesperis. Ant.} O virgo virginum quomodo fiet istud quia nec primam similem visa es nec habere sequentem filie Hierusalem quid me admiramini divinum est mysterium hoc quod cernitis. \V Ora pro nobis. {\color{Red} Oratio.}
\lettrine[lines=2]{\zallmancaps \color{Blue} D}{eus} qui de beate Marie virginis utero verbum tuum angelo nunciante carnem suscipere voluisti, presta supplicibus tuis ut qui vere eam genitricem dei credimus, eius apud te intercessionibus adiuvemur. Per eundem.
\newline {\color{Red} In Laudibus. Ant.} Spiritus sanctus in te descendet Maria ne timeas habebis in utero filium dei, alleluya. \V Elegit eam deus. {\color{Red} Oratio.} Deus qui de beate.
\newline \ul{Ab octavis Epiphanie usque ad festum Purificationis.}
\newline {\color{Red} In Vesperis. Ant.} O admirabile. \V Ora pro nobis. {\color{Red} Oratio.} Deus qui salutis.
\newline {\color{Red} In Laudibus. Ant.} Genuit puerpera. \V Elegit. {\color{Red} Oratio.} Deus qui salutis.
\newline \ul{A festo Purificationis usque ad Ramos Palmarum.}
\newline {\color{Red} In Vesperis. Ant.} Alma redemptoris. \V Ora pro nobis. {\color{Red} Oratio.} Concede nos.
\newline {\color{Red} In Laudibus. Ant.} Ave stella. \V Elegit. {\color{Red} Oratio.} Concede.
\newline \ul{Tempore Paschali.}
\newline {\color{Red} In Vesperis. Ant.} Regina celi letare alleluya quia quem meruisti portare alleluya, resurrexit sicut dixit alleluya, ora pro nobis deum alleluya. \V Ora pro nobis. {\color{Red} Oratio.} Concede. {\color{Red} Post Ascensionem.} Iam ascendit sicut, etc.
\newline {\color{Red} In Laudibus. Ant.} Beata dei genitrix Maria virgo perpetua, templum domini, sacrarium spiritus sancti, sola sine exemplo placuisti domino Ihesu Christo, ora pro populo, interveni pro clero: intercede pro devoto femineo sexu, alleluya alleluya. \V Elegit. \ul{Oratio ut supra.}
{\ubsubsection{Antiphone De Beata Virgine Dicende Post Completorium}{antiphone-de-beata-virgine-dicende-post-completorium-breviarium}}
\newline {\color{Red} Antiphona.}
\lettrine[lines=2]{\zallmancaps \color{Red} S}{alve} regina misericordie vita dulcedo et spes nostra salve. Ad te clamamus exules filii Eve ad te suspiramus gementes et flentes in hac lacrimarum valle. Eya ergo advocata nostra illos tuos misericordes oculos ad nos converte, et Ihesum benedictum fructum ventris tui nobis post hoc exilium ostende, O clemens, O pia, O dulcis Maria. {\color{Red} Tempore Paschali.} Alleluya.
\newline {\color{Red} Alia Antiphona.}
\lettrine[lines=2]{\zallmancaps \color{Blue} A}{ve} regina celorum, Ave domina angelorum, Salve radix sancta ex qua mundo lux est orta, gaude gloriosa super omnes speciosa, vale valde decora, et pro nobis semper Christum exora. {\color{Red} Tempore Paschali.} Alleluya.
\newline \ul{Post antiphonam dicatur,} \V Dignare me laudare te virgo sacrata. \R Da michi virtutem contra hostes tuos. \ul{Tempore Paschali dicatur cum} alleluya. \ul{Deinde oratio,} Concede nos. \ul{terminanda,} Per Christum.
\newline \ul{Predicte antiphone dicantur alternatim ad libitum post Completorium omni tempore cum versiculo et oratione predictis: preterquam in quarta et quinta et sexta feria ante Pascha.}
{\ubsubsection{Officium Cotidianum Beate Virginis}{officium-cotidianum-beate-virginis-breviarium}}
\lettrine[lines=2]{\zallmancaps \color{Red} O}{}\ul{\textsc{fficium} cotidianum beate virginis intermittitur, In vigilia Natalis Domini, In Matutinis et deinceps usque in crastinum post octavas Epyphanie exclusive. Et in quarta feria ante Cenam Domini in Vesperis et deinceps usque in secundam feriam post octavas Pasche exclusive. Similiter in vigilia Pentecostes in Vesperis et deinceps usque in secundam feriam post festum Trinitatis exclusive. Intermittitur eciam in omnibus festis duplicibus et totis duplicibus in Primis Vesperis et deinceps usque in crastinum diei festi exclusive. Intermittitur eciam quandocumque fit de ea officium in conventu. Siqui tamen in supradictis temporibus dicere ex devotione voluerint: per se et summisse illud dicant.}%typo: 2nd Intermittur
{\ubsubsection{Officium Defunctorum}{officium-defunctorum-breviarium}}
\lettrine[lines=2]{\zallmancaps \color{Blue} S}{}\ul{\textsc{ciendum} quod pro defunctis semper faciende sunt vigilie novem lectionum: a fratribus in conventu sive extra ubicumque fuerint, semel in ebdomada: nisi temporibus infra notatis in quibus intermitti possunt. In Vesperis autem et Matutinis et Laudibus antiphone, psalmi, cantica, versiculi lectiones et responsoria dicenda sunt sicut supra in commemoratione defunctorum notata sunt: excepto quod nonum responsorium dicatur tantum cum his tribus versibus scilicet,} \hyperlink{libera-me-domine}{Dies illa. Tremens.} \ul{et} Creator. \ul{Post ultimum versum responsorium ex integro repetatur.}%typo novum
\newline \ul{Intermittuntur autem predicte vigilie novem lectionum in omnibus Sabbatis et in vigilia cuiuslibet festi nisi trium fuerit lectionum. Item in festis duplicibus et totis duplicibus in utrisque Vesperis. Et in vigilia Natalis Domini et deinceps usque ad octavas Epyphanie exclusive. Item in quarta feria ante Cenam Domini et deinceps usque ad feriam secundam post octavas Pasche exclusive. Item in vigilia Ascensionis et deinceps usque ad octavam ipsius exclusive. Item in vigilia Pentecostes et deinceps usque ad Dominicam post Trinitatem exclusive. Item per omnes octavas. Si tamen presens defunctus fuerit: officium vigiliarum pro eo agi poterit, non tamen in die Natalis Domini: nec in die Cene, nec in Parascheve nec in Sabbato Sancto, nec in die Pasche, nec in die Pentecostes. Quod si anniversarium quodcumque ordinis evenerit: vel pro presenti defuncto vel ex alia qualibet causa infra ebdomadam quocumque die fuerint predicte vigilie faciende: sufficit quod ille fiant, et ad alias in ipsa ebdomada non teneantur fratres.}
\lettrine[lines=2]{\zallmancaps \color{Red} I}{}\ul{\textsc{n} vigiliis vero trium lectionum, psalmi, lectiones et responsoria dicantur hoc modo. In Dominica et in quarta feria dicantur, psalmi primi nocturni cum suis antiphonis et versiculo et lectionibus et responsoriis. In secunda et quinta feria dicantur psalmi secundi nocturni: cum suis antiphonis, et versiculo et lectionibus et responsoriis. In tercia vero et sexta feria dicantur psalmi tercii nocturni, similiter cum suis antiphonis, versiculo, lectionibus, et responsoriis. In ultimo autem responsorio dicatur tantum primus versus, scilicet,} \hyperlink{libera-me-domine}{Dies illa.} \ul{Post resumptionem vero, responsorium non repetatur. Quod si festum intervenerit: nichilominus in feriis singulis predictus ordo servetur.}
\newline \ul{Hoc autem ordine dicantur Orationes in vigiliis defunctorum. Pro presenti sive per uno defuncto: si episcopus fuerit: prima Oratio dicatur,} Deus qui inter apostolicos, \ul{si alius,} Inclina, \ul{si vero femina,} Quesumus domine, \ul{secunda,} Deus venie, \ul{tercia,} Fidelium.
\ul{Pro fratribus, familiaribus et benefactoribus, prima,} Deus venie, \ul{secunda,} Deus qui nos patrem, \ul{tercia,} Fidelium.
\ul{In omni anniversario, prima,} Deus indulgentiarum, \ul{secunda,} Deus venie, \ul{tercia,} Fidelium.
\ul{Si simul evenerint obitus et anniversarium dimittatur,} Deus venie, \ul{et prima Oratio dicatur pro presenti defuncto, secunda pro anniversario tercia non mutatur. In vigiliis autem trium lectionum dicatur prima,} Deus venie, \ul{secunda,} Deus qui nos patrem, \ul{tercia} Fidelium.
{\ubsubsection{Orationes in Officio Defunctorum}{orationes-in-officio-defunctorum-defunctorum-breviarium}}
\newline {\color{Red} Pro Episcopo Defuncto. Oratio.}
\lettrine[lines=2]{\zallmancaps \color{Blue} D}{eus} qui inter apostolicos sacerdotes, famulum tuum pontificali fecisti dignitate vigere, presta quesumus, ut eorum quoque perpetuo aggregetur consortio. Per.
\newline {\color{Red} Pro Viro Defuncto. Oratio.}
\lettrine[lines=2]{\zallmancaps \color{Red} I}{nclina} domine aurem tuam ad preces nostras, quibus misericordiam tuam supplices deprecamur ut animam famuli tui quam de hoc seculo migrare iussisti, in pacis ac lucis regione constituas et sanctorum tuorum iubeas esse consortem. Per.
\newline {\color{Red} Pro Femina Defuncta. Oratio.}
\lettrine[lines=2]{\zallmancaps \color{Blue} Q}{uesumus} domine pro tua pietate miserere anime famule tue, et a contagiis mortalitatis exutam in eterne salvationis partem restitue. Per.
\newline {\color{Red} In Omni Anniversario. Oratio.}
\lettrine[lines=2]{\zallmancaps \color{Red} D}{eus} indulgentiarum domine da animabus famulorum tuorum quorum anniversarium depositionis diem commemoramus refrigerii sedem, quietis beatitudinem luminis claritatem. Per.
\newline {\color{Red} Pro Fratribus Familiaribus et Benefactoribus Defunctis. Oratio.}
\lettrine[lines=2]{\zallmancaps \color{Blue} D}{eus} venie largitor et humane salutis actor, quesumus clementiam tuam, ut nostre congregationis fratres familiares et benefactores qui ex hoc seculo transierunt, beata Maria semper virgine intercedente cum omnibus sanctis ad perpetue beatitudinis consortium pervenire concedas. Per.
\newline {\color{Red} Pro Parentibus Defunctis. Oratio.}
\lettrine[lines=2]{\zallmancaps \color{Red} D}{eus} qui nos patrem et matrem honorare precepisti miserere clementer animabus parentum nostrorum, eorumque peccata dimitte, nosque eos in eterne claritatis gaudio fac videre. Per.
\newline {\color{Red} Pro Omnibus Fidelibus Defunctis. Oratio.}
\lettrine[lines=2]{\zallmancaps \color{Blue} F}{idelium} deus omnium conditor et redemptor, animabus famulorum famularumque tuarum remissionem cunctorum tribue peccatorum, ut indulgentiam quam semper optaverunt, piis supplicationibus consequantur. Qui vivis et regnas cum deo, et cetera.
{\ubsubsection{Ante Completorium}{ante-completorium-breviarium}}
\lettrine[lines=2]{\zallmancaps \color{Red} A}{}\ul{\textsc{nte} Completorium dicatur,} Iube domne benedicere. \ul{Et subiungatur benedictio,} Noctem quietam et finem perfectum tribuat nobis omnipotens et misericors dominus. Amen. \ul{Postea sequitur.} {\color{Red} Lectio.}
\lettrine[lines=2]{\zallmancaps \color{Blue} F}{ratres} sobrii estote et vigiliate, quia adversarius vester diabolus tanquam leo rugiens circuit querens quem devoret: cui resistite fortes in fide. Tu autem. \ul{Postea dicatur,} Adiutorium nostrum in nomine domini. \R Qui fecit celum et terram. \ul{Deinde dicto,} Pater noster, \ul{dicatur,} Confiteor. \ul{Postea Completorium inchoetur.}
{\ubsubsection{Ad Confiteor}{ad-confiteor-breviarium}}
\lettrine[lines=2]{\zallmancaps \color{Red} C}{onfiteor} deo et beate Marie et omnibus sanctis et vobis fratres, quia peccavi nimis cogitatione locutione, opere et omissione mea culpa precor vos orate pro me. {\color{Red} Responsio.}
\lettrine[lines=2]{\zallmancaps \color{Blue} M}{isereatur} tui omnipotens deus, et dimittat tibi omnia peccata tua liberet te ab omni malo, salvet et confirmet in omni opere bono, et perducat ad vitam eternam. \R Amen.%typo: obere
\newline \ul{Quando frater solus dicit Primam vel Completorium, nec habet cum quo dicat,} Confiteor: \ul{dicat sic,}
\lettrine[lines=2]{\zallmancaps \color{Red} C}{onfiteor} deo et beate Marie et omnibus sanctis quia peccavi nimis cogitatione, locutione, opere et omissione, mea culpa precor beatam Mariam et omnes sanctos orare pro me. \ul{Postea subiungat.}
\lettrine[lines=2]{\zallmancaps \color{Blue} M}{isereatur} michi omnipotens deus et dimittat michi omnia peccata mea liberet me, \ul{et cetera.}
{\ubsubsection{Ad Preciosa}{ad-preciosa-breviarium}}
\newline {\bfseries \color{Blue} P}reciosa est in conspectu domini. \R Mors sanctorum eius. {\color{Red} Oratio.}
\lettrine[lines=2]{\zallmancaps \color{Red} S}{ancta} Maria et omnes sancti intercedant pro nobis ad dominum, ut mereamur adiuvari ab eo: qui vivit et regnat per omnia secula seculorum. Amen. \V Deus in adiutorium meum intende. \R Domine ad adiuvandum me festina. \ul{Et dicatur ter. Deinde.} Gloria patri, \ul{et} Sicut erat, \ul{etc.} Kyrie eleyson. Christe eleyson. Kyrie eleyson. Pater noster. Et ne nos. Sed libera. \V Respice domine in servos tuos et in opera tua, et dirige filios eorum. \R Et sit splendor domini dei nostri super nos, et opera manuum nostrarum dirige super nos, et opus manuum nostrarum dirige. \ul{Deinde sequitur.} {\color{Red} Oratio.}
\lettrine[lines=2]{\zallmancaps \color{Blue} D}{irigere} et sanctificare digneris domine sancte pater omnipotens eterne deus hodie corda et corpora nostra in lege tua, et in operibus mandatorum tuorum, ut hic et in eternum te auxiliante semper salvi esse mereamur. Per Christum dominum nostrum. Amen.
\newline \ul{Postea dicto,} Iube, \ul{si de evangelio legendum sit, detur benedictio.} Divinum auxilium maneat semper nobiscum. Amen. \ul{Si vero de constitutionibus: detur benedictio.} Regularibus disciplinis, instruat nos magister celestis. Amen.
\newline \ul{Data benedictione sequitur lectio de evangelio vel de constitutionibus, fratres tamen qui extra conventum sunt, si evangelium in promptu non habuerint: dicant.} {\color{Red} Secundum Iohannem.}
\lettrine[lines=2]{\zallmancaps \color{Red} I}{n} principio erat verbum. Et verbum erat apud deum, et deus erat verbum. Hoc erat in principio apud deum. Omnia per ipsum facta sunt, et sine ipso factum est nichil. Tu autem.
\newline \ul{De constitutionibus vero dicatur.} {\color{Red} Lectio.}
\lettrine[lines=2]{\zallmancaps \color{Blue} Q}{uoniam} ex precepto regule iubemur habere cor unum et animam unam in domino: iustum est ut qui sub una regula et unius professionis voto vivimus uniformes in observantiis canonice religionis inveniamur. Tu autem. \ul{Finita lectione sequitur,} Commemoratio fratrum familiarium benefactorum defunctorum: ordinis nostri. \ul{Prelatus,} Requiescant in pace. Amen. \ul{Deinde sequitur.} {\color{Red} Ps.} \psalm{116} \ul{etc.} Gloria patri. Sicut erat. \ul{Deinde,} \V Ostende nobis domine misericordiam tuam. \R Et salutare tuum da nobis. Dominus vobiscum, vel Domine exaudi. Oremus. {\color{Red} Oratio.}
\lettrine[lines=2]{\zallmancaps \color{Red} A}{ctiones} nostras quesumus domine aspirando preveni et adiuvando prosequere: ut cuncta nostra operatio a te semper incipiat, et per te cepta finiatur. Per Christum dominum nostrum. Amen. \ul{Postea dicatur,} Adiutorium nostrum in nomine domini. \R Qui fecit celum et terram.
{\ubsubsection{Benedictiones Lectionum in Matutinis}{benedictiones-lectionum-in-matutinis-breviarium}}
\newline {\color{Red} In Primo Nocturno.}
\lettrine[lines=2]{\zallmancaps \color{Blue} B}{enedictione} perpetua, benedicat nos pater eternus. {\color{Red} Responsio.} Amen.
%pp142
\newline Unigenitus dei filius, nos benedicere et adiuvare dignetur.
\newline Spiritus sancti gratia illuminare dignetur sensus et corda nostra.
\newline {\color{Red} In Secundo Nocturno.}
\newline {\bfseries \color{Blue} D}eus pater omnipotens, sit nobis propitius et clemens.
\newline Ad gaudia paradisi, perducat nos misericordia Christi.
\newline Ignem sui amoris, accendat deus in cordibus nostris.
\newline {\color{Red} In Tercio Nocturno.}
\newline {\bfseries \color{Red} I}lle nos benedicat, qui sine fine vivit et regnat.
\newline {\color{Red} Quando omelia legitur.}
\newline {\bfseries \color{Blue} E}vangelica lectio, sit nobis salus et protectio.
\newline Divinum auxilium, maneat semper nobiscum.
\newline Ad societatem civium supernorum, perducat nos rex angelorum.
\newline {\color{Red} In Natali Domini Octava Benedictio.}
\newline Per evangelica dicta, deleantur nostra delicta.
\newline Intellectum sancti evangelii, aperiat nobis virtus spiritus sancti.
{\ubsubsection{Benedictiones in Festis Beate Virginis}{benedictiones-in-festis-beate-virginis-breviarium}}
\newline {\color{Red} In Primo Nocturno.}
\newline {\bfseries \color{Red} A}lma virgo virginum, intercedat pro nobis ad dominum.
\newline Sancta dei genitrix, sit nobis auxiliatrix.
\newline Nos cum prole pia, benedicat virgo Maria.
\newline {\color{Red} In Secundo Nocturno.}
\newline {\bfseries \color{Blue} A}b hoste maligno eripiat nos dei genitrix virgo.
\newline Ihesus Marie filius, sit nobis clemens et propicius.
\newline Stella Maria maris succurre piissima nobis.
\newline {\color{Red} In Tercio Nocturno.}
\newline {\bfseries \color{Red} O}ret voce pia, pro nobis virgo Maria. {\color{Red} Vel.} Evangelica lectio, \ul{si fuerit Omelia.}
\newline In tribulatione et angustia subveniat nobis pia virgo Maria.
\newline Ad societatem civium supernorum, perducat nos regina angelorum.
\newline \ul{In Sabbatis dicantur tres ultime quando de ea fit officium in conventu.}
% {\color{Red} .}
% {\color{Red} .}
% {\color{Red} .}
% {\color{Red} .}

% \lettrine[lines=2]{\zallmancaps \color{Red} O}{}
% \lettrine[lines=2]{\zallmancaps \color{Blue} I}{n}
% \lettrine[lines=2]{\zallmancaps \color{Red} O}{}

% \begin{responsory}[]
% \end{responsory}%

% \newline \ul{} \grecross \ {\color{Red} Lectio Prima.}
% \newline \ul{} \grecross \ {\color{Red} Lectio Secunda.}
% \newline \ul{} \grecross \ {\color{Red} Lectio Tertia.}
% \newline \ul{} \grecross \ {\color{Red} Lectio Quarta.}
% \newline \ul{} \grecross \ {\color{Red} Lectio Quinta.}
% \newline \ul{} \grecross \ {\color{Red} Lectio Sexta.}
% \newline \ul{} \grecross \ {\color{Red} Lectio VII. Secundum Matheum.}
% \newline \ul{} \grecross \ {\color{Red} Omelia Beati.}
% \newline \ul{} \grecross \ {\color{Red} Lectio VIII.}
% \newline \ul{} \grecross \ {\color{Red} Lectio IX.}

% \newline {\color{Red} Lectio Octava.}
% \newline {\color{Red} Lectio Nona.}

% \newline {\color{Red} Ad Mat. Invitat.}
% \begin{invitatory}[-invitatorium]
% {}
% \end{invitatory}
% \newline {\color{Red} Ymnus.} \hyperlink{}{.}

% \newline {\color{Red} In Primo Nocturno. Ant.} . {\color{Red} Ps.} \psalm{1} {\color{Red} Ant.} . {\color{Red} Ps.} \psalm{2} {\color{Red} Ant.} . {\color{Red} Ps.} \psalm{3} \V 
% \newline {\color{Red} In Secundo Nocturno. Ant.} . {\color{Red} Ps.} \psalm{1} {\color{Red} Ant.} . {\color{Red} Ps.} \psalm{2} {\color{Red} Ant.} . {\color{Red} Ps.} \psalm{3} \V 
% \newline {\color{Red} In Tertio Nocturno. Ant.} . {\color{Red} Ps.} \psalm{1} {\color{Red} Ant.} . {\color{Red} Ps.} \psalm{2} {\color{Red} Ant.} . {\color{Red} Ps.} \psalm{3} \V 

% \newline {\color{Red} In Laudibus. Ant.} . {\color{Red} Ps.} \psalm{92} {\color{Red} Ant.} . {\color{Red} Ps.} \psalm{99} {\color{Red} Ant.} . {\color{Red} Ps.} \psalm{62} {\color{Red} Ant.} . {\color{Red} Can.} \psalm{benedicite} {\color{Red} Ant.} . {\color{Red} Ps.} \psalm{148}
% \newline {\color{Red} In Laudibus. Ant.} . {\color{Red} Ant.} . {\color{Red} Ant.} . {\color{Red} Ant.} . {\color{Red} Ant.} .
% \newline {\color{Red} Ad Ben. Ant.} 
% \newline {\color{Red} Ad Mag. Ant.} 
% \newline {\color{Red} Ad Terciam. Cap.} 
% \newline {\color{Red} Ad Sextam. Cap.}
% \newline {\color{Red} Ad Nonam. Cap.}

% {\ubsubsection{Ad Vesperas}{-breviarium}}
% \newline {\color{Red} Oratio.}
% \newline {\color{Red} Ad Vesperas. Super psalmos. Ant.} 

% \hymn{}{}{
% \par {\bfseries \color{Red} }
% \par {\bfseries \color{Blue} }
% \par {\bfseries \color{Red} }
% \par {\bfseries \color{Blue} }
% \par {\bfseries \color{Red} }
% }

% \grecross \grealtcross \gothRbar \gothVbar \greheightstar \ \gresixstar

\end{multicols}\fancyhead[RO,LE]{Commune Sanctorum}

% \backmatter

% \thispagestyle{fancy}

% \phantomsection
% \addcontentsline{toc}{chapter}{Indices}
% \hypertarget{Indices}{}
% \fancyhead[C]{}
% \fancyhead[RE,LO]{Indices}
% Optional: change headers for index pages
% \renewcommand{\leftmark}{Indices}
% \renewcommand{\rightmark}{Index Hymnorum}

% Print each index
% \cprintindex{P}{Index Psalmorum et Canticorum}
% \newpage
% \cprintindex{H}{Index Hymnorum}
% \newpage
% \cprintindex{R}{Index Responsoriorum}%It does not index manual labels of responsories
% \newpage
% \cprintindex{S}{Index Responsoriorum Breve}
% \newpage
% \cprintindex{I}{Index Invitatorium}

\end{document}